% Federigo Enriques, Le superficie algebriche (1949).
% Complete editable reader edition. Compile with pdfLaTeX, twice (or latexmk -pdf).
% Source comments beginning % ?? and % [sic] are retained verbatim.
% Resolved doubts carry a separate EDITORIAL status comment and no active doubt note.
\documentclass[11pt,a4paper,twoside,openany,leqno]{book}
\usepackage[T1]{fontenc}
\usepackage[utf8]{inputenc}
\usepackage{lmodern}
\usepackage[sc]{mathpazo}
\AtBeginDocument{\DeclareFontShape{T1}{pplx}{b}{sc}{<->ssub*pplx/m/sc}{}}
\usepackage{amsmath,amssymb,mathtools}
\usepackage{microtype}
\usepackage{geometry}
\geometry{inner=30mm,outer=46mm,top=24mm,bottom=26mm,headheight=25pt,headsep=8mm,marginparwidth=35mm,marginparsep=6mm}
\usepackage[english]{babel}
\usepackage{array,longtable,booktabs}
\usepackage{ragged2e}
\usepackage{xcolor}
\definecolor{editiongray}{gray}{0.28}
\usepackage{marginnote,marginfix}
\usepackage{fancyhdr}
\usepackage{needspace}
\usepackage{titlesec}
\usepackage[unicode,pdfencoding=auto,psdextra,hyperfootnotes=false,hidelinks]{hyperref}
\usepackage{bookmark}
\hypersetup{pdftitle={Le superficie algebriche},pdfauthor={Federigo Enriques},pdfsubject={Edizione di lettura dalla stampa Zanichelli, Bologna 1949},pdfcreator={pdfLaTeX; edizione IA, novembre 2026},bookmarksopen=true,bookmarksopenlevel=0}
\newif\ifshownotes
\shownotesfalse % Change this ONE flag to \shownotesfalse for a reading copy without editorial marginal notes.
\setlength{\parindent}{1.3em}
\setlength{\parskip}{0pt plus .35pt}
\linespread{1.045}
\raggedbottom
\setlength{\emergencystretch}{2.5em}
\tolerance=1600
\clubpenalty=8000
\widowpenalty=8000
\displaywidowpenalty=5000
\allowdisplaybreaks[2]
\setlength{\abovedisplayskip}{10pt plus3pt minus3pt}
\setlength{\belowdisplayskip}{10pt plus3pt minus3pt}
\setlength{\abovedisplayshortskip}{5pt plus2pt minus2pt}
\setlength{\belowdisplayshortskip}{7pt plus2pt minus2pt}
\setlength{\LTpre}{4pt}
\setlength{\LTpost}{7pt}
\setlength{\tabcolsep}{4pt}
\setlength{\marginparpush}{6pt}
\setlength{\footnotesep}{8pt}
\interfootnotelinepenalty=150
\renewcommand{\thefootnote}{(\arabic{footnote})}
\pagestyle{fancy}
\fancyhf{}
\fancyhead[LE,RO]{\normalfont\fontsize{8}{9.5}\selectfont\parbox[b]{\textwidth}{\nouppercase{\leftmark}}}
\fancyfoot[LE,RO]{\small\thepage}
\renewcommand{\headrulewidth}{0.25pt}
\fancypagestyle{plain}{\fancyhf{}\fancyfoot[LE,RO]{\small\thepage}\renewcommand{\headrulewidth}{0pt}}
\titleformat{\subsection}{\normalfont\bfseries\large}{}{0pt}{}
\titlespacing*{\subsection}{0pt}{2.4ex plus1ex minus.4ex}{.9ex plus.3ex}
\newcommand{\gesperrt}[1]{\textls[100]{#1}}
\newcommand{\weq}[2]{\begin{equation*}\tag*{\ensuremath{#1}}#2\end{equation*}}
\makeatletter
\newif\ifedtable
\AddToHook{env/longtable/before}{\global\edtabletrue}
\AddToHook{env/longtable/after}{\global\edtablefalse\edFlush}
\AddToHook{env/tabular/before}{\global\edtabletrue}
\AddToHook{env/tabular/after}{\global\edtablefalse}
\let\edPending\@empty
\newcommand{\edQueue}[1]{\ifshownotes\g@addto@macro\edPending{#1}\fi}
\newcommand{\edFlush}{\ifshownotes\ifedtable\else\ifmmode\else\ifx\edPending\@empty\else\ifinner\else\let\edTemp\edPending\global\let\edPending\@empty\edTemp\fi\fi\fi\fi\fi}
\newcommand{\edNote}[3]{\marginpar{\normalfont\fontsize{7.1}{8.4}\selectfont\RaggedRight\color{editiongray}\textbf{#1}\, [#2]\quad #3\par}}
\newcommand{\printersignature}[1]{\edQueue{\marginpar{\normalfont\fontsize{6.8}{8}\selectfont\RaggedRight\color{editiongray}Segnatura: #1\par}}}
\newcommand{\weberanchor}[1]{\hypertarget{old:#1}{}\ifcsname wpAlias:#1\endcsname\hypertarget{old:\csname wpAlias:#1\endcsname}{}\fi}
\newcommand{\weberpage}[1]{\ifvmode\leavevmode\fi\weberanchor{#1}\begingroup\reversemarginpar\marginnote{\normalfont\fontsize{8}{9}\selectfont\color{editiongray}\ifcsname wpAlias:#1\endcsname[\csname wpAlias:#1\endcsname]\else#1\fi}\endgroup\ifmeasuring@\else\iffirstchoice@\ifcsname wpNotes:#1\endcsname\expandafter\edQueue\expandafter{\csname wpNotes:#1\endcsname}\fi\ifcsname wpDefer:#1\endcsname\else\edFlush\fi\fi\fi\ignorespaces}
\newcommand{\weberblank}[1]{\weberanchor{#1}}
\newcommand{\origfootnote}[3]{\begingroup\renewcommand\thefootnote{(#1)}\footnote[#1]{\hypertarget{fn:#2}{}\normalfont#3}\endgroup}
\newcommand{\origfootmark}[2]{\textsuperscript{\normalfont\hyperlink{fn:#2}{(#1)}}}
\newcommand{\origfoottext}[3]{\begingroup\renewcommand\thefootnote{(#1)}\footnotetext[#1]{\hypertarget{fn:#2}{}\normalfont#3}\endgroup}
\makeatother
\newcommand{\edchapter}[4]{\pdfbookmark[0]{#3}{ch.#1}\markboth{Chapter #1. #4}{Chapter #1. #4}\vspace*{12mm}\begin{center}{\large\scshape Chapter #1.}\par\vspace{7mm}{\Large\bfseries #2\par}\end{center}\vspace{8mm}}
\newcommand{\edfrontheading}[3]{\pdfbookmark[0]{#2}{#3}\markboth{#2}{#2}\vspace*{8mm}\begin{center}{\Large\bfseries#1}\end{center}\vspace{6mm}}
\newcommand{\edsection}[3]{\subsection*{#1}\pdfbookmark[1]{#2}{#3}}
\newcommand{\editiongap}{\par\bigskip\hypertarget{old:284}{}\hypertarget{old:285}{}\pdfbookmark[1]{Gap: pages 284–285}{gap.284.285}\noindent\fbox{\begin{minipage}{\dimexpr\linewidth-2\fboxsep-2\fboxrule\relax}\small\textbf{Gap in the source: pages 284–285.}\par\medskip Pages 284 and 285 are missing from the scan. The text of p.~283 breaks off in the middle of a word; the source resumes on p.~286, with § 8. No text has been reconstructed.\end{minipage}}\par\bigskip}
\newcommand{\sourcewarning}[1]{\par\medskip\noindent\begin{minipage}{\linewidth}\small\color{editiongray}\textbf{Note on the source.} #1\end{minipage}\par\medskip}
\newcommand{\editionpage}{\thispagestyle{plain}\pdfbookmark[0]{Note on the edition}{front.edizione}\markboth{Note on the edition}{Note on the edition}\vspace*{7mm}\begin{center}\Large\scshape Note on the edition\end{center}\bigskip
\noindent Federigo Enriques, \textit{Le superficie algebriche}, Bologna, Nicola Zanichelli Editore, 1949. English translation.\par\medskip
\noindent Scan: Maxim Leyenson.\par\medskip
\noindent Transcription AI-produced (OpenAI gpt-6-astra, November 2026). Typesetting and corrections of the Italian edition AI-produced (ChatGPT, Pro mode; further corrections: Claude, Anthropic). Translation AI-produced (Claude Opus 5.5, Anthropic, November 2026) from the Italian text as printed.\par\bigskip
\noindent The original work is in the public domain in the European Union: Federigo Enriques died in 1946 and Guido Castelnuovo, the author of the preface, in 1952. Elsewhere the situation may differ. To the extent that we hold rights in the AI-produced transcription, translation, corrections and typesetting, they are dedicated to the public domain under CC0 1.0. This edition is published for educational, non-profit purposes, and we claim no rights we do not hold. If you hold rights in this work and object to its publication here, please open a GitHub issue and we will take it down.\par\bigskip
\noindent\textbf{Reading guide.} The translation follows the Italian edition page by page and keeps Enriques's notation, numbering and formulas as printed. The figures in the inner margin give the pagination of 1949; those in square brackets are inferred from the sequence or from the table of contents where the number is not printed or not legible. The entries of the CONTENTS link to these places. The author's footnotes keep the numbering of the original pages. The editorial margin notes of the Italian edition (doubtful readings and possible misprints) are not repeated here; they are in the Italian edition and its \texttt{corrections.csv}.\par\medskip
\noindent\textbf{Limits of the source.} Pages 284--285 are missing. Passages covered or damaged in the scans, in particular on pp.~186 and 368, are marked without further reconstruction. The publisher's device on the title page is taken from the scan.\par}

\newsavebox{\publisherbox}\sbox{\publisherbox}{\hbox to 23mm{\vrule width0pt height26.25mm\pdfliteral{q 65.19685 0 0 74.40945 0 0 cm BI /W 255 /H 291 /BPC 8 /CS /DeviceGray /F /ASCIIHexDecode ID
FFFFFFFFFFFFFFFFFFFFFFFFFFFFFFFFFFFFFFFFFFFFFFFFFFFFFFFFFFFFFFFFFFFFFFFFFFFFFFFFFFFFFFFFFFFFFFFFFFFFFFFFFFFFFFFFFFFFFFFFFFFFFFFFFFFFFFFFFFFFFFFFFFFFFFFFFFFFFFFFFFFFFFFFFFFFFFFFFFFFFFFFFFFFFFFFFFFFFFFFFFFFFFFFFFFFFFFFFFFFFFFFFFFFFFFFFFFFFFFFFFFFFFFFFFFFFFFFFFFFFFFFFFFFFFFFFFFFFFFFFFFFFFFFFFFFFFFFFFFFFFFFFFFFFFFFFFFFFFFFFFFFFFFFFFFFFFFFFFFFFFFFFFFFFFFFFFFFFFFFFFFFFFFFFFFFFFFFFFFFFFFFFFFFFFFFFFFFFFFFFFFFFFFFFFFFFFFFFFFFFFFFFFFFFFFFFFFFFFFFFFFFFFFFFFFFFFFFFFFFFFFFFFFFFFFFFFFFFFFFFFFFFFFFFFFFFFFFFFFFFFFFFFFFFFFFFFFFFFFFFFFFFFFFFFFFFFFFFFFFFFFFFFFFFFFFFFFFFFFFFFFFFFFFFFFFFFFFFFFFFFFFFFFFFFFFFFFFFFFFFFFFFFFFFFFFFFFFFFFFFFFFFFFFFFFFFFFFFFFFFFFFFFFFFFFFFFFFFFFFFFFFFFFFFFFFFFFFFFFFFFFFFFFFFFFFFFFFFFFFFFFFFFFFFFFFFFFFFFFFFFFFFFFFFFFFFFFFFFFFFFFFFFFFFFFFFFFFFFFFFFFFFFFFFFFFFFFFFFFFFFFFFFFFFFFFFFFFFFFFFFFFFFFFFFFFFFFFFFFFFFFFFFFFFFFFFFFFFFFFFFFFFFFFFFFFFFFFFFFFFFFFFFFFFFFFFFFFFFFFFFFFFFFFFFFFFFFFFFFFFFFFFFFFFFFFFFFFFFFFFFFFFFFFFFFFFFFFFFFFFFFFFFFFFFFFFFFFFFFFFFFFFFFFFFFFFFFFFFFFFFFFFFFFFFFFFFFFFFFFFFFFFFFFFFFFFFFFFFFFFFFFFFFFFFFFFFFFFFFFFFFFFFFFFFFFFFFFFFFFFFFFFFFFFFFFFFFFFFFFFFFFFFFFFFFFFFFFFFFFFFFFFFFFFFFFFFFFFFFFFFFFFFFFFFFFFFFFFFFFFFFFFFFFFFFFFFFFFFFFFFFFFFFFFFFFFFFFFFFFFFFFFFFFFFFFFFFFFFFFFFFFFFFFFFFFFFFFFFFFFFFFFFFFFFFFFFFFFFFFFFFFFFFFFFFFFFFFFFFFFFFFFFFFFFFFFFFFFFFFFFFFFFFFFFFFFFFFFFFFFFFFFFFFFFFFFFFFFFFFFFFFFFFFFFFFFFFFFFFFFFFFFFFFFFFFFFFFFFFFFFFFFFFFFFFFFFFFFFFFFFFFFFFFFFFFFFFFFFFFFFFFFFFFFFFFFFFFFFFFFFFFFFFFFFFFFFFFFFFFFFFFFFFFFFFFFFFFFFFFFFFFFFFFFFFFFFFFFFFFFFFFFFFFFFFFFFFFFFFFFFFFFFFFFFFFFFFFFFFFFFFFFFFFFFFFFFFFFFFFFFFFFFFFFFFFFFFFFFFFFFFFFFFFFFFFFFFFFFFFFFFFFFFFFFFFFFFFFFFFFFFFFFFFFFFFFFFFFFFFFFFFFFFFFFFFFFFFFFFFFFFFFFFFFFFFFFFFFFFFFFFFFFFFFFFFFFFFFFFFFFFFFFFFFFFFFFFFFFFFFFFFFFFFFFFFFFFFFFFFFFFFFFFFFFFFFFFFFFFFFFFFFFFFFFFFFFFFFFFFFFFFFFFFFFFFFFFFFFFFFFFFFFFFFFFFFFFFFFFFFFFFFFFFFFFFFFFFFFFFFFFFFFFFFFFFFFFFFFFFFFFFFFFFFFFFFFFFFFFFFFFFFFFFFFFFFFFFFFFFFFFFFFFFFFFFFFFFFFFFFFFFFFFFFFFFFFFFFFFFFFFFFFFFFFFFFFFFFFFFFFFFFFFFFFFFFFFFFFFFFFFFFFFFFFFFFFFFFFFFFFFFFFFFFFFFFFFFFFFFFFFFFFFFFFFFFFFFFFFFFFFFFFFFFFFFFFFFFFFFFFFFFFFFFFFFFFFFFFFFFFFFFFFFFFFFFFFFFFFFFFFFFFFFFFFFFFFFFFFFFFFFFFFFFFFFFFFFFFFFFFFFFFFFFFFFFFFFFFFFFFFFFFFFFFFFFFFFFFFFFFFFFFFFFFFFFFFFFFFFFFFFFFFFFFFFFFFFFFFFFFFFFFFFFFFFFFFFFFFFFFFFFFFFFFFFFFFFFFFFFFFFFFFFFFFFFFFFFFFFFFFFFFFFFFFFFFFFFFFFFFFFFFFFFFFFFFFFFFFFFFFFFFFFFFFFFFFFFFFFFFFFFFFFFFFFFFFFFFFFFFFFFFFFFFFFFFFFFFFFFFFFFFFFFFFFFFFFFFFFFFFFFFFFFFFFFFFFFFFFFFFFFFFFFFFFFFFFFFFFFFFFFFFFFFFFFFFFFFFFFFFFFFFFFFFFFFFFFFFFFFFFFFFFFFFFFFFFFFFFFFFFFFFFFFFFFFFFFFFFFFFFFFFFFFFFFFFFFFFFFFFFFFFFFFFFFFFFFFFFFFFFFFFFFFFFFFFFFFFFFFFFFFFFFFFFFFFFFFFFFFFFFFFFFFFFFFFFFFFFFFFFFFFFFFFFFFFFFFFFFFFFFFFFFFFFFFFFFFFFFFFFFFFFFFFFFFFFFFFFFFFFFFFFFFFFFFFFFFFFFFFFFFFFFFFFFFFFFFFFFFFFFFFFFFFFFFFFFFFFFFFFFFFFFFFFFFFFFFFFFFFFFFFFFFFFFFFFFFFFFFFFFFFFFFFFFFFFFFFFFFFFFFFFFFFFFFFFFFFFFFFFFFFFFFFFFFFFFFFFFFFFFFFFFFFFFFFFFFFFFFFFFFFFFFFFFFFFFFFFFFFFFFFFFFFFFFFFFFFFFFFFFFFFFFFFFFFFFFFFFFFFFFFFFFFFFFFFFFFFFFFFFFFFFFFFFFFFFFFFFFFFFFFFFFFFFFFFFFFFFFFFFFFFFFFFFFFFFFFFFFFFFFFFFFFFFFFFFFFFFFFFFFFFFFFFFFFFFFFFFFFFFFFFFFFFFFFFFFFFFFFFFFFFFFFFFFFFFFFFFFFFFFFFFFFFFFFFFFFFFFFFFFFFFFFFFFFFFFFFFFFFFFFFFFFFFFFFFFFFFFFFFFFFFFFFFFFFFFFFFFFFFFFFFFFFFFFFFFFFFFFFFFFFFFFFFFFFFFFFFFFFFFFFFFFFFFFFFFFFFFFFFFFFFFFFFFFFFFFFFFFFFFFFFFFFFFFFFFFFFFFFFFFFFFFFFFFFFFFFFFFFFFFFFFFFFFFFFFFFFFFFFFFFFFFFFFFFFFFFFFFFFFFFFFFFFFFFFFFFFFFFFFFFFFFFFFFFFFFFFFFFFFFFFFFFFFFFFFFFFFFFFFFFFFFFFFFFFFFFFFFFFFFFFFFFFFFFFFFFFFFFFFFFFFFFFFFFFFFFFFFFFFFFFFFFFFFFFFFFFFFFFFFFFFFFFFFFFFFFFFFFFFFFFFFFFFFFFFFFFFFFFFFFFFFFFFFFFFFFFFFFFFFFFFFFFFFFFFFFFFFFFFFFFFFFFFFFFFFFFFFFFFFFFFFFFFFFFFFFFFFFFFFFFFFFFFFFFFFFFFFFFFFFFFFFFFFFFFFFFFFFFFFFFFFFFFFFFFFFFFFFFFFFFFFFFFFFFFFFFFFFFFFFFFFFFFFFFFFFFFFFFFFFFFFFFFFFFFFFFFFFFFFFFFFFFFFFFFFFFFFFFFFFFFFFFFFFFFFFFFFFFFFFFFFFFFFFFFFFFFFFFFFFFFFFFFFFFFFFFFFFFFFFFFFFFFFFFFFFFFFFFFFFFFFFFFFFFFFFFFFFFFFFFFFFFFFFFFFFFFFFFFFFFFFFFFFFFFFFFFFFFFFFFFFFFFFFFFFFFFFFFFFFFFFFFFFFFFFFFFFFFFFFFFFFFFFFFFFFFFFFFFFFFFFFFFFFFFFFFFFFFFFFFFFFFFFFFFFFFFFFFFFFFFFFFFFFFFFFFFFFFFFFFFFFFFFFFFFFFFFFFFFFFFFFFFFFFFFFFFFFFFFFFFFFFFFFFFFFFFFFFFFFFFFFFFFFFFFFFFFFFFFFFFFFFFFFFFFFFFFFFFFFFFFFFFFFFFFFFFFFFFFFFFFFFFFFFFFFFFFFFFFFFFFFFFFFFFFFFFFFFFFFFFFFFFFFFFFFFFFFFFFFFFFFFFFFFFFFFFFFFFFFFFFFFFFFFFFFFFFFFFFFFFFFFFFFFFFFFFFFFFFFFFFFFFFFFFFFFFFFFFFFFFFFFFFFFFFFFFFFFFFFFFFFFFFFFFFFFFFFFFFFFFFFFFFFFFFFFFFFFFFFFFFFFFFFFFFFFFFFFFFFFFFFFFFFFFFFFFFFFFFFFFFFFFFFFFFFFFFFFFFFFFFFFFFFFFFFFFFFFFFFFFFFFFFFFFFFFFFFFFFFFFFFFFFFFFFFFFFFFFFFFFFFFFFFFFFFFFFFFFFFFFFFFFFFFFFFFFFFFFFFFFFFFFFFFFFFFFFFFFFFFFFFFFFFFFFFFFFFFFFFFFFFFFFFFFFFFFFFFFFFFFFFFFFFFFFFFFFFFFFFFFFFFFFFFFFFFFFFFFFFFFFFFFFFFFFFFFFFFFFFFFFFFFFFFFFFFFFFFFFFFFFFFFFFFFFFFFFFFFFFFFFFFFFFFFFFFFFFFFFFFFFFFFFFFFFFFFFFFFFFFFFFFFFFFFFFFFFFFFFFFFFFFFFFFFFFFFFFFFFFFFFFFFFFFFFFFFFFFFFFFFFFFFFFFFFFFFFFFFFFFFFFFFFFFFFFFFFFFFFFFFFFFFFFFFFFFFFFFFFFFFFFFFFFFFFFFFFFFFFFFFFFFFFFFFFFFFFFFDF2FEF9F3FBFFFFFEFEFEFFFEFCFEF8FEFCFEF9FBFEFDFDFFFEFEFFFEFBFDFEFFFFFFFFFFFFFFFFFFFFFFFFFFFFFFFFFFFFFFFFFFFFFFFFFFFFFFFFFFFFFFFFFFFFFFFFFFFFFFFFFFFFFFFFFFFFFFFFFFFFFFFFFFFFFFFFFFFFFFFFFFFFFFFFFFFFFFFFFFFFFFFFFFFFFFFFFFFFFFFFFFFFFFFFFFFFFFFFFFFFFFFFFFFFFFFFFFFFFFFFFFFFFFFFFFFFFFFFFFFFFFFFFFFFFFFFFFFFFFFFFFFFFFFFFFFFFFFFFFFFFFFFFFFFFFFFFFFFFFFFFFFFFFFFFFFFFFFFFFFFFFFFFFFFFFFFFFFFFFFFFFFFFFFFFFFFFFFFFFFFFFFFFFFFFFFFFFFFFFFFFFFFFFFFFFFFFFFFFFFFFFFFFFFFFFFFFFFFFFFFFFFFFFFFFFFFFFFFFFFFFFFFFFFFFFFFFFFFFFFFFFFFFFFDEFF2F6F1F6F8F9FBFDFBFCFBF1FDF6F5FAFAF6F9FDFDFCFDFAF9FCFEFAFCFBFBFDFFFFFFFFFFFFFFFFFFFFFFFFFFFFFFFFFFFFFFFFFFFFFFFFFFFFFFFFFFFFFFFFFFFFFFFFFFFFFFFFFFFFFFFFFFFFFFFFFFFFFFFFFFFFFFFFFFFFFFFFFFFFFFFFFFFFFFFFFFFFFFFFFFFFFFFFFFFFFFFFFFFFFFFFFFFFFFFFFFFFFFFFFFFFFFFFFFFFFFFFFFFFFFFFFFFFFFFFFFFFFFFFFFFFFFFFFFFFFFFFFFFFFFFFFFFFFFFFFFFFFFFFFFFFFFFFFFFFFFFFFFFFFFFFFFFFFFFFFFFFFFFFFFFFFFFFFFFFFFFFFFFFFFFFFFFFFFFFFFFFFFFFFFFFFFFFFFFFFFFFFFFFFFFFFFFFFFFFFFFFFFFFFFFFFFFFFFFFFFFFFFFFFFFFFFFFFFFFFFFFFFFFFFFFFFFFFFFFFFFFFEFDFBECF7FCF6F4F5F9F9F9FAF9EFFCF8F1FCF9FCFCFDFBFAFBF8F6F9FEFEFEFAF9FBFDFFFFFFFFFFFFFFFFFFFFFFFFFFFFFFFFFFFFFFFFFFFFFFFFFFFFFFFFFFFFFFFFFFFFFFFFFFFFFFFFFFFFFFFFFFFFFFFFFFFFFFFFFFFFFFFFFFFFFFFFFFFFFFFFFFFFFFFFFFFFFFFFFFFFFFFFFFFFFFFFFFFFFFFFFFFFFFFFFFFFFFFFFFFFFFFFFFFFFFFFFFFFFFFFFFFFFFFFFFFFFFFFFFFFFFFFFFFFFFFFFFFFFFFFFFFFFFFFFFFFFFFFFFFFFFFFFFFFFFFFFFFFFFFFFFFFFFFFFFFFFFFFFFFFFFFFFFFFFFFFFFFFFFFFFFFFFFFFFFFFFFFFFFFFFFFFFFFFFFFFFFFFFFFFFFFFFFFFFFFFFFFFFFFFFFFFFFFFFFFFFFFFFFFFFFFFFFFFFFFFFFFFFFFFFFFFFFFFFFFEFAF8FDF4F9F9F6FAFBF6F8FBF6FEFEF1FAFEFDFBFCFBF6F5FCFEFEFBFDFCFAFDFFFAFBFFFFFFFFFFFFFFFFFFFFFFFFFFFFFFFFFFFFFFFFFFFFFFFFFFFFFFFFFFFFFFFFFFFFFFFFFFFFFFFFFFFFFFFFFFFFFFFFFFFFFFFFFFFFFFFFFFFFFFFFFFFFFFFFFFFFFFFFFFFFFFFFFFFFFFFFFFFFFFFFFFFFFFFFFFFFFFFFFFFFFFFFFFFFFFFFFFFFFFFFFFFFFFFFFFFFFFFFFFFFFFFFFFFFFFFFFFFFFFFFFFFFFFFFFFFFFFFFFFFFFFFFFFFFFFFFFFFFFFFFFFFFFFFFFFFFFFFFFFFFFFFFFFFFFFFFFFFFFFFFFFFFFFFFFFFFFFFFFFFFFFFFFFFFFFFFFFFFFFFFFFFFFFFFFFFFFFFFFFFFFFFFFFFFFFFFFFFFFFFFFFFFFFFFFFFFFFFFFFFFFFFFFFFFFFFFFFFFFFFEFAFEF9FDFBFCECF2FEF8F5FAF5FBFEFEF4FAF6F6FAF7F9FDFBFCFBFDFEF8F3FCFEFAFBFFFFFFFFFFFFFFFFFFFFFFFFFFFFFFFFFFFFFFFFFFFFFFFFFFFFFFFFFFFFFFFFFFFFFFFFFFFFFFFFFFFFFFFFFFFFFFFFFFFFFFFFFFFFFFFFFFFFFFFFFFFFFFFFFFFFFFFFFFFFFFFFFFFFFFFFFFFFFFFFFFFFFFFFFFFFFFFFFFFFFFFFFFFFFFFFFFFFFFFFFFFFFFFFFFFFFFFFFFFFFFFFFFFFFFFFFFFFFFFFFFFFFFFFFFFFFFFFFFFFFFFFFFFFFFFFFFFFFFFFFFFFFFFFFFFFFFFFFFFFFFFFFFFFFFFFFFFFFFFFFFFFFFFFFFFFFFFFFFFFFFFFFFFFFFFFFFFFFFFFFFFFFFFFFFFFFFFFFFFFFFFFFFFFFFFFFFFFFFFFFFFFFFFFFFFFFFFFFFFFFFFFFFFFFFFFFFFFFFFEFEFDF2FCFAF7FAFBFBFCFAEFE5FAFCFAFEF4DEF0FEF2F2F7F0F7F1FBFEFEFAFDFDFEFFFFFFFFFFFFFFFFFFFFFFFFFFFFFFFFFFFFFFFFFFFFFFFFFFFFFFFFFFFFFFFFFFFFFFFFFFFFFFFFFFFFFFFFFFFFFFFFFFFFFFFFFFFFFFFFFFFFFFFFFFFFFFFFFFFFFFFFFFFFFFFFFFFFFFFFFFFFFFFFFFFFFFFFFFFFFFFFFFFFFFFFFFFFFFFFFFFFFFFFFFFFFFFFFFFFFFFFFFFFFFFFFFFFFFFFFFFFFFFFFFFFFFFFFFFFFFFFFFFFFFFFFFFFFFFFFFFFFFFFFFFFFFFFFFFFFFFFFFFFFFFFFFFFFFFFFFFFFFFFFFFFFFFFFFFFFFFFFFFFFFFFFFFFFFFFFFFFFFFFFFFFFFFFFFFFFFFFFFFFFFFFFFFFFFFFFFFFFFFFFFFFFFFFFFFFFFFFFFFFFFFFFFFFFFFFFFFFFFFFFDF9FDFDFEF9F7EBF1FBF5F4BF3FF3F1D5EFFC65A8F2FEFDFDFEFEFEFEF7FAFDFEF6FAFFFFFFFFFFFFFFFFFFFFFFFFFFFFFFFFFFFFFFFFFFFFFFFFFFFFFFFFFFFFFFFFFFFFFFFFFFFFFFFFFFFFFFFFFFFFFFFFFFFFFFFFFFFFFFFFFFFFFFFFFFFFFFFFFFFFFFFFFFFFFFFFFFFFFFFFFFFFFFFFFFFFFFFFFFFFFFFFFFFFFFFFFFFFFFFFFFFFFFFFFFFFFFFFFFFFFFFFFFFFFFFFFFFFFFFFFFFFFFFFFFFFFFFFFFFFFFFFFFFFFFFFFFFFFFFFFFFFFFFFFFFFFFFFFFFFFFFFFFFFFFFFFFFFFFFFFFFFFFFFFFFFFFFFFFFFFFFFFFFFFFFFFFFFFFFFFFFFFFFFFFFFFFFFFFFFFFFFFFFFFFFFFFFFFFFFFFFFFFFFFFFFFFFFFFFFFFFFFFFFFFFFFFFFFFFFFFFFFFFFFCEEFDF1F6FCFE8189E0A971E5BA1E1A0516C5E85896F7F9F7FAEDEBFEF7EDFCFEF6FAFFFFFFFFFFFFFFFFFFFFFFFFFFFFFFFFFFFFFFFFFFFFFFFFFFFFFFFFFFFFFFFFFFFFFFFFFFFFFFFFFFFFFFFFFFFFFFFFFFFFFFFFFFFFFFFFFFFFFFFFFFFFFFFFFFFFFFFFFFFFFFFFFFFFFFFFFFFFFFFFFFFFFFFFFFFFFFFFFFFFFFFFFFFFFFFFFFFFFFFFFFFFFFFFFFFFFFFFFFFFFFFFFFFFFFFFFFFFFFFFFFFFFFFFFFFFFFFFFFFFFFFFFFFFFFFFFFFFFFFFFFFFFFFFFFFFFFFFFFFFFFFFFFFFFFFFFFFFFFFFFFFFFFFFFFFFFFFFFFFFFFFFFFFFFFFFFFFFFFFFFFFFFFFFFFFFFFFFFFFFFFFFFFFFFFFFFFFFFFFFFFFFFFFFFFFFFFFFFFFFFFFFFFFFFFFFFFFFFFFFFEFBEDFB719CE6ED872F140A1B21020213031B19131D1E308FE79F74EAFDF3F8FDF8FBFFFFFFFFFFFFFFFFFFFFFFFFFFFFFFFFFFFFFFFFFFFFFFFFFFFFFFFFFFFFFFFFFFFFFFFFFFFFFFFFFFFFFFFFFFFFFFFFFFFFFFFFFFFFFFFFFFFFFFFFFFFFFFFFFFFFFFFFFFFFFFFFFFFFFFFFFFFFFFFFFFFFFFFFFFFFFFFFFFFFFFFFFFFFFFFFFFFFFFFFFFFFFFFFFFFFFFFFFFFFFFFFFFFFFFFFFFFFFFFFFFFFFFFFFFFFFFFFFFFFFFFFFFFFFFFFFFFFFFFFFFFFFFFFFFFFFFFFFFFFFFFFFFFFFFFFFFFFFFFFFFFFFFFFFFFFFFFFFFFFFFFFFFFFFFFFFFFFFFFFFFFFFFFFFFFFFFFFFFFFFFFFFFFFFFFFFFFFFFFFFFFFFFFFFFFFFFFFFFFFFFFFFCFDFFFFFFFEFEFFFFFDF6FDD57A3C2D1904000002020002030B0300020201041A4065937EA77F96988A95CEF8FCFBFFFDF8FFFFFFF6FBFCFDFFFAF9FDFFFFFFFFFFFFFFFFFFFFFFFFFFFFFFFFFFFFFFFFFFFFFFFFFFFFFFFFFFFFFFFFFFFFFFFFFFFFFFFFFFFFFFFFFFFFFFFFFFFFFFFFFFFFFFFFFFFFFFFFFFFFFFFFFFFFFFFFFFFFFFFFFFFFFFFFFFFFFFFFFFFFFFFFFFFFFFFFFFFFFFFFFFFFFFFFFFFFFFFFFFFFFFFFFFFFFFFFFFFFFFFFFFFFFFFFFFFFFFFFFFFFFFFFFFFFFFFFFFFFFFFFFFFFFFFFFFFFFFFFFFFFFFFFFFFFFFFFFFFFFFFFFFFFFFFFFFFFFFFFFFFFFFFFFFFFFFFFFFFFFFFFFFFFFFFFFFFFFFFFFFFFFFFFFFFFFFFFFFFFFFFFFFFFFCFDFFFFFEFCFCFDFDF7FCF789965A1706060001020909061204000E060002000108172C1A251A2670697AC9F5B886CDD56EC4EEE9F8F7FBFCFCFDFCFDFFFFFFFFFFFFFFFFFFFFFFFFFFFFFFFFFFFFFFFFFFFFFFFFFFFFFFFFFFFFFFFFFFFFFFFFFFFFFFFFFFFFFFFFFFFFFFFFFFFFFFFFFFFFFFFFFFFFFFFFFFFFFFFFFFFFFFFFFFFFFFFFFFFFFFFFFFFFFFFFFFFFFFFFFFFFFFFFFFFFFFFFFFFFFFFFFFFFFFFFFFFFFFFFFFFFFFFFFFFFFFFFFFFFFFFFFFFFFFFFFFFFFFFFFFFFFFFFFFFFFFFFFFFFFFFFFFFFFFFFFFFFFFFFFFFFFFFFFFFFFFFFFFFFFFFFFFFFFFFFFFFFFFFFFFFFFFFFFFFFFFFFFFFFFFFFFFFFFFFFFFFFFFFFFFFFFFFFFFFFFFFFFFFFFEFEFFFEFDFBFBFDFDF9F1FA6D40230D04080E0A2C8A1CAEA519178D3A0202000003080300070C03222B748B9C8C6AA0581E425351D2F8FAF5F8FFFEFEFFFFFFFFFFFFFFFFFFFFFFFFFFFFFFFFFFFFFFFFFFFFFFFFFFFFFFFFFFFFFFFFFFFFFFFFFFFFFFFFFFFFFFFFFFFFFFFFFFFFFFFFFFFFFFFFFFFFFFFFFFFFFFFFFFFFFFFFFFFFFFFFFFFFFFFFFFFFFFFFFFFFFFFFFFFFFFFFFFFFFFFFFFFFFFFFFFFFFFFFFFFFFFFFFFFFFFFFFFFFFFFFFFFFFFFFFFFFFFFFFFFFFFFFFFFFFFFFFFFFFFFFFFFFFFFFFFFFFFFFFFFFFFFFFFFFFFFFFFFFFFFFFFFFFFFFFFFFFFFFFFFFFFFFFFFFFFFFFFFFFFFFFFFFFFFFFFFFFFFFFFFFFFFFFFFFFFFFFFFFFFFFFFFFFFFFFFFEFDFCFCFEEFA3B2A1340D02080541611F46C737F267989FD454060808001045640004395C1506637751352E3B1405050504AA8DB9ECFBFDFDFFFFFFFFFFFFFFFFFFFFFFFFFFFFFFFFFFFFFFFFFFFFFFFFFFFFFFFFFFFFFFFFFFFFFFFFFFFFFFFFFFFFFFFFFFFFFFFFFFFFFFFFFFFFFFFFFFFFFFFFFFFFFFFFFFFFFFFFFFFFFFFFFFFFFFFFFFFFFFFFFFFFFFFFFFFFFFFFFFFFFFFFFFFFFFFFFFFFFFFFFFFFFFFFFFFFFFFFFFFFFFFFFFFFFFFFFFFFFFFFFFFFFFFFFFFFFFFFFFFFFFFFFFFFFFFFFFFFFFFFFFFFFFFFFFFFFFFFFFFFFFFFFFFFFFFFFFFFFFFFFFFFFFFFFFFFFFFFFFFFFFFFFFFFFFFFFFFFFFFFFFFFFFFFFFFFFFFFFFFFFFFFFFFFFFFFFFFEFEFEFEFEFFF3B488340805060908327187A8E2A2F58BDBF3897E370608030C354D0200274712002D3D1E05020106050007054A2145B0F0F9FBFFFFFFFFFFFFFFFFFFFFFFFFFFFFFFFFFFFFFFFFFFFFFFFFFFFFFFFFFFFFFFFFFFFFFFFFFFFFFFFFFFFFFFFFFFFFFFFFFFFFFFFFFFFFFFFFFFFFFFFFFFFFFFFFFFFFFFFFFFFFFFFFFFFFFFFFFFFFFFFFFFFFFFFFFFFFFFFFFFFFFFFFFFFFFFFFFFFFFFFFFFFFFFFFFFFFFFFFFFFFFFFFFFFFFFFFFFFFFFFFFFFFFFFFFFFFFFFFFFFFFFFFFFFFFFFFFFFFFFFFFFFFFFFFFFFFFFFFFFFFFFFFFFFFFFFFFFFFFFFFFFFFFFFFFFFFFFFFFFFFFFFFFFFFFFFFFFFFFFFFFFFFFFFFFFFFFFFFFFFFFFFFFFFFFFFFFEFEFEFFFFFFFEFDEE7A080301453D1405329CF5FBF6F9F587D63A5935040F11050300070004010403040203040204010500081307010348BAFAFBFCFFFFFFFFFFFFFFFFFFFFFFFFFFFFFFFFFFFFFFFFFFFFFFFFFFFFFFFFFFFFFFFFFFFFFFFFFFFFFFFFFFFFFFFFFFFFFFFFFFFFFFFFFFFFFFFFFFFFFFFFFFFFFFFFFFFFFFFFFFFFFFFFFFFFFFFFFFFFFFFFFFFFFFFFFFFFFFFFFFFFFFFFFFFFFFFFFFFFFFFFFFFFFFFFFFFFFFFFFFFFFFFFFFFFFFFFFFFFFFFFFFFFFFFFFFFFFFFFFFFFFFFFFFFFFFFFFFFFFFFFFFFFFFFFFFFFFFFFFFFFFFFFFFFFFFFFFFFFFFFFFFFFFFFFFFFFFFFFFFFFFFFFFFFFFFFFFFFFFFFFFFFFFFFFFFFFFFFFFFFFFFFFFFFFFEFEFEFEFFFFFEFBF9E9640A0601584C18072C8DEBEDEAE099667B23030411393D07040006000700040E08040202060600000301581F0401057FF1F8FAFFFFFFFFFFFFFFFFFFFFFFFFFFFFFFFFFFFFFFFFFFFFFFFFFFFFFFFFFFFFFFFFFFFFFFFFFFFFFFFFFFFFFFFFFFFFFFFFFFFFFFFFFFFFFFFFFFFFFFFFFFFFFFFFFFFFFFFFFFFFFFFFFFFFFFFFFFFFFFFFFFFFFFFFFFFFFFFFFFFFFFFFFFFFFFFFFFFFFFFFFFFFFFFFFFFFFFFFFFFFFFFFFFFFFFFFFFFFFFFFFFFFFFFFFFFFFFFFFFFFFFFFFFFFFFFFFFFFFFFFFFFFFFFFFFFFFFFFFFFFFFFFFFFFFFFFFFFFFFFFFFFFFFFFFFFFFFFFFFFFFFFFFFFFFFFFFFFFFFFFFFFFFFFFFFFFFFFFFFFFFFFFFFFFFEFEFEFEFFFEFBF5F0DC560D060201010F0F5DCDB8C6A58505891E05101725565A0E0802000516131436382418170F05110007027E280401024DB4EAFCFFFFFFFFFFFFFFFFFFFFFFFFFFFFFFFFFFFFFFFFFFFFFFFFFFFFFFFFFFFFFFFFFFFFFFFFFFFFFFFFFFFFFFFFFFFFFFFFFFFFFFFFFFFFFFFFFFFFFFFFFFFFFFFFFFFFFFFFFFFFFFFFFFFFFFFFFFFFFFFFFFFFFFFFFFFFFFFFFFFFFFFFFFFFFFFFFFFFFFFFFFFFFFFFFFFFFFFFFFFFFFFFFFFFFFFFFFFFFFFFFFFFFFFFFFFFFFFFFFFFFFFFFFFFFFFFFFFFFFFFFFFFFFFFFFFFFFFFFFFFFFFFFFFFFFFFFFFFFFFFFFFFFFFFFFFFFFFFFFFFFFFFFFFFFFFFFFFFFFFFFFFFFFFFFFFFFFFFFFFFFFFFFFFFFFFCFDFEFEFDF8ECCF8B4B07000303052D30497E688A5203030F0C124C5E5E6159370406021C6060627E8C716769442B4F16060201000B08010853C7FFFFFFFFFFFFFFFFFFFFFFFFFFFFFFFFFFFFFFFFFFFFFFFFFFFFFFFFFFFFFFFFFFFFFFFFFFFFFFFFFFFFFFFFFFFFFFFFFFFFFFFFFFFFFFFFFFFFFFFFFFFFFFFFFFFFFFFFFFFFFFFFFFFFFFFFFFFFFFFFFFFFFFFFFFFFFFFFFFFFFFFFFFFFFFFFFFFFFFFFFFFFFFFFFFFFFFFFFFFFFFFFFFFFFFFFFFFFFFFFFFFFFFFFFFFFFFFFFFFFFFFFFFFFFFFFFFFFFFFFFFFFFFFFFFFFFFFFFFFFFFFFFFFFFFFFFFFFFFFFFFFFFFFFFFFFFFFFFFFFFFFFFFFFFFFFFFFFFFFFFFFFFFFFFFFFFFFFFFFFFFFFFFFFFFFFF9FAFEFDFCF7DE9722000900000607545F194FAC26980A005320549BB3D3EEC45E0A010744A8AFB08980A1CBE5C98C946C12070006412303305C99FFFFFFFFFFFFFFFFFFFFFFFFFFFFFFFFFFFFFFFFFFFFFFFFFFFFFFFFFFFFFFFFFFFFFFFFFFFFFFFFFFFFFFFFFFFFFFFFFFFFFFFFFFFFFFFFFFFFFFFFFFFFFFFFFFFFFFFFFFFFFFFFFFFFFFFFFFFFFFFFFFFFFFFFFFFFFFFFFFFFFFFFFFFFFFFFFFFFFFFFFFFFFFFFFFFFFFFFFFFFFFFFFFFFFFFFFFFFFFFFFFFFFFFFFFFFFFFFFFFFFFFFFFFFFFFFFFFFFFFFFFFFFFFFFFFFFFFFFFFFFFFFFFFFFFFFFFFFFFFFFFFFFFFFFFFFFFFFFFFFFFFFFFFFFFFFFFFFFFFFFFFFFFFFFFFFFFFFFFFFFFFFFFFFFFFFFEFAFAFDF9E7D4CE79000000020A04000001387715000100A886933305619B5101010A00000B02004B4F1B3081C4CA6CA61A0200248C380364D5F9FEFFFFFFFFFFFFFFFFFFFFFFFFFFFFFFFFFFFFFFFFFFFFFFFFFFFFFFFFFFFFFFFFFFFFFFFFFFFFFFFFFFFFFFFFFFFFFFFFFFFFFFFFFFFFFFFFFFFFFFFFFFFFFFFFFFFFFFFFFFFFFFFFFFFFFFFFFFFFFFFFFFFFFFFFFFFFFFFFFFFFFFFFFFFFFFFFFFFFFFFFFFFFFFFFFFFFFFFFFFFFFFFFFFFFFFFFFFFFFFFFFFFFFFFFFFFFFFFFFFFFFFFFFFFFFFFFFFFFFFFFFFFFFFFFFFFFFFFFFFFFFFFFFFFFFFFFFFFFFFFFFFFFFFFFFFFFFFFFFFFFFFFFFFFFFFFFFFFFFFFFFFFFFFFFFFFFFFFFFFFFFFFFFFFFFFFAFCFEFDF7DDA96E5A0207010000010000090F0001081201008F2B080A04020C1101010802000A08080702050B2B829513001206000000090870FBFCFFFFFFFFFFFFFFFFFFFFFFFFFFFFFFFFFFFFFFFFFFFFFFFFFFFFFFFFFFFFFFFFFFFFFFFFFFFFFFFFFFFFFFFFFFFFFFFFFFFFFFFFFFFFFFFFFFFFFFFFFFFFFFFFFFFFFFFFFFFFFFFFFFFFFFFFFFFFFFFFFFFFFFFFFFFFFFFFFFFFFFFFFFFFFFFFFFFFFFFFFFFFFFFFFFFFFFFFFFFFFFFFFFFFFFFFFFFFFFFFFFFFFFFFFFFFFFFFFFFFFFFFFFFFFFFFFFFFFFFFFFFFFFFFFFFFFFFFFFFFFFFFFFFFFFFFFFFFFFFFFFFFFFFFFFFFFFFFFFFFFFFFFFFFFFFFFFFFFFFFFFFFFFFFFFFFFFFFFFFFFFFFFFFFFFFAFBFAF5F7F0B74805000003412F120608080101174F04670A09120701082C33061E3F1C533127320F0302041B2C150B050001000008030030A4FAF5FFFFFFFFFFFFFFFFFFFFFFFFFFFFFFFFFFFFFFFFFFFFFFFFFFFFFFFFFFFFFFFFFFFFFFFFFFFFFFFFFFFFFFFFFFFFFFFFFFFFFFFFFFFFFFFFFFFFFFFFFFFFFFFFFFFFFFFFFFFFFFFFFFFFFFFFFFFFFFFFFFFFFFFFFFFFFFFFFFFFFFFFFFFFFFFFFFFFFFFFFFFFFFFFFFFFFFFFFFFFFFFFFFFFFFFFFFFFFFFFFFFFFFFFFFFFFFFFFFFFFFFFFFFFFFFFFFFFFFFFFFFFFFFFFFFFFFFFFFFFFFFFFFFFFFFFFFFFFFFFFFFFFFFFFFFFFFFFFFFFFFFFFFFFFFFFFFFFFFFFFFFFFFFFFFFFFFFFFFFFFFFFFFFFF2F9FEFEFDF7CC650F000B0000594914000128450F2293069610000202010E5B8B5E717981E37481672103040434662A01060B01020007030078EDFAFDFFFFFFFFFFFFFFFFFFFFFFFFFFFFFFFFFFFFFFFFFFFFFFFFFFFFFFFFFFFFFFFFFFFFFFFFFFFFFFFFFFFFFFFFFFFFFFFFFFFFFFFFFFFFFFFFFFFFFFFFFFFFFFFFFFFFFFFFFFFFFFFFFFFFFFFFFFFFFFFFFFFFFFFFFFFFFFFFFFFFFFFFFFFFFFFFFFFFFFFFFFFFFFFFFFFFFFFFFFFFFFFFFFFFFFFFFFFFFFFFFFFFFFFFFFFFFFFFFFFFFFFFFFFFFFFFFFFFFFFFFFFFFFFFFFFFFFFFFFFFFFFFFFFFFFFFFFFFFFFFFFFFFFFFFFFFFFFFFFFFFFFFFFFFFFFFFFFFFFFFFFFFFFFFFFFFFFFFFFFFFFFFFFFFFCFDFBF3F6FDE29130010905001F140405043C650807241B211102000604071B5EB8BA97E0AB4368473C150107131709000205005255163756AFFDFBFBFFFFFFFFFFFFFFFFFFFFFFFFFFFFFFFFFFFFFFFFFFFFFFFFFFFFFFFFFFFFFFFFFFFFFFFFFFFFFFFFFFFFFFFFFFFFFFFFFFFFFFFFFFFFFFFFFFFFFFFFFFFFFFFFFFFFFFFFFFFFFFFFFFFFFFFFFFFFFFFFFFFFFFFFFFFFFFFFFFFFFFFFFFFFFFFFFFFFFFFFFFFFFFFFFFFFFFFFFFFFFFFFFFFFFFFFFFFFFFFFFFFFFFFFFFFFFFFFFFFFFFFFFFFFFFFFFFFFFFFFFFFFFFFFFFFFFFFFFFFFFFFFFFFFFFFFFFFFFFFFFFFFFFFFFFFFFFFFFFFFFFFFFFFFFFFFFFFFFFFFFFFFFFFFFFFFFFFFFFFFFFFFFFFFFFFCFCF8F6F9ECA32A010201020401000101121E050F0B090B050400020706051C62BBC39272161A4746200C0805051139123E14B98B66ACD5E9F4F1FCFFFFFFFFFFFFFFFFFFFFFFFFFFFFFFFFFFFFFFFFFFFFFFFFFFFFFFFFFFFFFFFFFFFFFFFFFFFFFFFFFFFFFFFFFFFFFFFFFFFFFFFFFFFFFFFFFFFFFFFFFFFFFFFFFFFFFFFFFFFFFFFFFFFFFFFFFFFFFFFFFFFFFFFFFFFFFFFFFFFFFFFFFFFFFFFFFFFFFFFFFFFFFFFFFFFFFFFFFFFFFFFFFFFFFFFFFFFFFFFFFFFFFFFFFFFFFFFFFFFFFFFFFFFFFFFFFFFFFFFFFFFFFFFFFFFFFFFFFFFFFFFFFFFFFFFFFFFFFFFFFFFFFFFFFFFFFFFFFFFFFFFFFFFFFFFFFFFFFFFFFFFFFFFFFFFFFFFFFFFFFFFFFFFFFDFDFEFDFCF9E5AB6024070001000502000123390E07060D05044A1B01080E2C3C44B3D4542303051C413B18060502309927B014CF95BEE7F9FCF8F8FEFFFFFFFFFFFFFFFFFFFFFFFFFFFFFFFFFFFFFFFFFFFFFFFFFFFFFFFFFFFFFFFFFFFFFFFFFFFFFFFFFFFFFFFFFFFFFFFFFFFFFFFFFFFFFFFFFFFFFFFFFFFFFFFFFFFFFFFFFFFFFFFFFFFFFFFFFFFFFFFFFFFFFFFFFFFFFFFFFFFFFFFFFFFFFFFFFFFFFFFFFFFFFFFFFFFFFFFFFFFFFFFFFFFFFFFFFFFFFFFFFFFFFFFFFFFFFFFFFFFFFFFFFFFFFFFFFFFFFFFFFFFFFFFFFFFFFFFFFFFFFFFFFFFFFFFFFFFFFFFFFFFFFFFFFFFFFFFFFFFFFFFFFFFFFFFFFFFFFFFFFFFFFFFFFFFFFFFFFFFFFFFFFFFFFDFBFDFFFDFAF0D09C3205000300362A090134540900020B00076A2803030E3C4B3EB0CF49000005013F4519021D4783DB7AF05DE0D0F0F8FAFDFDFEFEFFFFFFFFFFFFFFFFFFFFFFFFFFFFFFFFFFFFFFFFFFFFFFFFFFFFFFFFFFFFFFFFFFFFFFFFFFFFFFFFFFFFFFFFFFFFFFFFFFFFFFFFFFFFFFFFFFFFFFFFFFFFFFFFFFFFFFFFFFFFFFFFFFFFFFFFFFFFFFFFFFFFFFFFFFFFFFFFFFFFFFFFFFFFFFFFFFFFFFFFFFFFFFFFFFFFFFFFFFFFFFFFFFFFFFFFFFFFFFFFFFFFFFFFFFFFFFFFFFFFFFFFFFFFFFFFFFFFFFFFFFFFFFFFFFFFFFFFFFFFFFFFFFFFFFFFFFFFFFFFFFFFFFFFFFFFFFFFFFFFFFFFFFFFFFFFFFFFFFFFFFFFFFFFFFFFFFFFFFFFFFFFFFFFFBFAFCF9F8FAEBC5913F000004054B4011001B2D000006050004331403000622456FB6B5280200020622220B032664919F9CF3CAF7FEFEFEFEFEFEFEFEFFFFFFFFFFFFFFFFFFFFFFFFFFFFFFFFFFFFFFFFFFFFFFFFFFFFFFFFFFFFFFFFFFFFFFFFFFFFFFFFFFFFFFFFFFFFFFFFFFFFFFFFFFFFFFFFFFFFFFFFFFFFFFFFFFFFFFFFFFFFFFFFFFFFFFFFFFFFFFFFFFFFFFFFFFFFFFFFFFFFFFFFFFFFFFFFFFFFFFFFFFFFFFFFFFFFFFFFFFFFFFFFFFFFFFFFFFFFFFFFFFFFFFFFFFFFFFFFFFFFFFFFFFFFFFFFFFFFFFFFFFFFFFFFFFFFFFFFFFFFFFFFFFFFFFFFFFFFFFFFFFFFFFFFFFFFFFFFFFFFFFFFFFFFFFFFFFFFFFFFFFFFFFFFFFFFFFFFFFFFFFFFFFFFFAF7FBF8F6F8E2AB7B410006030623250E02060700020A02020200020204060938899172071548301002030205255D766554A3F0FCFEFEFEFEFEFEFEFEFFFFFFFFFFFFFFFFFFFFFFFFFFFFFFFFFFFFFFFFFFFFFFFFFFFFFFFFFFFFFFFFFFFFFFFFFFFFFFFFFFFFFFFFFFFFFFFFFFFFFFFFFFFFFFFFFFFFFFFFFFFFFFFFFFFFFFFFFFFFFFFFFFFFFFFFFFFFFFFFFFFFFFFFFFFFFFFFFFFFFFFFFFFFFFFFFFFFFFFFFFFFFFFFFFFFFFFFFFFFFFFFFFFFFFFFFFFFFFFFFFFFFFFFFFFFFFFFFFFFFFFFFFFFFFFFFFFFFFFFFFFFFFFFFFFFFFFFFFFFFFFFFFFFFFFFFFFFFFFFFFFFFFFFFFFFFFFFFFFFFFFFFFFFFFFFFFFFFFFFFFFFFFFFFFFFFFFFFFFFFFFFFFFFFDF5FAFFF6F3EBAE622101090300020705030506060107010303030200050A072F745822082E925B11000404023791A7863E5FE7FCFEFEFEFEFEFEFEFEFFFFFFFFFFFFFFFFFFFFFFFFFFFFFFFFFFFFFFFFFFFFFFFFFFFFFFFFFFFFFFFFFFFFFFFFFFFFFFFFFFFFFFFFFFFFFFFFFFFFFFFFFFFFFFFFFFFFFFFFFFFFFFFFFFFFFFFFFFFFFFFFFFFFFFFFFFFFFFFFFFFFFFFFFFFFFFFFFFFFFFFFFFFFFFFFFFFFFFFFFFFFFFFFFFFFFFFFFFFFFFFFFFFFFFFFFFFFFFFFFFFFFFFFFFFFFFFFFFFFFFFFFFFFFFFFFFFFFFFFFFFFFFFFFFFFFFFFFFFFFFFFFFFFFFFFFFFFFFFFFFFFFFFFFFFFFFFFFFFFFFFFFFFFFFFFFFFFFFFFFFFFFFFFFFFFFFFFFFFFFFFFFFFFF9F9F9FAF3EFE38C1509080000000402030402071A2508372E333333372D182270AB6324042458370F0101011463BEB88AB1B3FBFEFEFEFEFEFEFEFEFEFFFFFFFFFFFFFFFFFFFFFFFFFFFFFFFFFFFFFFFFFFFFFFFFFFFFFFFFFFFFFFFFFFFFFFFFFFFFFFFFFFFFFFFFFFFFFFFFFFFFFFFFFFFFFFFFFFFFFFFFFFFFFFFFFFFFFFFFFFFFFFFFFFFFFFFFFFFFFFFFFFFFFFFFFFFFFFFFFFFFFFFFFFFFFFFFFFFFFFFFFFFFFFFFFFFFFFFFFFFFFFFFFFFFFFFFFFFFFFFFFFFFFFFFFFFFFFFFFFFFFFFFFFFFFFFFFFFFFFFFFFFFFFFFFFFFFFFFFFFFFFFFFFFFFFFFFFFFFFFFFFFFFFFFFFFFFFFFFFFFFFFFFFFFFFFFFFFFFFFFFFFFFFFFFFFFFFFFFFFFFFFFFFFFF5F8F2E8E4E6CA6405011001051A13050A0C142847803CCBBCC6BBC5D3B66067BA9D597706125261350C01024AB7BD71A9F9F3FEFEFEFEFEFEFEFEFEFEFFFFFFFFFFFFFFFFFFFFFFFFFFFFFFFFFFFFFFFFFFFFFFFFFFFFFFFFFFFFFFFFFFFFFFFFFFFFFFFFFFFFFFFFFFFFFFFFFFFFFFFFFFFFFFFFFFFFFFFFFFFFFFFFFFFFFFFFFFFFFFFFFFFFFFFFFFFFFFFFFFFFFFFFFFFFFFFFFFFFFFFFFFFFFFFFFFFFFFFFFFFFFFFFFFFFFFFFFFFFFFFFFFFFFFFFFFFFFFFFFFFFFFFFFFFFFFFFFFFFFFFFFFFFFFFFFFFFFFFFFFFFFFFFFFFFFFFFFFFFFFFFFFFFFFFFFFFFFFFFFFFFFFFFFFFFFFFFFFFFFFFFFFFFFFFFFFFFFFFFFFFFFFFFFFFFFFFFFFFFFFFFFFFFFEEEE9CC98938406030103031671551A02023CAABA65B9EDFEE9F5FBEDF3D57E331D1C2809588C8D661F05114FB2EBE1EFFDFCFCFDFEFEFEFEFEFEFEFEFFFFFFFFFFFFFFFFFFFFFFFFFFFFFFFFFFFFFFFFFFFFFFFFFFFFFFFFFFFFFFFFFFFFFFFFFFFFFFFFFFFFFFFFFFFFFFFFFFFFFFFFFFFFFFFFFFFFFFFFFFFFFFFFFFFFFFFFFFFFFFFFFFFFFFFFFFFFFFFFFFFFFFFFFFFFFFFFFFFFFFFFFFFFFFFFFFFFFFFFFFFFFFFFFFFFFFFFFFFFFFFFFFFFFFFFFFFFFFFFFFFFFFFFFFFFFFFFFFFFFFFFFFFFFFFFFFFFFFFFFFFFFFFFFFFFFFFFFFFFFFFFFFFFFFFFFFFFFFFFFFFFFFFFFFFFFFFFFFFFFFFFFFFFFFFFFFFFFFFFFFFFFFFFFFFFFFFFFFFFFFFFFFFFFDD38BAAC2640C061806000000200F020710196BE6D0FEFAFDFEFFEDF6F07E1201000003236BD6A0300500577E79C0FDFDF9FAF5FDFEFEFEFEFEFEFEFEFFFFFFFFFFFFFFFFFFFFFFFFFFFFFFFFFFFFFFFFFFFFFFFFFFFFFFFFFFFFFFFFFFFFFFFFFFFFFFFFFFFFFFFFFFFFFFFFFFFFFFFFFFFFFFFFFFFFFFFFFFFFFFFFFFFFFFFFFFFFFFFFFFFFFFFFFFFFFFFFFFFFFFFFFFFFFFFFFFFFFFFFFFFFFFFFFFFFFFFFFFFFFFFFFFFFFFFFFFFFFFFFFFFFFFFFFFFFFFFFFFFFFFFFFFFFFFFFFFFFFFFFFFFFFFFFFFFFFFFFFFFFFFFFFFFFFFFFFFFFFFFFFFFFFFFFFFFFFFFFFFFFFFFFFFFFFFFFFFFFFFFFFFFFFFFFFFFFFFFFFFFFFFFFFFFFFFFFFFFFFFFFFFFFF9EAAE7B8F74080C000000060A020E060000549762EDFEFAFBFEFFF9F07A170101030100086BA17D7C2700115BBEF0FAFEFEFDFCFDFEFEFEFEFEFEFEFEFFFFFFFFFFFFFFFFFFFFFFFFFFFFFFFFFFFFFFFFFFFFFFFFFFFFFFFFFFFFFFFFFFFFFFFFFFFFFFFFFFFFFFFFFFFFFFFFFFFFFFFFFFFFFFFFFFFFFFFFFFFFFFFFFFFFFFFFFFFFFFFFFFFFFFFFFFFFFFFFFFFFFFFFFFFFFFFFFFFFFFFFFFFFFFFFFFFFFFFFFFFFFFFFFFFFFFFFFFFFFFFFFFFFFFFFFFFFFFFFFFFFFFFFFFFFFFFFFFFFFFFFFFFFFFFFFFFFFFFFFFFFFFFFFFFFFFFFFFFFFFFFFFFFFFFFFFFFFFFFFFFFFFFFFFFFFFFFFFFFFFFFFFFFFFFFFFFFFFFFFFFFFFFFFFFBF9FEFBF6FEFDF8FDFEFCEBBB6416030700000403010704050206144FC1FCFDFBFDF5FBFDF0680B01115485011264E88D2D1004034CBEF8FEFFFEF8FDFEFEFEFEFEFEFEFEFEFFFFFFFFFFFFFFFFFFFFFFFFFFFFFFFFFFFFFFFFFFFFFFFFFFFFFFFFFFFFFFFFFFFFFFFFFFFFFFFFFFFFFFFFFFFFFFFFFFFFFFFFFFFFFFFFFFFFFFFFFFFFFFFFFFFFFFFFFFFFFFFFFFFFFFFFFFFFFFFFFFFFFFFFFFFFFFFFFFFFFFFFFFFFFFFFFFFFFFFFFFFFFFFFFFFFFFFFFFFFFFFFFFFFFFFFFFFFFFFFFFFFFFFFFFFFFFFFFFFFFFFFFFFFFFFFFFFFFFFFFFFFFFFFFFFFFFFFFFFFFFFFFFFFFFFFFFFFFFFFFFFFFFFFFFFFFFFFFFFFFFFFFFFFFFFFFFFFFFFFFFFFFFFFFFFCFBFEFEFBF2FDFDF3F0D2A198600F0D0204020005000A704A10044ECCF0ECFCEDF6FC5A42451C030507182F0003118B9A3007525976D3FDFDFEFEFAFEFFFFFFFFFFFFFFFFFFFFFFFFFFFFFFFFFFFFFFFFFFFFFFFFFFFFFFFFFFFFFFFFFFFFFFFFFFFFFFFFFFFFFFFFFFFFFFFFFFFFFFFFFFFFFFFFFFFFFFFFFFFFFFFFFFFFFFFFFFFFFFFFFFFFFFFFFFFFFFFFFFFFFFFFFFFFFFFFFFFFFFFFFFFFFFFFFFFFFFFFFFFFFFFFFFFFFFFFFFFFFFFFFFFFFFFFFFFFFFFFFFFFFFFFFFFFFFFFFFFFFFFFFFFFFFFFFFFFFFFFFFFFFFFFFFFFFFFFFFFFFFFFFFFFFFFFFFFFFFFFFFFFFFFFFFFFFFFFFFFFFFFFFFFFFFFFFFFFFFFFFFFFFFFFFFFFFFFFFFFFFFFFFFFFFEFEFFFFFDEFF4FDFBE36B2D2F2409030116694510010228220D528389ECF6FEFAFCF9C434080B0701060C07030C52A03C0C01205CA9E7FDFAFAFEFCFEFFFFFFFFFFFFFFFFFFFFFFFFFFFFFFFFFFFFFFFFFFFFFFFFFFFFFFFFFFFFFFFFFFFFFFFFFFFFFFFFFFFFFFFFFFFFFFFFFFFFFFFFFFFFFFFFFFFFFFFFFFFFFFFFFFFFFFFFFFFFFFFFFFFFFFFFFFFFFFFFFFFFFFFFFFFFFFFFFFFFFFFFFFFFFFFFFFFFFFFFFFFFFFFFFFFFFFFFFFFFFFFFFFFFFFFFFFFFFFFFFFFFFFFFFFFFFFFFFFFFFFFFFFFFFFFFFFFFFFFFFFFFFFFFFFFFFFFFFFFFFFFFFFFFFFFFFFFFFFFFFFFFFFFFFFFFFFFFFFFFFFFFFFFFFFFFFFFFFFFFFFFFFFFFFFFFFFFFFFFFFFFFFFFFFFFBF6FAFEF9FAFCDD8E7E36060602040B093229080A0202080524708D82F7FDF7F2654F1B020304040604020B4989440C00090763CCF3FDFBFCFEFDFDFEFFFFFFFFFFFFFFFFFFFFFFFFFFFFFFFFFFFFFFFFFFFFFFFFFFFFFFFFFFFFFFFFFFFFFFFFFFFFFFFFFFFFFFFFFFFFFFFFFFFFFFFFFFFFFFFFFFFFFFFFFFFFFFFFFFFFFFFFFFFFFFFFFFFFFFFFFFFFFFFFFFFFFFFFFFFFFFFFFFFFFFFFFFFFFFFFFFFFFFFFFFFFFFFFFFFFFFFFFFFFFFFFFFFFFFFFFFFFFFFFFFFFFFFFFFFFFFFFFFFFFFFFFFFFFFFFFFFFFFFFFFFFFFFFFFFFFFFFFFFFFFFFFFFFFFFFFFFFFFFFFFFFFFFFFFFFFFFFFFFFFFFFFFFFFFFFFFFFFFFFFFFFFFFFFFFFFFFFFFFFFFFFFFFBF6F7FBF8F2FEFDD64844230006070206000708010601000603032B69A6FAFCFAF6A223702F000406020207062F42070801367AB7E4FAFDFCFDFEFDFBFEFFFFFFFFFFFFFFFFFFFFFFFFFFFFFFFFFFFFFFFFFFFFFFFFFFFFFFFFFFFFFFFFFFFFFFFFFFFFFFFFFFFFFFFFFFFFFFFFFFFFFFFFFFFFFFFFFFFFFFFFFFFFFFFFFFFFFFFFFFFFFFFFFFFFFFFFFFFFFFFFFFFFFFFFFFFFFFFFFFFFFFFFFFFFFFFFFFFFFFFFFFFFFFFFFFFFFFFFFFFFFFFFFFFFFFFFFFFFFFFFFFFFFFFFFFFFFFFFFFFFFFFFFFFFFFFFFFFFFFFFFFFFFFFFFFFFFFFFFFFFFFFFFFFFFFFFFFFFFFFFFFFFFFFFFFFFFFFFFFFFFFFFFFFFFFFFFFFFFFFFFFFFFFFFFFFFFFFFFFFFFFFFFFFCF9FBFEF9E786E67D190C06020706010201030000080100070704349DCD96F8FEF9F9907D28030603020604093C5905090265D3F5F6FDFDFAFAFEFEFDFEFFFFFFFFFFFFFFFFFFFFFFFFFFFFFFFFFFFFFFFFFFFFFFFFFFFFFFFFFFFFFFFFFFFFFFFFFFFFFFFFFFFFFFFFFFFFFFFFFFFFFFFFFFFFFFFFFFFFFFFFFFFFFFFFFFFFFFFFFFFFFFFFFFFFFFFFFFFFFFFFFFFFFFFFFFFFFFFFFFFFFFFFFFFFFFFFFFFFFFFFFFFFFFFFFFFFFFFFFFFFFFFFFFFFFFFFFFFFFFFFFFFFFFFFFFFFFFFFFFFFFFFFFFFFFFFFFFFFFFFFFFFFFFFFFFFFFFFFFFFFFFFFFFFFFFFFFFFFFFFFFFFFFFFFFFFFFFFFFFFFFFFFFFFFFFFFFFFFFFFFFFFFFFFFFFFFFFFFFFFFFFFFFFFBF8FAF5FCE67B7C1901060302050606103C32090007064436100234818288F9F6FAF9942B0410403808070009396106090666D3FDFFFFFDFBF8FDFFFFFFFFFFFFFFFFFFFFFFFFFFFFFFFFFFFFFFFFFFFFFFFFFFFFFFFFFFFFFFFFFFFFFFFFFFFFFFFFFFFFFFFFFFFFFFFFFFFFFFFFFFFFFFFFFFFFFFFFFFFFFFFFFFFFFFFFFFFFFFFFFFFFFFFFFFFFFFFFFFFFFFFFFFFFFFFFFFFFFFFFFFFFFFFFFFFFFFFFFFFFFFFFFFFFFFFFFFFFFFFFFFFFFFFFFFFFFFFFFFFFFFFFFFFFFFFFFFFFFFFFFFFFFFFFFFFFFFFFFFFFFFFFFFFFFFFFFFFFFFFFFFFFFFFFFFFFFFFFFFFFFFFFFFFFFFFFFFFFFFFFFFFFFFFFFFFFFFFFFFFFFFFFFFFFFFFFFFFFFFFFFFFFFFFFFAF7FAF1EF99E51B05020101010519364354490D0102075D481101021F66A6F8F4FEFA841A0115554C08040304000F002C3A68CAFCFDFDFFFDF9FCFEFDFEFFFFFFFFFFFFFFFFFFFFFFFFFFFFFFFFFFFFFFFFFFFFFFFFFFFFFFFFFFFFFFFFFFFFFFFFFFFFFFFFFFFFFFFFFFFFFFFFFFFFFFFFFFFFFFFFFFFFFFFFFFFFFFFFFFFFFFFFFFFFFFFFFFFFFFFFFFFFFFFFFFFFFFFFFFFFFFFFFFFFFFFFFFFFFFFFFFFFFFFFFFFFFFFFFFFFFFFFFFFFFFFFFFFFFFFFFFFFFFFFFFFFFFFFFFFFFFFFFFFFFFFFFFFFFFFFFFFFFFFFFFFFFFFFFFFFFFFFFFFFFFFFFFFFFFFFFFFFFFFFFFFFFFFFFFFFFFFFFFFFFFFFFFFFFFFFFFFFFFFFFFFFFFFFFFFEFDFDFEFBFCFEFEFFFCFAFDE5539617010202050714427E87472304020401000806010138896FF5FCFEFAF541060B070000020609050204619EADDEFDFCFBFFFEFCFCFEFCFEFFFFFFFFFFFFFFFFFFFFFFFFFFFFFFFFFFFFFFFFFFFFFFFFFFFFFFFFFFFFFFFFFFFFFFFFFFFFFFFFFFFFFFFFFFFFFFFFFFFFFFFFFFFFFFFFFFFFFFFFFFFFFFFFFFFFFFFFFFFFFFFFFFFFFFFFFFFFFFFFFFFFFFFFFFFFFFFFFFFFFFFFFFFFFFFFFFFFFFFFFFFFFFFFFFFFFFFFFFFFFFFFFFFFFFFFFFFFFFFFFFFFFFFFFFFFFFFFFFFFFFFFFFFFFFFFFFFFFFFFFFFFFFFFFFFFFFFFFFFFFFFFFFFFFFFFFFFFFFFFFFFFFFFFFFFFFFFFFFFFFFFFFFFFFFFFFFFFFFFFFFFFFFFFFEFDFAFBFEF5F9FCFAFBF6F3FBEEA90C0402010D2827467E8086A8460D0100050203040203052BB6F5F2F9FCF56C1F2F2104080902020F1A3069DBFBF9FBFFFFFFFFFFFFFFFFFFFFFFFFFFFFFFFFFFFFFFFFFFFFFFFFFFFFFFFFFFFFFFFFFFFFFFFFFFFFFFFFFFFFFFFFFFFFFFFFFFFFFFFFFFFFFFFFFFFFFFFFFFFFFFFFFFFFFFFFFFFFFFFFFFFFFFFFFFFFFFFFFFFFFFFFFFFFFFFFFFFFFFFFFFFFFFFFFFFFFFFFFFFFFFFFFFFFFFFFFFFFFFFFFFFFFFFFFFFFFFFFFFFFFFFFFFFFFFFFFFFFFFFFFFFFFFFFFFFFFFFFFFFFFFFFFFFFFFFFFFFFFFFFFFFFFFFFFFFFFFFFFFFFFFFFFFFFFFFFFFFFFFFFFFFFFFFFFFFFFFFFFFFFFFFFFFFFFFFFFFFFFFFFFFFEF9FEFDFDFEFCF9F5F1F4FBF2CE3242340D082F8D807DA7754E0B060604000A050D0C040201042EB8C0EBD2F4C474A771243B2402040C74AC99E6F8F9FCFFFFFFFFFFFFFFFFFFFFFFFFFFFFFFFFFFFFFFFFFFFFFFFFFFFFFFFFFFFFFFFFFFFFFFFFFFFFFFFFFFFFFFFFFFFFFFFFFFFFFFFFFFFFFFFFFFFFFFFFFFFFFFFFFFFFFFFFFFFFFFFFFFFFFFFFFFFFFFFFFFFFFFFFFFFFFFFFFFFFFFFFFFFFFFFFFFFFFFFFFFFFFFFFFFFFFFFFFFFFFFFFFFFFFFFFFFFFFFFFFFFFFFFFFFFFFFFFFFFFFFFFFFFFFFFFFFFFFFFFFFFFFFFFFFFFFFFFFFFFFFFFFFFFFFFFFFFFFFFFFFFFFFFFFFFFFFFFFFFFFFFFFFFFFFFFFFFFFFFFFFFFFFFFFFFFFFFFFFFFFFFFFFFEF9FEF8F9FCF9FBFFFEFEF9F3FEF9D7DA2D1B5E8F7B96BF680900010100050400041F1D080B0700010AD216B357B1D67A80C3740A06267FEEFCF4F2F9FEFFFFFFFFFFFFFFFFFFFFFFFFFFFFFFFFFFFFFFFFFFFFFFFFFFFFFFFFFFFFFFFFFFFFFFFFFFFFFFFFFFFFFFFFFFFFFFFFFFFFFFFFFFFFFFFFFFFFFFFFFFFFFFFFFFFFFFFFFFFFFFFFFFFFFFFFFFFFFFFFFFFFFFFFFFFFFFFFFFFFFFFFFFFFFFFFFFFFFFFFFFFFFFFFFFFFFFFFFFFFFFFFFFFFFFFFFFFFFFFFFFFFFFFFFFFFFFFFFFFFFFFFFFFFFFFFFFFFFFFFFFFFFFFFFFFFFFFFFFFFFFFFFFFFFFFFFFFFFFFFFFFFFFFFFFFFFFFFFFFFFFFFFFFFFFFFFFFFFFFFFFFFFFFFFFFFFFFFFFFFFFFFFFFEFCFEF8F7FCFCFCFEFDF9F3FEFAD52FB474C1A96071816D5B0601011501001607006C77180000020E12000000AED6527D890704143478FEFEF3FAFEF5F8FFFFFFFFFFFFFFFFFFFFFFFFFFFFFFFFFFFFFFFFFFFFFFFFFFFFFFFFFFFFFFFFFFFFFFFFFFFFFFFFFFFFFFFFFFFFFFFFFFFFFFFFFFFFFFFFFFFFFFFFFFFFFFFFFFFFFFFFFFFFFFFFFFFFFFFFFFFFFFFFFFFFFFFFFFFFFFFFFFFFFFFFFFFFFFFFFFFFFFFFFFFFFFFFFFFFFFFFFFFFFFFFFFFFFFFFFFFFFFFFFFFFFFFFFFFFFFFFFFFFFFFFFFFFFFFFFFFFFFFFFFFFFFFFFFFFFFFFFFFFFFFFFFFFFFFFFFFFFFFFFFFFFFFFFFFFFFFFFFFFFFFFFFFFFFFFFFFFFFFFFFFFFFFFFFFFFFFFFFFFFFFFFFFEFDFEFEFEFDFDF8F4FAFEE9BADBCCAE2E91410F2C7B61090E0800011F1D030100162D200B0A020603010007040C131F736C07105EE8F6FEFEFDFAFDF7F9FFFFFFFFFFFFFFFFFFFFFFFFFFFFFFFFFFFFFFFFFFFFFFFFFFFFFFFFFFFFFFFFFFFFFFFFFFFFFFFFFFFFFFFFFFFFFFFFFFFFFFFFFFFFFFFFFFFFFFFFFFFFFFFFFFFFFFFFFFFFFFFFFFFFFFFFFFFFFFFFFFFFFFFFFFFFFFFFFFFFFFFFFFFFFFFFFFFFFFFFFFFFFFFFFFFFFFFFFFFFFFFFFFFFFFFFFFFFFFFFFFFFFFFFFFFFFFFFFFFFFFFFFFFFFFFFFFFFFFFFFFFFFFFFFFFFFFFFFFFFFFFFFFFFFFFFFFFFFFFFFFFFFFFFFFFFFFFFFFFFFFFFFFFFFFFFFFFFFFFFFFFFFFFFFFFFFFFFFFFFFFFFFFFEFEFCFDFEFEFBFBFDEBC99B44241D2B061407040A1711000101000E6A5D0C000280BD8340291B000300080902010008140E0C284782FDEBFCFEFBFDFEFEFFFFFFFFFFFFFFFFFFFFFFFFFFFFFFFFFFFFFFFFFFFFFFFFFFFFFFFFFFFFFFFFFFFFFFFFFFFFFFFFFFFFFFFFFFFFFFFFFFFFFFFFFFFFFFFFFFFFFFFFFFFFFFFFFFFFFFFFFFFFFFFFFFFFFFFFFFFFFFFFFFFFFFFFFFFFFFFFFFFFFFFFFFFFFFFFFFFFFFFFFFFFFFFFFFFFFFFFFFFFFFFFFFFFFFFFFFFFFFFFFFFFFFFFFFFFFFFFFFFFFFFFFFFFFFFFFFFFFFFFFFFFFFFFFFFFFFFFFFFFFFFFFFFFFFFFFFFFFFFFFFFFFFFFFFFFFFFFFFFFFFFFFFFFFFFFFFFFFFFFFFFFFFFFFFFFFFFFFFFFFFFFFFFFFEFCFAFAFDD6E2DF81386A8A02060700000102010101000200050419150204056C8C86AFC899492B2815231026160509080D50C2DAD8FAFEF0FAF6FAFFFFFFFFFFFFFFFFFFFFFFFFFFFFFFFFFFFFFFFFFFFFFFFFFFFFFFFFFFFFFFFFFFFFFFFFFFFFFFFFFFFFFFFFFFFFFFFFFFFFFFFFFFFFFFFFFFFFFFFFFFFFFFFFFFFFFFFFFFFFFFFFFFFFFFFFFFFFFFFFFFFFFFFFFFFFFFFFFFFFFFFFFFFFFFFFFFFFFFFFFFFFFFFFFFFFFFFFFFFFFFFFFFFFFFFFFFFFFFFFFFFFFFFFFFFFFFFFFFFFFFFFFFFFFFFFFFFFFFFFFFFFFFFFFFFFFFFFFFFFFFFFFFFFFFFFFFFFFFFFFFFFFFFFFFFFFFFFFFFFFFFFFFFFFFFFFFFFFFFFFFFFFFFFFFFFFFFFFFFFFFFFFFFFFEF8FDFDFCFD9C9A7E2137813004020100030100000204030000020003000003012063789EC389C5BF933FB3278D3F0405030C3F90A99AEAD4F3FDF9F9FCFFFFFFFFFFFFFFFFFFFFFFFFFFFFFFFFFFFFFFFFFFFFFFFFFFFFFFFFFFFFFFFFFFFFFFFFFFFFFFFFFFFFFFFFFFFFFFFFFFFFFFFFFFFFFFFFFFFFFFFFFFFFFFFFFFFFFFFFFFFFFFFFFFFFFFFFFFFFFFFFFFFFFFFFFFFFFFFFFFFFFFFFFFFFFFFFFFFFFFFFFFFFFFFFFFFFFFFFFFFFFFFFFFFFFFFFFFFFFFFFFFFFFFFFFFFFFFFFFFFFFFFFFFFFFFFFFFFFFFFFFFFFFFFFFFFFFFFFFFFFFFFFFFFFFFFFFFFFFFFFFFFFFFFFFFFFFFFFFFFFFFFFFFFFFFFFFFFFFFFFFFFFFFFFFFFFFFFCF9F8F9FDF9FEFBFEFED1EAD07623021A380302010300050A28150203020305000501002B1C06092674B9BFCDF8EB6DBBF25EC97114010115466654A0A0A5F1FDFDFBFCFAF9FFFFFDF7FEFEFFFFFFFFFFFFFFFFFFFDFAFFFEFEFEFEFCFDFDFDFFFFFEFDFEFDFEFFFFFFFFFFFFFFFFFFFFFFFFFFFFFFFFFFFFFFFFFFFFFFFFFFFFFFFFFFFFFFFFFFFFFFFFFFFFFFFFFFFFFFFFFFFFFFFFFFFFFFFFFFFFFFFFFFFFFFFFFFFFFFFFFFFFFFFFFFFFFFFFFFFFFFFFFFFFFFFFFFFFFFFFFFFFFFFFFFFFFFFFFFFFFFFFFFFFFFFFFFFFFFFFFFFFFFFFFFFFFFFFFFFFFFFFFFFFFFFFFFFFFFFFFFFFFFFFFFFFFFFFFFFFFFFFFFFFFFFFFFFFFFFFFFFFFFFFFFFFFAFCFDFBFAFAFCF6FDFEF2FC97B0C7590102040303070402060319603200070E141B0303050E543B12060457A0ACF4FDF3ACFBFBC2C1903605000B26351B4F539CC6D9F4FDFEF6F3FAFAFDF9FEFDFEFFFFFFFFFFFFFFFFFEFBFEFDFDF1D3EEFBFBFDFBF9FBFEFEFAFEFFFFFFFFFFFFFFFFFFFFFFFFFFFFFFFFFFFFFFFFFFFFFFFFFFFFFFFFFFFFFFFFFFFFFFFFFFFFFFFFFFFFFFFFFFFFFFFFFFFFFFFFFFFFFFFFFFFFFFFFFFFFFFFFFFFFFFFFFFFFFFFFFFFFFFFFFFFFFFFFFFFFFFFFFFFFFFFFFFFFFFFFFFFFFFFFFFFFFFFFFFFFFFFFFFFFFFFFFFFFFFFFFFFFFFFFFFFFFFFFFFFFFFFFFFFFFFFFFFFFFFFFFFFFFFFFFFFFFFFFFFFFFFFFFFFFFFFFFFFFF8FAFEFFFEFCFAFAFCF9F3F2BA725A2801020606020C03000903103E220B224256421D17040C352A0E0101265E92C7D0F5CDF0F9FB99642A040001030903021D575598D1E4F3F9F7FAF9FAF6FDFDFEFFFFFFFFFFFFFFFCF9FEF8F9FCD584CADEE8FBFAF6F9FEFEFBFEFFFFFFFFFFFFFFFFFFFFFFFFFFFFFFFFFFFFFFFFFFFFFFFFFFFFFFFFFFFFFFFFFFFFFFFFFFFFFFFFFFFFFFFFFFFFFFFFFFFFFFFFFFFFFFFFFFFFFFFFFFFFFFFFFFFFFFFFFFFFFFFFFFFFFFFFFFFFFFFFFFFFFFFFFFFFFFFFFFFFFFFFFFFFFFFFFFFFFFFFFFFFFFFFFFFFFFFFFFFFFFFFFFFFFFFFFFFFFFFFFFFFFFFFFFFFFFFFFFFFFFFFFFFFFFFFFFFFFFFFFFFFFFFFFFFFFFFFFFFFFCFBFBFCFEFBF8FEFCF8FEF2FA5B0D0502000316260B2D02040402061D4E7B897E63585C17010607010303062756616FB561EACFD2872D07030307182000050405036291B6E3FBFDFBFBF9F4FDFBFEFFFFFFFFFFFFFFFCF8FDF7F9FDBA4A90AED2F8FCFAFBFDFBFFFFFFFFFFFFFFFFFFFFFFFFFFFFFFFFFFFFFFFFFFFFFFFFFFFFFFFFFFFFFFFFFFFFFFFFFFFFFFFFFFFFFFFFFFFFFFFFFFFFFFFFFFFFFFFFFFFFFFFFFFFFFFFFFFFFFFFFFFFFFFFFFFFFFFFFFFFFFFFFFFFFFFFFFFFFFFFFFFFFFFFFFFFFFFFFFFFFFFFFFFFFFFFFFFFFFFFFFFFFFFFFFFFFFFFFFFFFFFFFFFFFFFFFFFFFFFFFFFFFFFFFFFFFFFFFFFFFFFFFFFFFFFFFFFFFFFFFFFFFFFFFFFFDF8F9FEFCFAFDF7DCFAE5D44F0D030302012666147905050A202367A4B6A2798995885414210E0305050405050619440ECD6B63621F0206040F3E550509000004354F90C5E3F4F8FBFDFAFDF9FEFFFFFFFFFFFFFFFEFCF9FCFAF1A148589BDCF9FAF9FEFCF7FFFFFFFFFFFFFFFFFFFFFFFFFFFFFFFFFFFFFFFFFFFFFFFFFFFFFFFFFFFFFFFFFFFFFFFFFFFFFFFFFFFFFFFFFFFFFFFFFFFFFFFFFFFFFFFFFFFFFFFFFFFFFFFFFFFFFFFFFFFFFFFFFFFFFFFFFFFFFFFFFFFFFFFFFFFFFFFFFFFFFFFFFFFFFFFFFFFFFFFFFFFFFFFFFFFFFFFFFFFFFFFFFFFFFFFFFFFFFFFFFFFFFFFFFFFFFFFFFFFFFFFFFFFFFFFFFFFFFFFFFFFFFFFFFFFFFFFFFFFFFFFFFEFDFBFBFEFDFBF9EF7BCC8A61280F06040401215F1A6B0B24185E5DA2A2767D8598AA709F405641210803080803000001035C0A0409081712020D354908010300000523586D8EC1DDF3FBF5FDFCFEFFFFFFFFFFFFFFFEFAF2F9F4D8937D5EB7F9FCF5F4FFFDF8FFFFFFFFFFFFFFFFFFFFFFFFFFFFFFFFFFFFFFFFFFFFFFFFFFFFFFFFFFFFFFFFFFFFFFFFFFFFFFFFFFFFFFFFFFFFFFFFFFFFFFFFFFFFFFFFFFFFFFFFFFFFFFFFFFFFFFFFFFFFFFFFFFFFFFFFFFFFFFFFFFFFFFFFFFFFFFFFFFFFFFFFFFFFFFFFFFFFFFFFFFFFFFFFFFFFFFFFFFFFFFFFFFFFFFFFFFFFFFFFFFFFFFFFFFFFFFFFFFFFFFFFFFFFFFFFFFFFFFFFFFFFFFFFFFFFFFFFFFFFFFFFFDFEFDFDFCFAF6E2E5274F14050513110502020B1802171B773466867577815E2B317692DD7355865616010107090104070000040C0311543C01030600020007010000080500245B91D8FBFDFCF8FEFFFFFFFFFFFFFFFEFBF7FEFDDEBDDEBEDEFCFBF9FAFEFCF8FFFFFFFFFFFFFFFFFFFFFFFFFFFFFFFFFFFFFFFFFFFFFFFFFFFFFFFFFFFFFFFFFFFFFFFFFFFFFFFFFFFFFFFFFFFFFFFFFFFFFFFFFFFFFFFFFFFFFFFFFFFFFFFFFFFFFFFFFFFFFFFFFFFFFFFFFFFFFFFFFFFFFFFFFFFFFFFFFFFFFFFFFFFFFFFFFFFFFFFFFFFFFFFFFFFFFFFFFFFFFFFFFFFFFFFFFFFFFFFFFFFFFFFFFFFFFFFFFFFFFFFFFFFFFFFFFFFFFFFFFFFFFFFFFFFFFFFFFFFFFFFFFEFEFAFAFAF7E69FB844050010085B43090A06287F208615E18B24414496A5300000167BECD18D8F420D02050600040F080201010000094C39010305020900071403030104043376ACEEFDF7F6FEFFFFFFFFFFFFFFFFFDF7F4FEFBEEE8F8F7F1F5FBFEFEFCF9FBFEFEFFFFFFFFFFFFFFFFFFFFFFFFFFFFFFFFFFFFFFFFFFFFFFFFFFFFFFFFFFFFFFFFFFFFFFFFFFFFFFFFFFFFFFFFFFFFFFFFFFFFFFFFFFFFFFFFFFFFFFFFFFFFFFFFFFFFFFFFFFFFFFFFFFFFFFFFFFFFFFFFFFFFFFFFFFFFFFFFFFFFFFFFFFFFFFFFFFFFFFFFFFFFFFFFFFFFFFFFFFFFFFFFFFFFFFFFFFFFFFFFFFFFFFFFFFFFFFFFFFFFFFFFFFFFFFFFFFFFFFFFFFFFFFFFFFFFFFFFFFFFFCFDF8F8FCF7C93E451A010801023C310703061E7EAFFAA6D46D19053683600707000844C6EBFCC59B49120303060000021500090B080200050905050A000B44611509050401245689DAFEFDFBFEFFFFFFFFFFFFFFFFFFFEFDFFFFFFFFF5FBFFFFFFFDFAF9FFFEFEFEFFFFFFFFFFFFFFFFFFFFFFFFFFFFFFFFFFFFFFFFFFFFFFFFFFFFFFFFFFFFFFFFFFFFFFFFFFFFFFFFFFFFFFFFFFFFFFFFFFFFFFFFFFFFFFFFFFFFFFFFFFFFFFFFFFFFFFFFFFFFFFFFFFFFFFFFFFFFFFFFFFFFFFFFFFFFFFFFFFFFFFFFFFFFFFFFFFFFFFFFFFFFFFFFFFFFFFFFFFFFFFFFFFFFFFFFFFFFFFFFFFFFFFFFFFFFFFFFFFFFFFFFFFFFFFFFFFFFFFFFFFFFFFFFFFFFFFFFFFFFFCFDFAFAD9884B13020007080100010907000032C6F4FEFCEE91230136AEBC46050C0C61E3F6FDFDF6D6AC53111000010C0600098C2A080402020F46680B558D815C1F4458260137A8EFFDFDEEFEFFFFFFFFFFFFFFFFFFFEFEFFFFFFFFFEFEFFFFFFFEFEFEFFFEFEFEFFFFFFFFFFFFFFFFFFFFFFFFFFFFFFFFFFFFFFFFFFFFFFFFFFFFFFFFFFFFFFFFFFFFFFFFFFFFFFFFFFFFFFFFFFFFFFFFFFFFFFFFFFFFFFFFFFFFFFFFFFFFFFFFFFFFFFFFFFFFFFFFFFFFFFFFFFFFFFFFFFFFFFFFFFFFFFFFFFFFFFFFFFFFFFFFFFFFFFFFFFFFFFFFFFFFFFFFFFFFFFFFFFFFFFFFFFFFFFFFFFFFFFFFFFFFFFFFFFFFFFFFFFFFFFFFFFFFFFFFFFFFFFFFFFFFFFFFFFFFFEFAFBF4B4605C340005007432110D0605486F2175FCFEEFEF7A832F0F3E666B480418445EB5F3FCFAECED8F1702050B0100000C4A2D0A115E62528B38051B7EB58C2034391C1120407ECFFEFCFDFEFFFFFFFFFFFFFFFFFFFFFFFFFFFFFFFFFFFFFFFFFFFFFFFFFFFFFFFFFFFFFFFFFFFFFFFFFFFFFFFFFFFFFFFFFFFFFFFFFFFFFFFFFFFFFFFFFFFFFFFFFFFFFFFFFFFFFFFFFFFFFFFFFFFFFFFFFFFFFFFFFFFFFFFFFFFFFFFFFFFFFFFFFFFFFFFFFFFFFFFFFFFFFFFFFFFFFFFFFFFFFFFFFFFFFFFFFFFFFFFFFFFFFFFFFFFFFFFFFFFFFFFFFFFFFFFFFFFFFFFFFFFFFFFFFFFFFFFFFFFFFFFFFFFFFFFFFFFFFFFFFFFFFFFFFFFFFFFFFFFFFFFFFFFFFFFFF8FAFDD87630411C00010B501703020103256797876CFEFBECA157150E227CC36A00041261D8F2FBFDFCFAD97714010203000A030F08110E333727431F3B9EBDEFD86E200904030530A2F3F3FBF7FEFFFFFFFFFFFFFFFFFFFFFFFFFFFFFFFFFFFFFFFFFFFFFFFFFFFFFFFFFFFFFFFFFFFFFFFFFFFFFFFFFFFFFFFFFFFFFFFFFFFFFFFFFFFFFFFFFFFFFFFFFFFFFFFFFFFFFFFFFFFFFFFFFFFFFFFFFFFFFFFFFFFFFFFFFFFFFFFFFFFFFFFFFFFFFFFFFFFFFFFFFFFFFFFFFFFFFFFFFFFFFFFFFFFFFFFFFFFFFFFFFFFFFFFFFFFFFFFFFFFFFFFFFFFFFFFFFFFFFFFFFFFFFFFFFFFFFFFFFFFFFFFFFFFFFFFFFFFFFFFFFFFFFFFFFFFFFFFFFFFFFFFFFFFFFFFFFCF6F4C95D0910090001031608080A514640B5EBD69DF6FEDA7C822C031158B6AE3F070333AFF2FDF9F6B36E3906000204000E01010407030A0B080C0B35A1E3FCC5440E01043E603F68DBFBFDF6FEFFFFFFFFFFFFFFFFFFFFFFFFFFFFFFFFFFFFFFFFFFFFFFFFFFFFFFFFFFFFFFFFFFFFFFFFFFFFFFFFFFFFFFFFFFFFFFFFFFFFFFFFFFFFFFFFFFFFFFFFFFFFFFFFFFFFFFFFFFFFFFFFFFFFFFFFFFFFFFFFFFFFFFFFFFFFFFFFFFFFFFFFFFFFFFFFFFFFFFFFFFFFFFFFFFFFFFFFFFFFFFFFFFFFFFFFFFFFFFFFFFFFFFFFFFFFFFFFFFFFFFFFFFFFFFFFFFFFFFFFFFFFFFFFFFFFFFFFFFFFFFFFFFFFFFFFFFFFFFFFFFFFFFFFFFFFFFFFFFFFFFFFFFFFFFFFFFFAC971230206030E0000020604044779A3EAFAF7E6FAF9F094691F062B6494A47A0F5391D8F2F7F3F7CA9E918C6A07010004010402085D460004080635BFF4FD9A5F451E0625381D42C6FCFEF9FEFFFFFFFFFFFFFFFFFFFFFFFFFFFFFFFFFFFFFFFFFFFFFFFFFFFFFFFFFFFFFFFFFFFFFFFFFFFFFFFFFFFFFFFFFFFFFFFFFFFFFFFFFFFFFFFFFFFFFFFFFFFFFFFFFFFFFFFFFFFFFFFFFFFFFFFFFFFFFFFFFFFFFFFFFFFFFFFFFFFFFFFFFFFFFFFFFFFFFFFFFFFFFFFFFFFFFFFFFFFFFFFFFFFFFFFFFFFFFFFFFFFFFFFFFFFFFFFFFFFFFFFFFFFFFFFFFFFFFFFFFFFFFFFFFFFFFFFFFFFFFFFFFFFFFFFFFFFFFFFFFFFFFFFFFFFFFFFFFFFFFFFFFFFFFFFFF8FAD7802A0303010803000002060452B5BA859DFBF5FEFBF1CF450A051850877C390D4662C3F3FAF9FCEFE2E2BD79670C010300020105433200041039AEE9FCFC8443401C020D120439D1FBECFCFEFFFFFFFFFFFFFFFFFFFFFFFFFFFFFFFFFFFFFFFFFFFFFFFFFFFFFFFFFFFFFFFFFFFFFFFFFFFFFFFFFFFFFFFFFFFFFFFFFFFFFFFFFFFFFFFFFFFFFFFFFFFFFFFFFFFFFFFFFFFFFFFFFFFFFFFFFFFFFFFFFFFFFFFFFFFFFFFFFFFFFFFFFFFFFFFFFFFFFFFFFFFFFFFFFFFFFFFFFFFFFFFFFFFFFFFFFFFFFFFFFFFFFFFFFFFFFFFFFFFFFFFFFFFFFFFFFFFFFFFFFFFFFFFFFFFFFFFFFFFFFFFFFFFFFFFFFFFFFFFFFFFFFFFFFFFFFFFFFFFFFFFFFFFFFFFFF7FAD7791E0405000202006221486CACEAA03E9DFDF6EFFEE89B6B230B20537E733D071651C6F0F6FCF9F8F9F8D893726A02020200030E483809133A7DB2F4FEFD902B170D0407070136D2FCEDFAFEFFFFFFFFFFFFFFFFFFFFFFFFFFFFFFFFFFFFFFFFFFFFFFFFFFFFFFFFFFFFFFFFFFFFFFFFFFFFFFFFFFFFFFFFFFFFFFFFFFFFFFFFFFFFFFFFFFFFFFFFFFFFFFFFFFFFFFFFFFFFFFFFFFFFFFFFFFFFFFFFFFFFFFFFFFFFFFFFFFFFFFFFFFFFFFFFFFFFFFFFFFFFFFFFFFFFFFFFFFFFFFFFFFFFFFFFFFFFFFFFFFFFFFFFFFFFFFFFFFFFFFFFFFFFFFFFFFFFFFFFFFFFFFFFFFFFFFFFFFFFFFFFFFFFFFFFFFFFFFFFFFFFFFFFFFFFFFFFFFFFFFFFFFFFFFFFFCFDDC7C1D0407000C56196F67B1DDF2CD848FE6FDF9EFFEE9846D260A3482A071460B45AEEDF5F5F9F7FAFFFCD88F7475060002040451753B12407A81A5F7FEFD9D2A03042936180234C9FBFBFBFEFFFFFFFFFFFFFFFFFFFFFFFFFFFFFFFFFFFFFFFFFFFFFFFFFFFFFFFFFFFFFFFFFFFFFFFFFFFFFFFFFFFFFFFFFFFFFFFFFFFFFFFFFFFFFFFFFFFFFFFFFFFFFFFFFFFFFFFFFFFFFFFFFFFFFFFFFFFFFFFFFFFFFFFFFFFFFFFFFFFFFFFFFFFFFFFFFFFFFFFFFFFFFFFFFFFFFFFFFFFFFFFFFFFFFFFFFFFFFFFFFFFFFFFFFFFFFFFFFFFFFFFFFFFFFFFFFFFFFFFFFFFFFFFFFFFFFFFFFFFFFFFFFFFFFFFFFFFFFFFFFFFFFFFFFFFFFFFEFEFEFEFEFEFEFEFEFEE6B874280302020D7A6A5CBBF0F8FAB6A2F3FAF7FAFEFCF2D6721C011E688845153C83C7F2FDFDF9FAFDFFF8C38CCB630900000C3F824D0B0C488C94EAFEFDD56C3A2C2C7CA37A5D79DCFAF7FEFFFFFFFFFFFFFFFFFFFFFFFFFFFFFFFFFFFFFFFFFFFFFFFFFFFFFFFFFFFFFFFFFFFFFFFFFFFFFFFFFFFFFFFFFFFFFFFFFFFFFFFFFFFFFFFFFFFFFFFFFFFFFFFFFFFFFFFFFFFFFFFFFFFFFFFFFFFFFFFFFFFFFFFFFFFFFFFFFFFFFFFFFFFFFFFFFFFFFFFFFFFFFFFFFFFFFFFFFFFFFFFFFFFFFFFFFFFFFFFFFFFFFFFFFFFFFFFFFFFFFFFFFFFFFFFFFFFFFFFFFFFFFFFFFFFFFFFFFFFFFFFFFFFFFFFFFFFFFFFFFFFFFFFFFFFFFFFFFEFEFEFEFEFEFEFEFEFCCD7C3D1901000B384DCDE2F8F8F8FAF4F4FDFAF3FEF9FAEFBA8F2B02113444241A627FA9EFFDF6FAFBFAF9F5EFF0EB9D4707000B76470C042A5C94EEF3FEF9B23D545371AFD4CFCAE3FBF7F0FEFFFFFFFFFFFFFFFFFFFFFFFFFFFFFFFFFFFFFFFFFFFFFFFFFFFFFFFFFFFFFFFFFFFFFFFFFFFFFFFFFFFFFFFFFFFFFFFFFFFFFFFFFFFFFFFFFFFFFFFFFFFFFFFFFFFFFFFFFFFFFFFFFFFFFFFFFFFFFFFFFFFFFFFFFFFFFFFFFFFFFFFFFFFFFFFFFFFFFFFFFFFFFFFFFFFFFFFFFFFFFFFFFFFFFFFFFFFFFFFFFFFFFFFFFFFFFFFFFFFFFFFFFFFFFFFFFFFFFFFFFFFFFFFFFFFFFFFFFFFFFFFFFFFFFFFFFFFFFFFFFFFFFFFFFFFFFFFFFEFEFEFEFEFEFEFEFEE8DECF933C05021A93E7F1FDFCFBFCF6F7FBFAFEFBFCF2FBD8563210031E719947163AB2EAFDFFF2F8FDF7F1F5FEFCD55CDA140102142C1A2D85D3EDFEFEF9F5F58A2706687E8E8C6EC0D7E6FEFCFEFFFFFFFFFFFFFFFFFFFFFFFFFFFFFFFFFFFFFFFFFFFFFFFFFFFFFFFFFFFFFFFFFFFFFFFFFFFFFFFFFFFFFFFFFFFFFFFFFFFFFFFFFFFFFFFFFFFFFFFFFFFFFFFFFFFFFFFFFFFFFFFFFFFFFFFFFFFFFFFFFFFFFFFFFFFFFFFFFFFFFFFFFFFFFFFFFFFFFFFFFFFFFFFFFFFFFFFFFFFFFFFFFFFFFFFFFFFFFFFFFFFFFFFFFFFFFFFFFFFFFFFFFFFFFFFFFFFFFFFFFFFFFFFFFFFFFFFFFFFFFFFFFFFFFFFFFFFFFFFFFFFFFFFFFFFFFFFEFEFEFEFEFEFEFEFEB9958342030109130FADE9FCF5FAFBF9FAFDFDFAF6FEFCF4F8D7B13600020B1005208CEBFCFCFCFEF8F9FBFEFCFBFCF2DDF2080209117A597B8C84FDFAF9EDFDFE96280454656E6D618F56CFFDF5FDFFFFFFFFFFFFFFFFFFFFFFFFFFFFFFFFFFFFFFFFFFFFFFFFFFFFFFFFFFFFFFFFFFFFFFFFFFFFFFFFFFFFFFFFFFFFFFFFFFFFFFFFFFFFFFFFFFFFFFFFFFFFFFFFFFFFFFFFFFFFFFFFFFFFFFFFFFFFFFFFFFFFFFFFFFFFFFFFFFFFFFFFFFFFFFFFFFFFFFFFFFFFFFFFFFFFFFFFFFFFFFFFFFFFFFFFFFFFFFFFFFFFFFFFFFFFFFFFFFFFFFFFFFFFFFFFFFFFFFFFFFFFFFFFFFFFFFFFFFFFFFFFFFFFFFFFFFFFFFFFFFFFFFFFFFFFFFFEFEFEFEFEFEFEFEFEEEA342120A0000010813BCE4E0F1FDFDF7F5F8F7FEFCF3FEDF360301010202000C387CEFFCF5F2FBFEFEFBFBFDFAF9D138E10C000301100C1F91FEFEFAF3FEF6F9A126060D10120F040127CEFCFDFEFFFFFFFFFFFFFFFFFFFFFFFFFFFFFFFFFFFFFFFFFFFFFFFFFFFFFFFFFFFFFFFFFFFFFFFFFFFFFFFFFFFFFFFFFFFFFFFFFFFFFFFFFFFFFFFFFFFFFFFFFFFFFFFFFFFFFFFFFFFFFFFFFFFFFFFFFFFFFFFFFFFFFFFFFFFFFFFFFFFFFFFFFFFFFFFFFFFFFFFFFFFFFFFFFFFFFFFFFFFFFFFFFFFFFFFFFFFFFFFFFFFFFFFFFFFFFFFFFFFFFFFFFFFFFFFFFFFFFFFFFFFFFFFFFFFFFFFFFFFFFFFFFFFFFFFFFFFFFFFFFFFFFFFFFFFFFFFEFEFEFEFEFEFEFEFEB083884E0105070323D93C0D1A92E8F2F3FFFFFCF5F3FEFEFECB300003070D0A2E93878DE8FCFEFEF5FAFEFAEBE5FEFBD60F000000060326A4FDFDFEFEFEFDF2FDFE922100669141040734CCFEF4FDFFFFFFFFFFFFFFFFFFFFFFFFFFFFFFFFFFFFFFFFFFFFFFFFFFFFFFFFFFFFFFFFFFFFFFFFFFFFFFFFFFFFFFFFFFFFFFFFFFFFFFFFFFFFFFFFFFFFFFFFFFFFFFFFFFFFFFFFFFFFFFFFFFFFFFFFFFFFFFFFFFFFFFFFFFFFFFFFFFFFFFFFFFFFFFFFFFFFFFFFFFFFFFFFFFFFFFFFFFFFFFFFFFFFFFFFFFFFFFFFFFFFFFFFFFFFFFFFFFFFFFFFFFFFFFFFFFFFFFFFFFFFFFFFFFFFFFFFFFFFFFFFFFFFFFFFFFFFFFFFFFFFFFFFFFFFFFFEFEFEFEFEFEFEFEFEF89E360A0200001AA2F5380F97E8FDFEFEFFFFFEFDF9FAEADB619F4806003097E1E6F3ECFCF9F6FDFEFEFAFAFDFDF8EDFCE40D0605011193F1F6F4FBF4F8EF988D782B0F040856772D0D2BE5FEFEFEFFFFFFFFFFFFFFFFFFFFFFFFFFFFFFFFFFFFFFFFFFFFFFFFFFFFFFFFFFFFFFFFFFFFFFFFFFFFFFFFFFFFFFFFFFFFFFFFFFFFFFFFFFFFFFFFFFFFFFFFFFFFFFFFFFFFFFFFFFFFFFFFFFFFFFFFFFFFFFFFFFFFFFFFFFFFFFFFFFFFFFFFFFFFFFFFFFFFFFFFFFFFFFFFFFFFFFFFFFFFFFFFFFFFFFFFFFFFFFFFFFFFFFFFFFFFFFFFFFFFFFFFFFFFFFFFFFFFFFFFFFFFFFFFFFFFFFFFFFFFFFFFFFFFFFFFFFFFFFFFFFFFFFFFFFFFFFFEFEFEFEFEFEFEFEFEAF4A0E02050500072DE237AC61ACE8F8F9F2F5FEF9F0FDEFE7B43F0B08080940917B292E90F2FCF5FAFFFAF8F7F7FEF9DF3FC415000013A2FDFDF7F4FDFEFAE38B21050502538752130127DAFEFBFEFFFFFFFFFFFFFFFFFFFFFFFFFFFFFFFFFFFFFFFFFFFFFFFFFFFFFFFFFFFFFFFFFFFFFFFFFFFFFFFFFFFFFFFFFFFFFFFFFFFFFFFFFFFFFFFFFFFFFFFFFFFFFFFFFFFFFFFFFFFFFFFFFFFFFFFFFFFFFFFFFFFFFFFFFFFFFFFFFFFFFFFFFFFFFFFFFFFFFFFFFFFFFFFFFFFFFFFFFFFFFFFFFFFFFFFFFFFFFFFFFFFFFFFFFFFFFFFFFFFFFFFFFFFFFFFFFFFFFFFFFFFFFFFFFFFFFFFFFFFFFFFFFFFFFFFFFFFFFFFFFFFFFFFFFFFFFFFEFEFEFEFEFEFEFEFEA63905010202041692E9D2F2D2E7F9F7F7FBFBFCFDF6D54B342D0A0907276865231B020B31A7F1FBFBF5F7FDFDF9F7FAF4B54F07020012A0FCFCF8F7FEF2A63F1F040704091C657E300624C6E4FCFEFFFFFFFFFFFFFFFFFFFFFFFFFFFFFFFFFFFFFFFFFFFFFFFFFFFFFFFFFFFFFFFFFFFFFFFFFFFFFFFFFFFFFFFFFFFFFFFFFFFFFFFFFFFFFFFFFFFFFFFFFFFFFFFFFFFFFFFFFFFFFFFFFFFFFFFFFFFFFFFFFFFFFFFFFFFFFFFFFFFFFFFFFFFFFFFFFFFFFFFFFFFFFFFFFFFFFFFFFFFFFFFFFFFFFFFFFFFFFFFFFFFFFFFFFFFFFFFFFFFFFFFFFFFFFFFFFFFFFFFFFFFFFFFFFFFFFFFFFFFFFFFFFFFFFFFFFFFFFFFFFFFFFFFFFFFFFFFBF9FEF9F4FEF7F7BE541603030A0A071F4DCDFAFDF6FBFDF9FAFCF9FCFFF9CE991501020A050E447451010E65A8DCF2FCF8F8FBFEFDFBF9F8EB62A12706000FA0F8F6F9FBF8EFD869080B0A4B56235482330C0A3946E5FCFFFFF4F8FCFCFDFFFFFFFFFFFFFFFFFFFFFFFFFFFFFFFFFFFFFFFFFFFFFFFFFFFFFFFFFFFFFFFFFFFFFFFFFFFFFFFFFFFFFFFFFFFFFFFFFFFFFFFFFFFFFFFFFFFFFFFFFFFFFFFFFFFFFFFFFFFFFFFFFFFFFFFFFFFFFFFFFFFFFFFFFFFFFFFFFFFFFFFFFFFFFFFFFFFFFFFFFFFFFFFFFFFFFFFFFFFFFFFFFFFFFFFFFFFFFFFFFFFFFFFFFFFFFFFFFFFFFFFFFFFFFFFFFFFFFFFFFFFFFFFFFFFFFFFFFFFFFFFFFFFFFFFFFFFFFFFFFDF6EEFCF8F6FAB986652706072B594B71A3E7F6FAFBFCF6F9FDFDF9F8F6C663580B0A10080809132A28010C3652BDF0FCFCFBFCFDFDFDFBF4EDAA62931D0310A0FAFAFBFCFBF1C244085169553A18202E1305000215C1D3F7EBF1F6FAFDF8F4FFFFFFFFFFFFFFFFFFFFFFFFFFFFFFFFFFFFFFFFFFFFFFFFFFFFFFFFFFFFFFFFFFFFFFFFFFFFFFFFFFFFFFFFFFFFFFFFFFFFFFFFFFFFFFFFFFFFFFFFFFFFFFFFFFFFFFFFFFFFFFFFFFFFFFFFFFFFFFFFFFFFFFFFFFFFFFFFFFFFFFFFFFFFFFFFFFFFFFFFFFFFFFFFFFFFFFFFFFFFFFFFFFFFFFFFFFFFFFFFFFFFFFFFFFFFFFFFFFFFFFFFFFFFFFFFFFFFFFFFFFFFFFFFFFFFFFFFFFFFFFFFFFFFFFFFFFFFFFFCF5F1F8FEF9D07C7A541B05081C5A8E687CC7EBFAF7F1F4FBFAF9FCF6D76A970D041B5D24040704030A09102461CFF7EEF9FDF6F6FBFBF9F9F8F69D6612030FA0FBFCFAF6FDF8D0996C7CBD95521C020204010308104D6CD9EDF6F9FAFCF9F7FFFFFFFFFFFFFFFFFFFFFFFFFFFFFFFFFFFFFFFFFFFFFFFFFFFFFFFFFFFFFFFFFFFFFFFFFFFFFFFFFFFFFFFFFFFFFFFFFFFFFFFFFFFFFFFFFFFFFFFFFFFFFFFFFFFFFFFFFFFFFFFFFFFFFFFFFFFFFFFFFFFFFFFFFFFFFFFFFFFFFFFFFFFFFFFFFFFFFFFFFFFFFFFFFFFFFFFFFFFFFFFFFFFFFFFFFFFFFFFFFFFFFFFFFFFFFFFFFFFFFFFFFFFFFFFFFFFFFFFFFFFFFFFFFFFFFFFFFFFFFFFFFFFFFFFFFFFFFFFCF8FBF7FAFECF8E5D1D0303040A4698475E7DBAEAF5F3FAF9F5F5F7F9E094C92D202F63350901092F5E554A95D3E7EFEDF9FCF6F7FCFAF6F8FDF4DF30050112A3F8F8FBF7FCF9F8F9D8BCBE944F1702040C160807060169C4F3FCF9FAFCFAFBFFFFFFFFFFFFFFFFFFFFFFFFFFFFFFFFFFFFFFFFFFFFFFFFFFFFFFFFFFFFFFFFFFFFFFFFFFFFFFFFFFFFFFFFFFFFFFFFFFFFFFFFFFFFFFFFFFFFFFFFFFFFFFFFFFFFFFFFFFFFFFFFFFFFFFFFFFFFFFFFFFFFFFFFFFFFFFFFFFFFFFFFFFFFFFFFFFFFFFFFFFFFFFFFFFFFFFFFFFFFFFFFFFFFFFFFFFFFFFFFFFFFFFFFFFFFFFFFFFFFFFFFFFFFFFFFFFFFFFFFFFFFFFFFFFFFFFFFFFFFFFFFFFFFFFFFFFFFFFFFFDFCFCF5EDC49C581803030729677840618385CCF6FFFCF4F6FAF7FDF4C97E867D826F6728001863B6B586C1B1A89BC9F0FAFCFDFDFCF6F6F8F08A7316011391E8F8FFFEFCF9FEF9F5FDC098561802062550190007076482C5F2F9FBFCF9FAFFFFFFFFFFFFFFFFFFFFFFFFFFFFFFFFFFFFFFFFFFFFFFFFFFFFFFFFFFFFFFFFFFFFFFFFFFFFFFFFFFFFFFFFFFFFFFFFFFFFFFFFFFFFFFFFFFFFFFFFFFFFFFFFFFFFFFFFFFFFFFFFFFFFFFFFFFFFFFFFFFFFFFFFFFFFFFFFFFFFFFFFFFFFFFFFFFFFFFFFFFFFFFFFFFFFFFFFFFFFFFFFFFFFFFFFFFFFFFFFFFFFFFFFFFFFFFFFFFFFFFFFFFFFFFFFFFFFFFFFFFFFFFFFFFFFFFFFFFFFFFFFFFFFFFFFFFFFFFFDFBFDFEF5B88379602001020B3167766938C79FCFF3FAFAF9FAFCFDF2CB7D6E879AD0A27F2F021F5692BEA56B3148298ED2F0FAFCFAF8F9FCF2F4817D19000957C2F8FBFCFEFEFBFCF9FDE3E1922C01052A5A3417060901517CC1F1FDFBF8F9FFFFFFFFFFFFFFFFFFFFFFFFFFFFFFFFFFFFFFFFFFFFFFFFFFFFFFFFFFFFFFFFFFFFFFFFFFFFFFFFFFFFFFFFFFFFFFFFFFFFFFFFFFFFFFFFFFFFFFFFFFFFFFFFFFFFFFFFFFFFFFFFFFFFFFFFFFFFFFFFFFFFFFFFFFFFFFFFFFFFFFFFFFFFFFFFFFFFFFFFFFFFFFFFFFFFFFFFFFFFFFFFFFFFFFFFFFFFFFFFFFFFFFFFFFFFFFFFFFFFFFFFFFFFFFFFFFFFFFFFFFFFFFFFFFFFFFFFFFFFFFFFFFFFFFFFFFFFFFFBF5F7FBE37C7E611800020302022080CD9BE4EBFAF7F5FAFCF9FAFAD88C707F041977631B02071D33427FAD3A04061094A6CCF1FCFAF6F9F6F4F4E72A0301012997ECF6FAFDF9F5FCFEFAD5A3591A02000526695D0B1706786D80C0F0FDF8F2FFFFFFFFFFFFFFFFFFFFFFFFFFFFFFFFFFFFFFFFFFFFFFFFFFFFFFFFFFFFFFFFFFFFFFFFFFFFFFFFFFFFFFFFFFFFFFFFFFFFFFFFFFFFFFFFFFFFFFFFFFFFFFFFFFFFFFFFFFFFFFFFFFFFFFFFFFFFFFFFFFFFFFFFFFFFFFFFFFFFFFFFFFFFFFFFFFFFFFFFFFFFFFFFFFFFFFFFFFFFFFFFFFFFFFFFFFFFFFFFFFFFFFFFFFFFFFFFFFFFFFFFFFFFFFFFFFFFFFFFFFFFFFFFFFFFFFFFFFFFFFFFFFFFFFFFFFFFFFFEF9F2F5C1765F150100070A090D0D1D7ADFEBFDFCF8FCFBF5F7FDF3BF738003010503622B06041C49503E3D000C02397876B0ECFDFDFEF7EDFAF998741306083F88CDFBFEFBF2FAFDFEF0B970381105102D7EE19E31691502707476C3F6F2FFFFFFFFFFFFFFFFFFFFFFFFFFFFFFFFFFFFFFFFFFFFFFFFFFFFFFFFFFFFFFFFFFFFFFFFFFFFFFFFFFFFFFFFFFFFFFFFFFFFFFFFFFFFFFFFFFFFFFFFFFFFFFFFFFFFFFFFFFFFFFFFFFFFFFFFFFFFFFFFFFFFFFFFFFFFFFFFFFFFFFFFFFFFFFFFFFFFFFFFFFFFFFFFFFFFFFFFFFFFFFFFFFFFFFFFFFFFFFFFFFFFFFFFFFFFFFFFFFFFFFFFFFFFFFFFFFFFFFFFFFFFFFFFFFFFFFFFFFFFFFFEFEFEFEFEFEFEFEFEFAF5F9F0DD7F0E03000302040A1F342C5DD4FCFDF4F9FCFAFCFAF9D3887A011100000302090400070B0707080D0B126D6FB2EDFDF9F9FBF7FFFEF26BEA2C030C6183ADFAFEFCFDFCF9FBFDDBA6632120589DE4E4805A9D2108016F4A84C2AFF6FBF3FFFEFDFEFEFEFFFEFEFFFEFDFDFFFFFFFFFFFFFFFFFFFFFFFFFFFFFFFFFFFFFFFFFFFFFFFFFFFFFFFFFFFFFFFFFFFFFFFFFFFFFFFFFFFFFFFFFFFFFFFFFFFFFFFFFFFFFFFFFFFFFFFFFFFFFFFFFFFFFFFFFFFFFFFFFFFFFFFFFFFFFFFFFFFFFFFFFFFFFFFFFFFFFFFFFFFFFFFFFFFFFFFFFFFFFFFFFFFFFFFFFFFFFFFFFFFFFFFFFFFFFFFFFFFFFFFFFFFFFFFFFFFFFFFFFFFFFFFEFAFBFEFCFCFDFEFCFBF4CF8E66061103000002263442644079E7FEF9FDFEEDF4FDFCFBC3480A00010000000B040001031B413935756B6A5DB8EFFDFDFDFFFFFFFFFEFCFBFA3301000E90FCFCFDFBFBFEFDF1FCC6440D062D77BBEBF7FDEE130606020F42794E38B4DF9AFEFDF2FCFBFDFEFCFCFEFCF6F6FEFEFEFFFFFFFFFFFFFFFFFFFFFFFFFFFFFFFFFFFFFFFFFFFFFFFFFFFFFFFFFFFFFFFFFFFFFFFFFFFFFFFFFFFFFFFFFFFFFFFFFFFFFFFFFFFFFFFFFFFFFFFFFFFFFFFFFFFFFFFFFFFFFFFFFFFFFFFFFFFFFFFFFFFFFFFFFFFFFFFFFFFFFFFFFFFFFFFFFFFFFFFFFFFFFFFFFFFFFFFFFFFFFFFFFFFFFFFFFFFFFFFFFFFFFFFFFFFFFFFFFFFFFFFFFEFEFCFEFBF8FBFEFEF8ECCF4708640A100500002E7B7C270221C0FEFEF8FDFEFEFEFCFBD46101060100080200010B2945485A9EB2DCF9ECF6F9FDFEFCFCFFFFFFFFFEFCFBFA921E131D8CF6FBF8FEFDFEFDF4FCDD6115022D433485F7EEF8850E00070358530D2A79858AFEFBFDFEFEFEF9F7F9FDFEFEFEFDF4FEFFFFFFFFFFFFFFFFFFFFFFFFFFFFFFFFFFFFFFFFFFFFFFFFFFFFFFFFFFFFFFFFFFFFFFFFFFFFFFFFFFFFFFFFFFFFFFFFFFFFFFFFFFFFFFFFFFFFFFFFFFFFFFFFFFFFFFFFFFFFFFFFFFFFFFFFFFFFFFFFFFFFFFFFFFFFFFFFFFFFFFFFFFFFFFFFFFFFFFFFFFFFFFFFFFFFFFFFFFFFFFFFFFFFFFFFFFFFFFFFFFFFFFFFFFFFFFFFFFFFFFFFFFFFFCFEFEFAFBFDFFFDF9D758267D1A03010000153A4B2F156182E0F6F1FBF0F7F9C0B7B31A0200040107010501011E525A4991D4FBFEF9FDFEF2EEF7FEFFFFFFFFFEFCFCFB941F020A58C6FBF8FBFBFEFDCFB4893A140B314D28248BF2A6EE2D0D08020D0C01387139B3EDF8FEFEFDFEFEFEFEFCF9FCFFFDF8FEFFFFFFFFFFFFFFFFFFFFFFFFFFFFFFFFFFFFFFFFFFFFFFFFFFFFFFFFFFFFFFFFFFFFFFFFFFFFFFFFFFFFFFFFFFFFFFFFFFFFFFFFFFFFFFFFFFFFFFFFFFFFFFFFFFFFFFFFFFFFFFFFFFFFFFFFFFFFFFFFFFFFFFFFFFFFFFFFFFFFFFFFFFFFFFFFFFFFFFFFFFFFFFFFFFFFFFFFFFFFFFFFFFFFFFFFFFFFFFFFFFFFFFFFFFFFFFFFFFFFFFFFFEFAFAFEFFFCF8F6EEC7A8B414091705070A23462B090F0E17688BE2FBFCFEFBF7FAB97C5A03030D0C641005060204080F25485458DCB0D4F9FEFBFBFDFFFFFFFFFFFEFCFCFB9B61213E7AC1FEFAF6FEFEFDC692601C03030C0F296771842089170404010008090A140937ACFDFEF2FBFCFDFFFFFDFCFCFCFEFDFEFFFFFFFFFFFFFFFFFFFFFFFFFFFFFFFFFFFFFFFFFFFFFFFFFFFFFFFFFFFFFFFFFFFFFFFFFFFFFFFFFFFFFFFFFFFFFFFFFFFFFFFFFFFFFFFFFFFFFFFFFFFFFFFFFFFFFFFFFFFFFFFFFFFFFFFFFFFFFFFFFFFFFFFFFFFFFFFFFFFFFFFFFFFFFFFFFFFFFFFFFFFFFFFFFFFFFFFFFFFFFFFFFFFFFFFFFFFFFFFFFFFFFFFFFFFFFFFFFFFFFFFFFEFDFDF8FBFAFAFCF9CA7F6E01040B000203304A240306040F36B8A0EEA0DBF4FAF6E68C180A010A166B110301000506050C1513185F4DA2C59887C6EDF6FEFFFFFFFEFDFDFCDCD78CACD7E7FFFDFAF6FEFDB65B1F05080601003EB28B20091C1102010C0702080B0E2B604DFFFFFBFDFBF8FAFCFCFEFDFBF6FFFFFFFFFFFFFFFFFFFFFFFFFFFFFFFFFFFFFFFFFFFFFFFFFFFFFFFFFFFFFFFFFFFFFFFFFFFFFFFFFFFFFFFFFFFFFFFFFFFFFFFFFFFFFFFFFFFFFFFFFFFFFFFFFFFFFFFFFFFFFFFFFFFFFFFFFFFFFFFFFFFFFFFFFFFFFFFFFFFFFFFFFFFFFFFFFFFFFFFFFFFFFFFFFFFFFFFFFFFFFFFFFFFFFFFFFFFFFFFFFFFFFFFFFFFFFFFFFFFFFFFFFFFFFEF6FDFCFDFCF5E8C2AA9E2300010B01020010170E0703010415743BD5AB90C9F3FBC581490401020B22050000080602080B04080418105C4C2A5F86BBEEFEFEFFFFFEFDFDFCFAFDDCEAFBFAFAF7FDFBFEFBD38B411105022B435EAC53560E0405070205033D430406256D25F8F4FBFEFDFCFDFBF8FAFEFBF2FFFFFFFFFFFFFFFFFFFFFFFFFFFFFFFFFFFFFFFFFFFFFFFFFFFFFFFFFFFFFFFFFFFFFFFFFFFFFFFFFFFFFFFFFFFFFFFFFFFFFFFFFFFFFFFFFFFFFFFFFFFFFFFFFFFFFFFFFFFFFFFFFFFFFFFFFFFFFFFFFFFFFFFFFFFFFFFFFFFFFFFFFFFFFFFFFFFFFFFFFFFFFFFFFFFFFFFFFFFFFFFFFFFFFFFFFFFFFFFFFFFFFFFFFFFFFFFFFFFFFFFFFFFFFEF9FEF8FAF8F7F1C28B6C05010904244D141E353E3C220202111E08748879CCF6FCBFA38E00060D020B00134720071C35343A61270A502F47295068B4F4FDFEFFFFFFFEFEFDF4F9F8F5F8FCFCF9FEFFFEFAEDB255140302304842602C5F110D56154E6229484D0608133B6DB8F3D9B9E6F6F9FDF8F6FFFCF4FFFFFFFFFFFFFFFFFFFFFFFFFFFFFFFFFFFFFFFFFFFFFFFFFFFFFFFFFFFFFFFFFFFFFFFFFFFFFFFFFFFFFFFFFFFFFFFFFFFFFFFFFFFFFFFFFFFFFFFFFFFFFFFFFFFFFFFFFFFFFFFFFFFFFFFFFFFFFFFFFFFFFFFFFFFFFFFFFFFFFFFFFFFFFFFFFFFFFFFFFFFFFFFFFFFFFFFFFFFFFFFFFFFFFFFFFFFFFFFFFFFFFFFFFFFFFFFFFFFFFFFFFFFFFEFAFDF8FAFAFCF7BA581E01010B00305C3B615F47422B06020A0300222F61CEF5FAE5B46507040B040900195C25031D4F7997A4626A7D5F6E724D78CFFAFDFDFEFFFFFEFEFDF6F7FBF8F9FCFDFEFCFEFBFAF9CD60130204131B1619132712567366B59E371C1E0507071979459499A1C9C4D8F8FCFAFFFDF8FFFFFFFFFFFFFFFFFFFFFFFFFFFFFFFFFFFFFFFFFFFFFFFFFFFFFFFFFFFFFFFFFFFFFFFFFFFFFFFFFFFFFFFFFFFFFFFFFFFFFFFFFFFFFFFFFFFFFFFFFFFFFFFFFFFFFFFFFFFFFFFFFFFFFFFFFFFFFFFFFFFFFFFFFFFFFFFFFFFFFFFFFFFFFFFFFFFFFFFFFFFFFFFFFFFFFFFFFFFFFFFFFFFFFFFFFFFFFFFFFFFFFFFFFFFFFFFFFFFFFFFFFFFFFEF7FCFCF8FBFBEE951B00000426471C21326040191612080500040051182CB5F3FCD28F4D0803000503020A251303165890959D82894BA66E793D83DEFCFDFDFEFFFFFEFEFEFEFCFCFDFFFDFDFEFDFDF6FBFCCA793D1002011D3915000D5E7980869C5F1A0004050002112E0124477A92A0CDF5FEFDFAFBFAFFFFFFFFFFFFFFFFFFFFFFFFFFFFFFFFFFFFFFFFFFFFFFFFFFFFFFFFFFFFFFFFFFFFFFFFFFFFFFFFFFFFFFFFFFFFFFFFFFFFFFFFFFFFFFFFFFFFFFFFFFFFFFFFFFFFFFFFFFFFFFFFFFFFFFFFFFFFFFFFFFFFFFFFFFFFFFFFFFFFFFFFFFFFFFFFFFFFFFFFFFFFFFFFFFFFFFFFFFFFFFFFFFFFFFFFFFFFFFFFFFFFFFFFFFFFFFFFFFFFFFFFFFFFFEFBFCF8F4FAFAED840203040C34700C1F372B100200030909040102792D0A79DDCD91835E020400010303000005041E617E5C5E622B1186472F167DE1FDFFFFFFFFFFFFFFFFFFFFFFFFFFFFFFFFFFFDF7FEFBC1885A1E0201325E210313877689783F11020127340002080A03060C3178C0EAF9FEFCF9FBFCFEFEFFFFFFFFFFFFFFFFFFFFFFFFFFFFFFFFFFFFFFFFFFFFFFFFFFFFFFFFFFFFFFFFFFFFFFFFFFFFFFFFFFFFFFFFFFFFFFFFFFFFFFFFFFFFFFFFFFFFFFFFFFFFFFFFFFFFFFFFFFFFFFFFFFFFFFFFFFFFFFFFFFFFFFFFFFFFFFFFFFFFFFFFFFFFFFFFFFFFFFFFFFFFFFFFFFFFFFFFFFFFFFFFFFFFFFFFFFFFFFFFFFFFFFFFFFFFFFFFFFFFFFFFFFFFFBF3FAFDFBF1880205080D0D1C0340642708040006090D351203211705208C817C792303000000000601000303133D6453161C04022216030166CBF5FFFFFFFFFFFFFFFFFFFFFFFFFFFFFFFFFFFDF7FFFFDE91370E0403172307050B29BD44832C02253F65752D010212060201298FE3FCFBFBFAFDFAF6FEFEFFFFFFFFFFFFFFFFFFFFFFFFFFFFFFFFFFFFFFFFFFFFFFFFFFFFFFFFFFFFFFFFFFFFFFFFFFFFFFFFFFFFFFFFFFFFFFFFFFFFFFFFFFFFFFFFFFFFFFFFFFFFFFFFFFFFFFFFFFFFFFFFFFFFFFFFFFFFFFFFFFFFFFFFFFFFFFFFFFFFFFFFFFFFFFFFFFFFFFFFFFFFFFFFFFFFFFFFFFFFFFFFFFFFFFFFFFFFFFFFFFFFFFFFFFFFFFFFFFFFFFFFFFFFF4FAFDF9F6DC7304040000070D020C1F180601040300227C3002020502002678671806040002000412050409090307181F162136340E2B2F30469CE9FFFFFFFFFFFFFFFFFFFFFFFFFFFFFFFFFEFAECFDE8D29B2802040200020500010FC42E49393497B9827D780100000307041E6ED3FBFBF7F6FBF2EAFFFFFFFFFFFFFFFFFFFFFFFFFFFFFFFFFFFFFFFFFFFFFFFFFFFFFFFFFFFFFFFFFFFFFFFFFFFFFFFFFFFFFFFFFFFFFFFFFFFFFFFFFFFFFFFFFFFFFFFFFFFFFFFFFFFFFFFFFFFFFFFFFFFFFFFFFFFFFFFFFFFFFFFFFFFFFFFFFFFFFFFFFFFFFFFFFFFFFFFFFFFFFFFFFFFFFFFFFFFFFFFFFFFFFFFFFFFFFFFFFFFFFFFFFFFFFFFFFFFFFFFFFFFFFEFCFBF7F8FAE87E161011000209030100000300070700031422040003190901021B18090E010504010E1F0D06010001010A182A8AC7E94079C6D9D3E2F7FFFFFFFFFFFFFFFFFFFFFFFFFFFFFFFFFDF9FEF8AAA07E210A0B040007110200011F208ECFD3E299838F24040201121D1029A4DFEDF8FEFBF9F8FDFDFEFFFFFFFFFFFFFFFFFFFFFFFFFFFFFFFFFFFFFFFFFFFFFFFFFFFFFFFFFFFFFFFFFFFFFFFFFFFFFFFFFFFFFFFFFFFFFFFFFFFFFFFFFFFFFFFFFFFFFFFFFFFFFFFFFFFFFFFFFFFFFFFFFFFFFFFFFFFFFFFFFFFFFFFFFFFFFFFFFFFFFFFFFFFFFFFFFFFFFFFFFFFFFFFFFFFFFFFFFFFFFFFFFFFFFFFFFFFFFFFFFFFFFFFFFFFFFFFFFFFFFFFFFEF2FDFEFDFA80787A0200010B031307020000000507010208000A15149C4E1A0902040600010A190120B8230000050A11346B7487FEF65B98EAFBF7FBFFFFFFFFFFFFFFFFFFFFFFFFFFFFFFFFFFFFFEFCF8EF962805000E100601030C010F02C047C36F9DDC8561860002022BAD1A0F1D65BFEEF3F4FEFBF3F3FEFFFFFFFFFFFFFFFFFFFFFFFFFFFFFFFFFFFFFFFFFFFFFFFFFFFFFFFFFFFFFFFFFFFFFFFFFFFFFFFFFFFFFFFFFFFFFFFFFFFFFFFFFFFFFFFFFFFFFFFFFFFFFFFFFFFFFFFFFFFFFFFFFFFFFFFFFFFFFFFFFFFFFFFFFFFFFFFFFFFFFFFFFFFFFFFFFFFFFFFFFFFFFFFFFFFFFFFFFFFFFFFFFFFFFFFFFFFFFFFFFFFFFFFFFFFFFFFFFFFFFFFFFEFEFEF6F0F2E3E993010300031C00060105070502090D02060010C61ABCEDD0550A0003050308010C00150A0601002C735F2890F0FDFCF7ECF3FBFEFEFEFFFFFFFFFFFFFFFFFFFFFFFFFFFFFFFFFFFCF1FD9176802B08547D4002000017C023D2F4F8FB7E1D0D0E13030F010D200D0234B1B185B0F1FBF3F8FEFEFFFFFFFFFFFFFFFFFFFFFFFFFFFFFFFFFFFFFFFFFFFFFFFFFFFFFFFFFFFFFFFFFFFFFFFFFFFFFFFFFFFFFFFFFFFFFFFFFFFFFFFFFFFFFFFFFFFFFFFFFFFFFFFFFFFFFFFFFFFFFFFFFFFFFFFFFFFFFFFFFFFFFFFFFFFFFFFFFFFFFFFFFFFFFFFFFFFFFFFFFFFFFFFFFFFFFFFFFFFFFFFFFFFFFFFFFFFFFFFFFFFFFFFFFFFFFFFFFFFFFFFFFEF8F3F1FAFEF98E1402070B49B604040D034B75380003050001102007C469C3DA711500010508010D000106020003071241808795F3FCF4F9FDF3F9FCFCFFFFFFFFFFFFFFFFFFFFFFFFFFFFFFFFFFFEFDFBF09D2E050013190E08010B1614BEF65D211B0C02040705030801060200080C245CABE9FEFCF8FBFCFEFFFFFFFFFFFFFFFFFFFFFFFFFFFFFFFFFFFFFFFFFFFFFFFFFFFFFFFFFFFFFFFFFFFFFFFFFFFFFFFFFFFFFFFFFFFFFFFFFFFFFFFFFFFFFFFFFFFFFFFFFFFFFFFFFFFFFFFFFFFFFFFFFFFFFFFFFFFFFFFFFFFFFFFFFFFFFFFFFFFFFFFFFFFFFFFFFFFFFFFFFFFFFFFFFFFFFFFFFFFFFFFFFFFFFFFFFFFFFFFFFFFFFFFFFFFFFFFFFFFFFFFFFEFDFDFAF2F982160804040C0D210A02040311180B000100010702020C229D6839190403020201030200050102000000060F151C31ACFEFD7F9EE5FAFCFEFFFFFFFFFFFFFFFFFFFFFFFFFFFFFFFFFEFCF8FD987A7822000306050305209C33242E10050201040409478E0208000F010503031257BEFAFDFDFDFBFEFFFFFFFFFFFFFFFFFFFFFFFFFFFFFFFFFFFFFFFFFFFFFFFFFFFFFFFFFFFFFFFFFFFFFFFFFFFFFFFFFFFFFFFFFFFFFFFFFFFFFFFFFFFFFFFFFFFFFFFFFFFFFFFFFFFFFFFFFFFFFFFFFFFFFFFFFFFFFFFFFFFFFFFFFFFFFFFFFFFFFFFFFFFFFFFFFFFFFFFFFFFFFFFFFFFFFFFFFFFFFFFFFFFFFFFFFFFFFFFFFFFFFFFFFFFFFFFFFFFFFFFFFFFFEAF8FB82825C010100020301020247665B663506080000040000030128168057020003050004000602227E3501000304030B082939BE669EDCF6FCFBFBFFFBF3F6FEF0FDFFFFE68EB0FCFBFCF9FCF8FEE1913A0B0101000304010C3E13AD2A0502016C5854BDCCB02A9E2291220802050541B3FAFCFEFDF5FAFEFAFDF7FAFFFFFFFAF4FFFFFEFBF5FCFFFFFEFEFBFAFCF9FEFFFFFFFFFFFFFEFDFDFFFFFFFFFFFFFFFFFFFFFFFFFFFFFFFFFFFFFFFFFFFFFFFFFFFFFFFFFFFFFFFFFFFFFFFFFFFFFFFFFFFFFFFFFFFFFFFFFFFFFFFFFFFFFFFFFFFFFFFFFFFFFFFFFFFFFFFFFFFFFFFFFFFFFFFFFFFFFFFFFFFFFFFFFFFFFFFFFFFFFFFFFFFFFFFFFFFFFEEEF4F7FCCC77260004090401092FA7BBD4D6D0C0B1AF913302060009070209311F020101062F900203000C3013020403010103030701321D2F47514E98B192CB7FF5E8F960D57ABC7291FBFDFCF9FDF6B2958864632500000203090F0802451302030790CFD3F7F4E7C0E99A48386E24020A1A346ECAFAF3FBFDF2FAFEFAF6FAFBF9FBF9FBFEFBFBFAFEFFFCFBFBFCF8F3FCFEFDFDFDFEFFFFFEFEFEFFFFFFFFFFFFFFFFFFFFFFFFFFFFFFFFFFFFFFFFFFFFFFFFFFFFFFFFFFFFFFFFFFFFFFFFFFFFFFFFFFFFFFFFFFFFFFFFFFFFFFFFFFFFFFFFFFFFFFFFFFFFFFFFFFFFFFFFFFFFFFFFFFFFFFFFFFFFFFFFFFFFFFFFFFFFFFFFFFFFFFFFFFFFFFFFFFFFFEF6FDFCEC8B2D0502070403060E41E3F3F9FDF9F5F4F7ECBB9C8F2A07040806040605000207194D01040008080100040400000201030B000C02040703293654A8255C655D13C847442B345B5F625B66989F312F2E3215010104060506011303020276A2D3F2F6FCFAFAF0D36EA1925917000409277BDCFCF6FFFFF2FBF8FAF9F9FBFCFEFBF1FCFAF9E284CCFAF7F7FBF8F9FEFBFDFDFDFDFDFEFFFFFFFFFFFFFFFFFFFFFFFFFFFFFFFFFFFFFFFFFFFFFFFFFFFFFFFFFFFFFFFFFFFFFFFFFFFFFFFFFFFFFFFFFFFFFFFFFFFFFFFFFFFFFFFFFFFFFFFFFFFFFFFFFFFFFFFFFFFFFFFFFFFFFFFFFFFFFFFFFFFFFFFFFFFFFFFFFFFFFFFFFFFFFFFFFFFFFFFFFFFEFDF8F2962C060503030B375E0746E6FDF5F9F9F7F9F7F3FBE9878A2B05050E0904040304060A020805000505050806050500010E54390D090505030A0D0924480E010D0007592003090D0B0A08010E2B520F020006040302305E2300050D070502516E7278B3F9FDFAF5DEA6F5E284260806051B496D84C0F7F9FAFCF1F8F6F7F9F6F8F8F6FCEB7FD13853707FC3FDF9F9FEFBFEFEFDFDFDFEFFFFFFFFFFFFFFFFFFFFFFFFFFFFFFFFFFFFFFFFFFFFFFFFFFFFFFFFFFFFFFFFFFFFFFFFFFFFFFFFFFFFFFFFFFFFFFFFFFFFFFFFFFFFFFFFFFFFFFFFFFFFFFFFFFFFFFFFFFFFFFFFFFFFFFFFFFFFFFFFFFFFFFFFFFFFFFFFFFFFFFFFFFFFFFFFFFFFFFFFFFFEFDF99D8633030204040A32570927A3F1F8FAF9FAFCF6F1F5EC957E2907025C374A470F042166050D0C001B5525080407090000104E330B0A0507023E430F070A0E0306070403020E504646724509080808080200060B08022A532000050D0605010207084CB8F4FDF9F1FBFDFCFCD697825D22020102134F7ECEA7DBA6CAABB9BBA6CB9BE289E1206C3152315EC8F8F7F6F4FDFFFFFEFDFDFEFFFFFFFFFFFFFFFFFFFFFFFFFFFFFFFFFFFFFFFFFFFFFFFFFFFFFFFFFFFFFFFFFFFFFFFFFFFFFFFFFFFFFFFFFFFFFFFFFFFFFFFFFFFFFFFFFFFFFFFFFFFFFFFFFFFFFFFFFFFFFFFFFFFFFFFFFFFFFFFFFFFFFFFFFFFFFFFFFFFFFFFFFFFFFFFFFFFFFFFFFFFFFFF88E8532000104070303050B1A89ECFAFAFCFEFCF9F6EFECA21F04054F847EC1BE4D0B2D72175B100120672905040406030104010701582B423E59493768716B5A111261235441957D7EC37C5330530B040500000A070101020200020F000100090B047DC3B6FAF2EFFCFCFCFEF3F8E4902F040300010A036E2DB0525D38474A376036C60A8A13011C5F3A5D9EC7F9FAF1FFFFFFFFFFFFFFFEFEFEFFFFFFFFFFFFFFFFFFFFFFFFFFFFFFFFFFFFFFFFFFFFFFFFFFFFFFFFFFFFFFFFFFFFFFFFFFFFFFFFFFFFFFFFFFFFFFFFFFFFFFFFFFFFFFFFFFFFFFFFFFFFFFFFFFFFFFFFFFFFFFFFFFFFFFFFFFFFFFFFFFFFFFFFFFFFFFFFFFFFFFFFFFFFFFFFFFFFFEFCEDB02D0101040907050300053AABF1F6F7FBFEFDFAE5C3C03C040C0DB887BFFAFDBA5F4503208A0506030000070B0A090903010D422BB4649EAB7961A0F5FBEAA26567E084D5A2B6ACACBE9CCCA0C86741120704030202040506000100030B020B0A05505659D0ADF8FBF6FAF7F7F9D1915421040207080014076F310503020205042685040701030302030C45AEF3F2FBFBFCFCFDFEFFFFFEFCFDFFFFFFFFFFFFFFFFFFFFFFFFFFFFFFFFFFFFFFFFFFFFFFFFFFFFFFFFFFFFFFFFFFFFFFFFFFFFFFFFFFFFFFFFFFFFFFFFFFFFFFFFFFFFFFFFFFFFFFFFFFFFFFFFFFFFFFFFFFFFFFFFFFFFFFFFFFFFFFFFFFFFFFFFFFFFFFFFFFFFFFFFFFFFFFFFFFFFFEF9B62D020402092B31010002185C85E8FDFDFCFDFCFAD68481340A4550C6F5F7FAFDE1B780020706040E070101001835280808083ACA94B58DB8DFEFF9F9F5F8D875F6FBAEE3EDEBCDD8F7F9F8FBFBF5F9C25A34130D0701060A080000020D4919030D0B03043A6629BBF3F5F3EDEEF8F7FADB7D1F080C0000000B03050105060608060A0509000108050B05166FB3ADCBFAF1F4F6FAFDFEFFFEFCFDFFFFFFFFFFFFFFFFFFFFFFFFFFFFFFFFFFFFFFFFFFFFFFFFFFFFFFFFFFFFFFFFFFFFFFFFFFFFFFFFFFFFFFFFFFFFFFFFFFFFFFFFFFFFFFFFFFFFFFFFFFFFFFFFFFFFFFFFFFFFFFFFFFFFFFFFFFFFFFFFFFFFFFFFFFFFFFFFFFFFFFFEFEFFFEFDFFFEFCF228010002020D56720401063CC6F9F8FCFAFCFBFAFDF0B95A110285DEFAFFFDFCFCAC41310200000C0400000408295F4F01011366B2B2BAE2D1DED0D3D4C9E6BBEAB6E95F9BBAA16C83D0F8F3ECEFFAF1F7F8973E25100019321601010B3CA04012030000030706042AC1EFBEF4FEF6F6F9FAD4761B00020000070601000204253307051109000501113819000000258CDAD1DAD8F6FAFDFEFDFCFDFFFFFFFFFFFFFFFFFFFFFFFFFFFFFFFFFFFFFFFFFFFFFFFFFFFFFFFFFFFFFFFFFFFFFFFFFFFFFFFFFFFFFFFFFFFFFFFFFFFFFFFFFFFFFFFFFFFFFFFFFFFFFFFFFFFFFFFFFFFFFFFFFFFFFFFFFFFFFFFFFFFFFFFFFFFFFFFFFFFFFEFBFCFEFAF6FEFCF1C1540D0201040401020208080F79F1F5FBFEFDF7F5F7CF8290380D26BCF3FDFCFCFCFAE99400000000100200072D200102050200010D0408753E674C595F51AB259D07B3BE22050369565CC1B9B8B4D9F8D4CEA7956C29054B96510D0010AC561439170403000001031262A062E0F7E3D4D9DADAC364121306073D1E09030A095678030816400300051785410405011E606C6C7F6AE9ECF3F8FAFDFEFFFFFFFFFFFFFFFFFFFFFFFFFFFFFFFFFFFFFFFFFFFFFFFFFFFFFFFFFFFFFFFFFFFFFFFFFFFFFFFFFFFFFFFFFFFFFFFFFFFFFFFFFFFFFFFFFFFFFFFFFFFFFFFFFFFFFFFFFFFFFFFFFFFFFFFFFFFFFFFFFFFFFFFFFFFFFFFFFFFFFAFAFCFBFAFCFCEDC343CB2F09050000000104134E66BEF8F6F8FDF9F9F1CCBBB81B122BA4F1EEF4FCFEFEFEAF1304070008010000326F450A060000030801201E0F0E060005053D870C000000010002120C03010001020D2BC68D72BA7624041B21378634000328E571AB442D230102000A429020AFBABD927E9271558FB6913B3D49D7593B18020E120C0F001BB001050801000000080C0401146B4E5EBDC2BED8F4FDFDFFFFFFFFFFFFFFFFFFFFFFFFFFFFFFFFFFFFFFFFFFFFFFFFFFFFFFFFFFFFFFFFFFFFFFFFFFFFFFFFFFFFFFFFFFFFFFFFFFFFFFFFFFFFFFFFFFFFFFFFFFFFFFFFFFFFFFFFFFFFFFFFFFFFFFFFFFFFFFFFFFFFFFFFFFFFFFFFFFFFFBFAFCFDFAF7FCE348C022010007080206132263D0C45FEEFCFCFDFEF2BC5D1006011BDD60B8F1FEFEFDFCF7E08E040D000E19140401120D0001051321212B0A794E2C1D1005050F0E01090F0006060100000A0A000209010C06626E1D0401004E7C3105060006A9BBDF7348B076000101020B06070A1209399B843E50A1E4EDEFE2B7C067A0580F0000010300031D0C01090200090500051C2F0B05160706061E7DDAFDFEFFFFFFFFFFFFFFFFFFFFFFFFFFFFFFFFFFFFFFFFFFFFFFFFFFFFFFFFFFFFFFFFFFFFFFFFFFFFFFFFFFFFFFFFFFFFFFFFFFFFFFFFFFFFFFFFFFFFFFFFFFFFFFFFFFFFFFFFFFFFFFFFFFFFFFFFFFFFFFFFFFFFFFFFFFFFFFFFFFFFFEFEFBFAFEFAF1EEC0201F090100000305074DD17AF9EDEEF0FBF4F8EBCD913308143DC6EAF7F6FCF1F1FAFEC23B00010716A421000609191A06001A685387935D3DA19C511E0D05194A01010C0100080303030303000100050B0A1B0A000F100D0D090B010000231BA1D5BCDE792230081201010000010207111C42606D8A6F7584271C7753180507242E020204074912020E341A0002255E5A0503282A04175199D8F3F9FFFFFFFFFFFFFFFFFFFFFFFFFFFFFFFFFFFFFFFFFFFFFFFFFFFFFFFFFFFFFFFFFFFFFFFFFFFFFFFFFFFFFFFFFFFFFFFFFFFFFFFFFFFFFFFFFFFFFFFFFFFFFFFFFFFFFFFFFFFFFFFFFFFFFFFFFFFFFFFFFFFFFFFFFFFFFFFFFFFFFFF9F7FCF9F4FCF069100702341F02010B2E79C0DCFEF6FAFCF4F8FDCF6E210C10368DF3E7FAFDF6F8F9F9FADF7E00031042270600042C60430304091D439AD0D3D4EDEFE0871D052D82523D3F11311005032A310C070107010C0003010003232D11001233081203011B907F9682736508050600000019060002040C173A6A4A2777280D1A0E0200014E700708000BA3110017763901050B1211020C5D640F052B599BDDFCFFFFFFFFFFFFFFFFFFFFFFFFFFFFFFFFFFFFFFFFFFFFFFFFFFFFFFFFFFFFFFFFFFFFFFFFFFFFFFFFFFFFFFFFFFFFFFFFFFFFFFFFFFFFFFFFFFFFFFFFFFFFFFFFFFFFFFFFFFFFFFFFFFFFFFFFFFFFFFFFFFFFFFFFFFFFFFFFFFFFFDFBFAFAF9FBFDF5940E0200713A0B090A408C9CFAF5FBFDFEDDD5D3BB7A28001B717FDBF1F4FCF5F6FBFBF8B22502010F96180505040E1C130001020D56669EC4E8FAF1F58F1103206FCBD0C560B16B454A6A674740230F1E0604000A090245532111286D1D09060002227758686C1C02020802020207020B0600060A0C1B1606221A050100100900162206020212330648192116050204050404031A23050368B2DBF5FEFFFFFFFFFFFFFFFFFFFFFFFFFFFFFFFFFFFFFFFFFFFFFFFFFFFFFFFFFFFFFFFFFFFFFFFFFFFFFFFFFFFFFFFFFFFFFFFFFFFFFFFFFFFFFFFFFFFFFFFFFFFFFFFFFFFFFFFFFFFFFFFFFFFFFFFFFFFFFFFFFFFFFFFFFFFFFFFFFFFFFEFAFAFDFDFBF4E5450700083016040A0A174EA9FDF9F9FDF7CE9F957941130F3F75B3EBFBFBFEFBFBF1F3DF7B0501050D3A0B020201030A0C06010607201B679BA2B0AEB5933C145CCCF9FEFACDF1D8C7C9A69CBFBC947795452D060404011A29180913290D0102004E19281920210204030200020202030505020203040A01000A060000080A0505050303184606050A7C1F000610252202030B1E330400047CDCFCF9F8FFFFFFFFFFFFFFFFFFFFFFFFFFFFFFFFFFFFFFFFFFFFFFFFFFFFFFFFFFFFFFFFFFFFFFFFFFFFFFFFFFFFFFFFFFFFFFFFFFFFFFFFFFFFFFFFFFFFFFFFFFFFFFFFFFFFFFFFFFFFFFFFFFFFFFFFFFFFFFFFFFFFFFFFFFFFFFFFFFFFFFFAFAFFFCE9BDAB0A020406030A0603081F5BA0F8FDF5FDD8C697725324021B66B8EDFCFEF3FCFFFFF7F8AD453D00030B000601000B08060703000D2D2606236A5A435C56665749BABFB3F8FEF7FDFAFBFDDFDAF8F8E4D9E9BE77221627200C070701050202103C4C7159491C00000103010100000701132F160504030807060005000205062A230B040404237B0202053A0F020315463F02070F444F03264195E4FFFBFAFFFFFFFFFFFFFFFFFFFFFFFFFFFFFFFFFFFFFFFFFFFFFFFFFFFFFFFFFFFFFFFFFFFFFFFFFFFFFFFFFFFFFFFFFFFFFFFFFFFFFFFFFFFFFFFFFFFFFFFFFFFFFFFFFFFFFFFFFFFFFFFFFFFFFFFFFFFFFFFFFFFFFEFDFCFBFDFEFFFDFDFAFBFFFBDF8C43020003000209050004328FC9F2FCF9FEC8A57D472717163A7AE3FFFEFEF4FAFDFFF8D86E3065050505000802000907010003051250490508281E09463B292B489E567CB3FBF8C3E9FAF9F9F6F8FDFCFAF2FCB1777C665E350C02010402030F6BBA7FACBC4B02020000041C380205002C5322071B2E2401050100050205044C3F0E0003174241060500050202010B262C1D241F453C1E61A5D2F3FDFDFEFFFFFFFFFFFFFFFFFFFFFFFFFFFFFFFFFFFFFFFFFFFFFFFFFFFFFFFFFFFFFFFFFFFFFFFFFFFFFFFFFFFFFFFFFFFFFFFFFFFFFFFFFFFFFFFFFFFFFFFFFFFFFFFFFFFFFFFFFFFFFFFFFFFFFFFFFFFFFFFFFFFFFEFDFCFBFDFEFEFCFDFCF9FAFDE2800100060002032B1B062059A9F5F5F5FEF4D380411606174D91B7E9FDF5FDFDFCF7FBF5BC531F3905070F0006020A29170101080D0F302D000B0404002F2B090019341B7949C0C18BB4E9F5F3F4F6FAFDFCF0FAD9D3E0A79D884410060407040476C0ABC3BE4B040200030F2B660200011F32150522433306000634190201033227071C2929650604050C360D0000071F36548274716D6FA1E4FBFDFCFDFFFFFFFFFFFFFFFFFFFFFFFFFFFFFFFFFFFFFFFFFFFFFFFFFFFFFFFFFFFFFFFFFFFFFFFFFFFFFFFFFFFFFFFFFFFFFFFFFFFFFFFFFFFFFFFFFFFFFFFFFFFFFFFFFFFFFFFFFFFFFFFFFFFFFFFFFFFFFFFFFFFFFCF9FFFEFAFAFEFEFBFDFDF7F4FBDC5D000709020501563A09439BCBF2F8F4FDF0D1591302163C82DDFCF7FCF1F9FEFFF5F7F9C15204010202150304001956300403060A060004341B080502020100010202144D084B527A5D9FC8DEF7FBF9FCFBFDF9F6F9FADBD4C8863912020B08067857D1846F32030001080E1441000203000203021131381A030B6B581908050607043F5A1B451230001271210200144C7394E6DDD3D7CFDDF6FDFAFBFEFFFFFFFFFFFFFFFFFFFFFFFFFFFFFFFFFFFFFFFFFFFFFFFFFFFFFFFFFFFFFFFFFFFFFFFFFFFFFFFFFFFFFFFFFFFFFFFFFFFFFFFFFFFFFFFFFFFFFFFFFFFFFFFFFFFFFFFFFFFFFFFFFFFFFFFFFFFFFFFFFFFFF9F1FFFEF7F6FDFFF9FAFAF9E8CDA71C060E29120502372A1C5AB6F9FCF7F9F1D2992301033773A0DAF9F6FEF9F6F6FBFEFAE3932F02070D250006000216372109050D1D19040B733A0B000218190A03132C090004000746162F569AD0D0CECCCEDEFBFEF9FDE2E4B3845D1D0D22130A4D0AAD395D390503030600010B000502000202030A2D4D2D05093E8432100A022522303E0102197805014D1D0401185AA0CEF6FAFFFEFAFCFEFDF9FAFDFDFFFFFFFFFFFFFFFFFFFFFFFFFFFFFFFFFFFFFFFFFFFFFFFFFFFFFFFFFFFFFFFFFFFFFFFFFFFFFFFFFFFFFFFFFFFFFFFFFFFFFFFFFFFFFFFFFFFFFFFFFFFFFFFFFFFFFFFFFFFFFFFFFFFFFFFFFFFFFFFFFFFCF8FDFCFEFCFBFBFBFAFBFECD7142060F02802A0504060A509BC2FDFAFBF2D159390E020D5DB5D1DCF6FAFCF1FBECF6FAFDB13E0205091C7902000307040200000324554E0B0A2A201002004A4D1F0A347405020C000003060300284E4E4C576C7BD1D8CEE5AFAD8956392641661B0C0006491F3A200306020B1F0002000D1E07050B1D141635240107034C2606000052510F000303114F0D00020007131F6CD5FAF1F2FDFEF9FEFEFBFBFDFBFBFFFFFFFFFFFFFFFFFFFFFFFFFFFFFFFFFFFFFFFFFFFFFFFFFFFFFFFFFFFFFFFFFFFFFFFFFFFFFFFFFFFFFFFFFFFFFFFFFFFFFFFFFFFFFFFFFFFFFFFFFFFFFFFFFFFFFFFFFFFFFFFFFFFFFFFFFFFFFFFFFFFEFDF9F6FCFEF3F2FAFDFAFBDC7A13020106590E0000011147ACFCFBF7FDF4A31101030F469EE5FBF9FFFEFCFDF6F9FDF4E9C45D0A0401145E04010E0200050300010A2432000503250F03012B2E10011C3E0005010500021C0E0208090210415F274F5361907382783E183D7C611800030701150B121B14113E7A020600347E1D01295A3A00050D00000879460B03032F320E020A15000106000B13256465A0FAFCFFFDF9FEFDFEFDF8FCFFFEFEFFFFFFFFFFFFFFFFFFFFFFFFFFFFFFFFFFFFFFFFFFFFFFFFFFFFFFFFFFFFFFFFFFFFFFFFFFFFFFFFFFFFFFFFFFFFFFFFFFFFFFFFFFFFFFFFFFFFFFFFFFFFFFFFFFFFFFFFFFFFFFFFFFFFFFFFFFFFFFFFFFFCF9FDF9F5FEFFFCF4F7F2EDC45603100A00000905011071B3CAF2FAFDF2C33A0C0E0A34A2E8F8FFFFFEFDFDFCFEE8F6FD93A47806030802000106000002040F0B060C05000100007552100B05070D0808100000011503117F3808050D10051E37000A374E3D2724295B6A403B08030E160E088B4F79AD62488A3E01080015410F0721372405000001000C3B100115130500040D408902090004187945ADD4D7F5EFFEFEFEFDFDFDFFFDFEFFFBFBFFFFFFFFFFFFFFFFFFFFFFFFFFFFFFFFFFFFFFFFFFFFFFFFFFFFFFFFFFFFFFFFFFFFFFFFFFFFFFFFFFFFFFFFFFFFFFFFFFFFFFFFFFFFFFFFFFFFFFFFFFFFFFFFFFFFFFFFFFFFFFFFFFFFFFFFFFFFFFFFFFFBF4FEFDFBFDFEE3ADCAAA5E2B0103000200000C02007C759BFEFEFEF6C53E00004F6A68C2FDFDFEFEFDF8FEF4FDFEFCFDC46C2604020600010D060005072C734A02020604100B07381D05050262620B000208010900000036170012130105050000022D380B0E07001D351B000321B209A826379379423179B60D0300000701000000000005070104000000100200000103020016410F00000097D953A8FEFCFCFDF9FAFDF7F5FDFEFEFEFEFAF9FFFFFFFFFFFFFFFFFFFFFFFFFFFFFFFFFFFFFFFFFFFFFFFFFFFFFFFFFFFFFFFFFFFFFFFFFFFFFFFFFFFFFFFFFFFFFFFFFFFFFFFFFFFFFFFFFFFFFFFFFFFFFFFFFFFFFFFFFFFFFFFFFFFFFFFFFFFFFFFFFFFEFCF6F8FDF3B15D201C22130303060000030006017F6D86D5F8FEF2BA390600030A54A8E7FAFAFDFCFDFCFEFDF5FDFEF9F7DC66030102060103169A1E000E2F20020006110100010602050502262C0B020726B3000007040D0B6A554D6905000105040309054E642D0403022E9022490C42110938220D0C1F3601020C000201000040693D06000100000000010001000403020708041AADB4C4E4F8CAE2FEFCFCFDFBFDFDFDFCFDFEFCFBFDFEFEFFFFFFFFFFFFFFFFFFFFFFFFFFFFFFFFFFFFFFFFFFFFFFFFFFFFFFFFFFFFFFFFFFFFFFFFFFFFFFFFFFFFFFFFFFFFFFFFFFFFFFFFFFFFFFFFFFFFFFFFFFFFFFFFFFFFFFFFFFFFFFFFFFFEFEFFFAF6F9FBFBFBF9FCF8FBA98B641F0505060500040701029AAAACD9C9D9F2F5FEEFA515022FB0ABC7EAFBFEFEFEFEFEFEFFFEFAFCF8FBFDA4330300030B0212053F0F0E07060702020618635D00020A0300060805020201114F0011A12009053D848C7D630103010C0002031F2710010000103611090C0C02000A0601040706010401000400050121291507000017773901793778355A5510012FA0B8F1F5F1FAFDF7FBFEFEFEFEFEFEFEFEFEFEFEFEFEFEFFFFFFFFFFFFFFFFFFFFFFFFFFFFFFFFFFFFFFFFFFFFFFFFFFFFFFFFFFFFFFFFFFFFFFFFFFFFFFFFFFFFFFFFFFFFFFFFFFFFFFFFFFFFFFFFFFFFFFFFFFFFFFFFFFFFFFFFFFFFFFFFFFFFFFFEFBFDFDFCFDFEFDFCFCF1B4A58E3D2E0F03020609000101109E58DCF3F9F5F8FCF6E4AF4803001B3E8BDEFBFFFFFFFFFFFFFFFEFDFBFDFDF89834090203030200040510030301030300000210382B00764E130102024A64160A040D0805410F02050D383B27300200030501466E290A04010102060803000201000B050002030307309D000007020104070704032476996A882CA746A05394BB9EA7B5EDEEFEFEF8FDFCFEFFFFFFFFFFFFFFFFFFFFFFFFFFFFFFFFFFFFFFFFFFFFFFFFFFFFFFFFFFFFFFFFFFFFFFFFFFFFFFFFFFFFFFFFFFFFFFFFFFFFFFFFFFFFFFFFFFFFFFFFFFFFFFFFFFFFFFFFFFFFFFFFFFFFFFFFFFFFFFFFFFFFFFFFFFFFFFFFFFFFFEFDFEFEF8F5F9FDFDEEAA778F440E0D06010005072E781F8BEBAEDEFAFFFEFBFBFCF1AA0D020A3091BDECFDFFFFFFFFFFFFFFFEFEFDF1FBF4CD560704010307000002020B0507041F554F0B040B09004420070001002B3A0B03030E030012040505070F0C09090209060001302E100301010101237022020704000406070308416B1A3E070007010703020302010D4CB8B4E0A6E3B9DEBACEE9ECF3F3FAEEF9FEFDFEFEFEFFFFFFFFFFFFFFFFFFFFFFFFFFFFFFFFFFFFFFFFFFFFFFFFFFFFFFFFFFFFFFFFFFFFFFFFFFFFFFFFFFFFFFFFFFFFFFFFFFFFFFFFFFFFFFFFFFFFFFFFFFFFFFFFFFFFFFFFFFFFFFFFFFFFFFFFFFFFFFFFFFFFFFFFFFFFFFFFFFFFFFFCFBFDF9F7FCFCF9E99A323C14020904060401347C73866978E6F1FCFBF5F7F3B56A4C0001073DD6EFF9FDFFFFFFFFFFFFFFFFFEEEF593CD9F34040301030B011193190702040117362D0C110001090C0A0D0601040C0A070D01188A080400020D5840020502010D3C6F070C0C0502040202041950120103020B6F350A05042E542E73060B030207020002031C5DA8DFEEFDEEFCF1F9F4F3F9FEFDFAF6F5FCFFFFFCFDFEFFFFFFFFFFFFFFFFFFFFFFFFFFFFFFFFFFFFFFFFFFFFFFFFFFFFFFFFFFFFFFFFFFFFFFFFFFFFFFFFFFFFFFFFFFFFFFFFFFFFFFFFFFFFFFFFFFFFFFFFFFFFFFFFFFFFFFFFFFFFFFFFFFFFFFFFFFFFFFFFFFFFFFFFFFFFFFFFFFFFFFFBF6FBFDFDFDECBA724708050405053D3847473F7A88642315E8F9FBFBF8F7F4C1590A000524A6F1FAFBFDFFFFFFFFFFFFFFFFFEF3F4A8828B450807040106080B6E14010003030709040A0E01090606090A08020304040909030F68070000050A4235050204070A2E5405014440433C1A0203060F000800010A53250402000914215701070702022B4E435593CCF8FBFCF6FEFEEFF9F7F9FAF7F9FAF8FDFEF9FDFCFDFEFFFFFFFFFFFFFFFFFFFFFFFFFFFFFFFFFFFFFFFFFFFFFFFFFFFFFFFFFFFFFFFFFFFFFFFFFFFFFFFFFFFFFFFFFFFFFFFFFFFFFFFFFFFFFFFFFFFFFFFFFFFFFFFFFFFFFFFFFFFFFFFFFFFFFFFFFFFFFFFFFFFFFFFFFFFFFFFFFFFFFDFCF8FBFCFAFAE08E394404060009026C998F5B1C315E190A02E6F9FBF7FBF8BE804E030A2A7694F3F9FAFDFFFFFFFFFFFFFFFEF7FEEFF19F874108062A5601070707020B0403050503184336090D41410E090602010206090003070B0106541E05070A0600020604093C48376D7A82571A0202010103120001020C040202040700050C0602060701274859A2F1F7F5FDF8F7F5F9F3FAF7FAFAF7F8FAF7FDF6F6FDFEFCFEFFFFFFFFFFFFFFFFFFFFFFFFFFFFFFFFFFFFFFFFFFFFFFFFFFFFFFFFFFFFFFFFFFFFFFFFFFFFFFFFFFFFFFFFFFFFFFFFFFFFFFFFFFFFFFFFFFFFFFFFFFFFFFFFFFFFFFFFFFFFFFFFFFFFFFFFFFFFFFFFFFFFFFFFFFFFFFFFFFFFF7FCFEFEFAF8FBDF893F5605072D60657E95641E030409195869EDFDFDF6F4C2823F0602154FC67FF1F7FBFEFFFFFFFFFFFFFFFDF5FEF6F8F26C0E0200316C06000802000E07000203011D4D400E0854490D0B0B342F0F060705030A0500087423050402060701082B4A61B8815553552B0402040410535D180200080403020507000309080300223E294784CAF8F7F3FBF1FEF4F8FDFCF7F9FBFCFDFAF8FCF3FAFEFFF9FCFFFFFFFFFFFFFFFFFFFFFFFFFFFFFFFFFFFFFFFFFFFFFFFFFFFFFFFFFFFFFFFFFFFFFFFFFFFFFFFFFFFFFFFFFFFFFFFFFFFFFFFFFFFFFFFFFFFFFFFFFFFFFFFFFFFFFFFFFFFFFFFFFFFFFFFFFFFFFFFFFFFFFFFFFFFFFFFFFFFFF5FCFAFCFCFAF3BD520B0B02277A9BA26421050203070A1791ABEEFCFDFCF29E805E0408278AECF8FCF8FBFDFFFFFFFFFFFFFFFDF9FEFDF7FB690A030202050B030003000103040402020504020703070102010D584C1625391005090102041409050602070C030E4C936AA767220A0501000002071A8587270401371B0401042338030501020341AA97C8F8F7F5FBFDFBF3FBFEFDFDFBF3F6F8FBFEFDFEFDF9FEF8FDF9FCFFFFFFFFFFFFFFFFFFFFFFFFFFFFFFFFFFFFFFFFFFFFFFFFFFFFFFFFFFFFFFFFFFFFFFFFFFFFFFFFFFFFFFFFFFFFFFFFFFFFFFFFFFFFFFFFFFFFFFFFFFFFFFFFFFFFFFFFFFFFFFFFFFFFFFFFFFFFFFFFFFFFFFFFFFFFFFFFFFFFFEFEF5F8FDEFC57722020F05577116020100020908000D320A05A0BAEEFAF6CE5A040E2663CCFEEFF8FBFEFFFFFFFFFFFFFFFFFDF5FBF8F9D44E04000303000307020B330D0014301C03060202000000040A04010706024778130200063F0A0335170504082027010D4E7E1D031E2C100000152910070D0B0000097C3D0802094E7C040204134570E0FBFCFEFFFFFDFCFDFDF7FEFCFCFEFFFFFEFEFDFCFEFCF9FEF1FCFFFFFFFFFFFFFFFFFFFFFFFFFFFFFFFFFFFFFFFFFFFFFFFFFFFFFFFFFFFFFFFFFFFFFFFFFFFFFFFFFFFFFFFFFFFFFFFFFFFFFFFFFFFFFFFFFFFFFFFFFFFFFFFFFFFFFFFFFFFFFFFFFFFFFFFFFFFFFFFFFFFFFFFFFFFFFFFFFFFFFFFFFCF9FDFCFCDB957168010201192507000100020B0D01299916010104BDDCE56603040D4C9D66FEF6FDF5FBFFFFFFFFFFFFFFFFFEFBF5F5FB7E1F1A1803000306010D4AA4200948926512000103020203050E030101010303010506000CBC180A933F04041A637804030B0005015F65230201458330010C040000361301000101020102142E39EDF9FCFBFDFFFFFFFFFFFFFFFFFFFFFFFFFFFFFFFFFFFFFFFFFFFFFFFFFFFFFFFFFFFFFFFFFFFFFFFFFFFFFFFFFFFFFFFFFFFFFFFFFFFFFFFFFFFFFFFFFFFFFFFFFFFFFFFFFFFFFFFFFFFFFFFFFFFFFFFFFFFFFFFFFFFFFFFFFFFFFFFFFFFFFFFFFFFFFFFFFFFFFFFFFFFFFFFFFFFFFFFFFFFFFFFFFFFFFFFFFFE9F6FFF8F8DF7A08090D088185160403000A0B0100081E05000000A262A6620607000F5EEFFEFEFEEFFAFFFFFFFFFFFFFFFFFDEFFDFDFDFB62786A020308000537B11E0B2B857F83590108100003230D02000105070602030903010B0E010101040204050106130B00020003000D160E080602070207010302983D070207060001083FEB4ADFF9FCFEFEFFFFFFFFFFFFFFFFFFFFFFFFFFFFFFFFFFFFFFFFFFFFFFFFFFFFFFFFFFFFFFFFFFFFFFFFFFFFFFFFFFFFFFFFFFFFFFFFFFFFFFFFFFFFFFFFFFFFFFFFFFFFFFFFFFFFFFFFFFFFFFFFFFFFFFFFFFFFFFFFFFFFFFFFFFFFFFFFFFFFFFFFFFFFFFFFFFFFFFFFFFFFFFFFFFFFFFFFFFFFFFFFFFFFFFFFFEF7FBFEF0EDDE84140F010F7C7114000705060D0A0036CD3307050002A0F16E0609041774FDF8F9FDFEFEFEFFFFFFFFFFFFFFFFFEFEFBFCF75A0A0D06080C000000000000001E7F5A03080F02070DA36A0F000904000B0E02071A050000050A0C030302060D060718070000060168752C000100040C010F000C0304050207090B161C59F4F4FDFEEFF7FEFFFFFFFFFFFFFFFFFFFFFFFFFFFFFFFFFFFFFFFFFFFFFFFFFFFFFFFFFFFFFFFFFFFFFFFFFFFFFFFFFFFFFFFFFFFFFFFFFFFFFFFFFFFFFFFFFFFFFFFFFFFFFFFFFFFFFFFFFFFFFFFFFFFFFFFFFFFFFFFFFFFFFFFFFFFFFFFFFFFFFFFFFFFFFFFFFFFFFFFFFFFFFFFFFFFFFFFFFFFFFFFFFFFFFFFFF4FBFAF8FAFEFDF4CE740B0005071204000C0D030A11022BCE141406010F8C5F63711261F2F7FDFCFCFEFBFDFFFFFFFFFFFFFFFFFEEEFBF85B2A1600000000020D00080009070F0B03000204020B071003220B050404649635000105000600000003100B00061F5A91080000000406131B110503050C010007001C0D050C080A6BF6F6F6F7FDF6F5FBFCFEFFFFFFFFFFFFFFFFFFFFFFFFFFFFFFFFFFFFFFFFFFFFFFFFFFFFFFFFFFFFFFFFFFFFFFFFFFFFFFFFFFFFFFFFFFFFFFFFFFFFFFFFFFFFFFFFFFFFFFFFFFFFFFFFFFFFFFFFFFFFFFFFFFFFFFFFFFFFFFFFFFFFFFFFFFFFFFFFFFFFFFFFFFFFFFFFFFFFFFFFFFFFFFFFFFFFFFFFFFFFFFFFFFFFFFFFFDFEEFF6FEFAD6812E1F01000D010105025C732D498E1B06191BB62D09010F151F1A0A5BFEFDFAFDF9F6FBFEFFFFFFFFFFFFFFFFFEF8FDFEE7DC630A0E0800010100181103000100030308357C758BE4E2D875866D270D5D8D1B0615030D030A021186620003080A12000509130801576B2A04062D993C0000180794578F7674EFFBF5F9FCEBFCFDFEFBFDFFFFFFFFFFFFFFFFFFFFFFFFFFFFFFFFFFFFFFFFFFFFFFFFFFFFFFFFFFFFFFFFFFFFFFFFFFFFFFFFFFFFFFFFFFFFFFFFFFFFFFFFFFFFFFFFFFFFFFFFFFFFFFFFFFFFFFFFFFFFFFFFFFFFFFFFFFFFFFFFFFFFFFFFFFFFFFFFFFFFFFFFFFFFFFFFFFFFFFFFFFFFFFFFFFFFFFFFFFFFFFFFFFFFFFFFF4FBFBFDFCFAF5E1B86D040003010C04031A1C09152B0F030FB242140300000204061F78F4F6FDFDF6FBFEFFFFFFFFFFFFFFFFFFFEFEEDF8FD91351002010000062E9F140009071611022C7B9597CE6B99EAE6889362141C2E0522A632CC340800061D17080603070104015B7F1D0021210C04010E2D0C0016A3372A9F7B89E4F8F2FCFDF9FAFEFBEFF8FEFFFFFFFFFFFFFFFFFFFFFFFFFFFFFFFFFFFFFFFFFFFFFFFFFFFFFFFFFFFFFFFFFFFFFFFFFFFFFFFFFFFFFFFFFFFFFFFFFFFFFFFFFFFFFFFFFFFFFFFFFFFFFFFFFFFFFFFFFFFFFFFFFFFFFFFFFFFFFFFFFFFFFFFFFFFFFFFFFFFFFFFFFFFFFFFFFFFFFFFFFFFFFFFFFFFFFFFFFFFFFFFFFFFFFFFFFBFDF8F3F9FDFCFAE896160600000A020504060A366D3609043D09992B020302000D5BD5F8F1FEF7F7FDF9FCFFFFFFFFFFFFFFFDF3F9F8F7FCCFAB5A0A020409020E3119030323654A2D83DFDCBB831C97EFF3DB7C2F0D406E121A62C9C8A937140600080A040E170802081F270E0607030101010104011206321112B2DAE1FDFBF9FBFEFBF5FEF5F7F5FAFFFFFFFFFFFFFFFFFFFFFFFFFFFFFFFFFFFFFFFFFFFFFFFFFFFFFFFFFFFFFFFFFFFFFFFFFFFFFFFFFFFFFFFFFFFFFFFFFFFFFFFFFFFFFFFFFFFFFFFFFFFFFFFFFFFFFFFFFFFFFFFFFFFFFFFFFFFFFFFFFFFFFFFFFFFFFFFFFFFFFFFFFFFFFFFFFFFFFFFFFFFFFFFFFFFFFFFFFFFFFFFFFFFFFFFFFBFBFAFAF8FAFDFAF0C36515080002100816140518597C150002034513010100192764F9F8F8F3FBFCF8FAFDFFFFFFFFFFFFFFFDF7FBEDF8FAA7663005010404000306090704102A46789FB7E99E7C4578DFF6F9902B0427420822A5ED6D4B8A2C07030303114763040612150C221D0A0C12110E0500040706020084CEEAFAFDFCFEFBF7FEFAF5F8FDFBFDFFFFFFFFFFFFFFFFFFFFFFFFFFFFFFFFFFFFFFFFFFFFFFFFFFFFFFFFFFFFFFFFFFFFFFFFFFFFFFFFFFFFFFFFFFFFFFFFFFFFFFFFFFFFFFFFFFFFFFFFFFFFFFFFFFFFFFFFFFFFFFFFFFFFFFFFFFFFFFFFFFFFFFFFFFFFFFFFFFFFFFFFFFFFFFFFFFFFFFFFFFFFFFFFFFFFFFFFFFFFFFFFFFFFFFFFFCFCFAFCFCF9FBFDEAB27C4C0A00050A04445F2E09223D0B02000201050100064E6D7CFDFAFCF2FEFCF9FDFFFFFFFFFFFFFFFFFEFDFEF4FCF9CBA6540502040900030D1D0902030D4395C3D0ECC37B272EAFF2FCC8570A355809116FF38F0F4B16080600051227370509537C629370375B726D73360505000A03069D92C1F9FCFCFDFAF9FEFBF8FDFCFDFEFFFFFFFFFFFFFFFFFFFFFFFFFFFFFFFFFFFFFFFFFFFFFFFFFFFFFFFFFFFFFFFFFFFFFFFFFFFFFFFFFFFFFFFFFFFFFFFFFFFFFFFFFFFFFFFFFFFFFFFFFFFFFFFFFFFFFFFFFFFFFFFFFFFFFFFFFFFFFFFFFFFFFFFFFFFFFFFFFFFFFFFFFFFFFFFFFFFFFFFFFFFFFFFFFFFFFFFFFFFFFFFFFFFFFFFFFFFDF8FBFFFCFCFDF0C68F430306050805336E571707010003120205050602053883C5FBFDFDF9FEF6FBF9FCFFFFFFFFFFFFFFFEFEFBF2F8FAFDFC82040408112D011569190301195993C2F5F8E76A081F85D1F6B9510E284F21001FCA850E040807060204070405061B76BAC5BF9E9AC5D4D3D797370C010501054186C2D5E4FBF8FCFEFEFCFEFFFBFCFEFFFFFFFFFFFFFFFFFFFFFFFFFFFFFFFFFFFFFFFFFFFFFFFFFFFFFFFFFFFFFFFFFFFFFFFFFFFFFFFFFFFFFFFFFFFFFFFFFFFFFFFFFFFFFFFFFFFFFFFFFFFFFFFFFFFFFFFFFFFFFFFFFFFFFFFFFFFFFFFFFFFFFFFFFFFFFFFFFFFFFFFFFFFFFFFFFFFFFFFFFFFFFFFFFFFFFFFFFFFFFFFFFFFFFFFFFFFCFAFDFFFFFEFCFAF4C55F161A170D05072D38130101000113030506070213166DF7F7FDFDF7FEF7FCF7FBFFFFFFFFFFFFFFFEFBFBF2F8FDFDFA9519050C1975010F591505011863A5CBF6FDE87C265F6FB2E9C0540C0329520D0C597411010803051C1A020404073B80A8EECFC6EAFBFBFBF7E7941821080301048C987BB5F5F9FDFCFFFEFDFDFCFDFEFFFFFFFFFFFFFFFFFFFFFFFFFFFFFFFFFFFFFFFFFFFFFFFFFFFFFFFFFFFFFFFFFFFFFFFFFFFFFFFFFFFFFFFFFFFFFFFFFFFFFFFFFFFFFFFFFFFFFFFFFFFFFFFFFFFFFFFFFFFFFFFFFFFFFFFFFFFFFFFFFFFFFFFFFFFFFFFFFFFFFFFFFFFFFFFFFFFFFFFFFFFFFFFFFFFFFFFFFFFFFFFFFFFFFFFFFEFEFEFEFCFCFDFBFCFAF0CB523C4013010102020203000305280200081A21605B80FBF7FBFDF5FBFEFEFCFDFFFFFFFFFFFFFFFEFAFBFAFCF6F8E8AC4B080A17610503030303011566C8F6FBFDF2BD879184C4EEFC9426124066150009660F04031C1942410501050A3175A5F1FDF8FBFDFDF9F7FFD83C88131F0D07412E48A2DFF9FCF9FCFEFBF9FDFFFFFFFFFFFFFFFFFFFFFFFFFFFFFFFFFFFFFFFFFFFFFFFFFFFFFFFFFFFFFFFFFFFFFFFFFFFFFFFFFFFFFFFFFFFFFFFFFFFFFFFFFFFFFFFFFFFFFFFFFFFFFFFFFFFFFFFFFFFFFFFFFFFFFFFFFFFFFFFFFFFFFFFFFFFFFFFFFFFFFFFFFFFFFFFFFFFFFFFFFFFFFFFFFFFFFFFFFFFFFFFFFFFFFFFFFFFFFDFDFEFCF8F9FCFBFDFAFAF68743411703020201010303060D75260919678ACBC8BEFAF1DEFDFEF9FAFEFCFEFFFFFFFFFFFFFFFFFEF8F7F9EFF6C1945B09040404070703050204146FD7FBF7FBFBF5FCACCBF5FBFCC35436694F08042C02020305563E393807030E171C2CA6F1FAF5F7FBFCFAF9FCE796DE2E6520040302468DB8F6FDFAFAFEFBF8FDFDFEFFFFFFFFFFFFFFFFFFFFFFFFFFFFFFFFFFFFFFFFFFFFFFFFFFFFFFFFFFFFFFFFFFFFFFFFFFFFFFFFFFFFFFFFFFFFFFFFFFFFFFFFFFFFFFFFFFFFFFFFFFFFFFFFFFFFFFFFFFFFFFFFFFFFFFFFFFFFFFFFFFFFFFFFFFFFFFFFFFFFFFFFFFFFFFFFFFFFFFFFFFFFFFFFFFFFFFFFFFFFFFFFFFFFFFFFFCFCFBFBFBFCFEFDFDFBF0F4901A0F0A05010304020814030A7E8C356ED4FBF9EADFDCC6A4F7E7FDF2FCF7FBFFFFFFFFFFFFFFFEFEFEFCFAFDFAA86D440301020B0006080103111D66CBFAF7F4F9F9FAF7FCFCF3F8C85538836B060E77070701075530060700020F3459035ED7F6FDF8FAFEFDFAF7F8EA958A9E20010A0802329CEAF9FEFEFFFFF8FDFAFCFFFFFFFFFFFFFFFFFFFFFFFFFFFFFFFFFFFFFFFFFFFFFFFFFFFFFFFFFFFFFFFFFFFFFFFFFFFFFFFFFFFFFFFFFFFFFFFFFFFFFFFFFFFFFFFFFFFFFFFFFFFFFFFFFFFFFFFFFFFFFFFFFFFFFFFFFFFFFFFFFFFFFFFFFFFFFFFFFFFFFFFFFFFFFFFFFFFFFFFFFFFFFFFFFFFFFFFFFFFFFFFFFFFFFFFFFFFFF4F9FFFCFCFCF9F9F5F8BD453D3F19091015152A4F1D1773E4A0E3F7FBF0B07E7F7793D297FBF9FDF7FBFFFFFFFFFFFFFFFEFCFEF4EFF9F7C994510900000403160300001B5E99CBEDF2FDFAEDFBFAFBF4FAFBC45F1B3358101082040104000A07020100010B31590E33AEE3F9FDFDFCF9F5F7FDEA6AF88A1B03010A09357DA0DDFDFBFEFEF9FDFBFDFFFFFFFFFFFFFFFFFFFFFFFFFFFFFFFFFFFFFFFFFFFFFFFFFFFFFFFFFFFFFFFFFFFFFFFFFFFFFFFFFFFFFFFFFFFFFFFFFFFFFFFFFFFFFFFFFFFFFFFFFFFFFFFFFFFFFFFFFFFFFFFFFFFFFFFFFFFFFFFFFFFFFFFFFFFFFFFFFFFFFFFFFFFFFFFFFFFFFFFFFFFFFFFFFFFFFFFFFFFFFFFFFFFFFFFFFEFEF3F9FDEBE8FEF9EFFBEAFED6CCCA554075887881998082FDFEFEFEEEF3AA764B072B6F868DFBFEFEFEFFFFFFFFFFFFFFFFFFFEFEF8FDF4F8BE47000102000418930F0802256696CAF9FDFAFFFEFDE8FAFEFDF3BF59407C7C740F7B05080300000C100702090F0800002D9FE8FBF6F6FBF9F5FCF7EEFDFEFD510B05020607338CB1FBF8FAFEF9FDFFFFFFFFFFFFFFFFFFFFFFFFFFFFFFFFFFFFFFFFFFFFFFFFFFFFFFFFFFFFFFFFFFFFFFFFFFFFFFFFFFFFFFFFFFFFFFFFFFFFFFFFFFFFFFFFFFFFFFFFFFFFFFFFFFFFFFFFFFFFFFFFFFFFFFFFFFFFFFFFFFFFFFFFFFFFFFFFFFFFFFFFFFFFFFFFFFFFFFFFFFFFFFFFFFFFFFFFFFFFFFFFFFFFFFFFFFFFF4FCF4F9E9BCBFF2FFFEFAFAFEEFFCC57BA4E8FEFDD99DEAFFFFFFF9FDFFFFAD683F0B031242FFFFFAFAFDFFFFFFFFFFFFFFFFFFFEF6FDFF99C7B747070402040008721B6D23022F9EC0B9FDFAFFFEF4FEF3FAF8FCC05132699DCA0F800400000A0C04040403060000000081C1B5DFFAFCFBFEF8F1FFFEFBFEFF2E0B05000607368C9DFFFFFFFFF9FDFFFFFFFFFFFFFFFFFFFFFFFFFFFFFFFFFFFFFFFFFFFFFFFFFFFFFFFFFFFFFFFFFFFFFFFFFFFFFFFFFFFFFFFFFFFFFFFFFFFFFFFFFFFFFFFFFFFFFFFFFFFFFFFFFFFFFFFFFFFFFFFFFFFFFFFFFFFFFFFFFFFFFFFFFFFFFFFFFFFFFFFFFFFFFFFFFFFFFFFFFFFFFFFFFFFFFFFFFFFFFFFFFFFFFFFFFFFFFEFEF8F7EAB87B6A86B1F1FDFCFCFAB3BCE8FDFEFEFCF7FDFFFEFAFDFBFDFEF7B67E5C275139FEF9FBFDFAFDFFFFFFFFFFFFFFFFFEFDFDFEA276824C0500000000001097721A01155393C5F4FAFEFCFBFEFAF4F9FDFAC68DAAD581791000000201000201000409000500072571A6DFFAFEFEFEFEFDFFFEFEFE7A7B2200060603082A7A95FEF9FEF9F2FAFFFFFFFFFFFFFFFFFFFFFFFFFFFFFFFFFFFFFFFFFFFFFFFFFFFFFFFFFFFFFFFFFFFFFFFFFFFFFFFFFFFFFFFFFFFFFFFFFFFFFFFFFFFFFFFFFFFFFFFFFFFFFFFFFFFFFFFFFFFFFFFFFFFFFFFFFFFFFFFFFFFFFFFFFFFFFFFFFFFFFFFFFFFFFFFFFFFFFFFFFFFFFFFFFFFFFFFFFFFFFFFFFFFFFFFFFFF9FCFDFDE596381B498FABF1FCF6FCF3F5FCFFFFFFFFFFFFFFFEFEFDF9F7FCFEECAA6F9CD7A0F6FDFEFDFDFEFFFFFFFFFFFFFFFEFDFEFCF8EEAC834508040004040302562807000C3093F2FDFBF7FDF8F8FDFDFEF9F5F2ECBF5B2D5E010C01000A050143450A3D6A02113093B0E5FAF6F9FFFFFFFFFFFFFFFF98D9B3340100070B0735937DFDEDFAFDF1F8FFFFFFFFFFFFFFFFFFFFFFFFFFFFFFFFFFFFFFFFFFFFFFFFFFFFFFFFFFFFFFFFFFFFFFFFFFFFFFFFFFFFFFFFFFFFFFFFFFFFFFFFFFFFFFFFFFFFFFFFFFFFFFFFFFFFFFFFFFFFFFFFFFFFFFFFFFFFFFFFFFFFFFFFFFFFFFFFFFFFFFFFFFFFFFFFFFFFFFFFFFFFFFFFFFFFFFFFFFFFFFFFFFFFFFFFFDFEFAFAE9B48081ACC3B2FBFEF1F6FDF8FBFFFFFFFFFFFFFFFEFBFAFCFDFBF8F7CBA9EDF9E6F8FEFFFEFCFEFFFFFFFFFFFFFFFEFAF5FCFDA885480F03040001050000251A62284D99CAF0F9FDFBFDFDFDFEFEFEFBFAC486866D0B1408030C73240912453D072C4A0004157AAAD7F6FCFDFFFFFFFFFFFFFFFFE9EBE4AE3E070803072675B1F6F9FEEDF5FDFFFFFFFFFFFFFFFFFFFFFFFFFFFFFFFFFFFFFFFFFFFFFFFFFFFFFFFFFFFFFFFFFFFFFFFFFFFFFFFFFFFFFFFFFFFFFFFFFFFFFFFFFFFFFFFFFFFFFFFFFFFFFFFFFFFFFFFFFFFFFFFFFFFFFFFFFFFFFFFFFFFFFFFFFFFFFFFFFFFFFFFFFFFFFFFFFFFFFFFFFFFFFFFFFFFFFFFFFFFFFFFFFFFFFFFFFCFEFCF9F6ECE0E3EFF0E0FFFFFAF9F8F9FCFFFFFFFFFFFFFFFEFDF3F6F9FBFCF9EDE9FAF2FCFDFEF8FDFCFEFFFFFFFFFFFFFFFEFAFAFDFEB54F100000052B66010002080A5282BDE6F4F8F8F7FDF9FFFFF5F9FDFEFECE693D560A07030008642457371516070D17040013588DADE8FAFBFFFFFFFFFFFFFFFFF4F7F99F2E0506070B0B3093FBF6FDF8F9FCFFFFFFFFFFFFFFFFFFFFFFFFFFFFFFFFFFFFFFFFFFFFFFFFFFFFFFFFFFFFFFFFFFFFFFFFFFFFFFFFFFFFFFFFFFFFFFFFFFFFFFFFFFFFFFFFFFFFFFFFFFFFFFFFFFFFFFFFFFFFFFFFFFFFFFFFFFFFFFFFFFFFFFFFFFFFFFFFFFFFFFFFFFFFFFFFFFFFFFFFFFFFFFFFFFFFFFFFFFFFFFFFFFFFFFFFFDFEFEFBFCFFFDF7F3FBF8F6FAFEFDF4FAFEFFFFFFFFFFFFFFFFFFF9FAF8FCFFF6F5F9F2F8FDF4FEF6F8FCFEFFFFFFFFFFFFFFFEFCECF9F9A075340503062956050101021B8AD5EFFAFAF7F9F2FEFCFEFAF5FBFEFEF9C459111C04040E00021E0F522E103E3E010302000B34417FDEFCFDFFFFFFFFFFFFFFFFF9F6FCAD823E04060A296C4CF7FAF7F7F7FCFFFFFFFFFFFFFFFFFFFFFFFFFFFFFFFFFFFFFFFFFFFFFFFFFFFFFFFFFFFFFFFFFFFFFFFFFFFFFFFFFFFFFFFFFFFFFFFFFFFFFFFFFFFFFFFFFFFFFFFFFFFFFFFFFFFFFFFFFFFFFFFFFFFFFFFFFFFFFFFFFFFFFFFFFFFFFFFFFFFFFFFFFFFFFFFFFFFFFFFFFFFFFFFFFFFFFFFFFFFFFFFFFFFFFFFFFFFCFBFCFEFCFAF9F6FDFDF4FBF9FDF4F8FDFFFFFFFFFFFFFFFEF2FCFDFAFCFBF6CFBBF4F5FDF5FEF9FAFDFFFFFFFFFFFFFFFFFDF7EEFBAA8985410802020D180A0400175D83D7FAFCF8F1F4F7FBFEF8FAFAFCFBFCFAA42E0406030D6F100100021A141340420002000401100D387FBCEAFFFFFFFFFFFFFFFFFAF6FCE3A07A440409276637FAFEFAF9F8FCFFFFFFFFFFFFFFFFFFFFFFFFFFFFFFFFFFFFFFFFFFFFFFFFFFFFFFFFFFFFFFFFFFFFFFFFFFFFFFFFFFFFFFFFFFFFFFFFFFFFFFFFFFFFFFFFFFFFFFFFFFFFFFFFFFFFFFFFFFFFFFFFFFFFFFFFFFFFFFFFFFFFFFFFFFFFFFFFFFFFFFFFFFFFFFFFFFFFFFFFFFFFFFFFFFFFFFFFFFFFFFFFFFFFFFFFFDF7F9FDFFFBF9FDFFFFFEFAFEF5FCF9F8FCFFFFFFFFFFFFFFFEF3FCFBFBFAF7F8A470DBB8F9FCFDFDFEFDFEFFFFFFFFFFFFFFFDF6FBF8958450190400011D54060200167189D8FBF7F7F6F4F6FAFEF4FCFEF6F6FCF8911902060A127711020004060708151A0801010601000610379CE8FFFFFFFFFFFFFFFFFAFBF8FCB07C470606091E22A4B7EBFCFAFCFFFFFFFFFFFFFFFFFFFFFFFFFFFFFFFFFFFFFFFFFFFFFFFFFFFFFFFFFFFFFFFFFFFFFFFFFFFFFFFFFFFFFFFFFFFFFFFFFFFFFFFFFFFFFFFFFFFFFFFFFFFFFFFFFFFFFFFFFFFFFFFFFFFFFFFFFFFFFFFFFFFFFFFFFFFFFFFFFFFFFFFFFFFFFFFFFFFFFFFFFFFFFFFFFFFFFFFFFFFFFFFFFFFFFFFFFDF8FCFDFBFAFBFCFCF9FFFEF8F9FAFDFCFDFFFFFFFFFFFFFFFEFEFEFCF6BD9FA463378068F7F7FAFEFAFBFDFFFFFFFFFFFFFFFEFCFBF7A0531A020001022B780100060245BEF5F8F5F6FAF9F5FEFEF5FEFEF2F7F7CE67130303175026010202060E3A2801060A030207253B0F2E4FA2E9FFFFFFFFFFFFFFFFFFF8F5FCE6995E3801041C4D2752DBFAFCFDFFFFFFFFFFFFFFFFFFFFFFFFFFFFFFFFFFFFFFFFFFFFFFFFFFFFFFFFFFFFFFFFFFFFFFFFFFFFFFFFFFFFFFFFFFFFFFFFFFFFFFFFFFFFFFFFFFFFFFFFFFFFFFFFFFFFFFFFFFFFFFFFFFFFFFFFFFFFFFFFFFFFFFFFFFFFFFFFFFFFFFFFFFFFFFFFFFFFFFFFFFFFFFFFFFFFFFFFFFFFFFFFFFFFFFFFFFFFFFFCF9FAFFFDFAF6FFFEF4FDFAFDFFFFFFFFFFFFFFFFFFFEFAFDFEF190583B120A1C33B5BAF8FEF4FBFEFEFEFEFFFEFDFDFDFEF7C4822C0C040102010B1D0501143A2BC1FBFBFBF7F4FCFCFBFAF9FEFCF8FCF5B44A0A05031B810B030004040B634A02040305020A45854E899AA9DDFCFCFFFEFEFFFFFFFCF4FBFFFBE9D883010A34820126ADEEFCFEFFFFFFFFFEFEFCFFFFFFFFFFFFFFFFFFFFFFFFFFFFFFFFFFFFFFFFFFFFFFFFFFFFFFFFFFFFFFFFFFFFFFFFFFFFFFFFFFFFFFFFFFFFFFFFFFFFFFFFFFFFFFFFFFFFFFFFFFFFFFFFFFFFFFFFFFFFFFFFFFFFFFFFFFFFFFFFFFFFFFFFFFFFFFFFFFFFFFFFFFFFFFFFFFFFFFFFFFFFFFFFFFFFFFFFFFFFFFFFFFFFFFFFFFFFFFFFFFFFFFFFFFFFFFFFFFFFFFFFFFFFFEF6FBFCF7E5D2670100030D297CFAFBF7FBFCFDFBF9FBFCFBFAFAFEF47C97490F23210400010511361B945DCBFBFDFCFAFAFEFCF6F6FDFEF8FCFEDBD586180303071D0308020107031613020406010205196DCAF3F0EDF6F8F7FCFDFBF9F7FEF8FDFAFDFAFAF3A62A080E11020D4BD7FBFBFCFDFEFEFDFCF8FFFFFFFFFFFFFFFFFFFFFFFFFFFFFFFFFFFFFFFFFFFFFFFFFFFFFFFFFFFFFFFFFFFFFFFFFFFFFFFFFFFFFFFFFFFFFFFFFFFFFFFFFFFFFFFFFFFFFFFFFFFFFFFFFFFFFFFFFFFFFFFFFFFFFFFFFFFFFFFFFFFFFFFFFFFFFFFFFFFFFFFFFFFFFFFFFFFFFFFFFFFFFFFFFFFFFFFFFFFFFFFFFFFFFFFFFFFFFFFFFFFFFFFFFFFFFFFFFFFFFFFFFFFFFEFDF8F7F6FBF68D1803153954ECFBF3F8FAFBFDFAF5F5FAFDFCFFFEEDD85D110673690500040D1CAA0841C9F3FDFCFCFDFEFDFBFCFDF9FCEEFAFD92967718080404030A0400040301000506030703010D005CBFDCF5FCF9F9FAF7FBFBEDE7FCFDFEF1F9FEF9F9EFA21D0002000050D5FAF8F7FBFDFCFCFEFDFDFEFFFFFFFFFFFFFFFFFFFFFFFFFFFFFFFFFFFFFFFFFFFFFFFFFFFFFFFFFFFFFFFFFFFFFFFFFFFFFFFFFFFFFFFFFFFFFFFFFFFFFFFFFFFFFFFFFFFFFFFFFFFFFFFFFFFFFFFFFFFFFFFFFFFFFFFFFFFFFFFFFFFFFFFFFFFFFFFFFFFFFFFFFFFFFFFFFFFFFFFFFFFFFFFFFFFFFFFFFFFFFFFFFFFFFFFFFFFFFFFFFFFFFFFFFFFFFFFFFFFFFFFFFEF7FDFBFDFDF9E2832667D3EEF5FDFEFDFDF4F6FCFDFDFCF5F4FEFCF464170101110E000102010C210F1FD5EAF6FEFEFEFDFAF7FEFDFCF7F6FCFDF29033100000000319141A040C0600070A0B0F01010100051B76DAF3F2F5F9FEFCFEFCFBF5F7FCFCFDFBF6F7F0DDA3000A010126C2F8FEFEFEFEFCFBFEFDF5FEFFFFFFFFFFFFFFFFFFFFFFFFFFFFFFFFFFFFFFFFFFFFFFFFFFFFFFFFFFFFFFFFFFFFFFFFFFFFFFFFFFFFFFFFFFFFFFFFFFFFFFFFFFFFFFFFFFFFFFFFFFFFFFFFFFFFFFFFFFFFFFFFFFFFFFFFFFFFFFFFFFFFFFFFFFFFFFFFFFFFFFFFFFFFFFFFFFFFFFFFFFFFFFFFFFFFFFFFFFFFFFFFFFFFFFFFFFFFFFFFFFFFFFFFFFFFFFFFFFFFFFFFFEF8FDFEFEF0E98B7ADC9C54FEFEF7FDEDEEF9FEFBF4F4FEFEFDF9EDF6DE49090D0D090600000306000022CCF5FBF6F6FBFCF6F1FBEDECFDFFFEFCF4F48D110303001AD806E32D060207061667880C110900000152CCF9FAF1F3F6ECFBECE4FDFEF8F9FDFEFEE88E5EB4060002090348C7F6F8F7F8F9F8FDFEFCFEFFFFFFFFFFFFFFFFFFFFFFFFFFFFFFFFFFFFFFFFFFFFFFFFFFFFFFFFFFFFFFFFFFFFFFFFFFFFFFFFFFFFFFFFFFFFFFFFFFFFFFFFFFFFFFFFFFFFFFFFFFFFFFFFFFFFFFFFFFFFFFFFFFFFFFFFFFFFFFFFFFFFFFFFFFFFFFFFFFFFFFFFFFFFFFFFFFFFFFFFFFFFFFFFFFFFFFFFFFFFFFFFFFFFFFFFFFFFFFFFFFFFFFFFFFFFFFFFFFFFFFFFFEFEF8F9F4F9FEF0EEF8EFF1FDFEF1FDFEFEF2F2FBF9F1F3FEFDFBFC4A1E1007017673130E020800000530CA6CAEF3F7ECEFFEFCFAFCEDFDFFFDF8FB9928040506092CF006D03EA14B0400060B050000000500000E186F7B82CCFCFDFDFDFDFAFEF8FCFEF2F3FCED9F4CB7050000033BC9FEFEFAF9FDFBF8FEFDFEFFFFFFFFFFFFFFFFFFFFFFFFFFFFFFFFFFFFFFFFFFFFFFFFFFFFFFFFFFFFFFFFFFFFFFFFFFFFFFFFFFFFFFFFFFFFFFFFFFFFFFFFFFFFFFFFFFFFFFFFFFFFFFFFFFFFFFFFFFFFFFFFFFFFFFFFFFFFFFFFFFFFFFFFFFFFFFFFFFFFFFFFFFFFFFFFFFFFFFFFFFFFFFFFFFFFFFFFFFFFFFFFFFFFFFFFFFFFFFFFFFFFFFFFFFFFFFFFFFFFFFFFFFFFFEF9FBF4EAF1FAF9FDFAEEF4FCFE9C809B9DC0F0FDFAF7F5FCF0E82F050703171502020201021000052EB5F1F7F8FAFDFBF7FAE4FCF5F8FD83A2781B538A0724AE3FD3F7E0F0510005070809000009006E891F01161861967FBEFAFEFEFCFAFDFBF7F9FDFAF1DF9F25000010033FC3EEF6FEFEFEFCF6FEFEFEFFFFFFFFFFFFFFFFFFFFFFFFFFFFFFFFFFFFFFFFFFFFFFFFFFFFFFFFFFFFFFFFFFFFFFFFFFFFFFFFFFFFFFFFFFFFFFFFFFFFFFFFFFFFFFFFFFFFFFFFFFFFFFFFFFFFFFFFFFFFFFFFFFFFFFFFFFFFFFFFFFFFFFFFFFFFFFFFFFFFFFFFFFFFFFFFFFFFFFFFFFFFFFFFFFFFFFFFFFFFFFFFFFFFFFFFFFFFFFFFFFFFFFFFFFFFFFFFFFFFFFFFFFFFEFF9F8F9FAFCFCFAF9F6FAF6FEF3AA727D505CBAFAFEFAF6EF56DA441304000303010200010A1FAC24B3F3FDFBFCFEFEFBF7FCF8F8FAF6FDD98C3306131C0844EEC3FAF6EB709B4E0102010200010B00172F5B6A2902131E49CCFCF9FEFDF2FCFEF6F8FDFEFCF9EBC4A40003001D2A6BC4F0F4F8FCF8FEFEFEFFFFFFFFFFFFFFFFFFFFFFFFFFFFFFFFFFFFFFFFFFFFFFFFFFFFFFFFFFFFFFFFFFFFFFFFFFFFFFFFFFFFFFFFFFFFFFFFFFFFFFFFFFFFFFFFFFFFFFFFFFFFFFFFFFFFFFFFFFFFFFFFFFFFFFFFFFFFFFFFFFFFFFFFFFFFFFFFFFFFFFFFFFFFFFFFFFFFFFFFFFFFFFFFFFFFFFFFFFFFFFFFFFFFFFFFFFFFFFFFFFFFFFFFFFFFFFFFFFFFFFFFFEFDF9FDFAF5F9F2EEFCFAF5FEFBF3F9AC3A26144ABDFAF9FEFCF8C8661C060104030405050901070A5DB5F2FBFCFDFEFCFBFDFDFDFBFCFEFCFAF8DC6A11010435B6F8F8EFFAF9CADC60050401020000020602051D250D000339A4EEFAEDFEFEFBF3F9FBFAFAFDFBF8F7C956981105060743B8FAFBFCFDFAFEF9FEFFFFFFFFFFFFFFFFFFFFFFFFFFFFFFFFFFFFFFFFFFFFFFFFFFFFFFFFFFFFFFFFFFFFFFFFFFFFFFFFFFFFFFFFFFFFFFFFFFFFFFFFFFFFFFFFFFFFFFFFFFFFFFFFFFFFFFFFFFFFFFFFFFFFFFFFFFFFFFFFFFFFFFFFFFFFFFFFFFFFFFFFFFFFFFFFFFFFFFFFFFFFFFFFFFFFFFFFFFFFFFFFFFFFFFFFFFFFFFFFFFFFFFFFFFFFFFFFFFFFFFFFFEFDFEFCF9F398565A93F7EFFEF9F9FEAF2B050651C2FBFCFDFAFDFBB12F0500020102040506080F23AAF4FDFDFDFEFFFDFDFFFFFEFAF9FEFFF7FCB24307010440E5FBFEF7F7FDF6F7C3500B0A020200000604030806023D5A436EDAFBF7FBFEFBF9F2FAFBFBFDFCFAF7E38E460A0400023AACF6FFFBF7FDFEF7FEFFFFFFFFFFFFFFFFFFFFFFFFFFFFFFFFFFFFFFFFFFFFFFFFFFFFFFFFFFFFFFFFFFFFFFFFFFFFFFFFFFFFFFFFFFFFFFFFFFFFFFFFFFFFFFFFFFFFFFFFFFFFFFFFFFFFFFFFFFFFFFFFFFFFFFFFFFFFFFFFFFFFFFFFFFFFFFFFFFFFFFFFFFFFFFFFFFFFFFFFFFFFFFFFFFFFFFFFFFFFFFFFFFFFFFFFFFFFFFFFFFFFFFFFFFFFFFFFFFFFFFFFFEFEF3FAF8EE65041364C38FFEFBFBF7DD82273292E3FBFAFEFCF486731C050502010101060A1F949B71EEFDFEFFFFFFFFFFFFFFFEFCFAF5FDF9F2C65910050646E2FBFBF9FAFBF9FDF6C28F51000202050242551603002A3B438DE2FDFDFCFCF5F8F7FCFDFDF8F6FDFCFBE389100300041D70CFFCFCF6F3FEFBFEFFFFFFFFFFFFFFFFFFFFFFFFFFFFFFFFFFFFFFFFFFFFFFFFFFFFFFFFFFFFFFFFFFFFFFFFFFFFFFFFFFFFFFFFFFFFFFFFFFFFFFFFFFFFFFFFFFFFFFFFFFFFFFFFFFFFFFFFFFFFFFFFFFFFFFFFFFFFFFFFFFFFFFFFFFFFFFFFFFFFFFFFFFFFFFFFFFFFFFFFFFFFFFFFFFFFFFFFFFFFFFFFFFFFFFFFFFFFFFFFFFFFFFFFFFFFFFFFFFFFFFFFFEFAF5FCA07433020A3D88A7F2FEF8F8FDDBA1AADFF9F8FBFDFBF69D862405040404070A0A0B248EE393F1F8FCFFFFFFFFFFFFFFFDF8FDF5F7F1B4833F0E032B7C8CF5FBF9FEFCFDFDFAFBE982000002040336440E070F0D04225D8DDBF5FEFCF8FAFDF3F6FCF7F3FAFEFDFBDE691501032077CFF4FAF8F4FCFFFFFFFFFFFFFFFFFFFFFFFFFFFFFFFFFFFFFFFFFFFFFFFFFFFFFFFFFFFFFFFFFFFFFFFFFFFFFFFFFFFFFFFFFFFFFFFFFFFFFFFFFFFFFFFFFFFFFFFFFFFFFFFFFFFFFFFFFFFFFFFFFFFFFFFFFFFFFFFFFFFFFFFFFFFFFFFFFFFFFFFFFFFFFFFFFFFFFFFFFFFFFFFFFFFFFFFFFFFFFFFFFFFFFFFFFFFFFFFFFFFFFFFFFFFFFFFFFFFFFFFFFFFFFEF9FDF4987337050A40A6F3F0F9FAFCFAF9FDFCFAF8F9FFFCF5FCFD9927060203033B490703258398EFFDF9FDFFFFFFFFFFFFFFFDF7FAFBF7BA8A3B0F080B367D94F6FCFAFCFBFBFCFAF6DF800B0601000302040509393C1907042183E5F3FEF4FCFCF1F4FCFCF9FAFEFBF4DF71150C031C60A0CCEEF9F9FBFBFEFFFFFFFFFFFFFFFFFFFFFFFFFFFFFFFFFFFFFFFFFFFFFFFFFFFFFFFFFFFFFFFFFFFFFFFFFFFFFFFFFFFFFFFFFFFFFFFFFFFFFFFFFFFFFFFFFFFFFFFFFFFFFFFFFFFFFFFFFFFFFFFFFFFFFFFFFFFFFFFFFFFFFFFFFFFFFFFFFFFFFFFFFFFFFFFFFFFFFFFFFFFFFFFFFFFFFFFFFFFFFFFFFFFFFFFFFFFFFFFFFFFFFFFFFFFFFFFFFFFFFFFFFFFFFCEEF8F67419478FBAEEFBF7F6FBF8FDFFFBF8F9FAFAFCF4EEF6351E0C0401024C58020833D782EAF8FCFEFFFFFFFFFFFFFFFFFDF3FBFBAE8D45140A2C6089FEF9F4FAFBFAF9FBF8ECBD80500E05150005070B063B4723091A3A40D4E6F8EFFAF7FBFDFBFCFFFFFEFBF3EF5E0B0F05001D5899D5FAF8FEF9FEFFFFFFFFFFFFFFFFFFFFFFFFFFFFFFFFFFFFFFFFFFFFFFFFFFFFFFFFFFFFFFFFFFFFFFFFFFFFFFFFFFFFFFFFFFFFFFFFFFFFFFFFFFFFFFFFFFFFFFFFFFFFFFFFFFFFFFFFFFFFFFFFFFFFFFFFFFFFFFFFFFFFFFFFFFFFFFFFFFFFFFFFFFFFFFFFFFFFFFFFFFFFFFFFFFFFFFFFFFFFFFFFFFFFFFFFFFFFFFFFFFFFFFFFFFFFFFFFFFFFFFFFFFFFF9F8FAEC9B74B6CA92F4F5FEF0FCFBFEFFFDFDFDF8F5FCEBB9DA2D5F2505040704000515378EF0FEF1F8FDFFFFFFFFFFFFFFFEFAEDF6F8FA781D040A46ADE9FDF6EFF8FEFEFCFAF9F0B4A4BC372E5A000A0306070503020028603BBAACE4EEF6F6F8FBF8F9FDFCFDFCFCF4D52B0A07032A5C87C2F6FEFDFEFEFFFFFFFFFFFFFFFFFFFFFFFFFFFFFFFFFFFFFFFFFFFFFFFFFFFFFFFFFFFFFFFFFFFFFFFFFFFFFFFFFFFFFFFFFFFFFFFFFFFFFFFFFFFFFFFFFFFFFFFFFFFFFFFFFFFFFFFFFFFFFFFFFFFFFFFFFFFFFFFFFFFFFFFFFFFFFFFFFFFFFFFFFFFFFFFFFFFFFFFFFFFFFFFFFFFFFFFFFFFFFFFFFFFFFFFFFFFFFFFFFFFFFFFFFFFFFFFFFFFFFFFFFEFCFCF4F3F9F8F8F8F5F8FAEDFCFEFDF3FAFFFAFAFDFFFFF8D7889E267536090C080202011C4C15A4FBF1FCFFFFFFFFFFFFFFFFFDF3CFCFEBEC6B19072B5168AFF9F9FDEFFBFDF4F9F4F6F3EAD5989FAB0A020501000408050107111F6E1FC2A3E2F1F3F6FBFFFBF7FAF7F9EFED6A3E1B0A0D1F5DB6E6F9FAFDFEFFFFFFFFFFFFFFFFFFFFFFFFFFFFFFFFFFFFFFFFFFFFFFFFFFFFFFFFFFFFFFFFFFFFFFFFFFFFFFFFFFFFFFFFFFFFFFFFFFFFFFFFFFFFFFFFFFFFFFFFFFFFFFFFFFFFFFFFFFFFFFFFFFFFFFFFFFFFFFFFFFFFFFFFFFFFFFFFFFFFFFFFFFFFFFFFFFFFFFFFFFFFFFFFFFFFFFFFFFFFFFFFFFFFFFFFFFFFFFFFFFFFFFFFFFFFFFFFFFFFFFFFFEF3F3FAFBF8FBF9F8FDFDFCF6F4FDFAE9F7FDFDF9F9FEFEFDFDEB762103000006040102073C860312F7F3FCFEFFFFFFFFFFFFFFFCED827A9BA953100247BABE77EDEBF9FCF9FBF4F5FBFBFAE7AC93F3B438090004010303030706000305118E0FBFB8DFFDFEFDFEFDFBFDFDEFEBC18D2A020000105AC7ECFDFEFEFFFFFFFFFFFFFFFFFFFFFFFFFFFFFFFFFFFFFFFFFFFFFFFFFFFFFFFFFFFFFFFFFFFFFFFFFFFFFFFFFFFFFFFFFFFFFFFFFFFFFFFFFFFFFFFFFFFFFFFFFFFFFFFFFFFFFFFFFFFFFFFFFFFFFFFFFFFFFFFFFFFFFFFFFFFFFFFFFFFFFFFFFFFFFFFFFFFFFFFFFFFFFFFFFFFFFFFFFFFFFFFFFFFFFFFFFFFFFFFFFFFFFFFFFFFFFFFFFFFFFFFFFFFFFCF4FAFBFCFFFCF8FAFEFEFDF7F9FEFAF6FAFCFCFDFBF5FEF9BD3E060001000001030200000228F8FFFBFDFEFEFFFEFEFFFEF5CAAE520A000603032280E5FEFEFEF7F7FDF0FAFAFCFFF8F5FCF9F9D26E0B0008030107050001041005010100CB3F9EF7FCF8F8FEFBF2FEFDCE63443C1503042471B7E5F6FEFEFFFFFFFEFEFFFEFEFFFFFFFFFFFFFFFFFFFFFFFFFFFFFFFFFFFFFFFFFFFFFFFFFFFFFFFFFFFFFFFFFFFFFFFFFFFFFFFFFFFFFFFFFFFFFFFFFFFFFFFFFFFFFFFFFFFFFFFFFFFFFFFFFFFFFFFFFFFFFFFFFFFFFFFFFFFFFFFFFFFFFFFFFFFFFFFFFFFFFFFFFFFFFFFFFFFFFFFFFFFFFFFFFFFFFFFFFFFFFFFFFFFFFFFFFFFFFFFFFFFFFFFFFFFFF3FBFDFBFCFFFEFCF9F5FCFDE7BCF2F9FBFEFAFAFEFEFBF4F37515020201000B0A000308000223F8FFFBFAF9FCFFFCFBFFFCE48C4326060000093073C3F8FBF5FEFEFEF5FCFEFEFEFFFEFEFEFEF3A95033060D04000707000204104C0E00008C32A7FAFEFBF8FAFBF0FDF2FDFCF89627030400104086DDFCFEFFFFFFFCFDFFFCFBFFFFFFFFFFFFFFFFFFFFFFFFFFFFFFFFFFFFFFFFFFFFFFFFFFFFFFFFFFFFFFFFFFFFFFFFFFFFFFFFFFFFFFFFFFFFFFFFFFFFFFFFFFFFFFFFFFFFFFFFFFFFFFFFFFFFFFFFFFFFFFFFFFFFFFFFFFFFFFFFFFFFFFFFFFFFFFFFFFFFFFFFFFFFFFFFFFFFFFFFFFFFFFFFFFFFFFFFFFFFFFFFFFFFFFFFFFFFFFFFFFFFFFFFFEF2FCF3FBFBFCFDFCFCF9FAE398764BC6CDEAFCFDFEF7EFFBFEF9BF3305005229000000001A4C0922F8B8E8FBFBFCFFFEFEFFFEFCF4DA4E090B05004D748EE1FEFEFDF2FAFDFEFFFFFFFFFFFFFFFFFEF6DF9D1A010008060E10030A277E2104060D03A9FEF7F6FCFBFEFCF5FCF5B8C7733C1F0D0825698CBAFDFBF5F5F7F7FBFFFEFEFFFFFFFFFFFFFFFFFFFFFFFFFFFFFFFFFFFFFFFFFFFFFFFFFFFFFFFFFFFFFFFFFFFFFFFFFFFFFFFFFFFFFFFFFFFFFFFFFFFFFFFFFFFFFFFFFFFFFFFFFFFFFFFFFFFFFFFFFFFFFFFFFFFFFFFFFFFFFFFFFFFFFFFFFFFFFFFFFFFFFFFFFFFFFFFFFFFFFFFFFFFFFFFFFFFFFFFFFFFFFFFFFFFFFFFFFFFFFFFFFFFFFFFFFEFDFEF6FBFBFCFDFCFCF1CD8159123C6360CAFBF9FBFDFCFEF1A5791A0200744A0B020628777E02128B64E1FEFDFDFFFDFAFBFEF9E68A13011708003395E2F7F6F2FEFDFEFEFFFFFFFFFFFFFFFFFFFFFEFBBA4928270F0209090001071A0A00011105609DD3F7F6FCF7F9FCF7CB829CA97C2C02010B2F73DAFEFEFDFDFDFDFEFCF8F9FFFFFFFFFFFFFFFFFFFFFFFFFFFFFFFFFFFFFFFFFFFFFFFFFFFFFFFFFFFFFFFFFFFFFFFFFFFFFFFFFFFFFFFFFFFFFFFFFFFFFFFFFFFFFFFFFFFFFFFFFFFFFFFFFFFFFFFFFFFFFFFFFFFFFFFFFFFFFFFFFFFFFFFFFFFFFFFFFFFFFFFFFFFFFFFFFFFFFFFFFFFFFFFFFFFFFFFFFFFFFFFFFFFFFFFFFFFFFFFFFFFFFFFFFFFFF8FAFEFEFEF6F2FBF3C36657170E2F81D9FAF9F8FAF8F8ED98942B07031B110A070533681F105B75CDF5FEFBFCFEFDF9F9FEF7C742040105023A69B0F7FEFCFBFEFAFDFEFFFFFFFFFFFFFFFFFBF0F4FEC97B68734418080101070A020906000417534095F8F3F9F8EEFEFCE57C362931382408050E2FB4FEFBFBF9FBFEFFFCFAFDFFFFFFFFFFFFFFFFFFFFFFFFFFFFFFFFFFFFFFFFFFFFFFFFFFFFFFFFFFFFFFFFFFFFFFFFFFFFFFFFFFFFFFFFFFFFFFFFFFFFFFFFFFFFFFFFFFFFFFFFFFFFFFFFFFFFFFFFFFFFFFFFFFFFFFFFFFFFFFFFFFFFFFFFFFFFFFFFFFFFFFFFFFFFFFFFFFFFFFFFFFFFFFFFFFFFFFFFFFFFFFFFFFFFFFFFFFFFFFFFFFFFFFFFFEF6FAFEF5FAFAECD0CCF3EAB7531E312CB8F5FDFDFBFBFCEFBB83380F0301000005040111220424CBD5F0F9FDF8F9FDFEFDFDFEF2A62B0D41280E3E6EA9E6FAFFFFFDF3FDFEFFFFFFFFFFFFFFFFFDF9FAFECA9A9CB9852D020000030503064C11010859334FC4FBFAF5F7F6FBC73B0801264F330301021374DDFCFDFCF6F9FAFBFCFFFFFFFFFFFFFFFFFFFFFFFFFFFFFFFFFFFFFFFFFFFFFFFFFFFFFFFFFFFFFFFFFFFFFFFFFFFFFFFFFFFFFFFFFFFFFFFFFFFFFFFFFFFFFFFFFFFFFFFFFFFFFFFFFFFFFFFFFFFFFFFFFFFFFFFFFFFFFFFFFFFFFFFFFFFFFFFFFFFFFFFFFFFFFFFFFFFFFFFFFFFFFFFFFFFFFFFFFFFFFFFFFFFFFFFFFFFFFFFFFFFFFFFFFFFEFAF6F2F6B6A6D3A781FAFDEEA56BB198DAFAEFF3FEFDFEDC6C60200144270801332F070B1D5A79FAFEF9FAFDFAFBFEFEFEFEFAE8AE390955393B639AD1F3F9F7FBFEF9FEFCFEFFFFFFFFFFFFFFFDF9F7FEEADCDFE5914C2D160800000107762103002818438FC2F5F9FEFAFCCC5316031122150003153E497BDEFBFEFBFCFCFAFBFFFFFFFFFFFFFFFFFFFFFFFFFFFFFFFFFFFFFFFFFFFFFFFFFFFFFFFFFFFFFFFFFFFFFFFFFFFFFFFFFFFFFFFFFFFFFFFFFFFFFFFFFFFFFFFFFFFFFFFFFFFFFFFFFFFFFFFFFFFFFFFFFFFFFFFFFFFFFFFFFFFFFFFFFFFFFFFFFFFFFFFFFFFFFFFFFFFFFFFFFFFFFFFFFFFFFFFFFFFFFFFFFFFFFFFFFFFFFFFFFFFFFFFFFFFEFCF8F6BC8255796B5DE7BBFCEDD0F2E7F1F6DAD0E4E3C2933B561E04682E070149470907258391F9FEF0FAFEFDFEFDFCFCFBEECD8F4304291A6CB4DEF3FAFCF9FBFDFBFEF9FDFFFFFFFFFFFFFFFEFAF9FDFCFBFDFAB07D4F1A0502080005381102010403447F8BC0F9F2FEFCECB4511302050A06021C5236609AD9F5FBFDFFFBFBFFFFFFFFFFFFFFFFFFFFFFFFFFFFFFFFFFFFFFFFFFFFFFFFFFFFFFFFFFFFFFFFFFFFFFFFFFFFFFFFFFFFFFFFFFFFFFFFFFFFFFFFFFFFFFFFFFFFFFFFFFFFFFFFFFFFFFFFFFFFFFFFFFFFFFFFFFFFFFFFFFFFFFFFFFFFFFFFFFFFFFFFFFFFFFFFFFFFFFFFFFFFFFFFFFFFFFFFFFFFFFFFFFFFFFFFFFFFFFFFFFFFFFFFFFFEFBFEFC9872311F2B48A599FEFCFEF8F8FDF8D8A4B2C172321A2F0E0334110200242404032C828DF9FEF1FAFCFDFEFDFBFCFAE9BE783B04042BA3EDFEF9F8FFFFFEFBFAFEF9FDFFFFFFFFFFFFFFFDF9FEF8FCFBFBFBC68E6E4311050802000101000100002276A3A4FAF2FAF9FCF498342A1C0809060C2A285A7ACCF3F8FCFEFCFCFFFFFFFFFFFFFFFFFFFFFFFFFFFFFFFFFFFFFFFFFFFFFFFFFFFFFFFFFFFFFFFFFFFFFFFFFFFFFFFFFFFFFFFFFFFFFFFFFFFFFFFFFFFFFFFFFFFFFFFFFFFFFFFFFFFFFFFFFFFFFFFFFFFFFFFFFFFFFFFFFFFFFFFFFFFFFFFFFFFFFFFFFFFFFFFFFFFFFFFFFFFFFFFFFFFFFFFFFFFFFFFFFFFFFFFFFFFFFFFFFFFFFFFFFFFEFCFDF8C98C52200B1F60D0FDF5FEFDFAFCE8CF947C7E3D000200050000000207050100042C89C3FAF9CBECF9FDFEFAF7FAF4CE93601E022E65BFEDFAFDFDFFFFFEFCFDFEFCFEFFFFFFFFFFFFFFFCF7FEF8FEFEFCF4C9A39C6A160506030001053045182A226EE1D1FDFBFAFAFDFCD1847E3D010507000218587DCBF6FBFCFEFDFCFEFFFFFFFFFFFFFFFFFFFFFFFFFFFFFFFFFFFFFFFFFFFFFFFFFFFFFFFFFFFFFFFFFFFFFFFFFFFFFFFFFFFFFFFFFFFFFFFFFFFFFFFFFFFFFFFFFFFFFFFFFFFFFFFFFFFFFFFFFFFFFFFFFFFFFFFFFFFFFFFFFFFFFFFFFFFFFFFFFFFFFFFFFFFFFFFFFFFFFFFFFFFFFFFFFFFFFFFFFFFFFFFFFFFFFFFFFFFFFFFFFFFFFFFFFFFFFCF9FCD3B55E214075FBFAF4FEF7E5D8B59578441E1800020936100400020B090201154B95FDF2EC61D7F9FEFBF7F6F6E8AC4D1E04015E72B9E5F2FAFFFFFFFFFFFFFFFFFFFFFFFFFFFFFFFFFEFDFFFEFCFDFEF5E2E1C47931100A010500068FA538574A5BC2D6FBF9FDFEFCF9E4C3B14C0304010502108B8BC2F3FEFCFCFEFEFEFFFFFFFFFFFFFFFFFFFFFFFFFFFFFFFFFFFFFFFFFFFFFFFFFFFFFFFFFFFFFFFFFFFFFFFFFFFFFFFFFFFFFFFFFFFFFFFFFFFFFFFFFFFFFFFFFFFFFFFFFFFFFFFFFFFFFFFFFFFFFFFFFFFFFFFFFFFFFFFFFFFFFFFFFFFFFFFFFFFFFFFFFFFFFFFFFFFFFFFFFFFFFFFFFFFFFFFFFFFFFFFFFFFFFFFFFFFFFFFFFFFFFFFFFEFBFEFCFAFCF87D56CDDBF5F5FFFFF4D1987B3F1E120102020117872E0C020202061E33559DC3FDEDB92BCFD4EDEBDBE9F7E7B63E040B014386D0F6F3F4FFFFFFFFFFFFFFFFFFFFFFFFFFFFFFFFFEFEFFFEF9FBFDFEF7F9E6C0822202010D0500A176232829254B97FDFAFBFFFFFAD29A8B3E1020180C0C35BBC7E5FAFFFAFAFDFEFFFFFFFFFFFFFFFFFFFFFFFFFFFFFFFFFFFFFFFFFFFFFFFFFFFFFFFFFFFFFFFFFFFFFFFFFFFFFFFFFFFFFFFFFFFFFFFFFFFFFFFFFFFFFFFFFFFFFFFFFFFFFFFFFFFFFFFFFFFFFFFFFFFFFFFFFFFFFFFFFFFFFFFFFFFFFFFFFFFFFFFFFFFFFFFFFFFFFFFFFFFFFFFFFFFFFFFFFFFFFFFFFFFFFFFFFFFFFFFFFFFFFFFFFFFEF3FBF7F2F9F58D7CF1F6FDF8FEFEF4D393491404020003060003281A0402070A0455B69D86FEF9E8362EAC79C8D995A3EAD5792202210D77DBE9F4FCFEFFFFFFFFFFFFFFFFFFFFFFFFFFFFFFFFFDFCFFFEFAF9FDFEFEFEFBCD6613070300110338110200000544BBFCFDFEFEFEFDE4B761182D6144011684E2F7FEFFFFFCFBFCFDFEFFFFFFFFFFFFFFFFFFFFFFFFFFFFFFFFFFFFFFFFFFFFFFFFFFFFFFFFFFFFFFFFFFFFFFFFFFFFFFFFFFFFFFFFFFFFFFFFFFFFFFFFFFFFFFFFFFFFFFFFFFFFFFFFFFFFFFFFFFFFFFFFFFFFFFFFFFFFFFFFFFFFFFFFFFFFFFFFFFFFFFFFFFFFFFFFFFFFFFFFFFFFFFFFFFFFFFFFFFFFFFFFFFFFFFFFFFFFFFFFFFFFFFFFFEF4F2FBFDFDF5DCDEFEFDFCFDFAF6C34B2B0D0A080003060000042203150A1111072293E5CDFAEAD5060154B5DEE4B17F8D800C080A8C4B87E6FEF6F8FEFFFFFFFFFFFFFFFFFFFFFFFFFFFFFFFFFDF8FFFFFFFDFDFDFCFCFDE8992E070F00020A020008040005298CFCFEFDFAFAFEFCE09A350F2119070843C0FCFEFEFEFFFEFCFBFCFFFFFFFFFFFFFFFFFFFFFFFFFFFFFFFFFFFFFFFFFFFFFFFFFFFFFFFFFFFFFFFFFFFFFFFFFFFFFFFFFFFFFFFFFFFFFFFFFFFFFFFFFFFFFFFFFFFFFFFFFFFFFFFFFFFFFFFFFFFFFFFFFFFFFFFFFFFFFFFFFFFFFFFFFFFFFFFFFFFFFFFFFFFFFFFFFFFFFFFFFFFFFFFFFFFFFFFFFFFFFFFFFFFFFFFFFFFFFFFFFFFFFFFFFEFEEDF9FEF8F5F8FDFCF4F1F5F9B66F912602010203080B00000EAC2085318D510C62DBFEFCFAFCDE060B14BFD9869379190F0300022770E9F9FCFEF8F9FFFFFFFFFFFFFFFFFFFFFFFFFFFFFFFFFDFAFEFEFEFEFFFCF8F8FCBE7784280600130805010707090D4AC7FEFDFDFCFDFEF7D44D0B06000103178AF6FDFBFBFEFFFEFDFDFFFFFFFFFFFFFFFFFFFFFFFFFFFFFFFFFFFFFFFFFFFFFFFFFFFFFFFFFFFFFFFFFFFFFFFFFFFFFFFFFFFFFFFFFFFFFFFFFFFFFFFFFFFFFFFFFFFFFFFFFFFFFFFFFFFFFFFFFFFFFFFFFFFFFFFFFFFFFFFFFFFFFFFFFFFFFFFFFFFFFFFFFFFFFFFFFFFFFFFFFFFFFFFFFFFFFFFFFFFFFFFFFFFFFFFFFFFFFFFFFFFFFFFFFFFEEDFBEFEBF6FDFEFEFEF0B876DB25131F0E0206476D39010000091C9FEBDD711A728188FEFEFDE53101299C5DB3CA5D0A030307040E7CE9F3FBFBF2F7FEFFFFFFFFFFFFFFFFFFFFFFFFFFFFFFFFFEFEFEFCFCFEFFFEFEFBFCECDDBD27070000B72C1307745D64EAFBFDFDFEFDFBFEF9E7A036000106080328B5FCFCFCFEFFFEFDFEFFFFFFFFFFFFFFFFFFFFFFFFFFFFFFFFFFFFFFFFFFFFFFFFFFFFFFFFFFFFFFFFFFFFFFFFFFFFFFFFFFFFFFFFFFFFFFFFFFFFFFFFFFFFFFFFFFFFFFFFFFFFFFFFFFFFFFFFFFFFFFFFFFFFFFFFFFFFFFFFFFFFFFFFFFFFFFFFFFFFFFFFFFFFFFFFFFFFFFFFFFFFFFFFFFFFFFFFFFFFFFFFFFFFFFFFFFFFFFFFFFFFFFFFFFFEFCF8FBFCFDFEFFF7F2FBE48F4C90160B030000233212020B0E0B8BEEF1F3D04687D8D5F5FBD741AA040F4FB9D084734E030A09000A81F5F7FCFEFAFCFFFFFFFFFFFFFFFFFFFFFFFFFFFFFFFFFFFFFFFEFCFCFEFFFFFFF9FAFCF9F5B9160F052E0F03002687817AFDFAFBFFFEF5F6FCFCEAA440042771677AF0FDFEFEFCFBFDFEFFFFFFFFFFFFFFFFFFFFFFFFFFFFFFFFFFFFFFFFFFFFFFFFFFFFFFFFFFFFFFFFFFFFFFFFFFFFFFFFFFFFFFFFFFFFFFFFFFFFFFFFFFFFFFFFFFFFFFFFFFFFFFFFFFFFFFFFFFFFFFFFFFFFFFFFFFFFFFFFFFFFFFFFFFFFFFFFFFFFFFFFFFFFFFFFFFFFFFFFFFFFFFFFFFFFFFFFFFFFFFFFFFFFFFFFFFFFFFFFFFFFFFFFFFFFFFFF63D4FDFBFAFAF9F0BA76BA1E370604010503070903031F84C4EAFCF6FBF8C8D9F8F9FBCB5FA83F0D03115BB061241D07021A78BEE4FCFDF8FCFAFBFFFFFFFFFFFFFFFFFFFFFFFFFFFFFFFFFFFFFFFEFEFEFEFCFBFCFDF9F7FCD5550D0301050501010A85D7D2FDFCF6F7FEFDFBFDF8CE66140010291836BEF9FEFFFCFCFDFFFBFAFFFFFFFFFFFFFFFFFFFFFFFFFFFFFFFFFFFFFFFFFFFFFFFFFFFFFFFFFFFFFFFFFFFFFFFFFFFFFFFFFFFFFFFFFFFFFFFFFFFFFFFFFFFFFFFFFFFFFFFFFFFFFFFFFFFFFFFFFFFFFFFFFFFFFFFFFFFFFFFFFFFFFFFFFFFFFFFFFFFFFFFFFFFFFFFFFFFFFFFFFFFFFFFFFFFFFFFFFFFFFFFFFFFFFFFFFCFAFBFDFCF9FBEE62F42EC190C2F4F8FBF8D9874B09080005395D28020203060D3858DDF1FAFEFCF5F9FEFFBF3920440909010113361D05060304123AB6E4FCFFF7F8FBFDFFFFFFFFFFFFFFFFFFFFFFFFFFFFFFFFFFFFFFFFFFFFFFFEFDFEFEFBF8FCEDBA9B160200060109072B9BF5F7FEFBF8FFFFF1FAC78A753D00060B1687E3FAFEFCFBFDF7F6FBFDFFFFFFFFFFFFFFFFFFFFFFFFFFFFFFFFFFFFFFFFFFFFFFFFFFFFFFFFFFFFFFFFFFFFFFFFFFFFFFFFFFFFFFFFFFFFFFFFFFFFFFFFFFFFFFFFFFFFFFFFFFFFFFFFFFFFFFFFFFFFFFFFFFFFFFFFFFFFFFFFFFFFFFFFFFFFFFFFFFFFFFFFFFFFFFFFFFFFFFFFFFFFFFFFFFFFFFFFFFFFFFFFFFFFFFFFFDFCFDFDFCFCFBE11C5F0E512BA7A3AEFBF1BC3C0B032979181E3A1A03050403031886E3F6FDFFFCFCFEFDF7D37309268707000603040402040806196083BCEEFDFCFCFDFAFBFFFFFFFFFFFFFFFFFFFFFFFFFFFFFFFFFFFFFFFFFFFFFFFFFFF9F8FDF7EDE27A0A0A0B04010465563199F9FAFEFAF8F3F6FAE5CAA54A030605144063C5F3FCFAFDFBFAF7F7FFFFFFFFFFFFFFFFFFFFFFFFFFFFFFFFFFFFFFFFFFFFFFFFFFFFFFFFFFFFFFFFFFFFFFFFFFFFFFFFFFFFFFFFFFFFFFFFFFFFFFFFFFFFFFFFFFFFFFFFFFFFFFFFFFFFFFFFFFFFFFFFFFFFFFFFFFFFFFFFFFFFFFFFFFFFFFFFFFFFFFFFFFFFFFFFFFFFFFFFFFFFFFFFFFFFFFFFFFFFFFFFFFFFFFFFFFFFFCFAF6F1E87890150D0E044A8AC7F3C14B070809254C1204090A0503080F091CBBF7FDFEFDFDFEFFFCF8F3AD0538DF07000401090A030204165FC3B9D6F2FBFBFCFEFCFCFFFFFFFFFFFFFFFFFFFFFFFFFFFFFFFFFFFFFFFFFFFFFFFFFFFBFBFFFAF4EEA0170B070102024135042E769DFBFCF3F7FDF4F8FAC04C020418583D61D9F8F5F7FDFBFAF8FAFFFFFFFFFFFFFFFFFFFFFFFFFFFFFFFFFFFFFFFFFFFFFFFFFFFFFFFFFFFFFFFFFFFFFFFFFFFFFFFFFFFFFFFFFFFFFFFFFFFFFFFFFFFFFFFFFFFFFFFFFFFFFFFFFFFFFFFFFFFFFFFFFFFFFFFFFFFFFFFFFFFFFFFFFFFFFFFFFFFFFFFFFFFFFFFFFFFFFFFFFFFFFFFFFFFFFFFFFFFFFFFFFFFFFFFFFFFCFAFAF7F2E88A7F0D0A035E85C5F5BF500801070C09000300060704070D0C1B6DDFFFFFFDFAFBFEFFFCFBFCD5799CDA0B17641D07204632164F89BEF5FAFDFDFBFBFCFEFFFFFFFFFFFFFFFFFFFFFFFFFFFFFFFFFFFFFFFFFFFFFFFFFFFFFBFBFBFAFDE98F741702070503030100041251FCFDECFDF8F5F9FDDE872903184A344CCBF4F6FAFEFAFAFDFFFFFFFFFFFFFFFFFFFFFFFFFFFFFFFFFFFFFFFFFFFFFFFFFFFFFFFFFFFFFFFFFFFFFFFFFFFFFFFFFFFFFFFFFFFFFFFFFFFFFFFFFFFFFFFFFFFFFFFFFFFFFFFFFFFFFFFFFFFFFFFFFFFFFFFFFFFFFFFFFFFFFFFFFFFFFFFFFFFFFFFFFFFFFFFFFFFFFFFFFFFFFFFFFFFFFFFFFFFFFFFFFFFFFFFFFFFFF9F8FAF9FCF9EC7D6E135483D3F8CA7F3A040706060305010509040123586E86D6FBFCFDFDFAFBFEFFFDFBFBF8D8A2D50B196D22001A4B38185590BDFAF9FBFEFEFDFCFDFFFFFFFFFFFFFFFFFFFFFFFFFFFFFFFFFFFFFFFFFFFFFFFFFFFFFBFBF9F7FBE291881501070803020105090D50F9FDF1F7F6FCFAFCDB8A31040607060FAEF1FDFCFDFCFDFDFCFFFFFFFFFFFFFFFFFFFFFFFFFFFFFFFFFFFFFFFFFFFFFFFFFFFFFFFFFFFFFFFFFFFFFFFFFFFFFFFFFFFFFFFFFFFFFFFFFFFFFFFFFFFFFFFFFFFFFFFFFFFFFFFFFFFFFFFFFFFFFFFFFFFFFFFFFFFFFFFFFFFFFFFFFFFFFFFFFFFFFFFFFFFFFFFFFFFFFFFFFFFFFFFFFFFFFFFFFFFFFFFFFFFFFFFFFFF9F8F9F8FEFDF9F8F779D387C4FBAC6F5A060A0B1C4B0E243B1405062A86E0FEFDF9FAFCFEFDFDFEFEFEFBF7F3D66D910A04010304011833324DBCFAF8F7F9FDFDFDFCFDFFFFFFFFFFFFFFFFFFFFFFFFFFFFFFFFFFFFFFFFFFFFFFFFFFFFFDFBF9FAF9F3EF810C010006090701070C237AACFEF8F4FDFCF8FBD37B270504010B47CCF9FAF7FBFCFDFCFBFFFFFFFFFFFFFFFFFFFFFFFFFFFFFFFFFFFFFFFFFFFFFFFFFFFFFFFFFFFFFFFFFFFFFFFFFFFFFFFFFFFFFFFFFFFFFFFFFFFFFFFFFFFFFFFFFFFFFFFFFFFFFFFFFFFFFFFFFFFFFFFFFFFFFFFFFFFFFFFFFFFFFFFFFFFFFFFFFFFFFFFFFFFFFFFFFFFFFFFFFFFFFFFFFFFFFFFFFFFFFFFFFFFFFFFFFCFBFCFBFBFEFCF5FFFFFDFCFEFFF0BD5D0F08040537890D4265241F311B38B1FBFDFCFCFEFFFFFFFEFEFEFEDFA8B08C060E07030508093B8CABAFF5FCFBFCFEFDFCFBFCFCFDFFFFFFFFFFFFFFFFFFFFFFFFFFFFFFFFFFFFFFFFFFFFFFFFFFFCF7F9FDFAFDF4AF1A02040133220623313C936EF9F4FDFFF7FBFDE8B5641E0A031574A4E0F3F7FCFBFCFCFCFFFFFFFFFFFFFFFFFFFFFFFFFFFFFFFFFFFFFFFFFFFFFFFFFFFFFFFFFFFFFFFFFFFFFFFFFFFFFFFFFFFFFFFFFFFFFFFFFFFFFFFFFFFFFFFFFFFFFFFFFFFFFFFFFFFFFFFFFFFFFFFFFFFFFFFFFFFFFFFFFFFFFFFFFFFFFFFFFFFFFFFFFFFFFFFFFFFFFFFFFFFFFFFFFFFFFFFFFFFFFFFFFFFFFFFFF7F5F7F9FFFEFCFEFDFEFAF8FAFDC6702801030A000206010401043B623D3293E8FDFCFEFEFDFEFEFEFEFFFCB73E3B011007060F3A3C4761B0FDFCFCF8F5F9FDFCFDFEFEFBFBFFFFFFFFFFFFFFFFFFFFFFFFFFFFFFFFFFFFFFFFFFFFFFFFFFFCFCFDFBFBFCE85119010B0365482D7D97517BFCF3FDFEFEF6FBFDC5998F590501031232AEF2FEFDFDFEFBF9FFFFFFFFFFFFFFFFFFFFFFFFFFFFFFFFFFFFFFFFFFFFFFFFFFFFFFFFFFFFFFFFFFFFFFFFFFFFFFFFFFFFFFFFFFFFFFFFFFFFFFFFFFFFFFFFFFFFFFFFFFFFFFFFFFFFFFFFFFFFFFFFFFFFFFFFFFFFFFFFFFFFFFFFFFFFFFFFFFFFFFFFFFFFFFFFFFFFFFFFFEFEFDFEFEFEFEFEFFFFFFFFFFFFFFFFFFFEF9F8FEFEFBFCFAFDF9FBF6F56A627300000004080006010302040F2B3955C2FCF6FEF9FAFCFBFDFEFEE6C8CAA0010005395DCDFAFAFBFFF9F3FAFFFFFFFFF9FAFDFFFFFFFFFFFFFFFFFFFFFFFFFFFFFFFFFFFFFFFFFFFFFFFFFFFFFFFFFEFCF8F6FFFCEED04F0A0100080E7C9883F4F9FDFEF7FBFEFBFDF9FBFEC9550302080585C4EAF7FCFCFDFEFFFFFFFFFFFFFFFFFFFFFFFFFFFFFFFFFFFFFFFFFFFFFFFFFFFFFFFFFFFFFFFFFFFFFFFFFFFFFFFFFFFFFFFFFFFFFFFFFFFFFFFFFFFFFFFFFFFFFFFFFFFFFFFFFFFFFFFFFFFFFFFFFFFFFFFFFFFFFFFFFFFFFFFFFFFFFFFFFFFFFFFFFFFFFFFFFFFFFFFFFEF9FCF7FAFCFAFDFCFFFFFFFFFFFFFFFFFFFAFAFDFCFEFCF6FDF8FCF2F4DB5800030403040F01000B02030A054FC5FBF9FEFBFEFDECFCFDF0F5FBF8920C150E060116D6F4FCFFFFFFFFFFFFFFFFFFFFFFFFFFFFFFFFFFFFFFFFFFFFFFFFFFFFFFFFFFFFFFFFFFFFFFFFFFFFFFFFFFFFFCF7F7FEFEFEE425CF2D070401076DCBD4E5FAF6FEFDFEFDF7FBFDFDFCCF6910030400162DA9E7FAFEFDFCFEFFFFFFFFFFFFFFFFFFFFFFFFFFFFFFFFFFFFFFFFFFFFFFFFFFFFFFFFFFFFFFFFFFFFFFFFFFFFFFFFFFFFFFFFFFFFFFFFFFFFFFFFFFFFFFFFFFFFFFFFFFFFFFFFFFFFFFFFFFFFFFFFFFFFFFFFFFFFFFFFFFFFFFFFFFFFFFFFFFFFFFFFFFFFFFFFFFFFFFFAFCFBFBFDFCFAFDFCFFFFFFFFFFFFFFFFFFFCFCFDFCFCFBF6FAFDEAF8E55B0806070005120211030002080A0B59CAFCFBFEFBF8FCFDFEEDF7DFE1F89354B50003082B38DCFEFFFFFFFFFFFFFFFFFFFFFFFFFFFFFFFFFFFFFFFFFFFFFFFFFFFFFFFFFFFFFFFFFFFFFFFFFFFFFFFFFFFFFEFEF6F1FBF4FCF54F2B0700090A080B185DE1FDFCF6FEF3FBFAFCF6F6F6C152040A059EE9FBFEFAFCFDFAFAFFFFFFFFFFFFFFFFFFFFFFFFFFFFFFFFFFFFFFFFFFFFFFFFFFFFFFFFFFFFFFFFFFFFFFFFFFFFFFFFFFFFFFFFFFFFFFFFFFFFFFFFFFFFFFFFFFFFFFFFFFFFFFFFFFFFFFFFFFFFFFFFFFFFFFFFFFFFFFFFFFFFFFFFFFFFFFFFFFFFFFFFFFFFFFFFFFFFFFFCFDFAFCFEFCFBFCF9FFFFFFFFFFFFFFFFFDFDFDFEFDFCFCFDF8F8FCF2699457040504131E0C161F020A4C7991BFEEFCF7F8FEFEFDFAFEFDFB8A93CC7C0E120D0D47CFE6F6FDFFFFFFFFFFFFFFFFFFFFFFFFFFFFFFFFFFFFFFFFFFFFFFFFFFFFFFFFFFFFFFFFFFFFFFFFFFFFFFFFFFFFFFFAF8FAF6FCF7FBB8CB3B100400000A569967F6FCEAFDFBFDF7FAF6F6C78CA35704000C91DCF7F3F8FEFBFAFFFFFFFFFFFFFFFFFFFFFFFFFFFFFFFFFFFFFFFFFFFFFFFFFFFFFFFFFFFFFFFFFFFFFFFFFFFFFFFFFFFFFFFFFFFFFFFFFFFFFFFFFFFFFFFFFFFFFFFFFFFFFFFFFFFFFFFFFFFFFFFFFFFFFFFFFFFFFFFFFFFFFFFFFFFFFFFFFFFFFFFFFFFFFFFFFFFFFFFEFEFDFCFDFCFBFDFCFFFFFFFFFFFFFFFFFDFDFEFDFDFBFDFEF8F8F8FBC353100308133F820E68870D011871BBE8FCFEFCFBFEF9F3FDFEFFCFCD71221401020719CE60E5FCFEFFFFFFFFFFFFFFFFFFFFFFFFFFFFFFFFFFFFFFFFFFFFFFFFFFFFFFFFFFFFFFFFFFFFFFFFFFFFFFFFFFFFF9F9FEFEFBFBFBED762F0A02000506031056CFF6FBFAF7F0FBFDFCF8F3EEDFB96D15001128B1F1FDFDFFFBF9FFFFFFFFFFFFFFFFFFFFFFFFFFFFFFFFFFFFFFFFFFFFFFFFFFFFFFFFFFFFFFFFFFFFFFFFFFFFFFFFFFFFFFFFFFFFFFFFFFFFFFFFFFFFFFFFFFFFFFFFFFFFFFFFFFFFFFFFFFFFFFFFFFFFFFFFFFFFFFFFFFFFFFFFFFFFFFFFFFFFFFFFFFFFFFFFFFFFFFFFF3FAF4F7FDFCFEFEFFFFFFFFFFFFFFFFFEFEFEFDFAF8F9FEF8FCF8C8541503000204143742B1263309299BE5FBFDF9F4F3FEFDFBFCD2C37DA45A070700000F5AF2D4F7F5FBFFFFFFFFFFFFFFFFFFFFFFFFFFFFFFFFFFFFFFFFFFFFFFFFFFFFFFFFFFFFFFFFFFFFFFFFFFFFFFFFFFFFFBFCFDFDFDFBFCE7AE4C150300010101143463FDFDFAF9FAF7FBFDFDF6F6DF976C3903046CDCEBF7FEF4F2FDFFFFFFFFFFFFFFFFFFFFFFFFFFFFFFFFFFFFFFFFFFFFFFFFFFFFFFFFFFFFFFFFFFFFFFFFFFFFFFFFFFFFFFFFFFFFFFFFFFFFFFFFFFFFFFFFFFFFFFFFFFFFFFFFFFFFFFFFFFFFFFFFFFFFFFFFFFFFFFFFFFFFFFFFFFFFFFFFFFFFFFFFFFFFFFFFFFFFFFFFFAFCF6F9FDFCFDFFFFFFFFFFFFFFFFFFFFFDFDFCFAFCFBF8FBC3F2811F042629090001104ECE619C42448DD2F1F8FCFBFAFDF8FEFE66407C5D1B02020F2F50D2F6FDF9F9FDFFFFFFFFFFFFFFFFFFFFFFFFFFFFFFFFFFFFFFFFFFFFFFFFFFFFFFFFFFFFFFFFFFFFFFFFFFFFFFFFFFFFFCFCFAF8FDFBFCF1759522040502040D436B42FCF4FDFDFEFCFBFBF5F4F8E1A96D4518022D80DBFBFEFAF6FAFFFFFFFFFFFFFFFFFFFFFFFFFFFFFFFFFFFFFFFFFFFFFFFFFFFFFFFFFFFFFFFFFFFFFFFFFFFFFFFFFFFFFFFFFFFFFFFFFFFFFFFFFFFFFFFFFFFFFFFFFFFFFFFFFFFFFFFFFFFFFFFFFFFFFFFFFFFFFFFFFFFFFFFFFFFFFFFFFFFFFFFFFFFFFFFFFFFFFFF6FCF3F7FCFBFAFCFBFFFFFFFFFFFFFFFFFFFBFBFAF8FFF9E5FD7DC995270D50470B15242077E5D2E4BFB5C8E6F7FBFEFFFFFEFAFBF6450C371C00031B64B7C3F9FBFAF8FCFEFFFFFFFFFFFFFFFFFFFFFFFFFFFFFFFFFFFFFFFFFFFFFFFFFFFFFFFFFFFFFFFFFFFFFFFFFFFFFFFFFFFFFDF8F6F8F7FCFAE9BE82170108060307213339E1F8F2FAF8F6FBF6F5FAFCDCA098932C020246C7F0F6FDFBFBFFFFFFFFFFFFFFFFFFFFFFFFFFFFFFFFFFFFFFFFFFFFFFFFFFFFFFFFFFFFFFFFFFFFFFFFFFFFFFFFFFFFFFFFFFFFFFFFFFFFFFFFFFFFFFFFFFFFFFFFFFFFFFFFFFFFFFFFFFFFFFFFFFFFFFFFFFFFFFFFFFFFFFFFFFFFFFFFFFFFFFFFFFFFFFFFFFFFFFD2F0D3E9F8F4F8FCFAFFFFFFFFFFFFFFFFFFFAFAF8F3FEF3C4F7915D4F170B30250B385F3CC4F8FEFAFBF9F8FCFFFDFCFFFEFCFBD9C8370E062234375AB0F9FBF6FEFCF5FCFFFFFFFFFFFFFFFFFFFFFFFFFFFFFFFFFFFFFFFFFFFFFFFFFFFFFFFFFFFFFFFFFFFFFFFFFFFFFFFFFFFFFFFFF9F8FCF7F9FCF1F8BA29060505020004164076F5F6FBEEDAEEF7F6F8FCECCED09D25030060CBF0F9FDFCFCFFFFFFFFFFFFFFFFFFFFFFFFFFFFFFFFFFFFFFFFFFFFFFFFFFFFFFFFFFFFFFFFFFFFFFFFFFFFFFFFFFFFFFFFFFFFFFFFFFFFFFFFFFFFFFFFFFFFFFFFFFFFFFFFFFFFFFFFFFFFFFFFFFFFFFFFFFFFFFFFFFFFFFFFFFFFFFFFFFFFFFFCFBFFFEFBFDFDFE75AF9BB6E5F7F7F7FDFFFFFFFFFFFFFFFFFFFAFAFAF4F7EB95C4690D0002030402082A5474F6FEFEFEF7FCFDFCFDFDEFD4CDF4D3685D1709236499926E91E9D4FAFDFAF7FDFFFFFFFFFFFFFFFFFFFFFFFFFFFFFFFFFFFFFFFFFFFFFFFFFFFFFFFFFFFFFFFFFFFFFFFFFFFFFFFFFFFFFFFFFDFDFEFCFAFFFEF9FD410B020201000832622BCBFACFAF84D3FBF9F4F7FDFDEEAE37271543B7F2FEFBFAFCFDFDFFFEF8FEFBFEFFFFFFFFFFFFFFFFFFFEFEFCFEFDFEFFECD8D0C9D2F5FFFFFEFEFEFFFFFFFFFFFFFFFFFFFFFFFFFFFFFFFFFFFFFFFFFFFFFFFFFFFFFFFFFFFFFFFFFFFFFFFFFFFFFFFFFFFFFFFFFFFFFFFFFFFFFFFFFFFFFFFFFAF6FAFCF7F9D4C21E466F6EBDFDF9F5FBFFFFFFFFFFFFFFFFFCFBFAFCFBF5E77E560B020522130604020833ADFFFEFEF8F5F8F8F7FBFFE1A18CBB74090605014E7EAEAC5449BAA2FEFCF9FDFBFDFFFFFFFFFFFFFFFFFFFFFFFFFFFFFFFFFFFFFFFFFFFFFFFFFFFFFFFFFFFFFFFFFFFFFFFFFFFFFFFFFFFFFFFFFFFFFFFFFFFFFFFA3E2B1201020307254A269AC1654330A8EFFCF8F6FBF7F7E4505D2E097FDFF9F8FAFEFDFCFBFDF2FEF9FDFFFFFFFFFFFFFFFFFFFEFDFBFDFCFDE8A36B5D6093F2F6FEFCFDFEFFFFFFFFFFFFFFFFFFFFFFFFFFFFFFFFFFFFFFFFFFFFFFFFFFFFFFFFFFFFFFFFFFFFFFFFFFFFFFFFFFFFFFFFFFFFFFFFFFFFFFFFFFFFFFFCFAF5F3FCF8EF686D1206463189E9F7FBF9FFFFFFFFFFFFFFFFF9FAFAEDD8CCC4650002060A572E0C0814263297FBFBFBF4F8EDE4E4E8EBD097785B1706051A286684706D4C276894FEFEFEFCF8FCFFFFFFFFFFFFFFFFFFFFFFFFFFFFFFFFFFFFFFFFFFFFFFFFFFFFFFFFFFFFFFFFFFFFFFFFFFFFFFFFFFFFFFFFFFFFFFFFFFFFFFF851662D0305070101071E83510E02247EC6EFFCFCF8F3FCE6475420083DB7F1FAFDFAFAFAF4FCF7FEFBFCFEFFFFFFFFFFFFFFFEFAFDFEFEF9E7B5540B01207FDDEFFAF9FCFEFFFFFFFFFFFFFFFFFFFFFFFFFFFFFFFFFFFFFFFFFFFFFFFFFFFFFFFFFFFFFFFFFFFFFFFFFFFFFFFFFFFFFFFFFFFFFFFFFFFFFFFFFFFFFFF9F8F8F4FBFAE403761A03010B62C5F1FAFBFFFFFFFFFFFFFFFFFDFDFCC974655F13000606044C2702083E550A50F4F6F6F9F6D6BDB4B1A98D6B5B23031E0D639595D36D3540022391FCFBFEF7FDFFFFFFFFFFFFFFFFFFFFFFFFFFFFFFFFFFFFFFFFFFFFFFFFFFFFFFFFFFFFFFFFFFFFFFFFFFFFFFFFFFFFFFFFFFFFFFFFFFFFFFFFFDAB813005050905050D1751080302266B97D4FCFCFCFCF9E566170601067BCBEDFBF4F3F9F3FFFEFDFCF9FCFFFFFFFFFFFFFFFDF5FBF7FBF2AA8450150422658EFBF6F6FAFDFFFFFFFFFFFFFFFFFFFFFFFFFFFFFFFFFFFFFFFFFFFFFFFFFFFFFFFFFFFFFFFFFFFFFFFFFFFFFFFFFFFFFFFFFFFFFFFFFFFFFFFFFFFFFFFFF9F7FEFAF1E701690E0402125FC8FEF8F2FFFFFFFFFFFFFFFFFDFDF69F3B5B080004000302020603103D400A4CDCEBFAF3FCCB846F715119040D000D8A25C9F7F2ED781D201F40B8FEF9FDF6FCFEFFFFFFFFFFFFFFFFFFFFFFFFFFFFFFFFFFFFFFFFFFFFFFFFFFFFFFFFFFFFFFFFFFFFFFFFFFFFFFFFFFFFFFFFFFFFFFFFFFFFFFFBFEA4380906040205010702001516205996D7F6F6FBFDFDF4AE2513040E3558A3E8FAFDFDFDFEFEF8FDFFFFFFFFFFFFFFFFFFFFFEFBEEFAF27546320D01021C75F6FCF7FDFFFFFFFFFFFFFFFFFFFFFFFFFFFFFFFFFFFFFFFFFFFFFFFFFFFFFFFFFFFFFFFFFFFFFFFFFFFFFFFFFFFFFFFFFFFFFFFFFFFFFFFFFFFFFFFFFEF5F1FBFEFBD7040F0500071257C2FCF9FDFFFFFFFFFFFFFFFFF7F9F6A342570105141012030A150B043A550B2E6890F2F4F8D75F27372A0600030B21D5A8E5EFF9FEC897959CB3F1FDFEFBF7FBFDFFFFFFFFFFFFFFFFFFFFFFFFFFFFFFFFFFFFFFFFFFFFFFFFFFFFFFFFFFFFFFFFFFFFFFFFFFFFFFFFFFFFFFFFFFFFFFFFFFFFFFFEFCFD540A0800020300020901366D2146BDF8FBFBFCF7F6E98F22100A0101022591F0F6FAFCFCFEF9FEFEFEFFFFFFFFFFFFFFFEFEFEFCF9BB430D0C0E080645EEF9F5FCFDFEFFFFFFFFFFFFFFFFFFFFFFFFFFFFFFFFFFFFFFFFFFFFFFFFFFFFFFFFFFFFFFFFFFFFFFFFFFFFFFFFFFFFFFFFFFFFFFFFFFFFFFFFFFFFFFFAFAF9F6FBDD6500000B060666A1D0FBF5F3FFFFFFFFFFFFFFFFFEFDC1521004030701921B010067480E3C5645780746E29FBE8C712F020001031772A2756370C5F4FEF3F5FCFBFEFBF4FEFEFEFBFDFFFFFFFFFFFFFFFFFFFFFFFFFFFFFFFFFFFFFFFFFFFFFFFFFFFFFFFFFFFFFFFFFFFFFFFFFFFFFFFFFFFFFFFFFFFFFFFFFFFFFFFEECF05A0500040500030F00092847100F6FD2F8FAFEFEFEF2B21B813A0A090A517E9DF9F1FEFCFEFEFFFAFDFFFFFFFFFFFFFFFEFCFDFEFECAB1540B00000231E5F9FEF4F2FCFFFFFFFFFFFFFFFFFFFFFFFFFFFFFFFFFFFFFFFFFFFFFFFFFFFFFFFFFFFFFFFFFFFFFFFFFFFFFFFFFFFFFFFFFFFFFFFFFFFFFFFFFFFFFFFEFAF8FAD35A051102008A4351BFFBF8F5FEFFFFFFFFFFFFFFFFFEF7DE840F03090112611B0D023E2A000206204D0815521E282C3C1A00000308279E9E0C941D4FC9E8F2FCFCFEFDF7F0FEFEFEFEFEFFFFFFFFFFFFFFFFFFFFFFFFFFFFFFFFFFFFFFFFFFFFFFFFFFFFFFFFFFFFFFFFFFFFFFFFFFFFFFFFFFFFFFFFFFFFFFFFFFFFFFF7FC8D81390B0C000002020D02010B020940B7FEFAF6FBFEFAF9CA732E1D09003B4160F5E5FEFAECFDFFFEFEFFFFFFFFFFFFFFFEFAFDFEEDA68A3E0D09000547DE62E1FBFCFDFFFFFFFFFFFFFFFFFFFFFFFFFFFFFFFFFFFFFFFFFFFFFFFFFFFFFFFFFFFFFFFFFFFFFFFFFFFFFFFFFFFFFFFFFFFFFFFFFFFFFFFFFFFFFFF9F9FCCE58090204000A532A7ADDFEFEFDFFFFFFFFFFFFFFFFFFF5F6C0520F06010006130401040E080205060B08040608030805060605051B5CADD53606481E0DA1EFFDFCFCF9F5FEFDFEF9F8FDFFFFFFFFFFFFFFFFFFFFFFFFFFFFFFFFFFFFFFFFFFFFFFFFFFFFFFFFFFFFFFFFFFFFFFFFFFFFFFFFFFFFFFFFFFFFFFFFFFFFFFFFFEFDB3D3C75301010605000500062766230D6ACEF6FBF9FEFDF387270B040607090D2B5D7FF8FDF5FDFCFCFEFFFFFFFFFFFFFFFEFDFEF7FBCE691902030100176CABF0F8FDFEFFFFFFFFFFFFFFFFFFFFFFFFFFFFFFFFFFFFFFFFFFFFFFFFFFFFFFFFFFFFFFFFFFFFFFFFFFFFFFFFFFFFFFFFFFFFFFFFFEF4FDFFF6FAFDF9F9CC5B3287090300050F0A52C6FBFAFFFFFFFFFFFFFFFFFAFDF8FAE07B367319000004000308030606010103040000000004010004042B8DDEE8DE30090B0D62CBF2FAFEFEF5EBFFFDF2FDFCFDFEFFFFFFFFFFFFFFFFFFFFFFFFFFFFFFFFFFFFFFFFFFFFFFFFFFFFFFFFFFFFFFFFFFFFFFFFFFFFFFFFFFFFFFFFFFFFFFFFFFFFFFF7FDF1E99D86580105060530A16A48AB4E67AEDFFDFDFBF9FEF9C43A0704020200466838B1FCFEF5FEF6FCFEFFFFFFFFFFFFFFFFFEFEFDFDB43A060305060006265EF8FCFEFFFFF9FAFAFAFFFCF7FFFFFFFFFFFFFFFFFFFFFFFFFFFFFFFFFFFFFFFFFFFFFFFFFFFFFFFFFFFFFFFFFFFFFFFFFFFFFCF7FEFBF4FCFBF8FBB8B1A2191A5D030C1006010F79DBFAFCFCF6F9FFFDEDF1FCFDFDFCF5C35B214C1300048A1808040001010406010500000D0B0202002473877FC2F8E9A6890E58C8F3FBFBFCFBF7F8FBF7FBFEFEFDFEFFFFFFFFFFFFFFFFFFFFFFFFFFFFFFFFFFFFFFFFFFFFFFFFFFFFFFFFFFFFFFFFFFFFFFFFFFFFFFFFFFFFFFFFFFFFFFFFFFFFFFFCF7F9F8BE797A6303040F1864747F739FD6EDF6F9F5FEF8EDF995230D0401000030481F7EF8F0FCFEFCFEFFFFFFFFFFFFFFFFFEFEFEF9FDD5AD5B140105337B0A20FC71E7FF9BC5FBFDF9F8FEFCFFFFFFFFFFFFFFFFFFFFFFFFFFFFFFFFFFFFFFFFFFFFFFFFFFFFFFFFFFFFFFFFFFFFFFFFFFFFFDF8EDF6F8FAFDF1F7914D46080A160C860F030116A0F3FBFCFCF7F5FEFBF6F8F1FBFEFEF9DA7E181105000461136E2906050333512562116708050703375B88AAB5D4F3F8EBE19DCDF4F9FCF9FAFBFBFCF8FAF9FCFAFDFEFFFFFFFFFFFFFFFFFFFFFFFFFFFFFFFFFFFFFFFFFFFFFFFFFFFFFFFFFFFFFFFFFFFFFFFFFFFFFFFFFFFFFFFFFFFFFFFFFFFFFFF6F9FCF7EBBA804101040D041953B9457DD9FBFDF7F4FDE69A7C993A0C01020708101531A9FBF5FEF4FDFCFDFFFFFFFFFFFFFFFDF4FEF7FCF1AE51140103214D0C098718697795DAF4F3FAF6F1FFFFFFFFFFFFFFFFFFFFFFFFFFFFFFFFFFFFFFFFFFFFFFFFFFFFFFFFFFFFFFFFFFFFFFFFFFFFFFFCF1F4F9FBFCFDB9904D150E0702137B7B7C10061AADFBFBFAFDFCF7F3FBF1F5F5F6FCFFCF8F4E0C02090A7E210D531E0D5282A8BE9BC991D3680907034A99BAD2EDF5F8FEFCFBEDF1FAFBF9F9FBFEFDDC9AFEFBF5F8FDFFFFFFFFFFFFFFFFFFFFFFFFFFFFFFFFFFFFFFFFFFFFFFFFFFFFFFFFFFFFFFFFFFFFFFFFFFFFFFFFFFFFFFFFFFFFFFFFFFFFFFFFF9FCF6F6FCEDB2512A610402021E6E1E3389C5CBD4F6FDEEB03D5A2E064A7A410B032D768FFEFDFEF8F3FAFEFFFFFFFFFFFFFFFEF7F5FAFDBB8C531A01020B1307021E09151559CBF7F6FDFBF4FAFFFFFFFFFFFFFFFFFFFFFFFFFFFFFFFFFFFFFFFFFFFFFFFFFFFFFFFFFFFFFFFFFFFFFFFFFFFFFBF2FEFAFCC0914B24110605041C7DD88E710E0819AEF9F8FAFAFDFCF5FEF4F5F9F8F8F49E5F4E19020916790B07170B094B82B8F4E8CB84886216052082D8F0F4F9F7FCFBFCFEF9F4F8F7F9FCFECF8E704AFAF5F7FCFEFFFFFFFFFFFFFFFFFFFFFFFFFFFFFFFFFFFFFFFFFFFFFFFFFFFFFFFFFFFFFFFFFFFFFFFFFFFFFFFFFFFFFFFFFFFFFFFFFFFFFFFFFFF7F5F6FBFBEBAE6E7D5F02020415050A385770A0C7C6C48424180D59737D3D07002B5C54F6F5FAF9F5F4FBFFFFFFFFFFFFFFFEFCEFF9FDBCB77B230001041E660A000C020015A8FBFEFCFAFDFBFFFFFFFFFFFFFFFFFFFFFFFFFFFFFFFFFFFFFFFFFFFFFFFFFFFFFFFFFFFFFFFFFFFFFFFFFFFFFDFBF7B5A04F1B06010004071F7EE6F2E17310040E76D8FAFBFBFAF7FEFBFEFCFAFDD4A65E464C1C00020E130C010C0805124C8FBED48964100E0F235EA0E2F9F9FAF7FCFAFAFAFBF5F7F7F9FDFA98171536F9F5FEFCF4FBFFFFFFFFFFFFFFFFFFFFFFFFFFFFFFFFFFFFFFFFFFFFFFFFFFFFFFFFFFFFFFFFFFFFFFFFFFFFFFFFFFFFFFFFFFFFFFFFFFFFFFFCF6FCF8FAFCFCEEB8756D0303030007030D0C25546157511B09094A837E73370800091429A2F4FAEFFCF3FAFFFFFFFFFFFFFFFEFBF6FAF7ECF1BB551002041F7A0B05070003089FF8FBFBF8FEFEFFFFFFFFFFFFFFFFFFFFFFFFFFFFFFFFFFFFFFFFFFFFFFFFFFFFFFFFFFFFFFFFFFFFFFFFFFFFFEF4BA431F070506022A30257BE6F9FDF29B17020125A3F8FDFFFBF2F9F7D1D2F6DB7932150B07030003000158180601050A4F696CC77A780100054AA6A5D6F9F8FAF8FEFEFDF8F9EDFAF9F9FCCF5E08205BB2FDFBFBEFFAFFFFFFFFFFFFFFFFFFFFFFFFFFFFFFFFFFFFFFFFFFFFFFFFFFFFFFFFFFFFFFFFFFFFFFFFFFFFFFFFFFFFFFFFFFFFFFFFFFFFFFFBFBFDF8F9FAFAFCF0A80A060C02000A060304070D100901000404AC8C8696480B0401031771ACFEF1FCF9FDFFFFFFFFFFFFFFFEFDFCF8F4F9FDD06C1A0706040F0008000305056ED5FAFDFAF8FDFFFFFFFFFFFFFFFFFFFFFFFFFFFFFFFFFFFFFFFFFFFFFFFFFFFFFFFFFFFFFFFFFFFFFFFFFFFFF5BF48090C02010502557981E2FBF3F8F0540B04093A86D4FEFDFEFCF4F8A686AC741E01000204020F3A0C03942C050301000D1345AD771004114B7DDDF3F8FEFCF5F2FDFCFFFEFDF7FBF2F7FCAC503249BF65FBF8F9FDFEFFFFFFFFFFFFFFFFFFFFFFFFFFFFFFFFFFFFFFFFFFFFFFFFFFFFFFFFFFFFFFFFFFFFFFFFFFFFFFFFFFFFFFFFFFFFFFFFFFFFFFFCF9FDF3F7FCFBFAE39F3B0316010102142F1A04050004060001027EAF4A2B1207070001259E5BFCFCFEFBFDFFFFFFFFFFFFFFFEFDF9F2FCFDFDB2470F07010203010003070101106FBCDAF8F4FFFFFFFFFFFFFFFFFFFFFFFFFFFFFFFFFFFFFFFFFFFFFFFFFFFFFFFFFFFFFFFFFEFDFDFDFDFEF6CA540A0000000422152086EBF8FDFBF1FCB6430F177F81AAFBF9FCFEFCF8DD882D0F050402010709278C220012020008061F1B01070F08010A2FBCF2F9F9FCFCFCFEFFF4D7EEFEFAFDFBF8FBF8F3EEAE71FBF4FEFBEEF9FDFDFEFEFEFEFDFFFFFEFDFBFEFFFFFEFEFEFDFCFEFEFEFEFEFEFFFFFFFFFFFEFEFEFEFEFEFBFDFEFFFCFDFFFDFDFFFEFDFEFEFEF7F8FDFBF9FBF9F0B669B9090000030535703C040300040B00062F160A080300040401030432C6F8F3FCF2F9FEFEFEFCFCFFFEFCFAF7FBFDFBFDE47F2F220C060206330D01000007071466CADFEFD9F0D7FAFFFDFEFEFEFFFFFFFFFFFFFFFFFFFFFFFFFFFFFFFFFFFFFFFFFFFDFCF7FAF8F5F9D66B0600010104017F698DEAFBFBFAF8FBF73DB6290127A5F8FAF9FDFCF6E96A736E1A1314070D170C04010503040201040D726305020104020C2A5DE3FFFDFDFCFBF3E8C9589BDCF0FAFEFAFCFCFBFAF8F8FCFFFEEFF2F5E7EAE5EFFDFAF9FEFEFBF8F2FCFDFEFAFAFDF8F4FDFDFCFEFBFDFEFCFEFEFEFCFCFDFBFDFCF1FAEDE7ECF8FDF8EDE9EDDFF6FDFBF8FDFEF4F7F7F5F8F8EDB926030200000503040300010404000BA1520E15121724261702020C3AFBEDF9EEF8FAFBFCF5F3FEFDF8FAFDF6F9FDFE99355E842000020ED015010100040301041C5C9C44A348DEE9F4FCFDFDFFFFFFFFFFFFFFFFFFFFFFFFFFFFFFFFFFFFFFFFFFF7FCFCFCFEF3EFD5570600080A050B1086E5F2F2FEFCFEFEFDD524010021A6F6F5FCF7F9F2990C0106056F6E23517E1D010A1710020F0800000202060E07010F51D9DCF8FEFBFCF3F2BC65884C081ABBF4FDE9F3FEFAF9F9F1FAFEFDFEF5C56C90619CFAF8FCFBFCEFFEFEFDF3FAF4F6FEFEFAF7FCF9FDF1FCFAF1F6F3F2F9FCF3E9F6EDF9EF7E59C6EFE4E5A4669654C9FDFEFCFEFAF8F1F5FDFEFBFAF6BD210200040107040002060501010000000D8E6098D6D99810092DE4DCE9EAEEF7F0F0F8FCFCFBFCF6FDF8FDFEFDFAF48F1C0100031000080C03000001110F00010208011048915EC9FBFCFDFFFFFFFFFFFFFFFFFFFFFFFFFFFFFFFFFFFFFFFFFFFEFEF9F7F9FBFEFEE06200000A02965492EFF8F6FEFBF1FAF7110C0101199FFCFCFCFAF5FCF2851B02070309070C1004080014ED20020000050402020307010C0A0167E6F9F8FCFEFEFBCC54000FBFF9EFF3FBF6F9FDFAF9FDFCFCE9F2D943954137AEFEFBFAFDFCFCF1F9E9F4FCF6F5F9F9F9FCFAF9F1FDFAFDFEE55C88F4D39566A061AF14087D832E02030747A35DC4E2FEFEFEFAFCFEEBEAF4ECFEF9DEB50F030F050201060700051612030A05040988E887150B0000010C131BC149A8699CF2E6EFFDFCFBF0F2FDFEFDFA9228040100030E0108000002060D7E7306030609050608040B52C4F9FEFFFFFFFFFFFFFFFFFFFFFFFFFFFFFFFFFFFFFFFFFFFFF7F9FCFDFBCE5F120800000501191279677BA3291B3BFDF7EADF2E0023A6F3F2FAF9FEF08314021507000000000100030B01130F02020A0E06000204082136A237CFF1FBFEFDF9F5F7D054090FB6FAFDFEEAF4FEFCFAFBFDF8F8F6DB5EB04A0A043DA4C18495FCE96A9377AA669A838BDDF0F2D345ECF32B1414180E0D19191C5FB0081A01000E0C030A0901070D0B2220CCF3ECF1EBF1F9FDFEF7F1FAD4201202010000060702307F5001030100000C9377877907060A030001020F25A86581977F97668B88FDFDEDEEFAFEF58817010910010000001CB0480D2485820001040748A9133AC5FAFEFFFFFFFFFFFFFFFFFFFFFFFFFFFFFFFFFFFFFFFFFFFFFFFAFAFBFAFBF6C0480D001306010202161116688B39CCFEFDEA3E0D1026A5FAFAFDFDE88E788D0D6466280A0F0E040100010000020006914908120D0001050B1C1066E2FDF8FAF9F6F9F0D1BF49ABF2EEF875A5E4F2FDFCF7CB57F3F6A8420D0B2F7A7123337F7E241B9558261219737C36282943AE242513AD34000A08000307101C06020207020502010100050A00030221252C32AB7393F0A46BEEFCEDC29A1D120500060E0A0E271B0000000B0804826628180201010001191001111E1013161219152375F4CEF3F797978B420A0238971D000E00A14D91CDDBDD89100C00040E1C1B0E2075D6FFFFFFFFFFFFFFFFFFFFFFFFFFFFFFFFFFFFFFFFFFFFF4FAFEFAF9F9F9F3BF51000501060E0A0A06518E749A56DAE9DF2503178FDCF2EDE3FDF6D7D88E0E231D0D0350692002040312060502012A1402050400010300000247CEF3FAF7F3F7F2C186B8B848C75BAB4798848C9A72E6E69648D35391280213231806111D1D0E032C160306081E1C06000315390C0C05300E0103030302010618160A0305000103010B0D0105080007020302032415203A649065E2BE342E08050100336745040408190607030200281E05080A01050415AD0F100701140D040B0F0D041A37403A90BD7477662204112E090A0400380F2B3BA4F7E1B87D1701090006043BAEEEFCFFFFFFFFFFFFFFFFFFFFFFFFFFFFFFFFFFFFFFFFFFFBFDFAF7FAFBD79DA3690F0900110966480D216988360F5145430C051D918E98A393D89BA8EC82170C0B04001F2A0D02061B8C1302000916080404030B0D010100011580D5E7E6E7E4B77D7E59440D38143816352421281B444B420645113710060501000103030E050004080101070404060205040404040902000002040706010A3DA870080102121005070E0A0403010A000B0219130402051E3518475E9D0F000E040416291C010421693E050B0904080D0A04040202030D4204060B1A1317161F55670A0203070022422429260D0205080903001F0A120D0A7C9585A7577A00030F010B1B68CBFBFFFFFFFFFFFFFFFFFFFFFFFFFFFFFFFFFFFFFFFFFFFFFCFBFAFCFEE8AF80833D0504080335280801274804070A0701000A0E343F726C2B5B687C937650561A040201030505020E480E0402056D3C07020B3F64010200003A7662545455533C26331403030103060304020305020105060306000302030507111E201C0E0E10010308020204050909020104240008150B020E0C0303010522563A050401465B2E38522A110A160B0E1A0A703516100B1009141025550F030802020201010309144420010604040255500C0002030303010C121D754A6E915648755005010111090E0F080402040903050505941A752E062F282037164203000800060A1D73D5FFFFFFFFFFFFFFFFFFFFFFFFFFFFFFFFFFFFFFFFFFFFFDFEFBF6F9FDE9B2852602020104120D02003561060400061C04050A0515382C000D30352B3D353B1202010506040D18010323070801462A0401092E45010000052A3D1004020001020000041502030F04150A131403020C1700070500010300061A5893917A3E4A420B010C040C0D050F201819159721237037115F48010200000504000308012B70604B453E56286D20708B2048376654546431772F7D06130E100402070B090505030C0102030002023A3B0D0D1903070001106C3BAA6C81986A582E37070F104C45664D1D17140509000004085F134E1E01060D0C05030A1C130000003A4B6DC3FFFFFFFFFFFFFFFFFFFFFFFFFFFFFFFFFFFFFFFFFFFBF9FDFBF7F9F9F1EBA22C0300020B5C3F090F42520606040B6E12020808151204091314000B270509100500050703265D0F0F740B06010106140F0B140A170004050B040004050715211504215E341B7A296E4D61591E132B75021C190103030607184C777666373E4610020608433F204078837038DC8082897C4A9271100B171200060002090303333E46473B522D591A61771D03597B82BFAC63C848CC3C713B56220626412C060005120E020300130D02040831590C0500040D5C369B58353434450D020F47505C7FC68B67785B2414020B040409030103111346561823408759131102326CA8E2FFFFFFFFFFFFFFFFFFFFFFFFFFFFFFFFFFFFFFFFFFFAF8F9FCFCF9F4EADF9C450109024681480C384A030708060D73230A1B28503602153B440B3473424A401201000D182C5A13136D0E0007160F4349415337590005010E03000A162E5D826D4377909283EB57BBB7C8C18C5F51D02D8266201C150F120F0401040501064E3F1609094641456E7D8D7F74827F7C73768087804D3F754C1D0A0B130D0F080201304A2832562B1E00150B04543C89F7CFACDC94DDABE58A6A250428472D0401134A2B0305004C2E06040832590E000004080112532804060107120C25568F88ACD5ACB1B3998156211F1B040C08060449498BC57F8CAAD58B551C121E65CDFFFFFFFFFFFFFFFFFFFFFFFFFFFFFFFFFFFFFFFFFFFFFDFEF7F9F9F3F6E9A86A2B01210D68390B053F53060702041711651E677877400114404F1244888384551602012C4A17000507020B040C4F274E7B788073670004101206022A4371ABC6BAA9C58FC9E1F497CAF5FDFEFFCF96E199F1987B7B57283C280003050102085D823D140503013A4907071A73040A006C215981A38170C3887A37435535491700020301216B858E750C5C1501021292FAF9FAF7FDF8FEFAB23C09000205050A11195B2E0509015E390A12170B0B020005000007020714130A05061544706CB8F3C58A828373A1CF978470712C020A08045989BBF5F6F8F5C37A6C0B2A699DDDFFFFFFFFFFFFFFFFFFFFFFFFFFFFFFFFFFFFFFFFFFFFFCFCF8F7FAF9F2D79A8B31044B2A080209050103020514117522DF4BBAB03B0C19130A170B10134847100C0E053A5B170C1B030506080166330D49472D6304040A294A1107386494B097819A898B978EDEEBEDF7F2F3F4EEF3EAF7D975E7F0924C64450E0A080E020006553C0B040704050500030800040000110A035480484292A8D585AFA57E972E000303012C6632C0971E821F0406067BE8E8E9F0F7EFFAD58C6F1F00060605213E170A050200010C0502313D100305020D001A150510492F0204143F94CFF2F2E4AF3B120D3EA6C39EF3E9E8820B0000050681E0D8DDE8EBD9610A012680BFEAFDFFFFFFFFFFFFFFFFFFFFFFFFFFFFFFFFFFFFFFFFFFFEFBFDF4F4FDE7CAC85B020076400D080A0F0B070D1A5718ED86DF89CB80233A613D31683723313F330D373715030007275311000010060703000A130D03010B042068130103337AB4AA5C2C18080044BFACD4F0BEC2C9ADD5A7E99DF5EBD481656564593E2E36100A090905010201080B06020409020E020102120B01050F3DA0E2BEA3EEC2A88A310C130F040102206D010F040E16131250878991BAF0FDF17D7B7223010403001E554103010600050B08034F591114000300037013021461420702124B86A2B6C1AADE81563D3980DEF7DBA1DBA0210301000A4F929B8F92AFC73F020309207FD6FAFFFFFFFFFFFFFFFFFFFFFFFFFFFFFFFFFFFFFFFFFFFBFBFCFBF7F2D0926D400E0803482715396A45294E44D2189497979D9988516081534D8D683954584B1D4E8A60170407227A1C010064230A0603022B4301000210110200001021255B846360230002868E7172C46E7C9158B12FD60AEB91E8FFFEF4F6EAB4754F353B00000400000234340E111101060800081205090B0A1A5395A0C5E5FDFEFED0833F30302113090306070100556F6E2C020103278AF4F1E26D7D1B0A05000304010305010800010B3F2F1104012B510411021283200706010003030012548D48846A8667929CC0CDC0D78F6A98626618140D09050302021F72907406072F67ADE1F9FFFFFFFFFFFFFFFFFFFFFFFFFFFFFFFFFFFFFFFFFFFFFAF9FCE5CED2B13B010A5318BC9D67BCF68567B03FE973030A0503030101040705051528330C031535324C4B21030A1A00000A0096230001030A4B6D0B0408030B0A0102013A54516B66190D080A00117C2B5141494C2E661F7C08895EE0FEF7F9FEF1D0A78A8A5C00051706020A5A91592E5008135202035021053A320602030D2DB8F2FCFBFDEC8F8DB1916E5B0817000F017DA1AF4904030237A2968C7E8D14000000070700071200030A070B02673D0B0003406F010501000000000001060601060E05000352525931578FCD943D58286E0B8D9D746D5E62310303194D72040014024C837BA0E1FFFFFFFFFFFFFFFFFFFFFFFFFFFFFFFFFFFFFFFFFFFFF1F8FED39571420201067C2298E7F7F8FBBEA7D79EE58A71676C676B273F6C3C0B2E4D4A5E3074424E471D0804001E541C0000000A030301010204050007286302000301090805052D5F2E0300000A010504040305090505050A0575EAF6F8FCF0F3FEFCE8AF6A1D0807373016040445573B65001CEC0961DBA175877F654638747FA9E3BFE9F6FBBEAAD8F6D69A67010001040607080A05000206080B0D0D060A010105030100000011592C01070509070200000409020700000C0303060100113C3F0001060403030504050670952F06080501124FDEF0F4EDA038273C03144C04000D06336D9DD8FFFFFFFFFFFFFFFFFFFFFFFFFFFFFFFFFFFFFFFFFFFFFAFBFEE09137070000001940D3F7FDFBF9F1F0DEADE182EDE2E9EBA07A89874C438B511D3C968F4856363C411E002A6E0F07040002000402053E4300020D1E7B0300374618000707050F1100075B2103014A25463F2F6E6B6F7375EBFDF9FBFBFBFDF9F7E5B0BDB65C279E66210701070B0A3563109070DCA699D2A17E804E3C909191566DC0C9EEF4F3F4F6D6A4DA7C0B0000080202053A3D26411A650114590E000000020201000113653401070501050200000205030A00040100000002001849470208030A0000010100033C49130302070872DDFAF7FCF8EFB0A0BE6E9ED2450201000569AEC2E0FFFEFDFEFCFEFEFEFFFFFFFFFFFFFFFFFFFFFFFFFBFCFEFEE29131000002140254919094D5C9A5A0878892D98CF5FEFEA3784520203D590F072B50290F3F3E4D4A1E0008180B0C610C02000006063F490200020623000427512001010B2D5F78797C7F79807BC5A4867D98DAE5D59BE4F59BEAFEF0C4CEF4E6C4C5C1C4E19877873A080000000026740E32DF952E1D68825024130F271E1B0D455550A6EBD7CAC4B4CDBCD462821F060000094576A0BE6D6A5E29B2304B1D0001284725030F130F532C0900070404344227660206075E21552302080D100D07180002000A42381504111D100205021170F2F1F9F0F6F5E9E7EAE2F0F7C34F12010338515182D4F7FCF9F7FEFEFEFFFFFFFFFFFFFFFFFFFFFFFFFDFDFCFDE18A270512114E3620221F25645629253A6B3D891E92A789783C0D01071215051915100C020F1115130B080806020A7A1502020D0703121303050206030203091E0C00070A1C598F827634B7A884908E507FC8A08D89919C9A8591967C60A3C280708E8996CB907F9644113C3E12080C200A0A7F2F10021423140305090301000010103C85925E7C9D8CB16BCDD69214010803495153B0ECC49ED47EE5915C1D06022147340703000A64280602064070807C6A7C7C7E5B656A7971280202020613393C070003493815030205040309125A91EEF9FDFAFAFAFBFCF6F5FDF7E6490A08030E112E82DCFCFCF9FAFEFCFEFFFFFFFFFFFFFFFFFFFFFFFFFEFEFBF9E2953102080D4A2F0E0A0204171103030E220C230E262E27290E0308070102090B0601050000000103040303050D132A0C0302060205040204020103000A01060601010406102432322B0780432A2A2C1953702C22306F29297434291D225C5B1F264C53545C3159461A073D411003040903012F0805010404010002060705060001053548381D374B475D27809432060A07044F41185797948A8C8494772E0907040C19120304040421120502024A7C5A6DA340858ED097CC96722E09040104123A5303000022150604030507051A6F897B88EDB39C9CD2CCBDE1C1DA99A1740C0A0301012380DDFDFAF8FDFEFCFEFFFFFFFFFFFFFFFFFFFFFFFFF7FCFDFBF3CA782300051D0D06452645233031223D1C4A115E054D5D1B00070D25352B603F022D572C392D0E2F392A391D41115C08457465223D3C28351B54635B5E604912202D3135322A16050D44410C000804041A1E04000D3300002C0804061C4432133A3D1B1813031B110607191704021D4E00000B0102000438320800010403040007051D150A070F12141F0C2C2701020602021E1A04112C323125362B251104030609213C28040C0001402A09061F3224416D0E32389782A77C3A110605060607182200030407030002060705031A733D2D2984432D2D716B527355862C3E5E0A04012A3B285AC9FBFCFBFCFEFAFDFFFFFFFFFFFFFFFFFFFFFFFFF9FDFCFDFBE097360001083A607A39535F9599888C36B865C85BBC8E64512D408DA283A3695E8D9D6F747143465E83A77EAA60CD5CB39CBC8472788FA485C8D2CECBBC7B62839C816278A0843E1EB3213B585B6066643C04184514595B451439243E65635F6B6B58291F5C1F364D27322E263239A60D5D0251123F217661355B310203020405012E34131E2D3031243C1155120401060439302A38193B0C481F400D174C462B244A745C3429381A854440582E092D483748060F2636382F0D02010105112D370003020306313110030503000A2A00011120070C061C2014160D2B0608230C03003645396CCFFDFFFEFCFDF8FDFFFFFFFFFFFFFFFFFFFFFFFFFFF9F8FCFBE0A1572203075E87511E1F83DFEBE5C57CF5CCFDCCF396B88D3957ACC9D3BF8FD5DFBC8C6167522F61B4C1D6F7D2F6CFF0999B955E5E97AFBDF9FBFEFDDB899CC8BE948CB0D2D6AF70E863A4CFCEA5969272597DAC68CFCD825C72719AA6B0C2987DAC885CB2816B7667716B81997CDD62CE56C166AB7EB19768824F0205020500023F6E5D7E7B565F637156871D040106085269677170A96EAA275D14226E7C7C899890718184AC7AC27CA6986A5C6D6A5C794C0E040905030605010304114A7C4F552C0A074B4A1800021C33144106551301000C0601060D23371934173E10030717306DBBEBF7F8FEFDFDFBFDFFFFFFFFFFFFFFFFFFFFFFFFFCF5F8FBFCF1CD893905082A5A25070387EDFDFBEADDFCF7F6FAF9C6B67D2C28649CCEE5D9FDFBCC66211E1D0F4C7B74E9FBFDF5FEF3A943311F172C3D68B4CCFAEFAE6D6686594F868E6995E3DBECDCF6FBFBA95735599BBEECDEFAF18C815285BDCEE4F8D18B969C6174C49A5A73625BABE7D9F1D6F1D4F4D8F0E2EAD6752A25000203000004206EB1CD8A3447827F883F0E010005092084A998D9E79F980D2607162C4993DEEDBD8A95A6B9D2F2D9F2A78799726BB17D7502080007050806020506052E86B28B4E1101201E0A01002E6D60B94F9C0E0000030303020B407E729C6A8B2C31170545AFF1FAF3F4FDFBFBFEFFFFFFFFFFFFFFFFFFFFFFFFFFFAFCFFFBF8F9DA770F02122B81480E2585D6F6FCFBF5FAFCECFDFFF48A813F0E4F92A9FAFCFDFEAA260601001A382735BEF1FFFDC1F55F0F01212408000E458DE6AF5B201416090C2728185AD5FEFEFEFCFAF8C24703183570F0C0F7BC843B0A2D72C6F8F1D9BE78311628B1C35122112A8FEEFCF5FEFCF9FEFEF8FCFFE67202040000070300173379BBAA55111051A4730E0100050C04076CB0BDE8A95224000A0101041666BDDDD3AC5B353093ECFDFCC16A363281E197460603000101000504000823518EBB59350A000303041B333556C1EF9F3F0A020407060200063386B8FAD5A0546A1D0640A2E9FCFEFEFDF8F8FDFEFFFFFFFFFFFFFFFFFFFFFFFFFEFCFDFBF0E2BF6A14030E815F22066D6193EBF7FBF9F7FEF6FEFAD9C06118082E7CC0E5FBFCE586290E1B1857661E1761C9FDFA53BF2B0020736D1F000943B7C8728B1A07141C1106032284D1F4FDFDF7FBE0862D0500043FBA68EB4F460E001E89DCF3D89D8A6C06206DD58824020155CAF6FEFDFEFAFDFEFEF2F7FCD663050D141F0E2F2D6BC0E5A7917B330139D5AC1603000A040604184092C2628901010B030D1D3B9DC09487711400052BB2EAF29F31024BC6CA74B02E0026090002051D262B70D8EAE1B156140503060341A29465E5E2650401040D092015001150927FFBFB9E2D1A06182D73CFF9FDFCFDFDFEFDFEFFFFFFFFFFFFFFFFFFFFFFFFFBF2F6FBC7879BB44708022D1201051C0F6AD2E4EFEDF6FBF6FEEA79A8460F0E00114797F9E8838E6F1965663A19073D8F61F1D8E144081683787D751102196CE0D1C323014C6E3D130C0D2B43D4F8F8F8E2738E631204050A4BDDF6CFA63A090649B2E6E5C55D110C8471A3702906006AE8FDFCF8FEFEFEFFFFFDF4EE84160213538348C8BEE2F8FBEBE69D320837B95C2C05020006190D0E5DC1D1922C060802092C6D64B0EFC0630B0003069B4BB2D5E76A0C1D70D7E3F7C820BC1E01081C8EB76E82FBF2E29E32050B111A20267AE0E4EDE0BF02041B211F8557040C2874DCF1FBAC260402389AD9F1FEFCFAF9F9FDFEFEFFFFFFFFFFFFFFFFFFFFFFFFFEFEF5F4C25112250E0007110500010005131E6BE8FDF9EBF0FDF4BD510F0003051B6EDBE8906B7B2E071418090B0F0D4FC8D534D03205956D9CF3E481030D56CC4DB02D09121553793E0226B2EFFEFEFE6E9363180404000333F3FAF152A13302033EACD2771903142098759B5A160668DFFDFDFDFEFBFEFEFEF5F0D9690C03030B1D6CC52D82E6F2F3B041101939C4D8D22107110F050907172E2E5EA002060F030D1C4FAFE6EDC243071500130DAB79A669125B9068FBF2F5E9DD340A032BA4EFF0F7F9DEA387540F0B67BDDAE0E7F1F8CD43251001BAE44F1812030151F0FBFDFDEE9E1B05042169C5FEF8F0FEFDFDFCFDFFFFFFFFFFFFFFFFFFFFFFFFFEF7FAFEC24F0301050000000000170D0207649373EDEAFEFDFCFD4809010407040619397FEEFB93170006030303013CC7FAFBE14507081E94E1F4E78300031B54B32D05000042805E17052CC1FBF5FBFEDF671B100900000440FEFDEBC2300D0230868B767501187EDCF0E6701D12096EEEFDF8F2FEFEFEF3F2FAF3747B7003050104101064C5E6F8FEFCAD350039C347150500ADDB4C9475010C51A514250000030100005496948E6307000007A9619E672207136CDEEFEEFCED3DAD2A19A6F7FBFEFDFCF6C8721F0F7875A7868BD7F7F2F8DFBF000B2422A85C0503031A7AF7EFFCB5170B03003EB6F8F1F4FEFEFDFDF8FDFFFFFFFFFFFFFFFFFFFFFFFFFEF1F8FEF8EEBF52010400009B56120203010E1D5FF5FBFCF03DC1330D051C7469001667D8F78575832100677C2C00051B50E9F21E01050CA0F8F7F5E4890D38EEEE9C2904030C1B140104289FFEFEFEF2711902026C880F03102246EBF8D4A732061448927F001384FCFD8F95521C667573D6FEFDFCF2FDFDFCFAF8D3721A050B010401011275E3FEFDA9400F030A44B02A03091EAFDF7315060950EDA20300004B8830379DE5BE51170D040700241815906710015FFCFDF9E8FCC638110534B1F3EEF7FBFEED8107011811171A6DD1F3FBF7DE4313010212210F000F6FD7D3FEFDFDE3931A000D1256C4FDF8F1FEFDF9FDFEFFFFFFFFFFFFFFFFFFFFFFFFF8FAFBFDF9C57A80B53B0600C55C1C0504020609133655FDF9B052AE3A000B191A0D0318508AE68B2A05657A7C662C12012BE245B62B050129ABF1962F28081071D82B070401080A03002F8F72BDF7F6FB5C0E0A06181F040301052FF5E4DD5B1403010A1716070827AEF9CC7D30091E6EC8F8FAF3FBFCFEF7FDFBF9E76B0606134D834DC5C7DCF8EEF6F08C2007022BC4260D0204BB711708060A5CC51B0308062CC6CCCBE8FCF3B13B0305010804020A1F130A638E67FCF8FBFEE02E041E95EBFCFAFDFCF7E6800E00050203031777E4F6F7DB14030009040703030A2988F9FCF2FCBF2C0A064499D6F2F8F8FCFDF5FDFEFFFFFFFFFFFFFFFFFFFFFFFFFFFBFDF6F6FBECD7E0F5C1330B2F1580490C01020407062FE04DE728301007060A0D061A5FBEC8A62B0A0374C5DC982B05000D46AD510F030078799AD472050E41C7B82F3A6F30060505020E2A2EB8FCF47B1E0302010609070700052EE7F5F7AB3000070700041F66C0E4E8F37119010369E3FFFCF7FEFFFFFAFDEBF5EA69001556926DC9F5F4FAF0F5FBB08977292B995E0803009069994C030F549462950006298A7CD0F9F3C59D9D6A0010893F1402020300021F6CC9FEFDFEF8F7A9241243B0F9F8F5FDFAED7A0F03006C5B1A015BD9F6D863961E04059138000203082B87FCFAFAE6891801123881D0F9FBFBFEFBF9FAFEFFFFFFFFFFFFFFFFFFFFFFFFFEFCFAFBFCFCFAF9FCE25C050008371E0803000102000E4419E31F040705010307050B2853B2930D074EB5EDB387701F0719A0F01F090200332E84AC2C00042190B61D1F2E4E5F2F040500021EBAF9FA4C09040207084971060D010F68E4FCCF4302030504020C3AB0FDF6FD6A1B040168E2FBFCFDFDF8FAFDFEF7EA917B5407307FC6F8FDFCF4F6FCFEDE9F5111123D140B000539947523000A66D9903F00051566878CA5A3B2A77A7F3D1BB34A813B0E050100025FF5F2FEFBFBE965121981DCFAF9FAE7A89D8A080A017CBEB3AACBF1F9EEA84D0C0104C24A6F5001031765FBFEFDFBA42002003CA0E7FBFAFBFEFBFCFDFEFFFFFFFFFFFFFFFFFFFFFFFFFFFDFEFEFDFDFDFCFDE65B0408070401080600020400000312DF1F08081F583F0808011148D9D258162D50613A2F3912257E73EE1B841F010608323700010D39C3BB1207014EADB7772D00000E4864E74209045E9C9CBEA3030C00052AE2F6E694340604071456A7D5FEF6FD5E5F491768E2FAFDFCF3FBFCFCFAFDF0B67430001774F8F5FEFDF1FBFDFEC86015031F7B8F21010402442A0604498691D77B000105203A2B31326D9D6732221FB1B4DF5F130200003FB7FCF8FCF2F56E8E231158C5FAF9FCEDC07A805D0F53796E5E9FB69EDFE07B8B020B03C04AA07101020A349BFCEDFD9F1D013598E2FBFDFCFDFEFDFEFFFFFFFFFFFFFFFFFFFFFFFFFFFFFFFFFFFFFFFFFFFFFFE850568C2F0D60460C0001070B02168F6F170F6996C4B77A460C0D44BAC08C4C100107050206091E7190EC9776246332040204133B86ABF8B310384D61BFF5B1605536127089DF510E5E75B8B9B2D2856E031E807CE2E598703A09051656BDFBFEF4B03F4837146BE4FBFBF3F7FEF6F6FAFDF9E469030A48BDF9F6FAF9F8F6FBFCA835020740B6851D0006020406070A75D1ADCB5C150001090202042482AD57031A65D2FAFEB3430E02012FA0FBFAF7F6F190731F116FD6FCF9F8FAF2BA775008453B4B41483F3082D28B6D04080AC4ACD16F07090009319CF0C15911012887DDFFFFFFFFFFFFFFFFFFFFFFFFFFFFFFFFFFFFFFFFFFFFFFFFFFFFFFFFFFFEF0A1C9F8A43DA46F1300022B60137686757D7CD0ABC6F8E5B17C758AB2BEB8832539353F3E3A677087E383EB8482D59731103E7983A0D5A1B42F7A7D9CAEB5B08F7C7E3CB992D49E27A8A1BEB881ADD77E7A807F78DFE296757B7D73737D7BBAF6F6A98B793A4475ABE7D598F08AEE8BDDA1CBAA6F444487ACF2C4ABEBD18EB8E4C76A395B96D0937E77752B5E334C79B4F3CD806A64061B611D050544B1BA60443E62D2EFFAA8843D110436AAF7F1F6F0F5EE91673BADF6F9FBF6F8FCF6B53E0C09063B451C03286E9ACC21080D6CDAF5F0723D480B0D339BD6765B13032C8BDFFFFFFFFFFFFFFFFFFFFFFFFFFFFFFFFFFFFFFFFFFFFFFFFFFFFFFFFFFFFDFBFEFDF8FDA3B96E1E0804346D1A80038E959190757E929FBCCEBECC908F89834B573F4F707A8E8B96DF08D0A097978A6F3E4B898178703EBF3A5862A78A738A6B4F935B6316748A67839598915356B538DD8B6B8E8F8D71395E9994908341468F968D958F797E7A74A4AB36D60C8E1173265E424A5D6A9658C770748A5F43768B8C847D838EAADAA1999573887E81C0C68C7C738C6907228661191E5E886B4555445182CAB16C80813A3272B9D2A2E099E99BE89885CDF0EDE5CBC3F2ECC24B49323F313B34396B858BD5746D6EE3A9D5EAA87F77492F7EA7AB385F122588DEFDFFFFFFFFFFFFFFFFFFFFFFFFFFFFFFFFFFFFFFFFFFFFFFFFFFFFFFFFFFFDFAFEF7F5FCF6C888663F1007090B06040903060A6141051648706B85080001484E1B02053A5B2609198D0686250700085F50020C3A770729860F043F514A4F2F072B5C410004000372360A05040234772E7A0E8C1B030003020404060A0000020005060304584C50574271248A0B0204040B0A0506103B5E3B60366E080454581E043C5B270A21741B000D6C256B3F50570430760E060700206925305839080305206C287233623674445152386A217D1E8E17BD45ACC0A79C8C7F7C99976F60847786525A707E806D97979DAB9EE53EABBDCA8341653572257208080960D1FAFCFFFFFFFFFFFFFFFFFFFFFFFFFFFFFFFFFFFFFFFFFFFFFFFFFFFFFFFFFFFBF2F7F5FCFDEA9BBBEBB55E392B40071600040F0602090D0600124E670000030D030101030707060904020000010808030D0C0000020204070B040B0B03090E070206060A010901040605090E070001030905060B4A150A0C090403010001010608000D0905060200000204030B0005000114080403070C0400000305030E0200060F0A010712050B0F00060C03020205060505000400120102000404010103080600000303020000000308070404000303001D7B2F7037050618515E021656B04E534F190C5152426F0C080400048C2A733B545102040A000205112A4E95DFF4F8FFFFFFFFFFFFFFFFFFFFFFFFFFFFFFFFFFFFFFFFFFFFFFFFFFFFFFFFFFFEFBE9FDFEF9FCFCF9FBF9F5C57DDF44500A060308090703010102060B01070508010806222B14110907090A010A05361D0200040B0D0D0E070013040404050402183113050C00070B0B16090606040A1941103F33DD6948202C260D070408070B1904060606020004040A1700030D040607321A0804080B0206040A02050906040C202413050511073C10020904090C080404080B100808030404090504080504080407060E090101000000070D06040D0B000203050B04000405010303080B0107000100000103030306030E001004000C070006090609020A1C3F7AE7FAFBFEFFFFFFFFFFFFFFFFFFFFFFFFFFFFFFFFFFFFFFFFFFFFFFFFFFFFFFFFFFFFFDFDFEFAFCF5FBF8FBFBFBFAF7FDFBFEF944391B090A261F060D1B19010A282E0D071010798F3130462F0D500A2D22B76A10072345392D0F3B0C3C1A25242B251C5AA25E2E400D39384E461A2B2716254DCF64F45BDFFAE08ABA8C261B1F1D234D1A1E3620311D0C24271C360A195A03321CC27E44261C2219371D251B2D14272817658F5D2D182429ED651645191C1A272926443811463D17091B241310191C16142929250F0903090806031B260612050400071531110303090E06080E02020905020C06050B080417330D04000506020A0B0400090E01040435D4F9FEFFF8FAFFFFFFFFFFFFFFFFFFFFFFFFFFFFFFFFFFFFFFFFFFFFFFFFFFFFFFFFFFFFFDFAFEFEECFBF9FCFEFFFFFCF8F9FFFDEAFDFF653D56B2983E559B9F1739B6A3494F3E3889A39FECEEB760F253E562BA763F4A97FEF8CD6EDC77C187997E88A6A7A28859B9EB62D1F9ECF887ADA08EA845F9FEFEEBFBF7FCFCFDFDDCAB9CAAA7FFFE84BD85B59E56748E89CA5649E30FE45ECDFAF29D9BBB87C780CA81B5A296B0A2AC9A94AA98D560D0CC52D46DBF8BA1AE8AE19C71F7A3363F5B976B2F588D6D67AB57BE59291D272C382174A41A3D3C31033E0FC252221C1F3035381B1A1E031B1F21132017182246BC2132053C41172A17050401010500030C78E7FDFEF3F8FFFFFFFFFFFFFFFFFFFFFFFFFFFFFFFFFFFFFFFFFFFFFFFFFFFFFFFFFFFFFEFEFEFEFEFEFEFEFEFFFFFEFEFEFFFEFEFEFFF7F5F6FBFAF5F6E2A3A0EAFAF0F3F6EDF0F9FAF5F6D76CE6FEF6FDF7FBF8F5F6F9FDFAFCF7F7F7FCE3C097CCF5F5F2C891C3FAF7FCFEFEFEF9FBFAF5C285F7FAFEF8F7FCFEFEFEFEFDFAFAFAFAFFFEF9FBF9FBFAF6A08BEFF9EFF6F540E7F6FCFEFEFAFAFBF9FCF8FCF8FBFAF9FBFAFAF7F6F8F3F5F4DF59DCF6F4FAF9F8F7F1BF69F4DA4FE8F5B68BC7F2CA96B0F2DA94B66FD57CBDECF39796BD81EDEFE851E471C786A89CAFF2EEDCA996A450AE88AD8CAB9DA1B182C56FDF4EE8CA59956D1F0401000311408ADCFBFEFCFBFDFFFFFFFFFFFFFFFFFFFFFFFFFFFFFFFFFFFFFFFFFFFFFFFFFFFFFFFFFFFFFFFFFFFFFFFFFFFFFFFFFFFFFFFFFFFFFFFFFFFFFFFFFFFFFFFFFAEBDBFDFEECFAFEFDFCFEF8F4FBF2D5FBFAFEFFFFFFFFFFFFFFFDFBFFFEF9FEFFFCECE1F2F8F8FEF1DDEFFEFFFFFFFFFFFFFFFFFEE9D3FCF9FDFDFCF8FDFFFFFFFFFFFFFFFFFFFFFFFFFFFFFFFFF3EEFAFEF9EBFDD4F9FEFFFFFFFFFFFFFFFFFFFFFFFFFFFFFFFFFFF9F7FEF9EFFBF9D2F4FAFCFEFFFDFCFAF1E5F8F7D6FBFFF6E3F1FEF5EDF4FDF7EBF6EDF4E8F5FDFDEEEFF0E3FDFCF2E7FBD4F1E6F0E4E6FCFEFAEFEBF3DDF4F1F6F1F5EEEEF3EDF1D8FCE9E965C47AA4611C3F4A5274DEF1FBFEFDFDFEFCFCFFFFFFFFFFFFFFFFFFFFFFFFFFFFFFFFFFFFFFFFFFFFFFFFFFFFFFFFFFFFFFFFFFFFFFFFFFFFFFFFFFFFFFFFFFFFFFFFFFFFFFFFFFFFFFFEFCF7FFFFF6FCFEFDFBFCFDFCFCF9F9FEFAFEFFFFFFFFFFFFFFFDFDFFFEFAFBFEFFFCFBFDF8F6FCFDFAFDFFFFFFFFFFFFFFFFFFFEFBF6F8F9F8FEFEFCFDFFFFFFFFFFFFFFFFFFFFFFFFFFFFFFFFFFFCFBFCFEF9FCF9F1FBFFFFFFFFFFFFFFFFFFFFFFFFFFFFFFFFFFF9F6FDFDF8F7FEF9FCFDFEFEFFFEFEFDFDFCFBF8FAF7FDFFFCFDF6F6FDFCF9FAFDFFFFFDFDFEF9FBFEFAF8FCFFFDF3FDF3F6FDFCFDFBF2F3FEFCFCFCF8FBFCFAFDFDF8F9FDFEFCFAF7F2FEFADEF2E2EE9389D0CCDFE6FCFEFAF3F8FDFFFCFCFFFFFFFFFFFFFFFFFFFFFFFFFFFFFFFFFFFFFFFFFFFFFFFFFFFFFFFFFFFFFFFFFFFFFFFFFFFFFFFFFFFFFFFFFFFFFFFFFFFFFFFFFFFFFFFDFAFAFDFEFDFEFCFAFCFBFCFFFDFAF7F3FDFEFFFFFFFFFFFFFFFEFDFDFEFDFDF9F9FDFBFAFBFCFCFCFFFFFFFFFFFFFFFFFFFFFFF9F7F9EDFDFDFEF9FDFEFFFFFFFFFFFFFFFFFFFFFFFFFFFFFFFFF8F6FDFAF8FCF1FAF9FDFFFFFFFFFFFFFFFFFFFFFFFFFFFFFFFFFFFDFBFBFBFDFCFAF2FDFEFFFFFFFFFFFEF9F1FCF5FBF9FDF3FAFFFCF8F6F9FDFBF5F9F1FDF1FAFDFDF3F1F8F6FDFAFBF6F2FCFAF7F3F9FCF8F9FBF4FAFCFEFEF1FAF7F8FDFDFCFDFDFEF9F5F6FBFCF5F8DEE1FAF7FAFDF8FEFCFBFDFEFFFAFBFFFFFFFFFFFFFFFFFFFFFFFFFFFFFFFFFFFFFFFFFFFFFFFFFFFFFFFFFFFFFFFFFFFFFFFFFFFFFFFFFFFFFFFFFFFFFFFFFFFFFFFFFFFFFFFDFBFCFAF9FEFDFBFCFEF9F7FBFAFAF3F9FFFFFFFFFFFFFFFFFFFEFBFAFDFEFEFCFCFEFDF9FBFFFDFCFFFEFBFFFFFFFFFFFFFFFFFCF9F7F4FFFEFDF6FDFFFFFFFFFFFFFFFFFFFFFFFFFFFFFFFFFFFBFAFDFBFBFAF8FCFEFFFFFFFFFFFFFFFFFFFFFFFFFFFFFFFFFFFCFCFEFAF6FAFEFCF9FEFFFEFDFDFDFEFFFCF8FBF6FCFBFEF9FCF9FBFCF8F9FFFCF2FCF8FBF8FCFFFBF6F9FBFAFDFCFCFBFAFEF7F9F9FAFDFAF8FDF9FAFDFDFCF8FAFAFCFDFBFBFBF6F9FAFAFAFDFEFAF7F9FCFBF7F9FDFBFEFEFFFDFBFCFBFCFFFFFFFFFFFFFFFFFFFFFFFFFFFFFFFFFFFF>
EI Q}\hss}}
\pdfxform\publisherbox\edef\publisherform{\the\pdflastxform}\newcommand{\publisherdevice}{\pdfrefxform\publisherform}
\newcommand{\weberfigure}[2]{\begin{center}\vspace{23mm}\leavevmode\hbox{\publisherdevice}\par\vspace{20mm}\end{center}}
\expandafter\def\csname ednote1\endcsname{\edNote{1}{s001R}{Dubbio. piccolo segno dopo «ZANICHELLI» -- macchia / punto / virgola.}}
\expandafter\def\csname ednote2\endcsname{\edNote{2}{IX}{Stampa. teoria-stessa -> teoria stessa, se il trattino è tipografico, unisce due parole intere.}}
\expandafter\def\csname ednote3\endcsname{\edNote{3}{IX}{Dubbio. piccolo segno fra «delle» e «lezioni» in «i corsi delle lezioni» -- macchia / punto stampato.}}
\expandafter\def\csname ednote6\endcsname{\edNote{4}{XI}{Dubbio. punto dopo «56» -- punto stampato / macchia.}}
\expandafter\def\csname ednote7\endcsname{\edNote{5}{XIV}{Stampa. Disgressione -> Digressione, probabile s soprannumeraria; a p. XIII è stampato «Digressione».}}
\expandafter\def\csname ednote8\endcsname{\edNote{6}{3}{Stampa. F -> f, la superficie di partenza è indicata con f nelle righe precedenti.}}
\expandafter\def\csname ednote10\endcsname{\edNote{7}{11}{Stampa. C -> \textbackslash{}\allowbreak bar\{C\}, la curva doppia di \textbackslash{}\allowbreak Phi è indicata con \textbackslash{}\allowbreak bar\{C\} nelle righe contigue.}}
\expandafter\def\csname ednote11\endcsname{\edNote{8}{12}{Stampa. C -> \textbackslash{}\allowbreak bar\{C\}, qui si parla della curva doppia di \textbackslash{}\allowbreak Phi, detta \textbackslash{}\allowbreak bar\{C\} nelle due alternative precedenti.}}
\expandafter\def\csname ednote12\endcsname{\edNote{9}{12}{Stampa. \textbackslash{}\allowbreak bar\{K\}\_\allowbreak 1 in «la curva \textbackslash{}\allowbreak bar\{K\}\_\allowbreak 1 proiettata» e in «di \textbackslash{}\allowbreak bar\{K\}\_\allowbreak 1» -> K\_\allowbreak 1, la curva su F è denotata K\_\allowbreak 1 nel periodo precedente.}}
\expandafter\def\csname ednote15\endcsname{\edNote{10}{19}{Stampa. C -> |C|, the first symbol denotes the system, otherwise written with vertical bars.}}
\expandafter\def\csname ednote16\endcsname{\edNote{11}{22}{Stampa. mn -> mn\textasciicircum{}2, the square on the cone order appears to be missing.}}
\expandafter\def\csname ednote18\endcsname{\edNote{12}{24}{Stampa. Mathematik -> Mathematics, apparent misprint in the English journal title.}}
\expandafter\def\csname ednote19\endcsname{\edNote{13}{27}{Stampa. contenuti -> contenute, agreement with "esse" (the curves).}}
\expandafter\def\csname ednote21\endcsname{\edNote{14}{39}{Stampa. \textbackslash{}\allowbreak omega\_\allowbreak 1 -> \textbackslash{}\allowbreak omega\_\allowbreak 2, the immediately preceding clause names the transformed curve \textbackslash{}\allowbreak omega\_\allowbreak 2.}}
\expandafter\def\csname ednote23\endcsname{\edNote{15}{41}{Dubbio. small mark after n in 3n-3 -- scanning speck or prime (3n'-3).}}
\expandafter\def\csname ednote27\endcsname{\edNote{16}{44}{Stampa. (\textbackslash{}\allowbreak partial\textbackslash{}\allowbreak varphi\textbackslash{}\allowbreak varphi\_\allowbreak 3) -> \textbackslash{}\allowbreak partial(\textbackslash{}\allowbreak varphi\textbackslash{}\allowbreak varphi\_\allowbreak 3), the opening parenthesis is printed before the derivative symbol.}}
\expandafter\def\csname ednote28\endcsname{\edNote{17}{45}{Stampa. -\textbackslash{}\allowbreak partial f/\textbackslash{}\allowbreak partial x\_\allowbreak 3 -> \textbackslash{}\allowbreak partial f/\textbackslash{}\allowbreak partial x\_\allowbreak 3, if the stroke is type, the preceding derivative column has no minus.}}
\expandafter\def\csname ednote29\endcsname{\edNote{18}{49}{Stampa. punt -> punti, the final i is absent before "base".}}
\expandafter\def\csname ednote30\endcsname{\edNote{19}{50}{Stampa. dalle -> delle, the phrase identifies the system of plane sections.}}
\expandafter\def\csname ednote33\endcsname{\edNote{20}{58}{Stampa. la molteplicità -> le molteplicità, the verb "si lasciano" is plural.}}
\expandafter\def\csname ednote35\endcsname{\edNote{21}{61}{Dubbio. relation sign -- \textbackslash{}\allowbreak doteq / =; the small mark above the bars may be a scan speck.}}
\expandafter\def\csname ednote36\endcsname{\edNote{22}{61}{Dubbio. relation sign -- \textbackslash{}\allowbreak doteq / =; the small mark above the bars may be a scan speck.}}
\expandafter\def\csname ednote40\endcsname{\edNote{23}{77}{Stampa. C\_\allowbreak 4 -> C\_\allowbreak 7, these residual curves were denoted C\_\allowbreak 7 earlier on the page.}}
\expandafter\def\csname ednote41\endcsname{\edNote{24}{77}{Stampa. C\_\allowbreak 4 -> C\_\allowbreak 7, repeated subscript discrepancy for the same curves.}}
\expandafter\def\csname ednote42\endcsname{\edNote{25}{79}{Stampa. O -> A, the point under discussion is called A in this sentence.}}
\expandafter\def\csname ednote43\endcsname{\edNote{26}{79}{Stampa. i volte A per -> i volte per A, the words are transposed in print.}}
\expandafter\def\csname ednote45\endcsname{\edNote{27}{84}{Stampa. dal 4º ordine -> del 4º ordine, the preposition appears to be a typesetting error.}}
\expandafter\def\csname ednote46\endcsname{\edNote{28}{84}{Stampa. di questa -> di questo, the pronoun appears to refer to the point O.}}
\expandafter\def\csname ednote47\endcsname{\edNote{29}{85}{Stampa. è in una cubica -> e in una cubica, a conjunction is required here.}}
\expandafter\def\csname ednote48\endcsname{\edNote{30}{87}{Stampa. grado -\textbackslash{}\allowbreak nu -> grado \textbackslash{}\allowbreak nu, the signed degree is denoted by \textbackslash{}\allowbreak nu on the preceding page and in the concluding formula.}}
\expandafter\def\csname ednote50\endcsname{\edNote{31}{87}{Stampa. -\textbackslash{}\allowbreak nu=-\textbackslash{}\allowbreak nu\_\allowbreak 0-\textbackslash{}\allowbreak nu\_\allowbreak 1+2=-1 -> \textbackslash{}\allowbreak nu=-\textbackslash{}\allowbreak nu\_\allowbreak 0-\textbackslash{}\allowbreak nu\_\allowbreak 1+2=-1, the initial sign conflicts with -\textbackslash{}\allowbreak nu=1 above and \textbackslash{}\allowbreak nu=-1 below.}}
\expandafter\def\csname ednote53\endcsname{\edNote{32}{91}{Stampa. \textbackslash{}\allowbreak varphi\_\allowbreak \{n-4\} -> \textbackslash{}\allowbreak varphi\_\allowbreak \{n-3\}, the section of degree n-4 together with a line has degree n-3.}}
\expandafter\def\csname ednote54\endcsname{\edNote{33}{92}{Stampa. del punto -> dal punto, the projection is taken from O.}}
\expandafter\def\csname ednote55\endcsname{\edNote{34}{94}{Dubbio. punctuation after “Castelnuovo” -- no punctuation or a faint period.}}
\expandafter\def\csname ednote56\endcsname{\edNote{35}{95}{Stampa. supeficie -> superficie, the letter r is missing.}}
\expandafter\def\csname ednote57\endcsname{\edNote{36}{98}{Stampa. di essa -> di esso, the stated antecedent is the masculine “fascio”.}}
\expandafter\def\csname ednote58\endcsname{\edNote{37}{99}{Stampa. e suoi poli -> e i suoi poli, the article is absent in print.}}
\expandafter\def\csname ednote59\endcsname{\edNote{38}{101}{Stampa. contenerle -> contenerli, the antecedent is “i punti doppi”.}}
\expandafter\def\csname ednote60\endcsname{\edNote{39}{102}{Stampa. contenerle -> contenerli, the antecedent is the masculine plural “punti doppi”.}}
\expandafter\def\csname ednote61\endcsname{\edNote{40}{102}{Stampa. contenerle -> contenerla, the antecedent is the singular curve C.}}
\expandafter\def\csname ednote62\endcsname{\edNote{41}{103}{Stampa. d’ordine n -> d’ordine m, the order of the curve C was denoted by m.}}
\expandafter\def\csname ednote63\endcsname{\edNote{42}{103}{Stampa. alle totalità -> alla totalità, the passage apparently refers to the single totality of these surfaces.}}
\expandafter\def\csname ednote64\endcsname{\edNote{43}{109}{Stampa. sezioni piane di s -> sezioni piane di F, the plane sections under discussion are those of the surface F.}}
\expandafter\def\csname ednote65\endcsname{\edNote{44}{109}{Stampa. -\textbackslash{}\allowbreak frac\{s(s-1)\}\{2\}n -> +\textbackslash{}\allowbreak frac\{s(s-1)\}\{2\}n, the sign disagrees with the preceding expression and the stated equality.}}
\expandafter\def\csname ednote66\endcsname{\edNote{45}{111}{Stampa. e dalla serie -> e della serie, the text is describing a sum of two series.}}
\expandafter\def\csname ednote68\endcsname{\edNote{46}{115}{Stampa. della superficie -> delle superficie, the following participle "costrette" is plural.}}
\expandafter\def\csname ednote70\endcsname{\edNote{47}{116}{Stampa. p\_\allowbreak a+\textbackslash{}\allowbreak pi\_\allowbreak \{h-3\}+\textbackslash{}\allowbreak delta\_\allowbreak \{h-2\}+\textbackslash{}\allowbreak delta\_\allowbreak \{h-1\} -> p\_\allowbreak a+\textbackslash{}\allowbreak pi\_\allowbreak \{h-3\}-1+\textbackslash{}\allowbreak delta\_\allowbreak \{h-2\}+\textbackslash{}\allowbreak delta\_\allowbreak \{h-1\}, the constant -1 in the preceding dimension formula is missing here.}}
\expandafter\def\csname ednote71\endcsname{\edNote{48}{118}{Stampa. f -> F, the surface is otherwise denoted by capital F.}}
\expandafter\def\csname ednote72\endcsname{\edNote{49}{119}{Stampa. se zione -> sezione, the word is split at a printed line end without a hyphen.}}
\expandafter\def\csname ednote73\endcsname{\edNote{50}{119}{Stampa. f -> F, the surface is otherwise denoted by capital F.}}
\expandafter\def\csname ednote76\endcsname{\edNote{51}{124}{Stampa. \textbackslash{}\allowbreak bar\{C\}=|C-\textbackslash{}\allowbreak theta| -> |\textbackslash{}\allowbreak bar\{C\}|=|C-\textbackslash{}\allowbreak theta|, the left-hand system delimiters are absent in the print.}}
\expandafter\def\csname ednote77\endcsname{\edNote{52}{124}{Dubbio. small marks after nu -- scan speckles with no punctuation (read), comma or full stop.}}
\expandafter\def\csname ednote79\endcsname{\edNote{53}{127}{Stampa. addittivo -> additivo, the printed word contains a doubled t.}}
\expandafter\def\csname ednote80\endcsname{\edNote{54}{130}{Stampa. delle serie completa -> della serie completa, the printed article does not agree with the singular noun and adjective.}}
\expandafter\def\csname ednote82\endcsname{\edNote{55}{133}{Stampa. RiemannRoch -> Riemann-Roch, the internal hyphen is absent in this occurrence.}}
\expandafter\def\csname ednote84\endcsname{\edNote{56}{141}{Stampa. opening « only -> closing » apparently omitted, the quotation is left open in print.}}
\expandafter\def\csname ednote86\endcsname{\edNote{57}{144}{Dubbio. the faint mark after the final 1 -- scan speck (not transcribed) or full stop.}}
\expandafter\def\csname ednote87\endcsname{\edNote{58}{145}{Stampa. non sia abbia -> non si abbia, an extra a is printed in sia.}}
\expandafter\def\csname ednote88\endcsname{\edNote{59}{146}{Stampa. e quindi l’indice di specialità della serie -> e quindi la serie, the printed subject repeats the index that the predicate assigns.}}
\expandafter\def\csname ednote89\endcsname{\edNote{60}{148}{Dubbio. "A tal" -- "A tal" or "A. tal"; a small intervening mark may be a scan speck.}}
\expandafter\def\csname ednote90\endcsname{\edNote{61}{149}{Stampa. r -> \textbackslash{}\allowbreak nu, the surrounding text denotes the degree by \textbackslash{}\allowbreak nu.}}
\expandafter\def\csname ednote92\endcsname{\edNote{62}{152}{Dubbio. "che occorre" -- "che occorre" or "che. occorre"; the small intervening point may be a scan speck.}}
\expandafter\def\csname ednote93\endcsname{\edNote{63}{154}{Stampa. "sia A\_\allowbreak 1, un numero" -> "sia A\_\allowbreak 1, ha un numero", the clause appears to omit its verb.}}
\expandafter\def\csname ednote94\endcsname{\edNote{64}{155}{Stampa. (CL') -> (DL'), the preceding paragraph places this residual series on D.}}
\expandafter\def\csname ednote95\endcsname{\edNote{65}{157}{Stampa. "rappresentato del fascio" -> "rappresentato dal fascio", the preposition appears to be a typographical error.}}
\expandafter\def\csname ednote96\endcsname{\edNote{66}{158}{Stampa. Le possibilità -> La possibilità, the following verb "costringe" is singular.}}
\expandafter\def\csname ednote97\endcsname{\edNote{67}{159}{Stampa. i sistema -> il sistema, the singular article is missing its final l.}}
\expandafter\def\csname ednote99\endcsname{\edNote{68}{160}{Stampa. n<\textbackslash{}\allowbreak pi-1 -> n\textbackslash{}\allowbreak leq\textbackslash{}\allowbreak pi-1, the preceding inequality gives only a non-strict bound when p=1.}}
\expandafter\def\csname ednote100\endcsname{\edNote{69}{161}{Dubbio. apparent dot after "dividere" -- scan speck (not transcribed), printed full stop.}}
\expandafter\def\csname ednote101\endcsname{\edNote{70}{164}{Dubbio. dots after "connesse" and "connessione" -- printed full stops, scan speckles.}}
\expandafter\def\csname ednote102\endcsname{\edNote{71}{164}{Stampa. \textbackslash{}\allowbreak delta\_\allowbreak 1\textbackslash{}\allowbreak delta\_\allowbreak 2 -> \textbackslash{}\allowbreak delta\_\allowbreak 1,\textbackslash{}\allowbreak delta\_\allowbreak 2, the three previously distinguished degrees are being enumerated.}}
\expandafter\def\csname ednote103\endcsname{\edNote{72}{165}{Dubbio. subscript on the second F in "F\_\allowbreak 1+F\_\allowbreak 2" -- 2, a dash-like damaged character.}}
\expandafter\def\csname ednote104\endcsname{\edNote{73}{165}{Stampa. m\_\allowbreak 1+m\_\allowbreak 2+ +4 -> m\_\allowbreak 1+m\_\allowbreak 2-4, the same adjoint order is printed with -4 elsewhere on this page.}}
\expandafter\def\csname ednote105\endcsname{\edNote{74}{168}{Dubbio. punctuation after "avremo" -- colon or an isolated scan speck.}}
\expandafter\def\csname ednote107\endcsname{\edNote{75}{169}{Stampa. "(r-1)\textasciicircum{}2, volte" -> "(r-1)\textasciicircum{}2 volte", the comma separates the multiplier from "volte".}}
\expandafter\def\csname ednote108\endcsname{\edNote{76}{170}{Stampa. "serie segata" -> "serie segate", agreement with the plural "delle serie".}}
\expandafter\def\csname ednote109\endcsname{\edNote{77}{170}{Stampa. "(C\_\allowbreak 2L)" -> "(C\_\allowbreak 2L\_\allowbreak 1)", the preceding count concerns the residual curve L\_\allowbreak 1.}}
\expandafter\def\csname ednote110\endcsname{\edNote{78}{172}{Stampa. "dato" -> "data", agreement with the subject "L’equivalenza".}}
\expandafter\def\csname ednote111\endcsname{\edNote{79}{173}{Stampa. "(\textbackslash{}\allowbreak pi-2)" -> "(2\textbackslash{}\allowbreak pi-2)", the preceding sentence gives the multiplicity as 2\textbackslash{}\allowbreak pi-2.}}
\expandafter\def\csname ednote112\endcsname{\edNote{80}{173}{Stampa. "P. Bonnesen" -> "T. Bonnesen", the main text gives the initial T.; "Académie Royal" -> "Académie Royale", adjective agreement in French.}}
\expandafter\def\csname ednote113\endcsname{\edNote{81}{175}{Stampa. "della formula (2)" -> "dalla formula (3)", the preposition is anomalous and the preceding formula for d is labelled (3).}}
\expandafter\def\csname ednote114\endcsname{\edNote{82}{175}{Stampa. "2n\textasciicircum{}2-7n+15+2\textbackslash{}\allowbreak pi(n-7)" -> "2n\textasciicircum{}2-7n+15+2\textbackslash{}\allowbreak pi(n-7)-p\textasciicircum{}\{(1)\}", the subsequent equation (6) retains the -p\textasciicircum{}\{(1)\} term.}}
\expandafter\def\csname ednote115\endcsname{\edNote{83}{175}{Stampa. "n+14" -> apparently "-7n+15-\textbackslash{}\allowbreak tau", comparison with the preceding expressions for (C\_\allowbreak jD) and a gives a different linear and constant part.}}
\expandafter\def\csname ednote116\endcsname{\edNote{84}{176}{Stampa. "\textbackslash{}\allowbreak delta,I,\textbackslash{}\allowbreak varrho,\textbackslash{}\allowbreak Pi,\textbackslash{}\allowbreak tau,t" -> "\textbackslash{}\allowbreak delta,I,d,\textbackslash{}\allowbreak varrho,\textbackslash{}\allowbreak Pi,\textbackslash{}\allowbreak tau,t", only six entries are printed after "i sette caratteri".}}
\expandafter\def\csname ednote117\endcsname{\edNote{85}{176}{Stampa. "(2)" -> "(3)", the expression for d on page 174 is labelled (3).}}
\expandafter\def\csname ednote119\endcsname{\edNote{86}{182}{Stampa. 3(n+6\textbackslash{}\allowbreak pi-k) -> 3(n+6\textbackslash{}\allowbreak pi)-k, comparison with the fourth formula A).}}
\expandafter\def\csname ednote120\endcsname{\edNote{87}{182}{Stampa. G. H. Zeuthen -> H. G. Zeuthen, the initials appear in the opposite order on page 180.}}
\expandafter\def\csname ednote122\endcsname{\edNote{88}{185}{Dubbio. the final vowel of delle before biaggiunte -- delle; dello.}}
\expandafter\def\csname ednote123\endcsname{\edNote{89}{185}{Stampa. D (twice) -> d, the branch curve is denoted d later in this paragraph, while D denotes the coincidence curve.}}
\expandafter\def\csname ednote124\endcsname{\edNote{90}{186}{Dubbio. the page number -- 184 is visible in the overlapped area; the underlying page number is not securely readable.}}
\expandafter\def\csname ednote125\endcsname{\edNote{91}{186}{Dubbio. covered beginnings, continuations after cor- and dirama-, and the first footnote mark -- no reliable reading of the covered portions is possible; only visible fragments are retained, and the note is attached here provisionally.}}
\expandafter\def\csname ednote126\endcsname{\edNote{92}{186}{Dubbio. the clipped suffix before si può -- mente; uente; the preceding text and the beginning of the line containing numerici are covered.}}
\expandafter\def\csname ednote127\endcsname{\edNote{93}{186}{Dubbio. clipped word beginnings -- po (possibly uopo); urve (possibly curve); esti (possibly questi); other covered text cannot be read.}}
\expandafter\def\csname ednote128\endcsname{\edNote{94}{187}{Stampa. -3 -> -3\textbackslash{}\allowbreak nu, substitution of the preceding expression for s leaves a factor \textbackslash{}\allowbreak nu.}}
\expandafter\def\csname ednote129\endcsname{\edNote{95}{187}{Stampa. i valore -> il valore, the final l of the article is absent.}}
\expandafter\def\csname ednote130\endcsname{\edNote{96}{191}{Stampa. \textbackslash{}\allowbreak S II, 8 -> Cap. II, \textbackslash{}\allowbreak S 8, the reference at the top of this page has that form.}}
\expandafter\def\csname ednote131\endcsname{\edNote{97}{191}{Stampa. 2 : 24 -> 2 \textbackslash{}\allowbreak cdot 24, the displayed result is 48.}}
\expandafter\def\csname ednote133\endcsname{\edNote{98}{194}{Stampa. m/2 -> m, the preceding text specifies 2m simple branch points.}}
\expandafter\def\csname ednote134\endcsname{\edNote{99}{199}{Stampa. cu sui -> su cui, apparent transposition of the initial consonants.}}
\expandafter\def\csname ednote135\endcsname{\edNote{100}{204}{Stampa. 1932.; -> 1932;, a period is printed immediately before the semicolon.}}
\expandafter\def\csname ednote137\endcsname{\edNote{101}{206}{Stampa. Toveremo -> Troveremo, the r after T is missing in print.}}
\expandafter\def\csname ednote138\endcsname{\edNote{102}{208}{Stampa. m = n - 2d - 3k -> m = n\textasciicircum{}2 - 2d - 3k, apparent omitted square exponent.}}
\expandafter\def\csname ednote139\endcsname{\edNote{103}{212}{Stampa. "sulla su" -> "sulla"; "dalle curva" -> "dalle curve", duplicated preposition and number disagreement.}}
\expandafter\def\csname ednote142\endcsname{\edNote{104}{212}{Stampa. con curva una L -> con una curva L, transposed words.}}
\expandafter\def\csname ednote143\endcsname{\edNote{105}{214}{Stampa. \textbackslash{}\allowbreak omega = 0 -> \textbackslash{}\allowbreak omega' = 0, the preceding formula uses \textbackslash{}\allowbreak omega'.}}
\expandafter\def\csname ednote147\endcsname{\edNote{106}{220}{Stampa. espressa, da -> espressa da, intrusive comma between the participle and its complement.}}
\expandafter\def\csname ednote148\endcsname{\edNote{107}{222}{Stampa. curve canonica -> curva canonica, the printed noun and adjective disagree in number.}}
\expandafter\def\csname ednote149\endcsname{\edNote{108}{226}{Stampa. un’altro -> un altro, apostrophe before a masculine form.}}
\expandafter\def\csname ednote151\endcsname{\edNote{109}{230}{Stampa. definite -> definita, apparently modifying the singular "involuzione".}}
\expandafter\def\csname ednote152\endcsname{\edNote{110}{230}{Stampa. D\_\allowbreak 3 -> D\_\allowbreak 9, the tricanonical curve is identified above with the branch curve D\_\allowbreak 9.}}
\expandafter\def\csname ednote153\endcsname{\edNote{111}{235}{Stampa. Ipotesi I) -> Ipotesi I, the closing parenthesis has no matching opening parenthesis.}}
\expandafter\def\csname ednote154\endcsname{\edNote{112}{237}{Dubbio. The short mark after "tacnodali" may be a hyphen or an artefact -- tacnodali-p; tacnodali p.}}
\expandafter\def\csname ednote155\endcsname{\edNote{113}{240}{Stampa. "riduzione." -> "riduzione,", the following lower-case subordinate clause continues the sentence.}}
\expandafter\def\csname ednote156\endcsname{\edNote{114}{242}{Dubbio. short mark between "un" and "punto" -- printed hyphen or ink/scanning artefact.}}
\expandafter\def\csname ednote159\endcsname{\edNote{115}{248}{Stampa. sisema -> sistema, the printed word omits "t".}}
\expandafter\def\csname ednote160\endcsname{\edNote{116}{251}{Stampa. siano -> siamo, the printed verb is inconsistent with the first-person plural construction "ci ... valsi".}}
\expandafter\def\csname ednote162\endcsname{\edNote{117}{256}{Stampa. auf der Flaechen -> auf den Flaechen, apparent German agreement error.}}
\expandafter\def\csname ednote163\endcsname{\edNote{118}{257}{Stampa. algebraics Varieties. -> algebraic Varieties »., apparent extra s and missing closing guillemet.}}
\expandafter\def\csname ednote165\endcsname{\edNote{119}{259}{Dubbio. character after d in d’ramazione -- apostrophe or a defective i.}}
\expandafter\def\csname ednote166\endcsname{\edNote{120}{260}{Stampa. OP and P -> OB and B, the point is called B in the preceding description of type II.}}
\expandafter\def\csname ednote167\endcsname{\edNote{121}{262}{Stampa. +s-1 -> +s+1, the same relation is printed with +1 on page 263.}}
\expandafter\def\csname ednote168\endcsname{\edNote{122}{263}{Stampa. P=2(p-1) -> P=2(p-1)+1, an apparent omitted +1 when compared with the plurigenus formulas.}}
\expandafter\def\csname ednote169\endcsname{\edNote{123}{263}{Stampa. p-1=n-3 -> p-1=n-3-s, inconsistent with p=n-2-s printed above.}}
\expandafter\def\csname ednote171\endcsname{\edNote{124}{263}{Stampa. 2s-3 -> 2(s-3), the contribution is written as 2(s-3) in the following total.}}
\expandafter\def\csname ednote172\endcsname{\edNote{125}{265}{Dubbio. mark before superficie in the second quoted title -- opening square bracket, opening parenthesis, or scan blemish.}}
\expandafter\def\csname ednote173\endcsname{\edNote{126}{267}{Dubbio. the low mark between "e" and "mercè" -- comma (transcribed) or scan speck.}}
\expandafter\def\csname ednote174\endcsname{\edNote{127}{274}{Stampa. C\_\allowbreak 7(O\textasciicircum{}5O\_\allowbreak 1\textasciicircum{}2) -> C'\_\allowbreak 7(O\textasciicircum{}5O\_\allowbreak 1\textasciicircum{}2), the preceding plane-image notation uses a prime.}}
\expandafter\def\csname ednote176\endcsname{\edNote{128}{275}{Stampa. fuori di esso -> fuori di essa, apparent agreement error with "curva doppia".}}
\expandafter\def\csname ednote177\endcsname{\edNote{129}{277}{Stampa. dal -> del, the preposition appears erroneous after "fa parte".}}
\expandafter\def\csname ednote178\endcsname{\edNote{130}{282}{Dubbio. An oblique mark follows B in "17 punti A e B" -- a scan speck or a prime on B.}}
\expandafter\def\csname ednote180\endcsname{\edNote{131}{286}{Stampa. varietà cubiche S\_\allowbreak 4 -> varietà cubiche di S\_\allowbreak 4, a preposition before the ambient-space symbol appears to be omitted.}}
\expandafter\def\csname ednote181\endcsname{\edNote{132}{287}{Stampa. delle -> dalle, the surfaces appear to be the agents of "segato".}}
\expandafter\def\csname ednote182\endcsname{\edNote{133}{290}{Stampa. F after "contenuto in" -> |F|, apparently missing system bars in a sentence about subsystems.}}
\expandafter\def\csname ednote185\endcsname{\edNote{134}{293}{Stampa. F -> |F|, apparently missing system bars, unlike the preceding analogous display.}}
\expandafter\def\csname ednote186\endcsname{\edNote{135}{293}{Stampa. n=p-1 -> n=p-2, inconsistent with the exponent p-2 printed at the end of the same item.}}
\expandafter\def\csname ednote188\endcsname{\edNote{136}{295}{Stampa. \textbackslash{}\allowbreak leq -> \textbackslash{}\allowbreak geq, the preceding sentence states n \textbackslash{}\allowbreak geq 2r.}}
\expandafter\def\csname ednote189\endcsname{\edNote{137}{295}{Stampa. 2p+1 -> 2p-1, the preceding displayed bound for the composed case has a minus sign.}}
\expandafter\def\csname ednote190\endcsname{\edNote{138}{295}{Stampa. s -> r, the base-point count is denoted r immediately above.}}
\expandafter\def\csname ednote191\endcsname{\edNote{139}{296}{Stampa. S\_\allowbreak \{p-2\} -> S\_\allowbreak \{p-1\}, the canonical space is denoted S\_\allowbreak \{p-1\} in the preceding discussion.}}
\expandafter\def\csname ednote192\endcsname{\edNote{140}{296}{Stampa. p=7 -> p=6, inconsistent with p\textasciicircum{}\{(1)\}=2p-3 and the printed value p\textasciicircum{}\{(1)\}=9.}}
\expandafter\def\csname ednote193\endcsname{\edNote{141}{296}{Stampa. p\textasciicircum{}\{(1)\}=4n-2 -> p\textasciicircum{}\{(1)\}=4n-1, inconsistent with the relation p\textasciicircum{}\{(1)\}=2p-3 stated on this page.}}
\expandafter\def\csname ednote194\endcsname{\edNote{142}{296}{Stampa. 4n-8 -> 4n-7 and 4n-4 -> 4n-3, both inconsistent with the relation p\textasciicircum{}\{(1)\}=2p-3 stated on this page.}}
\expandafter\def\csname ednote195\endcsname{\edNote{143}{297}{Stampa. la serie -> le serie, the following adjective and verb are plural.}}
\expandafter\def\csname ednote196\endcsname{\edNote{144}{299}{Stampa. S\_\allowbreak \{p-7\} -> S\_\allowbreak \{p-1\}, the preceding page places this normal surface in S\_\allowbreak \{p-1\}.}}
\expandafter\def\csname ednote197\endcsname{\edNote{145}{303}{Stampa. 3p\_\allowbreak 8 -> 3p\_\allowbreak 8z, the factor z present in the corresponding terms of the two preceding equations is absent.}}
\expandafter\def\csname ednote198\endcsname{\edNote{146}{305}{Stampa. p=1 -> p=2, this refers to the p=2 case just described.}}
\expandafter\def\csname ednote199\endcsname{\edNote{147}{305}{Dubbio. the conjunction after p\textasciicircum{}\{(1)\}=2 -- è / e; the apparent accent may be a stray mark.}}
\expandafter\def\csname ednote200\endcsname{\edNote{148}{306}{Stampa. F'\_\allowbreak 2 -> F'\_\allowbreak 3, the same composite surface is written F\_\allowbreak 3+F'\_\allowbreak 3 immediately above.}}
\expandafter\def\csname ednote201\endcsname{\edNote{149}{307}{Stampa. addittivo -> additivo, apparent doubled t.}}
\expandafter\def\csname ednote202\endcsname{\edNote{150}{309}{Stampa. (K\_\allowbreak 2 - |\textbackslash{}\allowbreak theta|)' -> |K\_\allowbreak 2 - \textbackslash{}\allowbreak theta|', the printed bars enclose only theta rather than the residual system named in the preceding clause.}}
\expandafter\def\csname ednote203\endcsname{\edNote{151}{310}{Stampa. C -> C', the preceding sentence describes the adjoint curves C'.}}
\expandafter\def\csname ednote204\endcsname{\edNote{152}{311}{Stampa. bicanoniche -> tricanoniche, the preceding discussion concerns the tricanonical cubic model.}}
\expandafter\def\csname ednote206\endcsname{\edNote{153}{318}{Stampa. la la -> la, the article is repeated across two printed lines.}}
\expandafter\def\csname ednote207\endcsname{\edNote{154}{318}{Stampa. genere 2 -> genere 1, this quartic is described as genus 1 above and on the following page.}}
\expandafter\def\csname ednote208\endcsname{\edNote{155}{319}{Stampa. \textbackslash{}\allowbreak lvert C\_\allowbreak \{14\}\textbackslash{}\allowbreak rvert-sega -> \textbackslash{}\allowbreak lvert C\_\allowbreak \{14\}\textbackslash{}\allowbreak rvert sega, an extraneous hyphen separates the subject from the verb.}}
\expandafter\def\csname ednote210\endcsname{\edNote{156}{321}{Stampa. orta -> porta, the initial p is absent from the visible word.}}
\expandafter\def\csname ednote211\endcsname{\edNote{157}{324}{Dubbio. the small mark following “cfr.” -- a scan speck or a second printed stop.}}
\expandafter\def\csname ednote212\endcsname{\edNote{158}{330}{Stampa. q >= -p -> q >= p, the immediately preceding statement asserts a family of dimension p with p\_\allowbreak a=-p.}}
\expandafter\def\csname ednote213\endcsname{\edNote{159}{332}{Stampa. fade -> falde, the surrounding text refers to the same hypersurface sheets.}}
\expandafter\def\csname ednote214\endcsname{\edNote{160}{333}{Stampa. e dir vero -> a dir vero, apparent misprint in the idiomatic phrase.}}
\expandafter\def\csname ednote215\endcsname{\edNote{161}{333}{Stampa. un questione -> una questione, missing feminine ending before questione.}}
\expandafter\def\csname ednote216\endcsname{\edNote{162}{334}{Stampa. (riducibile)con -> (riducibile) con, no separating space is printed.}}
\expandafter\def\csname ednote217\endcsname{\edNote{163}{335}{Stampa. r=p+n-\textbackslash{}\allowbreak pi+1 -> r=p\_\allowbreak a+n-\textbackslash{}\allowbreak pi+1, the parallel formula for r' uses p\_\allowbreak a.}}
\expandafter\def\csname ednote218\endcsname{\edNote{164}{340}{Dubbio. relation after the first p\_\allowbreak g -- printed equality crossed by a later stroke, or printed \textbackslash{}\allowbreak ne.}}
\expandafter\def\csname ednote220\endcsname{\edNote{165}{344}{Dubbio. damaged initial in "d’ordine" -- d or l.}}
\expandafter\def\csname ednote221\endcsname{\edNote{166}{344}{Stampa. Comptes, rendus -> Comptes rendus, apparently extraneous comma in the journal title.}}
\expandafter\def\csname ednote222\endcsname{\edNote{167}{346}{Dubbio. dot-like mark after "(gruppi di punti)" -- scan speck or printed full stop.}}
\expandafter\def\csname ednote223\endcsname{\edNote{168}{348}{Dubbio. the symbols in the final parenthesis -- Greek \textbackslash{}\allowbreak nu or Latin v; \textbackslash{}\allowbreak varrho or \textbackslash{}\allowbreak rho.}}
\expandafter\def\csname ednote224\endcsname{\edNote{169}{348}{Dubbio. the symbols in the final parenthesis -- Greek \textbackslash{}\allowbreak nu or Latin v; \textbackslash{}\allowbreak varrho or \textbackslash{}\allowbreak rho.}}
\expandafter\def\csname ednote225\endcsname{\edNote{170}{349}{Stampa. C\_\allowbreak 8 -> C\_\allowbreak 5, the cutting curves are called quintics C\_\allowbreak 5 on p. 348.}}
\expandafter\def\csname ednote226\endcsname{\edNote{171}{354}{Dubbio. the damaged n in “canonica” -- canonica or canoica.}}
\expandafter\def\csname ednote227\endcsname{\edNote{172}{357}{Stampa. indipendenti. rappresentazione -> indipendenti, rappresentazione, the printed full stop is followed by a lower-case continuation.}}
\expandafter\def\csname ednote228\endcsname{\edNote{173}{358}{Stampa. "XXIII,," -> "XXIII,", apparently duplicated comma after the volume number.}}
\expandafter\def\csname ednote229\endcsname{\edNote{174}{360}{Dubbio. mark after "A" in "A tal fine" -- scan mark or apostrophe.}}
\expandafter\def\csname ednote230\endcsname{\edNote{175}{361}{Dubbio. initial word -- e or è; the mark above it may be scan noise.}}
\expandafter\def\csname ednote231\endcsname{\edNote{176}{362}{Stampa. M'\_\allowbreak \{q-1\} -> K\_\allowbreak \{q-1\}, the image of the series was denoted K\_\allowbreak \{q-1\} earlier in this paragraph.}}
\expandafter\def\csname ednote232\endcsname{\edNote{177}{363}{Stampa. first M\_\allowbreak \{q-2\} -> M\_\allowbreak \{q-1\}, the preceding paragraph places this intersection on M\_\allowbreak \{q-1\}.}}
\expandafter\def\csname ednote233\endcsname{\edNote{178}{365}{Stampa. q\_\allowbreak g = 4 -> q = 4, this case is introduced as q = 4 in the preceding paragraph.}}
\expandafter\def\csname ednote234\endcsname{\edNote{179}{366}{Stampa. p\_\allowbreak a = 0 -> p\_\allowbreak g = 0, the converse concerns the geometric-genus condition stated immediately before.}}
\expandafter\def\csname ednote235\endcsname{\edNote{180}{366}{Dubbio. mark after "affermata" -- scan noise or comma.}}
\expandafter\def\csname ednote237\endcsname{\edNote{181}{368}{Dubbio. Only the ending "NO" is exposed; the full head is read from the matching page in scan s197L -- CAPITOLO NONO / another obscured head ending in NO.}}
\expandafter\def\csname ednote238\endcsname{\edNote{182}{368}{Dubbio. Several overlaid fragments are cropped -- g / p; qu / q; TI / UI; L / C; e F / F.}}
\expandafter\def\csname ednote239\endcsname{\edNote{183}{368}{Dubbio. The covered portions of this paragraph, including its opening, are restored solely from the matching exposed text in scan s197L -- identical wording / wording differing beneath the fold.}}
\expandafter\def\csname ednote240\endcsname{\edNote{184}{368}{Dubbio. The covered beginning and portions of the next line are read from the matching paragraph in scan s197L -- identical wording / wording differing beneath the fold.}}
\expandafter\def\csname ednote241\endcsname{\edNote{185}{368}{Stampa. In-text (1), footnote (3) -> in-text (3), the printed call and the number printed beside its footnote disagree.}}
\expandafter\def\csname ednote242\endcsname{\edNote{186}{372}{Dubbio. A diagonal blemish immediately follows "formeranno" -- "formeranno una" / "formeranno, una".}}
\expandafter\def\csname ednote243\endcsname{\edNote{187}{375}{Dubbio. The mark after "che" is read as a comma -- comma / scan blemish.}}
\expandafter\def\csname ednote244\endcsname{\edNote{188}{376}{Dubbio. small mark before "Cap." in the footnote -- scan speck or an additional full stop.}}
\expandafter\def\csname ednote245\endcsname{\edNote{189}{376}{Dubbio. small mark after "esplicito" -- scan speck or a printed full stop.}}
\expandafter\def\csname ednote246\endcsname{\edNote{190}{377}{Stampa. (t+1) -> (t-1), the preceding general formula has the factor t\_\allowbreak i-1.}}
\expandafter\def\csname ednote247\endcsname{\edNote{191}{378}{Stampa. -r+1 -> +r+1, the printed right-hand expression has the opposite sign for this r term.}}
\expandafter\def\csname ednote248\endcsname{\edNote{192}{378}{Stampa. è quindi -> e quindi, the accented verb interrupts the coordination of the two clauses introduced by "che".}}
\expandafter\def\csname ednote249\endcsname{\edNote{193}{382}{Stampa. II e 1 -> II e I, the hypotheses are labelled I and II.}}
\expandafter\def\csname ednote250\endcsname{\edNote{194}{383}{Stampa. delle serie -> della serie, the following adjective "s-bicanonica" is singular.}}
\expandafter\def\csname ednote251\endcsname{\edNote{195}{383}{Stampa. sK -> sK\textasciicircum{}\textbackslash{}\allowbreak pi, the following sentence compares the same system with |K\textasciicircum{}\textbackslash{}\allowbreak pi|.}}
\expandafter\def\csname ednote252\endcsname{\edNote{196}{383}{Stampa. para-s-canoniche -> para-s\textbackslash{}\allowbreak pi-canoniche, the preceding text calls C para-\textbackslash{}\allowbreak pi-canonica and then takes its s-fold multiple.}}
\expandafter\def\csname ednote253\endcsname{\edNote{197}{385}{Stampa. 341). -> 341))., the parenthesis opened before "Teorema" is not closed in print.}}
\expandafter\def\csname ednote254\endcsname{\edNote{198}{385}{Stampa. dato-con -> dato con, an extraneous hyphen joins these words.}}
\expandafter\def\csname ednote255\endcsname{\edNote{199}{386}{Stampa. 2(K+K\_\allowbreak 1+\textbackslash{}\allowbreak ldotp\textbackslash{}\allowbreak ldotp\textbackslash{}\allowbreak ldotp\textbackslash{}\allowbreak ldotp+K\_\allowbreak r) -> 2(K\_\allowbreak 1+\textbackslash{}\allowbreak ldotp\textbackslash{}\allowbreak ldotp\textbackslash{}\allowbreak ldotp\textbackslash{}\allowbreak ldotp+K\_\allowbreak r), the following virtual curve omits the unindexed K.}}
\expandafter\def\csname ednote256\endcsname{\edNote{200}{388}{Stampa. indipsndentemente -> indipendentemente, apparent substitution of s for e.}}
\expandafter\def\csname ednote257\endcsname{\edNote{201}{388}{Stampa. K\_\allowbreak \{\textbackslash{}\allowbreak pi-2\} -> K\_\allowbreak \{2\textbackslash{}\allowbreak pi-2\}, the curve has the latter subscript earlier in this paragraph.}}
\expandafter\def\csname ednote258\endcsname{\edNote{202}{389}{Stampa. K\_\allowbreak \{\textbackslash{}\allowbreak pi-2\} -> K\_\allowbreak \{2\textbackslash{}\allowbreak pi-2\}, the corresponding curve is given the latter subscript below.}}
\expandafter\def\csname ednote259\endcsname{\edNote{203}{389}{Stampa. g -> \textbackslash{}\allowbreak bar g, the same series carry an overbar later in the paragraph.}}
\expandafter\def\csname ednote260\endcsname{\edNote{204}{390}{Stampa. K\_\allowbreak \{\textbackslash{}\allowbreak pi-2\} in S\_\allowbreak \textbackslash{}\allowbreak pi -> K\_\allowbreak \{2\textbackslash{}\allowbreak pi-2\} in S\_\allowbreak \{\textbackslash{}\allowbreak pi-2\}, comparison with the same construction on the preceding pages.}}
\expandafter\def\csname ednote261\endcsname{\edNote{205}{390}{Stampa. sivuole (printed si- / vuole) -> si vuole, an anomalous line-end hyphen joins two separate words.}}
\expandafter\def\csname ednote262\endcsname{\edNote{206}{395}{Stampa. n=6 -> m=6, the order of the curves K\_\allowbreak m is denoted m on the preceding page.}}
\expandafter\def\csname ednote263\endcsname{\edNote{207}{396}{Dubbio. the first closed-path symbol -- \textbackslash{}\allowbreak nu, a faint \textbackslash{}\allowbreak gamma, or Latin v.}}
\expandafter\def\csname ednote264\endcsname{\edNote{208}{397}{Stampa. contratto -> contatto, an extra r is printed in the contact term.}}
\expandafter\def\csname ednote265\endcsname{\edNote{209}{400}{Stampa. 13-12p\_\allowbreak a-p\textasciicircum{}\{(1)\} -> 13+12p\_\allowbreak a-p\textasciicircum{}\{(1)\}, the following lower bound on p\_\allowbreak a requires a plus sign.}}
\expandafter\def\csname ednote266\endcsname{\edNote{210}{403}{Dubbio. punctuation after h\_\allowbreak i\textbackslash{}\allowbreak varrho\_\allowbreak i -- a full stop, a faint comma, or a scan speck.}}
\expandafter\def\csname ednote267\endcsname{\edNote{211}{405}{Stampa. (z-a)\textasciicircum{}\{r\_\allowbreak 1\} -> (z-a\_\allowbreak 1)\textasciicircum{}\{r\_\allowbreak 1\}, the preceding formula indexes the first branch value by 1.}}
\expandafter\def\csname ednote268\endcsname{\edNote{212}{405}{Stampa. 2n-2\textbackslash{}\allowbreak pi-2 -> 2n+2\textbackslash{}\allowbreak pi-2, compare the corresponding equality on printed p. 404.}}
\expandafter\def\csname ednote269\endcsname{\edNote{213}{408}{Dubbio. punctuation after "formula" -- a full stop or a scanning speck.}}
\expandafter\def\csname ednote270\endcsname{\edNote{214}{408}{Stampa. \textbackslash{}\allowbreak Sigma n(s-1) -> \textbackslash{}\allowbreak Sigma\textbackslash{}\allowbreak frac\{n\}\{s\}(s-1), the corresponding equality on page 410 contains the factor n/s.}}
\expandafter\def\csname ednote271\endcsname{\edNote{215}{408}{Stampa. \textbackslash{}\allowbreak geq1 -> \textbackslash{}\allowbreak geq2, the equality with 2n gives the stronger bound used below for t=3.}}
\expandafter\def\csname ednote272\endcsname{\edNote{216}{408}{Stampa. \textbackslash{}\allowbreak geq1 -> \textbackslash{}\allowbreak geq2, the apparently intended bound is the one used in the following t=3 calculation.}}
\expandafter\def\csname ednote273\endcsname{\edNote{217}{408}{Stampa. t-1 -> t-2, the next displayed calculation gives a bound of 1 when t=3.}}
\expandafter\def\csname ednote274\endcsname{\edNote{218}{409}{Stampa. s\_\allowbreak 1\textbackslash{}\allowbreak geq3 -> s\_\allowbreak 1=3, the following fourth case separately treats s\_\allowbreak 1\textbackslash{}\allowbreak geq4.}}
\expandafter\def\csname ednote275\endcsname{\edNote{219}{411}{Stampa. s=2 -> s\_\allowbreak 1=2, the other two multiplicities in this example are indexed 2 and 3.}}
\expandafter\def\csname ednote276\endcsname{\edNote{220}{411}{Stampa. \textbackslash{}\allowbreak Sigma\textbackslash{}\allowbreak frac\{s+11\}\{s\} -> \textbackslash{}\allowbreak Sigma\textbackslash{}\allowbreak left[\textbackslash{}\allowbreak frac\{s+11\}\{s\}\textbackslash{}\allowbreak right], the integer-part brackets used in the general formula and subsequent calculations are absent here.}}
\expandafter\def\csname ednote277\endcsname{\edNote{221}{411}{Stampa. 24 -> 6t, the preceding bound is 6 for each of t summands.}}
\expandafter\def\csname ednote278\endcsname{\edNote{222}{411}{Stampa. P\_\allowbreak \{12\}\textbackslash{}\allowbreak geq12t-47\textbackslash{}\allowbreak geq12(t-4)+1>1 -> P\_\allowbreak \{12\}\textbackslash{}\allowbreak geq6t-23=6(t-4)+1>1, using the preceding per-term bound with t summands.}}
\expandafter\def\csname ednote279\endcsname{\edNote{223}{412}{Stampa. s\_\allowbreak 1>4 -> s\_\allowbreak 1\textbackslash{}\allowbreak geq4, the displayed enumeration otherwise omits s\_\allowbreak 1=4.}}
\expandafter\def\csname ednote280\endcsname{\edNote{224}{415}{Stampa. pluricanonice -> pluricanoniche, the h is absent in print.}}
\expandafter\def\csname ednote281\endcsname{\edNote{225}{415}{Stampa. prima serie -> prima specie, the preceding discussion calls these transformations "di prima specie".}}
\expandafter\def\csname ednote282\endcsname{\edNote{226}{416}{Stampa. (x-b) -> (z-b), the other factors and the subsequent equations use z.}}
\expandafter\def\csname ednote284\endcsname{\edNote{227}{422}{Stampa. \textbackslash{}\allowbreak pi\textbackslash{}\allowbreak omega=\textbackslash{}\allowbreak pi -> \textbackslash{}\allowbreak tau\textbackslash{}\allowbreak omega=\textbackslash{}\allowbreak pi, the same left-hand side already equals \textbackslash{}\allowbreak omega\_\allowbreak 1 in this row and the reverse composition is listed above.}}
\expandafter\def\csname ednote285\endcsname{\edNote{228}{425}{Dubbio. punctuation after P\_\allowbreak \{12\}=1 -- comma, possibly followed by a period.}}
\expandafter\def\csname ednote288\endcsname{\edNote{229}{431}{Stampa. Schwartz -> Schwarz, the running text on this page spells the same name Schwarz.}}
\expandafter\def\csname ednote289\endcsname{\edNote{230}{431}{Stampa. f -> F, the surface is otherwise denoted by capital F in this paragraph.}}
\expandafter\def\csname ednote290\endcsname{\edNote{231}{432}{Stampa. birationelles -> birationnelles, apparent missing n in the French adjective.}}
\expandafter\def\csname ednote292\endcsname{\edNote{232}{434}{Stampa. è possiede -> e possiede, the conjunction is printed with an accent.}}
\expandafter\def\csname ednote293\endcsname{\edNote{233}{435}{Stampa. eccezio. i -> eccezioni, the printed word has a point and a gap where n is expected.}}
\expandafter\def\csname ednote294\endcsname{\edNote{234}{437}{Stampa. subgroup \textbackslash{}\allowbreak infty\textasciicircum{}p and group \textbackslash{}\allowbreak infty\textasciicircum{}q -> subgroup \textbackslash{}\allowbreak infty\textasciicircum{}q and group \textbackslash{}\allowbreak infty\textasciicircum{}p, the dimensions are reversed relative to the surrounding notation.}}
\expandafter\def\csname ednote295\endcsname{\edNote{235}{440}{Stampa. sopra F -> sopra \textbackslash{}\allowbreak Phi, the following sentence explicitly places the curves L on \textbackslash{}\allowbreak Phi.}}
\expandafter\def\csname ednote298\endcsname{\edNote{236}{443}{Stampa. aventi -> avente, apparent agreement error with the singular “un sistema bicanonico”.}}
\expandafter\def\csname ednote300\endcsname{\edNote{237}{444}{Stampa. dalla -> della, apparent preposition error after “conseguenza”.}}
\expandafter\def\csname ednote303\endcsname{\edNote{238}{446}{Stampa. Picard, ha -> Picard ha, apparently intrusive comma between subject and verb.}}
\expandafter\def\csname ednote305\endcsname{\edNote{239}{449}{Stampa. \textbackslash{}\allowbreak le0 -> \textbackslash{}\allowbreak ge0, apparently reversed inequality in view of the bound asserted on the next line.}}
\expandafter\def\csname ednote306\endcsname{\edNote{240}{449}{Stampa. p\_\allowbreak a=1 -> p\_\allowbreak a=-1, a minus sign apparently omitted in the limiting case.}}
\expandafter\def\csname ednote307\endcsname{\edNote{241}{451}{Stampa. Cap. XII -> Cap. VII, apparently inconsistent with the reference for the same family on the preceding page.}}
\expandafter\def\csname ednote308\endcsname{\edNote{242}{451}{Stampa. quale ordini -> quali ordini, adjective and plural noun do not agree.}}
\expandafter\def\csname ednote309\endcsname{\edNote{243}{453}{Stampa. d ordine -> d’ordine, apparently missing elision apostrophe.}}
\expandafter\def\csname ednote311\endcsname{\edNote{244}{455}{Stampa. "hyperellyptiques" -> "hyperelliptiques", apparent French spelling error.}}
\expandafter\def\csname ednote312\endcsname{\edNote{245}{456}{Stampa. "0)." -> "0.", unmatched closing parenthesis.}}
\expandafter\def\csname ednote313\endcsname{\edNote{246}{457}{Stampa. "distinti" -> "distinte", apparent agreement error referring to "superficie".}}
\expandafter\def\csname ednote314\endcsname{\edNote{247}{457}{Stampa. B') -> B''), the preceding distinction assigns the finite families to B''.}}
\expandafter\def\csname ednote315\endcsname{\edNote{248}{457}{Stampa. first B') and final B'') -> first B'') and final B'), these labels reverse the distinction stated above.}}
\expandafter\def\csname ednote316\endcsname{\edNote{249}{459}{Stampa. "pag. 463" -> "pag. 461", the Notizia reproduced in this packet begins on printed page 461.}}
\expandafter\def\csname ednote319\endcsname{\edNote{250}{463}{Dubbio. the small mark after "schiera" -- a comma (read here) or a scan speck.}}
\expandafter\def\csname ednote321\endcsname{\edNote{251}{463}{Stampa. final "P\_\allowbreak 4=0,1" -> "P\_\allowbreak 6=0,1", P\_\allowbreak 4 is repeated in the list of plurigenera.}}
\expandafter\def\csname ednote322\endcsname{\edNote{252}{464}{Dubbio. possible accent over the first o in "piuttosto" -- piuttosto / piuttósto.}}
\expandafter\def\csname wpAlias:[s002R]\endcsname{V}
\expandafter\def\csname wpAlias:[s004R]\endcsname{IX}
\expandafter\def\csname wpAlias:[s005R]\endcsname{XI}
\expandafter\def\csname wpAlias:[s008L]\endcsname{XVI}
\expandafter\def\csname wpAlias:[s015L]\endcsname{14}
\expandafter\def\csname wpAlias:[s015R]\endcsname{15}
\expandafter\def\csname wpAlias:[s030R]\endcsname{41}
\expandafter\def\csname wpAlias:[s043L]\endcsname{66}
\expandafter\def\csname wpAlias:[s043R]\endcsname{67}
\expandafter\def\csname wpAlias:[s054R]\endcsname{89}
\expandafter\def\csname wpAlias:[s060R]\endcsname{101}
\expandafter\def\csname wpDefer:136\endcsname{1}
\expandafter\def\csname wpDefer:145\endcsname{1}
\expandafter\def\csname wpAlias:[s095L]\endcsname{166}
\expandafter\def\csname wpAlias:[s095R]\endcsname{167}
\expandafter\def\csname wpDefer:206\endcsname{1}
\expandafter\def\csname wpAlias:[s120L]\endcsname{216}
\expandafter\def\csname wpAlias:[s120R]\endcsname{217}
\expandafter\def\csname wpAlias:[s130L]\endcsname{236}
\expandafter\def\csname wpAlias:[s130R]\endcsname{237}
\expandafter\def\csname wpDefer:259\endcsname{1}
\expandafter\def\csname wpAlias:[s146R]\endcsname{267}
\expandafter\def\csname wpAlias:[s174R]\endcsname{325}
\expandafter\def\csname wpAlias:[s198L]\endcsname{370}
\expandafter\def\csname wpAlias:[s198R]\endcsname{371}
\expandafter\def\csname wpDefer:409\endcsname{1}
\expandafter\def\csname wpAlias:[s226R]\endcsname{427}
\expandafter\def\csname wpDefer:445\endcsname{1}
\expandafter\def\csname wpNotes:[s001R]\endcsname{\csname ednote1\endcsname}
\expandafter\def\csname wpNotes:[s004R]\endcsname{\csname ednote2\endcsname\csname ednote3\endcsname}
\expandafter\def\csname wpNotes:[s005R]\endcsname{\csname ednote6\endcsname}
\expandafter\def\csname wpNotes:XIV\endcsname{\csname ednote7\endcsname}
\expandafter\def\csname wpNotes:3\endcsname{\csname ednote8\endcsname}
\expandafter\def\csname wpNotes:5\endcsname{}
\expandafter\def\csname wpNotes:11\endcsname{\csname ednote10\endcsname}
\expandafter\def\csname wpNotes:12\endcsname{\csname ednote11\endcsname\csname ednote12\endcsname}
\expandafter\def\csname wpNotes:19\endcsname{\csname ednote15\endcsname}
\expandafter\def\csname wpNotes:22\endcsname{\csname ednote16\endcsname}
\expandafter\def\csname wpNotes:24\endcsname{\csname ednote18\endcsname}
\expandafter\def\csname wpNotes:27\endcsname{\csname ednote19\endcsname}
\expandafter\def\csname wpNotes:34\endcsname{}
\expandafter\def\csname wpNotes:39\endcsname{\csname ednote21\endcsname}
\expandafter\def\csname wpNotes:[s030R]\endcsname{\csname ednote23\endcsname}
\expandafter\def\csname wpNotes:43\endcsname{}
\expandafter\def\csname wpNotes:44\endcsname{\csname ednote27\endcsname}
\expandafter\def\csname wpNotes:45\endcsname{\csname ednote28\endcsname}
\expandafter\def\csname wpNotes:49\endcsname{\csname ednote29\endcsname}
\expandafter\def\csname wpNotes:50\endcsname{\csname ednote30\endcsname}
\expandafter\def\csname wpNotes:55\endcsname{}
\expandafter\def\csname wpNotes:57\endcsname{}
\expandafter\def\csname wpNotes:58\endcsname{\csname ednote33\endcsname}
\expandafter\def\csname wpNotes:61\endcsname{\csname ednote35\endcsname\csname ednote36\endcsname}
\expandafter\def\csname wpNotes:68\endcsname{}
\expandafter\def\csname wpNotes:77\endcsname{\csname ednote40\endcsname\csname ednote41\endcsname}
\expandafter\def\csname wpNotes:79\endcsname{\csname ednote42\endcsname\csname ednote43\endcsname}
\expandafter\def\csname wpNotes:84\endcsname{\csname ednote45\endcsname\csname ednote46\endcsname}
\expandafter\def\csname wpNotes:85\endcsname{\csname ednote47\endcsname}
\expandafter\def\csname wpNotes:87\endcsname{\csname ednote48\endcsname\csname ednote50\endcsname}
\expandafter\def\csname wpNotes:[s054R]\endcsname{}
\expandafter\def\csname wpNotes:91\endcsname{\csname ednote53\endcsname}
\expandafter\def\csname wpNotes:92\endcsname{\csname ednote54\endcsname}
\expandafter\def\csname wpNotes:94\endcsname{\csname ednote55\endcsname}
\expandafter\def\csname wpNotes:95\endcsname{\csname ednote56\endcsname}
\expandafter\def\csname wpNotes:98\endcsname{\csname ednote57\endcsname}
\expandafter\def\csname wpNotes:99\endcsname{\csname ednote58\endcsname}
\expandafter\def\csname wpNotes:[s060R]\endcsname{\csname ednote59\endcsname}
\expandafter\def\csname wpNotes:102\endcsname{\csname ednote60\endcsname\csname ednote61\endcsname}
\expandafter\def\csname wpNotes:103\endcsname{\csname ednote62\endcsname\csname ednote63\endcsname}
\expandafter\def\csname wpNotes:109\endcsname{\csname ednote64\endcsname\csname ednote65\endcsname}
\expandafter\def\csname wpNotes:111\endcsname{\csname ednote66\endcsname}
\expandafter\def\csname wpNotes:115\endcsname{\csname ednote68\endcsname}
\expandafter\def\csname wpNotes:116\endcsname{\csname ednote70\endcsname}
\expandafter\def\csname wpNotes:118\endcsname{\csname ednote71\endcsname}
\expandafter\def\csname wpNotes:119\endcsname{\csname ednote72\endcsname\csname ednote73\endcsname}
\expandafter\def\csname wpNotes:121\endcsname{}
\expandafter\def\csname wpNotes:124\endcsname{\csname ednote76\endcsname\csname ednote77\endcsname}
\expandafter\def\csname wpNotes:127\endcsname{\csname ednote79\endcsname}
\expandafter\def\csname wpNotes:130\endcsname{\csname ednote80\endcsname}
\expandafter\def\csname wpNotes:133\endcsname{\csname ednote82\endcsname}
\expandafter\def\csname wpNotes:136\endcsname{}
\expandafter\def\csname wpNotes:141\endcsname{\csname ednote84\endcsname}
\expandafter\def\csname wpNotes:144\endcsname{\csname ednote86\endcsname}
\expandafter\def\csname wpNotes:145\endcsname{\csname ednote87\endcsname}
\expandafter\def\csname wpNotes:146\endcsname{\csname ednote88\endcsname}
\expandafter\def\csname wpNotes:148\endcsname{\csname ednote89\endcsname}
\expandafter\def\csname wpNotes:149\endcsname{\csname ednote90\endcsname}
\expandafter\def\csname wpNotes:150\endcsname{}
\expandafter\def\csname wpNotes:152\endcsname{\csname ednote92\endcsname}
\expandafter\def\csname wpNotes:154\endcsname{\csname ednote93\endcsname}
\expandafter\def\csname wpNotes:155\endcsname{\csname ednote94\endcsname}
\expandafter\def\csname wpNotes:157\endcsname{\csname ednote95\endcsname}
\expandafter\def\csname wpNotes:158\endcsname{\csname ednote96\endcsname}
\expandafter\def\csname wpNotes:159\endcsname{\csname ednote97\endcsname}
\expandafter\def\csname wpNotes:160\endcsname{\csname ednote99\endcsname}
\expandafter\def\csname wpNotes:161\endcsname{\csname ednote100\endcsname}
\expandafter\def\csname wpNotes:164\endcsname{\csname ednote101\endcsname\csname ednote102\endcsname}
\expandafter\def\csname wpNotes:165\endcsname{\csname ednote103\endcsname\csname ednote104\endcsname}
\expandafter\def\csname wpNotes:168\endcsname{\csname ednote105\endcsname}
\expandafter\def\csname wpNotes:169\endcsname{\csname ednote107\endcsname}
\expandafter\def\csname wpNotes:170\endcsname{\csname ednote108\endcsname\csname ednote109\endcsname}
\expandafter\def\csname wpNotes:172\endcsname{\csname ednote110\endcsname}
\expandafter\def\csname wpNotes:173\endcsname{\csname ednote111\endcsname\csname ednote112\endcsname}
\expandafter\def\csname wpNotes:175\endcsname{\csname ednote113\endcsname\csname ednote114\endcsname\csname ednote115\endcsname}
\expandafter\def\csname wpNotes:176\endcsname{\csname ednote116\endcsname\csname ednote117\endcsname}
\expandafter\def\csname wpNotes:178\endcsname{}
\expandafter\def\csname wpNotes:182\endcsname{\csname ednote119\endcsname\csname ednote120\endcsname}
\expandafter\def\csname wpNotes:183\endcsname{}
\expandafter\def\csname wpNotes:185\endcsname{\csname ednote122\endcsname\csname ednote123\endcsname}
\expandafter\def\csname wpNotes:186\endcsname{\csname ednote124\endcsname\csname ednote125\endcsname\csname ednote126\endcsname\csname ednote127\endcsname}
\expandafter\def\csname wpNotes:187\endcsname{\csname ednote128\endcsname\csname ednote129\endcsname}
\expandafter\def\csname wpNotes:191\endcsname{\csname ednote130\endcsname\csname ednote131\endcsname}
\expandafter\def\csname wpNotes:194\endcsname{\csname ednote133\endcsname}
\expandafter\def\csname wpNotes:199\endcsname{\csname ednote134\endcsname}
\expandafter\def\csname wpNotes:204\endcsname{\csname ednote135\endcsname}
\expandafter\def\csname wpNotes:205\endcsname{}
\expandafter\def\csname wpNotes:206\endcsname{\csname ednote137\endcsname}
\expandafter\def\csname wpNotes:208\endcsname{\csname ednote138\endcsname}
\expandafter\def\csname wpNotes:212\endcsname{\csname ednote139\endcsname\csname ednote142\endcsname}
\expandafter\def\csname wpNotes:214\endcsname{\csname ednote143\endcsname}
\expandafter\def\csname wpNotes:215\endcsname{}
\expandafter\def\csname wpNotes:[s120R]\endcsname{}
\expandafter\def\csname wpNotes:220\endcsname{\csname ednote147\endcsname}
\expandafter\def\csname wpNotes:222\endcsname{\csname ednote148\endcsname}
\expandafter\def\csname wpNotes:226\endcsname{\csname ednote149\endcsname}
\expandafter\def\csname wpNotes:229\endcsname{}
\expandafter\def\csname wpNotes:230\endcsname{\csname ednote151\endcsname\csname ednote152\endcsname}
\expandafter\def\csname wpNotes:235\endcsname{\csname ednote153\endcsname}
\expandafter\def\csname wpNotes:[s130R]\endcsname{\csname ednote154\endcsname}
\expandafter\def\csname wpNotes:240\endcsname{\csname ednote155\endcsname}
\expandafter\def\csname wpNotes:242\endcsname{\csname ednote156\endcsname}
\expandafter\def\csname wpNotes:244\endcsname{}
\expandafter\def\csname wpNotes:246\endcsname{}
\expandafter\def\csname wpNotes:248\endcsname{\csname ednote159\endcsname}
\expandafter\def\csname wpNotes:251\endcsname{\csname ednote160\endcsname}
\expandafter\def\csname wpNotes:253\endcsname{}
\expandafter\def\csname wpNotes:256\endcsname{\csname ednote162\endcsname}
\expandafter\def\csname wpNotes:257\endcsname{\csname ednote163\endcsname}
\expandafter\def\csname wpNotes:258\endcsname{}
\expandafter\def\csname wpNotes:259\endcsname{\csname ednote165\endcsname}
\expandafter\def\csname wpNotes:260\endcsname{\csname ednote166\endcsname}
\expandafter\def\csname wpNotes:262\endcsname{\csname ednote167\endcsname}
\expandafter\def\csname wpNotes:263\endcsname{\csname ednote168\endcsname\csname ednote169\endcsname\csname ednote171\endcsname}
\expandafter\def\csname wpNotes:265\endcsname{\csname ednote172\endcsname}
\expandafter\def\csname wpNotes:[s146R]\endcsname{\csname ednote173\endcsname}
\expandafter\def\csname wpNotes:274\endcsname{\csname ednote174\endcsname}
\expandafter\def\csname wpNotes:275\endcsname{\csname ednote176\endcsname}
\expandafter\def\csname wpNotes:277\endcsname{\csname ednote177\endcsname}
\expandafter\def\csname wpNotes:282\endcsname{\csname ednote178\endcsname}
\expandafter\def\csname wpNotes:283\endcsname{}
\expandafter\def\csname wpNotes:286\endcsname{\csname ednote180\endcsname}
\expandafter\def\csname wpNotes:287\endcsname{\csname ednote181\endcsname}
\expandafter\def\csname wpNotes:290\endcsname{\csname ednote182\endcsname}
\expandafter\def\csname wpNotes:291\endcsname{}
\expandafter\def\csname wpNotes:292\endcsname{}
\expandafter\def\csname wpNotes:293\endcsname{\csname ednote185\endcsname\csname ednote186\endcsname}
\expandafter\def\csname wpNotes:295\endcsname{\csname ednote188\endcsname\csname ednote189\endcsname\csname ednote190\endcsname}
\expandafter\def\csname wpNotes:296\endcsname{\csname ednote191\endcsname\csname ednote192\endcsname\csname ednote193\endcsname\csname ednote194\endcsname}
\expandafter\def\csname wpNotes:297\endcsname{\csname ednote195\endcsname}
\expandafter\def\csname wpNotes:299\endcsname{\csname ednote196\endcsname}
\expandafter\def\csname wpNotes:303\endcsname{\csname ednote197\endcsname}
\expandafter\def\csname wpNotes:305\endcsname{\csname ednote198\endcsname\csname ednote199\endcsname}
\expandafter\def\csname wpNotes:306\endcsname{\csname ednote200\endcsname}
\expandafter\def\csname wpNotes:307\endcsname{\csname ednote201\endcsname}
\expandafter\def\csname wpNotes:309\endcsname{\csname ednote202\endcsname}
\expandafter\def\csname wpNotes:310\endcsname{\csname ednote203\endcsname}
\expandafter\def\csname wpNotes:311\endcsname{\csname ednote204\endcsname}
\expandafter\def\csname wpNotes:312\endcsname{}
\expandafter\def\csname wpNotes:318\endcsname{\csname ednote206\endcsname\csname ednote207\endcsname}
\expandafter\def\csname wpNotes:319\endcsname{\csname ednote208\endcsname}
\expandafter\def\csname wpNotes:320\endcsname{}
\expandafter\def\csname wpNotes:321\endcsname{\csname ednote210\endcsname}
\expandafter\def\csname wpNotes:324\endcsname{\csname ednote211\endcsname}
\expandafter\def\csname wpNotes:330\endcsname{\csname ednote212\endcsname}
\expandafter\def\csname wpNotes:332\endcsname{\csname ednote213\endcsname}
\expandafter\def\csname wpNotes:333\endcsname{\csname ednote214\endcsname\csname ednote215\endcsname}
\expandafter\def\csname wpNotes:334\endcsname{\csname ednote216\endcsname}
\expandafter\def\csname wpNotes:335\endcsname{\csname ednote217\endcsname}
\expandafter\def\csname wpNotes:340\endcsname{\csname ednote218\endcsname}
\expandafter\def\csname wpNotes:344\endcsname{\csname ednote220\endcsname\csname ednote221\endcsname}
\expandafter\def\csname wpNotes:346\endcsname{\csname ednote222\endcsname}
\expandafter\def\csname wpNotes:348\endcsname{\csname ednote223\endcsname\csname ednote224\endcsname}
\expandafter\def\csname wpNotes:349\endcsname{\csname ednote225\endcsname}
\expandafter\def\csname wpNotes:354\endcsname{\csname ednote226\endcsname}
\expandafter\def\csname wpNotes:357\endcsname{\csname ednote227\endcsname}
\expandafter\def\csname wpNotes:358\endcsname{\csname ednote228\endcsname}
\expandafter\def\csname wpNotes:360\endcsname{\csname ednote229\endcsname}
\expandafter\def\csname wpNotes:361\endcsname{\csname ednote230\endcsname}
\expandafter\def\csname wpNotes:362\endcsname{\csname ednote231\endcsname}
\expandafter\def\csname wpNotes:363\endcsname{\csname ednote232\endcsname}
\expandafter\def\csname wpNotes:365\endcsname{\csname ednote233\endcsname}
\expandafter\def\csname wpNotes:366\endcsname{\csname ednote234\endcsname\csname ednote235\endcsname}
\expandafter\def\csname wpNotes:367\endcsname{}
\expandafter\def\csname wpNotes:368\endcsname{\csname ednote237\endcsname\csname ednote238\endcsname\csname ednote239\endcsname\csname ednote240\endcsname\csname ednote241\endcsname}
\expandafter\def\csname wpNotes:372\endcsname{\csname ednote242\endcsname}
\expandafter\def\csname wpNotes:375\endcsname{\csname ednote243\endcsname}
\expandafter\def\csname wpNotes:376\endcsname{\csname ednote244\endcsname\csname ednote245\endcsname}
\expandafter\def\csname wpNotes:377\endcsname{\csname ednote246\endcsname}
\expandafter\def\csname wpNotes:378\endcsname{\csname ednote247\endcsname\csname ednote248\endcsname}
\expandafter\def\csname wpNotes:382\endcsname{\csname ednote249\endcsname}
\expandafter\def\csname wpNotes:383\endcsname{\csname ednote250\endcsname\csname ednote251\endcsname\csname ednote252\endcsname}
\expandafter\def\csname wpNotes:385\endcsname{\csname ednote253\endcsname\csname ednote254\endcsname}
\expandafter\def\csname wpNotes:386\endcsname{\csname ednote255\endcsname}
\expandafter\def\csname wpNotes:388\endcsname{\csname ednote256\endcsname\csname ednote257\endcsname}
\expandafter\def\csname wpNotes:389\endcsname{\csname ednote258\endcsname\csname ednote259\endcsname}
\expandafter\def\csname wpNotes:390\endcsname{\csname ednote260\endcsname\csname ednote261\endcsname}
\expandafter\def\csname wpNotes:395\endcsname{\csname ednote262\endcsname}
\expandafter\def\csname wpNotes:396\endcsname{\csname ednote263\endcsname}
\expandafter\def\csname wpNotes:397\endcsname{\csname ednote264\endcsname}
\expandafter\def\csname wpNotes:400\endcsname{\csname ednote265\endcsname}
\expandafter\def\csname wpNotes:403\endcsname{\csname ednote266\endcsname}
\expandafter\def\csname wpNotes:405\endcsname{\csname ednote267\endcsname\csname ednote268\endcsname}
\expandafter\def\csname wpNotes:408\endcsname{\csname ednote269\endcsname\csname ednote270\endcsname\csname ednote271\endcsname\csname ednote272\endcsname\csname ednote273\endcsname}
\expandafter\def\csname wpNotes:409\endcsname{\csname ednote274\endcsname}
\expandafter\def\csname wpNotes:411\endcsname{\csname ednote275\endcsname\csname ednote276\endcsname\csname ednote277\endcsname\csname ednote278\endcsname}
\expandafter\def\csname wpNotes:412\endcsname{\csname ednote279\endcsname}
\expandafter\def\csname wpNotes:415\endcsname{\csname ednote280\endcsname\csname ednote281\endcsname}
\expandafter\def\csname wpNotes:416\endcsname{\csname ednote282\endcsname}
\expandafter\def\csname wpNotes:422\endcsname{\csname ednote284\endcsname}
\expandafter\def\csname wpNotes:425\endcsname{\csname ednote285\endcsname}
\expandafter\def\csname wpNotes:426\endcsname{}
\expandafter\def\csname wpNotes:430\endcsname{}
\expandafter\def\csname wpNotes:431\endcsname{\csname ednote288\endcsname\csname ednote289\endcsname}
\expandafter\def\csname wpNotes:432\endcsname{\csname ednote290\endcsname}
\expandafter\def\csname wpNotes:433\endcsname{}
\expandafter\def\csname wpNotes:434\endcsname{\csname ednote292\endcsname}
\expandafter\def\csname wpNotes:435\endcsname{\csname ednote293\endcsname}
\expandafter\def\csname wpNotes:437\endcsname{\csname ednote294\endcsname}
\expandafter\def\csname wpNotes:440\endcsname{\csname ednote295\endcsname}
\expandafter\def\csname wpNotes:442\endcsname{}
\expandafter\def\csname wpNotes:443\endcsname{\csname ednote298\endcsname}
\expandafter\def\csname wpNotes:444\endcsname{\csname ednote300\endcsname}
\expandafter\def\csname wpNotes:445\endcsname{}
\expandafter\def\csname wpNotes:446\endcsname{\csname ednote303\endcsname}
\expandafter\def\csname wpNotes:448\endcsname{}
\expandafter\def\csname wpNotes:449\endcsname{\csname ednote305\endcsname\csname ednote306\endcsname}
\expandafter\def\csname wpNotes:451\endcsname{\csname ednote307\endcsname\csname ednote308\endcsname}
\expandafter\def\csname wpNotes:453\endcsname{\csname ednote309\endcsname}
\expandafter\def\csname wpNotes:455\endcsname{\csname ednote311\endcsname}
\expandafter\def\csname wpNotes:456\endcsname{\csname ednote312\endcsname}
\expandafter\def\csname wpNotes:457\endcsname{\csname ednote313\endcsname\csname ednote314\endcsname\csname ednote315\endcsname}
\expandafter\def\csname wpNotes:459\endcsname{\csname ednote316\endcsname}
\expandafter\def\csname wpNotes:460\endcsname{}
\expandafter\def\csname wpNotes:461\endcsname{}
\expandafter\def\csname wpNotes:463\endcsname{\csname ednote319\endcsname\csname ednote321\endcsname}
\expandafter\def\csname wpNotes:464\endcsname{\csname ednote322\endcsname}
\begin{document}\frontmatter

% !PRINTED
\weberblank{[s001L]}
% head:
% scan: s001L; nessun numero di pagina stampato.

\clearpage\thispagestyle{empty}\vspace*{12mm}
% !PRINTED
\weberpage{[s001R]}
% scan: s001R; nessun numero di pagina stampato.
\begin{center}
{\Large\textbf{\gesperrt{FEDERIGO ENRIQUES}}}\par\vspace{23mm}

{\fontsize{26}{31}\selectfont\bfseries ALGEBRAIC\\SURFACES\par}
\end{center}

\weberfigure{}{Publisher’s device: a man in a cap strides forward scattering seeds; on the base the inscription «LABORAVI FIDENTER».}

\begin{center}
\textbf{\gesperrt{NICOLA ZANICHELLI PUBLISHER}} %

\textbf{\gesperrt{BOLOGNA 1949}}
\end{center}
\clearpage\editionpage\clearpage

\thispagestyle{empty}\vspace*{50mm}
% !PRINTED
\weberpage{[s002L]}
% scan: s002L; nessun numero di pagina stampato.
\begin{center}
\textbf{THE PUBLISHER, HAVING FULFILLED ITS OBLIGATIONS, WILL EXERCISE THE RIGHTS SANCTIONED BY LAW}
\end{center}

\noindent\rule{\linewidth}{0.4pt}
\begin{center}
\textbf{S.T.E.B. - Società Tipografica Editrice Bolognese - \textsc{vi} - 1949}
\end{center}
\clearpage

\clearpage\thispagestyle{plain}\markboth{Preface by G. Castelnuovo}{Preface by G. Castelnuovo}\pdfbookmark[0]{Preface by G. Castelnuovo}{front.castelnuovo}
% !PRINTED
\weberpage{[s002R]}
% scan: s002R; nessun numero di pagina stampato.
\textit{The volume which sees the light today contains, reworked and set in order, the greater part of the mathematical work of Federigo Enriques, namely the theory of algebraic surfaces, to which his name will always remain linked. This work, begun immediately after his graduation in \textup{1893} and pursued for more than half a century until his death, with brief interruptions when problems of philosophy and of the history of science occupied his universal mind, lies scattered in some fifty Notes and Memoirs published in various collections.}

\textit{The present volume will serve to bring out the high value of this work, which opened up a vast and fertile field, still almost unexplored when Enriques began his researches. It will enable scholars of algebraic geometry, and in particular university students, to find their bearings in this wide territory.}

\textit{The late Author had already felt the need for a systematic treatment of the geometry on algebraic surfaces after his call to the University of Rome, where he could devote himself exclusively to the advanced courses in mathematics and where the number of pupils was greater. Availing himself of the collaboration of Luigi Campedelli, then his assistant and mine, now professor at the University of Florence, he gathered the general part of the theory into a lithographed volume, which appeared in \textup{1932}; while a second part, concerning particular categories of surfaces, was published in print two years later by the Mathematical Seminar of the University of Rome, with the help of the same collaborator.}

\textit{These two volumes having sold out within a few years, the Author had to think of a second edition; and during the enforced leisure to which he was constrained by the racial laws, and then by the war and the persecutions, he meditated a reworking of the whole work, with the purpose of simplifying certain proofs or putting them beyond the reach of criticism, of completing or extending classifications of particular types of surfaces, and of taking into account the most recent results which he himself or other researchers had reached.}

\textit{The manuscript was fortunately complete when death overtook Federigo Enriques in the full lucidity of his mind.}

% !PRINTED VI
\weberpage{VI}
% scan: s003L.
\textit{The Author had reserved for himself the making of certain final corrections on the printer’s proofs, but unfortunately fate did not grant him the time for this task. It was taken on, with the finesse of their minds, with the memory of their master’s teachings, with the affection of sons, by his last disciples Pompilj and Franchetta. And I, who assisted them in their doubts on some obscure or less satisfactory point of the text, can say that the task was arduous and would have required, to be fully accomplished, researches not of months but of years. The Author himself takes care to warn, from the preface onwards, that the treatise, rather than expounding a doctrine already static and crystallized, aspires to arouse in the reader the desire to bring complements and improvements to various theories. And where the ground is less solid the Author puts the student on his guard.}

\textit{Of these points still in flux, the one which presents the hardest difficulty and the greatest interest in the general theory of surfaces forms the subject of Chap. IX, which deals with the continuous systems of algebraic curves (not contained in linear systems) that exist on every irregular surface. A brilliant intuition led Enriques in \textup{1904} to state and establish this characteristic property of irregular surfaces; subsequent fundamental researches have rested upon it. But a careful criticism showed, several years later, that Enriques’s proof is not satisfactory. The property is true, at least under some restriction, as follows from the researches of Henri Poincaré by transcendental methods. But all the attempts made afterwards by Enriques and by others to prove it by means of algebro-geometric considerations have come up against difficulties hitherto unsurmounted. This is explicitly stated in the aforesaid Chap. IX, in which the Author also gives suggestions on a path to be tried in order to reach the goal. I must confess that I do not see how that path could be turned into an irreproachable procedure.}

\textit{Other doubts, of far lesser importance, may be raised by some points of Chaps. VII, VIII and X, which deal with the classification and construction of very particular types of surfaces. Here it is worth recalling that Enriques’s fundamental researches can be divided into two groups:}

\textit{\textup{a)} the general theory of algebraic surfaces, and in particular the determination of the characters invariant under birational transformations: essentially, the \textup{genera;}}

\textit{\textup{b)} the classification of the surfaces of the lowest genera.}

\textit{This second group led the Author to the discovery of very interesting and unforeseen families of surfaces, and also furnished conceptual models which suggested the discovery of hidden properties belonging to the general theory. But all researchers know what pitfalls lurk in very minute classifications, where some hidden cases easily escape notice, and others, believed at}
% !PRINTED VII
\weberpage{VII}
% scan: s003R.
\textit{first to be realizable, turn out to be non-existent under a more scrupulous examination.}

\textit{Here, then, as the Author himself suggests, there presents itself a first field of inquiry for the student, who, after having become familiar with the delicate discussions which the subject entails, will then be in a position to face the gravest difficulties of the general theory, difficulties which concern not only the question to which we alluded above, but are met with whenever one wishes to bring into clear light the links between the algebraic and the transcendental (or the topological) treatment of algebraic surfaces. Naturally, for these latter researches the student will first have to study the transcendental theory in depth, reading the treatises of Picard-Simart, of Lefschetz, Zariski, Hodge, the above-mentioned Memoir of Poincaré, and the important algebro-transcendental works of Severi, in which are expounded researches that go beyond the programme of the present volume.}

\textit{Will there soon come the continuator of the work of the Italian and French schools, who will succeed in giving the theory of algebraic surfaces the perfection that the theory of algebraic curves has attained? I hope so, but I doubt it.}

\textit{I am led to entertain my doubts by the observation that mathematics has taken, in the present century, a direction very different from the one that prevailed in the last century. Imagination and intuition, which guided the research of those days, are today regarded with suspicion, for dread of the errors to which they may lead. Theories arose in answer to the need the mathematician felt to delineate and make precise objects of thought which were already present to his mind in a vague form. It was the exploration of a wide territory glimpsed from a distant summit. Thus were built, in the last century, those jewels which are called the theory of analytic functions, of elliptic and abelian functions, surfaces of minimal area, cubic surfaces.... Today, more than the ground to be explored, what is of interest is the road that leads to it, and this road is now strewn with artificial obstacles, now soars among the clouds.}

\textit{One is led to believe that this tendency is not a short-lived manifestation by comparison with the analogous phenomenon occurring in the plastic arts and in music. There too imagination is banished as a survival of the Romantic age; there too the content of the work of art interests less than the technique or the means of expression which the artist employs. Various illustrious critics think that this represents a decadence of art, and draw from it gloomy forebodings about the future of our civilization.}

\textit{It would be rash to extend these pessimistic judgements to the evolution that mathematics is undergoing. Nevertheless, when one compares the half-century now about to close with the corresponding period of the last century}
% !PRINTED VIII
\weberpage{VIII}
% scan: s004L.
\textit{in which there flourished names such as Gauss, Abel, Jacobi, Cauchy and so many other very great ones, one cannot conceal misgivings about the future of our science.}

\textit{One day, however, near or distant, the love for the great theories whose foundations were laid by our masters of the nineteenth century will be reborn. And on that day the treatise of Federigo Enriques will be read and meditated upon as the account of an exploration in a territory where many gems have already been gathered and many others await one who is worthy of discovering them.}

\begin{flushright}
\textsc{G. Castelnuovo}
\end{flushright}

Rome, January 1949.

\clearpage\thispagestyle{plain}
% !PRINTED
\weberpage{[s004R]}
% scan: s004R; nessun numero di pagina stampato.
\edfrontheading{\gesperrt{\textit{PREFACE}}}{PREFACE}{front.autore}

\textit{In this treatise we propose to expound the theory of algebraic surfaces, regarded particularly in its algebro-geometric aspect, as it has gradually matured over half a century, especially in the Italian geometric school. For many years we have worked on this exposition, revising and sometimes redoing, or at least reordering, the older proofs, so as to give the theory-itself the most rigorous and simplest arrangement.} %

\textit{This work of revision and ordering was carried out first of all through the courses of lectures we gave at the University of Rome, which were collected, in a first redaction, by Luigi Campedelli, now professor at the University of Florence, to whom we are glad to attest our gratitude once again. The lectures of Enriques-Campedelli were published in two volumes; the first part, with the title « Lezioni sulla teoria delle superficie algebriche, raccolte dal Dott. Luigi Campedelli », in a lithographed edition by the publishing house Cedam of Padua, in \textup{1932}, and the second part in \textup{1934}, with the title « Sulla classificazione delle superficie algebriche, particolarmente di genere zero, lezioni raccolte dal Dott. Luigi Campedelli », in the Rendiconti del Seminario Matematico della R. Università di Roma.} %

\textit{These volumes remain even today the foundation of our work of systematizing the theory, a work which has nevertheless been carried on since then, both from the university chair and in the retirement of recent years, after October \textup{1938}.}

\textit{We have taken up Campedelli’s exposition again, sometimes correcting or simplifying some point, but above all giving the theory itself a wider extension by making use of the contributions brought to the science by some young geometers. Among these it is dear to us to recall the names of other disciples of ours who, after our Florentine Colleague, have helped us in this reconstruction, going more deeply into some points or inviting us to reconsider them through the discussion of our ideas and our explanations: they are Prof. Giuseppe Pompilj --- soon called away to arms --- and Dr. Alfredo Franchetta (two young hopes of our science), who will be found cited in the course of these Lectures; to whom we wish publicly to attest our gratitude.}

% !PRINTED X
\weberpage{X}
% scan: s005L.
\textit{We need not explain at greater length the plan of the work, which already appears from the general table of contents. But it matters to us to point out that we have sought not only to give the theories expounded a logical arrangement, but also (as in earlier works) to offer a historical perspective of their development. In this way we wish to offer the reader, not the gift of something perfect that lets itself be contemplated from outside, but rather the view of an acquisition and of a progress whose reasons he must understand, and which he is invited to win back by himself and for himself, finding in the book an instrument of work. Precisely for this reason the inductive path of science, which starts from examples and particular cases, and at times also the errors which had to be overcome and corrected, are especially brought into relief.}

\textit{We add that, as regards the knowledge of preliminary theories (algebraic curves, rational surfaces), we refer particularly to the exposition given of them in other lectures of ours, which were collected in the four volumes of Enriques-Chisini « Teoria geometrica delle equazioni e delle funzioni algebriche »\origfootnote{1}{s005L-1}{\textup{Bologna, Zanichelli 1915, 1918, 1924, 1934.}} and in that of Fabio Conforto on « Le superficie razionali », likewise drawn up from our lectures\origfootnote{2}{s005L-2}{\textup{Bologna, Zanichelli, 1939.}}.}

\textit{The text which we offer the reader was composed and found its definitive form by the month of April \textup{1942}, when fascist persecution forbade, in our country, even the most innocent mathematical publications, and moreover publication abroad gave rise to difficulties which we were afterwards to experience. In it reference was made to some notes, in particular notes sent by us (as a corresponding member) to the Academy of Sciences of Madrid, about which we have not been able to learn whether they actually found a place in those academic proceedings.}

\textit{Today, having come out of the stormy period of the war, at the moment of handing the manuscript over to the printer, we still find it difficult to gain an exact notion of the scientific literature of the last three years; and on the other hand the commitments of new works prevent us from going back over the work done in order to re-examine it in the light of certain studies which, more or less approximately, we came to know only later. Therefore we preserve for the text of the present lectures the form given to it, as has been said, in \textup{1942}, save for completing some bibliographical information and indicating, with a few notes, the gaps which would require a new examination.}

\begin{flushright}
\textsc{F. Enriques}
\end{flushright}

Rome, September 1945.

\clearpage\thispagestyle{plain}
% !PRINTED
\weberpage{[s005R]}
% scan: s005R; nessun numero di pagina stampato.
\begingroup\small
\edfrontheading{\gesperrt{CONTENTS}}{CONTENTS}{front.indice}

\begin{longtable}{@{}p{0.79\linewidth}cr@{}}
\textit{Preface} by Prof. \textsc{G. Castelnuovo} \dotfill & Page & \hyperlink{old:V}{\textsc{v}} \\
\textit{Preface} by the Author \dotfill & » & \hyperlink{old:IX}{\textsc{ix}} \\
\textit{Introduction} \dotfill & » & \hyperlink{old:1}{1} \\
\hspace*{1em}\makebox[2em][r]{1. --} Review of elementary notions \dotfill & » & \hyperlink{old:1}{1} \\
\hspace*{1em}\makebox[2em][r]{2. --} Rational transformations \dotfill & » & \hyperlink{old:2}{2} \\
\hspace*{1em}\makebox[2em][r]{3. --} Surfaces in hyperspaces: resolution of singularities & » & \hyperlink{old:5}{5} \\
\hspace*{1em}\makebox[2em][r]{4. --} Surfaces with normal singularities \dotfill & » & \hyperlink{old:6}{6} \\
\hspace*{1em}\makebox[2em][r]{5. --} Note \dotfill & » & \hyperlink{old:8}{8} \\[1ex]
\textsc{Chapter} I. -- \textit{Linear systems of curves} \dotfill & » & \hyperlink{old:15}{15} \\ %
\hspace*{1em}\makebox[2em][r]{1. --} Linear pencils \dotfill & » & \hyperlink{old:15}{15} \\
\hspace*{1em}\makebox[2em][r]{2. --} Irrational pencils \dotfill & » & \hyperlink{old:17}{17} \\
\hspace*{1em}\makebox[2em][r]{3. --} Linear systems \dotfill & » & \hyperlink{old:18}{18} \\
\hspace*{1em}\makebox[2em][r]{4. --} Characteristic property of linear systems \dotfill & » & \hyperlink{old:19}{19} \\
\hspace*{1em}\makebox[2em][r]{5. --} Extension of Bertini’s theorems \dotfill & » & \hyperlink{old:20}{20} \\
\hspace*{1em}\makebox[2em][r]{6. --} Surfaces that are images of linear systems \dotfill & » & \hyperlink{old:22}{22} \\
\hspace*{1em}\makebox[2em][r]{7. --} Equivalent curves and complete linear systems \dotfill & » & \hyperlink{old:24}{24} \\
\hspace*{1em}\makebox[2em][r]{8. --} Sum and difference of linear systems: residue theorem & » & \hyperlink{old:26}{26} \\
\hspace*{1em}\makebox[2em][r]{9. --} Image surface of the sum system \dotfill & » & \hyperlink{old:30}{30} \\
\hspace*{1em}\makebox[2em][r]{10. --} Virtual characters \dotfill & » & \hyperlink{old:31}{31} \\
\hspace*{1em}\makebox[2em][r]{11. --} Exceptional curves \dotfill & » & \hyperlink{old:33}{33} \\
\hspace*{1em}\makebox[2em][r]{12. --} Note on reducible exceptional curves \dotfill & » & \hyperlink{old:37}{37} \\
\hspace*{1em}\makebox[2em][r]{13. --} Virtual characteristic series \dotfill & » & \hyperlink{old:40}{40} \\[1ex]
\textsc{Chapter} II. -- \textit{Covariant and invariant systems} \dotfill & » & \hyperlink{old:41}{41} \\
\hspace*{1em}\makebox[2em][r]{1. --} Jacobian curves \dotfill & » & \hyperlink{old:41}{41} \\
\hspace*{1em}\makebox[2em][r]{2. --} Jacobian system \dotfill & » & \hyperlink{old:42}{42} \\
\hspace*{1em}\makebox[2em][r]{3. --} Fundamental theorem \dotfill & » & \hyperlink{old:48}{48} \\
\hspace*{1em}\makebox[2em][r]{4. --} Canonical curves \dotfill & » & \hyperlink{old:50}{50} \\
\hspace*{1em}\makebox[2em][r]{5. --} Properties of canonical curves \dotfill & » & \hyperlink{old:52}{52} \\ %
\hspace*{1em}\makebox[2em][r]{6. --} Surface genus and linear genus \dotfill & » & \hyperlink{old:53}{53} \\
\hspace*{1em}\makebox[2em][r]{7. --} The exceptional curves as fixed parts of the canonical system \dotfill & » & \hyperlink{old:56}{56.} \\ %
\hspace*{1em}\makebox[2em][r]{8. --} Bicanonical and pluricanonical curves \dotfill & » & \hyperlink{old:60}{60} \\
\hspace*{1em}\makebox[2em][r]{9. --} Historical note \dotfill & » & \hyperlink{old:62}{62} \\
\end{longtable}

% scan: s006L
% !PRINTED XII
\weberpage{XII}
\begin{longtable}{@{}r@{\ }p{0.70\linewidth}@{\quad}c@{\ }r@{}}
\multicolumn{2}{@{}p{0.80\linewidth}}{\textsc{Chapter III.} – \textit{The adjoint surfaces} \dotfill } & Page & \hyperlink{old:67}{67} \\
1. – & Surfaces of order \(n-3\) adjoint to a surface of order \(n\) without proper multiple points \dotfill & » & \hyperlink{old:67}{67} \\
2. – & Adjoint surfaces of arbitrary order \dotfill & » & \hyperlink{old:69}{69} \\
3. – & Ruled surfaces \dotfill & » & \hyperlink{old:70}{70} \\
4. – & Characteristic property of canonical curves and of curves adjoint to a linear system \dotfill & » & \hyperlink{old:71}{71} \\
5. – & Isolated multiple points \dotfill & » & \hyperlink{old:72}{72} \\
6. – & Adjoints and subadjoints \dotfill & » & \hyperlink{old:73}{73} \\
7. – & Multiple points of the first kind \dotfill & » & \hyperlink{old:76}{76} \\
8. – & Double planes \dotfill & » & \hyperlink{old:77}{77} \\
9. – & The genus of the fundamental curves and its influence on the adjoints \dotfill & » & \hyperlink{old:80}{80} \\
10. – & Point singularities of higher kind \dotfill & » & \hyperlink{old:82}{82} \\
11. – & Bicanonical curves and biadjoint surfaces \dotfill & » & \hyperlink{old:89}{89} \\
12. – & Surfaces of genus \(p=0\) and bigenus \(P>0\) \dotfill & » & \hyperlink{old:92}{92} \\
13. – & Criteria of equivalence \dotfill & » & \hyperlink{old:96}{96} \\
\end{longtable}

\begin{longtable}{@{}r@{\ }p{0.70\linewidth}@{\quad}c@{\ }r@{}}
\multicolumn{2}{@{}p{0.80\linewidth}}{\textsc{Chapter IV.} – \textit{The arithmetic genus and the Riemann-Roch theorem for surfaces} \dotfill } & » & \hyperlink{old:101}{101} \\
1. – & Introduction \dotfill & » & \hyperlink{old:101}{101} \\
2. – & Conditions imposed on a surface by passing through a curve: postulation formulae \dotfill & » & \hyperlink{old:102}{102} \\
3. – & Dimension of the systems of surfaces adjoint to a given one \dotfill & » & \hyperlink{old:107}{107} \\
4. – & The arithmetic genus \dotfill & » & \hyperlink{old:110}{110} \\
5. – & Complements \dotfill & » & \hyperlink{old:114}{114} \\
6. – & Note on the conditions of regularity of a surface . & » & \hyperlink{old:117}{117} \\
7. – & Examples of irregular surfaces \dotfill & » & \hyperlink{old:120}{120} \\
8. – & The Riemann-Roch theorem for surfaces: systems larger than the canonical system \dotfill & » & \hyperlink{old:123}{123} \\
9. – & Elimination of the exceptional curves \dotfill & » & \hyperlink{old:127}{127} \\
10. – & Completeness of the series cut out by a small linear system on the curves of a large system \dotfill & » & \hyperlink{old:128}{128} \\
11. – & The Riemann-Roch theorem for arbitrary linear systems \dotfill & » & \hyperlink{old:131}{131} \\
12. – & Deficiency of the characteristic series \dotfill & » & \hyperlink{old:134}{134} \\
13. – & Historical note \dotfill & » & \hyperlink{old:137}{137} \\
14. – & Virtual curves \dotfill & » & \hyperlink{old:142}{142} \\
15. – & Regular and superabundant linear systems \dotfill & » & \hyperlink{old:147}{147} \\
16. – & Clarifications on the regularity of the system adjoint to a curve on a regular surface \dotfill & » & \hyperlink{old:150}{150} \\
17. – & Continuation: the system adjoint to a connected curve formed of multiple components on a regular surface \dotfill & » & \hyperlink{old:152}{152} \\
18. – & On the reducibility of the canonical system \dotfill & » & \hyperlink{old:156}{156} \\
19. – & On the regularity of the bicanonical system \dotfill & » & \hyperlink{old:158}{158} \\
\end{longtable}

% scan: s006R
% !PRINTED XIII
\weberpage{XIII}
\begin{longtable}{@{}r@{\ }p{0.70\linewidth}@{\quad}c@{\ }r@{}}
20. – & Note on the irreducibility of the bicanonical system \dotfill & Page & \hyperlink{old:162}{162} \\
21. – & Genera of a reducible surface \dotfill & » & \hyperlink{old:163}{163} \\
\end{longtable}

\begin{longtable}{@{}r@{\ }p{0.70\linewidth}@{\quad}c@{\ }r@{}}
\multicolumn{2}{@{}p{0.80\linewidth}}{\textsc{Chapter V.} – \textit{Numerical invariants and multiple planes} \dotfill } & » & \hyperlink{old:167}{167} \\
1. – & The Zeuthen-Segre invariant \dotfill & » & \hyperlink{old:167}{167} \\
2. – & Irrational pencil \dotfill & » & \hyperlink{old:172}{172} \\
3. – & Expression of the Zeuthen-Segre invariant by means of the genera \dotfill & » & \hyperlink{old:173}{173} \\
4. – & Cuspidal curves of a net \dotfill & » & \hyperlink{old:177}{177} \\
5. – & Historical account \dotfill & » & \hyperlink{old:180}{180} \\
6. – & Characters of a net and multiple planes \dotfill & » & \hyperlink{old:180}{180} \\
7. – & \((1,n)\) correspondence between two surfaces: double planes . & » & \hyperlink{old:183}{183} \\
8. – & Correspondence formulae \dotfill & » & \hyperlink{old:185}{185} \\
9. – & Existence theorem \dotfill & » & \hyperlink{old:192}{192} \\
10. – & Complements: limiting forms of the branch curve according to Chisini \dotfill & » & \hyperlink{old:199}{199} \\
11. – & The moduli of a class of algebraic surfaces \dotfill & » & \hyperlink{old:204}{204} \\
12. – & Digression on the completeness of the characteristic series of a complete system of plane curves with nodes and cusps \dotfill & » & \hyperlink{old:207}{207} \\
13. – & The moduli of regular surfaces of genus \(p>3\) . . & » & \hyperlink{old:209}{209} \\
14. – & Historical note and complements \dotfill & » & \hyperlink{old:213}{213} \\
\end{longtable}

\begin{longtable}{@{}r@{\ }p{0.70\linewidth}@{\quad}c@{\ }r@{}}
\multicolumn{2}{@{}p{0.80\linewidth}}{\textsc{Chapter VI.} – \textit{Regular surfaces: minimum of the genera and conditions of rationality} \dotfill } & » & \hyperlink{old:217}{217} \\
1. – & Lower bound of the linear genus \dotfill & » & \hyperlink{old:217}{217} \\
2. – & The bigenus and the plurigenera \dotfill & » & \hyperlink{old:220}{220} \\
3. – & Examples: double planes of linear genus \(p^{(1)}=1\) \dotfill & » & \hyperlink{old:224}{224} \\
4. – & Conditions of rationality of a surface \dotfill & » & \hyperlink{old:230}{230} \\
\end{longtable}

\begin{longtable}{@{}r@{\ }p{0.70\linewidth}@{\quad}c@{\ }r@{}}
\multicolumn{2}{@{}p{0.80\linewidth}}{\textsc{Chapter VII.} – \textit{Classification of the surfaces of linear genus \(p^{(1)}=1\)} \dotfill } & » & \hyperlink{old:237}{237} \\
1. – & Surfaces of genus 0 and bigenus 1, with bicanonical curve of order zero \dotfill & » & \hyperlink{old:237}{237} \\
2. – & The surfaces with all genera equal to 1 \dotfill & » & \hyperlink{old:247}{247} \\
3. – & Historical account \dotfill & » & \hyperlink{old:254}{254} \\
4. – & The surfaces of linear genus \(p^{(1)}=1\) with \(P>1\) . . & » & \hyperlink{old:257}{257} \\
\end{longtable}

\begin{longtable}{@{}r@{\ }p{0.70\linewidth}@{\quad}c@{\ }r@{}}
\multicolumn{2}{@{}p{0.80\linewidth}}{\textsc{Chapter VIII.} – \textit{Regular canonical and pluricanonical surfaces}} & » & \hyperlink{old:267}{267} \\
1. – & Introduction \dotfill & » & \hyperlink{old:267}{267} \\
2. – & Canonical surfaces of genus \(p=4\) and linear genus \(p^{(1)}=6{,}5\) \dotfill & » & \hyperlink{old:268}{268} \\
3. – & Surfaces with \(p=4\) and \(p^{(1)}=7\) \dotfill & » & \hyperlink{old:271}{271} \\
4. – & Surfaces with \(p=4\) and \(p^{(1)}=8\) \dotfill & » & \hyperlink{old:273}{273} \\
5. – & Surfaces with \(p=4\) and \(p^{(1)}=9\) \dotfill & » & \hyperlink{old:281}{281} \\
6. – & Surfaces with \(p=4\) and \(p^{(1)}>9\) \dotfill & » & \hyperlink{old:283}{283} \\
7. – & Canonical surfaces with \(p=5\) and \(p^{(1)}=9\) \dotfill & » & \hyperlink{old:284}{284} \\
\end{longtable}

% scan: s007L
% !PRINTED XIV
\weberpage{XIV}
\begin{longtable}{@{}r@{\ }p{0.70\linewidth}@{\quad}c@{\ }r@{}}
8. – & Surfaces with \(p=5\) and \(p^{(1)}=10\) \dotfill & Page & \hyperlink{old:286}{286} \\
9. – & The surfaces on a variety of 3 or more dimensions . & » & \hyperlink{old:289}{289} \\
10. – & Examples of hyperspatial canonical surfaces \dotfill & » & \hyperlink{old:292}{292} \\
11. – & Minimum value of the linear genus with respect to the surface genus \dotfill & » & \hyperlink{old:294}{294} \\
12. – & Upper bound of the linear genus \(p^{(1)}\) with respect to the surface genus \dotfill & » & \hyperlink{old:299}{299} \\
13. – & Canonical systems belonging to an involution \dotfill & » & \hyperlink{old:300}{300} \\
14. – & Surfaces of linear genus \(p^{(1)}=2\): first case \(p=2\) & » & \hyperlink{old:303}{303} \\
15. – & Surfaces with \(p^{(1)}=2\) and \(p=1\) \((P=3, P_3=5)\) \dotfill & » & \hyperlink{old:305}{305} \\
16. – & Surfaces with \(p^{(1)}=2\) and \(p=0\) \((P=2, P_3=4)\) \dotfill & » & \hyperlink{old:308}{308} \\
17. – & Surfaces of linear genus \(p^{(1)}=3\): first case \(p=3\) & » & \hyperlink{old:311}{311} \\
18. – & Surfaces with \(p^{(1)}=3\) and \(p=2\) \((P=5)\) \dotfill & » & \hyperlink{old:312}{312} \\
19. – & Surfaces with \(p^{(1)}=3\) and \(p=1\) \((P=4)\) \dotfill & » & \hyperlink{old:316}{316} \\
20. – & Surfaces with \(p=0\) and \(p^{(1)}=3\) \dotfill & » & \hyperlink{old:321}{321} \\
21. – & Simple and multiple pluricanonical surfaces \dotfill & » & \hyperlink{old:321}{321} \\
\end{longtable}

\begin{longtable}{@{}r@{\ }p{0.70\linewidth}@{\quad}c@{\ }r@{}}
\multicolumn{2}{@{}p{0.80\linewidth}}{\textsc{Chapter IX.} – \textit{Irregular surfaces and continuous systems of non-equivalent curves} \dotfill } & » & \hyperlink{old:325}{325} \\
1. – & Introduction \dotfill & » & \hyperlink{old:325}{325} \\
2. – & Arithmetic condition for a curve, on a surface of geometric irregularity \(q\), to belong to a continuous series \(\infty^q\) of non-equivalent curves \dotfill & » & \hyperlink{old:326}{326} \\
3. – & The Picard variety corresponding to an irregular surface \dotfill & » & \hyperlink{old:328}{328} \\
4. – & Characteristic property of irregular algebraic surfaces: fundamental theorem for \(p_g=0\) \dotfill & » & \hyperlink{old:328}{328} \\
5. – & Fundamental theorem for \(p_g>0\) \dotfill & » & \hyperlink{old:330}{330} \\
6. – & History of the theory of continuous systems \dotfill & » & \hyperlink{old:339}{339} \\
7. – & On various attempts at proof and extension of the fundamental theorem \dotfill & » & \hyperlink{old:347}{347} \\
8. – & The paracanonical system \dotfill & » & \hyperlink{old:354}{354} \\
9. – & Digression on abelian varieties \dotfill & » & \hyperlink{old:357}{357} \\ %
10. – & Continuation: the genus of abelian varieties \dotfill & » & \hyperlink{old:361}{361} \\
11. – & Irrational pencil on the surfaces of geometric genus zero \dotfill & » & \hyperlink{old:365}{365} \\
12. – & Note on the surfaces of irregularity 1 \dotfill & » & \hyperlink{old:369}{369} \\
\end{longtable}

\begin{longtable}{@{}r@{\ }p{0.70\linewidth}@{\quad}c@{\ }r@{}}
\multicolumn{2}{@{}p{0.80\linewidth}}{\textsc{Chapter X.} – \textit{The surfaces of geometric genus zero} \dotfill } & » & \hyperlink{old:371}{371} \\
1. – & Introduction \dotfill & » & \hyperlink{old:371}{371} \\
2. – & The surfaces containing a system of curves of genus \(\pi\) and of degree \(n>2\pi-2\) \dotfill & » & \hyperlink{old:372}{372} \\
3. – & Lemma on the reducible curves of a pencil \dotfill & » & \hyperlink{old:376}{376} \\
4. – & Surfaces of genus \(p_g=0\) with \(p_a<-1\) \dotfill & » & \hyperlink{old:378}{378} \\
5. – & Surfaces \(F\) with \(p_g=0\) and \(p_a=-1\), possessing an elliptic pencil of curves \(K\) of genus \(\pi>1\): lemma I \dotfill & » & \hyperlink{old:379}{379} \\
6. – & Lemma II: outline of the proof \dotfill & » & \hyperlink{old:381}{381} \\
7. – & First case: the genus \(\pi\) of the curves \(K\) is even \dotfill & » & \hyperlink{old:384}{384} \\
\end{longtable}

% scan: s007R
% !PRINTED XV
\weberpage{XV}
\begin{longtable}{@{}r@{\ }p{0.70\linewidth}@{\quad}c@{\ }r@{}}
8. – & Second case: the genus \(\pi\) of the curves \(K\) is odd . . & Page & \hyperlink{old:389}{389} \\
9. – & Lemma III \dotfill & » & \hyperlink{old:392}{392} \\
10. – & Conclusion: the surfaces \(F\) also possess a linear pencil of elliptic curves \dotfill & » & \hyperlink{old:393}{393} \\
11. – & Surfaces of genera \(p_g=0\) and \(p_a=-1\) with an elliptic pencil of curves of genus \(\pi=1\) \dotfill & » & \hyperlink{old:394}{394} \\
12. – & Elliptic surfaces \dotfill & » & \hyperlink{old:398}{398} \\
13. – & Construction of the elliptic surfaces of genus \(p_g=0\) & » & \hyperlink{old:401}{401} \\
14. – & Surfaces with zero plurigenera: characterization of the ruled surfaces \dotfill & » & \hyperlink{old:405}{405} \\
15. – & Historical note \dotfill & » & \hyperlink{old:412}{412} \\
16. – & Surfaces with pluricanonical curves of order zero: types with elliptic curves \(K\) of general modulus \dotfill & » & \hyperlink{old:414}{414} \\
17. – & Continuation: harmonic case \dotfill & » & \hyperlink{old:417}{417} \\
18. – & Equianharmonic case \dotfill & » & \hyperlink{old:421}{421} \\
19. – & Summary \dotfill & » & \hyperlink{old:424}{424} \\
\end{longtable}

\begin{longtable}{@{}r@{\ }p{0.70\linewidth}@{\quad}c@{\ }r@{}}
\multicolumn{2}{@{}p{0.80\linewidth}}{\textsc{Chapter XI.} – \textit{General classification of surfaces} \dotfill } & » & \hyperlink{old:427}{427} \\
1. – & Introduction \dotfill & » & \hyperlink{old:427}{427} \\
2. – & Surfaces admitting a continuous series of birational transformations into themselves: cases leading to the ruled surfaces \dotfill & » & \hyperlink{old:429}{429} \\
3. – & Elliptic and hyperelliptic surfaces \dotfill & » & \hyperlink{old:433}{433} \\
4. – & Characterization of elliptic and hyperelliptic surfaces by means of the values of the genera \dotfill & » & \hyperlink{old:438}{438} \\
5. – & Surfaces with genera \(p_g=1\) and \(p^{(1)}=1\) \dotfill & » & \hyperlink{old:447}{447} \\
6. – & Note on the geometric theory of hyperelliptic surfaces \dotfill & » & \hyperlink{old:450}{450} \\
7. – & General classification of algebraic surfaces . . & » & \hyperlink{old:455}{455} \\
8. – & Surfaces of negative arithmetic genus \dotfill & » & \hyperlink{old:459}{459} \\
9. – & Summary of the refined classification \dotfill & » & \hyperlink{old:463}{463} \\
\end{longtable}
% !PRINTED
\weberblank{[s008L]}
% scan: s008L; pagina bianca, senza numero di pagina né testatina stampati.

\endgroup

\clearpage\mainmatter\thispagestyle{plain}
% scan: s008R
% !PRINTED 1
\weberpage{1}

\edfrontheading{INTRODUCTION}{INTRODUCTION}{intro}

\edsection{\textbf{1. Review of elementary notions.}}{1. Review of elementary notions.}{sec.intro.1}

In these lectures we refer, in general, to an algebraic surface of ordinary space, which we suppose irreducible. The surface \(f\) will be defined by an equation in Cartesian coordinates
\[
f(x,y,z)=0
\]
or in homogeneous coordinates
\[
f(x_1,x_2,x_3,x_4)=0:
\]
\(f\) denotes in the first case a polynomial in the variables \(x,y,z\), and in the second case a homogeneous polynomial, that is, a form.

We assume as known the elements of the projective theory of surfaces, which include the notions of order, class, polar surfaces, etc.\origfootnote{1}{s008R-1}{Cf. \textsc{Enriques-Chisini}, \textit{Teoria geometrica delle equazioni e delle funzioni algebriche.} Book III, vol. II. Bologna, Zanichelli, 1918.}. And we recall some properties which will be used in what follows.

A \textit{point} \(O\) of the surface \(f\), which we suppose of order \(n\), is \textit{simple} for it if the lines issuing from \(O\) cut the surface in \(n-1\) points outside \(O\). The neighbourhood of a simple point of the surface can be represented one-to-one on a plane, to an order of approximation as high as one wishes, that is, taking account of the infinitesimals of order \(1,2,\ldots\), etc. In the first place it is clear that, to the first order of approximation, the surface can be replaced by the tangent plane there; but if one wishes to take account of the second order, supposing the Cartesian axes suitably chosen, one will instead replace it by the osculating paraboloid of the same order \(z=\varphi_2(xy)\), which can be projected one-to-one onto the plane \((x,y)\) from the point at infinity of the \(z\) axis. And so on, successively, for the neighbourhood of the third order, etc. Therefore, if a curve is given on the surface having an arbitrary singularity at \(O\), one will be able to
\printersignature{\textsc{Enriques F.} - \textit{Algebraic surfaces.}}
% scan: s009L
% !PRINTED 2
\weberpage{2}

regard this singularity as formed of infinitely near multiple points, just as happens in the plane\origfootnote{1}{s009L-1}{Cf. \textsc{Enriques-Chisini}, Op. cit., Book IV, vol. II.}.

The singular or multiple points of the surface \(f\) are defined as follows: a \textit{multiple point \(O\) of order \(r\)} \((r>1)\) is one such that the lines through \(O\) cut the surface in \(n-r\) residual points. One has to distinguish \textit{multiple curves} and \textit{isolated multiple points}, without excluding the consideration of notable multiple points belonging to a multiple curve, which present themselves as a superposition of the two cases, and on which we shall have occasion to give some elucidation later on. A multiple curve of order \(r\) can in general be regarded as the intersection of \(r\) surfaces or \textit{sheets} differentially distinct at a generic point of the curve, corresponding to the \(r\) tangent planes there. An isolated multiple point, on the other hand, appears as a rule conical, the lines having \(r+1\) coincident intersections with the surface forming a cone of order \(r\), generally irreducible.

But there may be particular cases. Thus two or more sheets of the surface along a multiple curve may touch each other or even fuse into a \textit{sheet of higher order}: the simplest example is offered by the cuspidal double curve. On the other hand, the tangent cone at an isolated \(r\)-fold point may break up or even reduce to a plane counted \(r\) times: the simplest example is offered by the uniplanar double point.

To master all the complications to which the singularities of the surface may give rise, it is convenient to define the \textit{infinitely near multiple points} and the \textit{curves} (\textit{proper} or \textit{infinitesimal}) formed by them, as is done by means of the branches of curves issuing from a point\origfootnote{2}{s009L-2}{Cf. \textsc{Enriques-Chisini}, Op. cit. Book IV, chap. IV.}. In this order of ideas a uniplanar point appears as a double point to which three double points in distinct directions are infinitely near. The simplest example of an infinitesimal double curve (line), infinitely near to an isolated multiple point, is the \textit{tacnode}: that is, a double point \(O\) such that the plane sections through \(O\) have, next to \(O\), another infinitely near double point. But we do not wish to dwell on this subject, for which we refer to the cited lectures of \textsc{Enriques-Chisini}. %

\edsection{\textbf{2. Rational transformations.}}{2. Rational transformations.}{sec.intro.2}

A rational transformation of the space to which the surface \(f\) belongs is defined by taking the coordinates \(X,Y,Z\) to be rational functions of \(x,y,z\):
\weq{(1)}{X=\frac{\varphi_1(xyz)}{\varphi_0(xyz)}\qquad Y=\frac{\varphi_2(xyz)}{\varphi_0(xyz)}\qquad Z=\frac{\varphi_3(xyz)}{\varphi_0(xyz)},}
% scan: s009R
% !PRINTED 3
\weberpage{3}

where the \(\varphi\) are polynomials which we may suppose to be of the same order. By means of this transformation, which we suppose non-degenerate, to every generic point of the space \((xyz)\) there corresponds one point of the space \((XYZ)\), and conversely to a point \((XYZ)\) there corresponds in general a certain finite number \(n\) \((\geq 1)\) of points \((xyz)\), which are obtained by solving the equations (1), in which the \(XYZ\) are taken as given and the \(xyz\) as unknowns: for \(n>1\) the groups of points thus obtained will form, in the space \((xyz)\), an involution \(I_n\) of order \(n\), that is, a series of groups of \(n\) points such that every point belongs in general to one group.

Now the transformation (1) will in general carry the \(\infty^2\) points of the surface \(f(x,y,z)=0\) into \(\infty^2\) distinct points, constituting a surface \(F(X,Y,Z)=0\), the transform of \(f\). To a generic point of \(f\) there will thus correspond \textit{one} point of \(F\), and conversely to a generic point of \(F\) there will as a rule correspond \textit{one} point of \(f\), the transformation turning out to be uniquely invertible for the points of the surface \(F\)\origfootnote{1}{s009R-1}{Hence the distinction between a birational transformation holding between two surfaces and a birational or Cremona transformation, extendable to the space containing them.}. This invertibility, and the particular circumstances in which it ceases to hold, require an explanation. Indeed we have seen that the transformation (1), considered between the two spaces \((xyz)\) \((XYZ)\), is not in general uniquely invertible, but on the contrary leads from the points \((XYZ)\) to the groups of \(n\) points of an involution \(I_n\) in the space \((xyz)\). But if \(f\) is chosen in a general way within this space, it will not belong to the involution \(I_n\), so that to a point of \(f\) there will be conjugate, in \(I_n\), \(n-1\) points outside the surface; and hence to a generic point of \(F\) there will correspond, by (1), one point of \(f\) and \(n-1\) points outside it, which describe the surface conjugate to it with respect to \(I_n\).

Only in the case where it belongs to the involution \(I_n\), or to an involution partially contained in it, will the transform of \(F\) reduce to a multiple surface, in \([1,n]\) or \([1,m]\) \((m<n)\) correspondence with \(f\). %

Now, referring to the general case, in which \(f\) and \(F\) are in one-to-one correspondence, it should nevertheless be noted that the transformation (1) admits, in general, some exceptions\origfootnote{2}{s009R-2}{One should bear in mind the analogous considerations relating to the transformation of curves: Cf. \textsc{Enriques-Chisini}, Op. cit., Book V, vol. III.}. Indeed the formulae of the transformation become indeterminate at all those points \((xyz)\) at which the \(\varphi\) vanish together (that is, at the base points of the transforming system determined by the surfaces \(\varphi_0=0,\ \varphi_1=0,\ \varphi_2=0,\ \varphi_3=0\)). An exception of this kind can occur first of all
% scan: s010L
% !PRINTED 4
\weberpage{4}

at a simple point of \(f\); this point is then transformed in general into a line or a curve of \(F\), a curve which is the transform of a simple point, and which will be called an \textit{exceptional curve} of the surface \(F\).

On the other hand, it may happen that a base point of the transforming system of the \(\varphi\) falls at a singular point of \(f\). And if there is a multiple curve of \(f\) on which the \(\varphi\) vanish simultaneously, this will in general be transformed into a simple curve, to every \(r\)-fold point of the curve there corresponding \(r\) distinct points, in relation to the \(r\) sheets of \(f\) passing through it. If instead a base point of the \(\varphi\) falls at an isolated multiple point \(O\) of \(f\), this point is in general transformed into a curve whose points correspond to the points infinitely near to \(O\).

These assertions break down when there occur multiple points and curves infinitely near to the singular points considered. In these cases the general effect of the transformation is to transform the infinitely near multiple points of the first order into proper multiple points, and in general to simplify the singularity, reducing from \(r\) to \(r-1\) the neighbourhood of the points considered, at least when it is supposed that the surfaces \(\varphi\) passing through a point \(O\) have no fixed tangent in common.

However, the transformation which succeeds in simplifying the singularities of a surface in general introduces new singularities, because there will in general be \(\infty^1\) pairs of points of \(f\) at which \(X,Y,Z\) take the same value again, so that to them there will correspond points of a double curve of \(F\).

This follows, in substance, from an easy count of constants. If \(X,Y,Z\) are to take the same value at the points \((x,y,z)\) and \((x',y',z')\), the following 5 equations in 6 unknowns must hold:
\[
\begin{aligned}
&\varphi_1(xyz)\varphi_0(x'y'z')-\varphi_1(x'y'z')\varphi_0(xyz)=0\\
&\varphi_2(xyz)\varphi_0(x'y'z')-\varphi_2(x'y'z')\varphi_0(xyz)=0\\
&\varphi_3(xyz)\varphi_0(x'y'z')-\varphi_3(x'y'z')\varphi_0(xyz)=0\\
&f(xyz)=0\\
&f(x'y'z')=0,
\end{aligned}
\]
which will in general admit \(\infty^1\) solutions \((xyzx'y'z')\), apart from the \(\infty^2\) trivial solutions
\[
x=x'\qquad y=y'\qquad z=z'\text{\origfootmark{1}{s010L-1}}.
\]\origfoottext{1}{s010L-1}{The first three equations define three hypersurfaces of the space \(S_6\), which have in common the \(S_3\) \((x=x',y=y',z=z')\) and, outside it, a certain three-dimensional variety \(V\); \(V\) is intersected by the cylindrical varieties \(f(xyz)=0\) and \(f(x'y'z')=0\) in a curve.}

In particular cases there may be \(\infty^1\) groups of more than two points of \(f\), to which there correspond the points of a multiple curve of \(F\). But
% scan: s010R
% !PRINTED 5
\weberpage{5}

there may also be curves of \(f\) on which the \(XYZ\) remain constant, to which there correspond isolated multiple points of \(F\). And as the singularities of \(f\) grow more complicated, and with the intervention of other particular circumstances, the genesis of higher singularities for \(F\) cannot be excluded.

\edsection{\textbf{3. Surfaces in hyperspaces: resolution of singularities.}}{3. Surfaces in hyperspaces: resolution of singularities.}{sec.intro.3}

The transformation of ordinary space which we considered above can be extended by defining a transformation which carries the surface \(f\) into a space \(S'_r\) of more dimensions: %
\[
X_1=\frac{\varphi_1(xyz)}{\varphi_0(xyz)}\qquad X_2=\frac{\varphi_2(xyz)}{\varphi_0(xyz)}\ldots\qquad X_r=\frac{\varphi_r(xyz)}{\varphi_0(xyz)};
\]
the surface \(f(xyz)=0\) will be transformed into a surface \(F\) of the space \(S_r\), as defined by the formulae written above, which on the other hand can also be defined in various ways as a partial intersection of varieties of hyperspace\origfootnote{1}{s010R-1}{Cf. \textsc{Enriques-Chisini}, Op., cit., Book I, \S\ 23.}. %

For this transformation too, much of what was said for a transformation of the space \(S_3\) is repeated; namely, that in general the transformation is one-to-one for the points of the two surfaces (and not for the spaces containing them), and that an ordinary multiple curve of \(f\) which is a base curve for the transforming system of the surfaces \(\varphi\) is in general resolved, becoming the locus of the \(r\) simple points which correspond to the \(r\)-fold points of \(f\) constituting it; and similarly for the isolated multiple points, which change into curves, etc.

But here, at least for \(r\geq5\), no new singular points are in general produced for \(F\), since precisely for \(r\geq5\) there exist in general no pairs of points of \(f\) at which \(X_1,X_2,\ldots,X_r\) take the same value again, while for \(r=4\) there is in general a finite number of such pairs, giving rise to as many double points of \(F\).

One thus perceives intuitively that it will always be possible to resolve the singularities of a surface \(f\) and to transform it into another of the space \(S_r\) entirely free of singularities, that is, endowed only with simple points. But this reduction, which is easy for surfaces endowed only with ordinary multiple curves and points (multiple curves with distinct sheets, and isolated multiple points with a quite general osculating cone), requires a more minute examination when \(f\) possesses infinitely near singularities, however complicated. The analysis needed for this purpose can, however, be decomposed into two parts: an analysis in the sense of differential geometry, where it is a matter of resolving the singularities defined in the neighbourhood of a point \(O\) of a curve of ordinary space, and in which one may make use of birational or Cremona transformations
% scan: s011L
% !PRINTED 6
\weberpage{6}
of the whole space $S_3$; and subsequently a critical examination of the new singularities which may be produced by the transformation itself. But this last step, which involves minute difficulties when it is a matter of reducing the singularities by birational transformations of $S_3$\origfootnote{1}{s011L-1}{Problem solved by \textsc{O. Chisini}, in Memorie Acc. Bologna, 1921.}, can be avoided thanks to a simple observation of \textsc{B. Levi}: indeed, if the surface $f(xyz) = 0$ is transformed into another $F(XYZ) = 0$, the surface of the 6-dimensional space which is the locus of the points $(x, y, z, X, Y, Z)$ will possess only singularities corresponding to the singularities common to $f$ and $F$.

It is not our purpose to stop here to explain the reduction of the singularities of surfaces as it was first rigorously obtained by \textsc{B. Levi}, and we shall confine ourselves to stating the fundamental result that is reached, referring for the proof to the principal memoirs which contain it and develop it in different ways\origfootnote{2}{s011L-2}{\textsc{B. Levi}, \textit{Risoluzione delle singolarità puntuali delle superficie algebriche.} Atti Acc. Torino, 1897 (cf. Ann. di Mat. 1897).

\textsc{O. Chisini}, \textit{La risoluzione delle singolarità delle superficie mediante trasformazioni birazionali dello spazio.} Memorie Acc. Bologna, 1921.

\textsc{G. Albanese}, \textit{Trasformazione birazionale di una superficie algebrica qualunque in un’altra priva di punti multipli.} Circ. Mat. Palermo, 1924.

\textsc{R. J. Walker}, \textit{Reduction of singularities of an algebraic surface.} Annals of Math., 1935.

\textsc{O. Zariski}, \textit{The reduction of the singularities of an algebraic surface.} Ann. of Math., 1939. Cf. Annals of Math., 1941.}.

\textit{A surface $f$ endowed with arbitrary singularities can always be transformed birationally into another one free of singularities, belonging to a space of five dimensions.} %

\edsection{4. Surfaces with normal singularities.}{4. Surfaces with normal singularities.}{sec.intro.4}

Given a surface with arbitrary singularities, let us begin by transforming it into an $F$ of the space $S_5$ entirely free of multiple points. The surface $F$, projected successively from two points onto ordinary space, will give us in the latter an image $\Phi$ of the given surface. We have to examine what the singularities of $\Phi$ will be when the projection of $F$ is made from two generic points. We wish to prove that $\Phi$ will possess only a nodal double curve (tangent planes generally distinct), endowed with triple points which are also triple for the surface.

It is clear first of all that this is the type of \textit{normal singularities} which a surface of ordinary space will in general have when it is obtained as the projection of a surface of a space of more dimensions free of singular points; this follows from a simple count of con%
% scan: s011R
% !PRINTED 7
\unskip\weberpage{7}%
\unskip stants. Observe indeed that the chords of a surface of the space $S_r$ are always $\infty^4$, and hence for $r = 5$ there is a finite number of chords passing through a point, while for $r = 4$ there are $\infty^1$ of them. Hence, projecting $F$ from a generic point of $S_5$, the projected surface $F'$ will acquire, in correspondence with the chords of $F$ passing through the centre, a finite number of double points; then, projecting $F'$ from a generic point of $S_4$, the projected surface will acquire an (apparent) double curve, on which the double points of $F'$ will come to lie. In other words, through two points of $S_5$, or through the line joining them, there will pass an $\infty^1$ series of planes which meet $F$ in two points and a finite number of planes which meet it in three points, so that the projection of $F$ made from those centres will indeed possess a nodal double curve endowed with triple points, triple at once for the double curve and for the surface.

However, one may suspect that, for a particular surface of $S_5$, even one free of multiple points, the projection made from two generic points of $S_5$ might come to acquire higher singularities, namely: a multiple curve instead of a double one, or multiple points that are more than triple for the double curve. To dispel the first doubt it suffices to note that the $\infty^4$ chords of a surface of any space of more than three dimensions cannot be trisecants, since otherwise the chords of a twisted curve which is a hyperplane section of the surface would also be trisecants, which is impossible\origfootnote{1}{s011R-1}{Cf. \textsc{Enriques-Chisini}, Op. cit., Book III, chap. IV, \S~43, vol. II, p. 289.}. From this it follows that no trisecants of $F$ can pass through a generic point of $S_5$, and that when $F$ is projected from a generic point into an $F'$ of $S_4$, only a finite number of trisecants can pass through a generic point of the latter.

It remains to prove that, if the centres of projection are chosen generically in $S_5$, it cannot happen that the trisecant planes through the said centres meet the surface $F$ in more than three points, and hence that the double curve of the projection $\Phi$ possesses, instead of triple points, points of higher multiplicity\origfootnote{2}{s011R-2}{Cf. \textsc{F. Enriques}, \textit{Sulle singolarità che nascono per proiezione di una superficie o varietà algebrica}, in Scritti matematici offerti a \textsc{Luigi Berzolari}, Pavia, 1936, p. 351.}. For this it suffices to exclude that all the planes which are required to meet the surface in three points should consequently come to meet it in four or more points. Indeed, if every plane trisecant to $F$ turns out to be quadrisecant, the same will happen for every twisted curve which is a hyperplane section of $F$ (and which cannot belong to an $S_3$), and hence the projection of the latter from one of its points onto $S_3$ will also have to give a curve every chord of which is trisecant, which we have already recalled to be impossible.

% scan: s012L
% !PRINTED 8
\weberpage{8}
In this way a fundamental theorem already recognized by \textsc{Noether} (1888) is rigorously proved:

\textit{A surface with arbitrary singularities can always be transformed into another one of ordinary space endowed only with normal singularities, that is, with a nodal double curve which in general possesses a finite number of triple points, which are also triple for the surface.} %

The said double curve may also possess a finite number of double points (\textit{crossings}) arising from those chords of the surface of $S_4$ projected onto $S_3$ which are double for the cone of the projecting chords.

It is also worth recalling that on the said \textit{nodal curve} there will in general be a finite number of \textit{cuspidal points} (pinch-points), at which the two linear sheets of the surface fuse into a sheet of the second order, and which turn out to be simple base points for the variable curves of intersection of the polar surfaces.

In what follows, since we have to develop a theory of surfaces in which one seeks the properties invariant under birational transformations, we shall always be able to refer to a model surface free of singularities in the space of 5 dimensions, or to a surface of ordinary space endowed with the normal singularities of the preceding statement. In the latter case a biplanar point of the double curve must be regarded as the superposition of two simple points, whose neighbourhoods correspond to the two sheets, just as in the invariant theory of curves a node is considered as the superposition of two points belonging to the two branches.

\edsection{5. Note.}{5. Note.}{sec.intro.5}

What has been said concerning the transformation of a surface into another of ordinary space endowed with normal singularities can be made more precise in the sense that: it is legitimate to suppose that the \textit{transform of any surface whatever has an irreducible double curve}, which possesses, as has been said, a certain number of triple points, triple alike for the surface and for its nodal curve.

We shall expound the proof of this property\origfootnote{1}{s012L-1}{Cf. \textsc{A. Franchetta}, \textit{Sulla curva doppia della proiezione di una superficie generale dello $S_4$ da un punto generico su uno $S_3$.} Rendic. R. Acc. d’Italia, 1941, see also the Rend. dell’Acc. dei Lincei, s. VIII, vol. II, fasc. 3, 1947.} assuming as known the principles of the theory of linear systems, which are expounded in the first sections of these lectures.

For this purpose it is convenient to start from a surface free of singularities in the space $S_5$, or better from an $F$ which is a generic projection of it into $S_4$, and therefore endowed only with improper double points
% scan: s012R
% !PRINTED 9
\weberpage{9}
that are \textit{ordinary}, each of which can be regarded as an apparent singularity, representing a pair of \textit{distinct} superposed points. We have to examine the singularities of a surface $\Phi$ of ordinary space, the projection of $F$ from a generic point $O$ of $S_4$: in which cases it can happen that the double curve of $\Phi$ is always reducible as the centre of projection $O$ varies.

This examination leads to establishing that: if a surface $F$, endowed with ordinary improper double points in $S_4$, is projected from a generic point into a $\Phi$ of $S_3$ whose double curve is reducible, then $F$ is a projection of the Veronese surface (the surface of the fourth order of $S_5$ represented on the plane by the $\infty^5$ system of conics).

To reach this result it is convenient to state first the following:

\textit{Lemma}: there exist no surfaces of $S_4$ possessing $\infty^3$ pairs of tangent planes meeting in a line, in such a way that the lines joining their points of contact fill the whole space.

Let $P$ be a generic point of $F$ and $\pi$ the tangent plane there. By hypothesis there exist $\infty^1$ planes tangent to $F$ and meeting $\pi$ in lines; their points of contact form a curve, possibly reducible, of which we consider a component $C$. Let $Q$ be a point of it, and $Q'$ the point infinitely near to $Q$ on $C$. The planes tangent to $F$ at $Q$ and $Q'$ cut $\pi$ in lines; hence each of them lies in an $S_3$ with $\pi$. But the tangent plane at $Q'$ contains $Q$; hence the said planes determine with $\pi$ a single $S_3$, which is tangent to $F$ at all the points of $C$. As $P$ varies, the curve $C$ describes a system (at least) $\infty^1$; and if, as has been supposed, the surface $F$ does not belong to an $S_3$, each of these curves $C$, having to lie in the $S_3$ determined by it and in the $S_3$ determined-by the infinitely near curve, is contained in a plane (it may reduce, in particular, to a line); the planes of two infinitely near curves $C$ lie in an $S_3$. One concludes that a surface for which the assumed hypothesis holds is ruled, or contains $\infty^1$ plane curves whose planes form a developable. In the first case the lines joining the points of contact of two planes meeting in lines are the generators of the ruled surface; in the second they are the lines of the planes of the developable; in neither case do they fill $S_4$. %

This being so, let there be given in $S_4$ a surface $F$, endowed only with a finite number of ordinary improper double points, whose projection $\Phi$ from a generic point $O$ possesses a reducible double curve; we say first of all that this curve must be connected, that is, formed of two or more parts having some point in common. Indeed the apparent double curves
\[
C = \Sigma K
\]
which correspond to the sets of these parts $K$ form a system

% scan: s013L
% !PRINTED 10
\weberpage{10}
rational (or at least unirational) whose elements correspond to the points $O$ of the space $S_4$, and therefore the $C$ belong to a linear system of curves on $F$\origfootnote{1}{s013L-1}{Theorem established for rational series of groups of points on a curve and extended to varieties by \textsc{Enriques}. Cf. \textsc{Enriques-Chisini}, op. cit., Book V, chap. I, \S~9, vol. III, p. 78 (cf. p. 485).}. If this system is irreducible, the reducible $C$ that belong to it must be connected\origfootnote{2}{s013L-2}{Principle of degeneration. Cf. \textsc{Enriques-Chisini}, Op. cit., Book V, chap. III, \S~36, vol. III, p. 405.}; if instead all the $C$ of the said linear system are reducible, they are composed of the curves $K$ of a pencil, which have no variable intersections with one another.

But this case cannot occur, because two curves $K$ containing two pairs of points of $F$ such as $AA_1$ and $AA_2$ certainly have a common point (at $A$); and note that they cannot coincide in one and the same curve, since, as $A_2$ varies, this curve would come to contain all the points of the surface.

We are therefore led to examine whether and how the surface $\Phi$, projection of $F$ from a generic point $O$, can have a double curve composed of two or more parts connected with one another.

A point of connection of such parts may be:

$1^\circ)$ a simple crossing, that is, a double point common to two simple branches, which will therefore be a tacnode for the surface, or

$2^\circ)$ a triple point, through which three branches of the double curve pass, two of which may a priori also belong to one and the same irreducible component.

$1^\circ)$ A point $\bar{P}$ common to two simple branches of the double curve $\bar{C}$ of $\Phi$ will in general be the projection of two points $P$ and $P'$ collinear with the centre $O$ in $S_4$, and one sees at once that the tangent planes to $F$ at $P$ and $P'$ belong to the same hyperplane, the one having as its trace the tangent plane to $\Phi$ at $\bar{P}$, and that they therefore meet in a line. But since the double point $\bar{P}$ of $\bar{C}$ is supposed to exist for every position of the centre $O$, the surface $F$ would have to possess $\infty^3$ pairs of incident tangent planes, and the lines joining the points of contact of these pairs of planes would have to fill the whole space; this is absurd by the lemma stated beforehand.

The case in which the points $P$ and $P'$ are infinitely near, thus falling at the points which a tangent through $O$ has in common with the surface $F$, requires a special examination. In this case the point $\bar{P}$, the crossing of the two branches of the double curve of $\Phi$, will still be a tacnode for $\Phi$, so that every plane through $\bar{P}$ will cut $\Phi$ in a curve having also a double point infinitely near to $\bar{P}$. It follows from this that
% scan: s013R
% !PRINTED 11
\weberpage{11}
the hyperplane projecting from $O$ the tacnodal plane of $\Phi$ must cut $F$ in a curve which has four intersections with every plane through the tangent $P\,P'$ and therefore possesses $P$ and $P'$ as double points; save for examining the particular cases in which the effective multiplicity of $P$ becomes greater than the virtual multiplicity 2, while the multiplicity of $P'$ is lowered instead. The hypothesis to which our discussion refers implies that, as $O$ moves in the space $S_4$, there is always through it a tangent, say $p$, meeting $F$ in two infinitely near points, which would be double for the section by the tangent hyperplane at $P$; hence each of the $\infty^3$ tangent hyperplanes to $F$ would have to touch the surface at two infinitely near points $P$ and $P'$. But it is easy to see that this property cannot pertain to a surface $F$ of $S_4$; indeed a generic projection of it onto $S_3$ would be such that each of its tangent planes would be tangent to it also at an infinitely near point, and it would therefore turn out to be a developable; then the surface of $S_4$ of which it is the projection would also have to be developable, and yet a developable of $S_4$ possesses only $\infty^1$, and not $\infty^3$, hyperplanes which touch it at two infinitely near points.

The conclusion of this discussion is to prove impossible the hypothesis we have made, namely that the projection of our surface $F$ from a generic point $O$ of $S_4$ is a $\Phi$ whose double curve consists of two components crossing at a point $\bar{P}$, image of two infinitely near points $P$ and $P'$ of $F$. However, it is necessary to examine the particular case alluded to above: that instead of a tangent hyperplane at $P$ whose intersection with $F$ has two double points $P$ and $P'$, one meets with a tangent hyperplane whose section has a triple point at $P$ and a simple point at $P'$, or else only a quadruple point at $P$. Of the two cases which present themselves here, the first gives rise to no exception in the preceding argument, because there must still be $\infty^3$ tangent hyperplanes to $F$ which touch it. also (virtually) at an infinitely near point. Instead, it is necessary to dwell a moment on the second hypothesis, in which one would a priori have no longer $\infty^3$, but $\infty^2$ hyperplanes cutting $F$ in curves endowed with a quadruple point. In fact this hypothesis cannot occur, because a surface of $S_4$ such that through every point of that space there is a tangent hyperplane cutting it in a curve endowed with a quadruple point, such a surface — we say — cannot belong to $S_4$, but lies entirely in a space $S_3$. For its hyperplane section curves are such that the osculating plane at any point has four-point contact with the curve itself, and therefore they are plane curves. %

$2^\circ)$ A point $\bar{P}$ which is triple for the double curve $C$ of $\Phi$ and a point of connection of irreducible components of it may be: %

% scan: s014L
% !PRINTED 12
\weberpage{12}
$a)$ a simple point common to three irreducible components of $\bar{C}$; or

$b)$ a simple point for one irreducible component of $\bar{C}$, and double for another.

$a)$ Under hypothesis $a)$ the point $\bar{P}$ will be the projection of three points of $F$, say $P_1$, $P_2$, $P_3$; and the components $\bar{K}_1$, $\bar{K}_2$, $\bar{K}_3$ of $C$ for which $\bar{P}$ is a triple point will be projections of three curves, respectively $K_1$, $K_2$, $K_3$, for which one may assume that the first passes through the points $P_2$ and $P_3$; the second through $P_1$ and $P_3$, and the third through $P_1$ and $P_2$. %

Let us recall that the surface $\Phi$, projection of $F$, must have a singularity similar to $\bar{P}$, however the centre of projection $O$ may vary in $S_4$, and in particular if it is made to move on the line $P_1P_2P_3$. Now, corresponding to this variation of $O$, it may happen that the apparent double curves $K_1K_2K_3$ on $F$ remain fixed, or vary with $O$. If one of them, e.g. $K_1$, remains fixed, one deduces that it is a conic: indeed, for the curve $\bar{K}_1$ always to be projected into the same curve, every line joining a point of $\bar{K}_1$ with a point of the line $P_2P_3$ must meet $K_1$ in another point, and hence $K_1$ turns out to belong to a plane through the line $P_2P_3$; a plane curve which, like this one, projects into a double curve (and not one of higher multiplicity) must be a conic, and therefore the surface $F$, possessing $\infty^2$ conics, will have to be a projection of a Veronese surface. %

We must now examine the hypothesis that the curve $K_1$ varies as $O$ varies on the line $P_2P_3$. In that case, for a suitable position of $O$, $K_1$ will come to pass also through the point $P_1$ of that line, and correspondingly one will have a surface $\Phi$ endowed with a particular triple point $\bar{P}$, at which two sheets, e.g. those corresponding to the neighbourhoods of $P_1$ and $P_2$, turn out to be tangent to each other. But this implies that the tangent planes to $F$ at $P_1$ and $P_2$ meet in a line. Under our hypotheses, therefore, $F$ will come to have $\infty^3$ pairs of mutually incident tangent planes, in such a way that the lines joining the points of contact of these pairs of planes fill the space $S_4$; and this is excluded by the lemma.

$b)$ Let us now examine the hypothesis in which the triple point $\bar{P}$ of $\Phi$ is simple for one component $\bar{K}_1$ and double for a component $\bar{K}_2$ of the double curve. $\bar{K}_1$ and $\bar{K}_2$ will be projections of two apparent double curves $K_1$ and $K_2$ of $F$, of which the first may be supposed to pass simply through $P_2$ and $P_3$, while the second will pass simply through $P_2$ and $P_3$ and doubly through $P_1$. Now, if the centre of projection $O$ is made to vary on the line $P_1P_2P_3$, it may happen that the curve $K_1$ remains fixed, or that it varies with $O$. If it remains fixed, one concludes as before that it is a plane curve, and hence a conic, whence it follows that $\Phi$ is the projection of a Veronese surface. If instead $K_1$ varies with %
% scan: s014R
% !PRINTED 13
\weberpage{13}
$O$, then for some position of the latter it comes to pass also through $P_1$: the singularity which $\Phi$ presents at $\bar{P}$ now turns out to be a triple point at which at least two sheets touch, and therefore the tangent planes to $F$ at two points such as $P_1$ and $P_2$ will turn out to meet in a line: $F$ possesses $\infty^3$ pairs of incident tangent planes, and we fall back into the case already excluded.

From all the foregoing it follows that \textit{the projection into the space $S_3$ of a surface of $S_5$ free from singularities possesses in general an irreducible double curve; the only exception is the case of a Veronese surface, which in fact projects into the Steiner surface, endowed with three double lines passing through a triple point.}
% !PRINTED
\weberblank{[s015L]}
% scan: s015L; pagina bianca, senza numero di pagina stampato.

\clearpage\thispagestyle{plain}
% !PRINTED
\weberpage{[s015R]}
% scan: s015R; nessun numero di pagina stampato.
\edchapter{I}{LINEAR SYSTEMS OF CURVES}{Chapter I. LINEAR SYSTEMS OF CURVES}{LINEAR SYSTEMS OF CURVES}

\edsection{1. Linear pencils.}{1. Linear pencils.}{sec.I.1}

Let $F$ be a surface which, for simplicity of exposition, we shall suppose to belong to ordinary space and to be endowed with the normal singularities defined in the preceding introduction.

Let us consider a rational function $t$ of the variable point on $F$ which does not reduce to a constant:
\weq{(1)}{t = \frac{\varphi(xyz)}{\psi(xyz)}.}

For every determinate value $t = t_0$ there will exist points of $F$ at which $t$ takes this value: their locus $C$ will be called a \textit{level curve} of the rational function $t$ on $F$. As $t$ varies, one will have on $F$ a simple infinity of level curves, which is said to constitute a \textit{linear pencil}. Among the $C$ of the pencil there figure the \textit{curve of zeros}, that is, the curve $\varphi = 0$ which corresponds to $t = 0$, and the \textit{curve of poles}, that is, the curve $\psi = 0$ for which $t = \infty$.

In case these two curves have a part $K$ in common, $K$, on which $t$ takes an indeterminate form, may be regarded as a fixed component of the $C$, or else be removed from all the $C$.

The definitions given above are perfectly analogous to those which refer to curves and to the series $g_n^1$ on them, and give rise to similar observations. The pencil of the $C$ defined by (1) is equally given by the rational function
\[
\tau = \frac{\alpha t + \beta}{\gamma t + \delta};
\]
the rational substitution has the sole effect of changing the value of $t$ which corresponds to each single curve $C$; hence all the $C$ appear on an equal footing in the pencil, and thus neither the curve of zeros nor that of poles has any particular significance for it. Given the
% scan: s016L
% !PRINTED 16
\weberpage{16}
pencil, a rational function corresponding to it is determined by the arbitrary choice of the curves $\tau=0$ and $\tau=\infty$, together with that of the unit curve $\tau=1$, this choice determining the preceding substitution. If for two functions $t$ and $\tau$ the curve of zeros and the curve of poles coincide, they differ by a multiplicative constant.

Let us note moreover that the rational function (1) is defined only with respect \textit{to the modulus $F$}. Thus two different functions $\frac{\varphi}{\psi}$ and $\frac{\varphi_1}{\psi_1}$ must be regarded as identical on $F$ when for all the points of $F$ one has:
\[
\frac{\varphi}{\psi}=\frac{\varphi_1}{\psi_1}; %
\]
that is,
\weq{(2)}{\varphi\psi_1-\psi\varphi_1\equiv 0 \qquad (\mathrm{mod.}\ F)}
or
\[
\varphi\psi_1-\psi\varphi_1=AF,
\]
where $A$ represents a polynomial.

The curves $C$ of a linear pencil all have the same order; if for one of them the order is apparently reduced, this means that the curve has split and contains a part at infinity. To eliminate the apparent reduction, when one does not wish to resort to a projective transformation of the surface, it suffices to make use of homogeneous coordinates $x_1, x_2, x_3, x_4$. Thus, for example, in (1), if the polynomial $\varphi$ was of lower degree than $\psi$, then on passing to homogeneous coordinates $\varphi$ is multiplied by a suitable power of $x_4$, and it thus becomes manifest that the curve of zeros comes to contain also the section of $F$ by the plane at $\infty$ ($x_4=0$), counted a certain number of times.

The definition of a linear pencil of curves $C$ has a simple projective meaning: indeed the equation
\weq{(3)}{\varphi(xyz)-t\psi(xyz)=0}
represents a pencil of surfaces which cut $F$ along the curves $C$; and (2) says that \textit{the same pencil of curves can be cut out on $F$ by different pencils of surfaces.}

The pencil of surfaces (3) will have a base curve $\Gamma$ which in general will not belong to $F$ and will meet $F$ in a finite number of points. These \textit{base points of the linear pencil} of curves $C$ are points of essential indeterminacy for the rational function (1).

Indeed the function, which at a base point takes the form $\frac{0}{0}$, possesses there a non-removable singularity: if in fact one tries to define its value there by considerations of limits, the value so defined by continuity will be different when one approaches the base point along different curves $C$.

% scan: s016R
% !PRINTED 17
\weberpage{17}
A base point for the curves $C$ of a linear pencil may have a certain multiplicity $i>1$ for the $C$, and in that case it will be called an \textit{$i$-fold base point of the pencil.}

Finally it may happen that the base curve $\Gamma$ of the pencil (3), or a part $K$ of it if it is reducible, lies on $F$; then (as has already been said) $K$ may be considered as a fixed component of the curves $C$ of the (\textit{reducible}) pencil, but one may also disregard $K$ and regard the pencil of the $C$ as consisting solely of the variable intersections of $F$ with the surfaces (3).

\edsection{2. Irrational pencils.}{2. Irrational pencils.}{sec.I.2}

A linear pencil of curves $C$ on $F$ is characterized geometrically by the following properties:

1) it is an (algebraic) series $\infty^1$ of \textit{index} 1, that is, such that through every generic point of $F$ there passes one $C$;

2) this series is \textit{rational}, that is, it can be put in algebraic one-to-one correspondence with a range of points on a line or with the series of values of a parameter $t$.

Indeed the $t$ so defined turns out to be a single-valued, and therefore rational, function of the points of $F$, and the curves $C$ are level curves for it.

If one calls a \textit{pencil} a series of curves $C$ which is of index 1, that is, which enjoys property 1), one can therefore say that \textit{every rational pencil is linear.}

But on a surface there may also be \textit{irrational pencils}, whose curves constitute the elements of a non-rational algebraic entity having a certain genus $p>0$. Thus, for example, the generators of a ruled surface (whose plane sections are) of genus $p$ form precisely an irrational pencil of lines.

One can construct an example of a surface having an irrational pencil of curves of genus $p$ by intersecting with an arbitrary variety of $S_4$ a cone of the second kind formed by the planes of that same $S_4$ which project from a line the points of a curve of genus $p$.

It is worth noting that a \textit{pencil of curves belonging to the surface $F$ is certainly rational if it possesses a base point at a point simple for the surface.}

Indeed, if the base point $O$ is simple for the curves $C$ of a pencil, the series of the $C$ is in one-to-one correspondence with the pencil of tangents at $O$. And if instead the point is $i$-fold for the $C$ ($i>1$), then the $C$ correspond in general to the groups of an involution of order $i$ in the pencil of the aforesaid tangents, and this involution is rational (Lüroth’s theorem)\origfootnote{1}{s016R-1}{Cf. \textsc{Enriques-Chisini}, Op. cit., Book II, chap. I, \S\ 3, vol I, p. 167.}.

\printersignature{\textsc{Enriques F.} - \textit{Superficie algebriche.}\quad 2}

% scan: s018L
% !PRINTED 18
\weberpage{18}
However, this argument may fail if the curves $C$ have fixed tangents at $O$, that is, if they possess as base points not only $O$ but also other points infinitely near to it. To eliminate the exception it suffices to observe that, proceeding along a branch of a $C$ up to a suitably high order $r$, one will find a first point $O_r$, belonging to the $r$-th neighbourhood of $O$, capable of describing freely the neighbourhood of the base point $O_{r-1}$, which, like the neighbourhood of $O$, constitutes a rational entity: if $O_{r-1}$ is a simple base point, the curves $C$ correspond one-to-one to the points $O_r$ of the said neighbourhood; if $O_{r-1}$ is $i$-fold and is followed by several free points, $O_r$, the $C$ correspond to the groups of an involution described by these groups of points.

\edsection{3. Linear systems.}{3. Linear systems.}{sec.I.3}

The definition of the linear pencil of curves on $F$ can be generalized by considering \textit{linear systems of dimension $r>1$}.

A linear system $\infty^r$ of curves $C$ is cut out on $F$ by a linear system of surfaces (or of hypersurfaces)
\weq{(1)}{\lambda_0\varphi_0+\lambda_1\varphi_1+\ldots+\lambda_R\varphi_R=0;}
the \textit{dimension} $r$ turns out to be less than $R$ when in the system (1) there are $\infty^h$ surfaces which contain $F$ as a part, and then $r=R-h-1$.

For multiply infinite linear systems $|C|$, as for the pencil, one defines the curves which are \textit{fixed components} of the $C$ and the \textit{base points} of $|C|$; a base point is a point common to all the $C$, where the $\varphi$ vanish simultaneously. %

A linear system $|C|$ can be defined intrinsically, independently of the projective character of the surface, just as was done for the linear pencil. To this end consider $r$ rational functions of the variable point of $F$ which all have the same polar curve and which we shall suppose represented by
\[
t_1=\frac{\varphi_1}{\varphi_0}\qquad t_2=\frac{\varphi_2}{\varphi_0}\qquad \ldots \qquad t_r=\frac{\varphi_r}{\varphi_0},
\]
assuming, for simplicity of exposition, that they are linearly independent on $F$, so that it is \textit{not possible} to find $r+1$ constants (not all zero) $a_0, a_1, \ldots, a_r$ for which one has identically
\[
a_0\varphi_0+a_1\varphi_1+\ldots+a_r\varphi_r\equiv 0 \qquad (\mathrm{mod.}\ F).
\]

Then the level curve of the function
\[
t_0=\lambda_1t_1+\lambda_2t_2+\ldots+\lambda_rt_r
\]
describes the linear system $|C|$ of dimension $r$.

In this the polar curve $\varphi_0$ seems to occupy a special position; but one can choose as such any curve whatever of the
% scan: s018R
% !PRINTED 19
\weberpage{19}
system, since to one and the same system $|C|$ there correspond infinitely many linear systems of functions, which are deduced from one another by linear substitutions.

The dimension $r$ of the linear system $|C|$ is defined geometrically by the following property: \textit{through $r$ generic points of the surface there passes one and only one curve of the system.}

A linear system $C$ will be called \textit{irreducible} if its generic curve $C$ is irreducible. For an irreducible linear system $|C|$ the genus of its generic curve will be called the \textit{effective genus}, and the number of variable intersections of two generic $C$ the \textit{effective degree}. For $r=1$ the degree is zero, since every point common to two curves of the pencil is a base point belonging to all the $C$. %

Between the dimension $r$ and the degree $n$ of a linear system $|C|$ there holds the inequality $r\leq n+1$; indeed the series cut out on any $C$ of $|C|$ by the remaining curves of the system, outside the base points, is a $g_n^{r-1}$, for which one has $r-1\leq n$.

The linear series $g_n^{r-1}$, defined on a generic curve $C$, is of fundamental importance in the study of the linear system and is called the \textit{characteristic series} of $|C|$\origfootnote{1}{s018R-1}{The characteristic series of a linear system was introduced by \textsc{C. Segre} and \textsc{G. Castelnuovo} around 1890. The name \textit{characteristic series} was used for the first time by \textsc{G. Castelnuovo} in: “\textit{Ricerche generali sopra i sistemi lineari di curve piane}”. Memorie dell’Acc. di Torino, ser. II, vol. XLII, 1891, reprinted on p. 137 ff. of the “Memorie scelte”, Zanichelli, Bologna, 1937.}.

The effective degree of a linear system of curves is also defined for a system of reducible curves which is free from fixed parts. Now one can prove that the degree is $n>0$ unless the system itself is a linear pencil; or unless its curves are composed of those of a pencil (rational or not).

Indeed, if $n=0$, all the curves $C$ of the system which pass through a generic point must coincide in a single one (dimension $r=1$); or else they must have a curve $K$ in common, and the set of the $K$ will constitute a pencil, rational or not, of which the $C$ are composed.

\edsection{4. Characteristic property of linear systems.}{4. Characteristic property of linear systems.}{sec.I.4}

We have noted that, $|C|$ being a linear system of dimension $r$ on the surface $F$, through $r$ generic points of $F$ there passes one and only one curve $C$. This property is characteristic for $r>1$\origfootnote{2}{s018R-2}{\textsc{F. Enriques}, \textit{Una questione sulla linearità ecc.} Atti Acc. Lincei, 1893. Cf. \textsc{C. Segre}, \textit{Introduzione alla geometria sopra un ente algebrico semplicemente infinito.} Ann. di Mat., ser. II, vol. 22 (1894).},
% scan: s020L
% !PRINTED 20
\weberpage{20}
whereas we have seen that for $r=1$ there may also be non-linear pencils.

The proof of the theorem is given as follows.

Suppose, for simplicity of exposition, that the curves of the system $|C|$ are irreducible, and consider first the case $r=2$. The $\infty^2$ curves cut one another two by two in the groups of an involution $I_n$ of order $n$ equal to the degree of $|C|$. Take these groups as the elements (points) of a new surface $F'$; then one will have on $F'$ a system $\infty^2$ of curves $C'$, transform of $|C|$, which cut one another two by two in one point. The curves $C'$ which pass through a point $A'$ form a pencil endowed with a base point, and therefore rational (cf. the preceding \S). Alongside this pencil consider another one, with a base point $B'$. The two pencils of $C'$ with centres $A'$ and $B'$ can be referred respectively to two pencils of lines $A$ and $B$ in the plane, by means of projectivities which both make the line $AB$ correspond to the curve $C'$ joining $A'$ and $B'$. From this reference there results a one-to-one correspondence between the points of $F'$ and the points of the plane, in which to the lines of the plane there correspond the curves $C'$, and to the pencils of lines the pencils of $C'$; this implies that the $C'$ on $F'$, and hence also the $C$ on $F$, form a linear system (\textit{net}). The theorem being thus established for nets, it will be extended to the case of linear systems $\infty^3$, and so on in general, passing from dimension $r$ to dimension $r+1$. If $|C|$ is a system $\infty^3$ of curves such that through three generic points there passes one $C$, the $C$ passing through a fixed point $A$ will form a net. Now the nets defined by two base points $A$ and $B$ can be referred projectively to two stars of planes $A'$ and $B'$ of ordinary space, in such a way that to the $C$ of the pencil common to the two nets there correspond, in both alike, the same planes common to the two stars. From this there results a one-to-one (homographic) correspondence between the system of the $C$ and the planes of the space $S_3$, in which to the planes of a pencil there will correspond the $C$ of a pencil, and to the planes of a star the $C$ of a net; this implies that the system $|C|$ is a linear system, Q.E.D.

\edsection{5. Extension of Bertini’s theorems.}{5. Extension of Bertini’s theorems.}{sec.I.5}

\textsc{E. Bertini} gave two fundamental theorems for linear systems of plane curves\origfootnote{1}{s020L-1}{Rendic. Istituto Lombardo, 1880.}: the first says that the curves of a linear system cannot have variable double or multiple points which do not belong to a fixed component; the second — which is a corollary of the first — asserts that whenever the variable parts of a linear system are reducible, they are composed of the curves of a pencil.

% scan: s020R
% !PRINTED 21
\weberpage{21}
These theorems extend easily to linear systems of curves belonging to an arbitrary surface\origfootnote{1}{s020R-1}{Cf. \textsc{F. Enriques}, \textit{Ricerche di Geometria ecc.} Memorie Acc. di Torino, 1893 and \textit{Introduzione alla geometria sopra una superficie algebrica.} Società Italiana delle Scienze (called “dei XL”), 1896. The theorem on reducibility is already noted in Noether, Math. Ann., VIII, 1874.}.

If one follows the proof indicated in the lectures of \textsc{Enriques-Chisini}\origfootnote{2}{s020R-2}{Op. cit., Book II, chap. I, \S\ 5, vol. I, p. 180.}, the extension is immediate.

If the generic curve $C$ of the linear system $|C|$ possesses a variable double or multiple point $O$, the curve $C$ has at least two intersections coinciding at $O$ with every infinitely near curve, and hence with the curves of all the pencils to which it belongs: this means that $O$ is a base point, or else a point belonging to a curve $K$ which is a fixed component of the $C$.

We therefore state that: \textit{the curves of a linear system on a surface cannot have variable double, or multiple, points} (outside the base points) \textit{which do not belong to a fixed component.}

Bertini’s second theorem follows as a corollary.

Let the curves $C$ of $|C|$ (having no fixed components) be reducible, for example into two (variable) parts $K$ and $K'$. If $K$ and $K'$ do not belong to one and the same pencil of which the curves $C$ are composed, there will be through a generic point $O$ at least one curve $K$ and one curve $K'$ distinct from $K$; consequently the $K$ and $K'$ must meet in at least one variable point, and the $C$ must have at least one double point capable of varying freely on the surface, which is absurd. Indeed, if the points common to the $K$ and $K'$ which compose a $C$ always fell at base points of the system $|C|$, all the $K$ and all the $K'$ would pass through those fixed points and would have no variable intersections. Nor can it happen that the complementary curves $K$ and $K'$ meet in points of a fixed curve, because this would turn out to be a fixed component of the $C$, which has been supposed not to exist.

We therefore state: \textit{if the generic curve of a linear system $|C|$ on $F$ is reducible:}

1) \textit{$C$ contains a fixed part, on removing which there remains an irreducible system,} or

2) \textit{the curves $C$ are composed of $r>1$ curves varying in a pencil, rational or not} (and on the latter, conceived as an $\infty^1$ entity, they form a $g_r$), or finally

3) \textit{the $C$ are composed, as in case 2), of the curves of a pencil, and moreover of fixed curves which are added to these.}

Cases 2) and 3) include all the linear systems of dimension $r>1$ and of degree $n=0$ (cf. \S\ 3).

% scan: s021L
% !PRINTED 22
\weberpage{22}
\edsection{6. Image surfaces of linear systems.}{6. Image surfaces of linear systems.}{sec.I.6}

Consider an irreducible linear system \(|C|\) of dimension \(r \geq 3\), and therefore of degree \(n > 0\). By referring \(|C|\) projectively to the system of hyperplanes of a space \(S_r\), one gives rise to a transformation of our surface \(F\) into an \(F'\) of order \(n\), on which the transforms of the \(C\) are cut out by the hyperplanes. Indeed the \(C\) can be regarded abstractly as the elements (hyperplanes) of a space \(S_r\), and hence the indicated projectivity makes a star of hyperplanes passing through a point \(P'\) of \(S_r\) correspond to the linear system \(\infty^{r-1}\) of the \(C\) passing through a point \(P\), and, as \(P\) varies on \(F\), the point \(P'\) varies, describing a surface \(F'\).

The single-valued correspondence between \(F\) and \(F'\) is, in general, uniquely invertible, because to the star with centre \(P'\) there will correspond a linear system of curves \(C\) endowed with a (variable) base point \(P\), and this system will not in general have other base points consequent upon the existence of \(P\).

In truth one cannot exclude that the conditions for the \(C\) to pass through a generic point of \(F\) entail as a consequence their passing through other points \(P_1P_2\ldots P_{m-1}\); then the groups analogous to \(PP_1\ldots P_{m-1}\) will form on \(F\) an \textit{involution of order \(m\) to which the \(C\) will belong}; and the surface \(F'\) of \(S_r\) determined above will turn out to be a surface of order \(\dfrac{n}{m}\), to be regarded as multiple of order \(m\).

That belonging to an involution is in fact a particular circumstance for a linear system \(|C|\) of dimension \(r \geq 3\) follows from a simple count of constants, because there are at least 3 linearly independent curves passing through \(P\), and in general there is no reason why one of these should pass through the intersection of the other two. On the other hand, a linear system \(\infty^2\) (that is, a net) always belongs to an involution of order \(n\), if it is of degree \(n > 1\); for through a point \(P\) there pass only two linearly independent \(C\), and all the other \(C\) through \(P\) form a pencil, which has their intersections as base points. Thus, then, a net \(|C|\) of degree \(n\) leads to a \textit{multiple plane of order \(n\)}, where the \(C\) have as images the \(n\)-fold lines, and on which special importance attaches to the \textit{branch curve}, the locus of the points of the plane to which there correspond groups of the involution endowed with a double point.

One can give an example of a linear system \(|C|\) of dimension \(r \geq 3\) belonging to an involution by considering the system of the \(C\) which are intersections of the surface \(F\) with the cones of a certain order \(n\) having their vertex at a fixed point of space; the system \(|C|\) will now belong to an involution of order \(m\) and will be of degree \(mn\), denoting by \(m\) the order of the surface \(F\). %

% scan: s021R
% !PRINTED 23
\weberpage{23}
Wishing to \textit{exclude membership in an involution} for a system \(|C|\), we shall say briefly that it is a \textit{simple system}.

Returning in general to the transformation of the surface \(F\) to which a linear system \(|C|\) of dimension \(r \geq 3\) gives rise, we shall note that the projective correspondence discussed above is translated into formulae as follows: if the system \(|C|\) is represented by
\[
\lambda_0\varphi_0 + \lambda_1\varphi_1 + \ldots + \lambda_r\varphi_r = 0,
\]
one must take as the coordinates of the points of \(S_r\):
\[
X_0 \equiv \varphi_0 \quad X_1 \equiv \varphi_1 \quad \ldots \quad X_r \equiv \varphi_r.
\]

We shall conclude by stating a theorem which extends to surfaces a well-known consideration relating to curves and to the linear series on them:

\textit{the geometry of birational transformations on surfaces, when one seeks the invariant relations with a simple linear system, \(\infty^3\) at least, is mirrored in the projective geometry of the image surface of the system.}

Thus the relation between two linear systems \(|C|\) and \(|K|\) (simple, \(\infty^3\) at least) on \(F\), of which the first contains the second, is mirrored in the fact that the image surface of \(|K|\) can be regarded as a projection of the image surface of \(|C|\) from a certain linear space. This space will in general cut the said surface in a group of points, if \(|K|\) is deduced from \(|C|\) by the imposition of those base points, and will instead cut the surface in an \(L\) if \(|K|\) is partially contained in \(|C|\) in such a way that by adding the fixed curve \(L\) to the \(K\) one obtains curves \(C\).

As concerns base points, it will be useful to bear in mind the following observations. If \(O\) is, on the surface \(F\), an \(i\)-fold base point of \(|C|\) (simple for \(F\)), %
it will be fundamental for the transformation which changes \(F\) into the \(\Phi\); %
having the \(C\) as hyperplane sections; and if the \(C\) have no fixed tangent at \(O\), to \(O\) there will correspond on \(\Phi\) \textit{an exceptional curve} of order \(i\). If instead, \(O\) being an \(i\)-fold base point for \(|C|\) with \(i \geq 1\), the \(C\) have a fixed tangent at \(O\) (that is, a simple base point \(O_1\) infinitely near to \(O\)), without higher contact, one can still say that to \(O\) there corresponds on \(\Phi\) an exceptional curve of order \(i\), but this will in general consist of a curve of order \(i-1\), whose points correspond to the variable points in the first-order neighbourhood of \(O\), and of a line, image of the neighbourhood of \(O_1\) (whose points belong to the second-order neighbourhood of \(O\)). In particular, if \(i=1\), the whole neighbourhood of the simple point \(O\) will have as its image on \(\Phi\) the neighbourhood of a conical double point, to which will be added the exceptional line, image of \(O_1\) (think, for example, of the plane representation of the quadric cone).

% scan: s022L
% !PRINTED 24
\weberpage{24}
\textsc{Note.} -- In general, if the simple point \(O\) of the surface \(F\) constitutes an arbitrarily complicated singularity for the curves \(C\) of \(|C|\), one can resolve this singularity by transforming \(F\) into a surface \(\Phi\) on which the neighbourhood of the point \(O\) (taken in its entirety) will be represented by an exceptional curve composed of as many parts as there are branches (linear or not) constituting the singularity of the \(C\) at \(O\); and the connection of these parts will mirror the composition of the said singularity by means of infinitely near multiple or simple points, and their relations of proximity\origfootnote{1}{s022L-1}{\textsc{Enriques-Chisini}, Op. cit., Book IV, vol. II, p. 327 ff.} %
in a manner which has been the object of the analysis of \textsc{E. Barber} and of \textsc{O. Zariski}\origfootnote{2}{s022L-2}{Cf. \textsc{Barber} and \textsc{Zariski}, \textit{Reducible exceptional curves of the first Kind.} American Journal of Mathematik, 1935 (p. 119). \textsc{P. Du Val}, American Journal of Math., 1936.}. But we shall have occasion to return to the subject. %

\edsection{7. Equivalent curves and complete linear systems.}{7. Equivalent curves and complete linear systems.}{sec.I.7}

The considerations that follow extend step by step to linear systems of curves on a surface what is said for equivalent groups of points and for complete linear series belonging to a curve\origfootnote{3}{s022L-3}{Cf. \textsc{Enriques-Chisini}, Op. cit., Book V, vol. III (1924).}.

On the surface \(F\) two curves (of the same order) are called \textit{equivalent} when they belong to one and the same linear system of dimension \(r \geq 0\): if they are distinct, this means that they belong to a linear pencil; the consideration of systems of dimension zero means that a curve is to be regarded as equivalent to itself.

The equivalence of two curves \(C\) and \(K\) will be expressed by writing
\[
C \equiv K \qquad \text{or also} \qquad C = K.
\]

Two equivalent curves \(C\) and \(K\) can always be considered as the curve of zeros and the curve of poles of a rational function \(t\) on \(F\); the two curves are interchanged if the function \(t\) is replaced by \(\dfrac{1}{t}\).

The relation of equivalence between curves on a surface enjoys the three fundamental properties of equality:
\begin{enumerate}
\item[1)] the \textit{reflexive} property, expressed by \(C \equiv C\),
\item[2)] the \textit{symmetric} property, if \(C \equiv K\) then also \(K \equiv C\),
\item[3)] the \textit{transitive} property, if \(C \equiv K\) and \(K \equiv L\) then also \(C \equiv L\).
\end{enumerate}

To prove this last property, construct a rational function \(t\) having as curve of zeros \(K\) and as curve of poles
% scan: s022R
% !PRINTED 25
\weberpage{25}
\(C\), and then a rational function \(\tau\) having \(K\) as curve of zeros and \(L\) as curve of poles; then the function \(\dfrac{\tau}{t}\) will have as curve of zeros \(C\) and as curve of poles \(L\). This proves precisely that \(C\) and \(L\) are equivalent.

From the aforesaid properties there follows the fundamental theorem:

\textit{On a surface \(F\) the totality of the curves equivalent to a given curve \(C\) constitutes a linear system, which is called complete.}

This system is obtained as the set of the level curves of the rational function which results from combining linearly the independent rational functions for which \(C\) is the curve of poles.

The \textit{complete} system \(C\) thus defined will in general be \textit{free from base points}, and in any case it will be understood, unless otherwise stated, to be \textit{virtually} free from base points, in the sense that will be made precise later; but one can also define the complete system relative to a \textit{group of assigned base points} on \(C\). Note that, on transforming the surface in such a way that the base points of \(C\) become (exceptional) curves, the transform of \(C\) must be regarded as stripped of these curves, which would otherwise appear as fixed components of the transformed system. And so the notion of the complete system relative to a group of base points reduces in general to that of the complete system without base points, on a transformed surface.

In any case, the irreducible complete system, with or without base points, is defined, starting from a linear system (\(\infty^2\) at least) contained in it, as the largest system of the same degree which contains the given one.

From the transitive property of equivalence it follows that the complete linear system \(|C|\), with or without assigned base points, is equally defined starting from any one of its curves.

It is obvious that the complete system \(|C|\) can be constructed as follows.

First construct an arbitrary linear system which contains \(C\) (totally), and then enlarge this system successively, as long as possible; since the dimension cannot exceed the degree increased by one, one will in this way arrive at a system \(|C|\) which will not be contained in another, larger linear system of the same order, and this will be the complete linear system defined by \(C\). Hence, however one proceeds in the construction of the successively larger systems that have been discussed, one will finally arrive at the same complete linear system.

To the concept of complete linear system there corresponds the concept of \textit{normal surface}, a surface of a space \(S_r\) which cannot be obtained as the projection of another of the same order belonging to a higher space: for the surface to be normal means that the
% scan: s023L
% !PRINTED 26
\weberpage{26}
(irreducible) system of its hyperplane sections is complete, and conversely.

Here, then, is the projective meaning of the theorem on complete systems. Consider, for example in ordinary space, a surface \(F\) of a certain order \(n\); it may happen that \(F\) is obtained as the projection of two surfaces of the same order of \(S_4\), \(F_1\) and \(F'_1\), the projections being made respectively from two points \(O\) and \(O'\). Now if \(F_1\) and \(F'_1\) are normal surfaces, they will turn out to be projective to each other, that is, projectively identical. If this is not so, it means that \(F_1\) can be regarded as the projection of a surface \(F_2\) of \(S_5\), and \(F'_1\) as the projection of an \(F'_2\) of \(S_5\), and so on; proceeding in any manner whatever from one surface to another of the higher space, one will finally arrive at a surface \(F_r\) of the space \(S_{r+3}\) which will be normal; and this \(F_r\) will always remain projectively defined, however one arrives at it through different series of surfaces (\(F_1, F_2, \ldots,\) or \(F'_1, F'_2, \ldots\)) which are deduced from one another by projection.

Examples relating to normal surfaces are already met with in the theory of rational surfaces\origfootnote{1}{s023L-1}{Cf. \textsc{Enriques-Conforto}, \textit{Le superficie razionali}, Bologna, 1939 (1945).}.

Thus a surface \(F_5\) of the \(5^\circ\) order of ordinary space, with elliptic plane sections, endowed with a double curve of the \(5^\circ\) order belonging to a quadric cone, and having a triple point, triple also for the surface, at the vertex of the said cone, can be regarded as the projection of several surfaces \(F_5\) of the same order of the space \(S_4\), not projective to one another; but these are projections of a single normal surface of \(S_5\), projectively defined. These assertions are justified by starting from the plane representation of \(F_5\) by means of a linear system \(\infty^3\) of cubics \(C_3\) with 4 simple base points \(A_1A_2A_3A_4\).

Indeed such a system \(\infty^3\) is contained in \(\infty^1\) linear systems \(\infty^4\) of \(C_3\) with the same base points \(A_1A_2A_3A_4\), which are, in the plane, homographically (and also birationally) distinct, and therefore represent surfaces \(F_5\) of \(S_4\) not projective to one another. On the other hand, the normal surface of \(S_5\) of which these \(F_5\) are projections (from an external point) is projectively defined by the complete system of the \(C_3\) through \(A_1A_2A_3A_4\).

\edsection{8. Sum and difference of linear systems: residue theorem.}{8. Sum and difference of linear systems: residue theorem.}{sec.I.8}

Curves composed of equivalent curves, that is, \textit{sums of equivalent curves, are equivalent.}

Indeed, let \(C_1\) and \(C_2\), \(K_1\) and \(K_2\) be two pairs of equivalent curves, and take as \(t\) a rational function having \(C_1\) as polar curve
% scan: s023R
% !PRINTED 27
\weberpage{27}
and \(C_2\) as curve of zeros; and as \(\tau\) a rational function having \(K_1\) as polar curve and \(K_2\) as curve of zeros; then the function \(\tau t\) will have \(C_1+K_1\) as polar curve and \(C_2+K_2\) as curve of zeros. Thus from
\[
C_1 \equiv C_2,\qquad K_1 \equiv K_2
\]
one deduces
\[
C_1+K_1 \equiv C_2+K_2.
\]

It follows from this that, given any two linear systems \(|C|\) and \(|K|\), there exists a well-determined complete system \(|C+K|\) which totally contains the composite curves \(C+K\); this will be called the \textit{complete sum system} of \(|C|\) and \(|K|\).

It is worth noting that if \(|C|\) and \(|K|\) are irreducible linear systems of dimension greater than zero, not coinciding in one and the same pencil, the complete system \(|C+K|\) is irreducible. Let us add that if the irreducible systems \(|C|\) and \(|K|\) possess assigned base points, these are always base points for \(|C+K|\). This is an obvious consequence of the extension of Bertini’s theorem. In particular the double, and in general the \(s\)-fold multiple system \(|sC|\), of the system \(|C|\) is irreducible, except in the case where \(|C|\) is a pencil.

With \(|C|\) and \(|K|\) again two systems not coinciding in one and the same pencil, consider the composite curves \(C+K\); we have said that they are contained in the complete system \(|C+K|\); but in general they may also be contained in systems of lower dimension, and in particular there will be a \textit{minimum sum system} of the two given ones. %

For example, take a surface \(F\) of order \(n\) of ordinary space, free from singularities; through a point \(O\) of it there pass \(\infty^2\) planes, which cut out a complete net; the minimum double of this net is the system \(\infty^5\) cut out by the quadric cones with vertex at \(O\); but this system, which has \(O\) as a double base point, is not complete, being contained in the system \(\infty^6\) cut out on \(F\) by the quadrics which touch it at \(O\).

The \textit{subtraction} of two linear systems, one of which partially contains the other, is defined as the \textit{inverse operation of addition}, and precisely as follows.

Let \(|C|\) be a linear system to which the split curve \(C \equiv L+K\) belongs; then the \textit{residual} system of the curve \(K\) with respect to \(|C|\) is the system of all the curves \(L\) which together with \(K\) constitute a curve of \(|C|\). By virtue of the properties of equivalent curves one can assert that:

if a particular curve \(C=L+K\) of the complete system \(|C|\) partially contains a \(K\), every other curve equivalent to \(K\) forms part of a \(C\). And hence one has the \textit{residue theorem}: \textit{the linear}
% scan: s024L
% !PRINTED 28
\weberpage{28}
\textit{system \(|L|\) residual to a \(K\) with respect to a system \(|C|\) is also residual to every other curve equivalent to \(K\).}

In short, the \textit{complete linear difference system} of \(|C|\) and \(|K|\):
\[
|L| = |C-K|,
\]
is defined by the relation
\[
|C| = |L+K|.
\]

\textsc{Note.} -- The theorem stated admits a projective complement which we shall prove further on. \textit{We shall see, in fact, that on a surface \(F\) of ordinary space with normal singularities, the (adjoint) surfaces of a given order which pass through the double curve cut out, outside it, a complete linear system.}

In particular, on a surface of ordinary space free from singularities, the surfaces \(\varphi\) of a given order cut out a complete system.

If, referring to the general case, one draws the \(\varphi\) of a certain order passing through a curve \(C\) of \(F\), and considers the residual intersections \(K\) of these adjoint surfaces with \(F\), \textit{every curve \(C\) of the complete system \(|C|\) will equally be residual to every \(K\)}, with respect to the system cut out by the \(\varphi\) (outside the double curve), and conversely every \(K\) is residual to every \(C\).

This is the form of the \textit{Restsatz} of \textsc{M. Noether}\origfootnote{1}{s024L-1}{\textit{Mathematische Annalen}, Bd. VIII. For curves cf.: \textsc{Brill} and \textsc{Noether}, \textit{Math. Ann.}, Bd. VII; \textsc{Enriques-Chisini}, Op. cit., Book V.}, which, analogously to what is said for the geometry on a curve (an observation of \textsc{G. Castelnuovo}, 1890), contains two distinct assertions: the first relating to the subtraction of complete systems, the other supplying a projective complement to it, concerning the completeness of the system cut out on a surface \(F\) of the space \(S_3\) by its adjoints \(\varphi\) of a given order.

It should now be observed that, while the operation of addition leads in general, as has been said, from irreducible linear systems to irreducible systems, this no longer holds for subtraction: when one subtracts from an irreducible system \(|C|\) a curve \(K\) (irreducible or not), the complete residual system \(|C-K|\) may turn out to be reducible. A simple example is obtained by considering a linear pencil \(|L|\), which for simplicity we suppose free from base points, and a linear system \(|K|\): if one adds \(|K|\) to \(|2L|\) one obtains a system \(|C|=|K+2L|\). And subtracting \(|K|\) from this, one has the double system of \(|L|\), a reducible system composed of the curves of the pencil \(|L|\).

It may also happen that, on subtracting a \(K\) from the irreduc%
% scan: s024R
% !PRINTED 29
\unskip\weberpage{29}%
\unskip{}ible system \(|C|\), one obtains a linear system containing a fixed part \(K'\), whose detachment from the \(C\) is a consequence of the detachment of \(K\); and similarly it may also happen that, on subtracting a curve \(K\) from \(|C|\), the residual system \(|L|\) has as a consequence some new base point \(O\) besides the base points of \(|C|\) which do not belong to \(K\). This case reduces to the preceding one by a transformation of the surface which changes \(O\) into an exceptional curve, because on the transformed surface that curve will appear as a fixed part of \(|L|\).

\textit{Remark.} -- The definition of the system \(|L|=|C-K|\) is bound up with a presupposition of existence, since the difference acquires an effective meaning only if the system \(|C|\) partially contains \(|K|\). But in what follows it will also be useful to speak of \textit{virtual curves}, which are defined by subtraction even when this is effectively impossible. Accordingly the symbolic relation
\[
|C-K| = |L-M|
\]
will be regarded as equivalent to
\[
|C+M| = |K+L|.
\]

Virtual curves are defined like the negative numbers in Peano’s theory of operators, that is, as possible addends. To add the virtual curve \(C-K\) to a system \(|D|\) will mean to add \(C\) to it and to subtract \(K\) from it, which will in general be possible for sufficiently large systems. In any case the equivalence of the virtual curves \(C-K\) and \(L-M\) means the equivalence of the linear systems
\[
|D+C-K| = |D+L-M|,
\]
where
\[
|D+(C-K)| = |D+C-K|: %
\]

However, one can give a more restricted concept of virtual curves, as curves having no effective existence but of which the double, or else a multiple, exists, as we shall have occasion to see in what follows. And this concept corresponds, rather than to the negative numbers, to that of the \textit{ideals} which are met with in the theory of algebraic fields.

\textit{Historical note.} -- The idea, if not the name, of “virtual curves”, that is, the idea that the curve \(C-K\), defined by means of the linear system which is the difference of two others, corresponds to a mathematical entity endowed with a certain reality, even when the subtraction does not lead to an effective curve, occurred to \textsc{Enriques} in connection with canonical curves: indeed it constitutes the principal motive of the new development given to the theory in the passage from the \textit{Ricerche} of 1893 to the \textit{Introduzione} of 1896, where what is sought is in substance the widest meaning assumed by the invariance of the said canonical curves, even when the ge%
% scan: s025L
% !PRINTED 30
\unskip\weberpage{30}%
\unskip{}nus \(p=0\). Of the difference curve, \(C-K\), the Author defines the virtual characters, degree and genus (cf. \S~10), and notes that, though it may be missing, it can have a double endowed with effective existence. The term “virtual”, applied to the curve as well as to its characters (which were so designated by \textsc{Enriques}), is found in a Note of \textsc{Severi}, \textit{Sulle curve algebriche virtuali appartenenti ad una superficie algebrica} (Rendic. Istituto Lombardo, 1905). The topological interpretation of virtual curves as two-dimensional cycles was noticed by \textsc{S. Lefschetz}.

\edsection{9. Image surface of the sum system.}{9. Image surface of the sum system.}{sec.I.9}

Let there be on the surface \(F\) two irreducible linear systems \(|C|\) and \(|K|\), of dimension \(r\) and \(s\) respectively, where \(r \geq 1\), and \(s \geq 1\). One can in general construct a transformed surface \(\Phi\), image of the sum system \(|C+K|\), in a space of \(R\) dimensions, where \(R > r+s\); the sum system in question may be the minimum sum system or also, if one wishes, the complete system.

In any case, there will be in \(S_r\) two linear series of spaces, of \(R-r-1\) dimensions and of \(R-s-1\) dimensions respectively, cutting the \(K\) and the \(C\), in such a way that the \(\infty^r\) hyperplanes passing through an \(S_{R-r-1}\) of the first series will cut out the \(C\), containing the \(S_{R-s-1}\) to which these belong, and conversely let the same be said for the hyperplanes which pass through an \(S_{R-s-1}\) of the second series. %

Fix an \(S_{R-r-1}\) of the first series cutting a curve \(C\) of \(|C|\), and an \(S_{R-s-1}\) of the second, cutting a curve \(K\) of \(|K|\); these will be contained in a hyperplane of \(S_R\), and will therefore have an \(S_{R-r-s-1}\) in common; projecting from this onto an \(S_{r+s}\), one will obtain in that space an image \(F\) of the given surface which will contain a particular \(K\), section of an \(S_{s-1}\), and a particular \(C\), section of an \(S_{r-1}\), and in which the hyperplanes through the \(S_{s-1}\) cut out the \(\infty^r\) curves \(C\), and the hyperplanes through the \(S_{r-1}\) the \(\infty^s K\).

Let us take in particular \(r=2\), \(s=1\): we shall be led to a surface of ordinary space on which the net of the curves \(C\) is cut out by the planes through a point \(O\), and the pencil of the \(K\) by the planes through a line \(a\) (not passing through \(O\)), and the surface so obtained will certainly be in one-to-one correspondence with the one from which we started, if \(|C|\) and \(|K|\) do not belong to one and the same involution.

The construction of the surface \(F\) can also be obtained directly, by setting up a projective correspondence between the net \(|C|\) and the star of planes \(O\), and between the pencil \(|K|\) and the pencil of planes \(a\). To the \(\overline{K}\) which is added to the \(C\) of the net there thus comes to correspond the neighbourhood of the multiple point \(O\), which will be precisely \(i\)-fold, if one denotes
% scan: s025R
% !PRINTED 31
\weberpage{31}
by \(i\) the number of intersections of a \(C\) and a \(K\); indeed the neighbourhood of \(O\) must have (like \(\overline{K}\)) \(i\) points in common with a plane section \(C\). On the other hand, to the \(\overline{C}\), which is added to the \(K\) of the pencil, there will correspond the line \(a\), multiple for \(F\), which will be precisely \(n\)-fold, denoting by \(n\) the degree of \(|C|\). The plane \(Oa\) will cut the surface \(F\), outside \(a\), in \(i\) exceptional lines issuing from \(O\). The surface \(F\) will therefore be of order \(n+i\), and one sees that it is the image of a particular system \(\infty^3\) contained in \(|C+K|\), which possesses \(i\) base points at the points common to \(C\) and \(K\): it is these points that give rise to the exceptional lines through \(O\) mentioned above.

If now one wishes to construct on \(F\) the sum system of the net \(|C|\) and the pencil \(|K|\), it is clear that one should have recourse to the system of the surfaces obtained by adding a plane through \(O\) and a plane through \(a\), and therefore to the linear system \(\infty^5\) of the quadrics through \(O\) and \(a\): one will thus obtain the minimum sum system \(|C+K|\).

\edsection{10. Virtual characters.}{10. Virtual characters.}{sec.I.10}

Let \(|C_1|\) and \(|C_2|\) be two irreducible linear systems on the surface \(F\), and denote by \(n_1\) and \(n_2\) respectively their effective degrees, by \(\pi_1\) and \(\pi_2\) their genera, and by \(i\) the number of variable intersections of a \(C_1\) with a \(C_2\). The \textit{degree of the system} \(|C_1+C_2|\) will be the number of intersections
\[
(C_1+C_2)^2 = C_1^2+C_2^2+2C_1C_2,
\]
and it will therefore be
\[
n = n_1+n_2+2i.
\]

The \textit{genus of the sum system} will be
\[
\pi = \pi_1+\pi_2+i-1,
\]
this being the expression for the genus of a curve which comes to split into two parts, an expression determined by \textsc{Noether} in accordance with the requirement of continuity\origfootnote{1}{s025R-1}{Cf. \textsc{Enriques-Chisini}, Op. cit., Book V, vol. III, p. 399.}.

These formulae lead one to define \textit{the virtual degree and genus of a curve \(C\)} on \(F\), independently of the possibility of extending the linear system \(\infty^0\) determined by it. To this end associate with \(C\) an irreducible system \(|K|\) such that the system \(|C+K|\) turns out to be irreducible; this can easily be done by passing through \(C\) the surfaces of a sufficiently high order and taking as curves \(K\) the residual curves. Now, denoting by \(n\) the degree of the \(C\) which is to be
% scan: s026L
% !PRINTED 32
\weberpage{32}
defined, by \(m\) the degree of \(K\), and by \(i\) the number of intersections of a \(C\) with a \(K\), the degree of \(|C+K|\) will be given by
\[
N=n+m+2i,
\]
from which one obtains
\[
n=N-m-2i.
\]

\textit{The virtual degree of \(C\)} thus defined \textit{is independent of the choice of the auxiliary system \(|K|\) that has been associated with it.} Indeed, let \(|K_1|\) and \(|K_2|\) be two auxiliary systems satisfying the conditions laid down, with degrees \(m_1\) and \(m_2\), and let \(i_1\) and \(i_2\) be the numbers of intersections \(CK_1\) and \(CK_2\), and \(j\) the number of variable intersections of a \(K_1\) and a \(K_2\).

One can compare the degree \(n_1\) defined with respect to \(|K_1|\) and the degree \(n_2\) defined with respect to \(|K_2|\) by evaluating the degree of the system
\[
|(C+K_1)+K_2|=|(C+K_2)+K_1|=|C+(K_1+K_2)|.
\]

One will therefore have
\[
\begin{aligned}
N&=(n_1+m_1+2i_1)+m_2+(2j+2i_2)\\
&=(n_2+m_2+2i_2)+m_1+(2j+2i_1),
\end{aligned}
\]
whence it follows that
\[
n_1=n_2.
\]

In an entirely similar way one can define the \textit{virtual genus} \(\pi\) of an isolated curve \(C\), and this character retains a meaning even when \(C\) is a multiple curve, or a reducible curve containing some multiple part.

Let us add that the definition of the virtual \textit{characters} (degree and genus) extends also to the case of a \textit{virtual curve} \(C-K\) defined as a difference independently of its effective existence (§ 8). For this purpose it suffices to add to \(C-K\) a curve \(L\) such that \(|M|=|C-K+L|\) is an irreducible linear system, and then to evaluate the characters of the difference \(|M-L|\), which turn out to be independent of \(|L|\).

Returning to effective curves, let us now note that the notion of the virtual characters of a curve \(C\) gives rise to some observations which tend to attribute to these characters a meaning invariant with respect to birational transformations of the surface.

Transform \(F\) into another surface \(\bar F\) in such a way that to a point \(O\) of the curve \(C\) there corresponds an exceptional curve \(o\) of \(\bar F\); then one may regard as the transform of \(C\) the composite curve \(\bar C+o\), or else only \(\bar C\), leaving out \(o\), which corresponds to a fundamental point of the transformation. The choice which can thus be made corresponds to a different way of considering the point \(O\) on
% scan: s026R
% !PRINTED 33
\weberpage{33}
the curve \(C\): if \(C\) is capable of belonging to an irreducible linear system of dimension \(>0\) with the base point \(O\), the transform of this system will be the likewise irreducible system \(|\bar C|\), in which \(o\) does not appear as a fixed part. But if instead \(|C|\) can be extended so as no longer to possess the base point \(O\), its transform will contain in particular the system \(o+|\bar C|\), where \(o\) appears as a fixed part. In other words, the exceptional curve \(o\) corresponding to \(O\) will or will not have to be regarded as a component of the transformed curve of \(C\), according as one intends to consider \(C\) as a curve of a linear system not having, or having, \(O\) as a base point. The circumstance that a system \(|C|\) without a base point \(O\) effectively exists can be regarded as accidental with respect to the above consideration. The \textit{virtual degree} of \(C\), having regard to the possible transformations of its points into exceptional curves, will therefore be defined \textit{in general} by supposing that it does \textit{not} have \textit{base points}, or else by declaring which \textit{base points} are understood to be assigned on the curve, apart from which every other point of \(|C|\) that turns out to be a \textit{base point} for an effectively constructed \(|C|\) must be regarded as \textit{virtually non-existent}. Let us add that if \(C\) possesses a multiple point (simple for the surface), of multiplicity \(\geq i\), this may be taken virtually as \(i\)-fold for the curve \(C\), even if in effect it possesses a higher multiplicity, and even if this turns out to be the effective multiplicity of the largest system \(|C|\) which is required to possess that point as an \(i\)-fold point.

In particular, then, if on the surface \(F\) the curve \(C\) possesses an \(i\)-fold point which is not to be regarded as an assigned base point for the system \(|C|\), it must be considered \textit{virtually non-existent}, and therefore the \textit{virtual genus} of \(C\) will be obtained by adding to the effective genus \(\pi\), \(\dfrac{i(i-1)}{2}\), which is precisely the number by which the genus of an irreducible curve diminishes when it comes to acquire an \(i\)-fold point.

Similarly the \textit{virtual degree} of \(C\) will be evaluated by adding \(i^2\) to the effective number of intersections of two \(C\) outside this point, since it is precisely by \(i^2\) that the number of intersections of two curves diminishes when they acquire a common \(i\)-fold point.

\edsection{11. Exceptional curves.}{11. Exceptional curves.}{sec.I.11}

From what has been said above, the neighbourhood of a (simple) point of the surface \(F\), when this point is a base point for a linear system of curves, can be regarded as an \textit{infinitesimal curve} which, under a transformation of the surface, changes into a curve properly so called (exceptional curve). Conversely, an exceptional curve is a
\printersignature{\textsc{Enriques F.} - \textit{Superficie algebriche.}\quad 3}

% scan: s027L
% !PRINTED 34
\weberpage{34}
curve traced on the surface \(F\) which can be changed into the neighborhood of a simple point.

Wishing to examine exceptional curves more closely, we shall first of all establish this distinction:

an \textit{irreducible} curve \(\omega\) on the surface \(F\) will be called an \textit{exceptional curve of the \(1^{\mathrm{st}}\) kind} if it is capable of being transformed into a simple point by means of a linear system \(|C|\) which has no base points on \(\omega\). On the other hand, a curve \(\omega\) transformable into a point by means of a linear system \(|C|\) which has some base point on it will be called an \textit{exceptional curve of the \(2^{\mathrm{nd}}\) kind}. %

Examples of exceptional curves of the first kind are the 27 lines belonging to a general cubic surface \(F_3\). Indeed \(F_3\) can be represented on the plane in such a way that to any one of these lines there corresponds the neighborhood of one of the 6 base points of the system of representative cubics. In the plane, on the other hand, there are no exceptional curves of the first kind. But the lines, the conics and in general the plane curves which are fundamental for a Cremona transformation constitute exceptional curves of the second kind: a line will be transformed into a point by means of a homaloidal net (for example of conics) having two base points on the line; and a conic will be transformed into a point by taking, for example, a homaloidal net of cubics with a double base point and four simple base points on the conic.

Already from these examples it follows that the distinction between exceptional curves of the first and of the second kind has only a projective character, or enjoys an \textit{invariance relative} to transformations without base points: indeed an \textit{exceptional curve of the second kind} \(\omega\) belonging to a surface \(F\) \textit{can always be changed into an exceptional curve of the first kind} on a transformed surface \(\Phi\). To this end consider the transforming linear system \(|C|\) which changes \(\omega\) into a point; by hypothesis it has a certain number \(s\) of base points, with certain multiplicities \(i_1,i_2,\ldots,i_s\), on \(\omega\), and the number of its intersections with \(\omega\) equals precisely the sum \(i_1+i_2+\ldots+i_s\). Now, if one adds to \(|C|\) a linear system \(|L|\) without base points, and for which \(\omega\) is not a fundamental curve, the system \(|C+L|\) will serve to transform \(\omega\) into a curve which will be fundamental for the transform of the system \(|C|\) (without base points on it) and which will therefore turn out to be an exceptional curve of the first kind.

Let us try to evaluate the genus and the degree of an irreducible exceptional curve of the first kind \(\omega\), referring to the simple point \(O\) whose neighborhood corresponds to \(\omega\) on a transformed surface. If on the latter one takes any linear system \(|C|\) for which \(O\) is not a base point, the imposition of the base point \(O\) diminishes the degree of \(|C|\) by one; moreover the curves \(C\) through \(O\) have one point in common with
% scan: s027R
% !PRINTED 35
\weberpage{35}
the neighborhood of \(O\) itself. Therefore, denoting by \(n\) and \(\pi\) the degree and the genus of \(|C|\), and by \(\nu\) and \(\varrho\) the degree and the genus of the infinitesimal curve which constitutes the neighborhood of \(O\), that is, of the exceptional curve \(\omega\), we shall have
\[
\begin{aligned}
n&=n-1+\nu+2\cdot1\\
\pi&=\pi+\varrho+1-1,
\end{aligned}
\]
and hence
\[
\nu=-1\qquad\text{and}\qquad\varrho=0.
\]
The formula \(\varrho=0\) follows, moreover, a priori, the neighborhood of \(O\) being rational.

Naturally the same computation can be repeated on the surface \(\Phi\), the transform of \(F\), which contains \(\omega\). Here one has a system \(|C|\) for which \(\omega\) is a \textit{fundamental curve} (without intersections with the variable curves \(C\)), and the curves \(C\) passing through a point of \(\omega\) break up into \(\omega\) and a residual curve \(C_1\) unisecant to \(\omega\): on writing down the degree and the genus of \(|C|=|\omega+C_1|\), one finds again the preceding formulas. %

We therefore state that: \textit{an irreducible exceptional curve of the \(1^{\mathrm{st}}\) kind has virtual genus \(\varrho=0\) and virtual degree \(\nu=-1\).}

Passing now to exceptional curves of the \(2^{\mathrm{nd}}\) kind, it suffices to recall that such a curve \(\omega\) can be transformed into an exceptional curve of the \(1^{\mathrm{st}}\) kind, provided one removes from it certain points, to be regarded as assigned base points of a transforming system \(|C|\); therefore \textit{the virtual genus and degree of an exceptional curve of the \(2^{\mathrm{nd}}\) kind will be}
\[
\varrho=0\qquad\text{and}\qquad\nu\geq0.
\]

For example, the degree of the generators of a ruled surface is
\[
\nu=0;
\]
the degree of a line of the plane is
\[
\nu=1,
\]
that of a conic
\[
\nu=4,\ \text{etc.}
\]

Let us now note that the requirement that the virtual \textit{genus} of an \textit{exceptional curve be} \(\varrho=0\) implies that an \textit{irreducible exceptional curve} \(\omega\) of the \(1^{\mathrm{st}}\) kind \textit{cannot} have \textit{double points} or multiple points. Indeed the virtual genus \(\varrho\) is evaluated by calculating the genus of the sum system \(|C|=|C_1+\omega|\), and therefore a possible double point of the component \(\omega\), not being a base point for \(|C|\), must be regarded as non-existent; and then, since already \(\varrho=0\), the curve \(\omega\) would turn out to be reducible.

Anticipating some notions which will be acquired later on, we wish to say here that exceptional curves are defined by their
% scan: s028L
% !PRINTED 36
\weberpage{36}
virtual characters found above. If \(\omega\) is an irreducible curve of genus \(\varrho=0\) and of degree \(\nu=-1\) on the surface \(F\), it can be transformed into a simple point by proceeding as follows. Take on \(F\), in an entirely general way, a linear system \(\infty^3|C|\) without base points on \(F\), of a certain genus \(\pi\) and of a certain degree \(n\), whose curves cut \(\omega\) in \(i\) points. Adding \(\omega\) to \(|C|\) one obtains a system \(|C+\omega|\) of curves \((i-1)\)-secant to \(\omega\), which will be of degree \(n+2i-1\) and of genus \(\pi+i-1\), so that the difference between the degree and twice the genus, \(n-2\pi\), is increased by one unit. And, as we shall recognize later on, this increase has an effective meaning, because the addition of \(\omega\) does in fact succeed in enlarging the system \(|C|\), the dimension of the irreducible system \(|C+\omega|\) turning out to be greater than that of \(|C|\). This being so, one can proceed by adding \(\omega\) again to the system obtained, and so on, until one arrives at a system \(|C+i\omega|\) without base points on \(F\), for which \(\omega\) turns out to be fundamental. By means of this system the surface \(F\) can be transformed into a surface \(\Phi\) on which to \(\omega\) there corresponds the neighborhood of a simple point \(O\). In conclusion, \(\omega\) is an exceptional curve of the \(1^{\mathrm{st}}\) kind.

The procedure of transformation used here leads to the following result:

if the surface \(F\) contains a certain number \(s\) of exceptional curves which cut the curves \(C\) in \(i_1,i_2,\ldots,i_s\) points respectively, where
\[
i_1+i_2+\ldots+i_s>2\pi-2-n
\]
one will construct on \(F\) an irreducible linear system of degree \(N\) and of genus \(\Pi\) for which
\[
N>2\Pi-2.
\]

We shall see later on that the existence of systems whose characters give rise to such an inequality serves to define a particular family of surfaces; indeed it characterizes the family of rational surfaces and of those referable to ruled surfaces; on a surface which does not belong to this particular family there holds, for every linear system of genus \(\pi\) and of degree \(n\), the fundamental inequality
\[
n\leq2\pi-2.
\]

A consequence of what has been said will be that: \textit{a surface not belonging to the family of ruled surfaces can always be transformed into another one entirely free of exceptional curves.}

At the same time it will be recognized that ``a surface which contains exceptional curves of the \(2^{\mathrm{nd}}\) kind belongs to the family of ruled surfaces''.

Indeed, let us first transform, as has been said, an excep%
% scan: s028R
% !PRINTED 37
\unskip\weberpage{37}%
tional curve of the \(2^{\mathrm{nd}}\) kind into a curve \(\omega\) of genus \(\varrho=0\) and of degree \(\nu\geq0\). Here, since \(\nu\geq0\), the procedure of adding \(\omega\) to a system \(|C|\) never comes to an end, because the number of intersections of the curves of the system with \(\omega\) increases, or at least does not diminish, so that one will certainly arrive at a linear system of genus \(\Pi\) and of degree \(N\) for which
\[
N>2\Pi-2.
\]

\edsection{12. Note on reducible exceptional curves.}{12. Note on reducible exceptional curves.}{sec.I.12}

When a surface \(F\) is transformed into another, \(\Phi\), by means of a linear system \(|C|\) which has on \(F\) distinct base points \(O,O_1,O_2,\ldots\), the exceptional curves which correspond to these points, and which are at the same time fundamental for the same system \(|L|\) (the transform of the plane sections of \(F\)), are irreducible curves disconnected from one another. But what happens when the points \(O,O_1,O_2,\ldots\) become infinitely near?

Suppose that the said points \(O,O_1,O_2,\ldots\) have for \(|C|\) certain multiplicities \(i,i_1,i_2,\ldots\), where \(i>i_1>i_2\ldots\) and the differences \(i-i_1,\ i_1-i_2,\ldots\) are as large as necessary. To these points there correspond in general on \(\Phi\) exceptional curves \(\Omega,\Omega_1,\Omega_2,\ldots\), of orders \(i,i_1,i_2,\ldots\) resp. But if the point \(O_1\) becomes infinitely near to \(O\), the curve \(\Omega\) of order \(i\) comes to break up into two parts \(\omega\) and \(\omega_1(=\Omega_1)\): the first, of order \(i-i_1\), corresponding to the neighborhood of the point \(O\) deprived of \(O_1\), and the second, of order \(i_1\), the limit of the exceptional curve corresponding to \(O_1\).

In general, if several points \(O_1O_2\ldots\) become infinitely near to a proper point \(O\), following one another on one or more branches, linear or superlinear, the exceptional curve which corresponds to the point \(O\) will break up into several parts connected with one another, comprising the exceptional curves which arise from the successive points. The matter can be made precise by referring to the characteristic case in which \(O,O_1,O_2,\ldots\) follow one another on a single branch, and by beginning with the simplest examples. If \(O,O_1,O_2,\ldots,O_s\) follow one another on a linear branch, the exceptional curve of order \(i\) which corresponds to \(O\) becomes a composite curve of order \((i-i_1)+(i_1-i_2)+\ldots+(i_{s-1}-i_s)+i_s\):
\[
\Omega=\omega+\omega_1+\omega_2+\ldots+\omega_s,
\]
while the curves arising from \(O_1O_2\ldots\) become
\[
\begin{aligned}
\Omega_1&=\omega_1+\omega_2+\ldots+\omega_s\\
\ldotp\quad&\ldotp\quad\ldotp\quad\ldotp\quad\ldotp\quad\ldotp\quad\ldotp\quad\ldotp\quad\ldotp\\
\Omega_{s-1}&=\omega_{s-1}+\omega_s\\
\Omega_s&=\omega_s.
\end{aligned}
\]
% scan: s029L
% !PRINTED 38
\weberpage{38}
And, according to the principle of continuity, these composite curves must be regarded as exceptional curves, whereas their parts (e.g. \(\omega_1\omega_2\ldots\omega_{s-1}\)) will not constitute \textit{whole} exceptional curves, but only \textit{components} of them. Indeed the component \(\omega\) will be in one-to-one correspondence not with the neighborhood of \(O\) (taken in its entirety), but with this neighborhood deprived of the point \(O_1\); and likewise \(\omega_1\) will be in correspondence with the neighborhood of \(O_1\) deprived of the point \(O_2\), etc. We can pass from the surface \(F\) to \(\Phi\) by means of several successive transformations, from \(F\) to \(F_1\), from \(F_1\) to \(F_2\), \(\ldots\), from \(F_{s-1}\) to \(F_s=\Phi\); the first of these transformations will change \(O\) into a curve \(\omega\), the second will change a point \(O_1\) of \(\omega\) into a curve \(\omega_1\), the third will change a point \(O_2\) of \(\omega_1\) (different from the point common to \(\omega\) and \(\omega_1\)) into a curve \(\omega_2\), etc. Therefore each \(\omega_i\) (for \(i<s\)) will be connected at one point with the next one, \(\omega_{i+1}\), and its degree (which is \(-1\) on \(F_i\)) will become \(-1-1=-2\) on \(F_{i+1}\) and on \(\Phi\). Consequently the \textit{composite exceptional curve} \(\Omega=\omega+\omega_1+\ldots+\omega_s\), corresponding to a series of \textit{points} \(O\,O_1\ldots O_s\), \textit{infinitely near and following one another on a linear branch}, will have the \textit{genus}
\[
\varrho=0+0+\ldots+0+(s-1)-(s-1)=0
\]
and the \textit{degree}
\[
\nu=-2-2\ldots-2-1+2(s-1)=-1.
\]
In short, one thus finds again the characters of the irreducible exceptional curves of which the composite curves \(\Omega,\Omega_1,\Omega_2,\ldots,\Omega_s\) appear as limit curves.

This coincidence of characters must still occur when the points \(O,O_1,O_2,\ldots,O_s\) come to follow one another on a superlinear branch; but in the definition of the limit curves we meet with a surprise: the components of the curves \(\Omega_i\) in general become multiple curves!

Let us refer, for simplicity, to the case of three points \(O,O_1,O_2\) following one another on a cuspidal branch: this case can be treated by regarding the cuspidal branch as the limit of a linear branch in which \(O_2\), varying in the neighborhood of \(O_1\), takes up the position of its satellite point. Here it appears that the curve of order \(i-i_1\), the transform of the neighborhood of \(O\) (deprived of \(O_1\)), comes to reduce to a curve \(\omega\) of order \(i-i_1-i_2\) and to the curve \(\omega_2\), of order \(i_2\), corresponding to the neighborhood of \(O_2\); consequently this last curve comes to appear twice in
\[
\omega+\omega_1+2\omega_2.
\]

To understand more clearly the reasons for what happens here, let us carry out three successive transformations of the surface \(F\),
% scan: s029R
% !PRINTED 39
\weberpage{39}

changing first \(O\) into \(\omega\), then again a point \(O_1\) of \(\omega\) into \(\omega_1\), and finally the point \(O_2\) common to \(\omega\) and \(\omega_1\) into \(\omega_2\): the position of the point \(O_2\) thus defined corresponds to the satellite point of \(O_1\) on the surface \(F\). But since \(O_2\) is a double point for the curve \(\omega+\omega_1\), the transformed exceptional curve \(\omega_1\) must be counted as double, in agreement with what has been said above. %
The multiplicities of the composite exceptional curves which correspond to the successive points of a superlinear branch have been determined by \textsc{E. Barber} and \textsc{O. Zariski}\origfootnote{1}{s029R-1}{\textsc{Barber} and \textsc{Zariski}, l. c. Ann. Journ. of Math., 1935.}, who have also pointed out the order of the connections of these curves and have thus illustrated, under a new aspect, the theory of singularities of \textsc{Enriques}\origfootmark{2}{s029R-2}. %
We shall confine ourselves to quoting the following general rule: \textit{the composite exceptional curve \(\Omega\), which corresponds to a series of points \(O_1,\ldots,O_s\) infinitely near to a proper point \(O\) and following one another on a superlinear branch, comprises \(s+1\) components, and the component \(\omega_i^{h_i}\), which occupies place \(i\), appears in it with the multiplicity \(h_i\) equal to the minimum order of the branch issuing from \(O\)\origfootnote{2}{s029R-2}{When one speaks of a branch issuing from the point \(O_1\) (or from \(O_2\), etc.) it is understood that the transformation which changes \(O_1\) into a proper point has first been carried out.} which contains the \(i\) points \(O_1,O_2,\ldots,O_i\).} %

The same rule applies to the composite exceptional curves \(\Omega_1,\Omega_2,\ldots,\Omega_s\)\origfootnote{3}{s029R-3}{Cf. \textsc{Enriques-Chisini}, \textit{Lezioni etc.}, Book I, vol. II.}. Moreover the curve \(\omega_i\) is \textit{connected} not always to the next one, \(\omega_{i+1}\), but --- in general --- to the curve \(\omega_{i+t}\) \textit{which follows it after \(t\) places}, corresponding to the circumstance that \textit{the points \(O_{i+1}\ldots O_{i+t}\) are proximate to \(O_i\)} (as is also evident from inspection of the graphical scheme of the branch \(O_{i+1}\ldots O_{i+t}\)). Let it be added that the said curve \(\omega_i\) turns out to be of \textit{degree} \(-1-t\), as appears from the fact that it represents the neighborhood of a point \(O_i\) from which \(t\) proximate points have been removed. Taking into account the multiplicities which belong to the curves \(\omega_i\) in \(\Omega_i\), the composite exceptional curves \(\Omega,\Omega_1,\ldots,\Omega_s\) always turn out to be of degree \(\nu=-1\) and of genus \(\varrho=0\), like the irreducible exceptional curves of which they are limits.

Finally we shall point out that the set of irreducible exceptional curves, disconnected from one another, whose limit is considered here as \(O_1,O_2,\ldots\) approach \(O\), tends to a curve
\[
\Omega+\Omega_1+\ldots+\Omega_s
\]
which turns out to be composed of the curves \(\omega,\omega_1,\omega_2,\ldots,\omega_s\), each counted with a multiplicity which is the sum of the multiplicities that it possesses with respect to the \(\Omega_i\) of which it forms part. We shall see later on the significance of this observation. %

% scan: s030L
% !PRINTED 40
\weberpage{40}

\edsection{13. Virtual characteristic series.}{13. Virtual characteristic series.}{sec.I.13}

We have explained how the virtual characters are determined, and in particular the degree \(n\) of a curve \(C\) traced on the surface \(F\), even if this is an \textit{isolated curve}, not contained in a system \(|C|\) of dimension \(r>0\). The procedure which leads to evaluating \(n\)\origfootnote{1}{s030L-1}{Cf. \textsc{F. Enriques}, \textit{Introduzione etc.}, 1896: nos. 15 and 16.} also leads to defining, for \(n>0\), the characteristic series of \(C\), even independently of the circumstance that it belongs to a larger system\origfootnote{2}{s030L-2}{\textsc{F. Severi}, \textit{Osservazioni sui sistemi continui di curve.} Atti Acc. Torino, 1904.}.

Indeed, associate with \(C\) an irreducible linear system \(|K|\) such that \(|C+K|\) turns out to be irreducible: the characteristic series of \(C\) will be defined by the curves of \(|C+K|\) which pass through the intersections of a curve \(C\) with a curve \(K\). This series, in fact, turns out to be independent of the choice of \(|K|\), as one sees by comparing the series defined in relation to two systems \(|C+K|\) and \(|C+L|\) and to the system \(|C+K+L|\). Note that, as we shall have occasion to explain later on, it will always be permissible to choose \(|K|\) in such a way that the curves of \(|C+K|\) cut out a complete series on \(C\), and thus one obtains on \(C\) the whole \textit{virtual characteristic series}.

The importance of this series shows itself above all in the study of continuous systems of curves \(C\) not consisting of equivalent curves; then the virtual characteristic series of \(C\) is cut out (wholly or in part) by the curves of the system infinitely near to it. %

\textsc{Note.} -- In this manner the characteristic series of a curve \(C\) belonging to a continuous system presented itself implicitly to \textsc{G. Castelnuovo}, as appears from a communication of his by letter to \textsc{Enriques}\origfootnote{3}{s030L-3}{Cf. Rendic. Circolo Matematico di Palermo, 25 Dec. 1898 (no. 4).}. Later \textsc{Castelnuovo} himself came to point out its significance explicitly, in relation to the (hyperelliptic) surface which represents the pairs of points of a curve of genus two, and drew to this example the attention of \textsc{Severi}, who was thereby led to the definition set forth above, as he himself mentions in a note at the foot of the first page of his paper. %

\clearpage\thispagestyle{plain}
% scan: s030R
% !PRINTED none
\weberpage{[s030R]}

\edchapter{II}{COVARIANT AND INVARIANT SYSTEMS}{Chapter II. COVARIANT AND INVARIANT SYSTEMS}{COVARIANT AND INVARIANT SYSTEMS}

\edsection{1. Jacobian curves.}{1. Jacobian curves.}{sec.II.1}

On the surface \(F\), which we may suppose to be free of any singularity in a hyperspace, let us consider an irreducible linear system \(|C|\) of dimension 2, that is, a \textit{net}. The locus of the points of the surface which are double for a curve of the net is an algebraic curve \(C_j\), which we shall call the \textit{Jacobian} of the net. The curve \(C_j\) can also be defined as the locus of the points of contact of two curves of the net; indeed the pencil determined by these contains a curve which has a double point at the point of contact. In the plane the Jacobian of a net of curves of order \(n\ (>1)\) is a curve of order \(3n-3\); %
and an \(i\)-fold base point of the net has in general the multiplicity \(3i-1\) for it; thus a simple base point \(O\) is, in general, double for the Jacobian, but if the curves of the net have a fixed tangent at \(O\), the Jacobian has a triple point there; with that fixed tangent\origfootnote{1}{s030R-1}{Cf. \textsc{Enriques-Chisini}, Op. cit., Book III, vol. II, chap. I, § 5.}. %
These properties remain unaltered also for nets given on any surface \(F\) whatever, since the determination of the multiplicities of a base point \(O\) for the Jacobian of a net \(|C|\) depends only on the differential characters of the curves \(C\) under consideration and of the surface which contains them; and this surface, to an order of approximation as high as one wishes, can be approximated by an osculating paraboloid at the point \(O\), which in its turn can be represented point by point on the plane. For the Jacobian curves of the nets extracted from a larger linear system \(|C|\) the following theorem holds:

\textit{The Jacobians of the nets contained in one and the same irreducible linear system \(|C|\) of dimension \(r>2\) are equivalent.}

Let us extract from \(|C|\) any two nets and compare their Jacobians, supposing first of all that the two nets have in common
% scan: s031L
% !PRINTED 42
\weberpage{42}
a pencil, and therefore belong together to a linear system \(\infty^3\)\origfootnote{1}{s031L-1}{This system is determined by the common pencil and by two curves taken respectively in the two nets, outside it.}. A generic point \(P\) of the surface \(F\) is double for a curve \(\overline{C}\) of the aforesaid system \(\infty^3\), and \(\overline{C}\), together with the pencil common to the two nets, determines a net of curves \(C\) whose Jacobian passes through \(P\). One sees in this way that the Jacobians of the \(\infty^1\) nets of curves \(C\) contained in the system \(\infty^3\) and having a pencil in common form a pencil (through every point of \(F\) there passes one of them); but this pencil is certainly linear, because it is rational, being in one-to-one correspondence with the pencil of the nets of a linear system \(\infty^3\) which have a pencil in common\origfootnote{2}{s031L-2}{A set to be likened to a pencil of planes of \(S_3\).}. Now if one takes within \(|C|\) (supposed to be of dimension \(r>3\)) two nets of curves having a curve \(\overline{C}\) in common, one can construct an auxiliary net which contains \(\overline{C}\) and has a pencil in common with each of the two nets: the Jacobian of the auxiliary net will turn out to be equivalent to the Jacobians of the two nets; and, by the transitive property of equivalence, these will therefore be equivalent to each other.

Finally, if one has within \(|C|\) two nets without a common curve (the system \(|C|\) having dimension \(r>4\)), one can compare their Jacobians by constructing an auxiliary net which has a curve in common with each of them. The Jacobians of the two nets will therefore turn out to be equivalent, Q.E.D.

\edsection{2. Jacobian system.}{2. Jacobian system.}{sec.II.2}

From what has been proved above it follows that:

\textit{The Jacobians of the nets of curves contained in an irreducible linear system \(|C|\) on the surface \(F\) belong to one and the same complete linear system, which will be called the Jacobian system \(|C_j|\) of \(|C|\).} The dimension of \(|C|\) is supposed for the present to be \(\geq 2\).

The preceding definition is perfect as it stands when the system \(|C|\) has no base points, it being understood that the system of the Jacobians \(C_j\) is to be completed without the imposition of any base point. If instead \(|C|\) has an \textit{\(i\)-fold base point}, it will be understood that this is assigned as a \textit{point of multiplicity \(3i-1\)} to the complete system \(|C_j|\), since, as we have said, \(3i-1\) is in general the multiplicity of the Jacobian of a net contained in \(|C|\).

This holds also if the base points are infinitely near; thus, for example, if \(|C|\) possesses a simple base point with a fixed tangent, that is, two simple base points \(O,O_1\), infinitely near,
% scan: s031R
% !PRINTED 43
\weberpage{43}
the Jacobians of the nets contained in it will actually have at \(O\) a triple point with the fixed tangent \(OO_1\)\origfootnote{1}{s031R-1}{Cf. \S~1.}, but the system of these curves will generally be contained in a larger system with two double base points \(O\) and \(O_1\), which will constitute our complete Jacobian system \(|C_j|\). Even if the enlargement is not actually possible, the singularity constituted by a triple point and by an infinitely near simple point must be regarded virtually as constituted by the two double points \(O\) and \(O_1\). We have supposed so far that the system \(|C|\) was irreducible; but it is clear how the definition is to be extended to the case in which \(|C|\) contains a fixed part \(K\) (which we shall suppose to be simple). Here every point of \(K\) appears as a base point with a fixed tangent for any net contained in \(|C|\); therefore the curve \(K\) will split off 3 times from the Jacobians of all the nets contained in \(|C|\). It must be added 3 times to these Jacobians, and the reducible system thus defined will in general be contained in the larger complete irreducible system \(|C_j|\). %

To tell the truth, what is thus proved is that \(K\) splits off \textit{at least} three times from the Jacobians of the nets of \(|C|\). To recognize that in fact it does not split off more than three times, one can reduce the matter from the surface to the plane by substituting for \(F\) a suitable osculating paraboloid at a point of the curve \(K\), or also a rational surface which has a contact of very high order with \(F\) along the whole curve \(K\)\origfootnote{2}{s031R-2}{Cf. \textsc{Enriques-Chisini}, Book XV, chap. IV, \S~39, vol. II, p. 634 ff.}. One thus recognizes that a curve \(K\) which is added as a \textit{fixed component} of a \textit{net} of irreducible curves \(C\) appears \textit{three times in the Jacobian} of this net; it cannot appear in it more than 3 times unless it is a part of the Jacobian of \(|C|\).

\textsc{Note.} -- One can give a direct analytic verification of the theorem at which we have arrived. To do this one needs to know how a complete linear system is cut out on a surface \(F\) (with ordinary singularities) of ordinary space, namely that ``the surfaces \(\varphi\) adjoint to \(f\), that is, passing simply through the double curve of \(f\), cut out complete systems on \(f\)''\origfootnote{3}{s031R-3}{Cf. the note to \S~4 and chap. III.}. For simplicity of exposition we shall suppose that \(f\) is a surface of order \(n\) entirely free of singularities, so that its complete intersections with the totality of the surfaces of a given order \(m\) form a complete system.

Given on \(f\) a net of irreducible curves \(|C|\) which are not complete intersections, to which a fixed component \(K\) is added, one can always suppose that the net \(|C|\) is residual to a curve \(L\) not containing \(K\) as a part, in such a way that the curves \(C+L\) are complete
% scan: s032L
% !PRINTED 44
\weberpage{44}
intersections with surfaces \(\varphi_i\) of order \(m\). Hence the net \(|K+C|\) is partially contained in the net represented by an equation of the type:
\[
\varphi(x_1x_2x_3x_4)[\lambda_1\varphi_1(x_1x_2x_3x_4)+\lambda_2\varphi_2(x_1x_2x_3x_4)+\lambda_3\varphi_3(x_1x_2x_3x_4)]=0,
\]
where, according to our hypotheses, \(\varphi_1,\varphi_2,\varphi_3\) will not vanish simultaneously on \(K\), which instead forms part of \(\varphi=0\), whose order will be denoted by \(s\). Now we have to write down the equation of the Jacobian of the net thus represented, and to recognize that the curve \(\varphi=0\), and in particular \(K\), which constitutes a part of it, appears in it 3 times as a component. The points of the Jacobian are points of contact between a surface of the system cutting out the net and \(f\). Writing down the conditions for contact, one obtains a system of linear equations, for whose compatibility the determinant must vanish
\[
D(x_1x_2x_3x_4)=
\begin{vmatrix}
\dfrac{\partial(\varphi\varphi_1)}{\partial x_1} &
\dfrac{\partial(\varphi\varphi_2)}{\partial x_1} &
\dfrac{\partial(\varphi\varphi_3)}{\partial x_1} &
\dfrac{\partial f}{\partial x_1} \\[6pt]
\dfrac{\partial(\varphi\varphi_1)}{\partial x_2} &
\dfrac{\partial(\varphi\varphi_2)}{\partial x_2} &
\dfrac{\partial(\varphi\varphi_3)}{\partial x_2} &
\dfrac{\partial f}{\partial x_2} \\[6pt]
\dfrac{\partial(\varphi\varphi_1)}{\partial x_3} &
\dfrac{\partial(\varphi\varphi_2)}{\partial x_3} &
\dfrac{\partial(\varphi\varphi_3)}{\partial x_3} &
\dfrac{\partial f}{\partial x_3} \\[6pt]
\dfrac{\partial(\varphi\varphi_1)}{\partial x_4} &
\dfrac{\partial(\varphi\varphi_2)}{\partial x_4} &
\dfrac{(\partial\varphi\varphi_3)}{\partial x_4} &
\dfrac{\partial f}{\partial x_4} %
\end{vmatrix}
=0
\]

Now, after carrying out the derivatives of \(\varphi\varphi_1,\varphi\varphi_2,\varphi\varphi_3\), one is led to decompose the determinant \(D\) into a sum of determinants: leaving aside those determinants which turn out to be identically zero, the others present the common factor \(\varphi^2\), on removing which there remain four determinants, one of which still contains the factor \(\varphi\). We need to prove that the sum of the other 3 determinants vanishes (simply) at the points of the common curve \(\varphi=0\). Indeed this sum is easily transformed into a single determinant. Let us write:
\begin{gather*}
\Phi=\varphi_1
\begin{vmatrix}
\dfrac{\partial\varphi}{\partial x_1} &
\dfrac{\partial\varphi_2}{\partial x_1} &
\dfrac{\partial\varphi_3}{\partial x_1} &
\dfrac{\partial f}{\partial x_1} \\[6pt]
\dfrac{\partial\varphi}{\partial x_2} &
\dfrac{\partial\varphi_2}{\partial x_2} &
\dfrac{\partial\varphi_3}{\partial x_2} &
\dfrac{\partial f}{\partial x_2} \\[6pt]
\dfrac{\partial\varphi}{\partial x_3} &
\dfrac{\partial\varphi_2}{\partial x_3} &
\dfrac{\partial\varphi_3}{\partial x_3} &
\dfrac{\partial f}{\partial x_3} \\[6pt]
\dfrac{\partial\varphi}{\partial x_4} &
\dfrac{\partial\varphi_2}{\partial x_4} &
\dfrac{\partial\varphi_3}{\partial x_4} &
\dfrac{\partial f}{\partial x_4}
\end{vmatrix}
\\[12pt]
% scan: s032R
% !PRINTED 45
\text{\weberpage{45}}
+\varphi_2
\begin{vmatrix}
\dfrac{\partial\varphi_1}{\partial x_1} &
\dfrac{\partial\varphi}{\partial x_1} &
\dfrac{\partial\varphi_3}{\partial x_1} &
\dfrac{\partial f}{\partial x_1} \\[6pt]
\dfrac{\partial\varphi_1}{\partial x_2} &
\dfrac{\partial\varphi}{\partial x_2} &
\dfrac{\partial\varphi_3}{\partial x_2} &
\dfrac{\partial f}{\partial x_2} \\[6pt]
\dfrac{\partial\varphi_1}{\partial x_3} &
\dfrac{\partial\varphi}{\partial x_3} &
\dfrac{\partial\varphi_3}{\partial x_3} &
\dfrac{\partial f}{\partial x_3} \\[6pt]
\dfrac{\partial\varphi_1}{\partial x_4} &
\dfrac{\partial\varphi}{\partial x_4} &
\dfrac{\partial\varphi_3}{\partial x_4} &
\dfrac{\partial f}{\partial x_4}
\end{vmatrix}
\\[12pt]
+\varphi_3
\begin{vmatrix}
\dfrac{\partial\varphi_1}{\partial x_1} &
\dfrac{\partial\varphi_2}{\partial x_1} &
\dfrac{\partial\varphi}{\partial x_1} &
\dfrac{\partial f}{\partial x_1} \\[6pt]
\dfrac{\partial\varphi_1}{\partial x_2} &
\dfrac{\partial\varphi_2}{\partial x_2} &
\dfrac{\partial\varphi}{\partial x_2} &
\dfrac{\partial f}{\partial x_2} \\[6pt]
\dfrac{\partial\varphi_1}{\partial x_3} &
\dfrac{\partial\varphi_2}{\partial x_3} &
\dfrac{\partial\varphi}{\partial x_3} &
\dfrac{\partial f}{\partial x_3} \\[6pt]
\dfrac{\partial\varphi_1}{\partial x_4} &
\dfrac{\partial\varphi_2}{\partial x_4} &
\dfrac{\partial\varphi}{\partial x_4} &
\dfrac{\partial f}{\partial x_4}
\end{vmatrix};
\end{gather*}\edFlush
interchanging in the first determinant the \(1^{\mathrm{st}}\) column with the \(3^{\mathrm{rd}}\), and in the second the \(2^{\mathrm{nd}}\) with the \(3^{\mathrm{rd}}\), and taking account of the signs, there appears here the expansion of the determinant of the \(5^{\mathrm{th}}\) order:
\[
\Phi=
\begin{vmatrix}
\varphi_1 & \varphi_2 & \varphi_3 & 0 & 0 \\[6pt]
\dfrac{\partial\varphi_1}{\partial x_1} &
\dfrac{\partial\varphi_2}{\partial x_1} &
\dfrac{\partial\varphi_3}{\partial x_1} &
\dfrac{\partial\varphi}{\partial x_1} &
\dfrac{\partial f}{\partial x_1} \\[6pt]
\dfrac{\partial\varphi_1}{\partial x_2} &
\dfrac{\partial\varphi_2}{\partial x_2} &
\dfrac{\partial\varphi_3}{\partial x_2} &
\dfrac{\partial\varphi}{\partial x_2} &
\dfrac{\partial f}{\partial x_2} \\[6pt]
\dfrac{\partial\varphi_1}{\partial x_3} &
\dfrac{\partial\varphi_2}{\partial x_3} &
\dfrac{\partial\varphi_3}{\partial x_3} &
\dfrac{\partial\varphi}{\partial x_3} &
-\dfrac{\partial f}{\partial x_3} \\[6pt] %
\dfrac{\partial\varphi_1}{\partial x_4} &
\dfrac{\partial\varphi_2}{\partial x_4} &
\dfrac{\partial\varphi_3}{\partial x_4} &
\dfrac{\partial\varphi}{\partial x_4} &
\dfrac{\partial f}{\partial x_4}
\end{vmatrix}.
\]

Subtracting from the \(1^{\mathrm{st}}\) row multiplied by \(m\) the \(2^{\mathrm{nd}}\) multiplied by \(x_1\), the \(3^{\mathrm{rd}}\) by \(x_2\), the \(4^{\mathrm{th}}\) by \(x_3\), the \(5^{\mathrm{th}}\) by \(x_4\), and taking account of Euler’s theorem, one can replace it by
\[
0,\ 0,\ 0,\ -s\varphi,\ -nf %
\]
since by \(n,m,s\) we have denoted respectively the orders of \(f\), of the \(\varphi_i\) and of \(\varphi\).

It is clear that the determinant in which this row appears vanishes when simultaneously \(\varphi=0\) and \(f=0\), because, when these terms vanish, the first row becomes identically zero, Q.E.D.

% scan: s033L
% !PRINTED 46
\weberpage{46}
Let us add that \(K\) cannot split off more than three times from the Jacobian of the net; or at least, if this happens, \(K\) forms part of the Jacobian of the net stripped of the fixed part \(K\). %

Indeed a further splitting off of \(K\) would depend on a contact between the surface \(f\) and \(\Phi\) along the whole curve \(K\). Now let us suppose that at a point \(P\) of this curve the tangent planes to \(f\) and to \(\Phi\) coincide. We must have
\[
\varrho\frac{\partial f}{\partial x_i}=\frac{\partial\Phi}{\partial x_i},
\qquad (i=1,2,3,4).
\]
Let us calculate \(\dfrac{\partial\Phi}{\partial x_i}\). We can differentiate by columns; in this way the derivative will be expressed as a sum of determinants obtained by substituting in one of the columns the derivatives with respect to \(x_i\). Now it is clear that, since we are at a point where \(\varphi=0,f=0\), the determinants obtained by differentiating the first three columns are zero. Hence there remains:
\[
\begin{aligned}
m\left(\frac{\partial\Phi}{\partial x_i}\right)_P
&=
\begin{vmatrix}
0 & 0 & 0 & -s\dfrac{\partial\varphi}{\partial x_i} & 0 \\[6pt]
\dfrac{\partial\varphi_1}{\partial x_1} &
\dfrac{\partial\varphi_2}{\partial x_1} &
\dfrac{\partial\varphi_3}{\partial x_1} &
\dfrac{\partial^2\varphi}{\partial x_1\partial x_i} &
\dfrac{\partial f}{\partial x_1} \\[6pt]
\cdots & \cdots & \cdots & \cdots & \cdots \\[6pt]
\dfrac{\partial\varphi_1}{\partial x_4} &
\dfrac{\partial\varphi_2}{\partial x_4} &
\dfrac{\partial\varphi_3}{\partial x_4} &
\dfrac{\partial^2\varphi}{\partial x_4\partial x_i} &
\dfrac{\partial f}{\partial x_4}
\end{vmatrix}
\\[12pt]
&+
\begin{vmatrix}
0 & 0 & 0 & 0 & -n\dfrac{\partial f}{\partial x_i} \\[6pt]
\dfrac{\partial\varphi_1}{\partial x_1} &
\dfrac{\partial\varphi_2}{\partial x_1} &
\dfrac{\partial\varphi_3}{\partial x_1} &
\dfrac{\partial\varphi}{\partial x_1} &
\dfrac{\partial^2 f}{\partial x_1\partial x_i} \\[6pt]
\cdots & \cdots & \cdots & \cdots & \cdots \\[6pt]
\dfrac{\partial\varphi_1}{\partial x_4} &
\dfrac{\partial\varphi_2}{\partial x_4} &
\dfrac{\partial\varphi_3}{\partial x_4} &
\dfrac{\partial\varphi}{\partial x_4} &
\dfrac{\partial^2 f}{\partial x_4\partial x_i}
\end{vmatrix}.
\end{aligned}
\]

Expanding the two determinants along the first row, one has:
\begin{gather*}
m\left(\frac{\partial\Phi}{\partial x_i}\right)_P
=\varrho m\left(\frac{\partial f}{\partial x_i}\right)_P
\\[10pt]
=s\frac{\partial\varphi}{\partial x_i}
\begin{vmatrix}
\dfrac{\partial\varphi_1}{\partial x_1} &
\dfrac{\partial\varphi_2}{\partial x_1} &
\dfrac{\partial\varphi_3}{\partial x_1} &
\dfrac{\partial f}{\partial x_1} \\[6pt]
\cdots & \cdots & \cdots & \cdots \\[6pt]
\dfrac{\partial\varphi_1}{\partial x_4} &
\dfrac{\partial\varphi_2}{\partial x_4} &
\dfrac{\partial\varphi_3}{\partial x_4} &
\dfrac{\partial f}{\partial x_4}
\end{vmatrix}
\\[10pt]
% scan: s033R
% !PRINTED 47
\text{\weberpage{47}}
-n\frac{\partial f}{\partial x_i}
\begin{vmatrix}
\dfrac{\partial\varphi_1}{\partial x_1} &
\dfrac{\partial\varphi_2}{\partial x_1} &
\dfrac{\partial\varphi_3}{\partial x_1} &
\dfrac{\partial\varphi}{\partial x_1} \\[6pt]
\cdots & \cdots & \cdots & \cdots \\[6pt]
\dfrac{\partial\varphi_1}{\partial x_4} &
\dfrac{\partial\varphi_2}{\partial x_4} &
\dfrac{\partial\varphi_3}{\partial x_4} &
\dfrac{\partial\varphi}{\partial x_4}
\end{vmatrix}
\qquad (i=1,2,3,4).
\end{gather*}\edFlush

Denoting briefly the first determinant of the right-hand side by \(A\), and the second by \(B\), this equation can be written
\[
m\varrho\frac{\partial f}{\partial x_i}
=s\frac{\partial\varphi}{\partial x_i}A
-n\frac{\partial f}{\partial x_i}B,
\qquad (i=1,2,3,4),
\]
which a priori offers two possibilities:

1) \(\dfrac{\partial f}{\partial x_i}=\lambda\dfrac{\partial\varphi}{\partial x_i}\), that is, the surfaces \(f=0\) and \(\varphi=0\) touch each other, which is contrary to the hypotheses; or

2) \(A=0\), and then, as one recognizes at once from the structure of the determinant \(A\), the point \(P\) lies on the Jacobian of the net.

The algebraic procedure which we have developed here can be illustrated by some simple geometric considerations, which we shall confine ourselves to sketching, referring to the simplest case, in which \(f\) is a surface of order \(n\) without singularities in ordinary space, \(\varphi_1,\varphi_2,\varphi_3\) are three linearly independent surfaces of order \(m>n\), and \(\varphi\) is a surface of order \(s\).

Let us observe that the surface \(f\), taken together with a surface \(\psi\) of order \(m-n\), can be combined linearly with \(\varphi_1,\varphi_2,\varphi_3\); one will then construct the Jacobian surface of the linear system \(\infty^3\) thus determined, that is, the locus of the double points of the surfaces of the system, which will be a surface \(D\) of order \((4m-4)\)\origfootnote{1}{s033R-1}{Here we have an immediate extension of the Jacobian of a net of plane curves. Cf. \textsc{Enriques-Chisini}, Op. cit., Book III, chap. I, \S~5, vol. II, p. 31.}. Now \(D\) will cut \(f\) in a curve of order \((4m-4)n\), which will turn out to be composed of the Jacobian of the net of curves \(|C|\) cut out by
\[
\lambda_1\varphi_1+\lambda_2\varphi_2+\lambda_3\varphi_3=0
\]
and of the curve \(L\), the intersection of \(f\) and \(\psi\); since the latter is of order \(n(m-n)\), the Jacobian of the net \(|C|\) will turn out to be of order \(n(3m+n-4)\).

This being so, let there be added to the net
\[
\lambda_1\varphi_1+\lambda_2\varphi_2+\lambda_3\varphi_3=0
\]
a surface \(\varphi\) of order \(s\), cutting \(f\) in a curve \(K\) of order \(ns\); adding \(\varphi\) also to \(f\), one will have a linear system \(\infty^3\) of sur%
% scan: s034L
% !PRINTED 48
\unskip\weberpage{48}%
faces of order \(m+s\) in which \(\varphi\) appears as a fixed part. Therefore this \(\varphi\) will split off 4 times from the Jacobian \(D\) of the system; its intersection with \(f\) (of order \(ns\)) will appear once in the curve of intersection of \(f\) with \(\psi+\varphi\), and hence will appear 3 times in the Jacobian curve of the net \(|K+C|\), which is now of order
\[
n[3(m+s)+n-4].
\]

\edsection{3. Fundamental theorem.}{3. Fundamental theorem.}{sec.II.3}

From the preceding observations there follows the

\textsc{Fundamental theorem:} \textit{Given on the surface \(F\) two irreducible linear systems \(|C|\) and \(|L|\), which we always suppose to be virtually free of base points and of dimension not less than \(2\), the Jacobian system \(|C+L|_j\) of their sum is expressed by}
\[
|C+L|_j=|C_j+3L|=|3C+L_j|.
\]

Indeed the system \(|C+L|\) contains the system \(|C|+L\), with a fixed curve \(L\), and the system \(C+|L|\), with a fixed curve \(C\), whose complete Jacobian systems are respectively the systems
\[
|C_j+3L|\quad\text{and}\quad |L_j+3C|,
\]
which therefore coincide in the system \(|C+L|_j\).

On the basis of the fundamental theorem one can now define the complete Jacobian system of a linear system \(|C|\) which has fewer than two dimensions, and which therefore reduces to a pencil, or even to an isolated curve \(C\). Indeed it will suffice to add to \(C\) an auxiliary system \(|L|\), at least \(\infty^2\), and to take
\[
|C_j|=|3C+L_j-3L|.
\]

The fundamental theorem thus extends to the Jacobians of linear systems which are reducible in any way whatever. Only one caution is needed, concerning the case in which one considers linear systems \(|C|\) and \(|L|\) with \textit{assigned base points} on the surface \(F\). If \(|C|\) possesses an \(i\)-fold base point \(O\) which is not a base point for \(|L|\), and \(|L|\) possesses an \(s\)-fold base point \(O'\) which is not a base point for \(|C|\), the Jacobian system \(|C+L|_j\) will have at \(O\) the (virtual) multiplicity \(3i-1\), and at \(O'\) the multiplicity \(3s-1\), whereas \(O\) will have the multiplicity \(3i\) for \(|3C+L_j|\), and \(O'\) the multiplicity \(3s\) for \(|3L+C_j|\); hence the equivalence between the two systems \(|C_j+3L|\) and \(|3C+L_j|\) will hold only when one adds to the two systems respectively the neighborhoods of the points \(O\) and \(O'\). That is, one will have:
\[
|3C+L_j+O|=|3L+C_j+O'|;
\]
% scan: s034R
% !PRINTED 49
\weberpage{49}
and the equivalence thus made precise has the meaning which results from our conventions: when the surface \(F\) is changed by transforming \(O\) and \(O'\) into exceptional curves, these are to be added respectively to the transforms of the two systems \(|C_j+3L|\) and \(|L_j+3C|\).

After this, let us return to the case, to which we shall generally refer in what follows, in which one is dealing with \textit{complete systems virtually free of base points} on the surface \(F\). The equivalence \(|C_j+3L|=|L_j+3C|\) can be reduced to %
\[
|C_j-3C|=|L_j-3L|=|K|,
\]
which in words means: \textit{if the Jacobian \(|C_j|\) of a system \(|C|\) contains the triple of the latter, then the Jacobian \(|L_j|\) of any other system \(|L|\) will also contain the triple of \(|L|\), and the system \(|K|\) obtained by the subtraction will turn out to be independent of the system from which one starts.} This system \(|K|\) will be defined as the (\textit{impure}) \textit{canonical system} of the surface \(F\).

But the symbolic equivalence which we have written does not lose its meaning even if \(|C_j|\) does not contain \(|3C|\), that is, under the hypothesis that the difference system \(|C_j-3C|\) has only a virtual existence. In any case the meaning of the equivalence between virtual systems is given here by the relation which expresses the fundamental theorem on Jacobians; there is nothing but a simple difference of language.

It is also convenient to reduce the equivalence of virtual systems indicated above to a simpler form. For this purpose we shall define the \textit{linear system \(|C'|\) adjoint to \(|C|\)} by means of the relation
\[
|C'|=|C_j-2C|,
\]
observing that \(|C'|\) (unlike \(|K|\)) will in general have an \textit{effective existence} for the systems \(|C|\) traced on any surface whatever: this is verified already in the plane for linear systems of curves of order \(m\) and of genus \(\pi>0\), since the curves \(C'\) are then the curves of order \(m-3\) passing \(i-1\) times through every \(i\)-fold base point of \(|C|\); and the fact will be verified later on for linear systems whose genus exceeds a certain limit.

With the introduction of the adjoint system, the canonical system is given by:
\[
|K|=|C'-C|=|L'-L|\quad .\,.\,.\,.
\]
and the equivalence of effective or virtual systems thus expressed means
\[
|C'+L|=|L'+C|
\]
which is the \textit{fundamental relation for adjoint systems} (without base points), a simple transform of the relation given above for Jacobian systems. %

\printersignature{\textsc{Enriques F.} - \textit{Superficie algebriche.}\quad 4}

% scan: s035L
% !PRINTED 50
\weberpage{50}
\edsection{4. Canonical curves.}{4. Canonical curves.}{sec.II.4}

We have seen how, starting from a linear system \(|C|\) given on a surface \(F\), one constructs the adjoint linear system \(|C'|\) and hence the canonical system \(|K|=|C'-C|\), which may have an effective or a virtual existence, and it has been recognized that this last system turns out to be independent of the choice of the system \(|C|\), supposed free of base points on \(F\). Even in the case where \(|C|\) possesses an \(i\)-fold base point \(O\), simple for \(F\), the system \(|C'-C|\) can still be regarded as identical with \(|K|\), if one disregards the infinitesimal curve constituting the neighborhood of \(O\), which, being \((i-1)\)-fold for \(|C'|\), would have to be added to the curves \(K\). The independence of \(|K|\) from the choice of \(|C|\) means that the canonical system \(|K|\) has an invariant meaning under birational transformations of the surface.

To explain the matter better, consider two surfaces \(F\) and \(\Phi\) of ordinary space, of orders \(n\) and \(m\) respectively, endowed with ordinary singularities, and suppose that a birational correspondence holds between them, adding at first the hypothesis that this correspondence has no fundamental points on either surface. The plane sections of the surface \(F\) will form a linear system \(|C|\) whose transform has no base points on \(\Phi\), and the plane sections \(L\) of \(\Phi\) will likewise form a linear system on \(\Phi\) whose transform has no base points on \(F\). Now one can construct the canonical system \(|K|\) on \(F\) either starting from the system \(|C|\) of the plane sections or starting from the system \(|L|\), and one will have
\[
|K|=|C'-C|=|L'-L|;
\]
the same system will be constructed on \(\Phi\), either starting from the system \(|L|\) of the plane sections or starting from the system \(|C|\), one again having %
\[
|K|=|L'-L|=|C'-C|.
\]
Therefore the curves of the canonical system \(|K|\) constructed on \(F\) starting from the system of plane sections are transformed into the curves of the canonical system of \(\Phi\) constructed starting from the system \(|L|\) of plane sections. It remains to see what this construction amounts to.

The canonical system can be constructed on \(F\) by constructing first of all the Jacobian system of \(|C|\). To this end we shall observe that a star of plane sections of \(F\), having as its center an arbitrary point \(O\) of space, constitutes a net whose Jacobian is the intersection of \(F\) with the polar surface of \(O\), of order \(n-1\), that is, its further intersection outside the nodal curve.

The Jacobians of the nets contained in the system \(|C|\) of plane sections will thus form a linear system \(\infty^3\) cut out on \(F\) by the polars, outside the nodal curve; but this system will have as simple base
% scan: s035R
% !PRINTED 51
\weberpage{51}
points the cuspidal points on the said nodal curve. That these points must be regarded as ``simple points'' of the surface, and therefore, according to our conventions, must be removed in order to obtain the complete Jacobian system \(|C_j|\), can easily be recognized. Indeed, if we refer to a transformed surface on which the nodal curve of \(F\) is represented by the curve of the neutral pairs with respect to the system \(|C|\), the cuspidal points will give on this curve base points of contact for a net of curves \(C\), and hence double base points for a pencil of curves \(C\): they are therefore points which belong simply to the Jacobians of all the nets contained in the system \(\infty^3\) of the curves \(C\).

This being so, the complete Jacobian system \(|C_j|\), free, at least virtually, of base points, will contain the whole system of the curves of section of \(F\) with the surfaces of order \(n-1\) passing through its double curve (\textit{adjoint} surfaces). If now 3 planes are removed from these surfaces, we shall have that the surfaces of order \(n-4\) (supposed to exist) passing through the double curve of \(F\) cut out canonical curves on \(F\).

Similarly the adjoint surfaces of order \(m-4\), passing through the double curve of \(\Phi\), will cut out canonical curves on \(\Phi\).

Hence the invariance of the canonical system under birational transformations will mean that ``the curves of section of \(F\) with the adjoint surfaces of order \(n-4\) are changed, by the birational transformation, into curves equivalent to the sections of \(\Phi\) with its adjoint surfaces of order \(m-4\)''.

The theorem takes on a more expressive form if one anticipates a result which will be established in the following chapter, namely that ``the adjoint surfaces of any order, and in particular those of order \(n-4\), cut out on the surface \(F\) of order \(n\) a complete linear system'', that is, in this case, the whole canonical system \(|K|\).

We shall therefore have: \textit{the} (canonical) \textit{curves of intersection of the surface \(F\) with its adjoints of order \(n-4\)} (outside the nodal curve) \textit{are changed, by a birational transformation, into the sections of \(\Phi\) with its adjoint surfaces of order \(m-4\).}

To tell the truth, it has been supposed that the birational correspondence between \(F\) and \(\Phi\) is without exceptions, that is, that the system \(|L|\) has no base points on \(F\) and \(|C|\) has no base points on \(\Phi\). But if, for example, \(|L|\) has an \(i\)-fold base point \(O\) on \(F\), the theorem will still hold with this modification, namely that: the \textit{exceptional curve} \(\omega\) corresponding to the point \(O\) of \(F\) will be added to the transforms of the curves \(|K|\), appearing as a \textit{fixed part of the impure canonical system constituted by the sections of \(\Phi\) with its adjoint surfaces of order \(m-4\).}

\textsc{Note.} -- It has been supposed that \(F\) and \(\Phi\) possess only ordinary singularities (that is, a double curve with triple points); but the
% scan: s036L
% !PRINTED 52
\weberpage{52}
theorem extends also if this restriction is dropped, provided that the \textit{adjoint surfaces} to a given surface are defined in relation to the higher singularities which it possesses. It suffices here to mention that an ordinary \(i\)-fold curve of \(F\) must be \((i-1)\)-fold for the adjoint surfaces, and a conical point of order \(s\), outside the multiple curve, must in general have the multiplicity \(s-2\) for these adjoints. But of this later.

\edsection{5. Properties of canonical curves.}{5. Properties of canonical curves.}{sec.II.5}

The canonical curves on a surface \(F\) (I mean the \textit{impure} canonical curves, sections of the surface \(F\) of order \(n\) with the adjoint surfaces of order \(n-4\), outside the double curve) enjoy a property which we wish to recognize explicitly here, and which, under the simplest conditions, also serves to characterize them, as we shall see in the following chapter.

Consider on \(F\) an irreducible linear system without base points, \(|C|\), of dimension \(r\geq 2\), of genus \(\pi\) and of degree \(n\). The Jacobian of a net of curves \(C\) cuts out on a curve \(C\) the group of the \(2n+2\pi-2\) double points of the characteristic \(g_n^1\), and this group belongs to the series which is the sum of the canonical series and of twice the \(g_n^1\). Removing \(2C\) from \(|C_j|\), one will have curves \(C'\) adjoint to \(|C|\) cutting out canonical groups on \(C\); and removing one more \(C\), one will have that the canonical curves \(K\) cut out on \(C\) groups residual to the characteristic series.

We shall therefore state that:

\textit{The adjoint curves \(C'\) cut out canonical groups on the curves of the system \(|C|\). The canonical curves \(K\) cut out on the curves \(C\) groups residual to the characteristic series.}

It follows from this that the adjoints \(C'\) cut a curve \(C\) of virtual genus \(\pi\) in \(2\pi-2\) points, and the canonical curves \(K\), effective or virtual, cut it in
\[
i=2\pi-2-n
\]
points.

\textit{This equality holds for any curve \(C\) whatever, even isolated, irreducible or reducible,} always referring to the virtual characters, defined on the supposition that \(C\) has no assigned base points on \(F\). Indeed, add to \(C\) a suitable curve \(L\) of genus \(\varrho\) and degree \(\nu\) which meets it in \(s\) points, in such a way that the complete linear system \(|C+L|\) is irreducible and of dimension \(>2\); it will suffice to take as the system \(|L|\) the system cut out by the surfaces of sufficiently high order passing through \(C\). Now the curves \(C'+L\) belonging to the system \(|(C+L)'|\) will cut a curve \(C+L\) in
\[
2(\pi+\varrho+s-1)-2
\]
% scan: s036R
% !PRINTED 53
\weberpage{53}
points; removing their intersections with \(L\), which are \(2\varrho-2+s\), and then the \(s\) intersections with \(L\), there will result
\[
(CC')=2\pi-2.
\]
Similarly the virtual number of intersections of \(C\) and \(K\) is always \((CK)=2\pi-2-n\).

\textsc{Note.} – Let us add that for an irreducible curve \(C\), without multiple points on \(F\), the group cut out on it by an adjoint \(C'\) will always be a canonical group, and the group cut out by a curve \(K\) will be residual to the characteristic series (effective or virtual). It will suffice to prove the assertion for a curve \(C\) which forms part of a composite curve \(C+L\) belonging to an irreducible net. Taking up again the reasoning given above, one will now find that a curve of
\[
|C_j|=|(C+L)_j-3L|
\]
cuts out on \(C\) a group belonging to the series which is the sum of the canonical series and of twice the virtual characteristic series of \(C\) itself.

\edsection{6. Surface genus and linear genus.}{6. Surface genus and linear genus.}{sec.II.6}

The invariance of the canonical system under birational transformations of the surface implies that its characters likewise have an invariant meaning with respect to such transformations. It is necessary, however, to distinguish:

1) \textit{absolute invariants}, which do not change under any transformation whatever, and

2) \textit{relative invariants}, which keep the same value under birational transformations without exceptions, but are altered by transformations possessing fundamental points, which change a (simple) point of one surface into an exceptional curve of the other.

The dimension of the canonical system, or better the number of linearly independent canonical curves, gives us precisely an absolute invariant, which is called the \textit{geometric surface genus} (\(p\) or \(p_g\)), or simply the ``genus'', of the surface.

Once it has been proved that the complete canonical system is cut out on a surface \(F\) of order \(n\) of ordinary space by its adjoints \(\varphi_{n-4}\) of order \(n-4\) (cf. Chap. III), one has that the genus of a general surface of order \(n\) (which is the number of linearly independent \(\varphi_{n-4}\)) is
\[
p=\frac{(n-3)(n-2)(n-1)}{6}.
\]

% scan: s037L
% !PRINTED 54
\weberpage{54}
Hence the genus of the plane, of quadrics and of cubic surfaces will be
\[
p=0.
\]
This is verified directly for the plane by observing that in it the system \(|C|\) of general curves of order \(n\geq 3\) contains its adjoint \(|C'|\), consisting of the curves of order \(n-3\); thus, since \(|C|\) is larger than \(|C'|\), it cannot happen that \(|C'|\) contains \(|C|\): \(C'-C\) has no effective existence, which means \(p=0\).

From the fact that the genus of the plane is zero it follows that all \textit{rational surfaces} (representable point by point on the plane) likewise have \textit{genus zero}. Among these are precisely the quadrics and the cubics (at least those which are not cones) mentioned above.

Another family of surfaces with \textit{zero} (geometric) \textit{genus} consists of \textit{the ruled surfaces}, and hence also of their transforms. The assertion is justified by observing that a generator of a ruled surface \(F\) cannot have an adjoint system, since this would have to cut out the canonical series on the generator. In general the system \(|C'|\) adjoint to a linear system \(|C|\) on the ruled surface cannot contain \(|C|\) because, the curves \(C\) cutting the generators in \(n\) points, the curves \(C'\) cut them in \(n-2\) points (cf. Chap. III, \S{} 3).

Returning to arbitrary surfaces, of genus \(p\geq 0\), the genus of the canonical curves (effective or virtual) furnishes a relative invariant, called the \textit{linear genus} \(p^{(1)}\), or also the \textit{relative} linear genus, which diminishes by 1 for every point of the surface \(F\) which is changed into an exceptional curve and, conversely, increases by 1 for every exceptional curve which is changed into a point. Indeed we have seen that, when a point \(O\) of the surface \(F\) is changed into an exceptional curve \(\omega\) of \(\Phi\), this \(\omega\) is added to the curves \(\bar{K}\), transforms of the canonical curves \(K\); now such a curve is not connected with the said \(K\), that is, it has with them a virtual number of intersections equal to zero, since the point \(O\) is not — at least virtually — a base point of \(|K|\), the said system \(|K|\) having no assigned base point on \(F\). Therefore, on adding to the \(\bar{K}\) the curve \(\omega\) of genus zero and of degree \(-1\), disconnected from them, the genus of these curves diminishes by 1.

Here there becomes apparent the advisability of considering, at least for surfaces of genus \(p>0\), where the canonical system has an effective existence, besides the relative linear genus \(p^{(1)}\), also the genus of the \textit{pure canonical system}, stripped of its fixed exceptional components, which will then be the genus of the canonical curves on a transformed surface free of exceptional curves, as will be made precise in the following \S{} 7; we shall thus obtain an absolute invariant, which we shall call the \textit{absolute linear genus} of the given surface. Its value will be
\[
\bar{p}^{(1)}=p^{(1)}+\sigma
\]
% scan: s037R
% !PRINTED 55
\weberpage{55}
denoting by \(\sigma\) the number of exceptional curves belonging to the surface.

We reserve for the next section the task of making precise that these curves are always exceptional curves of the \(1^{\mathrm{st}}\) kind (at least for \(p>0\)), not connected with one another, examining in particular the case of reducible exceptional curves.

Like the genus of the canonical system, its degree also furnishes a relative invariant of the surface, which we may denote by \(p^{(2)}\). When, through a transformation of the surface \(F\), an exceptional curve \(\omega\) of degree \(-1\), not connected with the curves \(K\), comes to be added to the \(\bar{K}\), \(p^{(2)}\) likewise diminishes by 1. And here too, besides the degree of the impure canonical system, we can define the degree of the pure canonical system, which will constitute an absolute invariant, defined at least for surfaces of genus \(p>0\). But \(p^{(2)}\) does not give us a new invariant character compared with the linear genus. Indeed \textit{there holds between the genus and the degree of the impure canonical system the relation.} %
\[
p^{(2)}=p^{(1)}-1.
\]
To prove it, it suffices to observe that the system adjoint to the canonical system \(|K|\) is
\[
|K'|=|2K|;
\]
hence the number of intersections of \(K\) with \(2K\) is \(2p^{(1)}-2\), and therefore
\[
(KK)=p^{(1)}-1.
\]
In the case \(p>0\) (and in all the cases in which it will be possible to define the absolute characters under discussion, that is, when the surface is transformable into another free of exceptional curves) \textit{the preceding relation still holds between the genus and the degree of the pure canonical system, which are}
\[
p^{(1)}+\sigma\quad\text{and}\quad p^{(2)}+\sigma.
\]

\textsc{Note.} – It will turn out later that the only surfaces for which it is not possible to define the absolute linear genus (by referring to a transform free of exceptional curves) are those which belong to the family of ruled surfaces (rational or not). Even for these, however, one succeeds in constructing an absolute invariant, namely the one given — according to \textsc{Castelnuovo}\origfootnote{1}{s037R-1}{\textit{Sul genere lineare di una superficie e sulla classificazione a cui esso dà luogo.} Rendic. Acc. Lincei, vol. VI, 1897. ``Memorie scelte''; p. 443.} — by the maximum of the relative linear genus, taken over all the possible transforms of the given surface.

Here it is opportune to calculate the linear genus of the plane and of the irrational ruled surfaces of genus \(p\) (\(>0\)), which correspond precisely to that maximum.

% scan: s038L
% !PRINTED 56
\weberpage{56}
The linear genus of the plane will be calculated by referring to a linear system \(|C|\) of plane quartics, of genus \(\pi=3\) and of degree \(n=16\): the genus \(\pi'=0\) of the adjoint system, which is the net of lines, is expressed by the formula
\[
\pi'=\pi+p^{(1)}+2\pi-2-n-1.
\]
Putting here \(\pi=3\), \(n=16\), \(\pi'=0\), one obtains that the linear genus of the plane is
\[
p^{(1)}=10.
\]

Let us pass to consider a ruled surface of order \(m\) and of genus \(p>0\). The double system of the plane sections (cut out by the quadrics) has the genus
\[
\pi=2p+m-1
\]
and the degree
\[
n=4m.
\]

The system adjoint to it consists of groups of \(m+2p-2\) lines\origfootnote{1}{s038L-1}{This can be verified in various ways; we refer to Chap. III.}, and hence has the genus
\[
\pi'=-(m+2p-2)+1=-m-2p+3.
\]

Hence from the formula
\[
\pi'=\pi+p^{(1)}+2\pi-2-n-1,
\]
one obtains that: \textit{the linear genus of ruled surfaces of genus \(p>0\)} is
\[
p^{(1)}=-8(p-1)+1.
\]

\edsection{7. Exceptional curves as fixed parts of the canonical system.}{7. Exceptional curves as fixed parts of the canonical system.}{sec.II.7}

It is useful to return and make precise, on some points, what concerns the exceptional curves belonging to a surface of genus \(p>0\), and their belonging, as fixed parts, to the canonical system \(|K|\).

That an exceptional curve of the first kind, \(\omega\), on the surface \(F\), belongs as a fixed component to the system \(|K|\) follows directly from the fact that it has genus \(\varrho=0\) and degree \(\nu=-1\). Indeed \(\omega\) must be intersected by the curves \(K\) in
\[
2\varrho-2-\nu=-1
\]
points; since this (virtual) number is negative, \(\omega\) must form part of every \(K\). One can say more: on the surface \(F\), of genus \(p>0\), there cannot be two irreducible exceptional curves of the first kind connected with each other, that is, having some point in common. Indeed, if
% scan: s038R
% !PRINTED 57
\weberpage{57}
there are two such curves, \(\omega\) and \(\omega'\), of genus 0 and of degree \(-1\), with \(i(\geq 1)\) common points, they must appear as fixed parts, to be counted a certain finite number of times, in all the canonical curves \(K\). But if the first splits off \(r\) times and the second \(s\) times, the number of intersections of \(\omega\) and \(\omega'\) with the residual curves will be resp.
\[
\begin{aligned}
(K\omega)-r(\omega\omega)-s(\omega\omega')&=-1+r-si,\\
(K\omega')-s(\omega'\omega')-r(\omega\omega')&=-1+s-ri.
\end{aligned}
\]
Now, of these two numbers, whose sum is negative, at least one is negative, so that the corresponding curve, \(\omega\) or \(\omega'\), would have to split off further from the curves \(K\), contrary to the supposition.

By the same reasoning one proves that the surface \(F\), of genus \(p>0\), cannot contain an exceptional curve \(\omega\) of the second kind. For this purpose it is useful to refer to a surface \(\Phi\), a transform of \(F\), on which \(\omega\) becomes a curve of degree \(\nu\geq 0\), always keeping the genus \(\varrho=0\). On this surface \(\omega\) must have in common with a canonical curve \(K\) the virtual number of points
\[
2\varrho-2-\nu=-2-\nu
\]
which is negative; \(\omega\) therefore splits off from the curves \(K\) as a fixed part; but it would have to split off infinitely many times, which is absurd. Indeed, if one supposes that it splits off \(s\) times, the number of intersections with the residual curves will be
\[
(K\omega)-s(\omega\omega)=-2-(s+1)\nu
\]
which is always a negative number.

We shall summarize the results obtained by stating the theorem:

\textit{On a surface of genus \(p>0\) there can be at most a finite number of irreducible exceptional curves of the first kind, not connected with one another, and there are never exceptional curves of the second kind.}

Let it be noted that when there are only irreducible exceptional curves of the \(1^{\mathrm{st}}\) kind, each of these, \(\omega\), splits off simply from the canonical system, remaining fundamental for the residual system \(|K-\omega|\). However, it cannot be excluded that \(\omega\) splits off further from \(|K-\omega|\), appearing twice as a fixed part of \(|K|\). This means that, on transforming \(\omega\) into a point \(O\), this point will turn out to be a \textit{base point common} to all the curves of the \textit{transformed canonical system}. Since, however, this base point is not assigned, it must be regarded as virtually non-existent; thus \(\omega\) too, on the given surface, must still be regarded as forming \textit{part of the pure canonical curves} which are obtained by splitting off once the exceptional curves which constitute fixed parts of the \textit{impure canonical system}. %

\textsc{Note I.} – Particular mention is deserved by the case in which there are \textit{reducible exceptional curves, arising from base points infinitely}

% scan: s039L
% !PRINTED 58
\weberpage{58}
\textit{near} on a transformed surface. Here the \textit{multiplicities of the exceptional curves} as fixed parts of the impure canonical system can be determined on the basis of the principle of continuity, since we know that the case in question arises as a limit when the base points of a transforming system approach one another indefinitely. Recalling the remark with which \S{} 12 of the preceding Chap. closes, we now point out that it has the following meaning: \textit{Each component \(\omega_i\) of a reducible exceptional curve} %
\[
\Omega=\omega+\omega_1+h_2\omega_2+\ldots+h_i\omega_i+\ldots+h_s\omega_s
\]
\textit{splits off from the impure canonical system a number of times equal to the sum of the multiplicities with which it appears in the exceptional curves of which it forms part}\origfootnote{1}{s039L-1}{\textsc{A. Franchetta}, \textit{Sulle curve eccezionali riducibili di prima specie.} Boll. dell’Unione Mat. It., Bologna, 1940.}.

For example, if one has an exceptional curve composed of three components, each connected with the next, corresponding to three infinitely near points following one another on a linear branch, these components give rise to three composite exceptional curves %
\[
\omega+\omega_1+\omega_2,\qquad \omega_1+\omega_2,\qquad \omega_2;
\]
then \(\omega\) will appear as a simple fixed part of the canonical system \(|K|\), \(\omega_1\) will appear in it counted twice, and \(\omega_2\) three times.

If instead the three aforesaid components come from successive points on a cuspidal branch (so that \(\omega\) and \(\omega_1\) are connected with \(\omega_2\) and not with each other), one will have three exceptional curves
\[
\omega+\omega_1+2\omega_2,\qquad \omega_1+\omega_2,\qquad \omega_2;
\]
hence \(\omega\) will again appear simply as a fixed part of the canonical system \(|K|\), \(\omega_1\) will appear in it doubly, and \(\omega_2\) as a quadruple curve.

The theorem that a surface \(F\) of genus \(p>0\) possesses in general only irreducible exceptional curves of the \(1^{\mathrm{a}}\) kind, not connected with one another, makes it evident that \(F\) itself can be transformed into another surface entirely free of exceptional curves, as \textsc{Enriques} indicated in his « Introduzione » of 1896. But the rigorous proof of this is obtained most simply following \textsc{Castelnuovo-Enriques} (1901), by recourse to the so-called Riemann-Roch theorem for surfaces, under broader conditions which include (as has already been mentioned) all the surfaces not belonging to the family of ruled surfaces. We therefore postpone its development until later. Here we wish only to note that the case in which one
% scan: s039R
% !PRINTED 59
\weberpage{59}
has an exceptional curve composed of \(s\) parts gives rise to no difficulty, because one of its components constitutes an irreducible exceptional curve, which can always be eliminated, thus reducing from \(s\) to \(s-1\).

\textsc{Note II.} – Let us bring out more explicitly what is included in the characteristic property of the canonical curves of cutting out, on a curve \(C\) of degree \(n\) and genus \(\pi\), a group of \(2\pi-2-n\) points. Precisely by virtue of this property the difference
\[
2\pi-2-n
\]
(which is in general positive, but which must be considered even if virtually negative) appears as an additive character for the systems of curves \(|C|\) on the surface; and its additive property can also be recognized directly from the formulas which give the degree and the genus of composite curves, independently of the effective or virtual meaning which here makes it evident. This said, one can assert that: \textit{On a surface of genus \(p>0\) an irreducible curve \(C\) of degree \(n\) and genus \(\pi\) for which}
\[
n>2\pi-2
\]
\textit{can only be an exceptional curve of the first kind:}
\[
\pi=0,\qquad n=-1.
\]
Indeed, if \(2\pi-2-n\) is less than zero, \(C\) must split off from all the canonical curves \(K\); but if one supposes that it splits off \(s\) times, the number of intersections of \(C\) with the residual curves will be
\[
2\pi-2-(s+1)n
\]
and therefore always negative, unless \(n=-1\), and consequently \(\pi=0\).

This theorem extends to the case of curves composed in any way whatever, by observing that a sum of negative or zero characters \(n_i-(2\pi_i-2)\) can never give a positive number. But one runs into a somewhat delicate analysis in the case of composite exceptional curves. The difficulty is avoided by referring to a transformed surface free of exceptional curves. Here is the result: \textit{On a surface of genus \(p>0\) there exist no curves (however composed) of genus \(\pi\) and degree \(n>2\pi-2\), apart from those which are obtained by adding disconnected exceptional curves to curves with \(n\leq 2\pi-2\)} (having with them a virtual number of common points equal to zero).

In particular one deduces that an exceptional curve of the first kind, however composed, belonging to a surface of genus \(p>0\),
% scan: s040L
% !PRINTED 60
\weberpage{60}
is characterized as « a connected curve of genus \(\varrho=0\) and degree \(\nu=-1\) »\origfootnote{1}{s040L-1}{With a slight restriction, namely excluding the case in which the curve belongs to a linear pencil of rational curves (with some fixed component), the theorem extends to all surfaces (\(p\geq 0\)). Cf. \textsc{A. Franchetta}, \textit{Sulla caratterizzazione delle curve eccezionali ecc.} Boll. Unione Mat. It., 1941.}.

\textsc{Exercise.} – The student is invited to resolve the following paradox: start from a surface \(F_4\) of the \(4^{\mathrm{o}}\) order (of genus \(p=1\)) containing a line \(a\); consider on \(F_4\) a cubic \(C\), a plane section through \(a\), and carry out a transformation of \(F_4\) into a surface \(\Phi\) by means of a linear system with a base point \(O\) on \(C\). We shall thus have on \(\Phi\) an elliptic curve \(C\), to which there is added an exceptional curve \(\omega\) meeting it in one point. Now consider the sum
\[
C+2\omega;
\]
it is a curve of virtual genus \(\varrho=0\) and of degree \(\nu=-1\), and yet it is not an exceptional curve!

\edsection{8. Bicanonical and pluricanonical curves.}{8. Bicanonical and pluricanonical curves.}{sec.II.8}

We have seen that the virtual existence of the canonical system
\[
|K|=|C'-C|
\]
constructed starting from a system \(|C|\) without base points on a surface \(F\) entails the existence of numerical characters such as \(p^{(1)}\) and \(p^{(2)}\), which are defined independently of the effective existence of \(|K|\). But the consideration of a virtual canonical system which has no effective existence takes on a more concrete meaning when it happens that the non-existent system gives rise to a double or a multiple that do effectively exist, and which can be defined directly by the relation
\[
|2K|=|2C'-2C|
\]
or
\[
|iK|=|iC'-iC|,
\]
since the fundamental theorem for adjoints gives in every case
\[
|iC'-iC|=|iL'-iL|.
\]
We also note that the bicanonical system can be defined by subtraction of the system \(|C|\) from its second adjoint
\[
|(C+C')'|=|C+C''|=|2C'|;\qquad |2K|=|C''-C|
\]

The simplest example of a surface of genus \(p=0\) to which there belongs a bicanonical curve (of order zero), and also the first
% scan: s040R
% !PRINTED 61
\weberpage{61}
that presented itself, is offered by the surface \(F_6\) of the sixth order which passes doubly through the edges of a tetrahedron, to which we shall have occasion to return later. The plane sections of \(F_6\) are sextics \(C\) of genus 4; the curves adjoint to them are twisted sextics \(C'\) cut out by the cubic surfaces which pass simply through the edges of the said tetrahedron; the curves \(C\) and \(C'\) are non-equivalent; but their doubles are equivalent, that is,
\[
|2C'|\doteq|2C|; %
\]
indeed the complete system \(|2C|=|2C'|\) is cut out on \(F_6\) by the surfaces of the \(6^{\mathrm{o}}\) order which likewise have the edges of our tetrahedron as double lines.

The \(F_6\) contains no exceptional curves, but if it is transformed into a surface \(\Phi\) in such a way that a point \(O\) of it changes into a curve \(\omega\), the exceptional curve \(\omega\) counted twice will by itself constitute the bicanonical curve of \(\Phi\).

The characters of the bicanonical system, and similarly those of the pluricanonical systems, furnish invariants of the surface \(F\). First of all, the dimension of the bicanonical system, or better this dimension increased by one, that is, the \textit{number of linearly independent bicanonical curves}, constitutes an absolute invariant which will be called the \textit{bigenus} \(P=P_2\). On the other hand, the genus and the degree of the bicanonical system will furnish relative invariants of the surface. But it is clear that these characters will be expressed by means of the (relative) linear genus \(p^{(1)}\).

Indeed, denoting by \(|K|\) the virtual canonical system of genus \(p^{(1)}\) and degree \(p^{(2)}=p^{(1)}-1\), the genus of the bicanonical system \(|2K|\) will be given by
\[
2p^{(1)}+p^{(2)}-1\doteq 3p^{(1)}-2 %
\]
and its degree will be
\[
4p^{(2)}=4(p^{(1)}-1).
\]

In an analogous way the \(i\)-canonical system of the surface \(F\), whether or not the pluricanonical systems of lower order have effective existence, will give us: an absolute invariant, which is the \(i\)-genus \(P_i\), that is, the number of linearly independent \(i\)-canonical curves, and relative invariants which express its genus and its degree. But these depend on the linear genus \(p^{(1)}\) of the surface; precisely, the genus of the \(i\)-canonical system is
\[
\frac{i(i+1)}{2}(p^{(1)}-1)+1,
\]
and its degree
\[
i^2(p^{(1)}-1).
\]

% scan: s041L
% !PRINTED 62
\weberpage{62}
The considerations developed in the preceding section concerning the exceptional curves on a surface of genus \(p>0\) now extend to the surfaces of genus \(p=0\) whose bigenus is \(P>0\), or, in general, which have some \(P_i>0\). Indeed, the effective existence of the bicanonical or \(i\)-canonical curves of \(|iK|=|K_i|\) always entails that a curve of genus \(\pi\) and degree \(n\) with \(n>2\pi-2\) must split off as a fixed part from the \(K_i\). In particular the exceptional curves of the first kind will appear as fixed parts of the \(K_i\), each counted in general \(i\) times. And, as in the case \(p>0\), there can be on the surface only exceptional curves of the \(1^{\mathrm{a}}\) kind not connected with one another. We do not pause to discuss in particular the composite exceptional curves, which do not give rise to new observations.

We point out only that, in consequence of what has been said for surfaces with a plurigenus \(P_i>0\) (even when \(p=0\)), there will always be defined the \textit{absolute linear genus}, which moreover is equal to \(p^{(1)}+\sigma\), denoting by \(p^{(1)}\) the relative linear genus and by \(\sigma\) the number of exceptional curves. On the other hand, the surface with \(p=0\) and with some plurigenus different from zero can be transformed into another free of exceptional curves, and the linear genus of the latter will give us precisely the absolute linear genus.

\edsection{9. Historical note.}{9. Historical note.}{sec.II.9}

The genus of an algebraic surface \(F\) presented itself for the first time to \textsc{A. Clebsch}\origfootnote{1}{s041L-1}{\textit{Comptes-rendus de l’Académie des Sciences de Paris}, 1868.} as the « number of double integrals of the first kind (that is, everywhere finite) belonging to \(F\) ». If \(F\) is supposed to be of order \(n\), in ordinary space, the algebraic integrands of the first kind are formed starting from the adjoint surfaces \(\varphi_{n-4}\), of order \(n-4\), just as the integrands of the first kind of a plane curve of order \(n\) are formed starting from the adjoint curves of order \(n-3\). It thus appears that these \(\varphi\), or rather the (canonical) curves which are their sections with \(F\) outside the double curve, have an invariant character under birational transformations of the surface. Their invariance received a first algebraic proof from \textsc{M. Noether}\origfootnote{2}{s041L-2}{\textit{Zur Theorie des eindeutigen Entsprechens algebraischer Gebilde.} Math-Annalen Bd. 2 (1869), Bd. 8 (1874).}. Then \textsc{Noether} himself noted that the canonical system \(|K|\) furnishes, besides the genus \(p\) (dimension increased by one), also a second invariant character of the surface, namely the genus of the \(K\), that is, the linear genus \(p^{(1)}\), defined at least for \(p>0\), and designated by him as « Kurvengeschlecht » in contrast %
% scan: s041R
% !PRINTED 63
\weberpage{63}
to « Flächengeschlecht ». As for the degree \(p^{(2)}\) of the canonical system, he saw that it is expressed in terms of \(p^{(1)}\) by the formula
\[
p^{(2)}=p^{(1)}-1,
\]
a consequence of the fundamental property of the \(K\) of cutting out on the generic curve of a linear system \(|C|\) groups residual to the characteristic series. The formula indicated, which \textsc{Noether} believed to be always valid for the effective characters, in truth holds, without exception, for the virtual characters defined by \textsc{Enriques}.

\textsc{F. Enriques}, taking up again after twenty years the general study of algebraic surfaces from the point of view of birational transformations\origfootnote{1}{s041R-1}{\textit{Ricerche di geometria sulle superficie algebriche.} Memorie della R. Accademia delle Scienze di Torino, 1893.}, turns first of all to ascertaining to what extent the fundamental property of the canonical curves \(K\), supposed to hold with respect to a particular system \(|C|\), suffices to characterize the said \(K\): this indeed happens (save for one exception) if \(|C|\) is an irreducible linear system, at least \(\infty^2\), free of proper fundamental curves (lowering the genus of the residual curves). From this the author derives a geometric proof of the invariance of the \(K\) with respect to birational transformations; a birational transformation of the surface is decomposed into successive transformations, each of which preserves the sections by the planes of a star. These and other topics treated in the cited « Ricerche » (which we shall have occasion to recall in what follows) are taken up again by \textsc{F. Enriques} in the « \textit{Introduzione alla Geometria sopra le superficie algebriche} » of 1896\origfootnote{2}{s041R-2}{Memorie della Società Italiana delle Scienze (detta dei XL) series III, t. X, 1896.}. As regards the canonical system \(|K|\), the most important thing is the extension of the invariance theorem to the case \(p_g=0\), in which \(|K|\) no longer has effective existence. Already in the «Ricerche» the author had considered the system \(|C'|\) \textit{adjoint} to a system \(|C|\), which — for \(p>0\) — presents itself as the sum \(|C'| = |C+K|\). Now he deepens the study of this system independently of the existence of \(|K|\) (that is, for \(p=0\)). The curves \(C'\) adjoint to \(|C|\) on a surface \(F\) (like the curves of order \(n-3\) adjoint to a plane curve of order \(n\)) cut out canonical groups on a generic \(C\). But this property alone does not suffice to characterize them if \(|C|\) possesses proper fundamental curves; one nevertheless succeeds in defining them in every case, in such a way as to satisfy, with regard to any two linear systems (free of base points), the functional condition %
\[
|C+D'|=|C'+D|.
\]

% scan: s042L
% !PRINTED 64
\weberpage{64}
This \textit{fundamental property of adjunction} contains within itself, for \(p>0\), the invariance of the canonical system:
\[
|K|=|C'-C|=|D'-D|,
\]
and constitutes the generalization of the invariance theorem itself, which still appears meaningful in the case in which \(|K|\), defined as a difference (and therefore a priori as \textit{virtual}), has no effective existence. The meaning in question is brought out by \textsc{Enriques} under several aspects:

1) The fundamental relation between the adjoint systems leads to defining uniquely the \textit{virtual characters} of the canonical system, \(p^{(1)}\) and \(p^{(2)}\), independently of its existence; and between them there holds, as has been said, the relation
\[
p^{(2)}=p^{(1)}-1.
\]

2) It may happen that, although \(|C|\) is not contained in \(|C'|\), the double or a multiple \(|sC|\) of it is contained in \(|sC'|\); this means that \textit{a virtual canonical curve, which has no effective existence, may give rise to an effectively existing multiple system.}

We have already mentioned the examples which illustrate this possibility. We shall see later the usefulness presented, particularly in problems of classification, by the new characters depending on the dimensions of the multiple systems of the canonical system, that is, the \textit{bigenus} and the \textit{plurigenera}.

The idea, which occurred some years later to the same author\origfootnote{1}{s042L-1}{Expounded in a course of lectures at the University of Bologna in the year 1897-98. Cf. the \textit{Programma} published in the Bollettino di bibliografia e storia delle Matematiche, no. 13 (1899).

For surfaces cf. \textit{Intorno ai fondamenti della geometria sopra le superficie algebriche.} Atti Accademia di Torino, 37, p. 19 (1901).}, of simply defining the adjoint systems and the canonical system by means of the Jacobians, was to bring a notable methodological advance to the theory of invariants thus outlined by \textsc{Enriques}. The idea is so simple, its implementation so easy, its scope so general with regard to the theory of varieties of any number of dimensions, that it has seemed to some (and not only to mathematicians who passed through the school of \textsc{Enriques}) to be a matter belonging to the common heritage of geometers, one which ought to have been known for a long time, at least as far as curves are concerned. But whoever rereads the expositions of the geometry on a curve by \textsc{C. Segre} and by \textsc{E. Bertini} of 1894\origfootnote{2}{s042L-2}{\textsc{C. Segre}, \textit{Introduzione alla Geometria sopra un ente algebrico semplicemente infinito.} Ann. di Mat., series II, t. 22.

\textsc{E. Bertini}, \textit{La geometria delle serie lineari sopra una curva piana secondo il metodo algebrico.} Annali di Matematica, s. II, t. 22.} in fact finds there no other demonstra%
% scan: s042R
% !PRINTED 65
\unskip\weberpage{65}%
tion of the invariance of the canonical series on a curve than that of \textsc{Brill} and \textsc{Noether}, connected with the theorem of \textsc{Riemann-Roch}, according to which this series is the only \(g_{2p-2}^{p-1}\) belonging to a curve of genus \(p\)\origfootnote{1}{s042R-1}{Cf. \textsc{C. Segre}, l. c., no. 75; \textsc{Bertini}, l. c., no. 21 \textit{b}) and \textit{g}).}. And if this were not so, if at that time the definition of the canonical series and the proof of its invariance based on the property of the Jacobian series of the \(g_n\) had been known, the more laborious development of the theory of the canonical curves of a surface, given by \textsc{Enriques} in the « Ricerche » of 1893 and in the « Introduzione » of 1896, would indeed be incomprehensible.

Let us add that the idea of the Jacobian systems can be implemented in different ways and thus gives rise to certain variants in the exposition of other authors.

The form adopted here, not far, moreover, from the one which appears in the «Fondamenti», corresponds to the development of the theory of curves expounded in the «Lezioni» by \textsc{Enriques-Chisini}\origfootnote{2}{s042R-2}{Vol. III. See also, for the analytic treatment: \textsc{F. Enriques}, \textit{Sul teorema d’invarianza della serie canonica \(g_{2p-2}^{p-1}\) ecc.} Memorie R. Accademia di Bologna, 1914.} and — if we are not mistaken — attains the maximum of simplicity.

\printersignature{\textsc{Enriques F.} - \textit{Superficie algebriche.}\quad 5}
% scan: s043L
% !PRINTED none
\weberblank{[s043L]}

\clearpage\thispagestyle{plain}
% scan: s043R
% !PRINTED none
\weberpage{[s043R]}
\edchapter{III}{ADJOINT SURFACES}{Chapter III. ADJOINT SURFACES}{ADJOINT SURFACES}

We have seen that on a surface \(F\) the curves \(C'\) adjoint to an irreducible linear system \(|C|\) cut out canonical groups on the latter. Now we wish to ascertain whether, and to what extent, this property suffices to characterize the curves \(C'\). The question is connected with another of a projective character: how the curves adjoint to the system of plane sections \(C\) are cut out on the surface \(F\), taken in ordinary space, by means of certain surfaces, which we shall call \textit{adjoint} to it, defined by the conditions of passing \(i-1\) times through every \(i\)-fold curve of \(F\) and by a suitable behavior at the isolated multiple points of \(F\).

\edsection{1. Surfaces of order \(n-3\) adjoint to a surface of order \(n\) free of proper multiple points.}{1. Surfaces of order n-3 adjoint to a surface of order n free of proper multiple points.}{sec.III.1}

Consider a surface \(F\) of a certain order \(n\), in ordinary space, and denote by \(C\) its plane sections. The surface \(F\) can possess two kinds of singularities:

1) one or more multiple curves, e.g. a double curve which corresponds, on a transform of \(F\), to a locus of pairs of points neutral for the transformed linear system of \(|C|\), or an \(i\)-fold curve corresponding to a locus of neutral groups of \(i\) points, etc.;

2) particularly notable multiple points, either isolated points or points of the multiple curve, whose singularity does not result from the passage of that curve through them.

As regards the multiple curve, we shall as a rule refer, for simplicity of exposition, to the case of ordinary multiple curves, and sometimes also of a nodal curve; but it will be easy to convince oneself that what we shall say holds equally for any curve of extraordinary singularity, when one takes account of the infinitely near multiple curves which define it.

% scan: s044L
% !PRINTED 68
\weberpage{68}
As regards the notable point singularities which must be taken into account in the considerations of this chapter, we shall make the distinction precise by recourse to the following criterion: a multiple point \(O\) of \(F\) will be called \textit{improper} if the genus of a plane section through \(O\) equals the genus of the sections by planes which do not pass through \(O\); on the other hand, a singular point of \(F\) will be called a \textit{proper multiple point} if it \textit{lowers the genus of the plane sections.}

If we refer to a transformed surface of \(F\) on which there corresponds to the multiple curve a locus of groups of points neutral for \(|C|\), then to an improper multiple point \(O\) of \(F\) there will in general correspond a \textit{neutral} group \(G\) of points (presenting only one condition to the \(C\) which are to contain them); thus, e.g., to a point which is triple for the surface \(F\) and for its double curve there corresponds in general a neutral set of three points; to a point \((i+1)\)-fold for \(F\) and for an \(i\)-fold curve of it, a neutral group of \(i+1\) points, etc.

It may happen, however, that by a transformation of \(F\) the point \(O\) gives rise to a curve \(\theta\) fundamental for the system \(|C|\), which may in general be composed of as many disconnected exceptional parts as there are points in the group \(G\) discussed above; and it is worth noting that this formation of \(\theta\) is indicated, in a characteristic way, by the property of being an \textit{improper fundamental curve}, not lowering the genus of the \(C\). Indeed this property entails that each connected part of \(\theta\) is of genus zero and has only one point in common with a residual curve with respect to \(|C|\) (or to a multiple of it), and thus this part is characterized as an exceptional curve. %

This premised, consider a surface \(F\) which \textit{is not ruled}, whose plane sections \(C\) are of order \(n\) and of genus \(\pi>1\); and on \(F\) a curve \(L\) which cuts the \(C\) in canonical groups. In a generic plane the group of points cut out by \(L\) will belong to a curve \(\varphi_{n-3}\) of order \(n-3\), adjoint to the plane section \(C\) of \(F\), constituting the complete intersection of the latter outside the multiple points (where an \(i\)-fold point of \(C\) absorbs \(i(i-1)\) intersections); \(C\) cannot have other intersections with the \(\varphi_{n-3}\) outside \(L\), unless it is reducible.

Take in space a generic line \(a\) and, in each plane through it, the \(\varphi_{n-3}\) adjoint to the plane section, determined by the group of points in which \(L\) meets the plane; as the plane varies through \(a\), the \(\varphi_{n-3}\) will describe a surface \(\Phi\), passing \(i-1\) times through every \(i\)-fold curve of \(F\)\origfootnote{1}{s044L-1}{The assertion can be verified analytically by reducing to the case of a surface described by the normal translation of a plane curve with an \(s\)-fold point (\(s=i-1\)), where it is evident that the \(s\)-fold point describes an \(s\)-fold line.}, which a priori may be supposed of order \(n-3+h\), containing the line \(a\) as an \(h\)-fold line. %

% scan: s044R
% !PRINTED 69
\weberpage{69}
Now if \(h>0\), the total intersection of \(F\) with \(\Phi\) will have to contain the points common to \(a\) and \(F\), which are outside \(L\) (and the multiple curve). If \(A\) is one of these points, there will be at least one plane, tangent to \(\Phi\) at \(A\), whose section with \(F\) will be intersected by the adjoint \(\varphi_{n-3}\) at the point \(A\) outside the canonical group lying on \(L\), and which therefore turns out to be reducible. But by a well-known theorem of \textsc{Kronecker-Castelnuovo}\origfootnote{1}{s044R-1}{Cf. \textsc{Enriques-Conforto}, Op. cit., Book II, chap. II, \S~7, p. 273.} the surface \(F\), having then \(\infty^2\) reducible plane sections, would have to be a \textsc{Steiner} surface with plane sections of genus zero, which contradicts our hypotheses. One concludes that the curves \(L\) cutting out canonical groups on the plane sections \(C\) are intersections of \(F\), outside the multiple curve, with adjoint surfaces \(\Phi_{n-3}\) of order \(n-3\). Hence the curves \(C'\) adjoint to \(|C|\) must be contained (totally or partially) in the system \(|L|\). Now the \(C'\) and the \(L\) are curves of the same order \(2\pi-2\), that is, they have the same number of intersections with the \(C\), and therefore they can differ only by fundamental curves of \(|C|\); this means that the system \(|C'|\) (where it is less ample than \(|L|\)) must be deduced from \(|L|\) by the imposition of some base point, which falls at a multiple point \(O\) of \(F\); but the \(C'\) are not subjected to having any base point at an improper multiple point \(O\), because the adjoints to the plane sections \(C\) through \(O\), which are still of genus \(\pi\), also meet these sections in \(2\pi-2\) points (at least virtually) outside \(O\).

What has been said can be clarified by referring to a transformed surface of \(F\) on which there corresponds to the point \(O\) a neutral group of simple points, \(G\): here the system \(|C'_1|\) adjoint to \(|C_1|\) (which is deduced from \(|C|\) by the imposition of simple points) does not differ from \(|C'|\). We shall therefore state the theorem:

\textit{On a surface \(F\) of order \(n\) in ordinary space, not ruled and free of proper multiple points, the curves adjoint to the plane sections \(C\) are characterized by the property of cutting out canonical groups on the said \(C\); and the complete system of these adjoints is cut out on \(F\) by the surfaces \(\Phi_{n-3}\) of order \(n-3\) (adjoint to it) passing \(i-1\) times through every \(i\)-fold curve of \(F\).}

\edsection{2. Adjoint surfaces of any order.}{2. Adjoint surfaces of any order.}{sec.III.2}

This result extends to the curves \(L\) which cut out on the plane sections \(C\) of \(F\) groups belonging to the series that is the sum or the difference of the canonical series and of a multiple, by any integer \(r\), of the characteristic series: under the aforesaid conditions (\(F\) not ruled and free of proper multiple points) the \(L\) are cut out on \(F\) (outside the multiple curve) by surfaces \(\Phi_{n-3+r}\), of order \(n-3+r\), adjoint
% scan: s045L
% !PRINTED 70
\weberpage{70}
to it. Suppose first that \(r>0\). It suffices to repeat the preceding reasoning, with just one caution: when \(r\geq3\), through the group \(\varphi\) cut out by \(L\) on a plane through the line \(a\) there pass infinitely many curves \(\varphi_{n-3+r}\) adjoint to the section \(C\) of that plane, which form a linear system of dimension \(\dfrac{(r-1)(r-2)}{2}\); and one should determine one of these curves by assigning, e.g., a suitable number of its points on \(a\), and then a certain number of tangents and of osculating curves. If instead \(r<0\), add to the curve \(L\) the curve \(H\) cut out on \(F\) by a generic surface \(\Psi\) of order \(|r|\). The curve \(L+H\) then cuts out on the plane sections of \(F\) groups of the canonical series, and therefore there exists a \(\Phi_{n-3}\), adjoint to \(F\), which passes through it. But the said \(\Phi_{n-3}\), having in common with \(\Psi\) the curve \(H\), of order \(|r|n > |r|(n-3)\), splits into \(\Psi\) itself and a \(\Phi_{n-3+r}\), again adjoint to \(F\) and passing through \(L\). %

From the preceding considerations there follows, in particular:

\textit{On a surface of order \(n\) in ordinary space, not ruled and not possessing proper multiple points, the surfaces \(\Phi_{n-4}\) passing \(i-1\) times through every \(i\)-fold curve of \(F\)} (that is, \textit{adjoint} to it) \textit{cut out the complete system \(|K|\) of the canonical curves} (\textit{impure}), as was already announced in Chap. II, \S~4; \textit{and in general the \(\Phi_{n-4+r}\) with \(r>0\) cut out completely on \(F\) the system adjoint to the \(r\)-fold system of the plane sections, which as a rule is cut out only in part by the system of all the surfaces of order \(r\).}

\textit{Moreover this system \(|(rC)'|=|C'+(r-1)C|=|K+rC|\) is characterized by the property of its curves of cutting out on the plane sections \(C\) of \(F\) groups belonging to the series that is the sum of the canonical series and of \(r\) times the characteristic series.}

\edsection{3. Ruled surfaces.}{3. Ruled surfaces.}{sec.III.3}

The preceding theorem has been established on the basis of two restrictive hypotheses:

1) that \(F\) is neither ruled nor a Steiner surface, and that

2) it does not possess proper multiple points.

We see at once that the restriction of not being ruled is truly essential. If \(F\) is a \textit{ruled surface of order \(n\) and of genus \(\pi>0\)} (which, for simplicity of exposition, we shall assume to be endowed only with a double curve), \textit{there exist no curves \(C'\) adjoint to the plane sections \(C\), nor adjoint surfaces \(\Phi_{n-3}\), but there are groups of \(2\pi-2\) generators cutting the \(C\) in canonical groups.}

Indeed, introducing the virtual canonical system \(|K|\), we can write
\[
|C'|=|C+K|;
\]
% scan: s045R
% !PRINTED 71
\weberpage{71}
the \(C\) are unisecant to the generators, the \(K\) cut them in \(-2\) points\origfootnote{1}{s045R-1}{The curves of the system \(|2C+K|\), being split into \(2\pi-2+n\) generators, cut them in zero points; hence the curves \(K\) cut the generators in \(-2\) points.}, hence the \(C'\) would have to cut the generators in
\[
1-2=-1
\]
points, hence they would have to contain them, which constitutes an absurdity.

On the other hand one also sees that there cannot exist surfaces \(\Phi_{n-3}\) of order \(n-3\) passing through the double curve of \(F\), since this curve meets the generators in \(n-2\) points.

One may add that on the ruled surface \(F\) there also exist groups of \(2\pi-2+(r-1)n\) generators cutting out on the plane sections \(C\) groups of the series that is the sum of the canonical series and of \(r\) times the characteristic series, but, except for \(r=2\), these groups do not constitute curves adjoint to \(|rC|\) and are not intersections of adjoint surfaces \(\Phi_{n-4+r}\), since the curves of \(|(rC)'|\) must meet the generators of \(F\) in \(r-2\) points (and not in zero).

By contrast, for \(r\geq2\), \textit{the surfaces \(\Phi_{n-4+r}\) passing through the double curve cut out on the ruled surface curves adjoint to the \(r\)-fold system of the plane sections.} This is proved by observing that these curves must cut out canonical groups on the curves of \(|rC|\), at least on those cut out on \(F\) by general surfaces of order \(r\), \(\Phi_r\). Indeed, on a \(\Phi_r\) the curves of intersection with the \(\Phi_{n-4+r}=\Phi_{r+n-4}\) turn out to be adjoint to the intersections with \(F\).

That, further, the \textit{system cut out on the ruled surface \(F\) by the \(\Phi_{n-4+r}\) is always complete} will be proved in what follows, when we learn how to evaluate its dimension.

\edsection{4. Characteristic property of the canonical curves and of the curves adjoint to a linear system.}{4. Characteristic property of the canonical curves and of the curves adjoint to a linear system.}{sec.III.4}

If we wish to state the results obtained in invariant form, we shall have to refer to a surface \(F\) and to an irreducible linear system \(|C|\) on it, and characterize with respect to \(|C|\) the canonical curves and the curves adjoint to the \(C\). But before setting out this statement it is well to free ourselves from a superfluous restriction. Indeed, we have supposed that the system \(|C|\) is not only irreducible but also \textit{simple, of dimension \(r\geq3\)}, so that \(\infty^3 C\) can appear as the plane sections of a transformed surface.

We can free ourselves from this restriction.

If \(|C|\) is (or contains) a net of irreducible curves of genus
% scan: s046L
% !PRINTED 72
\weberpage{72}
$\pi$ and degree $n$, we may suppose that this net is cut out on a surface $F$, of order $m>n$, by the planes of a star, with center at an $s$-fold point $O$ $(m=n+s)$. %

Then, starting from a curve $L$ which cuts out canonical groups on the $C$, and in relation to a pencil of planes whose axis $a$ passes through $O$, one will construct, as before, an adjoint surface $\Phi_{m-3}$ cutting out $L$ on $F$, which will pass $i-1$ times through every $i$-fold curve of $F$ and will have multiplicity $s-1$ at its $s$-fold point $O$; hence it will turn out that $L=C'$. %

The possibility of this construction, however, rests on two hypotheses:

1) that there do not exist $\infty^1$ reducible plane sections through $O$;

2) that $F$ does not possess proper multiple points (even infinitely near to $O$) or multiple lines through $O$, which might turn out to be of lower multiplicity for the $\Phi_{m-3}$ that we have constructed.

Looking at the matter from the invariant point of view, the two restrictions can be combined into a single one by requiring that the net of the $C$ shall not contain $\infty^1$ reducible curves. Indeed the presence of a fundamental curve for a linear system $\infty^r|C|$ means that $\infty^{r-1}$ reducible curves belong to it.

This said, we shall state the following theorem:

\textit{If on a surface $F$ there is given an irreducible linear system $|C|$, of dimension $r>1$, which does not contain $\infty^{r-1}$ split curves, the curves $C'$ adjoint to it are characterized by the property of cutting out canonical groups on the $C$. And the canonical curves of $F$ are defined by the property of cutting out on the $C$ groups residual to the characteristic series.}

The curves adjoint to a multiple system of $|C|$ can also be defined by the groups which they cut out on the $C$; the matter is clear without any need to insist on this point.

\edsection{5. Isolated multiple points.}{5. Isolated multiple points.}{sec.III.5}

Let us return to the projective consideration of the plane sections of a surface in ordinary space.

We wish to examine to what extent the results obtained extend to the case in which the surface $F$ possesses proper multiple points lowering the genus of the sections through them. For simplicity of exposition it will suffice to refer to the curves $C'$ adjoint to the system of plane sections $C$. It is therefore a question of examining to what extent these curves are characterized by the mere property of cutting out canonical groups on the plane sections $C$ of $F$. On the surface $F$ of order $n$, considered projectively, a curve $L$ cutting the $C$ in canonical groups is again, as was seen above, the intersection of a $\Phi_{n-3}$, of order $n-3$, passing $i-1$ times through
% scan: s046R
% !PRINTED 73
\weberpage{73}
every $i$-fold curve of $F$. But the system $|L|$ may fail to coincide with $|C'|$, while partially containing the latter: in any case one will have $|L|=|C'+\theta|$, where $\theta$ denotes a proper fundamental curve, or a set of proper fundamental curves, of $|C|$. And, speaking projectively, the surfaces $\Phi_{n-3}$ cutting out the $C'$ will be distinguished from the $\Phi_{n-3}$ cutting out the $L$ (subject only to passing $i-1$ times through every $i$-fold curve of $F$) by the condition of passing further through certain multiple points of $F$.

To clarify how things stand, let us suppose that $F$ possesses an \textit{isolated} $s$-fold point, that is, one outside the multiple curve. Consider the planes of space passing through a point $P$; the sections by these planes form a net of $|C|$, and, from the projective point of view, the Jacobian of the net is cut out by the polar surface of $P$, of order $n-1$, which in general passes $i-1$ times through every $i$-fold curve of $F$ and similarly with multiplicity $s-1$ through the $s$-fold point $O$. Note, however, that the neighborhood of the $s$-fold point $O$ corresponds (on a transformed surface) to a curve $\theta$ fundamental for $|C|$, which forms part (counted in general once) of the Jacobian $C_j$\origfootnote{1}{s046R-1}{By hypothesis $\theta$ has no variable intersections with the $C$. Take a generic point $A$ on $\theta$; the $C$ constrained to pass through $A$ split into $\theta$ and the residual curves $C_1$ of a pencil; among these there is a $C_1$ passing through $A$, which together with $\theta$ constitutes a curve of the net having a double point at $A$.}. Consequently, to obtain a whole Jacobian curve one must add to the intersection of the said polar the neighborhood of the point $O$; hence the system $|C_j|$ is in reality cut out by the surfaces of order $n-1$ passing $i-1$ times through every $i$-fold curve of the surface $F$ and with multiplicity $s-2$ through the $s$-fold point $O$. It follows that \textit{the curves $C'$ adjoint to the plane sections will be cut out by the surfaces $\Phi_{n-3}$ subject to the conditions of passing $i-1$ times through every $i$-fold curve of $F$ and of passing $s-2$ times through the isolated $s$-fold point.}

The same reasoning can be repeated for a proper $s$-fold point belonging to the multiple curve of $F$: the $s$-fold point is (in general and at least) $(s-2)$-fold for the surfaces which cut out on $F$ the adjoints $C'$ to the plane sections.

\edsection{6. Adjoints and subadjoints.}{6. Adjoints and subadjoints.}{sec.III.6}

The preceding analysis, relating to a surface $F$ with proper multiple points, leads us to distinguish, for such surfaces:

1) the surfaces subject to passing $i-1$ times through every $i$-fold curve of $F$, and

2) the surfaces which, besides passing $i-1$ times through every $i$-fold curve of $F$, are also subject to particular conditions of passage
% scan: s047L
% !PRINTED 74
\weberpage{74}
through its proper multiple points, in such a way that (for any order, however high) they give as intersections the curves adjoint to a multiple of the plane sections.

Reserving for the latter the name of \textit{adjoint surfaces}, we shall give to the former the name of \textit{subadjoint surfaces}, and accordingly the curves defined by the property of cutting out canonical groups on the $C$ will be called \textit{subadjoint curves} to the plane sections $C$. This said, there will be occasion to investigate, for each particular kind of proper multiple points, the further conditions which they impose on the subadjoint surfaces in order that these may become adjoint. In this connection it is useful to observe that the conditions relating to a multiple point, for the adjoints of a given order, may turn out to be in part dependent on the other conditions imposed by the multiple curves or by other multiple points; but one may admit from now on (and it will be clarified more precisely later) that this does not happen for surfaces of sufficiently high order. In any case, by evaluating the number of these conditions without taking account of the possible dependences we have discussed, we shall have a number which will always turn out to be the same\origfootnote{1}{s047L-1}{When one requires the surfaces $\varphi$ of sufficiently high order to possess a given point singularity, $O$, one is led to impose on the $\varphi$ a certain multiplicity at the point $O$ and at a finite number of points infinitely near to $O$ belonging to its successive neighborhoods, in consequence of which the $\varphi$ will possibly acquire the assigned multiplicities along infinitesimal curves of these neighborhoods. Thus the number of conditions imposed on the $\varphi$ will turn out to be independent of their order.} for adjoint surfaces of any order. %
If it actually corresponds to the number relative to the adjoint $\Phi_{n-4}$ of order $n-4$, it will express the \textit{diminution of the genus} of $F$ brought about by this multiple point. In any case, in saying that the multiple point diminishes the genus of the surface by so much, we shall be referring to the \textit{virtual diminution} of that genus, whose meaning we shall bring out later, recognizing that the virtual expression of the genus furnishes another invariant character of the surface (its \textit{arithmetic genus}), which may turn out to be less than the (\textit{geometric}) genus introduced earlier.

With the language we have introduced we shall say that:

\textit{A proper $s$-fold point $O$ of the surface $F$} (whether or not belonging to the multiple curve) \textit{is} (in general and at least) \textit{$(s-2)$-fold for the adjoint surfaces.}

In the simplest case, in which $O$ is isolated with an osculating cone of genus $\dfrac{(s-1)(s-2)}{2}$, the multiplicity of $O$ for the adjoint surfaces is precisely $s-2$, while there is no condition of passage for the subadjoints; therefore one can say that it imposes a priori
% scan: s047R
% !PRINTED 75
\weberpage{75}
$\dfrac{s(s-1)(s-2)}{6}$ linear conditions on the adjoint surfaces of order $n-4$, and thus \textit{diminishes} \textit{its genus} by $\dfrac{s(s-1)(s-2)}{6}$.

In particular a \textit{conical double point does not influence the genus} of the surface.

\textit{Examples.} -- The surface $F_4$ of the $4^\circ$ order without singularities is of (surface) genus $p=1$, possessing \textit{one} canonical curve of order zero, the section by the adjoint surface of order $n-4=0$. If the $F_4$ acquires a triple point, its genus diminishes by 1, thus becoming $p=0$; in agreement with the fact that the $F_4$ with a triple point is rational.

The $F_4$ with a conical double point is still of genus 1 and therefore certainly not rational. But there are rational surfaces of the $4^\circ$ order with one (single) double point\origfootnote{1}{s047R-1}{Surfaces determined by \textsc{M. Nöther}. Cf. \textsc{Enriques-Conforto}, Op. cit., Book I, chap. VIII.}: the first type of them is the $F_4$ endowed with a tacnode. From this example it already appears that « double points of a particular kind (among those with the osculating cone reduced to a double plane) can diminish the genus of the surface »: whence one draws the motive for a thorough analysis.

The surface $F_5$ of the $5^\circ$ order $(n=5)$ without singularities has genus $p=4$, possessing $\infty^3$ canonical curves, which are the sections by the planes of space (adjoint to $F_5$). If the $F_5$ acquires a quadruple point $(s=4)$, its genus diminishes by $\dfrac{s(s-1)(s-2)}{6}=4$, in agreement with the fact that it becomes rational.

The $F_5$ with a double line $a$ is of genus $p=2$, possessing as canonical curves the $\infty^1$ cubics which are the sections by the (adjoint) planes passing through $a$. If it acquires a triple point $O$ outside $a$, its genus reduces to 1. In this case there is \textit{one} adjoint plane, namely the plane $Oa$, which cuts $F_5$ in the \textit{impure canonical curve} consisting of \textit{three lines} through $O$. These lines are \textit{exceptional} for $F_5$; indeed one can recognize that the quadrics through $O$ and $a$ cut out on $F_5$ a linear system $|C|$ of genus 5 and of degree 8, by means of which $F_5$ is transformed into an $F_8$ of $S_5$ containing a plane cubic $C_3$, image of the point $O$, while the three lines of the impure canonical curve are transformed into three collinear points of $C_3$, simple for the surface $F_8$: and by projecting $F_8$ from the line which contains these three points of $C_3$ one obtains $F_5$ again. Since the hyperplane sections of $F_8$ are canonical curves of genus 5 in $S_4$, their system turns out to be its own adjoint, and one thus has the verification that the genus of the surface is equal to 1.

% scan: s048L
% !PRINTED 76
\weberpage{76}
\edsection{7. Multiple points of the first kind.}{7. Multiple points of the first kind.}{sec.III.7}

If we wish to investigate more closely the conditions imposed on the adjoint surfaces by a singular point of $F$, we must first of all distinguish the \textit{multiple points of the first kind} from those of higher kind. We shall call « of the first kind » a singular point $O$ of $F$, which we suppose to lower the genus of the plane sections $C$, when near it there is no other multiple point lowering the genus of the $C$ through $O$. On a suitable transform of $F$ the point $O$ becomes a curve $\theta$ fundamental for the system $|C|$, and what characterizes $\theta$ as a \textit{fundamental curve of the first kind} is that « there is no part of it which constitutes a fundamental curve of the residual system $|C-\theta|$ ».

The conditions imposed on the \textit{adjoint surfaces $\Phi$ by a multiple point of the first kind $O$} of $F$ are expressed simply by saying that: \textit{the $\Phi$ must be cut by the planes through $O$ in curves which, when a line through $O$ is added to them, become adjoint to the sections of $F$.}

The proof is given most simply by referring to the $\Phi_{n-4}$ of order $n-4$, and by supposing that the net of the sections $C$ of $F$ by the planes through $O$ does not contain $\infty^1$ reducible curves. For then it becomes evident that the $\Phi_{n-4}$ indicated above cut out on the $C$ groups residual to the characteristic series, a property which --- under the hypotheses made --- suffices to define the canonical curves. The reasoning is repeated analogously for the $\Phi_{n-4+r}$ with $r>0$, even if the $\Phi_{n-4}$ are lacking.

\textit{Examples.} -- Let there be a surface $F$ endowed with a \textit{tacnode} $O$, that is, with a point $O$ such that the plane sections through it have two double points infinitely near at $O$\origfootnote{1}{s048L-1}{Cf. \textsc{Enriques-Chisini}, Op. cit., Book IV.}. Applying the preceding criterion, one sees that the plane sections of the $\Phi_{n-4}$ through $O$, taken together with a line through $O$, must constitute curves adjoint to the sections of $F$. This entails that \textit{the said adjoints must pass simply through the tacnode $O$.} %

If instead one has an \textit{oscnode}, that is, a point $O$ such that the plane sections through it have three double points infinitely near at $O$, one finds in a similar way that \textit{the adjoints must pass simply through the oscnode and touch there the tangent plane to $F$.}

\textit{Exercise.} -- Let $F_5$ be a surface of the $5^\circ$ order endowed with a double line $a$ and a tacnode $O$. The adjoints of order $n-4=1$ to this surface reduce to the plane $Oa$, so that $p=1$.
% scan: s048R
% !PRINTED 77
\weberpage{77}
The intersection of $F_5$ with the plane $Oa$ (outside $a$) will consist of two exceptional curves: a conic and the line $b$, which touch each other at $O$.

These assertions will be verified by transforming $F_5$ into an $F_4$ of the $4^\circ$ order without singularities, containing a line $a'$: this surface (on which the system of plane sections is its own adjoint) is of genus $p=1$ and possesses a canonical curve of order zero. %

The aforesaid transformation will be carried out by means of a special quadratic transformation of the first kind of space, defined by the homaloidal system of the quadrics which pass through two incident lines and touch a given plane at a point of one of them: the base lines will be $a$ and the line $b$, intersection of the tacnodal plane at $O$ with the plane $Oa$, and the tacnodal plane at $O$ will be the fixed tangent plane referred to. As a result of this transformation the surface $F_5$ will become an $F_4$ of the $4^\circ$ order, and the quartic that is the further intersection of the tacnodal plane of $F_5$, outside the line $b$, will change into a line $a'$ of $F_4$. The exceptional conic of $F_5$ will be transformed into a point $O'$ of $F_4$ and the exceptional line $b$ into a point $B$; the transforms of the plane sections of $F_5$ will be curves $C_7$ of order 7, further intersections of $F_4$ with the quadrics through $a'$, having a double base point at $O'$ and a simple base point at $B$.

If one takes arbitrarily a surface $F_4$ containing a line $a'$, and on $F_4$ a generic point $O'$ (belonging to a cubic $K$, a plane section through $a'$), one can, conversely, transform $F_4$ into an $F_5$ endowed with a double line and a tacnode. For this purpose it is convenient to employ the homaloidal system of the quadrics which pass through $a'$ and through the line $b'$ tangent at $O'$ to the cubic $K$, and touch $F_4$ at $O'$. Here it must be observed that the imposition of a double point $O'$ on the curves $C_4$, partial intersections of the quadrics through the line $a'$, also entails as a consequence a simple base point at the tangential point of $O'$ on $K$. %
Thus the complete system of the $C_4$ without base points, which is of genus 6, of dimension 6 and of degree 10, gives rise to a superabundant system with one double and one simple base point, of genus 5, of degree 5 and of dimension 3. %

\edsection{8. Double planes.}{8. Double planes.}{sec.III.8}

Let us make a brief digression to dwell on a notable family of surfaces, which is illustrated by the preceding considerations. Consider the surface $F$ defined by the equation $z^2=f(xy)$, where $f$ is a curve of even order $2n$: by projection from the point at infinity $O$ of the $z$-axis, $F$ is represented on the \textit{double plane with branch curve $f$.} An $f$ of even order, taken arbitrarily and regarded with respect to the birational transformations of the plane, de%
% scan: s049L
% !PRINTED 78
\unskip\weberpage{78}%
fines a class of surfaces to which $F$ belongs; if $f$ is of odd order, the line at infinity of the plane is added to it as a branch line\origfootnote{1}{s049L-1}{Cf. \textsc{Enriques-Chisini}, \textit{Lezioni}, Book V, \S~12, vol. III, p. 90.}.

The surface $F$ is of order $2n$ and, if $f$ is free of singularities, possesses only one proper multiple point, at the point $O$. The analysis of this singularity reduces to the analysis of the singularity which the point at infinity of the $z$-axis presents for the hyperelliptic curve which is the section by a plane parallel to that axis, for example for the hyperelliptic curve $z^2=f(x)$, with $f$ of order $2n$. It is known\origfootnote{2}{s049L-2}{Cf. \textsc{Enriques-Chisini}, Op. cit., Book V, \S~12 note, vol. III, p. 93.} that this curve possesses at $O$ a $(2n-2)$-fold point with $n-1$ infinitely near double points proximate to it, on a branch of order $n-1$ tangent to the line at infinity of the plane.

Hence the surface $F$ possesses at $O$ a $(2n-2)$-fold singular point defined by the property that every plane section through it must have $n-1$ double points near $O$.

This said, let us consider the surfaces of order $2n-4$ adjoint to $F$. These are cones (and more properly cylinders) with vertex $O$, since, by the addition of a plane, they come to give as sections curves adjoint to the sections of $F$; but moreover such cones are subject to the condition of containing the double points of $F$ near $O$, and therefore the plane at infinity counted $n-1$ times splits off from them. We conclude that the surfaces of order $2n-4$ adjoint to $z^2=f(xy)$ consist of the plane at infinity counted $n-1$ times and of the cones of order $n-3$ with vertex $O$. \textit{The plane at infinity cuts $F$ in $2n$ exceptional lines. The canonical curves of the double plane having as branch curve a general curve of order $2n$ are represented on it by the curves $C_{n-3}$ of order $n-3$. The genus of the double plane is}
\[
p=\frac{n(n-3)}{2}+1=\frac{(n-1)(n-2)}{2}.
\]

Now note that a quadruple point of the branch curve $f$ is a tacnode for the surface $z^2=f(xy)$, the tacnodal plane being the plane $z=0$, where the lines through the origin have four intersections with the surface. We deduce from this that: \textit{a quadruple point of the branch curve diminishes the genus of the double plane by 1.}

Similarly a sextuple point of the branch curve gives rise to an oscnode of $F$ and therefore turns out to be double for the canonical $C_{n-3}$; in general a \textit{$2i$-fold point of the branch curve is $(i-1)$-fold for the canonical $C_{n-3}$ and thus diminishes the genus of the double plane by $\dfrac{i(i-1)}{2}$.}

% scan: s049R
% !PRINTED 79
\weberpage{79}
A multiple point $A$ of $f$, \textit{of odd order $2i+1$}, must \textit{be regarded}, in general, as \textit{of order $2i$}, inasmuch as the birational transformations of the plane having a fundamental point at $O$ change the neighborhood of this point into a curve which forms part of the transform of the branch curve. %
But, for the same reason, a multiple point of $f$ of odd order $2s+1$ which falls infinitely near to $A$ will have to be regarded as $(2s+2)$-fold; and thus, in particular, if one has two infinitely near multiple points $A$ and $A_1$, of the same odd order $2i+1$, the canonical $C_{n-3}$ will have to pass $i$ times through $A$ (instead of $i-1$) and $i-1$ times through $A_1$. %
The first example occurs for $i=1$: \textit{the singularity $[3,3]$, so defined, diminishes the genus of the double plane by 1}, just as happens for a 4-fold point. In \S~10 we shall find this result again by examining the singular double point which corresponds to it for the surface $F$.

\textit{Historical notice.} -- The notion of double planes goes back to \textsc{A. Clebsch}, « Ueber den Zusammenhang einer Klasse von Flächenabbildungen mit der Zweitheilung der Abelschen Funktionen »\origfootnote{1}{s049R-1}{Math. Annalen, Bd. III, 1861. Cf. also: Crelle’s Journal, 63.}, who represented on the simple plane the double planes with branch curve of order 2 and 4 (the latter arising --- according to \textsc{Geiser} --- from a projection of the cubic surface) and those with a branch sextic endowed with a quadruple point, arising, by projection, from a surface of the $4^\circ$ order with a double line. Then (1878) \textsc{M. Noether}\origfootnote{2}{s049R-2}{\textit{Ueber die Einzweideutigen Ebenentransformationen.} Sitzungsberichte d. Phys-med. Societät zu Erlangen, 10.} systematically investigated all the rational double planes, reducing them to three birationally distinct types, which --- by means of a $(1,2)$ transformation of \textsc{R. De Paolis} --- give rise to the types of the involutions of \textsc{E. Bertini}. Noether’s third type is the double plane with branch sextic endowed with two infinitely near triple points.

Later (1896) \textsc{Enriques}, undertaking the study of various classes of non-rational double planes\origfootnote{3}{s049R-3}{In particular in the memoir \textit{Sui piani doppi di genere uno.} Memorie della Società It. d. Scienze (detta dei XL), 1896.

The theory of rational double planes and of plane involutions of the $2^\circ$ order, elaborated and systematized in the teaching of \textsc{Enriques}, is expounded in the \textit{Lezioni} by \textsc{Enriques-Campedelli} (chap. 5, \S\S~57-63) and in the treatise \textit{Le superficie razionali} by \textsc{Enriques-Conforto} (Book II, chap. IV and V).}, indicated explicitly how the images of the canonical and pluricanonical curves are obtained on the double plane (cf. Chap. III, \S~11 and chap. V, \S~7).

% scan: s050L
% !PRINTED 80
\weberpage{80}
\edsection{9. The genus of the fundamental curves and the influence on the adjoints.}{9. The genus of the fundamental curves and the influence on the adjoints.}{sec.III.9}

Let us return to the consideration of the multiple points of the first kind of the surface $F$ to which there correspond \textit{irreducible fundamental curves} of the system $|C|$ of plane sections. Such a fundamental curve $\theta$ possesses a genus $\varrho$ and a degree $-s$ (virtual); $s$ will denote --- as we shall see --- the multiplicity of the point $O$ for $F$, and $\varrho$ the genus of the osculating cone (simple or multiple) at $O$ (account being taken of the tangents to the branches of the double curve which pass through $O$, and the other possible double generators of this cone being regarded as virtually non-existent).

We know that, owing to the presence of the multiple point $O$, the system $|C'|$ adjoint to the system of plane sections $C$ differs from the subadjoint system $|L|$ by the splitting off of $\theta$ counted a certain number of times, and we propose to find this number $x$, for which we shall have $|L|=|C'+x\theta|$.

Expressing that $\theta$ is a fundamental curve, without intersections with the curves of $|C|=|C_1+\theta|$, we shall have
\[
(\theta C_1)+(\theta\theta)=0
\]
and hence
\[
(\theta C_1)=s:
\]
the number $s$, essentially positive, is the number of intersections of $\theta$ with the curves $C_1$ residual to it with respect to $|C|$, that is, in general, the multiplicity of the point $O$ for $F$, as was mentioned above. This said, in order to determine the number $x$ defined above we may start from $|C'|$ and try to enlarge it successively by adding $\theta$ to it, as long as possible.

It is certain that this successive enlargement reaches a limit when the number of intersections of $\theta$ with the curves to which it is added becomes zero or negative. Therefore we must have:
\[
(\theta C')+x(\theta\theta)\geq 0.
\]

Let us write
\[
\begin{aligned}
|C|&=|C_1+\theta|,\quad |C'|=|C_1+\theta'|\\ %
(\theta C')&=(\theta C_1)+(\theta\theta')=s+2\varrho-2;
\end{aligned}
\]
the preceding inequality becomes
\[
2\varrho-2-(x-1)s\geq 0,
\]
which fixes an upper limit for the value of $x$:
\[
x\leq 1+\frac{2\varrho-2}{s}.
\]

% scan: s050R
% !PRINTED 81
\weberpage{81}
From the preceding inequality we obtain first of all that, if $\varrho=0$, then also $x=0$. Therefore the irreducible fundamental curves of the first kind and of genus zero of a linear system $|C|$ do not influence the definition of the adjoint curves. Or, in projective form: \textit{a multiple point of the first kind whose osculating cone} (simple or multiple) \textit{is irreducible of (virtual) genus zero does not influence the genus of the surface.}

In this class of multiple points there falls first of all the conical double point, of which we have spoken above, and then a triple point which is simple for the double curve of $F$, etc. If in the preceding inequality we put $\varrho=1$, we get $x\leq 1$. Here one can always take $x=1$. Indeed, let us refer to an $F$ endowed with a double curve on which there lies an $s$-fold point through which pass $\dfrac{(s-1)(s-2)}{2}-1$ branches of that curve; the $s$-fold point $O$ under discussion corresponds precisely to a fundamental curve of genus $\varrho=1$. Now the osculating cone to $F$ at $O$ is a cone of order $s$ which admits an adjoint cone of order $s-3$; this means that the subadjoint surfaces (of very high order), subject to passing through the double curve, will have multiplicity $s-3$ at $O$, while the adjoints must have multiplicity $s-2$ there.

The case of multiple points with $\varrho=1$ also includes the \textit{tacnode}, where the irreducible curve $\theta$ reduces to the neighborhood of $O$ on the tacnodal plane, counted twice. Indeed $\theta$ here appears represented on the double pencil of tangents with four branch lines, which correspond to the cuspidal plane sections issuing from the tacnode. One thus finds again the result established in the preceding section, that the tacnode diminishes the genus of the surface by 1.

In a similar way one sees that the \textit{oscnode} corresponds to a fundamental curve $\theta$ of genus $\varrho=2$, and here the value of the number $x$ turns out to be $x=2$, in agreement with what has been seen: the oscnode diminishes the genus of the surface by 3.

Now, referring to the case of an irreducible fundamental curve of the first kind of any genus $\varrho$ whatever, one may ask whether the system $|C'|$ can always be enlarged successively by adding $\theta$ to it, as long as $x$ satisfies the preceding inequality, as happens in the case of an isolated multiple point. The answer is negative: the upper limit of $x$ is not always a maximum that is effectively attained. To see this, an example suffices. Suppose that $F$ possesses a double curve passing with six branches through a sextuple point $O$, whose osculating cone will therefore be of genus $\varrho=4$. In general the tangents to these branches will not belong to a quadric cone, and hence a surface subadjoint to $F$ will have to have multiplicity 3 at $O$, having there an osculating cone of the $3^\circ$ order adjoint to the oscu%
\printersignature{\textsc{Enriques F.} - \textit{Superficie algebriche.}\quad 6}

% scan: s051L
% !PRINTED 82
\unskip\weberpage{82}%
lating cone of \(F\). For the adjoint surfaces the point \(O\) will have multiplicity \(4\), so that the system adjoint to the plane sections is deduced from the larger subadjoint system by the detachment of a single curve \(\theta\) corresponding to the neighbourhood of \(O\), and the value \(1\) is below the limit indicated for \(x\) by the preceding inequality:
\[
6-(x-1)6\geq 0,
\]
which gives, for the maximum value of \(x\), \(x=2\).

\edsection{10. Point singularities of higher kind.}{10. Point singularities of higher kind.}{sec.III.10}

The analysis of point singularities of higher kind, and of the influence which they have on the genus of the surface, reduces to the definition of the multiple points which may lie infinitely near to a proper singular point and which are not defined by the passage of singular curves, proper or infinitesimal.

If one imagines the point singularity possessed by a surface \(F\) at a point \(O\) to be completely resolved, the system of plane sections of \(F\), or equally a multiple of it (which we may take to be of dimension as large as we please), will give us on the transform of \(F\) a linear system \(|C|\) endowed with a fundamental curve which corresponds to the neighbourhood of the point \(O\). If the singularity of \(O\) is of the first kind, this fundamental curve is likewise of the first kind, so that, on detaching it from \(|C|\), no part of it remains fundamental for the residual system. On the other hand, one will define as a \textit{fundamental curve of the second kind} (corresponding to a multiple point \(O\) of the second kind) a fundamental curve which contains some parts constituting fundamental curves of the first kind for the residual system. And so, in general, the fundamental curves of \textit{kind \(s\)} will be defined inductively, their concept being reduced to that of the curves of kind \(s-1\).

From this definition it follows that the problem of determining the influence of a multiple point of kind \(s\) on the genus of the surface \(F\) reduces to the analogous problem relative to singularities of kind \(s-1\); and so the determination of the conditions which a fundamental curve of the system \(|C|\) imposes on the subadjoint curves that are to be adjoint to \(|C|\) is made to depend on the analogous determination with regard to the system obtained from \(|C|\) by detaching this curve.

Let us refer, for simplicity, to the case of a fundamental curve of the second kind of the system \(|C|\).

On the surface \(F\) which has the \(C\) as plane sections we have a multiple point \(O\) to which there are infinitely near, in the neighbourhood of the
% scan: s051R
% !PRINTED 83
\weberpage{83}
first order, some (proper) multiple points \(O_1, O_2, \ldots\ldotp\). On a transformed surface there will correspond to the point \(O\) a curve \(\theta\) fundamental for \(|C|\); to \(O_1\) a fundamental curve \(\theta_1\), etc. Now it is interesting to note that the curves \(\theta_1, \theta_2, \ldots\ldotp\) necessarily appear as parts of the fundamental curve \(\theta\), so that one will have in general
\[
\theta=\theta_0+\theta_1+\theta_2+\ldots\ldotp
\]

Indeed, the transformation to which the surface \(F\) has been subjected can be carried out in stages, taking first a transforming system of surfaces with the simple base point \(O\), so that the neighbourhood of \(O\) is changed into a curve \(\theta_0\) (simple or multiple) to which the singular points \(O_1, O_2, \ldots\ldotp\) will belong; and then transforming successively each of the points \(O_1, O_2, \ldots\ldotp\) into a curve which will in general turn out to be connected with \(\theta_0\). By virtue of these connections, when the curve \(\theta_0\) is detached from \(|C|\), the \(\theta_1, \theta_2, \ldots\ldotp\) also come to be detached as a consequence, these being likewise fundamental curves (without intersections with the \(C\)). In other words the sum
\[
\theta=\theta_0+\theta_1+\ldots\ldotp
\]
constitutes a \textit{monovalent} fundamental curve for \(|C|\), whose detachment from \(|C|\) entails only one condition. On the other hand the residual system \((C-\theta)\) will still possess the fundamental curves \(\theta_1, \theta_2, \ldots\ldotp\), so that it may be said that the curves
\[
\theta_0+2\theta_1+\theta_2+\ldots\ldotp
\]
and
\[
\theta_0+\theta_1+2\theta_2+\ldots\ldotp
\]
etc., will be \textit{bivalent} fundamental curves of \(|C|\), whose detachment from \(|C|\) entails \textit{two conditions}.

Let us consider, for example, the case of an (isolated) \textit{uniplanar double point}. This is a double point \(O\) at which the osculating cone reduces to a plane counted twice; the lines through \(O\) in this plane have in general three-point contact with the surface, except for three lines with four-point contact, which define three double points \(O_1, O_2, O_3\) infinitely near to \(O\)\origfootnote{1}{s051R-1}{Cf. e.g. \textsc{Enriques-Chisini}, Op. Cit., Book IV.}. We shall suppose for simplicity that the surface \(F\) to which the uniplanar point \(O\) belongs is a cubic surface \(F_3\), and we shall then illustrate what has been said by referring to its plane representation.

As a representative system of \(F_3\) one can take in the plane an \(\infty^3\) linear system of cubics \(C_3\) passing through three collinear base points \(O_1, O_2, O_3\) and through three other points infinitely near to
% scan: s052L
% !PRINTED 84
\weberpage{84}
these, \(O'_1, O'_2, O'_3\), where the line to which \(O_1, O_2, O_3\) belong corresponds to the uniplanar point \(O\) and to the part \(\theta_0\) of its neighbourhood. Let us transform our system of \(C_3\) by a Cremona transformation of the \(4^{\mathrm{o}}\) order, taking as transforming net a net of quartics with a triple base point, three simple base points \(O_1, O_2, O_3\) and three other simple base points. %
The system \(|C_3|\) will be transformed into another system of curves \(|C_9|\), with a sextuple base point \(P\), three collinear triple points \(A_1, A_2, A_3\), three other double points \(B_1, B_2, B_3\), and three simple points \(O_1, O_2, O_3\) which will belong respectively to the lines \(PB_1, PB_2, PB_3\). The line \(A_1A_2A_3\) counted twice represents the curve \(\theta_0\), while the \(PB_1, PB_2, PB_3\) are the images of \(\theta_1, \theta_2, \theta_3\). It appears at once that the group of lines \(\theta_0+\theta_1+\theta_2+\theta_3\) constitutes a monovalent fundamental curve for the system \(|C|\), since on detaching the line \(\theta_0\) one consequently detaches also the lines \(\theta_1, \theta_2, \theta_3\) which meet it, each at a point outside the base points, %
and consequently the said line \(A_1A_2A_3\) is detached a second time; hence the residual system of \(\theta=\theta_0+\theta_1+\theta_2+\theta_3\) is constituted by the quartics \(C_4\) which pass three times through \(P\), simply through \(B_1, B_2, B_3\) and again simply through \(A_1, A_2, A_3\). For this residual system \(\theta_0\) is no longer a fundamental line; on the other hand the lines \(PB_1, PB_2, PB_3\) are still fundamental lines.

The analysis of the singularity of the \(2^{\mathrm{a}}\) kind which a surface \(F\) possesses at an \textit{ordinary uniplanar point \(O\)} enables one to recognize that this point \textit{has no influence on the adjoint surfaces or on the genus}.

Indeed, consider the system \(|C|\) of the plane sections of \(F\) (which may, if need be, be replaced by a multiple of it). The plane sections \(C_1\) through \(O\) form a linear system which possesses only the fundamental curves \(\theta_1, \theta_2, \theta_3\), of the \(1^{\mathrm{a}}\) kind and of genus zero, corresponding to the double points \(O_1, O_2, O_3\) infinitely near to \(O\). Therefore these double points have no influence on the canonical curves, which remain defined by the property of cutting on the \(C_1\) groups residual to the characteristic series, that is, groups such that they belong to a curve of order \(n-3\) adjoint to a \(C_1\) and consisting of a \(\varphi_{n-4}\) and a line through \(O\). Hence the curves adjoint to the \(C_1\) are cut out by the adjoint surfaces \(\Phi_{n-3}\) which pass simply through \(O\), and, on adding the curve \(\theta\) which represents the neighbourhood of the latter, the curves \(C'\) adjoint to the plane sections of \(F\) will be cut out by \(\Phi_{n-3}\) which are no longer required to pass through \(O\). %

While the ordinary uniplanar double point \(O\) has no influence on the genus of a surface \(F\), there may be particularizations of this singularity which do have an influence on the genus. We have already seen that this happens for the tacnode, characterized by the property that
% scan: s052R
% !PRINTED 85
\weberpage{85}
the tangent lines through \(O\) have four-point contact with \(F\), and which therefore lowers by two, instead of by one, the genus of the plane sections through it. But more remarkable is the case of a uniplanar point \(O\) which lowers by only \(1\) the genus of the plane sections through it, and which possesses an infinitely near tacnode \(O_1\), in such a way that the genus of the plane sections through \(OO_1\) turns out to be lowered by \(3\) instead of by \(2\). In this case the system \(|C_1|\) of the plane sections through \(O\) possesses a fundamental curve \(\theta\) of genus \(1\), which has the tacnode \(O_1\) as its image, and this entails that the surfaces \(\Phi_{n-4}\) must pass simply through \(O_1\), that is, as a result of unloading, through \(O\). One therefore deduces that \textit{a uniplanar double point}, lowering by only \(1\) the genus of the plane sections through it, \textit{to which a tacnode is infinitely near, diminishes the genus of the surface by \(1\)}.

\textit{Note.} -- Noether pointed out two types of surfaces of the fourth order which turn out to be rational by virtue of possessing the singularity indicated above\origfootnote{1}{s052R-1}{Cf. \textsc{Enriques-Conforto}, \textit{Le superficie razionali}, Op. cit., pp. 196 and 204.}, and the rationality of these surfaces agrees with the fact that the aforesaid singular point diminishes by \(1\) the genus of the surface of the \(4^{\mathrm{o}}\) order, and so reduces it to zero. The two types of Noether surfaces give rise, respectively, to a plane representation:

1) with the representative system of \(\infty^3\) curves \(C_7\) of the \(7^{\mathrm{o}}\) order having a triple base point and \(9\) double base points on a fundamental cubic; and

2) with the representative system of \(\infty^3\) curves \(C_9\) of the \(9^{\mathrm{o}}\) order having \(8\) triple base points, one double and one simple, on a fundamental cubic.

In both cases the fundamental cubic of the representative system constitutes a \(\theta\) of the \(2^{\mathrm{a}}\) kind which is fundamental also for the residual system.

And it appears from the indicated representation that this \(\theta\) constitutes, on the surface \(F_4\), or rather on a transform in which the singularity is resolved without introducing new exceptional curves, a curve of genus \(1\) and degree \(-1\): indeed the plane cubic from which \(10\) base points are removed is precisely a curve of degree \(-1\)\origfootnote{2}{s052R-2}{Here it must be pointed out that type 2) of the Noether surfaces presents a projective complication, consisting in this, that the tacnodal plane of \(O_1\) passes through \(O\), %
so that the plane sections through \(OO_1\) contain a third fixed double point collinear with \(OO_1\) and split into the line \(OO_1\) and a cubic of genus \(1\) for which it is an inflexional tangent. %
Owing to this circumstance the analysis of the singularity must be carried out by having recourse, instead of to the system of plane sections, to a system of intersections with surfaces of higher order. Cf. \textsc{Enriques-Conforto}, Op. cit., pp. 206-7.}.

% scan: s053L
% !PRINTED 86
\weberpage{86}
In this regard the analysis of the singularity offered by the Noether surfaces \(F_4\) extends to the similar singularities of any surface whatever: \textit{the singularity of the \(2^{\mathrm{a}}\) kind constituted by a double point with an infinitely near tacnode corresponds in general to a fundamental curve of genus \(1\) and degree \(-1\) for the linear system \(|C|\) of the plane sections of the surface}\origfootnote{1}{s053L-1}{Cf. D. \textsc{Pedoe}, \textit{On the canonical System of certain algebraic Surfaces.} Proc. Phil. Soc. Cambridge, vol. 33 (1937).}.

The direct proof of the assertion (independent of the possibility of approximating the given surface by an \(F_4\) having the same singularity) can be given: first by showing that a fundamental curve of genus \(1\) and degree \(-1\) of a linear system \(|C|\) produces the said singularity on a surface having \(\infty^3\) \(C\) as plane sections; and conversely by starting from an \(F\) which possesses the said singularity, which is then resolved, and computing the characters of the fundamental curve which corresponds to it.

Assume that \(\theta\) is a fundamental curve for a system \(|C|\) of genus \(\pi\) and degree \(n\), a curve having in its turn genus \(\varrho=1\) and degree \(\nu=-1\); it appears at once that the number of intersections %
\[
(C-\theta)\theta=C\theta-\theta\theta=-\nu=1,
\]
that is, the curves of the residual system \(|C_1|=|C-\theta|\) have \textit{one} intersection, necessarily fixed, with \(\theta\) (\(C_1\theta=1\)). Consequently the genus \(x\) of \(|C_1|\) will be given by the equation
\[
\pi=x+1+1-1=x+1,
\]
that is
\[
x=\pi-1
\]
and its (virtual) degree will be computed from the formula
\[
n=CC=(C_1+\theta)(C_1+\theta)=C_1C_1+\theta\theta+2 %
\]
which gives
\[
C_1C_1=n-1
\]
(whereas the effective degree of \(|C_1|\), there being a base point on \(\theta\), will be \(n-2\)).

This being so, on the surface \(F\) which has \(\infty^3\) \(C\) as plane sections, \(\theta\) will have as its image a point \(O\), which lowers by \(1\) the genus of the plane sections through it and is therefore double for \(F\); and near \(O\) there will be a singular point \(O_1\) which lowers by two the genus of the (plane or surface) sections through it, and is therefore a tacnode.

Conversely, if one starts from a surface \(F\) which possesses a double point \(O\) with an infinitely near tacnode \(O_1\), to the point \(O\) there will cor%
% scan: s053R
% !PRINTED 87
\unskip\weberpage{87}%
respond, on a suitable transform of \(F\), a curve \(\theta\) fundamental for the system of plane sections \(C\), which a priori will be composed of two parts \(\theta_0\) and \(\theta_1\), the first of genus zero, corresponding to the neighbourhood of the \(1^{\mathrm{o}}\) order of \(O\) (the tacnode excluded), the other of genus \(1\), corresponding to the neighbourhood of \(O_1\): the curve
\[
\theta=\theta_0+\theta_1
\]
fundamental for \(|C|\), and necessarily connected, will represent the complete neighbourhood of the point \(O\).

Now one can evaluate the genus \(\varrho\) and the degree \(-\nu\) of the curve \(\theta\). %
Expressing that its detachment from the system \(|C|\) of genus \(\pi\) lowers the genus of \(C\) by \(1\), we shall have
\[
\begin{aligned}
|C|&=|C_1+\theta|\\
\pi&=(\pi-1)'+\varrho+x-1, %
\end{aligned}
\]
where \(x\geq 1\) denotes the number of the points of connection of the split \(C\) that are common to \(\theta\) and \(C_1\); and from this one deduces
\[
x=1 \qquad \text{and} \qquad \varrho=1.
\]

Hence, computing the number of intersections \(\theta C_1\) as a function of the degree of \(\theta\), one will have
\[
C_1\theta=(C-\theta)\theta=-\theta\theta=-\nu=1.
\]

After this let us evaluate the degree \(-\nu_1\) of \(\theta_1\) by writing that the passage of the sections \(C_1\) of \(F\) through \(O_1\) lowers their genus by \(2\); we shall have
\[
\begin{aligned}
|C_1|&=|C_2+\theta_1|\\
\pi-1&=\pi-3+1+y-1
\end{aligned}
\]
where \(y\) denotes the number of intersections
\[
C_2\theta_1=(C_1-\theta_1)\theta_1=-\theta_1\theta_1=\nu_1,
\]
and from this one obtains
\[
\nu_1=y=2.
\]

Now, denoting by \(-\nu_0\) the degree of \(\theta_0\), one will have
\[
-\nu=-\nu_0-\nu_1+2=-1, %
\]
and hence
\[
\begin{aligned}
\nu_0+\nu_1&=3,\\
\nu_0&=1.
\end{aligned}
\]

In conclusion the curve \(\theta_0\), of genus \(0\) and degree \(-1\), turns out to be an exceptional curve which can be changed into a simple point of the curve \(\theta_1\) of genus \(1\), which thus becomes a \(\theta\) of degree
\[
\nu=-1. %
\]
(preserving the genus \(\varrho=1\)).

% scan: s054L
% !PRINTED 88
\weberpage{88}
The singularity constituted by a double point to which a tacnode is near is met with in the study of double planes, when the branch curve possesses two infinitely near triple points \(O, O_1\), or, as one says, a point \([3,3]\). In this case consider, for example, the system of curves \(C\) having a certain genus \(\pi\) which is represented on the double plane by the \(\infty^5\) system of conics; the conics passing through \(O\) represent curves \(C_1\) of genus \(\pi-1\), but those through \(O\) and \(O_1\) represent curves of genus \(\pi-3\). Hence the point \(O\) is a double point of the surface \(z^2=f(xy)\) which lowers by only \(1\) the genus of the (plane or surface) sections through it, while the point \(O_1\) near to it lowers the genus of these sections by a further two, and is therefore a tacnode. Thus we find again that: \textit{a point \([3,3]\) belonging to the branch curve of order \(2n\) of a double plane} must be regarded as a simple base point for the curves of order \(n-3\) that are images of the canonical curves (cf. \S~7) and hence \textit{diminishes the genus of the double plane by \(1\)}.

Noether’s two \textit{rational} surfaces of the \(4^{\mathrm{o}}\) order which we recalled above are projected from their singular point into a double plane, with a branch curve of the \(6^{\mathrm{o}}\) order which possesses, for the two types respectively, a quadruple point and a point \([3,3]\): thus it is confirmed that the singularity of those surfaces diminishes by \(1\) the genus of the surfaces of the \(4^{\mathrm{o}}\) order.

\textit{Exercises.}

\(1^{\mathrm{o}})\) Show that the (isolated) singularity constituted by two infinitely near triple points \(OO_1\) imposes two simple base points on the adjoints and therefore has the same influence on the genus of the surface as two distinct triple points.

Then compute the genus of a surface \(F_5\) of the \(5^{\mathrm{o}}\) order with a double line and two infinitely near triple points \(OO_1\) off the line: one will find
\[
p=0.
\]

This deduction agrees with the fact that \(F_5\) is rational, since it possesses a pencil of rational quartic plane sections through \(OO_1\).

One can obtain the plane representation of the \(F_5\) indicated above by making use of a quadratic transformation of the \(2^{\mathrm{a}}\) kind, by means of the homaloidal system of the quadrics passing through the double line of \(F_5\) and through three other points of it, two of which are the infinitely near triple points \(O\) and \(O_1\). In this way \(F_5\) will be transformed into a surface of the \(4^{\mathrm{o}}\) order \(F_4\) with a double line.

\(2^{\mathrm{o}})\) Take a surface \(F_4\) of the \(4^{\mathrm{o}}\) order, free of singularities, containing a line \(a\): the linear system of plane sections, added to twice the pencil of the cubic sections \(K_3\) through \(a\), gives a linear system of curves \(C\) of genus \(3+1+1+(3-1)
% scan: s054R
% !PRINTED none
\text{\weberpage{[s054R]}}
+(3-1)=9\)\edFlush{}, of which one easily finds that the degree is \(16\) and the dimension is equal to \(9\). Impose on the \(C\) the condition of possessing a triple point at a point \(P\) of \(F_4\) off \(a\) which is a flex for the \(K_3\) passing through it. One thus obtains an \(\infty^3\) linear system of curves of genus \(6\) and degree \(7\), which possesses the fundamental curve \(K_3\), of genus \(1\) and degree \(-1\). By means of this system \(F_4\) can be transformed into a surface \(F_7\) of the \(7^{\mathrm{o}}\) order possessing a quadruple line, a triple line (which reduces the genus of the plane sections to \(6\)) and a double point with an infinitely near tacnode: one thus finds again, for \(F_7\), the genus of \(F_4\), which is \(p=1\). There is an adjoint surface of order \(7-4=3\), which reduces to the plane containing the quadruple and the triple line, counted twice, and to the plane projecting the isolated double point from the quadruple line, which cuts \(F_7\) in an exceptional cubic.

Find the inverse transformation of \(F_7\) into \(F_4\).

\textit{Note.} -- In what precedes we have made some observations in connection with interesting examples, but we by no means consider that we have exhausted the theory of the fundamental curves of linear systems in its relations with singularities. It is therefore necessary to examine and ponder a richer series of examples, and to recognize the complications to which they give rise. To go more deeply into this order of questions, we invite the student to consider certain singularities, even very simple ones, which may belong to the branch curve of a double plane: infinitely near quadruple points, series of \(4\) or \(6\) etc. infinitely near triple points, series of triple points infinitely near to quadruple points, etc. The task is to define these cases by characterizing the fundamental curve, for the system of plane sections of the surface \(F\), which corresponds to the said singular points, and the circumstances to which its detachment gives rise\origfootnote{1}{s054R-1}{Cf.: A. \textsc{Franchetta}, \textit{Sui punti doppi isolati delle superficie algebriche.} Rend. Acc. Lincei, 1945. Notes I and II.}. %

\edsection{11. Bicanonical curves and biadjoint surfaces.}{11. Bicanonical curves and biadjoint surfaces.}{sec.III.11}

Given in ordinary space a surface \(F\) of a certain order \(n\), we have learned to construct on it the canonical curves as intersections of surfaces \(\Phi_{n-4}\) of order \(n-4\) adjoint to \(F\), and more generally the curves adjoint to the \(r\)-fold system of the plane sections \(C\) of \(F\), by means of the adjoint surfaces \(\Phi_{n-4+r}\), which always cut out a complete linear system on \(F\). Now it will be easy to determine on \(F\) the bicanonical curves and also the curves second adjoint to the \(r\)-fold sys%
% scan: s055L
% !PRINTED 90
\unskip\weberpage{90}%
tem of \(|C|\), making use for this purpose of \textit{biadjoint surfaces} of \(F\), \(\psi_{2n-8}\) and \(\psi_{2n-8+r}\), of order \(2n-8\) and \(2n-8+r\) respectively.

Let us refer, for simplicity of exposition, to the case of genus \(p>0\), in which the adjoints \(\Phi_{n-4}\) and the canonical curves, their sections, actually exist. The behaviour of the biadjoint surfaces \(\psi_{2n-8}\) in relation to the singularities of \(F\) is as a rule defined by the behaviour of two \(\Phi_{n-4}\) which, taken together, make up precisely a \(\psi_{2n-8}\). However, it is necessary to define this behaviour more precisely, as we do with the observations which follow.

If \(F\) possesses only a double curve and no proper multiple points, the \(\Phi_{n-4}\) pass simply through this curve and the \(\psi_{2n-8}\) pass through it doubly: the linear system cut out by these \(\psi_{2n-8}\) is complete because it is deduced from the system cut out by the adjoints \(\Phi_{2n-8}\) (which we already know to be complete) by detaching the neighbourhood of the double curve.

But if \(F\) possesses a curve of multiplicity \(i>2\), the biadjoint surfaces \(\psi_{2n-8}\) (and similarly the \(\psi_{2n-8+r}\)) will be required to pass through this curve not with the multiplicity \(2i-2\), but with the multiplicity \(i\), having a suitable contact with the sheets of the surface \(F\) along it; for example, if \(i=3\) the \(\psi\) will possess the triple curve of \(F\) as a triple curve and will simply touch the sheets of \(F\). Indeed, for a surface \(\psi\) which passes four times through a triple curve of \(F\) one can always substitute a linear combination of \(\psi\) and \(F\) which passes three times through this curve, touching the sheets of \(F\). Of the \(\psi\) which are thus subjected to less restrictive conditions one can say that they certainly cut out on \(F\) a linear system cut out by the adjoints of the same order upon detaching, a certain number of times, the neighbourhood of the multiple curve. On the other hand the \(\psi\) passing four times through a triple curve of \(F\) will not in general give a complete system.

Similarly an isolated \(s\)-fold point of the surface \(F(s>2)\), which is in general \((s-2)\)-fold for the adjoint surfaces \(\Phi\), will have to be multiple for the biadjoint surfaces \(\psi\) not with the multiplicity \(2s-4\), but with the multiplicity \(s\), the condition being added of touching \(F\) at the point itself \(s-4\) times.

Thus, if \(O\) is an isolated triple point of \(F\), the \(\psi\) are required to pass doubly through it. If instead \(O\) is a quadruple point, the \(\psi\) are only required to have a quadruple point there; but if \(O\) is a quintuple point of \(F\), it will have to be only quintuple for the \(\psi\), which, however, will have to have the same osculating cone there.

Particular attention must be paid to the behaviour of the biadjoint surfaces at a tacnode of the surface \(F\). At first sight one might think that such a point, being simple for the adjoint surfaces \(\Phi_{n-4}\), must be double for the \(\psi_{2n-8}\);
% scan: s055R
% !PRINTED 91
\weberpage{91}
as for the surfaces composed of two \(\Phi_{n-4}\). But let us reflect on the way in which the behaviour of the surfaces \(\Phi_{n-4}\) at a tacnode \(O\) is established: a plane section through \(O\) possesses two infinitely near double points \(O\) and \(O_1\), and the section of a \(\Phi_{n-4}\) is such that together with a line through \(O\) it must constitute an adjoint \(\varphi_{n-4}\) to the section of \(F\). %
Therefore the \(\Phi_{n-4}\) is required a priori to pass through the point \(O_1\) infinitely near to \(O\), and only as a result of \textit{unloading} does it come to pass through the proper point \(O\). Similarly a plane through \(O\) will be cut by a biadjoint surface \(\psi_{2n-8}\) in a curve which a priori will be required to pass doubly through \(O_1\), and, by unloading, will come to pass simply through \(O\) and \(O_1\). Hence \textit{the conditions which a tacnode imposes on the biadjoint surfaces consist only in passing through this point and touching the tacnodal plane there}: thus the tacnode \textit{diminishes the bigenus by three}. These conditions can also be verified by noting that they entail, for the bicanonical curves which are sections of the \(\psi_{2n-8}\), the characteristic property of cutting on the generic curve of a linear system \(|C|\) (for example of the net of plane sections through \(O\)) groups of points belonging to the double of the series residual to the characteristic series.

Let us now turn our attention to the case of the Noetherian singularity constituted by a double point \(O\) with an infinitely near tacnode \(O_1\). Here too the passage of the \(\Phi_{n-4}\) through \(O\) is a consequence of the principle of unloading of singularities: a priori the \(\Phi_{n-4}\) ought to pass through the infinitely near point \(O_1\). Hence the \(\psi_{2n-8}\), required to pass twice through \(O_1\), will actually be determined by the condition of passing through the two infinitely near points \(O\) and \(O_1\). Therefore: \textit{a double point with an infinitely near tacnode in general requires the biadjoint surfaces to pass simply through the two infinitely near points \(O\) and \(O_1\)}, and hence \textit{diminishes the bigenus of the surface by only two}. It is easy to verify these results by referring to a double plane with branch curve \(f(xy)=0\) of a certain order \(2n\), that is, to the surface \(z^2=f(xy)\). We have seen that a quadruple point of \(f=0\) is a tacnode for the surface \(F\), and, in the plane, the canonical curves of \(F\) are represented by curves of order \(n-3\) passing simply through \(O\). Now the curves of order \(2n-6\) passing doubly through \(O\) will represent bicanonical curves of \(F\), although the system of these curves is not, in general, complete\origfootnote{1}{s055R-1}{A complete linear system of curves \(|C|\) on the double plane represents, on the surface \(F\), not a complete system, but rather a system of curves belonging to the involution \(I\) whose pairs correspond to the points of the double plane; and this can in general be extended to a larger system whose curves no longer belong to the involution \(I\). This is seen already in the case of the double plane with branch curve of the \(4^\circ\) order \(f_4=0\), which is obtained by projection, from the point at infinity \(O\) of the \(z\)-axis, of the surface of the \(4^\circ\) order \(F_4\) with equation \(z^2=f_4(xy)\), which has at \(O\) a tacnode with tacnodal plane at infinity. Here the \(\infty^5\) conics of the plane are images of the curves cut out on \(F_4\) by the quadric cones through \(O\), hyperelliptic curves which belong to the involution \(I\). But the \(\infty^5\) system of such curves is contained in the complete \(\infty^6\) linear system which is cut out on \(F_4\) by the quadrics touching the tacnodal plane at \(O\). %

As regards the possibility of extending the system of the bicanonical curves which have as images on the double plane curves of order \(2n-6\), the question depends on the behaviour of the biadjoint surfaces at the point at infinity of the \(z\)-axis, where \(F\) presents a known singularity. Suffice it to say that in general the \(\psi_{2n-8}\) will have at that point multiplicity \(2n-2\) and a suitable contact with the surface. It appears from this that, while the canonical curves of a surface representable on the double plane always belong to the involution whose pairs correspond to the points of this plane, the same does not in general happen for the bicanonical curves (Cf. Chap. V, \S~7).
}. Thus one succeeds in verifying that a tacnode diminishes the bigenus of \(F\) by \(3\).

% scan: s056L
% !PRINTED 92
\weberpage{92}
From here one passes to the case of the Noetherian singularity which corresponds to a point \([3,3]\) of the branch curve \(f\). The point \([3,3]\) is to be regarded as a double point \(O\) to which a 4-fold point \(O_1\) is infinitely near. The images of the bicanonical curves are therefore required to pass doubly through \(O_1\), that is --- as a result of unloading --- the images of the bicanonical curves will have to pass simply through \(O\) and \(O_1\), and hence \textit{a point \([3,3]\) of the branch curve of the double plane diminishes the bigenus of the surface by \(2\).}

\edsection{12. Surfaces of genus \(p=0\) and bigenus \(P>0\).}{12. Surfaces of genus p=0 and bigenus P>0.}{sec.III.12}

What has been said makes it easy to construct examples of surfaces of genus \(p=0\) and bigenus \(P>0\), which are therefore certainly not rational.

The first example presented itself to \textsc{Enriques} in 1896\origfootnote{1}{s056L-1}{\textit{Introduzione etc.}, l. c.}. It is the surface of the \(6^\circ\) order \(F_6\) which passes doubly through the edges of a tetrahedron. Taking this tetrahedron as fundamental for a system of projective coordinates, \(F_6\) is represented by an equation of the form:
\[
F_6=\varphi(x_2x_3x_4,\ x_3x_4x_1,\ x_4x_1x_2,\ x_1x_2x_3)+x_1x_2x_3x_4f(x_1x_2x_3x_4)=0
\]
where \(f\) and \(\varphi\) are quadratic forms, respectively in the \(x_1x_2x_3x_4\) and in the products of triples of these variables. Moreover, \(F_6\) belongs to the linear system which is the double of that constituted by the cubic surfaces passing simply through the edges of the tetrahedron.

% scan: s056R
% !PRINTED 93
\weberpage{93}
The \(F_6\) is of genus \(p=0\), because there exist no quadrics adjoint to it, that is, passing through the edges of the tetrahedron. But the bigenus of \(F_6\) is \(P=1\), because there exists a surface biadjoint to it of order \(2\cdot6-8=4\), namely our tetrahedron. This \(\psi_4\) cuts out on \(F_6\) \textit{a bicanonical curve of order zero.} The linear system \(|C|\) of the plane sections of \(F_6\) is of genus \(\pi=4\) and degree \(n=2\pi-2=6\); its characteristic series is a complete and non-special \(g_6\), whose double gives the double of the canonical series. The system \(|C'|\) adjoint to \(|C|\) is cut out on \(F_6\) by the cubic surfaces \(\Phi_3\) passing simply through the edges of the tetrahedron; like \(|C|\), it is a system of sextics of genus 4 and degree 6. And there is this noteworthy fact, that the doubles of the two systems \(|C|\) and \(|C'|\) coincide:
\[
|2C|=|2C'|;
\]
this coincidence expresses the concrete meaning of the existence of the bicanonical system of order zero \(|2C'-2C|\).

Since, moreover,
\[
|(2C)''|=|2C'|=|C+C''|
\]
one deduces that
\[
|C''-C|=|2C'-2C|=|2C-2C|
\]
and hence
\[
|C''|=|C|;
\]
in words: the system of the plane sections of \(F_6\) is in its turn the adjoint of its adjoint system, which is cut out by the surfaces \(\Phi_3\) passing through the edges of the tetrahedron.

Let us evaluate the linear genus \(p^{(1)}\) of \(F_6\) by computing the virtual genus of the curves of \(|C'-C|\) according to \S~10 of Chap. I. Denoting by \(\pi\) and \(n\) the genus and degree of \(|C|\), and by \(x\) and \(y\) the genus and degree of the virtual curve \(K=C'-C\) of order zero, which has \(i=0\) intersections with the \(C\), we shall write the genus and degree of \(|C+K|=|C'|\):
\[
\begin{aligned}
\pi'&=\pi+x+0-1=\pi\\
n'&=n+y+2\cdot0=n,
\end{aligned}
\]
whence
\[
x=1,\qquad y=0:
\]
the linear genus of \(F_6\) is \(p^{(1)}=1\).

A second example of a surface of genus \(p=0\) and bigenus \(P>0\) was pointed out by \textsc{Castelnuovo}, also in 1896\origfootnote{1}{s056R-1}{G. \textsc{Castelnuovo}, \textit{Sulle superficie di genere zero.} Memorie della Soc. It. delle Scienze (detta dei XL), series III, t. 10, 1896.}. It is the surface of the seventh order \(F_7\) which passes three times through a line \(r\), twice through a conic \(\Gamma\) (not meeting \(r\)), and moreover possesses three tacnodes whose tacnodal planes pass through \(r\).

% scan: s057L
% !PRINTED 94
\weberpage{94}
Such an (irreducible) surface does in fact exist and can be constructed as follows: first take a cubic surface \(F_3\) which passes simply through the line \(r\) and through the conic \(\Gamma\), and consider three planes \(\alpha_1\alpha_2\alpha_3\) passing through the line \(r\) and tangent to \(F_3\) at the points \(A_1,A_2,A_3\) (off \(r\)) respectively. It is easy to see that there exists an irreducible surface of the \(4^\circ\) order \(F_4\) passing doubly through the conic \(\Gamma\) and tangent to the planes \(\alpha_1\alpha_2\alpha_3\) at the points \(A_1,A_2,A_3\) respectively: this surface, taken together with the three planes \(\alpha_1\alpha_2\alpha_3\), constitutes a particular reducible \(F_7\) which passes twice through the conic \(\Gamma\), three times through the line \(r\), and possesses the tacnodes \(A_1,A_2,A_3\) with the tacnodal planes \(\alpha_1\alpha_2\alpha_3\). But a second split surface satisfying the same conditions is obtained by adding to \(F_3\) counted twice a plane through the line \(r\); hence, by forming the pencil of the two split \(F_7\), one will obtain an irreducible \(F_7\) endowed with the singularities indicated above.

There are no surfaces \(\Phi_3\) adjoint to \(F_7\): indeed a \(\Phi_3\), having to pass through the conic \(\Gamma\) and doubly through the line \(r\), which does not meet it, will a priori have to split, containing the plane of \(\Gamma\) as a part. But the residual quadric is required first of all to pass doubly through \(r\), and hence must split into two planes through \(r\), and must moreover pass through the tacnodes \(A_1,A_2,A_3\); this entails three conditions impossible to satisfy. One concludes that the genus of \(F_7\) is
\[
p=0.
\]

Now let us determine the surfaces \(\psi_{2\cdot7-8}=\psi_6\) biadjoint to \(F_7\). Among these there will certainly be found the surfaces composed of the plane of \(\Gamma\) counted twice and of four planes through the line \(r\), three of which are required to pass through the tacnodes \(A_1,A_2,A_3\) and therefore coincide with the planes \(\alpha_1\alpha_2\alpha_3\); thus one obtains \(\infty^1\) biadjoint surfaces \(\psi_6\) which contain as variable part a generic plane through \(r\) and cut out on \(F_7\) \(\infty^1\) bicanonical quartics of genus 1, forming a pencil without base points. It is clear that this bicanonical system of degree zero cannot be enlarged, and therefore the condition imposed on the biadjoint surfaces \(\psi_6\), of passing triply through the line \(r\) and touching the sheets of \(F_7\) through it, here entails as a consequence that the \(\psi_6\) possess \(r\) as a quadruple line.

Summing up, we shall say that: \textit{Castelnuovo’s surface \(F_7\) has genus \(p=0\) and bigenus \(P=2\), containing a pencil of bicanonical plane quartics of genus \(1\) and degree \(0\).} %

It is easy to determine the linear genus of \(F_7\), that is, the (virtual) genus \(p^{(1)}\) of its virtual canonical curves \(K\); indeed, the degree of \(|K|\) being equal to \(p^{(1)}-1\), the genus and the degree of \(|2K|\) will be respectively
\[
3p^{(1)}-2\qquad \text{and}\qquad 4p^{(1)}-4.
\]

% scan: s057R
% !PRINTED 95
\weberpage{95}
From \(3p^{(1)}-2=1\) and \(4p^{(1)}-4=0\) one likewise obtains that the \textit{linear genus} is \(p^{(1)}=1\).

For many years no other examples of surfaces of genus \(p=0\) and bigenus \(P>0\) presented themselves, and in particular the question remained unanswered whether among such surfaces there exist some for which the linear genus \(p^{(1)}>1\), possessing an irreducible bicanonical system of dimension \(>1\). This question is now answered by the examples actually constructed, simultaneously and independently of one another, by L. \textsc{Godeaux}\origfootnote{1}{s057R-1}{Cf. L. \textsc{Godeaux}, \textit{Sur une surface algébrique de genre zéro et de bigenre deux.} Rend. Acc. dei Lincei, vol. XIV, 1931.} and by L. \textsc{Campedelli}\origfootnote{2}{s057R-2}{L. \textsc{Campedelli}, \textit{Sui piani doppi con curva di diramazione del decimo ordine,} Ibidem, vol. XV, 1932.} in 1931-32.

We shall confine ourselves to reporting the example constructed by Campedelli, which concerns a double plane with branch curve of order 10 endowed with 6 points \([3,3]\) (a pair of infinitely near triple points). Such a curve \(f_{10}\) can actually be constructed, at least as a curve composed of three conics, bitangent in pairs, and of a quartic \(f_4\) which touches them at their 6 points of contact. Indeed, one will obtain 3 mutually bitangent conics by taking e.g. two concentric circles (which touch each other at the circular points of the plane) and an ellipse which touches the smaller circle at the ends of a diameter and the larger one at the ends of the perpendicular diameter. The 6 points of contact of the conics thus obtained offer 12 linear conditions to a quartic. Taking one of these quartics together with the three conics to constitute the branch curve \(f_{10}\) of a double plane, we shall have that:

\(1^\circ)\) there exists on the double plane no canonical curve, which would be a conic passing through the 6 singular points of \(f_{10}\), whereas in the example constructed one sees that these 6 points do not belong to a conic (circle); hence the genus of the double plane is
\[
p=0;
\]

\(2^\circ)\) there exists a net of quartics which, like the quartic \(f_4\), a component of the branch curve \(f_{10}\), pass through the 12 triple points of the latter: they represent \(\infty^2\) bicanonical curves of the surface, belonging to the involution whose pairs have the points of the plane as images, and one can see that they form the complete bicanonical system, so that the bigenus of the surface is %
\[
P=3.
\]

\textit{Note.} -- The proof that \(P=3\), and not \(P>3\), will result later from the expression of \(P\) by means of the linear genus, which here gives \(P=p+p^{(1)}=p^{(1)}=3\), where the value of \(p^{(1)}\) is found
% scan: s058L
% !PRINTED 96
\weberpage{96}
directly by observing that the virtual double conic constituting the canonical curve would possess 8 branch points: 6 at the triple points of \(f_{10}\) and two others outside them. If one does not wish to have recourse to the virtual canonical curve, it is also sufficient to observe that the degree of the bicanonical system is twice the degree of the system of quartics, so that one has
\[
4(p^{(1)}-1)=2\cdot4=8,\qquad p^{(1)}=3.
\]
If one does not wish to make use of the relation between the bigenus and the linear genus, one can nevertheless prove directly that, for our surface, the net of the double quartics that are images of the bicanonical curves is complete. For this it will suffice to observe that \(f_4\), a component of the branch curve, forms part of only one bicanonical curve, constituted by the curve itself counted twice, and that the \(\infty^2\) quartics through the triple points of \(f_{10}\) cut out on this \(f_4\) a complete and non-special \(g_4\): indeed, if this series were the canonical series of \(f_4\), the 12 triple points of \(f_{10}\) would come to belong to a cubic, which one easily sees to be an absurdity.

\edsection{13. Criteria of equivalence.}{13. Criteria of equivalence.}{sec.III.13}

The construction of the sub-adjoint and adjoint surfaces of a non-ruled surface \(F\) of ordinary space contains a \textit{criterion of equivalence}: two curves \(L\) and \(M\) which cut on the plane sections \(C\) of \(F\) groups of the series that is the sum (or difference) of the canonical series of \(C\) and of a multiple of the characteristic series are equivalent, up to fundamental curves of the system \(|C|\); they are equivalent outright if \(|C|\) possesses no proper fundamental curves. And one may observe that the equivalence of the two curves \(L\) and \(M\) can be inferred from the sole condition of cutting equivalent groups on the \(C\), even if the series to which these belong is not the sum of the canonical series and of a multiple of the characteristic series, when the system \(|C|\) has, as has been said, no proper fundamental curves.

Indeed, drawing through \(L\) a surface adjoint to \(F\), of sufficiently high order, one will have as further intersection a curve \(D\). Now \(L+D\) and \(M+D\) cut on a \(C\) equivalent groups of the series that is the sum of the canonical series and of a multiple of the characteristic series; therefore they are equivalent, and hence \(L\) and \(M\) also turn out to be equivalent.

The criteria of equivalence can be developed independently of the theory of adjoint and subadjoint surfaces, as has been done, in an interesting manner, by F. \textsc{Severi}\origfootnote{1}{s058L-1}{Cf. F. \textsc{Severi}, \textit{Il teorema d’Abel sulle superficie algebriche.} Ann. di Mat., (3), 12, 1905. \textit{Osservazioni varie di geometria sopra una superficie algebrica.} Atti Ist. Veneto, 65, part II, 1906. \textit{Alcune relazioni di equivalenza etc.} Ibidem, 70, 1911.}. Let us briefly set out
% scan: s058R
% !PRINTED 97
\weberpage{97}
these developments. First, let there be given on a surface \(F\) a linear \textit{pencil} of curves \(C\) endowed with at least one \textit{simple base point} \(O\) and \textit{free of reducible curves.} Then \textit{if two curves \(K_0\) and \(K_\infty\) not passing through \(O\) cut on the \(C\) equivalent groups of points, they are equivalent.}

Indeed, on any \(C\) of the pencil we can take the groups cut out by \(K_0\) and \(K_\infty\) as the group of zeros and the group of poles of a rational function \(\varphi\), which remains determined up to a multiplicative constant, and we can also determine this constant by requiring \(\varphi\) to take the value 1 at the point \(O\).

As \(C\) varies in its pencil, \(\varphi\) gives us a rational function of the points of the surface which has \(K_0\) as curve of zeros and \(K_\infty\) as curve of poles, so that these curves turn out to be equivalent, Q.E.D. For the deduction to be rigorous, however, it is necessary to prove that \(K_0\) exhausts the curve of zeros of our function, and likewise that \(K_\infty\) exhausts its curve of poles. To this end observe that in the opposite hypothesis one would have to add to \(K_0\) a curve which, having no variable intersections with the irreducible \(C\) of the pencil, would have to be composed of one or more \(C\); but then, on such a \(C\), the point \(O\) would come to constitute a point of indeterminacy of the function, where it would simultaneously have the values zero and 1, and this is impossible, since \(O\) belongs neither to the group of zeros nor to the group of poles which determine on \(C\) the subordinate rational function. %

The preceding criterion of equivalence can be generalized. \textit{Let \(|C|\) be a pencil of curves, rational or irrational, free of reducible curves: two curves \(K_0\) and \(K_\infty\) which cut equivalent groups on the \(C\) are equivalent to each other up to curves of the pencil}, that is, one has
\[
K_0+\Sigma C\equiv K_\infty+\Sigma'C.
\]

To prove this \textit{second criterion of equivalence}, we shall proceed as before to construct a rational function \(\varphi\) having \(K_0\) and \(K_\infty\) respectively as curve of zeros and as curve of poles. We know that \(\varphi\) is defined up to a multiplicative constant on every \(C\) of the pencil, and (denoting by \(t\) the parameter of the \(C\) of this pencil) we obtain in this way a function of the points of the surface, \(\lambda(t)\varphi\), where \(\lambda\) has to be determined. For this purpose it is not necessary to determine a point on each \(C\), as was done in the proof of the first criterion, but it suffices to determine on it a group of a certain number \(m\) of points, formed by the intersections which \(C\) itself has with an auxiliary curve \(L\) drawn on the surface; there will then be given the sum of the values
\[
\lambda\varphi_1+\lambda\varphi_2+\ldotp\,\ldotp\,\ldotp\,\ldotp+\lambda\varphi_m
\]
\printersignature{\textsc{Enriques F.} - \textit{Superficie algebriche.}\quad 7}
% scan: s059L
% !PRINTED 98
\weberpage{98}
which \(\varphi\) assumes at the points of this group. More simply, it will suffice to take, in place of \(\varphi\), the function
\[
f=\frac{\varphi}{\varphi_1+\varphi_2+\ldotp\,\ldotp\,\ldotp\,\ldotp+\varphi_m},
\]
which is a rational function of the points of the surface, independent of \(\lambda\).

It is obvious that \(K_0\) forms part of the curve of zeros of \(f\) and \(K_\infty\) forms part of the curve of poles, but these curves are not thereby exhausted: to \(K_0\) one must add a group \(\Sigma C\) formed by the \(C\) which pass through the points of intersection of \(L\) with \(K_\infty\); indeed, on these \(C\) the denominator of \(f\) becomes \(\infty\) and hence \(f\) is zero everywhere. Similarly, to \(K_\infty\) one must add a group \(\Sigma'C\) formed by the \(C\) which pass through the intersections of \(L\) with \(K_0\), as one sees by interchanging the roles of the curves \(K_0\) and \(K_\infty\), substituting \(\frac{1}{f}\) for \(f\).

We shall therefore have
\[
K_0+\Sigma C\equiv K_\infty+\Sigma'C.
\]
That \(K_0+\Sigma C\) actually exhausts the curve of zeros of \(f\) (and analogously for \(K_\infty+\Sigma'C\)) is proved by observing that this curve of zeros cannot contain, outside \(K_0\), anything but a certain number of curves without variable intersections with the \(C\), and therefore a group of (irreducible) \(C\); but on a \(C\) one cannot have identically \(f=0\) except when the sum \(\varphi_1+\varphi_2+\ldotp\,\ldotp\,\ldotp\,\ldotp+\varphi_m\), and hence one of the \(\varphi_i\), becomes infinite, which means that \(C\) itself must pass through a point common to \(L\) and \(K_\infty\).

Let us add two observations. The first is that when the pencil of the curves \(C\) is rational, its \(C\) are equivalent, and therefore there exists a difference \(\Sigma C-\Sigma'C\) consisting of a group of a certain number \(s\) of curves \(C\), to be taken positively or negatively, so that one will have %
\[
K_0\equiv K_\infty-sC
\]
or
\[
K_\infty\equiv K_0+sC.
\]

And if in particular the pencil of the \(C\) is endowed with base points, one has outright the equivalence
\[
K_0\equiv K_\infty
\]
when the equivalence between the two groups \((CK_0)\) and \((CK_\infty)\) holds with account taken also of the points which fall at the base points.

Indeed, from the fact that, for example,
\[
K_0\equiv K_\infty-sC
\]
% scan: s059R
% !PRINTED 99
\weberpage{99}
it follows that
\[
(CK_0)\equiv(CK_\infty)-s(CC);
\]
but we also have
\[
(CK_0)\equiv(CK_\infty),
\]
and hence, if the group \((CC)\) is not null, we must have \(s=0\).

Let us pass to the \textit{third criterion of equivalence: if on a surface \(F\) there is given an \(\infty^1\) system of index}\origfootnote{1}{s059R-1}{It is understood that through a generic point of the surface there pass \(\nu\) curves of the system.} \textit{\(\nu>1\) of curves which are all irreducible, two curves \(K_0\) and \(K_\infty\) which cut the \(C\) in equivalent groups are equivalent.}

To prove it, one will construct as before on each curve \(C\) a rational function \(\varphi\) which has its zeros at the intersections with \(K_0\) and its poles at the intersections with \(K_\infty\), and which remains determined up to a multiplicative constant \(\lambda\). Then one will take on \(F\) an auxiliary curve \(L\) having a certain number \(m\) of intersections with the \(C\), and --- for each \(C\) --- one will consider the sum of the values which \(\varphi\) takes at these \(m\) points:
\[
\varphi_1+\varphi_2+\ldotp\,\ldotp\,\ldotp\,\ldotp+\varphi_m.
\]
In this way there is determined on each \(C\) the rational function
\[
f=\frac{\varphi}{\varphi_1+\varphi_2+\ldotp\,\ldotp\,\ldotp\,\ldotp+\varphi_m}.
\]
After this we shall determine a rational function \(\psi\) of the points of the surface \(F\) by attributing to it, at a point \(P\), the sum of the values \(f_1,f_2,\ \ldotp\,\ldotp\,\ldotp\,\ldotp\ f_\nu\) which \(f\) takes on the \(\nu\) curves \(C\) passing through \(P\):
\[
\psi=f_1+f_2+\ldotp\,\ldotp\,\ldotp\,\ldotp+f_\nu.
\]
We can indicate the aforesaid values of the \(f_i\) by making use of a superscript; and then we shall write
\[
f_i=\frac{\varphi^{(i)}}{\varphi_1^{(i)}+\ldotp\,\ldotp\,\ldotp\,\ldotp+\varphi_m^{(i)}}.
\]

The rational function \(\psi\) evidently vanishes on the curve \(K_0\), where all the \(\varphi^{(i)}\) are equal to zero, and its poles must necessarily lie on \(K_\infty\); but it is not a priori evident that all the points of \(K_\infty\) are poles of the rational function \(\psi\), since a sum of infinities is not always infinite. What we can assert is that the curve \(\overline{K}_0\) of the zeros of \(\psi\) is of the type \(\overline{K}_0=K_0+A\), while, calling \(\overline{K}_\infty\) the curve of poles, one has \(K_\infty=\overline{K}_\infty+B\). Cutting the curves \(\overline{K}_0\) and \(\overline{K}_\infty\) with the generic \(C\), one has %
\[
(K_0,C)+(A,C)\equiv(K_\infty,C)-(B,C)
\]
% scan: s060L
% !PRINTED 100
\weberpage{100}
But by hypothesis \((K_0,C)\equiv(K_\infty,C)\), and hence the curves \(A\) and \(B\) would have no intersections with the generic curve \(C\), whereas our hypotheses exclude the existence of such curves. In conclusion, one has
\[
\overline{K}_0=K_0,\quad \overline{K}_\infty=K_\infty
\]
and the equivalence of the curves \(K_0\) and \(K_\infty\) is thus proved.

\textit{Remark.} -- In what precedes we have dealt with the criteria which allow one to recognize the equivalence of two (isolated) curves on the surface \(F\). Other, more significant, criteria are obtained in relation to a continuous series (at least \(\infty^1\)) of curves \(|C|\) given on \(F\). The first criterion of this kind was established by \textsc{Enriques} in 1895. A \textit{rational series} of curves \(|C|\) always consists of \textit{equivalent curves}, that is, it is contained in a linear system. The property is proved as for rational series of groups of points on a curve\origfootnote{1}{s060L-1}{Cf. \textsc{Enriques-Chisini}, Op. cit., Book V, chap. I, note to \S~10, vol. III, p. 73 (and p. 485).}, if one does not wish to reduce it to that case by means of the criteria given above. A second notable criterion was brought to light by \textsc{Severi} (1905)\origfootnote{2}{s060L-2}{\textit{Osservazioni etc.}, l. c.}. One can assert, in general, that \textit{a series \(\infty^1\)} (at least) \textit{of curves \(C\) consists of equivalent curves if the \(C\) cut equivalent groups on a single generic curve of a continuous system \(\{D\}\) of index \(\nu>1\).}

Indeed, the condition for an \(\infty^1\) series of groups of points on a curve to consist of equivalent groups (that is, to be contained in a linear series) can be expressed, following Castelnuovo, by means of an enumerative criterion, in which there appears the number of double points of the series itself (which must have the maximum value)\origfootnote{3}{s060L-3}{Cf. \textsc{Enriques-Chisini}, op. cit., Book V, chap. IV, 42 (vol. III, p. 482).}. Now, if this criterion is satisfied by the groups of intersection of the \(C\) with one \(D\), by reasons of continuity it will also be satisfied with respect to the other \(D\), and so one is brought back to the third criterion, given above.

Let us add that the criterion of equivalence for \(\infty^1\) series of groups of points on a curve, which, as we recalled above, was established by \textsc{Castelnuovo}, also extends to surfaces, as R. \textsc{Torelli} has shown. He writes an inequality satisfied by the number \(d\) of the curves endowed with a double point that belong to an \(\infty^1\) series (an inequality in which the Zeuthen-Segre invariant of the surface appears: cf. chap. V); the maximum of this number \(d\) corresponds to the case in which the said series consists of equivalent curves\origfootnote{4}{s060L-4}{Cf. R. \textsc{Torelli}, \textit{Sui sistemi algebrici di curve appartenenti ad una superficie algebrica.} Atti R. Accademia delle Scienze; Torino, 1907.}.

\clearpage\thispagestyle{plain}
% scan: s060R
% !PRINTED none
\weberpage{[s060R]}
\edchapter{IV}{THE ARITHMETIC GENUS AND THE RIEMANN-ROCH THEOREM\\FOR SURFACES}{Chapter IV. THE ARITHMETIC GENUS AND THE RIEMANN-ROCH THEOREM FOR SURFACES}{THE ARITHMETIC GENUS AND THE RIEMANN-ROCH THEOREM FOR SURFACES}

\edsection{1. Introduction.}{1. Introduction.}{sec.IV.1}

To understand the meaning of the developments which follow, it is well to recall the problems which are met with in the theory of curves and which find here their natural extension\origfootnote{1}{s060R-1}{Cf. e.g. \textsc{Enriques-Chisini}, Op. cit. Book V, vol. III.}.

After having defined the linear series on a curve and the elementary operations concerning them, and after having recognized the invariance of the canonical series, the question arises of evaluating the dimension \(r\) of the complete series determined on a curve of genus \(p\) by a group of \(n\) points. The first answer which is obtained here, namely \(r\geq n-p\), is given by the computation of the dimension of the system of the curves of sufficiently high order adjoint to the plane curve \(C\) of order \(m\) and genus \(p\). Since the system of these adjoints \(\varphi_{m-3+s}\) always cuts out on \(C\) (whatever \(s\) may be, positive or negative) a complete series, and since, for \(s\) sufficiently high, the double points of \(C\) offer to the curves of order \(m-3+s\) which must contain them %
\[
d=\frac{(m-1)(m-2)}{2}-p; %
\]
independent linear conditions, the series of degree \(n=2p-2+sm\) cut out on \(C\) by these adjoints turns out to be of dimension
\[
n-p.
\]

It is clear, then, that any series cut out on \(C\) by the \(\varphi_{m-3+s}\) passing through a group \(G\) of points which imposes on them conditions that are all independent will have dimension equal to the degree diminished by the genus \(p\). These results are then made more precise in the sense that the canonical series, of degree \(2p-2\), cut out on \(C\) by the
% scan: s061L
% !PRINTED 102
\weberpage{102}
adjoints \(\varphi_{m-3}\), always has the dimension \(p-1\), as results from the number of conditions offered by the double points of \(C\) to the curves of order \(m-3\) which must contain them, these conditions turning out to be independent of one another. Finally the Riemann-Roch theorem properly so called answers in the most general way the question of evaluating the dimension of the complete series determined on \(C\) by a group of \(n\) points: denoting by \(i\) (\(\geq 0\)) the index of speciality of the group, that is, the number of linearly independent canonical groups which contain it, the dimension of the said series turns out to be exactly %
\[
r=n-p+i.
\]
Now, wishing to extend this order of problems to surfaces, we shall be led first of all to seek, for the surface \(F\) of order \(n\) of ordinary space, the dimension of the linear system cut out by the adjoint surfaces \(\Phi_{n-3+s}\), for \(s\) sufficiently high, and to relate the result obtained to the computation of the surface genus of \(F\), of which we shall in this way have a virtual expression. Following this path we shall be led to evaluate the dimension of the linear systems which present themselves as adjoints of other systems, thus establishing a « Riemann-Roch theorem for adjoint systems », and then to examine the case of an arbitrary linear system, which no longer presents itself as the adjoint of another.

To evaluate the dimension of the system of the \(\Phi_{n-3+s}\) adjoint to a surface \(F\), it is necessary to recall the so-called postulation formulae of Cayley (1870) and of Noether (1871), which express the number of independent linear conditions imposed by a curve on the surfaces of sufficiently high order which are to pass through it, once or more times.

\edsection{2. Conditions imposed on a surface by passing through a curve:\\ postulation formulae.}{2. Conditions imposed on a surface by passing through a curve: postulation formulae.}{sec.IV.2}

Consider in ordinary space an (irreducible) twisted curve \(C\) of a certain order \(m\) and a certain genus \(p\), which is free of multiple points; the surfaces \(\Phi\) of any order \(n\) cut out on \(C\) a linear series \(g_{mn}^{r}\), of order \(mn\), whose dimension we denote by \(r\). If the series is complete and non-special, one can write
\[
r=mn-p.
\]

Under this hypothesis one at once finds the postulation of the curve \(C\) with respect to the surfaces of order \(n\), that is, the number of independent linear conditions which \(C\) offers to the \(\Phi\) which must contain it; indeed the number of these conditions will be %
\[
r+1=mn-p+1.
\]

% scan: s061R
% !PRINTED 103
\weberpage{103}
Our question would thus receive the simplest answer, were it not that two causes may modify the result indicated in opposite directions: the series \(g_{mn}^{r}\) may be special, and hence of dimension \(r>mn-p\), and it may be deficient, that is, have a dimension less than that of the complete series. Owing to the first circumstance the postulation of \(C\) for the \(\Phi\) may turn out to be \(>mn-p+1\), and owing to the second it may turn out to be smaller. But of the two disturbing causes one is eliminated at once by referring to surfaces \(\Phi\) of sufficiently high order \(n\); for it suffices to take \(n\) such that
\[
mn>2p-2.
\]
Hence one can assert that for the surfaces \(\Phi\) of sufficiently high order \(n\), the twisted curve \(C\), of order \(n\) and genus \(p\), has a postulation \(\leq nm-p+1\); this means that through \(C\) there pass \textit{at least} %
\[
\frac{(n+3)(n+2)(n+1)}{6}-nm+p-1
\]
linearly independent surfaces \(\Phi\).

One can make the result obtained more precise by proving that for \(n\) sufficiently high, \(n\geq m-2\), the said series \(g_{mn}^{r}\) is complete\origfootnote{1}{s061R-1}{The proof which follows is due to \textsc{G. Castelnuovo}, \textit{Sui multipli di una serie lineare di gruppi di punti appartenente ad una curva algebrica.} Rend. Circolo Mat. di Palermo, t. VII, 1893. Reprinted in “ Memorie Scelte „ p. 95.}. %
For this purpose let us project the curve \(C\) onto a plane from a generic point \(O\) of space: we shall thus have a curve of order \(m\) with \(d=\frac{(m-1)(m-2)}{2}-p\) double points, on which the adjoint curves of order \(\geq m-2\) will cut out a complete and non-special series. It follows from this that the cones of order \(m-2\) (or greater) having their vertex at \(O\) and containing the \(d\) chords of \(C\) which pass through \(O\) will cut out on \(C\), outside the points of incidence of these chords, a complete series. A fortiori the surfaces of order \(n\geq m-2\) passing through the same points will cut out on \(C\), outside these, a complete series. This series is enlarged if one extends the system of the aforesaid surfaces by referring to the totality of the surfaces \(\Phi_n\) of order \(n\) of space. %
Now, to prove that these cut out on \(C\) a complete series, it will suffice to recognize that the group \(G\) of the \(2d\) points of incidence of \(C\) with its chords through \(O\) offers independent linear conditions to the \(\Phi_n\), so that the addition of the points of \(G\) leads from a complete (non-special) series of dimension %
\[
mn-2d-p
\]
to a complete series of dimension
\[
mn-p.
\]

% scan: s063L
% !PRINTED 104
\weberpage{104}
For this purpose, referring for simplicity of exposition to the case \(n=m-2\), it will suffice to construct a surface of order \(m-2\) which passes through \(2d-1\) points of \(G\) and not through the remaining one; such a surface will be obtained by adding, to a cone of order \(m-3\) which passes through \(d-1\) chords of \(C\) issuing from \(O\) and not through the remaining one, a plane which passes through one of the two points of incidence of the remaining chord and not through the other. The existence of the aforesaid cone of order \(m-3\) results from the known theorem that the double points of a plane curve of order \(m\) (the projection of \(C\)) offer independent conditions to the curves of order \(m-3\) which are to be adjoint to it.

We therefore conclude: \textit{an irreducible twisted curve of order \(m\) and genus \(p\), without multiple points, offers}
\[
nm-p+1
\]
\textit{linear conditions to the surfaces of order \(n\geq m-2\) which must contain it.}

The computation of the postulation of a curve which we have indicated above extends, with slight modifications, to the case of twisted curves endowed with multiple points, and also of reducible ones. Let us first consider the case which is most important for us, in which we are dealing with a \textit{curve \(C\)} (of order \(m\) and genus \(p\)) \textit{which is endowed with a certain number \(t\) of triple points,} with tangents not belonging to the same plane, and which we wish to suppose still irreducible. Here one proves first that the \(t\) said points impose \(4t\) independent linear conditions on the surfaces \(\Phi\) of sufficiently high order \(n\) which are to possess them as double points, and then that the \(\Phi\) passing doubly through the \(t\) triple points of \(C\) cut out on \(C\) (outside these) a complete (non-special) series, so that \textit{the postulation of the curve \(C\) for the \(\Phi\) of sufficiently high order \(n\) turns out to be}
\[
4t+nm-6t-p+1=nm-2t-p+1.
\]

The first assertion is evident, since the assignment of a double point in an assigned position entails \(4\) linear conditions which the surface must satisfy; the second assertion is justified by the same reasoning as was carried out for the curves \(C\) without multiple points.

In a similar way one obtains the postulation of \(C\) in the case in which \textit{\(C\) possesses a certain number \(\tau\) of double points} (nodes): here too it is convenient first to make the \(\Phi\) pass also through these \(\tau\) points, which entails \(\tau\) linear conditions and diminishes by \(2\tau\) the number of intersections with \(C\); thus \textit{the postulation} of \(C\) for the \(\Phi\) of order \(n\) turns out to be
\[
\tau+nm-2t-2\tau-p+1=nm-2t-\tau-p+1.
\]

% scan: s063R
% !PRINTED 105
\weberpage{105}
Finally we wish to extend this formula to the case in which the curve \(C\) splits into several parts, respectively of orders \(m_1, m_2, \ldotp\ldotp\ldotp\ldotp\) and of genera \(p_1, p_2, \ldotp\ldotp\ldotp\ldotp\)

For simplicity of exposition let us refer to the case in which there are only two components \(C_1\) and \(C_2\), which may be connected through simple crossings, that is, double points of the composite curve \(C=C_1+C_2\), or also through any of the triple points of \(C\) that is double for one of the components and simple for the other. Let us observe first of all that the surfaces \(\Phi\) of sufficiently high order \(n\), passing doubly through the triple points of \(C\) and simply through its double points, will cut out complete (non-special) series on \(C_1\) as well as on \(C_2\); indeed, let the conditions relating to the passage through the double or triple points of \(C_2\) be dropped for the \(\Phi\); the \(\overline{\Phi}\) thus obtained will cut out a complete series on \(C_1\); but one can always suppose that through a group of this series there passes so large a system of surfaces \(\overline{\Phi}\) that one can always satisfy the omitted conditions and hence find a \(\Phi\) which passes through that group. We say more: not only are the series cut out by the \(\Phi\) on \(C_1\) and \(C_2\) each complete on its own account, but also the series which one has on the split curve turns out to be complete, since to any group of the series given on \(C_1\) one can associate a group on \(C_2\); in other words, when one considers the \(\Phi\) passing through a group of the series on \(C_1\) and through another, associated group of the series on \(C_2\), the one and the other being taken arbitrarily, it does not happen that these \(\Phi\) must consequently contain one or the other of the two curves. To see this, by the same procedure used in the case of irreducible \(C\), we can reduce ourselves to the plane, invoking the proposition that « on a reducible plane curve the adjoint curves of very high order always cut out a complete series ».

We have thus established that for the surfaces \(\Phi\) of sufficiently high order, passing, as has been said, through the points of connection of the components of \(C\), the postulation of the curve \(C\) is equal to the sum of the postulations of its components. It remains only to evaluate these postulations. For this purpose let us denote by \(m_1\) and \(m_2\) the orders and by \(p_1\) and \(p_2\) the genera of the said components. If one disregards the fixed points assigned to the \(\Phi\), the orders of the two series cut out on \(C_1\) and on \(C_2\) would be \(nm_1\) and \(nm_2\) respectively, and the dimensions \(nm_1-p_1\) and \(nm_2-p_2\). But a double point of \(C=C_1+C_2\) diminishes by \(1\) the dimension of the system of the \(\Phi\) and diminishes by \(2\) altogether the sum of the orders of the two series cut out on \(C_1\) and on \(C_2\), whether this point is a double point for one of the irreducible components of \(C\) or a simple point common to both. Similarly a triple point of the composite curve \(C\), having to be

% scan: s064L
% !PRINTED 106
\weberpage{106}
double for the \(\Phi\), lowers by \(4\) the dimension of the system of the \(\Phi\), and lowers altogether by \(6\) the sum of the orders of the series cut out by it on \(C_1\) and on \(C_2\), whether it is a double point for one and simple for the other, or a triple point for one of them. In short, one recognizes in this way that the \(t\) triple points of the curve \(C\) and its \(\tau\) double points have the same influence as in the case of the irreducible curve \(C\). Hence the postulation of \(C\) will be expressed by the same formula written above, in which the \(t\) triple points of \(C\) and its \(\tau\) double points enter, provided one makes the modifications that we indicate here. In place of the series of order \(x=mn-6t-2\tau\) cut out by the \(\Phi\) on the irreducible curve \(C\), one will have to consider the two series cut out on \(C_1\) and \(C_2\), whose orders we denote by \(x_1\) and \(x_2\); and, in place of the dimension \(x-p\) of that series, one will have the sum of the dimensions of the two series given on \(C_1\) and \(C_2\), that is
\[
x_1-p_1+x_2-p_2.
\]
But, whereas in the preceding case the postulation of \(C\) was given by \(x-p+1\), now this postulation will be obtained by adding the postulations of \(C_1\) and \(C_2\), and will therefore be
\[
x_1-p_1+1+x_2-p_2+1=x-(p_1+p_2)+2.
\]

In conclusion, the postulation of the curve \(C=C_1+C_2\), for the surfaces of sufficiently high order \(n\), will be expressed by
\[
mn-2t-\tau-(p_1+p_2)+2.
\]

This formula reduces to the one we wrote for irreducible \(C\) if, in place of the genera \(p_1\) and \(p_2\) of the components \(C_1\) and \(C_2\) of \(C\), one introduces the genus \(p\) of the composite curve, computed, following Noether, under the hypothesis that the components have no points of connection\origfootnote{1}{s064L-1}{These points, instead, have been removed by imposing them as base points on the \(\Phi\).}:
\[
p=p_1+p_2-1.
\]

Let us state the general result that one reaches in the case of \textit{a curve of order \(m\) composed of \(h\) irreducible curves of genera \(p_1, p_2, \ldotp\ldotp\ldotp\ldotp p_h\), which is endowed with \(t\) triple points (with generic tangents) and \(\tau\) double points; the postulation of the curve for the surfaces of sufficiently high order \(n\) is}
\[
nm-2t-\tau-p+1
\]
\textit{where}
\[
p=p_1+p_2+\ldotp\ldotp\ldotp\ldotp+p_h-h+1
\]
\textit{denotes the genus of the composite curve under the hypothesis that its parts are not connected with one another.}

% scan: s064R
% !PRINTED 107
\weberpage{107}
\edsection{3. Dimension of the systems of surfaces adjoint to a given one.}{3. Dimension of the systems of surfaces adjoint to a given one.}{sec.IV.3}

Let us refer to a surface \(F\) of a certain order \(n\), belonging to ordinary space and endowed with normal singularities: a nodal double curve with triple points having non-coplanar tangents and, if one wishes, also double ones\origfootnote{1}{s064R-1}{A priori one may suppose that the given surface \(F\) possesses an irreducible double curve without crossings: \(\tau=0\) (Cf. \textit{Introduction, etc.}, \S\ 5). Nevertheless it may be useful to develop the formulae also in the case \(\tau>0\).}. Let us evaluate the number \(N_m\) of the linearly independent surfaces adjoint to \(F\), of a sufficiently high order \(m=n-4+s\). Denoting these surfaces again by \(\Phi_m\), and denoting by \(d\) the order of the double curve of \(F\), by \(\varrho\) its genus, by \(t\) the number of triple points and by \(\tau\) the number of double points, we shall have that the number \(N_m\) of the \(\Phi_m\) will be
\[
N_m=\binom{m+3}{3}-(dm-2t-\tau-\varrho+1);
\]
hence the difference \(N_m-N_{m-1}\) will be given by %
\[
N_m-N_{m-1}=\binom{m+2}{2}-d.
\]

The meaning of this relation is that \textit{the adjoint surfaces \(\Phi_m\) of a sufficiently high order cut out on a plane the complete linear system of the curves \(\varphi_m\) adjoint to the section of \(F\).}

Indeed one can deduce the dimension \(N_{m-1}-1\) of the system \(\lvert\Phi_{m-1}\rvert\) from the dimension \(N_m-1\) of the system of the \(\Phi_m\) by detaching a plane from them; and the detaching of the plane entails precisely \(r+1\) linear conditions, when \(r\) is the dimension of the system cut out by the \(\Phi_m\). In this case, therefore, one finds.
\[
\begin{aligned}
r+1&=\binom{m+2}{2}-d,\\
r&=\frac{m(m+3)}{2}-d,
\end{aligned}
\]
which is the dimension of the complete system of the \(\varphi_m\) adjoint to the plane section of \(F\).

In what follows we shall say that the number \(N_m\) defined by the formula written above expresses in every case the \textit{virtual number of the surfaces \(\Phi_m\) adjoint to \(F\)}: we know that for \(m\) sufficiently high the virtual number is equal to the effective number of the \(\Phi_m\), whereas for the lower values of \(m\) there may be a difference. Nevertheless it is important to observe that: \textit{for \(m\geq n-4\) the effective number of the surfaces adjoint to \(F\) is always \(\geq N_m\).}

% scan: s065L
% !PRINTED 108
\weberpage{108}
To prove this, let us start from a value of the order \(m\) such that the effective number of the \(\Phi_m\) is equal to the virtual number \(N_m\), and such that \(m\geq n-3\).

From the number of the \(\Phi_m\) adjoint to \(F\) one deduces the number of the \(\Phi_{m-1}\) by subtracting the dimension \(r\) of the system cut out by the \(\Phi_m\) on a plane, increased by one. If this system is the complete system of the \(\varphi_m\) adjoint to the plane section, the effective number of the- \(\Phi_{m-1}\) again turns out equal to the virtual number; if instead there is a difference, the system of the \(\varphi_m\), sections of the \(\Phi_m\) on a plane, must then fail to be complete, and then %
\[
r<\frac{m(m+3)}{2}-d,
\]
and hence the effective number of the \(\Phi_{m-1}\) turns out greater than \(N_{m-1}\). This reasoning can be applied as long as \(m\geq n-3\), because under this hypothesis the complete system of the curves adjoint to the plane section of \(F\), of order \(n\), is regular, that is, it has the dimension \(\frac{m(m+3)}{2}-d\) and not greater, the double points of the section curve of \(F\) imposing independent conditions on the curves \(\varphi_m\) that must pass through them. If instead one takes \(m=n-4\), the system of the \(\varphi_{n-4}\) adjoint to the section of \(F\) turns out superabundant, and so the preceding reasoning is no longer applicable.

We shall call \textit{numerical} or \textit{arithmetic genus} \(p_a\) of the surface \(F\) the virtual number of linearly independent canonical curves, that is, the number \(N_{n-4}\) of the adjoint surfaces \(\Phi_{n-4}\); the surface genus of \(F\) itself, which we defined in the preceding chapter, will be distinguished, where necessary, from the arithmetic genus by designating it as the \textit{geometric genus}: \(p_g\).

By what we have seen above one can affirm that: \textit{the arithmetic genus of the surface \(F\) is in every case less than or equal to the geometric genus:}
\[
p_a\leq p_g.
\]

We shall shortly point out the invariance of the arithmetic genus with respect to birational transformations of the surface \(F\). Meanwhile let us write the value of this character by means of the postulation formulae established above:
\[
p_a=\binom{n-1}{3}-(n-4)d+2t+\tau+\varrho-1.
\]

Consequently the number of the linearly independent surfaces \(\Phi_{n-4+s}\) adjoint to \(F\) can be expressed as
% scan: s065R
% !PRINTED 109
\weberpage{109}
follows:
\[
\begin{aligned}
N_{n-4+s}&=\binom{n-1+s}{3}-(n-4+s)d+2t+\tau+\varrho-1\\
&=p_a+\binom{n-1+s}{3}-\binom{n-1}{3}-sd.
\end{aligned}
\]

Hence the linear system formed by the intersections of these adjoints \(\Phi_{n-4+s}\) with \(F\) will in general have the dimension:
\[
N_{n-4+s}-1-\binom{s-1}{3}
=p_a+\binom{n-1+s}{3}-\binom{n-1}{3}-\binom{s-1}{3}-sd-1,
\]
where \(\binom{s-1}{3}\) (which vanishes for \(s=1,2,3\)) is the number of linearly independent surfaces of order \(s-4(\geq 0)\), that is, the dimension of the system of the surfaces \(\Phi_{n-4+s}\) that pass through one and the same intersection with the surface \(F\) of order \(n\). The dimension of the system of the sections of the \(\Phi_{n-4+s}\), which we have written above, can be expressed in a very simple form by introducing the genus \(\pi_s\) of the curves cut out on \(F\) by the surfaces of order \(s\). Let us first compute this \(\pi_s\) by means of the genus \(\pi\) (\(=\pi_1\)) of the plane sections of \(s\); we shall have %
\[
\pi_s=s\pi+\frac{s(s-1)}{2}n-s+1,
\]
and, substituting for \(\pi\) its value
\[
\pi=\frac{(n-1)(n-2)}{2}-d,
\]
we shall obtain
\[
\begin{aligned}
\pi_s&=s\left[\frac{(n-1)(n-2)}{2}-d\right]+\frac{s(s-1)}{2}n-s+1\\
&=\frac{s(n-1)(n-2)}{2}+\frac{s(s-1)}{2}n-s+1-sd.
\end{aligned}
\]

But one has
\[
\begin{aligned}
\frac{s(n-1)(n-2)}{2}-\frac{s(s-1)}{2}n-s+1&{}\\ %
&=\binom{n-1+s}{3}-\binom{n-1}{3}-\binom{s-1}{3},
\end{aligned}
\]
hence
\[
N_{n-4+s}-1-\binom{s-1}{3}=p_a+\pi_s-1.
\]

We conclude that \textit{the linear system cut out on \(F\) by the surfaces \(\Phi_{n-4+s}\) of order \(n-4+s\), which is the system adjoint to the system}
% scan: s066L
% !PRINTED 110
\weberpage{110}
\textit{\(|C_s|\) cut out by the surfaces of order \(s\), for \(s\) sufficiently high, has the dimension}
\[
p_a+\pi_s-1,
\]
having denoted by \(\pi_s\) the genus of \(|C_s|\), the \(s\)-fold multiple of the system of plane sections.

\edsection{4. The arithmetic genus.}{4. The arithmetic genus.}{sec.IV.4}

Let us again refer to a surface \(F\) of ordinary space, having a certain order \(n\) and normal singularities. We propose to evaluate the dimension of the system \(|C'|\) adjoint to an irreducible linear system \(|C|\) belonging to \(F\). Let us recall first of all that the adjoint curves \(C'\) cut out canonical groups on a \(C\); by detaching \(C\) from \(|C'|\) one obtains the canonical system \(|K|\), of dimension \(p_g-1\); hence through the canonical group of intersection of a \(C'\) with a \(C\) there must pass \(\infty^{p_g}\) curves \(C'\). It follows from this that the dimension of the system \(|C'|\) is given in every case by
\[
r=p_g+\pi-1-\delta(C)
\]
where \(\delta(C)\) denotes the deficiency of the series cut out by the \(C'\) on a \(C\); in particular one has \(\delta(C)=0\) and
\[
r=p_g+\pi-1
\]
if \(|C'|\) cuts out on \(C\) the complete series \(g^{\pi-1}_{2\pi-2}\).

A priori the deficiency \(\delta(C)\) depends on the choice of the system \(|C|\); but we can recognize a fundamental inequality to which one is led by considering the sum of two irreducible linear systems (at least \(\infty^1\))\origfootnote{1}{s066L-1}{If they are two pencils, it is understood that they are distinct, so that \(|C+L|\) also turns out irreducible.}, \(|C|\) and \(|L|\).

One has in fact
\[
\delta(C)\leq\delta(C+L).
\]

We shall suppose for simplicity that the systems \(|C|\) and \(|L|\) are free of base points on the surface \(F\). Let us start from the fundamental relation
\[
|(L+C)'|=|L+C'|=|L'+C|.
\]

On the basis of this relation the system \(|C'|\) will be obtained from \(|(L+C)'|\) by detaching the curve \(L\), and this entails a number of linear conditions which is given by the dimension of the series cut out by \(|(L+C)'|\) on \(L\), increased by one.

% scan: s066R
% !PRINTED 111
\weberpage{111}
Now the series indicated on \(L\) is the sum of the canonical series \((LL')\) and of the series cut out by \(|C|\), which is certainly of order greater than zero, the systems \(|C|\) and \(|L|\) being both of dimension greater than or equal to one. %

This established, let us denote by \(\pi\) the genus of \(|C|\), by \(\varrho\) the genus of \(|L|\), and by \(i\) the number of the intersections \((LC)\). Then the genus of \(|L+C|\) will be
\[
\pi'=\pi+\varrho+i-1,
\]
and hence the dimension of \(|(L+C)'|\) will be
\[
p_g+\pi'-1-\delta(L+C).=p_g+\pi+\varrho+i-2-\delta(L+C). %
\]

The dimension of \(|C'|\), on the other hand, is
\[
p_g+\pi-1-\delta(C).
\]

The difference between the two dimensions is given, as has been said, by the dimension of the non-special series \(g\) cut out by \(|(L+C)'|\) on \(|L|\), increased by one. If the series \(g\) is complete, and hence of dimension
\[
2\varrho-2+i-\varrho=\varrho+i-2,
\]
it follows at once that
\[
\delta(C)=\delta(L+C);
\]
if instead \(g\) is deficient, its deficiency expresses the difference \(\delta(L+C)-\delta(C)\). In every case one has:
\[
\delta(C)\leq\delta(L+C).
\]

The reasoning still holds in the case where the two systems \(|C|\) and \(|L|\) have base points on \(F\). The base points common to the two systems, or the base points of \(|C|\) that do not belong to \(|L|\), have a priori no influence. If there is a base point of \(|L|\) that does not belong to \(|C|\), and therefore, at least virtually, not to \(|C'|\) either, the neighbourhood of the point (regarded as an infinitesimal curve) would have to be added to \(|C'|\), which does not increase its dimension; thus the preceding conclusion extends to this case as well.

Let us observe moreover that the hypothesis made on \(|L|\), of belonging to an irreducible system at least \(\infty^1\), contains more than is necessary: in reality \(L\) may also be an isolated curve composed of several parts \(L_1, L_2, \ldots\), provided that none of these is fundamental for the system \(|C|\), and provided that by adding \(L_1, L_2, \ldotp\ldotp\ldotp\ldotp\) successively to \(|C|\) one obtains irreducible linear systems of higher dimension. Indeed one will be able to detach \(L\) from \(|(L+C)'|\) by successively detaching its parts \(L_1, L_2 \ldots\), and if these are not fundamental for \(|C|\), the series cut out on them, whose dimension we need, will always be non-special, since they contain the canonical series.

% scan: s067L
% !PRINTED 112
\weberpage{112}
Now let us recall that we have computed the dimension of the system cut out on \(F\) by the adjoint surfaces for \(s\) sufficiently high: this system, which is the adjoint of the \(s\)-fold system of the plane sections, will here be conveniently denoted by \(|L_s|\). Denoting by \(\pi_s\) the genus of \(|L_s|\), we found that the dimension of the system \(|L'_s|\) is
\[
p_a+\pi_s-1.
\]

On the other hand it is expressed in terms of the geometric genus \(p_g\) and of the deficiency of the series cut out by \(|L'_s|\) on \(L_s\), being equal to
\[
p_g+\pi_s-1-\delta(L_s).
\]

From this one deduces that
\[
\delta(L_s)=p_g-p_a,
\]
\(s\) being, as has been said, a sufficiently high number.

This established, consider on \(F\) any irreducible system \(|C|\), at least \(\infty^1\), of a certain genus \(\pi\), which we shall suppose as usual to be free of base points. It is clear that for \(s\) sufficiently high there are surfaces of order \(s\) passing through a \(C\), and hence \(|C|\) is contained in \(|L_s|\), which indeed presents itself as the sum of \(|C|\) and of another irreducible system \(|M|=|L_s-C|\).

Therefore one deduces from this that
\[
\delta(C)\leq\delta(L_s).
\]

This conclusion extends easily to the case where \(|C|\) has base points, with certain multiplicities \(i_1, i_2, \ldots\). Indeed, referring to a sufficiently high value of \(s\), one can always require the system \(|L_s|\) to possess the same base points as \(|C|\) with the same multiplicities, thus constructing a system \(|\bar{L}_s|\) which contains \(|C|\), and indeed is such that the residual \(|M|=|\bar{L}_s-C|\) is an irreducible linear system. Now it is legitimate to suppose that the base points of \(|\bar{L}_s|\) impose \(\frac{i(i-1)}{2}\) independent linear conditions on the adjoint curves \(L'_s\), which must pass through an \(i\)-fold point with multiplicity \(i-1\). We shall thus have a linear system \(|\bar{L}_s|\) with the same base points as \(|C|\), for which one again has
\[
\delta(\bar{L}_s)=p_g-p_a.
\]
And one will deduce as before
\[
\delta(C)\leq\delta(\bar{L}_s).
\]

In this way we have proved that: \textit{on the surface \(F\) the deficiency of the} (canonical) \textit{series that the system adjoint to an irreducible linear system \(|C|\) cuts out on a generic \(C\) has a maximum value,}
% scan: s067R
% !PRINTED 113
\weberpage{113}
\textit{which constitutes an intrinsic character of the surface and which equals the difference \(p_g-p_a\) between its geometric genus and its arithmetic genus.} This \textit{maximum} is attained for a class of linear systems belonging to \(F\), which in every case includes the \textit{sufficiently large systems}, that is, those which contain within themselves a multiple of very high order of another irreducible system. Indeed this class is defined in such a way that « every irreducible linear system \textit{greater} than another one of the class (that is, containing it partially) belongs to the same class ».

A priori there can only be a second, restricted class of linear systems \(|C|\), for which the deficiency of the (canonical) series cut out on \(C\) by \(|C'|\) remains below the maximum.

From what has been said it follows that \textit{the arithmetic genus \(p_a\) of a surface}, defined directly by the postulation formulae (on the given surface or on a transform with normal singularities), constitutes a character \textit{invariant under birational transformations} of the surface itself, not only in the case where \(p_a=p_g\), but also when \(p_a<p_g\).

It will nevertheless be convenient to distinguish the two cases, designating as \textit{regular surfaces} those for which \(p_a=p_g\), and as \textit{irregular} (of irregularity \(p_g-p_a\)) those for which \(p_a<p_g\).

For both kinds of surfaces the result obtained contains the extension (from series on a curve to systems on a surface) of the Riemann-Roch theorem, for the moment limited to the linear systems \(|C'|\) that are adjoint to other irreducible systems \(|C|\) (of genus \(\pi\) and degree \(n\)). The dimension of \(|C'|\), or at least a lower bound for it, is computed by means of the genus of \(|C|\) from the formula:
\[
r'\geq p_a+\pi-1.
\]

But it is easy to express the same \(r'\) by means of the characters, genus \(\pi'\) and degree \(n'\), of the system \(|C'|\). Indeed — denoting by \(p^{(1)}\) the (relative) linear genus of the surface \(F\), and supposing, for simplicity of exposition, that \(|C|\) and \(|C'|=|C+K|\) are free of base points on \(F\) we shall have
\[
\begin{aligned}
\pi'&=\pi+p^{(1)}+2\pi-2-n-1=3\pi-3-n+p^{(1)}\\
n'&=n+(p^{(1)}-1)+4\pi-4-2n=4\pi-5-n'+p^{(1)}, %
\end{aligned}
\]
and hence
\[
\pi-1=n'-\pi'+1.
\]

Thus the dimension of the system \(|C'|\) will be
\[
r'\geq p_a+n'-\pi'+1,
\]
\printersignature{\textsc{Enriques F.} - \textit{Superficie algebriche.}\quad 8}
% scan: s068L
% !PRINTED 114
\weberpage{114}
and similarly, if \(|C|\) itself is the adjoint of another irreducible system, its dimension will be
\[
r\geq p_a+n-\pi+1.
\]

We shall see later how this \textit{fundamental inequality}, which constitutes the \textit{Riemann-Roch theorem for adjoint linear systems}, extends, with a small modification, to all the linear systems belonging to a surface of arithmetic genus \(p_a\). Hence we shall call \textit{regular} every complete non-special linear system whose dimension is expressed by the preceding \textit{equality}, and \textit{superabundant} a system for which the \textit{inequality} holds strictly.

\edsection{5. Complements.}{5. Complements.}{sec.IV.5}

In the preceding developments we have referred to a surface \(F\) of ordinary space endowed with normal singularities: under this hypothesis we have determined the postulation of the adjoint surfaces that are required to pass simply through the double curve of \(F\). But there is no difficulty in extending our considerations so as to contemplate also surfaces with higher singularities. If one wishes to follow the path indicated in the preceding sections, one will be led first of all to extend the postulation formulae, determining (with Noether) the number of conditions that a curve of order \(m\) and genus \(p\) imposes on the surfaces of sufficiently high order \(n\) that must contain it as an \(s\)-fold curve, which is\origfootnote{1}{s068L-1}{Cf. also \textsc{L. Campedelli}, \textit{Sulla postulazione di una curva \(i\)-pla.} (Rend. Circolo Matematico di Palermo, to. 55, 1931).}:
\[
\frac{s(s+1)}{6}[(3n-4s+4)m-(2s+1)(p-1)].
\]\origfootnote{2}{s068L-2}{To be modified in the presence of multiple points of the curve.}

And similarly one will have to evaluate the number of conditions imposed on the surfaces of a given order by the possession of an isolated multiple point, which for an \(s\)-fold point is
\[
\binom{s+2}{3}.
\]

But, as the singularities of \(F\) become more complicated, the computation of the postulation on which the dimension of the systems of adjoint surfaces depends becomes more difficult, as already appears from what we have said regarding the virtual influence on the genus of the proper multiple points of \(F\). It is therefore useful to point out a quicker path, independent of the nature of the singularities of the surface \(F\). We shall de%
% scan: s068R
% !PRINTED 115
\unskip\weberpage{115}%
note in every case by \(n\) the order of \(F\), by \(\pi\) the genus of its plane sections, by \(n_s\) and \(\pi_s\) the characters of the \(s\)-fold system cut out on \(F\) by the surfaces of order \(s\).

Let us observe first of all that, the arithmetic genus \(p_a\) being supposed determined (in relation to a transform of \(F\) with normal singularities), we know a priori the dimension of the systems of surfaces adjoint to \(F\), and we can affirm that the adjoint surfaces \(\Phi_{n-3+s}\), for \(s\) sufficiently large, cut out on \(F\) a complete linear system of dimension
\[
p_a+\pi_s-1,
\]
which in turn cuts out a complete series of order \(2\pi-2+sn\) on the plane section \(C\) of \(F\). Therefore the \(\Phi_{n-3+s}\) cut out on every plane the complete system of the curves of order \(n-3+s\) adjoint to the plane section \(C\).

This established, one sees that \textit{the dimensions of the linear systems of the surfaces \(\Phi_{n-3+s}\) adjoint to \(F\), for \(s\) sufficiently large, form an arithmetic progression of the third order having as its difference}
\[
\pi+sn+\frac{s(s-3)}{2}.
\]

This progression, extended to the smaller values of \(s\), serves to define in every case the \textit{virtual expression} of the number of linearly independent adjoint surfaces of a given order, and in particular the \textit{arithmetic genus} \(p_a\).

In this sense one may say that the aforesaid arithmetic progression, which is moreover a consequence of the postulation formulae, takes the place of the actual knowledge of these formulae\origfootnote{1}{s068R-1}{Independently of the circumstance that the base group assigned to the \(\Phi\) consists of the singularities of a surface \(F\) to which the \(\Phi\) must be adjoint, one can prove that the numbers expressing the dimension of the system of the surfaces \(\Phi\) of sufficiently high order constrained to pass with certain multiplicities through the points and the curves of an assigned base group always give rise to an analogous arithmetic progression of the third order, the knowledge of which serves, in many cases, to make up for the actual knowledge of the postulation formulae. Cf. \textsc{G. Castelnuovo}, \textit{Alcuni risultati sui sistemi lineari di curve appartenenti ad una superficie algebrica.} Mem. Soc. It. delle scienze (detta dei XL), series III, t. 10, 1896. Reprinted in the “Memorie scelte„ p. 335.}. %

Let us return to the consideration of the systems of the surfaces \(\Phi_{n-3+s}\) adjoint to \(F\).

Let us denote by \(h\geq0\) the smallest number such that for \(s\geq h\) these systems cut out on a plane the totality of the adjoint curves of the same order; this means that, unless \(h=0\), the system of the \(\Phi_{n-3+h-1}\) (which is of order greater than or equal to \(n-3\)) will cut out on a plane a non-complete system of adjoint curves, with the de%
% scan: s069L
% !PRINTED 116
\unskip\weberpage{116}%
ficiency \(\delta_{h-1}>0\). For \(s<h-1\) the systems of the \(\Phi_{n-3+s}\) (\(s=0, 1, \ldotp\ldotp\ldotp\ldotp\ h-2\)) will cut out on the same plane systems of adjoint curves, complete or not, which a priori may be supposed to have the deficiencies
\[
\delta_0\geq0,\quad \delta_1\geq0,\quad \ldotp\ldotp\ldotp\ldotp,\quad \delta_{h-2}\geq0. %
\]

Now we have seen that, in correspondence with the deficiency \(\delta_{h-1}\), and as compared with its virtual expression, the dimension of the system of the adjoint surfaces of order \(n-3+h-2\) increases precisely by \(\delta_{h-1}\), as does that of the system of curves cut out by them on \(F\), which will be
\[
p_a+\pi_{h-2}-1+\delta_{h-1}.
\]

Similarly, if \(\delta_{h-2}>0\), there will be an increase, in an analogous sense, by \(\delta_{h-2}\) in the dimension of the system cut out by the \(\Phi_{n-3+h-3}\), which will be
\[
p_a+\pi_{h-3}+\delta_{h-2}+\delta_{h-1}. %
\]

Finally the dimension of the canonical system, cut out by the \(\Phi_{n-4}\), will turn out to be
\[
p_g-1=p_a-1+\delta_0+\delta_1+\ldotp\ldotp\ldotp\ldotp+\delta_{h-1}.
\]

That is, the \textit{irregularity \(p_g-p_a\) of the surface equals the sum of the deficiencies of the systems of curves cut out on a plane by the surfaces adjoint to \(F\), of order}
\[
n-3+s\qquad \text{for}\qquad s=1,2,\ldotp\ldotp\ldotp\ldotp
\]

We can make the preceding statements somewhat more precise, proving first of all that \textit{if the adjoint surfaces of an order greater than or equal to \(n-2\) cut out on a plane the complete system of the adjoint curves of the same order, the same happens for all the linear systems of adjoint surfaces of higher order.} (We shall then say how the statement extends to \(s=0\), that is, to the case where the adjoint surfaces of order \(n-3\) cut out on a plane the complete system of the adjoints of this order).

To prove the theorem stated, we may pass from the systems of the surfaces \(\Phi\) to the systems of curves cut out by them on \(F\), and hence to the series that these systems cut out on a plane section \(C\) of \(F\). And it will suffice to establish that: « if the curves of the system \(|sC+C'|\), for \(s\geq1\), cut out on \(C\) the complete series, the series cut out on \(C\) by the curves of \(|(s+1)C+C'|\) is certainly complete as well ».

We shall suppose first that \(s>1\) and — for simplicity of exposition — let us take \(s=2\). The hypothesis is then that the curves of \(|2C+C'|\), sections of the adjoint \(\Phi_{n-1}\), cut out on \(C\) the complete series \(g\) of order \(2\pi-2+2n\) determined by the adjoint curves \(\varphi_{n-1}\). And evidently it suffices to show that the minimal sum of this \(g\) and of the series cut out by the lines is necessarily complete.

% scan: s069R
% !PRINTED 117
\weberpage{117}
In short, it is a matter of establishing (with \textsc{Castelnuovo}) the following:

\textit{Lemma on the completeness of the sum of two series belonging to a curve.}

On a curve of genus \(\pi\), add to a complete and non-special series \(g_m\) of a certain order \(m\) a series \(g_n\), at least \(\infty^1\), without fixed points, which is contained in \(g_m\) and gives as residual with respect to it a non-special series; the minimal sum of \(g_m\) and \(g_n\) is necessarily complete.

Indeed, add to \(g_m\), which is of dimension \(m-\pi\), two groups \(G\) and \(G'\) of \(g_n\), to be regarded as fixed: one will thus have two linear series of order \(n+m\) and of dimension \(m-\pi\), which have in common the series formed by the \(2n\) fixed points constituting the group \(G+G'\) and by the groups of \(|g_m-g_n|\), which is a series of order \(m-n\), supposed non-special, and therefore of dimension
\[
m-n-\pi.
\]

Therefore our two series, which have the groups \(G\) and \(G'\) as fixed parts, will be contained in a series of dimension at least
\[
m+n-\pi
\]
which therefore turns out to be complete, Q.E.D.

This lemma applies immediately to the case \(h\geq2\). But, with a slight modification, the reasoning extends to the case \(h=1\) (not, however, to \(h=0\)). It is a matter of proving that, by adding to the series \(g_{2\pi-2+n}\) cut out on \(C\) by the adjoint curves \(\varphi_{n-2}\) the series cut out by the lines, one always obtains, as minimal sum, the complete series cut out by the \(\varphi_{n-1}\).

Indeed, add to \(g_{2\pi-2+n}\) two fixed groups \(G\) and \(G'\) cut out by two lines \(a\) and \(a'\) of the plane; one will have two linear series of dimension \(\pi-2+n\), which have in common the canonical \(g^{\pi-1}_{2\pi-2}\) (of dimension not \(\pi-2\), but \(\pi-1\)): the minimal sum of the two series with the two fixed groups \(G\) and \(G'\) is now a series of order \(2\pi-2+2n\) and of dimension \(\pi-2+n-1\) instead of \(\pi-2+n\). This series, which is not complete, is cut out on \(C\) by the system of all the \(\varphi_{n-1}\) that pass through the point common to the lines \(a\) and \(a'\). Now the smallest linear system containing within itself two systems of \(\varphi_{n-1}\) with a different base point is the complete system of all the \(\varphi_{n-1}\). It follows therefore that the minimal sum of the series \(g_{2\pi-2+n}\) and of the series cut out by the lines of the plane is necessarily the complete series cut out by the \(\varphi_{n-1}\).

\edsection{6. Note on the conditions of regularity of a surface.}{6. Note on the conditions of regularity of a surface.}{sec.IV.6}

We have seen that, given a surface \(F\) of order \(n\) of ordinary space, which for simplicity we wish to suppose endowed with
% scan: s070L
% !PRINTED 118
\weberpage{118}
normal singularities, if the \(\Phi_{n-2}\) adjoint to it, of order \(n-2\), cut out on a plane the complete system of the curves adjoint to the section of \(F\), the same happens for the adjoint surfaces \(\Phi_{n-2+r}\) of order \(>n-2\); in this case, therefore, the irregularity \(p_g-p_a\) of the surface is equal to the deficiency of the (canonical) series cut out on the plane sections \(C\) of \(F\) by the \(\Phi_{n-3}\). In particular, if the \(\Phi_{n-3}\) cut out on the said \(C\) the complete canonical series, the irregularity \(p_g-p_a=0\), that is, the surface \(F\) is regular.

But the criterion of regularity, thus established by \textsc{Enriques}, can be reduced to another, simpler one.

\textit{The surface \(F\) of order \(n\) of ordinary space is regular \((p_a=p_g)\) if the surfaces \(\Phi_{n-3}\) of order \(n-3\) adjoint to it cut out on a plane section the complete canonical series.}

We shall give a brief sketch of the proof that \textsc{Castelnuovo}\origfootnote{1}{s070L-1}{\textit{Alcune proprietà fondamentali dei sistemi lineari di curve tracciati sopra una superficie algebrica.} Annali di Mat., series II, t. 25 (1897). Cf. \textit{Memorie scelte.} XXIII.} has given of this theorem. It is convenient first of all to establish the

\textit{Lemma.} – If a group of \(m<n\) points of our surface \(F\), chosen among the \(n\) points cut out by a line \(a\), belongs to a special series \(g_m^h\), of dimension \(h\geq1\), on a plane section \(C\) through \(a\), the same happens for every other section of \(F\) by any plane through \(a\).

It is an immediate consequence of the hypothesis that the \(\Phi_{n-3}\) cut out on the plane sections \(C\) of \(F\) the complete canonical series. Indeed, by the Riemann-Roch theorem, a group of the said \(g_m^h\) on \(C\) presents precisely \(m-h\) conditions to the canonical groups of \(C\), or equally to the \(C'\), which must contain it: but the number of such conditions is independent of the choice of a particular \(C\) of the pencil cut out on \(F\) by the planes through \(a\).

On the basis of the aforesaid lemma we are in a position to recognize that « if the \(\Phi_{n-3}\) cut out the complete canonical series on the plane sections \(C\) of \(f\), the characteristic series of the complete system \(|C|\) turns out complete ». %

We know that the complete system \(|C|\) can be constructed on \(F\) by means of the linear system of all the adjoints \(\Phi_{n-2}\) that pass through the section \(C'\) of a \(\Phi_{n-3}\) (residue theorem). It will therefore suffice to show that these \(\Phi_{n-2}\) through a \(C'\) cut out on a plane the complete series of the curves \(\varphi_{n-2}\) of order \(n-2\) adjoint to the section of \(F\).

To this end, let us try to construct a surface \(\Phi_{n-2}\) adjoint to \(F\) as the locus of a curve \(\varphi_{n-2}\) of the same order \(n-2\), adjoint to the section \(C\) by a plane \(\alpha\) varying in a pencil which has as its axis a generic line of space, \(a\), not meeting \(C'\). We shall require the \(\varphi_{n-2}\):
% scan: s070R
% !PRINTED 119
\weberpage{119}
1) to pass through the \(2\pi-2\) points of a canonical group, sec tion of \(C'\); %

2) to pass through \(s-1\) points, \(A_1, A_2, \ldotp\ldotp\ldotp\ldotp A_{s-1}\), chosen among the \(n\) points common to \(C\) and to the line \(a\), \(s\) denoting the dimension of the characteristic series of \(C\) made complete, that is, the dimension of the complete series to which belongs the \(g_n\) cut out on \(C\) by the lines of the plane \(\alpha\);

3) and finally to touch a plane \(\gamma\) passing through a point \(A_1\) of this group of \(s-1\) points and not containing \(a\).

Conditions 1) and 2) determine, in a generic plane through \(a\), a pencil of curves \(\varphi_{n-2}\), to which belongs the curve composed of the line \(a\) and of the \(\varphi_{n-3}\) passing through the canonical group of \(C\), section of \(C'\). Adding condition 3), \textit{one} curve \(\varphi_{n-2}\) remains determined in the plane; and, as this plane varies, the \(\varphi_{n-2}\) thus constructed will describe a surface \(\Phi_{n-2+h}\) adjoint to \(F\), which a priori may be of an order \(n-2+h>n-2\), passing \(h\) (\(>0\)) times through the line \(a\). For our purpose it is necessary to prove that one must have \(h=0\). We shall therefore suppose \(h>0\) and reduce this hypothesis to absurdity.

To this end, observe that, since \(h>0\), there will be, at a point \(A_s\) of \(a\), \(h\) planes tangent to \(\Phi_{n-2+h}\), any one of which, \(\bar{\alpha}\), will cut this surface in a curve composed of the line \(a\) counted \(h\) times and of a \(\varphi_{n-2}\) passing through the \(s\) points \(A_1A_2\ldotp\ldotp\ldotp\ldotp A_{s-1}A_s\). Now one can place \(A_s\) at one of the \(n-(s-1)\) intersections of the surface \(F\), or of the curve \(\bar{C}\), section of \(\bar{\alpha}\), with \(a\), outside \(A_1A_2\ldotp\ldotp\ldotp\ldotp A_{s-1}\): since the \(\bar{\varphi}_{n-3}\), section of \(\bar{\alpha}\), does not pass through \(A_1\) (\(a\) not meeting \(C'\)), the group of points \(A_1A_2\ldotp\ldotp\ldotp\ldotp A_s\) must belong not merely to one, but to a pencil of curves \(\varphi_{n-2}\) adjoint to the section \(\bar{C}\) of \(F\) with \(\bar{\alpha}\), and hence this group must belong to a special \(g_s^1\), of dimension (at least) equal to 1. But this cannot happen on a particular plane \(\bar{\alpha}\) through \(a\) without its happening for all the planes \(\alpha\). Therefore one deduces that \(h=0\), and so we succeed in constructing a surface \(\Phi_{n-2}\) adjoint to \(F\) which satisfies the conditions of:

1) passing through a curve \(C'\), section of \(f\) with an adjoint \(\Phi_{n-3}\); %

2) passing through \(s-1\) of the \(n\) points of intersection of \(F\) with a line \(a\);

3) touching a plane \(\gamma\) through one of these points.

By varying the points that figure in condition 2) one can obtain \(\infty^{s+1}\) \(\Phi_{n-2}\) through \(C'\), which cut out, on a generic plane \(\alpha\), \(\infty^s\) curves \(\varphi_{n-2}\), that is, the totality of the adjoint curves \(\varphi_{n-2}\) passing through the canonical group \(G_{2\pi-2}\), section of \(C'\). This means that the \(\Phi_{n-2}\) through \(C'\) cut out on \(F\) a linear system \(|C|\) of dimension \(s+1\), which has complete characteristic series, Q.E.D.

To establish Castelnuovo's theorem it now suffices to point out that: « on a curve \(C\) of genus \(\pi(>0)\) the minimal sum of the canonical
% scan: s071L
% !PRINTED 120
\weberpage{120}
series \(g_{2\pi-2}^{\pi-1}\) and of another complete series \(g_n^s\), without fixed points and not composed with an involution, is always complete ».

This follows from observing that a group \(G_n\) of the second series presents \(s\) independent conditions to a group of the same \(g_n^s\) that must contain it, and \(n-s\) conditions to a group of the canonical series, hence\origfootnote{1}{s071L-1}{Cf. \textsc{Castelnuovo}, \textit{Sui multipli di una serie lineare di gruppi di punti appartenente ad una curva algebrica.} Rendic. Circolo di Palermo, 1893. \textit{Memorie scelte}, VIII.} at least
\[
s+n-s-1=n-1
\]
conditions to a group of the minimal sum series \(g_{2\pi-2+n}=g_{2\pi-2}^{\pi-1}+g_n^s\); therefore the dimension of the latter exceeds the dimension of the canonical series by (at least) \(n-1\), and is thus
\[
(n-1)+\pi-1=n+\pi-2:
\]
this means that the minimal series \(g_{2\pi-2}^{\pi-1}+g_n^s\) has the dimension of the complete series \(g_{2\pi-2+n}^{\pi-2+n}\), that is, it coincides with it.

\textit{Note.} – Castelnuovo's theorem, which expresses the conditions of regularity of the surface in relation to the system \(|C|\) of its plane sections, extends easily to irreducible linear systems \(|C|\), even non-simple ones, as the author himself has pointed out. We shall therefore have:

\textit{The necessary and sufficient condition for a surface to be regular \((p_a=p_g)\) is that the system \(|C'|\) adjoint to an irreducible linear system \(|C|\), of genus \(\pi>0\), at least \(\infty^2\), cut out on the generic \(C\) the complete canonical series, that is, that it be of dimension \(p_g+\pi-1\).}

\textit{Every complete irreducible linear system belonging to a regular surface has complete characteristic series.}

We shall find this last theorem again later as a particular case of a property concerning the deficiency of the characteristic series of the complete linear systems belonging to irregular surfaces.

\edsection{7. Examples of irregular surfaces.}{7. Examples of irregular surfaces.}{sec.IV.7}

The preceding developments are illustrated by some examples. For the general surface of order \(n\), without singularities, the adjoint surfaces \(\Phi_{n-4}\) constitute the totality of the surfaces of order \(n-4\); the effective number of linearly independent ones among these is equal to the virtual number
\[
p_g=p_a=\binom{n-1}{3}=\frac{(n-1)(n-2)(n-3)}{6}\text{\origfootmark{2}{s071L-2}}.
\]\origfoottext{2}{s071L-2}{If one wishes to apply the general postulation formulae, one must regard the double curve of order \(n=0\) as being of genus \(\varrho=1\).}

% scan: s071R
% !PRINTED 121
\weberpage{121}
Thus surfaces that are general of their order are regular; in particular planes, quadrics and cubic surfaces that are not cones, and hence also all rational surfaces, have the surface genus
\[
p_g=p_a=0.
\]

Now, if one requires a surface of order \(n\) to possess a double curve of relatively small order, one likewise obtains regular surfaces; if instead one wants the surface to possess a double curve of rather high order, one finds a priori a number of conditions exceeding the number of parameters on which the surface of the given order depends; therefore it is difficult, by this route, to obtain examples of irregular surfaces.

But let us consider a cone \(F\) of an order \(n>2\), and of genus \(p>0\), and let \(O\) be its vertex (an \(n\)-fold point); here there are certainly no adjoint surfaces of order \(n-4\), which would have to pass through the point \(O\) with multiplicity \(n-2>n-4\), so that (as we already know)
\[
p_g=0.
\]

For the same reason there are also no surfaces \(\Phi_{n-3}\) adjoint to \(F\), which would have to cut out on it \(p_a+p\) linearly independent curves adjoint to the plane sections, so that
\[
p_a+p\leq 0,
\]
that is,
\[
p_a\leq -p.
\]

But on the other hand there do exist adjoint cones \(\Phi_{n-2}\) (with \(O\) \((n-2)\)-fold) which cut out on a plane the complete system of the curves of order \(n-2\) adjoint to the section of \(F\). Therefore our results entitle us to say that the arithmetic genus
\[
p_a=-p.
\]

One concludes that \textit{ruled surfaces of genus \(p\) have arithmetic genus \(-p\) and irregularity}
\[
p_g-p_a=p.
\]

This theorem is proved by transforming the ruled surface into a cone; but even without carrying out this transformation one can repeat the preceding reasoning, based on the facts that: \(1^\circ)\) there exist on the surface no curves adjoint to the plane sections \(C\); and \(2^\circ)\) the curves of \(|C+C'|\) (sections of the adjoints \(\Phi_{n-2}\)) cut out on a \(C\) the complete series.

The example of ruled surfaces has been known since the first studies of A. \textsc{Cayley}\origfootnote{1}{s071R-1}{. \textit{On the deficiency of certain surfaces.} Math. Annalen, Bd. III.} %
% scan: s072L
% !PRINTED 122
\weberpage{122}
(1871) and of M. \textsc{Noether}, who recognized that the virtual expression of the genus becomes negative in this case. Subsequently the first example of an irregular surface for which \(p_g>0\) was indicated by G. \textsc{Castelnuovo}\origfootnote{1}{s072L-1}{Cf. G. \textsc{Castelnuovo}, \textit{Osservazioni intorno alla geometria sopra una superficie.} Rend. Ist. Lombardo, 1891. Note I. “ Memorie scelte „ no. XVII.}.

Let a surface \(F_6\) of order 6 be required to possess 8 triple points \(A_1,A_2,\ldots,A_8\) in arbitrarily assigned position in space. The \(F_6\) depend on 83 constants, and the imposition of a triple point entails 10 linear conditions; hence the \(F_6\) passing triply through \(A_1,\mathinner{\ldotp\ldotp\ldotp\ldotp},A_8\) will be \(\infty^3\), and they will split into the triples of quadrics (belonging to a pencil) passing through those 8 points. But instead of taking \(A_1,A_2,\mathinner{\ldotp\ldotp\ldotp\ldotp},A_8\) in arbitrary position, suppose that they constitute the base group of a net of quadrics \(\Phi_2\). In this case, by linearly combining the triples of \(\Phi_2\), one will obtain a linear system of irreducible \(F_6\), of dimension 9\origfootnote{2}{s072L-2}{One can prove that the system of the \(F_6\) with the eight triple points cannot have dimension \(r>9\), by evaluating \(r\) by means of the dimension of the system that these \(F_6\) cut out on a quadric of the net.}. %

Now the \(F_6\) so constructed turn out to be of genus \(p_g=3\), having as adjoints of order \(6-4=2\) the quadrics \(\Phi_2\) of the net under consideration. But the 8 base points \(A_1,A_2,\mathinner{\ldotp\ldotp\ldotp\ldotp},A_8\) impose on the \(\Phi_2\) seven independent linear conditions instead of eight, which is the virtual number of these conditions. Therefore the arithmetic genus of the \(F_6\) is %
\[
p_a=p_g-1=2.
\]

It is interesting to note that the one constructed is the first case of a series of analogous examples. In general one can construct irreducible surfaces of order \(2n\) \((n\geq 3)\) possessing 8 \(n\)-fold points at the base points of a net of quadrics: these \(F_{2n}\) turn out irregular, of genus
\[
p_g=\frac{n(n-1)}{2}
\]
and
\[
p_a=n-1.
\]

Here it is of interest to observe that the quadrics of the net cut out on an \(F_{2n}\) a linear system composed of elliptic quartics, and these quartics form on \(F_{2n}\) an irrational pencil of genus
\[
p_g-p_a=\frac{(n-1)(n-2)}{2};
\]
indeed the genus of the pencil equals the genus of the osculating cone of \(F_{2n}\) at one of the points \(A\), which is a base point for the pencil of the quartics themselves.

% scan: s072R
% !PRINTED 123
\weberpage{123}
Let us add that the series of the aforesaid examples, indicated by \textsc{Castelnuovo}, can be generalized, as L. \textsc{Campedelli} has shown (1936)\origfootnote{1}{s072R-1}{L. \textsc{Campedelli}, in Atti R. Accademia di Torino, 1936.}.

The construction of an irreducible \(F_{2n}\) with 8 \(n\)-fold points depends in fact on that of a linear pencil of elliptic curves of an order \(m\) dividing \(n\) \((n=ms)\), possessing 8 \(s\)-fold base points on a quadric: this group of base points always belongs to an elliptic quartic, and gives on it a group arising from the division by \(n\) of the series cut out on the curve by all the surfaces of order \(2n\). The pencil of elliptic curves in question can be projected onto the plane from one of its base points into a Halphen pencil of curves \(C_{3m}\) with 9 \(m\)-fold base points.

Starting from the pencil indicated on a quadric, one obtains irregular surfaces \(F_{2n}\), with 8 \(n\)-fold points, which always possess an irrational pencil of genus \(p_g-p_a\).

\textit{Note.} – On the basis of the theorem of \S~6, which affirms the completeness of the characteristic series of a complete linear system on a regular surface, one recognizes in general that « the possession of an irrational pencil (of genus \(\pi>0\)) of curves implies that the surface is irregular ».

More precisely, we shall find that the surface has irregularity \(p_g-p_a\geq\pi\). But of this later.

\edsection{8. The Riemann-Roch theorem for surfaces: systems ampler than the canonical system.}{8. The Riemann-Roch theorem for surfaces: systems ampler than the canonical system.}{sec.IV.8}

We have recognized that a linear system of degree \(n\) and genus \(\pi\), which is the adjoint of another irreducible system on a surface of arithmetic genus \(p_a\), has the dimension
\[
r\geq p_a+n-\pi+1,
\]
and this fundamental inequality, even under the restrictive conditions supposed here, constitutes an extension to surfaces of the well-known Riemann-Roch theorem relating to curves. It is convenient first of all to account for the meaning of the restriction that \(|C|\) be the adjoint system of another irreducible system. For simplicity of exposition let us assume that \(|C|\) is free of base points on the surface \(F\), and that the same holds for the system \(|L|\) of which it is the adjoint. Denoting by \(|K|\) the canonical system we shall have
\[
|C|=|L+K|.
\]

% scan: s074L
% !PRINTED 124
\weberpage{124}
Thus the condition imposed on \(|C|\), of being the adjoint of another system, translates into that of containing the canonical system, giving rise to a residual system \(|L|\) of order greater than zero, or, as we shall say, of being \textit{ampler} than the canonical system.

A system \(|C|\) ampler than the canonical system \(|K|\) can always be regarded as the adjoint of the system
\[
|L|=|C-K|,
\]
but one cannot say a priori that \(|L|\) is irreducible. Indeed there are cases in which, even when the dimension of \(|C|\) is increased by multiplying it by an arbitrarily large number, the system \(|L|\) always turns out reducible, containing fixed components not connected with the variable components. This happens indeed if \(|C|\) is the system of plane sections of a surface \(F\) which possesses a proper multiple point influencing the genus, for example an isolated multiple point of order \(\nu\geq 3\); %
which represents a fundamental curve \(\theta\) of \(|C|\), although here it is easy to recognize that the inequality in question (dimension of \(|C|\geq p_a+n-\pi+1\)) holds a fortiori. Indeed, denoting by \(\bar{n}\) and \(\bar{\pi}\) the characters of the system \(\bar{C}=|C-\theta|\), %
we shall have that \(|\bar{C}|\) is the adjoint of the irreducible system \(|\bar{L}|=|L-\theta|\):
\[
|\bar{C}|=|\bar{L}'|,
\]
and hence the dimension of \(|\bar{C}|\) is
\[
\bar{r}\geq p_a+\bar{n}-\bar{\pi}+1;
\]
now the degree of \(\theta\) is \(-\nu\) and its genus \(\varrho=\frac{(\nu-1)(\nu-2)}{2}\geq 1\), hence we shall have
\[
\begin{aligned}
n&=\bar{n}-\nu\\ %
\bar{\pi}&=\bar{\pi}+\varrho+\nu-1 %
\end{aligned}
\]
and consequently the dimension of \(|C|\), \(r>\bar{r}\), will be
\[
r>p_a+n-\pi+1.
\]

But, in order to proceed rigorously with the greatest generality, we shall prove that every linear system \(|C|\) ampler than the canonical system \(|K|\) has the dimension \(r\geq p_a+n-\pi+1\), resorting to other considerations which hold equally for reducible or irreducible systems \(|C|\).

Given on the surface \(F\) a curve \(C\), irreducible or reducible, we can always suppose that the surfaces adjoint to \(F\), \(\Phi_{n-4+s}\) (of sufficiently high order \(n-4+s\)), that are constrained to pass through \(C\) do not in consequence pass through another curve or through points
% scan: s074R
% !PRINTED 125
\weberpage{125}
outside \(C\) and the double curve\origfootnote{1}{s074R-1}{This follows from observing that in the postulation formulae relating to a composite curve (or to a curve and a group of points) each component brings an effective number of conditions.}. This established, the curve \(C\) will be contained in the system \(|L'_s|\) cut out on \(F\) by the \(\Phi_{n-4+s}\), and will give with respect to it an irreducible residual system (without base points)
\[
|D|=|L'_s-C|.
\]

Now, denoting by \(r'_s\) the dimension, and by \(\pi'_s\) and \(n'_s\) the genus and the degree of \(|L'_s|\), we have found
\[
r'_s\geq p_a+n'_s-\pi'_s+1
\]
(indeed in this case precisely the sign of equality holds); by detaching a curve \(D\), one will find that the residual system \(|C|=|L'_s-D|\) must have the dimension
\[
r\geq p_a+n-\pi+1,
\]
provided one knows that the series cut out by \(|L'_s|\) on \(D\) is non-special, as can be affirmed in the case where \(|C|\) is ampler than the canonical system \(|K|\), whose curves cut out on \(D\) groups residual to the characteristic series of \(D\).

Let us carry out the computation by introducing the characters genus \(\varrho\) and degree \(\nu\) of \(D\) and the number \(t\) of the intersections of \(C\) and \(D\). The series cut out by \(|L'_s|\) on \(D\) will have the order \(\nu+t\) and (if it is complete) the dimension \(\nu+t-\varrho\); hence the detaching of \(D\) lowers the dimension \(r'_s\) of \(|L'_s|\) by at most \(\nu+t-\varrho+1\); it follows that
\[
r\geq r'_s-(\nu+t-\varrho+1).
\]

On the other hand
\[
\begin{aligned}
n'_s&=n+\nu+2t\\
\pi'_s&=\pi+\varrho+t-1
\end{aligned}
\]
and hence
\[
n'_s-\pi'_s=n-\pi+\nu-\varrho+t+1;
\]
consequently
\[
r\geq p_a+n-\pi+1.
\]

We have thus proved that: \textit{for irreducible or reducible linear systems of degree \(n\) and genus \(\pi\) ampler than the canonical system, the dimension is always}
\[
r\geq p_a+n-\pi+1.
\]

This inequality constitutes the so-called Riemann-Roch theorem for surfaces, which will later be extended — in the way that we shall see — to arbitrary linear systems.

% scan: s075L
% !PRINTED 126
\weberpage{126}
Meanwhile, referring to the \textit{systems ampler than the canonical system}, we shall distinguish:

1) the \textit{regular} systems, for which the equality holds
\[
r=p_a+n-\pi+1
\]

2) and the \textit{superabundant} systems, for which instead
\[
r>p_a+n-\pi+1.
\]

We have already seen that certain regular systems exist, for such are the adjoints of the multiples of very high order of an irreducible system \(|L|\).

The result that we have obtained leads at once to noteworthy consequences.

Let \(|C|\) be an \textit{irreducible linear system} of genus \(\pi\) and degree \(n\), \textit{ampler than the canonical system and regular} on the surface \(F\), and let there be added to it an irreducible curve \(D\) of genus \(\varrho\) and degree \(\nu\), which intersects the \(C\) in a certain number \(m\) \((>0)\) of points.

The degree and the genus of \(|C+D|\) are respectively
\[
n+\nu+2m
\]
and
\[
\pi+\varrho+m-1;
\]
and hence their difference is
\[
(n-\pi)+(\nu-\varrho)+m+1>n-\pi
\]
if
\[
m\geq\varrho-\nu.
\]

One deduces that \textit{the system \(|C+D|\) has, under this hypothesis, dimension greater than}
\[
p_a+n-\pi+1
\]
which is the dimension of \(|C|\); and therefore it is certainly \textit{irreducible}.

Moreover this system \(|C+D|\) will be \textit{regular} and must cut out on \(D\) a complete (non-special) series, since otherwise by detaching \(D\) one would obtain a superabundant system \(|C|\).

Note that this enlargement of \(|C|\) would be impossible for \(m=0\), that is, if the curve \(D\) were fundamental for \(|C|\); in this case one must therefore have \(\nu-\varrho<0\), which follows a priori from observing that the fundamental curve \(D\), of genus \(\varrho\geq 0\), must have negative degree: \(\nu<0\), \(-\nu\) being the number of intersections of \(D\) with its residual curves (of \(|C-D|\)).

If in the preceding observation one interchanges the roles of \(|C|\) and \(|D|\), one is led to affirm: \textit{given on the surface \(F\) a linear system} (virtually free of base points) \(|C|\), \textit{reducible in any manner,}
% scan: s075R
% !PRINTED 127
\weberpage{127}
\textit{one can always add to it a regular irreducible linear system \(|D|\) so as to obtain an irreducible system \(|D+C|\) that is likewise regular.}

It suffices to take as the system \(|D|\) the multiple \(|L_s|=|sL|\) of a regular system \(|L|\) without fundamental curves, and therefore having with the \(C\) (or with their components) a number of intersections that grows large with \(s\). In case the curves \(C\) are reducible, e.g., \(C=C_1+C_2+\mathinner{\ldotp\ldotp\ldotp\ldotp}\), one will first add a multiple of \(|L|\) to \(C_1\), and then the irreducible system (as ample as one wishes) so obtained to \(C_2\), and so on, until one arrives at an irreducible system \(|sL+C|\).

\edsection{9. Elimination of exceptional curves.}{9. Elimination of exceptional curves.}{sec.IV.9}

Let us take up again the considerations of the preceding section regarding the addition of a curve to a regular linear system \(|C|\) on a surface \(F\). And let us take as the curve to be added to \(|C|\) an exceptional curve \(\omega\) of genus \(\varrho=0\) and of degree \(\nu\geq -1\). If it is not fundamental for \(|C|\), and hence meets the \(C\) in \(m>0\) points, by adding it to \(|C|\) one will always obtain a linear system of greater dimension, irreducible and regular like \(|C|\). And the curves of this sum system will have \(m+\nu\) intersections with it. Therefore \(\omega\) can be added successively to \(|C|\), and one will obtain an unlimited sequence of systems of increasing dimension if \(\nu\geq 0\), that is, in the case where \(\omega\) is an exceptional curve of the second kind. If instead \(\omega\) is of the first kind \((\nu=-1)\), the sequence of systems obtained will stop at the system \(|C+m\omega|\), for which \(\omega\) becomes fundamental.

Now consider, for every linear system \(|C|\) of degree \(n\) and genus \(\pi\) belonging to the surface \(F\), the difference \(2\pi-2-n\) (the number of intersections of its curves with the virtual canonical curves), which gives an \textit{additive character} \(a(C)\) of the systems on the same \(F\). This difference decreases by at least one each time \(\omega\) is added to \(|C|\) or to one of the systems successively obtained in this way, and hence, if the sequence of these is unlimited, it certainly ends by becoming negative, so that there will be on \(F\) systems of genus \(\pi\) and of degree \(n>2\pi-2\). %

This established, we wish to suppose that \textit{the surface \(F\) does not contain systems for which the difference \(-a(C)=n-(2\pi-2)>0\).} Under this hypothesis it will not contain exceptional curves of the second kind and can contain only a finite number \(s\) of exceptional curves of the first kind.

For simplicity we shall speak of these assuming that they are irreducible exceptional curves, with the proviso of pointing out afterwards what must be said in the excluded case (which corresponds to exceptional curves that are images of infinitely near points).

% scan: s076L
% !PRINTED 128
\weberpage{128}
Assume that the (irreducible regular) system \(|C|\) is free of base points on \(F\), not only virtually but also effectively. If, as has been said above, the \(C\) cut the exceptional curve \(\omega\) in \(m\) points, \(\omega\) will be a fundamental curve (of degree \(-1\), and therefore lowering by \(1\) the degree of the residual curves) for the system \(|C+m\omega|\). And on transforming the surface \(F\) into another \(F'\) by means of this system, \(\omega\) will be changed into a simple point \(O\) of \(F'\).

Now each of the other \(s-1\) exceptional curves on \(F\) will be without connection with \(\omega\), and will therefore give on \(F'\) again an exceptional curve of the first kind (or a simple point of the surface); otherwise one would obtain on \(F'\) exceptional curves of degree \(\geq 0\) (that is, of the second kind), and hence one would construct on the surface \(F'\), and also on \(F\), systems \(|C|\) for which the character
\[
-a(C)>0.
\]

Thus the transformation of \(F\) into \(F'\) succeeds in decreasing the number of exceptional curves belonging to the surface, so that after \(s\) transformations we shall obtain a surface \(\Phi\) entirely free of exceptional curves.

To tell the truth, we have simplified our discussion by supposing that the exceptional curves on \(F\) are all irreducible. But if there is on \(F\) an exceptional curve of the first kind composed of \(h\) parts (given in a certain order \(\omega,\,\omega' \ldotp\ldotp\ldotp\ldotp \omega^{(h-1)}\)), the last of these components always constitutes an exceptional curve of degree \(\nu=-1\), and when this curve is reduced to a point by the transformation that changes \(F\) into \(F'\), the component preceding it becomes an exceptional curve (of the first kind) of degree \(\nu=-1\). Thus the exceptional curve composed of \(h\) parts (with which \(h\) exceptional curves of \(F\) are formed, one irreducible and \(h-1\) reducible) becomes on \(F'\) an exceptional curve composed of \(h-1\) parts. And finally it disappears after \(h\) transformations, reducing (on \(\Phi\)) to a group of \(h\) infinitely near simple points. %

We conclude therefore: \textit{A surface can in general be transformed into another one free of exceptional curves. This transformation fails only for the surfaces} (of genus \(p_g=0\), bigenus \(P=0\), etc.) \textit{on which there exist linear systems \(|C|\) for which the character} \(-a(C)=n-(2\pi-2)>0\), that is, as we shall see, for \textit{the surfaces that belong to the family of ruled surfaces.}

\edsection{10. Completeness of the series cut out by a small linear system on the curves of a large system.}{10. Completeness of the series cut out by a small linear system on the curves of a large system.}{sec.IV.10}

The extension of the Riemann-Roch theorem to arbitrary linear systems (even to the « special » ones, which are less ample than the
% scan: s076R
% !PRINTED 129
\weberpage{129}
canonical system) will be given in the simplest way by relying on the completeness of the series that a « \textit{small} » linear system cuts out on the generic curve of a « \textit{normally large} » system. « Small » and « large » are relative terms which are defined as follows: one regards as « large » with respect to a given linear system \(|C|\) (virtually free of base points on the surface \(F\)) a linear system \(|D|\) which contains \(|C|\) in such a way that the difference \(|D-C|\) equals the multiple \(|sL|\), according to an integer \(s\) as large as one wishes, of an irreducible linear system \(|L|\) without base points, or else constitutes a system « ampler » than \(|sL|\). And \(|D|\) will be said to be \textit{normally} large with respect to \(|C|\) if the preceding condition is satisfied with regard to a system \(|L|\) \textit{free of fundamental curves}, so that, by choosing \(s\) sufficiently large, one can arrange that the system \(|B|=|D-C|\) be irreducible\origfootnote{1}{s076R-1}{This is no longer true if \(|L|\), and hence \(|L_1|=|D|\), possesses a fundamental curve \(\theta\) which is not fundamental for \(|C|\), since then \(\theta\) splits off from \(|D-C|\) as a fixed part. Hence the fundamental lemma also no longer holds for \(|C|\) in comparison with this (\textit{not} normally) large system \(|D|\): for example the canonical system \(|K|\) no longer cuts out, in general, on a \(D\) the complete series (cut out by \(|K+\theta|\)).}. %

We wish to prove the \textit{Fundamental lemma.}

\textit{A complete linear system \(|C|\), irreducible or reducible, cuts out a complete series on the irreducible curve of a linear system \(|D|\) which is normally large with respect to it.}

To establish this lemma, take on a \(D\) an arbitrary group \(G\) of the complete series \(CD\), cut out by \(|C|\), and associate with it a group \(G'\), whose existence will be proved below, satisfying the two following properties:

1) \(G+G'\) is a group of the characteristic series (generally not complete) cut out on \(D\) by the other curves of \(|D|\);

2) \(G'\) is a group of the series cut out on \(D\) by an irreducible curve \(B\) of \(|D-C|\).

Then \(G+G'\) will constitute the base group of a pencil of curves \(D\), and with respect to this pencil there will be an irreducible fundamental curve \(B\) passing through \(G'\); consequently the said pencil will contain a curve \(D\) split into \(B\) and a \(C\) through \(G\); every group \(G\) of the complete series \(CD\) on \(D\) will therefore be a group cut out by a \(C\).

It is only necessary to recognize that, \(|D|\) being very large with respect to \(|C|\), there actually exists for every group \(G\) of the complete series \(CD\) a group \(G'\) associable with it which satisfies properties 1) and 2). To this end it is convenient to show that the two linear series, a priori not complete, \(DD-G\) and \(BD\), cut out on \(D\) respectively by the \(D\) through \(G\) and by the \(B\), and contained in the same series
\printersignature{\textsc{Enriques F.} - \textit{Superficie algebriche.}\quad 9}

% scan: s077L
% !PRINTED 130
\weberpage{130}
complete, certainly have in common an effective linear series, described by the groups \(G'\). And therefore it suffices to establish that the deficiencies of these two series are less than their dimensions, because then, as one easily sees, the sum of the dimensions of the two series is greater than the dimension of the complete series. %

Now this last condition is certainly satisfied, since one recognizes that, for \(s\) large and tending to infinity, the dimension of the characteristic series of \(|D|\), on \(D\), and hence also the dimension of the series \(DD-G\), becomes infinite of the \(2^\circ\) order, like \(\dfrac{s^2}{2}\,n\) (\(n\) the degree of \(|L|\)), and the same happens for the dimension of the series \(BD\); on the other hand the defect of the characteristic series \(DD\), and hence also that of the series \(DD-G\), even if it may increase indefinitely with \(s\), cannot a priori become infinite of order higher than the first.

We shall limit ourselves to carrying out this simple computation in the case where
\[
|D|=|L_s| \qquad |C|=|L|.
\]

Denoting by \(n\) and \(\pi\) the degree and genus of \(|C|\), and by \(n_s\) and \(\pi_s\) the degree and genus of \(|L_s|\), the dimension \(r_s\) of \(|L_s|\) is given by the Riemann-Roch theorem:
\[
r_s\geq p_a+n_s-\pi_s+1,
\]
where
\[
n_s=s^2n,\qquad \pi_s=s\pi+\frac{s(s-1)}{2}n-s+1,
\]
and hence
\[
r_s\geq \frac{s^2}{2}n-s\pi-\frac{s}{2}n+s-1;
\]
therefore the effective dimension of the characteristic series \(|DD|=(L_sL_s)\) is %
\[
r_s-1\geq \frac{s^2}{2}n+\ldotp\ldotp\ldotp\ldotp
\]
neglecting the terms of order \(s\). And since the group \(G=CL_s\) comprises \(sn\) points, the dimension of the effective series \(DD-G\) will also be
\[
\frac{s^2}{2}n+\ldotp\ldotp\ldotp\ldotp,
\]
neglecting the terms of order lower than the \(2^\circ\).

On the other hand, the dimension of the series cut out on \(D=L_s\) by the curves \(B\) of \(|L_{s-1}|=|L_s-C|\) will also be
\[
\frac{s^2}{2}n+\ldotp\ldotp\ldotp\ldotp
\]

Now one can evaluate the defect of these series, or at least an upper bound for it, which depends on the defect of the character%
% scan: s077R
% !PRINTED 131
\unskip\weberpage{131}%
\unskip istic series \(DD\), and hence on its index of speciality. If this series is complete (of dimension \(r_s-1=p_a+n_s-\pi_s\)) its index of speciality will be \(p_a\); if instead it has the defect \(d\), its index of speciality will be
\[
i=d+p_a.
\]

In any case this index \(i\) will be less than the order of the residual series, that is, than the number of intersections of \(D=L_s\) with a canonical curve \(K\):
\[
DK=s(2\pi-2-n),
\]
and consequently the defect of the aforesaid series \(DD\), and hence also the defects of \(DD-G\) and of \(BD\), cannot become infinite with \(s\) of an order higher than the first. \hfill Q.E.D.

The same computation extends, with slight modifications, to the case in which, \(|L|\) being a system without fundamental curves, one again has \(|D|=|L_s|=|sL|\) and \(|L|=|B+C|\), or else in which
\[
|D|=|L_s+E|,
\]
and \(|C|\) coincides with \(|L|\) or is partially contained in it:
\[
|L|=|B+C|.
\]

One likewise arrives at the conclusion that the complete system \(|C|\) cuts out a complete series on the curve \(D\) of a normally large system.

\edsection{11. The Riemann-Roch theorem for arbitrary linear systems.}{11. The Riemann-Roch theorem for arbitrary linear systems.}{sec.IV.11}

In place of the linear system \(|C|\) considered in the preceding section, take the canonical system \(|K|\), supposing therefore \(p_g>0\), and construct an irreducible linear system \(|D|\) which is normally large with respect to \(|K|\). We may assert that \(|K|\) cuts out, on a generic curve \(D\) of \(|D|\), the complete series residual to the characteristic series.

This assertion also holds for \(p_g=0\), in the sense that « on a surface of genus \(p_g=0\), the \textit{characteristic series} of a system \(|D|\), normally large, is a \textit{non-special} series ». To prove it, take an auxiliary linear system \(|C|\) and consider, \textit{together} with it, its adjoint system \(|C'|\). If \(|D|\) is sufficiently large with respect to \(|C|\) and to \(|C'|\), both these systems cut out complete series on a generic \(D\): say \(g\) and \(g'\). If the characteristic series \(DD\) is special, and hence contained in the canonical series \(DD'\), where \(|D'|=|D+C'-C|\), the series \(g\) will be contained in \(g'\), and precisely a group \(G\) of \(g\) will form part of a group \(G'=G+G_1\) of \(g'\). It is a matter of recognizing that, under this hypothesis, \(|C'|\) also contains \(|C|\), and hence \(p_g>0\), contrary to the hypothesis.

% scan: s078L
% !PRINTED 132
\weberpage{132}
For this purpose it suffices to observe that the curve \(C\) passing through \(G\) (and, in the case of reducibility, each component of it) has with \(D\) (normally large) a number of intersections greater than \(CC'\), and hence forms part of the \(C'\) through \(G+G_1\).

Thus we have established that \textit{on a surface of genus \(p_g\geq 0\), the characteristic series of a normally large linear system} (in comparison with the canonical system) \textit{has the index of speciality \(p_g\).}

Now consider on the surface \(F\) a linear system \(|C|\) of any speciality, of index \(i\) \((\geq 0)\), and a system \(|D|\) normally large with respect to it. We wish to recognize that « the index of speciality of the series cut out on a \(D\) by the system \(|D+C|\) is equal to the index of speciality of \(|C|\), that is, to the number \(i\) of linearly independent canonical curves \(K\) of which a \(C\) forms part ».

Indeed, since the series cut out by \(|K|\) on \(D\) is complete, through a group \(G\) of the series \(CD\) on a \(D\) there pass as many linearly independent canonical curves \(K\) as the index of speciality of the series cut out on the same \(D\) by \(|D+C|\); each of these \(K\) has with \(C\) (or, in the case of reducibility, with any component of \(C\)) a number of intersections \(CD>CK\), and hence contains \(C\) entirely.

Now we are in a position to deduce the dimension \(r\) of the system \(|C|\) of degree \(n\), genus \(\pi\) and index of speciality \(i\), which is obtained by detaching a \(D\) from the system \(|D+C|\); for this one must subtract from the dimension of this latter system, which is regular, the dimension of the series cut out on \(D\) by \(|D+C|\), whose index of speciality \(i\) is known, but which may fail to be complete. In this way one finds a lower bound for \(r\): precisely, one finds, as for the systems larger than the canonical system: \(r\geq p_a+n-\pi+1\) when \(i=0\), that is, for all \textit{non-special} systems \(|C|\), and in general
\[
r\geq p_a+n-\pi+1-i
\]
\textit{for special systems of index \(i\).}

\textit{This inequality constitutes the most general expression of the Riemann-Roch theorem for surfaces.}

Since it has been recognized for complete linear systems \(|C|\), irreducible or reducible, without base points on the surface \(F\), one deduces that it is valid for any linear system whatever, endowed with assigned base points: indeed an \(i\)-fold base point imposed on the curves \(C\) diminishes the degree by \(i^2\), the genus by \(\dfrac{i(i-1)}{2}\), and the dimension by
\[
\frac{i(i+1)}{2}=i^2-\frac{i(i-1)}{2}
\]
at most.

% scan: s078R
% !PRINTED 133
\weberpage{133}
The linear systems \(|C|\) for which the dimension \(r\) turns out equal to \(r=p_a+n-\pi+1-i\) will be called \textit{regular}; \textit{superabundant} if
\[
r>p_a+n-\pi+1-i.
\]

The Riemann-Roch theorem for curves of genus \(p\) gives us the dimension of the complete special series \(g_n^r\), expressed in terms of the dimension of its residual series with respect to the canonical \(g_{2p-2}^{p-1}\). This theorem has as an obvious complement the observation that the residual series of the \(g_n^r\) has the order \(2p-2-n\). For special linear systems belonging to a surface, too, one can complete the theorem of RiemannRoch, by expressing the characters genus and degree of the residual system by means of those of the given system. Indeed, let \(|C_1|\) and \(|C_2|\) be two linear systems residual to each other with respect to the canonical system \(|K|\). Let \(n_1\) and \(\pi_1\) be the characters, genus and degree, of the first system; \(n_2\) and \(\pi_2\) the characters of the second; and denote by \(m\) the number of intersections of \(|C_1|\) and \(|C_2|\). Introducing the linear genus \(p^{(1)}\) of the surface, we can express the genus and the degree of \(|K|\) by means of the formulae %
\[
\begin{aligned}
p^{(1)}&=\pi_1+\pi_2+m-1\\
p^{(1)}-1&=n_1+n_2+2m.
\end{aligned}
\]

But one can compute \(m\) by taking into account that \(|K|\) cuts out on \(C_1\) a group of
\[
2\pi_1-2-n_1=n_1+m
\]
points, so that
\[
m=2\pi_1-2-2n_1.
\]
We shall therefore have
\[
\begin{aligned}
p^{(1)}&=\pi_1+\pi_2+2\pi_1-2-2n_1-1=3\pi_1+\pi_2-2n_1-3\\
p^{(1)}-1&=n_1+n_2+4\pi_1-4-4n_1=4\pi_1-3n_1+n_2-4,
\end{aligned}
\]
from which one obtains
\[
\begin{aligned}
\pi_2&=p^{(1)}-3\pi_1+2n_1+3\\
n_2&=p^{(1)}+3n_1-4\pi_1+3.
\end{aligned}
\]

\textit{These formulae serve to express the characters of the system \(|C_2|\), residual to the special system \(|C_1|\) with respect to the canonical system, by means of the characters of \(|C_1|\).}

Note that the dimensions of the two systems \(|C_1|\) and \(|C_2|\), residual to each other, are given by
\[
\begin{aligned}
r_1&=p_a+n_1-\pi_1-r_2+\omega_1\\
r_2&=p_a+n_2-\pi_2-r_1+\omega_2,
\end{aligned}
\]
% scan: s079L
% !PRINTED 134
\weberpage{134}
where \(\omega_1\) and \(\omega_2\) denote the respective \textit{superabundances}. Hence from the preceding formulae one obtains
\[
\omega_1=\omega_2.
\]
that is: \textit{two special systems residual to each other with respect to the canonical system have the same superabundance.}

\edsection{12. Deficiency of the characteristic series.}{12. Deficiency of the characteristic series.}{sec.IV.12}

On the surface \(F\) of arithmetic genus \(p_a\) and geometric genus \(p_g\geq p_a\) consider an irreducible linear system \(|C|\) of genus \(\pi\), degree \(n\) and dimension \(r\), which for simplicity we suppose non-special. On the generic curve \(C\) we have two linear series, residual to each other with respect to the canonical series; they are: the characteristic series \(g_n^{r-1}\), cut out by the other curves \(C\), and the series of order \(2\pi-2-n\) and dimension \(p_g-1\) cut out by the canonical curves \(K\). If these series have respectively the deficiencies \(d\) and \(d'\), they are contained in two complete series of dimension
\[
r-1+d \qquad \text{and} \qquad p_g-1+d'.
\]

Applying the Riemann-Roch theorem on the curve, one has
\[
r-1+d=n-\pi+p_g-1+d'+1
\]
that is
\[
r=n-\pi+p_g+d'-d+1.
\]
Hence if \(d'=0\)
\[
r=p_g-d+n-\pi+1.
\]

Now, if \(|C|\) is a regular linear system, normally large with respect to the canonical system, its dimension is
\[
r=p_a+n-\pi+1
\]
and on the other hand we know that the canonical curves \(K\) cut out on \(C\) the complete series, so that one has \(d'=0\).

One deduces that
\[
d=p_g-p_a.
\]

In words: \textit{the characteristic series of a complete irreducible linear system, normally large, which is regular}\origfootnote{1}{s079L-1}{It will be seen later that « every normally large linear system is regular ».}, \textit{has the deficiency}
\[
d=p_g-p_a.
\]

% scan: s079R
% !PRINTED 135
\weberpage{135}
This being established, we wish to prove that the characteristic series of any irreducible linear system \(|C|\) belonging to the surface \(F\) always has a deficiency
\[
d\leq p_g-p_a.
\]

To obtain this result we shall establish the following

\textit{Lemma}: If a linear system
\[
|C|=|C_1+C_2|
\]
which is the sum of two irreducible systems cuts out the complete series on a curve of the system \(|C_1|\), the deficiency \(d_1\) of the characteristic series of \(|C_1|\) is less than or equal to that, \(d_2\), of \(|C_2|\):
\[
d_1\leq d_2.
\]

Denote by \(n_1,\ \pi_1,\ n_2,\ \pi_2\) the characters of \(|C_1|\) and \(|C_2|\), by \(r_1\) and \(r_2\) their respective dimensions, and by \(s\) the number of intersections of a \(C_1\) with a \(C_2\); finally let \(r\) be the dimension of \(|C|\). By hypothesis the series cut out by \(|C|\) on a \(C_1\) is complete. Hence the curves \(C\) passing through the group \(G\) of the \(s\) points common to a \(C_1\) and a \(C_2\) will also cut out on \(C_1\) a complete series of dimension \(r_1-1+d_1\).

On the other hand the dimension of the complete series cut out by \(C\) on \(C_1\) is
\[
r-r_2-1
\]
and hence the dimension of the series cut out on \(C_1\) by the curves \(C\) through \(G\) will be
\[
r-r_2-1-(s-\varepsilon),
\]
where \(s-\varepsilon\) denotes the number of linearly independent conditions which the group \(G\) of \(s\) points presents to the curves \(C\) that are to contain it. Thus we shall have
\[
r_1-1+d_1=r-r_2-1-(s-\varepsilon),
\]
that is
\[
d_1=r-(r_1+r_2)-(s-\varepsilon).
\]

This equality also holds independently of the hypothesis that \(|C|\) cuts out the complete series on \(|C_1|\), provided one replaces \(d_1\) by the deficiency \(d'_1\) of the characteristic series of \(|C_1|\) with respect to that cut out by the curves \(C\) through \(G\) \((d'_1\leq d_1)\); applying the same reasoning to the system \(|C_2|\) one finds:
\[
d_2\geq d'_2=r-(r_1+r_2)-(s-\varepsilon).
\]

From the comparison of the expressions of \(d_1\) and \(d_2\) there results in any case
\[
d_1\leq d_2 \qquad\qquad \text{Q.E.D.}
\]

% scan: s080L
% !PRINTED 136
\weberpage{136}
This being established, given on the surface \(F\) any irreducible linear system \(|C|\), at least \(\infty^1\), of degree \(n\) and genus \(\pi\), construct on \(F\) another irreducible linear system \(|D|\) which is regular and sufficiently large that the deficiency of its characteristic series equals \(p_g-p_a\).

\ifshownotes\clearmargin\fi\edFlush
We know (\S\ 8) that by adding \(|C|\) to \(|D|\) (since the number of intersections \((CD)\) is greater than \(n-\pi\)) one again obtains a regular system \(|C+D|\), which cuts out on \(C\) the complete series. Therefore the deficiency \(d\) of the characteristic series of \(|C|\) must be
\[
d\leq p_g-p_a.
\]

In words: \textit{the characteristic series of an irreducible linear system on a surface of genera \(p_a\) and \(p_g\geq p_a\) has the deficiency}
\[
d\leq p_g-p_a.
\]

Thus \textit{the irregularity} \(p_g-p_a\) of a surface is also \textit{defined} as the \textit{maximum value of the deficiency of the characteristic series} for the irreducible linear systems belonging to it.

\textit{On a regular surface} \((p_g=p_a)\) \textit{every complete irreducible linear system has a complete characteristic series.} Or, from the projective point of view: a regular surface, normal in a certain space, has normal curves as hyperplane sections. If the sections of a normal surface are not normal, one deduces
\[
p_a<p_g.
\]

We can illustrate what has been said by the example of \textit{ruled surfaces.}

We know that the ruled surfaces (with sections) of genus \(p\) have the geometric genus \(p_g=0\) and the arithmetic genus \(p_a=-p\). Therefore the normal ruled surfaces of order \(n\) must in general belong to spaces of dimension \(r=n-2p+1\); and indeed \textsc{C. Segre}\origfootnote{1}{s080L-1}{\textsc{C. Segre}, \textit{Courbes et surfaces réglées.} Math. Annalen, 1892.}, as early as 1890, recognized that a non-special ruled surface of order \(n\) and genus \(p\), which is not a cone, belongs precisely to a space \(S_{n-2p+1}\); on the other hand a \textit{special} ruled surface, whose plane or hyperplane sections are special curves of index \(i\), is normal in a space of \(n-2p+1+i\) dimensions (always excluding the case of the cone, in which the system of sections turns out superabundant); thus the characteristic series of the system of hyperplane sections has in every case the deficiency \(p\).

\textit{Remark.} -- The theorem on the deficiency of the characteristic series of complete linear systems led \textsc{G. Castelnuovo} to recognize that: \textit{every surface containing an irrational pencil of curves, of genus \(p\) \((>0)\), is irregular, of irregularity}
\[
p_g-p_a\geq p. %
\]

% scan: s080R
% !PRINTED 137
\weberpage{137}
This theorem falls within another, more general one, established by \textsc{F. Enriques}\origfootnote{1}{s080R-1}{\textit{Una proprietà delle serie continue di curve appartenenti ad una superficie algebrica regolare.} Rendic. Circolo Mat. Palermo, t. XIII, 1899.}: \textit{a surface which contains an irreducible algebraic system \(\{C\}\) of non-equivalent curves is irregular}; more precisely, \textit{if \(\{C\}\) contains \(\infty^d\) \((d>0)\) non-equivalent curves, the irregularity of the surface is}
\[
p_g-p_a\geq d.
\]

\textsc{Enriques} recognized the deficiency of the (canonical) series cut out on a \(C\) by the adjoint system \(|C'|\). On that occasion \textsc{Castelnuovo} --- to whom he had communicated the result --- observed the deficiency of the characteristic series of the complete linear system to which a \(C\) belongs. The proof takes its simplest form in the exposition of \textsc{F. Severi}\origfootnote{2}{s080R-2}{\textit{Osservazioni sui sistemi continui di curve appartenenti ad una superficie algebrica.} Atti R. Accad. Scienze di Torino, 1904.}. Indeed it suffices to note that the complete continuous system \(\{C\}\) will be formed of \(\infty^d\) linear systems \(|C|\) having, in general, a certain dimension \(r\): hence the characteristic series of \(|C|\) on a generic \(C\) will have the dimension \(r-1\) and will be contained in the characteristic series of the continuous system, of dimension \(r-1+d\), that is, in the series cut out on \(C\) by the infinitely near curves of \(\{C\}\).

Let it be explicitly noted that the theorem relating to the irrational pencil falls within the one thus established: it suffices to note that, by subtracting from a sufficiently large linear system the groups of \(p\) curves of a pencil of genus \(p\), one obtains a continuous system formed of \(\infty^p\) non-equivalent linear systems.

\edsection{13. Historical note.}{13. Historical note.}{sec.IV.13}

The consideration of the arithmetic genus of a surface presented itself from the very first studies of \textsc{A. Cayley}, \textsc{H. G. Zeuthen} and \textsc{M. Noether}. After \textsc{Noether} had defined the (geometric) genus \(p\) of a surface \(F_n\) of order \(n\) as the number of linearly independent adjoint surfaces \(\Phi_{n-4}\), it was natural to seek an expression of \(p\) by means of the characters of the double curve of \(F_n\); and for this purpose the postulation formulae of \textsc{Cayley} were at hand. This same mathematician\origfootnote{3}{s080R-3}{\textsc{A. Cayley}, \textit{On the deficiency of certain Surfaces.} Math. Annalen, Bd. III (1871).} had occasion to note that, in the case of ruled surfaces of genus \(p\), the calculation just mentioned leads to the negative value \(-p\), whereas the geometrically defined genus cannot, by its nature, descend below zero and, in this case, is precisely zero.

% scan: s081L
% !PRINTED 138
\weberpage{138}
Meanwhile \textsc{Zeuthen}, from a study of the correspondences between algebraic surfaces, was led to prove (albeit with some restriction as regards the singularities) that « the arithmetic (virtual) expression of the genus, according to \textsc{Cayley}, remains invariant under birational transformations of the surface, and this independently of its (presumed) geometric meaning and of the proof of the invariance of the canonical system furnished by \textsc{Noether} ».

These results were interpreted by \textsc{Noether}\origfootnote{1}{s081L-1}{\textit{Mathematische Annalen.} Bd. VIII (1875).} in the sense that the genus of a surface $F_n$ of order $n$ can be defined in two ways: \textit{geometrically}, as the effective number of linearly independent adjoint surfaces $\Phi_{n-4}$, and \textit{numerically}, by means of the postulation formulae; one thus introduces the \textit{geometric genus} and the \textit{arithmetic genus} of the surface, both a priori invariant under birational transformations, but which are presumed to coincide with each other, except that the arithmetic genus can become negative (for ruled surfaces) and therefore seems to absorb the other into a character of more general meaning.

But the presumption of \textsc{Noether} that, for non-negative $p_a$, one must have $p_a=p_g$, came to fall as a result of the new examples of non-regular surfaces which were added to the example of ruled surfaces. As has been said, the first example of these was brought forward by G. \textsc{Castelnuovo} in a Note of the Istituto lombardo of 1891. Other examples emerged from the researches of E. \textsc{Picard} and of G. \textsc{Humbert} on surfaces possessing integrals of total differentials of the first kind, and in particular on hyperelliptic surfaces, for which $p_g=1$ and $p_a=-1$.

Subsequent studies have brought to light large classes of irregular surfaces, as was indicated in the Remark of § 7.

Meanwhile the researches begun in 1893 and continued in 1896 (in the cited « Introduzione ») led F. \textsc{Enriques} to recognize the geometric meaning of the difference $p_g-p_a$ in relation to the deficiency of the (canonical) series which the system $|C'|$ adjoint to an irreducible linear system $|C|$ cuts out on the generic curve of $|C|$: whence results the entirely general proof of the invariance of $p_a$, the invariance of $p_g$ being already established. For \textsc{Enriques} the irregularity $p_g-p_a$ is the maximum value of the deficiency indicated, while it remains a priori possible that there exist systems $|C|$ for which the said deficiency is $<p_g-p_a$. On this point E. \textsc{Picard} completed the theory, proving (by transcendental means) that the system $|C'|$ adjoint to an irreducible simple system $|C|$ (at least $\infty^3$) is always regular, and therefore $p_g-p_a$ is defined simply as the deficiency of the series cut out on the plane section of a
% scan: s081R
% !PRINTED 139
\weberpage{139}
surface $F_n$ by the adjoint surfaces $\Phi_{n-3}$\origfootnote{1}{s081R-1}{E. \textsc{Picard}, \textit{Sur quelques questions se rattachant à la connexion linéaire dans la théorie des fonctions algébriques de deux variables indépendantes.} See Crelle, Bd. 129, 1905. Cf. \textsc{Picard} et \textsc{Simart}, \textit{Traité}, t. II, p. 437. An algebro-geometric proof (and also an extension of the conditions of validity) of the theorem was given by \textsc{Severi}, in the Note \textit{Sulla regolarità del sistema aggiunto} (Lincei, 8 Nov. 1908), based however on the theorem of the completeness of the characteristic series of a continuous system}. As regards this theorem, which for $p_g=p_a$ includes the theorem of Castelnuovo set forth in § 6, see the developments of the following section.

The extension of the Riemann-Roch theorem to surfaces is first found sketched by M. \textsc{Noether} in 1886\origfootnote{2}{s081R-2}{\textit{Extension du théorème de Riemann-Roch aux surfaces algébriques.} Comptes Rendus de l’Académie des Sciences, Paris, t. 103.}. The author considers a complete irreducible linear system $|C|$ of dimension $r$, degree $n$ and genus $\pi$, having a certain index of speciality $i\geq 0$, on a surface of (geometric) genus $p$. The characteristic series of $|C|$, cut out by $|C|$ on a generic $C$, is a $g_n^{r-1}$ residual to the series $g_{2\pi-2-n}^{p-i}$ marked out by the canonical system $|K|$; and \textsc{Noether}, without declaring it explicitly, postulates that the two series (cut out by complete systems) are complete; under this hypothesis the Riemann-Roch theorem relative to the curve $C$ furnishes the equality
\[
r=p+n-\pi+1-i.
\]

The question is taken up again by \textsc{Enriques} as early as the « Ricerche » of 1893, in which he succeeds in proving that for \textit{regular} surfaces of genus $p>1$ the characteristic series of a complete system is always complete, if such is the characteristic series of the canonical system; under this hypothesis (which will later turn out to involve a superfluous restriction) there is thus proved the inequality
\[
r\geq p+n-\pi+1,
\]
and the author points out that it cannot be converted into an equality, that is, that there certainly exist \textit{superabundant} systems, the presence of fundamental curves of genus $\varrho>0$ being sufficient for this purpose. %

In the « Introduzione » of 1896, with the new definition of the arithmetic genus, \textsc{Enriques} gave, for all surfaces, regular and non-regular, the extension of the Riemann-Roch theorem to adjoint systems, which means — substantially — to the systems larger than the canonical system.

To extend the theorem to arbitrary linear systems it seemed necessary to seek the maximum of the deficiency of the characteristic series and to prove that it is given by $p_g-p_a$. This result
% scan: s082L
% !PRINTED 140
\weberpage{140}
was in fact achieved by G. \textsc{Castelnuovo} in two memoirs\origfootnote{1}{s082L-1}{\textit{Alcuni risultati sui sistemi lineari di curve appartenenti ad una superficie algebrica.} Memorie della Società Italiana delle Scienze, detta dei XL (1896, vol. X). \textit{Alcune proprietà fondamentali dei sistemi lineari di curve tracciati sopra una superficie algebrica.} Annali di Matematica, series II, vol. XXV. Reproduced in “ Memorie scelte „, XXII and XXIII.} of the years 1896 and 1897.

The procedure used for this purpose is based on a suitable extension of the idea which leads to finding the simplest conditions of regularity of a surface, as we have set forth in § 6. These are methods which still seem susceptible of fruitful application, but they turn out rather laborious.

A notable simplification was brought into this order of questions by F. \textsc{Severi}, in two Notes of 1903 and 1905\origfootnote{2}{s082L-2}{\textit{Sulla serie caratteristica di un sistema lineare di curve appartenente ad una superficie algebrica.} Rendiconti R. Accademia dei Lincei, series V, vol. XII (1903). \textit{Sul teorema di Riemann-Roch e sulle serie continue di curve appartenenti ad una superficie algebrica.} Atti Accademia delle Scienze di Torino, vol. XL (1905).}; particularly in the second Note the author succeeds in inverting the order of the deductions, first justifying the Riemann-Roch theorem and then drawing from it the maximum of the deficiency of the characteristic series of irreducible systems.

The starting point of these developments is the well-known $Af+B\varphi$ theorem of Noether, extended to surfaces in the following form:

Every surface $f_m(xyz)=0$, of order $m$, not containing the plane $z=0$, which passes through the points of the group $G$ common to the plane $z=0$ and to the curve $D_s$ of order $ns$, the complete intersection of the surface $f_n$ of order $n$, devoid of proper multiple points, with a surface $f_s$ of order $s$, has an equation of the form
\[
f_m=Af_n+Bf_s+zf_{m-1}=0,
\]
where $f_{m-1}$ is a polynomial of order $m-1$ and $A$ and $B$ are polynomials in $x$, $y$, of degree $m-n$ and $m-s$ respectively.

Assuming outright that the surface $f_n$ is endowed with ordinary singularities, it follows from this theorem that the surfaces $f_m$ of any order, passing through the double points of a $D_s$, cut out on $D_s$ the complete series: indeed one proves that the proposition is true for $m-1$ if it is true for $m$.

From this \textsc{Severi} derives the

\textit{Lemma.} – For $s>n-4$ the surfaces $\Phi_{n-4}$ adjoint to $f_n$ cut out on the curve $D_s$ (the complete intersection of $f_n$ and of $f_s$) a complete series.

Indeed the surfaces of order $n-4$ passing through the double points
% scan: s082R
% !PRINTED 141
\weberpage{141}
of $D_s$ come to contain the double curve of $f_n$ entirely, and hence turn out to be adjoint to it.

As a consequence of the preceding lemma one has that « the canonical curves $K$, and then also the curves of a complete (special) system $|C|$ contained in $|K|$, cut out on the curves $D_s$ the complete series. And therefore, by detaching $D_s$ from the complete system $|D_s+C|$, one obtains for the dimension of $|C|$ the inequality which expresses the Riemann-Roch theorem. %

On the other hand the same lemma stated above, once it has been verified that the complete system $|D_s|=|sD|$ (without proper fundamental curves) is regular (cf. § 15), implies that the characteristic series of a $D_s$ has the deficiency $p_g-p_a$.

From this one deduces (as is done here in the preceding section) that the characteristic series of any irreducible linear system $|C|$, at least $\infty^1$, has a deficiency $\leq p_g-p_a$.

The reasoning which leads to this result is the very one that \textsc{Enriques} had employed in the « Ricerche » to pass from the characteristic series of the canonical system to that of any other linear system, on a regular surface, and that \textsc{Castelnuovo} had taken up again in the most general form, as an entirely secondary means, to justify the theorem on the deficiency of the characteristic series in the particular cases which escaped his direct proof (non-simple systems).

The new proof of the Riemann-Roch theorem for arbitrary linear systems, as it is offered in this treatise\origfootnote{1}{s082R-1}{Cf. F. \textsc{Enriques}, \textit{Sur l’extension du théorème de Riemann-Roch aux systèmes linéaires de courbes appartenant à une surface algébrique.} Bull. des Sciences Mathématiques, 1940 (Cf. in particular no. 5).}, is at least as simple in conception as that of \textsc{Severi}. It has, in our eyes, the advantage of indicating the deep reason for the fact that the canonical curves $K$ (or the curves of a special system) cut out a complete series on the $D_s$: a reason which lies not in the $D_s$ being complete intersections of the given surface $F_n$ with an $F_s$, but rather in $|K|$ being very small with respect to the system $|D_s|$. Moreover our procedure affords the most direct proof of the theorem in view, remaining within the same order of ideas in which the preliminary case of adjoint systems, or of systems larger than the canonical system, is treated, without any need to resort to some new argument, as seemed necessary when following the paths of \textsc{Castelnuovo} and of \textsc{Severi}.

And further, since the Riemann-Roch theorem for a linear system $|C|$ is obtained from the consideration of the series which a system $|L|=|C+D|$ cuts out on a curve $D$, which will be detached
% scan: s083L
% !PRINTED 142
\weberpage{142}
from $|L|$, it remains justified \textit{for the system} $|C|=|L-D|$, however \textit{reducible} this may be.

Now the consideration relating to the detaching of a $D$ from $|L|$ is already met with in the « Ricerche » of \textsc{Enriques} (of 1893)\origfootnote{1}{s083L-1}{L. c., IV, 4.}, and therefore one may hold that \textsc{Enriques} himself, and \textsc{Castelnuovo} with him, knew very early, at least implicitly, the possibility of extending the Riemann-Roch theorem to reducible systems. In keeping with their inductive turn of mind, they were to make this knowledge explicit on the occasion of concrete cases. The extension to reducible systems is formulated in a joint memoir of theirs of 1900\origfootnote{2}{s083L-2}{\textsc{Castelnuovo-Enriques}, \textit{Sopra alcune questioni fondamentali nella teoria delle superficie algebriche.} Annali di Mat., t. VI, series III, no. 4.}.

The Riemann-Roch theorem also affords notable criteria for the effective existence of \textit{virtual curves}, as we shall say in the following section.

\edsection{14. Virtual curves.}{14. Virtual curves.}{sec.IV.14}

Let us bear in mind that the calculation of the dimension of a linear system $|C|=|A-B|$, residual to a $B$ with respect to a regular and non-special system $|A|$, was obtained by evaluating the dimension of the series cut out by $|A|$ on this curve $B$; it appears from this that the Riemann-Roch theorem, which expresses the result of this calculation, must also afford a criterion for the effective existence of the virtual curves of $|A-B|$. This is first of all evident in the case in which $|A|$ cuts out on a $B$ a non-special series.

The formula which gives the dimension of the system adjoint to another of genus $\pi$:
\[
r\geq p_a+\pi-1
\]
is a first application of this criterion, which already presented itself in the « Introduzione » of \textsc{Enriques} of 1896. A second application is the formula which gives the dimension of the $i$-canonical system on a surface of linear genus $p^{(1)}$:
\[
P_i-1\geq p_a+\frac{i(i-1)}{2}(p^{(1)}-1)
\qquad (i>1)
\]
as it is met with in a Note of \textsc{Castelnuovo} of 1897\origfootnote{3}{s083L-3}{\textit{Sul genere lineare di una superficie etc.} Rend. R. Acc. dei Lincei, 1897; reproduced in « Memorie scelte », XXIV.} and in the memoir of \textsc{Castelnuovo-Enriques} of 1900\origfootnote{4}{s083L-4}{\textit{Sopra alcune questioni fondamentali della teoria delle superficie.} Ann. di Matematica, 1900.}: indeed the
% scan: s083R
% !PRINTED 143
\weberpage{143}
$i$-canonical system $|iK|$ can be deduced from a regular system $|A|=|B+iK|$, the $i$-th adjoint of a regular system $|B|$, by detaching $|B|$ from it, and $B+iK$ does indeed cut out a non-special series on a $B$. %

A third application of the Riemann-Roch theorem as a criterion of effective existence of virtual curves is given by \textsc{Severi} (1905), and contemplates in general the case where $|A|$ cuts out on a curve $B$ a series of any speciality: as we shall say further on.

If we wish to examine the question of the existence of the curves of a difference system $|C|=|A-B|$, it is legitimate to suppose that the two systems $|A|$ and $|B|$ which define the virtual curves $C$ are both \textit{irreducible, regular and without fundamental curves}, and also without base points, on a surface $F$ which is itself without singularities; for one can always add to the said systems $|A|$ and $|B|$ another system which is regular and without fundamental curves.

This being established, let us observe that:

1) The effective existence of the virtual curves of $|C|=|A-B|$ requires first of all that the degree of the minuend system be greater than or equal to that of the subtrahend
\[
A^2\geq B^2.
\]

Indeed if $A^2<B^2$, the curves $C$ must have a negative number of intersections with the curves of the system $|A+B|$, since $(A+B)(A-B)=A^2-B^2$, and therefore they would have to be fixed components of this system, which, by the hypotheses admitted, is irreducible.

2) In the second place, for curves of $|A-B|$ to exist it is necessary that the number of intersections of the curves $A$ and $B$ lie between the degrees of the two systems $|A|$ and $|B|$:
\[
B^2\leq AB\leq A^2.
\]

If this were not so, the curves $C$ would have a negative number of intersections with the curves $B$ or with the curves $A$, since one has respectively: for $AB<B^2$
\[
B(A-B)<0
\]
and for $AB>A^2$
\[
A(A-B)<0.
\]

The conditions 1) and 2) are \textit{necessary} both for the effective existence of $|C|=|A-B|$ and for that of a multiple $|i(A-B)|$. Moreover, if
\[
A^2>B^2
\]
it is necessary that
\[
B^2<AB<A^2,
\]
otherwise the curves $C$ (or the $iC$) — which are not, in the absolute sense, curves of order zero — would have zero intersections with the curves $B$ or with the curves $A$ and
% scan: s084L
% !PRINTED 144
\weberpage{144}
would therefore be fundamental curves for the corresponding system (which we have supposed not to possess such curves).

To the two necessary conditions stated above there corresponds the \textit{sufficient condition} for the existence of a multiple of $|C|$:

\textit{If the virtual degree of $|C|=|A-B|$ is positive and $A^2>AB>B^2$, the curves of the multiple system $|iC|$, for $i$ sufficiently large, always have effective existence.} %

Indeed, consider the series cut out by the system $|iA-(i-1)B|$ on a curve $B$ (supposing that such a system effectively exists); this series is of order
\[
AB+(i-1)(AB-B^2),
\]
which, for $i$ sufficiently large. comes to exceed the order of the canonical groups of $B$; therefore the said series is non-special, and one can calculate the dimension of the residual system %
\[
|iA-(i-1)B-B|=|iA-iB|
\]
by means of the formula
\[
r_i\geq p_a+n_i-\pi_i+1,
\]
where
\[
\begin{aligned}
n_i&=i^2n\\
\pi_i&=i\pi+\frac{i(i-1)}{2}n-i+1 %
\end{aligned}
\]
($n=(A-B)^2=$ virtual degree of $|C|$; $\pi=$ virtual genus of $|C|$).

From the preceding formula there results $r_i>0$ for sufficiently high values of $i$, and hence the effective existence of the curves of $|iC|$.

However, in the preceding argument the existence of the system $|iA-(i-1)B|$ was supposed a priori; this supposition is justified for the values of $i$ which make $r_i>0$, since under such conditions one can calculate successively the virtual dimensions of $|iA-B|$, $|iA-B-B|=|iA-2B|$ ...., up to $|iA-(i-1)B|$, which all turn out a fortiori greater than $r_i$: indeed for the said values of $i$, the dimension of the series cut out by $|iA-(i-2)B|$ on $B$ will a fortiori be greater than that of the canonical series, and hence non-special, etc. %

The criterion of existence of virtual curves which we have thus proved applies to the pluricanonical system of a surface for which the linear genus is $p^{(1)}>1$, and hence the virtual degree of the canonical system is
\[
p^{(2)}=p^{(1)}-1>0.
\]

One thus finds the formula already written above
\[
P_i\geq p_a+\frac{i(i-1)}{2}(p^{(1)}-1)+1.
\]

% scan: s084R
% !PRINTED 145
\weberpage{145}
Note that the criterion does not lead to any decisive answer concerning the existence of the multiple $|iC|$ of a virtual system of degree $n=0$, and thus does not allow one to assert a priori the existence of pluricanonical curves for $p^{(1)}=1$; although this existence is later proved, by means of other considerations, for surfaces without exceptional curves (not belonging to the family of ruled surfaces).

\edFlush
The case in which the degree of $|C|$ is $n<0$ remains a fortiori undecided; one cannot say whether or not the two systems $|A|$ and $|B|$, or their multiples $|iA|$ and $|iB|$, for $i$ sufficiently large, give rise to an effective difference.

To complete this discussion we wish further to consider the case in which the systems $|A|$ and $|B|$ have the same degree:
\[
A^2=B^2.
\]

In this case the difference $|A-B|$, or a multiple of it, cannot have effective existence if the number of intersections of the curves $A$ and $B$ is different from $A^2$; one thus has the necessary condition of existence
\[
AB=A^2=B^2,
\]
and the virtual degree of $|C|=|A-B|$ is
\[
n=(A-B)^2=0.
\]

Now the curve $C$ cannot have effective existence except as a curve of order zero, when $A\equiv B$. On the other hand the curve of $|iC|$ will have effective existence, again as a curve of order zero, when
\[
iA\equiv iB.
\]

Here there is occasion to observe that if, for no value of $i$, one has $iA\equiv iB$, one can construct on the surface (which we may suppose to have the curves $B$ as plane or hyperplane sections) an unlimited series of systems $|iA-(i-1)B|$, all formed of curves of the same order, and consequently contained in one or more continuous series of non-equivalent curves, which implies that the surface is irregular: %
\[
p_g>p_a.
\]

One may therefore assert that:

\textit{On a regular surface two linear systems of curves $|A|$ and $|B|$ of the same degree, which intersect in}
\[
AB=A^2=B^2
\]
\textit{points, are aliquot parts of one and the same multiple system:}
\[
|iA|=|iB|,
\]
\printersignature{\textsc{Enriques F.} - \textit{Superficie algebriche.}\quad 10}
% scan: s085L
% !PRINTED 146
\weberpage{146}
so that the $i$-fold multiple of their difference is the effective curve of order zero: $iA-iA$.

If one disregards the regularity of the surface, one can assert in general that: « if two curves $A$ and $B$ of the same degree $n$ meet in $n$ points, two multiples of them belong to one and the same continuous series of curves, or, as one says, are \textit{algebraically equivalent} »\origfootnote{1}{s085L-1}{Cf. F. \textsc{Severi}, Math. Annalen, 62 (1906), p. 194. A. \textsc{Todd}, Proc. of the Cambridge Phil. Soc. 23 March, 1939. For the relations with a theorem on the base due to W. V. D. \textsc{Hodge}, J. London Math. Soc. 12, 1937, cf. B. \textsc{Segre}, Ann. di Mat., 1937 and J. \textsc{Bronowski}, J. London Math. Soc., 1938.}.

The criterion of effective existence of $|C|=|A-B|$ which, as we said, was indicated by \textsc{Severi}\origfootnote{2}{s085L-2}{\textit{Sulle curve algebriche virtuali appartenenti ad una superficie algebrica.} Rendic. Ist. Lombardo II, 38 (1905).}, consists in the simple \textit{formal extension of the Riemann-Roch theorem to virtual curves}, such as one obtains by defining the \textit{index of speciality} $i$ of $|C|$ as « the number of linearly independent curves of the system residual to $|C|$ itself with respect to the (impure) canonical system $|K|$ », that is, as the dimension of the system
\[
|K-C|=|B+K-A|
\]
increased by one.

Denoting by $n$ and $\pi$ the virtual characters of $|C|$, the inequality
\[
r\geq p_a+n-\pi+1-i,
\]
affords a lower bound for the dimension $r$ of $|C|$, and hence assures the effective existence of the system when
\[
p_a+n-\pi+1-i\geq 0.
\]

To justify what has been said within our order of ideas (and limiting ourselves, for simplicity of discourse, to the hypothesis $i>0$) it suffices to bear in mind the lemma which led us to the Riemann-Roch theorem for special systems, observing that the system $|B+K-A|$ can be regarded as a small system with respect to the normally large system $|sB|$ for $s$ very high; therefore $|B+K-A|$ will cut out on the curve of the said system a complete series of dimension $i-1$, and hence the index of speciality of the series cut out on that same curve by $|A+(s-1)B|$ will be precisely the index of speciality $i$: from this one deduces the formula written above, which gives a lower bound for the dimension $r$ of %
\[
|A+(s-1)B-sB|=|A-B|.
\]

% scan: s085R
% !PRINTED 147
\weberpage{147}
Here it is appropriate to add the observation that the index of speciality $i$ of a virtual curve $C$ can turn out as large as one wishes, unlike the index of an effective curve, which is always
\[
i\leq p_g.
\]

Indeed the residual of $C$ with respect to $|K|$,
\[
|B+K-A|,
\]
need not a priori be contained in $|K|$: on the contrary, its being contained in it expresses, in another way, the necessary and sufficient condition for a special virtual curve to have effective existence.

\edsection{15. Regular and superabundant linear systems.}{15. Regular and superabundant linear systems.}{sec.IV.15}

On a surface of arithmetic genus $p_a$ let there be given a complete irreducible linear system of genus $\pi$, degree $n$ and index of speciality $i$: its dimension satisfies the inequality which constitutes the Riemann-Roch theorem
\[
r\geq p_a+n-\pi+1-i,
\]
or
\[
r=p_a+n-\pi+1-i+\omega
\]
with
\[
\omega\geq 0;
\]
the system has been called \textit{regular} if $\omega=0$, and \textit{superabundant} if $\omega>0$; in the latter case $\omega$ denotes its \textit{superabundance}.

There is occasion to inquire whether there exist criteria which allow one to evaluate the superabundance of a system $|C|$, or at least to decide whether or not $|C|$ is regular.

The most important result, in this order of ideas, concerns the regularity of the system adjoint to an irreducible linear system $|C|$, which we have proved for a sufficiently large multiple of $|C|$ ($\infty^3$ or $\infty^2$ at least) and which, as we said, was recognized by Picard (§ 13) for the system adjoint to $|C|$ itself.

Here we can establish very simply the \textit{theorem of Picard}, provided however that a superfluous restriction is added:

\textit{The system $|C'|$, adjoint to an irreducible linear system $|C|$ of dimension $r\geq p_g-p_a$, is certainly regular} (of dimension $p_a+\pi-1$).\origfootnote{1}{s085R-1}{Cf. \textsc{Enriques}, Bull. de Sc. Math., l. c. 1940.}

The proof rests on a

\textit{Lemma.} – Let $|C|$ and $|D|$ be two irreducible linear systems, say for simplicity of dimension greater than one. One can recog%
% scan: s086L
% !PRINTED 148
\unskip\weberpage{148}%
nize that « \(\lvert C\rvert\) cuts out a complete series on the curve of the sum system \(\lvert C+D\rvert\) if the dimension of \(\lvert D\rvert\) exceeds \(p_g-p_a\) ». Indeed, consider on the curve \(\bar{L}\) of \(\lvert L\rvert=\lvert C+D\rvert\) a group \(G\) of the complete series containing the intersection groups of the curves \(C\). Since the characteristic series of \(\lvert L\rvert\) has the deficiency \(\delta\leq p_g-p_a\), among the \(\infty^r\) groups \(LD\) there will be \(\infty^{r-(p_g-p_a)}\) which, associated with \(G\), constitute groups of the characteristic series of \(\lvert L\rvert\): let \(G'\) be one of these groups. Now through \(G+G'\) there passes a pencil of curves \(L\), and, with respect to this, the curve \(D\) through \(G'\) is fundamental: detaching it, there remains a curve of \(\lvert C\rvert\) which passes through \(G\). Thus, then, the complete series of the groups \(G\) is cut out on \(\bar{L}\) by the system \(\lvert C\rvert\).

This being established, take on a surface an irreducible linear system \(\lvert C\rvert\), of a certain genus \(\pi\) and of dimension greater than \(p_g-p_a\); we wish to prove that its adjoint \(\lvert C+K\rvert\) (\(K\) denoting the canonical system) is regular, of dimension \(p_a+\pi-1\).

For this purpose associate with \(\lvert C\rvert\) an irreducible linear system \(\lvert D\rvert\) of dimension greater than \(p_g-p_a\), whose adjoint is regular; and construct the system %
\[
\lvert L\rvert=\lvert C+D+K\rvert.
\]

On a curve \(L\), both the system \(\lvert C+K\rvert\) and the system \(\lvert D+K\rvert\) will cut out complete series, which must be residual to each other with respect to the canonical series. Since one knows the dimension of the series cut out by \(\lvert D+K\rvert\), which is equal to the dimension of the system \(\lvert D+K\rvert\), one will deduce the dimension of the series cut out by \(\lvert C+K\rvert\), which is that cut out by this same system.

Indeed, let \(\varrho\) and \(m\) be the genus and the degree of \(\lvert D+K\rvert\), and \(s\) the number of intersections of a \(C\) with a \(D+K\): the dimension of the series cut out by \(\lvert D+K\rvert\) on \(L\) is
\[
p_a+m-\varrho+1
\]
and its order is \(m+s\), while the genus of \(L\) is
\[
\bar{\pi}=\pi+\varrho+s-1.
\]

Hence the Riemann-Roch theorem tells us that the dimension of the series cut out on the same \(L\) by \(\lvert C+K\rvert\) is
\[
p_a+m-\varrho+1-(m+s-\bar{\pi})-1=p_a+\pi-1.
\]

Another criterion of regularity (which proves useful in questions of the postulation of a surface with respect to the varieties \(V_{n-1}\) of a space \(S_n\) of \(n>3\) dimensions) concerns the multiples of a linear system: \textit{The multiples of very high order of an irreducible system \(\lvert C\rvert\), at least \(\infty^2\), without fundamental curves, are regular.}

% scan: s086R
% !PRINTED 149
\weberpage{149}
Consider a certain value of \(h\) for which the system \(\lvert hC\rvert\) partially contains the canonical system \(\lvert K\rvert\), so that one has
\[
\lvert hC-K\rvert=\lvert D\rvert.
\]

We know that, for \(s\) sufficiently high, the system adjoint to \(\lvert sC\rvert\), that is, \(\lvert sC+K\rvert\), is regular. Now one will have
\[
\lvert(s+h)C\rvert=\lvert sC+K+D\rvert,
\]
where \(D\) will be a curve composed, in general, of several components, simple or multiple, of positive or negative degree. This being established, it is a matter of showing that the regular system \(\lvert sC+K\rvert\) is successively enlarged into regular systems (according to the criterion of § 8) when the irreducible components of \(D\) are added to it.

Indeed, denoting by \(L\) an irreducible component of \(D\), with characters \(\varrho\) and \(\nu\), where in any case \(\varrho\geq0\), one sees that the number of intersections of \(L\) with the curves of \(\lvert sC+K\rvert\) to which it is added (where \(s\) is as high as one wishes) is
\[
m=st+2\varrho-2-r>\varrho-\nu, %
\]
where \(t>0\) because \(L\) is not a fundamental curve for \(\lvert C\rvert\). The analogous inequality always holds for the number of intersections of the curves of \(\lvert sC+K+L\rvert\) with a component of \(D\), and so on successively in relation to the sum of \(\lvert sC+K\rvert\) with all the irreducible components of \(\lvert D\rvert\).

In the last analysis the system \(\lvert sC+K+D\rvert\), which is the multiple system \(\lvert(s+h)C\rvert\), for \(s+h\) very high, turns out regular, Q.E.D.

It will suffice to explain the matter by referring to the first significant case, in which it is a question of adding \(\lvert sC+K+L\rvert\) to \(L\) itself, which is supposed to enter as a double component in a curve \(D\) and to be of genus \(\varrho\geq0\) and of negative degree \(\nu=-\nu'\). We say then that, in this case, the number of intersections of \(L\) with the curves of \(\lvert sC+K+L\rvert\) is
\[
st+2\varrho-2+\nu'-\nu'=st+2\varrho-2>\varrho+\nu',
\]
for \(s\) sufficiently high in comparison with \(\nu'=-\nu\).

The theorem established above can be stated by saying that:

\textit{Every normally large linear system} (with respect to the canonical system) \textit{is regular.}

It is worth pointing out explicitly that the \textit{property does not hold} for large systems which are \textit{not} so \textit{normally}.

The simplest example is given by the linear systems which possess proper fundamental curves:

\textit{If the linear system \(\lvert C\rvert\) has a proper fundamental curve \(\theta\)}
% scan: s087L
% !PRINTED 150
\weberpage{150}
\textit{of genus \(\varrho>0\), the multiples of \(\lvert C\rvert\), however large, have the superabundance}
\[
\omega\geq\varrho\text{\origfootmark{1}{s087L-1}}.
\]\origfoottext{1}{s087L-1}{Cf. \textsc{Enriques}, \textit{Ricerche etc.}, 1893.}

Indeed the detaching of a \(\theta\), of degree \(-\nu\), from \(\lvert sC\rvert\) diminishes by 1 the effective dimension of the system, by \(\nu\) its degree and by \(\nu+\varrho-1\) its genus, and hence by \(\varrho-1\) its virtual dimension; this means that the superabundance \(\omega\) of \(\lvert sC\rvert\) exceeds by \(\varrho\) that, \(\omega'\), of \(\lvert sC-\theta\rvert\): %
\[
\omega=\omega'+\varrho.
\]

Note that, in general (for \(\varrho>1\)), also \(\omega'>0\), and therefore the genus \(\varrho\) of the fundamental curve is only a lower bound for the superabundance of \(\lvert sC\rvert\). Instructive here is the simple projective example of a surface \(F\) of order \(n\) (with plane sections \(C\)) endowed with an isolated multiple point of order \(i>3\), on which the surfaces \(\varphi_s\) (of sufficiently high order \(s\)) cut out the linear system \(\lvert sC\rvert\).

One sees at once not only that the general surfaces \(\varphi_s\) cut out on \(F\) a superabundant system (which has the fundamental curve \(\theta\) represented by the neighbourhood of \(O\)), but also that the systems cut out on \(F\) by the surfaces \(\varphi_s\) passing through \(O\) with the multiplicities \(1,\,2,\,\ldotp\ldotp\ldotp\ldotp\,i-3\) are superabundant, while that cut out by the adjoint surfaces \(\varphi_s\) which pass \(i-2\) times through \(O\) is regular. Whence one easily derives the calculation of the superabundance \(\omega\) of \(\lvert sC\rvert\). %

\edsection{16. Refinements concerning the regularity of the system adjoint to a curve\\ on a regular surface.}{16. Refinements concerning the regularity of the system adjoint to a curve on a regular surface.}{sec.IV.16}

We wish to examine more deeply the question of the regularity or superabundance of the linear systems larger than the canonical system, referring to the simple case of \textit{regular surfaces}, with \(p_g=p_a=p>0\).

If \(\lvert C\rvert\) is a linear system larger than the canonical system \(\lvert K\rvert\), one has
\[
\lvert C-K\rvert=\lvert L\rvert,
\]
where \(\lvert L\rvert\) denotes an effective system, and hence
\[
\lvert C\rvert=\lvert L+K\rvert=\lvert L'\rvert.
\]

Thus \(\lvert C\rvert\) can be considered as the system adjoint to an effective curve \(L\). Now \textit{if \(L\) is irreducible}, of a genus \(\pi>0\), \textit{its adjoint system is certainly regular}, of dimension \(p+\pi-1\), since it cuts out on \(L\) a (complete) series which cannot exceed the canonical series \(g^{\pi-1}_{2\pi-2}\).

% scan: s087R
% !PRINTED 151
\weberpage{151}
The same conclusion --- \textit{regularity of the system adjoint to \(L\)} --- extends to the case in which the curve \(L\) is a \textit{connected curve split in any manner into simple parts}.

\textit{Connected curve} formed of simple components means « a curve which, however divided into two parts, gives rise to curves having some (connecting) point in common », or also, topologically, a « curve or Riemann surface on which there can always be drawn a continuous path joining any two of its points ».

For the connected curves formed of simple components, thus defined, there holds the

\textit{Lemma of geometric composition of connected curves.} A connected curve \(L\) formed of \(s\) simple components can always be obtained by adding to a connected part formed of \(s-1\) components \(L_1+L_2+\ldotp\ldotp\ldotp\ldotp+L_{s-1}\) an \(s\)-th component \(L_s\) which has some (connecting) point in common with it.

This is proved by complete induction: supposing that the given \(L\) contains a connected part formed of \(h\) simple irreducible components, with \(h<s\), this part must be connected with the remaining one, and therefore with at least one of its components; by adding this component to it one obtains a connected part of \(L\), formed of \(h+1\) components.

As a \textit{Corollary}: given a connected curve \(L\) formed of \(s\) simple components, one can take these in an order \(L_1,\,L_2,\,\ldotp\ldotp\ldotp\ldotp,\,L_s\) such that in
\[
L=L_1+L_2+\ldotp\ldotp\ldotp\ldotp+L_s
\]
all the parts \(L_1+L_2+\ldotp\ldotp\ldotp\ldotp+L_h\) (\(h=1,\,2,\,\ldotp\ldotp\ldotp\ldotp,\,s\)) are connected.

By virtue of this property it is easy to recognize that « the adjoint system \(\lvert L'\rvert\) is regular », as was stated above. Indeed, denote by \(\pi_h\) the genus of \(L_1+\ldotp\ldotp\ldotp\ldotp+L_h\), by \(\bar{\pi}_h\) the genus of \(L_h\), and by \(i_h\) the number of intersections \((L_1+\ldotp\ldotp\ldotp\ldotp+L_{h-1})L_h\). We shall have that the genus of
\[
L=L_1+\ldotp\ldotp\ldotp\ldotp+L_s=(L_1+\ldotp\ldotp\ldotp\ldotp+L_{s-1})+L_s
\]
will be
\[
\pi_s=\bar{\pi}_s+\pi_{s-1}+i_s-1,
\]
and the series cut out by \(\lvert L'\rvert\) on \(L_s\) will be of order
\[
n_s=2\bar{\pi}_s-2+i_s
\]
and of dimension
\[
r_s\leq\bar{\pi}_s-2+i_s;
\]
hence the detaching of \(L_s\) from \(\lvert L'\rvert\) will involve
\[
r_s+1\leq\bar{\pi}_s+i_s-1
\]
conditions.

% scan: s088L
% !PRINTED 152
\weberpage{152}
It follows, then: if the system \(\lvert(L_1+\ldotp\ldotp\ldotp\ldotp+L_{s-1})'\rvert\) adjoint to \(\lvert L_1+\ldotp\ldotp\ldotp\ldotp+L_{s-1}\rvert\) is regular, one has precisely %
\[
r_s+1=\bar{\pi}_s+i_s-1
\]
and the system \(\lvert L'\rvert\) adjoint to \(\lvert L\rvert=\lvert L_1+\ldotp\ldotp\ldotp\ldotp+L_s\rvert\) is also regular. The theorem stated holds for an \(L\) formed of \(s\) components (simple and connected with one another) if it holds for a curve formed of \(s-1\) components. And thus it can be regarded as proved in general, since \(\lvert L'_1\rvert\) is regular.

Contrary to what happens for connected curves, the systems adjoint to disconnected split curves are not regular, but superabundant. More precisely:

\textit{The system \(\lvert L'\rvert\) adjoint to a curve \(L\) composed of \(s\) simple irreducible parts without common points has the superabundance \(s-1\).}

This is proved by the same inductive reasoning used for connected curves. Take, for example, a curve \(L\) (of genus \(\pi\)) composed of two parts \(L_1\) and \(L_2\), of genera \(\pi_1\) and \(\pi_2\), without common points (\(\pi=\pi_1+\pi_2-1\)). Then the adjoint system \(\lvert L'\rvert=\lvert(L_1+L_2)'\rvert\) cuts out on \(L_1\) the complete canonical series \(g^{\pi_1-1}_{2\pi_1-2}\), and the residual of \(L_1\) with respect to it --- which is \(\lvert L'_2\rvert\) --- cuts out on \(L_2\) the canonical series \(g^{\pi_2-1}_{2\pi_2-2}\); this latter system \(\lvert L'_2\rvert\) is regular, hence \(\lvert L'\rvert\) is superabundant, of dimension
\[
p+\pi=p+(\pi_1+\pi_2-1)-1+\omega
\]
with
\[
\omega=1.
\]

After this it is also easy to evaluate the superabundance of the system \(\lvert(L+\Sigma s_i\theta_i)'\rvert\) adjoint to a disconnected curve \(L+\Sigma s_i\theta_i\), formed of an irreducible curve \(L\) of degree \(>0\), say for simplicity of a curve varying in a linear system \(\infty^2\), and of several irreducible curves \(\theta_i\) of degree \(\nu_i<0\), fundamental for this system \(\lvert L\rvert\), each counted a certain number \(s_i\) of times.

\edsection{17. Continued: the system adjoint to a connected curve formed of\\ multiple components on a regular surface.}{17. Continued: the system adjoint to a connected curve formed of multiple components on a regular surface.}{sec.IV.17}

The results indicated do not exhaust the question of the regularity or superabundance of the system adjoint to a curve \(L\) on a regular surface, when \(L\) is composed in any manner of simple and multiple parts. The difficulty met with here is due, in fact, to the concept of the « connected curve », which must be extended precisely to the case of an arbitrary \(L\), composed as above. %

We designate as an \textit{(arithmetically) connected curve} a curve \(L\), formed in any manner of simple and multiple irreducible components,
% scan: s088R
% !PRINTED 153
\weberpage{153}
when any two parts \(L_1\) and \(L_2\) into which it is divided always have a (virtual) number of common points
\[
L_1L_2>0:
\]
it is not excluded that some part of \(L_1\) has a negative number \(-n\) of points in common with \(L_2\) (containing at least one multiple component of \(L_2\), of negative degree), provided that the remaining part then has more than \(n\) points in common with the same \(L_2\).

The consideration of connected curves, formed in any manner with multiple parts, gives rise to difficulties which have been set forth by A. \textsc{Franchetta}\origfootnote{1}{s088R-1}{\textit{Sulla caratterizzazione delle curve eccezionali riducibili di prima specie}, in Boll. dell’Unione Mat., Bologna, 1941.}, and which it is fitting to point out here.

Contrary to what intuition suggests in the case of simple parts (having with others a number \(\geq0\) of common points), it is \textit{not} possible to say in general that by adding to a connected curve \(L\) a component \(C\) which meets it in \(n>0\) points one always obtains a connected curve \(L+C\).

Indeed \textsc{Franchetta} adduces \textit{examples of disconnected curves formed of \(s\) components, decomposable into a connected curve formed of \(s-1\) components and a component which meets it in a positive number of points.}

In consequence of this, the proof of the lemma (\textit{of geometric composition}) which we established in the preceding section for connected curves formed of simple components also breaks down: one cannot say in general that an arithmetically connected curve formed of \(s\) components is always obtained by adding an irreducible component to a connected curve formed of \(s-1\) components.

Nevertheless there hold for arithmetically connected curves some properties, brought to light by \textsc{Franchetta}\origfootnote{2}{s088R-2}{l. c.}, which take the place of the preceding ones, as they are stated in the following theorems:

\textit{Theorem I.} -- A curve \(L\) (formed in any manner of simple or multiple irreducible components) is arithmetically connected if every connected part of it is connected with the remaining part (that is, has with it a virtual number of common points \(>0\)).

Indeed, if \(L=A+B\) is not connected, so that \(AB\leq0\), either \(A\) is connected, and one then has a connected part of \(L\) which intersects the residual part in a number \(\leq0\) of points, or else it is possible to decompose \(A\) into two parts \(A_1\) and \(A_2\) not connected with each other:
\[
A=A_1+A_2,\qquad A_1A_2\leq0;
\]

% scan: s089L
% !PRINTED 154
\weberpage{154}
while from \(AB\leq0\) it follows that at least one of the two curves, say \(A_1\), has a non-positive number of intersections with \(B\), but in this case we have %
\[
L=A_1+(A_2+B)
\]
with
\[
A_1(A_2+B)=A_1A_2+A_1B\leq0,
\]
that is, there is a part \(A_1\) of \(A\), within \(L\), disconnected from the residual part. Now either \(A_1\) will be connected, or it will contain a part \(A_{11}\) disconnected from the residual curve \((L-A_{11})\), and so on. But the sequence of successive parts \(A,\,A_1,\,A_{11},\,\ldotp\ldotp\ldotp\ldotp\), whose number of components keeps decreasing, will end with a connected curve, formed of one or more irreducible components: the last part thus obtained would be a curve connected in itself and disconnected from the residual one.

\textit{Theorem II.} -- From an (arithmetically) connected curve \(L\) one can remove a connected part \(B\) (formed of one or more components) in such a way that the residual part \(L-B\) is still connected.

Among all the possible connected parts of \(L\), let us choose one, \(B\), that has the least (positive) number of points in common with the residual part \(A=L-B\); we claim that \(A\) is also connected. If this were not so, we would have a division into disconnected parts
\[
\begin{aligned}
A&=A_1+A_2:\\
A_1A_2&\leq0,
\end{aligned}
\]
where one may suppose that \(A_1\) is, in itself, connected (theor. I). But, in that case, \(B\) must have a positive number of intersections with each of the parts \(A_1\) and \(A_2\) of \(A\), for otherwise we would have \((B+A_1)A_2\leq0\) or \((B+A_2)A_1\leq0\), and therefore \(L\) would be disconnected; we shall thus have
\[
BA_1>0\qquad\text{and}\qquad BA_2>0,
\]
hence
\[
BA_1<BA=B(A_1+A_2),
\]
and a fortiori
\[
(B+A_2)A_1<BA;
\]
thus \(A_1\), a connected part of \(L\), would have with the residual curve a smaller number of intersections than \(B\), contrary to our hypothesis.

The theorems of \textsc{Franchetta} take the place, up to a certain point, of the intuitive properties that characterize connected curves formed of simple components, as we shall see in what follows.

However, it does not seem easy to deduce from them, without exception, the regularity of the system \(\lvert L'\rvert\) adjoint to an arbitrary connected curve \(L\) (with multiple components of negative degree). One may attempt a somewhat indirect route, which is closely related to our developments on the Riemann-Roch theorem, but which nevertheless remains subordinate to a working hypothesis that will be brought to light below.

% scan: s089R
% !PRINTED 155
\weberpage{155}
Take on the surface \(F\) a regular linear system \(\lvert C\rvert\) which is a multiple of very high order of an irreducible system, at least \(\infty^2\), without fundamental curves, so that, in the first place, \(\lvert D\rvert=\lvert C+L\rvert\) turns out to be irreducible: one may assume that both \(\lvert L\rvert\) and its adjoint system \(\lvert L'\rvert=\lvert L+K\rvert\) cut complete series on a curve \(\bar{D}\) of this normally large system \(\lvert C+L\rvert\). Now one will recognize that \(\lvert C\rvert\) also cuts a complete series on such a curve \(\bar{D}\). Indeed, let \(G\) be a group of the complete series \(CD\) and \(\Gamma\) a group residual to \(G\) with respect to the characteristic series of \(\lvert D\rvert\), namely the section of \(D\) with \(L\); \(G+\Gamma\) constitutes the base group of a pencil of curves \(D\) which have no intersection with the components of \(L\) outside \(\bar{D}\).

Now if a \(D\) of this pencil is required to contain a point of \(L\) outside \(\bar{D}\), then, owing to the connectedness of \(L\), the whole connected curve will split off from \(D\), leaving a \(C\) through \(G\), Q.E.D.

The preceding proof rests on the assumption that, \(L\) being arithmetically connected, when a \(D\) of the aforesaid pencil is required to contain a point of \(L\) outside \(\Gamma\), \(L\) splits off completely, that is, each of its components splits off with the multiplicity with which it appears in \(L\).

This property appears to us to be true, but we do not possess an exhaustive proof of it. We shall introduce it as a working hypothesis, it being thus understood that the preceding result, and the consequences we shall draw from it, are subordinate to the validity of this hypothesis.

Now, having recognized that the regular system \(\lvert C\rvert\) cuts on the curve \(D\) of \(\lvert D\rvert=\lvert C+L\rvert\) a complete series \(g\), which is residual to the one cut by \(\lvert L'\rvert=\lvert L+K\rvert\), we can evaluate the dimension of the latter (and hence of \(\lvert L'\rvert\)), which equals the index of speciality of \(g\) diminished by one. Denoting by \(n\) the degree of \(\lvert C\rvert\), by \(\varrho\) its genus, and by \(m\) the number of intersections \(CL\), the order of the series \(g\) will be
\[
n+m
\]
and its dimension
\[
r=p+n-\varrho+1,
\]
while the genus of \(D\) is
\[
\pi+\varrho+m-1.
\]

Hence the index of speciality \(i\) of \(g\) is given by:
\[
n+m-(\pi+\varrho+m-1)+i=p+n-\varrho+1,
\]
and from this it follows that the dimension of the series \((CL')\) residual to \(CD\) is %
\[
i-1=p+\pi-1.
\]

Therefore \(\lvert L'\rvert\) has this same dimension and is regular.

% scan: s090L
% !PRINTED 156
\weberpage{156}
In conclusion: \textit{On a regular surface of genus \(p>0\), the system adjoint to an arithmetically connected curve is always regular.}

\edsection{18. On the reducibility of the canonical system.}{18. On the reducibility of the canonical system.}{sec.IV.18}

The theorem established in the preceding section leads to notable applications concerning the bicanonical system and the pluricanonical systems. To understand their meaning it is well to go back to the question of the reducibility of the canonical system.

Noether believed he could prove that the canonical system of a surface of genus \(p>0\) can be reducible only in two cases:

1) owing to the presence of exceptional curves (of genus \(\varrho=0\) and degree \(\nu=-1\)) which constitute fixed parts of the impure canonical system;

2) because the (absolute) linear genus is \(p^{(1)}=1\), in which case the pure canonical curves are in general composed of the elliptic curves of a pencil.

Thus for \(p^{(1)}>1\) the pure canonical system would always be irreducible. But Noether’s deduction is in error, and there exist actual \textit{examples of reducible} (pure) \textit{canonical curves} \(K\), belonging to surfaces of genus \(p>0\) and linear genus \(p^{(1)}>1\):

1) The first example is offered by the double plane with branch curve \(C_{2n}\) of order \(2n\), endowed with a \((2n-6)\)-fold point \(O\) (\(n>4\)) and with two points \([3,3]\), \(A_1\) and \(A_2\), collinear with \(O\) on a line \(a\)
\[
(p=2n-7,\qquad p^{(1)}=4n-15).
\]

The canonical curves \(K\) have as (double) images the curves \(C_{n-3}\) having \(O\) as an \((n-4)\)-fold point and passing simply through \(A_1\) and \(A_2\), and therefore split into \(C_{n-4}\) having \(O\) as an \((n-5)\)-fold point and the fixed part \(a\). Now this line represents an (irreducible) curve \(\theta\), a fixed component of the \(K\), which is not exceptional, having genus \(\varrho=0\) and degree \(\nu=-2\): as also appears from the fact that it has two intersections with the residual components of the \(K\), because the number of these intersections is:
\[
\begin{aligned}
&2\varrho-2-2\nu=2\\
&(\varrho=0,\qquad \nu=-2).
\end{aligned}
\]

2) A second example is offered by the double plane with branch curve \(C_{10}\), of order 10, endowed with 3 points \([3,3]\) on a line \(a\) which is not part of \(C_{10}\). A simple count of constants gives us a linear system of \(C_{10}\) with 3 assigned 3-fold points on \(a\) and as many 3-fold points infinitely near these in given directions, a system of
% scan: s090R
% !PRINTED 157
\weberpage{157}
dimension 29, which therefore does not reduce to the \(\infty^{26}\) one formed by \(a\) counted 2 times and by curves \(C_8\) with 3 tacnodes and given tacnodal tangents. Now, on the double plane we have constructed, the images of the canonical curves are the conics split into the fixed line \(a\) and a variable line. The corresponding fixed part \(\theta\) of the canonical system \(\lvert K\rvert\) is a curve of genus \(\varrho=1\) and degree \(\nu=-1\); the genera of the surface \(F\) are
\[
p=3,\qquad p^{(1)}=6.
\]

The degree \(\nu=-1\) of \(\theta\) will be computed by noting that \(\theta\) --- the fixed part of \(\lvert K\rvert\) --- has 2 intersections with the residual curves, so that
\[
2\varrho-2-2\nu=-2\nu=2,\qquad \nu=-1.
\]

On the other hand, the complete curve corresponding to \(a\) on \(F\) is a curve of genus 4 composed of 3 elliptic curves of degree \(-1\) corresponding to the points \([3,3]\) and of our \(\theta\), which has one point in common with each of them: the total degree of this complete curve is
\[
-3-1+2.3=2.
\]

3) A third example is offered by the double plane with branch curve \(C_{12}\), of order 12, with 4 points \([3,3]\) lying on a line \(a\) (which is not part of \(C_{12}\)); by a count of constants one proves the existence of irreducible \(C_{12}\) not containing \(a\). Now the canonical curves \(K\) are given on the double plane by cubics, split into the line \(a\) and a variable conic. To the line \(a\) there corresponds a \(\theta\), a fixed component of the \(K\), which has genus \(\varrho=1\) and degree \(\nu=-2\); indeed \(\theta\) has 4 intersections with the residual curves, so that %
\[
2\varrho-2-2\nu=-2\nu=4,\qquad \nu=-2\text{\origfootmark{1}{s090R-1}}. %
\]\origfoottext{1}{s090R-1}{We further cite the example of an irregular surface with \(p_a=2\), \(p_g=3\), \(p^{(1)}=11\), such as is represented on the double plane with branch curve \(C_{12}\) composed of 3 quartics \(C_4\) (of a pencil) which touch one another at 8 points of a conic \(C_2\); the images of the canonical curves split into the conic \(C_2\), to which there corresponds a fixed \(\theta\) of genus 3, and a variable line. The surface contains an irrational (elliptic) pencil of curves, represented by the pencil of the \(C_4\), in agreement with a general theorem of M. \textsc{De Franchis} on irregular double planes.}

The complete curve corresponding to \(a\) on \(F\) is a curve of genus 5 composed of 4 elliptic curves of degree \(-1\) corresponding to the points \([3,3]\) and of our \(\theta\), which has one point in common with each of them: the total degree of this complete curve is
\[
-4-2+2.4=2.
\]

% scan: s091L
% !PRINTED 158
\weberpage{158}
\textit{Note.} -- In the present state of our knowledge it is not known whether there exist surfaces with \(p^{(1)}>1\) whose canonical system is reducible in such a way that its variable parts are composed of the curves \(k\) of a pencil\origfootnote{1}{s091L-1}{Actual examples have recently been constructed by \textsc{G. Pompilj}: \textit{Alcuni esempi di superficie algebriche a sistema canonico puro degenere}, Rend. Acc. Lincei, s. 8, v. IV, 1948.}. Regarding this possibility we can only prove the following proposition:

\textit{If the pure canonical curves of a surface of genus \(p>1\)} (apart from fixed parts) \textit{are composed of the curves \(k\) of a pencil of genus \(\pi\) and degree \(n\), one has in every case}
\[
\pi\leq p^{(1)},\qquad n\leq p^{(2)}=p^{(1)}-1.
\]

Denote by \(\theta\) the possible fixed components of the \(K\):
\[
|K|=\Sigma\theta+|sk|
\]
and let \(\varrho\) and \(\nu\) be the characters of a \(\theta\). Let us express the degree of \(|K|\) as the sum of the intersection numbers of its components \(k\) (irreducible or not) and \(\theta\); we shall have
\[
s(2\pi-2-n)+\Sigma(2\varrho-2-\nu)=p^{(1)}-1.
\]

Now, referring to a transformed surface free of exceptional curves, we shall have for every \(\theta\)
\[
2\varrho-2-\nu\geq0,
\]
hence
\[
2\pi-2-n\leq p^{(1)}-1.
\]

On the other hand, a curve \(k\) has, with the residual curve \((s-1)k+\Sigma\theta\), the virtual number of intersections
\[
2\pi-2-2n\geq0.
\]

From this one deduces
\[
n\leq\pi-1,
\]
and
\[
\pi\leq p^{(1)},\qquad n\leq p^{(1)}-1,\qquad \text{Q.E.D.}
\]

\edsection{19. On the regularity of the bicanonical system}{19. On the regularity of the bicanonical system}{sec.IV.19}

The possibility that the pure canonical system \(|K|\) of a surface with \(p>0\) and \(p^{(1)}>1\) be reducible compels us to examine more deeply the question “whether the bicanonical system and the pluricanonical systems are regular”. For this it is necessary to prove that, however the \(K\) may be split into simple and multiple parts, they are always (for \(p^{(1)}>1\)) connected; hence the adjoint system is regular, %
\[
|2K|=|K'|,
\]
% scan: s091R
% !PRINTED 159
\weberpage{159}
and similarly, the curves of \(|2K|\) being also connected, the tricanonical system is regular %
\[
|3K|=|(2K)'|\quad\text{etc.}
\]

In words: \textit{On a regular surface of genus \(p>0\) and of linear genus \(p^{(1)}>1\), the bigenus is exactly}
\[
P=p+p^{(1)}
\]
\textit{and the \(i\)-genus}
\[
P_i=p+\frac{i(i-1)}{2}p^{(1)}.
\]

To prove this important theorem it is convenient to recognize successively that, for \(p>0\) and \(p^{(1)}>1\): %

1) the variable parts of the bicanonical system are irreducible;

2) the canonical curves \(K\), however reducible, are always connected, with \textit{order of connection} \(\geq2\) (that is, any part of \(K\) has at least 2 intersections with the residual part).

3) consequently the pure bicanonical system is regular, of dimension \(p+p^{(1)}-1\) and not greater.

These deductions extend a fortiori to the pluricanonical systems of order \(>2\).

1) The first proposition is easily proved. First of all, if \(p\geq2\) and \(p^{(1)}>1\), the variable parts of the canonical curves are irreducible or formed of the irreducible curves \(C\), of genus \(\pi>1\), belonging to a pencil; hence the curves \(C'\) adjoint to a \(C\) form a linear system of dimension \(p+\pi-1\) (at least), cutting on \(C\) the complete canonical series \(g_{2\pi-2}^{\pi-1}\); two \(C'\) have at least \(2\pi-2\) variable intersections, and therefore the \(C'\) are not composed of the curves of a pencil; a fortiori the bicanonical system \(|2K|\), which partially contains \(|C'|\), cannot have reducible variable parts composed of the curves of a pencil.

The theorem extends to the case \(p=1\). In this case, if the variable parts of the bicanonical curves are composed of the irreducible curves \(C\) of a pencil, first of all --- since \(p^{(1)}>1\) --- the genus of these must be \(\pi>1\)\origfootnote{1}{s091R-1}{Indeed we shall see later that if a surface possesses a pencil of elliptic curves and its geometric genus is positive, then its linear genus is necessarily \(p^{(1)}=1\).}. This being so, if the components \(C\) that enter into the curves of \(|2K|\) are \(s\) in number (\(P=s+1\) denoting the bigenus of the surface), then, since \(p=1\), there will be \(s\) curves \(C_{\frac{1}{2}}\) whose doubles belong to the pencil \(|C|\):
\[
K=\theta+sC_{\frac{1}{2}}.
\]

% scan: s092L
% !PRINTED 160
\weberpage{160}
Now the system \(|C'_{\frac{1}{2}}|\) adjoint to a curve \(C_{\frac{1}{2}}\) and contained in \(|2K|\) would have to have as variable parts the curves \(C\) of \(|C|\); but this is impossible, because its dimension is not one but (at least) \(p+\pi-1=\pi>1\)\origfootnote{1}{s092L-1}{Prop. 1) can also be proved in another way, by showing first of all that one must have \(\pi=2\) and hence necessarily \(p>1\): If the variable parts of the bicanonical curves \(2K\) are composed of \(s\) (\(>1\)) curves \(C\) of a linear, complete pencil, then, denoting by \(\theta\) the fixed parts of \(|2K|\), we shall write
\[
|2K|=\Sigma\theta+|sC|.
\]

Now let \(n\) and \(\pi\) denote the degree and the genus of \(|C|\), and \(\nu\) and \(\varrho\) those of a \(\theta\), and let us evaluate the additive character of the \(2K\) which expresses the number of their intersections with a \(K\); we shall have
\[
2K\cdot K=2p^{(1)}-2=\Sigma(2\varrho-2-\nu)+s(2\pi-2-n).
\]

But, referring, as is permissible, to a surface without exceptional curves:
\[
2\varrho-2-\nu\geq0,
\]
and therefore
\[
s(2\pi-2-n)\leq2p^{(1)}-2.
\]
Hence, in order that the pencil \(|C|\) be complete, its dimension must be \(r=1\); therefore, if \(|C|\) is not contained in the canonical system \(|K|\):
\[
p+n-\pi+1\leq1
\]
and
\[
n<\pi-1; %
\]
and this inequality holds a fortiori if \(|C|\) is contained in \(|K|\), because in that case \(|2C|\) is part of the adjoint system
\[
|C'|=|C+K|.
\]

This being so, if \(\pi>2\), we shall have
\[
\begin{aligned}
s(\pi-1)&\leq2p^{(1)}-2\\
\pi-1&\geq2\\
p^{(1)}-1&\geq s
\end{aligned}
\]
and consequently the dimension of \(|2K|\) is
\[
\geq p+p^{(1)}-1>s,
\]
in contradiction with the hypothesis that \(|2K|\) is formed by the groups of \(s\) curves of the pencil \(|C|\) and hence has dimension \(s\).

Having thus established our proposition in the case \(\pi>2\), it remains to discuss the hypothesis \(\pi=2\).

In this case the pencil of the curves \(C\) of genus \(\pi=2\) allows one to represent the surface on a double plane with branch curve \(f_{2m}\) of a certain even order \(2m\), endowed with a \((2m-6)\)-fold or \((2m-5)\)-fold point \(O\) and with other elementary singularities: 4-fold points and points \([3,3]\).

In this plane the bicanonical curves have as images curves \(\Phi_{2m-6}\) adjoint to \(f_{2m}\), which are composed of a fixed \(\Phi_{2m-6-s}\) and of \(s\) lines through \(O\). Since among the bicanonical curves there is at least one canonical curve counted twice, the \(\Phi_{2m-6-s}\) must be the double of a curve \(\Psi_{m-3-h}\) (\(s=2h\)) which, together with \(h\) lines through \(O\), constitutes the image of a canonical curve. Since \(h\geq1\), one deduces \(p>1\), a case already exhausted by the first considerations of the text!}.

% scan: s092R
% !PRINTED 161
\weberpage{161}
2) To prove the connectedness of the canonical curves \(K\), it will suffice to prove that the bicanonical curves of \(|2K|\) are connected, and for this it is a matter of reducing to absurdity the hypothesis that these are disconnected. Now the bicanonical system \(|2K|\) will consist of a fixed part \(L\) (however composed) and of irreducible variable parts \(C\):
\[
|2K|=L+|C|.
\]

If the curves of this system are disconnected, they can be divided into two parts having a number \(\leq0\) of common points: one of these parts will contain a \(C\) and a (fixed) part \(E\) of \(L\), while the other will be a fixed \(D\): %
\[
L=D+E,\qquad (C+E)D\leq0.
\]

Now one may assume that \(D\) is connected; if it were not, it could be decomposed into two parts \(D_1\) and \(D_2\) with a zero or negative number of common points, and one of these --- say, for example, \(D_1\) --- would have
\[
D_1(C+E+D_2)\leq0
\]
intersections with the residual curve; and, continuing in this way, one would in any case end up finding a connected part of \(L\), and hence of \(C+L\), having a number \(\leq0\) of intersections with the residual curve.

Assuming then that \(D\) is connected, let \(\varrho\) denote its genus and \(\nu\) its degree: since \(D\) is connected, we shall have
\[
\varrho\geq0.
\]\origfootnote{1}{s092R-1}{This proposition is proved, by complete induction, for the curves \(D\) composed of \(s\) irreducible components, assuming it proved for curves composed of fewer than \(s\) components. Indeed \(D\) can be decomposed into two parts \(D_1\) and \(D_2\), each connected in itself and connected with each other, with \(n>0\) common points; hence, denoting by \(\varrho_1\) and \(\varrho_2\) the genera of these curves, we shall have
\[
\begin{aligned}
\varrho&=\varrho_1+\varrho_2+n-1,\\
\varrho_1&\geq0,\quad \varrho_2\geq0,\quad \varrho\geq0;
\end{aligned}
\]
\hfill Q.E.D.}

Let us write the number of intersections of \(D\) with the residual curve with respect to the bicanonical system \(|2K|\):
\[
D(2K-D)=D(C+E)=4\varrho-4-3\nu\leq0.
\]

From this one deduces
\[
\nu\geq\frac{4}{3}(\varrho-1).
\]

Now, if \(\varrho=0\), since \(\nu\) is an integer, we shall have
\[
\nu\geq-1:
\]
\printersignature{\textsc{Enriques F.} - \textit{Superficie algebriche.}\quad 11}
% scan: s093L
% !PRINTED 162
\weberpage{162}
hence either \(\nu\geq0\), which entails \(p=0\), or \(D\) is an exceptional curve of the surface (\(\varrho=0,\ \nu=-1\)). One may suppose that the surface has been transformed so as not to contain such curves, and therefore one may discard the hypothesis \(\varrho=0\) and take
\[
\varrho\geq1.
\]

But if \(\varrho=1\), since we cannot have \(\nu>0\) (otherwise \(p=0\)), we shall have \(\nu=0\), and then one finds that the bicanonical and pluricanonical curves of \(|iK|\) will have zero intersections with \(D\), while \(D\) itself always has zero intersections with the residual curves of \(|iK-D|\), \(|iK-2D|\), etc.; from this it follows that the curves of \(|iK|\) are composed of the elliptic curves of a pencil and the linear genus is \(p^{(1)}=1\), contrary to our hypothesis (\(p^{(1)}>1\)).

In conclusion, if the bicanonical curves are disconnected, a part \(D\) of them, connected in itself and having a number \(\leq0\) of intersections with the residual curves, must be taken to be of genus \(\varrho>1\) and of degree
\[
\nu\geq\frac{4}{3}(\varrho-1)>\varrho-1.
\]

But this conclusion is absurd, as it conflicts with the hypothesis that \(D\) is part of the fixed part \(L\) of the curves of \(|2K|\), since the dimension of the complete system \(|D|\), of index of speciality \(p\), is, by the Riemann-Roch theorem,
\[
p+\nu-(\varrho-1)-p>0.
\]

Therefore the curves of the bicanonical system \(|2K|\), and hence also the curves \(K\) of the canonical system, are connected: any part of a \(K\), of genus \(\varrho\) and degree \(\nu\), will have an even number \(>0\) of points in common with the residual part:
\[
2\varrho-2-2\nu\geq2,
\]
and thus the \textit{order of connection of the \(K\)} is \(\geq2\).

3) From the fact that the pure canonical curves \(K\) are connected (for \(p>0\), \(p^{(1)}>1\)) one deduces that the bicanonical system \(|2K|=|K'|\) is regular. Similarly all the pluricanonical curves turn out to be connected, and hence the corresponding systems are again regular
\[
|K_i|=|\{(i-1)K\}'|.
\]

\edsection{20. Note on the irreducibility of the bicanonical system.}{20. Note on the irreducibility of the bicanonical system.}{sec.IV.20}

Having proved that on a regular surface with \(p>0\) and \(p^{(1)}>1\) \textit{the pure bicanonical system} is regular, one asks whether it is always irreducible, and --- since we have excluded the reducibility of its variable parts --- it is a matter of \textit{proving that it has no fixed parts} (apart from the exceptional curves).

% scan: s093R
% !PRINTED 163
\weberpage{163}
But the proof runs into difficulties of a delicate order when the canonical curves are supposed to be formed in an arbitrary way of simple and multiple components. We shall confine ourselves to proving the theorem in the case where \textit{the} (pure) \textit{canonical curves} \(K\) \textit{are formed of simple components.}

In this case we know that a \(K\) remains connected when one removes from it an (irreducible) component \(K_1\) having the least number \(n\) (\(\geq2\)) of points in common with the residual curve \(K_2\).

Now, denoting by \(\pi_1\) and \(\pi_2\) the genera of \(K_1\) and \(K_2\) respectively, we shall have for the genus of \(K\):
\[
\pi=\pi_1+\pi_2+n-1.
\]

Hence the splitting off of \(K_1\) from \(|K|\) diminishes by \(r_1+1=\pi_1+n-1>0\) the dimension of the regular system adjoint to \(|K|\), \(r_1\) denoting the dimension of the series cut by \(|K'|\) on \(K_1\). This means that \(K_1\) cannot be a fixed part of \(|K'|=|K_1+K'_2|\), since otherwise, by removing this part, one would obtain the system \(|K'_2|\) with the same dimension.

Now --- once the component \(K_1\) is removed --- the curves of \(|K_2|\) remain connected, but no longer (in general) with order of connection \(\geq2\): they may have \textit{handles}, that is, components having \textit{one} point in common with the residual curve, which --- when they are of genus 0 --- will constitute fixed parts of \(|K'_2|\). However, these handles are not, in any case, fixed parts of \(|K'|\); and this because otherwise the series \(g_{2\pi_1-2+n}\) cut by \(|K'|\) on \(K_1\) would have fixed points and would not be complete, whereas it must have dimension
\[
r_1=\pi_1-2+n.
\]

We have therefore recognized that, since the \(K\) are connected curves with order of connection \(\geq2\), the adjoint system \(|K'|\) cannot contain as fixed parts either \(K_1\) or the handles of the residual curve \(K_2\). Now, by stripping a \(K_2\) of these handles, one will have a connected curve \(\overline{K}_2\), with order of connection \(\geq2\), and one will prove similarly that neither a component removable without breaking the connectedness of the curve, nor the handles of the residual curve, are fixed parts of \(|\overline{K}'_2|\). Proceeding in this way, it follows that none of the components of \(K\) can be a fixed part of \(|2K|=|K'|\). \hfill Q.E.D.

\edsection{21. Genera of a reducible surface.}{21. Genera of a reducible surface.}{sec.IV.21}

To the developments of this chapter we wish to add some indications concerning reducible surfaces. If a surface \(F\) of a certain order \(m\), which we wish to suppose endowed with ordinary
% scan: s094L
% !PRINTED 164
\weberpage{164}
singularities, comes to split into two surfaces \(F_1\) and \(F_2\), of order \(m_1\) and \(m_2\) respectively, the question is to determine the numerical characters of \(F\) as functions of those of \(F_1\) and \(F_2\) and of the curve common to them.

It is well to make precise that the limit surface \(F_1+F_2\) will have a double curve, the limit of that of \(F\), which will be composed of three parts: a part of a certain order \(\delta_1\), double for \(F_1\) and not belonging to \(F_2\), a second part, whose order we shall denote by \(\delta_2\), double for \(F_2\) and not belonging to \(F_1\), and finally a third part --- say of order \(\delta\) --- belonging to \(F_1\) and \(F_2\). Apart from this there will be a curve \(C\) of a certain genus \(\pi\), the intersection of \(F_1\) and \(F_2\), which establishes the connection between the two surfaces \(F_1\) and \(F_2\), whose sum is thought of as the limit of \(F\).

Now the problem is to determine the numerical characters of the reducible surface \(F=F_1+F_2\) as functions of the analogous characters of the two component surfaces, and of those of the common curve \(C\) through which they are connected.\origfootnote{1}{s094L-1}{That there must exist a connecting curve. of order greater than zero follows from the principle of degeneration applied to the plane section curves of \(F\). Cf. \textsc{Enriques-Chisini}, Op. cit.}. %

Let us speak first of the arithmetic genus. We shall denote simply by \(p\) (equal to \(p_a\)) the arithmetic genus of \(F\); by \(p_1\) and \(p_2\) respectively the arithmetic genera of \(F_1\) and \(F_2\), while \(\pi\) has been called the genus of the curve \(C\).

It is clear that one can compute \(p\) by making use of the postulation formulae relative to the double curves of \(F_1\) and \(F_2\), or rather to the curves of order \(\delta_1\delta_2\) and \(\delta\) which constitute the limit of the double curve of \(F\) when the latter degenerates into \(F_1\) and \(F_2\). In this way one can arrive at the formula %
\[
p=p_1+p_2+\pi,
\]
which constitutes the extension to surfaces of Noether’s formula giving the genus of a reducible curve\origfootnote{2}{s094L-2}{In general, for varieties of any dimension: cf. \textsc{F. Severi}, \textit{Fondamenti per la geometria sulle varietà algebriche.} Rendic. Circolo Mat. di Palermo, 1909.}.

But we can justify the formula without carrying out the computation indicated above.

Indeed the surfaces \(\Phi_{m_1+m_2-4}\) of order \(m_1+m_2-4\) adjoint to \(F_1+F_2\) (for which \(C\) is not to be regarded as a double curve) will cut on \(F_1\) the system adjoint to \(|C|\), which has the virtual dimension:
\[
p_1+\pi-1,
\]
and for each of these curves there will in general be infinitely many surfaces \(\Phi_{m_1+m_2-4}\); the virtual dimension of this system of \(\Phi_{m_1+m_2-4}\)
% scan: s094R
% !PRINTED 165
\weberpage{165}
will be \(p_2\), there being \(\infty^{p_2-1}\) \(\Phi_{m_1+m_2-4}\) split into \(F_1\) and a surface \(\Phi_{m_2-4}\) adjoint to \(F_2\). In this way one finds the dimension of the system of the \(\Phi_{m_1+m_2-4}\) adjoint to \(F_1+F_2\):
\[
p-1=p_1+p_2+\pi-1.
\]

To this argument one might object that the surfaces \(\Phi\) adjoint to \(F_1+F_2\), by the fact of being obliged to pass through a curve of order \(\delta_2\) not belonging to \(F_1\), might cut on \(F_1\) a system of dimension lower than that of the system \(|C'|\) adjoint to \(|C|\). But, since we are dealing with \textit{virtual} dimensions, this circumstance has no effect on our computation, for which an analogous computation on adjoint surfaces of higher order could be substituted. %

We therefore conclude that the \textit{arithmetic genus of a reducible surface whose components of genera \(p_1\) and \(p_2\) are connected by means of a common curve of genus \(\pi\), is}
\[
p=p_1+p_2+\pi.
\]

Let us now set ourselves to evaluate the linear genus \(p^{(1)}\) of \(F_1+F_2\).

For this purpose it is necessary to know the linear genera \(p_1^{(1)}\) and \(p_2^{(1)}\) of \(F_1\) and \(F_2\), the genus of the connecting curve \(C\), which we have already denoted by \(\pi\), and the degrees \(n_1\) and \(n_2\) which this curve possesses on \(F_1\) and on \(F_2\) respectively.

\textit{The linear genus of the reducible surface is}
\[
p_1^{(1)}+p_2^{(1)}+8\pi-n_1-n_2-9.
\]

Indeed, as we have observed, the surfaces of order \(m_1+m_2+4\) adjoint to \(F_1+F_2\) cut on \(F_1\) curves \(C'\) adjoint to \(C\), which are of genus %
\[
p_1^{(1)}+3\pi-3-n_1;
\]
similarly the curves which the same \(\Phi\) cut on \(F_2\) are of genus
\[
p_2^{(1)}+3\pi-3-n_2.
\]

Now a curve of the first system on \(F_1\) and a curve of the second on \(F_2\) will together constitute the intersection of \(F_1+F_2\) with a \(\Phi_{m_1+m_2-4}\) when they have in common a canonical group of \(2\pi-2\) points constituting their intersections with \(C\), so that the genus of the curves composed of these two will be
\[
p^{(1)}=p_1^{(1)}+p_2^{(1)}+8\pi-n_1-n_2-9;
\]
\hfill Q.E.D.
% scan: s095L
% !PRINTED none
\weberblank{[s095L]}

\clearpage\thispagestyle{plain}
% scan: s095R
% !PRINTED none
\weberpage{[s095R]}
\edchapter{V}{NUMERICAL INVARIANTS AND MULTIPLE PLANES}{Chapter V. NUMERICAL INVARIANTS AND MULTIPLE PLANES}{NUMERICAL INVARIANTS AND MULTIPLE PLANES}

\edsection{1. The Zeuthen-Segre invariant.}{1. The Zeuthen-Segre invariant.}{sec.V.1}

In what precedes we have introduced the invariants of surfaces that are defined in relation to the canonical system, following a procedure which appears as a natural extension of the one that leads to the definition of the genus of a curve. But this latter procedure can be extended in different directions, and one thus succeeds in defining other numerical characters of surfaces which likewise have an invariant meaning and which can be used with advantage in the further developments of the theory, although they are not essentially new characters, since they can be expressed by means of the arithmetic genus \(p_a\) and the linear genus \(p^{(1)}\), as we shall have occasion to explain.

On the surface \(F\) consider a linear pencil of irreducible curves of genus \(\pi\geq0\), having a certain number \(n\) of effective base points (simple or multiple), and denote by \(\delta\) the number of curves of the pencil endowed with a double point, outside the base points.

\textit{The expression}
\[
I=\delta-n-4\pi,
\]
\textit{formed with the number of base points, the genus, and the number of curves endowed with a double point of a linear pencil, does not depend on the arbitrary choice of the pencil, and therefore constitutes a} (relative) \textit{invariant of the surface.}

The invariant meaning of this character was first recognized by \textsc{Zeuthen}\origfootnote{1}{s095R-1}{\textsc{H. G. Zeuthen}, \textit{Etudes géométriques etc.} Math. Ann., IV, 1871.} with reference to a pencil of plane sections of a surface. \textsc{C. Segre}\origfootnote{2}{s095R-2}{\textsc{C. Segre}, \textit{Intorno ad un carattere delle superficie etc.} Atti Acc. Scienze di Torino, 1896.} proved directly the equality of the analogous characters relative to any two pencils belonging to the surface. And therefore \(I\) is usually designated as the \textit{Zeuthen-Segre invariant.}

% scan: s096L
% !PRINTED 168
\weberpage{168}
Let us set out Segre’s proof, comparing two pencils in entirely generic relation to each other, \(|C_1|\) and \(|C_2|\), taken on the surface.

The first pencil will have a certain number \(n_1\) of base points, which we wish to suppose simple and distinct; the second will similarly have a certain number \(n_2\) of base points, different from the former and likewise simple and distinct. We shall denote by \(\pi_1\) and \(\pi_2\) respectively the genera of the generic curves of the two pencils, and by \(s\) the number of intersections \((C_1C_2)\). And we shall consider the curve \(L\), locus of the points of contact of the curves \(C_1\) and \(C_2\), that is, the Jacobian of the two pencils.

A curve \(C_1\) meets \(L\) at the points of the Jacobian group of the series \(g_s^1\) cut by \(|C_2|\) on \(C_1\), and therefore in \(2s+2\pi_1-2\) points, apart from the \(n_1\) base points, which also belong to \(L\). Similarly a \(C_2\) cuts \(L\) in \(2s+2\pi_2-2\) variable points.

Now on the curve \(L\), whose genus we shall denote by \(\pi\), we shall have two \(\infty^1\) series, of order \(2s+2\pi_1-2\) and \(2s+2\pi_2-2\) respectively, which we shall denote by \(g_1\) and \(g_2\). Let us consider the Jacobian group of \(g_1\), which is formed of
\[
4s+4(\pi_1-1)+2\pi-2
\]
points.

Observe that the double points of the series \(g_1\) on \(L\) are given by:

1) the \(\delta_1\) double points of the pencil \(|C_1|\); indeed one of these points is double for the group cut on \(L\) by the corresponding \(C_1\);

2) the \(\tau\) points where a \(C_1\) and a \(C_2\) have three-point contact, because there one has two infinitely near points of a \(C_1\) common to \(L\);

3) and finally the \(n_2\) base points of the pencil \(C_2\); indeed one of these points appears as a fixed point for the series cut by the \(C_2\) on the \(C_1\) passing through it, and hence counts twice in the Jacobian of this series, so that it constitutes a simple contact of \(C_1\) with \(L\).

This being so, we shall have: %
\[
4s+4(\pi_1-1)+2\pi-2=\delta_1+\tau+n_2. %
\]

Similarly, interchanging \(|C_1|\) with \(|C_2|\), one will obtain
\[
4s+4(\pi_2-1)+2\pi-2=\delta_2+\tau+n_1.
\]

Subtracting, one deduces
\[
\delta_1-n_1-4\pi_1=\delta_2-n_2-4\pi_2,
\]
\hfill Q.E.D.

\textit{Remark.} -- We have supposed that the base points of each of the two pencils \(|C_1|\) and \(|C_2|\) were simple and distinct, but we can free ourselves from these restrictions.

Suppose for example that a base point \(O\) of \(|C_2|\) is a multiple point of order \(r\) greater than one, with variable tangents. Then it is
% scan: s096R
% !PRINTED 169
\weberpage{169}
easy to show that this point will have multiplicity \(2r-1\) for the Jacobian \(L\): one verifies indeed that \(L\), completed by two curves \(\overline{C}_1\) and \(\overline{C}_2\) belonging respectively to the two pencils, constitutes the Jacobian curve of the net containing the two pencils
\[
\overline{C}_2+|C_1|\qquad \text{and}\qquad \overline{C}_1+|C_2|,
\]
and therefore it will have multiplicity \(3r-1\) at \(O\); removing \(C_2\), it remains that \(L\) has multiplicity \(2r-1\) at \(O\). Now when, in the preceding argument, one seeks the multiplicity with which \(O\) appears to belong to the group of points of section of \(L\) with the \(C_1\) passing through \(O\), one will find that this multiplicity is \(2r\), and hence that \(O\) belongs \(2r\) times to the Jacobian group of the \(g_s^1\) on \(L\); consequently the point \(O\) always comes to appear once in the count of the order of the Jacobian series of \(g_1\).

As regards the case in which one of our pencils possesses infinitely near base points, one can say that our formula always holds, in accordance with a requirement of continuity, provided one takes into account the curves of the pencil that come to acquire a double point (diminishing the genus \(\pi\)) at a base point of the pencil itself, or at a point infinitely near to it. It will suffice to refer to the case in which the pencil \(|C_1|\) has a simple base point \(O\) at which its curves touch one another: then the pencil contains a curve having a double point at \(O\) and a second curve (infinitely near to the preceding one) which has a double point infinitely near to \(O\): these two curves must appear in the count of \(\delta\), which will therefore be equal to the number of curves of the pencil endowed with a double point outside the base points, increased by 2.

After this we can evaluate the equivalence of \textit{a curve \(C_1\) endowed with an \(r\)-fold point}, outside the base points, by showing that it \textit{takes the place in general of \((r-1)^2\) curves endowed with a double point.} Compare the pencil \(|C_1|\), to which the curve endowed with the \(r\)-fold point belongs, with a second generic pencil \(|C_2|\). First of all, as has already been observed above, the \(r\)-fold point, \(P\), will be \((r-1)\)-fold for the Jacobian of the two pencils \(|C_1|\) and \(|C_2|\). And on each branch of this Jacobian the point \(P\) will appear as \((r-1)\)-fold for the series cut by the \(C_2\); so that, repeating the preceding argument, \(P\) will come to count \((r-1)^2\), times in the computation of the invariant \(I\). %

Instead of a curve endowed with a multiple point, there may be found in a pencil \(|C|\) a curve that contains a whole double or multiple component: for example a
\[
C^{(0)}=2\chi+\psi,
\]
where \(\chi\) has a certain genus \(\varrho\geq0\) and the residual curve \(\psi\) (which may also be absent) meets \(\chi\) in a certain number \(i\geq0\) of points.
% scan: s097L
% !PRINTED 170
\weberpage{170}
It is not difficult to evaluate the equivalence of this curve \(2\chi+\psi\), that is, the number of curves \(C\) endowed with a double point that are absorbed by it.

This computation is found worked out in a memoir by \textsc{Castelnuovo-Enriques}\origfootnote{1}{s097L-1}{\textsc{Castelnuovo-Enriques}, \textit{Sopra alcune questioni fondamentali etc.} Ann. di Mat., 1901.}. We shall set it out here in the simplest case, in which the complication of infinitely near base points does not intervene, referring therefore to the hypothesis that the pencil possesses no base points, which is the case in which it will suit us later to apply the formula in relation to irrational pencils.

Comparing the linear pencil \(|C|=|C_1|\) with another generic pencil \(|C_2|\), one will observe first of all that \(\chi\) splits off from the Jacobian \(L\) of the two pencils, and the curve \(L_1\), the residual part of the latter, will pass through the \(i\) points common to \(\chi\) and \(\psi\), which are triple for the curve \(C^{(0)}\). Following a procedure analogous to the one followed in the general case, one must find the order of the series cut on \(L_1\) by the curves of the pencils \(|C_1|\) and \(|C_2|\). %

In the first place, the order of the series \(g_{m_1}^1\) cut by the curves \(C_1\) on \(L_1\) is
\[
m_1=2s+2\pi_1-2\qquad \text{with}\qquad s=(C_1C_2),
\]
since that is how many \(C_2\) touch a generic \(C_1\).

To determine the order of the series \(g_{m_2}^1\) cut on \(L_1\) by the \(C_2\), one must bear in mind that the Jacobian group of the series \(g_s^1\) cut by the \(C_1\) on a generic \(C_2\) also includes the group of the points of intersection of that same \(C_2\) with \(\chi\). Denoting by \(\sigma\) the number of these points \((C_2\chi)\), one finds that the number of curves \(C_1\) properly touching a generic \(C_2\) is
\[
m_2=2s+2\pi_2-2-\sigma,
\]
and this is also the number of intersections \((C_2L)\). Now, denoting by \(P\) the genus of the Jacobian \(L_1\), the Jacobian group of the series \(g_{m_1}^1\) on \(L_1\) will be composed of \(2m_1+2P-2\) points, and this same group is formed by: %

1) the \(\bar{\delta}\) double points of the remaining nodal curves belonging to the pencil \(|C_1|\);

2) the \(\tau\) three-point contacts of a \(C_1\) with a \(C_2\);

3) the \(n_2\) base points of the pencil \(|C_2|\);

4) the \(i\) points common to \(\chi\) and \(\psi\), each of which, being triple for the group cut by \(C^{(0)}=2\chi+\psi\) on \(L_1\), must be counted twice;

% scan: s097R
% !PRINTED 171
\weberpage{171}
5) the \(2\sigma+2\varrho-2\) double points of the series cut by the pencil \(|C_2|\) on \(\chi\), since \(L_1\) passes through these points, which are double for the group cut by \(C^{(0)}\).

We shall therefore have
\[
2m_1+2P-2=\bar{\delta}+\tau+n_2+2i+2\sigma+2\varrho-2.
\]

On the other hand, for the series cut by the curves \(C_2\) on \(L_1\) one has simply
\[
2m_2+2P-2=\delta_2+\tau,
\]
the pencil \(|C_1|\) having no base points.

Eliminating \(2P-\tau-2\) from these two relations and substituting the values of \(m_1\) and \(m_2\), one obtains
\[
\bar{\delta}-4\pi_1+2\varrho-2+2i=\delta_2-n_2-4\pi_2.
\]

But if the pencil \(|C_1|\) contained only curves having isolated double points, and \(\delta_1\) denoted their number, one would have
\[
I=\delta_1-4\pi_1.
\]

From the comparison it follows that \textit{a curve \(C^{(0)}\) endowed with a double component of genus \(\varrho\) which intersects the residual curve in \(i\) points counts for}
\[
\delta_0=2\varrho-2+2i
\]
\textit{curves endowed with a double point}, in the computation of the Zeuthen-Segre invariant, at least under the hypothesis that the \textit{pencil has no base points.}

The preceding argument extends to the case in which the pencil \(|C|=|C_1|\) contains a curve
\[
C^{(0)}=h\chi+\psi,
\]
containing a component of multiplicity \(h\) (greater than 1), free of double points, which meets the residual curve \(\psi\) in \(i\) points.

In this case one will have
\[
\delta_0=(h-1)(2\varrho-2)+hi.
\]

More generally, suppose that the pencil \(|C|\) contains a reducible curve
\[
C^{(0)}=h_1\chi_1+h_2\chi_2+\ldotp\ldotp\ldotp\ldotp+h_t\chi_t
\]
consisting of \(t\) simple or multiple components \(\chi\) \((h_r\geq1)\), and that the curves \(\chi_r\) and \(\chi_s\) have \(i_{rs}\) points in common. The equivalence of this %
% scan: s098L
% !PRINTED 172
\weberpage{172}
curve for the purposes of computing the Zeuthen-Segre invariant is given by \textit{Godeaux’s formula}\origfootnote{1}{s098L-1}{L. \textsc{Godeaux}, \textit{Sur le calcul de l’invariant de Zeuthen-Segre.} Mémoires de la Société des Sciences du Hainaut, t. 66, 1920. Cf. Bull. de l’Académie Royale de Belgique, 13 Oct. 1934. See also L. \textsc{Campedelli}, \textit{Sul computo dell’invariante di Zeuthen-Segre per una superficie algebrica.} - \textit{Ancora sul computo etc.} Rend. Acc. Lincei, series 6, t. XIX, 1934.}. %
\[
\delta_0=\Sigma(h_r-1)(2\varrho_r-2)+\Sigma i_{rs}(h_r+h_s-1) %
\]

This formula is capable of covering all cases, when account is taken of two modifications:

1) If \textit{one of the curves, \(\chi_1\), possesses a double point} or a multiple one, \(\delta_0\) must be duly increased: if it is a double point and \(\varrho_1\) is taken as the \textit{effective genus}, \(2h_1-1\) must be added to \(\delta_0\); if instead \(\varrho_1\) represents \textit{the virtual genus} of \(\chi_1\), \(\delta_0\) is to be increased by one unit only.

2) In case the pencil \(|C|\) has infinitely near base points on a multiple component \(\chi_1\) of \(C^{(0)}\), these points can be eliminated by means of a transformation of the surface that changes them into exceptional curves: this amounts to adding to the \(\chi_r\) one or more neighbourhoods of base points. \textit{With this convention Godeaux’s formula also covers the case of infinitely near base points.}

\edsection{2. Irrational pencil.}{2. Irrational pencil.}{sec.V.2}

The Zeuthen-Segre invariant relative to a surface \(F\) can be computed not only in relation to the linear pencils traced on \(F\), but also in relation to irrational pencils belonging to \(F\).

Precisely, the \textit{invariant \(I\) is expressed, in relation to a pencil of genus \(\varrho>0\) consisting of curves of genus \(\pi\) and containing \(\Delta\) curves endowed with a double point, by the formula}
\[
I=\Delta+4(\varrho-1)(\pi-1)-4\ \text{\origfootmark{2}{s098L-2}}.
\]\origfoottext{2}{s098L-2}{Cf. \textsc{Castelnuovo-Enriques}, l. c.}

This formula can be proved directly, by taking up again the comparison of two pencils \(|C_1|\) and \(|C_2|\) and replacing one of the two by the irrational pencil \(\{C\}\) under discussion here. But one can also reduce this proof to the case of two linear pencils, by replacing \(\{C\}\) by a linear pencil consisting of the groups of \(n\) curves \(C\) of a \(g_n^1\), where \(n\) is a sufficiently large integer, for example equal to \(\varrho+1\). Indeed the circumstance that one of the two pencils is composed of reducible curves does not invalidate the arguments developed in the preceding section.

% scan: s098R
% !PRINTED 173
\weberpage{173}
Thus the groups of \(n\) curves \(C\) of the aforesaid \(g_n^1\) will form a linear pencil of curves, of genus
\[
\Pi=n\pi-n+1
\]
without base points; and in this pencil there will be:

1) \(\Delta\) curves containing a \(C\) with a double point, and

2) \(2n+2\varrho-2\) curves containing a \(C\) of genus \(\pi\) counted twice, each of which is equivalent to \(2\pi-2\) curves endowed with a double point;

hence one will have
\[
I=\Delta+(\pi-2)(2n+2\varrho-2)-4\Pi=\Delta+4(\varrho-1)(\pi-1)-4, %
\]
\hfill Q.E.D.

\edsection{3. Expression of the Zeuthen-Segre invariant by means of the genera.}{3. Expression of the Zeuthen-Segre invariant by means of the genera.}{sec.V.3}

It has been said that the \textit{Zeuthen-Segre invariant \(I\) is a relative invariant of the surface \(F\)}; indeed it \textit{increases by one when a point \(O\) of \(F\) is changed into an exceptional curve.} The assertion is verified in two ways: if one takes a point \(O\) of \(F\) which is not a base point and is simple for a curve of the pencil \(|C|\), then, on changing \(O\) into an exceptional curve, the number \(\delta\) of curves endowed with a double point will increase by one; if instead \(O\) falls at a base point of the pencil \(|C|\), the transformation diminishes the number of base points by one.

Let us now observe that a transformation of the surface \(F\) substituting an exceptional curve for a point makes the invariant \(I\) and the linear genus \(p^{(1)}\) of the surface vary in opposite senses, since we know that the latter decreases by one unit. Therefore the sum
\[
I+p^{(1)}
\]
will constitute an \textit{absolute invariant} of the surface. But it is not a new invariant, since \textsc{Noether} proved that it is expressed in terms of the arithmetic genus by the formula
\[
I+p^{(1)}=12p_a+9\ \text{\origfootmark{1}{s098R-1}}.
\]\origfoottext{1}{s098R-1}{\textsc{Noether}, op. c., Math. Ann., VIII.}

This relation is connected with a group of formulae which bring to light various invariant numerical characters of surfaces. We shall prove it here following the route indicated by T. \textsc{Bonnesen}\origfootnote{2}{s098R-2}{P. \textsc{Bonnesen}, \textit{Sur les séries linéaires triplement infinies de courbes algébriques sur une surface algébrique.} Bull. de l’Académie Royal des Sciences et des Lettres de Danemarque, 1906, no. 4.}. %

% scan: s099L
% !PRINTED 174
\weberpage{174}
Let us refer to an irreducible, simple, \(\infty^3\) linear system of curves \(|C|\), which we may regard as the plane sections of a surface \(F\) endowed with a double curve \(D\); we shall suppose for simplicity that \(D\) is irreducible and endowed only with \(t\) triple points which are also triple for \(F\). The invariant \(I\) of the surface \(F\) can be computed in relation to a pencil of plane sections \(C\), and is given by
\weq{(1)}{I=\delta-n-4\pi,}
where \(n\) is the \textit{order} of \(F\), \(\pi\) the genus of the \(C\), and \(\delta\) denotes the \textit{class} of the surface \(F\). The class of \(F\) is also given by the effective degree of the system \(|C_j|\) cut on \(F\) by the polars \(\Phi_{n-1}\); this system possesses as simple base points the \(\tau\) cuspidal points of \(F\) on \(D\), so that \(\delta+\tau\) equals the degree of
\[
|C_j|=|3C+K|
\]
computed without regard to the aforesaid \(\tau\) base points. Introducing the linear genus \(p^{(1)}\) (genus of the impure canonical curves \(K\)) we shall therefore have
\weq{(2)}{\delta+\tau=3n+12(\pi-1)+p^{(1)}-1.}

Now certain characters (projective characters of the surface \(F\)) are connected with the \(\infty^3\) system \(|C|\), and in particular the following:

\(\alpha\)) the order \(d\) of the double curve \(D\), which depends on the degree \(n\) and on the genus \(\pi\) of \(|C|\):
\weq{(3)}{d=\frac{(n-1)(n-2)}{2}-\pi;}

\(\beta\)) the genus \(\varrho\) of the double curve \(D\);

\(\gamma\)) the number \(t\) of triple points of \(D\), which are also triple for the surface \(F\); moreover, calling \(\Delta\) the curve which corresponds to the double curve \(D\) on a model \(F'\), free of singularities, of the surface \(F\):

\(\delta\)) the (effective) genus \(\Pi\) of the curve of neutral pairs of \(|C|\), a curve \(\Delta\) which is represented on the double \(D\) with \(\tau\) branch points at the cuspidal points. It should be noted that the curve \(\Delta\) of neutral pairs appears to possess \(t\) triads of double points corresponding to the \(t\) triple points of \(D\), and therefore the \textit{virtual genus} of \(\Delta\) is
\[
\Pi+3t.
\]

Among the aforesaid characters, besides (1) and (2), there hold various relations, from which it appears that, besides the two characters \(n\) and \(\pi\) of the system \(|C|\), only two characters of the surface are capable of assuming independent and, up to a certain point, arbitrary values.

First of all we observe that the curve \(\Delta\), a covariant of the \(\infty^3\) system \(|C|\), turns out to be, on the surface \(F'\) (the transform of \(F\) in which
% scan: s099R
% !PRINTED 175
\weberpage{175}
the double curve \(D\) is resolved), a partial curve of the system \(|(n-4)C|\); more precisely, adding a canonical curve \(K\) to \(\Delta\), one obtains
\[
|\Delta+K|=|(n-4)C|.
\]

Therefore the intersections of \(\Delta\) with the curves of the adjoint system \(|\Delta+K|\) will form a canonical group on the curve \(\Delta\) of virtual genus \(\Pi+3t\), and one will therefore have
\weq{(4)}{2(\Pi+3t)-2=2(n-4)d,}
where \(d\) is expressed by formula (2). %

On the other hand, since the double curve \(D\) is in \((1,2)\) correspondence with \(\Delta\) and there are \(\tau\) branch points on it, we shall have, by a well-known formula of Zeuthen,
\weq{(5)}{2(2\varrho-2)+\tau=2\Pi-2.}

Finally we can obtain a sixth relation by evaluating the number of intersections of a \(C_j\), the Jacobian of a net of \(|C|\), with \(\Delta\), or, what is the same, with \(D\). For this purpose we observe that
\[
|C_j|=|3C+K|\qquad \text{and}\qquad |\Delta+K|=|(n-4)C|;
\]
consequently the number of intersections \((C_j\Delta)=(C_jD)\) is given by
\[
(3C+K)[(n-4)C-K]=2n^2-7n+15+2\pi(n-7). %
\]

But this same number \((C_jD)\) can be computed directly on the surface \(F\), by seeking the number \(\alpha\) of intersections of \(D\) with the second polar \(\Phi_{n-2}\) of a generic point \(O\); indeed the \(C_j\), intersection of the polar of \(O\), meets the double curve \(D\): at the \(\tau\) cuspidal points and at the \(\alpha\) points common to \(D\) and to \(\Phi_{n-2}\) (outside the triple points of \(D\)), at each of which the tangent plane to a sheet of \(F\) passes through \(O\). So we shall have %
\[
(C_jD)=\alpha+\tau,
\]
where
\[
\alpha=(n-2)d-3t.
\]

Consequently we can write
\weq{(6)}{2n^2+n+14+2\pi(n-7)-p^{(1)}=(n-2)d-3t.} %

To the six relations that we have written it is convenient to add a seventh, which introduces the genus \(p_a\): one has
\[
\begin{aligned}
p_a&=\frac{(n-1)(n-2)(n-3)}{6}-(n-4)\left[\frac{(n-1)(n-2)}{2}-\pi\right]+2t+\varrho-1=\\
&=\frac{1}{6}(n-1)(n-2)(9-2n)+\pi(n-4)+2t+\varrho-1.
\end{aligned}
\]

% scan: s100L
% !PRINTED 176
\weberpage{176}
We have found seven equations linking the eleven characters: \(n,\pi,\delta,I,d,\varrho,\Pi,\tau,t,p_a,p^{(1)}\).

Among these, \(n\) and \(\pi\) can be taken, up to a certain point, arbitrarily, since they depend on the choice of a pencil \(|C|\) on the surface \(F\); hence it is clear a priori that the seven characters
\[
\delta,\quad I,\quad \varrho,\quad \Pi,\quad \tau,\quad t %
\]
can be expressed as functions of the said \(n\) and \(\pi\) and of the invariants \(p_a\) and \(p^{(1)}\) of the surface. In particular \(I\) must be expressed as a function of \(p_a\) and \(p^{(1)}\): a priori one might suppose that this function also depends on \(n\) and \(\pi\), but in fact it must turn out to be independent of them, since it is an invariant character of \(F\).

It only remains to carry out the computation, and since (1) and (2) already contain an explicit expression of \(I\) and of \(d\), it is a matter of solving five linear equations, in which \(n\), \(\pi\), \(p_a\) and \(p^{(1)}\) appear as parameters. One thus obtains: %
\[
\begin{aligned}
\delta&=n+4\pi-p^{(1)}+12p_a+9\\
\varrho&=\frac{1}{2}(n^2-7n+48)+(n-12)\pi+9p_a-2p^{(1)}\\
\Pi&=n^2-6n+2(n-10)\pi+12p_a-3p^{(1)}+36\\
t&=\frac{1}{6}(n^3-9n^2+26n-78)-(n-8)\pi-4p_a+p^{(1)}\\ %
\tau&=2n+8\pi-12p_a+2p^{(1)}-22.
\end{aligned}
\]

From the first of these formulae one derives Noether’s expression for the Zeuthen-Segre invariant:
\[
I=12p_a-p^{(1)}+9.
\]

\textit{Note.} -- The four formulae giving \(\varrho\), \(\Pi\), \(t\) and \(\tau\) bring to light four relative invariants constructed with the characters of an \(\infty^3\) linear system \(|C|\) belonging to the surface. Particularly noteworthy is the last formula, which contains \(n\) and \(\pi\) linearly; it defines the invariant
\[
2n+8\pi-\tau=12p_a-2p^{(1)}+22.
\]

Combining this formula with the first, one will obtain another very simple relative invariant \origfootnote{1}{s100L-1}{Cf. \textsc{Bonnesen}, l. c.}, namely
\[
2\delta-\tau=36p_a-4p^{(1)}+40.\ \text{\origfootmark{2}{s100L-2}}
\]\origfoottext{2}{s100L-2}{Cf. \textsc{Bonnesen}, l. c.}

% scan: s100R
% !PRINTED 177
\weberpage{177}
This invariant is constructed from two characters inherent in an \(\infty^3\) linear system \(|C|\) traced on the surface: \(\delta\) denotes, in general, the number of curves \(C\) of a pencil endowed with a double point (outside the base points), and \(\tau\) indicates the number of pencils contained in the \(\infty^3\) system \(|C|\) which possess a double base point (outside the base points of \(|C|\)).

\textit{Remark.} -- The expression of the Zeuthen-Segre invariant by means of the arithmetic genus and the linear genus leads us to recover the \textit{linear genus of the plane and of ruled surfaces}, which we computed in Chap. II, \S\ 6.

First of all, one can compute the invariant \(I\) of the plane by referring to a pencil of conics, for which one has:

the genus \(\pi=0\),

the number of base points \(n=4\),

and the number of curves endowed with a double point \(\delta=3\).

We shall have
\[
I=\delta-n-4\pi=-1;
\]
hence from the formula
\[
I=12p_a-p^{(1)}+9
\]
since \(p_a=0\), one derives
\[
p^{(1)}=10.
\]

Let us now pass to the irrational ruled surfaces of genus \(p_a=-p\) \((p>0)\). Here one can compute the value of \(I\) in relation to the pencil of generating lines, which is of genus \(p\) and for which
\[
\pi=0,\quad \Delta=0,\quad n=0.
\]

We shall have
\[
I=\Delta+4(p-1)(\pi-1)-4=-4p.
\]

Now from Noether’s formula one deduces
\[
12p_a-p^{(1)}+9=-4p
\]
and, since \(p_a=-p\),
\[
p^{(1)}=-8(p-1)+1.
\]

The same result is reached by referring, instead of to the irrational pencil of generators, to a linear pencil of plane sections.

\edsection{4. Cuspidal curves of a net.}{4. Cuspidal curves of a net.}{sec.V.4}

Let us take up again the linear system \(\infty^3|C|\) and its characters considered in the preceding section. The formulae established per%

\printersignature{\textsc{Enriques F.} - \textit{Superficie algebriche.}\quad 12}

% scan: s101L
% !PRINTED 178
\unskip\weberpage{178}%
mit us to determine the characters belonging to a net of curves $|C|$, and in particular the number $\chi$ of the curves of a net that are endowed with a cusp.

For this purpose let us refer to the surface $F$ which has the curves $C$ as plane sections. The locus of the points of $F$ for which the tangent plane section possesses a cusp is the \textit{parabolic curve} $C_h$ of $F$, the section of $F$ itself by the Hessian surface $H$. And hence the $\chi$ plane sections endowed with a cusp that pass through a point $O$ of space will correspond to the intersections of $C_h$ with the curve $C_j$, the Jacobian of the net cut out by the planes through $O$, apart from the cuspidal points of $F$, which belong simply to $C_j$ and, as we shall see, are double for the parabolic curve $C_h$.

To evaluate $\chi$ it is therefore necessary:

1) to determine the behaviour of the Hessian along the nodal double curve of the surface $F$;

2) further, to find the multiplicities which the curve $C_h$ possesses at the $\tau$ cuspidal points of $F$.

The first computation is carried out simply by replacing the given surface by a surface, which for the moment we again denote by $F$, having at the origin a biplanar double point, with the tangent planes $x$ and $y$; we write, in homogeneous coordinates $x$, $y$, $z$, $u$:
\[
F=u^{n-2}xy+u^{n-3}\varphi_3(x\,y\,z)+\ldotp\ldotp\ldotp\ldotp+\varphi_n(x\,y\,z).
\]

The Hessian of $F$ is
\[
H=
\begin{vmatrix}
\dfrac{\partial^2 F}{\partial x^2} &
\dfrac{\partial^2 F}{\partial x\partial y} &
\dfrac{\partial^2 F}{\partial x\partial z} &
\dfrac{\partial^2 F}{\partial x\partial u} \\[6pt]
\dfrac{\partial^2 F}{\partial y\partial x} &
\dfrac{\partial^2 F}{\partial y^2} &
\dfrac{\partial^2 F}{\partial y\partial z} &
\dfrac{\partial^2 F}{\partial y\partial u} \\[6pt] %
\dfrac{\partial^2 F}{\partial z\partial x} &
\dfrac{\partial^2 F}{\partial z\partial y} &
\dfrac{\partial^2 F}{\partial z^2} &
\dfrac{\partial^2 F}{\partial z\partial u} \\[6pt]
\dfrac{\partial^2 F}{\partial u\partial x} &
\dfrac{\partial^2 F}{\partial u\partial y} &
\dfrac{\partial^2 F}{\partial u\partial z} &
\dfrac{\partial^2 F}{\partial u^2}
\end{vmatrix}
\]

Carrying out the computation one finds that the behaviour of the Hessian at the origin of coordinates $(x=y=z=0,\ u=1)$ is determined by the Hessian of the approximating surface of the third order, and thus $H$ possesses a quadruple point at the biplanar point of $F$. Returning now to our surface $F$ endowed with a nodal curve, one deduces that its Hessian $H$, of order $4n-8$, passes four times through the said nodal curve. Consequently the parabolic curve $C_h$ will belong to the linear system
\[
|4K+8C|,
\]
denoting as usual by $K$ the canonical curves of the surface.

% scan: s101R
% !PRINTED 179
\weberpage{179}
Now one sees that the $\tau$ cuspidal points of $F$ (simple for the $C_j$) are double for the parabolic curve $C_h$. Indeed, referring to a transformed surface on which the nodal curve of $F$ is resolved, a cuspidal point of $F$ will give on that surface a double base point for a pencil of $|C|$, in which there will therefore be \textit{two} curves $C$ endowed with a cusp.

This granted, the number of cuspidal curves of a net within $|C|$ will be given by
\[
\chi=(C_j\cdot C_h)-2\tau.
\]

But, since
\[
|C_j|=|3C+K|,\qquad |C_h|=|8C+4K|,
\]
we shall have
\[
(C_j\cdot C_h)=24n+20(2\pi-2-n)+4p^{(1)}-4,
\]
and, recalling the expression for $\tau$ from the preceding section,
\[
\tau=2n+8\pi-12p_a+2p^{(1)}-22
\]
it follows that
\[
\chi=24(p_a+\pi).
\]

One thus arrives at the relation discovered by \textsc{Zeuthen}: \textit{the number of cuspidal curves of a net of genus $\pi$ is in general $24(p_a+\pi)$.}

To tell the truth, we have proved this formula by referring to a net of curves $C$ that is totally contained in an $\infty^3$ system. But it is easy to extend the reasoning by referring to an $\infty^3$ linear system $|L+C|$ that contains $|C|$ partially. And thus the formula will be justified in every case\origfootnote{1}{s101R-1}{The proof of the formula given in the text can also be presented in a slightly different form, in which no further use is made of the properties of the Hessian. For this purpose it is convenient to associate with the surface $F$ its dual $F'$, which is the transform of $F$ obtained by making the Jacobians $C_j$ of the nets of $|C|$ correspond to the plane sections. If one now constructs the dual surface of $F'$, one must fall back upon $F$. This means that the Jacobian curves of the nets of $C_j$ on $F$, let us say the $C_{jj}$, must be identified with the plane sections $C$, up to a fixed component, which is recognized to be the parabolic curve $C_h$.

Indeed, consider the polars with respect to $F$ of the points of a plane $\alpha$, which form a net. It is clear first of all that the plane section $C$ by $\alpha$ belongs to the Jacobian $C_{jj}$ of this net, because the polar of a point $P$ of $C$ cuts $F$ in a curve which passes doubly through $P$. Moreover, if one takes a point $A$ on the parabolic curve $C_h$, the tangent plane at $A$ meets $\alpha$ in a line $a$, and the polars of the points of $a$ cut $F$ in the curves of a pencil which has at $A$ a base point with a fixed tangent, among which there is one endowed with a double point at $A$; therefore the curve $C_h$ forms part of the second Jacobian $C_{jj}$.

From the relation
\[
|C_{jj}|=|C_h+C|
\]
it therefore follows that
\[
|C_h|=|8C+4K|.
\]
}.

% scan: s102L
% !PRINTED 180
\weberpage{180}
\edsection{5. Historical note.}{5. Historical note.}{sec.V.5}

The relation between the number of cuspidal curves of a net and the arithmetic genus was discovered by \textsc{H. G. Zeuthen}\origfootnote{1}{s102L-1}{\textsc{H. G. Zeuthen}, \textit{Études géométriques etc.} Math. Ann., Bd. IV, 1871. Cf. the memoir of \textsc{Noether}, also cited, in the Math. Ann., VIII.}. Assuming the invariance of the arithmetic genus to be known (which \textsc{Zeuthen} and \textsc{Noether} recognized, albeit with some restriction relating to the singularities of the surface), \textsc{Zeuthen} thus came to prove that the number $\chi-24\pi$, computed starting from a net (of plane sections), turns out to be independent of it, and therefore constitutes an intrinsic character of the surface. In 1902 \textsc{F. Severi} gave a direct proof of this theorem, by an entirely elementary procedure\origfootnote{2}{s102L-2}{\textsc{F. Severi}, \textit{Il genere aritmetico ed il genere lineare in relazione alle reti di curve tracciate sopra una superficie algebrica.} Atti Acc. di Torino, vol. XXXVII, 1902.}.

The proof developed in the text is found in \textsc{Bonnesen}, l. c.\origfootnote{3}{s102L-3}{Cf. also \textsc{M. Pannelli}, \textit{Sui sistemi lineari triplamente infiniti di curve tracciati sopra una superficie algebrica.} Rend. Circ. Mat. Palermo, t. XX, 1905. The consideration relating to the Jacobian of a net of Jacobians, which we indicated in a preceding note, is also met with in \textsc{E. Caporali}.}.

We cannot close this historical sketch, which completes the information given in the preceding sections, without recalling an order of ideas in which one seeks the functional geometric interpretation of the enumerative formulae discussed above. It is well known how this point of view was introduced into the theory of curves, with the theorem of \textsc{Painlevè-Castelnuovo} (interpretation of Zeuthen's formula relating to $[1,n]$ correspondences) and with the consideration of the Jacobian series of a linear series, in which \textsc{Enriques} recognized the simplest way of defining the canonical series and proving its invariance. %

To find an analogous interpretation of enumerative geometry on surfaces, \textsc{Severi}, in successive memoirs beginning in 1932, has sought to construct a theory of series of groups of points, which he calls series \textit{of equivalence}, and which should constitute a natural extension of the theory of linear series on a curve.

For an extensive bibliography on the subject see the book of \textsc{F. Severi}: ``Serie, sistemi di equivalenza e corrispondenze algebriche sulle varietà algebriche'', Rome, ed. Cremonese, 1942.

\edsection{6. Characters of a net and multiple planes.}{6. Characters of a net and multiple planes.}{sec.V.6}

The formulae established in the preceding section also allow us to express the characters of a net of curves $|C|$ by means of two
% scan: s102R
% !PRINTED 181
\weberpage{181}
of them, e.g. the genus $\pi$ and the degree $n$, and of the invariants of the surface: $p_a$ and $p^{(1)}$.

Let $|C|$ be a net of irreducible curves (of genus $\pi$ and degree $n$) without base points on the surface $F$; to it there will in general belong the following numerical characters:

1) the number $N$ of the \textit{intersections of the curves $C$ with the Jacobian of the net $C_j$}:
\[
N=2n+2\pi-2;
\]

2) the number $\delta$ of the \textit{curves $C$ of a pencil endowed with a double point}, which we know to be expressed by
\[
\delta=n+4\pi+12p_a-p^{(1)}+9;
\]

3) the \textit{genus} of the said \textit{Jacobian} $C_j$ $(=3C+K)$
\[
P=p^{(1)}+9\pi-9;
\]

4) the number $k$ of the \textit{three-point contacts} of the $C$ of the net (base points of a pencil of osculating $C$);

5) the number $i$ of the \textit{curves} of the net $|C|$ \textit{endowed with a cusp};

6) the number $d$ of the \textit{double contacts} of the $C$ of the net (base points of a pencil of bitangent $C$);

7) the number $t$ of the $C$ of the net endowed with two double points.

One may suppose that the net of the $C$ is cut out on $F$ by a star of planes through a point $O$, which will in general be multiple for $F$ and will represent an irreducible curve $L$ (which, added to the $C$, gives the system of plane sections).

Then $F$ is projected from $O$ onto an \textit{$n$-fold multiple plane}, whose points correspond to groups of $n$ points of $F$ belonging to the involution $I$ defined by the net $|C|$. And the aforesaid characters take on the following meaning:

1) $N$ is the \textit{order of the branch curve $D$} of the $n$-fold plane (the curve which is the locus of the points of the plane to which there correspond groups of $I$ having a coincidence);

2) $P$ is the \textit{genus of the branch curve $D$};

3) $\delta$ is the \textit{class} of the said \textit{curve $D$};

4) $k$ is the number of lines through $O$ that have a three-point contact with the surface $F$, and hence the \textit{number of cusps} of $D$;

5) $i$ is the number of cuspidal curves, plane sections of $F$ through $O$, and hence the \textit{number of inflexions} of the \textit{branch curve $D$};

6) $d$ is the number of lines through $O$ bitangent to $F$, and hence the \textit{number of nodes of $D$};

% scan: s103L
% !PRINTED 182
\weberpage{182}
7) $t$ is the number of planes through $O$ bitangent to $F$, and hence the \textit{number of double tangents} (with distinct contacts) \textit{of $D$}.

These characters are connected by the \textit{Plücker formulae} relating to the curve $D$, which in general is irreducible and possesses no other singularities, point or tangential.

In this way one succeeds in establishing the following relations:
\weq{A)}{
\left\{
\begin{aligned}
N&=2n+2\pi-2\\
P&=p^{(1)}+9\pi-9\\
\delta&=n+4\pi+12p_a-p^{(1)}+9\\
k&=3(n+6\pi+p^{(1)}-4p_a-11)\\
i&=24(\pi+p_a)\\
d&=2[(n+\pi)^2-5n-17\pi-2p^{(1)}+6p_a+24]\\
t&=\frac{1}{2}[(n+4\pi+12p_a-p^{(1)})^2+15n-6\pi+132p_a-17p^{(1)}+74].
\end{aligned}
\right.
}

\textit{These formulae express in general the characters of an irreducible net $|C|$ without base points on the surface $F$, or the Plückerian characters of the branch curve $D$ of the corresponding multiple plane, by means of the genus and the degree, $\pi$ and $n$, of $|C|$ and of the invariants $p_a$ and $p^{(1)}$ of the surface.}

These invariants are expressed by means of the characters of $D$ by the formulae:
\weq{B)}{
\left\{
\begin{aligned}
p^{(1)}&=P-9\pi+9\\
p_a&=\frac{i}{24}-\pi\\
&\text{or also}\\
4p_a&=n+P-3\pi-\frac{k}{3}-2\\
(\pi&=\frac{N}{2}-n+1).
\end{aligned}
\right.
}

\textit{Remark.} – Among the formulae $A)$ the fourth and the fifth are of particular importance, since they bring out two invariant characters of the net $|C|$ containing $n$ and $\pi$ linearly, namely $3(n+6\pi-k)$ and $24\pi-i$. %

But above all the latter of these two formulae is important, since it gives us the absolute invariant $p_a$. As has been said, this relation between $p_a$ and the number of cuspidal curves of a net was discovered by \textsc{G. H. Zeuthen}. And \textsc{F. Severi} has proved in a direct and elementary way that the expression $24\pi-i$ is the same for all nets without base points belonging to the surface $F$. One thus has a new proof of the invariance of $p_a$. %

% scan: s103R
% !PRINTED 183
\weberpage{183}
\edsection{7. $(1,n)$ correspondence between two surfaces: double planes.}{7. (1,n) correspondence between two surfaces: double planes.}{sec.V.7}

Let there be given a $(1,n)$ correspondence between two surfaces $f$ and $F$: to a point of $F$ there corresponds \textit{one} point on $f$, while to the points of $f$ there correspond on $F$ the groups of an \textit{involution} $I$ of order $n$ (every point of $F$ belongs to \textit{one} group of $I$). For simplicity we add the hypothesis that there are \textit{no} \textit{fundamental} points, that is, points of $f$ or $F$ to which there corresponds a curve on the other surface.

To an irreducible linear system $|c|$ given on $f$ there will correspond on $F$ a linear system $|C|$ belonging to the involution $I$. To the Jacobian $|c_j|$ of a net contained in $|c|$ there corresponds on $F$ a curve which forms part of the Jacobian of the net of the homologous $|C|$. But to complete the Jacobian of the latter it is necessary to add to it the curve of coincidences $D$, which corresponds to the branch curve on $f$; indeed, a point of $D$ is a base point of a pencil of $|C|$ within the aforesaid net, which has a fixed tangent, and therefore contains a $C$ endowed with a double point. This granted, recalling the definition of the canonical curves $(|K|=|C_j-3C|)$, one deduces the theorem:

\textit{Under a $(1,n)$ correspondence between the surfaces $f$ and $F$, the transforms of the canonical curves of $f$, added to the curve homologous to the branch curve, constitute canonical curves of $F$}\origfootnote{1}{s103R-1}{Cf. \textsc{Enriques}, \textit{Ricerche etc.}}. %

The theorem must be modified when there is on $f$ \textit{a fundamental point} to which there corresponds on $F$ a curve (or a curve and a group of points); \textit{this fundamental curve must be added to the images of the canonical curves of $f$ and to the curve of coincidence} in order to form a canonical curve of $F$.

The theorem stated holds also when the \textit{canonical curves} of $f$ are \textit{virtual}. One can apply it to the problem of the canonical and pluricanonical curves of a \textbf{double plane}.

Let there be given a surface $F$ containing a rational involution $I$ of the $2^\circ$ order, by means of which it is represented on the double plane
\[
z^2=f_{2n}(x,y)
\]
where the branch curve $f_{2n}$, of order $2n$, may have singular points, whose multiplicity is to be evaluated according to the conventions of Chap. III, § 7. On the double plane one takes the curve $\varphi_m$ of a linear system having no base points outside the singular points of $f_{2n}$; to $\varphi_m$ there corresponds on $F$ a curve having as adjoints the transforms of the $\varphi_{m-3}$ adjoint to $\varphi_m$ in the plane, added to the curve of coincidence homologous to $f_{2n}$. If, instead of a first adjoint, one considers a second adjoint $\varphi_{m-6}$, then the curve of coincidence counted twice must be added to its transform, and so on.

% scan: s104L
% !PRINTED 184
\weberpage{184}
Indeed, on the plane the virtual canonical curves $K$ are defined by the subtraction
\[
|K|=|\varphi_{m-3}-\varphi_m|,
\]
and therefore canonical curves on $F$ will be obtained by adding the $K$ to the curve of coincidence $f_{2n}$; by adding to these the homologues of the $\varphi_m$ one will therefore have the adjoints to these latter curves.

Let us apply the result obtained, taking in place of the $\varphi_m$ the branch curve $f_{2n}$ and the second adjoints to it: $f_{2n-6}$. To have on $F$ the second adjoints to the curve of coincidence counted twice, which corresponds to $f_{2n}$, it will suffice to add this curve counted twice to the transforms of the $f_{2n-6}$. Subtracting from the system thus obtained the same curve of coincidence counted twice, one will have the bicanonical curves of $F$, among which are those, belonging to the involution $I$, that have the $f_{2n-6}$ as images.

Hence: \textit{the curves second adjoint to the branch curve of the double plane are images of bicanonical curves} belonging to the involution $I$. This deduction agrees with what was seen in Chap. III, § 7: that the canonical curves of the surface are represented by the $f_{n-3}$ belonging to the half of the system second adjoint to the branch curve.

The reasoning extends to the pluricanonical curves: \textit{one obtains tricanonical curves of the double plane by adding the third adjoints of the branch curve to this same curve}, and so on.

But let us return to the bicanonical curves, in order to observe that, also in another way, \textit{one obtains bicanonical curves of the double plane by adding the branch curve to the adjoints of the curves $f_{n-3}$, images of the canonical curves}; it suffices to recall that the bicanonical system cuts groups of the canonical series on the curves of the canonical system.

Now note that the bicanonical system is transformed into itself by the involution $I$, while its curves do not in general belong to this involution. Hence we shall have to regard this system as a linear space of $P-1$ dimensions ($P$ denoting the bigenus of $F$) on which $I$ induces an involutory homography with two spaces of united points $S_r$ and $S_t$, where
\[
P-1=r+t+1.
\]

One of the two spaces, say e.g. $S_r$, will correspond to the system of the $f_{2n-6}$, second adjoints of $f_{2n}$, while the other space $S_t$ will correspond to the system formed by the branch curve $f_{2n}$ taken together with the $f_{n-6}$, adjoints of the images of the canonical curves. Hence, denoting by $\pi=t+1$ the genus of the $f_{n-3}$, we shall have
\[
P=r+\pi+1.
\]

\textit{This formula expresses the bigenus of the double plane} as a function of the genus $(r+1)$ of the curves $f_{2n-3}$ adjoint to the branch curve
% scan: s104R
% !PRINTED 185
\weberpage{185}
and of the genus $\pi$ of the curves that are halves of the bi-adjoints $f_{2n-6}$, that is, of the genus of the involution of the $2^\circ$ order which $I$ induces on a canonical curve. %

What has been said extends easily to the plurigenera. We shall confine ourselves to pointing out that if the plurigenera must all be zero, all the successive adjoints to the branch curve will be lacking, and therefore this curve can be transformed birationally into a curve of order $2n$ with a $(2n-2)$-fold point, or into a quartic, or into a sextic with two infinitely near triple points, or into a particular case of these\origfootnote{1}{s104R-1}{Theorem of \textsc{Castelnuovo-Enriques}, 1900. Cf. \textsc{Enriques-Conforto}, \textit{Le superficie razionali}, § 27, p. 489.}.

Therefore one can assert that:

\textit{A double plane having all its genera zero is rational or else referable to a hyperelliptic ruled surface.}

\edsection{8. Correspondence formulae.}{8. Correspondence formulae.}{sec.V.8}

The theorem which we gave in the preceding section regarding the transformation of the canonical curves in a $(1,n)$ correspondence between two surfaces $f$ and $F$ says first of all that, in every case, the genus $P_g$ of $F$ cannot be less than that of $f$:
\[
P_g\geq p_g.
\]

Moreover, the same theorem leads at once to a relation between the linear genera $P^{(1)}$ and $p^{(1)}$ of the two surfaces.

It is convenient to premise an observation concerning the branch curve $D$ of the correspondence on $f$. It will in general be endowed with a certain number of nodes, to which there correspond on $F$ groups of $n$ points with two distinct coincidences. Moreover $D$ will possess a certain number of cusps, to each of which there corresponds a group of $n$ points of which three have come to coincide. This granted, by a well-known theorem of Zeuthen, to an $n$-fold canonical curve on $f$ which meets the branch curve $d$ in $s$ points there corresponds a curve of genus %
\[
n(p^{(1)}-1)+\frac{s}{2}+1.
\]

This curve meets the curve of coincidence $D$ in $s$ points, so that, denoting by $\varrho$ the effective genus\origfootnote{2}{s104R-2}{It is assumed that the curve $d$ presents no singularities other than the nodes and the cusps, which have the meaning explained above with respect to our correspondence.} of the latter, we shall have
\[
P^{(1)}=n(p^{(1)}-1)+\varrho+\frac{3}{2}s.
\]

% scan: s105L
% !PRINTED 186
\weberpage{186}
\sourcewarning{The upper left part of the scan is covered. In the following three paragraphs only the legible fragments are preserved, without filling in the gaps. The marker of the first footnote is placed provisionally. The continuous text resumes with ``Let us therefore denote''.}
\textit{connects the linear genera of the two surfaces in cor-} In it one must put $\varrho=1$ and $s=0$ \textit{when there is lacking} to an $f$ without singularities) \textit{the branch cur-}\origfootnote{1}{s105L-1}{\textsc{Enriques}, \textit{Ricerche etc.}, 1893.} %

ly one can write a second relation arithmetic $P_a$ and $p_a$ of $F$ and $f$\origfootnote{2}{s105L-2}{Cf. \textsc{Severi}, \textit{Sulle relazioni che legano i caratteri invarianti di due superficie algebriche.} Rend. Ist. Lombardo, 1903.}. %

pose one will take a pencil of curves $|c|$ on $f$ and the pencil orresponding curves $|C|$ on $F$. By means of the characters of ese two pencils one will evaluate the Zeuthen-Segre invariants $i$ and $I$ of the two surfaces and write down the relation by which they are connected. %

Let us therefore denote by $m$ the number of base points of $|c|$, by $\pi$ the genus of the $c$, and by $\delta$ the number of the $c$ endowed with a double point; we shall have
\[
i=\delta-m-4\pi.
\]

Now the number $M$ of base points of the pencil $|C|$ and the genus $\Pi$ of the $C$ are given respectively by
\[
M=mn
\]
and
\[
\Pi=n(\pi-1)+1+\frac{t}{2}
\]
where $t$ denotes the number of intersections of a $c$ with the branch curve $d$ on $f$. As for the number $\Delta$ of the curves $C$ endowed with a double point, among these one will find first of all the $\delta$ curves which correspond to the $c$ endowed with a double point, each of which, having $n$ double points, must count for $n$. But there will in general be $\Delta-n\delta$ further curves, which are obtained from the $c$ that touch the curve $d$, with the exception of those which have an improper contact, passing through a cusp of $d$, to which there corresponds on $F$ a point in which three of the $n$ corresponding points come together.

Therefore, denoting by $\varrho$ the effective genus of the branch curve $d$, and by $\chi$ the number of its cusps, we shall have
\[
\Delta=n\delta+2t+2\varrho-2-\chi,
\]
hence we shall have the Zeuthen-Segre invariant of $F$ expressed by
\[
I=\Delta-M-4\Pi=n\delta+2\varrho-2-\chi-nm-4n(\pi-1)-4.
\]

% scan: s105R
% !PRINTED 187
\weberpage{187}
This granted, recalling Noether's relation
\[
\begin{aligned}
i&=12p_a-p^{(1)}+9\\
I&=12P_a-P^{(1)}+9,
\end{aligned}
\]
one will deduce
\[
24(P_a+1)=24n(p_a+1)+6\varrho-6+3s-2\chi.
\]

\textit{This relation connects the arithmetic genera of the two surfaces $f$ and $F$ in} $(1,n)$ \textit{correspondence}, under the simplifying hypothesis that there exist \textit{no} \textit{fundamental points} of the correspondence on $f$. Here too, \textit{if there is no} branch curve, one will take
\[
\varrho=1,\qquad s=0.
\]

The formulae which connect the arithmetic genera and the linear genera of two surfaces $f$ and $F$ in $(1,n)$ correspondence can also be written in another form, by introducing, in place of the number $s$, which denotes the number of intersections of the branch curve $d$ of $f$ with the canonical curves $k$, the degree $\nu$ of the curve $d$: a virtual degree, which is expressed by means of $s$ and of the virtual genus of $d$, $\varrho+\varepsilon+\chi$, where $\varepsilon$ denotes the number of nodes of $d$ and $\chi$ the number of cusps. We shall have
\[
s=2(\varrho+\varepsilon+\chi)-2-\nu
\]
and hence \textit{the formulae which connect the characters of $f$ and of $F$ will become}:
\[
\begin{aligned}
P^{(1)}&=n(p^{(1)}-1)+4\varrho+3(\varepsilon+\chi)-\frac{3}{2}\nu-3\\
24(P_a+1)&=24n(p_a+1)+12\varrho-12+6\varepsilon+4\chi-3. %
\end{aligned}
\]

The first formula can equally be deduced by comparing the degrees of the canonical system on $f$ and on $F$. Recalling that one obtains a canonical curve $K$ of $F$ by adding to the transform of a $k$ of $f$ the curve of coincidence $D$, which meets it in $s$ points, and denoting by $x$ the degree of $D$, we shall have:
\[
P^{(1)}-1=n(p^{(1)}-1)+x+2s.
\]

This formula agrees with the preceding one, if one computes the value of $x$, which is %
\[
x=\frac{\nu}{2}-(\varepsilon+\chi).
\]

In particular, if one has on $f$ a branch curve, or a component of it, of degree $\nu$ which is entirely free of double points, the curve of coincidence corresponding to it on $F$ has degree $\dfrac{\nu}{2}$. This appears clearly especially in the case of a $(1,2)$ correspondence, $n=2$, where the branch curve must in general be regarded as
% scan: s106L
% !PRINTED 188
\weberpage{188}
free of double points: to this curve of degree \(\nu\) there corresponds a double curve of degree \(2\nu\), whose component therefore has degree \(\frac{\nu}{2}\).

It is interesting to compare the formulae obtained here with those which give the arithmetic genus and the linear genus of multiple planes, which we gave in \S~6.

We must recover these formulae, taking into account that the (relative) linear genus of the plane is \(p^{(1)} = 10\)\origfootnote{1}{s106L-1}{Cf. Chap. II, \S~6; Chap. V, \S~3.}, while its arithmetic genus is \(p_a = 0\).

For this purpose let us take up again the notation used in the formulae \(B)\) of \S~6, thus putting \(P\) in place of \(\varrho\), \(k\) in place of \(\chi\), \(d\) in place of \(\varepsilon\) and \(N^2\) in place of \(\nu\), where
\[
N = 2n + 2\pi - 2.
\]

With these substitutions, our formulae give
\[
P^{(1)} = 9n + 4P + 3(d + k) - \frac{3}{2}N^2 - 3 %
\]
and since
\[
P = \frac{(N - 1)(N - 2)}{2} - (d + k),
\]
one deduces
\[
P^{(1)} = P - 9\pi + 9;
\]
which is the first of our formulae \(B)\).

Similarly the formula giving \(24(P_a + 1)\) becomes
\[
24(P_a + 1) = 24n + 12(P - 1) + 6d + 4k - 3N^2
\]
whence
\[
4P_a = n + P - 3\pi - \frac{k}{3} - 2
\]
which is the second of our formulae \(B)\).

\textit{In this way the formulae giving the characters of a multiple plane have been recovered.}

\textit{Examples.} -- Let us take up again the notation of this section, in order to apply the formulae obtained to some interesting examples.

Consider the \textit{Kummer surface}, the singular surface of the quadratic complex of lines of space, which is defined as a surface \(f_4\) of the \(4^\circ\) order endowed with 16 double points\origfootnote{2}{s106L-2}{Cf. \textsc{Enriques-Chisini}, \textit{Lezioni etc.}, Book VI, \S~40, vol. IV, p. 246.}. It is well known that \(f_4\), which is a surface dual to itself, possesses 16 double tangent planes, each of which touches it along a conic to which 6 double points belong. This follows at once from the representation of \(f_4\) on a double plane with branch sextic split into 6 lines, such as is obtained by projecting \(f_4\) from one of its double points. Now consid%
% scan: s106R
% !PRINTED 189
\unskip\weberpage{189}%
er 3 double tangent planes of \(f_4\) passing through a double point \(A\); these will intersect two by two in lines which project from \(A\) three double points, \(B\), \(C\), \(D\); hence it follows that the tetrahedron \(ABCD\), whose vertices are double for \(f_4\), contains simply the remaining 12 double points of \(f_4\), three on each face.

This granted, denoting by \(T(x, y, z) = 0\) the equation of this tetrahedron, let us construct in the space \(S_4(x\,y\,z\,u)\) the surface \(F\) with equation
\[
\begin{aligned}
&u^2 = T(x\,y\,z)\\
&(f_4(x, y, z) = 0),
\end{aligned}
\]
which is represented on the double \(f_4\).

It is important to observe that on this double surface there is no branch curve of order greater than zero. But there are 16 infinitesimal branch curves, constituted by the neighbourhoods of the 16 double points, through which \(T\) passes an odd number of times. From this it follows that our double \(f_4\) is irreducible; and moreover this also results from Comessatti's criterion\origfootnote{1}{s106R-1}{We recall that according to Comessatti's criterion the necessary and sufficient condition for an irreducible curve \(C\), tangent to the branch curve \(D_{2m}\) wherever it meets it, to be the image of a split curve is that the group of points of contact belong to the series cut out on \(C\) by the curves \(D_m\), halves of the branch curve.}, on observing that on each double quartic plane section of \(f_4\) there are 8 apparent critical points which fall at the points of contact of a quadrilateral and do not belong to a conic. The \((1, 2)\) correspondence which subsists between \(f_4\) and \(F\) leads from the system of plane sections of \(f_4\), which has genus \(\pi = 3\) and degree \(n = 4\), to a linear system of genus 5 and degree \(8 = (2 \cdot 5 - 2)\) on \(F\). One deduces that \(F\), like the double \(f_4\), is of genus
\[
p_g = 1
\]
and has, like the latter, a canonical curve of order zero, and hence all its plurigenera equal to one.

A priori the (absolute) linear genus of \(F\) will therefore be \(P^{(1)} = 1\). But let us evaluate it with our formula:
\[
P^{(1)} = n(p^{(1)} - 1) + 4\varrho + 3(\varepsilon + \chi) - \frac{3}{2}\nu - 3.
\]

For this purpose observe that the branch curve \(d\) on \(f_4\) consists of the 16 neighbourhoods of the 16 double points, each of which is an infinitesimal curve of genus zero and degree \(-2\) \((\varepsilon = \chi = 0)\). We shall therefore find
\[
P^{(1)} = -4 \cdot 15 + \frac{3}{2} \cdot 32 - 3 = -15.
\]

% scan: s107L
% !PRINTED 190
\weberpage{190}
This paradoxical conclusion is cleared up if one bears in mind that to each branch curve of genus zero and degree \(-2\) there corresponds a priori a curve of genus zero and degree \(-1\), which is exceptional, and which in our case appears as the neighbourhood of a simple point of \(F\). To construct the canonical curve of \(F\) we added to the transform of the canonical curve of order zero of \(f_4\) the transform of the branch curve; consequently we evaluated \(P^{(1)}\) taking into account the neighbourhoods of 16 simple points as if these were exceptional curves of \(F\) for the purposes of the computation of the linear genus. In other words, we evaluated the linear genus of a transformed surface of \(F\) with 16 exceptional curves. Eliminating these curves one finds, as must be the case:
\[
P^{(1)} = 1.
\]

Proceeding further to evaluate the arithmetic genus of \(F\), one will have
\[
24(P_a + 1) = 96 - 12 \cdot 15 - 12 + 3 \cdot 32 = 0
\]
that is
\[
P_a = -1.
\]

\textit{Note.} -- As was shown by \textsc{F. Klein}\origfootnote{1}{s107L-1}{Cf. \textsc{Enriques-Chisini}, l. c.}, the Kummer surface \(f_4\) can be represented as an involution on the surface \(\Phi\) whose elements are the pairs of points of a curve \(C\) of genus 2; precisely, one takes as conjugate on \(C\) the pairs of points which are residual to each other with respect to a \(g_4^2\). In this way one defines on \(\Phi\) an involution \(I\) which has 16 points of coincidence, corresponding to the 16 pairs of \(C\) arising from the bisection of the \(g_4^2\). Now, \(I\) being represented by \(f_4\), it follows that \(\Phi\) is represented by this \textit{double Kummer surface, with \(16\) branch points, which fall at its \(16\) double points.} Therefore one deduces that the surface \(\Phi\) is equivalent to the \(F\) whose equation we wrote above. And thus \textit{the hyperelliptic surface representing the set of pairs of points of a curve of genus \(2\) has the genera}
\[
p_g = 1,\qquad p_a = -1
\]
\textit{and all its plurigenera equal to \(1\).} We shall have occasion to return to this surface later\origfootnote{2}{s107L-2}{For the construction indicated above of the hyperelliptic surface as a double Kummer surface, with branch curve defined by the tetrahedron \(T\), cf. \textsc{Enriques-Severi}, \textit{Mémoire sur les surfaces hyperelliptiques.}, Acta Math., 1909.}. %

Another interesting example, with regard to our theory, is obtained by considering the surface of genus \(p_a = p_g = 0\) and of bigenus
% scan: s107R
% !PRINTED 191
\weberpage{191}
\(P = 1\), formed by the sextic \(f_6\) which passes doubly through the edges of a tetrahedron \(T\) (cf. Chap. II, \S~8). Let us construct the surface \(F\) defined by the equations
\[
u^2 = T(x, y, z),\qquad f_6(x, y, z) = 0,
\]
which is represented on the double \(f_6\). Since \(T\) and \(f_6\) have in common the 6 edges of the tetrahedron \(T\) counted 2 times, on \(f_6\) there is no branch curve of order \(> 0\). But neither are there infinitesimal branch curves, that is, branch points: a priori these could fall only at the vertices of \(T\); but one of these vertices \(A\) arises on \(f_6\) from the superposition of 3 simple points, whose neighbourhoods belong to the face planes of the tetrahedron passing through \(A\). Now, since \(A\) is \(4\)-fold for the section of \(f_6\) by a face \(\alpha\), the points infinitely near to \(A\) on \(\alpha\) will not be branch points of our double \(f_6\). This granted, we observe that the curves \(C\), plane sections of \(f_6\), are irreducible double curves (Comessatti's criterion) and hence represent on \(F\) curves \(C\) of genus \(\pi = 7\) and of degree \(N = 2 \cdot 6 = 2\pi - 2\). Since the double of \(|C|\) coincides with the double of the adjoint system: %
\[
|2C'| = |2C|,
\]
one deduces that on \(F\)
\[
|C'| = |C|
\]
and therefore the geometric genus of \(F\) is
\[
P_g = 1.
\]

\(F\) possesses a canonical curve of order zero, hence its linear genus
\[
P^{(1)} = 1
\]
and all its plurigenera \(P_i = 1\).

Let us compare the characters of \(F\) with those of the double surface \(f_6\), which has
\[
p_a = p_g = 0,\qquad p^{(1)} = 1
\]
(cf. \S~II, 8). Our formulae give us %
\[
\begin{aligned}
&n = 2,\qquad \varrho = 1,\qquad \nu = 0,\qquad \varepsilon = \chi = 0,\\
&\qquad P^{(1)} = 2(p^{(1)} - 1) + 4\varrho - 3 = 1
\end{aligned}
\]
and
\[
24(P_a + 1) = 2 \cdot 24(p_a + 1) = 48 %
\]
whence
\[
P_a = 1.
\]

\textit{The sextic \(f_6\) passing doubly through the edges of a tetrahedron, which --- taken as simple --- is a regular surface of genus zero and of}
% scan: s108L
% !PRINTED 192
\weberpage{192}
\textit{bigenus one with a bicanonical curve of order zero, can be regarded as a double surface with genera equal to \(1\) and with a canonical curve of order zero}\origfootnote{1}{s108L-1}{Cf. \textsc{F. Enriques}, \textit{Un’osservazione relativa alle superficie di bigenere uno.} Rendic. Acc. Bologna, 1908.}.

\edsection{9. Existence theorem.}{9. Existence theorem.}{sec.V.9}

Let us return to the consideration of multiple planes. When a surface \(F\) with given genera \(p_a\) and \(p^{(1)}\) is represented on a multiple plane of a certain order \(n\), corresponding to a net of curves on it of genus \(\pi\) and degree \(n\), we have learned to evaluate the characters of the branch curve of the \(n\)-fold plane by means of \(n\), \(\pi\), \(p_a\) and \(p^{(1)}\) (cf. \S~8).

The formulae thus obtained already indicate that the branch curve of an \(n\)-fold plane, for \(n > 2\), cannot be taken arbitrarily. To be sure, given any plane curve
\[
f(x, y) = 0
\]
one can construct the (\textit{cyclic}) \(n\)-fold plane
\[
z^n = f(xy)
\]
which has precisely \(f = 0\) as branch curve.

But the branch curve of this \(n\)-fold plane does not consist of the irreducible \(f\) alone, but of \(f\) counted \(n - 1\) times, to which there is further added the line at infinity, counted once or more times, when the order \(m\) of \(f\) is not a multiple of \(n\)\origfootnote{2}{s108L-2}{To see this it suffices to refer to an \(n\)-fold line, or better to an \(n\)-sheeted Riemann surface, constructed by taking in the plane of the complex variable \(m\) branch points, to each of which there corresponds the cyclic substitution among the branches \((1\,2\,3\ldots n)\). A circuit encircling all the \(m\) branch points gives rise to a permutation among the branches expressed by \((1\,2\,3\ldots n)^m\), and this is the substitution corresponding to a circuit around the point at infinity, which therefore actually turns out to be a branch point when the said substitution is not the identity, that is, when \(m\) is not a multiple of \(n\).}. Thus, then, the branch curve of a cyclic \(n\)-fold plane is by no means an arbitrarily given irreducible curve.

Since the branch curve of an \(n\)-fold plane for \(n > 2\) cannot be given arbitrarily, there arises the problem of ``determining the conditions for a plane curve \(f(xy) = 0\) to be the branch curve of an \(n\)-fold plane''. These conditions will lead to an \textit{existence theorem} for multiple planes or for algebraic functions of two variables, a theorem which has a sense analogous to, but in a certain
% scan: s108R
% !PRINTED 193
\weberpage{193}
respect different from, the existence theorem for \(n\)-fold lines or for algebraic functions of one variable given by \textsc{Riemann}\origfootnote{1}{s108R-1}{Cf. e.g. \textsc{Enriques-Chisini}, \textit{Lezioni.} Book V, chap. III, \S~33 (vol. III, p. 355).}. Indeed, in Riemann's case the group of branch points can be given arbitrarily, and corresponding to it there is a finite number of classes of algebraic functions (or of \(n\)-fold lines), each of which depends on assigning the substitutions among the branches in relation to a system of loops encircling the branch points; here, on the other hand, as has been said, it is a question of finding the particular conditions which the branch curve must satisfy for the existence of the \(n\)-fold plane, and, at least in general, when these conditions are satisfied, there will be only one class\origfootnote{2}{s108R-2}{Cf.: \textsc{O. Chisini}, \textit{Sulla identità birazionale delle funzioni algebriche di due variabili dotate di una medesima curva di diramazione.} Rend. Ist. Lombardo di Scienze e Lettere, 1944.} of surfaces represented on the plane, that is, a unique \(n\)-fold plane, to which there corresponds a determined system of substitutions in relation to the loops encircling the branch curve, of which we speak below.

Nevertheless, the algebraic meaning of the question remains analogous to that of the construction of an algebraic function of a single variable whose branch points are given. When the surface \(F(xyz) = 0\) is projected from the point at infinity \(Z\) of the \(z\)-axis onto the plane \(z = 0\), one has in general in this plane a curve, the locus of the critical points of \(z(xy)\), which is represented by equating to zero the discriminant of this function,
\[
\Delta(xy) = 0
\]
that is, the resultant of the two equations
\[
F(xyz) = 0,\qquad \frac{\partial F}{\partial z} = 0.
\]

However, this curve \(\Delta = 0\) in general decomposes into two parts,
\[
\Delta = \Delta_1 \cdot \Delta_2^2,
\]
one of which, \(\Delta_2 = 0\), is an apparent critical curve, the projection of the double curve of \(F\), while the other, \(\Delta_1 = 0\) (the so-called \textit{principal part of the discriminant}), properly constitutes the branch curve of the algebraic function \(z(xy)\). Our problem consists in ``assigning the conditions for a plane curve to be the branch curve of an algebraic function of two variables'' and hence in the ``construction of this algebraic function with given branch curve''. The formulation of the problem itself can be made precise by assuming that the plane curve \(f\), which is required to be the
\printersignature{\textsc{Enriques F.} - \textit{Superficie algebriche.}\quad 13}
% scan: s109L
% !PRINTED 194
\weberpage{194}
branch curve of an \(n\)-fold plane, is a curve of even order \(2m\)\origfootnote{1}{s109L-1}{If it were of odd order, the line at infinity would necessarily be added to it.}, the projection from a point \(O\) of the apparent contour \(C\) of a surface \(F\), so that \(f\) possesses a certain number \(d\) of nodes and a certain number of cusps, which correspond to apparent double points of the said curve \(C\): besides these double points, \(f\) may possess nodes and cusps corresponding to isolated double points (respectively conical points and biplanar points) of \(F\); but one can leave these double points of \(f\) out of account, considering them as virtually non-existent. %

This premised, let us consider in the plane of the given curve \(f\) a pencil of lines, whose centre, for simplicity of discussion, we may suppose not to belong to \(f\); and let it be, for example, the pencil \(y = tx\) with centre at the origin \(O\). For a generic value of \(t\), one has an \(n\)-fold line on which \(2m\) simple branch points are marked; in the representative plane of the complex variable \(x\) one takes \(2m\) loops issuing from the point \(O\) and encircling these \(2m\) points; one can assign arbitrarily \(2m\) transpositions on the branches corresponding to the said loops, with the condition that the group generated by them be transitive and that their product be the identity: in this way one finds a finite number of irreducible curves, in general birationally distinct, which are represented on the \(n\)-fold line \(y = tx\) and which are all of genus
\[
\pi = \frac{m}{2} - n + 1; %
\]
these curves are defined up to birational transformations, that is, they correspond to determined classes of curves or algebraic functions, generally distinct.

The number \(r\) of the said curves \(C_1\,C_2\,\ldotp\ldotp\ldotp\ldotp\,C_r\) was determined by \textsc{Hurwitz}, but here it is not necessary to know it. It is important instead to bear in mind that, when \(t\) is made to vary describing a closed path in the plane \(\tau\) of this complex variable, the curves \(C_1,\ C_2,\ \ldotp\ldotp\ldotp\ldotp,\ C_r\) (or the classes of curves to which they belong) will in general be permuted among themselves: this happens in the first place because the \(2m\) branch points, intersections of the line \(y = tx\) with \(f\), are permuted among themselves, and also because the system of loops with which one constructs, e.g., the \(n\)-sheeted Riemann surface of \(C_1\) may be reduced by continuity to a new system of loops, which can be brought back to the old one by suitably changing the relative transpositions. Now if \(f\) is the branch curve of an \(n\)-fold plane, the projection of a rationally determined surface \(F\), one of the said curves, say e.g. \(C_1\), being a section of \(F\), must be rationally
% scan: s109R
% !PRINTED 195
\weberpage{195}
determined, and therefore its Riemann surface must not be interchanged with those of \(C_2,\ \ldotp\ldotp\ldotp\ldotp\,C_r\), but must always return to \(C_1\) for every closed circuit described by the point \(t\) in the plane of this complex variable.

It is a question, in the first place, of establishing the \textit{conditions of invariance of} (the class containing) \(C_1\) or \textit{of its Riemann surface, with respect to all the possible closed paths described by \(t\) in the plane \(\tau\).} These conditions will certainly be necessary conditions for the existence of an \(n\)-fold plane having \(f\) as branch curve.

Let us state precisely the sense of our investigation. First of all, choose an initial value of \(t\), say e.g. \(t = 0\), and in the plane of the complex variable \(x\) assign an ordered system of loops issuing from \(O\) and encircling the branch points \(A_1\,A_2\,\ldotp\ldotp\ldotp\ldotp\,A_{2m}\) (intersections of \(f\) with the line \(y = 0\)); then, with reference to a given naming of the branches \(1,\ 2,\ \ldotp\ldotp\ldotp\ldotp\,n\) of the algebraic function \(z(x)\), take note of the transpositions which correspond to the said loops \(OA_i\), thus succeeding in defining the \(n\)-sheeted Riemann surface representing our \(C_1\) (the one whose conditions of invariance with respect to the closed circuits described by \(t\) in the plane \(\tau\) are being sought).

Having defined the initial Riemann surface, among those which represent the multiple lines \(y = tx\), we shall make \(t\) vary by continuity (starting from the point \(t = 0\)), describing a closed path in the plane \(\tau\) of this variable; while \(t\) varies in this way, \(C_1\) also varies by continuity, and its Riemann surface can be constructed with sheets superposed on the same plane of the complex variable \(x\), since (in the language of algebraic geometry) the multiple line \(y = tx\) --- which has as branch points its intersections with \(f\) --- can in general be replaced by its projection onto the \(x\)-axis made from the point at infinity of the \(y\)-axis.

Now, corresponding to a continuous variation of \(t\) in the plane \(\tau\), excluding the positions of \(t\) to which there correspond two coincident points \(A_r\) and \(A_s\), we shall see in the plane of the complex variable \(x\) the points \(A_1\,A_2\,\ldotp\ldotp\ldotp\ldotp\,A_{2m}\) move continuously without ever coinciding, describing arbitrary curves (closed or open), and it is permissible to let the loops \(OA_i\) which encircle them vary, always continuously, in such a way that they --- preserving the order in which they follow one another --- remain open and free curves, not crossing one another.

After this, consider in the plane \(\tau\) a closed circuit described by the point \(t\), starting from the origin \(t = 0\). If this circuit does not enclose any \textit{critical point} to which there correspond two coincident points \(A_i\), it can be deformed by continuity until it reduces to a null cycle; therefore, corresponding to it, the ordered system of loops and the relative substitutions which define the Riemann surface \(C_1\) must return into itself: in short, the said cycle gives rise to no condition for the invariance of \(C_1\).

% scan: s110L
% !PRINTED 196
\weberpage{196}
Therefore the conditions of invariance of \(C_1\) will reduce to those which correspond to elementary circuits or loops of the plane \(\tau\) encircling a critical point \(t\) for which two points \(A_i\) coincide.

Now such circuits can belong to three different kinds:

1) circuits which encircle a point \(T\) to which there corresponds a simple tangent of the curve \(f\) drawn from the external point \(O\);

2) circuits which encircle a point \(D\) to which there corresponds a line through \(O\) going to a node of the curve \(f\); and

3) circuits which encircle a point \(Q\) to which there corresponds a line through \(O\) going to a cusp of \(f\).

For each of these three cases it is easy to determine the conditions of invariance of the Riemann surface \(C_1\), beginning with the simple hypothesis in which the two points \(A_1A_2\) that come to coincide at the critical point (\(T\) or \(D\) or \(Q\)) belong to loops \(OA_1\) and \(OA_2\) that are contiguous and \textit{honestly near}: by this we mean that in the limit (when \(A_1\) and \(A_2\) are superposed) the two curves \(OA_1\) and \(OA_2\), having in common the end point \(A_1 = A_2\), bound a surface in which none of the other points \(A_i\) falls, so that these curves can be brought infinitely near to each other without crossing other loops of the variable Riemann surface, and hence also without changing the substitutions on the branches of \(z(x)\) which correspond to them.

Let us examine in succession the three cases defined above:

1) If \(A_1\) and \(A_2\) come to coincide at a point \(T\), on the limiting multiple line \(OT\) the two branch points \(A_1\) and \(A_2\) are lost, merging into an apparent critical point, the image of a double point of \(C_1\) (and after the circuit of \(t\) they are interchanged with each other); this means that the two transpositions relating to the two (honestly near) loops \(OA_1\) and \(OA_2\) must be \textit{identical}. Conversely, if this happens, for example if the same transposition \((12)\) corresponds to \(OA_1\) and \(OA_2\), the Riemann surface \(C_1\) remains invariant under the loop of the plane \(\tau\) encircling \(T\).

2) If \(A_1\) and \(A_2\) come to coincide at a point \(D\), on the limiting multiple line \(OD\) the two branch points \(A_1\) and \(A_2\) are superposed at a point which is the image of two distinct points of \(C_1\). This means that the transpositions relating to the two (honestly near) loops \(OA_1\) and \(OA_2\) must bear on distinct pairs of branches, that is, they must be permutable and in general \textit{disjoint}, like \((12)\) and \((34)\)\origfootnote{1}{s110L-1}{They will instead be identical, and therefore still permutable, if the point \(D\) is a virtually non-existent node of \(f\).}. Conversely, if this happens, the Riemann surface \(C_1\) remains invariant under the loop of the plane \(\tau\) encircling \(D\).

% scan: s110R
% !PRINTED 197
\weberpage{197}
3) Finally, if \(A_1\) and \(A_2\) come to coincide at a point \(Q\), on the limiting multiple line \(OQ\) the two branch points \(A_1\) and \(A_2\) are superposed at \(Q\), and as a result of the circuit of \(t\) they are interchanged; but the point \(Q\) becomes the image of an inflexion of \(C_1\), so that the transpositions relating to \(A_1\) and \(A_2\) must give as their product a cyclic substitution of the \(3^\circ\) order, bearing therefore on two pairs of branches having one branch in common; that is, they must be \textit{concatenated}, like \((12)\) and \((23)\).

The conditions 1), 2) and 3) ensure the invariance of the Riemann surface \(C_1\) with respect to the circuits encircling the points \(T\), \(D\) and \(Q\) in the plane \(\tau\), when the two points \(A_rA_s\) which come to coincide at these points are the end points of loops \(OA_r\) and \(OA_s\) which become honestly near. The case in which \(OA_r\) and \(OA_s\) do not become honestly near can be brought back to the preceding one: in this case, indeed, it happens that the two aforesaid curves, which come to have the common end point \(A_r = A_s\), will form a closed curve bounding a surface within which some other of the points \(A_i\) falls; but one can always bring the curve \(OA_s\) infinitely near to \(OA_r\), making it vary by continuity across the said points \(A_i\): this means that as a result of this reduction, the order of the loops being changed, the transposition corresponding to the loop \(OA_i\) also changes\origfootnote{1}{s110R-1}{Cf. e.g. \textsc{Enriques-Chisini}, \textit{Lezioni}, Book V, \S~3, vol. III, p. 18.}. %

We can summarize the result obtained by saying that: \textit{the conditions of invariance of a Riemann surface \(C_1\), with regard to its variations in the pencil of \(n\)-fold lines \(y = tx\) with branch points on the curve \(f\), relate to the circuits which \(t\) performs in the plane \(\tau\), encircling respectively the critical points \(T\), \(D\) and \(Q\)} (points which correspond to the proper tangents and to the lines going to the nodes and to the cusps of the branch curve \(f\)); \textit{they are expressed by saying that the transpositions corresponding to the branch points \(A_r\) and \(A_s\) which come to coincide at one of those critical points must be respectively identical, or permutable} (and in general \textit{disjoint}) \textit{or concatenated, when the system of loops which defines the Riemann surface is transformed in such a way that the loops encircling the said points, \(OA_r\) and \(OA_s\), become honestly near.}

The conditions expressed above are \textit{necessary conditions} for the curve \(f\), endowed with nodes and cusps, to be the branch curve of an \(n\)-fold plane \((n > 2)\). We wish to show that they are also \textit{sufficient} for the existence of such an \(n\)-fold plane.

For this purpose consider the pencil of planes through the \(z\)-axis having as traces, in the plane \((z = 0)\) of the curve \(f(xy) = 0\), the lines \(y = tx\). In each of these planes we want to construct a well-defined curve which, by projection from the point at infinity of the \(z\)-axis, let us call it \(Z\), gives rise to the \(n\)-fold line \(C_1\) with branch points at \(A_1\,A_2\ldotp\ldotp\ldotp\ldotp\,A_{2m}\)
% scan: s111L
% !PRINTED 198
\weberpage{198}
intersections with the curve $f$. The curve $C_1$, as was said above, is defined as a Riemann surface in a way which is invariant with respect to the closed circuits of $t$ in the plane; one can therefore construct an irreducible projective image of $C_1$, that is, a curve $C^{(n+h)}$, of a certain sufficiently high order $n+h$, which passes $h$ times through the point $Z$ and touches the lines $ZA_i$. On such a curve $C^{(n+h)}$ one can also impose the condition of passing through $n$ distinct points $O_1\,O_2\,\ldotp\ldotp\ldotp\ldotp\,O_n$ of the line $OZ$, and moreover of satisfying a certain number of linear conditions, depending on $h$; one will thus find a finite number of irreducible curves $C^{(n+h)}$, and among these a finite number $s$ of curves represented by the Riemann surface $C_1$, to which there will correspond $s$ equally branched algebraic functions $z_1(x),\,z_2(x),\,\ldotp\ldotp\ldotp\ldotp\,z_s(x)$\origfootnote{1}{s111L-1}{Cf. \textsc{Enriques-Chisini}, \textit{Lezioni}. Book V, § 33, vol. III, p. 355.}; then the sum
\[
z(x)=z_1(x)+z_2(x)+\ldotp\ldotp\ldotp\ldotp+z_s(x),
\]
will likewise correspond to an irreducible curve $C$ represented on the Riemann surface $C_1$ and rationally defined. Let us make this $C$ vary with the parameter $t$, and we shall obtain a surface $F$ represented on the $n$-fold plane $z=0$ with the branch curve $f$. It is obvious that $f$ forms part of the branch curve of the $n$-fold plane obtained by projecting $F$ from the point $Z$; one might however suspect that to this curve there is added a certain number of branch lines issuing from $O$; but if this were so, some plane through the line $OZ$ would have to touch $F$ along a curve, and this is impossible if the points $O_1O_2\ldots O_n$ are distinct, as has been supposed.

In conclusion: \textit{The conditions of existence of an $n$-fold plane $(n>2)$ with assigned branch curve $f(xy)=0$, endowed with nodes and cusps,} (and not passing through the origin $O=[00]$) \textit{are expressed by means of the conditions of invariance of one of the Riemann surfaces represented on the multiple line $y=tx$, in relation to the closed paths which the parameter $t$ can describe, starting from the origin $O$ in the plane of its own complex variable, encircling the points which correspond to the tangents drawn to $f$ from $O$ or to the lines joining $O$ to its nodes or to its cusps.}

It is essential to observe that the conditions required for the existence of a multiple plane having a certain branch curve $f$ are \textit{qualitative conditions}, which cannot vary when $f$ itself varies continuously. Hence the existence theorem which we have established gives rise to the following:

\textit{Corollary. -- Given in the plane a continuous system of curves $f$, endowed with nodes and cusps, if one of these curves is the branch curve of an $n$-fold plane, the same happens for all the curves of the system.}

% scan: s111R
% !PRINTED 199
\weberpage{199}
\edsection{10. Complements: limiting forms of the branch curve according to Chisini.}{10. Complements: limiting forms of the branch curve according to Chisini.}{sec.V.10}

The existence theorem for multiple planes, with its corollary, set forth in the preceding section, was established by \textsc{F. Enriques} in the memoir ``Sulla costruzione delle funzioni algebriche di due variabili possedenti una data curva di diramazione'', published in the Annali di Matematica s. IV, t. I (1923).

Interesting developments are, moreover, connected with this theorem. First of all \textsc{O. Zariski}\origfootnote{1}{s111R-1}{\textsc{O. Zariski}, \textit{On the problem of existence of algebraic functions of two variables possessing a given branch curve}. Am. J. Math., vol. 51 (1929). \textit{On the linear connection index of the algebraic surfaces $z^n=f(xy)$}. Proc. Mat. Accad. Sc. U. S. A., vol. 15 (1929). \textit{On the irregularity of cyclic multiple planes}. Ann. of Math., II s., vol. 32 (1931). Cf. \textit{Algebraic surfaces}, Berlin, Springer, 1935.} transformed Enriques's conditions of existence into a topological and group-theoretic form, from which he deduces in particular that a curve without cusps cannot be the branch curve of a multiple plane of order $n>2$, except in the case of the curve counted $n-1$ times, which corresponds to a cyclic $n$-fold plane. %

But, whether in Enriques's form or in Zariski's, the existence theorem for multiple planes encounters considerable difficulties in its practical application, giving rise to the further problem of determining the branch curves satisfying the conditions laid down. An important contribution to this problem was made by \textsc{O. Chisini}\origfootnote{2}{s111R-2}{\textsc{O. Chisini}, \textit{Un teorema di esistenza dei piani multipli}. Notes I and II, Atti Acc. Lincei, VI, 19 (1934). \textit{Un più generale teorema di esistenza dei piani multipli}. Atti Acc. Lincei, VI, 27 (1938). \textit{Altre curve di diramazione dei piani $n$-pli}. Atti Acc. Lincei, VI, 29 (1939).}, with the investigation of the \textit{limiting forms} contained in a continuous system of branch curves endowed with a certain number of nodes and cusps.

To understand \textsc{Chisini}'s investigation, consider a surface $F$ of a certain order $m+h$, with an $h$-fold point $Z$, which comes to split into two parts $f_1$ and $f_2$ meeting the lines through $Z$ in $n$ and $s$ points respectively $(m=n+s)$. Here the branch curve of the $m$-fold plane onto which $F$ is projected from $Z$ will split into the two branch curves of the $n$-fold plane and of the $s$-fold plane corresponding to $f_1$ and to $f_2$, and into the double curve, projection of the curve common to $f_1$ and $f_2$, which becomes an apparent critical curve. This observation leads one to investigate whether one can construct the branch curve of an $m$-fold plane, with $m=n+s$, by adding to the branch curves of the two $n$-fold and $s$-fold planes a double curve, which must behave suitably with respect to them, %
% scan: s112L
% !PRINTED 200
\weberpage{200}
so that the composite curve, constructed as has been said, represents the limit curve of a continuous system of branch curves of $m$-fold planes, endowed with a certain number of nodes and cusps.

What has been said serves to explain the genesis of \textsc{Chisini}'s idea. However, the general search for the double curves to be added to the sum of the branch curves of an $n$-fold plane and of an $s$-fold plane still gives rise to considerable difficulties. To overcome them, \textsc{Chisini} restricts himself to the case in which, e.g., $s=1$, that is, in which one has only the branch curve of an $n$-fold plane, from which, by the addition of a suitable double curve, one wishes to derive the branch curve of an $(n+1)$-fold plane. It will suffice to explain Chisini's construction by referring to a simple example.

Let us set ourselves the task of constructing the system of branch curves of the quadruple planes which arise by projection of the general surface $F_4$ of the $4^\circ$ order. These branch curves are $C_{12}$ of order 12 with $d=12$ nodes and $k=24$ cusps.

If we make the surface $F_4$ pass through the centre of projection $Z$, the branch curve $C_{12}$ of the quadruple plane degenerates into the branch curve $C_{10}$ of a triple plane, endowed with $k'=18$ cusps (and no longer with nodes), and into a line $C_1$, a double tangent of $C_{10}$, counted twice. In its turn $C_{10}$ is capable of degenerating into a sextic $C_6$, entirely general, the branch curve of a double plane, and into a conic tangent to it at six points, counted twice; and this corresponds to the condition that $F_4$ be required to pass doubly through the point $Z$. Finally, if the double point $Z$ of $F_4$ becomes triple, the corresponding branch curve $C_6$ of the aforementioned double plane degenerates into a cubic counted twice.

It is now useful to examine what the singularities of the branch curve of the multiple plane become under the successive degenerations indicated above.

First of all, when the aforesaid $C_{12}$ degenerates into $C_{10}+2C_1$, each of the two points of contact $P_1$ and $P_2$ of $C_1$ and $C_{10}$ appears as the limit of 3 cusps of $C_{12}$, and each of the other 6 points common to $C_1$ and $C_{10}$ turns out to be the limit of two nodes of $C_{12}$. The assertion is justified by referring to the surface $F_4$ which comes to pass through the centre of projection $Z$; one must take into account that the cusps of $C_{12}$ are obtained by projecting from $Z$ the intersections of $F_4$ with the $1^{\mathrm{a}}$ and the $2^{\mathrm{a}}$ polar of the point $Z$, and that the nodes of $C_{12}$ come from the bitangents to $F_4$ passing through $Z$, which now become the six tangents drawn through $Z$ and touching elsewhere the quartic section of $F_4$ by the tangent plane at $Z$, each counted twice.

What has been observed regarding the degeneration of the singularities of $C_{12}$ when it becomes $C_{10}+2C_1$ still holds with re%
% scan: s112R
% !PRINTED 201
\unskip\weberpage{201}%
gard to $C_{10}$ degenerating into $C_6+2C_2$; that is to say, the 6 points of contact of $C_2$ with $C_6$ each come to absorb 3 cusps of $C_{10}$.

We have observed how the branch curve $C_{12}$ of a quadruple plane representing an $F_4$ can be made to degenerate successively into a $C_{10}+2C_1$, and then into a $C_6+2C_2+2C_1$, and finally also, if one wishes, into a $2C_3+2C_2+2C_1$, where $C_{10}$ is the branch curve of a triple plane, bitangent to the line $C_1$, and $C_6$ (the branch curve of a double plane) is tangent at six points to the conic $C_2$. Now this analytic process of reduction can be inverted, thus leading to a synthetic process of construction of multiple planes, which constitutes precisely \textit{Chisini's existence theorem}. Indeed, the author goes back from the limiting forms to more general forms of the multiple planes, recovering the triples of cusps and the pairs of double points which are absorbed by the limiting singularities. Thus, starting from the double plane with branch sextic $C_6$ and particularizing $C_6$ so that it admits a conic $C_2$ tangent to it at six points, one will regard the curve $C_6+2C_2$ as a particular case of a $C_{10}$, for which Enriques's conditions, explained in the preceding section, will be satisfied, and which will therefore turn out to be the branch curve of a triple plane. And then, adding further to $C_{10}$ a bitangent $C_1$, one will prove similarly that $C_{10}+2C_1$ (or $C_6+2C_2+2C_1$) is a limiting form of a $C_{12}$ which constitutes the branch curve of a quadruple plane representing a surface $F_4$ with plane sections of genus 3. Needless to say, the case to which we have referred is only a particular, though characteristic, example of a method which leads the author to the recognition of the conditions of existence for a wider family of multiple planes. We do not propose to set forth here \textsc{Chisini}'s more general developments, but shall confine ourselves to clarifying his constructive method, and also the delicate questions to which it gives rise, by referring to a notable family of surfaces which includes the quartics $F_4$, and to which we shall have to return in the sequel of these Lectures: we mean the surfaces with genera $p_a=p_g=1$ and $p^{(1)}=1$ possessing a canonical curve of order zero. %

In accordance with the formulae of § 6, one obtains surfaces belonging to the aforesaid family by starting from the $n$-fold multiple planes, with $n=2\pi-2$, defined by a branch curve $C_{3n}=C_{6\pi-6}$ of order $2n+2\pi-2=3n=6\pi-6$, endowed with $k(\pi)=24(\pi-2)$ cusps and $d(\pi)=18\pi^2-78\pi+84$ nodes: indeed, such an $n$-fold plane (if it exists) has the genera $p_a=p^{(1)}=1$ and also $p_g=1$, possessing a canonical curve of order zero; and the images of the lines of the plane belong to a linear system of genus $\pi$ and degree $n=2\pi-2$, and hence of dimension $\pi$. Thus the $n$-fold plane

% scan: s113L
% !PRINTED 202
\weberpage{202}
can be made to arise by projection, from external points, of a surface $F_n$ of order $n=2\pi-2$ in $S_\pi$, with canonical hyperplane sections of genus $\pi$.

Now here we shall attempt to prove the existence of $n$-fold planes corresponding to the branch curves $C_{3n}$ indicated above, and hence of the $F_n=F_{2\pi-2}$ of $S_\pi$ with canonical sections of genus $\pi$, for all values of $\pi=3,4,5\ldots$. Already in the preceding discussion we indicated how one succeeds in constructing the branch curve $C_{12}$ of a quadruple plane (representing an $F_4$) corresponding to $\pi=3$, starting from the double plane with branch sextic $C_6$ tangent at six points to a conic $C_2$. Now we shall proceed similarly to construct the sextuple plane corresponding to $\pi=4$, starting from a $C_{12}$ tangent at six points to a conic $C_2$: we shall ask whether $C_{12}+2C_2$ is the degenerate form of a $C_{16}$, the branch curve of a quintuple plane, and whether, by then adding to it a bitangent line $C_1$ counted twice, one arrives at a $C_{18}$, the branch curve of a sextuple plane. If this happens, whereas the $C_{12}$ indicated above possesses $k(\pi)=k(3)=24$ cusps and $d(\pi)=d(3)=12$ nodes, the $C_{16}$ will have in addition 18 cusps (3 cusps corresponding to each of the 6 contacts of $C_{12}$ with $C_2$) and 24 nodes (2 nodes arising from each of the 12 simple intersections of $C_{12}$ with $C_2$), and thus in all 42 cusps and 36 nodes; then, by adding to the $C_{16}$ the bitangent $C_1$ counted twice, one will obtain a $C_{18}$ having in addition another 6 cusps and another 24 nodes, and thus %
\[
k(4)=48\ \text{cusps}
\]
and
\[
d(4)=60\ \text{nodes}.
\]

In general, in order to continue this constructive procedure in correspondence with the successive values of $\pi$, it is useful to anticipate a notion which will be established in the following section, namely that the surfaces $F_n$ (with genera $p_a=p_g=p=1$, $p^{(1)}=1$) depend on at least $9p-2p^{(1)}+12=19$ essential parameters or moduli, invariant under birational transformations. Consequently, when there is given a multiple plane defined by a branch curve $C_{3n}=C_{6\pi-6}$ with the aforesaid characters, one will be able, by disposing of one parameter, to specialize the $C_{3n}$ so that it possesses a conic $C_2$ tangent to it at six points. And, in order to follow Chisini’s method, we are always led to ask ourselves whether $C_{3n}+2C_2$ can be regarded as the limit form of a $C_{3n+4}$, of order $3n+4$, which possesses, besides the cusps and the nodes of the $C_{3n}$, 18 cusps and $12n-24=24(\pi-2)$ nodes, corresponding respectively to the 6 points of contact of $C_2$ with $C_{3n}$ and to the $6n-12$ further intersections of these curves; and whether the $C_{3n+4}$ so constructed turn out to be branch curves of $(n+1)$-fold planes. %

If this can be done (in such a way that the conditions
% scan: s113R
% !PRINTED 203
\weberpage{203}
of invariance of the preceding section are satisfied), the $C_{3n+4}$ so constructed will turn out to be branch curves of $(n+1)$-fold planes. And, as we shall say, by adding to them a suitable bitangent line $C_1$, one will then obtain the branch curve of an $(n+2)$-fold plane, representing an $F_{n+2}$ with canonical sections of genus $\pi+1$ in the space $S_{\pi+1}$.

But the question expressed above does not always admit an affirmative answer. For $C_{3n}+2C_2$ to be effectively the limit form of a $C_{3n+4}$ satisfying our requirements, it is necessary that its singularities resolve into nodes and cusps in the manner indicated, and that for the curve $C_{3n+4}$ so constructed there be verified the conditions of invariance which allow one to regard it as the branch curve of an $(n+1)$-fold plane. Now the totality of these conditions is not satisfied for an arbitrary conic tangent at six points to $C_{3n}$; but, referring, for example, to the case $n=4$, it can be satisfied by reducing the $C_{12}$ to Chisini’s limit form $C_{12}=2C_2+2C_1+C_6$ and then assuming that there is a second conic tangent at six points to $C_6$, distinct from $C_2$. From this it appears how the difficulty can be overcome for the smallest values of $n$, but it is not apparent that this can be done for the larger ones.

In any case, if one has succeeded in seeing the curve $C_{3n}+2C_2$ as the limit form of a $C_{3n+4}$ which is the branch curve of an $(n+1)$-fold plane, it will still be necessary to add to it a bitangent line $C_1$, chosen so that the curve $C_{3n}+2C_2+2C_1$ is the limit form of the branch curve of an $(n+2)$-fold plane. But here (without any need to resort to considerations analogous to those which, in the case $n=4$, led us to the choice of the $C_2$ tangent at six points to $C_{12}$) the difficulty of choosing $C_1$ is overcome at once on the basis of this simple observation: that the $(n+1)$-fold plane defined by the branch curve $C_{3n+4}$ has genera one and therefore appears as representing an $F_{n+2}$ $(n=2\pi-2)$ with canonical sections of genus $\pi+1$, projected from one of its simple points; to this point there corresponds the exceptional line, bitangent to $C_{3n+4}$, which is the required $C_1$.

What we have said serves to explain Chisini’s constructive method and, in some measure, to test its significance by the examination of a concrete problem. As for this problem --- concerning the existence and the classification of the surfaces $F_n$ with genera one --- we shall give its direct solution in § 3 of Chap. VII, where one will also find a historical account of the course of these researches.

To the preceding we add some bibliographical indications. Whereas in the citations made we were referring to $n$-fold planes with $n>2$, one may list some special researches on the case of double planes: \textsc{F. Enriques}, “Sui piani doppi di genere 1”, Memorie della Società It. delle Scienze, detta dei XL, Rome, 1896; “Sui piani doppi di genere lineare $p^{(1)}=1$”. R. Accademia dei Lincei, 1898.
% scan: s114L
% !PRINTED 204
\weberpage{204}
\textsc{L. Campedelli}, “Sopra i piani doppi con tutti i generi uguali all’unità”. Rend. del Sem. mat. di Padova, vol. XI, 1940; “Le superficie con i generi uguali all’unità, rappresentabili in infiniti modi sul piano doppio”, Rend. del Sem. Mat. della R. Università di Roma, s. V, vol. I, 1940 and “La classificazione dei piani doppi con tutti i generi uguali all’unità”. Atti del $2^\circ$ Congresso dell’Unione Mat. Ital., Bologna, 1940; “Sui piani doppi con curva di diramazione dell’ottavo ordine”. R. Accademia dei Lincei, 1932; “Sui piani doppi con curva di diramazione del $10^\circ$ ordine”, ibidem, 1932.; “Sopra alcuni piani doppi notevoli con curva di diramazione del $10^\circ$ ordine”, ibidem, 1932. %

Somewhat akin to the case of double planes is the study of those particular quadruple planes which depend on Abelian equations, that is, on the extraction of two quadratic radicals: \textsc{P. Libois}, Rend. Acc. Lincei, 1934.

The following works also refer to noteworthy classes of $n$-fold multiple planes, with $n>2$: \textsc{B. Segre}, “Sulla caratterizzazione delle curve di diramazione dei piani multipli generali”, Reale Acc. d’Italia, vol. I, 1930. --- \textsc{G. Pompilj}, “Sulla rappresentazione algebrica dei piani tripli”, Rend. Sem. mat. della R. Università di Roma, 1939; “Osservazioni sui piani tripli”. R. Istituto Lombardo, 1941; “Sui piani multipli diedrici”. Rend. Sem. mat. della R. Università di Roma, 1939.

\edsection{11. The moduli of a class of algebraic surfaces.}{11. The moduli of a class of algebraic surfaces.}{sec.V.11}

The consideration of multiple planes with simple branch curve, on which an algebraic surface can in general be represented, together with the related existence theorem which we explained in the preceding section, makes it possible to say something regarding the number of \textit{moduli}, or parameters, on which there depends the determination of a \textit{class} of surfaces within a continuous family to which there correspond certain values of the characters $p$ and $p^{(1)}$. Let us assume for simplicity that we are dealing with regular surfaces $p=p_g=p_a$, having some plurigenus $P_i>0$, and therefore transformable so as to eliminate the exceptional curves, so that $p^{(1)}$ becomes for them the \textit{absolute linear genus}.

On a surface $F$ satisfying the aforesaid conditions, let us choose arbitrarily an irreducible linear system $|C|$ without base points, of genus $\pi$ and degree $n\leq2\pi-2$, which we shall suppose regular and non-special, hence of dimension
\[
r=p+n-\pi+1.
\]

A generic net contained in $|C|$ allows us to represent $F$ on a multiple plane of order $n$, endowed with a branch
% scan: s114R
% !PRINTED 205
\weberpage{205}
curve $D$ of order $2n+2\pi-2$, which will possess a certain number $d$ of nodes and a certain number $k$ of cusps, where these numbers are expressed in terms of $p$ and $p^{(1)}$ by means of the formulae of § 6. Now, if we let the curve $D$, of order $2n+2\pi-2$, vary within a continuous system of curves still having $d$ nodes and $k$ cusps, we have seen that this curve is the branch curve of an $n$-fold plane, to which there corresponds a surface having the same characters $p$ and $p^{(1)}$ as $F$. However, not all the curves of this system, of dimension
\[
s\geq(n+\pi-1)(2n+2\pi+1)-d-2k, %
\]
will correspond to surfaces birationally distinct from $F$; first of all, the $\infty^8$ homographic transforms of $D$ lead to the same $F$ and to the same net of curves $|C|$ on it; secondly, there are in the $\infty^r$ system $|C|$ infinitely many nets, each of which depends on $3r-6$ parameters. We shall therefore have at least
\[
(n+\pi-1)(2n+2\pi+1)-d-2k-8-(3r-6) %
\]
surfaces which must be regarded as birationally distinct from $F$. Indeed, if $\infty^1$ curves $D$ not corresponding to nets contained in $|C|$ give rise to birationally identical surfaces, this means that $F$ possesses a continuous series of non-equivalent nets of curves, in which a net of $C$ is contained, and this contradicts the hypothesis that the surface is regular $(p_a=p_g)$.

To tell the truth, with regard to this assertion one may raise the doubt that a net of $C$, say $|C|_2$, might present itself as the limit of a net of curves $|L|_2$ of higher order, endowed with base points, just as the net of the lines of the plane appears as the limit of a net of conics with three base points. But the doubt expressed is easily excluded for the surfaces $F$ to which we are referring, on which the degree of a linear system of curves cannot exceed twice the genus minus two, and which can therefore be regarded as free of exceptional curves. Indeed, if a net $|L|_2$ of effective genus $\pi$ and effective degree $n$ possesses some base point of order $i\geq1$, the $L$ will have with the canonical curves $K$ a virtual number of intersections $2\pi-2-n-\Sigma i^2$, and in the limit they will have to split into curves $C$ augmented by a fixed component $\theta$, which will have with the $K$ the negative number of intersections: $-\Sigma i^2$. This $\theta$ will therefore have a certain genus $\varrho$ and a degree
\[
\nu=2\varrho-2+\Sigma i^2>2\varrho-2; %
\]
and its presence on the surface $F$ contradicts the hypothesis made, according to which the surface has been transformed so as not to possess exceptional curves. Here, for the correctness of our count (the assertion that there will be $\infty^{3r-6+8}$ identical multiple planes), it is useful to add that our $F$ cannot admit a continuous series of birational transfor%
% scan: s115L
% !PRINTED 206
\unskip\weberpage{206}%
mations into itself; for it would follow either that the complete system $|C|$ would have to belong to a wider series of non-equivalent curves, which contradicts the hypothesis $p_a=p_g$, or that $|C|$ would be transformed into itself, and this contradicts the hypothesis that the surface $F$ does not belong to the family of ruled surfaces\origfootnote{1}{s115L-1}{Indeed, the surfaces which possess a continuous infinity of projective transformations are rational or ruled. Theorem of \textsc{Enriques-Fano}. Cf. \textsc{Enriques-Conforto}, \textit{Le superficie razionali}, Op. cit., chap. VI, § 37, p. 496.}.

\edFlush
We are now in a position to write down the number $M$ of moduli on which the class of the surfaces $F$ depends, or at least a lower bound for this number. We shall in fact have
\[
M\geq(n+\pi-1)(2n+2\pi+1)-d-2k-8-(3r-6),
\]
that is,
\[
M\geq(n+\pi-1)(2n+2\pi+1)-d-2k-3r-2.
\]

In this formula let us substitute for $d$ and $k$ their expressions in terms of $p$ and $p^{(1)}$ given by the formulae of § 6:
\[
\begin{aligned}
d&=2[(n+\pi)^2-5n-17\pi-2p^{(1)}+6p_a+24]\\
k&=3(n+6\pi+p^{(1)}-4p_a-11);
\end{aligned}
\]
moreover, let us substitute for $r$ its value
\[
r=p+n-\pi+1.
\]

We shall thus find %
\[
M\geq9p-2p^{(1)}+12,
\]
or, more precisely,
\[
M=9p-2p^{(1)}+12+\omega,
\]
where $\omega\geq0$ denotes the superabundance of the system of plane curves $D$; that is to say, in the continuous system of the curves $D$ of order $2n+2\pi-2$ endowed with $d+k$ double points, the conditions imposed in order that $k$ double points become cusps are not $k$ independent conditions, but depend only on $k-\omega$ conditions.

For the present we conclude:

\textit{Regular surfaces with some plurigenus $P_i>0$, of surface genus $p$ and absolute linear genus $p^{(1)}$, depend, in general, on}
\[
M=9p-2p^{(1)}+12+\omega
\]
\textit{moduli, with $\omega\geq0$.}

For surfaces with $p_a<p_g$, the preceding count must be modified in relation to the characteristic property of these sur%
% scan: s115R
% !PRINTED 207
\unskip\weberpage{207}%
faces which we shall establish later on, namely the property of containing series of $\infty^{p_g-p_a}$ non-equivalent linear systems. Hence for \textit{irregular surfaces the number of moduli will be}
\[
M=10p_a-p_g-2p^{(1)}+12+\omega
\]
with $\omega\geq0$\origfootnote{1}{s115R-1}{Save for the meaning of $\omega$, the formula holds a fortiori for surfaces which admit a continuous group of transformations into themselves.}.

Let us return to regular surfaces $(p_a=p_g=p)$. For surfaces with a canonical or bicanonical curve of order zero ($p=1$, $p^{(1)}=1$, and $p=0$, $p^{(1)}=1$) one will have respectively
\[
9p-2p^{(1)}+12=19
\]
and
\[
9p-2p^{(1)}+12=10,
\]
and we shall verify that in fact the number of moduli is precisely, for the former,
\[
M=19,
\]
and for the latter
\[
M=10;
\]
hence
\[
\omega=0.
\]

But in general one will find
\[
\omega>0.
\]

\edsection{12. Digression on the completeness of the characteristic series of a complete system of plane curves endowed with nodes and cusps.}{12. Digression on the completeness of the characteristic series of a complete system of plane curves endowed with nodes and cusps.}{sec.V.12}

In order to express the count of the moduli of a surface by a more significant formula, it is advisable to seek the meaning of the character $\omega$ which appears in the formula of the preceding section. It is easy to recognize that “the superabundance $\omega$ of a continuous system of plane curves $\{D\}$ of a certain order $n$, endowed with $d$ nodes and $k$ cusps, is less than or equal to the index of speciality of its characteristic series on a generic $D$, that is, of the series cut on it by the curves of the same order which pass simply through the $d$ nodes and the $k$ cusps and touch the cuspidal tangents at these points”.

Indeed, let us recall that, given a continuous series $\{D\}$ of plane curves $D$ of a certain order $n$, endowed with $d$ nodes and $k$ cusps, the curves of $\{D\}$ infinitely near to a given $D=D_0$ have with it 2 intersections absorbed at each node and 3 intersections absorbed at each cusp\origfootnote{2}{s115R-2}{Observation of De \textsc{Jonquières} and \textsc{Beck} (Math. Annalen, Bd. 14, 1878). Cf. \textsc{Enriques-Chisini}, Book II, chap. II (vol. I, p. 331, cf. vol. II, p. 286).}. This means that the characteristic series,
% scan: s116L
% !PRINTED 208
\weberpage{208}
cut on \(D_0\) by the infinitely near \(D\), is contained in the complete series cut on the said \(D_0\) by the curves of the same order passing through its \(d\) nodes and its \(k\) cusps and tangent at these points to the cuspidal tangents. Now, if the continuous system \(\{D\}\) has the superabundance \(\omega\), its dimension is not
\[
r = \frac{n(n+3)}{2} - d - 2k,
\]
but
\[
r + \omega = \frac{n(n+3)}{2} - d - 2k + \omega,
\]
and hence the dimension of the series \(g_m\), of order \(m = n - 2d - 3k\), cut on \(D_0\), of genus %
\[
\pi = \frac{n(n-3)}{2} + 1 - d - k,
\]
by the infinitely near \(D\) is \(\ge r + \omega = m - \pi + \omega\), that is, this characteristic series is special, with index of speciality at least equal to \(\omega\). At first sight one would be tempted to say that \(\omega\) is certainly equal to the said index of speciality, that is, that “the characteristic series of the complete system \(\{D\}\), consisting of the plane curves of order \(n\) with \(d\) nodes and \(k\) cusps, is complete”. Indeed, it is easy to verify that the following \textit{differential property} is satisfied: in the pencil determined by a \(D\) and by an adjoint curve \(\varphi_n\) of it which touches its cuspidal tangents at the \(k\) cusps, the curve infinitely near to \(D\) possesses \(d\) nodes and \(k\) cusps near to those of \(D\) itself.

But this differential property does not entail that each of the aforesaid curves with \(d\) nodes and \(k\) cusps, near to \(D\), belongs to a continuous series of curves endowed with the same number of nodes and cusps. Indeed, if we refer to the variety \(V\) which represents the elements (curves) of the system \(D\), this \(V\) will appear to us as the intersection of a variety \(W\) of the space \(S_R\) of \(R = \frac{n(n+3)}{2}\) dimensions, representing the system of all plane curves of order \(n\) endowed with \(d + k\) nodes, with \(k\) hypersurfaces or linear hypersurface sheets of the said space \(S_R\); and it might happen, e.g., that these \(k\) sheets have at every common point the same tangent space of \(R - (k - 1)\) dimensions, while intersecting in a variety of \(R - k\) (and not of \(R - k + 1\)) dimensions: in this case a point of \(V\) (representing a curve of \(\{D\}\)) would have not \(\infty^{r-1}\) but \(\infty^r\) neighboring points belonging to \(V\) (that is, images of curves of \(\{D\}\)).

% scan: s116R
% !PRINTED 209
\weberpage{209}
Although certain indications and the examples that we know tend to exclude this doubt, we have not succeeded in proving what we believe to be true, and we must confine ourselves to introducing a

\textit{Working hypothesis:} The complete continuous system of the plane curves \(D\) endowed with \(d\) nodes and \(k\) cusps has its characteristic series complete; therefore \textit{the superabundance} \(\omega\) of this system \(\{D\}\) \textit{equals the index of speciality} of the series itself, cut on a \(D\) by the infinitely near curves.

This hypothesis leads to some noteworthy consequences, and has a certain heuristic value in relation to the developments which follow.

\edsection{13. The moduli of regular surfaces of genus \(p > 3\).}{13. The moduli of regular surfaces of genus p > 3.}{sec.V.13}

We have seen how the number of moduli on which a class of algebraic surfaces depends (with some plurigenus greater than zero, that is, not belonging to the family of ruled surfaces) is expressed in general, when these surfaces are represented on multiple planes with branch curves \(D\) of a certain order \(n\), endowed with a certain number of nodes and cusps. The formula obtained in \S~11 contains, besides the characters, surface genera \(p_a\), \(p_g\) and absolute linear genus \(p^{(1)}\), also a non-negative numerical character \(\omega\), which is the superabundance of the continuous system of the \(D\), and — according to our \textit{working hypothesis} (\S~12) — the index of speciality of the characteristic series of \(\{D\}\). %

Now the meaning admitted for \(\omega\) is susceptible of another interpretation, and leads us to write a more significant formula as regards the number of moduli of regular surfaces of genus \(p > 3\).

For this purpose let us consider a surface \(F = F_n\) of ordinary space, of a certain order \(n\), endowed with ordinary singularities, and the representation of \(F\) on an \(n\)-fold plane which is obtained by projecting \(F\) itself from a generic external point, \(O\). In this representation the branch curve \(D\) arises by projection from the Jacobian curve \(C_j\) of the net of plane sections \(C\) by planes through \(O\). Now on the plane curve \(D\) the complete characteristic series of the continuous system \(\{D\}\), consisting of the curves of the same order with the same number of nodes and cusps, has precisely the index of speciality \(\omega\). It is a matter of finding this series, of speciality \(\omega\), again on the Jacobian curve \(C_j\). For this purpose let us observe that \(C_j\) is the section of \(F\) by the polar of the point \(O\), that is (denoting by \(K\) the canonical curves of \(F\)), it is a curve of the linear system \(|3C + K|\) cut on \(F\) by the surfaces \(\varphi_{n-1}\) of order \(n - 1\) adjoint to it which pass through the \(\tau\) cuspidal points of \(F\), whose group, belonging to the double curve of \(F\), will be denoted by \(G\). The other curves \(C_j\), passing through the \(\tau\) cuspidal points, cut on \(C_j\), outside them, a linear series which,
\printersignature{\textsc{Enriques F.} - \textit{Superficie algebriche.}\quad 14}
% scan: s117L
% !PRINTED 210
\weberpage{210}
as we shall see at once, is projected from \(O\) into the series cut on the branch curve \(D\) by its first polars. Indeed, if \(\overline{C}_j\) denotes the section of \(F\) by the polar surface of a point \(\overline{O}\), the tangent plane to \(F\) at a point \(P\) common to \(C_j\) and \(\overline{C}_j\) passes through the line \(O\overline{O}\): hence the tangent to the curve \(D\) at the point \(P'\), the projection of \(P\) from \(O\), passes through the trace of the line \(O\overline{O}\) on the multiple plane. %

It follows from this that the characteristic series of \(\{D\}\) on the curve \(D\), the projection of \(C_j\), arises by projection from the series which is cut on \(C_j\) by the curves of the linear system \(|L| = |4C + K|\) passing through \(G\) (or, in other words, by the adjoint surfaces \(\varphi_n\) passing through the \(\tau\) points of the said \(G\)). Hence the character \(\omega\), which denotes the index of speciality of the aforesaid characteristic series of \(\{D\}\), can also be interpreted as “the index of speciality of the series cut on the curve \(C_j\), which is a curve of \(|3C + K|\) through \(G\), by the curves of the linear system \(|L| = |4C + K|\) likewise passing through \(G\)”. %

This \textit{projective interpretation of the character \(\omega\)}, with respect to the surface \(F\), is equally valid for regular and irregular surfaces. And it leads to an evaluation of the number of moduli on which a class of surfaces \(F\) depends, within the widest continuous family of surfaces containing it, in a way which agrees very well with intuition.

Indeed, let us consider the continuous system of all the surfaces \(F\) having the same order \(n\) and a double curve with the same characters, on which there lies the same number \(\tau\) of cuspidal points: the surfaces of such a system \(\{F_n\}\), to which \(F\) belongs, determine on one of them, and in particular on the given \(F\), a \textit{characteristic linear system}, consisting of the intersections of \(F\) with the infinitely near surfaces of the system, and this linear system (which a priori one might suspect of not being complete) is cut on \(F\) by the adjoint surfaces \(\varphi_n\) of order \(n\) tangent to the cuspidal planes of \(F\) at the points of the group \(G\).

Our existence theorem, supplemented by the working hypothesis of \S~12, assures us that the continuous system \(F_n\) will cut on an \(F\) the \textit{complete characteristic linear system}; and, conversely, this presupposition \textit{is equivalent to the said working hypothesis}, and provides the means of computing the number of moduli on which the class of \(F\) depends within \(\{F_n\}\). Naturally, in order to have surfaces with \textit{general moduli}, it will be advisable to start from a projective model \(F\) whose existence within a continuous family of surfaces with the given characters can be assured a priori; for example from an \(F\) which corresponds to a regular linear system of degree \(n\) and genus \(\pi\) suitably high, for which the inequality \(n > \pi - 1\) is satisfied, or else — save in the case of small genus or of some analogous exception — from a canonical surface, having as plane sections \(\infty^3\) curves \(K\). Let us see what the count of the moduli leads to, referring to the two choices of
% scan: s117R
% !PRINTED 211
\weberpage{211}
\(F\) indicated above and supposing, for simplicity of discussion, that the surface is regular.

First of all, with the first choice one recovers the formula for the moduli of \S~11. Indeed, it is a matter of evaluating the dimension \(R\) of the characteristic system of \(\{F_n\}\), and then the number \(S\) which indicates on how many parameters the \(F_n\) of the said continuous system that are birationally identical to a given \(F\) depend; we shall then have
\[
M = R - S.
\]

To calculate \(R\), let us first evaluate the dimension of the system \(|L| = |4C + K|\) without base points, which we may legitimately suppose to cut on a \(C_j\) a complete and non-special series, and then let us subtract from this dimension \(\tau - \omega\), that is, the number of independent conditions which the group of the \(\tau\) cuspidal points of \(F\) imposes on the curves \(L\) required to contain it. The dimension of \(|L|\) will be
\[
p + 4\pi + 6n - 4,
\]
and hence the dimension of the characteristic system of \(\{F_n\}\) will be
\[
R - 1 = p + 4\pi + 6n - 4 - \tau + \omega,
\]
where
\[
\tau = 2n + 8\pi - 12p + 2p^{(1)} - 22,
\]
so that
\[
R = 13p + 4n - 4\pi - 2p^{(1)} + 19 + \omega.
\]

On the other hand, \(S\) is calculated by adding 15 (the number of parameters of a space homography) to the number of parameters which define an \(\infty^3\) linear system within the complete linear system \(|C|\) of dimension \(r\). One thus obtains
\[
S = 4r - 12 + 15 = 4r + 3
\]
that is,
\[
S = 4p + 4n - 4\pi + 7.
\]
Therefore one finds
\[
M = R - S = 9p - 2p^{(1)} + 12 + \omega,
\]
which is the formula given in \S~11.

But, taking into account the meaning of \(\omega\) (in accordance with our working hypothesis of \S~12), one can give a more significant formula, by showing that for \(p > 3\) one has in general
\[
\begin{aligned}
&\qquad\qquad \omega \ge p,\\
&M = 10p - 2p^{(1)} + 12 + \omega'\\
&\qquad\qquad\qquad\qquad\qquad (\omega' \ge 0).
\end{aligned}
\]

% scan: s118L
% !PRINTED 212
\weberpage{212}
For this purpose let us resume the preceding calculation, referring to a surface \(F_n\) of order \(n = p^{(1)} - 1\) whose plane sections are canonical curves \(K\). Here the following modifications are needed:

1) the curves \(L\) of \(|L| = |5K|\) cut on a \(C_j = K_j\) the complete canonical series; therefore the calculation of \(R\), if we again denote by \(\omega\) the index of speciality of the series cut on \(K_j\) by the curves of \(|5K|\) which pass through the group \(G\), will again give us %
\[
R - 1 = p + 4\pi + 6n - 4 - \tau + \omega
\]
from which it follows that the \(\tau\) cuspidal points impose not \(\tau - \omega\) but \(\tau - (\omega - 1)\) independent conditions on the \(L\) required to contain them:

2) the system \(|K|\), unlike the regular non-special system \(|C|\) considered before, has the dimension
\[
\begin{aligned}
&\qquad p - 1 = p + n - \pi.\\ %
&(n = p^{(1)} - 1,\qquad \pi = p^{(1)}),
\end{aligned}
\]
instead of
\[
p + n - \pi + 1;
\]
hence the number of surfaces \(F_n\) birationally identical to a given one turns out to be
\[
4p + 4n - 4\pi + 3 %
\]
instead of
\[
4p + 4n - 4\pi + 7.
\]

Therefore one will find that the number of moduli of \(F_n\) is given by
\[
M = 9p - 2p^{(1)} + 16 + \theta
\]
where \(\theta = \omega - 4\) denotes the index of speciality of the series cut on a \(K_j\) by the curves of \(|5K|\) passing through the group \(G\) of the cuspidal points which belong to the surface \(F\), the image of a linear system chosen within \(|K|\).

Now one can evaluate \(\omega\), or at least a lower bound for it, by observing that the group \(G\) of the cuspidal points is a group of points belonging to the curve \(K_j\), the Jacobian of a net of curves \(K\), and is the partial intersection with a curve \(L\) of the system %
\[
|L| = |5K|
\]
which is adjoint to
\[
|K_j| = |4K|.
\]

The group \(G_\tau\) will therefore be a special group, and will impose on the curves of \(|5K|\) required to contain it precisely \(\tau - \theta\) conditions, where \(\theta\) is the dimension of the complete series to which \(G_\tau\) belongs on \(K_j\) (Riemann-Roch theorem).

% scan: s118R
% !PRINTED 213
\weberpage{213}
But a lower bound for \(\theta\) is obtained by noting that the net having \(K_j\) as its Jacobian belongs to \(\infty^{p-4}\) linear systems \(\infty^3\) of curves \(K\) contained in the canonical system. Each of these \(\infty^3\) systems defines on \(K_j\) a group \(G_\tau\), which is the group of the points that are double for \(\infty^1\) curves of the system. Therefore it results that
\[
\theta \ge p - 4
\]
and finally
\[
M = 10p - 2p^{(1)} + 12 + \omega' \qquad (\omega' \ge 0).
\]

Let us state the result obtained, bearing in mind that there enters into it the working hypothesis of \S~12, of which a proof would need to be given:

\textit{The number of moduli on which there depends a class of regular surfaces of genus \(p > 3\), with irreducible simple canonical system \((p^{(1)} > 5)\), is probably}
\[
M \ge 10p - 2p^{(1)} + 12.
\]

In the following section we shall say how this formula extends to the case of irregular surfaces.

\edsection{14. Historical note and complements.}{14. Historical note and complements.}{sec.V.14}

The first attempt to determine the number \(M\) of moduli on which a class of algebraic surfaces depends is due to M. \textsc{Noether}\origfootnote{1}{s118R-1}{M. \textsc{Noether}, \textit{Anzahl der Moduln einer Classe algebraischer Flächen.} Sitzungsberichte der Akademie, Berlin, 1888.} (1888), who regards as valid, under broad conditions, the formula
\[
M = 10p - 2p^{(1)} + 12 \qquad (p = p_a = p_g),
\]
precisely the one which, according to our working hypothesis of \S~12, would give in general a lower bound for \(M\) for regular surfaces with \(p > 3\), corresponding to \(\omega' = 0\). \textsc{Noether} arrived at this sort of divination through an intuition of genius. He started from the formula which gives the postulation of a space curve \(C\) (endowed with triple points) for the surfaces \(F_n\) of order \(n\) which are required to pass doubly through it. If the number of conditions thus imposed on \(F_n\) is negative, \(-R\), he presumes that \(R\) indicates the number of conditions for the said curve \(C\) to be double for an \(F_n\); hence the number of moduli relative to the \(F_n\) with \(C\) double will be found by evaluating the number \(N\) of parameters on which \(C\) depends:
\[
M = N - R.
\]

% scan: s119L
% !PRINTED 214
\weberpage{214}
The question of the moduli of surfaces was taken up again by \textsc{Enriques} in two papers, of 1908 and 1912\origfootnote{1}{s119L-1}{F. \textsc{Enriques}, \textit{Sui moduli delle superficie algebriche.} Rend. Lincei, June 1908; \textit{Sui moduli di una classe di superficie algebriche e sul teorema d’esistenza per le funzioni algebriche di due variabili.} Att. R. Accademia di Torino, 1912. Cf. Annali di Mat., 1923-24.}. The evaluation is made by considering — in ordinary space — the continuous system of the surfaces \(F_n\) of order \(n\) endowed with a double curve with given characters, and by assuming that the characteristic system, cut on an \(F_n\) by the infinitely near \(F_n\), is complete: the author did not notice at the time that this presupposition implicitly contained a hypothesis which is not sufficiently proved by ascertaining the corresponding differential property; the gap was pointed out only by \textsc{Zariski} (1935)\origfootnote{2}{s119L-2}{O. \textsc{Zariski}, \textit{Algebraic Surfaces.} Berlin, 1935 (cf. V, 8).} after an analogous gap had been found by \textsc{Severi} (1921) in the proof of the fundamental theorem which expresses the characteristic property of irregular surfaces (cf. Chap. IX, \S~6).

In any case, \textsc{Enriques}, relying on the implicit hypothesis mentioned above, found first of all the formula of \S~11:
\[
M = 10p_a - p_g - 2p^{(1)} + 12 + \omega \qquad (\omega \ge 0)
\]
valid for all regular and irregular surfaces which do not belong to the family of ruled surfaces (that is, with some plurigenus \(P_i > 0\)); and secondly the more expressive formula which should hold, in general, for regular surfaces (with irreducible simple canonical system) whose genus is \(p > 3\) \((p^{(1)} \ge 6)\):
\[
M = 10p - 2p^{(1)} + 12 + \omega'; \qquad (\omega' \ge 0)
\]
this formula, as we have said, reduces for \(\omega = 0\) to that of Noether. %
In the second Note (of 1912) \textsc{Enriques} drew from the formula for the moduli some insight into the existence theorem for multiple planes, observing that the conditions for a plane curve endowed with a certain number of nodes and cusps to be the branch curve of a multiple plane are of a purely arithmetical nature, so that one can say that “if a curve \(D\) is the branch curve of an \(n\)-fold plane, the same must happen for all the curves of the widest continuous system which are endowed with the same number of nodes and cusps”. The author pointed out on that occasion that such arithmetical conditions could be determined by the study of the monodromy group, as he afterwards carried out in the memoir in the Annali di Matematica of 1923, and as has been expounded here in \S~9.

Now it is clear that, starting — as was done above in \S~11
% scan: s119R
% !PRINTED 215
\weberpage{215}
— from the existence theorem for multiple planes, one succeeds in giving the first formula for the moduli, which (the family of ruled surfaces being excluded) is absolutely general, and thus remains immune from Zariski’s criticism. That criticism affects instead the second, more expressive formula (for \(p > 3\)), which in this exposition has been made to depend on the working hypothesis of \S~12.

The question of the moduli of algebraic surfaces was taken up again by B. \textsc{Segre}\origfootnote{1}{s119R-1}{B. \textsc{Segre}, \textit{Sui moduli delle superficie algebriche irregolari.} Rendic. Acc. Lincei, vol. XIX (1934).} in 1934, especially with the aim of finding a more expressive formula for irregular surfaces. In this research, however, there remains the presupposition (of the completeness of the characteristic system of a continuous system of surfaces \(F_n\)) which we have reduced to the working hypothesis of \S~12; for \textsc{Segre} had not noticed the difficulty brought to light later by Zariski’s criticism.

Here we note that, on the basis of the said presupposition, the reasoning developed above for the general case of regular surfaces with \(p > 3\) extends at once to the case of irregular surfaces (bearing in mind the characteristic property of these surfaces, which has been mentioned). In this way one would arrive at the formula
\[
M = 12p_a - 2p_g - 2p^{(1)} + 12 + \omega \qquad (\omega \ge 0)
\]

But \textsc{Segre} follows another path, which leads him to another, more significant formula, so much so that already in the case of regular surfaces he is led to write \(\omega \ge p - 1\), and hence
\[
M \ge 10p - 2p^{(1)} + 11,
\]
a formula which (save for the restrictions mentioned above) should hold for all surfaces with some non-zero plurigenus, without any need to assume \(p > 3\), irreducible canonical curves, etc.

For that matter, this result falls within the more general formula which gives the moduli of irregular surfaces. Indeed, B. \textsc{Segre} finds \(\omega \ge 2p_g - p_a - 1\), and hence
\[
M = 9p_a + p_g - 2p^{(1)} + 11 + \theta %
\]

To close, we report a formula communicated to us by Prof. G. \textsc{Castelnuovo}, who has found the following inequality for the character \(\omega\):
\[
\omega \ge 2p_a + p_g + p^{(1)} - I - 3.
\]
% scan: s120L
% !PRINTED none
\weberblank{[s120L]}

\clearpage\thispagestyle{plain}
% scan: s120R
% !PRINTED none
\weberpage{[s120R]}
\edchapter{VI}{REGULAR SURFACES: MINIMUM OF THE GENERA\\AND CONDITIONS OF RATIONALITY}{Chapter VI. REGULAR SURFACES: MINIMUM OF THE GENERA AND CONDITIONS OF RATIONALITY}{REGULAR SURFACES: MINIMUM OF THE GENERA AND CONDITIONS OF RATIONALITY}

In this chapter we wish, for simplicity, to consider regular surfaces of genus
\[
p_g = p_a = p,
\]
although many results which refer to them will extend without difficulty to irregular surfaces, as we shall often have occasion to point out.

\edsection{1. Lower bound for the linear genus.}{1. Lower bound for the linear genus.}{sec.VI.1}

Let us consider a surface \(F\) of genus \(p > 0\), which we may suppose free of exceptional curves (Chap. IV, \S~9). By hypothesis \(F\) possesses canonical curves, which we shall at first suppose to be of order greater than zero. We wish to prove that the linear genus of \(F\) (which is here the absolute linear genus) is
\[
p^{(1)} > 0.
\]

We must recognize that the hypothesis \(p^{(1)} \le 0\) leads to an absurdity. Let then \(p^{(1)} \le 0\), and hence the degree of the canonical system \(|K|\) be
\[
p^{(2)} = p^{(1)} - 1 < 0.
\]

We shall thus have (at least) one canonical curve \(K\) of negative degree, which may consist of several simple or multiple parts:
\[
K = s_1K_1 + s_2K_2 + \mathinner{\ldotp\ldotp\ldotp\ldotp} + s_rK_r,
\]
where the \(K_i\) cannot be curves of virtual genus \(\varrho_i = 0\) and degree \(\nu_i = -1\), which would be exceptional curves of the first kind on \(F\). Let us intersect one of the components of \(K\), for example \(K_1\), with \(K\). The number of intersections will be %
\[
2\varrho_1 - 2 - \nu_1. %
\]

% scan: s121L
% !PRINTED 218
\weberpage{218}
This number cannot be negative. Indeed, if
\[
2\varrho_1 - 2 - \nu_1 < 0,
\]
then, since \(\varrho_1 \geq 0\), it follows that
\[
\nu_1 \geq -1
\]
and hence (the values \(\nu_1 = -1\), \(\varrho_1 = 0\) being excluded),
\[
\nu_1 \geq 0.
\]

But this is absurd, because a component \(K_1\) of degree \(\nu_1 \geq 0\) cannot have a negative number of intersections with \(K = s_1K_1 + s_2K_2 + \ldots\). Now, having proved that the components \(K_i\) of \(K\) must have a non-negative number of intersections with \(K\), one easily deduces that \(K\) itself must have a degree \(p^{(2)} \geq 0\). %

Indeed, denoting by \(n_{ij}\) the number of intersections of two components \(K_i\) and \(K_j\) of \(K\), the degree of \(K\) will be expressed by
\[
p^{(2)} = \Sigma s_i^2\nu_i + 2\Sigma s_i s_j n_{ij}.
\]

The expression for \(p^{(2)}\) now appears to us as a quadratic form in the \(s_i\), and, bearing in mind Euler’s theorem on homogeneous functions, one recognizes that it must be essentially non-negative, because its partial derivatives, which represent twice the number of intersections \((K \cdot K_j)\), are
\[
\frac{\partial p^{(2)}}{\partial s_i} = 2(s_i\nu_i + \Sigma s_j n_{ij}) \geq 0.
\]

The inequality \(p^{(1)} > 0\) still holds in the case where the surface \(F\) possesses a single canonical curve of order zero \((p = 1)\). Indeed, let us take on \(F\), which is free of exceptional curves, an irreducible linear system \(|C|\) without base points, of genus \(\pi\) and degree \(n\); in this case we shall have
\[
n = 2\pi - 2,
\]
the number of intersections being \((C \cdot K) = 2\pi - 2 - n = 0\). And, denoting by \(|C'|\) the adjoint system of \(|C|\), of genus \(\pi'\) and degree \(n'\), we shall have
\[
|C'| = |C|.
\]

Now, evaluating the genus \(\pi' = \pi\) of \(|C'|\), one finds
\[
\pi' = \pi = \pi + p^{(1)} + 2\pi - 2 - n - 1
\]
and hence
\[
p^{(1)} = 1.
\]

The theorem proved \((p^{(1)} \geq 1)\) holds also for surfaces of genus \(p = 0\) which have a plurigenus \(P_i > 0\). Indeed, one
% scan: s121R
% !PRINTED 219
\weberpage{219}
can repeat the reasoning, putting in place of the canonical curve \(K\) the \(i\)-canonical curve \(iK\), which, if \(p^{(1)} \leq 0\), would have to have with each of its components (of genus \(\varrho\) and degree \(\nu\)) a non-negative number of intersections, expressed by
\[
i(2\varrho - 2 - \nu).
\]

Also in the case \((P_i = 1)\) in which the \(i\)-canonical curve is of order zero, and hence a system \(|C|\) coincides with its \(i\)-th adjoint, one concludes, as in the case \(p = 1\), that
\[
p^{(1)} = 1.
\]

Finally, we point out that the proof of our theorem is entirely independent of the hypothesis that \(F\) is regular. Let us therefore state, without restrictions, that:

\textit{A surface of genus \(p_g > 0\), or also of genus \(p_g = 0\) having a plurigenus \(P_i > 0\), has absolute linear genus}
\[
p^{(1)} \geq 1.
\]

We wish to offer a second proof of the theorem stated above, which depends on a lemma that will receive important applications in what follows. Let there be given a surface \(F\), which we wish to suppose free of exceptional curves, of linear genus \(p^{(1)} \leq 0\), and let us take on it an irreducible regular linear system \(|C|\), of genus \(\pi\) and degree \(n \leq 2\pi - 2\). The \(C\) meet a virtual canonical curve \(K\) in
\[
2\pi - 2 - n \geq 0
\]
points. The curves \(C'\) of the adjoint system \((|C'| = |C + K|)\), on the other hand, meet the \(K\) in
\[
2\pi - 2 - n + p^{(1)} - 1 < 2\pi - 2 - n
\]
points. Thus, in passing from a linear system to its adjoint, the number of intersections of the curves of the system with the canonical curve keeps decreasing. But this entails that the process of adjunction itself must come to an end, or else that one arrives at linear systems of genus \(\pi\) and degree \(N > 2\pi - 2\); but the existence of such a system \(|L|\) again entails of necessity that the series of the adjoints \(|L'|\), \(|L''|\), etc. must come to an end, because the number of intersections \((L \cdot L) > (L' \cdot L) > (L'' \cdot L)\ldots\).

Hence the hypothesis \(p^{(1)} \leq 0\) entails that, starting from a linear system \(|C|\), one obtains only a limited series of successive adjoints, and therefore that a system \(|C|\) cannot be contained in any of its adjoint systems \(|C^{(i)}|\), that is, that all the plurigenera of the surface are \(P_i = 0\).

% scan: s122L
% !PRINTED 220
\weberpage{220}
Here two observations are needed. The first is that the hypothesis of a surface \(F\) free of exceptional curves is introduced only for the purpose of reductio ad absurdum; if this hypothesis is rejected, one must admit straightaway that the surface \(F\) possesses a system with \(n > 2\pi - 2\), starting from which the process of adjunction dies out.

The second observation refers to the possible reducibility of the system \(|C'|\) adjoint to the irreducible system \(|C|\) from which we started. If the surface \(F\) is regular and \(|C|\) is free of fundamental curves, it is clear that \(|C'|\) cannot be reducible, and so in particular it cannot reduce to the system \(|C|\) plus a fixed canonical curve. But, independently of the irreducibility of the adjoint systems successive to \(|C|\), it really suffices to arrive at a system \(|L|\), reducible in any way whatever, of genus \(\pi\) and degree \(N > 2\pi - 2\); every multiple of this system will always give a linear system which satisfies a fortiori the same inequality. It will therefore be easy to construct an irreducible system which likewise satisfies an analogous inequality, by adding \(|sL|\) to a suitable irreducible linear system, a point on which it is not worth dwelling further here.

Meanwhile we conclude by affirming the result stated above.

\edsection{2. The bigenus and the plurigenera.}{2. The bigenus and the plurigenera.}{sec.VI.2}

Let us refer to a surface \(F\) such that for every linear system of genus \(\pi\) the degree is \(n \leq 2\pi - 2\), a surface which, as we have seen (Chap. IV, \S~9), can be transformed so as to eliminate its exceptional curves. Starting from an irreducible linear system \(|C|\), which may be the system of hyperplane sections of an \(F\) without singularities, let us try to construct the bicanonical system of \(F\). To this end one must detach a \(C\) from the second adjoint system \(|C''|\). Now, if one supposes the difference to be
\[
d = 2\pi - 2 - n \geq 0,
\]
the \(C''\) will cut on \(C\) a series of order \(2\pi - 2 + d\) which is certainly non-special, and therefore the detachment of \(C\) from \(|C''|\) involves at most \(\pi - 1 + d\) conditions. Hence, calculating the dimension of \(|C''|\) by the Riemann-Roch theorem, one finds that the dimension of the bicanonical system \(|C'' - C|\) will be expressed by %
\[
P - 1 \geq p_a + p^{(1)} - 1;
\]
we shall therefore have
\[
P \geq p_a + p^{(1)}.
\]

This result is independent of the effective existence of the canonical curves (of order \(d \geq 0\)), that is, it holds for \(p_g \geq 0\). In particular \textit{for regular surfaces \((p = p_a = p_g)\), free of ex}%
% scan: s122R
% !PRINTED 221
\weberpage{221}%
\textit{ceptional curves, with virtual canonical curves of order \(d\) \((\geq 0)\), the bigenus is always}
\weq{1)}{P \geq p + p^{(1)}.}

If there are irreducible canonical curves of genus \(p^{(1)} > 1\), on which the bicanonical system then cuts the complete canonical series, one has exactly:
\weq{2)}{P = p + p^{(1)}.}

In place of the equality 2) one has \textit{in general} the inequality
\weq{3)}{P \geq p + 1 \qquad \text{for} \qquad p^{(1)} = 1,}
since the \textit{canonical system}, of degree \(p^{(1)} - 1 = 0\), then splits into the \textit{elliptic curves \(k\) of a pencil} (linear on the regular \(F\)): if a canonical curve is composed of \(p - 1\) curves \(k\) of such a pencil, the bicanonical curves are formed by \(2p - 2\) curves \(k\) and therefore give a linear system of dimension
\[
P - 1 = 2p - 2;
\]
if some fixed curve is added to the \(p - 1\) curves \(k\), twice this curve may constitute a curve \(k\) capable of varying in their pencil; and therefore it may turn out that
\[
P - 1 > 2p - 2.
\]

The preceding results undergo some modification in the case \(p = 1\), when there is \textit{an effective canonical curve of order \(d = 0\)}: one then has, in place of 1):
\[
P = 1
\]
and
\[
p^{(1)} = 1 \qquad (P < p + p^{(1)}).
\]

Indeed, if \(d = 0\), the system \(|C|\) considered on the surface \(F\) is of degree
\[
n = 2\pi - 2,
\]
and, since \(p = 1\), one has
\[
|C| = |C'| = |C''|;
\]
hence the bicanonical system coincides with the \textit{system of order zero} which constitutes the canonical system and which has the dimension
\[
p - 1 = P - 1 = 0;
\]
the genus of this system \(|K|\) is
\[
p^{(1)} = 1,
\]
% scan: s123L
% !PRINTED 222
\weberpage{222}
because, by adding a \(K\) of order zero to a \(C\), one has a composite curve of genus
\[
\pi = \pi + p^{(1)} + 0 - 1.
\]

Thus \textit{the regular surfaces of genus \(p = 1\) which possess a canonical curve of order zero are distinguished from those with a canonical curve of order greater than zero by having: the former} %
\[
P = 1,
\]
\textit{and the latter}
\[
P \geq p + p^{(1)} \geq 2.
\]

For \textit{a regular surface without exceptional curves \((n \leq 2\pi - 2)\), of genus \(p = 0\), with virtual canonical curve of order \(d = 0\) \((d = 2\pi - 2 - n)\), there still holds, as for \(d > 0\),} the \textit{inequality} 1), \textit{that is} %
\[
P \geq p^{(1)},
\]
\textit{since one has}
\[
P = 1,\qquad p^{(1)} = 1.
\]

Indeed, \(|C'|\) being different from \(|C|\), \(|C''|\) will also be distinct from \(|C'|\) and will cut on a \(C\) a non-special series \(g_{2\pi-2}\); consequently the detachment of \(C\) from \(|C''|\) will involve at most \(\pi - 1\) linear conditions: one deduces that it involves exactly \(\pi - 1\) conditions, with the result that
\[
|C''| = |C|,\qquad P = 1,\qquad p^{(1)} = 1.
\]

\textit{Remark.} -- It is important to observe that the preceding results, in particular the fundamental \textit{inequality}
\weq{1)}{P \geq p + p^{(1)},}
absolutely \textit{fail} for surfaces containing linear systems \(|C|\) of genus \(\pi\) and degree \(n > 2\pi - 2\), that is, having virtual canonical curves of negative order
\[
-d = n - (2\pi - 2).
\]

Indeed, the system \(|C''|\) will now cut on a \(C\) a series of order \(2\pi - 2 - d\) which may be special.

Let the case of the plane serve as an example, in which one takes the linear system \(|C_7|\) formed by the general curves of the \(7^\circ\) order. The adjoint curves are the quartics \(C_4\), and the second adjoints are the lines \(C_1\), cutting on a \(C_7\) a series \(g_7^2\), a special series with index of speciality
\[
i = 10.
\]

Here one can calculate the genus (linear genus \(p^{(1)}\)) of the virtual canonical system
\[
|K| = |C_1 - C_4|.
\]
% scan: s123R
% !PRINTED 223
\weberpage{223}
Since the \(K\) cut the \(C_1\) in \(-d = -3\) and the \(C_4\) (of genus 3) in \(-12\) points, one will have
\[
0 = p^{(1)} + 3 - 12 - 1
\]
and hence we find again that the \textit{(relative) linear genus of the plane is:}
\[
p^{(1)} = 10.
\]

Meanwhile, since it is not possible to detach \(|C|\) from \(|C_1|\), the bicanonical system is entirely lacking, and therefore
\[
P = 0.
\]

The formulae which express the bigenus by means of the surface genus and the linear genus extend to the plurigenera. When the surface \(F\), free of exceptional curves, possesses \textit{virtual canonical curves of order \(d > 0\)} \((d = 2\pi - 2 - n)\), \textit{one has}
\[
P_i \geq p_a + \frac{i(i - 1)}{2}(p^{(1)} - 1) + 1
\]
and, for a regular surface \((p_a = p_g = p)\) of \textit{linear genus \(p^{(1)} > 1\), with irreducible \((i - 1)\)-canonical curves:}
\[
P_i = p + \frac{i(i - 1)}{2}(p^{(1)} - 1) + 1.
\]

\textit{The equality becomes an inequality in the case \(p^{(1)} = 1\)}, where the \textit{pluricanonical curves}, of degree \(i^2(p^{(1)} - 1) = 0\), are all composed of the elliptic curves of one and the same pencil:
\[
P_i \geq p + 1.
\]

Exceptions to the fundamental inequality
\[
P_i \geq p + \frac{i(i - 1)}{2}(p^{(1)} - 1) + 1
\]
are the surfaces of bigenus \(P = P_2 = 1\) which possess a virtual canonical curve of order \(d = 0\) and an effective bicanonical curve, also of order zero. In this case the series cut on a \(C\) by the third adjoint \(|C'''| = |C'|\) is the canonical series (of dimension \(\pi - 1\)); hence its detachment from \(|C'''|\), involving \(\pi\) conditions, is not possible, so that
\[
P_3 = 0.
\]

Hence \textit{the regular surfaces of genus zero and bigenus one with bicanonical curve of order zero are distinguished from those with bicanonical curve of order greater than zero by having:}

\textit{the former, the trigenus}
\[
P_3 = 0,
\]
% scan: s124L
% !PRINTED 224
\weberpage{224}
\textit{while the latter have}
\[
P_3 \geq 1.
\]

For the former surfaces one moreover has
\[
P_{2i} = 1,\qquad P_{2i+1} = 0,
\]
while for the latter all the plurigenera are essentially positive. What is more, they come to exceed unity even in the case \(p^{(1)} = 1\). Indeed, by combining a bicanonical curve counted three times with a tricanonical curve counted twice, one obtains a \textit{pencil of 6-canonical curves}, so that
\[
P_6 \geq 2.
\]

\edsection{3. Examples: double planes of linear genus \(p^{(1)} = 1\).}{3. Examples: double planes of linear genus p\textasciicircum{}(1) = 1.}{sec.VI.3}

An interesting exemplification of the formulae of the preceding section is offered by the double planes of linear genus \(p^{(1)} = 1\)\origfootnote{1}{s124L-1}{Cf. \textsc{F. Enriques}, \textit{Sui piani doppi di genere lineare \(p^{(1)} = 1\).} Rendic. R. Acc. Lincei, April and May 1898.}. It is proved there that, with the exception of the double plane with an entirely general branch sextic
\[
z^2 = D_6(xy)
\]
which has a canonical curve of order zero (genus and plurigenera \(p = P_i = 1\)), the other double planes, for which \(p^{(1)} = 1\) and some one of the plurigenera is \(P_i > 1\), contain a pencil (rational or not) of elliptic curves \(L\), transformed into itself by the involution \(I\) whose pairs correspond to the points of the plane. And hence one is led to two families of surfaces representable on such double planes:

I) surfaces with a linear pencil of (irreducible) elliptic curves \(L\) transformed into themselves by the involution \(I\), which induces on them a linear involution \(g_2^1\); these surfaces are represented on double planes with branch curve \(D_{2n}\) of order \(2n\) \((n > 3)\) endowed with a \((2n - 4)\)-fold or \((2n - 3)\)-fold point \(O\), and possibly with other singularities lowering the genus: quadruple points and points \([3, 3]\); the lines through \(O\) are the images on the double plane of the curves \(L\);

II) surfaces with a pencil of irreducible elliptic curves \(L\), which may be curves of a rational pencil, transformed into themselves by the involution \(I\), which induces on them an elliptic involution \(\gamma_2^1\), or else curves of a pencil (rational or not) associated in pairs by \(I\) in a rational involution of the pencil; these surfaces are represented on double planes in such a way that to the \(L\) there correspond the
% scan: s124R
% !PRINTED 225
\weberpage{225}
elliptic curves of a Halphen pencil\origfootnote{1}{s124R-1}{Cf. e.g., \textsc{Enriques-Conforto}, \textit{Le superficie razionali.} (Zanichelli - Bologna).} (curves \(C_{3s}\) with 9 \(s\)-fold points), and therefore their branch curves (supposed free of exceptional curves) are composed of the curves of the said pencil and possibly also of the cubic passing through its nine base points.

Among the double planes of category I, noteworthy are those with branch curve \(D_{2n}\) endowed with pairs of infinitely near triple points (points \([3, 3]\)) collinear with the \((2n - 4)\)-fold point \(O\).

For \(n = 4\) we have the double plane with branch curve \(D_8\) endowed with a 4-fold point \(O\) and with two infinitely near triple points \(A\) and \(A_1\), collinear with \(O\); the \(\mathrm{D}_8\) is composed of a \(D_7\) with \(O^3A^2A_1^2\) and of the line \(r = OA\). The genus of this double plane is %
\[
p_g = p_a = p = 1,
\]
there being \textit{a canonical curve} represented by the line \(r = OA\) (subject to passing simply through \(O\) and \(A\)). There is some difficulty in determining the linear genus \(p^{(1)}\), since to the double line \(r = OA\) there seems to correspond only a pair of coincident disconnected rational curves. The difficulty is resolved by recognizing that the surface \(F\) represented by projection on the double plane possesses an impure canonical curve, consisting of two infinitely near exceptional curves whose points correspond to the generic points of \(r\), and of an infinitesimal curve represented by the neighborhood of the point \(A_1\); which, under a transformation of \(F\) into a surface without singularities, changes into a curve of order greater than zero and of genus \(p^{(1)} = 1\). In accordance with this finding, \(F\) will have the bigenus
\[
P = p + p^{(1)} = 2,
\]
there being a pencil of elliptic bicanonical curves represented on the double plane by the lines through \(O\): on adding the exceptional curves, one will have as images of the impure bicanonical curves the pairs of lines formed by \(r\) and by a variable line through \(O\). What has been said becomes clear on considering the \(\infty^3\) linear system of the conics through \(AA_1\), which leads to transforming \(F\) into a double quadric cone without exceptional curves: the images of these double conics, with 10 branch points, are hyperelliptic curves of genus 4, forming a linear system of degree 4 \((< 2 \cdot 4 - 2)\); these are cut in 2 (conjugate) points by the canonical curve which corresponds to the neighborhood of \(A_1\), and in 4 points, consisting of two pairs of points conjugate in the \(g_2^1\), by the bicanonical curves; therefore the latter are
\printersignature{\textsc{Enriques F.} - \textit{Superficie algebriche.}\quad 15}

% scan: s125L
% !PRINTED 226
\weberpage{226}
represented on the cone by the plane sections of a pencil, and on the plane by the lines through \(O\). Let us state precisely that the said double quadric cone will possess a branch curve \(D_{10}\) of the \(10^\circ\) order, of which the \(D_7\) can be regarded as the projection from a triple point \(P\) to which another triple point \(P_1\), which projects into \(O\), is infinitely near. The line \(PP_1\) lies in the tangent plane to the cone along the generator \(k\) through \(P\), and the planes through it cut out on the cone the images of the bicanonical curves. %

Alongside the double plane mentioned above, which has a branch curve \(D_8 = rD_7\) with \(O^4A^3A_1^3\), one may consider the double plane with branch curve \(D_8\) passing through \(O^4\) and through two pairs of infinitely near triple points collinear with \(O\): \(A^3A_1^3\) and \(B^3B_1^3\). This \(D_8\) splits into the two lines \(r = OA\) and \(s = OB\) and a \(D_6(O^2A^2A_1^2B^2B_1^2)\). The genus of such a double plane is
\[
p = 0,
\]
and the bigenus
\[
P = 1,
\]
there being an impure bicanonical curve represented by the pair of lines \(r\) and \(s\). This curve reduces to a pair of exceptional curves counted twice. On evaluating the trigenus of the double plane one would find \(P_3 = 0\). This follows indirectly from the observations below.

If the double plane is transformed into a surface \(F\) without exceptional curves, it appears that the bicanonical curve of \(F\) is of order zero: one can construct \(F\) by means of the \(\infty^4\) linear system of the cubics through \(OAA_1BB_1\), which are images of curves of genus 4 and degree 6; in this way one is led to a double cubic surface (normal in \(S_4\)) with branch curve of the \(10^\circ\) order.

We shall see later that \(F\) can also be transformed into the surface \(F_6\) of the \(6^\circ\) order which passes doubly through the edges of a tetrahedron.

Let us now consider the double planes (still of category I) with branch curve \(D_{10}\) of the \(10^\circ\) order, endowed with a 6-fold point \(O\) and with two pairs of infinitely near triple points \(AA_1\) and \(BB_1\) collinear with \(O\): the \(D_{10}\) is composed of the two lines \(r = OA\) and \(s = OB\) and of a \(D_8(O^4A^2A_1^2B^2B_1^2)\). Here one has an impure canonical curve represented by the two exceptional lines \(r\) and \(s\), and by a residual curve (which becomes of order \(> 0\) on a transformed surface without singularities) formed of two elliptic components which correspond to the neighborhoods of the points \(A_1\) and \(B_1\), so that
\[
p_g = p_a = p = 1,\qquad p^{(1)} = 1.
\]

% scan: s125R
% !PRINTED 227
\weberpage{227}
The bigenus \(P\) must exceed unity, and indeed one has a net of bicanonical curves which have as images on the double plane the pairs of lines through \(O\); one therefore concludes that
\[
P = 3 \qquad (> p + p^{(1)}).
\]

Alongside the double plane just mentioned, consider the one with branch curve \(D_{10}\) endowed with a 6-fold point \(O\) and with 3 pairs of infinitely near 3-fold points \(AA_1\), \(BB_1\), \(EE_1\), collinear with \(O\). The surface \(F\) so defined possesses no canonical curve, there being no conic passing doubly through \(O\) and simply through \(A\), \(B\), \(E\); hence %
\[
p_a = p_g = p = 0.
\]

Instead one has elliptic bicanonical curves represented by the lines through \(O\), to which are added the exceptional images of the branch lines
\[
r = OA,\qquad s = OB,\qquad t = OE;
\]
therefore
\[
p^{(1)} = 1
\]
and
\[
P = 2 \qquad (> p^{(1)}).
\]

Among these bicanonical curves there are three double curves (corresponding to the three lines \(r\), \(s\), \(t\), or better to the neighborhoods of \(A_1\), \(B_1\) and \(E_1\)), which, taken together simply, constitute a tricanonical curve of the surface \(F\):
\[
P_3 = 1 \qquad (= p^{(1)}).
\]

The double of the said tricanonical curve belongs to the triple of the bicanonical system; one thus obtains the linear system of the sexticanonical curves of \(F\), which have as images the triples of lines through \(O\). We therefore have:
\[
P_6 = 4.
\]

In category II, too, we find examples of interest for our theory. Let us take a double plane with branch curve \(D_{6n}\) \((n > 1)\) consisting of an irreducible elliptic curve with 9 \(2n\)-fold points, belonging --- as is well known --- to a Halphen pencil\origfootnote{1}{s125R-1}{Every (double) curve of the Halphen pencil represents two distinct curves \(L\) not belonging to the involution \(I\), since this happens for the \(D_{6n}\); moreover this agrees with the criterion of \textsc{Comessatti}, since the \(D_{6n}\) can be reduced by continuity to a double \(\varphi_3^n\). Cf. e.g. \textsc{Enriques-Chisini}, Book V, chap. IV, \S~38 (vol. III, p. 445).}. The surface \(F\) so defined possesses a canonical curve having as %
% scan: s126L
% !PRINTED 228
\weberpage{228}
image a curve of order \(3n-3\) subject to passing \(n-1\) times through the 9 base points of the pencil, which reduces to the cubic \(\varphi_3\) through the 9 points counted \(n-1\) times; let us say to the \(\varphi_3^{n-1}\). We therefore have
\[
p_a=p_g=p=1,\qquad p^{(1)}=1.
\]

Now the branch curve \(D_{6n}\), taken together with the \(\varphi_3^{n-2}\) adjoint to the \(\varphi_3^{n-1}\) (to be counted twice), will form a bicanonical curve of \(F\) (cf. Chap. V, § 7), while another bicanonical curve is given by the \(\varphi_3^{2n-2}\) to be counted twice. One deduces that the curves of the Halphen pencil defined by the \(D_{6n}\), taken together with the fixed \(\varphi_3^{n-2}\), represent a pencil of bicanonical curves, consisting therefore of irreducible elliptic curves \(L\)\origfootnote{1}{s126L-1}{Every curve \(L\) of this pencil is the image of \textit{two} distinct elliptic curves, in accordance with the criterion of \textsc{Comessatti}, because the group of apparent critical points cut out on \(L\) by the \(D_{6n}\) is equally cut out by the \(D_3^{2n}\), double of the \(D_3^n\). Cf. \textsc{Enriques-Chisini}, l. c.} (\textit{not} belonging to the involution \(I\)) and of a fixed part, which is an elliptic curve of multiplicity \(n-2\), whose component belongs to the involution \(I\) and is equivalent to an \((n-1)\)-th part of \(L\). Thus we shall have
\[
P=2\qquad (=p+p^{(1)}).
\]

One then easily finds that the tricanonical curves are given by the pairs of curves \(L\), so that
\[
P_3=3.
\]

Alongside the preceding one we further find in category II the double plane with branch curve \(D_{6n}=D_{6n-3}+D_3\), composed of an irreducible Halphen curve \(D_{6n-3}\), with 9 \((2n-1)\)-fold points, and of the cubic \(D_3\) which passes through these 9 points.

The surface \(F\) so defined possesses \textit{one} canonical curve, represented by the \(D_3\) counted \(n-1\) times, let us say \(D_3^{n-1}\), so that
\[
p_a=p_g=p=1,\qquad p^{(1)}=1.
\]

The bicanonical curves transformed into themselves by the involution \(I\) are given: by the double \(D_3^{n-1}\), counted twice, and by the branch curve \(D_{6n-3}+D_3\) taken together with the adjoint of the image of the canonical curve, which is the double \(D_3^{n-2}\). Hence one has a pencil of elliptic bicanonical curves represented by the curves of order \(6n-3\) of the Halphen pencil, which are simple\origfootnote{2}{s126L-2}{See preceding note.}, taken together with the multiple curve represented by the double \(D_3^{n-2}\):
\[
P=2.
\]

% scan: s126R
% !PRINTED 229
\weberpage{229}
We shall not dwell on the double planes of category II having a branch curve composed in any way of the curves of a Halphen pencil: these lead in general to irregular surfaces with
\[
p_a=1,\qquad p_g>1,\qquad p^{(1)}=1. %
\]

But we wish to point out a double plane with \(p^{(1)}=1,\ p=0,\ P=1\), endowed with a bicanonical curve of order greater than zero, which possesses a pencil of elliptic curves (not self-conjugate in \(I\)) represented by the curves of a Halphen pencil, and which nevertheless falls outside category II because of the presence of a line forming part of the branch curve, the image of an exceptional curve. It is a double plane which had already presented itself to \textsc{Enriques}; it has a branch curve \(D_{10}\) of order 10 and is met with in the classification of these planes by L. \textsc{Campedelli}\origfootnote{1}{s126R-1}{L. \textsc{Campedelli}, \textit{Sui piani doppi con curva di diramazione del 10\textsuperscript{o} ordine} and \textit{Sopra alcuni piani doppi notevoli con curva di diramazione del 10\textsuperscript{o} ordine}. R. Acc. Lincei, 1932.}: the \(D_{10}\) is composed of a \(D_9+r\), where the \(D_9\) is an (irreducible) elliptic Halphen curve with 9 triple points, six of which, \(AA_1,\ BB_1,\ EE_1\), are infinitely near in pairs, while the other three (distinct) lie on the line \(r\).

The surface \(F\) represented on this double plane possesses no canonical curves, since a conic passing through the three triple points of \(D_9\) on \(r\) and through the three points \(A,\ B,\) and \(E\) would split into the \textit{exceptional curve} which corresponds to the line \(r\)\origfootnote{2}{s126R-2}{It is easy to see that the curve corresponding to this branch line of genus 0 and degree \(-2\) has on \(F\) genus 0 and degree \(-1\).} and a line through \(A,\ B\) and \(E\), which instead are not collinear.

We therefore have
\[
p=0.
\]

On the other hand, a double line through \(A,\ B\) and \(E\), if it existed, would have 4 branch points (three at the said three points and one on \(r\)); hence the absolute linear genus of \(F\) (the genus of the canonical curves deprived of the exceptional curve) is
\[
p^{(1)}=1.
\]

Now on \(F\) there is \textit{one} irreducible bicanonical curve\origfootnote{3}{s126R-3}{The irreducibility of the double cubic \(D_3\) follows from the criterion of \textsc{Comessatti}, bearing in mind the construction of the Halphen pencil to which the \(D_9\) belongs.}, which is represented on the double plane by the cubic \(D_3\) passing through the 9 triple points of \(D_9\), to which is added the exceptional curve represented by the line \(r\). Therefore the bigenus of \(F\) is
\[
P=1.
\]

% scan: s127L
% !PRINTED 230
\weberpage{230}
In accordance with our theorem one must have the trigenus \(P_3\ge1\) and the sextigenus \(P_6\ge2\).

To verify this, observe that (the \(D_3\) possessing an adjoint curve of order 0) a tricanonical curve of \(F\) will be given by the branch curve \(D_9\) to which the exceptional curve \(r\) is added. This tricanonical curve is isolated, that is, not capable of varying in a pencil, since otherwise, having no points in common with the bicanonical curve represented by \(D_3\), it would have to contain it, and one would have \(p_g>0\). This means that the double elliptic curves of the 9\textsuperscript{o} order, of the Halphen pencil defined by the \(D_9\), are images of irreducible curves (self-conjugate in the involution \(I\) defined on \(F\) by the double plane). %
This deduction can also be confirmed directly by the criterion of Comessatti. Finally we have \textit{one} tricanonical curve \(D_3\) and a pencil of irreducible sexticanonical curves represented by the curves of the Halphen pencil, to which belong the tricanonical curve counted twice and the bicanonical curve counted three times. We therefore have: %
\[
P_3=1,\qquad P_6=2.
\]

\edsection{4. Conditions for the rationality of a surface.}{4. Conditions for the rationality of a surface.}{sec.VI.4}

A rational surface, that is, one that can be put into point-by-point correspondence with the plane, must, like the plane, have genus \(p=p_a=p_g=0\), and together with the genus all its plurigenera must vanish. The examples we have adduced (in particular the surface of the 6\textsuperscript{o} order which passes doubly through the edges of a tetrahedron) show that a regular surface of genus \(p=0\) can have bigenus \(P>0\), so that the vanishing of the genus does not suffice to furnish the conditions for the rationality of a surface. Now it is extremely interesting to recognize that these conditions can be expressed by equating to zero not only the genus but also the bigenus, that is, by writing
\[
p_a=P=0
\]
(from which \(p_g=0\) necessarily follows as well). This theorem is due to G. \textsc{Castelnuovo}, who gave it in a memoir of 1896\origfootnote{1}{s127L-1}{G. \textsc{Castelnuovo}, \textit{Sulle superficie di genere zero}. Soc. It. delle Scienze (detta dei XL), cf. \textit{Memorie scelte}, p. 307.}.

The proof of our theorem can be reduced to the rationality criteria of Noether and of Clebsch-Noether, relating to surfaces with a linear pencil of rational curves and to the types of double planes determined by them. For this purpose one starts from an irreducible linear system \(\lvert C\rvert\) on the given surface and constructs suc%
% scan: s127R
% !PRINTED 231
\unskip\weberpage{231}%
cessively its adjoint system \(\lvert C'\rvert\), the second adjoint \(\lvert C''\rvert\), and so on. It will be recognized that, since \(p_a=P=0\), the series of these successive adjoints comes to an end, and the last adjoint system that is at least \(\infty^1\) turns out to be composed of rational or elliptic curves, and therefore leads to one of the cases of rationality contemplated in the criteria of Noether and of Clebsch-Noether discussed above.

The discussion can be carried out in a very simple form by distinguishing two hypotheses: 1\textsuperscript{o}) the given surface \(F\) contains an irreducible linear system \(\lvert C\rvert\) of genus \(\pi\) and of degree \(n>2\pi-2\), which may be supposed to be the system of plane or hyperplane sections of an \(F\) entirely free of singularities; 2\textsuperscript{o}) for every irreducible linear system given on the surface \(F\) the inequality \(n\le2\pi-2\) holds, \(n\) and \(\pi\) denoting the degree and the genus of the system.

\textit{Hypothesis I.} -- It is clear that, by constructing the successive adjoints \(\lvert C'\rvert,\ \lvert C''\rvert\), etc., starting from the system \(\lvert C\rvert\) given on \(F\), one obtains a finite series of systems; in other words, the process of adjunction comes to an end. Indeed the curves \(C',\ C'',\ \ldots\ldotp\) will have a decreasing order
\[
2\pi-2,\qquad 2\pi-2-d,\qquad 2\pi-2-2d,\qquad \ldots\ldotp
\]
where
\[
d=n-(2\pi-2).
\]

We therefore wish to consider the last adjoint system \(\lvert C^{(t)}\rvert\) of dimension \(r_t\ge1\).

Let us suppose first that the adjoint systems \(\lvert C'\rvert,\ \lvert C''\rvert,\ \ldots\ldotp,\ \lvert C^{(t)}\rvert\) of the aforesaid series are irreducible. The system \(\lvert C^{(t)}\rvert\), the first of the series which has no adjoint of dimension greater than zero, consists of curves of genus zero or one. If the system \(\lvert C^{(t)}\rvert\) is formed of curves of genus zero, the theorem of Noether\origfootnote{1}{s127R-1}{Cf. e.g. \textsc{Enriques-Conforto}, Op. cit.} on surfaces containing a linear pencil of rational curves allows one to deduce at once that \(F\) is rational.

If instead the system \(\lvert C^{(t)}\rvert\) is of genus 1, it is convenient to distinguish the hypothesis that its dimension is \(r_t\ge2\), or \(r_t=1\). If \(r_t\ge2\), the complete system of elliptic curves \(\lvert C^{(t)}\rvert\), which, by the regularity of the surface, has complete characteristic series, will have degree \(n_t=r_t\), and will therefore contain within itself a net of degree 2, by means of which the surface \(F\) is represented on the double plane with branch quartic: from this it follows\origfootnote{2}{s127R-2}{Theorem of Clebsch. Cf. \textsc{Enriques-Conforto}, Op. cit.} that our surface is rational.

If instead \(r_t=1\), the system \(\lvert C^{(t-1)}\rvert\) which precedes \(\lvert C^{(t)}\rvert\) in the series of adjoints must be of genus two, and may be supposed
% scan: s128L
% !PRINTED 232
\weberpage{232}
of dimension \(r_{t-1}\ge2\), because otherwise one could go back to a preceding system, still of genus two. Now one can see that the system \(\lvert C^{(t-1)}\rvert\) is not simple: all the \(C^{(t-1)}\) passing through a point \(P\) of the surface \(F\) consequently pass through a point \(P'\) which is conjugate to \(P\) in the \(g_2^1\) on it; indeed the conjugate of \(P\) on a \(C^{(t-1)}\) varying through \(P\) must be fixed, otherwise it would describe a curve \(C^{(t)}\), or a component of \(C^{(t)}\), of genus zero, contrary to our hypotheses.

This established, since \(F\) is a regular surface, the system \(C^{(t-1)}\) adjoint to \(C^{(t-2)}\) will be a regular system, of dimension \(p_{t-2}-1\), \(p_{t-2}\) denoting the genus of \(\lvert C^{(t-2)}\rvert\), and it will have complete and non-special characteristic series. Therefore this series, which we have seen to be composed with the \(g_2^1\), will be precisely the double of this \(g_2^1\). Hence the system \(\lvert C^{(t-1)}\rvert\) will be an \(\infty^3\) system of curves of genus 2, belonging to an involution of pairs of points, by means of which our surface \(F\) is represented on a double quadric \(Q\), with branch curve of the sixth order \(D_6\). Since the generators of \(Q\) cannot possess an odd number of branch points, one deduces that \(Q\) is a cone, whose vertex constitutes a branch point to be added to the \(D_6\), which cuts the generators in three points. The double cone so defined can be represented on Noether's double plane with branch sextic endowed with two infinitely near triple points, and from this it follows that the double plane is rational\origfootnote{1}{s128L-1}{Theorem of Noether. Cf. \textsc{Enriques-Conforto}, Op. cit.}.

We have assumed that the series of successive adjoint systems \(\lvert C'\rvert,\ \lvert C''\rvert,\ \ldots\ldotp\) is formed of irreducible systems. If, in running through this series, one meets a system \(C^{(h)}\) endowed with a fixed curve \(\theta\), and whose variable part is irreducible, belonging to a system of dimension \(r_h\ge2\), one can in general continue the construction of our series of adjoint systems by putting in place of \(\lvert C^{(h)}\rvert\) the irreducible system \(\lvert C^{(h)}-\theta\rvert\), stripped of the fixed part \(\theta\). And if this system is of genus 0 or 1, and hence (by the completeness of the characteristic series) of degree \(n_h=r_h-1\) or \(n_h=r_h\), one will conclude (as before) that the surface \(F\) is rational.

It remains to examine more closely the case in which the system \(\lvert C^h-\theta\rvert\), stripped of the fixed part \(\theta\), is a pencil or is composed of the curves \(L\) of a pencil which, by the regularity of the surface, will be \textit{linear}.

In the first place, if the curves \(L\) of the pencil are rational, the well-known theorem of Noether allows us to assert the rationality of the surface.

% scan: s128R
% !PRINTED 233
\weberpage{233}
Assume instead that the curves \(L\) have genus \(\varrho>0\). Let us observe first of all that \(\theta\), the fixed part of \(\lvert C^h\rvert\), must be a fundamental curve for the irreducible system \(\lvert C^{h-1}\rvert\), since \(\lvert C^h\rvert\) cuts out on a \(C^{h-1}\) the complete canonical series, which is free of fixed points.

This established, the curves \(L\) will meet the \(C^{h-1}\) in \(\nu>1\) variable points (since from the hypothesis \(\nu=1\) one would deduce that the \(L\) are rational), and the complete canonical series cut out by \(\lvert C^h-\theta\rvert\) on a \(C^{h-1}\) will be composed of the groups of the involution of order \(\nu\) cut out by the pencil of the \(L\); hence one deduces that the \(C^{h-1}\) are hyperelliptic curves, with \(\nu=2\): the pencil of the \(L\), of which the curves of \(\lvert C^h-\theta\rvert\) are composed, is a linear pencil of bisecants.

Now let us recall that the complete system \(\lvert C^{h-1}\rvert\) adjoint to the irreducible system \(\lvert C^{h-2}\rvert\) (on the regular surface \(F\) of genus \(p=0\)) is a regular system (of dimension \(\pi_{h-2}-1\)), and therefore has complete and non-special characteristic series. This series is composed with the \(g_2^1\), because otherwise the conjugate of a point \(P\) on a \(C^{h-1}\) through \(P\) would describe a rational curve, forming the \(L\), or a component of the \(L\), through \(P\). Consequently the system \(\lvert C^{h-1}\rvert\) is not simple, but composed with an involution of the 2\textsuperscript{o} order, and the degree \(n_{h-1}\) of this system of hyperelliptic curves will be twice the dimension diminished by one unit:
\[
n_{h-1}=2(r_{h-1}-1);
\]
while its genus \(\pi_{h-1}\) will be
\[
\pi_{h-1}=r_{h-1}-1.
\]

We can arrive at a more precise determination of these characters by means of the following consideration. If one detaches from the system \(\lvert C^{h-1}\rvert\) a curve \(L\) (on which the \(C^{h-1}\) cut out the pairs of a \(g_2^1\)), one obtains a linear system of dimension \(r_{h-1}-2\), which must have as adjoint system the adjoint system of \(\lvert C^{h-1}\rvert\) diminished by the \(L\), and therefore the system constituted by all the groups of \(\pi_{h-1}-2\) curves \(L\); hence the detachment of the \(L\) from \(\lvert C^{h-1}\rvert\) diminishes its genus by one unit only. But on the other hand the \(C^{h-1}\) which come to split into the \(L\) and a residual curve must be connected, and hence the \(L\) must meet the residual curves in \(i\ge1\) points; thus, \(\pi_{h-1}-1\) being the genus of these curves, one will have
\[
\pi_{h-1}=(\pi_{h-1}-1)+\varrho+i-1
\]
from which (since \(\varrho>0\)) there follows
\[
\varrho=1,\qquad i=1.
\]

% scan: s129L
% !PRINTED 234
\weberpage{234}
Therefore the residual curves of an \(L\) with respect to \(\lvert C^{h-1}\rvert\) cut the \(L\) in a point which is necessarily fixed, and common to the \(L\) of our pencil, and this entails that these residual curves split into several \(L\) of the pencil itself.

Now, by means of the system \(\lvert C^{h-1}\rvert\), let us transform the surface \(F\) into a double surface of order \(m=r_{h-1}-1\) in the space of \(m+1=r_{h-1}\) dimensions. This surface, on which the curves \(L\) are represented by lines, will be a cone having a certain vertex \(O\), and must possess a branch curve \(D\) of order %
\[
2+2\pi_{h-1}=2m+2.
\]

If this \(D\) is to meet the generators of the surface in at least 3 variable points (in order that the genus of the \(L\) be \(\varrho>0\)), one must have \(m=2\); indeed a hyperplane drawn through \(O\) cuts our cone in \(m\) generators, and therefore one must have
\[
3m\le2m+2.
\]

In this way we have come to conclude that the system \(\lvert C^{h-1}\rvert\) leads to representing our surface \(F\) on a double quadric cone with branch sextic, which, as we have already recalled, is rational (Noether).

\textit{Hypothesis II.} -- We must still examine the second hypothesis: every irreducible linear system traced on the surface \(F\) has genus \(\pi\) and degree \(n\) satisfying the inequality
\[
n\le2\pi-2.
\]

Let us assume this hypothesis, which we shall reduce to an absurdity. According to the assumption, the surface \(F\) may be supposed transformed so as to contain no exceptional curves, and, if one wishes, free of singularities in a suitable space. Its hyperplane sections constitute a linear system \(\lvert C\rvert\), of which we denote the degree by \(n\) and the genus by \(\pi\), where, as has been said,
\[
n\le2\pi-2.
\]

Then the bigenus \(P\) of \(F\) can be evaluated by increasing by one unit the dimension of the bicanonical system, which is obtained by detaching a \(C\) from the second adjoint system \(\lvert C''\rvert\); and, as we have seen, one finds that for a regular surface of surface genus \(p=p_a=p_g=0\) and of linear genus \(p^{(1)}\), the bigenus is
\[
P\ge p^{(1)}.
\]

Consequently, if \(P=0\), one will have
\[
p^{(1)}\le0.
\]

% scan: s129R
% !PRINTED 235
\weberpage{235}
Then consider the series of the systems successively adjoint to \(\lvert C\rvert,\ \lvert C'\rvert,\ \lvert C''\rvert,\ \ldots\ldotp\); the number of intersections \((2\pi-2-n)\) which \(\lvert C\rvert\) has with the virtual canonical curves increases, on passing to the adjoint system, by \(p^{(1)}-1\), that is, in reality it decreases by at least one unit. Therefore at least one of the two following facts must occur:

1\textsuperscript{o}) either the difference \(n-(2\pi-2)\) formed with the characters of a system \(\lvert C\rvert\) becomes positive for some one of the successive adjoints, which directly contradicts our hypothesis;

2\textsuperscript{o}) or the series of the said adjoints is finite.

In that case, just as happened in Hypothesis I), one arrives at a last linear system composed of the rational curves of a linear pencil, or else at a pencil of elliptic curves which presents itself as adjoint to the hyperelliptic curves of genus 2 of a linear system, at least \(\infty^2\), of degree 4. %
By the same reasoning carried out in Hypothesis I, one deduces from this the rationality of the surface \(F\). But it is not even necessary to repeat that reasoning; it suffices to point out that, in contradiction to our Hypothesis II, we have succeeded in constructing on \(F\) a linear system of genus \(\pi\) and degree \(n>2\pi-2\), so that the hypothesis itself is reduced to an absurdity. Here it is opportune to observe that when one has on \(F\) a system \(\lvert L\rvert\), reducible in any way, which contains no exceptional components, of genus \(\pi\) and degree \(n>2\pi-2\), by adding a multiple of it to a regular system one can always obtain an irreducible linear system whose characters satisfy the same inequality. We shall conclude by stating the theorem: \textit{the necessary and sufficient conditions for a surface to be rational are expressed by the vanishing of the arithmetic genus and of the bigenus:}
\[
p_a=P=0;
\]
where the condition \(P=0\) entails a priori that \(p_g=0\) as well. Thus, in other words, one may say that the \textit{rational surfaces} are characterized as « \textit{regular surfaces of bigenus zero} ».

This theorem belongs to G. \textsc{Castelnuovo}, who proved it substantially by the method followed here, in which only simplifications of details have been introduced.
% scan: s130L
% !PRINTED none
\weberblank{[s130L]}

\clearpage\thispagestyle{plain}
% scan: s130R
% !PRINTED none
\weberpage{[s130R]}
\edchapter{VII}{CLASSIFICATION OF SURFACES\\OF LINEAR GENUS \(p^{(1)}=1\).}{Chapter VII. CLASSIFICATION OF SURFACES OF LINEAR GENUS p\textasciicircum{}(1)=1.}{CLASSIFICATION OF SURFACES OF LINEAR GENUS \(p^{(1)}=1\).}

\edsection{1. Surfaces of genus zero and bigenus 1, with bicanonical curve of order zero.}{1. Surfaces of genus zero and bigenus 1, with bicanonical curve of order zero.}{sec.VII.1}

We have seen (Chap. VI, § 2) that among the regular surfaces of genus \(p=0\) and bigenus \(P=1\), those endowed with a bicanonical curve of order zero are characterized, with respect to birational transformations, by the value of the trigenus, which for them is
\[
P_3=0,
\]
or of the sextigenus, which is
\[
P_6=1
\]
(instead of \(P_6>1\)).

We now wish to deepen the study of these surfaces \(F\) by reducing them to a projectively defined type, which will be the double plane with branch curve \(D_8\) formed by a sextic with two tacnodes \(M\) and \(N\) and with a double point at the intersection of their tacnodal tangents-\(p\) and \(q\), to which are added these two lines \(p\) and \(q\). We shall then recognize that this double plane can be transformed into a surface of the 6\textsuperscript{o} order \(F_6\), passing doubly through the edges of a tetrahedron, which can therefore equally be taken as the type of our family of surfaces (\(p=0,\ P=1,\ P_3=0\)). %

Let us recall first of all the fundamental properties of our surfaces \(F\), referring to a model of them free of exceptional curves:

1\textsuperscript{o}) every curve \(C\) of genus \(\pi\) traced on \(F\) has degree
\[
n=2\pi-2;
\]

2\textsuperscript{o}) the dimension of the complete linear system \(\lvert C\rvert\) (virtually free of base points) is
\[
r\ge\pi-1;
\]
% scan: s131L
% !PRINTED 238
\weberpage{238}
\(3^\circ\)) for every irreducible system \(|C|\) of genus \(\pi > 1\) one has exactly
\[
r = \pi - 1
\]
that is, the system is regular;

\(4^\circ\)) the linear system \(|C|\) has as adjoint a system \(|C'|\), distinct from it, with the same characters; the relation between the two systems is reciprocal, \(|C|\) being also the adjoint of \(|C'|\): the two systems are two different halves of the same double system
\[
|L| = |2C| = |2C'|;
\]

\(5^\circ\)) to the surface \(F\) there may belong \textit{elliptic curves}, that is, irreducible curves of virtual genus \(\pi = 1\), and hence of degree \(n = 0\). These curves can be \textit{isolated}, that is, constituting a regular system, of dimension \(r = 0\), or else belonging to a pencil (a superabundant system, of dimension \(r = 1\)). If \(|C|\) is a pencil of elliptic curves on \(F\), the adjoint system \(|C'|\) is constituted by two isolated elliptic curves \(\dfrac{C}{2}\) and \(\left(\dfrac{C}{2}\right)'\), adjoint to each other, whose doubles belong to the pencil. Indeed the adjoint curve \(C'\), having no variable intersections with the \(C\) of the pencil, and being distinct from a \(C\), must consist of parts of curves \(C\); and on the other hand \(|2C'|\) must be equivalent to a pair of curves \(C\).

Conversely, if in the pencil \(|C|\) there is a half curve \(\dfrac{C}{2}\), its adjoint \(\left(\dfrac{C}{2}\right)'\) will also be a half curve of the pencil, and the curves adjoint to \(\dfrac{C}{2} + \left(\dfrac{C}{2}\right)'\) will be the \(2\dfrac{C}{2} = C\). And hence, by property \(4^\circ\)), to a pencil of elliptic curves \(|C|\) there can belong only a single pair of double elliptic curves, formed with two halves of the pencil, which are adjoint to each other;

\(6^\circ\)) a curve \(\theta\) which is not a component of the elliptic curves of a pencil \(|C|\) necessarily cuts the \(C\) in an even number of points. Indeed this number is double the number of intersections of \(\theta\) with the curves \(\dfrac{C}{2}\);

\(7^\circ\)) an irreducible curve \(\theta\) of (virtual) genus zero, and hence of degree \(-2\), on the surface \(F\), cannot possess any adjoint curve.

Indeed, if there existed an adjoint \(\theta'\), it would have to meet \(\theta\) in \(-2\) points and would hence coincide with \(\theta\); but from \(\theta = \theta'\) it would follow that the genus \(p = 1\), contrary to the hypothesis.

In other words, the existence of \(\theta'\) entails that \(\theta\) is of genus \(> 0\);

% scan: s131R
% !PRINTED 239
\weberpage{239}
\(8^\circ\)) on the surface \(F\), given a pencil of elliptic curves \(|C|\), there cannot exist an irreducible curve \(\theta\) of genus zero which is bisecant to them.

Indeed the system \(C + \theta\), which would be a pencil of curves of virtual genus \(\varrho = 2\), would have to admit an adjoint system \(|L|\) of dimension 1, containing the composite curve \(\theta + \dfrac{C}{2} + \left(\dfrac{C}{2}\right)'\); instead one can recognize that this curve cannot vary in a pencil. To this end observe that the \(\dfrac{C}{2}\) must meet the \(L\) in a single point, since
\[
\frac{C}{2}L = \frac{C}{2}\theta + \frac{C}{2}\frac{C}{2} + \frac{C}{2}\left(\frac{C}{2}\right)'
\]
and
\[
\frac{C}{2}\frac{C}{2} = 0;
\]
hence the \(L\) cut \(\dfrac{C}{2}\), and similarly \(\left(\dfrac{C}{2}\right)'\), in a fixed point, and therefore — if there are \(\infty^1\) of them which do not reduce to \(\theta\) itself increased by the \(C' \doteq \dfrac{C}{2} + \left(\dfrac{C}{2}\right)'\) — they turn out to be devoid of intersections with the \(C\), and hence composed of the said \(C\) and of a curve \(\theta'\) adjoint to \(\theta\): but the existence of \(\theta'\) would entail that the genus of \(\theta\) is \(> 0\), contrary to the hypothesis. %

Let us bear in mind propositions 1) .... 8).

Since \(|C|\) is an irreducible linear system of dimension \(\pi - 1\), consider the double system \(|L|\), which has degree \(4n = 8\pi - 8\), genus \(4\pi - 3\), and dimension \(4\pi - 4\). Within the system \(|L|\) there are two \(\infty^{2\pi-2}\) systems, say \(\Sigma\) and \(\Sigma'\), constituted, the first by the pairs of curves \(C\), and the second by the pairs of curves \(C'\), each of which has the index
\[
i = \frac{1}{2}\binom{2\pi - 2}{\pi - 1}.
\]

Denoting by \(S_{4\pi-4}\) the space whose points represent the elements (curves) of the system \(|L|\), there will be in it two varieties representing \(\Sigma\) and \(\Sigma'\), of dimension \(2\pi - 2\) and of order \(i\), which will have in common \(i^2\) points, or else an infinity of points. Thus one finds curves \(L\) which are simultaneously decomposable into two curves \(C\) and into two curves \(C'\); these \(L\) are necessarily split, and in general they allow one to construct, within a system \(|C|\) of genus \(\pi > 1\), a system partially contained in it, of genus less than \(\pi\).

To be precise, one will find a curve composed of two \(C\):
\[
C = C_1 + C_2,\qquad C = C_3 + C_4
\]
% scan: s132L
% !PRINTED 240
\weberpage{240}
decomposed in such a way that a part of the first, together with a part of the second, must give a curve \(C'\), say for example
\[
C_1 + C_3 = C',\qquad C_2 + C_4 = C',
\]
so that, since
\[
C_1 + C_3 = (C_1 + C_2)'
\]
one will deduce
\[
C_3 = C'_2
\]
and similarly
\[
C_4 = C'_1.
\]

This implies that the curves \(C_1\), \(C_2\), and similarly the \(C_3\) and \(C_4\), are of genus greater than or equal to one.

We shall distinguish two cases:

\textit{a}) among the two curves \(C_1\) and \(C_2\) there is one, e.g. \(C_1\), of genus \(\pi_1 > 1\).

\textit{b}) the two curves \(C_1\) and \(C_2\) are of genus 1.

In case \textit{a}) one can assert that \(C_1\) belongs to a complete linear system \(|C_1|\) of dimension \(\pi_1 - 1 \geq 1\). Since this system is partially contained in \(|C|\), its dimension \(\pi_1 - 1 < \pi - 1\), and hence the genus \(\pi_1 < \pi\). %

Moreover, the variable parts of the curves \(C_1\) must be irreducible, and still of genus \(\varrho > 1\) (\(\varrho < \pi\)). Indeed, suppose that these parts are composed of \(s\) curves \(K\) of a linear pencil (\(s > 1\)). If the \(K\) were of genus greater than 1, the complete system which is the \(s\)-fold multiple of \(|K|\) would turn out to be of dimension greater than \(s\), and hence irreducible.

Let instead the genus of the \(K\) be equal to 1. In this hypothesis the groups of \(sK\) are only parts of the curves \(C_1\):
\[
|C_1| = \theta + |sK|.
\]

And in order that the genus of \(C_1\) be \(\pi_1 > 1\), it is necessary that \(\theta\), or a component \(\theta_1\) of it, meet the \(K\) in some point (a curve not intersecting the \(K\) turning out to be part of a \(K\), and hence of genus less than or equal to 1). The number of these intersections being even, and therefore greater than or equal to two, the addition of \(\theta_1\) to \(|sK|\) certainly succeeds in enlarging the system \(|sK|\) if \(\theta_1\) is of genus 1, and also if \(\theta_1\) is of genus 0, since in that hypothesis the number of intersections with the \(K\) must be greater than 2 (cf. prop. 8) and hence greater than or equal to 4.

In conclusion, the system of the \(C_1\), or else a complete system contained in it, constitutes an irreducible linear system of genus greater than one.

Therefore, the repeated application of our process of reduction. where case \textit{b}) is not met with, leads to finding a linear system partially contained in \(|C|\) and of genus two. %

% scan: s132R
% !PRINTED 241
\weberpage{241}
Let us place ourselves instead in case \textit{b}), supposing therefore that among the curves of a system \(|C|\) of genus \(\pi\) there is some curve split into two curves of genus 1,
\[
C = C_1 + C_2.
\]

Since \(C_1\) and \(C_2\) are curves of genus 1 and of degree 0, their doubles \(2C_1\) and \(2C_2\) will give rise to two pencils of elliptic curves, whose variable parts we shall denote by \(K_1\) and \(K_2\); and, as before, it is easy to recognize that neither of these two pencils can possess fixed parts which are not components of the curves of the pencil itself. Hence the curves of the two pencils \(|K_1|\) and \(|K_2|\) must intersect in a number of points greater than zero, in order that the degree of \(|C_1 + C_2|\) turn out to be greater than zero. Now two half curves of \(K_1\) and \(K_2\), say \(\dfrac{K_1}{2}\) and \(\dfrac{K_2}{2}\), constitute two isolated elliptic curves having a certain number \(s\) of intersections, such that
\[
1 \leq s \leq \pi - 1.
\]

In conclusion one may hold that, if our process of reduction does not lead to constructing a pencil of genus 2, one will always arrive at determining two isolated elliptic curves, which we again call \(C_1\) and \(C_2\), with a certain number \(s \leq \pi - 1\) of points in common (\(s > 0\)). Strictly speaking, one cannot assert that these \(C_1\) and \(C_2\) are without common parts; it might happen, for example, that \(C_1\) is split into two curves of genus zero with two common points, \(C_1 = \theta + \theta_1\), and similarly that \(C_2 = \theta + \theta_2\); then \(C_1\) and \(C_2\) would have in common the component \(\theta\); but in the deductions founded on the use of virtual characters, it is « as if » the two curves \(\theta + \theta_1\) and \(\theta + \theta_2\) were irreducible with
\[
s = (\theta + \theta_1,\ \theta + \theta_2)
\]
common points.

The process of reduction which we have put into operation, starting from an irreducible linear system on the surface \(F\), leads to constructing on \(F\): either an \(\infty^1\) linear system \(|C|\) of genus 2, or a pair of elliptic curves \(C_1\) and \(C_2\) with \(s(\geq 1)\) common points.

In the first case the curves \(C\) must be held to be irreducible, because their variable parts, constituting a pencil, are a priori irreducible and of genus 2: if indeed these were curves \(K\) of genus 1, \(|C|\) would possess a fixed part \(\theta\) cutting the \(K\) in only 2 common points, which is impossible.

Therefore there will be found within \(|C|\) curves split into two isolated elliptic components \(C_1\) and \(C_2\), with \textit{one} common point.

In the second case, when one has on \(F\) two isolated elliptic curves \(C_1\) and \(C_2\) with \(s > 1\) common points, one can find on \(F\) a pair of elliptic curves intersecting in \(s' < s\) (\(s' \geq 1\)) points, and hence also
\printersignature{\textsc{Enriques F.} - \textit{Superficie algebriche.}\quad 16}
% scan: s134L
% !PRINTED 242
\weberpage{242}
a pair of elliptic curves with \textit{one}-common point. The possibility of this new reduction is demonstrated as follows. %

If \(s > 1\), the linear system \(|C| = |C_1 + C_2|\) is of genus \(\pi = s + 1 > 2\). Hence among the pairs of curves \(2C\) there will in general be \(\dfrac{1}{4}\dbinom{2\pi - 2}{\pi - 1}^{2} \geq 9\) split curves. If these contain components of genus greater than or equal to two, one succeeds, by the process already explained, in constructing a pair of elliptic curves with a number of intersections \(s' < s\). If instead all the said split curves are formed of elliptic components, it is necessary to develop the following considerations. %

Precisely, one must admit that, besides the pair of split curves which constitutes
\[
C_1 + C_2 + C'_1 + C'_2,
\]
there is another pair of curves split into elliptic components
\[
C_3 + C_4 + C'_3 + C'_4,
\]
where \(C_3\) and \(C_4\) are distinct from \(C_1\), \(C_2\), \(C'_1\) and \(C'_2\): the existence of \(C_3 + C_4\) is deduced from admitting that the two varieties \(\Sigma\) and \(\Sigma'\) representing the pairs of curves \(C\) and \(C'\) have some other point in common, \textit{distinct} from the two points which correspond to the pairs
\[
C_1 + C_2 \qquad C'_1 + C'_2
\]
and
\[
C_1 + C'_2 \qquad C_2 + C'_1;
\]
and indeed we shall see shortly that the said pairs of isolated elliptic curves \(C_1 + C_2\) and \(C'_1 + C'_2\) cannot correspond to intersections of \(\Sigma\) and \(\Sigma'\) which fall infinitely near.

Now then, since \(C_3 + C_4\) is a split curve of the system \(|C|\) different from \(C_1 + C_2\) and from \(C_1 + C'_2\), it will be necessary that \(C_3\) and \(C_4\) cut each of the two curves \(C_1\) and \(C_2\) in some point, for otherwise they would contain as a part \(C_1\) itself, or else one of the two curves \(C_2\) or \(C'_2\); it follows that, for example, \(C_3\) must meet each of \(C_1\) and \(C_2\) in a number of points less than \(s\); so that we have obtained, as desired, a pair of isolated elliptic curves \(C_1\) and \(C_3\) with \(s' < s\) points in common.

It remains to justify that our system \(|C|\) cannot contain two infinitely near pairs of isolated elliptic curves \(C_1 + C_2\) and \(C_3 + C_4\), where, e.g., \(C_3\) is infinitely near to \(C_1\) and \(C_4\) to \(C_2\). Indeed the curves \(2C_1\) and \(2C_3\) would have to define a pencil of elliptic curves, of which \(C_1\) would turn out to be a fixed component, so that \(C_1\) itself
% scan: s134R
% !PRINTED 243
\weberpage{243}
would come to vary in a pencil of elliptic curves, whereas it was supposed isolated.

From what precedes one draws the conclusion that our surface \(F\) possesses a pair of isolated elliptic curves \(C_1\) and \(C_2\) having one point in common, in such a way that their adjoints, \(C'_1\) and \(C'_2\), are incident to \(C_2\) and to \(C_1\). The composite curve \(C_1 + C_2 + C'_2\) is therefore of genus 3, and defines a net of curves \(|K|\) of this genus: the \(K\) cut \(C_2\) and \(C'_2\) each in a single point, which must be fixed; hence the \(K\) have two simple base points, which fall precisely at the points common to \(C'_1\) and \(C_2\), and to \(C_1\) and \(C'_2\), since \(C'_1\) forms part of \(\infty^1\) curves \(K\), having as residual with respect to \(|K|\) the pencil of elliptic curves \(|2C_2|\): the assertion follows from observing that the curves \(C'_1 + 2C_2\) are second adjoints to the \(K = C_1 + C_2 + C'_2\), and hence are equivalent, belonging to the same system \(|K|\).

Therefore the \(K\) will turn out to be hyperelliptic curves, cutting one another two by two in two variable points; these \(K\) then cut out on \(C_1\) the pairs of a \(g^1_2\), there being only one residual curve of \(C_1\) with respect to \(|K|\), which is \(C_2 + C'_2\). It may be added that the base points of the \(K\) are, on any \(K\) whatever, two double points of the corresponding \(g^1_2\).

By means of the net of hyperelliptic curves \(|K|\) of genus 3, one can now represent the surface \(F\) on a double plane, which will have a branch curve \(D_8\) of order 8.

This \(D_8\) will include two branch lines, \(p\) and \(q\), corresponding to the base points of \(|K|\); and the point \(O = pq\) will be the image of the elliptic curve \(C'_1\), and hence quadruple for the \(D_8\), or double for its component \(D_6\) which remains on removing the lines \(p\) and \(q\). The lines of the pencil with center \(O\) will be, in the double plane, the images of the elliptic curves of the pencil \(|2C_2| = |2C'_2|\); to \(C_2\) and \(C'_2\) there will correspond two tacnodes of the \(D_6\) with the tacnodal tangents \(p\) and \(q\), that is, two points \([3,3]\) of the branch curve \(D_8\). These singularities of the \(D_8\) do in fact suffice to characterize it as the branch curve of a double plane of genus \(p_a = p_g = 0\) and \(P = 1\), with (pure) bicanonical curve of order zero, as is recognized on the basis of Chap. V; \S\ 7.

We conclude therefore that \textit{the regular surfaces of genus zero and bigenus \(1\), with bicanonical curve of order zero — such as are characterized by the value \(P_3 = 0\) of the trigenus, or \(P_6 = 1\) of the sextigenus — form a single family, which has as its type the double plane with branch curve}
\[
D_8 = D_6 + p + q,
\]
\textit{consisting of a sextic with two tacnodes and with a double point at the point of intersection of their tacnodal tangents, to which these same tangents are added.}

% scan: s135L
% !PRINTED 244
\weberpage{244}
This \textit{family} of surfaces depends on 10 \textit{moduli}. Indeed the construction of the \(D_6\) with 4 assigned double points, and hence of the \(D_8\), depends on
\[
\frac{6 \cdot 9}{2} - 5 \cdot 3 = 12
\]
constants, so that, adding the 6 constants which depend on the choice of the singular points of the \(D_6\), one has \(\infty^{18}\) curves \(D_8\). From the 18 parameters which belong to them one must remove 8, corresponding to the \(\infty^8\) homographic \(D_8\), which naturally define the same surface \(F\).

Consider the sextic surface \(F_6\) of ordinary space passing doubly through the edges of a tetrahedron \(T\). Since \(F_6\) is a regular surface with \(p = 0\), \(P = 1\), \(P_3 = 0\), it must be possible to represent it on a double plane with branch curve \(D_8\) endowed with the singularities mentioned above. One actually obtains a representation of \(F_6\) on a double plane by means of the skew projection from a pair of opposite edges \(a\) and \(b\) of the tetrahedron \(T\), since the lines incident to \(a\) and \(b\) meet \(F_6\) in two simple points. But this construction leads in general to a double plane with branch curve \(D_{10}\), which can be brought back to our \(D_8\) only by a further transformation of the plane. It is therefore useful to particularize the said construction, taking a representative plane which passes through an edge \(c\) of the tetrahedron, incident to \(a\) and \(b\). Then the lines of the double plane appear as images of the curves \(K\) cut out on \(F_6\) by the quadrics through \(a\), \(b\), \(c\). These \(K\) are hyperelliptic curves of genus 3, among which there is a curve split into the neighborhoods of 3 double lines of the tetrahedron \(T\): \(c\), \(d\), and \(e\), \(d\) and \(e\) being the edges of \(T\) that are the further sections of \(F_6\) with the planes \(ca\) and \(cb\). Note that the components of this split curve \(K\), that is, the neighborhoods of the 3 double lines \(c\), \(d\), and \(e\), are elliptic curves connected like the \(C_1\), \(C_2\), \(C'_2\) which we used above to construct the net of the \(K\) that led us to the double plane with branch curve \(D_8\). %

One thus succeeds in verifying in the simplest way the conclusions which are drawn a priori for \(F_6\) from the property of its being a regular surface of genus \(p = 0\) and bigenus \(P = 1\), with bicanonical curve of order zero.

Now observe that there are in space \(\infty^{13}\) sextics \(F_6\) passing doubly through the edges of a given tetrahedron \(T\) (whose equation we wrote in \S\ 11 of Chap. III), and that there are \(\infty^3\) homographies with their fixed points at the vertices of \(T\); it follows from this that the \(F_6\) depend on 10 parameters or projective invariants, just like the double planes with branch curve \(D_8\) on which we have learned to represent them. That these projective invariants give as many moduli of the surfaces \(F_6\) with respect to birational transformations is deduced
% scan: s135R
% !PRINTED 245
\weberpage{245}
from the fact that the system of plane sections is complete and, the surface being regular, cannot be contained in a continuous series of analogous systems, although indeed there may be on \(F\) a discontinuous series of such systems\origfootnote{1}{s135R-1}{The actual existence of such a series follows from the property of \(F_6\) of possessing a discontinuous infinity of birational transformations into itself. These can be defined in relation to the \(\infty^1\) elliptic quartics \(K\) cut out by the quadrics passing through 4 edges \(a\,b\,d\,e\) of \(T\) which form a skew quadrilateral. On each of these curves one obtains a birational transformation by adding to a point \(P\) the virtual group of order zero, the difference of the pairs \(G_d - G_e\) that are the sections by two incident edges \(d\) and \(e\): where it is useful to note that these pairs \(G_d\) and \(G_e\) cannot be equivalent, nor give rise to equivalent multiples, since it would follow that \(d\) and \(e\) themselves, or their multiples, would turn out to be equivalent, which contradicts the fact that \(d\) and \(e\) have a point in common. Therefore the transformation defined on each \(K\) of our pencil, carrying a point \(P\) to the point \(P' = P + G_d - G_e\), cannot be cyclic, and thus will give rise to an infinite series of powers, which again give birational transformations of the surface into itself.}.

It follows from this that the \(F_6\), on a par with and in place of the corresponding double plane with branch curve \(D_8\), can be taken as the general type of the regular surfaces of genus \(p = 0\) and of bigenus \(P = 1\), with bicanonical curve of order zero (\(P_3 = 0\)).

The problem of constructing the transform \(F_6\) starting from the double plane with branch curve \(D_8\) can be solved by determining in this plane a cubic \(K_3\) passing through the 5 double points of \(D_6\), which touches \(D_6\) in 4 points and which is the image of a split elliptic curve on the surface \(F\) represented by our double plane. That among the cubics \(K_3\) adjoint to \(D_6\) and quadritangent to it there actually exist images of split curves can be proved by using the criterion of Comessatti, already recalled several times. It is a matter of proving that there is a \(K_3\) quadritangent to \(D_6\) whose points of contact lie on another cubic likewise adjoint to \(D_6\). And for this one will start from an arbitrary cubic \(K_3\) and choose on it two points \(P\) and \(Q\) having the same tangential \(O\); then one will choose a group of 4 points \(A_1A_2A_3A_4\), the section of \(K_3\) with another cubic \(K'_3\) passing through \(O\) and tangent at \(P\) and \(Q\) to our \(K_3\); one then verifies that there are \(\infty^5\) curves of the \(6^\circ\) order \(D_6\) passing doubly through \(O\) and having at \(P\) and \(Q\) two tacnodes with the respective tacnodal tangents \(p = PO\) and \(q = QO\), and tangent to \(K_3\) at \(A_1\), \(A_2\), \(A_3\), \(A_4\): as the \(K_3\) and the points which we have chosen on it for the preceding construction vary, the \(D_6\) comes to vary in the complete system of sextics endowed with two tacnodes and with a double point at the point of intersection of their tacnodal tangents \(p = PO\) and \(q = QO\), and tangent to \(K_3\) at \(A_1\), \(A_2\), \(A_3\), \(A_4\). One deduces that, given a priori such a \(D_6\), there exist cubics \(K_3\) adjoint and quadritangent to it,
% scan: s136L
% !PRINTED 246
\weberpage{246}
which are images of two distinct elliptic curves on the surface $F$.

This established, consider on $F$ the curves of genus 2 which have as image the lines through one of the two tacnodes of $D_6$, say through $P$, which we may suppose distinct from $O$; by adding to these curves of genus two a bisecant elliptic curve corresponding to our $K_3$, one will define an $\infty^3$ linear system of genus 4 and degree 6 which leads to representing $F$ on the sextic $F_6$ which passes doubly through the edges of a tetrahedron. The images of the plane sections of this $F_6$ will be, on the double plane, $\infty^3$ curves of the $6^{\mathrm{o}}$ order $K_6$, having in common with $D_6$ the double points $O$, $P$, $Q$, and $P'$, $Q'$ infinitely near to the latter; having moreover a variable double point, and being 8-tangent to the $D_6$. %

The meaning of the construction carried out on the double plane can also be explained by saying that: on the surface $F$, which we have recognized to contain a simple quadrilateral of isolated elliptic curves, $C_1+C_2+C'_1+C'_2$, there exist two other isolated elliptic curves, $C_3$ and $C'_3$ (corresponding to the above-named $K_3$), incident to the sides of the quadrilateral. Three elliptic curves incident two by two, like $C_1$, $C_2$ and $C_3$, constitute a curve of genus 4 belonging to an $\infty^3$ linear system, which leads to transforming $F$ into an $F_6$ in which the said curves have as images three double lines, sides of a trilateral, and their adjoints, $C'_1 C'_2 C'_3$, likewise give rise to three double lines, edges of a tetrahedron of which the trilateral forms part, passing through the vertex opposite to its face.

Summing up the result obtained, we shall state the theorem\origfootnote{1}{s136L-1}{F. \textsc{Enriques}, \textit{Sopra le superficie algebriche di bigenere uno.} Memorie Soc. It. delle Scienze, detta dei XL, 1906.}:

\textit{The regular surfaces of genus $p=0$ and $P=1$, with bicanonical curve of order zero — such as are characterized by the value $P_3=0$ of the trigenus (or $P_6=1$ of the sextigenus) — can be transformed into sextics $F_6$ passing doubly through the edges of a tetrahedron.}

\textit{Remark.} – It is useful to observe some particular, noteworthy cases of the sextic $F_6$. %

First of all it may happen that the 4 planes of the tetrahedron, whose edges are double lines of $F_6$, come to pass through a point. The tetrahedron then reduces to a polyhedral angle, whose vertex becomes quadruple for $F_6$. This case was observed by \textsc{Castelnuovo}. By projection from the quadruple point, $F_6$ is represented on a double plane with branch curve $D_{10}$, which, by a quadratic transformation, can be brought back to the $D_8$ whose type was defined above.

% scan: s136R
% !PRINTED 247
\weberpage{247}
Another possible degeneration of the tetrahedron is obtained by imagining that two opposite edges become infinitely near skew lines, and that at the same time the pairs of faces through these lines also come to coincide. One thus obtains an $F_6$ which possesses a triple line and a double line infinitely near to it and skew to it, as well as two other double lines, skew to each other, incident to the triple line. This $F_6$ corresponds to a double plane with a branch curve $D_8=p+q+D_6$, where one of the tacnodes of the $D_6$ is infinitely near to the double point $O$.

\edsection{2. Surfaces with all genera equal to 1.}{2. Surfaces with all genera equal to 1.}{sec.VII.2}

We have seen that the regular surfaces of genus $p=1$ with canonical curve of order zero are characterized, as against others likewise having linear genus $p^{(1)}=1$, by their bigenus being $P=1$ instead of $P>1$. All the further plurigenera of the surface also turn out to be $P_i=1$.

Taking as model a surface $F$ free of exceptional curves, every linear system $|C|$ of genus $\pi$ without base points on it has degree $n=2\pi-2$ and is the adjoint of itself, so that for an irreducible $|C|$ the characteristic series is the canonical series.

\textit{On the surface $F$ $(p=P=1)$ a curve of virtual genus $\pi\geq0$, irreducible, or reducible provided it is connected, belongs to a complete linear system of dimension $r=\pi$} (cf. Chap. IV, § 17).

On the other hand, a disconnected curve belongs, in general, to a linear system of dimension greater than its genus $\pi$; in particular: \textit{the groups of $s$ $(>1)$ elliptic curves of a pencil lying on $F$ belong to a (reducible) linear system of genus 1 and of dimension $s$.} The linear systems containing as fixed part $\theta$ a curve of genus zero are also superabundant.

Let us note further that: if to a linear system $|C|$ without fixed parts lying on $F$ one adds a fixed curve $\theta$ intersecting the variable components of the $C$ in more than one point, the linear system in general is enlarged. The only case in which this enlargement does not take place is that in which one adds to the groups of $s$ elliptic curves of a pencil a curve of genus zero unisecant to such curves.

To classify our surfaces of genus 1 $(p=P=1)$, we shall take into consideration on a given $F$ the irreducible linear systems of minimum genus $\pi>1$ which belong to it. We may suppose first of all that the genus of the said minimum system $|C|$ is $\pi=2$, and then $\pi=3$, $\pi=4$, \ldots.

The hypothesis $\pi=2$ entails that the degree of the system is $n=2$ and its dimension $r=2$, so that $F$ will be represented on a double
% scan: s137L
% !PRINTED 248
\weberpage{248}
plane with branch sextic $D_6$. And in fact the double plane with branch curve $D_6$ has genus %
\[
p=p_a=p_g=1,
\]
and bigenus $P=1$, there being on it only one canonical curve, of order zero.

For $\pi=3$ one is led to the general surface of the $4^{\mathrm{o}}$ order $F_4$, which indeed has genus $p=1$ and possesses a canonical curve of order zero.

For $\pi=4$ we shall be led to consider a surface $F_6$ of $S_4$, such as is obtained as the complete intersection of a quadric $Q_2$ and a cubic variety $V_3$. Indeed the $\infty^{14}$ quadrics of $S_4$ cut out on $F_6$ the double system of the hyperplane sections, which is of genus
\[
4+4+6-1=13,
\]
and hence also of dimension 13; therefore there is a quadric $Q_2$ which contains $F_6$. Similarly one proves that through $F_6$ there pass $\infty^5$ irreducible cubic varieties $V_3$, so that $F_6$ turns out to be defined as the complete intersection of a $Q_2$ with a $V_3$. On the other hand it is clear that the intersection surface so defined has as hyperplane sections canonical curves, whose linear system therefore turns out to be an $\infty^4$ system of genus 4, adjoint to itself, whence it follows %
\[
p=p_a=p_g=1,\qquad P=1.
\]

For $\pi=5$ one is led to a surface $F_8$ of order 8 of $S_5$, which one easily sees to be the base surface of a net of quadrics.

It is easy to prolong the series of surfaces of genus 1 which we have constructed for the first values of $\pi$, e.g. by considering the intersections of quadrics and of rational three-dimensional varieties with elliptic section curves. But, without dwelling on these examples, let us answer the general question which presents itself here:

\textit{The regular surfaces of genus $p=1$ and bigenus $P=1$} (surfaces endowed with a canonical curve of order zero) \textit{can be classified on the basis of the minimum value of the genus $\pi$ $(>1)$ of the irreducible curves which belong to them, and they are therefore distributed into infinitely many families having as type surfaces $F_n$ of order $n=2\pi-2$, with canonical sections of genus $\pi$, belonging to the space $S_\pi$.}

\textit{There exist families of surfaces $F_n$ corresponding to all the values of $\pi=2,\,3,\,4,\,5,\,\ldots\ldotp$, each family depending exactly on 19 projective invariants which express as many moduli of them with respect to birational transformations. Surfaces corresponding to different values of $n$ or of $\pi$ are generally not transformable into one another by birational transformations.}

% scan: s137R
% !PRINTED 249
\weberpage{249}
The actual existence of surfaces $F_n$ for all integer values of $\pi$, or for the even values of $n=2\pi-2$, is recognized by observing that there exist at least $\infty^{18}$ particular surfaces $F_n$ which are obtained as the projection of an $F_{n+2}$ endowed with a double point. To prove this it suffices to observe that there exist, among the $\infty^{19}$ $F_n$, $\infty^{18}$ particular $F_n$ which contain a conic. For the moment let us justify this assertion by a simple count of constants: a priori the existence of a hyperplane quadritangent to $F_n$ will impose $x-\pi+4$ conditions, if the existence of one of these entails as a consequence the existence of $\infty^x$ similar hyperplanes. But if the section of a hyperplane quadritangent to $F_n$ splits into a conic and a residual curve of order $n-2$ (which must have 4 points in common with the conic), the existence of the aforesaid hyperplane will bring with it that of an $\infty^{\pi-3}$ system of similar hyperplanes, such as are the hyperplanes through the plane of the conic, and one will therefore have $x=\pi-3$. Therefore, since $x-\pi+4=1$, there will not in general be hyperplanes quadritangent to $F_n$ intersecting the surface in a split curve of which a conic forms part, and the surfaces for which such hyperplanes exist, that is, the $F_n$ containing a conic, will have to satisfy one condition.

Now, starting from an $F_n$ containing a conic $C$, consider on it the system that is the sum of $C$ and of the hyperplane sections; this will be a system of genus $\pi+1$ and dimension $\pi+1$, with respect to which $C$ is a fundamental curve; therefore this system will lead to the existence of an $F_{n+2}$ of $S_{\pi+1}$ with canonical section curves, endowed with a double point $O$ corresponding to $C$; this surface is projected from $O$ into the $F_n$ from which we started, or into a surface projective to it\origfootnote{1}{s137R-1}{As one sees, the preceding reasoning is suggested by the attempt to construct the $F_n$ with $n=2\pi-2$ inductively (passing from $\pi$ to $\pi+1$) by means of the constructive procedure of the multiple planes of \textsc{Chisini} (cf. V, 10). But the delicate question which one meets along that path (concerning the choice of a suitable conic sextitangent to the branch curve) is resolved here by the use of the theorem of \textsc{Riemann-Roch} for surfaces with genera one.}.

But the family of surfaces $F_{n+2}$ endowed with a double point, depending only on 18 absolute invariants, will be contained in a family of $F_{n+2}$ depending on 19 invariants, and therefore in general free of double points.

Indeed the existence of this more general family of $F_n$, that is, the related count of moduli, is justified (as in § 11 of chap. V) by representing the surface itself on an $n$-fold plane and calculating the dimension of the system of curves (with a given number of nodes and cusps) defined by the branch curve $C$. Here, since we are dealing with an $F_n$ endowed with a double point, the $C$ will have a double point to

% scan: s138L
% !PRINTED 250
\weberpage{250}
regard as virtually non-existent, which, as one understands a priori, must disappear for the most general curve of the system indicated above.

Nevertheless our assertions give rise to some critical doubt of a delicate nature. It is necessary to prove that:

$1^{\mathrm{o}}$) the $F_n$, which (according to § 11 of Chap. V) possess at least 19 moduli, possess 19 moduli and no more; %

$2^{\mathrm{o}}$) the $F_n$ (or the $F_{n+2}$) containing a conic actually possess 18 moduli and no more, that is, that the condition required for the possession of the conic is not identically satisfied for the most general $F_n$ of its family.

These two assertions are justified on the basis of the count of the projective invariants pertaining to a rational surface with canonical sectional curves of genus $\pi$, endowed with a triple point, belonging to the space $S_\pi$. Such surfaces can be represented on the plane by means of a linear system $\infty^\pi$, of genus $\pi$ and of degree $2\pi-2$, possessing a fundamental cubic on which there must lie 12 base points of the system, constituting a $G_{12}$\origfootnote{1}{s138L-1}{The theorem to which we appeal here is found in P. \textsc{Du Val}, \textit{On rational surfaces whose prime sections are canonical curves.} (Proc. of the Math. Society of London, 1932). It is also found virtually established in the exposition of \textsc{Enriques-Conforto}, \textit{Le superficie razionali}, Book I, § 41, p. 184 and § 46, p. 220. In this last place (cf. also § 46, p. 223) it is proved that a linear system $|C|$ can be transformed so that it possesses no fundamental curves splitting off from the adjoint system, or else into a linear system consisting of curves of a certain order $n$ endowed with two base points of multiplicities $r$ and $s$, with $n=r+s$; or else of curves of a certain order $n$, with an $(n-r)$-fold base point and other $r$-fold base points infinitely near to it in distinct directions. In these last two cases the system $|C|$, supposed of genus $\pi$, has degree $n>2\pi-2$; therefore the systems $|C|$ with canonical characteristic series all fall under the statement of p. 184.}; one therefore has exactly 12 absolute invariants, corresponding to the 11 independent points of the $G_{12}$ and to the absolute invariant of the cubic. Since the imposition of a triple point on a surface $F_n$ entails 6 conditions (and no more), one deduces that the $F_n$ of $S_\pi$ cannot form a family depending on more than 19 moduli. The same consideration shows that the general $F_{n+2}$, depending on 19 invariants, cannot possess a double point, since, by requiring this point to become triple, one thereby subjects the surface to 6 conditions and no more. From this it also follows that the $F_n$ possessing a conic, arising as they do from the projection of an $F_{n+2}$ with a double point, depend on 18 and not on 19 moduli. %

It remains to justify the assertion that the surfaces $F_n$ $(n=2\pi-2)$ corresponding to different values of $\pi$ $(>1)$ are in general birationally inequivalent to one another. This can be done inductively on the basis of a simple
% scan: s138R
% !PRINTED 251
\weberpage{251}
count of constants, like the one we made use of earlier to show that an $F_n$ does not in general possess a conic. By a similar count one indeed recognizes that « \textit{on a surface $F_n$ of $S_\pi$} (corresponding to the minimum value $\pi$ of the genus of the curves lying on it) \textit{there are in general no linear systems of curves other than the system of hyperplane sections and its multiples} ». %

First of all one asks whether, among the $\infty^\pi$ hyperplanes of $S_\pi$, there can be hyperplanes $m$-tangent to $F_n$ such that their sections with $F_n$ split into two curves of genera $\pi_1$ and $\pi_2$, with $m$ common points: $\pi=\pi_1+\pi_2+m-1$. If such a hyperplane exists, there exist $\infty^x$ of them, with $x=\pi_1+\pi_2$; therefore the number of parameters on which the hyperplanes sought depend will be $\pi-m-x=-1$: this means that the $F_n$ possessing such split curves are subject to one condition. In a similar way one recognizes that the section of an $F_n$ with a hypersurface of any given order cannot in general split except into two curves which are in their turn complete intersections of two hypersurfaces; in this last case the count of constants again gives one condition to be satisfied by the surface $F_n$, but it is obvious that this condition results here as a consequence of the fact that the $F_n$ contains the system of hyperplane sections, and is therefore identically satisfied.

Now it must be pointed out that, as has already been said above (in dealing with the $F_n$ that contain a conic), the result to which our count of constants leads leaves standing a doubt which already deprives of value Noether’s proof that the surface of the $4^{\mathrm{o}}$ order of ordinary space contains in general only curves that are complete intersections.

Nevertheless the matter can be established rigorously, and also simply, on the basis of the principle of continuity, by reducing an $F_n$ to the limiting case of the rational surface endowed with a triple point. We shall confine ourselves to indicating briefly how this proof can be carried out in the elementary case of Noether’s theorem on surfaces of the $4^{\mathrm{o}}$ order: the extension of the reasoning to the $F_n$ of genus one presents no difficulty.

Here, then, a quartic surface $F_4$ is made to vary continuously until it reduces to a general monoid $\Phi_4$ possessing a triple point $O$: the monoid, which contains 12 lines $a_1a_2a_3\ldots\ldotp a_{12}$ through $O$, can be represented on the plane, by projection from $O$, by means of a linear system of curves of the $4^{\mathrm{o}}$ order $C_4$ passing through the 12 base points $A_1 A_2\ldots\ldotp A_{12}$, traces of the lines $a_i$: these points lie on a fundamental cubic $C_3$ of genus one.

The proof that we propose to carry out aims at establishing that: while the $F_4$ tends continuously to the general monoid $\Phi_4$, a complete irreducible linear system of curves $C$, traced on
% scan: s139L
% !PRINTED 252
\weberpage{252}
$F_4$, reduces in the limit to a linear system of curves $\overline{C}$, complete intersections of $\Phi_4$ with surfaces of a certain order $d$ which do not contain its triple point $O$. Then, denoting by $\pi$ the genus of the $\overline{C}$, and hence also of the $C$, and $2\pi-2$ being therefore the degree and $\pi$ the dimension of $|C|$, it will follow that the said $C$ on $F_4$ must have $2\pi-2$ common points with the intersections of the surfaces $F_d$ of order $d$, which are likewise of genus $\pi$ and of degree $2\pi-2$; since a $C$ cannot contain a series $g_{2\pi-2}$ of dimension $\pi$, one deduces that the system $|C|$ on $F_4$ coincides with the linear system of the complete intersections of the $F_d$.

To develop the proof thus indicated, we observe that, when the $F_4$ reduces continuously to $\Phi_4$, a system $|C|$ of genus $\pi$ and degree $2\pi-2$, lying on $F_4$, will a priori reduce to a system $\infty^\pi$ of \textit{connected} limit curves formed by the intersections of $\Phi_4$ with surfaces passing through $O$ a certain number of times, augmented by the lines $a_i$, counted respectively $h_i$ times $(h_i\geq0)$, and by the neighborhood of $O$ counted $s$ $(\geq0)$ times. Hence the images of such curves on the representative plane of $\Phi_4$ will be, a priori, curves of a certain order $n$, belonging to a complete linear system
\[
|C_n|=C_m(A_1^{r_1} A_2^{r_2}\ldots A_{12}^{r_{12}})+sC_3+\Sigma h_iA_i,
\]
where the $C_m$ (of order $m\leq n$) belong to an irreducible linear system $\infty^\pi$. And since the $C_m$ have genus $\pi$ and degree $2\pi-2$, we shall have
\[
\begin{aligned}
n^2-\Sigma(r_i+s-h_i)^2&=2\pi-2\\
n(n-3)-\Sigma(r_i+s-h_i)(r_i+s-h_i-1)&=2\pi-2,
\end{aligned}
\]
and hence
\[
3n=\Sigma(r_i+s-h_i):
\]

This means that for the system $|C_n|$ the cubic $C_3$ is \textit{virtually fundamental}, that is, it has zero intersections with the $C_n$.

It is convenient to discuss first of all the hypothesis:

1) the $C_3$ is also effectively fundamental for the complete system $|C_m|$, that is %
\[
|C_n|=|C_m|,\qquad s=0,\qquad h_i=0.
\]

Under this hypothesis we easily prove that $|C_n|$ is a multiple of the system $C_4$ $(A_1A_2\ldots\ldotp A_{12})$, and therefore the corresponding system on $\Phi_4$ consists of the complete intersections of this monoid with the surfaces $F_d$ of order $d=\dfrac{n}{4}$.

Indeed, denoting by $G$ the group cut out on the elliptic cubic
% scan: s139R
% !PRINTED 253
\weberpage{253}
$C_3$ by a line, we shall have on the cubic the following equivalences:
\[
\begin{aligned}
\Sigma r_iA_i&\equiv nG\\
\Sigma A_i&\equiv4G;
\end{aligned}
\]
multiplying the second relation by $r_1$ and subtracting the first, one finds:
\[
\mathop{\Sigma}\limits_{i=2}^{i=12}(r_i-r_1)A_i\equiv(n-4r_1)G\,;
\]
but this last relation cannot hold, the group of the 11 points $A_2A_3\ldots\ldotp A_{12}$ being arbitrarily chosen, unless $n=4r_1$ and
\[
r_1=r_2=\ldots\ldotp=r_{12}=r;
\]
so that the $C_n$ are curves of order $4r$ belonging to the system $|rC_4|$, \hfill Q.E.D.

If, instead, hypothesis 1) is not satisfied, we shall have two possibilities, which have to be excluded.

2) First, that the $C_3$ does not form part of the $C_n$, but only the neighborhood of some point $A_i$:
\[
s=0,\qquad h_i>0.
\]

Now if the $C_3$ is fundamental for $|C_m|$, this irreducible linear system is a multiple, according to a certain number $r=r_1=r_2=\ldots\ldotp=r_{12}$, of $C$, and — if it is not to be enlarged when the neighborhoods of the points $A_i$ are added to it — the point to be added must be a single one, for example $A_1$, and it must have multiplicity 1 for $C_m$, hence $r=1$. But it is obvious that a curve $C$ passing through the 12 points $A_i$, with one of these added to it, cannot present itself as the limit of a curve traced on $F_4$ when this surface reduces to the monoid $\Phi_4$. %

If, instead, one supposes that the $C_3$ is not fundamental for $|C_m|$, this linear system is regular (and not superabundant), hence it is always enlarged if the neighborhood of some point $A_i$ is added to it. And this goes against the hypothesis that the $|C_n|$ is complete.

3) Finally, it is convenient to discuss the hypothesis that the cubic $C_3$ appears a certain number $s>0$ of times as a fixed part of the $C_n$. In this case the $|C_m|$, which does not possess the $C_3$ as a fundamental curve, is regular, and hence it is enlarged if one adds to it the neighborhood of a point $A_i$ which is not a base point for it: thus only a certain number $t$ of points $A_i$ can figure among the points constituting the sum $\Sigma h_iA_i$, and each of them $h_i\leq s$ times. Denoting by $l=C_mC_3$
% scan: s140L
% !PRINTED 254
\weberpage{254}
the number of variable intersections of a $C_m$ with the cubic $C_3$, we shall now have
\[
l+ts-3s\geq0.
\]

Hence one will be led to distinguish two cases possible a priori: either $l+t\geq4$, from which inequality there follows the enlargement of the linear system $C_m+C_3+\Sigma A_i$; or $l+t<4$, which entails $s\leq1$; in this last case the dimension of $|C_m|$ would turn out to be 1 less than the genus of $|C_n|$: and therefore $|C_n|$ could not be the limit system of a linear system lying on the $F_4$, whose dimension equals the genus.

Hypotheses 2) and 3) being thus excluded, there remains only hypothesis 1), from which, as we have said, one infers that the curves lying on a general $F$ are only its complete intersections with other surfaces of order $d=1,\,2,\,\ldots$.

\textit{Remark.} – As has been said, the preceding reasoning extends to the case of the surfaces $F$ of genus 1. However, one thus arrives at recognizing, not that all the curves lying on a general $F$ with 19 moduli are intersections of it with hypersurfaces, but rather that all these curves, and in particular the hyperplane sections of $F$, when they are not the irreducible curves of minimum genus lying on $F$, constitute multiple systems of one and the same system of curves of genus $\varrho$ and of degree $\nu$; one will therefore have
\[
n=s\nu=2\pi-2=s(2\varrho-2).
\]

Conversely, going back from a rational surface with a triple point to a surface of genus one with canonical hyperplane sections, one recognizes (at least inductively) that when
\[
n=s\nu\qquad \text{and}\qquad 2\pi-2=s(2\varrho-2),
\]
there exist at least two families of surfaces $F_n$ of genus 1: an $s$-fold multiple surface and a simple surface, both depending on 19 moduli and birationally inequivalent to each other.

\edsection{3. Historical notice.}{3. Historical notice.}{sec.VII.3}

The first examples of surfaces $F_n$ of genus 1 containing the same number of moduli as the surfaces of the $4^{\mathrm{o}}$ order, and birationally inequivalent to them, were indicated by \textsc{Castelnuovo}\origfootnote{1}{s140L-1}{Cf. F. \textsc{Enriques}, \textit{Ricerche etc.}, III, 6.} (1893). In 1896, \textsc{Enriques}\origfootnote{2}{s140L-2}{F. \textsc{Enriques}, \textit{Sui piani doppi di genere uno.} Memorie della Soc. It. d. Scienze (detta dei XL), 1896.} characterized, from the point of view of birational transformations, the family of these regular surfaces of genus 1 with ca%
% scan: s140R
% !PRINTED 255
\unskip\weberpage{255}%
nonical curve of order zero, by means of the value of the bigenus $P=1$. Later \textsc{Enriques}\origfootnote{1}{s140R-1}{F. \textsc{Enriques}, \textit{Le superficie di genere uno}, Rendic. Acc. di Bologna, 1908.} and \textsc{Severi}\origfootnote{2}{s140R-2}{F. \textsc{Severi}, \textit{Le superficie algebriche con curva canonica d’ordine zero.} Atti Ist. Veneto (vol. LXV, 2), 1909.}, independently of each other, succeeded in proving the existence of families of surfaces $F_n$ of $S_\pi$, of even order $n=2\pi-2$, for all values of $\pi$. Both took as their starting point the count of moduli made by \textsc{Enriques} (1908 Cf. Chap. V, § 11). Since the $F_n$ depend on 19 moduli (at least), the existence of a family of $F_n$ with that many moduli follows, for \textsc{Enriques}, from the fact that there are particular hyperelliptic $F_n$, with 3 moduli, endowed with 16 double points: from this one also argues that the number of moduli belonging to the $F_n$ must be 19 and not greater; the proof of this requires establishing that the possession of 16 double points characterizes the hyperelliptic $F_n$, as is known to happen for $n=4$ (Kummer surface).

\textsc{Severi} recognizes by algebro-geometric means the existence of the $F_n$ $(n=2\pi-2)$ for any value of $\pi$, by constructing the particular $F_n$ with 4 double points, which are obtained as transforms of quartics $F_4$ possessing 4 curves of genus (effective and virtual) zero, without common points. Admitting that the possession of such curves gives rise, for the $F_4$, to 4 independent conditions (which implies Noether’s theorem that « the curves lying on a general $F_4$ are complete intersections » and something more), he deduces from this that the general $F_n$ cannot contain more than 19 moduli, and hence depend on exactly 19 moduli.

It is appropriate to note here that the construction of $F_n$ can be carried out starting from quartics or from other particular surfaces with genera one, in many different ways, which however require some care, as is hinted at in the « Lezioni » of Enriques-Campedelli (p. 311). For example, one obtains $F_n$ with $n=2\pi-2$, where $\pi$ is of the form $\pi=2+3s$, starting from a double plane with a branch sextic which is the envelope of a series $\infty^1$ of cubics: for this purpose it suffices to add to the curves $C$, curves of genus two which are images of the lines of the plane, $s$ elliptic curves $K$, images of the said cubics, thus constructing the system $C+sK$, which is precisely of genus $\pi=2+3s$\origfootnote{3}{s140R-3}{If $\pi$ is of the form $\pi=3+3s$, one will proceed similarly starting from an $F_4$ containing a line; and if $\pi=4+3s$, one will instead start from a particular $F_6$ containing a plane cubic.}.

Since one starts from particular surfaces having exactly 18 moduli, and to these correspond all the $F_n$ containing a plane cubic, which are therefore subject to \textit{one} condition (and no more), one thus deduces that the $F_n$ depend on precisely 19 moduli, and not on more.

% scan: s141L
% !PRINTED 256
\weberpage{256}
Meanwhile the questions under discussion have been clarified by \textsc{A. Franchetta}\origfootnote{1}{s141L-1}{\textsc{A. Franchetta}, « \textit{Sulle superficie regolari di genere uno con curva canonica d’ordine zero} ». Istituto Lombardo, 1942.}, in the way that we have set out in the text. The author, having sought to connect the proof of the existence of the infinite series of the $F_n$ with Chisini’s existence theorem, is led to give this proof the simplest form, passing from the particular $F_n$ containing a conic to the $F_{n+2}$. \textsc{Franchetta} then proves that the $F_n$ (for a given value of $n$) cannot depend on more than 19 moduli, by having recourse to the rational $F_n$ with a triple point: precisely as was seen above.

That the $F_n$ corresponding to different values of $n$ are in general birationally inequivalent to one another depends, as we have said, on the theorem that « on a general $F_n$ there are no curves other than complete intersections with hypersurfaces ». One easily acquires a notion of the fact thus stated on the basis of Noether’s theorem concerning the quartics $F_4$, and of the count of constants on which it is founded. But with this one does not yet have a rigorous proof. \textsc{Enriques} and \textsc{Severi} thought that such a proof could be given (for $n\geq 4$) by making a general $F_n$ reduce continuously to a hyperelliptic $F_n$ with generic moduli, making use of the property, recognized by transcendental means\origfootnote{2}{s141L-2}{\textsc{Enriques} and \textsc{Severi}, \textit{Mémoire sur les surfaces hyperelliptiques.} Acta Math., vol. XXXII, no. 27.}, that on them there are no curves other than the complete intersections with hypersurfaces and the curves of contact with tangent hypersurfaces. And \textsc{Severi} has carried out this plan in an elegant way (there remains in his argument a small gap, which is easily filled). The algebro-geometric proof offered here, which proceeds by reduction to the limiting case of the rational $F_n$ with a triple point, is due to \textsc{A. Franchetta}, in his note already cited above.

Here we shall give some bibliographical indications concerning the various proofs of Noether’s theorem on the surfaces of the $4^\circ$ order and on the general surfaces of higher order.

\textsc{M. Noether}, « Zur Grundlegung der algebraischen Raumcurven ». Abhandlungen der Akademie der Wissenschaften, Berlin, 1883.

\textsc{K. Rohn}, « Die Raumcurven auf der Flaechen vierter Ordnung ». Berichten der koenigl. Saechs. Gesellschaft der Wissenschaften, Leipzig, 1897. %

\textsc{F. Severi}, « Le superficie algebriche con curva canonica d’ordine zero ». Atti Ist. Veneto (LXV, 2), 1909.

\textsc{G. Fano}, « Sulle varietà algebriche che sono intersezioni complete di più forme ». R. Acc. di Torino, vol. XLIV, 1909.

% scan: s141R
% !PRINTED 257
\weberpage{257}
\textsc{S. Lefschetz}, « On certain numerical Invariants of algebraics Varieties. (Trans. of Am. Math. Society, vol. 22, 1921); « Concerning Noether’s theorem on curves traced on non singular surfaces ». Bull. Am. Math. Society, vol 29, 1923. %

\textsc{A. Franchetta}, « Sulle superficie regolari di genere uno con curva canonica d’ordine zero ». Istituto lombardo, 18 June 1942.

\textsc{B. D’Orgeval}, « Sur les surfaces algébriques dont tous les genres sont 1 ». Gauthier-Villars, Paris, 1945.

To anyone pursuing a complete classification of the surfaces with genera 1, there now presents itself a new question, namely whether the surfaces $F_n$ with $n=2\pi-2$, with canonical sectional curves in an $S_\pi$, form, within this space, a single family, or several different families, each depending on 19 absolute invariants or moduli. In this connection, in agreement with the final remark of the preceding section, one may note that already in the space $S_5$ there are 2 types of $F_8$, namely the simple surface $F_8$, intersection of three quadrics, and the double Veronese surface $F_4$, with branch curve of order 12, which is the transform of a double plane with a branch sextic. Analogously, in the space $S_9$ one finds two birationally inequivalent types of $F_{16}$, one of which is derived from a surface of the $4^\circ$ order of $S_3$. From this one is led to recognize in $S_6$ two types of particular $F_{10}$ endowed with a double point, which are projected from this point respectively into the $F_8$ and into the double $F_4$ indicated above. But \textsc{P. du Val}\origfootnote{1}{s141R-1}{\textsc{P. du Val}, \textit{Superficie di genere uno che non sono base per un sistema di quadriche} and \textit{Osservazioni sulle superficie di genere uno che non sono base per un sistema di quadriche.} Rend. R. Acc. dei Lincei, February-March, 1932.} has shown that these two $F_{10}$ belong to one and the same family of general $F_{10}$, without double points; and this is in agreement with the general result established by \textsc{B. d’Orgeval}, in his volume already cited, who has indeed proved that the surfaces of $S_\pi$ with canonical sectional curves form a single family. %

\edsection{4. Surfaces of linear genus $p^{(1)}=1$ with $P>1$.}{4. Surfaces of linear genus p\textasciicircum{}(1)=1 with P>1.}{sec.VII.4}

Let there be given a surface $F$, of (absolute) linear genus $p^{(1)}=1$ and of surface genus $p>1$. Its $\infty^{p-1}$ canonical curves $K$ must form a linear system of elliptic curves, of degree $p^{(1)}-1=0$, which will therefore be composed of the elliptic curves of a pencil (without base points). The same may be said of the bicanonical or $i$-canonical curves of $F$, which must form a linear system of genus
\[
\frac{i(i+1)}{2}(p^{(1)}-1)+1=1,
\]
\printersignature{\textsc{Enriques F.} - \textit{Superficie algebriche.}\quad 17}
% scan: s142L
% !PRINTED 258
\weberpage{258}
and of degree
\[
i^2(p^{(1)}-1)=0.
\]

Therefore: \textit{the canonical or pluricanonical curves of a surface of linear genus $p^{(1)}=1$ are in general reducible, being composed of the elliptic curves of a pencil.}

The existence of a pencil of elliptic curves, and of a linear pencil for regular surfaces, therefore follows as a consequence of the fact that $p^{(1)}=1$, whenever the genus of the surface is $p>1$ or one of its plurigenera $P_i$ turns out to be
\[
P_i>1.
\]

But this inequality will always hold for regular, non-rational surfaces possessing some canonical or bicanonical curve of order greater than zero, since under this hypothesis it has been seen (Chap. VI, \S\ 2) that for $p>1$ one has
\[
P\geq 2,
\]
and for $p=0$:
\[
P\geq 1,\qquad P_6\geq 2.
\]

Since, on the other hand, it has also been found that the surfaces for which $p=0$ and $P=1$, with bicanonical curve of order zero ($P_6=1$), always possess linear pencils of elliptic curves, it is legitimate to assert that:

\textit{The regular (non-rational) surfaces of (absolute) linear genus $p^{(1)}=1$ possess a linear pencil of elliptic curves, with the exception of those which have all their genera equal to $1$} (surfaces $F_n$ with $n=2\pi-2$, with canonical sections in $S_\pi$, characterized by $p=P=1$). %

To classify the surfaces $F$ of linear genus $p^{(1)}=1$, with $P>1$, containing a linear pencil of elliptic curves $C$, it is convenient first of all to point out that they possess \textit{two invariant characters}, capable of assuming \textit{arbitrary integer values}, namely:

1) the genus $p$ ($=p_a$), which is defined by the number $S$ of the curves $C$ of the pencil endowed with a double point, the Zeuthen-Segre invariant being $I=S-4=12p+8\ (=12p_a-p^{(1)}+9)$,

2) and the \textit{determinant} $d$, which is defined as the minimum number of points that can be determined rationally on a $C$, as a function of the parameter on which the $C$ itself depends.

To understand the meaning of this last character it is convenient to recall that, in general, on an elliptic curve $C$ of order $n$, in the plane or in any space $S_r$, one cannot construct \textit{rationally} (as functions of the coefficients of its equation) any series other than the multiples of the $g_n$ cut out by the lines of the plane or by the hyperplanes of $S_r$. And thus it is not possible in general to lower the order
% scan: s142R
% !PRINTED 259
\weberpage{259}
$n$ of a variable elliptic $C$ by operating rationally on the parameters on which its variation depends\origfootnote{1}{s142R-1}{Cf. \textsc{F. Enriques}, \textit{Sulle superficie algebriche con un fascio di curve ellittiche.} Rendic. Acc. Lincei, 1912. In this note the deduction of no. 3 is vitiated by an error of calculation. Cf. also \textsc{Enriques-Chisini}, \textit{Lezioni}, vol. III, p. 345.}.

\edFlush
Therefore a series $\infty^1$ of elliptic curves $C$, of order $n$, taken in space, will in general form a surface $F$ with a pencil of $C$ which will not be susceptible of any transformation lowering the order of the $C$.

Now we can say that: the surfaces $F$ with a pencil of elliptic curves $C$, of \textit{determinant} $d>1$, can be transformed so that the $C$ become \textit{curves of order} $d$, and thus --- under the hypothesis of regularity --- into surfaces $F_m$ of a certain order $m$, of ordinary space, with an $(m-d)$-fold line (for $d=2$ into double planes); and no transformation of these lowering the order of the $C$ is possible.

To justify the assertion, note that --- by the terms of the definition --- the $F$ will contain a curve $L$ that is $d$-secant to the curves $C$, and no other curves cutting the $C$ in a smaller number of points. Now, adding a $C$ to the $L$, one will increase its degree by $2d$ and its genus by $d$, so that the dimension of $|L+C|$, given by the Riemann-Roch theorem, will exceed that of $|L|$ by $2d>d>1$, and hence the curves $L+C$, $L+2C$, $\ldotp\ldotp\ldotp\ldotp\ L+sC$ will be contained in larger irreducible linear systems, $d$-secant to the $C$. For $d=2$ one obtains a double plane, with a certain branch curve $D_{2n}$, on which the $C$ are represented by the lines through a $(2n-4)$-fold point of $D_{2n}$. %

To deepen the study of our surfaces of linear genus $p^{(1)}=1$ (regular, non-rational surfaces, with $P>p\geq 1$), it is convenient to study in particular those of \textit{determinant} $d=1$.

A surface $F$ of this family contains a linear pencil of elliptic curves $C$ and a \textit{curve $\theta$ unisecant} to the $C$. By means of a linear system $\infty^3$ contained in $|3\theta+sC|$ one can transform $F$ into a surface $F_m$, of a certain order $m$, on which the $C$ are cubics cut out by the planes through an $(m-3)$-fold line $a$; these $C$, varying in the pencil of the said plane sections, will have a rationally distinguished flex, which may be supposed fixed, at a point $A$, $(m-2)$-fold for $F_m$, lying on $a$; and one may also suppose that $a$ is the corresponding inflectional tangent of the $C$. Hence the $F_m$ is represented, by projection from $A$, on a double plane with branch curve $D_{2n}$, of a certain even order $2n$, endowed with a $(2n-3)$-fold point $O$ (trace of $a$). The $D_{2n}$ may have other singular points lowering the genus of the double plane: 4-fold points or $[3,3]$ points. But the presence of such singularities in general allows one to lower the order of $D_{2n}$ by quadratic transformations
% scan: s143L
% !PRINTED 260
\weberpage{260}
of the plane. If one wants the $D_{2n}$ \textit{to be of minimum order}, one must assume that it belongs to one of the two following types:

I) $D_{2n}$ with a $(2n-3)$-fold point $O$, without other singularities affecting the genus; and

II) $D_{2n}$ with a $(2n-3)$-fold point $O$ and a 4-fold point $B$, made up of the line $OB$ and of a $D_{2n-1}$ with $O$ $(2n-4)$-fold and $B$ 3-fold.

The double planes of type I are of genus
\[
p=n-2,
\]
there being $\infty^{p-1}$ canonical curves represented by the groups of $p-1$ lines through $O$. The $i$-canonical curves are given by the groups of $i(p-1)$ lines through $O$, hence
\[
P_i=i(p-1)+1.
\]

For the double planes of type II, there are $\infty^{n-4}$ canonical curves represented by the groups of $p-1=n-4$ lines through $O$, to which is added the exceptional branch line\origfootnote{1}{s143L-1}{On the line $OP$ there are three singular points of $D_{2n}$, namely $P$, $O$ and the point infinitely near to $O$; hence the degree of the branch line on the plane is $-2$, and the corresponding curve on the surface consists of two curves of degree $-1$ (and of genus 0).}, hence the genus %
\[
p=n-3.
\]

Similarly the $i$-genus is
\[
P_i=i(p-1)+1.
\]

Let us conclude for the time being:

\textit{The surfaces $F$ of absolute linear genus $p^{(1)}=1$ and of determinant $d=1$, with $P\geq 1$, can be represented on double planes of type I,} with branch curve $D_{2n}$ endowed with a $(2n-3)$-fold point $O$, \textit{or of type II,} that is, with branch $D_{2n}$ endowed with a $(2n-3)$-fold point $O$ and, outside it, a 4-fold point $B$: these double planes give rise to as many distinct families as there are values of the genus, which is, for type I:
\[
p=n-2
\]
and for type II
\[
p=n-3,
\]
and for them one always has
\[
P_i=i(p-1)+1.
\]

Every family of our $F^1$ contains infinitely many classes of birationally distinct surfaces, each of which is characterized by
% scan: s143R
% !PRINTED 261
\weberpage{261}
a certain number of parameters or moduli, which it is easy to evaluate. The $D_{2n}$ of type I for which the $(2n-3)$-fold point $O$ is fixed depends on
\[
\frac{2n(2n+3)}{2}-\frac{(2n-3)(2n-2)}{2}=8n-3
\]
parameters, and --- since there are $\infty^6$ homographies leaving the point $O$ fixed --- it possesses
\[
8n-9
\]
absolute invariants. Thus the $F^1$ \textit{of type I}, for each value of the genus $p=n-2$, depend on
\[
8p+7
\]
moduli.

Moreover, the $D_{2n}$ with $O$ $(2n-3)$-fold and $B$ 4-fold depend on
\[
8n-13
\]
parameters and possess
\[
8n-13-6=8n-19
\]
invariants with respect to the $\infty^6$ quadratic transformations which carry the two pencils of lines $O$ and $B$ into themselves. Hence, since $p=n-3$, there are
\[
8p+5
\]
moduli.

These results contradict, at first sight, what was established in \S\ 11 of Chap. V, namely that, in general, a class of regular surfaces of genera $p$ and $p^{(1)}$ depends on at least
\[
9p-2p^{(1)}+12
\]
moduli; since for $p^{(1)}=1$ one has
\[
9p-2p^{(1)}+12=9p+10
\]
which exceeds $8p+7$ and $8p+5$.

But the explanation of the paradox lies in this, that \textit{the surfaces $F^1$ of linear genus $p^{(1)}=1$ and of determinant $1$, with $P>1$, of a given genus $p$, are included in families $F^2$, of determinant $2$, having the same genus $p$.}

Indeed the double plane of type I is included as a particular case in the double plane $\mathrm{I}'$ with branch $D_{2n}$ endowed with a $(2n-4)$-fold point $O$, which is likewise of genus
\[
p=n-2,
\]
and depends on
\[
10n-12
\]
invariants, that is, on
\[
10p+8
\]
% scan: s144L
% !PRINTED 262
\weberpage{262}
moduli; and for
\[
n\geq 4,\qquad p>1,
\]
this number is always
\[
\geq 9p+10\qquad (=9p-2p^{(1)}+12).
\]

Similarly the double plane of type II is included, as a particular case, in the double plane $\mathrm{II}'$ with branch $D_{2n}$ endowed with a $(2n-4)$-fold point $O$ and a 4-fold point $B$, which is likewise of genus
\[
p=n-3
\]
and depends on
\[
10n-22
\]
invariants with respect to $\infty^6$ quadratic transformations, that is, on
\[
10p+8
\]
moduli.

If now, after having discussed the surfaces $F^1$ of linear genus $p^{(1)}=1$ and of determinant $d=1$ ($P>1$), we pass to examine the $F^2$ of determinant $d=2$ (with a linear pencil of elliptic curves $C$), we come upon something unexpected: the classification of such surfaces depends not only on the genus $p$, but also on a new character, which is the number $s$ of the curves of the pencil $|C|$ that reduce to double elliptic curves; on the other hand this character can be expressed through the bigenus $P$ by means of the formula
\[
P=2(p-1)+s-1. %
\]

To classify the surfaces $F^2$, let us start from the observation that a surface with a linear pencil of elliptic curves $C$, on which there exist curves bisecant to the $C$, contains an irreducible linear system of the type $(\theta+hC)$ and hence can be represented (on a double rational ruled surface and) on a double plane with branch curve $D_{2n}$, of a certain order $2n$, endowed with a $(2n-4)$-fold point $O$. The $D_{2n}$ may have, in general, other elementary singularities lowering the genus $p$: 4-fold points or $[3,3]$ points. But if one requires that the order of the $D_{2n}$ cannot be lowered by quadratic transformations of the plane, one will have the following \textit{irreducible types}:

$\mathrm{I}'$) $D_{2n}$ endowed with a $(2n-4)$-fold point without other singularities lowering the genus:
\[
p=n-2;
\]

$\mathrm{II}'$) $D_{2n}$ endowed with a $(2n-4)$-fold point and with a 4-fold point
\[
p=n-3;
\]

III) $D_{2n}$ endowed with a $(2n-4)$-fold point $O$ and with $s$ $(\geq 1)$ $[3,3]$ points: $A_1A_2\ldotp\ldotp\ldotp\ldotp A_s$. The $D_{2n}$ contains $s$ lines through $O$ and a remaining
% scan: s144R
% !PRINTED 263
\weberpage{263}
part $D'_{2n-s}$ endowed with $s$ tacnodes for which the said lines are the tacnodal tangents:
\[
p=n-2-s.
\]

\textit{The surfaces $F$ of linear genus $p^{(1)}=1$ and of determinant $d=2$} ($P>p\geq 1$) give rise, for each value of the genus $p$, \textit{not only} to the \textit{double planes of types} $\mathrm{I}'$ and $\mathrm{II}'$ \textit{for which}
\[
p=n-2,\qquad n-3,
\]
and
\[
P=2(p-1), %
\]
but also to \textit{infinitely many families of double planes of type III, depending on another integer}
\[
\begin{aligned}
s&=1,\,2,\,3\ \ldotp\ldotp\ldotp\ldotp: \\
p&=n-2-s,
\end{aligned}
\]
which \textit{can be distinguished by giving}, besides the genus $p$, \textit{the bigenus}
\[
P=2(p-1)+s+1.
\]

Indeed, on a double plane of type III, the canonical curves will be represented by groups of $p-1=n-3$ variable lines through $O$ (double lines, images of irreducible elliptic curves $C$), to which are to be added the $s$ lines $OA_1\ldotp\ldotp\ldotp\ldotp OA_s$ (images of fixed curves $\dfrac{C}{2}$), or rather the neighborhoods of the $[3,3]$ points: $A_1\ldotp\ldotp\ldotp\ldotp A_s$ (cf. Chap. III, \S\ 8, 11 and Chap. V, \S\ 7); and the bicanonical curves will be given by the groups of $2(p-1)+s$ variable lines through $O$. Let us add that for a double plane of type III \textit{the successive plurigenera} are expressed by %
\[
P_{2i}=2i(p-1)+si+1,
\]
and
\[
P_{2i+1}=(2i+1)(p-1)+si+1.
\]

One can also evaluate the number $M$ of \textit{moduli} on which a class of double planes of type \textit{III} depends. Indeed the construction of a $D'_{2n-s}$ with an assigned $(2n-s-4)$-fold point $O$ and given tacnodes $A_1\ldotp\ldotp\ldotp\ldotp A_s$, with tacnodal lines $OA_1\ldotp\ldotp\ldotp\ldotp OA_s$, depends on $5(2n-s)-6-6s$ parameters, %
while the determination of the group of the $s+1$ points $O$, $A_1$, $\ldotp\ldotp\ldotp\ldotp$, $A_s$ brings to light $2s-3$ projective invariants. %
Therefore our $D'_{2n-s}$ and the $D_{2n}$ will depend on $5(2n-s)-6-6s+2(s-3)=10n-9s-12$ invariants or moduli. And since %
\[
p=n-2-s,
\]
we shall have
\[
M=10p+s+8,
\]
or, introducing the bigenus $P$:
\[
M=8p+P+9.
\]

% scan: s145L
% !PRINTED 264
\weberpage{264}
For $p=0$, $P=1$, bicanonical curve of order zero, one finds again the 10 moduli of the typical double plane already encountered, with branch $D_8$.

The results obtained with regard to the classification of the surfaces of linear genus $p^{(1)}=1$ and of determinant 2 lead inductively to stating in general, for the surfaces $F^d$ of any determinant $d$, the following theorem:

\textit{In the classification of the regular surfaces of linear genus $p^{(1)}=1$ ($P>1$) and of determinant $d>1$ (with a linear pencil of elliptic curves $C$ which may be supposed of order $d$) one has to take into account not only the genus $p$, but also other integer characters corresponding to the divisors $r$ of $d$ ($1<r\leq d$), each of which characters, $s_r$, indicates the number of the elliptic curves $\dfrac{C}{r}$ which, counted $r$ times, belong to the pencil $|C|$.}

\textit{The various families of surfaces $F^d$}, which correspond to these characters, \textit{can also be distinguished by the values of the plurigenera}: corresponding to each of the $s_r$ curves $\dfrac{C}{r}$, the plurigenus of order $i$, $P_i$, increases (over $i(p-1)+1$) by $\left[\dfrac{i(r-1)}{r}\right]$, the greatest integer contained in $\dfrac{i(r-1)}{r}$, and hence the $r$-genus $P_r$ increases by $r-1$, and, in general, $P_{ir}$ by $i(r-1)$, etc.

To justify the statement, observe that if, on an $F^d$, there belongs to the pencil of the $C$ an elliptic curve $\dfrac{C}{r}$ which --- taken $r$-fold --- constitutes a $C$, this does not influence the number of the $C$ endowed with a double point, for the purposes of the computation of the Zeuthen-Segre invariant, and therefore it influences neither the arithmetic genus $p_a$ of the surface $F^d$ nor its geometric genus $p_g=p_a$; the $\dfrac{C}{r}$, counted $r-1$ times, comes to be added as a fixed component to the variable parts of the canonical curves $K$ of $F^d$ (composed of the $C$). Indeed, let there be constructed on $F^d$ a net of curves $|L+C|$, containing the pencil of split curves $L+|C|$; it is easy to see that the Jacobian of this net contains $r-1$ times, as a fixed part, every $\dfrac{C}{r}$, and from this it follows that the $\dfrac{(r-1)C}{r}$ will figure, as a fixed component, in the canonical curves $K$.

The problem of the classification of the surfaces of linear genus $p^{(1)}=1$ may be regarded as substantially solved by the preceding theorem. But to this it is worth adding an interesting relation which links our surfaces $F^d$, of determinant $d>1$, to the $F^1$ of
% scan: s145R
% !PRINTED 265
\weberpage{265}
determinant 1, and offers a means of studying the $F^d$. Consider, on an $F^d$, the corresponding pencil of elliptic curves $C$, which may be supposed of order $d$. To a $C$ there corresponds in general a variety representing the $g_d^{d-1}$ that belong to it; and if one regards these $g_d^{d-1}$ as elements, « points », of a curve, one has an elliptic curve $c$ birationally identical with the $C$, on which \textit{one} point is rationally given, corresponding to the series cut out on the $C$ by the planes or by the hyperplanes of the space containing it. Thus to the surface $F^d$, locus of the curves of the pencil $|C|$, one can associate an $F^1$, locus of the curves $c$, which are birationally identical with the homologous $C$. This assertion would seem, at first sight, to contradict the fact that the $F^1$ and the $F^d$ are birationally distinct surfaces; but one should note that between a $C$ and the homologous $c$ there is given not \textit{one} birational transformation, but rather a group of $d$ transformations, which are not rationally distinguishable as functions of the parameter of the $C$ itself, and thus between $F^1$ and $F^d$ there will be an algebraic correspondence $[d,d]$. Indeed, if a point $P$ is taken on $c$, a series $g_d^{d-1}$ is determined on the $C$, within which a group $G_d$ of $d$ points can be fixed rationally. Now if one associates with $P$ a point $P'$ among the $d$ points of the $G_d$ ($G_d=P'+G_{d-1}$), there is defined a one-to-one correspondence between $C$ and $c$, in which to a point $A$ of $C$ there corresponds on $c$ a point $A'$ for which the group $A'+G_{d-1}$ belongs to the series $g_d^{d-1}$ homologous to $A$.

We shall therefore say:

\textit{To every surface $F^d$ of linear genus $p^{(1)}=1$ and of determinant $d>1$ ($P>1$), containing a linear pencil of elliptic curves $C$, one can associate a surface $F^1$ of determinant $1$, with a pencil of curves birationally identical with the $C$, which is in $[d,d]$ correspondence with $F^d$.}

To the curves $C$ of $F^d$ endowed with a double point there will correspond in particular curves $c$ endowed with a double point on $F^1$; hence \textit{the two surfaces} $F^1$ and $F^d$ will have the same Zeuthen-Segre invariant and the \textit{same genus} $p_a$.

\textit{Notice.} -- In this section we have set out, clarified, and at some points corrected and completed, the results on the classification of the surfaces of linear genus $p^{(1)}=1$ obtained by \textsc{Enriques}, especially in the Notes of the years 1898, 1914: « \textit{Sui piani doppi di genere lineare $p^{(1)}=1$} ». Rendic. Acc. Lincei, $1^\circ$ sem. 1898, « \textit{Sulla classificazione delle [superficie algebriche e particolarmente sulle superficie di genere lineare $p^{(1)}=1$} ». Rendic. Acc. Lincei, Feb. 1914. %

What has been said extends easily to irregular surfaces ($p_g>p_a$), as soon as the characteristic property which concerns them (relating to non-linear continuous series of curves) has been established (Chap. IX, \S\ 2 and 3).

Precisely, in a Note « \textit{Intorno alle superficie algebriche di}
% scan: s146L
% !PRINTED 266
\weberpage{266}
\textit{genere lineare \(p^{(1)}=1\)} »\origfootnote{1}{s146L-1}{R. Acc. di Bologna, Dec. 1906.}, \textsc{Enriques} has established that: « \textit{the irregular surfaces of linear genus \(p^{(1)}=1\) and of arithmetic genus \(p_a\geq0\) \((p_g>p_a)\) contain an irrational pencil, of genus \(p_g-p_a\), of elliptic curves} ». Such surfaces can be classified in the same way as the regular ones, with the character \(p_g-p_a\) added for them.

Outside these there remain only irregular surfaces of linear genus \(p^{(1)}=1\), which have the arithmetic genus \(p_a=-1\), and which can be classified completely as «elliptic and hyperelliptic surfaces, with \(\infty^1\) and \(\infty^2\) birational transformations into themselves » (Chap. XI).

\clearpage\thispagestyle{plain}
% scan: s146R
% !PRINTED none
\weberpage{[s146R]}
\edchapter{VIII}{REGULAR CANONICAL AND PLURICANONICAL SURFACES}{Chapter VIII. REGULAR CANONICAL AND PLURICANONICAL SURFACES}{REGULAR CANONICAL AND PLURICANONICAL SURFACES}

\edsection{\textbf{1. Introduction.}}{1. Introduction.}{sec.VIII.1}

Surfaces of genus \(p_g>3\) possess a (pure) canonical system \(|K|\) of dimension \(p_g-1\geq3\), and by means of this they can as a rule be transformed into \textit{canonical surfaces} of the space of \(p_g-1\) dimensions \(S_{p_g-1}\), having the curves \(K\) as hyperplane sections. In this chapter we shall examine the construction of some canonical surfaces, referring to the case of \textit{regular surfaces}: \(p_g=p_a=p\). %

The canonical surfaces of \(S_{p-1}\), having (absolute) linear genus \(p^{(1)}\) \((>1)\), will in general be of order
\[
p^{(2)}=p^{(1)}-1
\]
and, in their projective properties, will mirror the invariant properties of the whole class of the transformed surfaces.

However, exceptions arise as regards the construction of canonical surfaces: first of all the surfaces of \textit{linear genus} \(p^{(1)}=1\), on which the canonical curves split into the elliptic curves of a pencil. This is the \textit{essential exception}, which takes us back to the family of surfaces characterized by the possession of a pencil of elliptic curves, studied in the preceding chapter: here not only the canonical surface is altogether lacking, but also the pluricanonical surfaces of which we speak below, all the pluricanonical curves being likewise composed of the elliptic curves of the pencil.

Other exceptions, or cases in which the general discussion undergoes at least some modification, are:

1) the surfaces for which the genus \(p\leq3\);

2) those on which the canonical system, though irreducible and \(\infty^3\) at least \((p>3)\), belongs to an involution \(I_n\) of order \(n>1\), so that the canonical surface reduces to a multiple (\(n\)-fold) surface;

% scan: s147L
% !PRINTED 268
\weberpage{268}
3) the surfaces of linear genus \(p^{(1)}>1\) for which the (pure) \textit{canonical curves} \(K\) are in any way reducible; as was seen in Chap. IV, § 18, there are actual examples in which the \(K\) contain fixed (non-exceptional) parts of genus zero and \(>0\); it remains doubtful whether there exist surfaces on which the variable parts of the \(K\) are composed of the irreducible curves \(K\) of a pencil.

If on a surface \(F\) the pure canonical system \(|K|\) possesses a fixed component, exceptional or not, and has genus \(p>3\), having irreducible variable components, the \(\infty^{p-1}\) canonical curves will always lead to a canonical surface, but this will be of order equal to the effective degree of the canonical system \(\bar{p}^{(2)}<p^{(2)}=p^{(1)}-1\). If instead the variable parts of the canonical curves are reducible, there will no longer be a canonical surface. In any case, in this case too, as in the preceding ones, for \(p^{(1)}>1\), one will be able to obtain a projective model of the class by having recourse to the \textit{bicanonical or tricanonical surfaces}, etc., whose plane sections are bicanonical or tricanonical curves respectively. But of this later.

In the following sections we shall refer first of all to regular surfaces with irreducible (pure) canonical curves (without base points) of genus \(p>3\), forming a simple linear system, \(\infty^3\) at least, of genus \(p^{(1)}\) and degree \(p^{(1)}-1\). The first question that presents itself is this: do there exist in one and the same space \(S_{p-1}\), that is, for the same value of the genus \(p\), canonical surfaces of different orders, corresponding therefore to different values of \(p^{(1)}\)? In other words: are the two characters \(p\) and \(p^{(1)}\) independent of each other? The answer, given by \textsc{Noether} as early as 1870, is affirmative.

The independence of \(p^{(1)}\) and \(p\) is proved at once by actual examples. It suffices to observe that a triple point and a tacnode, imposed on the surfaces \(F\) of a given order and with given singularities, equally diminish \(p\) by one unit, but diminish \(p^{(1)}\) by 3 and by 2 respectively.

\edsection{\textbf{2. Canonical surfaces of genus \(p=4\) and linear genus \(p^{(1)}=6,5\).}}{2. Canonical surfaces of genus p=4 and linear genus p\textasciicircum{}(1)=6,5.}{sec.VIII.2}

In this section and in the first ones following it, we propose to classify the regular surfaces of surface genus \(p=4\) and of linear genus \(p^{(1)}>1\), by constructing the corresponding canonical surfaces, which belong to ordinary space \(S_3\). We shall suppose \textit{as a rule} that we are dealing with surfaces with \textit{pure canonical curves \(K\) forming a simple irreducible system} (free of base points on a model without exceptional curves), while reserving some notes for the cases of reducibility; and so we shall look first of all for the simple canonical surfaces, which correspond to the hypothesis of a system \(|K|\) not belonging to an involution \(I_n\) of order \(n>1\).

% scan: s147R
% !PRINTED 269
\weberpage{269}
For \(p=4\) and \(p^{(1)}>1\), a canonical surface \(F\), supposed simple, must have order
\[
p^{(2)}=p^{(1)}-1\geq5.
\]

The simple type corresponding to the smallest value of \(p^{(1)}\) is therefore: \textit{the general surface of the \(5^\circ\) order, \(F_5\), a canonical surface with}
\[
p=4,\qquad p^{(1)}=6.
\]

It is clear that a surface with \(p=4\) and \(p^{(1)}=6\) whose canonical curves are irreducible can be reduced birationally to the projective model of the \(F_5\). As examples we cite the surfaces indicated by Noether:

\textit{a}) The surface of the \(7^\circ\) order \(F_7\), with a triple conic and a triple point outside it. By a quadratic transformation of the first kind the \(F_7\) is brought back to a particular \(F_5\), on which a point \(O\) and a conic are given; the point \(O\) corresponds to the exceptional line in which the plane of the triple conic cuts \(F_7\).

\textit{b}) The \(F_7\) with a double sextic \(C_6\) of genus 3. In the usual plane representation of a cubic surface \(F_3\) passing through \(C_6\), the \(C_6\) itself appears to have as its image a \(C'_6\) \((1^3\,2^2\,3^2\,4^2\,5^2\,6)\), whose behavior with respect to the fundamental points \(1,2,3,4,5,6\) is indicated by the preceding symbol; hence the residual variable intersections of two \(F_3\) through \(C_6\) are twisted cubics, represented by \(C_3(2,3,4,5,6^2)\), forming a homaloidal net; one deduces that these same \(F_3\) form the transforming system of a Cremona transformation of space, which reduces our \(F_7\) to the \(F_5\).

\textit{c}) The surface \(F_7\) with a double quintic \(C_5\) of genus 0. Here too the \(F_3\) through \(C_5\) form a homaloidal system, which however has as base curve the \(C_5\) itself completed by a line meeting it in 4 points, which is an exceptional line of the \(F_7\); and by means of this transforming system of \(F_3\) the \(F_7\) is reduced by a Cremona transformation to the \(F_5\).

Now we wish to prove with full rigor and generality that: \textit{every surface \(F\) with genera \(p=4\) and \(p^{(1)}=6\), such that the variable part of the canonical curves is irreducible, can be transformed into a canonical quintic \(F_5\).}

Let us denote by \(|k|\) the variable part of the canonical system, which by hypothesis is not composed of the curves of a pencil. Forming the projective image of this system, one obtains a surface \(F_n\) whose order \(n\) \((\leq5)\) equals the effective degree of \(|k|\).

If \(n=5\), the \(F_n\) is the quintic, and it turns out a posteriori that
\[
|k|=|K|.
\]

If \(n<5\), one cannot have a simple canonical surface (which would be of genus \(p\leq1\)), nor a multiple plane (which would en%
% scan: s148L
% !PRINTED 270
\unskip\weberpage{270}%
tail \(p=3\)), but there remains possible — a priori at least — only a double quadric \(F_2\): this case would in fact be met with if there were an exceptional fixed component of the pure canonical curves, that is, a base point of \(|K|\) on a model of \(F\) free of exceptional curves. The curves \(k\) being of genus \(\pi\leq p^{(1)}\) \((=6)\), the branch curve on the \(F_2\) will be a \(C_{2\pi+2}\) of order \(2\pi+2\), and we shall prove first of all that one certainly does not have
\[
\pi<5,
\]
and therefore only the following cases are possible a priori:
\[
\pi=6,5.
\]

Indeed one cannot have \(\pi<5\), that is, \(\pi\leq4\), since the canonical curves \(K\) would have to cut the \(k\) of genus \(\pi\) in groups residual to the characteristic series \(g_n^2=g_4^2\), and hence in not more than 2 points, whereas one must have \(kK>kk=4\). But if \(\pi=5\), one has on the double \(F_2\) a branch curve of order 12; and if \(\pi=6\), a curve of order 14. The first hypothesis leads to a double plane with branch curve \(C_{12}\) endowed with two 6-fold points, which has as canonical curves the conics through the two singular points (completed by the exceptional line joining them), and hence the absolute linear genus \(p^{(1)}=5\) instead of \(p^{(1)}=6\). The second hypothesis leads to a double plane with branch \(C_{14}\) endowed with an 8-fold point \(O\) and a 6-fold point \(O_1\), or with two infinitely near 7-fold points \(OO_1\), which has as canonical curves the \(\infty^3\) cubics through \(O^2\) and \(O_1\) (completed by the exceptional line \(OO_1\)), and is therefore of genus \(p=5\) instead of \(p=4\).

We have said that \(p^{(1)}=6\) is the minimum value that the linear genus \(p^{(1)}\) can assume for a surface of genus \(p=4\) that admits a simple canonical surface as projective model. But this is no longer so if the canonical system \(|K|\) belongs to an involution \(I_2\) of order 2, thus admitting a base point.

One then has as \textit{canonical surface} the \textit{double quadric} with branch curve of order 12, which leads, as has been said, to the \textit{double plane with branch curve} \(C_{12}\) of order 12 endowed with two 6-fold points. This double plane with \(p=4\) and \(p^{(1)}=5\) corresponds to the \textit{minimum value of \(p^{(1)}\) for \(p=4\)} (in agreement with the general inequality \(p^{(1)}\geq2p-3\), which we shall establish later). This follows from the preceding analysis, because a smaller value of \(p^{(1)}\) would have to lead to a simple canonical surface of order 3 or 2, which would be of genus \(p=0\), or to a triple plane, which would have genus \(p=3\).

And the same analysis also excludes the possibility that a hypothetical reducibility of the pure canonical system (that is, the presence of fixed components, exceptional or not) might lead to other types of surfaces with \(p=4\) and \(p^{(1)}=5\).

% scan: s148R
% !PRINTED 271
\weberpage{271}
We shall therefore state that: \textit{Every surface with \(p=4\) of minimum linear genus \(p^{(1)}=5\) can be transformed into the double plane with branch \(C_{12}\) endowed with two sextuple points.}

As an example, consider the surface \(F_6\) of the \(6^\circ\) order with two double lines \(a\) and \(b\) skew to each other: the adjoint quadrics cut out on such an \(F_6\) the canonical curves \(K\), which belong to the involution cut out by the lines meeting \(a\) and \(b\).

\edsection{\textbf{3. Surfaces with \(p=4\) and \(p^{(1)}=7\).}}{3. Surfaces with p=4 and p\textasciicircum{}(1)=7.}{sec.VIII.3}

Let us proceed to construct the canonical model for the surfaces of genera \(p=4\) and \(p^{(1)}=7\). We shall have, under the apparently most general hypotheses, surfaces \(F_6\) of order 6 of ordinary space, whose adjoint quadrics must split into a variable plane and a fixed plane, which meets \(F_6\) in a double cubic.

Thus \textit{for \(p=4\) and \(p^{(1)}=7\) one has the canonical surface of the \(6^\circ\) order \(F_6\) with a double plane cubic.}

Every surface whose genera are \(p=4\) and \(p^{(1)}=7\), with a simple canonical system, can be transformed birationally into the \(F_6\). Thus, for example, the surfaces of the \(6^\circ\) order \(F_6\) with 6 triple points, whose canonical curves (intersections of the quadrics through the 6 points) are precisely of genus 7 and degree 6.

\textit{Note.} – But if one wants to study more deeply the family of the surfaces for which \(p=4\) and \(p^{(1)}=7\), it is convenient to examine those for which the canonical system \(|K|\) (which is supposed irreducible apart from possible fixed components) is not simple, so that the canonical surface reduces to a double cubic \(F_3\), or to a quadric \(F_2\), double or triple. The principal aim of this examination is to answer the question « whether all the surfaces for which \(p=4\) and \(p^{(1)}=7\) belong to one and the same family represented by the canonical \(F_6\), or whether there are several distinct families of surfaces with the said characters ».

As regards the double cubics \(F_3\), it is easy to convince oneself that they are included as a particular case in the family of the \(F_6\).

First of all the \(F_3\) cannot be ruled, since it would not be normal and \(p=5\) would result; hence it will have elliptic plane sections and will possess a branch curve \(C_{12}\) of order
\[
2p^{(1)}-2=12.
\]

This \(C_{12}\) must be the complete intersection of the \(F_3\) with a surface of the \(4^\circ\) order; otherwise one would find on \(F_3\) a pencil of conics meeting the \(C_{12}\) in 6 points, which would be images of curves of genus 2 meeting the canonical curves \(K\) in 2 points, so that the \(K\)
% scan: s149L
% !PRINTED 272
\weberpage{272}
themselves would turn out to be hyperelliptic and the plane sections of \(F_3\) rational, which has been excluded.

Hence the double \(F_3\) can be represented on the \textit{double plane with branch curve \(C_{12}\) of order 12 endowed with 6 quadruple points}, which — in fact — has the genera \(p=4\) and \(p^{(1)}=7\), and on which the canonical curves have as images the cubics through the 6 points counted twice.

Let us prove that \textit{the said double plane is included as a particular case in the family of the \(F_6\) with a double plane cubic}; considering the pencil of surfaces determined by the \(F_6\) and by one of its adjoints \(F_3\) counted twice:
\[
F_3^2+\lambda F_6=0,
\]
for \(\lambda=0\) one obtains in this pencil the double cubic surface \(F_3^2\), with branch curve cut out by \(F_6=0\).

To prove this it suffices to recognize that a plane section \(f_6\) of \(F_6\), belonging to a pencil \(f_3^2+\lambda f_6=0\), reduces by continuity (for \(\lambda=0\)) to the double cubic \(f_3\) with its branch points at the 12 simple base points of the said pencil. Now — if we suppose the \(f_6\) entirely general (that is, without double points) — it constitutes an envelope of class 30 which, in the limit for \(\lambda=0\), must reduce to the cubic envelope \(f_3=0\) counted twice (which is of class \(2\cdot6=12\)), augmented by the 18 pencils having as centers the branch points of the said double cubic; and it is evident that the points of contact of the tangents drawn from any point \(P\) to the \(f_6\), for \(\lambda=0\), come to fall at the 18 intersections of \(f_6\) and \(f_3\)\origfootnote{1}{s149L-1}{Cf. \textsc{Enriques-Chisini}, \textit{Lezioni}, vol. III.}. If, further, the \(f_6\) has 3 collinear double points, taken as base of the pencil \(f_3^2+\lambda f_6=0\), the 3 pencils with those centers, counted 2 times, are to be removed from the 18 mentioned above, there remaining therefore a branch \(G_{12}\) of the double \(f_3\).

While the canonical surfaces consisting of double cubics with branch \(C_{12}\) are included as a particular case in the more general type of the \(F_6\) with a double plane curve of the \(3^\circ\) order, on the contrary those canonical surfaces which consist of double quadrics form a distinct family.

Without dwelling on a more minute examination, let us say that there is a type of surface \(F\), with \(p=4\) and \(p^{(1)}=7\), on which the canonical curves are hyperelliptic with two base points (taken among the double points of the \(g_2^1\)), and these lead to canonical double quadrics representable on the \textit{double plane with branch curve \(C_{10}\) of order 10, endowed with two \([3,3]\) points}.

And these double planes actually furnish \textit{a second family of surfaces with \(p=4\) and \(p^{(1)}=7\), which are not included among the simple canonical \(F_6\)}, because they depend on 39 moduli (\(39=65
% scan: s149R
% !PRINTED 273
\text{\weberpage{273}}
-18-8\)\edFlush{}), whereas we have seen that the \(F_6\) depend on 38 moduli.

Greater difficulties are presented by the examination of the canonical surfaces, with \(p=4\) and \(p^{(1)}=7\), which reduce to triple quadrics. The surfaces of this family are characterized by the property that the pure canonical curves \(K\) have no base points and contain a \(g_3^1\), so that \(|K|\) turns out to belong to a rational involution of the \(3^\circ\) order \(I_3\). \textit{The triple canonical quadric \(Q_2^3\)}, studied by \textsc{G. Pompilj}\origfootnote{1}{s149R-1}{\textsc{G. Pompilj}, \textit{Sulle superficie algebriche le cui curve canoniche posseggono una \(g_3^1\).} Rendiconti Istituto Lombardo, 1940.}, \textit{is a cone and possesses a branch curve \(C_{18}\) of order 18.}

One recognizes that the \textit{triple canonical quadrics \(Q_2^3\) are not included as particular cases in the family of the simple \(F_6\) with a double plane cubic, but rather constitute a third family of surfaces with \(p=4\) and \(p^{(1)}=7\), distinct from the preceding ones.}

To prove this it will suffice to verify that a plane curve \(C_6\) with 3 double points on a straight line, a section of \(F_6\), cannot reduce by continuity to a triple conic \(C_2^3\), a section of \(Q_2^3\), endowed with 18 branch points. To this end one will observe first of all that the system \(\infty^{27}\) of the plane sextics \(C_6\) does not contain within itself the system \(\infty^{29}\) of the triple conics of the same genus with 24 branch points. One cannot exclude, however, that among the \(\infty^{27}\) \(C_6\) there appear particular \(C_2^3\) with 24 critical points; but if one wants to find within this family a \(C_2^3\) with 18 distinct branch points, constituting a degenerate \(C_6\) with 3 double points on a straight line, it is necessary that 6 of the 24 critical points mentioned above coincide in pairs in 3 \textit{collinear} apparent critical points, and this is impossible for an \textit{irreducible \(C_2^3\)}.

\edsection{\textbf{4. Surfaces with \(p=4\) and \(p^{(1)}=8\).}}{4. Surfaces with p=4 and p\textasciicircum{}(1)=8.}{sec.VIII.4}

Let us pass to the next case: \(p=4,\ p^{(1)}=8\). The (simple) canonical surface \(F_7\) corresponding to these characters must in general possess a double curve \(C_7\) which, counted twice, constitutes the complete intersection of the \(F_7\) with the adjoint quadric \(\Phi_2\). The \(C_7\), by a formula already seen\origfootnote{2}{s149R-2}{Cf. chap. V, § 3.}, possesses a triple point which turns out to be double for the adjoint \(\Phi_2\), which is therefore a cone. This is in agreement with the fact that the \(C_7\) is the curve of contact of the \(\Phi_2\) with a surface of order 7, and therefore belongs to the half system of the linear system \(|C_{14}|\) consisting of the intersections of \(\Phi_2\) with the surfaces of the \(7^\circ\) order of space. But from the plane representation of the quadric it appears that the said system \(|C_{14}|\) cannot in general be halved; in order that it can be, the quadric must be
\printersignature{\textsc{Enriques F.} - \textit{Superficie algebriche.}\quad 18}

% scan: s150L
% !PRINTED 274
\weberpage{274}
a cone, whose representative system consists of the conics through two infinitely near points \(O,\ O_1\): the images of the \(C_{14}\) are then curves of the same order \(C'_{14}\) \((O^7O_1^7)\), and the half of the system \(|C'_{14}|\) is a priori a system \(C'_7(O^4O_1^3)\) or a more particular system contained in it. But the \(C'_7(O^4O_1^3)\) which do not have at \(O\) a multiplicity greater than 4 are images of curves \(C_7\) of \(\Phi_2\), curves of contact of the latter with certain \(F_7\). These cannot be double for an irreducible \(F_7\), because the generators of the cone \(\Phi_2\) would have eight intersections with such an \(F_7\). It is therefore necessary that the double curve of one of our \(F_7\) belong to the linear system which in the representative plane of \(\Phi_2\) has as image \(C_7\) \((O^5O_1^2)\), and hence that it pass three times through the vertex of the cone and have genus \(\varrho=4\). %

We have proved that the canonical \(F_7\) \((p=4,\ p^{(1)}=8)\), where it exists, possesses as double curve a \(C_7\) of genus \(\varrho=4\), endowed with a triple point \(V\). This curve belongs to the half of the system cut on \(\Phi_2\) by the surfaces of the \(7^\circ\) order which pass 3 times through \(V\) (or behave like those which have a triple point there). Now it is a matter of recognizing the effective existence of such \(F_7\). This can be inferred inductively from the count of moduli; if one assumes the more expressive formula of § 13, Chap. V:
\[
M\geq10p-2p^{(1)}+12=36,
\]
one is led to believe that every \(C_7\) (defined as above on a quadric cone with vertex \(V\) triple for it) is a double curve for \(\infty^{26}\) canonical surfaces \(F_7\). Indeed the \(C_7\) on the quadric cone are \(\infty^{17}\), and so, with respect to the \(\infty^7\) projective transformations of the cone, they give rise to \(\infty^{10}\) projectively distinct curves. On the other hand, the surfaces \(F_7\) having one and the same double curve \(C_7\) are \(\infty^{26}\), since by splitting off the cone one has \(\infty^{25}\) residual \(F_5\) passing through the \(C_7\).

The preceding conclusion is reached only inductively. This is so both because the formula for the moduli that has been used depends on a hypothesis not rigorously established, and because that formula serves to establish the extent of a continuous family of surfaces with given \(p\) and \(p^{(1)}\) only when the existence of some surface of the family is presupposed as given. Thus one must at least recognize the effective existence of a particular surface \(F_7\) (with the genera \(p=4,\ p^{(1)}=8\)) endowed with a double \(C_7\): one succeeds, for example, in constructing such an \(F_7\) starting from a sextic \(F_6\) with 5 triple points and a tacnode, for which precisely \(p=4,\ p^{(1)}=8\). But more directly one can prove that every \(C_7\) (belonging to the series considered on a quadric cone) is the double curve of a (canonical) surface. \(F_7\), by the method indicated %
% scan: s150R
% !PRINTED 275
\weberpage{275}
by \textsc{A. Franchetta}\origfootnote{1}{s150R-1}{\textsc{A. Franchetta}, \textit{Su alcuni esempi di superficie canoniche}, Seminario Mat. di Roma, 1939.}, and the count of moduli is thereby confirmed for these same \(F_7\).

The method to which we resort consists in this: one recognizes that a \(C_7\) is always a double curve for one (and indeed for \(\infty^6\)) rational surface of the seventh order \(F'_7\) which in addition possesses an assigned quadruple point; by linearly combining this particular \(F'_7\) with a degenerate \(F_7\) passing doubly through \(C_7\), one obtains the required canonical \(F_7\).

To carry out this program, observe first of all that there exist rational surfaces of the \(7^\circ\) order \(F'_7\) endowed with a double curve of the \(7^\circ\) order and with an isolated quadruple point. Such a surface is represented on the plane by the system of the curves of the eleventh order \(C_{11}(A^5B_1^3\ldotp\ldotp\ldotp\ldotp B_9^3D_1\ldotp\ldotp\ldotp\ldotp D_8)\), where the 18 points \(A,\ B\) and \(D\) belong to a fundamental curve of the seventh order \(K_7\) \((A^3B_1^2\ldotp\ldotp\ldotp\ldotp B_9^2D_1\ldotp\ldotp\ldotp\ldotp D_8)\). The \(C_{11}\) are curves of genus 8, and the characteristic series of the \(\infty^3\) system is self-residual, the system \(|C_{11}|\) being equivalent to the sum of the fundamental curve \(K_7\) and its adjoint. Precisely for this reason the represented surface is an \(F'_7\) with plane sections of genus 8, and hence endowed with a double curve \(C_7\), which moreover possesses a quadruple point corresponding to the \(K_7\). %

Now the \(F'_7\) thus constructed depend on a certain number of projective absolute invariants, which it is easy to evaluate: this number is 19. Moreover one can calculate the dimension of the system of the \(F'_7\) possessing one and the same double curve \(C_7\) and one and the same quadruple point \(O\) outside it; this dimension exceeds by one that of the characteristic system cut on an \(F'_7\) by the other \(F'_7\) of the same system, and therefore equals 6. %
This being so, by letting the quadruple point vary, one will have (at most) \(\infty^9\) \(F'_7\) with the same double curve \(C_7\), so that — there being \(\infty^{19}\) \(F'_7\) — one will have (at least) \(\infty^{10}\) curves \(C_7\) double for an \(F'_7\). But the \(C_7\) thus obtained are precisely those curves of the \(7^\circ\) order and of genus 4, lying on a quadric cone and passing triply through the vertex, which we defined above. By linearly combining an \(F'_7\) with the \(\infty^{25}\) degenerate \(F_7\) consisting of the quadric cone containing \(C_7\) and of an \(F_5\) through \(C_7\), one obtains \(\infty^{26}\) canonical surfaces \(F_7\) passing doubly through the assigned \(C_7\). Thereby it is proved that: \textit{there exist \(\infty^{36}\) projectively distinct canonical surfaces of the \(7^\circ\) order \(F_7\), whose double curves are the \(\infty^{10}\) curves of the \(7^\circ\) order and of genus 4 which are defined by means of the bisection of the linear system cut on a quadric cone by the surfaces of the \(7^\circ\) order passing triply through the vertex.}

\textit{The curves \(C_7\) can be characterized as hyperel}%
% scan: s151L
% !PRINTED 276
\unskip\weberpage{276}%
\textit{liptic space curves of genus $4$, endowed with a triple point}: indeed the hyperelliptic space curves of the $7^{\mathrm{o}}$ order and of genus 4 are $\infty^{28}$, and the possession of a triple point constitutes \textit{one} condition for them. This appears a priori from the fact that the $g_7^3$ possessing a neutral triple on a hyperelliptic curve of genus 4 are defined by containing the double of the $g_2^1$, and therefore depend on three arbitrary constants. %

We have explained the general construction of the canonical surfaces $F_7$ endowed with a double curve $C_7$: a hyperelliptic curve of genus 4, endowed with a triple point. But suppose that, instead of asking for the construction of an $F_7$ with double $C_7$, one asks for that of a canonical $F_7$ possessing some multiple curve of multiplicity greater than 2. Then the very simple example immediately presents itself of the $F_7$ endowed with a triple line $a$ and with two double conics in planes through it. When the triple line and the two double conics are assigned, there are $\infty^{33} F_7$: indeed, splitting off the planes of the two conics diminishes the number of parameters by 2, and there remains the dimension of the system of the $F_5$ passing simply through a line and through two conics in planes through it, which is 31.

Now observe that the curve composed of a line and of two conics in planes through it possesses an absolute invariant, which is the cross-ratio of the 4 points on the line. One will therefore have $\infty^{34}$ projectively distinct canonical surfaces $F_7$ possessing a triple line and two double conics. The question arises: do these $F_7$ belong, as particular surfaces, to the family of the canonical $F_7$ with double $C_7$, or do they constitute a distinct continuous series of surfaces, depending on only 34 instead of 36 moduli? If the formula giving the moduli were established without exception in its more expressive form cited above, the question would at once be resolved a priori by adopting the first alternative. In any case one can recognize directly in what way the $F_7$ with a triple line and two double conics are specializations of the $F_7$ with double $C_7$.

It is a matter of showing that the curves of ordinary space $C_2 + C'_2 + 3a$, consisting of two generic conics and of the line of intersection of their planes, belong as a particular case to the family of our $C_7$. Since the said $C_2 + C'_2 + 3a$ satisfy the three conditions of possessing a triple point and of being hyperelliptic (that is, of containing a $g_2^1$), it suffices to show that these composite curves belong to the family of the $C_7$ of genus 4. Now the $C_7$ of genus 4 of ordinary space can be obtained as projections of curves $C_8$ of the space $S_4$ from a simple point $O$. Among these $C_8$ there are curves split into two quartics $C_4$, $C'_4$ of the $1^{\mathrm{a}}$ species passing doubly through $O$ and having there the same tangent plane $\alpha$; these are projected from $O$ into a curve composed of the triple line $a$, trace of the plane $\alpha$, and of the two conics $C_2$, $C'_2$, projec%
% scan: s151R
% !PRINTED 277
\unskip\weberpage{277}%
tions of $C_4$, $C'_4$. Thus it appears that the curve $C_2 + C'_2 + 3a$ can be regarded as a specialization of our hyperelliptic $C_7$ endowed with a triple point. One can see directly that the conditions required for this degeneration of the $C_7$ are 9 in number, and therefore all the composite curves $C_2 + C'_2 + 3a$ belong to the family of our $C_7$. As we have already seen, this entails that the $F_7$ with a triple line and two double conics appear within the family of the general canonical surfaces with $p = 4$ and $p^{(1)} = 8$, satisfying 2 particular conditions.

The canonical surfaces $F_7$ which we have learned to construct correspond to the values of the genera $p = 4$ and $p^{(1)} = 8$ ($p^{(1)}$ the \textit{absolute} linear genus) when the canonical curves $K$ on a transformed surface $F$ without exceptional curves form a simple irreducible system free from base points. If instead one supposes that the system $|K|$ has on $F$ a (simple) \textit{base point}, this system will have the effective degree $p^{(2)} - 1 = p^{(1)} - 2 = 6$. Hence the corresponding canonical surface will reduce to a sextic $F_6$, on which the base point of the canonical system must have as image an exceptional line counted twice: this line forms part simply of the impure canonical system and then appears as a fixed component of the pure canonical system. %

Now it is easy to construct effectively a surface of the $6^{\mathrm{o}}$ order $F_6$ possessing an assigned double conic $C$ and, in the plane $\gamma$ of the latter, a tacnode $O$ whose tacnodal plane $\alpha$, different from $\gamma$, is likewise assigned. All the $F_6$ satisfying the indicated conditions are obtained by linear combination of two particular $F_6$. The first is the $F_6$ consisting of a quadric through $C$ counted twice and of the plane $\alpha$, also counted twice. The second is an $F_6$ composed of the plane $\gamma$ and of an $F_5$ which passes simply through $C$ and has at $O$ a biplanar point, touching with one of its sheets the assigned tacnodal plane $\alpha$.

The $F_6$ thus constructed in fact possesses a fixed component of the canonical curves. This is the line $a$ along which the plane $\alpha$ touches $F_6$, and which therefore appears twice in the intersections of $F_6$ with the $\infty^3$ adjoint quadrics, consisting of a variable plane and of the fixed plane $\alpha$. Denoting by $\varrho = 0$ and $-\nu$ the genus and the degree of this line $a$, the genus and the degree of $2a$ will be respectively: %
\[
2 \cdot 0 - \nu - 1 = -(\nu + 1) \qquad \text{and} \qquad -4\nu;
\]
one can therefore evaluate the degree $\bar{p}^{(2)}$ and the (virtual) genus $\bar{p}^{(1)}$ of the impure canonical system on $F_6$ (the latter is its relative genus):
\[
\begin{aligned}
\bar{p}^{(2)} &= 6 - 4\nu + 4 = 10 - 4\nu\\
\bar{p}^{(1)} &= 8 - (\nu + 1) + 2 - 1 = 8 - \nu \qquad (p^{(1)} = 8);
\end{aligned}
\]
% scan: s152L
% !PRINTED 278
\weberpage{278}
and from the relation
\[
\bar{p}^{(2)} = \bar{p}^{(1)} - 1
\]
one then derives
\[
\begin{aligned}
10 - 4\nu &= 7 - \nu,\\
\nu &= 1.
\end{aligned}
\]

Thus it is proved a posteriori that the line $a$ is effectively an exceptional line for $F_6$ ($-\nu = -1$), so that, owing to its presence, the relative linear genus
\[
\bar{p}^{(1)} = 7 = p^{(1)} - 1
\]
and the absolute linear genus $p^{(1)} = 8$. Moreover this line, removed once from the impure canonical curves, still remains a fixed component of the pure canonical curves, and thus corresponds to a simple base point of $|K|$ on a transformed surface $F$\origfootnote{1}{s152L-1}{One can also verify in another way that the $F_6$ is a surface of absolute linear genus $p^{(1)} = 8$ whose pure canonical curves possess the exceptional line as a fixed component. Indeed it suffices to transform the $F_6$ quadratically by means of the quadrics passing through the conic $C$ and through a point $A$ of it. One thus obtains a surface of the $7^{\mathrm{o}}$ order with a triple conic, which in addition possesses a double point (of \textsc{Nöther}'s type) $O$, with an infinitely near tacnode. This singular point (image of an elliptic fundamental curve of degree $-1$) gives rise to a simple base point for the canonical curves cut by the quadrics, which are the variable parts of the adjoint cubic surfaces.}.

Let us now count the number of parameters on which our $F_6$ depend. The $F_5$ passing through the conic $C$ and having a biplanar point $O$ with a sheet tangent to the plane $\alpha$ depend on 37 parameters; therefore the $F_6$ with assigned double conic, tacnode and tacnodal plane will be $\infty^{38}$.

The singularities assigned for the $F_6$ can be transformed projectively into other singularities given arbitrarily in space (since they possess no absolute invariant), and indeed those singularities are transformed into themselves by $\infty^3$ homologies. Hence one sees that: \textit{our $F_6$, canonical surfaces with $p = 4$ and $p^{(1)} = 8$ corresponding to a canonical system endowed with a base point, depend on $35$ absolute invariants or moduli.}

Here too the number of moduli is less than that (36) which belongs to the canonical surfaces $F_7$. One is therefore led to ask whether the said $F_6$ constitute a specialization of the $F_7$, corresponding to the degeneration of the latter into an $F_6$ and a plane.

We shall confine ourselves to indicating in what way this can come about. If one adds the plane $\gamma$ to our $F_6$, one obtains a surface $F_7$ which possesses the triple conic $C$ and in addition two infinitely near double lines coinciding with $a$; the same $F_7$ can be regarded as endowed
% scan: s152R
% !PRINTED 279
\weberpage{279}
only with the triple conic $C$ and the double line $a$. For this one regards as virtually non-existent the double line infinitely near to $a$, which we wish to consider as a line of connection between the $F_6$ and the plane $\gamma$. Thus the composite $F_7$ will have the (arithmetic) genus of the $F_6$, which would have to be increased by zero (genus of the plane) and by zero (genus of the line of connection). It is a matter of recognizing that the degenerate $F_7$ defined in this way can be considered as a particular case of a canonical $F_7$ with a double curve $C_7$ of order 7 and genus 4 (hyperelliptic and endowed with a triple point), account being taken of the point singularity which our construction comes to assign to it at the point $O$. For this purpose it will be useful to consider the said $C_7$ as the projection from a point $P$ of a $C_8$ of order 8 and genus 4 of the space $S_4$. One then lets the $C_8$ degenerate into a curve $C_6$ of genus 3 belonging to a quadric cone with vertex $P$ and having at $P$ a double point, and into two lines $a$ and $o$ meeting each other, the first of which is a chord of $C_6$ while the other passes through $P$. The projection of this degenerate $C_8$ from $P$ into ordinary space can be regarded as the limit curve of a variable $C_7$ which comes to degenerate into a triple conic $C$ and a line $a$ lying in the plane of $C$; this degenerate curve is to be completed by adding to it an infinitesimal curve consisting of the neighborhood of the point $O$, trace of the line $o$, in a plane through $a$. One understands in this way how also the \textit{possession of the tacnode of the $F_6$ can result as a limit of the possession of the double curve $C_7$ which defines the family of the $F_7$.}

\textit{Note.} -- To go deeper into the study of the surfaces with $p = 4$ and $p^{(1)} = 8$ it is necessary to examine the cases in which the canonical system $|K|$, irreducible in its variable components, belongs to an involution, giving rise to a multiple canonical surface. One is led to examine canonical surfaces consisting of double cubics $F_3$, or of double or triple quadrics $F_2$. Without dwelling on this examination, we shall confine ourselves to indicating some notable families which are obtained in this way.

First of all one finds a double $F_3$ which corresponds to the case in which the $K$, belonging to a rational involution $I_2$ (considered on a transform without exceptional curves), have a base point. The $F_3$, as in the preceding section, is recognized to have elliptic plane sections (that is, not to be ruled), and its branch curve now turns out to be a curve of order $2p^{(1)} - 2 = 14$. Representing this $F_3$ on the plane, one obtains a \textit{double plane with branch curve $C_{12}$ of order $12$, endowed with $5$ $4$-fold points and a $[3,3]$ point}. In fact this double plane has as (double) images of the canonical curves the cubics through the 6 singular points of $C_{12}$, and therefore
\[
p = 4,\qquad p^{(1)} = 8.
\]

% scan: s153L
% !PRINTED 280
\weberpage{280}
Now one can recognize that \textit{the family of surfaces with the genera $p = 4$ and $p^{(1)} = 8$ defined by the said double planes belongs to that of the simple canonical $F_6$}, which correspond to the case of a canonical system endowed with a base point. Indeed, consider the cubics adjoint to one of the said $F_6$, that is, passing through its double conic $C_2$ and through its tacnode $O$ (in the plane of $C_2$); among these surfaces take one, $F_3$, which touches the tacnodal plane of $F_6$ at $O$, and construct the pencil
\[
F_3^2 + \lambda F_6 = 0.
\]

In this pencil, for $\lambda = 0$, one sees the $F_6$ degenerate continuously into the $F_3^2$, with a branch curve of order 14, the section of the $F_6$ outside $C_2$. The line which, counted twice, constitutes the further section of $F_6$ by the plane of $C_2$ comes to form part of the branch curve.

A second family of multiple canonical surfaces is that of the \textit{double quadrics} referable to the type of the \textit{double plane with branch curves $C_{12} = r + C_{11}$ comprising an exceptional line $r$ and a $C_{11}$ of order $11$ endowed with two infinitely near $4$-fold points $AA_1$ on $r$ and with two $[3,3]$ points $B$ and $D$, one of which, $B$, also belongs to the line $r$} (the infinitely near triple point $B_1$ being outside $r$). The double images of the canonical curves $K$ are here the conics through $A$ and $D$, completed by the exceptional line $AA_1D$; therefore
\[
p = 4 \qquad \text{and} \qquad p^{(1)} = 8,
\]
and on a model of the surface without exceptional curves the canonical curves will have 3 base points. The question whether these double planes belong to the family defined by the canonical surfaces $F_7$ or to that of the $F_6$ (simple pure canonical system with a base point or an exceptional fixed component) seems bound to receive a negative answer. This is so although the number of moduli on which these double planes depend is 33, that is, less than that of the said $F_7$ (which have 36) and of the $F_6$ (which have 35). Indeed one does not see how a surface composed of a double quadric and of two planes could possess the characteristic singularities of the $F_6$.

We shall not dwell further on the multiple canonical surfaces with the genera $p = 4$ and $p^{(1)} = 8$, and we leave aside the problem of the canonical triple quadrics with the said genera. We shall only say that an investigation by \textsc{G. Pompilj} (which we shall have occasion to cite later) succeeds in excluding the existence of this type.

Finally we shall confine ourselves to pointing out that, for a complete classification of the surfaces with $p = 4$ and $p^{(1)} = 8$, it would still be necessary to examine further multiple canonical surfaces. These are the ones which might arise under the hypothesis of reducible canonical curves with a fixed part, whose variable components are of genus $< p^{(1)}$.

% scan: s153R
% !PRINTED 281
\weberpage{281}
\edsection{5. Surfaces with $p = 4$ and $p^{(1)} = 9$.}{5. Surfaces with p = 4 and p\textasciicircum{}(1) = 9.}{sec.VIII.5}

The canonical surfaces $F_8$ of genus $p = 4$ and (absolute) linear genus $p^{(1)} = 9$ with a simple irreducible canonical system without base points can be constructed by the same method used above for $p^{(1)} = 7,8$. The $F_8$ must possess, according to our formulas,
\[
t = 4
\]
triple points, and an adjoint surface $\Phi_3$ having 4 double points at these. One must consider the $F_9$ passing triply through the 4 double points of $\Phi_3$ and halve the linear system of the curves cut by them on $\Phi_3$.

Let us resort to the usual representation of the $\Phi_3$ on a plane, in which the plane sections have as images curves $C_3$ passing through three pairs of infinitely near points, $PP_1$, $QQ_1$ and $RR_1$, on a conic. The half system we are seeking is then given by the curves of the $9^{\mathrm{o}}$ order,
\[
C_9(P^4P_1Q^4Q_1R^4R_1).
\]

These curves are, on the $\Phi_3$, $\infty^{21}$ curves of the $12^{\mathrm{o}}$ order $C_{12}$ of genus 10, and a priori they constitute possible double curves of canonical surfaces $F_8$. Now we observe that a $C_{12}$ which is a double curve for an irreducible $F_8$ will also be double for $\infty^{13} F_8$, since, on splitting off the adjoint $\Phi_3$, there remain $\infty^{12}$ surfaces of the $5^{\mathrm{o}}$ order $F_5$ passing through the same $C_{12}$. Since the cubic surfaces with 4 double points are $\infty^{15}$, all projectively identical with one another, our $C_{12}$ belong to a family depending on 21 absolute invariants. If we admit that each $C_{12}$ is a double curve for one (and hence for $\infty^{13}$) irreducible $F_8$, one will have $\infty^{34} F_8$ projectively distinct. That is, the number of moduli of our canonical surfaces will be given exactly by the formula
\[
M = 10p - 2p^{(1)} + 12 = 40 - 18 + 12 = 34.
\]

The reasons set out in Chap. V, \S\ 13, tend to establish the aforesaid expression of $M$ at least as a lower limit for the number of moduli of surfaces of genus $p > 3$ (with simple canonical system). They lead one to believe that in fact all the $C_{12}$ defined above are double curves for actual irreducible $F_8$ (each for $\infty^{13}$ of these). The result thus inductively obtained can be confirmed by the method of \textsc{Franchetta}; in this way the count of moduli according to the formula written above is also confirmed, at least plausibly, for this class of surfaces. First of all one recognizes that \textit{there actually exist irreducible canonical surfaces $F_8$ having}
% scan: s154L
% !PRINTED 282
\weberpage{282}
\textit{a double $C_{12}$}, and this at least when this $C_{12}$ is a double curve for a rational surface $F'_8$ also endowed with an isolated quadruple point. The $F_8$ which are associated in this way with a rational $F'_8$ are obtained by linear combination of the $F'_8$ with the surfaces composed of the subadjoint $\Phi_3$ and of an $F_5$ passing through the same $C_{12}$. As for the effective existence of the rational $F'_8$, it is proved by resorting to their representative system in the plane. This is a linear system $\infty^3$ of curves of genus $p^{(1)} = 9$ with self-residual characteristic series, such as can be defined for example by means of the curves of the eleventh order $C_{11}\,(A_1^3\ldotp\ldotp\ldotp\ldotp A_{12}^3\ B_1\ldotp\ldotp\ldotp\ldotp B_5)$, where the 17 points $A$ and $B$ belong to a fundamental curve of the seventh order $K_7\,(A_1^2\ldotp\ldotp\ldotp\ldotp A_{12}^2\ B_1\ldotp\ldotp\ldotp\ldotp B_5)$. %

Having thus proved that irreducible canonical surfaces $F_8$ exist, one is also led to admit that the \textit{number of moduli} on which they depend is exactly
\[
M = 10p - 2p^{(1)} + 12 = 34,
\]
and hence that all the $C_{12}$ obtained from the bisection problem explained above are effectively double curves of canonical surfaces $F_8$. This assertion is justified at least if one admits that imposing a quadruple point on our $F_8$ entails 17 conditions for them, and not fewer. Indeed, in this way one recovers the number $(34 - 17 = 17)$ of the absolute invariants of the rational $F'_8$, as it can be evaluated from the plane representation.

\textit{Note.} -- For $p = 4$ and $p^{(1)} = 9$ too, one meets multiple canonical surfaces, which can give rise to various cases. We shall confine ourselves to pointing out two cases. First of all, there is the double $F_4$ with plane sections of genus 2 (that is, endowed with a double line), with branch curve of order 12, which can be represented on the \textit{double plane with branch curve $C_{14}$ of order $14$, endowed with a $6$-fold point $O$ and with $8$ $4$-fold points, $A_1A_2\ldotp\ldotp\ldotp\ldotp A_8$}. Here the canonical curves have as (double) images the quartics through $O^2A_1A_2\ldotp\ldotp\ldotp\ldotp A_8$, so that one has precisely
\[
p = 4 \qquad \text{and} \qquad p^{(1)} = 9.
\]

It is easy to recognize that the said double plane, or the corresponding $F_4^2$, belongs, as a \textit{particular case, to the family of surfaces with $p = 4$ and $p^{(1)} = 9$} defined by the simple $F_8$. This follows, with complete certainty, from the count of moduli, because the projective invariants of the branch $C_{14}$ of the said double plane number
\[
119 - 19 - 64 - 8 = 28,
\]
and this number turns out to be less than the number of moduli of the family to which the $F_4^2$ belong, which is certainly
\[
M \geq 9p - 2p^{(1)} + 12 = 30.
\]

% scan: s154R
% !PRINTED 283
\weberpage{283}
A second notable case in which --- for $p = 4$ and $p^{(1)} = 9$ --- one meets a multiple canonical surface is that of the double cubic $F_3$, with branch curve of order 16, which can be represented on the \textit{double plane with branch $C_{12}$ endowed with $4$ $4$-fold points, $A_1A_2A_3A_4$, and with $2$ $[3,3]$ points: $B_1$ and $B_2$}. Here the images of the canonical curves $K$ are given by the cubics through $A_1A_2A_3A_4B_1B_2$; and one has two base points of $|K|$ corresponding to $B_1$ and $B_2$. Our double planes depend on 32 invariants or moduli. We shall not stop to investigate whether or not they belong to the family defined by the simple $F_8$, to which $10p - 2p^{(1)} + 12 = 34$ moduli belong.

\edsection{6. Surfaces with $p = 4$ and $p^{(1)} > 9$.}{6. Surfaces with p = 4 and p\textasciicircum{}(1) > 9.}{sec.VIII.6}

The difficulties which already appear in the construction of the canonical surfaces $F_8$ with $p = 4$ and $p^{(1)} = 9$ go on increasing when one passes to the higher values of $p^{(1)}$. The method we have followed leads first of all to the bisection of a linear system on a surface $\Phi_{n-5}$ ($n = p^{(1)} - 1$), a problem which no longer lends itself to so simple a solution when non-rational surfaces are concerned. This case already presents itself for the $F_9$ with $p^{(1)} = 10$, for here one meets an adjoint surface $\Phi_4$ of order 4, with
\[
t = 10
\]
double points, and it is a matter of halving the linear system cut on the $\Phi_4$ by the surfaces $F_9$ of order 9 passing triply through these 10 double points.

The problem which thus presents itself for $p^{(1)} = 10, 11,\ldotp\ldotp\ldotp\ldotp$ gives rise to the following observations: %

$1^{\mathrm{o}}$) One cannot in general assert that the bisection of the system cut on a $\Phi_{n-5}$ with $t$ double points by the $F_n$ passing triply through them is possible, nor even that it gives a uniquely determined system.

$2^{\mathrm{o}}$) In any case one would still have to decide whether the hypothetical curves $C$ of the aforesaid half system, all of them or in part, are effectively double curves for irreducible surfaces $F_n$.

3) From observations $1^{\mathrm{a}}$) and $2^{\mathrm{a}}$) it follows that the (simple) canonical surfaces of ordinary space ($p = 4$, $p^{(1)} = n + 1$), instead of constituting a single continuous system, may give rise, for the values of $n > 8$, to several distinct families (perhaps possessing the same number of moduli $M = 10p - 2p^{(1)} + 12 = 50 - 2n$). This cir-

\editiongap

% scan: s155L
% !PRINTED 286
\weberpage{286}
\edsection{8. Surfaces with $p = 5$ and $p^{(1)} = 10$.}{8. Surfaces with p = 5 and p\textasciicircum{}(1) = 10.}{sec.VIII.8}

Passing now to the surfaces $F$ with $p = 5$ and $p^{(1)} = 10$ (with irreducible, simple canonical system), we shall first of all give an example in which it is easy to construct the corresponding canonical surface $F_9$, of order 9 ($= p^{(2)} = p^{(1)} - 1$), in the space $S_4$. Indeed, consider in $S_3$ the $F = F_6$ of the $6^{\mathrm{o}}$ order with 5 triple points $A_1A_2\ldotp\ldotp\ldotp\ldotp A_5$, for which one has precisely $p = 5$ and $p^{(1)} = 10$. The quadrics of $S_3$ passing simply through the 5 aforesaid points are the surfaces $\Phi_2$ ($= \Phi_{n-4}$) adjoint to the $F_6$ ($= F_n$), cutting the canonical curves on $F_6$. Now the system of such $\Phi_2$ represents a cubic variety $V_3^3$ of the space $S_4$. This is a particular $V_3^3$ (studied by \textsc{Segre} and \textsc{Castelnuovo}) which is endowed with 10 double points corresponding to the 10 edges of the pentahedron $A_1A_2\ldotp\ldotp\ldotp\ldotp A_5$. On this $V_3^3$ the cubic varieties of $S_4$ (which in general do not pass through any of the double points) cut surfaces $F_9$ of the $9^{\mathrm{o}}$ order, having as images in $S_3$ the $F_6$ passing triply through $A_1A_2\ldotp\ldotp\ldotp\ldotp A_5$. Therefore the $F_6$ of $S_3$ can be transformed into a canonical surface $F_9$ of $S_4$, complete intersection of two cubic varieties of this space, one of which is a $V_3^3$ endowed with 10 double points.

But the $F_6$ with 5 triple points from which we started constitute only a particular family of surfaces with $p = 5$ and $p^{(1)} = 10$. It is clear that every other surface with these characters will in general give rise to a canonical $F_9$ of $S_4$, complete intersection of two cubic varieties. Indeed, on such an $F_9$ one easily recognizes that the $\infty^{34}$ cubic varieties $S_4$ cut a complete linear system (adjoint to the system cut by the quadrics) of dimension 32. %

On the other hand, all the $F_9$ of $S_4$ that are complete intersections of two cubic varieties form a continuous system with the same characters $p = 5$ and $p^{(1)} = 10$. One can therefore say that the \textit{general type of the simple surfaces corresponding to the genera $p = 5$ and $p^{(1)} = 10$ is the surface $F_9$ of $S_4$, complete intersection of two cubic varieties}, which depends on
\[
2 \cdot 34 - 2 - 24 = 42
\]
absolute invariants and therefore possesses $M$ moduli:
\[
M = 10p - 2p^{(1)} + 12 = 42.
\]

It is interesting to see how the said canonical surface can be constructed by the method explained in the preceding sections, starting from a projection of it of the same order. This projection is a non-normal canonical surface $F'_9$ of the space $S_3$, endowed with a double curve $C_{18}$ of order 18 belonging to a surface of the $4^{\mathrm{o}}$ order $\Phi_4$. For this purpose it is useful to note that the projection $F'_9$ of an $F_9$, made from an external point $O$, must possess $t = 6$ triple points, $B_1B_2\ldotp\ldotp\ldotp\ldotp B_6$, which
% scan: s155R
% !PRINTED 287
\weberpage{287}
correspond to the 6 lines through $O$ of the cubic variety passing through $F_9$ and through $O$ itself\origfootnote{1}{s155R-1}{It is known and obvious that on a $V_3^3$ of $S_4$, and through a point $O$ of it, there are 6 lines. These belong to the surface section of $V_3^3$ by the space $S_3$ tangent at $O$, and therefore lie on the quadric cone osculating this surface section at $O$.}, and hence belong to a conic $C_2$\origfootnote{2}{s155R-2}{A priori one might believe that the surfaces of the $4^{\mathrm{o}}$ order with 6 double points on a conic constitute a system contained in another, wider system of $\Phi_4$ with 6 double points no longer lying on a conic. In reality the surfaces of the $4^{\mathrm{o}}$ order with 6 double points are distributed into several continuous systems $\infty^{28}$, one of which, not enlargeable, consists precisely of the $\Phi_4$ with the 6 double points on a conic. That there exist $\Phi_4$ with 6 double points not lying on a conic is seen at once. One fixes 5 double points in generic position, $A_1$, $A_2,\ldotp\ldotp\ldotp,$ $A_5$, and then imposes a double point (in unassigned position) on the $\Phi_4$ of the linear system of dimension $34 - 20 = 14$ thus defined. In this way one obtains $\infty^{28}\Phi_4$ with 6 double points not lying in a plane. But, on the other hand, suppose one requires a $\Phi_4$ to possess a tangent plane along a conic $C$. Then the $\infty^4$ sextics cut on the $\Phi_4$ by the quadrics through $C$ (among which there are $\infty^2$ quadrics split into the plane of $C$ and another variable plane) will meet $C$ in 6 fixed points, which turn out to be double for $\Phi_4$. Now it is easy to see that imposing a tangent plane along a conic entails precisely 6 conditions for $\Phi_4$ (which are expressed by requiring that $\Phi_4$ contain a conic self-residual with respect to the system of plane sections). One therefore deduces that the $\Phi_4$ with 6 double points on a conic likewise form a continuous system $\infty^{28}$ of surfaces of the $4^{\mathrm{o}}$ order with 6 double points. %

We believe that on a surface $\Phi_4$ endowed with 6 double points not lying on a conic it is not generally possible to bisect the system cut by the surfaces of the $9^{\mathrm{o}}$ order passing triply through the 6 points. The problem would have to be examined by representing the $\Phi_4$ on a double plane, by projection from one of its double points.}.

Now one sees that the double $C_{18}$ of $F'_9$ must belong to a system which is the half of that cut on the adjoint $\Phi_4$ by the surfaces of $9^{\mathrm{o}}$ order passing triply through the 6 points $B_1B_2\ldotp\ldotp\ldotp\ldotp B_6$, which are double for the said $\Phi_4$. The bisection of this system $|C_{36}|$ is easily obtained by observing that the $C_{18}$ of the required half system can be defined as curves belonging to a sum. This is the sum of the system $|C_{16}|$, cut on $\Phi_4$ by the general surfaces of the $4^{\mathrm{o}}$ order passing through $B_1B_2\ldotp\ldotp\ldotp\ldotp B_6$, and of the conic $C_2$ along which the plane of $C_2$ touches $\Phi_4$. %

The $\Phi_4$ is a regular surface of genus 1, with canonical curve of order zero, and therefore every linear system without base points (at simple points) on it has dimension equal to the genus $\pi$. Hence the system $|C_{16}|$, of degree $64 - 2 \cdot 6 = 52$, has genus equal to the dimension $26 + 1 = 27$. Adding to the $|C_{16}|$ the conic $C_2$ of genus 0, which meets the $C_{16}$ (outside $B_1B_2\ldotp\ldotp\ldotp\ldotp B_6$) in 2 points, one obtains $C_{18}$ of genus
\[
\pi = 27 + 0 + 2 - 1 = 28.
\]

% scan: s156L
% !PRINTED 288
\weberpage{288}
Hence the dimension of the system $|C_{18}|$ is also 28, and it exceeds by 1 that of the system $|C_{16}|$ to which the fixed conic $C_2$ is added: this means that the system $|C_{18}|$ is irreducible. Thus we have obtained on the $\Phi_4$ with 6 double points $B_1B_2\ldotp\ldotp\ldotp\ldotp B_6$ lying on a conic, $\infty^{28}$ curves $C_{18}$, generally irreducible, endowed with 6 triple points at $B_1B_2\ldotp\ldotp\ldotp\ldotp B_6$, which are a priori « possible » double curves for non-normal canonical surfaces $F'_9$. Let us admit that among these $C_{18}$ there are $\infty^x$ $(x\leq 28)$ through which actual irreducible surfaces $F'_9$ pass doubly. Each of these $C_{18}$ then belongs, as a double curve, to $\infty^5 F'_9$ that are irreducible, there being already one, and hence $\infty^4 F'_9$ split into the $\Phi_4$ and an $F_5$ through it. One will therefore have $\infty^{x+5} F'_9$ with double curve $C_{18}$ belonging to a given $\Phi_4$ (endowed with 6 double points on a conic). But there are $\infty^{14}\Phi_4$ with the same 6 double points and, by changing these, one obtains $\infty^{28}\Phi_4$, so that one will have, in the space $S_3$, $\infty^{28+x+5}F'_9$. Among these there are $\infty^{15}$ projectively identical with a given one, and $\infty^{19}$ birationally identical with it, corresponding to projections of $F_9$ out of $S_4$ from $\infty^4$ different centers. Therefore the number of birationally distinct $F'_9$ and $F_9$, that is, the number of moduli of the said canonical surfaces, turns out to be:
\[
M=28+x+5-19=x+14,
\]
and since on the other hand one knows that
\[
M=42,
\]
one deduces
\[
x=28.
\]

This means that « \textit{all the curves $C_{18}$ obtained by the bisection indicated above are double curves of non-normal canonical surfaces $F'_9$, actually existing and generally irreducible} ».

We conclude by proposing to the student some

\textit{Questions.} -- Whether, among the surfaces with simple canonical system, for $p=5$ and $p^{(1)}=10$, there exist or not surfaces on which this system possesses a base point, leading therefore to canonical surfaces $F_8$ of $S_4$, or to $F'_8$ of $S_3$, with double curve $C_{11}$ of order 11. And whether such surfaces belong as particular cases to the family of the $F_9$, or of the $F'_9$, with 42 moduli, which we have described. Whether, on the other hand, there belong as particular cases to the family of the $F_9$, or of the $F'_9$, surfaces in which the canonical system (possessing a base point) is composed with a rational involution $I_2$. These are $F_8$ which reduce to double $F_4$, representable on the double plane with branch curve $C_{12}$ of order 12, endowed with 4 4-fold points and a $[3,3]$ point: note that this double plane has exactly the characters $p=5$ and $p^{(1)}=10$, and that it depends on only 41 moduli.

% scan: s156R
% !PRINTED 289
\weberpage{289}
\textit{Note.} -- Passing to the case $p=5$ and $p^{(1)}=11$, observe now the \textit{double plane with branch $C_{12}$ endowed with 3 4-fold points and 2 $[3,3]$ points, which has the genera $p=5$ and $p^{(1)}=11$ and the (non-simple) canonical system with 2 base points}. This double plane, which leads to a double canonical quartic with branch $C_{18}$, depends on 40 moduli. Suppose the simple canonical surfaces $F_{10}$ of $S_4$ (with $p=5$ and $p^{(1)}=11$) likewise depend on $M=10p-2p^{(1)}+12=40$ moduli. This means that, as is to be considered probable, one will have for $p=5$ and $p^{(1)}=11$ \textit{two distinct families} of surfaces: the $F_{10}$ and the double $F_4$ indicated. Here it is worth pointing out that simple canonical $F_{10}$ of the space $S_4$ certainly exist; particular $F_{10}$ (endowed with a double line) are obtained starting from an $F_6$ of ordinary space with 4 triple points and 1 tacnode. %

\edsection{9. Surfaces on a variety of 3 or more dimensions.}{9. Surfaces on a variety of 3 or more dimensions.}{sec.VIII.9}

Instead of pursuing the analysis of the preceding sections, we prefer to indicate how notable examples of hyperspatial canonical surfaces can be obtained by direct construction, in relation to the geometry of varieties of more than 2 dimensions.

When one passes from surfaces to varieties of 3 or more dimensions, one meets three orders of properties:

$1^\circ$) properties of a general nature, which present themselves as a natural extension of those of curves and surfaces, and which it is therefore legitimate to regard as known; some of these are already indicated by \textsc{M. Noether} in his fundamental memoir in the Math. Ann. Bd. 8;

$2^\circ$) properties more or less foreseen or foreseeable as an extension of those of surfaces, which may however require less evident, and sometimes even difficult, proofs;

$3^\circ$) entirely new and unexpected properties which have no counterpart whatever in the theory of surfaces\origfootnote{1}{s156R-1}{As an example we recall the results of \textsc{Fano} and \textsc{Enriques} on questions of rationality; non-rational involutions of space; various types of varieties containing a congruence of rational curves, etc.

Among the works on the geometry of varieties we cite: \textsc{F. Severi}, \textit{Fondamenti per la geometria sulle varietà algebriche} (Rend. del Circ. Mat. di Palermo, vol. XXVIII, 1909). \textsc{A. Rosenblatt}, \textit{Varietà algebriche a tre e più dimensioni} (Atti del Congresso Internazionale dei Matematici, 1928), which is a complete report on the results achieved in this field. Extensive bibliographical indications are also to be found in it. \textsc{B. Segre}, \textit{Quelques résultats nouveaux dans la géométrie sur une $V_3$ algébrique} (Mémoires de l'Académie R. de Belgique 2\textsuperscript{e} s., t. XIV, 1936).}.

\printersignature{\textsc{Enriques F.} - \textit{Superficie algebriche.}\quad 19}

% scan: s157L
% !PRINTED 290
\weberpage{290}
We shall take into consideration only the first order of problems. These concern the general theory of the linear systems of surfaces belonging to a variety of 3 dimensions, the adjoint systems, the canonical system, and the invariants to which it gives rise.

For simplicity of exposition we refer to a variety of 3 dimensions $V$, free from singularities, in a suitable hyperspace, and we consider on it linear systems of surfaces free from base points or base curves. A surface $F$ on $V$ belongs in general to a complete linear system $|F|$, and one can operate on such systems --- as on linear systems of curves on a surface --- by \textit{addition and subtraction}.

If the system $|F|$ has dimension at least equal to three, every triply infinite system contained in $F$ will possess a \textit{Jacobian surface} $F_j$, locus of the double points of its surfaces. The $F_j$ relative to the various systems totally contained in $|F|$ are \textit{equivalent} to one another and determine a complete linear system: the \textit{Jacobian system} of $|F|$, which we shall denote by $|F_j|$. %

Analogously to what happens for systems of curves on a surface, there holds the \textit{fundamental theorem} expressed by the relation:
\[
|(F+\Phi)_j|=|F_j+4\Phi|=|4F+\Phi_j|,
\]
and the difference
\[
|F_j-4F|
\]
defines the \textit{canonical system} of $V$\origfootnote{1}{s157L-1}{It is hardly necessary to point out that the canonical surfaces, thus defined with respect to birational transformations on the variety $V$, are not to be confused with the \textit{projectively canonical} surfaces. The latter are defined, in a space $S_{p-1}$, by the property that their hyperplane sections are canonical curves (of the same surface).}, effective or virtual, while the \textit{adjoint system} of $F$ is given by
\[
|F'|=|F_j-3F|.
\]

If the variety $V_n$ of order $n$ belongs to the space of 4 dimensions $S_4$, it will in general possess a double surface $\Gamma$, whose points correspond to the \textit{neutral pairs} of the system of hyperplane sections on a transformed variety. Unless there are more complicated singularities, the canonical system is cut by the \textit{adjoint varieties} $\Phi_{n-5}$ of order $n-5$, defined simply by passing through the double surface $\Gamma$. The adjoints $\Phi_{n-4}$, on the other hand, will cut on the $V_n$ the surfaces adjoint to the hyperplane sections.

When the variety $V_n$ of $S_4$ possesses higher singularities, that is, a multiple surface of order $i>2$, or an $i$-fold curve $(i>2)$, or an isolated $i$-fold point $(i>3)$, the varieties $\Phi$ \textit{adjoint} to
% scan: s157R
% !PRINTED 291
\weberpage{291}
$V_n$ will be defined in general by the conditions of passing $i-1$ times through every $i$-fold surface, $i-2$ times through an $i$-fold curve and $i-3$ times through an isolated $i$-fold point.

On a variety $V$ of 3 dimensions (whether or not it belongs to a space $S_4$) the surfaces $F'$ of the system adjoint to a system $|F|$ cut a generic $F$ in canonical curves; and, conversely, this property characterizes the adjoint surfaces $F'$ when the system $|F|$ has no fundamental surfaces.

The characters of the canonical system of a variety furnish as many \textit{invariant characters} of the $V$ itself, already met by \textsc{Noether}. They are:

$1^\circ$) the dimension of the canonical system, increased by one, that is, the number of linearly independent canonical surfaces, which is the \textit{principal geometric} genus $P$ (or $P_g$) of $V$;

$2^\circ$) the geometric genus $p_g$ and the linear genus $p^{(1)}$ of the canonical surfaces, which are defined respectively as the \textit{surface genus $p_g$ and linear genus $p^{(1)}$} of the variety $V$.

The genus $x$ of the curve $K$, intersection of two canonical surfaces, and the degree $y$ of the canonical system (number of points common to 3 canonical surfaces) do not constitute new characters of $V$. They are in fact expressed in terms of the preceding ones by means of the formulas
\[
\begin{aligned}
2x+y-1&=p^{(1)}\\
4y&=p^{(1)}-1=p^{(2)},
\end{aligned}
\]
which express the fact that $|2K|$ is the canonical system of the canonical surface. There are therefore only three independent characters, as has already been noted by \textsc{Noether}.

Alongside the geometric characters one can define for $V$ the arithmetic characters: the \textit{principal arithmetic genus} $P_a$ and the \textit{surface arithmetic genus} $p_a$\origfootnote{1}{s157R-1}{Cf. \textsc{Severi}, l. c.}.

From what has been said it follows that, if one has in $S_4$ a surface $F$, complete intersection of two varieties $V^m$ and $V^n$ of orders $m$ and $n$, free from singularities on the $F$, the canonical system is cut by the varieties of order $m+n-5$\origfootnote{2}{s157R-2}{Cf. \textsc{Severi}, \textit{Su alcune questioni di postulazione.} Rend. Circ. Mat. di Palermo, vol. XVII, 1903.}; here we have the extension of the well-known theorem of Noether for space curves\origfootnote{3}{s157R-3}{Cf. \textsc{Enriques-Chisini}, op. cit., Book V, chap. V, \S~41 (vol. III, p. 528).}.

If the $F$ is not the complete intersection of $V^m$ and $V^n$, but these meet further in another surface $F^\alpha$, then on the $F$ the canonical system is cut by the varieties of order $m+n-5$ passing through $F^\alpha$. %

% scan: s158L
% !PRINTED 292
\weberpage{292}
Similarly, suppose that in the space of $r$ dimensions one has $r-2$ hypersurfaces (varieties of $r-1$ dimensions), of orders $n_1,n_2,\ldotp\ldotp\ldotp\ldotp,n_{r-2}$, free from singularities, meeting in a composite surface $F+F^\alpha$. Then on the $F$ the canonical curves are cut by the hypersurfaces of order $n_1+n_2+\ldotp\ldotp\ldotp\ldotp+n_{r-2}-r-1$ passing through $F^\alpha$. %

\edsection{10. Examples of hyperspatial canonical surfaces.}{10. Examples of hyperspatial canonical surfaces.}{sec.VIII.10}

The preceding notions lead to simple examples of projectively canonical surfaces in hyperspaces. First of all, consider in the space $S_4$ the surface which is the complete intersection of two varieties $V^m$ and $V^n$ without singular points. This will be a canonical surface if
\[
m+n-5=1,
\]
since, by what has been said, the hyperplanes will then cut the canonical curves on it. Here one recovers the two types already indicated above: the surface $F_8$, intersection of a quadric and a variety $V_3^4$, for which $p=5$ and $p^{(1)}=9$; and the surface $F_9$, intersection of two cubic varieties, for which $p=5$ and $p^{(1)}=10$.

But, more generally, the complete intersections of $r-2$ hypersurfaces, of orders $n_1,n_2,\ldotp\ldotp\ldotp\ldotp,n_{r-2}$, will be canonical surfaces of the space $S_r$ when these orders $(\geq 2)$ are such that one has
\[
n_1+n_2+\ldotp\ldotp\ldotp\ldotp+n_{r-2}-r-1=1,
\]
that is,
\[
n_1+n_2+\ldotp\ldotp\ldotp\ldotp+n_{r-2}=r+2.
\]

This relation, taking into account that one must have $n_i\geq 2$, furnishes some notable canonical surfaces, such as the following.

In the space $S_5$ one has the surface $F_{12}$ with $p=6$ and $p^{(1)}=13$, complete intersection of two quadrics and a cubic variety. And in the space $S_6$ one has the $F_{16}$ with $p=7$ and $p^{(1)}=17$, complete intersection of 4 quadrics.

These examples belong to more general families of canonical surfaces. These can be defined as the intersection of a normal variety of 3 dimensions of a hyperspace, with elliptic or projectively canonical curve sections, with a hypersurface of order 3 or 2 respectively.

For the $V_3^n$ of the space $S_r$, with $n=2r-4$, with projectively canonical curve sections\origfootnote{1}{s158L-1}{Varieties of this type are the cones projecting a surface with all genera equal to 1 from a point outside the space containing it. For the first values of the order these cones are contained in wider families of varieties without singular points. In any case the reasoning of the text seems extendable to the cones as well, with slight modifications.}, which we suppose free from singular points,
% scan: s158R
% !PRINTED 293
\weberpage{293}
it will be noted that their hyperplane sections $F$ are surfaces of genus 1 with canonical curve of order zero. Therefore the system $|F'|$ adjoint to $F$ consists only of the surface of order zero. Consequently the system adjoint to $|2F|$ is
\[
|(2F)'|=|F|;
\]
hence the surfaces $F$, that is, the hyperplanes of the space $S_r$, cut canonical curves on the surfaces of $|2F|$, that is, on the intersections of the quadrics.

As for the $V_3^n$ of the space $S_{n+1}$ $(n>3)$ with elliptic curve sections\origfootnote{1}{s158R-1}{The varieties of 3 dimensions with elliptic curve sections are classified in the note of \textsc{F. Enriques}, \textit{Sui sistemi lineari di superficie algebriche le cui intersezioni variabili sono curve ellittiche.} Rendic. Acc. Lincei, May-June 1894. Cf. Math. Annalen, Bd. 46.}, which are rational, having normal sections, it will suffice to note the following. The intersections of the quadrics, and hence the surfaces of the system $|2F|$, double of the system $|F|$ of hyperplane sections, are surfaces with genera equal to 1 and with canonical curve of order zero. Indeed, the surfaces of $|2F|$ have as hyperplane sections curves of order $2n$, of genus $n+1$, belonging to spaces of $n$ dimensions, and hence canonical curves. It follows from this that the system $|2F|$ possesses as its only adjoint surface the surface of order zero, and therefore the system adjoint to $|3F|$ is
\[
|(3F)'|=F, %
\]
so that the surfaces $F$, that is, the hyperplanes of the space $S_{n+1}$, cut on the intersection of $V_3^n$ with a cubic variety its canonical curves.

In conclusion, we have obtained \textit{two series of canonical surfaces of genus $p$ in the space $S_{p-1}$}:

$1^\circ$) \textit{the surfaces $F_{3p-6}$, of order $3p-6$ $(r=p-1,n=p-1)$, with $p^{(1)}=3p-5$, complete intersections of a cubic hypersurface with a rational normal variety with elliptic curve sections, $V_3^{p-2}$}: %

$2^\circ$) \textit{and the surfaces $F_{4p-12}$, of order $4p-12$, with $p^{(1)}=4p-11$ $(r=p-1,n=2p-6)$, complete intersections of a quadric with a variety $V_3^{2p-6}$ with projectively canonical curve sections}.

We shall see later other classes of canonical surfaces defined on the variety of 3 dimensions with rational curve sections. Here we shall not dwell further on the construction of the canonical surfaces of hyperspaces. We confine ourselves to pointing out recent works, connected in some way with the subject, in which this examination is pursued, with the construction of other interesting examples: \textsc{G. Fano}, « Superficie algebriche e varietà a 3 dimensioni a curve sezioni canoniche », Atti Acc. Lincei, 1936; « Su alcune varietà al%
% scan: s159L
% !PRINTED 294
\unskip\weberpage{294}%
gebriche a 3 dimensioni aventi curve sezioni canoniche » in « Scritti matematici offerti a L. Berzolari », p. 329 (1936); « Sulle varietà algebriche a tre dimensioni a curve sezioni canoniche », Mem. Acc. d'Italia, 1937.

\textsc{L. Godeaux}, « Sur la construction d'une surface canonique », Bull. de la Soc. Math. de Liège 5, (1936); « Sur une surface canonique appartenant à la variété de Segre représentant les couples de points de deux plans », Bull. de l'Académie Belgique, 22, (1936); « Sur les surfaces algébriques dont les sections hyperplanes sont des courbes canoniques », I, II, III, Bull. de la Soc. Math. de Liège, 6, (1937), « Construction d'une surface canonique du septième ordre », Bull. de la Soc. des Sciences de Liège, 1944; « Construction d'une surface canonique du huitième ordre », Bull. des Sciences math., 1944; « Construction d'une surface canonique du neuvième ordre », Bull. de l'Acad. de Belgique, 1944.

\textsc{P. Bourniat}, « Note sur quelques surfaces canoniques et sur des surfaces de genre un », Bull. de l'Acad. de Belgique 5, 22, (1936); « Note sur quelques surfaces canoniques », Acad. de Belgique 3, 22, (1936); « Sur des hypersurfaces et variétés canoniques », Acad. de Belgique (5), 23, (1937).

\textsc{E. A. Maxwell}, « Regular canonical surfaces of genus three and four », Proc. Cambridge phil. Soc., 33, (1937.)

\edsection{11. Minimum value of the linear genus with respect to the surface genus.}{11. Minimum value of the linear genus with respect to the surface genus.}{sec.VIII.11}

Already from the construction of the canonical surfaces it appears clear that the (absolute) linear genus $p^{(1)}$ of a surface cannot descend below a certain minimum value, corresponding to a given surface genus $p$, (or $p_g$). This holds at least if the canonical curves are irreducible, which excludes, for $p>2$, the surfaces with $p^{(1)}=1$. We now propose to prove, without any exception whatever, that \textit{for surfaces with $p^{(1)}>1$ the minimum value of $p^{(1)}$ with respect to $p$ or $p_g$ $(>2)$ is given by}
\[
p^{(1)}=2p-3.
\]

In other words one has
\[
p^{(1)}\geq 2p-3,
\]
or, what is the same,
\[
p\leq \frac{p^{(1)}+3}{2}:
\]
\textit{the maximum of $p$ with respect to $p^{(1)}$ $(>1)$ is $\dfrac{p^{(1)}+3}{2}$.}

This inequality is established by \textsc{Noether}, for the case of irreducible (pure) canonical curves, in his memoir in the Math.
% scan: s159R
% !PRINTED 295
\weberpage{295}
Ann. Bd. 8 (1875). If the canonical curves $K$ are supposed irreducible, it is a consequence of the theorem of Clifford, according to which the order of a special series $g_n^r$ on any curve is $n\geq 2r$\origfootnote{1}{s159R-1}{Cf. \textsc{Enriques-Chisini}, Op. c., Book V, chap. I, \S~11 (vol. III, p. 87).-}. Indeed, for the characteristic series of the system $|K|$ the degree is $n=p^{(1)}-1$ and the dimension $r=p-2$, so that %
\[
p^{(1)}-1\leq 2p-4. %
\]

But the same inequality is valid even when the canonical $K$ are admitted to be reducible, provided that the variable part splits into curves of genus $\pi\geq 3$. First of all, if their variable components $k$ are irreducible, their degree is $n\leq p^{(1)}-1$ (Chap. IV, \S~18), while the dimension of $|k|$ is the same as that of $|K|$, that is, $p-1$; one therefore has
\[
n\geq 2p-4
\]
and a fortiori
\[
\begin{aligned}
p^{(1)}-1&\geq 2p-4\\
p^{(1)}&\geq 2p-3.
\end{aligned}
\]
If, then, the variable $k$ are supposed composed of $p-1$ curves of a pencil, of genus $\pi\geq 3$, the genus of $|k|$ $(\leq p^{(1)})$ will be
\[
\geq (p-1)(\pi-1)+1;
\]
one deduces
\[
p^{(1)}\geq 2p-1.
\]
Now \textit{if for a surface the genus $p$ $(>2)$ attains the maximum $p=\dfrac{p^{(1)}+3}{2}$ $(p^{(1)}>2)$, the canonical curves are irreducible and hyperelliptic.}

Indeed, if the $K$ are irreducible, for the characteristic series $g_n^r$ on a $K$ $(n=p^{(1)}-1,r=p-2)$ one must have $n=2r$, and this entails precisely that the $K$ is hyperelliptic\origfootnote{2}{s159R-2}{Cf. \textsc{Enriques-Chisini}, Op. c., Book V, \S~12, vol. III, p. 97.}. If, on the other hand, the $K$ were reducible, the reasoning would be repeated for their variable components, which are certainly irreducible (otherwise $p^{(1)}\geq 2p+1$). The $K$ thus turn out to be hyperelliptic and lead to a double canonical surface, referable to the double plane which we indicate below, for which it turns out a posteriori that there are no fixed components of $|K|$. %

It is noteworthy that, for our surfaces with maximum $p$, there is thus also excluded the case in which the pure canonical curves $K$ contain exceptional fixed components, that is, in which $|K|$ possesses some base point (on a model free from exceptional curves). Indeed, if there are $r$ base points, one finds
\[
p^{(1)}\geq 2p-3+s. %
\]

% scan: s160L
% !PRINTED 296
\weberpage{296}
The canonical surface corresponding to a maximum $p$ $(p^{(1)}=2p-3)$ is a normal surface with rational curves as hyperplane sections, of order $p-2$, in $S_{p-2}$. Therefore it can be represented on the plane by means of a linear system of curves belonging to one of the following types\origfootnote{1}{s160L-1}{Cf. \textsc{Enriques-Chisini}, Book V, chap. II, \S~22, vol. III, p. 194.}: %

I) net of lines,

II) system $\infty^5$ of conics,

III) system of curves of order $n$ with a base point $O$ of multiplicity $n-1$, and possibly also another simple base point distinct from $O$, or several simple base points infinitely near to $O$. %

The double planes whose canonical systems have such systems of rational curves as images are respectively:

I) the double plane with branch curve $C_8$ of the $8^\circ$ order: $p=3$, $p^{(1)}=3$;

II) the double plane with branch curve $C_{10}$ of the $10^\circ$ order; $p=7$, $p^{(1)}=9$; %

III) the double plane with branch curve $C_{2n+6}$, of order $2n+6$, having a $2n$-fold point $O$: $p=2n+1$, $p^{(1)}=4n-2$; and %

III$'$) particular cases of type III) where the branch curve (which can be supposed to be of order not lowerable by quadratic transformations) possesses, besides $O$, a 6-fold or 4-fold point, with $p=2n-2$, $p^{(1)}=4n-8$ or respectively $p=2n$, $p^{(1)}=4n-4$. %

Indeed, as has been said, those types of double planes are to be excluded for which the characteristic series of the pure canonical system $|K|$ has some fixed point (base point of $|K|$): thus, e.g., a point $[5,5]$ or a point $[3,3]$.

Meanwhile we state that: \textit{the surfaces of surface genus $p>2$ with $p^{(1)}>1$, for which the absolute linear genus $p^{(1)}$ attains the minimum, $p^{(1)}=2p-3$, are representable on a double plane belonging to one of the types} I), II) III) and III$'$)\origfootnote{2}{s160L-2}{Cf. \textsc{F. Enriques}, \textit{Sopra le superficie algebriche di cui le curve canoniche sono iperellittiche.} Rendic. Acc. Lincei, series V, vol. V, p. 191 (1896).}.

For Noether's inequality
\[
p\leq \frac{p^{(1)}+3}{2}
\]
one can substitute another, more expressive inequality, which holds for surfaces with simple canonical system $(p>3)$, and which was
% scan: s160R
% !PRINTED 297
\weberpage{297}
brought to light by \textsc{Castelnuovo}\origfootnote{1}{s160R-1}{\textsc{G. Castelnuovo}, \textit{Osservazioni intorno alla geometria sopra una superficie.} Note II. (Rendic. Ist. Lomb., series II, vol. XXIV, 1891). «Memorie Scelte», no. XVIII.}. \textit{When the canonical system is irreducible and simple, the surface genus ($p$ or $p_g>2$) cannot exceed the value $\dfrac{p^{(1)}+6}{3}$}; one therefore has
\[
p\leq \frac{p^{(1)}+6}{3},\qquad p^{(1)}\geq 3p-6.
\]

To prove this inequality it is convenient to start from a lemma of Bertini concerning the self-residual series $g_n^r$ which can belong to a curve of genus $\pi$. %
On the curve containing such a $g_n^r$ let us take two groups $G$ and $G'$ of the series having in common a group $\Gamma$ of $r-1$ generic points. If the series is simple (that is, not composed with an involution), $G$ and $G'$ will have no other points in common besides the $r-1$ points of $\Gamma$, and the group $\Gamma$ cannot belong to a linear series $g_{r-1}^s$ of dimension $s>0$. This being so, the Riemann-Roch theorem says. that the group $G$ of the $g_n^r$ imposes $n-r$ conditions on the canonical groups which are to contain it\origfootnote{2}{s160R-2}{Cf. e.g. \textsc{Enriques-Chisini}, l. c., Book V, chap. I, vol. III, p. 85.}. %
Moreover the group $G'$ will impose $n-r-(r-1)=n-2r+1$ conditions (a priori not independent) on the canonical groups which already contain $r-1$ of its points, and therefore on those which have already been constrained to contain the $G$. Hence the number of conditions for a canonical group to contain $G+G'-\Gamma$ is less than or equal to
\[
n-r+(n-2r+1)=2n-3r+1.
\]

Now the canonical groups satisfying these conditions cannot be more than one, that is, they coincide with the $G+G'$ in which the points of $\Gamma$ are double. Otherwise the $\Gamma$ (which consists of $r-1$ generic points of the curve, with $r<\pi$) would have to belong to a series $g_{r-1}^s$ with $s>0$. One deduces that
\[
2n-3r+1\geq \pi-1,
\]
$\pi-1$ being the number of conditions which suffice to determine the $G+G'$ within the canonical series.

But since the $g_n^r$ is self-residual,
\[
n=\pi-1,
\]
hence
\[
\pi-3r\geq 0.
\]

% scan: s161L
% !PRINTED 298
\weberpage{298}
Let us apply this formula to the characteristic series, which is simple, of the canonical system on a surface of genera \(p\) and \(p^{(1)}\), and we shall have
\[
n=\pi-1=p^{(1)}-1;\qquad r=p-2,
\]
and hence
\[
p^{(1)}-3p+6\geq 0,\qquad \text{Q.E.D.}
\]

\textsc{Castelnuovo}, after having written this inequality, which he deduces precisely from Bertini's formula, gives actual examples of surfaces with irreducible simple canonical system for which the genus \(p\) attains the maximum with respect to \(p^{(1)}\). Indeed he succeeds in determining the \textit{entire family of surfaces for which this maximum is attained}. The first examples are:

1°) the surface of the 5° order of \(S_3\), for which \(p=4\) and \(p^{(1)}=6\);

2°) the canonical surface \(F_8\) of the space of four dimensions, intersection of a quadric with a variety of the 4° order, for which, as we have seen, \(p=5\) and \(p^{(1)}=9\); %

3°) the surface obtained in \(S_6\) on the three-dimensional cone of the 4° order which projects a Veronese surface from a point \(O\). One obtains it by intersecting this cone with a hypersurface of the 4° order passing through a two-dimensional quadric cone, with vertex \(O\), lying on it:
\[
p=7,\qquad p^{(1)}=15.
\]

We leave it to the student to verify that the \(F_{14}\) thus constructed is effectively a canonical surface.

\textit{For \(p>7\) the simple canonical surfaces, of order \(p^{(1)}-1\), whose linear genus has the minimum value \(p^{(1)}=3p-6\), are obtained as intersections of a simply infinite variety of planes \(V_3^{p-3}\), of order \(p-3\), in \(S_{p-1}\) with a hypersurface of the 4° order, \(V_{p-2}^4\), passing through \(p-5\) planes of the \(V_3^{p-3}\).}

We shall confine ourselves to indicating how one justifies the fact that a surface \(F\) of order \(3p-7\) obtained in this way is effectively a canonical surface of \(S_{p-1}\). For this purpose it is convenient to recall how the variety \(V_3^{p-3}\) with rational hyperplane sections\origfootnote{1}{s161L-1}{The varieties of three dimensions with rational curve sections are: rational simply infinite series of planes, or cones of the 4° order projecting the Veronese surfaces from a point, such as are represented on ordinary space by quadrics tangent at a point. Cf. \textsc{F. Enriques}, \textit{Sui sistemi lineari di superficie algebriche ad intersezioni variabili iperellittiche}. Acc. Lincei, 1893. Cf. Math. Ann., Bd. 46.} is represented in space by projection from \(p-4\) of its points. The representative system consists of \(\infty^{p-1}\) surfaces of order \(p-3\) passing \(p-4\) times through a line \(r\) and simply through \(p-4\) other lines \(t\) meeting it, which are the traces of the planes through the centers of projection.

% scan: s161R
% !PRINTED 299
\weberpage{299}
One can suppose that the surface \(F\) contains the aforesaid \(p-4\) centers of projection. It then has as image in the representative \(S_3\) a surface \(F_{2p-3}\) containing \(r\) as a \((p-4)\)-fold line and, as triple lines, the \(p-4\) lines \(t\) meeting it. It is therefore evident that the surfaces \(\Phi_{2p-7}\) adjoint to the \(F_{2p-3}\) split into the \(p-4\) planes \(rt\) and a residual surface \(\Phi_{p-3}\) of order \(p-3\) which passes simply through the \(p-4\) lines \(t\). Thus it appears that the hyperplane sections of the normal surface \(F\) of \(S_{p-7}\) are the canonical curves, Q.E.D. %

\edsection{12. Upper limit of the linear genus \(p^{(1)}\) with respect to the surface genus.}{12. Upper limit of the linear genus p\textasciicircum{}(1) with respect to the surface genus.}{sec.VIII.12}

The very construction of the canonical surfaces, as well as the count of moduli, suggests that the linear genus \(p^{(1)}\) cannot exceed a certain maximum with respect to \(p\) (\(p_a\) or \(p_g\)). Although one has not yet succeeded in fixing the maximum actually attained by \(p^{(1)}\), it should be said that a first \textit{inequality} in this direction was written by \textsc{A. Rosenblatt}\origfootnote{1}{s161R-1}{\textsc{A. Rosenblatt}, \textit{Sur quelques inégalités etc.} Comptes rendus Ac. des Sciences, Paris, t. 154 (1912).}: %
\[
p^{(1)}\leq 16p_a+27.
\]

More recently \textsc{F. Severi}\origfootnote{2}{s161R-2}{Rendic. Circolo Mat. di Palermo, t. LVI, p. 79 (1932).}, relying on a transcendental result of \textsc{Hodge}, arrived at the \textit{inequality}:
\[
p^{(1)}\leq 2p_g+8p_a+11-\varrho,
\]
where \(\varrho\) denotes the base number: the number of independent continuous systems of curves traced on the surface free from exceptional curves.

The question is still open from the geometric point of view, even in the case of regular surfaces. Let us take up again the calculation of the moduli which (for \(p_a=p_g=p>3\) and for simple canonical surfaces) has led us to the formula
\[
M=10p-2p^{(1)}+12+\omega,
\]
one has to seek the upper limit, and possibly the maximum actually attained, of the superabundance \(\omega\). It is in fact a matter of recognizing the minimum number of conditions which the group of cuspidal points of a surface of order \(n\) imposes on the intersections of the adjoint surfaces of the same order that are to contain them.

% scan: s162L
% !PRINTED 300
\weberpage{300}
\edsection{13. Canonical systems belonging to an involution.}{13. Canonical systems belonging to an involution.}{sec.VIII.13}

The constructions of the preceding sections bring out the fact that the canonical system of a surface of genus \(p>3\), even when supposed irreducible, can belong to an involution \(I_n\) of an order \(n\geq 2\), in such a way that the corresponding canonical surface turns out to be multiple. Here it is worth saying something about this eventuality.

The simplest case, which we have met in the preceding examples, corresponds to the hypothesis that the surface contains a rational involution of the 2nd order \(I_2\), that is, that it can be represented on the double plane with an arbitrary branch curve \(C_{2n}\) of even order \(2n\): in this case --- as was already observed by \textsc{Noether} --- the canonical curves (having as images double curves of order \(n-3\), which form half of the second adjoint system to \(C_{2n}\)) always belong to the involution \(I_2\); one therefore has as canonical surface a double rational surface (possibly itself reducible to a multiple surface).

Particularly noteworthy is the fact that the possession of an involution follows necessarily from the hypothesis that the linear genus \(p^{(1)}\) is sufficiently small with respect to \(p\): \(p^{(1)}<3p-6\), since it has been recognized that, if the canonical system is irreducible and simple,
\[
p^{(1)}\geq 3p-6.
\]

However, the converse proposition does not hold, since there exist, for example, double planes with \(p^{(1)}>3p-6\); for \(n\) sufficiently large such is the double plane with branch curve \(C_{2n}\) entirely free of singularities, for which
\[
p=\frac{n(n-3)}{2}+1,\qquad p^{(1)}=2(n-3)^2+1
\]
where the ratio \(\frac{p^{(1)}}{p}\) for \(n=\infty\) tends to 4.

Double canonical surfaces always correspond to surfaces possessing an involution of the second order \(I_2\), but this \textit{involution is not necessarily rational}. Some examples of surfaces whose canonical curves belong to a non-rational involution \(I_2\) have already been pointed out by others, whose reference we do not recall. In our line of thought we can construct an interesting example of a double canonical surface, not referable to a double plane, starting from the canonical surface \(F_{10}\) with \(p=4\) and \(p^{(1)}=11\), whose double curve possesses \(t=20\) triple points. This \(F_{10}\) will admit an adjoint surface of the 5th order \(\Phi_5\), endowed with 20 conical points. In the pencil
\[
\lambda F_{10}+\Phi_5^2=0
\]
% scan: s162R
% !PRINTED 301
\weberpage{301}
one will see \(F_{10}\) reduce by continuity to \(\Phi_5^2\), on which there will be no proper branch curve, but rather \(t=20\) conical points constituting as many infinitesimal branch curves. The quintic \(\Phi_5\), with genera \(p=4\) and \(p^{(1)}=6\), thus constitutes a \textit{double canonical surface} with the genera \(p=4\) and \(p^{(1)}=11\). Let it suffice for us to have mentioned this case, which we point out to the willing reader as a subject of study.

The first examples of surfaces whose canonical system belongs to an involution \(I_n\) of order \(n>2\) have been pointed out by \textsc{B. Segre}\origfootnote{1}{s162R-1}{\textsc{B. Segre}, \textit{Sulle superficie algebriche aventi il sistema canonico composto con un’involuzione}. Rendic. Acc. Lincei, ser. VI, vol. XVI, p. 316 (1932).}. The simplest is given by the surface of the 6th order \(F_6\) possessing two oscnodes \(A\) and \(B\): \(p=4\), \(p^{(1)}=9\). The \(\infty^3\) quadrics adjoint to \(F_6\) touch the oscnodal planes at \(A\) and \(B\), and therefore meet two by two in conics tangent to these same planes, which lie in planes through \(A\) and \(B\): the said conics form, in space, a congruence of the 1st order and cut out on the sextic \(F_6\) an involution \(I_4\), with which the canonical system turns out to be composed. The corresponding \textit{canonical surface} is a \textit{quadruple quadric}.

One is also led to a quadric, and more precisely to a 4-fold canonical quadric cone, with the same characters \(p=4\) and \(p^{(1)}=11\), by the \textit{double plane with branch curve} \(C_{18}\), of order 18, with 8 6-fold points, since here the canonical curves have as double images the sextics of genus two passing doubly through the 8 points, which belong, as is well known, to an involution \(I_2\) (Bertini's 3rd type).

We leave it to the student to compare the two examples and to recognize whether they correspond to distinct types of surfaces and whether or not they fall within one and the same wider family of surfaces.

The second example of \textsc{B. Segre} is the \textit{sextic \(F_6\) endowed with an oscnode \(A\) and with a triple point \(B\) near which there lies an infinitesimal double line}: \(p=4\), \(p^{(1)}=7\). It is clear that here too the adjoint quadrics cut one another two by two in the conics of a congruence of the 1st order, which meet the sextic in triples of points of an involution \(I_3\), with which the canonical system is composed. The sextic \(F_6\) therefore leads to a triple canonical quadric, which has, as has been said,
\[
p=4\qquad \text{and}\qquad p^{(1)}=7.
\]

This case of \textsc{B. Segre} comes to realize an observation of \textsc{Castelnuovo}, who, as early as 1892, had noted that, whenever the canonical curves of a surface contain a series \(g_3^1\), the characteristic series of the system must be composed with it, and hence the canonical system will belong to a rational involution of the 3rd order \(I_3\).

% scan: s163L
% !PRINTED 302
\weberpage{302}
The problem of ``\textit{determining all the surfaces \((p>2)\) whose canonical curves} (irreducible, without base points) \textit{contain a \(g_3^1\) and which are therefore transformed into triple canonical surfaces with rational sections}'' has recently been solved by \textsc{G. Pompilj} in the Note\origfootnote{1}{s163L-1}{Rendic. Ist. Lombardo, 27 March 1941.} ``Sulle superficie algebriche le cui curve canoniche posseggono una \(g_3^1\)''. There are \textit{three types of canonical surfaces} which meet these conditions:

1) a triple plane with branch curves \(C_{12}\) of order 12 (already pointed out in the ``Lezioni'' of Enriques-Campedelli)\origfootnote{2}{s163L-2}{There it is noted that such a triple plane with \(p=3\) and \(p^{(1)}=4\) (which the count of moduli shows to be a particular case of the family) arises from the surface of the 5th order \(F_5\) endowed with a tacnode, by projection from this singular point. One can similarly start from a canonical surface \(F_6\) of \(S_3\) which acquires a triple point.

This consideration is easily extended; thus from the canonical surfaces \(F_7\), \(F_8\ldots\) of \(S_3\) which acquire a triple point, one obtains the construction of particular multiple canonical planes, quadruple or quintuple planes or in general \((p^{(1)}-1)\)-fold planes, of genus \(p=3\) and \(p^{(1)}=5,6\ldots\)

On the other hand, by making the said canonical surfaces acquire two triple points, one obtains surfaces with \(p=2\) of linear genus \(p^{(1)}=2,3,\ldots\), where the pencil of canonical curves possesses \(p^{(1)}-1\) base points.}, which has
\[
p=3\qquad \text{and}\qquad p^{(1)}=4;
\]

2) a triple quadric cone with branch curves \(C_{18}\) of order 18, for which
\[
p=4\qquad \text{and}\qquad p^{(1)}=7,
\]
a type within which falls (as a particular case?) the example of \textsc{B. Segre};

3) a cubic cone, rational normal in \(S_4\), triple, with branch curve \(C_{24}\) of order 24, for which
\[
p=5\qquad \text{and}\qquad p^{(1)}=10.
\]

These three types are represented analytically by the following equations:

1) triple plane
\[
z^3+3p_4z+q_6=0,
\]
where \(p_4(xy)\) and \(q_6(xy)\), equated to zero, give general curves of order 4 and 6 respectively: branch curve
\[
4p_4^3+q_6^2=0;
\]

2) triple plane
\[
z^3+3p_6z+q_9=0,
\]
% scan: s163R
% !PRINTED 303
\weberpage{303}
\(p_6\) and \(q_9\) curves of order 6 and 9, \(p_6\) passing through two infinitely near triple points \(O^3\) and \(O_1^3\), and \(q_9\) through \(O^5\) and \(O_1^4\); branch curve
\[
4p_6^3+q_9^2=0;
\]

3) triple plane
\[
z^3+3p_8+q_{12}=0; %
\]
\(p_8\) and \(q_{12}\) curves of order 8 and 12, passing through a point \(O\) and through two points \(O_1\) and \(O_2\) proximate to it, as is indicated by the symbols \(p_8(O^5O_1^2O_2^2)\) and \(q_{12}(O^8O_1^4O_2^4)\); branch curve
\[
4p_8^3+q_{12}^2=0.
\]

We shall not dwell on the analysis of \textsc{Pompilj}, which is based on the theory of triple planes established by himself\origfootnote{1}{s163R-1}{\textsc{G. Pompilj}, \textit{Sulla rappresentazione algebrica dei piani tripli}. Rendic. Seminario mat. Università di Roma, ser. 4, vol. III. \textit{Osservazioni sui piani tripli}. Rendic. Istituto Lombardo, 27 February 1941.} and which can also be carried out in another way\origfootnote{2}{s163R-2}{Cf. \textsc{A. Franchetta}, ``\textit{Sulle superficie le cui curve canoniche posseggono una \(g_3^1\)}''. Boll. dell’U. M. I., 1942.}. But we wish to indicate how one justifies the observation of \textsc{Castelnuovo} that ``if the canonical curves of a surface contain a \(g_3^1\), the characteristic series is necessarily composed with it''.

To prove the assertion it suffices to consider a curve of genus \(\pi>3\) which contains a \(g_3^1\) and a self-residual series of dimension \(r\geq 1\), within which we may take a \(g_n^1\) with \(n=\pi-1\): we shall reduce to absurdity the hypothesis that the \(g_n^1\) is not composed with the \(g_3^1\) (apart from fixed points). Indeed, if the \(g_n^1\) is not composed with the \(g_3^1\), one can transform the said curve into a \(C_m\) of order \(m=n+3\) on which the \(g_3^1\) and the \(g_n^1\) are cut out respectively by the lines through two points \(A\) and \(B\): \(A\) \(n\)-fold and \(B\) 3-fold. Now a pair of lines through \(B\) would have to form part of a curve of order \(m-3\) adjoint to \(C_m\), and hence there would have to exist a curve of order \(m-5\) passing \(n-1=m-4\) times through \(A\): which is absurd.

\edsection{14. Surfaces of linear genus \(p^{(1)}=2\): first case \(p=2\).}{14. Surfaces of linear genus p\textasciicircum{}(1)=2: first case p=2.}{sec.VIII.14}

When, for a surface \(F\), the canonical system is of dimension less than 3, that is, the genus is \(p<4\), or if this (pure) system is reducible, provided the linear genus is \(p^{(1)}>1\), or also if it belongs to an involution, in all these cases, in which one no longer has a simple canonical surface of order \(p^{(1)}-1\), one can usefully take as projective model of the class of \(F\) the \textit{bicanonical surface} of \(S_{P-1}\) (\(P=p+p^{(1)}\)), whose hyperplane sections are the bicanonical curves, or the \textit{tricanonical surface}, etc. Let us show the
% scan: s164L
% !PRINTED 304
\weberpage{304}
advantage that can be drawn from this consideration, by determining in this way the regular surfaces of lowest linear genus: \(p^{(1)}=2,3\).

Suppose that the (absolute) linear genus of \(F\) is \(p^{(1)}=2\), and that its surface genus is \(p_a=p_g=p\). In accordance with the relation which gives the maximum of \(p\) with respect to \(p^{(1)}\), we shall have:
\[
p\leq 2;
\]
indeed, the canonical curves cannot form a net of irreducible curves, which would be of degree \(p^{(2)}=p^{(1)}-1=1\), nor can they be reducible, containing as parts the curves of a pencil, which would have to be of genus \(<2\) and hence elliptic, thus entailing \(p^{(1)}=1\). We shall therefore have to examine three cases:
\[
p=2,\qquad p=1,\qquad p=0.
\]
And corresponding to the three indicated values of \(p\), the bigenus will be
\[
P=p+p^{(1)}=4,3,2,
\]
and the trigenus
\[
P_3=p+3p^{(1)}-2=6,5,4.
\]

\textit{First case}: \(p^{(1)}=2\), \(p=2\), \(P=4\). Here one has a (linear) pencil of canonical curves \(K\), of genus 2, and an \(\infty^3\) linear system of bicanonical curves, of genus \(3p^{(1)}-2=4\), cutting out on the \(K\) the \(g_2^1\) and hence belonging to a rational involution \(I_2\). One will therefore have, as projective model of the class of the given surfaces \(F\), a bicanonical surface of order \(4(p^{(1)}-1)=4\); but this will reduce to a double quadric with branch curve \(C_{10}\) of order 10, cutting the conics that are its plane sections in 10 points; since, then, the generators must possess an even number of branch points, such a quadric will be a cone, on which the vertex \(V\) is added to \(C_{10}\) as a branch point: which, moreover, corresponds to the base point of the pencil \(|K|\) on \(F\), a point which (being self-residual) must fall, for every \(K\), at a double point of the \(g_2^1\).

Summing up:

\textit{The surfaces with genera \(p=2\), \(p^{(1)}=2\), admit as model the bicanonical surface constituted by the double quadric cone with branch curve of order 10, complete intersection with a surface of the 5th order which does not pass through the vertex, to which this same vertex is added as a branch point.} The said cone can also be represented on the \textit{double plane with branch curve of order 10 endowed with two infinitely near 5-fold points}.

\textit{Remark.} -- To the preceding analysis one may object that it supposes the canonical curves \(K\) irreducible. A priori the \(K\) could be composed of the curves \(k\) varying in a (linear) pencil and of
% scan: s164R
% !PRINTED 305
\weberpage{305}
other fixed curves \(\theta\): \(K=nk+\Sigma\theta\). But, in any case, the \(k\) could not be of genus \(\pi=0\) or \(\pi=1\), since it would follow that the surface is rational (\(p=0\)) or that it has \(p^{(1)}=1\). The \(k\), whose genus cannot exceed \(p^{(1)}\) (cf. Chap. IV, \S~18), will therefore be of genus 2.

Now the \(\infty^3\) system \(|k'|\) adjoint to \(|k|\) --- which is contained in \(|2K|\) --- will have degree 4 and will belong to a rational involution \(I_2\), leading to a model of the surface constituted by a double quadric cone with branch curve \(C_{10}\) of order 10: this is the general model of the bicanonical surfaces (\(p=1\), \(p^{(1)}=2\)) defined above, and it turns out a posteriori that there are no fixed parts of the canonical pencil \(|K|\), nor of the bicanonical system. %

\edsection{15. Surfaces with \(p^{(1)}=2\) and \(p=1\) (\(P=3\), \(P_3=5\)).}{15. Surfaces with p\textasciicircum{}(1)=2 and p=1 (P=3, P\_3=5).}{sec.VIII.15}

Here there is a canonical curve \(K\) of genus \(p^{(1)}=2\) and \(\infty^2\) bicanonical curves of genus \(3p^{(1)}-2=4\) meeting two by two in \(4(p^{(1)}-1)=4\) points. %
If \(K\) is supposed irreducible, the bicanonical curves of \(|2K|\) will also be irreducible, and they cannot be hyperelliptic, since otherwise the characteristic series of the net, formed by the pairs of a \(g_2^1\), would not be complete, as the regularity of the given surface \(F\) requires. This being so, we shall have on \(F\) an \(\infty^4\) tricanonical linear system \(|3K|\), also formed of irreducible curves, which --- cutting out the canonical series on a bicanonical curve, and since \(p>0\) --- will necessarily be simple. It leads to a tricanonical surface \(F_9\) of order \(9(p^{(1)}-1)=9\) in the space \(S_4\), with hyperplane sections of genus \(6p^{(1)}-5=7\).

The direct construction of the \(F_9\) of \(S_4\) seems to present some difficulty. On the other hand it is easy to characterize a projection \(F_6\) of it into ordinary space. For this purpose observe that \(F_9\) possesses a triple line, corresponding to the canonical curve \(K\), on which the tricanonical curves cut out a \(g_3^1\). From a point of this triple line \(F_9\) is projected into a surface of the 6th order \(F_6\) of \(S_3\), which a priori must contain 3 exceptional lines. The \(F_6\) will possess a singular double point, \(O\), trace of the triple line of \(F_9\), and such that the neighborhood of \(O\) corresponds to the curve \(K\) of genus 2. Now the plane sections of \(F_6\) will be of genus 7, that is, sextics endowed with 3 double points, the plane sections through \(O\) of genus 4; the singularity of the double point \(O\), corresponding to a fundamental curve of genus 2, will lower the genus of the plane sections by 3, and hence these will have no double points away from \(O\). One deduces that \(F_6\) possesses 3 double lines through \(O\) (which cannot coincide in a triple line, which would entail a pencil of elliptic plane sections and hence \(p^{(1)}=1\)); and these lines \(a\), \(b\) and \(c\)
\printersignature{\textsc{Enriques F.} - \textit{Superficie algebriche.}\quad 20}
% scan: s165L
% !PRINTED 306
\weberpage{306}
lie in the plane\origfootnote{1}{s165L-1}{Otherwise the point \(O\) would turn out to be triple.} tangent to \(F_6\) at \(O\) (a plane which, counted twice, constitutes its osculating cone): the singularity on every plane section of \(F_6\) through \(O\) consists of two branches having a contact of the 5th order (six-point contact).

An \(F_6\) endowed with the indicated singularities possesses an adjoint quadric which splits into the plane of the three double lines counted twice; in order that there be three exceptional lines, the said double lines will have to be cuspidal.

The \(F_6\) that we have constructed has geometric genus \(p_g=1\), possessing an adjoint quadric constituted by the plane \(a\,b\,c\) counted twice; it also has arithmetic genus \(p_a=1\), that is, it is regular; indeed, for this surface, whose plane sections are of genus 7, there exist \(\infty^7\) (\(7=p_a+\pi-1\)) adjoint surfaces of the 3rd order, namely those defined by the conditions of passing through the lines \(a\), \(b\), \(c\), and of having a contact of the 4th order with the sheets of \(F_6\) at \(O\); and finally it has absolute linear genus \(p^{(1)}=2\), because its bigenus equals \(P=p+p^{(1)}=3\), there being \(\infty^2\) bicanonical curves cut out by the planes through \(O\), to which the plane \(a\,b\,c\) counted three times is to be added: the value of \(p^{(1)}\) also follows from the order of \(F_6\), with tricanonical plane sections, taking account of the 3 exceptional lines:
\[
6=9(p^{(1)}-1)-3=9-3\ \text{\origfootmark{2}{s165L-2}}.
\]\origfoottext{2}{s165L-2}{The classification of surfaces with \(p^{(1)}=2\), \(p=1,2\), is given in the Note of \textsc{F. Enriques}, \textit{Le superficie algebriche di genere lineare \(p^{(1)}=2\)}. Rendic. Acc. Lincei, February 1897.}
One can construct an \(F_6\) with the singularities indicated above starting from a cubic surface \(F_3\) which contains three lines \(a\), \(b\), \(c\) lying in a plane and passing through a point \(O\). There exist \(\infty^4\) cubic surfaces \(F'_3\) which touch \(F_3\) at \(O\), giving rise on every plane section to a six-point contact: these \(F'_3\) consequently touch \(F_3\) along the three lines \(a\), \(b\), \(c\), as appears from the plane representation, and further cut \(F_3\) in a plane cubic, being characterized by these properties. Now a composite surface \(F_3+F'_3\) furnishes a particular sextic which possesses \(a\), \(b\), \(c\) as tacnodal lines, and presents at \(O\) the singularity indicated above (on every plane section through \(O\), two branches with six-point contact). Let us combine linearly \(F_3+F'_2\) with the degenerate sextic constituted by the plane \(a\,b\,c\) counted three times and by a cubic cone with vertex \(O\); in this way one will obtain an \(F_6\) which will have the required singularity at \(O\), and for which the lines \(a\), \(b\), \(c\) turn out to be cuspidal. %

One verifies at once that the \(F_6\) thus constructed possesses an \(\infty^7\) system of adjoint cubic surfaces \(\Phi_3\), defined by the conditions of passing
% scan: s165R
% !PRINTED 307
\weberpage{307}
through the three lines \(a\), \(b\), \(c\), and of touching \(F_3\) simply at \(O\), in such a way that the section of \(\Phi_3\) by a plane through \(O\) comes to contain 5 of the 6 successive double points of the section of \(F_6\). And from this it follows that the surface \(F_6\), whose plane sections are of genus 7, is regular, since the said \(\Phi_3\) cut out, on a plane section of \(F_6\), the complete system of its adjoint curves.

Meanwhile let us summarize the result obtained by stating that: \textit{The regular surfaces with genera \(p=1\) and \(p^{(1)}=2\) can be transformed into sextics \(F_6\) endowed with three cuspidal lines lying in a plane and passing through a singular double point \(O\), such that on every plane section through it there are two branches with contact of the 5th order.}

\textit{Remark.} -- In the preceding analysis it has been supposed that the canonical curve \(K\) is irreducible, and hence that the net of bicanonical curves \(|2K|\) has no fixed parts. Let us indicate briefly how we can free ourselves from this restrictive hypothesis, and to this end let us examine whether, for \(p=1\), \(p^{(1)}=2\), there can be a bicanonical system
\[
|2K|=|k_2|+\Sigma\theta
\]
having some fixed component \(\theta\), whose variable parts are irreducible curves \(k_2\) (cf. Chap. IV, \S~19). We denote by \(\pi\) and \(n\) the genus and degree of a \(k_2\), and by \(\varrho\) and \(\nu\) the genus and degree of a \(\theta\). Then the additive character which expresses the virtual number of intersections of the bicanonical curves with the canonical ones is given by %
\[
2p^{(1)}-2=2=2\pi-2-n+\Sigma(2\varrho-2-\nu).
\]
And, supposing the exceptional curves of the surface to have been eliminated (and hence, for \(\varrho=0\), \(\nu\leq -2\)), we shall have for all the components \(\theta\):
\[
2\varrho-2-\nu\geq 0,
\]
and hence
\[
2\pi-2-n\leq 2.
\]
Now if some \(\theta\) exists, since \(K\), and hence also the curves of \(|2K|\), are connected with order of connection \(\geq 2\), the number of intersections of a \(k_2\) with \(\Sigma\theta\) (residual with respect to the bicanonical system) will be
\[
k_2\cdot\Sigma\theta=4\pi-4-3n\geq 2.
\]

From this it follows that
\[
\begin{aligned}
4\pi-4-2n&=(4\pi-4-3n)+n\leq 4,\\
n&\leq 2\qquad \text{and hence}\qquad n=2(n>1),\\ %
2\pi-2&\leq n+2=4;
\end{aligned}
\]
that is, the net \(|k_2|\) consists of curves of genus \(\pi\leq 3\) and of degree \(n=2\). But this conclusion is absurd because (in the most
% scan: s166L
% !PRINTED 308
\weberpage{308}
favorable hypothesis: \(\pi = 3\)) it entails as a consequence the possibility of representing the given surface \(F\) on a double plane with branch curve \(C_8\) of order 8, which --- since \(p = 1\) --- must possess two elementary singularities, quadruple points or points \([3, 3]\), whence it follows that
\[
p^{(1)} = 1.
\]

Therefore the bicanonical system of an \(F\) of genus \(p = 1\) with \(p^{(1)} = 2\) possesses no fixed parts, and the classification of these surfaces turns out to be complete.

\edsection{16. Surfaces with \(p^{(1)} = 2\) and \(p = 0\) (\(P = 2,\ P_3 = 4\)).}{16. Surfaces with p\textasciicircum{}(1) = 2 and p = 0 (P = 2, P\_3 = 4).}{sec.VIII.16}

One is led to a \textit{tricanonical surface} of ordinary space \(S_3\): a surface \(F_9\) of order 9, with plane sections of genus 7 and hence \textit{endowed with a double curve \(C_{21}\) of order \(21\)}, which does not belong to any surface \(\Phi_5\), but is nevertheless double for \(\infty^4\) biadjoint surfaces \(\Phi_{10}\) of order 10, which cut out on it the twisted sextics of genus 4 of a pencil, which are bicanonical curves.

It is not easy to construct \(F_9\) directly and to define its singularities precisely, but one can at least prove its \textit{existence}, starting from particular surfaces of the family of which it provides the projective model. For this purpose one can start from the double plane of \textsc{Campedelli}, endowed with branch curves \(C_{10}\) of order 10 with a 4-fold point \(A\) and 5 points \([3, 3]\), for which
\[
p = 0 \qquad \text{and} \qquad p^{(1)} = 2,
\]
there being a pencil of bicanonical curves of genus \(3p^{(1)} - 2 = 4\), represented by the (double) quartics which pass doubly through \(A\) and simply through the 5 pairs of infinitely near triple points of \(C_{10}\) forming the singularities \([3, 3]\). The double plane thus defined represents a surface \(F\) on which the bicanonical curves are not hyperelliptic, and hence the tricanonical curves form a simple system (see the Remark further on). By means of this system \(F\) can be transformed into the tricanonical \(F_9\) indicated above, or rather into a particular case of it.

Indeed, our double planes depend on 4 moduli, there being 4 projective invariants of the branch curve \(C_{10}\):
\[
4 = 65 - 8 - 5 \cdot 9 - 8;
\]
whereas the number of moduli belonging to the family of the surface \(F_8\), with \(p = 0\) and \(p^{(1)} = 2\), is certainly %
\[
M \geq 9p - 2p^{(1)} + 12 = 8.
\]

% scan: s166R
% !PRINTED 309
\weberpage{309}
\textit{Remark.} -- We shall add that the tricanonical surface corresponding to a regular surface \(F\) with \(p = 0\) and \(p^{(1)} = 2\) is certainly simple, at least if the bicanonical curves of \(F\) are irreducible.

Let us indicate briefly how the assertion can be justified.

First of all it is convenient to establish that from the irreducibility of the bicanonical curves \(K_2\) it follows that the tricanonical curves \(K_3\) of the adjoint system \(|K'_2| = |K_3|\) are irreducible and without base points. Since the \(K'_2\) cut out on a \(K_2\) the complete canonical series, if they are reducible they must contain a fixed (irreducible) part \(\theta\) fundamental for the pencil \(|K_2|\), and then the tricanonical system stripped of \(\theta\) will constitute the adjoint to the curve \(K_2 - \theta\), so that \((K_2 - |\theta|)'\) would have the same dimension as \(|K'_2|\), that is, dimension 3: it is therefore a matter of recognizing that this is impossible because \(K_2 - \theta\) is a connected curve of genus \(< 4\) (Cf. Chap. IV, § 17). For this purpose one will realize the \(K_2\) as canonical sextics of ordinary space, projected onto the plane into sextics with two triple points, and one will study their possible degenerations: precisely in this way one sees that the detachment of an irreducible \(\theta\) from \(|K_2|\) necessarily leads to a connected curve of genus \(\leq 2\). %

A similar analysis serves to exclude that the irreducible \(|K_3|\) has some base point on the given surface \(F\): in the first place there cannot be a base point which is simple for the curve \(K_3\) passing through it, etc. This being so, since \(|K_3|\) is an irreducible linear system without base points, of degree 9, if it belongs to an involution \(I_n\) of order \(n\), \(n\) must be a divisor of 9, and hence \(n = 3\). Under this hypothesis one is led to examine a rational cubic surface \(F_3\) with elliptic sections, triple, with branch curves \(C_{12}\) of order 12, which would have to constitute the tricanonical surface of the \(F\) of genus \(p_a = p_g = 0\) and of linear genus \(p^{(1)} = 2\). But it is easy to convince oneself that such a triple cubic \(F_3\), with genera \(p = 0\) and \(p^{(1)} = 2\), \textit{does not exist}. Let us say first of all that our \(F_3\), which counted 3 times would constitute the tricanonical surface of a surface \(F\) with \(p = 0\) and \(p^{(1)} = 2\), must possess at least one (singular) double point. Indeed, if \(F_3\) is supposed free of double points, there will be on it two homaloidal nets of twisted cubics, complementary (with respect to the system of sextics cut out by the quadrics): \(|C_3|\) and \(|L_3|\); and one sees at once that these must represent two nets of curves, \(|C|\) and \(|L|\) respectively, of degree 3 and genus 4. If this were not so, since the \(C_3 + L_3\) cut the branch curve \(C_{12}\) of the triple \(F_3\) in 24 points, one of the two nets, for example \(|C|\), would be of genus \(\pi < 4\); but if \(\pi < 3\), all the plurigenera of \(F\) turn out to be zero; if instead \(\pi = 3\), the bicanonical curves of \(F\), which --- by hypothesis --- form a pencil, would cut out on a \(C\) the pairs of a \(g_2^1\), which is incompatible with the fact that the \(C\) contain a \(g_3^1\) without fixed points.

% scan: s167L
% !PRINTED 310
\weberpage{310}
This being so, since the \(C\) and the \(L\) are curves of genus 4, one recognizes that on them the characteristic \(g_3^1\) is self-residual (with respect to the canonical series). If this were not so, e.g. for the \(C\), one proves that the genus of the virtual canonical curves \(K\), which have 3 points in common with the \(C\), must equal \(p^{(1)} = 1\), instead of \(p^{(1)} = 2\). Indeed, take as model of an \(F\) a simple \(F_n\), of a certain sufficiently high order \(n\), on which the \(C\) are cut out by the planes through an \((n - 3)\)-fold point \(O\): the curves of the adjoint system \(|C'| = |C + K|\) belonging to \(F\) will be projected from \(O\) into plane curves of the \(6^{\mathrm{th}}\) order and, in particular, the \(C'\) through a group of the characteristic \(g_3^1\) of \(C\)(which are \(\infty^1\)) into sextics with a triple point: this means that the genus of the \(C\) equals %
\[
7 = 4 + p^{(1)} + 3 - 1,
\]
and hence
\[
p^{(1)} = 1.
\]

It thus remains established that each of the two nets of curves \(|C|\) and \(|L|\) on \(F\) possesses a self-residual characteristic series. Consequently the bicanonical curves \(K_2\) (of \(|2K|\)) cut out on a \(C\) groups of the double series of the said \(g_3^1\), and hence are partly contained in the double system \(|2C|\), and so similarly in \(|2L|\).

Since the complete \(|2C|\) has dimension 6 and the minimal \(|2C|\) dimension 5, there is \textit{one} curve of the pencil \(|K_2|\) belonging to the minimal \(|2C|\), and this is characterized by belonging to the involution \(I_3\) whose triples correspond to the points of \(F_3\), and hence it must also belong to the minimal system \(|2L|\).

Now represent the cubic \(F_3\) point by point on a plane \(\alpha\): the plane sections of \(F_3\) will have as images the cubics through 6 points \(A_1A_2\ldotp\ldotp\ldotp\ldotp A_6\), not lying on a conic; the curves of the minimal system \(|2C|\) may be supposed represented by the conics of \(\alpha\); and those of the minimal system \(|2L|\) by \(C_{10}\) of order 10 passing 4 times through the 6 points \(A_i\). Hence there would have to be in the plane \(\alpha\) a conic, passing through some point \(A_i\), which at the same time forms part of a \(C_{10}\) and meets these same \(C_{10}\) in 2 variable points: but such a conic does not exist. Thus the hypothesis of an \(F_3\) (triple image of a tricanonical \(F\) with \(p = 0\) and \(p^{(1)} = 2\)) possessing no double points is rejected.

Suppose on the other hand that there exists an \(F_3\) endowed with a double point \(O\) (however singular), whose plane sections are images of the tricanonical curves \(K_3\) of \(F\).

The plane sections \(C_3\) of \(F_3\) through \(O\) (as was said before for any twisted cubics belonging to this surface) will be images of curves \(C\) of genus 4, and the bicanonical curves \(K_2\) of \(F\) will have as images on \(F_3\) simple sextics, meeting the said \(C_3\) in 6 variable points, and hence not passing through \(O\). Con%
% scan: s167R
% !PRINTED 311
\unskip\weberpage{311}%
sequently the quadricanonical system adjoint to \(|K_3|\), that is
\[
|2K_2| = |K'_3| = |4K|
\]
will consist in general of curves not passing through \(O\); and to \(O\) there will correspond on \(F\) a curve \(\theta\) fundamental for the said \(|K_3|\). But from this one deduces that by removing \(\theta\) from this system one obtains the adjoint to \(|C|\):
\[
\begin{aligned}
|C'| &= |K'_3 - \theta|\\
(|K_3 - \theta| &= |C|),
\end{aligned}
\]
and hence the genus of the \(C\) must turn out equal to that of the \(K_3\) diminished by one unit:
\[
7 - 1 = 6;
\]
an absurd consequence, because the genus of the \(C\) equals, as has been said,
\[
\pi = 4!
\]

Thus the discussion is exhausted: \textit{triple bicanonical \(F_3\) of genus \(p = 0\) and \(p^{(1)} = 2\) do not exist.} \hfill Q.E.D. %

\edsection{17. Surfaces of linear genus \(p^{(1)} = 3\): first case \(p = 3\)\origfootnote{1}{s167R-1}{The classification of surfaces of linear genus \(p^{(1)} = 3\), for \(p = 1, 2, 3\), which is here deepened in various respects and extended to the case \(p = 0\), is given by \textsc{F. Enriques} in the Note \textit{Sulle superficie algebriche di genere lineare \(p^{(1)} = 3\)}. Rendic. Lincei, May 1897. Only the simple type for \(p^{(1)} = 3\) and \(p = 2\), of § 18, is missing.}.}{17. Surfaces of linear genus p\textasciicircum{}(1) = 3: first case p = 3.}{sec.VIII.17}

Let us examine the (regular) surfaces of linear genus \(p^{(1)} = 3\). Their surface genus
\[
p \leq \frac{p^{(1)} + 3}{2}
\]
cannot exceed the value 3 (cf. § 11). We shall therefore have a priori 4 cases:
\[
p = 3,\, 2,\, 1,\, 0.
\]

\textit{First case:} \(p^{(1)} = 3,\ p = 3\). The canonical system \(|K|\), supposed irreducible, is a net of degree \(p^{(2)} = p^{(1)} - 1 = 2\), which leads to representing the surface on the \textit{double plane with branch curve \(C_8\) of order \(8\)}.

The hypothesis that \(|K|\) is reducible into a fixed component \(\Sigma\theta\) and a variable part \(|k|\) does not lead to any new case. Let \(\pi\) and \(n\) denote the characters of the \(k\), and \(\varrho\) and \(\nu\) those of a \(\theta\). Computing the additive character which indicates the number of intersections of the canonical curves with these same curves, we shall have:
\[
2\pi - 2 - n + \Sigma(2\varrho - 2 - \nu) = 2.
\]

% scan: s168L
% !PRINTED 312
\weberpage{312}
Now the number of intersections of an (irreducible) \(k\) with the rest of the canonical curves, \(\Sigma\theta\), is
\[
2\pi - 2 - 2n \geq 0,
\]
and the degree of the system \(|k|\) is \(n \geq 2\); hence
\[
2\pi - 2 - n \geq 2.
\]

On the other hand, if one admits that every exceptional curve has been eliminated from the surface, one will have for each \(\theta\):
\[
2\varrho - 2 - \nu \geq 0;
\]
consequently
\[
\begin{aligned}
2\pi - 2 - n = 2,\qquad &2\varrho - 2 - \nu = 0,\\
n = 2 \qquad \text{and}\qquad &\pi = 3.
\end{aligned}
\]

Therefore the \(k\) form a net which leads to transforming the given surface into a double plane with branch curve \(C_8\), and the \(\theta\) turn out a posteriori not to exist.

\edsection{18. Surfaces with \(p^{(1)} = 3\) and \(p = 2\) (\(P = 5\)).}{18. Surfaces with p\textasciicircum{}(1) = 3 and p = 2 (P = 5).}{sec.VIII.18}

Supposing that the canonical pencil \(|K|\) is irreducible, apart from the fixed exceptional components which correspond to its base points, the \(\infty^4\) bicanonical curves, cutting out on every \(K\) the complete canonical series \(g_4^2\), will also form an irreducible system, of genus \(3p^{(1)} - 2 = 7\) and degree \(4(p^{(1)} - 1) = 8\), free of base points. And this system must be simple if the canonical curves of the pencil \(|K|\) are not hyperelliptic, as it is convenient to suppose in the first place. Hence one is led to a simple bicanonical \(F_8\) of \(S_4\), with hyperplane sections of genus \(3p^{(1)} - 2 = 7\), on which the \(K\) are plane quartics having two points in common and hence lying in planes through a line \(a\), which does not lie on \(F_8\). The hyperplanes through \(a\) cut \(F_8\) in pairs of quartics \(K\); it follows that the planes of the \(K\) form a quadric cone of the \(2^{\mathrm{nd}}\) species \(Q\), which has \(a\) as its axis. Now project \(F_8\) onto the space \(S_3\) from a generic point \(P\) of \(Q\): one will have in \(S_3\) a surface \(F'_8\) with two infinitely near quadruple lines \(o\) and \(o'\), issuing from a point \(O\). This means that \(F_8\) is the complete intersection of the quadric cone \(Q\) and a hypersurface of the \(4^{\mathrm{th}}\) order. But (as we have already seen) the intersection of a quadric and a quartic in \(S_4\) is in general a canonical surface with genera \(p = 5\) and \(p^{(1)} = 9\): to have a bicanonical \(F_8\) (\(p = 2,\ p^{(1)} = 3\)) one must impose on the said intersection suitable singularities which lower its genus, and precisely --- since the genus 9 of the hyperplane sections must be lowered %
% scan: s168R
% !PRINTED 313
\weberpage{313}
by 2 --- it will be convenient to impose on \(F_8\) a double curve \(C_2\) of the \(2^{\mathrm{nd}}\) order. Now, referring to \(F'_8\), the projection of \(F_8\) into \(S_3\) from a point of the quadric cone which contains it, one sees that \(C'_2\), the projection of \(C_2\), cannot consist of two lines, incident or not, which do not lie in a plane through \(o\), and hence it must be a conic lying in a plane through \(o\). In conclusion the \textit{type of the simple bicanonical surface, for \(p = 2\) and \(p^{(1)} = 3\), is an \(F_8\) of order \(8\) in the space \(S_4\) belonging to a quadric cone \(Q\) and having as projection from a point of \(Q\) into ordinary space an \(F_8\) endowed with two infinitely near quadruple lines \(o\) \(o'\), incident, and with a double conic, lying in another plane through \(o\).}

One constructs \(F'_8\) by combining linearly: a quartic with two double lines \(o\) and \(o'\) passing through the conic \(C'_2\), counted twice, and a degenerate \(F'_8\) composed of the plane \(\omega = o\,o'\) counted 4 times and of the plane \(\gamma\) of the conic \(C'\) counted twice, besides a generic quadric. And one verifies at once that the surface \(F'_8\) endowed with the singularities that result from this is precisely our bicanonical surface, which is not normal. Indeed the surfaces adjoint to it will reduce: 1) owing to the presence of the 2 4-fold lines \(o\) and \(o'\): to cones tangent along \(o\) (to which the plane \(\omega = oo'\) counted twice is added); and 2) owing to the double conic \(C'_2\) (whose plane \(\gamma\) splits off from the said cones) to planes through \(o\) cutting out on \(F'_8\) the quartics \(K\); whence follows \(p = 2,\ p^{(1)} = 3\). Moreover there is a biadjoint \(\Phi_7\), \(\Phi_7 = \omega^5\gamma^2\), which cuts \(F'_8\) only in the singular curves, so that the plane sections of \(F'_8\) turn out to be bicanonical curves.

\textit{Every surface with \(p = 2\) and \(p^{(1)} = 3\), with canonical curves} (irreducible) \textit{that are not hyperelliptic, can be transformed into the simple bicanonical \(F_8\) of \(S_4\).}

A particular example is given by the \textit{double plane with branch curve \(C_{12}\) of order \(12\), endowed with \(8\) \(4\)-fold points \(A_1A_2\ldotp\ldotp\ldotp\ldotp A_8\)}. Here the canonical curves are represented by the cubics \(C_3\) through \(A_1A_2\ldotp\ldotp\ldotp\ldotp A_8\), which --- taken as double, with the branch points at the intersections with \(C_{12}\) --- constitute, in general, non-hyperelliptic curves of genus 3. In effect there are \(\infty^4\) bicanonical curves having as images on the plane curves of order 12 passing 4 times through \(A_1A_2\ldotp\ldotp\ldotp\ldotp A_8\) and 8-tangent to the branch curve \(C_{12}\); and these represent a simple complete \(\infty^4\) linear system. Within the said complete system there are \(\infty^3\) bicanonical curves which have as images the sextics \(C_6\) passing doubly through the 4-fold points of \(C_{12}\), and which form a linear system, not complete, belonging to an involution of the \(4^{\mathrm{th}}\) order.

In the preceding example one sees a particular bicanonical \(F_8\) which --- from a certain external point --- is projected into a quadruple cone of the second order. Another noteworthy example is given by an \(F_8\) which can be projected from an external point \(O\) into a \textit{quartic \(F_4\)
% scan: s169L
% !PRINTED 314
\weberpage{314}
of genus \(1\), double, with plane branch curve of the \(4^{\mathrm{th}}\) order, \(F_4\) being moreover endowed with a plane tangent along a conic, on which there lie \(6\) conical branch points.}

Indeed, if \(F_8\) is transformed into itself by an involution \(I_2\), induced by a harmonic homology with center \(O\), it will have as projection from \(O\) an \(F_4\), on which the canonical curves \(K\) have as images the plane sections of genus 3 of a pencil, and which therefore must be free of double curve: since, moreover, \(I_2\) pairs off the canonical curves of the pencil \(|K|\) on \(F_8\), in this pencil there must be two united \(K\), one of which is projected into the branch curve \(C_4\) of \(F_4^2\), the other (hyperelliptic) into the conic \(C\), section by a tangent plane \(\alpha\), and on this there lie (besides the two intersections with \(C_4\)) 6 branch points, which are the conical points of \(F_4\) on \(\alpha\). Conversely, take an \(F_4\) with a plane \(\alpha\) tangent along a conic \(C\), and construct \(F = F_4^2\) by extracting \(\sqrt{\alpha\beta}\) on \(F_4\), so as to possess a branch quartic \(C_4\), \(\beta = 0\), to which the apparent critical conic \(C\) is added. It is easy to convince oneself that the \(\infty^1\) sections by the planes passing through the line common to the planes \(\alpha\) and \(\beta\) constitute double quartics, images of curves that split like the section by the plane \(\beta\); and hence --- taken simply --- they represent \(\infty^1\) canonical curves of \(F\), so that for the said \(F\) one finds \(p_g = 2\). One can also verify that \(p_a = 2\), either by showing a priori the regularity of the surface, or by computing \(p_a\) by means of the formulae established for the correspondence \([1, 2]\) between \(F_4\) and \(F\): where one must take into account that the 6 double points of \(F_4\) on the conic \(C\) constitute 6 infinitesimal branch curves (to be added to the quartic section by the plane \(\beta\)), to which correspond exceptional curves or simple points of the surface \(F\).

\textit{But a second, entirely distinct type}, which does not fall within that of the simple bicanonical \(F_8\) (\(p = 2,\ p^{(1)} = 3\)), is given by the \textit{bicanonical surface constituted by a Segre surface \(F_4\)}\origfootnote{1}{s169L-1}{Cf. \textsc{F. Enriques-F. Conforto}: ``\textit{Le superficie razionali}'', p. 144.} (with two double points) \textit{taken as double}, or by the \textit{double plane with branch curve \(C_{10}\) of order \(10\), endowed with \(2\) \(4\)-fold points and with \(2\) points \([3, 3]\)}\origfootnote{2}{s169L-2}{Or, if one prefers, with a \(C_{10}\) constituted by a line \(r\) and by a \(C_9\) with three 3-fold points on \(r\) and, outside it, two infinitely near 3-fold points.}.

Here the canonical curves are represented by \(\infty^1\) conics of a pencil (each to be taken as double with \(6 + 2 = 8\) branch points) and the bicanonical ones by \(\infty^4\) quartics with two double base points and two contacts:
\[
p = 2,\qquad p^{(1)} = 3.
\]

One easily proves that \textit{every surface \(F\) with \(p = 2\) and \(p^{(1)} = 3\) which leads in \(S_4\) to a double bicanonical surface has hyper%
% scan: s169R
% !PRINTED 315
\unskip\weberpage{315}%
elliptic canonical curves and can be represented on the \(F_4^2\) of \(S_4\) or on the aforesaid double plane.}

First of all, if the bicanonical system \(|2K|\) of \(F\) is not simple, this means that its curves passing through a point \(A\) must consequently pass through (at least) another point \(A'\); but this cannot be outside the \(K\) which passes through \(A\), and if it lies on the \(K\) (supposed irreducible) this means that the canonical series \(g_4^2\) marked out by \(|2K|\) on \(K\) is composite and \(K\) is hyperelliptic.

Now if the \(K\) are hyperelliptic, the bicanonical surface into which \(F\) is transformed will reduce to a double \(F_4\), and in order that \(F_4\) be normal in \(S_4\) it must have elliptic hyperplane sections, that is, be a Segre surface, complete intersection of two quadrics of \(S_4\). On \(F_4^2\), whose double hyperplane sections are of genus \(3p^{(1)} - 2 = 7\), there will be a branch curve \(C_{12}\) of order 12. Let us add that our \(F_4\) of \(S_4\) is not the most general intersection of two quadrics, but is endowed with two double points, which fall at the base points of the pencil of canonical conics (images of two double points of the \(g_2^1\) on the \(K\)). The point-by-point representation of \(F_4\) on the plane leads to a representative system of the hyperplane sections constituted by curves of the \(4^{\mathrm{th}}\) order with two double points and two simple base points of contact: these are the images of the bicanonical curves of a double plane, with \(p = 2\) and \(p^{(1)} = 3\), whose branch curve \(C_{10}\) possesses two 4-fold points and two points \([3, 3]\).

We have asserted that the double Segre \(F_4\) constitute a \textit{second type} of surfaces with \(p = 2\) and \(p^{(1)} = 3\), not falling as a particular case within the family which has as its type the simple bicanonical \(F_8\). To prove this one can resort to a count of moduli: both the simple bicanonical \(F_8\) and the \(F_4^2\) depend on \(24 = 9p - 2p^{(1)} + 12\) moduli, and therefore they form two distinct families. But this also appears in a direct way, because if one tries to reduce by continuity \(F'_8\), the projection of \(F_8\) into \(S_3\), to \(F'_4\), the projection of the Segre \(F_4\), counted twice, one arrives not at a quartic endowed with an irreducible double conic (as happens as a rule for our second type), but at a quartic with two infinitely near double lines: therefore there really are \textit{two distinct families of surfaces with \(p = 2\) and \(p^{(1)} = 3\)}, having a particular subfamily in common.

\textit{Remark.} -- Here too one can complete the analysis by examining the hypothesis that the \textit{canonical system \(|K|\) is reducible.}

First of all we may assert that the curves \(k\), variable components of the \(K\), must in any case be irreducible, because otherwise, being composed of the curves of a linear pencil (regularity of the surface), they would yield at least \(\infty^2\) \(K\) and \(p > 2\).

% scan: s170L
% !PRINTED 316
\weberpage{316}
Now one may suppose that the \(K\) possess, besides the \(k\), some fixed components \(\theta\), not exceptional
\[
|K| = \Sigma\theta + |k|.
\]

Denoting by \(\pi\) and \(n\) the characters of \(|k|\) and by \(\varrho\) and \(\nu\) those of a \(\theta\), we shall have
\[
2\pi - 2 - n + \Sigma(2\varrho - 2 - \nu) = p^{(1)} - 1 = 2,
\]
while \(\pi \geq 2\).

Moreover, referring to a surface without exceptional curves, one will have for every \(\theta\)
\[
2\varrho - 2 - \nu \geq 0.
\]

Since, further, a \(k\) has \(2\pi - 2 - 2n\) intersections with the residual curve \(\Sigma\theta\):
\[
2\pi - 2 - 2n \geq 0.
\]

From these inequalities there result three possibilities for the values of \(\pi\) and \(n\):
\weq{1)}{\pi = 2,\ n = 0;}
\weq{2)}{\pi = 2,\ n = 1;}
\weq{3)}{\pi = 3,\ n = 2.}

But hypothesis 1) is discarded at once. The system \(|k'|\) adjoint to \(|k|\) would be an \(\infty^3\) system of curves of degree
\[
0 + 2 \cdot 2 + 2 = 6,
\]
formed of hyperelliptic curves (bisecant to the \(k\)), and hence would lead to a double cubic ruled surface, which is not normal, so that \(|k'|\) (which has dimension \(p + 2 - 1 = 3\)) would have to be contained in a complete system of dimension \(> 3\): which is impossible.

Hypothesis 2) (\(\pi = 2,\ n = 1\)) is also easily discarded.

For the adjoint system \(|k'|\), which is a system of hyperelliptic curves, of genus 4 and degree 7, endowed with a simple base point, would again lead to a non-normal double cubic ruled surface, whence its dimension would be \(> 3\), which is impossible.

There remains hypothesis 3) (\(\pi = 3;\ n = 2\)). But under this hypothesis the system \(|2k|\) leads us precisely to the double surface \(F_4\) to which \(|K'| = |2K|\) led under the supposition that the \(K\) are irreducible. One therefore falls back on the same double plane and, a posteriori, one recognizes that the \(K\) are irreducible: \(|K| = |k|\).

\edsection{19. Surfaces with \(p^{(1)} = 3\) and \(p = 1\) (\(P = 4\)).}{19. Surfaces with p\textasciicircum{}(1) = 3 and p = 1 (P = 4).}{sec.VIII.19}

Supposing that the canonical curve \(K\) is irreducible, the bicanonical system \(|2K|\), which is \(\infty^3\), will also be irreducible and free of base points on
% scan: s170R
% !PRINTED 317
\weberpage{317}
\(K\). If it is simple it leads to transforming the given surface into an \(F_8\) of order 8, in the space \(S_3\), with plane sections of genus \(3p^{(1)} - 2 = 7\), and hence endowed with a double curve \(C_{14}\) of order 14. This \(C_{14}\) belongs to \textit{one} biadjoint surface \(\Phi_7\), of order \(2 \cdot 8 - 9 = 7\), which together with a plane constitutes a \(\Phi_8\) cutting out a bicanonical curve on \(F_8\); indeed \(C_{14}\) is the complete intersection of \(F_8\) and \(\Phi_7\left(\dfrac{7 \cdot 8}{4} = 14\right)\). On the other hand there must exist \textit{one} adjoint quartic \(\Phi_4\), which cuts out on \(F_8\), outside \(C_{14}\), the canonical \(K\) of order 4 (\(4 \cdot 8 - 2 \cdot 14 = 4\)). This being so, it is easy to characterize the curve \(C_{14}\): it is, in general, the double curve of the surface \(\Phi_7\) with elliptic sections, a surface possessing an adjoint quartic \(\Phi_4\) and hence rational.

Therefore \textit{the simple bicanonical surface \(F_8\) with genera \(p = 1\) and \(p^{(1)} = 3\) is defined, in general, by the possession of a double curve \(C_{14}\), of order \(14\), which is the double curve of a rational surface \(\Phi_7\) of the \(7^{\mathrm{th}}\) order with elliptic sections}\origfootnote{1}{s170R-1}{For surfaces with elliptic sections cf. e.g. \textsc{Enriques - Conforto}. \textit{Le superficie razionali}. Book II, chap. III.}. To construct \(F_8\) it will suffice to combine linearly the quartic \(\Phi_4\) adjoint to \(\Phi_7\), counted twice, and \(\Phi_7\) itself, taken together with a generic plane.

One can evaluate the number of moduli on which \(F_8\) depends: one finds that \(C_{14}\) (or \(\Phi_7\)) possesses 12 projective invariants and is double for \(\infty^4 F_8\), so that there are \(\infty^{16} F_8\) that are projectively distinct: the number of moduli of \(F_8\) equals:
\[
M = 16 = 10p - 2p^{(1)} + 12.
\]

Now the multiple bicanonical surfaces will be seen to fall, as particular cases, within the family of the simple bicanonical \(F_8\) with \(p = 1\) and \(p^{(1)} = 3\); these --- at least under the hypothesis of an irreducible canonical curve --- are only double quartics \(F_4^2\), and give rise to two types of these.

If on a surface \(F\) with \(p = 1\) and \(p^{(1)} = 3\) there is an irreducible canonical curve \(K\), and the bicanonical system \(|2K|\), which cuts out on \(K\) the complete \(g_4^2\), is not simple, it will be necessary that it belong to an involution of the \(2^{\mathrm{nd}}\) order \(I_2\) for which \(K\) is: 1) a double curve, or: 2) a hyperelliptic curve transformed into itself. And since in no case can \(|2K|\) have base points, one deduces that ``the only multiple bicanonical surfaces for \(p = 1\) and \(p^{(1)} = 3\) are double quartics \(F_4^2\)''.

Let us study in succession the two cases to which we are led by the two hypotheses 1) and 2) indicated above, regarding the belonging of \(|2K|\) to an involution \(I_2\).

% scan: s171L
% !PRINTED 318
\weberpage{318}
First: if hypothesis 1) is assumed, the bicanonical surface into which the given \(F\) is transformed will be a double quartic \(F_4\) with a plane branch quartic \(K\). And in order that \(K\) have genus \(p^{(1)} = 3\), and also in order that \(F_4^2\) turn out to be of genus \(p = 1\), it will be necessary that the plane sections of \(F_4\) have no double points, and that \(F_4\) itself be of genus 1. One can construct a double \(F_4\) of genus 1, representing an \(F\) with \(p = 1\) and \(p^{(1)} = 3\), by forming the pencil of a bicanonical \(F_8\) (with the said genera) and of its adjoint \(F_4\) counted twice; in this pencil
\[
\lambda F_8 + F_4^2 = 0
\]
one sees the generic surface (which is an \(F_8\)) reduce by continuity (for \(\lambda = 0\)) to the double \(F_4\). The \(F_4\) possesses as double points the triple points of \(F_8\), and these are added as branch points (infinitesimal curves) to the proper branch curve, which is the the quartic \(K\), simple intersection with \(F_8\); the remaining part, which is \(C_{14}\), curve of contact of \(F_4\) and of a \(\Phi_7\), constitutes an apparent critical curve. %

Now the number of triple points of \(C_{14}\), double curve of \(F_8\) (and also of \(\Phi_7\)), can be evaluated by the well-known formula, which gives
\[
t = 10.
\]

This number is also found in another way from the correspondence formula which links the arithmetic genera of the two surfaces of genus 1, \(F_4\) and \(F\), in correspondence \([1, 2]\).

From what precedes it results that there exists, as a particular case of the bicanonical \(F_8\) with \(p = 1\) and \(p^{(1)} = 3\), a double \(F_4\) of genus 2 possessing as branch elements a quartic plane section and 10 conical points. %
However, one might think that the double \(F_4\) constructed above represent only a subfamily of a more general type of double bicanonical surfaces with \(p = 1\) and \(p^{(1)} = 3\), which do not fall within the family of the \(F_8\). To resolve this doubt one must study our \(F_4^2\) more deeply. They can be defined by extracting on \(F_4\) the square root of a polynomial of the \(8^{\mathrm{th}}\) order, which can be reduced to \(\alpha\Phi_7\), \(\alpha = 0\) being the plane of the branch curve \(K\) and \(\Phi_7\) the surface of the \(7^{\mathrm{th}}\) order which has \(C_{14}\) (apparent critical curve) as double curve. For this polynomial \(\alpha\Phi_7\) one can substitute another polynomial \(\alpha F_7\), letting the surface vary freely within the system of the \(F_7\), of order 7, which pass triply through the 10 double points of \(F_4\) and touch it along a curve of the complete system \(\lvert C_{14}\rvert\). Indeed one can enlarge this system by adding to it twice the neighborhoods of the 10 conical points, thus taking, in place of a surface \(F_7\) which passes simply through these 10 points and touches \(F_4\) along the apparent critical curve of the \(14^{\mathrm{th}}\) %
% scan: s171R
% !PRINTED 319
\weberpage{319}
order (which we again call \(C_{14}\)). And it is then permissible to add to or subtract from such an \(F_7\) a surface of even order \(2n\) which touches \(F_4\) along a curve \(C_{4n}\), thus substituting for \(C_{14}\) an apparent critical curve \(C_{14} \pm C_{4n}\). More precisely, we wish to prove that one can remove from the system \(\lvert C_{14}\rvert\), cut out on \(F_4\) by the tangent \(F_7\) which pass through its 10 conical points, two quartic plane sections (each counted twice), so as to define \(F_4^2\) by the extraction of a \(\sqrt{\alpha F_3}\); where \(F_3\), equated to zero, represents a cubic tangent to \(F_4\) along a sextic \(C_6\) of genus 3; this last may therefore be taken as apparent critical curve of \(F_4^2\).

For this purpose let us observe that the linear system of the \(C_{28}\) cut out on \(F_4\) by the \(F_7\) through 10 conical points has degree
\[
4 \cdot 7^2 - 2 \cdot 10,
\]
and hence the half-system of the \(C_{14}\) cut out by the tangent \(F_7\) is of degree
\[
7^2 - 5 = 44,
\]
and therefore has genus and dimension \(\dfrac{44 + 2}{2} = 23\). Now the system \(\lvert C_{14}\rvert\)-cuts out on a plane quartic \(C_4\) a series \(g_{14}^{11}\), so that by removing \(C_4\) from \(\lvert C_{14}\rvert\) one will obtain a system \(\lvert C_{10}\rvert\) of dimension and genus \(23 - 12 = 11\). And since this \(\lvert C_{10}\rvert\) cuts out on another quartic plane section a \(g_{10}^7\), by removing this curve too one will have an \(\infty^3\) linear system of sextics \(C_6\) of genus 3, curves of contact of cubics passing through 10 conical points of \(F_4\). %

In conclusion: \textit{the bicanonical surface with \(p = 1\) and \(p^{(1)} = 3\), represented on a double surface of the \(4^{\mathrm{th}}\) order \(F_4\) of genus 1, is defined in relation to an \(F_4\) which is touched along a sextic \(C_6\) of genus 3 by a surface of the \(3^{\mathrm{rd}}\) order \(F_3\), and which therefore has 10 conical points on this \(C_6\); \(F_4^2\) is obtained by extracting \(\sqrt{\alpha F_3}\) on \(F_4\)}, where \(\alpha\) denotes the plane of the branch quartic.

It remains to verify that the \textit{said \(F_4^2\) falls as a particular case within the family of the simple bicanonical \(F_8\) (\(p = 1,\ p^{(1)} = 3\)).}

For this it suffices to compute the number of invariants or moduli on which our \(F_4^2\) depend. First of all it will be observed that ``the existence of a sextic of genus 3, \(C_6\), self-residual with respect to the system cut out by the cubic surfaces, entails for \(F_4\) 10 conditions and the possession of 10 double points (in general conical) on \(C_6\)''. Indeed, the existence of a \(C_6\) of genus 3 on \(F_4\) entails, to begin with, \textit{one} condition for \(F_4\); this condition being satisfied, there will be \(\infty^6\) cubics \(F_3\) cutting out on \(F_4\) two sextics \(C_6\) of the same genus 3, each residual to the other; and two such sextics --- forming together a curve \(C_{12}\) of genus 19 --- will have 14 points in common; but in order that the residual \(C_6\) be equi%
% scan: s172L
% !PRINTED 320
\unskip\weberpage{320}%
\unskip valent, it is necessary that among the 14 points mentioned there be 4 variable points and 10 base points, which will necessarily fall at 10 double points of \(F_4\); and the existence of a base point, a fixed point for the series \(g_{14}^{11}\) cut out by the \(C_6\) of a system on a residual \(C_6\), entails \textit{one} condition, but the existence of 10 base points having as residual series the canonical \(g_4^2\) entails only 9 conditions. In all, therefore, the conditions which \(F_4\) must satisfy in order to possess a self-residual \(C_6\) of genus 3 are 10 in number, and they consequently entail 10 double points of \(F_4\) on \(C_6\).

This being so, the \(F_4\) possessing a self-residual \(C_6\) depend on \(34 - 10 = 24\) parameters; to these are to be added the 3 parameters on which the choice of a plane section \(K\) depends, and from them are to be subtracted the 15 parameters of a homography of space: it remains that our \(F_4^2\) depend on. %
\[
M = 24 + 3 - 15 = 12
\]
projective invariants or moduli. And since
\[
M < 9p - 2p^{(1)} + 12 = 15;
\]
one deduces that the family of the said \(F_4^2\) is contained in a wider family of simple bicanonical \(F_8\), and precisely in the one we have defined above.

There is another type of bicanonical \(F_4^2\), corresponding to an \(F\) with genera \(p = 1\) and \(p^{(1)} = 3\) and with a hyperelliptic (irreducible) canonical curve, and it is obtained by taking as double a quartic \(F_4\) with sections of genus 2 (that is, endowed with a double line) with a branch curve \(C_{12}\) of order 12; and this \(F_4^2\) too falls as a particular case within the family of the simple bicanonical \(F_8\).

Representing \(F_4\) on a plane, one obtains a \textit{double plane with branch curve \(C_{10}\) of order 10, endowed with a 4-fold point \(O\) and with 4 points \([3, 3]\): \(A_1A_2A_3A_4\)}. In fact, on this double plane there is \textit{one} canonical curve \(K\) of genus \(p^{(1)} = 3\), which has as image the conic through \(OA_1A_2A_3A_4\), and a system of bicanonical curves, represented doubly by \(\infty^3\) quartics \(C_4\) with \(O\) double and tangent at \(A_1A_2A_3A_4\), which form the representative system of the plane sections of \(F_4\).

One recognizes that \textit{this double plane falls as a particular case within the family of the simple bicanonical \(F_8\)} (with \(p = 1\) and \(p^{(1)} = 3\)), because its branch curve \(C_{10}\) depends on
\[
65 - 8 - 36 - 8 = 13
\]
projective invariants or moduli, and one has
\[
13 < 9p - 2p^{(1)} + 12 = 15.
\]

We can summarize the results of the discussion carried out by stating that: \textit{all surfaces with \(p = 1\) and \(p^{(1)} = 3\)}, at least if the ca%
% scan: s172R
% !PRINTED 321
\unskip\weberpage{321}%
\unskip nonical curve is irreducible, \textit{fall within the more general family of the simple bicanonical \(F_8\) of ordinary space.} It also seems that the possible reducibility of the canonical curve does not give rise to any new case. But we shall not dwell on a minute examination of the question.

\edsection{20. Surfaces with \(p = 0\) and \(p^{(1)} = 3\).}{20. Surfaces with p = 0 and p\textasciicircum{}(1) = 3.}{sec.VIII.20}

Here the bicanonical curves, forming a net of degree \(4(p^{(1)} - 1) = 8\), lead in general to a multiple plane of order 8; therefore it is convenient to seek the type of the given surfaces by resorting to the \textit{tricanonical surface, which is an \(F_{18}\) of order 18 in the space of \(3p^{(1)} - 2 = 7\) dimensions.}

The high characters of this surface make it difficult to construct it directly in \(S_7\), or to define precisely a projection \(F'_{18}\) of it in \(S_3\): which will have to have plane sections of genus \(6p^{(1)} - 2 = 16\), and will therefore possess, in general, a double curve \(C_{120}\) of order 120. But we can at least prove the \textit{existence of the \(F_{18}\)}, because a particular case of the family is known, namely a double plane of Campedelli, defined by a branch curve \(C_{10}\) of order 10 with 6 points \([3, 3]\), not lying on a conic.

It is easy to see that the said double plane leads to a simple tricanonical surface \(F_{18}\). But this is not the most general \(F_{18}\), because the branch curve \(C_{10}\) named above depends on 3 projective invariants or moduli, and one has
\[
3 < 9p - 2p^{(1)} + 12 = 6.
\]

Let us add that, at least if the net of bicanonical curves of \(F\) (a regular surface) is irreducible, the \textit{corresponding tricanonical surface (\(p = 0\) and \(p^{(1)} = 3\)) is certainly simple.} Indeed a multiple tricanonical surface (of order \(\dfrac{18}{2} = 9\)) could arise only under the hypothesis that the bicanonical curves \(K_2\) of \(F\) are hyperelliptic, the (complete) canonical series which the tricanonical curves cut out on a \(K_2\) then turning out to be composite. Under this hypothesis the characteristic series of the bicanonical net on a generic curve \(K_2\) will be a \(g_8^2\) composed with the \(g_2^1\) belonging to \(K_2\) itself, and hence contained in a complete \(g_8^4\); but (if the surface is to be regular) this entails \(P = 6\) instead of \(P = 3\). %

\edsection{21. Simple and multiple pluricanonical surfaces.}{21. Simple and multiple pluricanonical surfaces.}{sec.VIII.21}

In the preceding sections we have met several types of surfaces which give rise to double bicanonical surfaces, in particular the following three general types:

\printersignature{\textsc{Enriques F.} - \textit{Superficie algebriche.}\quad 21}

% scan: s173L
% !PRINTED 322
\weberpage{322}
I) For \(p = 3\) and \(p^{(1)} = 3\), the double plane with branch curve \(C_8\) of order 8;

II) For \(p = 2\) and \(p^{(1)} = 2\), the double plane with branch curve \(C_{10}\) of order 10 endowed with two infinitely near quintuple points;

III) For \(p = 2\) and \(p^{(1)} = 3\), the double plane with branch curve \(C_{12}\) of order 12, endowed with two quadruple points and with two points \([3, 3]\).

For types I) and II) not only the bicanonical but also the tricanonical surface turns out to be a double surface, inasmuch as the bicanonical curves (of which the tricanonical ones are the adjoints) are hyperelliptic curves. %

Now we propose to prove that these types I) and II) are the only types of surfaces, with \(p > 0\) and \(p^{(1)} > 1\), whose tricanonical surfaces are multiple surfaces.

For this purpose let us note that, since the bicanonical curves, or their variable parts, are irreducible curves (cf. \S~19, Chap. IV), the adjoints and hence the tricanonical curves, or their variable parts, will certainly form a simple system if the said bicanonical curves are not hyperelliptic (and hence belonging to a rational involution \(I_2\)); if therefore a regular tricanonical surface (for \(p > 0\) and \(p^{(1)} > 1\)) is not simple, it must be representable on a double plane in such a way that the (hyperelliptic) bicanonical curves have as images the rational curves of a linear system at least \(\infty^2\), to which some fixed component may possibly be added.

It is convenient to examine in succession the three cases in which the images of the said bicanonical curves reduce to:

1) the \(\infty^2\) lines of the plane;

2) the \(\infty^5\) conics of the plane;

3) or, finally, to the curves \(C_n\) of a certain order \(n\) (\(> 1\)) forming a linear system of dimension \(2n - s\) with an \((n - 1)\)-fold base point \(O\) and without other base points (\(s = 0\)); or with another simple base point \(O_1\) (\(s = 1\)), or with \(s\) (\(> 1\)) simple base points \(O_1O_2\ldots O_s\) infinitely near to \(O\) (in distinct directions).

Hypothesis 1) --- which gives \(P = p + p^{(1)} = 3\), \(p = 1\) and \(p^{(1)} = 2\) --- is impossible, as appears from the examination of surfaces with these characters which we carried out in \S~15.

Hypothesis 2) leads to the double plane with branch curve \(C_8\), that is, to type I) with \(p = 3\), \(p^{(1)} = 3\), met in \S~17.

Hypothesis 3) is realized for \(n = 2\) when one takes as branch curve of the double plane a \(C_8\) of order 8 with two infinitely near triple points \(O\) and \(O_1\): the images of the bicanonical curves are the conics through \(O\) and \(O_1\).

% scan: s173R
% !PRINTED 323
\weberpage{323}
But the double plane so defined falls within the more general type where one takes as branch curve a \(C_{10}\) of order 10 with two infinitely near 5-fold points, which is type II) for \(p = 2\) and \(p^{(1)} = 2\): when the \(C_{10}\) also has a double point, it is reduced by a quadratic transformation to a \(C_8\) with a point \([3, 3]\).

To discuss in general the possibilities to which hypothesis 3) gives rise, suppose that the bicanonical curves of our double plane contain a fixed part of order \(d\), to be added to the \(C_n\); in order that the addition of this part should not enlarge the system \(\lvert C_n\rvert\), it is necessary that it consist of lines fundamental for \(C_n\), counted \(h_1h_2\ldots h_s\) times respectively, that is, of lines \(OO_i\), and hence that it pass with multiplicity \(d\) through \(O\) and with multiplicity \(h_i\) through \(O_i\). Thus the branch curve of the double plane will be a \(C_{2m}\) of order \(2m\) which has the curves \(C_n + \Sigma h_i \cdot OO_i\), of order \(n + d = n + \Sigma h_i\), as second adjoints; and we shall have
\[
2m = n + d + 6.
\]

Besides this it will suffice to bear in mind that \(C_{2m}\) must possess at \(O\) the multiplicity \(n - 1 + d + 2 = n + d + 1 = 2m - 5\). Now on a double plane with branch curve \(C_{2m}\) endowed with a \((2m - 5)\)-fold point \(O\), the images of the bicanonical curves are in general curves of order \(2m - 6 = n + d\) passing through \(O\) with multiplicity
\[
2m - 8 = n + d - 2;
\]
in order that this multiplicity increase by 1, becoming \(n + d - 1\) (as happens for the \(C_n + \Sigma h_i \cdot OO_i\)), it is necessary that near \(O\) there be another \((2m - 5)\)-fold point of \(C_{2m}\); but the sum
\[
2m - 5 + 2m - 5 \leq 2m,
\]
and hence
\[
m = 5,\qquad 2m = 10
\]
or (in particular)
\[
m = 4,\qquad 2m = 8;
\]
thus one finds type II) again.

In conclusion we shall state the theorem:

\textit{The tricanonical surface of a regular surface with genera \(p > 0\) and \(p^{(1)} > 1\) is always a simple surface, of order \(\leq 9(p^{(1)} - 1)\), with the exception of two types I) and II) of surfaces referable to the double plane, which correspond to the values \(p = 3\) and \(p^{(1)} = 3\), and \(p = 2\) and \(p^{(1)} = 2\).}

For these two types I) and II) both the bicanonical and the tricanonical surface turn out to be double: the bicanonical surface is, for I), a double Veronese surface (\(F_4\) of \(S_5\)) and, for II), a double quadric cone of \(S_3\); the tricanonical surface is, in the first case, a double surface \(F_9\), of order 9 with elliptic sections in \(S_9\), in the second case a double surface \(F_6\), of order 6 with elliptic sections in \(S_6\).

% scan: s174L
% !PRINTED 324
\weberpage{324}
Since, then, the tricanonical curves of our surfaces are in no case hyperelliptic, it follows (for \(p > 0\)) that their adjoints do not belong to an involution, that is: \textit{on a regular surface, for \(p > 0\) and \(p^{(1)} > 1\), the quadricanonical curves} (or their variable components), \textit{and hence also the \(i\)-canonical curves for \(i > 4\), always form a simple system.}

This proposition, which asserts the simplicity of the pluricanonical surfaces (\(p^{(1)} > 1\)) for \(i > 3\), will in all likelihood extend to the case \(p = 0\); but this requires a more difficult analysis, on which we shall not dwell (cf., for \(p^{(1)} = 2\), the discussion of \S~16). %

\clearpage\thispagestyle{plain}
% scan: s174R
% !PRINTED none
\weberpage{[s174R]}
\edchapter{IX}{IRREGULAR SURFACES AND CONTINUOUS SYSTEMS\\OF NON-EQUIVALENT CURVES}{Chapter IX. IRREGULAR SURFACES AND CONTINUOUS SYSTEMS OF NON-EQUIVALENT CURVES}{IRREGULAR SURFACES AND CONTINUOUS SYSTEMS OF NON-EQUIVALENT CURVES}

\edsection{1. Introduction.}{1. Introduction.}{sec.IX.1}

Let us recall some remarks that we made in Chap. IV, \S~12.

And first of all we point out that, when speaking of curves and of linear or continuous systems of curves \(C\) on a surface \(F\), we shall as a rule understand that \textit{irreducible curves without multiple points} are meant\origfootnote{1}{s174R-1}{If one wishes to consider also curves \(C\) with multiple points, it is convenient to regard these points as assigned base points for the linear and continuous systems referred to.}.

We have seen that, on a regular surface, a complete linear system of curves \(\lvert C\rvert\) has a complete characteristic series. And a surface to which there belongs a continuous system of curves \(C\) formed of \(\infty^d\) distinct linear systems, that is, a system \(\{C\}\) containing \(\infty^d\), with \(d > 0\), non-equivalent curves, is irregular, of irregularity
\[
p_g - p_a \geq d.
\]

The assertion is immediately justified for continuous systems \(\{C\}\) of irreducible curves \(\infty^{r+d}\), formed of \(\infty^d\) linear systems having a certain (generic) dimension \(r\): for the characteristic series cut out on a generic \(C\) by the infinitely near curves of \(\{C\}\) will have the dimension
\[
r + d - 1,
\]
and hence the characteristic series of the linear system \(\lvert C\rvert\) will have a deficiency (at least) equal to \(d\); it is therefore enough to invoke the theorem of Castelnuovo (Chap. IV, \S~12), by which
\[
d \leq p_g - p_a.
\]

Let us add that this inequality holds also for continuous systems \(\{C\}\) formed of reducible curves; because, in the first place,
% scan: s175L
% !PRINTED 326
\weberpage{326}
one may disregard any fixed components forming part of the \(C\), and in the second place, if the variable \(C\) are reducible, by subtracting them from a sufficiently large linear system \(\lvert L\rvert\) one can obtain a continuous system \(\{L - C\}\) which --- apart from fixed components --- will consist of irreducible curves and will contain at least --- like \(\{C\}\) --- \(\infty^d\) non-equivalent curves.

The principal aim of this Chapter is to set out the arguments which lead to the inversion of the preceding theorem, in the sense that: on an irregular surface, of genera \(p_a\) and \(p_g > p_a\), and hence of irregularity \(q = p_g - p_a\), there always exist continuous systems of curves formed of \(\infty^q\) distinct linear systems.

The possession of continuous systems of non-equivalent curves will thus constitute the \textit{characteristic property of irregular surfaces.} And the theorem which states it and makes it precise as above, with regard to the value of \(p_g - p_a\), may be called the \textit{fundamental theorem} in the theory of irregular surfaces.

But before tackling the proof of the fundamental theorem, in view of the delicate difficulties which it presents, it is convenient to place ourselves under the more general hypotheses in which it is set aside, and thus to examine what can be said of surfaces containing continuous systems \(\infty^q\) of non-equivalent curves when one knows only that
\[
0 \leq q \leq p_g - p_a,
\]
as was recognized above. For this purpose we shall distinguish two irregularities of the surface, which are a priori not always equal: \textit{the arithmetic irregularity}
\[
p_g - p_a
\]
and \textit{the geometric irregularity} \(q\), that is, the maximum dimension of a continuous series of non-equivalent curves lying on the surface:
\[
q \leq p_g - p_a.
\]

This number \(q\) denotes, in any case, an invariant character of the surface itself: all the known examples of irregular surfaces (surfaces with an irrational pencil of curves, surfaces of the pairs of points of a curve, etc.) give \(q > 0\), and indeed \(q = p_g - p_a\).

\edsection{2. Arithmetic condition for a curve, on a surface of geometric irregularity \(q\), to belong to a continuous series \(\infty^q\) of non-equivalent curves.}{2. Arithmetic condition for a curve, on a surface of geometric irregularity q, to belong to a continuous series ∞\textasciicircum{}q of non-equivalent curves.}{sec.IX.2}

Let \(q > 0\) be the geometric irregularity of a surface \(F\): this means that there belongs to \(F\) a continuous series \(\infty^q\) \(\{D\}\), and not a larger series, of non-equivalent curves.

% scan: s175R
% !PRINTED 327
\weberpage{327}
We ask: can an arithmetic criterion be given for a curve \(C\) on \(F\) to belong likewise to a series \(\infty^q\) of non-equivalent curves, that is, to a series formed of \(\infty^q\) distinct complete linear systems?

The answer is furnished by the computation of the virtual dimension of the system determined by the curve \(C\) (of genus \(\pi\), degree \(n\), and index of speciality \(i\)):
\[
r = p_a + n - \pi + 1 - i.
\]

Indeed, let \(\lvert C\rvert\) be a linear system of curves, irreducible or reducible, of degree \(n\), of genus \(\pi\), and of index of speciality \(i\), which is \textit{arithmetically effective}, that is, such that its virtual dimension
\[
p_a + n - \pi + 1 - i \geq 0;
\]
with it we shall associate another regular linear system \(\lvert D\rvert\) (irreducible, etc.), sufficiently large, so that \(\lvert C + D\rvert\) too is regular. The system \(\lvert D\rvert\) may be supposed to be contained in a continuous system \(\{D\}\) formed of \(\infty^q\) non-equivalent linear systems, one of which --- different from \(\lvert D\rvert\) --- we shall denote by \(\lvert\overline{D}\rvert\).

Now the system \(\lvert C\rvert\) can be constructed by detaching a curve \(D\) from \(\lvert C + D\rvert\), and its dimension turns out to be \(\geq p_a + n - \pi + 1 - i\), since \(\lvert C + D\rvert\) cuts out on the said \(D\) a series (complete or not) which has precisely the index of speciality \(i\) (cf. Chap. IV, \S~11).

Let us consider the series cut out by the same system \(\lvert C + D\rvert\) on a \(\overline{D}\). It is not possible that, for a generic choice of \(\overline{D}\) in \(\{D\}\), the said series should have an index of speciality greater than \(i\), because --- by varying \(\overline{D}\) continuously --- one can reduce it to \(D\), and in this passage to the limit the index of speciality of the series cannot decrease. Therefore the dimension of the linear system \(\lvert\overline{C}\rvert = \lvert C + D - \overline{D}\rvert\) will turn out (according to the well-known computation which yields the Riemann-Roch theorem) to be \(\geq p_a + n - \pi + 1 - i\): thus, as \(\overline{D}\) varies, one will obtain \(\infty^q\) distinct complete linear systems forming an irreducible series \(\{C\}\).

Let us state the theorem thus obtained:

\textit{On a surface of geometric irregularity \(q > 0\), an arithmetically effective linear system, that is, one of degree \(n\), genus \(\pi\) and index of speciality \(i\) such that}
\[
p_a + n - \pi + 1 - i \geq 0,
\]
\textit{always belongs to a series \(\infty^q\) of non-equivalent linear systems.}

Here one sees the simplest definition of the geometric irregularity (and then of the difference \(p_g - p_a\), in relation to the fundamental theorem which yields the characteristic property of irregular surfaces).

% scan: s176L
% !PRINTED 328
\weberpage{328}
\edsection{3. The Picard variety corresponding to an irregular surface.}{3. The Picard variety corresponding to an irregular surface.}{sec.IX.3}

In the preceding section we have seen how, on a surface of geometric irregularity \(q\ (>0)\), the \(\infty^q\) linear systems forming a certain continuous system \(\{|C|\}\) or \(\{C\}\) are linked to those of another analogous continuous system \(\{D\}\): there is a one-to-one correspondence between the elements (linear systems) of \(\{C\}\) and those of \(\{D\}\), the linear systems residual to each other with respect to a \(|C+D|\) being homologous. This remark (which is due to \textsc{Castelnuovo}) can be made precise by saying: \textit{The elements (linear systems) of the continuous series \(\infty^q\) belonging to a surface of geometric irregularity \(q\ (>0)\) can be regarded as the points of one and the same Abelian variety \(V\)}, which \textsc{Castelnuovo} chose to name after the French geometer, as the \textit{Picard variety} corresponding to the surface.

Indeed, the operation \(+D-\overline{D}\) by which one passes from a linear system \(|C|\) to another linear system \(|\overline{C}|\) contained in \(\{C\}\) can be interpreted as a point transformation of the variety \(V\) into itself; it is obvious that all the analogous transformations are permutable and form a group \(\infty^q\), simply transitive without exceptions, within the said variety \(V=V_q\) (of \(q\) dimensions): which for precisely this reason is characterized as an \textit{Abelian variety}, and a non-degenerate one.

The meaning of the theorem obtained is clarified by the following considerations.

Let us pose the problem of determining all the curves of a given order \(n\) belonging to an algebraic surface \(F\).

A priori these curves, which — in order to belong to \(F\) — are subjected to algebraic conditions, will be distributed into a certain number of families and will therefore constitute, in general, algebraic continuous systems.

Such systems will be linear systems if \(F\) is regular. If, on the other hand, \(F\) is of geometric irregularity \(q\) \((p_g>p_a)\), these systems — regarded as varieties having distinct linear systems as elements — will be, not algebraic varieties of arbitrary type, but Abelian varieties \(V_q\), or at least they will form Abelian varieties together with their multiples.

Some reminders, which it will be useful to give precisely about Abelian varieties, will help us to understand better what this implies (cf. § 9).

\edsection{4. Characteristic property of irregular algebraic surfaces: fundamental theorem for \(p_g=0\).}{4. Characteristic property of irregular algebraic surfaces: fundamental theorem for p\_g=0.}{sec.IX.4}

The proof of the fundamental theorem on irregular surfaces, which we stated in § 2, is given in the sim%
% scan: s176R
% !PRINTED 329
\unskip\weberpage{329}%
plest way for surfaces of genus \(p_g=0\), and requires instead more delicate considerations for \(p_g>0\). It is therefore convenient to distinguish these two cases.

Let us first suppose that \(p_g=0\), and hence \(p_a=-p<0\).

Consider, on our surface \(F\), a complete linear system \(|C|\) which is regular, non-special, of genus \(\pi\) and of degree \(n\), and hence of dimension
\[
r=-p+n-\pi+1.
\]

And with it associate another linear system \(|D|\), which is also regular, of genus \(\varrho\) and of degree \(\nu\), and hence of dimension
\[
r'=-p+\nu-\varrho+1,
\]
and such that the sum \(|C+D|\) too is a regular system (cf. Chap. IV, § 8).

Denoting by \(m\) the number of intersections of \(C\) and \(D\), we shall write the genus and the degree of \(|C+D|=|E|\):
\[
\begin{aligned}
\Pi &= \pi+\varrho+m-1\\
N &= n+\nu+2m,
\end{aligned}
\]
and hence its dimension:
\[
R=-p+N-\Pi+1=-p+(n-\pi)+(\nu-\varrho)+m+2,
\]
that is,
\[
R=r+r'+p+m.
\]

Let us recall that, on the basis of the \textit{principle of degeneration or of splitting}\origfootnote{1}{s176R-1}{Cf. \textsc{Enriques-Chisini}, Bk. V, Chap. III, 35 (vol. III, p. 405).}, an irreducible curve endowed with a certain number \(m\) of double points cannot reduce, by continuity, to a split curve whose components, if they are free of multiple points, do not have at least \(m+1\) common points (that is, with one more double point, through which the two components of the split curve are connected). From this principle one deduces that, by requiring the curves \(E\) to possess \(m\) double points, one will not be able to obtain a continuous system \(\overline{E}\) of dimension \(r+r'+p>r+r'\), formed of irreducible curves, which contains within itself the system of the curves of \(|E|=|C+D|\) split into a curve of \(|C|\) and a curve of \(|D|\); rather, one will come to define a system of which at least a part (of the same dimension) contains within itself a continuous system of curves split into a curve of the order of \(C\) and a curve of the order of \(D\), these two components thus describing two non-linear continuous systems, say \(\{C\}\) and \(\{D\}\), complementary to each other.

% scan: s177L
% !PRINTED 330
\weberpage{330}
To each curve \(C\) of the larger continuous system \(\{C\}\) there will correspond, as residual with respect to \(|E|\), a linear system of curves \(|D|\), and conversely; so that it is easy to evaluate the dimension of the continuous system \(\{C\}\), which must be (at least) \(r+p\), and similarly that of \(\{D\}\), which must be \(r'+p\). Indeed, one can say that \(\{C\}\) will consist of a series (at least) \(\infty^p\) of non-equivalent linear systems having dimension equal to \(r\), because the dimension of the complete linear systems forming \(\{C\}\), which is a priori \(\ge r\), cannot be \(>r\), seeing that these linear systems reduce by continuity to the \(|C|\) from which we started, which has exactly the dimension \(r\).

But in the case where the geometric genus of \(F\) is \(p_g=0\), we know that \(F\) cannot contain a continuous system of curves formed of \(\infty^d\) distinct linear systems where \(d>p_g-p_a=p\); therefore one concludes that, on irregular surfaces of geometric genus \(p_g=0\), there exist continuous systems of \(\infty^p\) non-equivalent curves, and not systems formed of a larger series of such curves.

We shall summarize the results obtained by stating the \textit{Restricted fundamental theorem: A surface of negative arithmetic genus \(p_a=-p\) certainly contains continuous systems formed of \(\infty^p\) non-equivalent curves, and, if the geometric genus \(p_g=0\), it does not possess larger systems of non-equivalent curves; in other words the geometric irregularity is}
\[
q\ge -p %
\]
and for \(p_g=0\) \textit{it equals the arithmetic irregularity}
\[
p_g-p_a=p.
\]

\edsection{5. Fundamental theorem for \(p_g>0\).}{5. Fundamental theorem for p\_g>0.}{sec.IX.5}

Let us try to extend the fundamental theorem to surfaces of geometric genus \(p_g>0\) and of arbitrary arithmetic genus \(p_a<p_g\). For this purpose let us again consider, on the surface \(F\) (supposed to be of genera \(p_a\) and \(p_g\)), two regular non-special linear systems \(|C|\) and \(|D|\) with the characters \(n,\pi\) and \(\nu,\varrho\), such that the sum \(|E|=|C+D|\) too is regular, cutting out on a \(D\) a complete and non-special series \(g_{m+\nu}\).

Denoting by \(N\) and \(\Pi\) the characters of \(|E|\), one sees (as above in the case \(p_g=0\)) that the linear system \(|E|\) will have the dimension
\[
\begin{aligned}
& R=p_a+N-\Pi+1=p_a+(n-\pi)+(\nu-\varrho)+m+2\\
& \qquad\qquad\qquad\qquad(m=CD),
\end{aligned}
\]
that is,
\[
R=r+r'+m-p_a,
\]
% scan: s177R
% !PRINTED 331
\weberpage{331}
where \(r\) and \(r'\) are the dimensions of \(|C|\) and \(|D|\). Now, if the linear system \(|C|\) is to be contained in a continuous system \(\{C\}\) of dimension \(r+q=r+p_g-p_a\) (formed of \(\infty^q\) distinct linear systems), it is necessary that the imposition on the \(E\) on \(F\) of \textit{\(m\) double points}, from which there results \textit{the splitting of such curves} into a \(C\) of \(\{C\}\) and a \(D\) of the residual system \(\{D\}\), should entail — not \(m\) independent conditions (or \(m-p_a\), as would appear from the existence of \(|C|\) and of \(|D|\)) — but \textit{only \(m-p_g\) conditions}. If this happens, the continuous system \(\{C\}\), of dimension \(p_g+n-\pi+1\), will have, on a \(C=C_0\), a complete characteristic series, and will hence not be capable of enlargement.

Thus, in order to show that « the imposition on the \(E\) on \(F\) of \(m\) double points, from which there results the splitting of such curves into curves of \(\{C\}\) and \(\{D\}\), cannot entail more than \(m-p_g\) conditions », and hence that the geometric irregularity \(q\) of \(F\) equals the arithmetic irregularity \((q=p_g-p_a)\), one is led to examine the characteristic series of the continuous system \(\{C\}\); it seems that the proof sought will be attained if one recognizes that this series is complete.

And for this purpose it will suffice to consider the curves \(E\) of \(|E|=|C+D|\) which pass through the group \(G\) of the \(m\) points common to the two components \(C_0\) and \(D_0\) of a split \(E\), \(E_0=C_0+D_0\); these \(E\) cut out on \(C_0\) the complete characteristic series, and one of them, together with \(E_0\), determines a linear pencil of curves which touch one another at the points of \(G\): it follows from this that the curve of the pencil infinitely near to \(E_0\) also possesses \(m\) double points near to those of \(G\)\origfootnote{1}{s177R-1}{Cf. e.g. \textsc{Enriques-Chisini}, \textit{Lezioni}, Bk. II, § 5, vol. I, p. 183 and Bk. III, Chap. II, § 20, vol. II, p. 167.}, and therefore, if the principle of degeneration is applicable also to these infinitely near curves, it must split into two curves \(C_1\) and \(D_1\), the first of which — near to \(C_0\) — cuts out on \(C_0\) a group (arbitrarily assignable) of the characteristic series.

Have we thus attained the proof (for arbitrary \(p_g\)) of the fundamental theorem?

No, the appearance of a proof must not deceive us: the theorem in question concerns an \textit{integral property}, and we have established only a \textit{differential property} which is a consequence of it, but which conversely does not necessarily imply the former. Indeed, what has been proved is at most that « the characteristic series of \(C_0\) is cut out completely by curves \(C_1\) infinitely near to it », but not that the \(C_1\) really belong, as curves tending to the limit, to a continuous system of proper curves \(\{C\}\).

There is here a critical point of the proof, whose historical significance we shall bring out later (§ 6): the condition by which « \(m-p_g\) double points imposed on the \(E\) must entail as a consequence a further
% scan: s178L
% !PRINTED 332
\weberpage{332}
\(p_g\) double points » has been verified not « in the large » but « in the small », up to infinitesimals of higher order, for the curves infinitely near to \(E_0\).

The differential property « in the small » follows, as has been said, from the integral, or « in the large », property which one wishes to establish, but the latter is not in its turn a consequence of the former.

To clarify the matter better, consider a variety \(V_r\) of dimension \(r\) which is defined in a space \(S_R\) \((R>r)\) as the intersection of a certain number of hypersurfaces (or of differentially distinct linear sheets of one and the same algebraic hypersurface): to evaluate the dimension \(r\) one is led to look for the number \(s\) of linearly independent points in the neighborhood of the first order of a generic point \(P\) of \(V_r\); but it is clear that the varieties of which \(V_r\) is defined as the intersection may have a certain contact at every point of it, and in that case the number of points near to \(P\) mentioned above will turn out to be \(s>r\).

For example, for \(r=1\), \(V_1\) may be the curve of contact of two or more surfaces of the space \(S_3\), a curve to be counted twice; then the number of its independent points near to a point \(P\) is no longer \(r=1\), but \(s=2\).

From this clarification there also emerges the requirement which must be satisfied in order to fill the gap that remains in our proof of the general fundamental theorem for \(p_g>0\), by passing from the property recognized in the small to the property in the large, of which we spoke above.

Consider the \(V_r\) of \(S_R\), defined as the intersection of \(m\) hypersurface sheets which correspond to the curves of \(|E|=|C+D|\) endowed with \(m\) double points (which constitute the largest continuous system to which there belong the \(E\) split into a curve \(C\) of the order of \(C_0\) and a \(D\) of the order of \(D_0\)): the said \(m\) sheets must not all touch one another along \(V_r\), or at a generic point \(P\) of it. This means that there cannot be a linear branch of curve issuing from \(P\) which has a contact of finite order \(s\) \((\ge 1)\) with the aforesaid sheets and not with \(V_r\); in other words: if such a branch has \(h<s+1\) \((>0)\) successive points in common with \(V_r\), it will also have in common with this variety the next point, the \((h+1)\)—th. %

Or, returning from \(V_r\) to our surface \(F\), one must show that for every curve \(C_1\) infinitely near to \(C_0\) in the neighborhood of the first order there exists a series of infinitely near curves \(C_1C_2\ldotp\ldotp\ldotp\ldotp\) in the neighborhoods of successive order \(1,\ 2,\ 3\,\ldotp\ldotp\ldotp\ldotp\); a \textit{regular} series (that is, one assimilable to a linear branch). The curve \(C_1\) forms part of a split curve \(E_1=C_1+D_1\) near to \(C_0+D_0\); \(C_2\) will form part of a successive curve \(C_2+D_2\), and so on. %

To meet the requirement thus posed, and hence to complete the general dem%
% scan: s178R
% !PRINTED 333
\unskip\weberpage{333}%
onstration of the fundamental theorem, an idea presents itself quite simply\origfootnote{1}{s178R-1}{The part which follows was left incomplete by the Author, and hence the arguments developed in it present numerous gaps; nevertheless it has been thought appropriate to reproduce it, because it contains ideas which perhaps, suitably completed, may provide the starting point for a systematization of the theory. (Note by the revisers).}: since two infinitely near non-equivalent curves \(C_0\) and \(C_1\) are given on \(F\), the operation \(C_1-C_0\), applied successively to a curve, e.g. to \(C_1\), will serve to define the series of the successive curves \(C_2C_3\,\ldotp\ldotp\ldotp\ldotp\), or that of the linear systems determined by it. Here, to tell the truth, a bold intuition comes into play, namely it is assumed that « one can operate on the surface \(F\) by addition and subtraction of curves, even when infinitely near curves are concerned ». We wish to justify this intuition. %

Let us start from a somewhat analogous question, which refers to curves. If, on a curve \(L\) of genus \(p>0\), one considers two non-equivalent groups of \(p\) points, \(G\) and \(G'\), the operation \(+G'-G\), applied successively starting from \(G\), leads to a (generally) unlimited series of groups %
\[
G,\quad G',\quad G'',\quad G''',\quad \ldotp\ldotp\ldotp\ldotp,
\]
which correspond to the homologues of \(G\) under the powers of a transformation of the first kind \(T\) of the Jacobian variety \(V\) relative to \(L\). Therefore the series of the successive groups \(G'',\,G'''\ldotp\ldotp\ldotp\ldotp\) does not cease to be defined\origfootnote{2}{s178R-2}{In the precise sense of the differential conditions which define infinitely near points. Cf. \textsc{Enriques-Chisini}, \textit{Lezioni}, Bk. IV, vol. II.} when \(G'\) becomes infinitely near to \(G\): in this case the transformation \(T\) becomes an infinitesimal transformation generating the continuous group of \(V\), in the sense of \textsc{S. Lie}, and defines an analytic series \(t\), a trajectory of the Jacobian group of \(L\).

To pass from here to the consideration of infinitely near curves on a surface, it is convenient first of all to justify the principle of splitting for the curve infinitely near to a \(C_0+D_0\) which has (as in the construction indicated above) \(m=C_0D_0\) double points near to the points common to \(C_0\) and \(D_0\); it is a matter, then, of showing that such a curve is a \(C_1+D_1\), formed of two parts near to \(C_0\) and \(D_0\) respectively.

For this purpose one may suppose that \(C_0\) and \(D_0\) are non-special normal curves belonging to spaces of very high dimension with respect to their genus: one can indeed reduce to this case by means of a transformation of the given surface \(F\). Then it is easy to verify that, in the space of \(F\), the irreducible curves \(E\) of order \(N\) and genus \(\Pi\) which constitute the family defined by the split \(C+D\) form a variety having a certain dimension \(R\), within which there is a variety
% scan: s179L
% !PRINTED 334
\weberpage{334}
of dimension \(R-m\) consisting of the (split) \(E=\overline{E}\) with \(m\) double points, which include the family of the \(C+D\), so that the splitting of the \(E\) \textit{in the space} of \(F\) (we do not say on the surface) depends precisely on \(m\) independent conditions, and not on fewer.

To prove this, project the split curve \(C_0+D_0\) with \(m\) double points (common to \(C_0\) and \(D_0\)) from a suitable linear space onto an \(N\)-fold multiple line (reducible) with \(2N+2\Pi-2\) critical points. %

This \(N\)-fold line can be regarded as a Riemann surface with \(N\) sheets which — for a certain choice of loops, taken in a given order, and for a suitable naming of the branches — can be defined by a system of transpositions on the branches, similar to the one that corresponds to the Lüroth-Clebsch theorem\origfootnote{1}{s179L-1}{Cf. \textsc{Enriques-Chisini}, \textit{Lezioni}, Bk. V, Chap. I, § 3, vol. III, p. 26; Chap. III § 33, vol. III, pp. 359-376.}; thus (\(n,\pi\) and \(\nu,\varrho\) being the characters of \(C_0\) and \(D_0\) respectively) one will have:

1) first of all \(2n+2\pi-2\) loops, to each of which there corresponds an even number of similar transpositions \((12)\ (12)\,\ldotp\ldotp\ldotp\ldotp\,((n-1)n),\ ((n-1)n)\);

2) then \(m\) loops (which go to the \(m\) apparent critical points, images of the points common to \(C_0\) and \(D_0\)), to each of which there corresponds the product of two transpositions
\[
(n(n+1))\cdot(n(n+1))=1;
\]

3) and finally \(2\nu+2\varrho-2\) loops, to each of which there corresponds an even number of similar transpositions
\[
\begin{aligned}
((n+1)(n+2)),&\quad ((n+1)(n+2))\,\ldotp\ldotp\ldotp\ldotp\\
((n+\nu-1)(n+\nu)),&\quad ((n+\nu-1)(n+\nu)).
\end{aligned}
\]

Now, by the existence theorem\origfootnote{2}{s179L-2}{Cf. \textsc{Enriques-Chisini}, \textit{Lezioni}, Bk. V, Chap. III, § 33 loc. cit.}, one can make the multiple line thus defined vary by continuity, separating the two branch points which coalesce in one of the apparent critical points. It therefore appears that the split curve \(\overline{E}=C_0+D_0\), in the space of \(F\), is capable of varying in a certain variety \(\infty^R\ V\) of irreducible curves \(E\), and within this there lies a variety \(\infty^{R-m}\ \overline{V}\) of split curves with \(m\) double points, so that the curves infinitely near to \(C_0+D_0\) with \(m\) double points constitute precisely the linear neighborhood of this curve-element on \(\overline{V}\). %

Finally, then, we have proved the \textit{fundamental lemma}: « every curve \(E\) with \(m\) double points, infinitely near to \(C_0+D_0\) (in the space of \(F\)), belongs to a (regular) series \(\infty^1\) of split curves which includes \(C_0+D_0\), and must therefore be regarded as effectively
% scan: s179R
% !PRINTED 335
\weberpage{335}
split into two improper curves \(C_1\) and \(D_1\), of the order of \(C_0\) and of \(D_0\) respectively, near to these curves ».

The fundamental lemma thus established means that « the curve infinitely near to \(C_0\) which forms part of a curve near to \(C_0+D_0\) with \(m\) double points, such as we have learned to construct on the surface \(F\), \textit{depends} only on \(C_0\) and \textit{not} on the choice \textit{of a \(D_0\)}, which can be associated with it in infinitely many different ways within a regular system \(|D|\) ».

On this basis we wish to recognize that, after the said curve \(C_1\), there is further, in the neighborhood of \(2^\circ\) order of \(C_0\), a curve \(C_2\), and then a curve \(C_3\), etc.

Let us be allowed here to speak of \textit{infinitely near curves}, which are presumed to be defined by differential conditions in the same way as points on a curve, as if they were \textit{proper curves}; we reserve the right to justify later the use we make of this language.

On the surface \(F\), which contains the linear systems \(|C|\) and \(|D|\) considered above, construct a normally large linear system \(|L|\) of a certain genus \(P\), such that \(|C|\) and \(|D|\) cut out on a generic \(L\), and also on every (irreducible) \(L\) endowed with a double point, \textit{complete series} of order less than \(P\)\origfootnote{1}{s179R-1}{Cf. Chap. IV, § 10.}. Let us consider in \(|L|\) a generic pencil of curves \(L\), having a certain group \(\Gamma\) of base points: we say first of all that « the infinitely near non-equivalent curves \(C_0\) and \(C_1\) cut out non-equivalent groups on the \(L\) of the pencil ».

This assertion follows from the fact that \(|C_0|\) cuts out a complete series on every \(L\): thus — since \(M=CL>C_0^2\) — there are as many equivalent curves infinitely near to \(C_0\) as there are groups on \(L\) equivalent to \(C_0L\) and near to it; and if \(C_1\) cut out on an \(L\) a group equivalent to \(C_0L\), \(C_1\) too would have to be equivalent to \(C_0\).

Now we can slightly modify our assumptions and make them precise, assuming that \(C_0\) belongs, within \(\{C\}\), to a regular system of dimension \(r=p+n-\pi+1\), so that the curve itself is determined by \(r\) simple base points forming a subgroup \(A\) of \(\Gamma\) (and similarly, where necessary, that \(|D_0|\) too is of dimension \(r'=p_a+\nu-\varrho+1\), so that the curve \(D_0\) remains determined on imposing \(r'\) simple base points forming a subgroup \(B\) of the same \(\Gamma\)). Then in the base group \(\Gamma\) of the pencil of the \(L\) one will choose, outside \(A\) and \(B\), a further \(P-M\) points (which may be reduced to a single point to be counted \(P-M\) times), forming a certain group \(S\), which, added to the group of the \(M\) points of section of \(C_0\) with \(L\), will give on the latter a non-special group \(G\) of \(P\) points, and, added to the group cut out by \(C_1\), will similarly give a group of \(P\) points \(G_1\), infinitely near %
% scan: s180L
% !PRINTED 336
\weberpage{336}
and non-equivalent to \(G\). The correspondence between \(G\) and \(G_1\) (an infinitesimal transformation of the group of the Jacobian variety of \(L\)) generates on the \(L\) of the pencil an analytic series \(t\) of groups of \(P\) points, to which there belongs a group \(G_2\) infinitely near to \(G\) in the neighborhood of \(2^\circ\) order of \(G\), that is, successive to \(G_1\). As \(L\) varies in its pencil, the group \(G_2\), which turns out to be determined algebraically and rationally\origfootnote{1}{s180L-1}{The differential expression which defines \(G_1\), from which that of \(G_2\) can be derived, is to be found computed by \textsc{Chisini}. Cf. \textsc{Enriques-Chisini}, \textit{Lezioni}, Bk. VI, Chap. II, § 28, vol. IV, p. 143.}, will describe on the surface \(F\) a curve \(C_2\) infinitely near to \(C_0\), in its neighborhood of \(2^\circ\) order, and successive to \(C_1\), to which are to be added the infinitesimal curves, not connected with them, which surround the points of the group \(A\).

To justify the assertion made, that as \(L\) varies in the pencil \(G_2\) generates a curve \(C_2\) near to \(C_1\) and of the same order, one must recognize that the locus of \(G_2\), defined up to irreducible curves \(L\) of the pencil, does not contain any of the points of its base group \(\Gamma\) outside \(A\), and therefore that on \textit{any} \(L\) of our pencil (chosen generically within \(|L|\)) the group of points \(G_2\) (near to \(G\) in the neighborhood of \(2^\circ\) order) cannot be special.

For this purpose we observe that:

1) on a curve \(L\) of genus \(P\), the group \(G_2\) cannot be special unless the groups \(G\) and \(G_1\) are also special; for a group infinitely near to a non-special group is certainly non-special;

2) since the group of \(M\) points \(LC_0\) (with \(A\) removed) does not belong to a linear series of dimension \(>0\), by adding to it a generic group of \(P-M\) points (the group \(S\)) one obtains a non-special group\origfootnote{2}{s180L-2}{Here we have a simple consequence of the Riemann-Roch theorem. Cf. e.g. \textsc{Enriques-Chisini}, \textit{Lezioni}. Bk. V, Chap. I, vol. III.}: to have a special group of \(P\) points one must subject \(S\) to \textit{one} condition, and the choice of \(S\) on \(L\) must be subjected to at least \textit{one other condition} in order that the group \(G=S+LC\) too be special; thus, then, the choice of a group \(S\) of \(P-M\) points of \(L\) which, added to \(LC_0\) and to \(LC_1\), gives rise to two special groups \(G\) and \(G_1\), entails at least \textit{two} conditions which \(S\) must satisfy.

This said, let us take up again the construction outlined previously of a pencil of curves \(L\) in \(|L|\): one will take first of all a particular curve \(\overline{L}\) of \(|L|\) (irreducible, of genus \(P\)), and on it a group of \(r\) points imposed as base points on \(|C|\), and then a group \(S\) of \(P-M\) points to be added to those of the group \(LC_0\) (from which \(A\) has been removed), and finally, on the same \(\overline{L}\), a group \(\Gamma\) of the charac%
% scan: s180R
% !PRINTED 337
\unskip\weberpage{337}%
teristic series \(g_N^{R-r}\) of \(\overline{L}\), containing \(A\). If the group of \(P\) points \(G=S+LC_0\) is to be special on \(\overline{L}\) and is to have, on the same curve, a near group \(G_1=S+LC_1\) which is also special, the choice of the group \(S\) must satisfy, as we have said, at least two conditions; and if, instead, the analogous conditions are to be satisfied on a curve \(L\) of the pencil having \(\Gamma\) as base group, since one has here one parameter at one's disposal, it will nevertheless be necessary for the choice of \(S\), and hence that of \(\Gamma\), to satisfy at least one condition on \(\overline{L}\); and so we shall have a pencil of \(L\) which is not chosen generically within \(|L|\). %

In this way the assertion is proved: the group \(G_2\) constructed, as above, on the \(L\) of a generic pencil of \(|L|\) describes a curve near to \(C_1\) and of the same order (not passing through any of the base points of \(\Gamma\)). And it is worth pointing out that our reasoning does not break down on account of the circumstance that within a pencil of \(L\) there are, in general, (irreducible) curves of effective genus \(P-1\), endowed with a double point \(X\); for it is legitimate to regard these curves as of virtual genus \(P\) (taking the double point \(X\) as virtually non-existent). This means that a group \(S\) of \((P-1)+1\) points on such an \(L\) must be regarded as non-special if it does not belong to a linear series \(\infty^1\) having a neutral pair \(X_1X_2\) corresponding to \(X\). Thus, in this case too, the speciality of \(G\) and \(G_1\) entails at least two conditions which the choice of the group \(S\) must satisfy.

In the preceding discussion we speak of the curves \(C_1\ C_2\,\ldotp\ldotp\ldotp\ldotp\) infinitely near to \(C_0\) as if they were proper curves. To justify the use of this language, it is convenient to show that the said curves can really be defined by means of a \textit{passage to the limit} from proper curves. This happens first of all for the curve \(C_1\) infinitely near to \(C_0\) in the neighborhood of the first order: for we have seen that \(C_1\) forms part of the reducible curve near to \(C_0\), as limit for \(\lambda=0\), within a pencil \(C_0+D_0+\lambda E\) \((|E|=|C+D|)\) which has as base points the \(m\) points of the group \(C_0D_0\). It is quite true that the curve of this pencil is, for \(\lambda\ne 0\), irreducible and tends to \(C_0+D_0\); but, according to the fundamental lemma established above, for infinitesimal \(\lambda\) this curve acquires \(m\) base points and splits into \(C_1+D_1\): this means that the analytic series \(t\ \infty^1\) of groups of points determined by \(G\) and \(G_1\) on the \(L\) of the said pencil does not depend at all on the choice of \(D_0\) or of the system \(|D|\) which has been associated with \(|C|\).

Now one can construct a proper curve whose variation defines \(C_2\), by considering on a generic \(L\) of the said pencil:

1) the group \(G\) cut out by \(C_0\), to which the points of \(A\) are added;

2) the group \(G'\) consisting (besides \(A\)) of the points of \(C_0+D_0+\lambda E\) near to those of \(G\);

3) the group of \(P\) points \(G''\) which corresponds to \(G'\) in the transformation \((GG')\) of the Jacobian variety of \(L\).

\printersignature{\textsc{Enriques F.} - \textit{Superficie algebriche.}\quad 22}

% scan: s181L
% !PRINTED 338
\weberpage{338}
The locus described by \(G''\), as \(L\) varies in the pencil, serves precisely to define the curve \(C_2\) successive to \(C\) and \(C_1\).

To tell the truth, as long as \(\lambda\ne 0\), the said locus is an irreducible curve which tends to a curve larger than \(C_0\), namely to \(C_0+D_0\); in fact a point of \(C_0+D_0+\lambda E\) near to \(C_0\) is liable, as \(L\) varies in the pencil, to be interchanged with a point near to \(D_0\), so that the curve defined by the variable \(G''\) will be generated by the sum of \(G''\) and of an analogous \(\bar{G}''\) defined with respect to \(D_0\) and near to this \(D_0\), and the sum \(G''+\bar{G}''\) will describe a curve \(E_2\) cutting out on the \(L\) groups of the linear series cut out by \(|E|=|C+D|\), or (a priori) of a multiple of this series, whose limit — for \(\lambda=0\) — is \(C_0+D_0\)\origfootnote{1}{s181L-1}{The justification of the assertion reduces to this: let there be given on a curve \(L\) of genus \(P\) two groups of \(P\) points \(G\) and \(\bar{G}\), and also two others \(G'\) and \(\bar{G}'\) such that \(G+\bar{G}=G'+\bar{G}'\), and let the groups \(G''=G'+|G'-G|\) and \(\bar{G}''=\bar{G}'+|\bar{G}'-\bar{G}|\) be constructed; then \(G''+\bar{G}''=G'+\bar{G}'=G+\bar{G}\).}.

But, on \(L\), the groups \(G\,G'\,G''\) belong to an analytic series \(\infty^1\) of groups of \(P\) points which — for \(\lambda\to 0\) — when they become infinitely near, depends only on \(C_0\) and \(C_1\) and \textit{not on} \(D_0\) or on \(|D|\), and which therefore really serves to define, in the neighborhood of \(2^\circ\) order of \(C_0\), the curve \(C_2\), the irreducible part of \(E_2\) for \(\lambda\to 0\), whose existence we wished to establish.

We have constructed on \(F\) a curve \(C_2\) near to \(C_1\) and successive to it, in the neighborhood of \(2^\circ\) order of \(C_0\), which forms part of a split \(E_2=C_2+D_2\) of the linear system \(|E|=|C+D|\). Similarly, starting from \(C_2\), one will construct a successive curve \(C_3\), near to \(C_0\) in the neighborhood of \(3^\circ\) order, etc. Thus the demonstrative requirement declared above is met; that is, it is established that near to the split \(E_0=C_0+D_0\) there are, successively in all the neighborhoods of finite order, split curves of the system \(|E|=|C+D|\), forming a regular series (assimilable to a linear branch of curve); and hence the fundamental theorem is proved in general.

On a surface of geometric genus \(p_g\geq 0\) and of arithmetic genus \(p_a<p_g\), every regular non-special linear system of irreducible curves \(|C|\) is contained in a continuous system \(\{C\}\) which cuts out on a \(C\) the complete characteristic series and is therefore formed of \(\infty^q\) (\(q=p_g-p_a\)) distinct linear systems. Hence: \textit{the geometric irregularity of a surface always equals the arithmetic irregularity:}
\[
q=p_g-p_a.
\]

And (\S~2) \textit{on a surface of irregularity \(q\) every curve \(C\) with the characters \(n\), \(\pi\) and \(i\) which is arithmetically effective, that is, such that}
\[
p_a+n-\pi+1-i\geq 0,
\]
\textit{belongs to a continuous series \(\infty^q\) of non-equivalent curves.}

% scan: s181R
% !PRINTED 339
\weberpage{339}
\edsection{6. History of the theory of continuous systems.}{6. History of the theory of continuous systems.}{sec.IX.6}

The theory of continuous systems of curves belonging to irregular surfaces has a fundamental significance, not only for the algebro-geometric study of surfaces, but also for the (transcendental) doctrine of the simple integrals (of total differentials) belonging to them, since it serves to connect these two different aspects of the science of surfaces. Therefore, placing ourselves at the general point of view which corresponds to its genesis, we wish to deepen the understanding of the general theory of these continuous systems by offering a historical perspective on it, in which one may also see the significance of the errors through which the theory itself has had to pass: a thing all the more essential if one considers that the development of the theory has not yet arrived at a definitive systematization.

The memoir of \textsc{Castelnuovo-Enriques} of 1900 « Sopra alcune questioni fondamentali nella teoria delle superficie algebriche »\origfootnote{1}{s181R-1}{\textit{Annali di Matematica}, t. VI, s. II, p. 165.}, while on the one hand it fruitfully concluded a series of efforts, on the other posed new problems, such as that of the determination of the surfaces with vanishing plurigenera (and hence of the surfaces of genus \(p_g=0\) and of negative arithmetic genus). To go more deeply into such problems, and in general to advance the knowledge of surfaces, also in relation to the transcendental theory of the French school, the study of continuous systems of non-equivalent curves presented itself as an instrument.

The general consideration of such systems goes back to a Note by \textsc{G. Humbert} of 1893\origfootnote{2}{s181R-2}{\textit{Sur une propriété d’une classe de surfaces algébriques.} (Comptes rendus, 1893).}: the author proves that the possession of such a system of curves entails the existence of simple integrals of the first kind (Picard integrals of total differentials) attached to the surface. \textsc{C. Segre} drew the attention of his colleagues to this Note, and shortly afterwards the subject became one of study for the Italian geometric school. Already in 1896 \textsc{Enriques} made an indirect contribution to it\origfootnote{3}{s181R-3}{\textit{Un’osservazione relativa alla rappresentazione parametrica delle curve algebriche.} Rendic. Circolo mat. di Palermo, 1895.} by showing that « on a surface a rational series of curves is always contained in a linear system ».

In 1899 \textsc{Enriques} began the geometric study of surfaces containing continuous systems \(\{C\}\) of non-equivalent curves\origfootnote{4}{s181R-4}{\textit{Una proprietà delle serie continue di curve appartenenti ad una superficie algebrica regolare.} Rendic. Circolo mat. di Palermo, 1899.}, establishing that they are irregular. The author shows that the system adjoint
% scan: s182L
% !PRINTED 340
\weberpage{340}
to a linear system \(|C|\) contained in \(\{C\}\) cannot cut out on a \(C\) the complete canonical series. When he had communicated the result to \textsc{Castelnuovo} before publication, the latter replied by offering him a second proof of the theorem, in which it is shown that the characteristic series of \(|C|\) is necessarily deficient. And in the Note he published, \textsc{Enriques} gives, after his own, the proof of his Colleague.

Subsequently (in 1901) \textsc{Enriques} himself — aiming to recognize the link between the geometrically defined irregular surfaces (\(p_g-p_a>0\)) and those which possess simple integrals of the first kind — proved that « the surfaces possessing \(q\) simple integrals of the first kind with \(2q\) periods contain continuous systems of non-equivalent curves\origfootnote{1}{s182L-1}{\textit{Sur les surfaces algébriques admettant des intégrales de différentielles totales de première espèce.} Annales de Toulouse (2), 3, p. 77.} ».

On the other hand, \textsc{Castelnuovo} succeeded in proving that « a surface of genus \(p_g=0\) and \(p_a<0\) containing a continuous system of non-equivalent curves possesses in consequence also an irrational pencil of curves » (\S~11). It was a property whose significance depended on possible further progress of the theory, and \textsc{Castelnuovo} confined himself to communicating it to \textsc{Enriques}, who was to make it public later (1904), when he was able to deduce from it that « every surface with \(p_g=0\) and \(p_a<0\) always possesses such an irrational pencil ». %

Meanwhile, in 1902, \textsc{F. Severi}, recently graduated from the University of Turin, came, on the advice of his teacher \textsc{C. Segre}, as assistant to the chair of geometry of the University of Bologna, then held by \textsc{Enriques}, with the aim of entering into closer relations with those who cultivated the theory of surfaces. \textsc{Enriques} at once drew the young geometer's attention to the subject of continuous systems, and in particular suggested to him the study of a particular family of surfaces containing such systems, namely the surfaces which represent the pairs of points of a curve, to which \textsc{Severi} indeed devoted a work published in 1903\origfootnote{2}{s182L-2}{\textit{Sulle superficie che rappresentano le coppie di punti di una curva algebrica.} Atti della R. Accademia di Torino, January 1903.

It should be said that, independently of \textsc{Severi}, \textsc{M. De Franchis} also published, in the meantime, a study \textit{Sulla varietà \(\infty^2\) delle coppie di punti di due curve o di una curva algebrica.} Rendic. Circolo mat. di Palermo, 1903.}. %

The study of the surfaces representing the pairs of points of a curve offered examples apt to illuminate questions of a general order: there, in particular, \textsc{Severi} came upon the problem of the base, which was later to form the object of some of his most im%
% scan: s182R
% !PRINTED 341
\unskip\weberpage{341}%
portant studies. Meanwhile it gave occasion for the introduction of the more extended concept of the characteristic series.

The idea was in germ in the proof which, as we have said, \textsc{Castelnuovo} had communicated to \textsc{Enriques} concerning the deficiency of the characteristic series of a \(|C|\) contained in a \(\{C\}\) of non-equivalent curves, because there one sees the deficient series of \(|C|\) being enlarged by the limit of the characteristic series of an infinitely near linear system. But more precisely, \textsc{Castelnuovo} brought out this concept explicitly with reference to the (hyperelliptic) surfaces which represent the pairs of points of a curve \(L\) of genus \(2\); on such surfaces there is a continuous system of curves \(C\) of genus two, which is of degree two and contains within itself the images of the series of pairs of points of \(L\) with a fixed point: now, a curve \(C\) of this system is isolated, so that there is lacking the linear system which ought to cut out on it the characteristic series, which (since \(p_g=1\) and the canonical curve is of order zero) would be the canonical \(g_2^1\); but on the said \(C\) this same \(g_2^1\) is cut out by the infinitely near curves of \(\{C\}\).

\textsc{Castelnuovo} having communicated this remark to \textsc{Severi}, the latter was led to develop it in general, pointing out first of all that the characteristic series of a \(C\) can be defined, independently of the existence of a larger linear system containing \(C\), as the difference
\[
(C+D)\,C-DC,
\]
where \(|D|\) is any auxiliary linear system (such that \(|C+D|\) cuts out a complete series on \(C\)). Indeed this difference series is independent of the choice of the auxiliary system \(|D|\)\origfootnote{1}{s182R-1}{\textsc{Severi}, \textit{Osservazioni sui sistemi continui di curve appartenenti ad una superficie algebrica.} Atti della R. Accademia di Torino, February, 1904.}. As one sees, we have here the geometric or functional interpretation of the procedure by which \textsc{Enriques} had already defined the virtual degree of the curve \(C\), regarding precisely \(|C|\) as the difference of two linear systems of curves.

The characteristic series of a curve \(C\) having been defined in this way, independently of the existence of a larger linear system \(|C|\), there is no difficulty in recognizing that « if \(C\) belongs in any way to a continuous system \(\{C\}\), the curves of the latter infinitely near to \(C\) cut out on it groups of the characteristic series ». And hence the proof of the theorems of Enriques recalled above was reduced to its simplest expression.

Ever since it had been recognized that the possession of a continuous system of non-equivalent curves entails the irregularity of the sur%
% scan: s183L
% !PRINTED 342
\unskip\weberpage{342}%
face (\(p_a<p_g\)), there naturally arose the demand to invert the theorem, that is, to establish that the possession of such continuous systems constitutes the characteristic property of irregular surfaces. Indeed this demand took on a more precise form: for geometers who only by force of examples had departed from the hypothesis implicitly admitted by \textsc{Noether} concerning the completeness of the characteristic series of complete linear systems were bound to ask themselves whether the characteristic series of a larger continuous system would not instead be complete. But the point was to find a means of discussing the question; for this purpose, in fact, a criterion was needed which would make it possible in some way to count the conditions for the splitting of a curve contained in a linear system. \textsc{Enriques} (1904) found such a criterion in the \textit{principle of degeneration or of splitting}: if an irreducible curve \(C\), varying continuously, comes to split into two components, these must have in common at least one point of connection which is not the limit of a variable double point of \(C\).

From the topological aspect there is here an evident truth: because by virtue of continuity the Riemann surface of \(C\), which is connected, must remain connected in every passage to the limit leading to a split \(C\)\origfootnote{1}{s183L-1}{Only recently has attention been drawn to a remark of \textsc{Noether} which anticipates the said principle. In the memoir \textit{Ueber die reductiblen algebraischen Curven} (Acta Mathematica, 1896), having developed the theory of series of groups of points (cut out by linear systems of curves or surfaces) on a split curve, plane or skew, he points out that for skew curves (!) the theory does not come within that relating to irreducible curves if the components do not have at least one common point.

The principle of splitting can also be established, instead of from a topological point of view, by algebro-geometric means. Cf. e.g., \textsc{Enriques-Chisini}, Book V, chap. III, \S~36, vol. III, p. 405.}.

With the help of the principle of splitting \textsc{Enriques} believed he could prove that « complete continuous systems have complete characteristic series »\origfootnote{2}{s183L-2}{\textit{Sulla proprietà caratteristica delle superficie algebriche irregolari.} Rendic. Acc. di Bologna, December 1904.} and thus that a regular non-special linear system (of characters \(n\) and \(\pi\)) on a surface of irregularity \(q=p_g-p_a\) belongs to a continuous system of dimension \(p_g+n-\pi+1\), composed therefore of \(\infty^q\) non-equivalent linear systems. For this purpose he had recourse to the representation of the surface on a multiple plane, which — for the purposes of the proof itself — is not at all essential, so much so that the reasoning can be repeated in a substantially identical manner on the surface itself, as was observed by \textsc{Severi}\origfootnote{3}{s183L-3}{\textit{Intorno alla costruzione dei sistemi completi non lineari.} Rendic. Circolo mat. di Palermo, 1905.}.

% scan: s183R
% !PRINTED 343
\weberpage{343}
It should be said at once that the reasoning of \textsc{Enriques} contained an error, as is explained in the developments of \S~5: the author had not noticed the critical point encountered here, and passed, without noticing it, from a differential property to an integral property. The difficulty, that is, the insufficiency of the principle of splitting for \(p_g>0\), has also been illustrated in the said \S~5 by pointing out that the true meaning of this principle consists in this: that « the splitting of the curve of a linear system \(|C+D|\) into two curves of the order of \(C\) and of \(D\) with \(m\) common points entails (not more than) \(m\) conditions »; whereas for \(p_g>0\) one must recognize that these conditions are, in general, only \(m-p_g\). To prove that this happens for the curves infinitely near to a split curve \(C+D\) of our linear system means precisely to establish a differential property which one is not authorized to take in an integral sense. The deduction turns out to be correct only for \(p_g=0\) (and \(p_a<0\)), in which case indeed there is no need at all to resort to the consideration of the characteristic series.

But the error, or the gap, in the reasoning of \textsc{Enriques} was not noticed by geometers for fifteen years, so much so that one must come to 1921 to find the criticism formulated by \textsc{Severi}; and by that time the fundamental theorem, concerning the existence of series \(\infty^q\) of non-equivalent curves, had been proved in another way, by a transcendental procedure, by \textsc{Poincaré} (1910)\origfootnote{1}{s183R-1}{\textsc{H. Poincaré}, \textit{Sur les courbes tracées sur les surfaces algébriques.} Annales de l’Ecole Normale Sup., p. 3 (27).}. In 1904 the reasoning of \textsc{Enriques} seemed to be confirmed also by the exposition of it given by \textsc{Severi} himself, who preserved — as we have mentioned — its essential scheme\origfootnote{2}{s183R-2}{Cf. the memoir cited: \textit{Intorno alla costruzione dei sistemi completi non lineari che appartengono ad una superficie irregolare.} Rendic. del Circolo mat. di Palermo, 1905.

Later, in the Note \textit{Nuovi contributi alla teoria dei sistemi continui di curve appartenenti ad una superficie algebrica} (Rendic. Lincei, 2 and 16 April 1916), in no. 1, \textsc{Severi} set out again the proof of \textsc{Enriques}, indeed reverting to its original form (representation of the surface on the multiple plane).}.

And therefore the theorem, thus accepted by geometers, became the starting point of important developments.

These proceed essentially in two directions:

1) \textsc{Enriques} applies himself first of all to the classification of surfaces of genus \(p_g=0\), and later to the more general problems of the classification of surfaces, which will constitute the subject of the two following chapters X and XI. Fortunately these developments are
% scan: s184L
% !PRINTED 344
\weberpage{344}
based, or can be based, for the most part on the restricted fundamental theorem.

This is clear a priori for Chap. X, relating to surfaces of genus \(p_g=0\), but it holds also for the results on surfaces of arithmetic genus \(p_a<0\) which we set out in Chap. XI, where, however, the proofs must be modified at some points. The exceptions are the theorems concerning surfaces with \(p_a=0\) and \(p_g=1\): that these surfaces cannot have the canonical curve of order zero, and that they always possess a pencil of elliptic curves. %

2) \textsc{Castelnuovo} and \textsc{Severi}, on the other hand, turned to developing the consequences of the fundamental theorem with regard to the transcendental theory: where it was a matter of making precise, from the quantitative point of view, the link which the said theorem brought to light between irregular surfaces and those possessing simple integrals of the first kind. In this connection we have already mentioned the first attempt made by us to invert the remark of \textsc{Humbert}, namely the theorem that surfaces possessing \(q\) simple integrals of the first kind with \(2q\) periods contain continuous systems of non-equivalent curves and are therefore irregular. This theorem would take on the widest significance if Picard's transcendental theory allowed one to know a priori that the number of the said integrals of the first kind attached to an arbitrary surface is always equal to half the number of the periods.

In the absence of this, already before the Note of \textsc{Enriques} « Sulla proprietà caratteristica.... », \textsc{Severi} had investigated in another way the question of surfaces possessing simple integrals of the first kind, proving that « the possession of integrals of the second kind (and therefore of the first) entails that the characteristic series of a complete linear system is deficient, and hence also that the surface is irregular »\origfootnote{1}{s184L-1}{\textit{Sulle superficie algebriche che posseggono integrali di Picard della 2\textsuperscript{a} specie.} Rendic. Lincei, 1904. Cf. Math. Annalen, 1905.}. Although this proof serves to establish a fundamental fact and contains an interesting analysis relating to the singularities on the polar curve of the integral, it does not seem that it can be completed so as to find the number of integrals of the first kind as a function of the irregularity without the addition of some essential element\origfootnote{2}{s184L-2}{The analysis of the question leads \textsc{Severi} to a result \textit{Sulla-differenza fra i numeri degli integrali di Picard della prima e della seconda specie appartenenti ad una superficie irregolare.} %
Atti R. Accad. Torino, 1905. The same result was reached, by his own procedure, by \textsc{Picard} (Comptes, rendus, 1905): Cf. the \textit{Traité} of \textsc{Picard} and \textsc{Simart}, II, p. 417.}. %

On the other hand, the consideration of the continuous systems of \(\infty^q\) non-equiv%
% scan: s184R
% !PRINTED 345
\unskip\weberpage{345}%
alent curves allows \textsc{Castelnuovo}\origfootnote{1}{s184R-1}{\textit{Sugli integrali semplici appartenenti ad una superficie irregolare.} Rendic. Lincei, May-June, 1905.} and \textsc{Severi}\origfootnote{2}{s184R-2}{\textit{Il teorema d’Abel sulle superficie algebriche.} Annali di Mat., August, 1905.}, independently of each other, to make the remark of \textsc{Humbert} precise from the quantitative aspect, by proving that « irregular surfaces, of irregularity \(q=p_g-p_a\), possess exactly \(q\) simple integrals of the first kind, with \(2q\) periods ». The proof of this capital theorem, which the consideration of the linear equation of Fuchs used by \textsc{Picard} does not succeed in reaching, will later be seen carried out in another way by \textsc{Poincaré}, in the memoir cited, in which he also finds again, as we have said, the existence of continuous systems with \(\infty^q\) non-equivalent curves.

For the proof of the theorem on integrals mentioned above, \textsc{Castelnuovo} introduces the important concept of what he called the \textit{Picard variety} attached to an irregular surface (cf. \S~3). This notion makes it possible in general to construct the continuous system defined by a curve \(C\), starting from any other continuous system \(\{D\}\) belonging to the surface. \textsc{Enriques} observed that the construction can be made precise on the basis of the Riemann-Roch theorem, deducing from it that « a linear system of curves of characters \(n\), \(\pi\) and \(i\), \textit{arithmetically effective}, that is, such that \(p_a+n-\pi+1-i\geq 0\), always belongs to a series \(\infty^q\) of non-equivalent linear systems » (cf. \S~2). The remark was communicated by the author to \textsc{Severi}, and since the latter had just then succeeded in simplifying considerably the proof of the Riemann-Roch theorem, \textsc{Enriques} left it to him to publish it\origfootnote{3}{s184R-3}{\textsc{Severi}, \textit{Sul teorema di Riemann-Roch e sulle serie continue di curve appartenenti ad una superficie algebrica.} R. Accad. Torino, 1905.}.

Let us now come to the criticism which discovered the weak point in the proof of Enriques's theorem on continuous systems. We have said that it was made by \textsc{Severi} in 1921. It is to be found in a series of Notes « Sulla teoria degli integrali semplici di prima specie appartenenti ad una superficie algebrica »\origfootnote{4}{s184R-4}{Rendic. Lincei, 1921.}, in which there is set out, with simplifications at some points, the proof of \textsc{Poincaré} of 1910, which comprises together the existence of the continuous systems and that of the integrals of the first kind\origfootnote{5}{s184R-5}{The proof of \textsc{Poincaré}, further simplified by \textsc{Lefschetz}, is set out by \textsc{O. Zariski}, \textit{Algebraic Surfaces} (Berlin, Springer, 1935). Cf. VII, 4, 5 (pp. 128, 133).}. The author observes that (thanks to \textsc{Poincaré}) at least the essential part of the theorem (the fundamental theorem as it is stated here in \S~5) remains safe. Later there will be renewed attempts to extend this theorem by proving
% scan: s185L
% !PRINTED 346
\weberpage{346}
that the characteristic series of any complete continuous system whatever is always complete; but of this later (\S~7).

\textsc{Enriques}, in several Notes of 1936-1937, on « Curve infinitamente vicine sopra una superficie algebrica »\origfootnote{1}{s185L-1}{Rendic. Lincei, April 1936 and Rendic. Seminario Matematico della R. Università di Roma, 1936.} and on « La proprietà caratteristica delle superficie algebriche irregolari e le curve infinitamente vicine »\origfootnote{2}{s185L-2}{Rendic. Lincei, October 1937.}, attempted to fill the gap pointed out in the proof of the fundamental theorem, trying to prove that to a curve \(C\) (when the characteristic series of \(C\) has the deficiency \(q=p_g-p_a\)) there are infinitely near not only \(\infty^q\) non-equivalent curves in the neighborhood of the first order, but also curves successive to these in the neighborhood of the second order, etc. The fundamental idea of this proof (which already appears in the « Lezioni » of \textsc{Enriques} and \textsc{Campedelli}) consists in this: that « on an Abelian variety one defines, with \textsc{Lie}, an infinitesimal transformation generating a group \(\infty^1\), starting from two infinitely near points ».

Thus, in particular, in the Jacobian variety constituted by the groups of \(p\) points of a curve of genus \(p\), one succeeds in defining an analytic series \(\infty^1\) of elements (groups of points) starting from two infinitely near groups. This construction extends to continuous systems of non-equivalent curves on a surface, when one presupposes the existence of a system whose elements form the (Abelian) Picard variety. The considerations developed by \textsc{Enriques} aim at justifying such a construction, passing from infinitely near curves in the neighborhood of the first order to curves in the neighborhood of the second order and so on, independently of the said presupposition; indeed, with the aim of thereby arriving at a justification of it. The starting point is thus that, near to a split curve \(C+D\), there are, in the neighborhood of the first order, curves \(\bar{C}+\bar{D}\) split into two components, one of the order of \(C\) and another of the order of \(D\), having the same number \(m\) of common points. Here, the author points out, it is first of all necessary to justify the implicit application which is made of the principle of splitting in the infinitesimal; but this offers no difficulty, inasmuch as one can recognize that, going outside the surface, one can construct, in space, a series of split curves which contains \(C+D\) and \(\bar{C}+\bar{D}\). %

The \textit{reality} of the splitting for infinitely near curves having thus been justified, and hence the effective existence of a non-equivalent \(\bar{C}\) infinitely near to \(C\), the construction of \textsc{Enriques} yields a series of operations on \(C+D\) and \(\bar{C}+\bar{D}\) which comes to depend only on \(C\) and \(\bar{C}\), and leads precisely to prolonging into the successive neighborhoods of \(C\) the series of curves defined by \(\bar{C}\).

% scan: s185R
% !PRINTED 347
\weberpage{347}
On the occasion of the Note of 1937, this proof was contested by \textsc{B. Segre}.

Nevertheless \textsc{Enriques} (while rebutting the objections of \textsc{Segre}) willingly granted preference to another proof based simply on the principle of continuity, which was proposed by \textsc{Segre} himself, while reserving the discussion of whether it was well founded: of this interesting attempt and of the developments it has received at the hands of other geometers we speak expressly in the following \S~7.

Only after the failure of the aforesaid attempt have we taken up again, as is natural, the criticized proof, which — with some refinements and clarifications — has been set out in the preceding \S~5.

\edsection{7. On various attempts to prove and to extend the fundamental\\theorem.}{7. On various attempts to prove and to extend the fundamental theorem.}{sec.IX.7}

The proof of the fundamental theorem, which expresses the characteristic property of irregular surfaces, has been attained in general (for \(p_g\geq 0\)) by means of a consideration of infinitesimals which takes the place of a differential calculation not carried out. Considerations of this nature, even while remaining within the field of algebraic analysis, are of a delicate order and arouse the distrust of mathematicians who are not sufficiently familiar with them, or at least leave the impression that some intuitive element has been introduced, in contrast with the demands of logical rigor.

This impression was all the more understandable in view of some rather rapid or less justified step in the Notes cited, whose exposition\origfootnote{1}{s185R-1}{In particular following the path indicated in the Notes of 1936.} has been taken up again here, in an effort to make precise and to clarify the most delicate points. For these reasons it appeared, and may still appear, desirable to furnish other proofs of the same theorem which would compel the assent of all, without doubts or hesitations.

We have already said that a proof of this kind, using the integrals belonging to the surface, was given by \textsc{poincaré} (1910), and today it can be seen, for example, in the simplified exposition of \textsc{Lefschetz} reported by \textsc{Zariski}. But, within the sphere of the algebro-geometric theory, special mention is deserved by the attempt which, as we have mentioned, was undertaken by \textsc{B. Segre}\origfootnote{2}{s185R-2}{\textit{Un teorema fondamentale della geometria sulle superficie algebriche e il principio di spezzamento.} Annali di Mat., 1938.}. %

The guiding idea of this author can be expressed as follows. Consider the linear system \(|E|=|C+D|\) which contains the split curve \(E_0=C_0+D_0\), endowed with \(m\) double points which form the group \(C_0D_0\); and impose on the curves \(E\), near to \(E_0\), \(m-p_g\)

% scan: s186L
% !PRINTED 348
\weberpage{348}
double points (near $m-p_g$ points of the said group); one wishes to prove that the curves $\bar{E}$ so defined necessarily split into a curve of the order of $C_0$ and a curve of the order of $D_0$: and for this purpose one should reduce to absurdity the hypothesis that the said $\bar{E}$ are irreducible. Accordingly, on an $\bar{E}$ with $m-p_g$ assigned double points, construct the characteristic series $g$ of the continuous system $\{\bar{E}\}$, made complete, and pass to the limit as $\bar{E}\to E_0$; the series $g$ has dimension
\[
2p_g+(n-\pi)+(\nu-\varrho)+1, %
\]
but, in the limit, it appears to reduce to the series $g_c+g_d$ formed with the groups of the characteristic series of $C_0$ and of $D_0$ (made complete), to which are added $p_g$ fixed points of the group $C_0D_0$ (the new double points acquired by the reducible $E_0$); now, the dimension of the limit series $g_c+g_d$ is
\[
2p_g+(n-\pi)+(\nu-\varrho), %
\]
one unit less than that of $g$!

The author asserted that this decrease in the dimension of a series is absurd, since it conflicts with the principle of continuity. But \textsc{Enriques}, having made an analogous attempt on his own (in the simplest case $p_g=1$), had been brought to a halt by an embarrassing paradox, and accordingly set out the difficulty he had encountered\origfootnote{1}{s186L-1}{\textsc{Enriques}, \textit{Sulla proprietà caratteristica delle superficie algebriche irregolari.} Rendic. Lincei, June, 1938.}. Consider, in the plane, the system of curves $C_8$ of order 8 with 15 double points; among these $C_8$ there are the curves split into two quartics $C_4+\bar{C}_4$, which have in common 15 points $A_1\ldotp\ldotp\ldotp\ldotp A_{15}$, limits of the double points of $C_8$, and another, sixteenth, point $O$. Now on a variable $C_8$ the canonical series $g=g_{10}^5$ is cut out by the quintics $C_5$ through the 15 double points; in the limit as $C_8\to C_4+\bar{C}_4$ the said canonical series reduces to the series composed of the groups of the two canonical series $g_4^2$ on $C_4$ and on $\bar{C}_4$, which is cut out by the $C_5$ through $A_1\ldotp\ldotp\ldotp\ldotp A_{15}$, which consequently have the common fixed point $O$; thus it seems that in the limit the aforesaid series $g$ reduces to a series of lower dimension. The doubt arises that the conservation of dimension, which seemed to follow from the principle of continuity, breaks down!

It is to \textsc{Segre}'s credit that he cleared up the difficulty by explaining the paradox. The principle of the conservation of the dimension of a linear series does indeed follow from a requirement of continuity which, once stated explicitly, proves irrefutable; but when one seeks the limit of the series $g$, as $C_8\to C_4+\bar{C}_4$, one must consider not only what the said $g$ becomes while a $C_5$ adjoint to $C_8$ tends to a generic $C_5$ through $A_1\ldotp\ldotp\ldotp\ldotp A_{15}$, but also the limit of
% scan: s186R
% !PRINTED 349
\weberpage{349}
this series corresponding to a variation by which $C_8$ and $C_5$ tend together to two degenerate curves containing $C_4$ (or $\bar{C}_4$) as a component, where the intersection groups of $C_8$ and $C_5$ become indeterminate. In other words: the paradoxical example brings to light the necessity of taking into account, in passing to the limit of the series $g$, the \textit{entire limit}!

B. \textsc{Segre} explains the import of this requirement well by the example of a hyperbola which, varying, reduces in the limit to the pair of its asymptotes: the $g_2^1$ cut out on the hyperbola by the lines through the center reduces to the fixed point which is the center, but also to the two symmetries with respect to the center itself which arise on the asymptotes. Similarly, in \textsc{Enriques}'s example, taking into account that the curves $C_8$ passing through $A_1\ldots A_{15}$ on $C_4$ do not consequently pass through $O$, the limit of the $g=g_{10}^5$ is given by two series of dimension 5, which are obtained by associating groups of $C_4$ and of $\bar{C}_4$: to speak of only one of these series, it is constructed by associating with the canonical groups of $\bar{C}_4$ the groups of a $G_6^3$ on $C_4$, which no longer has the fixed point $O$, but is defined by setting $g_6^3=g_4^2+2O$. %

Having thus explained the embarrassing paradox, B. \textsc{Segre} recognized a gap in his own proof of the fundamental theorem, and set about filling it by seeking a criterion capable of determining the limit of the intersection of two curves which, varying on a surface, come to acquire a common component. He alluded to this infinitesimal criterion in a Note ``Sur un théorème fondamental de géométrie sur les surfaces algébriques'' of 9 January 1939\origfootnote{1}{s186R-1}{\textit{Comptes-rendus}, 9 Janvier 1939.}, which was accompanied by another Note of \textsc{Enriques}\origfootnote{2}{s186R-2}{\textit{Sur la propriété caractéristique des surfaces algébriques irrégulières.} Comptes-rendus, 3 Janvier 1939.} which greeted with pleasure the success of his Colleague.

Meanwhile B. \textsc{Segre} left Italy, and it does not appear that any more extensive publication of his on the subject has come out up to now.

Subsequently \textsc{Severi}, in a memoir of the Accademia d’Italia\origfootnote{3}{s186R-3}{\textit{La teoria generale dei sistemi continui di curve sopra una superficie algebrica.} Memorie dell’Accademia d’Italia, vol. XIX, 1941}, offers a critical revision of the entire theory of continuous systems of curves on algebraic surfaces. Unfortunately this exposition is marred by some errors which, in part, have been pointed out by \textsc{Enriques} in a Note ``Sui sistemi continui di curve appartenenti ad una superficie algebrica'', and which must be re-examined here. %

To satisfy the demonstrative requirement relating to \textsc{Segre}'s reasoning, as explained above, \textsc{Severi} takes over from him — in a somewhat modified form — a
% scan: s187L
% !PRINTED 350
\weberpage{350}
lemma of differential geometry, from which it should follow that the characteristic series $g$ of the continuous system $\{\bar{E}\}$ on $\bar{E}$ (made complete), on passing to the limit as $\bar{E}\to E_0=C_0+D_0$, would necessarily have as fixed points the $p_g$ points $P$ of the group $C_0D_0$ which are not limits of the $m-p_g$ double points of $\bar{E}$, and this in whatever way the passage to the limit is made, even when the curves cutting out $g$ on $\bar{E}$ tend to contain $C_0$ as a component, so that the $\mathit{lim}\,g$ turns out to be indeterminate. The Segre-Severi lemma, which would allow this conclusion to be deduced, expresses a property of differential geometry according to which the said points $P$ would be limits of points belonging to groups of the series $G$. This differential property would in its turn descend from the one, recognized by \textsc{Enriques} in the original proof of the fundamental theorem, that ``the curves infinitely near to $C_0$ cut out on it the complete characteristic series''; in this way the gap in the proof of that theorem would be filled.

Now \textsc{Enriques}, who had already found some difficulty in reconstructing \textsc{Segre}'s infinitesimal criterion from the indications supplied to him by the author, did not succeed in understanding this step, and therefore declared the said lemma ``obscure and unconvincing''\origfootnote{1}{s187L-1}{Commentarii Mathematici Helvetici: received by the editors on 20 March 1942.}. But he believed that the aim could nevertheless be achieved by having the series $\bar{g}$ on the variable $\bar{E}$ cut out by curves that remain irreducible in the passage to the limit in question.

In this he was mistaken, and \textsc{Severi} has rightly exposed his error\origfootnote{2}{s187L-2}{\textit{Intorno ai sistemi continui di curve sopra una superficie algebrica.} Commentarii Mathematici Helvetici (30 November 1942). The editors of the Commentarii did not allow \textsc{Enriques} to reply to this article, and therefore he could not at that time either acknowledge his own error or point out the errors into which, for his part, his critic had fallen.}. However, the recognition of this error casts doubt on the whole proof proposed by \textsc{Segre} and \textsc{Severi}. It is not merely a matter of correcting the lemma criticized in \textsc{Severi}'s Memoir, as the author himself has done\origfootnote{3}{s187L-3}{Note presented to the R. Acc. d’Italia, 16 January 1942.} following an example of \textsc{Bompiani}; but rather of asking whether a differential property such as is expressed in the aforesaid lemma really suffices to exclude that the limit of $g$ may give rise on $C_0$ to a series without fixed points. Whoever analyzes the question more deeply will easily be persuaded that this doubt cannot be removed, and indeed that there is here a radical insufficiency of the lemma, namely that ``this Segre-Severi lemma does not answer the purpose for which it is introduced''. We shall not dwell on this, having reason to believe that \textsc{Severi} himself has by now recognized the insufficiency of his
% scan: s187R
% !PRINTED 351
\weberpage{351}
reasoning. At any rate it is worth having mentioned here an idea which — if not today — might tomorrow find a fruitful application.

We have mentioned that \textsc{Severi} (and before him B. \textsc{Segre}) attempted to extend the fundamental theorem on irregular surfaces, establishing that, on any surface whatever, ``the characteristic series of an irreducible continuous system without multiple points (however special or superabundant) is always complete''.

However, the author has now recognized, through examples, that the more extended result which he had in view admits some exceptions\origfootnote{1}{s187R-1}{And \textit{the uniqueness} of the continuous system to which an irreducible curve belongs, which the author had believed he had proved by the considerations developed at the end of p. 378 and at the beginning of p. 379 of the memoir \textit{La teoria generale etc.}, of 1941, also admits exceptions.}.

But let the choice of the continuous systems of irreducible curves $\{C\}$ be restricted as far as necessary, so that the property holds \textit{by hypothesis}: let $\{C\}$, then, be a continuous system $\infty^{d+1}$ which has on a $C$ the complete characteristic series of dimension $d$; it is worth specifying various circumstances that may arise.

The system $\{C\}$ will contain an infinity of linear systems of non-equivalent curves if $d+1$ exceeds the dimension of the complete linear system $|C|$, that is, if this system, of dimension
\[
r=p_a+n-\pi+1-i+\omega \qquad (\text{with }\omega\geq0)
\]
has a deficient characteristic series, of deficiency
\[
q-\varepsilon \qquad (q=p_g-p_a).
\]

Now it is useful to distinguish the abnormal \textit{elementary cases} which, by their combinations, give rise to all the possible cases. A priori it may happen:

1) that the system $|C|$ is regular and special, with index of speciality $i$:
\[
0<i<p_g,
\]
in the preceding notation:
\[
\omega=\varepsilon=0;
\]

2) that $|C|$ is non-special and superabundant, with superabundance $\omega>0$, while its characteristic series has the maximum deficiency $q=p_g-p_a$: $i=0$, $\varepsilon=0$.

3) that $|C|$ is non-special and superabundant, with superabundance $\omega=\varepsilon$, because it has a characteristic series of deficiency $q-\varepsilon$: $i=0$, $\varepsilon>0$.

% scan: s188L
% !PRINTED 352
\weberpage{352}
In \textit{case} 1) (if it is possible) the dimension of $|C|$ is
\[
r=p_a+n-\pi+1-i,
\]
and that of $\{C\}$ is $r+q$; hence \textit{the continuous system is formed of $\infty^q$ linear systems} (in general) \textit{with the same index of speciality $i$.}

In \textit{case} 2), the dimension of $|C|$ being $r=p_a+n-\pi+1+\omega$, that of $\{C\}$ is again $r+q$; hence $\{C\}$ \textit{is formed of $\infty^q$ linear systems with the same superabundance.} %

But in \textit{case} 3), the dimension of $|C|$ being $r=p_a+n-\pi+1+\varepsilon$, that of $\{C\}$ is $r+q-\varepsilon$, and one cannot say a priori whether, for example, $\{C\}$ \textit{is formed of $\infty^q$ regular linear systems having in general the dimension $r-\varepsilon$, among which is found the superabundant $|C|$ of dimension $r$, or of a series $\infty^{q-\varepsilon}$ of linear systems all with the same superabundance $\omega=\varepsilon$.} %

We have thus named only the two extreme cases among those which, under hypothesis 3), appear a priori possible. The actual possibility that these various hypotheses occur, and in particular the first (which a deeply hidden error had led to denying), has been recognized in the cited critical Note of \textsc{Enriques}\origfootnote{1}{s188L-1}{Commentarii Mathematici Helvetici, 20 March 1942.}, where the study of the example which is briefly sketched here is carried out in depth.

Consider the ruled surface $F_5$ of the $5^\circ$ order which has a triple line and a double line: $F_5$ has geometric genus $p_g=0$, arithmetic genus $p_a=-p=-2$ and irregularity $q=2$.

On $F_5$ the system $|C|$ of plane sections is superabundant, with superabundance $\varepsilon=1$, having degree $n=5$, genus $\pi=2$, and dimension $r=p_a+n-\pi+1+\varepsilon=3$, while the deficiency of the characteristic series is $q-\varepsilon=1$.

The linear system $|C|$ is contained in a continuous system $\{C\}$ of dimension 4, and one can see that this is formed not of $\infty^1$ linear systems $\infty^3$, but of $\infty^2$ regular linear systems $\infty^2$, among which is found the said $|C|$ of dimension 3. To verify this assertion one adds to the system $|C|$ a pair of generators $a\ b$, and by means of the linear system so obtained (which is $\infty^6$) one constructs a surface $F_9$ of order 9, a normal ruled surface in the space $S_6$, of which $F_5$ appears as the projection from a pair of incident generators $a'\ b'$. The surface $F_9$ has a double point $O$ at the point of intersection of $a'$ and $b'$, but does not possess a double curve; therefore it does not have infinitely many pairs of incident generators, and so the $\infty^2$ linear systems constituting $\{C\}$, that is, the systems cut out on $F_9$ by the hyperplanes through a pair of generators, are, in general, of dimension 2, \hfill Q.E.D.

% scan: s188R
% !PRINTED 353
\weberpage{353}
On the other hand one can point out on $F_9$ a continuous system $\{L\}$ formed of $\infty^1$ superabundant linear systems, all of dimension $3\ (=2+1)$: such is the continuous system cut out by the hyperplanes passing through the double point $O$ and through a generator of the ruled surface. And it should be noted that this example must be regarded as falling under case 3), and not under case 2), because the characteristic series of $\{L\}$ is a non-special series $g_5^3$.

\textit{Remark.} -- The $\infty^1$ linear systems $|L|$ forming $\{L\}$ are — on the ruled surface $F_5$ — the residuals of the generators with respect to the linear system $\infty^5$
\[
|C+a+b-A-B|,
\]
which has two base points, $A$ and $B$, on $a$ and $b$ respectively, which constitute a neutral pair for the system
\[
|C+a+b|.
\]

Now, if $a$ and $b$ are changed into two other generators $\bar{a}$ and $\bar{b}$, the pair $A\ B$ will no longer be neutral for $|C+\bar{a}+\bar{b}|$, and hence the subtraction of a generator from these linear systems will lead to $\infty^2$ linear systems $|\bar{L}|=|C+\bar{a}+\bar{b}-A-B|$.

The two systems, $\{L\}$ formed of $\infty^1$ linear systems and $\{\bar{L}\}$ formed of $\infty^2$ linear systems, are distinct, that is, $\{L\}+\{\bar{L}\}$ is reducible into two varieties $\infty^4$ as a system of curves, and must instead be regarded as irreducible as a variety $\infty^2$ of linear systems $|\bar{L}|$ which contains within itself the variety $\infty^1$ of the $|L|$. This circumstance is expressed by saying that the said system $\{L\}$ is \textit{exorbitant}. The phenomenon of exorbitance (of a linear system or of a continuous series of linear systems) was observed by \textsc{Rosenblatt} and \textsc{Albanese} and then by \textsc{Severi}\origfootnote{1}{s188R-1}{\textit{Nuovi contributi etc.} Lincei, 1916.}.

Other examples of exorbitant systems are given by the system of plane sections of a quite general cone of order $n\geq3$.

For $n=3$ the linear system $\infty^3$ of plane sections is contained in the series of the $\infty^1$ linear systems ($\infty^2$) formed by the triples of generators of the cone, to which the neighborhood of the vertex is to be added.

For $n=4$ the linear system $\infty^3$ of plane sections of the cone is contained in the series of the $\infty^3$ linear systems ($\infty^1$) formed by the quadruples of generators of the cone, to which the neighborhood of the vertex is to be added.

For $n\geq5$ the linear system $\infty^3$ of plane sections of the cone is contained in the series of the $\infty^n$ linear systems ($\infty^0$) formed by the groups of $n$ generators of the cone, to which the neighborhood of the vertex is always added.

\printersignature{\textsc{Enriques F.} - \textit{Superficie algebriche.}\quad 23}

% scan: s189L
% !PRINTED 354
\weberpage{354}
\edsection{8. The paracanonical system.}{8. The paracanonical system.}{sec.IX.8}

According to the fundamental theorem (\S~4) and the numerical condition of \S~2, the canonical system $|K|$ of a surface $f$, of arithmetic genus $p_a\geq0$ and geometric genus $p_g>p_a$, will always be contained in a continuous system $\infty^q$ ($q=p_g-p_a$) of distinct linear systems, of dimension
\[
r\geq p_a+p^{(1)}-1-p^{(1)}+1=p_a
\]

It was believed that the dimension of this \textit{paracanonical system} $\{K\}$ could be determined precisely by determining the deficiency of the characteristic series of the canonical system $|K|$, at least under the hypothesis that \textit{the pure canonical curves are irreducible}, and without multiple points, of genus $p^{(1)}>1$\origfootnote{1}{s189L-1}{\textsc{Enriques}, \textit{Des courbes paracanoniques appartenant à une surface algébrique irrégulière.} Soc. Royale des Sciences de Liège (19 October, 1939).}.

The said deficiency would be $q=p_g-p_a$, and hence the paracanonical system (granting that it has on a $K$ the complete characteristic series) would have dimension $2p_g-p_a-1$. But the \textit{deduction} is \textit{erroneous}\origfootnote{2}{s189L-2}{It seems that the proof of the same theorem offered by \textsc{Severi} in the cited memoir \textit{La teoria generale dei sistemi continui etc}, of 1941, also falls substantially into the same error.}, as is recognized by taking up the author's reasoning again and bringing to light an objection of a delicate nature which invalidates it.

To this end take on the surface $F$ a regular and non-special linear system $|C|$, of degree $n$ and genus $\pi$, whose adjoint $|C'| = |C+K|$ will have dimension %
\[
p_a+\pi-1
\]
and will cut out on a $C$ a deficient canonical series, of dimension %
\[
\pi-1-q.
\]

Denoting by $G=CK$ the group of $2\pi-2-n$ points common to a $C$ and a $K$, one wishes to compute the number of conditions which $G$ imposes on the curves of $|C'|$ that are to contain it. And for this one assumes that the $C'$ containing $G$ cut out on $C$ the characteristic series of $|C|$ made complete, that is, the characteristic series $g$ of $\{C\}$, cut out on the said $C$ by the $\bar{C}$ (in general non-equivalent) infinitely near to it. One would like to justify the assumption by saying that, among the residual curves of a $\bar{C}$ with respect to $|C'|$, there would be at least one curve $\bar{K}$ near to $K$, so that $C+K$ and $\bar{C}+\bar{K}$ would determine a
% scan: s189R
% !PRINTED 355
\weberpage{355}
pencil of curves $C'$ through $G$, passing through an arbitrarily given group of $g$.

But precisely here lies the error! If the paracanonical linear systems have a dimension lower than that, $p_g-1$, of the canonical system $|K|$, a paracanonical system $|\bar{K}|$ near to $K$ does not possess a curve $\bar{K}$ near to an arbitrary $K$: the limit of a $|\bar{K}|$ is a smaller linear system, contained in the latter.

We have explained the error that invalidates the preceding reasoning, from which it would follow that the deficiency of the canonical series of $|K|$ is $q=p_g-p_a$. Now one can see from actual examples that the conclusion itself is far from being verified. Indeed it seems \textit{probable} that ``on a surface of irregularity $q=p_g-p_a>1$, ($p_a\geq0$) with irreducible canonical curves $K$, without multiple points, \textit{the $\infty^q$ paracanonical linear systems are, in general, regular} and hence the continuous system $\{K\}$ is $\infty^{p_g}$: the characteristic \textit{series} of $|K|$ (at least if this system is not exorbitant) would therefore have \textit{deficiency} 1''.

Let us see how this is borne out for the surface representing the pairs of points of a quartic $L$ of genus $q=3$.

This surface $F$ has geometric genus $p_g=3$ and arithmetic genus $p_a=0$\origfootnote{1}{s189R-1}{Cf. F. \textsc{Severi}, \textit{Sulle superficie che rappresentano le coppie di punti di una curva algebrica.} Atti R. Acc. di Torino, 25 January 1903.}: the canonical curves of $F$ are given by the series of pairs of points of $L$ extracted from the quadruples of the special $g_4^1$, and are therefore $\infty^2$; the paracanonical curves, on the other hand, are given by the series of pairs of points extracted from the quadruples of the non-special $g_4^1$, and therefore form a continuous system $\infty^3$, consisting of $\infty^3$ non-equivalent curves, within which lies the superabundant system $|K|$.

One can easily verify that the paracanonical systems $|K|$ are, in general, of dimension $p_a=0$ and not $>0$; otherwise the characteristic series $g_6^x$ of $\{K\}$ on a $K$ (of genus $p^{(1)}=7$) would have dimension
\[
x\geq3,
\]
whereas a curve of genus greater than 4, not hyperelliptic, such as $K$ is, cannot contain a $g_6^3$.

What has been proved for the surface representing the pairs of points of a quartic $L$ of genus $q=3$ seems to extend to the surfaces $F$ representing the pairs of points of a curve $L$ of genus $q=4,\,5\ldots$.

For $q=4$, $L$ being realized by a sextic of $S_3$, the canonical curves $K$ of $F$ will be given by the ruled surfaces cut out on the congruence of the chords of $L$ by the linear complexes of $S_3$, and therefore one will have
\[
p_g=6;
\]
% scan: s190L
% !PRINTED 356
\weberpage{356}
on the other hand there will be paracanonical curves $\bar{K}$ of $F$ corresponding to the series of pairs of points of $L$ extracted from its non-special $g_6^1$. If, as we believe, it is verified that to the $g_6^1$ of a non-special $g_6^2$ there correspond the curves $\bar{K}$ of a complete linear system $|\bar{K}|$, the paracanonical continuous system $\{\bar{K}\}$ of $F$ will have dimension
\[
p_g=p_g-p_a+p_a=4+2=6
\]
containing within itself $\infty^4$ regular linear systems of dimension
\[
p_a=2.
\]

However, the canonical system $\infty^5$ will not be contained in this total paracanonical system $\{\bar{K}\}$, but must be \textit{exorbitant}, because the ruled surface of the chords of $L$ which belong to a special linear complex cannot be the limit of a ruled surface whose generators join the pairs of a (non-special) $g_6^1$ of $L$ itself.

Also for $q=5$, $L$ being realized in $S_4$ by a curve of order 8 (base of a net of quadrics), it seems that the surface $F$, representing the pairs of points of $L$, possesses a paracanonical system $\{K\}$ formed of $\infty^q=\infty^5$ regular linear systems of dimension
\[
p_a=5 \qquad (p_g=10),
\]
if it is verified that the $g_8^1$ extracted from non-special $g_8^3$ form a (non-linear) system contained in a complete linear system $\infty^5$.

The examples cited, whose study needs to be carried further, suggest various problems concerning irregular surfaces, in particular ``whether the paracanonical system $\{K\}$ of a surface with $p_a\geq0$ and $p_g>p_a$ has dimension $p_g$ and consists, in general, of $\infty^{p_g-p_a}$ regular linear systems $\infty^{p_a}$''.

But at least the case in which the absolute linear genus is $p^{(1)}=1$, and hence the canonical curves, for $p_g>2$, are reducible, being composed of the elliptic curves of a pencil, gives rise to an exception.

In this case one is dealing with surfaces $F$ of irregularity $q=p_g-p_a$ possessing a pencil of elliptic curves, of genus $q$; in fact one can prove that the continuous system $\{C\}$ to which an arbitrary curve $C$ belongs cuts out on the $k$ (components of the canonical curves $K$) a series of equivalent groups: if this were not so, one would find, besides the pencil $|k|$, a second elliptic pencil of curves, and one would deduce that the $k$ are curves of constant modulus and that $p_a=-1$ (cf. Chap. X). %

This being so, if $p_a>0$, the canonical curves of $F$ will be composed of $p_g-1>q-1$ curves $k$, forming within the pencil a non-special series: the paracanonical systems will therefore likewise be of dimension $p_g-1$, giving rise to a superabundant continuous system of dimension %
\[
2p_g-p_a-1.
\]

% scan: s190R
% !PRINTED 357
\weberpage{357}
On the other hand, for $p_a=0$, the canonical system $|K|$ will be the linear system $\infty^{p_g-1}$ formed by the $2q-2=2p_g-2$ curves of the canonical groups within the pencil $|k|$, and the paracanonical linear systems will have dimension $p_g-2$, giving rise to a continuous system $\{K\}$ of dimension
\[
2p_g-p_a-2.
\]

\textit{Remark.} -- Here it is worth adding that for $p^{(1)}=1$ and $p_a=0$ the case of irregular surfaces (of genus $p_g=1$) with pure canonical curve of order zero does not occur. Indeed \textit{for a surface $F$ of genus $p_g=1$ with canonical curve of order zero, one cannot have $p_a=0$}\origfootnote{1}{s190R-1}{Cf. \textsc{Enriques}, \textit{Intorno alle superficie algebriche di genere lineare $p^{(1)}=1$.} Rendic. Accademia Bologna, 9 December 1906 (no. 1).}. For a linear system $|C|=|C'|$ of genus $\pi$ and degree $n=2\pi-2$ on $F$ would give rise to a continuous series $\infty^1$ of linear systems of dimension $\pi-1$, and one of these, different from $|C|$ and not containing $|C|$, would have to cut out on a $C$ a non-special series of degree $2\pi-2$ and dimension $\pi-1$.

\edsection{9. Digression on Abelian varieties.}{9. Digression on Abelian varieties.}{sec.IX.9}

An Abelian variety (of rank 1) of $q$ dimensions, $V_q$, is an (algebraic) variety which admits a parametric representation by means of functions (Abelian, that is) $2q$ times periodic, of $q$ independent variables. representation one-to-one in the prismatoid of the periods: if these functions, as is supposed, are not degenerate, $V_q$ admits a simply transitive group $\infty^q$ of permutable transformations, without exception: this means that, given \textit{any} two (simple) points of $V_q$, there exists \textit{one} birational transformation of $V_q$ which changes the first point into the second. %
The simplest example of an Abelian variety is the variety of the groups of $p(=q)$ points, $G_p$; of a curve $C$ of genus $p$, that is, the \textit{Jacobi variety} corresponding to the curve\origfootnote{2}{s190R-2}{Cf., e.g., \textsc{Enriques-Chisini}, \textit{Lezioni}, Book VI, chap. III, \S~36. vol. IV, p. 223.}. %

This $V_p$ admits a simply transitive mixed group of birational transformations into itself, which is generated by the $\infty^p$ involutory transformations of the second kind defined by the $g_{2p}^p$ belonging to the curve, and by the $\infty^p$ permutable transformations of the first kind, which constitute a continuous group $\Gamma_p$ and are defined as products of an even number of transformations of the second kind. The said Jacobi variety, $V_p$, contains an \textit{exceptional variety}, $V_{p-1}$, corresponding to the series of special groups $G_p$ on the curve; this $V_{p-1}$ is transformed into a $V_{p-2}$ of dimension $p-2$ by the
% scan: s191L
% !PRINTED 358
\weberpage{358}
transformations of the first or of the second kind, for example by the involution between groups \(G_p\) defined on \(C\) by a \(g^p_{2p}\)\origfootnote{1}{s191L-1}{Cf. \textsc{Enriques-Chisini}, \textit{Lezioni}, Book VI, chap. II, § 27 (vol. IV, p. 143).}.

When among the \(p\) independent integrals of the first kind of the curve \(C\) of genus \(p\) there are \(q<p\) (linearly independent) ones reducible with \(2q\) periods, within the Jacobi \(V_p\) corresponding to the curve there is found a family \(\infty^{p-q}\) of Abelian varieties \(V_q\) constituting a system of imprimitivity for \(\Gamma_p\), each admitting a simply transitive algebraic group \(\infty^q\) of birational transformations into itself, and also, by the Picard-Poincaré theorem, an analogous family \(\infty^q\) of Abelian varieties \(V_{p-q}\)\origfootnote{2}{s191L-2}{Cf. e.g., \textsc{Enriques-Chisini}, \textit{Lezioni}, Book VI. Castelnuovo's geometric interpretation. Chap. III, 4, 38 (vol. IV, p. 234).}.

\textit{Abelian varieties} \(V_q\) of this type, contained imprimitively within a Jacobi \(V_p\), are the \textit{Picard varieties} attached to irregular algebraic surfaces\origfootnote{3}{s191L-3}{Cf. \textsc{R. Torelli}, \textit{Sulle serie algebriche semplicemente infinite di gruppi di punti appartenenti a una curva algebrica}, Rend. Circ. Mat. di Palermo, t. XXXVII, 1914 (§ II, no. 11).}: indeed, on an irregular surface \(F\), of irregularity \(q\), the non-equivalent curves of a continuous system \(\infty^q\) cut out, on a curve of genus \(p\) (\(>q\)), a series formed of \(\infty^q\) non-equivalent groups of points, or of \(\infty^q\) linear series, which constitutes precisely an Abelian variety \(V_q\) within the Jacobi \(V_p\) corresponding to the curve.

The preceding observation is of general scope, because one proves\origfootnote{4}{s191L-4}{Cf. \textsc{G. Castelnuovo - F. Enriques}, \textit{Sur les intégrales simples....}, Ann. de l’école normale sup., s. 3, t. XXIII,, 1906.}, though by transcendental methods, that every Abelian variety (of rank \(1\)) is a Picard variety attached to some irregular algebraic surface; so that one can at once state \textsc{Torelli}'s theorem: \textit{There always exist Jacobian varieties containing a given Abelian variety imprimitively.} %

If, however, one does not wish to have recourse to the result of \textsc{Castelnuovo-Enriques}, obtained by transcendental means, one can develop the following considerations, which lead to a theorem less precise than that of \textsc{Torelli}, but sufficient for the subsequent developments of our exposition.

Adopting the most general geometric point of view, we shall \textit{define} as an \textit{Abelian variety} of \(q\) dimensions, \(V_q\), a \(V_q\) which possesses a \textit{group \(\infty^q\) \(\Gamma\) of birational transformations into itself, simply transitive without exception}, such, therefore, that there exists \textit{one} non-degenerate transformation of \(\Gamma\) under which \textit{any two} (simple) \textit{points of \(V_q\)} correspond to each other.

% scan: s191R
% !PRINTED 359
\weberpage{359}
The property of the group \(\Gamma\) of being simply transitive without exception on \(V_q\) entails as a consequence that: \textit{there exists on \(V_q\) no linear system of dimension \(>0\) of varieties \(V_{q-1}\) which is transformed transitively into itself by \(\infty^r\) (\(r>0\)) transformations of \(\Gamma\).} In other words: a variety \(V_{q-1}\) on \(V_q\) which is not transformed into itself by a continuous subgroup of \(\Gamma\) cannot be transformed by \(\Gamma\) into the varieties of a linear system \(\lvert V_{q-1}\rvert\).

Let there be, on the Abelian variety \(V_q\), a linear system \(\lvert V_{q-1}\rvert\) of (irreducible) varieties, of dimension \(r>0\), which is transformed transitively into itself by an (algebraic) subgroup \(\infty^r\) of \(\Gamma\). One can represent \(\lvert V_{q-1}\rvert\) by means of a linear space \(S_r\), in which there is given a transitive continuous group \(\infty^r\) of homographies. Now, within this group, there will be found a continuous subgroup \(\Gamma'\), of dimension \(h\geq1\), which leaves invariant a line \(a\) of \(S_r\), permuting its points in \(\infty^1\) ways. Hence there will be in the \(\Gamma\) of \(V_q\) a subgroup \(\Gamma_h\) of dimension \(h\), which interchanges among themselves in \(\infty^1\) ways the varieties \(\overline{V}_{q-1}\) of a linear pencil within \(\lvert V_{q-1}\rvert\), and carries a generic point of a \(\overline{V}_{q-1}\) into the points of a \(\overline{V}_h\) formed of \(\infty^1\overline{V}_{h-1}\) belonging to the \(\overline{V}_{q-1}\) of the said pencil. But in this pencil there is (at least) one \(\overline{V}_{q-1}\) invariant under \(\Gamma_h\) (corresponding to a united point on the united line \(a\) of \(S_r\)), and hence \(\overline{V}_h\) must contain a \(\overline{V}_{h-1}\) invariant under \(\Gamma_h\): associating two points of this \(\overline{V}_{h-1}\), one will have not one but \(\infty^1\) transformations of \(\Gamma_h\), and this contradicts the supposition that the \(\Gamma\) of \(V_q\) is simply transitive without exceptions.

Therefore on \(V_q\) there are no linear systems \(\lvert V_{q-1}\rvert\) invariant under a continuous subgroup of transformations of \(\Gamma\), but there may be \(\lvert V_{q-1}\rvert\) invariant under a finite group, a subgroup of \(\Gamma\) itself.

Now, on \(V_q\) one can choose a variety \(V_{q-1}\), of \(q-1\) dimensions, which is not transformed into itself by any (non-identical) transformation of the group \(\Gamma\). Indeed, \(V_q\) being projectively realized in a suitable space \(S_r\), a generic hyperplane section of it will certainly not be changed into itself by any transformation of \(\Gamma\).

This being so, the series \(\infty^q\) of the transforms of the said \(V_{q-1}\), say \(\{V_{q-1}\}\), will constitute an algebraic entity identical with the group \(\Gamma\), and therefore also with the variety \(V_q\). In general, in the said series there may be found a finite number \(s\) of \(V_{q-1}\) equivalent to the one from which we started. Hence, cutting the \(V_{q-1}\) of the said series by a generic algebraic curve drawn on \(V_q\), one will obtain a series comprising \(\infty^q\) non-equivalent groups of points, which can be represented by an Abelian variety \(W_q\) (contained in the Picard variety of the given \(V_q\)), which, for \(s=1\), turns out to be identical with the family \(\{V_{q-1}\}\) and with \(V_q\) itself. If instead \(s>1\), the indicated family, and \(V_q\), are rep%
% scan: s192L
% !PRINTED 360
\unskip\weberpage{360}%
resented on \(W_q\) taken as multiple according to the number \(s\), since to \(sV_{q-1}\) that are equivalent there will correspond a single point of \(W_q\). Referring to this second case, \(s>1\), there will therefore be between \(W_q\) and \(V_q\) a correspondence \((1,s)\), by virtue of which there is defined on \(V_q\) an involution of order \(s\), \(I'\), invariant under the transformations of \(\Gamma\) and therefore free of double points, whose locus would give rise to an exception to the simple transitivity of \(\Gamma\).

In conclusion: \textit{an Abelian variety \(V_q\) can always be represented on another Abelian variety, \(W_q\), simple or multiple without branch points, which is contained imprimitively in the Jacobi variety corresponding to a curve of genus \(p(\geq q)\).}

Conversely: \textit{every irreducible variety \(V_q\) which can be represented on an Abelian variety \(W_q\), simple or multiple without branch points, contained imprimitively in a Jacobi variety, is itself an Abelian variety.}

To prove this theorem consider the correspondence \((s,s)\) in which to a group \((A,B\,\ldotp\ldotp\ldotp\ldotp)\) of the involution \(I'\) there corresponds another group of conjugate points \((A',B'\,\ldotp\ldotp\ldotp\ldotp)\); a correspondence defined by the transformation of the Abelian group of \(W_q\) which carries its homologous points into each other: \(\overline{A}\) and \(\overline{A}'\). It is a matter of proving that this correspondence \((s,s)\) splits into \(s\) birational transformations, and thus, for example, that the point \(A'\) depends on \(A\) not only algebraically but also uniquely, and hence rationally. To this end we shall make the point \(A\) vary continuously on \(V_q\), describing a closed cycle \(\gamma\), and we shall recognize that \(A'\) also moves continuously, likewise describing a closed cycle, and therefore cannot be interchanged, e.g., with \(B'\), but must unfailingly return to \(A'\). %

If the cycle described by \(A\) is a null cycle \(\gamma\), our assertion is obvious, since it cannot contain within itself any branch point of the correspondence \((s,s)\), the involution \(I_s\) being free of double points. But the matter can also be proved for a non-null cycle \(\gamma\) on the Riemann space corresponding to the variety \(V_q\). And the proof results from a requirement of continuity, if one takes into account that the correspondence \((A\,A')\) between the cycle \(\gamma\) and the path \(\gamma'\) described without ambiguity by the homologous point \(A'\) can vary continuously until it reduces to the identity, if \(A'\) is made to vary until it coincides with \(A\); and in the limit the path \(\gamma'\), supposed open, cannot reduce to a closed path, the points conjugate in the involution \(I_s\) always remaining distinct.

It thus follows that the variety \(V_q\) possesses a simply transitive group of birational transformations into itself, there being precisely \textit{one} transformation of it which carries a point \(A\) into another point \(A'\). In particular there will be a birational transformation which carries a point \(A\) into a point \(B\) conjugate to it in the involution \(I_s\),
% scan: s192R
% !PRINTED 361
\weberpage{361}
and thus the said involution will be generated by a finite group of \(s\) transformations: \(G_s\). %

Now, of the birational transformations \((A\,A')\) which we have found on \(V_q\) (and which correspond to the permutable transformations of the \(s\)-fold variety \(W_q\)) one cannot say, a priori, that they are permutable. Since they transform into itself the involution generated by the group \(G_s\), whose groups of points correspond to the points of \(W_q\), one can only say that the transform of one of them, \(\Pi\), by means of another, \(T\),
\[
T\Pi T^{-1}=\Pi'
\]
differs from \(\Pi\) by a cyclic transformation \(\Omega\) of \(G_s\):
\[
\Pi'=\Pi\Omega.
\]

But \(\Omega\) must remain unchanged when \(T\) varies in the continuous group \(\Gamma\), which contains the identity, and therefore
\[
\Omega=1.
\]

One therefore deduces that \(V_q\) possesses a simply transitive group \(\infty^q\) of permutable transformations into itself, that is, it is an Abelian variety.

\textit{Note.} -- The transcendental theory of Abelian varieties \(V_q\) results from the classical works of \textsc{Jacobi}, \textsc{Riemann}, and then of \textsc{Poincarè}, \textsc{Picard}; for (hyperelliptic) surfaces --- case \(q=2\) --- the study was carried further by \textsc{G. Humbert} and then also by the Italian school. %

\edsection{10. Continuation: the genus of Abelian varieties.}{10. Continuation: the genus of Abelian varieties.}{sec.IX.10}

We wish to complete what has been said about Abelian varieties by evaluating their (geometric, principal) genus.

We shall begin with the Jacobi varieties, proving that: \textit{The Jacobi variety, \(V_p=J_p\), corresponding to a curve of genus \(p\), has genus one.}

We shall prove that on the \(J_q\) of the \(q\)-tuples of points of a curve \(C\) of genus \(p\) (\(q\leq p\)) the variety \(K\) of the \(q\)-tuples taken from the individual groups of a \(g^{q-1}_{2p-2}\) contained in the canonical series of \(C\) is a canonical variety. Take first \(q=2\) (\(p\geq2\)). On the surface \(J'_2\) of the (unordered) pairs of points of the curve \(C\), there will be \(\infty^1\) curves \(M\), irreducible and birationally identical with \(C\), each of which corresponds to the series of pairs with a fixed point; so that the \(M\) constitute, on the surface, a system of index two. %

Let us take a \(g^1_{2p-2}\) contained in the canonical series of \(C\), and let \(K\) be the image curve of the pairs taken from this series \(\infty^1\); let us deter%
% scan: s193L
% !PRINTED 362
\unskip\weberpage{362}%
mine the intersections of \(K\) with a curve \(M\). To this end consider the \(M\) which corresponds to the series of pairs with a fixed point \(A\) (a generic point of \(C\)), and let
\[
A,\quad A_1,\quad A_2\quad \ldotp\ldotp\ldotp\ldotp\quad A_{2p-3}
\]
be the group of the \(g^1_{2p-2}\) which contains \(A\). The points corresponding to the pairs
\[
(A\,A)\,(A\,A_1)\,\ldotp\ldotp\ldotp\ldotp\,(A\,A_{2p-3})
\]
belong to \(M\) and constitute a canonical group of it, and on the other hand all these points, except the first, \((A\,A)\), also belong to \(K\). Now the group of these points, common to \(K\) and \(M\), is residual to the characteristic series of \(M\) with respect to its canonical series. Indeed, consider, besides \(M\), the curve \(M'\) representing the pairs of \(C\) which have in common another point \(B\): \(M\) and \(M'\) intersect at the point homologous to the pair \((A\,B)\), and this tends to the point \((A\,A)\) when \(B\) comes infinitely near to \(A\), so that \((A\,A)\) will constitute a group of the characteristic series on the curve \(M\). Now, since \(K\) cuts out on every curve of the system \(M\) a group residual to the characteristic series, by the \textit{third criterion of equivalence} (cf. Chap. III, § 13) \(K\) itself will be a canonical curve of \(J_q\), \hfill Q.E.D.

Having thus proved our theorem for \(q=2\), we shall prove it for \(q=3\), and so on successively by complete induction. Indeed, supposing the theorem true for the variety of the groups of \(q-1\) points of \(C\) (\(q\leq p\)), one recognizes that it is true also for the \(J_q\) of the \(q\)-tuples.

To this end consider on \(J_q\) the series of the \(\infty^1\) varieties \(M_{q-1}\), irreducible and birationally identical, each of which represents the \(q\)-tuples with a fixed point: this series is of index \(q\). Let \(K_{q-1}\) be the variety of the \(q\)-tuples taken from the individual groups of a \(g^{q-1}_{2p-2}\) contained in the canonical series of \(C\), and let us determine the intersection of \(K_{q-1}\) with \(M_{q-1}\). Let \(g^{q-2}_{2p-3}\) be the series residual to the point \(A\) with respect to the series considered. Adding the fixed point \(A\) to the groups of \(q-1\) points extracted from the groups of the \(g^{q-2}_{2p-3}\), one obtains a variety \(K_{q-2}\) which is common to the \(M_{q-1}\) corresponding to the point \(A\) and to the \(M_{q-1}\) image of the \(g^{q-1}_{2p-2}\). %
Now if, on \(M_{q-1}\), we add to \(K_{q-2}\) the variety \(M_{q-2}\) of the \((q-1)\)-tuples having the fixed point \(A\) (that is, the variety of the groups of \(q\) points of which the point \(A\) forms part twice), we shall obtain an (impure) \textit{canonical} variety of \(M_{q-1}\); and this by the theorem hypothetically assumed as the basis of our induction, because the said variety corresponds to the series of the groups of \(q-1\) points extracted from the individual groups of the series
\[
g^{q-2}_{2p-3}+A,
\]
a series with the fixed point \(A\), contained in the canonical series of \(C\).

% scan: s193R
% !PRINTED 363
\weberpage{363}
Therefore \(K_{q-1}\) cuts out on \(M_{q-2}\) a variety residual to a canonical variety with respect to \(M_{q-2}\); hence, by the third criterion of equivalence (which easily extends from the case of curves to that of varieties of any dimension), it will suffice to prove that \(M_{q-2}\) constitutes a variety of the characteristic system of \(M_{q-1}\). And indeed the \(M_{q-1}\) relative to the fixed point \(A\) and the \(M'_{q-1}\) relative to the fixed point \(B\) have in common the variety \(\overline{M}_{q-2}\) of the \(q\)-tuples with the two fixed points \(A\) and \(B\), and when \(B\) tends to \(A\), \(M'_{q-1}\) becomes infinitely near to \(M_{q-1}\) and \(\overline{M}_{q-2}\) tends precisely to \(M_{q-2}\). %

The stated theorem being thus established, put \(p=q\); we shall have: \textit{on the Jacobi variety \(J_p\), representing the groups of \(p\) points of a curve of genus \(p\), the variety \(K_{p-1}\) of the special \(p\)-tuples is a canonical variety} (a priori impure). But it is precisely an \textit{exceptional variety}, which can be transformed into a variety of dimension \(p-2\) on the variety, birationally identical with \(J_p\), representing the non-equivalent groups of \(p\) points of the curve \(C\) of genus \(p\). This transformation is obtained, moreover, by making use of the involutory correspondence which exists, on \(C\), between the groups of \(p\) points residual with respect to a \(g^p_{2p}\).

Therefore it follows that the \textit{Jacobi variety \(J_p\) has genus one.}

The theorem established for the Jacobi varieties \(J_p=V_p\) now extends easily to more general Abelian varieties, and it will suffice to indicate the proof briefly. First of all, referring to a simple \(V_q\) contained imprimitively in a \(V_p\) (\(q<p\)), this \(V_q\) will form part of a family of varieties of the same dimension, a family of index one; and hence the (impure) canonical variety of \(V_p\) will cut out on each such \(V_q\) \textit{a} canonical variety. Secondly, the theorem extends further to the case of a \(V_q\) which is represented on a multiple variety \(W_q\), contained in \(V_p\), without branch varieties (cf. § 9). Indeed the transformation \((1,s)\) which leads from \(W_q\) to \(V_q\) changes the canonical variety of \(W_q\) into a canonical variety of \(V_q\)\origfootnote{1}{s193R-1}{Theorem of Painlevé-Castelnuovo, established for \(q=1\) in \textsc{Enriques-Chisini}, \textit{Lezioni}. Book V, chap. I, § 9 (vol. III, p. 72). For \(q=2\) cf. Chap. V, § 7 of this treatise. The theorem extends without difficulty to all values of \(q\).}.

In conclusion we can state that: \textit{Abelian varieties \(V_q\) (\(q\geq1\)) are always of genus one, and with pure canonical variety of order zero.}

\textit{Historical note.} -- The evaluation of the genus of the Jacobi varieties \(V_p\), and in general of Abelian varieties, is immediate from the transcendental point of view, since one at once constructs a \(p\)-fold integral of the first kind attached to \(V_p\). The geometric proof of the
% scan: s194L
% !PRINTED 364
\weberpage{364}
theorem which we have set forth above was given by \textsc{Severi} in a Note of 1911\origfootnote{1}{s194L-1}{\textit{Sulle superficie e varietà algebriche irregolari di genere geometrico nullo.} R. Acc. Lincei, 1911.}.

Another geometric proof of the same theorem was indicated by \textsc{B. Segre}\origfootnote{2}{s194L-2}{\textit{Determinazione di certi gruppi covarianti di due o più serie lineari.} Rendic. Circolo Mat. Palermo, 1932.}, and consists in showing that the linear series \(g^{p-1}_{2p}\) contained in a series \(g^p_{2p}\) on a curve \(C\) of genus \(p\) have as their images, on the corresponding Jacobi variety, the varieties \(V_{p-1}\) of a linear system whose characteristic system, added to the intersection with the image of the variety of special groups, constitutes the canonical system on a \(V_{p-1}\). We shall confine ourselves to explaining the matter in the simplest case, \(p=2\), noting that the extension of the reasoning to \(p>2\) is indeed less simple than one might wish. %

Let us consider, then, a curve \(C\) of genus \(p=2\), and on it a generic linear series \(g^2_4\): the \(g^1_4\) contained in this series give rise, on the Jacobi surface \(V_2\), to the curves \(L\) of a net, free of fundamental curves, which will have a variable intersection \(A\) with the exceptional curve \(V_1\), image of the canonical \(g^1_2\) of \(C\), or a base point \(B=V_0\) at the point into which that curve is transformed by the involution defined by the said \(g^2_4\). The canonical curves \(K\) of \(V_2\) are defined a priori, on the surface, by the property of cutting out on a generic \(L\) groups residual to the characteristic series of \(\lvert L\rvert\), increased by the intersection \(A\) of \(L\) with the exceptional curve \(V_1\), or by the base point \(B\) (transform of \(V_1\)).

It is therefore a matter of proving that the series \(A+B+g^2_6\) on an \(L\) is (totally) contained in the canonical series \(g^4_8\). And the assertion derives essentially from the geometric interpretation of Zeuthen's formula which links the two canonical series of the curves \(C\) and \(L\) in \((2,3)\) correspondence\origfootnote{3}{s194L-3}{This step becomes quite simple (for \(p=2\)) if one introduces the auxiliary curve \(M\) which represents the pairs of points of \(L\) complementary with respect to a \(g^1_4\) contained in the \(g^2_4\).}.

Thus, then, the net of the curves \(L\) on \(V_2\), or rather on a transform of it without exceptional curves, turns out to be of genus \(\pi=5\), of effective degree \(6\) and of virtual degree \(n=8\), (\(=2\pi-2\)), having two simple base points, with a canonical characteristic series. The possession of such a net \(\lvert L\rvert\), without proper fundamental curves, implies precisely that \(V_2\) is a surface of genus \(p_g=1\) with canonical curve of order zero, \hfill Q.E.D. %

% scan: s194R
% !PRINTED 365
\weberpage{365}
\edsection{11. Irrational pencil on surfaces of geometric genus zero.}{11. Irrational pencil on surfaces of geometric genus zero.}{sec.IX.11}

From the circumstance that an Abelian variety \(V_q\) of dimension \(q\geq2\) has geometric genus equal to one, one deduces that ``\textit{every non-exceptional surface \(F\) belonging to an Abelian variety \(V_q\) has geometric genus \(p_g>0\)}''.

For \(q=2\) the matter is evident, \(F\) coinciding with \(V_2\). To prove the theorem in general one must observe that for \(q>2\) ``\(F\) on \(V_q\) forms part of a continuous system of analogous surfaces, of dimension at least \(q-2\)''.

Indeed, applying to \(F\) the \(\infty^q\) transformations of the group \(\Gamma\) of \(V_q\), one will have, in general, \(\infty^q\) surfaces analogous to \(F\): an exception is the case in which \(F\) admits a continuous series of transformations into itself; but in any case (by the simple transitivity of the group \(\Gamma\)) the transformations of \(F\) will be at most \(\infty^2\), and hence one will have at least \(\infty^{q-2}\) surfaces transformed from \(F\).

This being so, the property of \(F\) of being of non-zero geometric genus (as we have stated) is a consequence of the general property that: On a variety \(V_q\) of dimension \(q\) which has geometric genus \(P>0\), every surface \(F\) which belongs to a continuous family of dimension at least \(q-2\) has genus \(p_g>0\).

Let us prove this proposition referring first to the case \(q=3\), and then to the case \(q=4\), etc.

Supposing \(q=3\), the canonical surface which exists by hypothesis on \(V_q\) (even if it consists of a single surface, of order zero) intersects \(F\) along a curve which belongs to the system residual to the series of the characteristic system of \(F\) with respect to its canonical system\origfootnote{1}{s194R-1}{Cf. VIII, § 9.}; this suffices to prove the existence of the latter, that is, that the genus of \(F\) is \(p_g>0\).

If \(q_g=4\), the surfaces \(F\), which are at least \(\infty^2\), can be distributed in at least \(\infty^1\) simply infinite continuous systems, which therefore give rise to a series \(\infty^1\), at least, of varieties \(V_3\): now, since the genus of \(V_4\) is greater than zero, a canonical variety of \(V_4\) will cut out on each of the aforesaid \(V_3\) special surfaces, residual to the characteristic system of \(V_3\) with respect to the canonical system, and hence these \(V_3\) will be of genus \(P>0\). From this it follows that the \(\infty^1\) surfaces which constitute them are also of genus \(p_g>0\). %

By an analogous procedure, passing from \(q\) to \(q+1\), the matter is proved for a variety \(V_q\) of any dimension.

The theorem just established, that the surfaces belonging to an Abelian variety are of genus \(p_g>0\), now allows us to make precise, for
% scan: s195L
% !PRINTED 366
\weberpage{366}
surfaces of genus \(p_g=0\), the characteristic property of irregular surfaces, by proving that: \textit{Irregular surfaces of genus \(p_g=0\) possess an irrational pencil of curves, of genus \(-p_a=p\). Conversely: every irregular surface of genus \(p_a=0\) which possesses an irrational pencil of genus \(p\) (\(>0\)) is an irregular surface of arithmetic genus \(-p\).} %

Let \(F\) be an irregular surface and --- without making for the moment any hypothesis about the values of its characters --- consider on it a simply infinite continuous system of curves \(\{C\}\), \textit{of index} \(i>2\), which we may suppose \textit{simple}, that is, such that the \(i\) curves \(C\) passing through a generic point of the surface do not consequently have other points in common.

Denote by
\[
C_i=\Sigma C
\]
the curve which is the sum of the \(i\) curves of \(\{C\}\) issuing from a point of the surface \(F\).

Three cases may arise:

\textit{a}) the \(\infty^2\) curves \(C_i\) are all equivalent on the surface \(F\);

\textit{b}) the said \(C_i\) are distributed in \(\infty^1\) simply infinite systems of equivalent curves;

\textit{c}) the \(C_i\) are all non-equivalent to one another, or else each of them turns out to be equivalent to a finite number, \(s\geq1\), of curves.

Now we shall prove that \textit{in case a) all the curves \(C\) are equivalent to one another, that is, the series \(\{C\}\) is contained in a linear system.}

To this end it suffices to prove that the \(\infty^1\) curves \(C\) cut out a series of equivalent groups on an arbitrary curve \(M\) drawn on \(F\), and to invoke, for example, the first criterion of equivalence. Indeed, the property asserted for the series cut out on the curve \(M\) derives from a known criterion of equivalence of Severi-Castelnuovo relating to the series \(\infty^1\) of groups of points belonging to a curve\origfootnote{1}{s195L-1}{Cf. III, 13. See also \textsc{Enriques-Chisini}, \textit{Lezioni}. Book V, chap. IV, § 42 (vol. III, p. 483).}. %

Secondly we shall prove that \textit{in case b) one can construct on the surface \(F\) an irrational pencil of curves.}

Indeed, let a point \(A\) move on \(F\) in such a way that the curves \(C_i\) issuing from it remain equivalent; since the condition imposed is algebraic, the point \(A\) will describe an algebraic curve \(K\), and it is clear that the curves \(K\) so defined on \(F\), since no two of them can intersect, must form a pencil (free of base points that are simple for the surface). Supposing, for simplicity of discussion, that the \(K\) are irreducible (for otherwise it would suffice to consider in
% scan: s195R
% !PRINTED 367
\weberpage{367}
place of the \(K\) their irreducible components), we shall now prove that the pencil \((K)\) is irrational.

Indeed, the same consideration made above shows that the curves \(C\) of \(\{C\}\) must cut out, on any \(K\), \(\infty^1\) equivalent groups of points. Hence, by the \textit{second criterion of equivalence}\origfootnote{1}{s195R-1}{Cf. III, 13.}, the \(C\) will be equivalent up to curves of the pencil \((K)\); that is --- denoting by \(C\) and \(C_0\) any two curves of \(\{C\}\) --- one will have
\[
C+\Sigma K=C_0+\Sigma'K,
\] %
where the two summations \(\Sigma\) and \(\Sigma'\) contain the same number of terms. Consequently, if the curves \(C\) are not all equivalent, the pencil \((K)\) cannot be linear, for otherwise, by adding and subtracting the same number of curves \(K\), one would pass from a \(C\) to an equivalent \(C_0\).

An objection can be raised to this reasoning, founded on the second criterion of equivalence, under the hypothesis that the pencil \((K)\) contains reducible curves, such as
\[
K=K_1+K_2\qquad \text{and}\qquad \overline{K}=\overline{K}_1+\overline{K}_2;
\]
since then one might have, for example,
\[
C_0=C+K_1-\overline{K}_1+\ldotp\ldotp\ldotp\ldotp
\] %

But the doubt raised here is eliminated by the consideration that the curves \(C_0\) and \(C\) belong to one and the same continuous system \(\{C\}\), and as \(C_0\) varies in \(\{C\}\) the components of the reducible \(K\), being finite in number, remain fixed; hence, reducing \(C_0\) to \(C\) by continuity, one necessarily obtains
\[
C_0\to C,\qquad K_1=\overline{K}_1.
\] %

Finally, placing oneself in \textit{case c)}, one proves at once that the \textit{genus of the surface must be \(p_g\geq1\).} Indeed, representing the curves of \(\{C\}\) by the points of a curve \(L\), the \(\infty^2\) curves \(C_i\) define on \(L\) itself a series \(\gamma'_i\) which can be regarded as a doubly infinite variety, simple or multiple, contained in the Jacobi variety which corresponds to the curve \(L\), representing the non-equivalent linear series of \(L\) itself. %

The preceding discussion leads one to recognize that, for an irregular surface \(F\) of genus \(p_g=0\), every series \(\{C\}\) of index \(i\), consisting of non-equivalent curves, chosen on it, gives rise to case \textit{b)}, so that there is found on the surface an irrational pencil of curves \(K\).

% scan: s196L
% !PRINTED 368
\weberpage{368}
\sourcewarning{The scan shows overlapping sheets. For the following two paragraphs the transcription states that it recovered the covered portions from a corresponding scan, s197L, which is not included in the material supplied. The text as received is kept, but these readings cannot be verified here. The limitation ends before ``Historical note''.}
One then recognizes that the genus of \((K)\) equals the arithmetic genus with its sign changed: \(p=-p_a\). For, if one had \(p<-p_a\), by virtue of the second criterion of equivalence one would succeed in constructing on \(F\) a series \(\{C\}\) of non-equivalent curves cutting out on each \(K\) non-equivalent groups of points, and hence (case \(b\) occurring) a second irrational pencil of curves; but the possession of two irrational pencils on \(F\), \((K)\) and \((M)\), would imply that this surface must have geometric genus \(p_g>0\): the matter appears evident if the two pencils are formed of mutually unisecant curves\origfootnote{1}{s196L-1}{Indeed one obtains a canonical curve of \(F\) by adding to the curves of \((M)\) which pass through the points of a canonical group of a \(K\) the curves of \((K)\) which pass through the points of a canonical group of an \(M\).}; while the case in which the curves of the two pencils meet in a certain number \(s>1\) of points leads to representing \(F\) on a multiple (\(s\)-fold) surface with two pencils of mutually unisecant curves, whence it follows, a fortiori\origfootnote{2}{s196L-2}{Cf. V, 7.}, that its genus is \(p_g>0\). %

In conclusion the characteristic property of irregular surfaces of genus \(p_g=0\), which entails the possession of an irrational pencil of genus \(p=-p_a\), is proved, in accordance with the preceding statement. %

\textit{Historical note.} – The theorem concerning the existence of an irrational pencil on irregular surfaces of genus \(p_g=0\) was stated by \textsc{Enriques} as a consequence of the characteristic property of irregular surfaces, and this in his cited note in the Rendic. dell’Accademia di Bologna in 1904. Indeed, as early as 1900 \textsc{Castelnuovo} had had occasion to observe and to communicate to his colleague that ``from the existence of a continuous system of non-equivalent curves on a surface of genus \(p_g=0\), one deduces the existence of an irrational pencil''. This property, which \textsc{Enriques} recalls in his Note, mentioning its author, was justified by transcendental means, as follows.

Let \(\{C\}\) be a continuous system \(\infty^1\) of non-equivalent curves on the surface \(F\) of genus \(p_g=0\). Then, as was first pointed out by G. \textsc{Humbert}\origfootnote{1}{s196L-3}{Comptes rendus, 1893; Journal del Math., s. 4, t. X, 1894.}%
, one constructs on \(F\) simple integrals (of total differentials) of the first kind, by summing the values which the Abelian integrals of the first kind pertaining to the series \(\{C\}\) take in correspondence with the \(i\ (>1)\) curves \(C\) issuing from one and the same point \(A\) of \(F\). And these simple integrals \(I_1\,I_2\,\ldots.\) will necessarily be functions of one another; otherwise — as was
% scan: s196R
% !PRINTED 369
\weberpage{369}
indicated by \textsc{Noether}\origfootnote{1}{s196R-1}{\textit{Ueber die totalen algebraischen differentialausdrücke.} Math. Annalen, Bd. 29, 1887. Cf. \textsc{Picard} et \textsc{Simart}, op. c., vol. I, Chap. V, \S\ 15.} — one would deduce the existence of a double integral of the first kind, and hence \(p_g>0\).

This being so, \textsc{Castelnuovo} observes that the curves \(I=\mathrm{const.}\) will be composed of the algebraic curves of an irrational pencil, which are defined by making \(A\) move on the surface in such a way that the \(i\) curves \(C\) issuing from a point describe equivalent groups of curves. This reasoning, especially in the form it takes through the considerations of \textsc{Castelnuovo}'s Note ``Sugli integrali semplici appartenenti ad una superficie irregolare''\origfootnote{2}{s196R-2}{R. Acc. Lincei, 21 May, 4 and 18 June, 1905. Cf. Memorie scelte, p. 473 ff. (see in particular no. 10, p. 490).}, was translated into the geometric one set forth above by F. \textsc{Severi}\origfootnote{3}{s196R-3}{\textit{Sulle superficie e varietà algebriche irregolari di genere geometrico nullo.} Rendic. Lincei, s. V, vol. XX, 1911.}.

\edsection{12. Note on surfaces of irregularity one.}{12. Note on surfaces of irregularity one.}{sec.IX.12}

In view of possible applications, it is worth pointing out that the considerations developed in the preceding section also lead to a characterization of surfaces of irregularity \(q=p_g-p_a=1\), for which the following theorem holds: \textit{Surfaces \(F\) of irregularity \(q=1\) possess an elliptic pencil of curves} (of genus \(\pi\geq0\)). \textit{On them every continuous system of curves is obtained from a linear system by adding to it the difference of two curves of the said pencil.}

Let us observe that the Picard variety of \(F\) is an elliptic curve. Now take on \(F\) any system \(\infty^1\) of non-equivalent curves, \(\{C\}\), of index \(\nu>1\): through every point \(P\) of \(F\) there will pass \(\nu\) curves \(C\), and, as \(P\) varies, one will obtain \(\infty^2\) groups of \(\nu\) curves; hence, starting from a particular one of these points, \(P=\overline{P}\), there will be \(\infty^1\) points \(P\) through which there passes a group of \(\nu\) curves \(C\) equivalent to the group of the curves passing through \(\overline{P}\). The locus of these points \(P\) will be a curve \(K\) through \(\overline{P}\), and all the analogous \(K\) will form on \(F\) an elliptic pencil of curves \((K)\). The first part of the theorem being established in this way, to prove the second it will suffice to note that the curves \(C\) of the continuous system \(\{C\}\) will cut out equivalent groups on every \(K\), and consequently (by the second criterion of equivalence) two non-equivalent \(C\) will differ from each other only by the difference of two \(K\).

\printersignature{\textsc{Enriques F.} - \textit{Superficie algebriche.}\quad 24}
% scan: s198L
% !PRINTED none
\weberblank{[s198L]}

\clearpage\thispagestyle{plain}
% scan: s198R
% !PRINTED none
\weberpage{[s198R]}
\edchapter{X}{SURFACES OF GEOMETRIC GENUS ZERO}{Chapter X. SURFACES OF GEOMETRIC GENUS ZERO}{SURFACES OF GEOMETRIC GENUS ZERO}

\edsection{1. Introduction.}{1. Introduction.}{sec.X.1}

The characteristic property of irregular surfaces of genus \(p_g=0\), which we established in \S\S\ 4 and 11 of the preceding chapter, makes it possible, as we shall see, to characterize — completely — the surfaces of genus \(p_g=0\), and — in particular — to assign, by means of the vanishing of the plurigenera, the necessary and sufficient conditions which define the family of ruled surfaces.

The development of our deductions will proceed in the following order.

First of all we shall deal with surfaces which contain systems of curves of genus \(\pi\) and degree \(n>2\pi-2\), proving that they are transformable into ruled surfaces, rational or irrational; the rational case corresponds (as is well known) to the regularity of the surface.

Secondly we shall set forth a lemma concerning the reducible curves of an irrational pencil, which will enable us to state in a simpler form the second criterion of equivalence, as it has to be used in the developments that follow.

From this we shall deduce, to begin with, that surfaces of genus \(p_g=0\) and \(p_a<-1\) can at once be referred to ruled surfaces, of genus \(p_a=-p\). There then remain to be studied the surfaces with genera \(p_g=0\) and \(p_a=-1\), which give rise to:

1) surfaces transformable into elliptic ruled surfaces, and

2) surfaces containing an elliptic pencil of curves of genus \(\pi\geq1\).

We shall treat separately the two cases which arise, for \(\pi>1\) and \(\pi=1\), proving that in either case the surface contains a second linear pencil of elliptic curves, birationally identical.

Finally we shall study in depth the family of (\textit{elliptic}) surfaces defined by the aforesaid properties, we shall give their con%

% scan: s199L
% !PRINTED 372
\unskip\weberpage{372}%
struction, we shall evaluate their plurigenera and recognize that they admit an elliptic group, simply infinite, of birational transformations into themselves, etc. We shall end the chapter by characterizing the surfaces with vanishing plurigenera and those with plurigenera equal to zero or to one.

\edsection{2. Surfaces containing a system of curves of genus \(\pi\) and of degree \(n>2\pi-2\).}{2. Surfaces containing a system of curves of genus π and of degree n>2π-2.}{sec.X.2}

Let \(F\) be a surface which contains a linear system of curves \(|C|\), at least \(\infty^2\), irreducible, for which \(n>2\pi-2\). The curves \(C'\), \(C''\) etc., successive adjoints of \(|C|\), will cut the \(C\) in a decreasing number of points:
\[
2\pi-2,\qquad 4\pi-4-n,\qquad 6\pi-6-2n\,\ldots.
\]
and will therefore form a finite sequence; in other words, the operation of adjunction, starting from the linear system \(|C|\), must terminate after a finite number of steps, giving rise to a last adjoint system, \(|C^t|\), of dimension \(\geq0\), which no longer possesses any adjoint curve of order \(>0\). %

We already know one case in which adjunction terminates in this way: it is the case of the regular surfaces of genus \(p_g=p_a=0\), with bigenus \(P=0\), in which, as is known, the surface \(F\) is rational\origfootnote{1}{s199L-1}{Castelnuovo's theorem cf. Chap. V, \S\ 4.}. Setting this case aside, we must examine the hypothesis that the surface on which our system \(|C|\) lies is irregular, of genus \(p_g=0\) and \(p_a=-p\leq-1\), and therefore contains an irrational pencil, of genus \(p\), of curves \(K\). If the \(K\) are of genus \(\varrho=0\), \(F\) is at once transformable into a ruled surface of genus \(p\)\origfootnote{2}{s199L-2}{Enriques's theorem cf. e.g. \textsc{Enriques-Conforto}, \textit{Le superficie razionali}, Book II, chap. I, \S\ 4-6 (p. 253).}. We must examine, and reduce to absurdity, the hypothesis that the genus of the \(K\) is instead \(\varrho\geq1\). And it will be useful to distinguish the two cases:
\weq{1)}{p>1}
\weq{2)}{p=1.}

\textit{First case:} \(p>1\).

Denote by \(s\) the number of intersections of the \(C\) with the \(K\), and observe that the \(C+pK-p\overline{K}\) give rise to an \(\infty^p\) family of distinct linear systems cutting equivalent series on the \(K\); from this there follows
\[
s>1.
\]

Now, since the (virtual) canonical curves of \(F\) cut the \(K\), of genus \(\varrho\), in \(2\varrho-2\) points, the \(C'\) will cut the same \(K\) in \(s+2\varrho
% scan: s199R
% !PRINTED 373
\text{\weberpage{373}}
-2\)\edFlush{} points, the \(C''\) will cut them in \(s+2\,(2\varrho-2)\) points, and so on. Finally the curves of the last adjoint system \(|C^t|\) will cut the \(K\) in a number of points which will be, a fortiori, \(m\geq2\). But this conclusion at once appears absurd if the \(C^t\) are irreducible, because, since they possess an irrational involution \(\gamma_m^1\) of genus \(p\), cut by the \(K\), they must have genus (effective and virtual) greater than or equal to \(m(p-1)+1>p\), and hence possess an adjoint system of dimension \(>0\).

The conclusion (which is a reduction to absurdity of the hypothesis from which we started) extends to the case in which \(C^t\) is formed of reducible curves, since in any case a component of \(C^t\) must cut the \(K\) in at least \(2\) points: this assertion follows from observing that the \(C^t\) belong to a continuous series of non-equivalent curves, which is obtained by adding to them the difference of two curves of the pencil \((K)\): \(K-\overline{K}\).

\textit{Second case:} \(p=1\).

Proceed as in the preceding case: starting from the system \(|C|\) with the characters \(\pi\) and \(n>2\pi-2\), one will arrive at a last adjoint system \(|C^t|\), consisting of curves cutting the \(K\) in \(m\geq2\) points; but here it seems that the system adjoint to \(C^t\) may actually be lacking, if the genus (effective and virtual) of \(C^t\) is equal to \(1\), like that of the pencil of the \(K\). However, our hypothesis implies that the curves \(C^t\), supposed irreducible (or their components), belong to an elliptic family, \(\infty^1\), of non-equivalent elliptic curves; and since the curves of this family cannot form a second irrational pencil (different from \(|K|\)), we must accordingly admit that the curves \(C^t+K-\overline{K}\) give rise to a continuous system \(\{C^t\}\) of curves of degree \(\nu\geq1\).

But from this one infers that the surface \(F\) must possess \(\infty^1\) rational curves and hence must be birationally equivalent to an elliptic ruled surface, in contradiction with our hypotheses. Indeed, if \(\nu=1\), the elliptic family \(\infty^1\) of mutually unisecant curves \(C^t\) will contain \(\infty^1\) equivalent pairs of curves, and (by Lüroth's theorem)\origfootnote{1}{s199R-1}{Cf. \textsc{Enriques-Chisini}, \textit{Lezioni.} Book II, chap. I, \S\ 3 (vol. I, p. 169) and Book V, chap. I, \S\ 7 (vol. III, p. 48).} their points of intersection will generate \(\infty^1\) rational curves on \(F\). If instead \(\nu>1\), the (irreducible) curves of \(C^t\) with \(\nu-1\) fixed points will give rise to an \(\infty^1\) family of mutually unisecant elliptic curves, so that we fall back into the preceding case.

Thus the fundamental theorem is established:

\textit{Surfaces containing an irreducible linear system, at least \(\infty^2\), of genus \(\pi\) and degree \(n>2\pi-2\), are transformable into ruled surfaces, rational or irrational.}

% scan: s200L
% !PRINTED 374
\weberpage{374}
And from this, as we announced in \S\ 9 of Chap. IV, it follows that:

\textit{The surfaces belonging to the family of ruled surfaces are the only surfaces containing infinitely many exceptional curves; and all surfaces outside this family can be transformed into surfaces free of exceptional curves.}

As regards the \textit{fundamental theorem}, it should be added that it \textit{holds} under the widest conditions, \textit{whatever the dimension of the system \(|C|\) between whose characters the inequality \(n>2\pi-2\) holds, provided that the \(C\) themselves do not contain exceptional curves as components.}

To prove this, observe that the character \(a(C)=n-(2\pi-2)\) is an additive character of the linear systems to which it refers, and therefore it cannot be positive for the sum of two systems without being positive for one of the summands. Thus it will suffice to prove that:

\textit{Any surface \(F\) on which there lies a curve \(C\), irreducible and not exceptional, of genus \(\pi\) and degree \(n>2\pi-2\), certainly belongs to the family of ruled surfaces.}

Indeed the multiple of the said curve \(C\) by a sufficiently large number \(s\) will have genus and degree
\[
\pi_s=s\pi+\frac{s(s-1)}{2}n-s+1,\qquad n_s=s^2n;
\]
and hence it will have dimension
\[
r\geq p_a+n_s-\pi_s+1>0,
\]
constituting an irreducible system (for \(n>0\)) for which
\[
n_s>2\pi_s-2.
\]

And let it be explicitly noted that what has been said holds also for \(\pi=0\), because, since it is excluded that \(C\) is an exceptional curve, and hence that \(n=-1\), we shall have
\[
n\geq0;
\]
and indeed, in the case \(n=0\), which alone needs to be examined, the multiple \(|sC|\), for \(s\) very large, will always turn out to be composed with the rational curves of a pencil, whence it follows that \(F\) is transformable into a ruled surface.

\textit{Historical note.} – The theorem that the surfaces on which adjunction terminates, and therefore those which contain linear systems with characters \(n\) and \(\pi\) for which \(n>2\pi-2\), belong to the family of ruled surfaces, constitutes the fundamental result of the Memoir of \textsc{Castelnuovo-Enriques}, « Sopra alcune questioni fondamentali nella teoria delle superficie algebriche », published in the Annali di Matematica, 1901. %

% scan: s200R
% !PRINTED 375
\weberpage{375}
The original proof of this theorem did not make use of the property of irregular surfaces of genus \(p_g=0\) of containing an irrational pencil of curves (a property discovered, as has been said, in 1904). In its place there intervened an inequality which derives from comparing Noether's expression for the Zeuthen-Segre invariant with the Riemann-Roch theorem. Let us briefly indicate the idea of this line of proof.

If (on a surface \(F\) which must be of genus \(p_g=0\)), starting from a linear system \(|C|\), adjunction terminates, a last adjoint system \(|C^t|\) being reached, the genus of this system must be \(\pi\leq-p_a=p\).

Let us suppose, for simplicity, that \(|C^t|\) is irreducible, at least \(\infty^1\); denoting its degree by \(n\), we shall have that the Zeuthen-Segre invariant of the surface equals %
\[
\delta-n-4\pi=-12p-\dot{p}^{(1)}+9,\qquad (p=-p_a)\text{\origfootmark{1}{s200R-1}} %
\]\origfoottext{1}{s200R-1}{Here \(\dot{p}^{(1)}\) denotes the relative linear genus of the surface.}
while the Riemann-Roch theorem gives us its dimension
\[
r\geq n-\pi+1-p.
\]
From this one deduces
\[
5\pi+r\geq11p+p^{(1)}-8+\delta\geq11p+p^{(1)}-8.
\]

Then an appropriate discussion, in which one must distinguish the hypothesis \(p^{(1)}\leq1\) from the hypothesis \(p^{(1)}>1\), leads to finding on the surface:

1) either a linear system of curves of genus \(\pi\) and of dimension
\[
r\geq3\pi-5,
\]
whence it follows that the surface is rational or birationally equivalent to a ruled surface of genus \(\pi\)\origfootnote{2}{s200R-2}{This is a consequence of the argument used by \textsc{Enriques} in the Note \textit{Sulla massima dimensione dei sistemi lineari di curve di dato genere appartenenti ad una superficie algebrica.} Atti Acc. di Torino, 1894.};

2) or else a linear system of curves split into rational curves, whence it likewise follows that the surface is birationally equivalent to a ruled surface.

This sketch gives barely an idea of the path followed: one then needs a very minute analysis, in relation to various complicating circumstances, and a special examination of the cases in which \(p<3\).

The considerably simpler proof of our fundamental theorem, based on the irrational pencil belonging to an
% scan: s201L
% !PRINTED 376
\weberpage{376}
irregular surface of genus \(p_g=0\), was outlined in the « Lezioni » of \textsc{Enriques-Campedelli}\origfootnote{1}{s201L-1}{Cf. Chap. III.} published in the Rendiconti del Seminario Matematico di Roma (1934). That proof has been developed here in a simpler form, avoiding the appeal, made there, to the theorem of Castelnuovo-Enriques on double planes with vanishing plurigenera. Its place is taken by the elementary consideration used above, that an \(\infty^1\) series of equivalent pairs of elements, on an elliptic \(\infty^1\) entity, is a \(g_2^1\) and is therefore certainly rational. %

Finally we shall mention that the \textit{extension of the fundamental theorem} to the case of \textit{a single curve}, not exceptional, of genus \(\pi\ (\geq 0)\) and degree \(n>2\pi-2\), was brought to light by L. \textsc{Campedelli} in the Note « Intorno alle superficie algebriche su cui esistono curve di genere \(\pi\) e di grado \(n>2\pi-2\) » (Rendic. Lincei, \(1933_2\)), and the author himself has pointed out how the explicit use of this property serves to simplify our theory at several points. %

\textit{Note.} – The proposition which we used to prove the fundamental theorem of this section also deserves special notice. \textit{The surfaces on which adjunction terminates belong to the family of ruled surfaces.}

And therefore, recalling the developments given at the end of § 1 of Chap. VI, we can assert that: \textit{The surfaces not belonging to the family of ruled surfaces, and therefore reducible to a model without exceptional curves, have absolute linear genus}
\[
p^{(1)}\geq 1.
\]

In Chap. VI, § 1, it had only been recognized that for \(p^{(1)}\leq 0\), since adjunction terminates, all the plurigenera would have to be \(P_i=0\); here one sees further that the surface would have to be birationally equivalent to a ruled surface (in contradiction with the hypothesis that the exceptional curves can be eliminated).

\edsection{3. Lemma on the reducible curves of a pencil.}{3. Lemma on the reducible curves of a pencil.}{sec.X.3}

In view of the developments that follow, we must specify the \textit{valence} of a curve split in any way, belonging to a rational or irrational pencil, for the purpose of counting the curves endowed with a double point, whose number appears in the expression of the Zeuthen-Segre invariant (Chap. V, § 2, 3). We shall refer, for simplicity, to surfaces without exceptional curves, which therefore do not belong to the family of ruled surfaces, and we shall prove that: \textit{With regard to the computation of the invariant}
% scan: s201R
% !PRINTED 377
\weberpage{377}
\textit{of Zeuthen-Segre relative to a surface which possesses a pencil \(|C|\), rational or irrational, of irreducible curves of genus \(\pi>0\), without base points, every split curve of the pencil has valence \(\geq 0\).}

Consequently the number, \(\delta\) or \(\Delta\), of the curves of the pencil endowed with a double point must always be taken to be \(\geq 0\). \textit{If \(\delta=0\) or \(\Delta=0\), the pencil cannot contain any split curve, except multiple elliptic curves (of degree zero), in the case where the curves of the pencil are elliptic \((\pi=1)\).}

To justify these assertions, we must recall \textit{Godeaux's formula}, which (Chap. V, § 2) we said expresses the valence of a particular curve \(\bar{C}\) of the pencil, split in any way into multiple parts. Let us denote by \(\bar{\delta}\) the valence of a particular \(C=\bar{C}\), which we shall suppose split into multiple components of genus \(\varrho_1,\ \varrho_2\,\ldots.\):
\[
\bar{C}=t_1\theta_1+t_2\theta_2+\ldots.
\]
meeting one another in
\[
s_{hk}=\theta_h\theta_k
\]
points. We shall have:
\[
\bar{\delta}\geq \Sigma(t_i-1)(2\varrho_i-2)+\Sigma s_{hk}(t_h+t_k-1),
\]
the sign \(>\) referring to possible complications, for example to the case in which one of the components \(\theta_i\) possesses a double point, or to that in which some of the points of connection common to two of the components, \(\theta_h\) and \(\theta_k\), come to coincide in a point common to \(r>2\) components\origfootnote{1}{s201R-1}{The inequality \(\bar{\delta}\geq 0\) was recognized by \textsc{Castelnuovo-Enriques} (Annali di Mat., 1901) in the simplest case, in which \(\bar{C}\) is split into a multiple irreducible curve \(\theta_1\) of genus \(\varrho_1\) and another, simple, component \(\theta_2\) which meets it in \(s_{12}=\theta_1\theta_2\ (>0)\) points. The general formula written above is put forward, by extending the argument used in that case, by L. \textsc{Godeaux}, \textit{Sur le calcul de l'invariant de Zeuthen-Segre.} (Mém. de la Soc. des Sciences de Hainart, 1920) and later also by L. \textsc{Campedelli}, \textit{Sul computo dell'invariante di Zeuthen-Segre; Ancora sul computo etc.} Lincei, 1934. The latter author has derived from it the proof that \(\bar{\delta}\geq 0\), and more precisely that one can have \(\bar{\delta}=0\) only when the pencil \((C)\) is of genus \(\pi=1\) and \(C\) is a multiple elliptic curve of genus \(\varrho=1\): %
\[
\bar{\delta}=(t+1)(2\varrho-2)=0. %
\]

This is the property stated in our lemma, which also appears in the \textit{Lezioni} of \textsc{Enriques-Campedelli} (Sem. Mat. di Roma, 1934) and which now receives a considerably simpler proof due to A. \textsc{Franchetta}.}.

Let us observe that, since the surface is free of exceptional curves, we have, for every \(i\),
\[
s_{ii}=\theta_i\theta_i\leq 2\varrho_i-2;
\]
% scan: s202L
% !PRINTED 378
\weberpage{378}
hence
\[
\bar{\delta}\geq \Sigma s_{ii}(t_i-1)+\Sigma s_{hk}(t_h+t_k-1);
\]
moreover, since for \(h\neq k\) we have \(s_{hk}\geq 0\),
\[
\begin{aligned}
\Sigma s_{hk}(t_h+t_k-1)&=\Sigma s_{hk}(t_h-1)+\Sigma s_{hk}(t_k-1)+\Sigma s_{hk}\geq\\
&\geq\Sigma s_{hk}(t_h-1)+\Sigma s_{hk}(t_k-1).
\end{aligned}
\]

We therefore have
\[
\bar{\delta}\geq\Sigma s_{ii}(t_i-1)+\Sigma s_{hk}(t_h-1)+\Sigma s_{hk}(t_k-1).
\]

Let us introduce the curve
\[
C^*=\theta_1+\theta_2+\ldots
\]
and let us denote by \(\varrho\) and \(\nu\) its virtual characters. The expression on the right-hand side of the preceding inequality equals the number of virtual intersections \(C^*(\bar{C}-C^*)\); therefore the theorem will be proved if we show that \(C^*(\bar{C}-C^*)\geq 0\), that is, if we reduce to absurdity the hypothesis
\[
C^*(\bar{C}-C^*)=-i\qquad(i>0).
\]

Noting that \(C^*\bar{C}=C^*C=0\), the hypothesis adopted would entail
\[
\nu=C^*C^*=i>0;
\]
hence the curve \(rC^*\) would determine a linear system \(|rC^*|\) of dimension \(d_r\) which, for \(r\) sufficiently large, would turn out to be
\[
d_r\geq p_a+r^2i-r\varrho-\frac{r(r-1)}{2}i-r+1=p_a+\frac{r^2i}{2}-\frac{r}{2}(2\varrho-2-i)+1, %
\]
that is, infinite of the second order as \(r\) tends to infinity.

On the other hand the curve \(r\bar{C}\) determines a linear system \(|r\bar{C}|\), composed with the curves of the pencil; its dimension \(D_r\) is given, for \(r\) sufficiently large, by \(D_r=r-p\), where \(p\) is the genus of the pencil; that is, it is infinite of the first order as \(r\) tends to infinity. But it is clear that the system \(|rC^*|\) is partially contained in the system \(|r\bar{C}|\), and hence that \(d_r\leq D_r\). The hypothesis from which we started is therefore absurd. %

It follows moreover that one can have \(\bar{\delta}=0\) only if \(\Sigma s_{hk}=0\), and this is possible only if \(C^*\) consists of a single component \(\theta\), and if \(\theta\cdot\theta=2\varrho-2=0\), that is, if \(\varrho=1\).

\edsection{4. Surfaces of genus \(p_g=0\) with \(p_a<-1\).}{4. Surfaces of genus p\_g=0 with p\_a<-1.}{sec.X.4}

Let \(F\) be a surface of geometric genus \(p_g=0\) and of negative arithmetic genus \(p_a=-p<-1\). We know (Chap. IX, § 11) that it
% scan: s202R
% !PRINTED 379
\weberpage{379}
must contain a pencil of genus \(p\) of curves \(K\), having genus \(\pi\geq 0\). Let us prove that we must have \(\pi=0\), and therefore \(F\) birationally equivalent to a ruled surface of genus \(p\), by reducing to absurdity the hypothesis \(\pi>0\).

Let us therefore assume as hypothesis that \(\pi>0\), and hence that \(F\) (not birationally equivalent to a ruled surface) can be reduced to a surface without exceptional curves; its absolute linear genus will be (§ 2, Note)
\[
p^{(1)}\geq 1.
\]
Let us now write the formula which gives the Zeuthen-Segre invariant of \(F\) in relation to the irrational pencil \(K\). We shall have (Chap. V, § 2)
\[
\Delta-4+4(p-1)(\pi-1)=12p_a-p^{(1)}+9
\]
and hence, for \(\pi>0\),
\[
\Delta\leq 13-12p-p^{(1)}.
\]
But from this there follows, for \(p>1\), \(\Delta<0\), which is in contradiction with the lemma of our § 3.

We therefore conclude that \(\pi=0\). Thus: \textit{Surfaces of genus \(p_g=0\) and \(p_a=-p<-1\) are birationally equivalent to ruled surfaces of genus \(p\).}

\edsection{5. Surfaces \(F\) with \(p_g=0\) and \(p_a=-1\), possessing an elliptic pencil of curves \(K\) of genus \(\pi>1\): lemma I.}{5. Surfaces F with p\_g=0 and p\_a=-1, possessing an elliptic pencil of curves K of genus π>1: lemma I.}{sec.X.5}

The analysis of the surfaces of genera \(p_g=0\) and \(p_a=-p=-1\) leads to the examination of surfaces \(F\) which possess an elliptic pencil of curves \(K\) of genus \(\pi\geq 1\); and it is convenient to distinguish the two cases which present themselves as possible when \(F\) does not belong to the family – of ruled surfaces, namely: %
\[
\pi>1\qquad\text{and}\qquad\pi=1.
\]

Here we discuss the case \(\pi>1\).

First of all, observe that our surface \(F\), which we may assume to be without exceptional curves, has absolute linear genus \(p^{(1)}=1\). Indeed, computing the Zeuthen-Segre invariant for the elliptic pencil \((K)\)\origfootnote{1}{s202R-1}{Cf. V, 2 and 3.} one finds
\[
I=\Delta-4+4(p-1)(\pi-1)=\Delta-4=-12-p^{(1)}+9,
\]
so that, since \(\Delta\geq 0\), one deduces \(p^{(1)}\leq 1\), and hence (§ 2, Note)
\[
p^{(1)}=1.
\]

It is also worth noting that for our pencil \((K)\) we consequently have
\[
\Delta=0,
\]
% scan: s203L
% !PRINTED 380
\weberpage{380}
and therefore (§ 3) the pencil itself (formed of curves of genus greater than one) contains no curve \(K\) endowed with a double point, or split in any way.

This being so, let us proceed to establish a series of lemmas which will then lead us to a definitive theorem

\textsc{Lemma I.} – \textit{The surface \(F\) possesses at least one parabicanonical elliptic curve, cutting the curves \(K\) of the pencil in groups of the bicanonical series \(g_{4\pi-4}^{3\pi-4}\)}; and similarly \textit{a para-\(i\)-canonical elliptic curve}, for \(i>2\).

Indeed, consider the linear system \(|K''|\), second adjoint of a curve \(K\) of our pencil. Since \(p^{(1)}=1\), and the virtual canonical curves cut the \(K\) in groups of \(2\pi-2\) points, the genus of the \(K'\) is
\[
\pi'=\pi+1+2\pi-2-1=3\pi-2
\]
and hence the dimension of \(|K''|\) is
\[
r\geq-p+\pi'-1=3\pi-4,
\]
that is, at least equal to that of the bicanonical series of \(K\). Now, if the series cut by \(|K''|\) on \(K\) is not complete, through a group of it there passes at least a pencil of \(K''\), and one deduces the existence of a curve \(K''\) split into the \(K\) and a bicanonical curve, which is of genus
\[
3p^{(1)}-2=1:
\]
this constitutes an irreducible elliptic curve, or else a curve composed of elliptic curves without common points.

In the case where \(|K''|\) cuts the complete bicanonical series on \(K\), consider the continuous system \(\{K''\}\) (second para-adjoint), formed of \(\infty^1\) linear systems, which are obtained, starting from a particular \(|K''|\), by adding to it the difference of two curves \(K_1-K_2\) of the pencil \(K\).

Let us note that all the linear systems forming \(\{K''\}\) will cut the same bicanonical series on a \(K\), and similarly the groups cut by any two non-equivalent curves of a para-adjoint system (of order \(1,\ 2,\ 3,\ \ldots.\)) will always turn out to be equivalent.

It follows from this that there will be (at least) \(\infty^1\) curves of \(\{K''\}\) passing through one and the same bicanonical group of \(K\), and hence there will be a curve of this system split into a \(K\) and a residual parabicanonical curve: the latter, like the virtual bicanonical curve, will be of genus
\[
3p^{(1)}-2=1
\]
and will therefore consist of one or more irreducible elliptic curves without common points.

% scan: s203R
% !PRINTED 381
\weberpage{381}
In a perfectly analogous way one thus proves that the para-adjoint system of order \(i=3,\ 4\ldots.\), \(\{K^i\}\), of dimension \((2i-1)\pi-2i\), contains (at least) \(\infty^1\) curves passing through a group of the \(i\)-canonical series chosen on a \(K\), and therefore also a curve split into the \(K\) and a \textit{para-\(i\)-canonical curve}, elliptic or formed of irreducible elliptic curves without common points.

But one cannot assert that all these elliptic curves, for example the parabicanonical and the paratricanonical curve, are formed of \textit{distinct} irreducible components: for they might all reduce to a multiple paracanonical curve, that is, to an elliptic curve whose \textit{order} (number of intersections with the \(K\)) is \(2\pi-2\).

\edsection{6. Lemma II: outline of the proof.}{6. Lemma II: outline of the proof.}{sec.X.6}

\textit{On the surface \(F\)} (of genera \(p_g=0\) and \(p_a=-1\), with a pencil of curves of genus \(\pi>1\)) \textit{there exist at least two elliptic curves, not formed of the same components, whose multiples constitute two parapluricanonical curves}, cutting the \(K\) in equivalent groups of points.

We have already constructed a parabicanonical curve (or also a parapluricanonical curve of higher order), say \(C\), which is a curve of genus one, formed of irreducible components, simple or multiple, without common points. Let us now see how to construct a second curve, \(D\), of the same order \(4\pi-4\), which is not parabicanonical, but cuts on every \(K\) a non-bicanonical series whose multiple by \(s\) is the \(2s\)-canonical series cut by the curve \(sC\).

More precisely, in the case in which the genus \(\pi\) of the \(K\) is even, and therefore has the divisor \(s=2\), we propose to construct an elliptic curve \(D\) which, while having the same order \(4\pi-4\) as \(C\), cuts on every \(K\) a group which is not bicanonical, but belongs to a semiquadricanonical series of this curve. And we shall succeed in this in the simplest case, when the following two simplifying hypotheses are satisfied, which we shall discuss in the subsequent sections: I) on every curve \(K\) one can rationally determine a semibicanonical (but not canonical) series \(\bar{g}\); II) within the aforesaid \(\bar{g}\) one can rationally determine a group.

Under these hypotheses we shall first of all be able to determine rationally on a \(K\) a semibicanonical series \(\bar{g}=g_{2\pi-2}^{\pi-2}\) (different from the canonical series \(g=g_{2\pi-2}^{\pi-1}\)), then to determine in it a group \(G_{2\pi-2}\), and thus to construct a sufficiently large linear system of curves \(|L|\) which cuts on every \(K\) the corresponding series \(\bar{g}\). Then the linear system \(|L'|\), adjoint to \(|L|\), will cut on each \(K\) the semiquadricanonical series \(\bar{g}+g\). Now, comparing \(|2L'|\) and \(|2K''|\) (\(|K''|\) the second adjoint of a \(K\)) with respect to the second criterion of equivalence (made precise here by the lemma of § 3), we shall recognize that
% scan: s204L
% !PRINTED 382
\weberpage{382}
they differ only by a certain number \(r\) of curves \(K_1\,\ldots.\,K_r\) of the pencil \((K)\), to be counted twice:
\[
|2L'|=|2K''+2(K_1+\ldots.+K_r)|,
\]
and we shall deduce from this that the system \(|L'-(K_1+\ldots.+K_r)|\) has the same characters (degree, genus and hence virtual dimension \(3\pi-4\)) as \(|K''|\). Therefore it gives rise to a continuous system of dimension \(3\pi-3\) (one unit higher than that of the series cut on a \(K\)), from which it is possible to split off a \(K\): the residual curve is the semiparaquadricanonical \(D\) which we proposed to construct: \(C\) and \(D\) are certainly non-equivalent, and are not even para-equivalent, whereas \(2C\) and \(2D\) are para-equivalent.

Should the hypotheses II and 1 stated above not be fulfilled, the argument is repeated \textit{mutatis mutandis} in the following sense: %

1) When hypothesis II fails (I holding), instead of a \(\bar{g}\), a semibicanonical (non-canonical) series, one will construct on \(K\) a series \(\bar{\bar{g}}\), multiple of the said \(\bar{g}\) by an odd number \(i\), for which the said hypothesis is satisfied. Then in place of the system \(|L|\), or rather of its adjoint \(|L'|\), discussed above, one will succeed in constructing a linear system \(|M|\) which cuts on every \(K\) the series \(\bar{\bar{g}}\). Then the comparison of \(|2M|\) and \(|2K^i|\) (where \(|K^i|\) denotes the \(i\)-th adjoint of a \(K\)) will make us recognize that the first system, diminished by a certain number of curves \(K_1\,\ldots.\,K_r\) of the pencil \((K)\) counted twice, turns out to be para-equivalent to the second. From this one deduces that the system \(|M-(K_1+\ldots.+K_r)|\) has the same characters as \(|K^i|\); and hence (§ 5) that it belongs to a continuous system in which there is a curve split into a \(K\) and a residual elliptic curve \(D\).

The curve \(D\) has the same order as \(C=|K^i-\bar{K}|\), but the two curves are neither equivalent nor para-equivalent; rather, they have para-equivalent doubles, which cut equivalent groups on the \(K\).

2) In the absence of hypothesis I, one will pass from the surface \(F\) to an \(F^*\) with the same characters, represented on the \(m\)-fold \(F\) \((m>1)\) without branch curve, in such a way that the simplifying hypothesis (not verified for \(F\)) holds for \(F^*\). One will then construct on \(F^*\) a semiparapluricanonical elliptic curve \(D\) (of order \(4\pi-4\) or \(i(2\pi-2)\)), to which there will correspond on \(F\) a curve \(D\), likewise elliptic, (of order \(m(4\pi-4)\), or \(im(2\pi-2)\)). This curve \(D\) is not para-equivalent to the parapluricanonical curve \(C\) of the same order (that is, to the parabicanonical \(C=|K''-\bar{K}|\) which is constructed on \(F\) according to § 5); but the doubles \(2D\) and \(2C\) are para-equivalent, cutting equivalent groups on the \(K\).

In the case in which \(\pi\) is an odd number, denoting by \(s\) a prime (odd) divisor of it, we propose to construct on \(F\) an
% scan: s204R
% !PRINTED 383
\weberpage{383}
elliptic curve \(D\) of order \(4\pi-4\), which cuts on every \(K\) a group, not bicanonical, belonging to a series which is the \(s\)-th part of the \(s\)-bicanonical series. And we shall succeed in this in the simplest case, in which two hypotheses I and II are satisfied, analogous to those which arise in the first case of \(\pi\) even and \(s=2\). %

Indeed, there being given on every \(K\), according to these hypotheses, a non-special series \(\bar{g}=g_{2\pi-2}^{\pi-2}\), \(s\)-th part of the \(s\)-canonical series, and in it, rationally, a group of points, one will construct a sufficiently large linear system of curves \(|L|\), which cuts on each \(K\) the said series \(\bar{g}\). Then the adjoint system \(|L'|\) (of order \(4\pi-4\)) will cut on the \(K\) the series \(\bar{g}+g\), \(s\)-th part of the \(s\)-bicanonical series. And one will recognize that \(|sL'|\) differs from \(|sK''|\) (\(|K''|\) the second adjoint of a \(K\)) by a certain number \(r\) of curves \(K_1\,\ldots.\,K_r\) of the pencil \((K)\), to be counted \(s\) times:
\[
|sL'|=|sK''+s(K_1+\ldots.+K_r)|.
\]

In consequence of this relation, the system \(|L'-(K_1+\ldots.+K_r)|\) will have the characters (degree, genus and dimension \(3\pi-4\)) of \(|K''|\), and will be contained in an \(\infty^{3\pi-3}\) continuous system, to which there belongs a curve split into \(K\) and a \(D\): the latter is a curve of genus 1 and degree 0, whose \(s\)-fold multiple turns out to be para-equivalent to the \(s\)-fold multiple of the parabicanonical curve \(C\).

But, as in the first case (\(\pi\) even and \(s=2\)), it will be necessary to examine how the argument is modified if the hypotheses I and II, mentioned above, are not verified.

1) When hypothesis II does not hold, instead of a \(\bar{g}\), of order \(2\pi-2\), different from the canonical series \(g\) but the \(s\)-th part of the \(s\)-canonical series, one will take into consideration on \(K\) a series \(\bar{\bar{g}}=(2\pi-2)\bar{g}-(\pi-2)g\), of order \(\pi(2\pi-2)\), not pluricanonical, which depends rationally on the \(\bar{g}\) supposed given according to hypothesis I; of this \(\bar{\bar{g}}\) one will succeed in determining \textit{one} group. In this way one will construct, in place of \(|L|\), or rather of \(|L'|\), a linear system \(|M|\) which cuts on each \(K\) the corresponding series \(\bar{\bar{g}}\). Then one will prove that
\[
|sM|=|sK+s(K_1+\ldots.+K_r)| %
\]
and therefore
\[
|M-(K_1+\ldots.+K_r)|
\]
has the same characters as \(|K^\pi|\), and contains a curve split into a \(K\) and a residual elliptic curve \(D\): the curve \(D\) has the same order \(\pi(2\pi-2)\) as the para-\(\pi\)-canonical curve \(C=|K^\pi-\bar{K}|\), but the two curves are neither equivalent nor para-equivalent; only their multiples, \(sC\) and \(sD\), constitute para-\(s\)-canonical curves, cutting equivalent groups on the \(K\). %

% scan: s205L
% !PRINTED 384
\weberpage{384}
2) If, further, hypothesis I does not hold, one will in any case succeed in constructing rationally on \(K\) not one, but a group of \(m>1\) series \(\bar{g}\) (or \(\bar{\bar{g}}\)), and one will then obtain a surface \(F^*\), with the same characters as \(F\), represented on the \(m\)-fold \(F\) without branch curve; on this \(F^*\) the elliptic curves \(K^*\), corresponding to the \(K\), each contain a determinate series \(\bar{g}\) (or \(\bar{\bar{g}}\)). Then (hypothesis I being verified for \(F^*\)) one will succeed in constructing on \(F^*\) an elliptic curve \(D^*\), an aliquot part of a parapluricanonical curve, to which there will correspond on \(F\) the required elliptic curve \(D\).

We have thus explained the general outline of the proof of our lemma II; let us now carry out this proof in detail, distinguishing the two cases which present themselves: \(\pi\) even and \(s=2\), \(\pi\) odd and \(s\) a prime divisor of it (different from 2).

\edsection{7. First case: the genus \(\pi\) of the curves \(K\) is even.}{7. First case: the genus π of the curves K is even.}{sec.X.7}

Let us first of all recall that the bisection of a linear series \(g_{2n}^{2n-\pi}\) with \(n>\pi\), on a curve of genus \(\pi\), leads to \(2^{2\pi}\) series which are always distinct\origfootnote{1}{s205L-1}{Cf. e.g. \textsc{Enriques-Chisini}, \textit{Lezioni}, Book V, chap. III, § 35 (vol. III, p. 394). Cf. also Book VI, chap. III, § 34 (vol. IV, p. 199).}.

In particular the double of the canonical series on a \(K\), that is, the bicanonical series \(g_{4\pi-4}^{3\pi-4}\), will give rise to \(2^{2\pi}\) half-series, one of which is the canonical series \(g=g_{2\pi-2}^{\pi-1}\), while the others are \(2^{2\pi}-1\) semibicanonical series \(\bar{g}=g_{2\pi-2}^{\pi-2}\), distinct from \(g\) and distinct from one another, which we here take into consideration. Since \(K\), varying in its pencil, remains irreducible and never acquires double points, the said series \(\bar{g}\) never come to coincide; nevertheless they may be interchanged among themselves by a circuit of the curve \(K\) in its own (irrational) pencil, and it will be convenient to suppose in general that on a \(K\) one can define rationally — not \textit{one} series — but a group of \(m(\geq 1)\) series \(\bar{g}=g_{2\pi-2}^{\pi-2}\), distinct from one another.

Here, as we have announced, we shall introduce two simplifying hypotheses, which will later be justified relatively to the consequences deduced from them.

\gesperrt{Hypothesis I.} – That \(m=1\), that is, that for every \(K\) it is possible to determine uniquely, and therefore \textit{rationally, one series} \(\bar{g}\), semibicanonical (and not canonical).

\gesperrt{Hypothesis II.} – That in a series \(\bar{g}\) which is assumed as rationally given, one can determine uniquely, and therefore \textit{rationally, one group}, which, as \(K\) varies on \(F\), will give rise to a
% scan: s205R
% !PRINTED 385
\weberpage{385}
curve \(L\) cutting the \(K\) in semibicanonical groups of \(\bar{g}\). This second hypothesis might seem evident, since knowledge of the series \(\bar{g}\) on \(K\) leads in general to transforming \(K\) into a projectively defined \(K_{2\pi-2}\) of the space \(S_{\pi-2}\), on which the hyperplanes cut the series \(\bar{g}\), and thus this series is given \textit{linearly}, so that a hyperplane of \(S_{\pi-2}\) serves to cut out \textit{one} group of it. But a priori it does not follow that a series \(\bar{g}=g^{\pi-2}\) rationally given on \(K\) can be linearly defined in a unique way. For, among the infinitely many curves \(K_{2\pi-2}\), images of the series \(\bar{g}\), projective to one another, the choice of a particular \(K_{2\pi-2}\), such as one succeeds in determining uniquely when one starts from a group of \(\bar{g}\) itself, will in general depend on some numerical irrationality which has to be adjoined to the field of rationality given with \(K\) and with \(\bar{g}\), and such irrationalities will then give — as \(K\) varies — algebraic irrationalities\origfootnote{1}{s205R-1}{To clarify the difficulty, consider a plane quartic \(K_4\) defined as the envelope of an \(\infty^1\) series of index 2 of quadritangent conics \(K_2\); on \(K_4\) there is thus rationally defined the \(g_4^1\) of the quadruples of points of contact, which however is not defined linearly by means of a determinate \(\infty^1\) linear system of conics with 4 base points. Abstractly the \(g_4^1\) can be regarded as a rational curve representable on a line; but to determine the representation one must give a point of the said rational curve, which entails the adjunction of a quadratic irrationality. (Noether's theorem, cf. \textsc{Enriques-Chisini}, \textit{Lezioni.} Book V, chap. III, § 32 (vol. III, p. 341). %
}. %

On the basis of hypotheses I and II, let us first construct on \(F\) a curve \(L\) which cuts on each \(K\) of our elliptic pencil a group of \(2\pi-2\) points of the semibicanonical series \(\bar{g}\) (which has been assumed as given). The curves \(L'\) adjoint to \(L\) will cut on the \(K\) groups of a semiquadricanonical series, \(\bar{g}+g\), different from the bicanonical one, and will form a linear system \(|L'|\) which may be taken to be as large as one wishes, by adding to it — where necessary — a certain number of curves \(K\). Now compare the linear system \(|2L'|\) with \(|2K''|\) (the double of the second adjoint of a \(K\)). Since the two systems cut equivalent groups on every \(K\), by the second criterion of equivalence, made precise by the lemma of § 3, they must differ only by curves of the pencil \((K)\). But on this elliptic pencil the difference of two groups of elements (the first of which is very large) will turn out to be equivalent to a group of elements to be taken positively, so that we can write
\[
|2L'|=|2K''+K_1+K_2+\ldots.|. %
\]

Let us distinguish two cases which present themselves as possible a priori, namely that the number of the \(K_1\ K_2\,\ldots.\) is an even number, \(2r\), or that it is an odd number \(2r+1\).

\printersignature{\textsc{Enriques F.} - \textit{Superficie algebriche.}\quad 25}

% scan: s206L
% !PRINTED 386
\weberpage{386}
In the first case, by making the group of the curves \(K_1\,K_2\,\ldotp\ldotp\ldotp\ldotp\) vary within the linear series \(g_{2r}^{2r-1}\) which it determines in the elliptic pencil, one can arrange that the group itself reduces to a group counted twice:
\[
\begin{aligned}
&2K_1+2K_2+\ldotp\ldotp\ldotp\ldotp+2K_r\\
&(K_1=K_{r+1},\quad K_2=K_{r+2}\,\ldotp\ldotp\ldotp\ldotp).
\end{aligned}
\]

In the second case one can reduce the aforesaid group to another of \(r+1\) curves, among which \(r\) curves are counted twice:
\[
2K_1+2K_2+\ldotp\ldotp\ldotp\ldotp+2K_r+K_0.
\]

Hence we must have
\[
|2L'-2K''-2(K+K_1+\ldotp\ldotp\ldotp\ldotp+K_r)|=|K_0|; %
\]
and therefore the characters (\(\pi\) and \(n=0\)) of \(K_0\) are expressed by means of those (genus \(\varrho\) and degree \(\nu\)) of the virtual curve
\[
L'-K''-(K_1+\ldotp\ldotp\ldotp\ldotp+K_r)
\]
whose double is equivalent to \(K_0\).

We shall therefore have
\[
\begin{aligned}
n&=4\nu=0\\
\pi&=2\varrho+\nu-1=2\varrho-1,
\end{aligned}
\]
and consequently \(\pi\) must be odd. Since \(\pi\) has been supposed even, the second case examined as possible cannot occur, and we shall have
\[
|2L'-2(K_1+\ldotp\ldotp\ldotp\ldotp+K_r)|=|2K''|.
\]

Now the linear system \(|2L'-2(K_1+\ldotp\ldotp\ldotp\ldotp+K_r)|\) admits two non-equivalent halves:
\[
|L'-(K_1+\ldotp\ldotp\ldotp\ldotp+K_r)|
\]
and
\[
|K''|.
\]

Therefore the first system will have the same characters, genus and degree, as the second, which are:
\[
4\pi-4+\pi+1-1=5\pi-4
\]
and
\[
2\cdot(4\pi-4)=8\pi-8.
\]

Its dimension will therefore be
\[
3\pi-4,
\]
which is the dimension of the complete series cut by it on a \(K\). Consequently a group of \(3\pi-4\) points of this series will belong to \(\infty^1\) curves of the continuous system, \(\{L'-(K_1+\ldotp\ldotp\ldotp\ldotp+K_r)\}\)
% scan: s206R
% !PRINTED 387
\weberpage{387}
formed of \(\infty^1\) non-equivalent linear systems, and there will be at least one curve of this continuous system split into a \(K\) and a residual curve \(D\): the genus of \(D\) turns out to be equal to that of \(C=|K''-\bar K|\), that is, equal to \textit{one}. And — as was required — the elliptic curve \(D\), of the same order as the parabicanonical \(C\), is not equivalent (but rather semiequivalent) to it, and is therefore \textit{essentially distinct} from it, that is, it certainly contains some distinct (elliptic) component. %

But we must discuss the hypotheses I and II, which figure as presuppositions of the preceding deduction. Let us begin with the

\begin{center}
\textit{Discussion of hypothesis II.}
\end{center}

Hypothesis II, verified for \(\pi=2\)\origfootnote{1}{s206R-1}{In this case the series \(\bar g\) reduces to a single pair, which is therefore given with \(\bar g\).}, is no longer strictly true for the higher values of the even number \(\pi\) \((\pi=4,\,6\,\ldotp\ldotp\ldotp\ldotp)\). But, if one cannot in general cut out on \(K\) \textit{one} group \(G\) of \(\bar g\) or of the semibicanonical series \(\bar g+g\) (which is obtained from the \(\bar g\) supposed rationally given), we prove nevertheless that one succeeds in determining a group formed of \(i=\pi-1\) groups \(G\) (\(i\) odd), which belongs to a series — let us say \(\bar{\bar g}=i\bar g\) — not pluricanonical, but half of the \(2i\)-canonical series. Then one will be able to construct on \(F\) a curve \(M\) which cuts every \(K\) in a group of the corresponding \(\bar{\bar g}\), and, by operating on \(|M|\) as we previously operated on \(|L|\), or rather on \(|L'|\), that is, by splitting off a certain number of curves \(K\), one will obtain an elliptic semipara-\(2i\)-canonical curve \(D\) (essentially distinct from \(iC\)) whose double is para-equivalent to the bicanonical curve \(C\) counted \(i\) times.

Thus our argument based on hypothesis II can be repeated with slight modifications (substituting the series \(\bar{\bar g}\) for \(\bar g\) or for \(\bar g+g\)), provided one proves what is stated above: that one can determine \textit{one} group formed of \(i=\pi-1\) (an odd number) groups \(G\) of \(\bar g\) (and hence, linearly, \(\bar{\bar g}\)). For this purpose consider in the space \(S_{\pi-2}\) the family of all the mutually projective curves \(K_{2\pi-2}\) which are images of the semibicanonical series \(\bar g=g_{2\pi-2}^{\pi-2}\) given on \(K\). And let us try to construct two quadrics \(Q\) and \(Q'\), or a pencil of quadrics, in that space, which is (projectively) a covariant of \(K_{2\pi-2}\): then there will be determined, as a covariant of the pencil, the self-conjugate \(i\)-gon, with \(i=\pi-1\), common to \(Q\) and \(Q'\), whose faces will cut on \(K_{2\pi-2}\) \(\pi-1\) groups of \(\bar g\).

The construction of the aforesaid pencil of \(Q\) and \(Q'\) is carried out uniquely, and therefore rationally, starting from the \(K\) on which \(\bar g\) is given: this appears at once if \(K_{2\pi-2}\) does not lie on any quadric of \(S_{\pi-2}\), because the minimal double series of the series \(\bar g\), which is a bicanonical series
% scan: s207L
% !PRINTED 388
\weberpage{388}
(complete or not) cut on \(K_{2\pi-2}\) by the quadrics of its space, is defined linearly on the corresponding bicanonical curve \(K_{4\pi-4}\) of \(S_{3\pi-4}\) (itself also linearly given), by means of the hyperplanes of a linear system \(\alpha\) well determined by the curve \(K_{4\pi-4}\). Thus a pencil of quadrics of \(S_{\pi-2}\) is determined rationally as a function of \(\bar g\) by a pencil of hyperplanes of \(S_{3\pi-4}\) chosen in \(\alpha\) by imposing passage through a suitable group of points of \(S_{3\pi-4}\), fixed independently of \(K_{4\pi-4}\): that is, by a series \(g_{4\pi-4}^1\) contained in the bicanonical \(g_{4\pi-4}^{3\pi-4}\) (and more precisely in the minimal double series of \(\bar g\)). %

But if \(K_{2\pi-2}\) lies on \(\infty^r\) quadrics of \(S_{\pi-2}\) (\(r\geq1\)), the preceding argument breaks down, because to a group of the minimal double series of \(\bar g\), defined as the section of a hyperplane of \(\alpha\) on \(K_{4\pi-4}\), there no longer corresponds one quadric of \(S_{\pi-2}\), but rather an \(\infty^{r+1}\) linear system of quadrics; and similarly to an \(\infty^t\) linear system of hyperplanes of \(S_{3\pi-4}\) chosen in \(\alpha\) there corresponds an \(\infty^{r+t+1}\) linear system of quadrics of \(S_{\pi-2}\), cutting on \(K_{\pi-2}\) a series of smaller dimension \(g_{4\pi-4}^t\). Nevertheless, by putting %
\[
t=3\pi-4-2=3\pi-6,
\]
one will have in \(S_{\pi-2}\) a linear system of quadrics \(|Q|\), 2 units lower than the total system: now there is a harmonic family of envelope quadrics associated\origfootnote{1}{s207L-1}{Given, in a linear space of any number of dimensions, a locus-quadric \(a_x^2=0\) and an envelope quadric \(\alpha_u^2=0\), there is (according to \textsc{Battaglini} and \textsc{Rosanes}) a simultaneous linear invariant of them, the \textit{harmonizing invariant} \(a_\alpha^2\), which by its vanishing defines the relation of harmony between a locus quadric and the associated linear system of envelope quadrics. Cf. \textsc{Enriques-Chisini}, \textit{Lezioni}. Book III chap. I, \S~9 (vol. I, p. 61).} with \(|Q|\), which defines a conjugate \(i\)-gon (\(i=\pi-1\)) and hence a group of \(i\) groups of \(\bar g\) \hfill Q.E.D.

\begin{center}
\textit{Discussion of hypothesis I.}
\end{center}

Hypothesis II having been justified (at least in the sense that, if it is not verified for the series \(\bar g\), it will be for \(\bar{\bar g}\)), let us now see how hypothesis I too is relatively justified, in this sense: that — if it does not hold for our surface \(F\) — it will hold for an \(F^*\) with the same characters, represented on \(F\) counted \(m\) \((>1)\) times, without branch curve, so that it will be possible to construct on \(F^*\) an elliptic curve \(D^*\), to which there will correspond on \(F\) a semiparapluricanonical \(D\).

Suppose that on each \(K\) of \(F\) one is able to cut out uniquely not \textit{one} semibicanonical series \(\bar g\), but a group of \(m\) \((>1)\) such series \(\bar g\), interchangeable among themselves by a circuit of \(K\) in the
% scan: s207R
% !PRINTED 389
\weberpage{389}
pencil \((K)\). And, for simplicity of exposition, suppose \(\pi>2\), and hence \(\pi\geq4\), and accept with regard to each of these \(\bar g\) hypothesis II, namely that a linear image of it can be constructed as a rational function of it, that is, that one succeeds in defining \textit{one} curve \(K_{\pi-2}\) of \(S_{\pi-2}\), chosen among the infinitely many corresponding projective curves. %

The algebraic \(\infty^1\) variety which has as its elements our series \(g\) on the \(K\) (associated with these same \(K\)) can be regarded abstractly as a series \(\infty^1\Sigma\) of spaces \(S_{\pi-2}\), contained in a space of higher dimension \(S_h\). Since two of the \(m\) series \(\bar g\) on a \(K\) never come to coincide, the said series, \(\Sigma\), contains an elliptic involution \(\gamma_m^1\) of spaces, without double elements, and is therefore itself elliptic. %

Now in each space \(S_{\pi-2}\) of \(\Sigma\) one can construct (by hypothesis) \textit{one} curve \(K^*=K_{2\pi-2}\), image of the homologous \(\bar g\) chosen on the corresponding \(K\): the locus of the curves \(K^*\) (birationally identical with the \(K\)) will be a surface \(F^*\), represented on the \(m\)-fold \(F\) (without branch curve) in such a way that to each \(K\) of \(F\) there will correspond on \(F^*\) a group of \(m\) curves \(K^*\), always distinct from one another. On \(F^*\) the curves \(K^*\), birationally identical with the \(K\), form \textit{an elliptic pencil} \((K^*)\).

It is easy to see that the surface \(F^*\) has, like \(F\), arithmetic genus \(p_a=-1\) and linear genus \(p^{(1)}=1\), in agreement with the formula which gives the Zeuthen-Segre invariant of \(F^*\) in relation to the elliptic pencil of the \(K\), which are free of double points.

This being so, one can say of \(F^*\) what — under hypothesis I — was said of \(F\), since on \(F^*\), for each curve \(K^*\) (of genus \(\pi\)), there is given \textit{one} semibicanonical series \(\bar g\), rationally and linearly defined by the hyperplanes of \(S_h\). One deduces that \(F^*\) contains an elliptic semiparabicanonical curve \(D^*\), to which there will correspond on \(F\) a curve \(D\), likewise elliptic, semipara-\(2m\)-canonical.

\textit{Remark.} -- In the preceding discussion the hypothesis \(\pi=2\) has been excluded, a case in which, however, the discussion itself could be repeated with some modification. But for this case let it suffice to point out that, if the \(K\) are of genus \(\pi=2\), on each of them there is determined the group of the 6 double points of the \(g_2^1\), always distinct from one another, which — as \(K\) varies — will describe a semiparasexticanonical elliptic curve

\edsection{8. Second case: the genus \(\pi\) of the curves \(K\) is odd.}{8. Second case: the genus π of the curves K is odd.}{sec.X.8}

We shall denote by \(s\) an (odd) prime number dividing \(\pi\), and we shall consider on an arbitrary \(K\) one of the non-special series \(\bar g = g_{2\pi-2}^{\pi-2}\) which are obtained by dividing the \(s\)-canonical series by \(s\). Let us introduce, provisionally, two simplifying hypotheses analogous to I and II of the preceding section.

% scan: s208L
% !PRINTED 390
\weberpage{390}
\textit{Hypothesis I}: that one can determine uniquely, that is, rationally, on the variable \(K\) of the pencil \((K)\), a series \(\bar g\).

\textit{Hypothesis II}: that, a \(\bar g\) being rationally given on every \(K\), one can determine rationally \textit{one} group, which is thus uniquely defined as \(K\) varies in \((K)\); in other words, that the rationally given \(\bar g\) can be defined linearly (giving rise to \textit{one} \(K_{\pi-2}\) in \(S_\pi\)). %

On the basis of these hypotheses one can construct on \(F\) a curve \(L\) which cuts each \(K\) in a group of \(\bar g\); and \(|L|\), which is defined up to summands equal to \(K\), will constitute a linear system as large as one wishes. %

Let us denote by \(|L'|\) the linear system, of order \(4\pi-4\), adjoint to \(|L|\), and by \(|K''|\) the second adjoint of a \(K\), which is of the same order.

Let us compare, with respect to the usual criterion of equivalence, the two systems \(|sL'|\) and \(|sK''|\); they will differ by a certain number of curves of the pencil \((K)\), to be added to the smaller system \(|sK''|\); indeed it is legitimate to assume that the curve to be added to this lesser system is formed of a certain number of \(s\)-fold curves
\[
sK_1,\quad sK_2,\quad \ldotp\ldotp\ldotp\ldotp\quad sK_r
\]
and possibly also of a curve \(tK_0\), \(t\)-fold, of the said pencil, where \(t<s\):
\[
|sL'|=|sK''+s(K_1+\ldotp\ldotp\ldotp\ldotp+K_r)+tK_0|.
\]

But one proves that
\[
t=0,
\]
because the virtual curve
\[
\left(\frac{t}{s}K_0\right)=|L'-K''-(K_1+\ldotp\ldotp\ldotp\ldotp+K_r)|
\]
would have to have an integral genus \(\varrho\), whereas the genus of \(tK_0\) (of degree 0) is
\[
t\pi-t+1=s\varrho-s+1
\]
so that
\[
t(\pi-1)=s(\varrho-1):
\]
this relation cannot hold with \(\varrho\) an integer, because \(s\) (a prime divisor of \(\pi\)) is prime to \(\pi-1\) and to \(t\).

We shall therefore have
\[
|sL'|=|sK''+s(K_1+\ldotp\ldotp\ldotp\ldotp+K_r)|
\]
and
\[
|sL'-s(K_1+\ldotp\ldotp\ldotp\ldotp+K_r)|=|sK''|.
\]

% scan: s208R
% !PRINTED 391
\weberpage{391}
Hence
\[
|L'-(K_1+\ldotp\ldotp\ldotp\ldotp+K_r)|\quad \text{and}\quad |K''|
\]
will have the same characters (degree \(8\pi-8\) and genus \(5\pi-4\)) and therefore the same virtual dimension \(3\pi-4\), which is the dimension of the series cut by them on a \(K\). One deduces that in the continuous system \(\infty^{3\pi-3}\{L'-(K_1+\ldotp\ldotp\ldotp\ldotp+K_r)\}\) there is a curve split into a \(K\) and a residual curve \(D\), which (like the parabicanonical \(C = |K''-\bar K|\)) will be elliptic, of order \(4\pi-4\); and \(sD\), like \(sC\), will be a para-\(2s\)-canonical curve, the two curves cutting the \(K\) in equivalent (\(2s\)-canonical) groups. %

When the hypotheses II and I are not fulfilled, the preceding argument can nevertheless be repeated with some modification.

\textit{If hypothesis II does not hold}, that is, if one cannot determine uniquely a group of the non-special series \(\bar g=g_{2\pi-2}^{\pi-2}\), \(s\)-th part of the \(s\)-canonical series, which is supposed rationally given according to hypothesis I, it will be convenient to consider on \(K\), in place of the non-special series \(\bar g\), \(s\)-th part of the \(s\)-canonical series, the series
\[
\bar{\bar g}=(2\pi-2)\bar g-(\pi-2)g,
\]
of which a group is defined rationally by taking the sum of the groups \(\bar G_{2\pi-2}\) of \(\bar g\) which have as a \((\pi-2)\)-fold point a point of a (canonical) group \(G_{2\pi-2}\) of \(g\), and then removing from it the \(G_{2\pi-2}\) counted \(\pi-2\) times: it is essential to note that \(\bar{\bar g}\), whose \(s\)-fold multiple is an \(s\pi\)-canonical series, is not a \(\pi\)-canonical series, because otherwise from the relations
\[
(2\pi-2)\bar g-(\pi-2)g=\pi g
\]
that is
\[
(2\pi-2)\bar g=(2\pi-2)g
\]
and
\[
s\bar g=sg,
\]
since \(2\pi-2\) and \(s\) are numbers prime to each other, one would obtain
\[
\bar g=g\origfootmark{1}{s208R-1}.
\]\origfoottext{1}{s208R-1}{It suffices to solve the equation of indeterminate analysis \((2\pi-2)x-sy=1\); one will then deduce
\[
[(2\pi-2)x-sy]\bar g=[(2\pi-2)x-sy]g.
\]
}

Now, in place of \(|L|\), or rather of \(|L'|\), we shall construct a sufficiently large linear system \(|M|\) which cuts on each \(K\) the series \(\bar{\bar g}\) (constructed as a function of the \(\bar g\) which is supposed rationally given). Let us compare \(|sM|\) with the linear system \(|sK^\pi|\), multiple of \(|K^\pi|\),
% scan: s209L
% !PRINTED 392
\weberpage{392}
the \(\pi\)-th adjoint of a \(K\); denoting by \(K_0,\,K_1\,\ldotp\ldotp\ldotp\ldotp\,K_r\) particular curves of the pencil \((K)\), we shall have
\[
|sM|=|sK^\pi+s(K_1+\ldotp\ldotp\ldotp\ldotp+K_r)+tK_0|,
\]
where \(t=0\), and therefore
\[
|sM|=|sK^\pi+s(K_1+\ldotp\ldotp\ldotp\ldotp+K_r)|.
\]

From this one deduces that the linear system
\[
|M-(K_1+\ldotp\ldotp\ldotp\ldotp+K_r)|
\]
has the same characters — degree, genus and virtual dimension — as \(|K^\pi|\), and therefore belongs to a continuous system whose dimension exceeds by 1 that of the series \(\bar{\bar g}\), within which there will therefore be a curve split into a \(K\) and a residual elliptic curve \(D\), of order \(\pi(2\pi-2)\), whose multiple by \(s\) will constitute a para-\(s\pi\)-canonical curve.

Finally let us discuss what can be done if for the surface \textit{hypothesis I does not hold}: this means that on a variable curve \(K\) we are not able to determine rationally \textit{one non-special series} \(\bar g\), among those whose \(s\)-fold multiple constitutes the \(s\)-canonical series, but only a group of \(m\) \((>1)\) series \(\bar g\), interchangeable among themselves by a circuit of \(K\) in the pencil \((K)\). As we have already seen in the first case (\(\pi\) even and \(s=2\)), it is then convenient to pass from the surface \(F\) to an \(F^*\), with the same characters, represented on the \(m\)-fold \(F\) without branch curve, for which \(F^*\) hypothesis II is verified. Then one will construct on \(F^*\) an elliptic curve \(D^*\), of a certain order \((2\pi-2)\), or more generally \(\pi(2\pi-2)\), to which there will correspond on \(F\) a \(D\), of order \(m\pi(2\pi-2)\), likewise elliptic: and the curve \(D\), the \(s\)-th part of a para-\(sm\pi\)-canonical curve, will certainly contain some component distinct from those of a para-\(m\pi\)-canonical \(C\), constructed on \(F\) according to \S~5.

\edsection{9. Lemma III.}{9. Lemma III.}{sec.X.9}

\textit{On a surface \(F\), of genera \(p_g=0\) and \(p_a=-1\), possessing an elliptic pencil of curves of genus \(\pi>1\), let there be a curve \(E\) which cuts the \(K\) in groups of points equivalent to those cut by a para-\(i\)-canonical curve \(C\); if, apart from possible components in common with \(C\), the curve \(E\) has no points in common with \(C\), it is, like \(C\), an elliptic para-\(i\)-canonical curve.}

Compare \(E\) and \(C\) with respect to the second criterion of equivalence (made precise here by the lemma of \S~3): they will be curves equi%
% scan: s209R
% !PRINTED 393
\unskip\weberpage{393}%
valent up to the difference \(K_1-K_2\) of two curves of the pencil \((K)\):
\[
|E|=|C+K_1-K_2|.
\]

Hence \(E\) will have the same genus \textit{one} and the same degree \textit{zero} as \(C\), and will be, like it, a para-\(i\)-canonical curve.

\edsection{10. Conclusion: the surfaces \(F\) also possess a linear pencil of elliptic curves.}{10. Conclusion: the surfaces F also possess a linear pencil of elliptic curves.}{sec.X.10}

Lemmas I, II and III, established in the preceding sections, now lead us to a theorem which concludes the analysis of our surfaces \(F\).

\textit{The surfaces \(F\), of genus \(p_g=0\) and \(p_a=-1\), having an elliptic pencil of curves \(K\) of genus \(\pi>1\), also possess a linear pencil, without base points, of elliptic curves.}

Indeed, let us construct on \(F\) (lemma II) two elliptic para-\(i\)-canonical curves \(C\) and \(D\), cutting the \(K\) in equivalent groups of \(n=i(2\pi-2)\) points. On each \(K\) the said curves determine a linear series \(g_n^1\), which will possess a certain number
\[
m\leq2n+2\pi-2
\]
of groups \(G_n\) endowed with a double point, apart from those which may be absorbed by the section groups of \(C\) and of \(D\). The locus of the said groups \(G_n\) is a curve \(E\) which has (with respect to the \(K\)) the same order \(mn\) as \(mC\) and \(mD\), and cuts the \(K\) in groups equivalent to the sections of these curves; moreover, as \(K\) varies in the pencil \((K)\), it can never happen that a group \(G_n\) comes to coincide with the section group of \(C\), because this would entail the existence of double points for the elliptic involution cut by the \(K\) on an elliptic component of the said \(C\). Hence the curve \(E\) has zero intersections with \(mC\), and consequently it too is a para-\(mi\)-canonical curve, like \(mC\) and \(mD\).

This being so, let \(F\) be represented on an \(n\)-fold elliptic ruled surface, \(\Phi\), in such a way that the generators correspond to the \(K\), each of them representing the \(g_n^1\) defined, as above, by the sections of \(C\) and \(D\). The surface \(\Phi\) can be realized as a ruled surface of a certain space \(S_r\), on which there are two minimal directrices corresponding to \(C\) and \(D\), lying in two complementary spaces \(\gamma_1\) and \(\gamma_2\), of \(s\) and \(r-s-1\) dimensions; for example one can map \(\Phi\) onto a cone whose vertex corresponds to \(C\) and of which \(D\) is a hyperplane section. And on \(\Phi\) there will be an elliptic curve \(E\), having no points in common with \(\gamma_1\) and \(\gamma_2\). Now the \(\infty^1\) homographies of \(S_r\) which have \(\gamma_1\) and \(\gamma_2\) as spaces of united points will transform the image of \(E\) into a rational \(\infty^1\) system of elliptic curves of degree 0, and therefore into a linear pencil of elliptic curves, without base points;
% scan: s210L
% !PRINTED 394
\weberpage{394}
to this pencil of the \(n\)-fold \(\Phi\) there will similarly correspond on \(F\) a linear pencil of elliptic curves \(C\), without base points, whose existence is thus proved, \hfill Q.E.D.

It is not excluded that the \(C\) are reducible: in any case their variable \textit{irreducible} components will form a \textit{linear pencil, which we again denote by} \(|C|\).

Let us add that the computation of the Zeuthen-Segre invariant for this irreducible linear pencil \(|C|\) gives us
\[
\delta-4=12p_a-p^{(1)}+9=-4,
\]
and hence \(\delta=0\): this means that the elliptic curves \(C\), which, varying in \(|C|\), never acquire a double point, all have the same modulus, that is, they are birationally identical.

\edsection{11. Surfaces of genera \(p_g=0\) and \(p_a=-1\) with an elliptic pencil of curves of genus \(\pi=1\).}{11. Surfaces of genera p\_g=0 and p\_a=-1 with an elliptic pencil of curves of genus π=1.}{sec.X.11}

The theorem which asserts the existence of a linear pencil of elliptic curves on the surfaces with \(p_g=0\) and \(p_a=-1\) possessing an elliptic pencil of curves \(K\) of genus \(\pi>1\) now extends to the surfaces \(F\) with the same characters (\(p_g=0,\ p_a=-1\) and \(p^{(1)}=1\)) which possess an elliptic pencil of curves \(K\) of genus \(\pi=1\).

It is useful to start from the observation that, since the Zeuthen-Segre invariant is
\[
\Delta-4=12p_a-p^{(1)}+9=-4,
\]
we must have
\[
\Delta=0;
\]
and hence the elliptic \(K\), which, varying in \((K)\), never acquire a double point, are birationally identical with one another.

Let us transform \(F\), on which the \(K\) are realized projectively as curves of order \(m\), into a surface of a suitable hyperspace, which we again call \(F\), in such a way that the \(K\) become normal elliptic curves \(K_m\), of the same order \(m\), each lying in a space \(S_{m-1}\): each \(K\) of the original \(F\) will be, in general, a projection of the corresponding \(K_m\). Now two \(K_m\), normal elliptic curves with the same modulus, will be projective to each other in a certain number \(n\) of different ways: precisely, \(n=2m^2\) if the said curves are of general modulus, and instead \(n=4m^2\) or resp. \(n=6m^2\), if they are harmonic or equianharmonic. Thus between the two aforesaid \(K_m\) there is rationally defined the set \(\Gamma_n\) of \(n\) hyperspatial homographies; keeping fixed one of the two \(K_m\), and on it a point \(P\), there is then defined a curve \(C\) which passes through \(P\) and meets the other \(K_m\),
% scan: s210R
% !PRINTED 395
\weberpage{395}
the variable one, in the \(n\) points corresponding to \(P\) with respect to the \(n\) homographies of \(\Gamma_n\).

It seems, at first sight, that — the \(K_m\) being identical with one another — the aforesaid homographies always remain distinct from one another and make correspond to a point \(P\), fixed as above, \(n\) points that are likewise distinct; thus the curve \(C\), containing an elliptic involution \(\gamma_n^1\) without double points, would at once turn out to be elliptic and, as \(P\) varies, would give us a linear pencil \(|C|\) of elliptic curves, directrices of the elliptic pencil.

These curves may be reducible; in that case we shall again denote by \(C\) their irreducible components, which are still elliptic, without common points, and we shall likewise again denote by \(n\) the number of their intersections with the \(K=K_m\). We shall thus have, on the surface of irregularity \(p_g-p_a=1\), possessing the elliptic pencil of elliptic curves \((K)\), another pencil \(|C|\) of elliptic curves, directrices of the \(K\), a pencil which is necessarily rational\origfootnote{1}{s210R-1}{Cf. chap. IX, \S~11.}.

But to the argument made to prove that the \(C\) are of genus one, an objection can be raised which falls only in the face of a thorough examination, in which account is taken of the fact that we are dealing with surfaces of genus \(p_a=-1\). Indeed it is not excluded a priori that the homographies existing between two variable \(K_m\) come to degenerate, and so — supposing for example \(n=6\) — that there is a pencil of sextics \(K_6\), normal, in \(S_5\), within which there are sextics degenerating into double plane cubics \(K_3^2\): then between a generic \(K_6\) and the \(K_3\) there will exist, no longer 36, but 18 degenerate homographies, which yield as many \([1,\,2]\) correspondences between the two curves. %

Nor is it excluded a priori that among the \(K_m\) of an elliptic pencil, generally projective to one another, there are curves endowed with a cusp, or split into several parts, simple or multiple; but this case cannot arise on our surface of genus \(p_a=-1\), (and \(p^{(1)}=1\)) where we have seen that the number of the \(K_m\) endowed with a double point must be
\[
\Delta=13+12p_a-p^{(1)}=0.
\]

Since \(p_a=-1\), it is therefore only a matter of removing the exception in which the pencil \((K_m)\) contains multiple elliptic curves which, with regard to the Zeuthen Segre invariant, have valence zero in the computation of \(\Delta\) (cf. \S~3).

For this purpose observe that, in relation to the two pencils \((K)\) and \(|C|\), our surface \(F\) can be represented on an \(n\)-fold cubic cylinder, \(\Phi\), where the \(C\) correspond to the normal plane sections \(c\) of \(\Phi\) and the \(K\) to its generating lines: there will be on \(\Phi\) a branch cur%

% scan: s211L
% !PRINTED 396
\unskip\weberpage{396}%
formed of a certain number of plane sections \(c\) corresponding to multiple curves \(C=sC_s\) and (possibly) also of a certain number of generators \(k\), corresponding to curves \(K\) that degenerate into multiple elliptic curves \(K=rK_r\). If \(p_a=-1\), and hence (for the \(K\)) \(\Delta=0\), the multiple cylinder \(\Phi\) provides the representation of a surface \(F\), which may be supposed free of singularities, on which a \(C_s\) and a \(K_r\) meet in \(\frac{n}{rs}\) points (simple points of the surface), which correspond to the point \(A=ck\) of \(\Phi\): indeed, if some one, \(A'\), of the said points of intersection of \(C_s\) and \(K_r\) were double or multiple for \(F\), the degenerate limit curve of a \(K\) would be, not (as has been supposed) \(rK_r\), but \(rK_r\) augmented by the neighbourhood of \(A'\), so that (\S\ 3) it would follow that \(\Delta>0\). But now one proves, on the contrary, that if to the point \(A\) of the \(n\)-fold cylinder \(\Phi\) there corresponds a point of the surface \(F\), this must be a singular (double or multiple) point of \(F\), so that the hypothesis of a \(K\) degenerating into \(rK_r\) is reduced to an absurdity. Indeed, consider a plane \(\bar{\alpha}\), oblique to the generators of \(\Phi\), which passes through \(A\), and a variable plane \(\alpha\) which does not pass through \(A\) but tends to \(\bar{\alpha}\). The section of \(\Phi\) by the plane \(\alpha\) is a multiple curve \(l\), \(n\)-fold, on which there are two branch points \(A_1\) and \(A_2\), intersections of \(c\) and \(k\). On the Riemann surface of the multiple \(l\) let us describe a closed path \(\nu\) encircling \(A_1\) and \(A_2\), which may be regarded as the sum of two elementary loops encircling \(A_1\) and \(A_2\). To the first loop there corresponds a substitution on the branches equivalent to the product of %
\[
(s-1)\frac{n}{s}
\]
transpositions, and to the second a substitution equivalent to
\[
(r-1)\frac{n}{r}
\]
transpositions. Hence to the closed path \(\gamma\), which in the limit --- for \(\alpha\to\bar{\alpha}\) --- becomes a loop encircling \(A\), there corresponds a substitution formed of at least \(n\) transpositions: this means that the curve \(L\) homologous to \(l\) cannot have simple points corresponding to \(A\), since in the most unfavourable case this curve would give rise to a branch of order \(n\), to which there would correspond, on the Riemann surface \(l\) with \(n\) sheets, a substitution on the branches formed of \(n-1\) transpositions!

We therefore conclude:

\textit{The surfaces \(F\) of genera \(p_g=0\) and \(p_a=-1\) which possess an elliptic pencil of elliptic curves \(K\) also possess a linear pencil of elliptic curves \(C\), directrices of the \(K\).}

\textit{Note.} -- In the reasoning developed above use has been made of the fruitful concept of the singularities of curves set forth by \textsc{O. Chisini}
% scan: s211R
% !PRINTED 397
\weberpage{397}
in the Memoir ``Le singolarità di un ramo superlineare di curva piana definite mediante un prodotto di sostituzioni''\origfootnote{1}{s211R-1}{Atti del R. Istituto Veneto di Scienze, Lettere ed Arti, t. XXX, part \(2^{\mathrm{a}}\), 1921. This memoir, which is later than the publication of the second volume of the \textit{Lezioni} of \textsc{Enriques-Chisini}, brings a further important contribution to the theory developed there.}.

It is worth pointing out explicitly that, if the hypothesis \(p_a=-1\) is dropped, there are in fact surfaces possessing an elliptic pencil of elliptic curves \(K\) to which reducible curves belong: namely surfaces \(F\) for which \(p_a\geq0\) and \(p_g=p_a+1\), containing a linear pencil of curves \(C\), directrices of \((K)\), of genus \(\pi>1\). If we give the name \textit{elliptic} (for a reason that will be explained later) to the surfaces with \(p_g=0\) and \(p_a=-1\) studied in the preceding sections, these new surfaces \(F\), which present some notable analogies with them, may be called \textit{paraelliptic}. The principal analogy is that on the paraelliptic surfaces, as on the elliptic ones with \(p_a=-1\), there are \textit{two pencils of birationally identical curves}: \textit{an elliptic pencil} of elliptic curves \((K)\) and \textit{a linear pencil} of curves which here are not elliptic but of genus \(\pi>1\). However, whereas the elliptic surfaces admit --- as we shall see --- a group \(\infty^1\) of birational transformations into themselves, this property does not hold for the paraelliptic surfaces.

We confine ourselves here to indicating two examples of paraelliptic surfaces, constructed starting from a multiple cubic cylinder \(f(xy)=0\):
\weq{1)}{\begin{aligned}
&u^2=(z-a)(z-b)(z-c)(z-d)\,\psi(xy)\\
&\qquad f(xy)=0\quad(p_a=3,\quad p_g=4);
\end{aligned}}
\(\psi=0\) represents a quadric cylinder parallel to \(f\), which cuts \(f=0\) in 6 branch generators, to be added to the 4 branch cubics
\[
z=a,\quad z=b,\quad z=c,\quad z=d;
\]

\weq{2)}{\begin{aligned}
&u^3=(z-a)(z-b)(z-c)\,\psi(xy)\\
&f(xy)=0\quad(p_a=0,\quad p_g=1);
\end{aligned}}
\(\psi=0\) denotes a cubic cylinder which has with \(f\) a three-point contact along three generators, and touches it along another generator \(k\), cutting it further in a generator \(k'\): the branch curve of the triple cylinder \(f\) therefore consists of the three cubics %
\[
z=a,\quad z=b,\quad z=c,
\]
and of the two generators
\[
k\qquad\text{and}\qquad k'.
\]
Finally, we point out to young researchers this interesting subject
% scan: s212L
% !PRINTED 398
\weberpage{398}
of study: \textit{to determine and classify all paraelliptic surfaces}; and we refer to the Notes on pp. 123 and 134 of the ``Lezioni'' of \textsc{Enriques-Campedelli} in the Rendiconti del Seminario Matematico della R. Università di Roma (1934): from these the reader will be able to gather some useful observations and suggestions.

\edsection{12. Elliptic surfaces.}{12. Elliptic surfaces.}{sec.X.12}

Let us summarize the results obtained in the preceding sections: \textit{Every irregular surface \(F\) of genus \(p_g=0\) which does not belong to the family of ruled surfaces} (rational or not) \textit{has arithmetic genus \(p_a=-1\) and contains an elliptic pencil \((K)\) of curves of genus \(\pi\geq1\), and a linear pencil of elliptic curves \(|C|\)}; hence it can be represented on a \textit{cubic cylinder \(\Phi\) of multiplicity equal to the number \(n=CK\), with a branch curve consisting of plane sections normal to \(\Phi\).}

Conversely, an \textit{\(n\)-fold cubic cylinder \(\Phi\)}, with a branch curve consisting of normal plane sections (fixed so as to satisfy the conditions for the existence of an irreducible surface \(F\)), does \textit{not} represent \textit{in all cases a surface of genus \(p_g=0\)}, because the images of the plane sections of \(\Phi\) may be reducible elliptic curves whose components constitute a pencil \((C)\), not linear, \textit{of genus \(\varrho\geq1\)}.

Here it appears that our surfaces, with \(p_g=0\) and \(p_a=-1\), not birationally equivalent to ruled surfaces, fall within a \textit{more general family of surfaces \(F\)}, with \(p_a=-1\), which possess an \textit{elliptic pencil of curves \(K\)} of genus \(\pi\geq1\) and a \textit{pencil of genus \(\varrho\geq0\) of elliptic curves \(C\)}, for which surfaces \(F\) \textit{the geometric genus will be \(p_g=\varrho\)}.

Indeed, the existence of the two pencils \((C)\) and \((K)\) which characterize the said more general surfaces \(F\) implies, since \(p_a=-1\), that the irregularity of \(F\) is
\[
q\geq\varrho+1
\]
and hence
\[
p_g\geq\varrho;
\]
but, on the other hand, the canonical curves of \(F\), formed of curves \(C\), must cut every \(K\) in groups of the \(g_{2\pi-2}^{\pi-1}\) composed with the involution of genus \(\varrho\) cut by \((C)\), which is free of double points\origfootnote{1}{s212L-1}{Theorem of Painlevé-Castelnuovo. Cf. \textsc{Enriques-Chisini}, \textit{Lezioni}. Book V, chap. I, \S\ 9 (vol. III, p. 72).}; and therefore they must belong to the canonical series \(g_{2\varrho-2}^{\varrho-1}\) within the pencil \((C)\). There will therefore result:
\[
p_g\leq\varrho
\]
and consequently
\[
p_g=\varrho,
\]
\hfill Q.E.D.

% scan: s212R
% !PRINTED 399
\weberpage{399}
The more general surfaces \(F\), characterized --- as has been said above --- by having arithmetic genus \(p_a=-1\) and by possessing two pencils \((K)\) and \((C)\) of irreducible curves, namely

1) an elliptic pencil of curves \(K\) of a certain genus \(\pi\geq0\) (for which, from \(p_a=-1\), there follows \(\Delta=0\));

2) and a pencil of genus \(p_g\) of elliptic curves \(C\), intersecting one another in \(n=CK\geq1\) points, will be called \textit{elliptic surfaces}: and, more precisely, \textit{proper} elliptic surfaces those, for \(\pi\geq1\), which are not birationally equivalent to ruled surfaces, and \textit{improper} those, for \(\pi=0\), which reduce to ruled surfaces of genus 1.

The said name ``elliptic'' derives from a parametric representation of the surfaces in question, which is connected with the group \(\infty^1\) of their transformations into themselves, of which we speak later, and which was pointed out by \textsc{P. Painlevé}\origfootnote{1}{s212R-1}{\textsc{P. Painlevé}, \textit{Leçons sur la théorie analytique des équations différentielles}. (Paris, Hermann, 1897).}.

With regard to this parametric representation the number \(n=CK\) receives the name of \textit{determinant} of the surface.

The parametric representation mentioned above is translated into the following geometric property: \textit{Every elliptic surface \(F\) admits an elliptic group \(\infty^1\) of birational transformations into itself, which has as trajectories the curves of the pencil \((C)\).}

Let us prove the theorem by showing that there exists a birational transformation of \(F\) --- which leaves all the curves \(C\) invariant --- in which two points \(P_o\) and \(P'_o\), arbitrarily chosen on one of these curves, \(C_o\), correspond to each other. %

Indeed, consider another curve \(C\) of the same pencil \((C)\) and choose a point \(P\) on it. If the curves \(C\), birationally identical to one another, are elliptic curves of general modulus, there are --- as is well known --- \textit{two} transformations of \(C_o\) into \(C\) which change \(P_o\) into \(P\)\origfootnote{2}{s212R-2}{Cf. e.g. \textsc{Enriques-Chisini}, \textit{Lezioni}. Book V, chap. III, \S\ 27 (vol. III, p. 262).}; and among these transformations only one\origfootnote{3}{s212R-3}{Cf. \textsc{Enriques-Chisini}, loc. cit., chap. IV, \S\ 39 (vol. III, p. 446).} (rationally determined) makes correspond to each other the groups cut by the curves \(K\), which form on each of them an elliptic involution \(\gamma_n^1\). Therefore, the transformation of the first kind \((P_oP'_o)\) being given on \(C_o\), there is also rationally determined along with it \textit{one} transformation of the first kind of \(C\), in which to a point \(P\) there will correspond a point \(P'\), in such a way that \((PP')\) will reduce to \((P_oP'_o)\) when \(C\) reduces to \(C_o\). %

The proof extends without difficulty to the cases in which the curves \(C\) are harmonic or equianharmonic elliptic curves.

% scan: s213L
% !PRINTED 400
\weberpage{400}
It is also worth observing that the property established for elliptic surfaces is characteristic:

\textit{Every surface \(F\) which admits an elliptic group \(\infty^1\) \(\Gamma\) of birational transformations into itself is an elliptic surface, of arithmetic genus \(p_a=-1\), proper or improper, possessing an elliptic pencil of curves \(K\) of genus \(\pi(\geq0)\) and a pencil of genus \(p_g\) of elliptic curves \(C\).}

Indeed, \(F\) must possess in the first place a pencil \((C)\), of a certain genus \(\varrho\geq0\), formed by the trajectories of the group \(\Gamma\), which will be elliptic curves identical with the group \(\Gamma\) itself. But, in the second place, an elliptic pencil \((K)\) will be constructed on \(F\) by the following construction: consider on \(F\) an (irreducible) curve \(L\) and the elliptic series \(\infty^1\) \(\{L\}\) consisting of its transforms by the transformations of the group \(\Gamma\); the curves \(L\) of the said series will certainly be non-equivalent, since they must cut on a \(K\) non-equivalent groups of points, which correspond to one another in a transformation of the first kind; however, among the groups of curves \(L\) passing through a point \(P\) variable on \(F\), there will be \(\infty^1\) that are equivalent to one another: this means that there is on \(F\) an elliptic pencil of curves on which the \(L\) of the series \(\{L\}\) cut equivalent groups. (Cf. Chap. IX, \S\ 11).

Having thus constructed the two pencils \((C)\) and \((K)\) on \(F\), one readily determines the characters of the surface. First of all, the computation of the Zeuthen-Segre invariant for \((K)\) gives us
\[
\Delta=13-12p_a-p^{(1)}\geq0 %
\]
and hence
\[
p_a\geq-1\qquad\qquad(p^{(1)}\geq1).
\]

On the other hand, the presence of an elliptic pencil \((K)\) and of a pencil \((C)\) of genus \(\varrho\) implies that (there exists on \(F\) a series of \(\infty^{\varrho+1}\) non-equivalent curves, and hence that)
\[
\begin{aligned}
&p_g-p_a\geq\varrho+1\\
&p_g\geq\varrho.
\end{aligned}
\]

Finally, a consideration already made above serves to establish that the canonical curves of \(F\) --- which for \(\varrho>0\) are formed of \(2\varrho-2\) curves \(C\) --- must give within the pencil \((C)\) a series contained in the canonical series \(g_{2\varrho-2}^{\varrho-1}\); from the comparison there follows
\[
p_g=\varrho,\qquad p_a=-1.
\]
\hfill Q.E.D.

\textit{Remark.} -- From what we have said one deduces that on an elliptic surface \(F\) of determinant \(n>1\) there is an \textit{involution, \(I_n\), formed by the groups of \(n\) points common to the \(C\) and \(K\), which is}
% scan: s213R
% !PRINTED 401
\weberpage{401}
\textit{generated by an abelian group \(\Gamma_n\) of birational transformations} (permutable), namely by the transformations of the group \(\Gamma\) of the surface which leave the \(C\) and the \(K\) fixed. If the groups of points of \(I_n\) are taken as elements (points) of a new surface \(f\), one obtains an elliptic surface of determinant 1, on which \(F\) is represented in a multiple way, with multiplicity \(n\); when \(p_g=0\), this improper elliptic surface reduces to an elliptic ruled surface, and, if one wishes, to a cubic cylinder \(f(xy)=0\).

Here it is worth pointing out that the \textit{elliptic surfaces \(f\) of determinant 1} are defined in general as representative of the \textit{variety \(\infty^2\) of the pairs of points belonging to an elliptic curve \(C\) and to a curve \(K\) of genus \(\pi\) \((\geq0)\).}

Starting from these surfaces, by means of the introduction of algebraic irrationalities suitable for resolving the abelian group \(\Gamma_n\), one will be able to construct the elliptic surfaces of determinant \(n>1\). For this purpose one can say, a priori, that there will be needed \textit{a single radical} or \textit{two non-superposed radicals}, according as the said group \(\Gamma_n\) is \textit{cyclic or abelian with base two}. For \(\Gamma_n\), which is a finite group of transformations into itself of an elliptic curve \(C\), cannot have a base \(>2\)\origfootnote{1}{s213R-1}{Cf. \textsc{Enriques-Chisini}, \textit{Lezioni}. Book V, chap. III, \S\ 28 (vol. III, p. 271) and chap. IV, \S\ 39 (vol. III, pp. 446-451).}. In the following section, confining ourselves in particular to the hypothesis \(p_g=0\), and then again, more precisely, in \S\ 16, we explain the meaning of this construction, which serves to establish the \textit{existence} of elliptic surfaces, and the conditions relating to it.

\edsection{13. Construction of the elliptic surfaces of genus \(p_g=0\).}{13. Construction of the elliptic surfaces of genus p\_g=0.}{sec.X.13}

To construct the (proper) elliptic surfaces \(F\) of genus \(p_g=0\), having any determinant \(n>1\), we start from the cubic cylinder \(f(xy)=0\), which realizes --- as has been said --- the (improper) elliptic surface of genus \(p_g=0\) and of determinant 1.

By means of the two pencils of irreducible curves \(|C|\) and \((K)\) which characterize it, the surface \(F\) can be represented on an \(n\)-fold cylinder \(f\), with a branch curve composed of a certain number of normal plane sections \(c\), in such a way that the \(K\) have as images the generators \(k\) of the cylinder, and the \(C\) (trajectories of the group \(\infty^1\) of transformations of \(F\)) correspond to the normal plane sections \(c\) of the same cylinder. The branch curve of the \(n\)-fold cylinder will consist of a certain number of curves \(c\), to which there will correspond multiple curves of the type \(C=sC_s\).

\printersignature{\textsc{Enriques F.} - \textit{Superficie algebriche.}\quad 26}

% scan: s214L
% !PRINTED 402
\weberpage{402}
Hence the representation of the elliptic surface \(F\) on the \(n\)-fold cylinder \(f\) must satisfy the following conditions:

1) that the curve \(K\) defined by an \(n\)-fold \(k\) be irreducible and contain an involution \(g_n\) of the cyclic type or of the abelian type (properly so called) with base 2;

2) that the curves \(C\) corresponding to the normal plane sections \(c\) of \(f\), \(n\)-fold curves without branch points, be in general irreducible.

We observe indeed that, if this condition 2) is not satisfied, the surface \(F\) represented by the \(n\)-fold cylinder will no longer be an (elliptic) surface of determinant \(n\), but will have a determinant \(<n\), and will turn out, in general, to be of genus \(p_g>0\). This happens, for example, for the surface
\[
\left\{
\begin{aligned}
u^2&=(z-a)(z-b)(z-e)(z-d)\,\psi^2(xy)\\
f(xy)&=0
\end{aligned}
\right.
\]
which is an irreducible surface \(F\) represented on the double cylinder, with branch curve composed of the plane sections \(z=a\), \(z=b\), \(z=d\), \(z=e\).

Indeed, to a generic normal section \(c\):
\[
z=c\qquad(c=\text{const.})
\]
there corresponds on the said \(F\) a pair of distinct elliptic curves
\[
u=\pm\sqrt{(c-a)(c-b)(c-d)(c-e)}\,\psi,
\]
so that the irreducible curves \(C\) which compose the images of the \(c\) now form a pencil (no longer linear, but elliptic) of curves unisecant to the \(K\); and therefore the geometric genus of \(F\) is:
\[
p_g=1.
\]

Here it is convenient to recall the criteria for an \(n\)-fold elliptic curve, without branch points, to be irreducible\origfootnote{1}{s214L-1}{Cf. \textsc{Enriques-Chisini}, \textit{Lezioni}, Book V, chap. IV, \S\ 39 (vol. III, p. 446).}. There are two possible cases: the \textit{cyclic} case and the \textit{abelian} case (base 2). In the first case the \(n\)-fold curve is constructed, starting e.g. from a cubic \(f(xy)=0\), by extracting on it an \(n\)-th radical bearing on a polynomial \(\varphi(xy)=0\),
\[
\left\{
\begin{aligned}
u^n&=\varphi(xy)\\
f(xy)&=0;
\end{aligned}
\right.
\]
in order that the multiple curve have no branch points, it is necessary that the \textit{critical points}
\[
f=0,\qquad\varphi=0,
\]
% scan: s214R
% !PRINTED 403
\weberpage{403}
be \textit{apparent}, and therefore that \(\varphi\), taken e.g. of order \(n\), have three \(n\)-point contacts with \(f\). Hence the \textit{irreducibility} of the said multiple curve depends on the fact that ``the series \(g_3^2\) to which the triad of apparent critical points belongs is different from the series cut on \(f\) by the lines of the plane''.

In the abelian case (in which \(n=pq\)) the construction of the \(n\)-fold curve requires the extraction of two non-superposed radicals of index \(p\) and \(q\), bearing respectively on two polynomials \(\varphi(xy)\) and \(\psi(xy)\), which may legitimately be supposed of order \(p\) and \(q\): in order that the multiple curve have no branch points, the curves \(\varphi\) and \(\psi\) must have with \(f\) three points of contact of order \(p-1\) and \(q-1\) respectively, which thus form two triads of apparent critical points; the irreducibility of the said multiple curve requires that the said triads do not belong to the \(g_3^2\) cut on \(f\) by the lines of the plane, and moreover that they satisfy a certain condition of dissimilarity.

Let us develop the construction indicated in the \textit{cyclic case}, which is met with in particular when the determinant \(n\) is a prime number. We have then that on the surface \(F\) the involution \(I_n\) formed by the groups \(CK\) is generated by a cyclic group of transformations \(\Gamma_n\), of order \(n\).

Starting from the \(n\)-fold cubic cylinder
\[
f(xy)=0,
\]
let us fix the normal plane branch sections
\[
z=a_1,\quad z=a_2,\quad \ldotp\ldotp\ldotp\ldotp\quad z=a_t;
\]
the \(n\) points of the group \(G_n\) which corresponds to a generic point of \(f\) can be given by the values of a radical
\[
u=\sqrt[n]{(z-a_1)^{r_1}(z-a_2)^{r_2}\,\ldotp\ldotp\ldotp\ldotp\,(z-a_t)^{r_t}\cdot\varphi(xy)},
\]
satisfying the following conditions:

1) For the exponents \(r_i<n\) one will have
\[
r_1+r_2+\ldotp\ldotp\ldotp\ldotp+r_t\equiv0\qquad(\mathrm{mod.}\ n),
\]
otherwise a branch point at infinity would arise on the generators of the cylinder \(f\).

2) If \(h_i\) is the greatest common divisor of \(n\) and \(r_i\), so that
\[
n=h_i s_i,\qquad r_i=h_i\varrho_i. %
\]
in the expression of \(u\) there appears the factor
\[
\sqrt[n]{(z-a_i)^{r_i}}=\sqrt[s_i]{\varepsilon_i(z-a_i)^{\varrho_i}},
\]
where \(\varepsilon_i\) is an \(h_i\)-th root of unity.

% scan: s215L
% !PRINTED 404
\weberpage{404}
Thus on the \(n\)-fold generator of the cylinder \(f\), the point
\[
z=a_i
\]
is a branch point to which there corresponds on the homologous \(K\) a group \(G_n\) (of the \(g_n\) cut by the \(C\)) consisting of \(h_i\) groups of \(s_i\) coincident points. In other words, the \textit{point \(z=a_i\)} is for the multiple \(k\) a \textit{branch point of order \(s_i\)}, with
\[
s_i=\frac{n}{h_i},
\]
and hence one has
\[
\sum\frac{n}{s}(s-1)=2n+2\pi-2
\]
(the sum being extended to the values of \(s=s_1,\ s_2\,\ldotp\ldotp\ldotp\ldotp\,s_t\)), \(\pi\) \((\geq1)\) denoting the genus of the \(K\).

3) Each of the numbers \(s_i(i=1,2,\ \ldotp\ldotp\ldotp\ldotp\ t)\) is a divisor of the least common multiple of the remaining ones.

This property follows in general from the abelian character of \(\Gamma_n\). Indeed, on the multiple line \(k\), there being no branchings at infinity, the product of the cyclic substitutions on the branches of the algebraic function \(u(xyz)\), relative to the branch points \(a_i\), is equivalent to the identity:
\[
S_t\,\ldotp\ldotp\ldotp\ldotp\,S_2S_1=1
\]
so that, for example,
\[
S_1^{-1}=S_t\,\ldotp\ldotp\ldotp\ldotp\,S_3S_2,
\]
that is
\[
S_1^{s_1-1}=S_t\,\ldotp\ldotp\ldotp\ldotp\,S_3S_2.
\]
But the substitution \(S_1^{s_1-1}\) is cyclic of period \(s_1\), whereas the product \(S_t\,\ldotp\ldotp\ldotp\ldotp\,S_3S_2\) is cyclic with a period which is a divisor of the least common multiple of \(s_2,s_3,\ldotp\ldotp\ldotp\ldotp,s_t\)\origfootnote{1}{s215L-1}{Cf. e.g. \textsc{L. Bianchi}, \textit{Lezioni sulla teoria dei gruppi di sostituzioni e delle equazioni algebriche secondo Galois} (Pisa, 1900), chap. III, \S\ 30.}, since one has
\[
S_t^{s_t}\,\ldotp\ldotp\ldotp\ldotp\,S_3^{s_3}S_2^{s_2}
=(S_t\,\ldotp\ldotp\ldotp\ldotp\,S_3S_2)^{s_2s_3\,\ldotp\ldotp\ldotp\ldotp\,s_t}=1.
\]

4) Since the \(n\)-fold cylinder \(f\) possesses no branch generators, the polynomial \(\varphi(xy)\) which appears under the radical, in the expression of \(u\), must represent a cylinder having with \(f\) an \(n\)-point contact along every common generator. And in order that the curves \(C\) represented by the \(z=\text{const.}\) be irreducible, it will be necessary that the group of the apparent critical generators be not equivalent to a multiple of the triads of generators cut by the planes
\[
\alpha x+\beta y+\gamma=0.
\]

% scan: s215R
% !PRINTED 405
\weberpage{405}
We may even take \(\varphi\) of order \(n\), so that \(\varphi=0\) touches \(f=0\) along three generators (not lying in a plane) with \(n\)-point contact. In fact, one easily recognizes that, if \(\varphi\) is made to vary continuously, the surface \(F\) remains birationally identical with itself, the curves \(C\) and \(K\), represented respectively by the plane sections \(z=\text{const.}\) and by the generators of \(f\), being preserved. Hence, if in place of the cylinder \(\varphi\) one takes a cylinder \(\psi\) of order \(nh\), with \(h>1\), we may substitute for \(\psi(xy)\) a polynomial split into a simple factor \(\varphi\) of order \(n\) and an \(n\)-fold factor \(\varphi_1^n\) with \(\varphi_1\) of order \(h-1\):
\[
\psi=\varphi_1^n\varphi;
\]
thus there will result
\[
u=\varphi_1\sqrt[n]{\varphi(xy)\cdot(z-a_1)^{r_1}\,\ldotp\ldotp\ldotp\ldotp\,(z-a_t)^{r_t}},
\]
from which --- changing \(u\) into \(u\varphi_1\) --- one obtains precisely the expression which it was desired to establish.

In conclusion: \textit{as type of the elliptic surfaces of determinant \(n\) belonging to the cyclic case, one has the surface}
\[
\left\{
\begin{aligned}
u&=\sqrt[n]{(z-a)^{r_1}\,\ldotp\ldotp\ldotp\ldotp\,(z-a_t)^{r_t}\cdot\varphi(xy)}\\ %
f(xy)&=0
\end{aligned}
\right.
\]
\textit{where the exponents \(r_i\) satisfy the conditions stated above and \(\varphi(xy)=0\) is a cylinder of order \(n\) which touches \(f\), with \(n\)-point contact, along three generators not belonging to a plane.}

The construction of the surfaces \(F\), starting from an \(n\)-fold cubic cylinder \(f(xy)=0\), is repeated with few modifications in the case in which the group is \textit{abelian} (with base 2). Except that, instead of a single radical, there will now be, in the \textit{expression of \(u\), two non-superposed radicals}. We shall not dwell on this construction, which we develop later (\S\ 16) with reference to the family of the surfaces \(F\) for which the \(K\) as well as the \(C\) are elliptic curves. Here it suffices to note that again, as in the cyclic case, one will have the equality
\[
\sum\frac{n(s-1)}{s}=2n-2\pi-2, %
\]
and always \textit{each of the numbers \(s_i\) must divide the least common multiple of the remaining ones.}

\edsection{14. Surfaces with vanishing plurigenera: characterization of ruled surfaces.}{14. Surfaces with vanishing plurigenera: characterization of ruled surfaces.}{sec.X.14}

We have proved that ``the irregular surfaces of geometric genus \(p_g=0\), not belonging to the family of ruled surfaces, are elliptic surfaces, of arithmetic genus \(p_a=-1\)''. We now wish to determine the plurigenera of these elliptic surfaces and to verify that these
% scan: s216L
% !PRINTED 406
\weberpage{406}
characters cannot all be zero, as happens for the ruled surfaces (rational or not). In this way, recalling Castelnuovo's conditions of rationality, given by the vanishing of the arithmetic genus and of the bigenus (Chap. VI, § 4): $p_a=P_2=0$, the family of ruled surfaces will be defined by the vanishing of the plurigenera, and more precisely by the following theorem: \textit{A necessary and sufficient condition for a surface to be birationally equivalent to a ruled surface is that its quadrigenus and its sextigenus vanish:}
\[
P_4=P_6=0.
\]

Let us consider a (properly) elliptic surface $F$, of genus $p_g=0$ and of determinant $n\geq2$, represented on the $n$-fold cubic cylinder $f(xy)=0$, in the manner defined above, in relation to its two characteristic pencils $(K)$ and $|C|$. One can determine the virtual canonical curve of $F$ by adding to the transform of the canonical curve of the cylinder $f$ the curve corresponding to the branch curve given on this multiple cylinder (cf. Chap. V, § 7).

Now the virtual canonical curve of $f$ is obtained by subtracting the linear system $2c$; consisting of two plane sections normal to the cylinder, from its adjoint system, which consists of the three lines of section with the plane at infinity; and since these lines are images of exceptional curves on $F$, one will only have to add to the transform of the branch curve the virtual curve $-2C$ which corresponds to $-2c$. Therefore, denoting by $z=a_i$ the components of the branch curve, to each of which there corresponds a multiple curve $C=s_iC_i$, the pure canonical curve of $F$ will be given by %
\[
D-2C=\Sigma(s-1)C_s-2C.
\]

Since the curves $C$ homologous to the $c$ are irreducible, one sees at once that the subtraction indicated by the preceding formula is not possible, and thus it is confirmed that the aforesaid irreducibility implies that the genus of $F$ is $p_g=0$.

Indeed, the linear system $|D|$, defined by $D=\Sigma(s-1)C_s$, has dimension $x=0$ and therefore cannot contain within itself the system $|C|$, and still less $|2C|$. To justify this assertion, note that the curve $D$, being formed of elliptic components of degree $0$, must a priori define a linear system composed of the curves of the pencil $|C|$ and with components that are aliquot parts of such curves:
\[
|D|=|\Sigma(s-1)C_s|=|xC+\Sigma h_sC_s|.
\]

From this one derives
\[
|xC|=|\Sigma k_sC_s|,
\]
with
\[
0\leq k_s<s-1,
\]
% scan: s216R
% !PRINTED 407
\weberpage{407}
and it is clear that such a relation cannot hold for $x>0$ (the generic $C$ being irreducible), because the \textit{principle of degeneration} would be contradicted, inasmuch as a variable curve of $(C)$ (irreducible and therefore connected) would come to degenerate into parts $C_s$ not connected with one another.

Now from the formula written above let us pass to the expression of the pluricanonical curves and hence to the computation of the plurigenera of $F$.

In general the $m$-canonical curves will be given by
\[
|m\Sigma(s-1)C_s-2mC|=|m(t-2)C-m\Sigma C_s|,
\]
where the sum is extended to the various values of $s$ which correspond to the $t$ multiple curves of the pencil $|C|$:
\[
s=s_1,\ s_2;\ \ldotp\ldotp\ldotp\ldotp,\ s_t. %
\]

If one denotes by $x=P_m-1$ the dimension of the complete $m$-canonical system, this system will be composed of $x$ variable (irreducible) curves $C$ and of fixed components $C_s$, each to be counted a certain number of times
\[
h<s.
\]

We shall therefore have
\[
|m(t-2)C-m\Sigma C_s|=|xC+\Sigma hC_s|,
\]
with $0\leq h<s$, the sum being extended to the values
\[
s=s_1,\ s_2,\ \ldotp\ldotp\ldotp\ldotp,\ s_t
\]
and to the corresponding values of
\[
h=h_1,\ h_2,\ \ldotp\ldotp\ldotp\ldotp,\ h_t.
\]

Now from the preceding relation, since $sC_s=C$, one derives
\[
|(m(t-2)-x)C|=\left|\Sigma\frac{m+h}{s}C\right|:
\]
this means that the fractions $\dfrac{m+h}{s}$ are equal to integers:
\[
\frac{m+h}{s}=\varrho,
\]
each of these integers ($\varrho=\varrho_1,\ \varrho_2,\ \ldotp\ldotp\ldotp\ldotp,\ \varrho_t$) expressing the value by excess of the fraction $\dfrac{m}{s}$, that is, the greatest integer
\[
\left[\frac{m+s-1}{s}\right]
\]
contained in $\dfrac{m+s-1}{s}$.

% scan: s217L
% !PRINTED 408
\weberpage{408}
In conclusion we shall have
\[
x=m(t-2)-\Sigma\left[\frac{m+s-1}{s}\right],
\]
and hence \textit{the plurigenus of order $m(\geq2)$ is given by the formula.} %
\[
P_m=1+m(t-2)-\Sigma\left[\frac{m+s-1}{s}\right],
\]
where one must take $P_m=0$ when its expression turns out to be negative.

Let us here put $m=2$, taking into account that — for $s=s_1,s_2,\ldotp\ldotp\ldotp\ldotp,s_t$ — one always has
\[
s\geq2;
\]
we shall have
\[
\left[\frac{s+1}{s}\right]=1,\qquad
\Sigma\left[\frac{s+1}{s}\right]=t,
\]
and hence \textit{the bigenus is:}
\[
P_2=t-3.
\]

Now, for $t>3$, one certainly has $P_2>0$ and therefore also $P_4>0$, $P_6>0$, etc. If instead $P_2=0$, the hypothesis $t=2$, which leads to an $F$ birationally equivalent to a ruled surface, being excluded, one will have
\[
t=3.
\]

But, denoting by $n$ the determinant of $F$, one has the equality
\[
\Sigma n(s-1)=2n+2\pi-2, %
\]
from which (since $s_i\geq2$, $\pi\geq1$)
\[
\Sigma\left(1-\frac{1}{s}\right)\geq1, %
\]
that is
\[
\begin{aligned}
&t-\Sigma\frac{1}{s}\geq1\\ %
&\Sigma\frac{1}{s}\leq t-1. %
\end{aligned}
\]

For $t=3$ one will therefore have
\[
\Sigma\frac{1}{s}=\frac{1}{s_1}+\frac{1}{s_2}+\frac{1}{s_3}\leq1.
\]

Let us suppose that the (positive integer) numbers $s_1$, $s_2$, $s_3$ are in non-decreasing order:
\[
s_3\geq s_2\geq s_1;
\]
% scan: s217R
% !PRINTED 409
\weberpage{409}
the preceding inequality gives us
\[
s_1=2,\qquad s_2=3,\qquad s_3\geq6
\]
or
\[
s_1=2\qquad s_3\geq s_2\geq4,
\]
or else
\[
s_1\geq3,\qquad s_3\geq s_2\geq3, %
\]
or finally
\[
s_3\geq s_2\geq s_1\geq4.
\]

\edFlush
In the first and in the second case:
\[
\left[\frac{s_1+2}{s_1}\right]=2,\qquad
\left[\frac{s_2+2}{s_2}\right]=\left[\frac{s_3+2}{s_3}\right]=1
\]
and hence
\[
P_3=0,
\]
while in the third and fourth case:
\[
\left[\frac{s_1+2}{s_1}\right]=\left[\frac{s_2+2}{s_2}\right]=\left[\frac{s_3+2}{s_3}\right]=1,
\]
and
\[
P_3=1.
\]

On the other hand, computing the quadrigenus
\[
P_4=5-\Sigma\left[\frac{s+3}{s}\right], %
\]
we shall have:

in the first and third case
\[
P_4=0,
\]

in the second case
\[
P_4=1
\]

and in the fourth
\[
P_4=2.
\]

And for the sextigenus
\[
P_6=7-\Sigma\left[\frac{s+5}{s}\right]
\]
one will find:

in the first and third case
\[
P_6=1,
\]

in the second case
\[
P_6=0
\]

and in the fourth
\[
P_6=2.
\]

% scan: s218L
% !PRINTED 410
\weberpage{410}
For our purpose it is sufficient to point out that when $P_2=P_4=0$ (first and third case of the preceding discussion) there results
\[
P_6\geq1.
\]

In conclusion, for an elliptic surface not birationally equivalent to a ruled surface, at least one of the two plurigenera $P_4$ or $P_6$ will be different from zero:
\[
P_4+P_6\geq1,\qquad\qquad\text{Q.E.D.}
\]

\textit{Remark. -- The conditions for a surface to be birationally equivalent to a ruled surface can also be expressed by setting equal to zero the genus of order $12$:}
\[
P_{12}=0.
\]

Indeed, for the elliptic surfaces, which have $P_4\geq1$ or $P_6\geq1$, one certainly has
\[
P_{12}\geq1.
\]

It is also worth pointing out that one will have
\[
P_{12}=1
\]
for the \textit{elliptic surfaces} (with virtual canonical curve of order zero) possessing an \textit{elliptic pencil of elliptic curves $K$}, and instead
\[
P_{12}>1
\]
for the \textit{elliptic surfaces} possessing an \textit{elliptic pencil of curves $K$ of genus $\pi>1$}, whose virtual canonical curves cut the $K$ in $2\pi-2$ points.

Indeed, the difference between the two kinds of surfaces $F$ consists in this:

in the case of the $K$ of genus $\pi=1$ one has the \textit{equality}
\[
2n=\Sigma\frac{s-1}{s}n
\]
that is
\weq{1)}{t-2=\Sigma\frac{1}{s},}
whereas in the case of the $K$ of genus $\pi>1$ there holds the equality
\[
2n+2\pi-2=\Sigma\frac{s-1}{s}n
\]
which entails the \textit{inequality}
\weq{2)}{t-2>\Sigma\frac{1}{s}.}

That in case 1) one has
\[
P_{12}=1
\]
and not $P_{12}>1$ follows at once, not only from the formula for $P_{12}$,
% scan: s218R
% !PRINTED 411
\weberpage{411}
but also from observing that on the surface $F$ possessing two pencils of elliptic curves there cannot be (pure) pluricanonical curves of order $>0$.

Let us now prove that, if the inequality 2) holds, one necessarily has
\[
P_{12}>1
\]
and not $P_{12}=1$. This proof is not a simple deduction from the formula for $P_{12}$, but requires in addition that account be taken of a condition for the existence of elliptic surfaces (pertaining to the abelian character of the group $\Gamma_n$) already pointed out in § 13, namely that ``each of the integers $s_i$ must divide the least common multiple of the remaining ones''\origfootnote{1}{s218R-1}{It is precisely in this way that the inequality $P_{12}>1$ is established in § 15 of the Lezioni of \textsc{Enriques-Campedelli} (Seminario mat. di Roma, 1934). But in the original Note of \textsc{Enriques}; \textit{Sulla classificazione delle superficie algebriche etc.} (Rendic. Lincei, February 1914) the matter is not clarified, so that it is natural to suppose that it is a simple deduction from the general formula which gives $P_{12}$. Now this deduction would be erroneous, as \textsc{H. Geppert} has observed in his Bericht \textit{Die Klassification der algebraischen Flaechen} (Jahresbericht der Mathematische Vereinigung, 1931); indeed, taking, for example, $s=2$, $s_2=s_3=5$; one finds $P_{12}=1$. %

Therefore \textsc{Geppert} proposes to consider $P_{24}$ in place of $P_{12}$. But in reality for the elliptic surfaces with pluricanonical curves of order $>0$ one already has $P_{12}>1$. The exceptional case pointed out does not correspond to an actual elliptic surface, because $2$ is not a divisor of $5$!}.

Let us write the formula
\[
P_{12}=12t-23-\Sigma\frac{s+11}{s}, %
\]
and discuss the cases that can occur for $t>4$, for $t=4$ and for $t=3$.

For $t>4$, that is $t\geq5$, one has
\[
\left[\frac{s_i+11}{s_i}\right]\leq6
\]
and therefore certainly
\[
\begin{aligned}
&\Sigma\left[\frac{s_i+11}{s_i}\right]\leq24,\\ %
&P_{12}\geq12t-47\geq12(t-4)+1>1. %
\end{aligned}
\]

For $t=4$ the minimum value of $P_{12}$:
\[
P_{12}=12(t-4)+1=1,
\]
% scan: s219L
% !PRINTED 412
\weberpage{412}
is attained only when all four of the numbers $s_i$ are equal to $2$:
\[
s_1=s_2=s_3=s_4=2,
\]
so that the equality 1) holds, because only in this case will one have for each of them
\[
\left[\frac{s_i+11}{s_i}\right]=6;
\]
if instead one of these numbers exceeds $2$, as is required by the inequality 2), for example if
\[
s_4\geq3,
\]
one has
\[
\left[\frac{s_4+11}{s_4}\right]\leq4
\]
and hence
\[
P_{12}\geq3.
\]

Finally let $t=3$, and let us write all the integer solutions of the inequality
\[
\begin{aligned}
&\frac{1}{s_1}+\frac{1}{s_2}+\frac{1}{s_3}<1\\
&(s_1\leq s_2\leq s_3),
\end{aligned}
\]
satisfying the condition that each of the $s_i$ divide the common multiple of the remaining ones. One will have
\[
\left\{
\begin{aligned}
s_1&=2,&\qquad s_2&=4,&\qquad s_3&=8,\\
s_1&=2,& s_2&=5,& s_3&=10,\\
s_1&=2,& s_2&\geq6,& s_3&\geq s_2,\\
s_1&=3,& s_2&\geq3,& s_3&>3,\\
s_1&>4,& s_2&\geq s_1,& s_3&\geq s_2. %
\end{aligned}
\right.
\]

In all these cases one has
\[
\begin{aligned}
&\Sigma\left[\frac{s+11}{s}\right]<12\\
&P_{12}>1,
\end{aligned}
\]
\hfill Q.E.D.

\edsection{15. Historical note.}{15. Historical note.}{sec.X.15}

The classification of the irregular surfaces of genus $p_g=0$ ($p_a<0$), the proof that those which are not birationally equivalent to ruled surfaces are elliptic surfaces, containing an elliptic pencil of curves
% scan: s219R
% !PRINTED 413
\weberpage{413}
admitting a group $\infty^1$ of transformations into themselves, finally $K$ of genus $\pi\geq1$ and a second linear pencil of elliptic curves $C$, the characterization of ruled surfaces by the vanishing of the quadrigenus and of the sextigenus, constitute the subject of the memoir of \textsc{F. Enriques} ``Sulle superficie algebriche di genere geometrico zero'', published in the Rendiconti del Circolo matematico di Palermo of 1905. However (apart from some minor inaccuracy) the proof of these results left a gap, as regards lemma II of §§ 6, 7, 8; for the existence of a second elliptic curve essentially distinct from the parabicanonical curve $C$ was inferred simply from the fact that the curves of the system $\{K''\}$, second para-adjoint to a $K$, passing through a bicanonical group of the latter, must form a system $\infty^1$ of index $>1$: in fact it is not excluded that the reducible curves of this system which contain the said $K$ as a part reduce to a single one to be counted several times. %

To fill this gap the same author devised the reasoning which rests on the construction of curves cutting the $K$ in groups of the semiquadricanonical series or, in general, of an aliquot part of a multiple of the canonical series. This procedure was set forth by him in a Note of the Accademia dei Lincei, of 1930\origfootnote{1}{s219R-1}{\textsc{Enriques}, \textit{Sopra le superficie algebriche trasformabili in rigate.} (July, 1930).} and then by \textsc{Enriques-Campedelli}, in the ``Lezioni sulla classificazione delle superficie algebriche particolarmente di genere zero''. (Seminario matematico di Roma, 1934).

But in the new proof too there remains a weak point of a delicate nature, because it is not legitimate to assume a priori, as is done provisionally with hypothesis II of §§ 6, 7 and 8, that, a linear series being given in, or adjoined to, the field of rationality of a curve $K$, one can rationally determine \textit{one} group of it. Fortunately we have succeeded in overcoming the difficulty by recognizing that hypothesis II is verified for some multiple of the said series, so as to justify the procedure of construction of essentially distinct parapluricanonical curves.

We must further add a remark concerning the case in which our surfaces possess an elliptic pencil of elliptic curves $K$ (birationally identical); for the construction of the trajectories $C$, or rather the proof that they too are elliptic, to be rigorous, one must remove the doubt that some of these $K$ might degenerate into multiple curves (proof that, for $p_a=-1$, there cannot be paraelliptic surfaces): this
% scan: s220L
% !PRINTED 414
\weberpage{414}
clarification — which is based on the analysis of the singularity arising from the crossing of two branch curves (\textsc{Chisini}) is to be found in the ``Lezioni'' of \textsc{Enriques-Campedelli} of 1934 (§ 12).

Finally, as regards the effective construction of elliptic surfaces, the original Memoir of \textsc{Enriques} of 1905 considered only the cyclic case. The first examples relating to the abelian case (with base two) were brought to light by \textsc{G. Bagnera} and \textsc{M. De Franchis} in their study ``Sopra le superficie algebriche che danno le coordinate del punto generico esprimibili con funzioni meromorfe quadruplamente periodiche di due parametri''\origfootnote{1}{s220L-1}{R. Accademia dei Lincei, $1907_1$.}.

Finally, the general construction of all elliptic surfaces (which is reduced to the solution of abelian equations with group of base two) was developed by \textsc{O. Chisini} in the Notes on ``Le superficie ellittiche il cui determinante è un numero composto'' of 1921\origfootnote{2}{s220L-2}{R. Accademia dei Lincei, $1921_2$.}.

\edsection{16. Surfaces with pluricanonical curves of order zero: types with\\elliptic curves $K$ of general modulus.}{16. Surfaces with pluricanonical curves of order zero: types with elliptic curves K of general modulus.}{sec.X.16}

In § 14, and in particular in the Remark which closes it, we have proved that:

\textit{The properly elliptic surfaces of genera $p_g=0$ and $p_a=-1$ are distinguished from the ruled surfaces by having the 12-genus.}
\[
P_{12}>0,
\]
\textit{instead of}
\[
P_{12}=0.
\]

\textit{More precisely, the elliptic surfaces with $p_g=0$ and $p_a=-1$ give rise to two families:}

1) \textit{surfaces with virtual canonical curve of order zero, characterized by the possession of an elliptic pencil of curves $K$ of genus $\pi=1$;}

2) \textit{surfaces with virtual canonical curve of order $2\pi-2$, which contain an elliptic pencil of curves $K$ of genus $\pi>1$.}

\textit{The former} (on which every pure linear system of genus $p$ has degree $n=2p-2$) \textit{have}
\[
P_{12}=1,
\]
% scan: s220R
% !PRINTED 415
\weberpage{415}
\textit{while the latter} (for which $n<2p-2$) \textit{possess linear systems of dimension $r>0$ of pluricanonical curves} (of genus $p^{(1)}=1$) \textit{and have} %
\[
P_{12}\geq2.
\]

Now we wish to examine more closely the classification of the surfaces 1) for which $P_{12}=1$, showing that they give rise to a few well-determined types, whose equations we can write down.

Let us refer to the elliptic surfaces $F$ of family 1); it is convenient to distinguish the cases in which:

\textit{The elliptic curves $K$} (birationally identical to one another) \textit{have general modulus}, or \textit{the $K$ are harmonic}, or \textit{they are equianharmonic}.

Let us begin here with the \textit{general case}: we shall prove that the surface $F$ has \textit{determinant} $n=2$ (cyclic case) or $n=4$ (abelian case).

For this purpose one must consider the involution $g_n^1$ cut on an elliptic curve $K$ by the trajectories $C$. This involution is generated by an abelian group $\Gamma_n$ and, being rational (rather than of genus $1$), it must contain some involutory transformation $I$ of the second kind ($g_2^1$). Now, if the group $\Gamma_n$ does not reduce to the $\Gamma_2$ formed by an $I$ and the identity, then (since the products of two $I$ are transformations $\pi$ of the first kind) $\Gamma_n$ must contain as many transformations $\pi$ as $I$, say $m$ of each: $n=2m$. More precisely, the $\pi$ permutable with an $I$ will be involutory, and hence the abelian group $\Gamma_n$ generated by the said $\pi$ and $I$ will be the dihedral $\Gamma_4$ formed by two involutions $I$, one involution $\pi$ and the identity.

Hence our surface $F$ with a pencil of \textit{elliptic curves $K$ of general modulus will be of determinant $n=2$ or $n=4$, and the relevant $g_n^1$ defined on a $K$ will belong to one of the two following types:}

\textit{a)} $n=2,\quad g_n^1=g_2^1,$

\textit{b)} $n=4,\quad g_n^1=g_4^1$ \textit{abelian and dihedral, generated by two permutable transformations of the second kind, whose product is a $\gamma_2^1$ of the first series.} %

In the first case one may take as \textit{type of the surface $F$ of determinant $n=2$ the one which is defined by the equations:}
\weq{\mathrm{I}_a}{
\left\{
\begin{aligned}
u&=\sqrt{(z-a_1)(z-a_2)(z-a_3)(z-a_4)\varphi(xy)}\\ %
f(xy)&=0,
\end{aligned}
\right.
}
where $f$ is a cubic cylinder and $\varphi$ a quadric cylinder which has with it a two-point contact along three generators, not lying in a plane.

In the second case it is a matter of constructing, in the plane $(u\,z)$, the elliptic
% scan: s221L
% !PRINTED 416
\weberpage{416}
curve \(K\) containing a dihedral \(g_4^1\), starting from the corresponding representation on the quadruple line \(k\) (\(u=0\)).

This construction can be explained in relation to the general theory of abelian equations\origfootnote{1}{s221L-1}{Cf. e.g. L. \textsc{Bianchi}, op. cit., chap. IV, \S~76.}, or by referring more especially to the notions which concern elliptic curves\origfootnote{2}{s221L-2}{Cf. \textsc{Enriques-Chisini}, \textit{Lezioni.} Book V, chap. IV, \S~39 (vol. III).}, as we shall do here.

Let us observe that our dihedral \(g_4^1\) (on \(K\)) possesses 8 double points, which are distributed in two tetrads: a tetrad of double points of the first \(g_2^1\), consisting of two pairs of conjugate points of the second
\[
A_1\,A_2 \qquad \text{and} \qquad B_1\,B_2,
\]
and a tetrad of double points of the second \(g_2^1\), consisting of two pairs of the first
\[
D_1\,D_2 \qquad \text{and} \qquad E_1\,E_2.
\]

Hence the quadruple line \(k\) will have 4 (double) branch points:
\[
a\,b\,d\,e,
\]
and the field of rationality of \(K\) can be defined by putting
\[
u=\sqrt{\varphi\cdot(z-a)(x-b)}\pm\sqrt{\psi\cdot(z-d)(z-e)}, %
\]
(with \(\varphi\) and \(\psi\) constant with respect to \(z\) and \(u\)) in such a way that to the two branch points \(a\) and \(b\) there correspond the pairs of double points \(A_1\,A_2\) and \(B_1\,B_2\), and to the points \(d\) and \(e\) there correspond the pairs of double points \(D_1\,D_2\) and \(E_1\,E_2\).

That the equation written above does in fact represent the elliptic \(K\) corresponding to our quadruple \(k\) is proved by the fact that the curve represented possesses a \(g_4^1\) with 8 double points, and moreover --- freeing the equation itself of radicals --- one has an equation of the \(4^{\mathrm{o}}\) degree in \(z\) and \(u\), which is of the \(2^{\mathrm{o}}\) degree in \(z\), and defines a quartic having a double point at the point at infinity of the \(z\)-axis and another double point infinitely near to it. Let us add that the two transformations of the second kind, that is the \(g_2^1\) each endowed with 4 double points, of which the dihedral \(g_4^1\) is formed, can be represented by means of equations
\[
\begin{aligned}
u'&=\frac{\varphi\cdot(z-a)(z-b)-\psi\cdot(z-e)(z-d)}{u}\\
u''&=\frac{\psi\cdot(z-d)(z-e)-\varphi\cdot(z-a)(z-b)}{u},
\end{aligned}
\]
% scan: s221R
% !PRINTED 417
\weberpage{417}
as follows from the identity:
\[
\begin{aligned}
&\frac{\varphi\cdot(z-a)(z-b)-\psi\cdot(z-d)(z-e)}
{\sqrt{\varphi\cdot(z-a)(z-b)}+\sqrt{\psi\cdot(z-d)(z-e)}}\\
&=\sqrt{\varphi\cdot(z-a)(z-b)}-\sqrt{\psi\cdot(z-d)(z-e)}.
\end{aligned}
\]

Therefore the equations of the \textit{elliptic surface \(F\) of determinant \(n=4\), with \(K\) of general modulus, can be reduced to the type:}
\weq{\mathrm{I}_b}{
\left\{
\begin{aligned}
u&=\sqrt{(z-a)(z-b)\cdot\varphi(xy)}+\sqrt{(z-d)(z-e)\cdot\psi(xy)}\\
f(xy)&=0,
\end{aligned}
\right.
}
where \(f\) is a cubic cylinder and \(\varphi\) and \(\psi\) are two quadric cylinders which touch \(f\) along two triads of generators, not lying in a plane and not equivalent.

\edsection{17. Continued: Harmonic case.}{17. Continued: Harmonic case.}{sec.X.17}

In the harmonic case, besides the \(g_2^1\) and the \(g_4^1\) of the general case, there may be on the elliptic curve \(K\) a cyclic or abelian \(g_n^1\) whose groups of \(n\) points are transformed into themselves by a singular transformation of the fourth order. Indeed, the harmonic elliptic curves possess, besides the ordinary transformations of the first and second kind, \(\pi\) and \(I\), two continuous series \(\infty^1\) of singular transformations of the \(4^{\mathrm{o}}\) order, \(\omega\) and \(\tau\), the four kinds of transformations being connected with one another by the following relations\origfootnote{1}{s221R-1}{Cf. e.g. \textsc{Enriques-Chisini}, \textit{Lezioni}, Book V, chap. III, \S~27 (vol. III, p. 260).}: %
\[
\begin{aligned}
\pi\omega&=\omega_1 & \omega_1\omega&=I & I\omega&=\tau & \tau\omega&=\pi\\
\pi I&=I_1 & I_1I&=\pi & \omega I&=\tau & \tau I&=\omega\\
\pi\tau&=\tau_1 & \tau_1\tau&=I & I\tau&=\omega & \omega\tau&=\pi\\
\pi_1\pi&=\pi_2 & I\pi&=I_1 & \omega\pi&=\omega_1 & \tau\pi&=\tau_1.
\end{aligned}
\]

One can classify the various types of cyclic or abelian \(g_n^1\) belonging to a harmonic elliptic curve \(K\), starting from the observation that, if the relevant group \(\Gamma_n\) is not generated by ordinary transformations alone (corresponding to \(K\) of general modulus), it must contain an invariant subgroup \(\Gamma_4\) generated by a singular transformation \(\omega\) or \(\tau\). Now this \(\Gamma_4\) gives rise to a \(g_4^1\) with two quadruple points \(A\) and \(B\) and a pair of double points \(D_1\) and \(D_2\) (forming a group of the \(g_4^1\)), which leads to representing \(K\) on a quadruple line \(k\) (\(u=0\)) with two quadruple branch points: \(z=a\) and \(z=b\), and a double branch point: \(z=d\).
\printersignature{\textsc{Enriques F.} - \textit{Superficie algebriche.}\quad 27}
% scan: s222L
% !PRINTED 418
\weberpage{418}
If the \(g_n^1\) is not exhausted by the said \(g_4^1\) (that is, if \(n>4\)), the group \(\Gamma_n'\) must contain the transformations which correspond to the projectivities of the quadruple line transforming the point \(d\) into itself and interchanging the two points \(A\) and \(B\): and these transformations generate an abelian \(\Gamma_8\), formed by 4 singular transformations of the \(4^{\mathrm{o}}\) order, two \(\omega\) and two \(\tau\), by 2 involutions of the second kind \(I\), and by 2 transformations \(\pi\) of the first kind: an involution \(\gamma_2^1\) and the identity. %

In conclusion, the \(g_n^1\) not generated by ordinary transformations alone, on the \textit{harmonic elliptic curve \(K\)}, will belong to one of the two following types:

\textit{a}) \(n=4\), cyclic \(g_4^1\) generated by a singular transformation of the \(4^{\mathrm{o}}\) order: \(\omega\) or \(\tau\); and

\textit{b}) \(n=8\), \(g_8^1\) generated by an abelian group \(\Gamma_8\) containing 4 cyclic singular transformations of the \(4^{\mathrm{o}}\) order (two \(\omega\) and two \(\tau\)) and 4 ordinary transformations: two \(I\) and two \(\pi\), (that is, a \(\gamma_2^1\) and the identity).

In the \textit{first case} the elliptic \(K\) represented on the quadruple \(k\), with the branch points \(a\), \(b\) and \(d\), is realized in the plane (\(z\,u\)) by writing, for example:
\[
u=\sqrt[4]{\varphi\cdot(z-a)(z-b)(z-d)^2}
\]
with \(\varphi\) constant (with respect to \(z\) and \(u\)).

The equation written corresponds to a certain choice of the substitutions on the branches of the algebraic function \(u(z)\), in relation to the branch curves \(z=a,\ b,\ d\); if instead the inverse substitutions are taken, we shall have the \textit{similar type}, but birationally distinct:
\[
u=\sqrt[4]{\varphi\cdot(z-a)^3(z-b)^3(z-d)^2}.
\]

Hence \textit{the harmonic elliptic surfaces of determinant \(n=4\), cyclic case, can be reduced to the type represented by the equations:}
\weq{\mathrm{II}_a}{
\left\{
\begin{aligned}
u&=\sqrt[4]{(z-a)(z-b)(z-d)^2\cdot\varphi(xy)}\\
f(xy)&=0,
\end{aligned}
\right.
}
or to a \textit{similar type}: where \(f\) is a cubic cylinder and \(\varphi\) a cylinder of the fourth order touching \(f\) with four-point contact along three generators, not lying in a plane. The similar types, which differ, as has been said, in the choice of the substitutions on the branches of \(u\), although birationally distinct, nevertheless have the same characters.

In the properly abelian case we can again realize in the plane (\(uz\)) a type of the elliptic curve \(K\) which contains a \(g_8^1\) with two groups each consisting of two quadruple points and one group consisting of four double points: for this it is a matter of constructing an
% scan: s222R
% !PRINTED 419
\weberpage{419}
8-fold line with three branch points
\[
z=a,\qquad z=b,\qquad z=d
\]
in such a way that to the first two there corresponds a cyclic substitution on the branches of the \(4^{\mathrm{o}}\) order, while to the third there corresponds a simple substitution of period 2.

This construction is carried out according to the general principles of the theory of abelian equations\origfootnote{1}{s222R-1}{Cf. e.g., L. \textsc{Bianchi}, \textit{Lezioni sulla teoria dei gruppi di sostituzioni e delle equazioni algebriche secondo Galois.} (Pisa, 1900) chap. III, \S~30.}, and moreover the formula thus obtained is justified a posteriori\origfootnote{2}{s222R-2}{Cf. O. \textsc{Chisini}, \textit{Le superficie ellittiche il cui determinante è un numero composto.} Rendic. Lincei, s. V, vol. XXX, \(1921_2\).} by recognizing the substitutions on the branches which correspond to the branch points, of which we spoke above.

We say that the curve \(K\) can be reduced to the type represented by the equations
\[
u=\sqrt{(z-b)(z-d)\cdot\varphi}+\sqrt[4]{(z-a)(z-b)^3\cdot\psi}
\]
or to a \textit{similar type} (which differs from it by a different choice of the aforesaid substitutions): where \(\varphi\) and \(\psi\) denote constants (with respect to \(u\) and \(z\)).

To justify the assertion, let us distinguish the 8 branches of \(u\) by two indices:
\[
\begin{aligned}
&u_{11}\qquad u_{12}\qquad u_{13}\qquad u_{14}\\
&u_{21}\qquad u_{22}\qquad u_{23}\qquad u_{24},
\end{aligned}
\]
making the first index correspond to the two values \(+\sqrt{}\) and \(-\sqrt{}\) of the quadratic radical which enters into the expression of \(u\), and the second to the four values \(+\sqrt[4]{},\ i\sqrt[4]{},\ -\sqrt[4]{},\ -i\sqrt[4]{}\) of the radical of the \(4^{\mathrm{o}}\) order.

To the branch point \(z=a\) there will correspond the substitution on the branches
\[
S_a=(u_{11}\,u_{12}\,u_{13}\,u_{14})(u_{21}\,u_{22}\,u_{23}\,u_{24}).
\]

And to the point \(z=b\) there will likewise correspond a cyclic substitution of the \(4^{\mathrm{o}}\) order, namely
\[
S_b=(u_{11}\,u_{24}\,u_{13}\,u_{22})(u_{12}\,u_{21}\,u_{14}\,u_{23}).
\]

Whereas to the point \(z=d\) there will correspond the substitution of period 2:
\[
S_d=(u_{11}\,u_{21})(u_{12}\,u_{22})(u_{13}\,u_{23})(u_{14}\,u_{24}).
\]

One thus obtains three substitutions which, by multiplication, generate an abelian group of order 8, having the structure of our \(\Gamma_8\) indicated above.

% scan: s223L
% !PRINTED 420
\weberpage{420}
It may be added that the aforesaid substitutions satisfy the condition of existence\origfootnote{1}{s223L-1}{Cf. e.g., \textsc{Enriques-Chisini}, \textit{Lezioni.} Book V, chap. IV, \S~38 (vol. III, p. 429).} of the algebraic function \(u(z)\), (not branched at the point \(z=\infty\)), that is, that
\[
S_d\,S_b\,S_a=1;
\]
it is precisely in view of this condition that the exponent 3 (instead of 1) had to be assigned to the term \(z-b\) which appears under the radical of the \(4^{\mathrm{th}}\) order.

In this way the function \(u(z)\) realizes the required elliptic curve \(K\) with abelian \(g_8^1\), and indeed furnishes its \textit{type}, although there may also be another \textit{similar type}, not transformable into it by birational transformations, such as is given by
\[
u=\sqrt{(z-a)(z-b)^3\cdot\varphi}+\sqrt[4]{(z-a)(z-d)\cdot\psi}\,\text{\origfootmark{2}{s223L-2}}:
\]\origfoottext{2}{s223L-2}{From the projective point of view the curve \(K\) representing the function \(u(z)\) will have: corresponding to the branch point \(z=a\), a tangent (parallel to the \(u\) axis) with two four-point contacts, and corresponding to \(z=b\), a tacnode at which two linear branches touch with four-point contact (which thus accounts for \textit{two} quadruple points of the \(g_8^1\)), finally, corresponding to \(z=d\), a quadruple tangent, with four simple contacts.}
other types which apparently present themselves reduce to those indicated here, by simple interchanges of the names of the branch curves.

It is worth adding that the two substitutions on the branches of \(u\) in the neighborhood of the branch points \(z=a\) and \(z=b\) extend to birational transformations, of the same order 4, among the points of the curve \(K\) whose equation
\[
\theta(zu)=0
\]
is obtained by eliminating the radicals which appear in the expression of \(u\); and the same may be said of the substitution of period 2 relating to the point \(z=d\).

Indeed one can prove, for example, that \(u_{12}\) is a rational function of \(u_{11}\) and of \(z\); this follows from considering the two equations satisfied for one and the same value of \(z\):
\[
\theta(uz)=0,\qquad\theta(u'z)=0
\]
and their resultant
\[
R(uu')=0:
\]
this, together with \(\theta(uz)=0\), defines \(u'(uz)\), as a rational function.

% scan: s223R
% !PRINTED 421
\weberpage{421}
Now we can write the \textit{equations of an elliptic surface, of harmonic type, of determinant \(n=8\):}
\weq{\mathrm{II}_b}{
\left\{
\begin{aligned}
u&=\sqrt[2]{(z-a)(z-b)^3\cdot\varphi(xy)}+\sqrt[4]{(z-b)(z-d)\cdot\psi(xy)},\\
f(xy)&=0
\end{aligned}
\right.
}
where \(f\) represents a cubic cylinder and \(\varphi\) and \(\psi\) may be taken to be two cylinders, the one of the \(2^{\mathrm{nd}}\) and the other of the \(4^{\mathrm{th}}\) order, the first tangent to \(f\) along 3 generators and the second tangent to it with four-point contact along another 3 generators. In order that the surface \(F\) be precisely of determinant \(n=8\) (irreducibility of the curves \(C\) which correspond on it to \(z=\text{const.}\)) one must suppose that the two triples of generators of contact of \(f\) and \(\varphi\) and of \(f\) and \(\psi\) do not lie in planes, so that for \(z=\text{const.}\) the two curves
\[
u=\sqrt[2]{(z-a)(z-b)^3\cdot\varphi(xy)},\qquad f(xy)=0
\]
and
\[
u=\sqrt[4]{(z-b)(z-d)\cdot\psi(xy)},\qquad f(xy)=0,
\]
turn out irreducible, and further that the two curves
\[
u=\sqrt[2]{(z-a)(z-b)^3\cdot\varphi(xy)}
\qquad\text{and}\qquad
u=\sqrt[4]{(z-b)(z-d)\cdot\psi(xy)}
\]
be birationally distinct\origfootnote{1}{s223R-1}{Cf. \textsc{Enriques-Chisini}, \textit{Lezioni.} Book V, chap. IV, \S~33 (vol. III, p. 452).}: this entails that the triple \(g_1\,g_2\,g_3\) of the generators of four-point contact of \(\psi\) with \(f\) be \textit{dissimilar} from the triple \(g'_1\,g'_2\,g'_3\) of the generators of contact of \(\varphi\) and \(f\), that is, that, counted twice, it constitute a sextuple not equivalent to the one obtained by adding to \(g'_1\,g'_2\,g'_3\) three generators of \(f\) which lie in a plane.

\edsection{18. Equianharmonic case.}{18. Equianharmonic case.}{sec.X.18}

In the equianharmonic case, if the elliptic curve \(K\) contains a \(g_n^1\) which is not generated by mere ordinary transformations (\(n=2,\ 4\)), it must be generated by a group \(\Gamma_n\) of transformations, cyclic or abelian (with a basis of two), which contains some singular transformation of the \(3^{\mathrm{rd}}\) or of the \(6^{\mathrm{th}}\) order. Indeed the equianharmonic elliptic curves possess, besides the ordinary transformations of the first and of the second kind, \(\pi\) and \(I\), two continuous \(\infty^1\) series of cyclic singular transformations of the \(3^{\mathrm{rd}}\) order, \(\omega\) and \(\tau\), and another two series of cyclic singular transformations of the \(6^{\mathrm{th}}\) order, \(\sigma\) and \(\varrho\)\origfootnote{2}{s223R-2}{\textsc{Enriques-Chisini}, \textit{Lezioni.} Book V, chap. III, \S~27 (vol. III, pp. 260-63).}. And the transfor%
% scan: s224L
% !PRINTED 422
\unskip\weberpage{422}%
mations of the various types are connected with one another by the following relations:
\[
\left\{
\begin{aligned}
\pi\sigma&=\sigma_1 & \sigma_1\sigma&=\tau & \tau\sigma&=I & I\sigma&=\omega & \omega\sigma&=\varrho & \varrho\sigma&=\pi\\
\pi\tau&=\tau_1 & \tau_1\tau&=\omega & \omega\tau&=\pi & \sigma\tau&=I & I\tau&=\varrho & \varrho\tau&=\sigma\\
\pi I&=I_1 & I_1I&=\pi & \sigma I&=\omega & \omega I&=\sigma & \tau I&=\varrho & \varrho I&=\tau.\\ %
\pi\omega&=\omega_1 & \omega_1\omega&=\tau & \pi\omega&=\pi & \sigma\omega&=\varrho & \varrho\omega&=I & I\omega&=\sigma\\ %
\pi\varrho&=\varrho_1 & \varrho_1\varrho&=\omega & \omega\varrho&=I & I\varrho&=\tau & \tau\varrho&=\sigma & \sigma\varrho&=\pi.
\end{aligned}
\right.
\]

Now for the abelian \(g_n\) not generated by ordinary transformations on the \textit{equianharmonic curve \(K\)}, \textit{three cases} will be possible:

\textit{a}) \(n=3\), cyclic \(g_3^1\) generated by a singular transformation of the \(3^{\mathrm{rd}}\) order, \(\omega\) or \(\tau\);

\textit{b}) \(n=6\), cyclic \(g_6^1\), generated by a singular transformation of the \(6^{\mathrm{th}}\) order, \(\sigma\) or \(\varrho\);

\textit{c}) \(n=9\), abelian \(g_9^1\) (with a basis of two) whose group contains three cyclic singular transformations of the \(3^{\mathrm{rd}}\) order, permutable with one another, say \(\omega\,\omega_1\,\omega_2\), another three cyclic transformations as above \(\tau=\omega^2,\ \tau_1=\omega_1^2,\ \tau_2=\omega_2^2\), and finally three ordinary transformations of the first kind \(\pi,\ \pi^2,\ \pi^3=1\).

This classification can easily be justified starting from the observation that the group \(\Gamma_n\) must contain an invariant cyclic subgroup \(\Gamma_3\), generated by a singular transformation \(\omega\) or \(\tau\). Since \(\Gamma_3\) generates a cyclic \(g_3^1\) with three triple points, we are led to examine the transformations of \(K\) which correspond to the projectivities transforming into itself a triple line of the type
\[
u^3=(z-a)(z-b)(z-d);
\]
these are:

1) the involutions, e.g. \((ab)(dd)\), which have as double point a point of the triple \((abd)\) and which interchange the other two points; and

2) the cyclic projectivities \((abd)\) or \((adb)\).

Now the addition to \(\Gamma_3\) of a transformation of type 1) succeeds in enlarging \(\Gamma_3\) itself into a cyclic \(\Gamma_6\) (with a sextuple point, two triple points and three double points) which is not contained in a larger abelian group; while the addition of a transformation \((abd)\) of type 2) leads to an abelian \(\Gamma_9\), containing three singular transformations of the third order, which in its turn is not contained in a larger abelian group.

We can easily construct the equianharmonic curves \(K\) which correspond to the aforesaid abelian \(g_n^1\), for
\[
n=3,\qquad 6,\qquad 9,
\]
and hence the corresponding elliptic surfaces.

% scan: s224R
% !PRINTED 423
\weberpage{423}
We shall have, as \textit{type of the equianharmonic elliptic surfaces of determinant \(n=3\)} (cyclic case), \textit{the surfaces}:
\weq{\mathrm{III}_a}{
\left\{
\begin{aligned}
u&=\sqrt[3]{(z-a)(z-b)(z-d)\cdot\varphi(xy)}\\
f(xy)&=0,
\end{aligned}
\right.
}
where \(f\) is a cubic cylinder and \(\varphi\) another cylinder of the same order, which touches \(f\) with three-point contact along three generators not lying in a plane.

And we shall have, as \textit{type of the equianharmonic elliptic surfaces of determinant \(n=6\)} (cyclic case), \textit{the surfaces}:
\weq{\mathrm{III}_b}{
\left\{
\begin{aligned}
u&=\sqrt[6]{(z-a)(z-b)^2(z-d)^3\cdot\varphi(xy)}\\
f(xy)&=0,
\end{aligned}
\right.
}
where \(f\) denotes a cubic cylinder and \(\varphi\) another cylinder of the \(6^{\mathrm{th}}\) order which touches \(f\) with six-point contact, along three generators not belonging to a plane.

In the properly abelian case, in which we are dealing with the equianharmonic elliptic surfaces \(F\) of determinant \(n=9\), one can succeed in constructing a type of \(F\), in a manner analogous to the one followed for the harmonic case (\(n=8\)).

For this purpose one must construct the equianharmonic elliptic curve \(K\) of the plane (\(uz\)) which is represented on a multiple line of order 9 (\(u=0\)), image of our abelian \(g_9^1\), generated by two singular transformations of the \(3^{\mathrm{rd}}\) order, permutable with one another. As type one may take the curve
\[
u=\sqrt[3]{(z-a)^2(z-d)\cdot\varphi}+\sqrt[3]{(z-b)(z-d)^2\cdot\psi}
\]
with \(\varphi\) and \(\psi\) constant.

To prove that the \(u(z)\) thus defined does in fact correspond to an equianharmonic elliptic curve containing the abelian \(g_9^1\) indicated above, we shall have to compute the substitutions on the branches of \(u\) with respect to the branch points
\[
z=a,\qquad z=b,\qquad z=d.
\]

For this purpose we shall denote the branches by two indices, the first of which will take the values 1, 2, 3 in relation to the values of the first cubic radical
\[
\sqrt[3]{},\qquad \varepsilon\sqrt[3]{},\qquad \varepsilon^2\sqrt[3]{}
\qquad
\left(\varepsilon=e^{\frac{2\pi i}{3}}\right)
\]
which appears in the expression of \(u(z)\), while the second will similarly take three values in relation to those of the second cubic radical.
% scan: s225L
% !PRINTED 424
\weberpage{424}
We shall have, with respect to the three branch points, the following substitutions:

for \(z=a\)
\[
S_a=(u_{11}\,u_{31}\,u_{21})(u_{12}\,u_{32}\,u_{22})(u_{13}\,u_{33}\,u_{23});
\]

for \(z=b\)
\[
S_b=(u_{11}\,u_{12}\,u_{13})(u_{31}\,u_{32}\,u_{33})(u_{21}\,u_{22}\,u_{23})
\]

and for \(z=d\)
\[
S_d=(u_{11}\,u_{23}\,u_{32})(u_{12}\,u_{21}\,u_{33})(u_{13}\,u_{22}\,u_{31}).
\]

The third substitution is the inverse of the product of the other two, so that the condition of existence of the multiple line \(u=0\), not branched at the point at infinity \(z=\infty\)\origfootnote{1}{s225L-1}{Cf. \textsc{Enriques-Chisini}, \textit{Lezioni.} Book V, chap. IV, \S~38 (vol. III, p. 429).}, is satisfied: \(S_dS_bS_a=1\).

Hence one recognizes that the said substitutions generate an abelian group \(\Gamma_9\), which gives rise to a \(g_9^1\), likewise abelian, on an equianharmonic elliptic curve.

Thus one is led to prove that the \textit{equianharmonic elliptic surfaces of determinant \(n=9\)} (abelian case) \textit{give rise to the type}
\weq{\mathrm{III}_c}{
\left\{
\begin{aligned}
u&=\sqrt[3]{(z-a)^2(z-d)\cdot\varphi(xy)}+\sqrt[3]{(z-b)(z-d)^2\psi(xy)}\\
f(xy)&=0,
\end{aligned}
\right.
}
where \(f\) is a cubic cylinder and \(\varphi\) and \(\psi\) may likewise be taken to be cubic cylinders osculating \(f\) along two non-equivalent triples of generators, which do not lie in one and the same plane.

\edsection{19. Summary.}{19. Summary.}{sec.X.19}

We can summarize the analysis of the elliptic surfaces of genus \(p_g=0\) with virtual canonical curve of order zero by stating the following theorem:

\textit{The elliptic surfaces of genus \(p_g=0\) and \(p_a=-1\) with virtual canonical curve of order zero, characterized by the value of \(P_{12}=1\),} that is, the surfaces with \(p_a=-1\) possessing an elliptic pencil of elliptic curves \(K\) and a second linear pencil of elliptic curves \(C\), \textit{are distributed into seven families of determinant \(n=2,\ 4;\ 4,\ 8;\ 3,\ 6,\ 9\), represented by one of the types of equations \(I_a I_b II_a II_b III_a III_b III_c\), or else by a similar type.}

The types \(I_a\) and \(I_b\) (resp. \textit{cyclic} type and \textit{abelian} type) correspond to elliptic curves \(C\) of \textit{general modulus}; \(II_a\) and \(II_b\) (\textit{cyclic} type and \textit{abelian} type)
% scan: s225R
% !PRINTED 425
\weberpage{425}
to \textit{harmonic} curves \(C\), and finally \(III_a\) and \(III_b\) to a \textit{cyclic} type with \textit{equianharmonic} \(C\), and \(III_c\) \textit{to an abelian equianharmonic type.}

The aforesaid elliptic surfaces are also distinguished, up to a certain point, by the values of the plurigenera lower than \(P_{12}\), which can be computed by the formula of \S~14. The surfaces \(I_a\) and \(I_b\) (determinant 2 and 4) have the characters:
\[
P_2=1,\qquad P_4=1,\qquad P_{12}=1;
\]
the surfaces of harmonic type \(II_a\) and \(II_b\) (determinant 4 and 8) have:
\[
(P_2=0),\qquad P_4=1,\qquad P_6=0,\qquad P_{12}=1;
\]
and finally the surfaces of the equianharmonic type \(III_a\) (determinant 3) have
\[
P_2=0,\qquad P_3=1,\qquad P_4=1,\qquad P_{12}=1, %
\]
while \(III_b\) and \(III_c\) (determinant 6 and 9) have
\[
P_2=0,\qquad P_3=0,\qquad P_4=1,\qquad P_{12}=1.
\]

\textit{Historical note.} -- The elliptic surfaces with virtual canonical curve of order zero, already observed by \textsc{Enriques} in his study of surfaces of genus zero (1905), present themselves as particular hyperelliptic surfaces (more precisely, as involutions on the Jacobi surface which corresponds to a degenerate curve of genus two), and in this sense they were met with by G. \textsc{Bagnera} and M. \textsc{De Franchis} in the memoir of 1907 « Sopra le superficie algebriche che hanno le coordinate del punto generico esprimibili con funzioni meromorfe quadruplamente periodiche di due parametri »\origfootnote{1}{s225R-1}{Rendic. Lincei, ser. V, vol. XVI, \(1907_1\).}.

The said authors discover that such surfaces give rise to a finite number of types, and they succeed in classifying them by assigning their hyperelliptic parametric representations, which correspond to our 7 typical equations.

Independently of \textsc{Bagnera} and \textsc{De Franchis}, but after them, \textsc{Enriques} and \textsc{Severi} also found this result again in another way, likewise starting from the parametric representation of our surfaces by means of hyperelliptic functions\origfootnote{2}{s225R-2}{F. \textsc{Enriques} and F. \textsc{Severi}, \textit{Intorno alle superficie iperellittiche irregolari.} Rendic. Lincei, ser. V, vol. XVII, \(1908_1\) (cf. by the same authors ibidem, vol. XVI, \(1907_1\)). \textit{Mémoire sur les surfaces hyperelliptiques.} Acta Mathem., t. 32 and 33 (1909). See in particular \S\S~56 and 57.}.

In this treatment we find particularly noteworthy the idea of constructing the proper hyperelliptic or abelian surface as a multiple elliptic surface without branch curves, computing the characters which serve to define it (\(p_a=-1,\ p_g=1,\)
% scan: s226L
% !PRINTED 426
\weberpage{426}
\(P_4=1\) according to a theorem of Picard-Enriques). For example, referring to the simplest elliptic surface, of type \(I_a\), one will have an \(F_m\) of a certain order \(m\) in \(S_3\), possessing a biadjoint surface \(\varphi_{2n-8}\), and one will obtain a proper hyperelliptic surface (\(p_a=-1\), \(p_g=1\)) by extracting \(\sqrt[2]{\varphi_{2n-8}}\) on \(F_m\). %

The present treatment of the problem, that is, the algebraic classification of our elliptic surfaces (which, in the constructive part, makes use of the cited work of O. \textsc{Chisini}), formed the subject of a Note by \textsc{Enriques} « Sulle superficie ellittiche di genere zero »%
\origfootnote{3}{s226L-1}{Rendic. Lincei, ser. VI, vol. XIX, 18 February 1934.} of 1934, and is subsequently set forth in the « Lezioni » of \textsc{Enriques-Campedelli} of the Mathematical Seminar of Rome (1934).

\clearpage\thispagestyle{plain}
% scan: s226R
% !PRINTED none
\weberpage{[s226R]}
\edchapter{XI}{GENERAL CLASSIFICATION OF SURFACES}{Chapter XI. GENERAL CLASSIFICATION OF SURFACES}{GENERAL CLASSIFICATION OF SURFACES}

\edsection{1. Introduction.}{1. Introduction.}{sec.XI.1}

The reader who has followed the developments of this treatise, if — as often happens among mathematicians — he is above all a lover of general truths, may have drawn from it the impression that the author has given too much room to examples and particular cases, letting himself be guided somehow by the feeling of curiosity of the naturalist who collects in a museum the most diverse types of animals or plants or minerals.

But just as the museum succeeds in giving an idea of the richness of the forms of life, and thus leads to general problems of biology, so the collection of examples, in this field of mathematics, also takes on an essential significance with regard to the heuristic or historico-constructive aspect of science. We can illustrate its value by repeating the words with which G. \textsc{Castelnuovo} gave an account of our efforts, made in common, to dispel the obscurities that we met at the beginnings of the theory of surfaces%
\origfootnote{1}{s226R-1}{\textit{La geometria algebrica e la scuola italiana.} Lecture delivered at the International Congress of Mathematicians, Bologna, September 1928.}.

« We had constructed, in an abstract sense of course, a great number of models of surfaces of our space or of higher spaces; and we had distributed these models, so to speak, into two display cases. One contained the regular surfaces, for which everything proceeded as in the best of all possible worlds; analogy made it possible to carry over to them the most salient properties of plane curves. But when we tried to verify these properties on the surfaces of the other display case, the irregular ones, the troubles began, and exceptions of every kind presented themselves. In the end the assiduous study of our models had led us to divine certain properties which had to hold, with suitable modifications, for the surfaces of both
% scan: s227L
% !PRINTED 428
\weberpage{428}
display cases; we then put these properties to the test by constructing new models. If they withstood the test, we sought, as the last phase, their logical justification ».

But now we can say more: the determinations of particular families or classes of surfaces over which we have lingered do not yield mere examples, whose significance may be superseded in the inductive process of science; for they draw together into a synthesis which achieves what may be called the principal aim of this algebro-geometric theory; namely, the general classification of algebraic surfaces.

In this connection it should be noted that algebraic geometry is, in substance, a higher arithmetic, in which the very criterion of generality has only a relative meaning. There is here a very profound observation, which goes back to \textsc{Poncelet} and to \textsc{Noether}: from a purely logical point of view it cannot be said, for example, that the family of surfaces not referable to ruled surfaces is wider than that of ruled surfaces. If one proceeds to classify surfaces (of ordinary space or of a hyperspace) according to their order \(n\), the more general case appears to be that of surfaces not referable to ruled surfaces whose hyperplane sections are of genus \(\pi\) with
\[
n\leq 2\pi-2;
\]
but if instead one proceeds to classify surfaces according to the genus \(\pi\) of their hyperplane sections, the more general case will appear to be that in which the order
\[
n>2\pi-2,
\]
that is, the case which corresponds to surfaces referable to ruled surfaces.

We do not wish to draw paradoxical conclusions from what has been said. But if we continue to regard the study of surfaces not belonging to the family of ruled surfaces as \textit{more general} than that of ruled surfaces, this judgment, rather than a statement of fact, will be a \textit{value judgment}, in which one weighs — so to speak — the problems relating to the one and to the other type, the wealth of cases to which they give rise, and the possibility of subordinating the one family to the other from more significant points of view. From this we should draw the following lesson: that every family or class of surfaces which presents characters of its own distinguishing it from the rest constitutes an object of study equally worthy of being pursued, and that the definition of such a family in terms of integer invariant characters often carries an interest of a general order. For example, the definition of the family of ruled surfaces by the vanishing of the quadrigenus and of the sexigenus constitutes not so much a particular result concerning this family as a significant theorem of the
% scan: s227R
% !PRINTED 429
\weberpage{429}
general theory; one which teaches us the existence of canonical or pluricanonical curves (of order \(\geq 0\)) on any surface not referable to a ruled surface.

Our classification, gathering into a synthesis the results of the preceding analyses, succeeds in distinguishing 4 families of surfaces, which have essentially different properties and which are defined by the value of \(P_{12}\) and — subordinately — of \(p^{(1)}\). They are:

\textit{A}) the ruled surfaces (existence of linear systems of genus \(\pi\) and degree \(n>2\pi-2\), infinitely many non-eliminable exceptional curves, and also — as we shall see — continuous series of birational transformations not forming a group of finite dimension);

\textit{B}) the surfaces (without exceptional curves) with pluricanonical curves of order zero (pure linear systems of degree \(n=2\pi-2\));

\textit{C}) the surfaces (\(n<2\pi-2\) for \(\pi>1\)) having canonical or pluricanonical curves composed of the elliptic curves of a pencil (\(p^{(1)}=1\));

\textit{D}) the surfaces (\(n<2\pi-2\)) with \(p^{(1)}>1\), which give rise to canonical or pluricanonical surfaces furnishing a model whose projective properties mirror the invariant properties of the class of surfaces.

Each of these families constitutes an object of study interesting in itself, but bringing to light the characteristic properties which distinguish it from the others has, as we have said, a value for the general theory.

Now, in order that the classification take on its proper significance, we must extend to surfaces of genus \(p_g>0\) some observations made for the case \(p_g=0\); and for this purpose it is useful first of all to study more deeply the surfaces which admit a continuous series of birational transformations into themselves.

\edsection{2. Surfaces which admit a continuous series of birational transformations\\into themselves: cases which lead to ruled surfaces.}{2. Surfaces which admit a continuous series of birational transformations into themselves: cases which lead to ruled surfaces.}{sec.XI.2}

We know that the surfaces referable to ruled surfaces possess continuous series of birational transformations into themselves; in particular, on a rational surface, or on the plane, one knows%
\origfootnote{1}{s227R-1}{\textsc{Enriques}, 1893. Cf. \textsc{Enriques-Chisini}, \textit{Lezioni.} Book V; chap. II, § 28 (vol. III, p. 200).} the types of continuous groups of Cremona transformations, and on the other hand there are also continuous series of transformations, such as the quadratic ones, which do not generate by multiplication continuous groups of finite dimensions. And similar series can also be recognized on a ruled surface of genus \(p>0\); where, however, there will be series
% scan: s228L
% !PRINTED 430
\weberpage{430}
of transformations operating on the surface in a non-transitive way, if \(p>1\).

With some additional restriction, the preceding observations can be inverted, that is, the possession of a continuous infinity of birational transformations makes it possible to characterize the family of ruled surfaces. This can be done from four points of view:

1) every surface which possesses a \textit{rational} \(\infty^1\) \textit{group of transformations into itself} belongs to the \textit{family of ruled surfaces}.

2) every surface which possesses a \textit{continuous group} (algebraic or transcendental) \textit{of transformations into itself} which leaves transitively%
\origfootnote{1}{s228L-1}{It is understood that it permutes the elements (curves) of \(|L|\) in infinitely many ways.} invariant a \textit{linear system}, \(|L|\), at least \(\infty^1\), of curves, belongs to the \textit{family of ruled surfaces}.

3) every surface \(F\) which possesses a \textit{continuous group} of transformations into itself, \textit{of dimension} \(r\geq 2\), operating \textit{intransitively} on the points of \(F\), or \textit{of dimension} \(r\geq 3\), operating \textit{transitively} on it, belongs to the \textit{family of ruled surfaces}.

4) every surface which possesses a \textit{continuous series} of transformations into itself, \textit{not generating a group} of finite dimension, belongs to the \textit{family of ruled surfaces}.

Theorem I) follows directly from the Noether-Enriques theorem (1898)\origfootnote{2}{s228L-2}{Cf. e.g. \textsc{Conforto}, \textit{Le superficie razionali.} Book II, chap. I, §§ 4, 5, 6 (pp. 245, 260, 262).} concerning surfaces with a pencil of genus \(p\geq 0\) of rational curves; %
because the trajectories of the rational \(\infty^1\) group form. precisely a pencil of rational curves. %
But it must be said that the result in question was obtained by P. \textsc{Painlevé} before \textsc{Enriques} had extended to the case \(p>0\) the theorem given by \textsc{Noether} for \(p=0\): indeed the circumstance that the rational curves \(C\) of a pencil are trajectories of an \(\infty^1\) group makes it possible, easily, to construct a directrix curve unisecant to the \(C\), by means of which the \(C\) themselves are rationally transformed into lines.

Theorem 2) states, in substance, that « a surface transformed into itself by a continuous group of homographies is rational or referable to a ruled surface », and in this projective form it was given, with some restriction, by \textsc{Enriques} (1893), and has received its most complete proof from \textsc{Fano} (1895, 96, 97)\origfootnote{3}{s228L-3}{Cf. \textsc{Conforto}, \textit{Le superficie razionali.} Op. cit. Book II, chap. II, § 37 (page 496).}.

Nevertheless it is worth giving a new proof of the theorem itself, one that also covers the case of a group \(\Gamma\) for which a linear system \(|L|\) (at least \(\infty^2\)) belonging to an
% scan: s228R
% !PRINTED 431
\weberpage{431}
involution remains invariant, and the case of a \(\Gamma\) which leaves invariant a linear pencil \(|L|\); indeed the proof which we propose will reduce the general theorem precisely to this last case.

We take as our starting point the observation (of \textsc{Picard}) that « a continuous group of birational transformations of an algebraic surface into itself is always contained in a \textit{continuous algebraic group} »%
\origfootnote{1}{s228R-1}{Indeed the conditions for a surface \(F\) to be transformed into itself by birational transformations of a given order are algebraic, and therefore serve to define an algebraic group, continuous or mixed. In this second hypothesis the mixed group contains within itself a continuous group.}.

If the surface \(F\) admits a continuous algebraic group of transformations into itself which leaves transitively invariant a linear pencil of curves \(|L|\), one may suppose that the curves \(L\) are permuted by \(\Gamma\) according to an \(\infty^1\) group, since otherwise it would suffice to put in place of \(\Gamma\) the subgroup of its transformations which leave fixed one or two particular curves \(L\).

This being so, we shall have to distinguish the following hypotheses:

\textit{a}) the group \(\Gamma\) is \(\infty^1\); then its trajectories, admitting \(\infty^1\) transformations into themselves, are curves \(K\) (a priori rational or elliptic)\origfootnote{2}{s228R-2}{Theorem of \textsc{Schwartz}. Cf. \textsc{Enriques-Chisini}, \textit{Lezioni.} Book V, chap. III, § 31 (vol. III, p. 309).} on which the \(L\) cut out an invariant rational involution, and therefore they are rational\origfootnote{3}{s228R-3}{Curves which admit infinitely many transformations into themselves leaving invariant a \(g^1\) are reduced, by means of a multiple of this series, to curves (of a certain space) with infinitely many projective transformations, and therefore they are rational. Cf. e.g. \textsc{Enriques-Chisini}, op. cit. Book V, chap. III.}: it follows that \(F\) is referable to a ruled surface. %

\textit{b}) the group \(\Gamma\) is at least \(\infty^1\), and operates transitively on the points of the surface. Then the curves of the pencil \(|L|\), transformed into themselves by \(\infty^1\) transformations of \(\Gamma\), are rational or elliptic.

If they are rational, \(F\) can be transformed into a ruled surface. If instead the \(L\) are elliptic, consider the elliptic \(\infty^1\) series \(\{D\}\) which is constituted by the transforms of a curve \(D\), chosen in an entirely general way, on \(F\): through every point \(P\) of \(f\) there pass a certain number \(i\) of curves \(D\) of \(\{D\}\), and \(P\) can vary on a curve \(K\) in such a way that the groups of \(i\,D\) issuing from it turn out equivalent; therefore the \(K\) thus defined must form an elliptic pencil \((K)\), without base points\origfootnote{4}{s228R-4}{Cf. Chap. IX, 12.}. %

Hence, among the transformations of our \(\Gamma\) permuting the elements (curves) of \((K)\), there will be \(\infty^1\) transformations, forming a subgroup \(\Gamma'\), which will have the \(K\) as trajectories. A priori, by \textsc{Schwarz}'s theorem, these \(K\) would turn out rational or elliptic; but the se%
% scan: s229L
% !PRINTED 432
\unskip\weberpage{432}%
\unskip cond case is excluded by observing that the transformations of \(\Gamma'\) must leave invariant the linear series \(g_n^1\) cut out on a \(K\) by the \(L\) of \(|L|\). It follows that \(F\) is referable to a ruled surface, \hfill Q.E.D.

Having thus established theor. 2) for the case in which the linear system \(|L|\), invariant under the group \(\Gamma\), is a pencil, it is easy to extend the proof to the case in which \(|L|\) is a linear system, \(\infty^r\), with \(r>1\). Indeed the transformations of \(\Gamma\) (which may be supposed to be a continuous algebraic group) operate on the elements (curves) of \(|L|\) as homographies of an \(S_r\), likewise forming an algebraic group \(\Gamma'\), and the subgroups of dimension \(\geq 1\) formed by all the transformations of \(\Gamma'\) permutable with a given \(\omega\) will also be algebraic. Now such a subgroup, say \(\overline{\Gamma}'\), leaves fixed the same united points as \(\omega\), and hence also a united line not entirely made up of united points: to which there corresponds a linear pencil of curves \(L\), transitively invariant under \(\overline{\Gamma}'\). We thus fall back into the case treated above.

Theorem 3)%
\origfootnote{1}{s229L-1}{Cf. \textsc{Castelnuovo-Enriques}, \textit{Sur les surfaces algébriques admettant un groupe continu de transformations birationelles en elles-mêmes.} (Comptes rendus, Paris, 29 July 1895).} is proved as follows. %

Suppose first that the group \(\Gamma\), \(\infty^r\), of the transformations of \(F\) operates intransitively on the surface, that is to say, transforms into themselves the curves \(C\) of a pencil on it. If \(\Gamma\) is (as may be supposed) an algebraic \(\infty^r\) group with \(r\geq 2\), one will obtain an algebraic \(\infty^1\) subgroup of it, \(\Gamma'\), by requiring the transformations of \(\Gamma\) to leave fixed \(r-1\) generic points of \(F\): say \(O, O'\ldots\).

Now the group \(\Gamma'\) — an algebraic \(\infty^1\) series of elements with infinitely many transformations into itself — will be, a priori, rational or elliptic. But it cannot be elliptic, because it contains an invariant series \(g_n^1\), of a certain order \(n>1\), which corresponds to the series of the points of \(F\) infinitely near to \(O\). Therefore it is rational and (in accordance with theor. 1) \(F\) is referable to a ruled surface.

To tell the truth, the preceding reasoning would give rise to an exception if the transformations of \(\Gamma'\) left fixed all the points of \(F\) infinitely near to \(O\); but in that case one will repeat the same argument with regard to the successive neighborhood of a point \(O_1\) near to \(O\) in the neighborhood of the first order, and similarly — where necessary — for the first neighborhood of a point \(O_s\) near to \(O\) in the neighborhood of order \(s\), which is not entirely made up of united points for \(\Gamma'\).

In an analogous way one proves the second case of the theorem, where one has an algebraic group \(\Gamma\), at least \(\infty^3\), which operates transitively on the points of \(F\). Referring to the simplest hypothesis, in which \(\Gamma\) has dimension 3, one will obtain an algebraic \(\infty^1\) subgroup, \(\Gamma'\), by requiring the transformations of \(\Gamma\) to leave fixed a point \(O\), and
% scan: s229R
% !PRINTED 433
\weberpage{433}
— as in the preceding case — one will recognize that \(\Gamma'\) is a \textit{rational} group.

If instead \(\Gamma\) has a dimension \(r>3\), one will similarly obtain a subgroup \(\Gamma'\) of it, rational and \(\infty^1\), by requiring the transformations of \(\Gamma\) to leave fixed a point \(O\) of \(F\), and then also a point \(O_1\) infinitely near to \(O\), etc.

Finally let us prove theorem 4), which answers a more recondite question, raised by E. \textsc{Picard}%
\origfootnote{1}{s229R-1}{\textsc{Castelnuovo-Enriques}, \textit{Sopra alcune questioni fondamentali nella teoria delle superficie algebriche.} Annali di Mat., 1901 (no. 19).}.

If the surface \(F\) admits a continuous series \(\Gamma\) — let us say \(\infty^1\) — of transformations into itself which does not generate a group (of finite dimension), this means that the products of the transformations of \(\Gamma\), taken \(r\) at a time, will in general constitute the transformations of an \(\infty^r\) series, \(\Gamma_r\), where \(r\) can be made arbitrarily large. Hence the transformations of \(\Gamma_r\) will change a curve \(C\), chosen in the most general way on \(F\), into the curves \(C_r\) of a continuous system \(\{C_r\}\) of dimension \(r\). If the curve \(C\) belongs to a linear system without base points, of genus \(\pi\) and degree \(n\), \(\{C_r\}\) will be formed of linear systems \(|C_r|\) of the same genus and degree, having a certain number of base points of multiplicity \(i\) (\(i=1,\,2\,\ldots.\)); and the virtual degree and the virtual genus of a \(|C_r|\) will be given respectively by %
\[
N_r=n+\Sigma h_i^2
\]
and
\[
\Pi_r=\pi+\Sigma\frac{h_i(h_i-1)}{2}
\]

Therefore we shall have
\[
N_r-2\Pi_r=n-2\pi+\Sigma h_i.
\]

But, since the degree \(N_r\) of \(\{C_r\}\) (which a priori is \(\geq r-1\)) grows with \(r\) as much as we please, we can choose \(r\) so that
\[
\Sigma h_i^2>(2\pi-n)^2
\]
and hence
\[
\Sigma h_i\geq\sqrt{\Sigma h_i^2}>2\pi-n.
\]

Thus it appears that the surface \(F\) possesses curves of genus \(\Pi_r\) and of degree \(N_r>2\Pi_r-2\), and therefore\origfootnote{2}{s229R-2}{Cf. Chap. X, § 2.} it is referable to a ruled surface. \hfill Q.E.D.

\edsection{3. Elliptic and hyperelliptic surfaces.}{3. Elliptic and hyperelliptic surfaces.}{sec.XI.3}

The analysis of the preceding section now allows us to «determine all the surfaces \(F\), not referable to ruled surfaces, which admit a continuous group of transformations into themselves ».

\printersignature{\textsc{Enriques F.} - \textit{Superficie algebriche.}\quad 28} %

% scan: s230L
% !PRINTED 434
\weberpage{434}
First of all one may suppose that this group \(\Gamma\) (not contained in a larger group) is \textit{algebraic}; then the group will give rise to one of the following cases:

1) \(\Gamma\) operates intransitively on the points of the surface \(F\), and it is then (§ 2) an elliptic \(\infty^1\) group which changes into themselves the elliptic curves \(C\) of a pencil, having a certain genus \(\varrho\geq 0\);

2) or else \(\Gamma\) is an \(\infty^2\) group of transformations which operates transitively on the points of \(F\).

\textit{First case.} – In the first case \textit{the surface \(F\) which admits an algebraic \(\infty^1\) group of birational transformations is an elliptic surface, and has arithmetic genus \(p_a=-1\), geometric genus \(p_g=\varrho\), and linear genus \(p^{(1)}=1\).}

Indeed the \(\infty^1\) transformations of \(\Gamma\) carry a curve \(L\), chosen in a general way on \(F\), into the \(\infty^1\) curves of an elliptic series \(\{L\}\); through every point \(P\) there will pass a certain number \(i\) of curves \(L\), and \(P\) will be able to move in \(F\) on a curve \(K\) in such a way that the group of the aforesaid \(i\,L\) remains equivalent to itself; hence the curves \(K\) will form on \(F\) an elliptic pencil \((K)\), on the curves of which the \(L\) cut out equivalent groups: for this reason the curves \(K\) are essentially distinct from the \(C\), trajectories of \(\Gamma\).

Now the surface \(F\), which contains a pencil of genus \(\varrho\) of elliptic curves \(C\) and an elliptic pencil of curves \(K\) (of genus \(>0\)), is elliptic (Chap. X, § 12) and possesses the aforesaid characters, which can easily be evaluated. %

To this end let us compute the number \(\Delta\) of the curves of \((K)\) endowed with a double point and express it by means of the Zeuthen-Segre invariant. Since \((K)\) is an elliptic pencil, one will find
\[
\Delta-4=12p_a-p^{(1)}+9.
\]

And since the curves of \((K)\) are transformed into one another by the transformations of \(\Gamma\), which leave none of them fixed:
\[
\Delta=0;
\]
hence
\[
\begin{aligned}
&12\,p_a-p^{(1)}+13=0,\\
&\quad p_a=-1,\qquad p^{(1)}=1.
\end{aligned}
\]

Finally, if \(\varrho=0\) it follows that \(p_g=0\), and if \(\varrho>0\) the groups of the \(g_{2\varrho-2}^{\varrho-1}\) within the pencil \((C)\) will constitute the canonical curves of \(F\) (Chap. X, § 12); therefore we shall have, in every case,
\[
p_g=\varrho.
\]

\textit{Note.} – The simplest example of elliptic surfaces of genus \(p_g>0\) is that of the surfaces \(F\) of determinant 1 representing the pairs of points of an elliptic curve and of a curve of genus
% scan: s230R
% !PRINTED 435
\weberpage{435}
\(p_g\), which, on the other hand, can be represented on a multiple cubic cylinder \(\Phi\) (cf. Chap. X, § 13). The elliptic surfaces of genus \(p_g>0\) and of determinant \(n>1\) can be constructed starting from the multiple \(F\) of determinant 1, or also from a multiple cubic cylinder. In substance one repeats, with few modifications, the considerations developed for the case \(p_g=0\).

It therefore appears that \textit{the elliptic surfaces for \(p_g>0\) give rise to infinitely many families corresponding to the following integer characters}, independent up to a certain point:

\textit{a}) the \textit{geometric genus} \(p_g\), genus of the pencil of elliptic curves \((K)\);

\textit{b}) the \textit{determinant} \(n\), number of the intersections of a curve \(K\) with a curve of the elliptic pencil \((C)\) (consisting of the directrix curves of the \(K\));

\textit{c}) the number of the \textit{multiple elliptic curves} (\(K=sK_s\)) which belong to the \textit{pencil} \((K)\) and the relative orders of multiplicity \((s_i)\): as for \(p_g=0\), one can relate these numbers to the plurigenera of the surface%
\origfootnote{1}{s230R-1}{Here interesting classificatory problems present themselves: for example « to determine the elliptic surfaces corresponding to the first values of \(\pi\) ». For \(\pi=2\), \(p_g=0\), we cite the Memoir of L. \textsc{Campedelli}, \textit{Intorno alle superficie ellittiche con un fascio di curve di genere due.} Rendic. Seminario Mat. di Padova, 1935.}.

\textit{d}) the genus \(\pi\) of the curves of the elliptic pencil \((C)\): \origfootmark{1}{s230R-1}
\[
\pi=1+n(p_g-1)+\Sigma\frac{(s_i-1)n}{2s_i};
\]

\textit{Second case.} – Having exhausted the discussion of case 1), let us now examine case 2), which concerns the surfaces that admit a continuous group, \(\infty^2\), \(\Gamma\) of transformations into themselves.

First of all we note that, if the transformations of \(\Gamma\) are not all permutable with one another, those transformations which are permutable with a given one will form an algebraic \(\infty^1\) subgroup of \(\Gamma\), and hence we shall be brought back to particularizations of case 1).

Therefore one may suppose that \(\Gamma\) is an \(\infty^2\) group, transitive on \(F\), consisting of permutable transformations. And it is further legitimate to add that it is transitive \textit{without exceptions} on the surface \(F\) without exceptional curves: otherwise there would be on \(F\) at least one united point \(O\) or one united curve \(\chi\), and, by fixing a point near to \(O\) or a point of the curve \(\chi\), one would obtain within \(\Gamma\) an algebraic \(\infty^1\) subgroup with a united point, whence \(F\) could be transformed into a ruled surface (§ 2). %

Finally it is easy to recognize that our \(\Gamma\) operates in a \textit{simply} transitive way on the points of \(F\), that is, that there cannot be a
% scan: s231L
% !PRINTED 436
\weberpage{436}
finite number $n>1$ of transformations of $\Gamma$ in which to a point $A$ there corresponds one and the same point $A'$ of the surface. Indeed this hypothesis is reduced to absurdity as follows.

Granted that the hypothesis is verified, consider, alongside $F$, the surface $F'$ (which provides the \textit{parametric} representation of the group $\Gamma$ in the sense of S. \textsc{Lie}) whose points correspond, without exception, to the elements (transformations) of $\Gamma$; since the transformations of $\Gamma$ operate on one another by multiplication, $\Gamma$ will operate on $F'$ in a simply transitive way: $F'$ is therefore a hyperelliptic, that is, abelian, surface (Chap. IX, \S~9). On the other hand there is between $F$ and $F'$ a $(1,n)$ correspondence in which to the points of $F$ there correspond on $F'$ the groups of an involution $I_n$; and so $F'$ is represented on the $n$-fold $F$ without branch points. It follows from this that the involution $I_n$ on $F'$ is generated by a finite abelian group $G_n$ of transformations of $\Gamma$, invariant under the $\infty^2$ transformations of $\Gamma$ itself; and therefore to the $n$ transformations of $\Gamma$ which make two groups of the involution $I_n$ correspond to each other there will correspond a single transformation of $F$ in which the homologous points $A$ and $A'$ correspond, contrary to the hypothesis, from which we started, that there were $n$ transformations of $\Gamma$ carrying $A$ into $A'$.

In conclusion we can say that \textit{the surfaces $F$, not referable to ruled surfaces, which admit an $\infty^2$ group of birational transformations} (case 2) are \textit{hyperelliptic surfaces}, that is, two-dimensional \textit{abelian} varieties, according to the definition of Chap. IX, \S~9.

Of these surfaces we have already evaluated the genus (Chap. IX, \S~10), which is
\[
p_g=1,
\]
and it is easy to find the other characters.

First of all, the \textit{hyperelliptic surface $F$} (without exceptional curves) \textit{cannot possess canonical curves $K$ of order $>0$.}

Indeed the $\infty^2$ group of the transformations of $F$, leaving $K$ fixed, would no longer carry its points into other points of the surface, and there would thus be an exception to the transitivity of $\Gamma$. Here it is also useful to point out that, a priori, $F'$ cannot possess a canonical system of dimension $p_g-1>0$, because, this linear system being transformed into itself by the transformations of $\Gamma$, $F$ itself would turn out to be reducible to a ruled surface (\S~2).

Now, since the surface $F$ possesses a canonical curve of order zero, its (absolute) linear genus will be
\[
p^{(1)}=1.
\]

It remains to evaluate its arithmetic genus $p_a$. For this purpose we shall have to prove that the irregularity of $F$ is
\[
q=p_g-p_a=2,
\]
so that
\[
p_a=-1.
\]

% scan: s231R
% !PRINTED 437
\weberpage{437}
Let us begin by proving that « the Jacobi variety $V_p$ corresponding to a curve $C$ of genus $p$, which is at once seen to have an \textit{irregularity} $\geq 0$, has precisely the irregularity \textit{$p$ and not higher}, that is, that there cannot be in $V_p$ an $\infty^r$ series, with $r>p$, of non-equivalent varieties $V_{p-1}$ ».

This proposition follows from the third criterion of equivalence (Chap. III, \S~13), which extends easily from surfaces to varieties of more dimensions. On the basis of this criterion it is legitimate to assert that two varieties $V_{p-1}$ belonging to $V_p$ which are non-equivalent must cut out non-equivalent $V_{p-2}$ on the $V_{p-1}$ which corresponds to the series of the groups of $p$ points of $C$ with a fixed point $A_1$, and hence also non-equivalent $V_{p-3}$ on the $V_{p-2}$ which corresponds to the groups of $p$ points of $C$ with two fixed points $A_1$ and $A_2$, etc.

Therefore the stated proposition will be proved once it is recognized that the $V_2$ corresponding to the series of the groups of $p$ points of $C$ with $p-2$ fixed points $A_1 A_2 \ldotp\ldotp\ldotp\ldotp A_{p-2}$, that is, the surface representing the pairs of points of the curve $C$ of genus $p$, has precisely the irregularity $p$ and not $>p$.

And for this it suffices to observe that an $\infty^r$ series of non-equivalent curves on such a surface gives rise on $C$ to an $\infty^r$ series of non-equivalent groups of points, so that one certainly has
\[
r\leq p \qquad\qquad \text{Q.E.D.}
\]

Once the proposition on the irregularity of a Jacobi variety $V_p$ has been established, one deduces that « an abelian variety of $q$ dimensions, $V_q$, contained imprimitively in the Jacobi variety $V_p$ $(q<p)$, that is, transformed into itself by an $\infty^p$ subgroup of the $\infty^q$ group of $V_p$, also has the irregularity $q$ ». %

Since a $V_{q-1}$ of $V_q$ is transformed by the relative abelian group $\Gamma_q$ into the $V_{q-1}$ of an $\infty^q$ family, the irregularity of $V_q$ is certainly $\geq q$, and it will suffice to prove that it cannot be $>q$.

Indeed, if a $V_p$ contains such a $V_q$, so that the group $\Gamma_p$ of the $\infty^p$ permutable transformations of $V_p$ contains an algebraic subgroup $\Gamma_q$, $\infty^q$, which transforms into themselves the varieties $V_q$ of an $\infty^{p-q}$ family of imprimitivity, then $\Gamma_p$ will also contain an algebraic subgroup $\Gamma_{p-q}$, $\infty^{p-q}$, which will leave invariant the $V_{p-q}$ of an $\infty^q$ family\origfootnote{1}{s231R-1}{This theorem (already stated in Chap. IX, \S~9) derives, as has been said, from the geometric interpretation of a theorem of \textsc{Picard-Poincaré} concerning reducible abelian integrals on a curve, an interpretation due to \textsc{Castelnuovo}. Cf. \textsc{Enriques-Chisini}, \textit{Lezioni}. Book VI. Chap. III. \S~38 (vol. IV, pp. 237-241).}.

Now, if one of the aforesaid $V_q$ within $V_p$ possesses a continuous $\infty^r$ system of non-equivalent $V_{q-1}$, these, by means of the $\infty^{p-q}$ transfor%
% scan: s232L
% !PRINTED 438
\unskip\weberpage{438}%
mations of $\Gamma_{p-q}$, will generate a continuous system of dimension $r+p-q$ of non-equivalent $V_{p-1}$, so that we shall have
\[
\begin{aligned}
r+p-q&\leq p,\\
r&\leq q.
\end{aligned}
\]

Therefore: $V_q$, contained in $V_p$, has the irregularity $q$ and not $>q$.

In particular, for $q=2$, an abelian $V_2$, that is, a hyperelliptic surface $F$ which belongs imprimitively to a Jacobi $V_p$ $(p\geq 2)$, being transformed into itself transitively by $\infty^2$ transformations of the relative $\Gamma_p$, will have the irregularity
\[
p_g-p_a=2
\]
and hence
\[
p_a=-1.
\]

But (Chap. X, \S~9) any hyperelliptic surface $F'$ can always be reduced to such an $F$, simple or multiple, contained imprimitively in $V_p$, and if one has an $F$ multiple according to a certain number $n>1$, this is, in any case, without branch curve; therefore one can evaluate the arithmetic genus of $F'$ by using Severi's formula, which connects the arithmetic genera $p_a=-1$ and $p'_a$ of the two surfaces $F$ and $F'$ in $(1,n)$ correspondence (Chap. V; \S~28); since the branch curve of the multiple $F$ is absent, this formula is written %
\[
24(p'_a+1)=24n(p_a+1),
\]
and therefore gives
\[
p'_a=p_a=-1.
\]

In conclusion \textit{a hyperelliptic surface has the characters}
\[
p_g=1,\qquad p^{(1)}=1,\qquad p_a=-1.
\]

Moreover, since it possesses only a canonical curve of order zero, its \textit{plurigenera will all be equal to one.}

\edsection{4. Characterization of elliptic and hyperelliptic surfaces by means of the values of the genera.}{4. Characterization of elliptic and hyperelliptic surfaces by means of the values of the genera.}{sec.XI.4}

The values of the genera $(p_a,p_g,p^{(1)},P_4)$, which we have found to belong to the elliptic and hyperelliptic surfaces, conversely make it possible to define these surfaces (not referable to ruled surfaces), and hence the entire family of surfaces possessing a continuous group of birational transformations into themselves.

Let us first prove that \textit{a surface $F$ with the genera}
\[
p_a=-1,\qquad p_g>1,\qquad p^{(1)}=1,
\]
\textit{is an elliptic surface.}

% scan: s232R
% !PRINTED 439
\weberpage{439}
Indeed the $\infty^{p_g-1}$ canonical curves of $F$ will be composed of the elliptic curves $K$ of a pencil $(K)$, having a certain genus $\varrho$; and a priori the $\infty^q$ $(q=p_g-p_a)$ non-equivalent curves $L$ of a complete $\infty^q$ system, constructed on $F$, will cut out on the $K$ $\infty^q$ groups which are equivalent or else belong to a series formed of $\infty^1$ linear series: in the first case the $L$ of $\{L\}$ are obtained from one another by addition and subtraction of curves $K$ of $(K)$, and one deduces
\[
\varrho=q=p_g+1;
\]
in the second case the system $\{L\}$ is formed of $\infty^1$ continuous $\infty^{q-1}$ systems, each of which is deduced from one of its curves by adding and subtracting curves $K$ of $(K)$, so that
\[
\varrho=p_g;
\]
and besides $(K)$ --- as we have done several times before --- one constructs on $F$ a second irrational, elliptic pencil of directrix curves of the $K$. Consequently the surface $F$ turns out to be elliptic (Chap. X, \S~12).

Now we see that it is precisely this second case that occurs, and not the first.

For this purpose let us compute the Zeuthen-Segre invariant of $F$ with respect to the pencil of elliptic curves $K$: we shall have
\[
\Delta-4=12p_a-p^{(1)}+9=-4,
\]
and therefore
\[
\Delta=0;
\]
whence it is inferred that the $K$ of $(K)$, never acquiring double points, are birationally identical with one another; in $(K)$ there will be found a certain number of multiple elliptic curves with the same modulus
\[
s_1K_1,\quad s_2K_2,\quad \ldotp\ldotp\ldotp\ldotp
\]
where
\[
s_i>1.
\]

This being so, the procedure explained in \S~11 of Chap. X (for the case of a pencil of genus $\varrho=1$) makes it possible to construct a pencil (a priori rational or elliptic) of directrix curves $C$, of genus $\pi>\varrho$, without base points. On a $C$ the canonical curves of $F$ will cut out the canonical series $g_{2\pi-2}$ (in general not complete). Since the $K$ cut out on a $C$ an involution of genus $\varrho$, which has as $s_i$-fold points the intersections with the $K_i$, the canonical curves of $F$ will be formed by the groups of the $g^{\varrho-1}_{2\varrho-2}$ within the pencil $(K)$ of genus $\varrho$, to which the curves $(s_i-1)K_i$ are added. And from this there results the genus of $F$
\[
p_g=\varrho \qquad\qquad \text{Q.E.D.}
\]

% scan: s233L
% !PRINTED 440
\weberpage{440}
Also for $p_g=0$ and $p_a=-1$, $p^{(1)}=1$, as for $p_g>1$ and $p^{(1)}=1$, the surface $F$ turns out to be elliptic (Chap. X, \S~12), except that it may be improperly elliptic, reducing to an elliptic ruled surface $(P_{12} = 0)$\origfootnote{1}{s233L-1}{Cf. X, 14.}. %

But the case $p_g=1$ gives rise to an exception: the surfaces whose genera have the values
\[
p_a=-1,\qquad p_g=1,\qquad p^{(1)}=1,
\]
are distributed into two families:

1) surfaces with elliptic canonical curve of order greater than zero;

2) surfaces with pure canonical curve of order zero.

The surfaces of the first family are elliptic, that is, they admit an elliptic $\infty^1$ group of transformations into themselves, while the surfaces of the second family are hyperelliptic, and not in general elliptic: this means that they admit an abelian $\infty^2$ group, $\Gamma$, of transformations into themselves, but only for particular moduli does there exist in $\Gamma$ an algebraic (elliptic) $\infty^1$ subgroup.

To justify what has been said, concerning cases 1) and 2), we must establish a

\textit{Lemma on the surface of irregularity $q=2$:} « a surface $F$ of irregularity $q=p_g-p_a=2$, not containing a pencil of genus two of curves, can be represented on a hyperelliptic surface $\Phi$, simple or multiple ».

To prove this lemma, consider on $F$ a (simple) continuous system $\{C\}$, formed of $\infty^2$ non-equivalent curves, and regard the elements (curves) of $\{C\}$ as « points » of a surface $\Phi$: which (as the Picard variety of $F$) will be a hyperelliptic surface. Between the surfaces $F$ and $\Phi$ there will exist a correspondence, in which to the points of $F$ (that is, to the $\infty^1$ systems of curves $C$ passing through them) there will correspond the curves $L$ of an $\infty^2$ system $\{L\}$ on $F$. We say that the $\infty^2$ $L$ on $\Phi$ cannot be equivalent to one another. Otherwise the $\infty^1$ $C$ of $\{C\}$ passing through a point $O$ of $F$ would cut out equivalent groups on another $C=\bar{C}$, and, by the third criterion of equivalence (Chap. III, \S~13), they would turn out to be equivalent to one another. %

To justify the assertion we must consider on the surface $F$ the series of the groups of points $G$ which are the sections of a particular $C=\bar{C}$, not passing through $O$, with the $\infty^1$ $C$ through $O$; if we denote by $i$ the index of this series, there will be $i$ groups $G$ having in common a point $P$ of $\bar{C}$ and $i$ groups $G$ having in common another point $P'$ of the same curve; and, by virtue of a criterion of \textsc{Severi-Castelnuovo}\origfootnote{2}{s233L-2}{Cf. \textsc{Enriques-Chisini}, \textit{Lezioni}. Book V, chap. IV, \S~42 (vol. III, p. 483)-}, it is necessary %
% scan: s233R
% !PRINTED 441
\weberpage{441}
and sufficient to prove that the two sets of groups $G$, $P(G)$ and $P'(G)$, thus defined on $\bar{C}$, are equivalent. For this purpose consider on $\Phi$ the curve $\bar{L}$ which corresponds to the point $O$ (or rather to the series $C(O)$ of the $C$ through $O$) and the point $A$ which corresponds to the curve $\bar{C}$: if, as we are supposing, the $\infty^1$ $L$ through $A$ are equivalent on $\Phi$, they will cut out on $\bar{L}$ equivalent groups of points, and so in particular the groups cut out by the two $L$ --- say $L_P$ and $L_{I'}$ --- corresponding to the points $P$ and $P'$ of $F$ will be equivalent on the said curve. But --- returning to the surface $F$ --- this means that the two groups of curves $C$ through $O$ and $P$ and through $O$ and $P'$, cutting out on $\bar{C}$ the groups of points $P(G)$ and $P'(G)$, are equivalent \textit{within the series of the $C$ through $O$}, and so form part of a rational $\infty^1$ series of groups of curves $C$; recalling that a rational series of groups of points on a curve always consists of equivalent groups\origfootnote{1}{s233R-1}{Cf. \textsc{Enriques-Chisini}, \textit{Lezioni}. Book V, chap. I, \S~10, Note (vol. III, p. 78).}, one deduces that $P(G)$ and $P'(G)$, sections of $\bar{C}$ with these groups of $C$, are equivalent \textit{on the said $\bar{C}$.} %

Now, having proved that the $\infty^2$ curves $L$ of $\Phi$ (corresponding to the points of $F$) cannot all be equivalent, we must still examine the hypothesis that every $L$ of $\{L\}$ forms part of an $\infty^1$ series of equivalent curves $L$, and recognize that, under this hypothesis, the surface $F$ contains a pencil of genus two. Indeed, under the said hypothesis there will be found on $F$ a pencil of curves $K$ to whose points there correspond equivalent curves $L$, and one can establish that the $C$ of a series $C(O)$, passing through an arbitrary point $O$ of the surface, and hence all the $C$, cut out equivalent groups on a $K$: for this purpose it suffices to repeat the reasoning made above, showing that the groups of curves common to $C(O)\ C(P)$ and $C(O)\ C(P')$ are equivalent (within the series $C(O)$), and hence that the groups of points $P(G)$ and $P'(G)$, sections of $K$ with the $C$ of $C(O)$, are also equivalent on the curve $K$.

In this way one verifies that the $\infty^2$ curves $C$ of $F$, non-equivalent to one another, cut out equivalent groups on every $K$; hence the continuous $\infty^2$ system $\{C\}$ is constructed starting from a linear system $|C|$, or from a curve $C$, by adding and subtracting curves $K$ of the pencil $(K)$: which therefore turns out to be of genus ($\varrho\geq 2$ and hence) $p_g-p_a=2$.

Thus the stated \textit{Lemma} is justified.

Now one can add that: \textit{the surfaces of irregularity $q=2$, with $p_g>0$, containing a pencil of genus two of curves} (of genus $\pi>0$) \textit{have $p_a\geq 0$.}

Indeed, if a surface $F$ of arithmetic genus $p_a=-1$ is given, which contains a pencil of genus two of curves $K$, let us prove that for it we have: $p_g>1$.

% scan: s234L
% !PRINTED 442
\weberpage{442}
Computing the Zeuthen-Segre invariant with respect to the pencil $(K)$, one finds
\[
\Delta-4+8(\pi-1)=12p_a-p^{(1)}+9=-3-p^{(1)}
\]
and hence
\[
p^{(1)}=1,\qquad \Delta=0,\qquad \pi=1;
\]
therefore the elliptic curves of the pencil $(K)$ turn out to be birationally identical with one another. Then (as in Chap. X, \S~11) one can construct a pencil (which a priori would be rational or elliptic and will then turn out to be elliptic) of curves $C$, without base points, directrices of the $K$: hence\origfootnote{1}{s234L-1}{. Cf. the observation in \S~8 of Chap. IX.} the canonical curves of $F$ will be composed of the canonical groups of the $K$ of $(K)$, and possibly also of fixed parts (more precisely of $(s-1)C_s$, where $sC_s=C$ is an $s$-fold multiple curve of the pencil $(C)$); one deduces %
\[
p_g\geq 2 \qquad\qquad \text{Q.E.D.}
\]

On the basis of the lemma on surfaces of irregularity $q=2$, and of the observation that the surfaces with $q=2$ containing a pencil of genus two have $p_g\geq 2$, and hence $p_a\geq 0$, we can say that a surface $F$ with $p_a=-1$ and $p_g=1$ can be represented on a hyperelliptic surface $\Phi$, simple or multiple, which will have a certain branch curve $D$. To $D$ there corresponds on $F$ the canonical curve; therefore, if the linear genus of $F$ is $p^{(1)}=1$, one of the two following cases will occur, which will bring us back to cases 1) and 2) mentioned above:

1) $D$ is an elliptic curve of order greater than zero, or

2) $D$ is a curve of order zero.

1) In the first case the $\infty^2$ transformations into itself of the hyperelliptic surface $\Phi$, belonging to the continuous group $\Gamma$, change $D$ (of degree zero) into the curves of a pencil, and each $D$ of this pencil (or each irreducible component of $D$) is invariant under an elliptic $\infty^1$ subgroup of $\Gamma$. It follows from this that the multiple surface $\Phi$ admits an elliptic $\infty^1$ group of transformations, that is, it is an elliptic surface, Q.E.D.

2) If instead on the multiple surface $\Phi$ (deprived of exceptional curves) there is no branch curve, $F$ is, like $\Phi$, an abelian, that is, hyperelliptic, surface (Chap. IX, \S~9); and thus case 2) occurs.

The two cases 1) and 2) can be distinguished by observing that, in case 1) (because $p_g=1$ and not $p_g>1$), to the elliptic curves of the pencil $(D)$ of $\Phi$, or rather to their irreducible components, there correspond on
% scan: s234R
% !PRINTED 443
\weberpage{443}
$F$ the curves $K$ of an elliptic pencil of elliptic curves, in which the number of the curves endowed with a double point is $\Delta=0$.

Therefore in this pencil $(K)$, formed of birationally identical curves, there will be at most a certain number of multiple curves $D_i$, each to be counted $s_i$ $(>1)$ times, and the particular curve $D$ which constitutes the branch curve of the multiple $\Phi$ will be composed of a certain number $r$ of these curves:
\[
D=\Sigma(s_i-1)D_i.
\]

It follows from this that in the pencil $(K)$ on $F$ the canonical curve, corresponding to $D$, is
\[
\Sigma\frac{s_i-1}{s_i}K.
\]

Hence, if $r>1$, there will be on $F$ a bicanonical system of dimension $\geq r-1$:
\[
|rK+\Sigma\frac{s_i-2}{s_i}K|,
\]
having as fixed components the curves which correspond to the $D_i$ of multiplicity $s_i>2$, counted $s_i-2$ times; and, in any case, if $r=1$, there will be on $F$ a quadricanonical system, at least $\infty^1$, of which the variable components are pairs of curves $K$, forming a $g^1_2$ within the elliptic $(K)$. %

In conclusion, in case 1), the elliptic surface $F$ $(p_a=-1,\ p_g=1,\ p^{(1)}=1)$ has the quadrigenus
\[
P_4>1,
\]
while the hyperelliptic surface (with the same $p_a=-1$, $p_g=1$ and $p^{(1)}=1$, but with canonical curve of order zero) will have all its plurigenera equal to one and, in particular, the \textit{quadrigenus}:
\[
P_4=1.
\]

Let us summarize the results obtained:

\textit{The surfaces, not referable to ruled surfaces} ($P_4+P_6>0$ or $P_{12}>0$), \textit{which admit a continuous group of birational transformations into themselves, are characterized globally by the values of the genera:}
\[
p_a=-1,\qquad p^{(1)}=1;
\]
\textit{they are always either elliptic surfaces (with an $\infty^1$ group), and $p_g>1$ or else $p_g=1$ and $P_4>1$; or hyperelliptic surfaces (with an $\infty^2$ group) if}
\[
p_g=P_4=1.
\]

% scan: s235L
% !PRINTED 444
\weberpage{444}
That hyperelliptic surfaces are not in general also elliptic, that is, that their group $\infty^2$ does not generally contain any algebraic (elliptic) subgroup $\infty^1$, can be seen by referring to the Jacobi surface that corresponds to a curve $C$ of genus two (the surface $F$ representing the pairs of points of $C$), or similarly to the abelian involutions on $F$ or to the multiple surfaces $F$ without branch curve.

The hyperelliptic $F$ thus constructed contains no elliptic curves, provided that $C$ neither contains an elliptic involution $\gamma_n^1$, of a certain order $n>1$, nor degenerates into a pair of elliptic curves; in the particular case in which the said $F$ contains an elliptic curve $C$, this will be carried by the transformations of the group $\Gamma$ of $F$ into the curves of an elliptic pencil $(C)$, and (in accordance with the Picard-Poincaré theorem on reducible integrals, interpreted geometrically by \textsc{Castelnuovo}) there will be on $F$ a second elliptic pencil $(K)$, likewise of elliptic curves: these are constructed as directrices of the pencil $(C)$ in the manner indicated in \S~11 of Chap. X. Therefore $F$ can be regarded in two different ways as an elliptic surface, admitting two elliptic groups $\infty^1$ of transformations, subgroups of the $\Gamma$ $\infty^2$, which have as trajectories the curves $C$ and the curves $K$. Finally we add that the said $F$, twice elliptic, of determinant $n>1$ (where the curves $C$ and the curves $K$ meet in $n$ points), can be represented on the particular hyperelliptic surface (and twice elliptic, of determinant \textit{one}) which is the Jacobi surface corresponding to the curve of genus two degenerate into two elliptic curves, a surface to be taken as $n$-fold. %

\textit{Remark.} -- The conditions with which we have defined the family of elliptic and hyperelliptic surfaces ($p_a=-1\quad p^{(1)}=1$) contain something redundant:

\begin{center}
\textit{the condition $p^{(1)}=1$ is a consequence of $p_a=-1$.} %
\end{center}

We shall prove this, for elliptic surfaces with arbitrary $p_g$, in \S~8; here the assertion will be justified as far as the surfaces with $p_g=1$ are concerned, that is, essentially for the hyperelliptic case.

Let us start from the conclusion reached earlier that: « the surfaces $F$ with genera $p_a=-1$ and $p_g=1$ can be represented on a hyperelliptic surface $\Phi$, simple or multiple ». If we have a multiple $\Phi$, there will be on it a branch curve $D$ of genus $p^{(1)}$ equal to the linear genus of $F$, and it is a matter of ascertaining that from the hypothesis $p^{(1)}>1$ there follows
\[
p_a\geq 0.
\]

Let us prove this in the simplest hypothesis, in which the branch curve $D$ of $\Phi$ is a simple curve, of genus $\varrho=p^{(1)}$, endowed with a certain number $\delta$ of nodes and a certain number $\tau$ of cu%
% scan: s235R
% !PRINTED 445
\unskip\weberpage{445}%
sps, the case to which the formula of Severi refers\origfootnote{1}{s235R-1}{To complete the proof it is therefore necessary to extend this formula to the case in which the said branch curve is composed in any way whatever of multiple parts.

That in this way one arrives at the stated result we are assured \textit{a priori} by a consideration (which rests in part on the theory of the simple integrals of the first kind belonging to a surface), as kindly communicated to us by G. \textsc{Castelnuovo}. It is this: for irregular surfaces not birationally equivalent to ruled surfaces, denoting by $I$ the Zeuthen-Segre invariant, one has
\[
I+4=13+12p_a-p^{(1)}\geq 0.
\]

Indeed, for $p_a<0$ (that is, $p_a=-1$) there follows
\[
p^{(1)}=1.
\]
}, which connects the arithmetic genera $p'_a$ and $p_a$ of the two surfaces $\Phi$ and $F$ in $(1,n)$ correspondence (cf. Chap. V, \S~8). The said formula gives us
\[
24(p_a+1)\geq 24n(p'_a+1)+6(\varrho-1)-2\tau,\origfootmark{2}{s235R-2}
\]\origfoottext{2}{s235R-2}{The case in which the inequality holds corresponds to the hypothesis in which there are on $D$ some fundamental points for the correspondence $(1,n)$.}
where
\[
p'_a=-1,\qquad \varrho>1,
\]
that is
\[
24(p_a+1)\geq 6(\varrho-1)-2\tau.
\]

\edFlush
But, on the hyperelliptic surface $\Phi$, the curve $D$ belongs at least to a system $\infty^2$, and possesses a characteristic series (cut out by the infinitely near curves) of degree %
\[
2\delta+3\tau\leq 2(\varrho+\delta+\tau-1).
\]

Therefore we shall have
\[
\tau\leq 2(\varrho-1) %
\]
and
\[
24(p_a+1)>0,
\]
that is
\[
p_a\geq 0. \qquad\qquad \text{Q.E.D.}
\]

Summing up: \textit{The surfaces with genera}
\[
p_a=-1\qquad \text{and}\qquad p_g=1
\]
\textit{have linear genus}
\[
p^{(1)}=1
\]
\textit{and, where they possess a canonical curve of order zero (which happens for $P_4=1$), are hyperelliptic surfaces.}

\textit{Historical note.} -- The analysis of the surfaces, not birationally equivalent to ruled surfaces, that admit a continuous group of transformations into themselves, be%
% scan: s236L
% !PRINTED 446
\unskip\weberpage{446}%
gins with the study of hyperelliptic surfaces made by \textsc{Picard} from the transcendental point of view\origfootnote{1}{s236L-1}{E. \textsc{Picard}, \textit{Mémoire sur la théorie des fonctions algébriques de deux variables indépendantes.} Journal de Mathématiques, 1889. (Chap. III). Cf. vol. II of the Traité of \textsc{Picard} and \textsc{Simart}, Paris 1906.}.

A brilliant application of the simple integrals of total differentials, which he himself introduced into science, leads the author to characterize the surfaces of the indicated kind that possess a group \(\infty^2\) of permutable transformations, by means of:

1) the existence of two simple integrals of the first kind;

2) the existence of a double integral of the first kind (which, according to \textsc{Clebsch} and \textsc{Noether}, means \(p_g=1\));

3) and the absence of a (pure) canonical curve of order greater than zero.

Condition 3) (which \textsc{Picard} noticed but, misled by an error of \textsc{Noether}, considered superfluous) entails the invertibility of the two simple integrals and hence the parametric representation of the surface by quadruply periodic (hyperelliptic) functions of two independent variables. %

From this result \textsc{Enriques}\origfootnote{2}{s236L-2}{\textit{Sulle superficie algebriche che ammettono un gruppo continuo di trasformazioni birazionali in se stesse.} Rendic. Circolo Mat. di Palermo, 1905.} derived the characterization of hyperelliptic surfaces by means of the values of the genera \(p_a=-1\), \(p_g=P_4=1\).

Later \textsc{Severi}\origfootnote{3}{s236L-3}{\textit{Sulle superficie algebriche che ammettono un gruppo continuo permutabile a due parametri di trasformazioni birazionali.} Atti Istituto Veneto, 1908.} translated the proof of Picard's theorem into geometric form, deducing directly the existence of the group \(\infty^2\) of transformations from the conditions \(p_a=-1\), \(p_g=1\) and canonical curve of order zero, as has also been done here.

The case of surfaces with an elliptic group \(\infty^1\), already encountered by \textsc{Picard}, is exhausted, from the transcendental point of view, by the analysis of \textsc{Painlevé}\origfootnote{4}{s236L-4}{P. \textsc{Painlevé}, \textit{Leçons sur la théorie analytique des équations différentielles.} Stockholm, 1895 (p. 270).}, who defines elliptic surfaces by means of a parametric representation by functions elliptic in one parameter and algebraic in another, and thus succeeds in classifying all surfaces, not birationally equivalent to ruled surfaces, possessing a continuous group of transformations into themselves. The geometric translation of this result and the classification of the whole family of surfaces by means of the values of the genera was achieved by \textsc{Enriques} in the cited Memoir of 1905\origfootnote{5}{s236L-5}{\textit{Sulle superficie algebriche che ammettono un gruppo continuo etc.} l. c. In part already in the preceding Memoir, likewise published in the Rendiconti del Circolo Mat. di Palermo: \textit{Sulle superficie di genere zero} (1905).}.

% scan: s236R
% !PRINTED 447
\weberpage{447}
\edsection{5. Surfaces with genera \(p_g=1\) and \(p^{(1)}=1\).}{5. Surfaces with genera p\_g=1 and p\textasciicircum{}(1)=1.}{sec.XI.5}

The results established in the preceding sections now allow us to classify the surfaces of geometric genus \(p_g=1\) having linear genus \(p^{(1)}=1\). These surfaces can be distributed into three families:

1) Regular surfaces of genus \(p_a=p_g=1\) with canonical curve of order zero (characterized by the value of the bigenus \(P_2=1\)), which give rise to infinitely many classes of surfaces (with 19 moduli) with canonical hyperplane sections of genus \(\pi\), of order \(n=2\pi-2\), in the space \(S_\pi(\pi=2,\ 3\ldotp\ldotp\ldotp\ldotp)\) (Cf. Chap. VII, \S~2).

2) Hyperelliptic surfaces with genera \(p_a=-1\) and \(p_g=1\), likewise possessing a canonical curve of order zero (characterized by the value of the quadrigenus \(P_4=1\)), which again give rise to infinitely many classes of normal surfaces with canonical section curves of genus \(\pi\), of order \(n=2\pi-2\), in the space \(S_{\pi-2}\) (these classes depending, as we shall say, on three moduli).

3) Surfaces possessing a (rational or elliptic) pencil of elliptic curves: this property belongs to the surfaces 1) and 2) only for particular values of their moduli.

In order to prove that « \textit{the three families of surfaces} 1), 2), 3) \textit{exhaust the surfaces having genera \(p_g=1\) and \(p^{(1)}=1\)} » it is appropriate partly to recall observations already made earlier in this treatise, and partly to recognize some facts not yet observed. These are the following propositions:

\textit{a}) the surfaces of geometric genus \(p_g=1\) have arithmetic genus
\[
p_a\ge -1,
\]
that is
\[
p_a=1 \qquad \text{or} \qquad p_a=0 \qquad \text{or} \qquad p_a=-1;
\]

\textit{b}) the (regular) surfaces for which \(p_g=p_a=1\), with canonical curve of order zero, are defined by the value of the bigenus \(P_2=1\) (Chap. VII, \S~2); on the other hand, the surfaces of this family that have linear genus \(p^{(1)}=1\) and possess a (pure) canonical curve of order greater than zero, having \(P_2>1\), will also contain a (linear) pencil of elliptic curves, bicanonical or components of the bicanonical ones;

\textit{c}) the surfaces for which \(p_g=1\), with canonical curve of order zero, have arithmetic genus \(p_a=1\) or else \(p_a=-1\) (Chap. IX, \S~8) and, if \(p_a=-1\), also have linear genus \(p^{(1)}=1\), and are
% scan: s237L
% !PRINTED 448
\weberpage{448}
hyperelliptic surfaces or elliptic surfaces, with a pencil of elliptic curves, according as \(P_4=1\) or \(P_4>1\) (\S~4);

\textit{d}) the surfaces for which \(p_g=1\), \(p_a=0\) and \(p^{(1)}=1\) (which by \textit{b}) must contain a canonical curve of order greater than zero) possess an elliptic pencil of paracanonical curves, and therefore have bigenus \(P_2>1\).

Among these propositions, \textit{b}) and \textit{c}) recall knowledge already acquired; therefore it will suffice to prove \textit{a}) and \textit{d}), which bring something new.

Since prop. \textit{a}) is already established for elliptic and hyperelliptic surfaces, it will suffice to prove it under the hypothesis that the surface \(F\) (of genus \(p_g=1\)) does not possess a continuous group of transformations into itself (later on, \S~7, this proposition will appear as a particular case of a more general theorem according to which the arithmetic genus of surfaces not birationally equivalent to ruled surfaces is always \(p_a\ge-1\)).

Let \(F\) be a surface of genus \(p_g=1\) that does not admit a continuous group of transformations, and let \(q=p_g-p_a\) be its irregularity: we want to reduce to absurdity the hypothesis
\[
q>2,
\]
that is
\[
p_a<-1.
\]

To this end, assuming \(q>2\), consider the Picard variety \(V_q\) corresponding to a continuous system \(\infty^q\) of non-equivalent curves \(C\), taken on our \(F\). As we have seen (Chap. IX, \S~11), provided \(F\) does not contain an irrational pencil of curves \(K\), a series \(\infty^1\) of curves \(C\) of \(\{C\}\) leads to the construction of a surface, simple or multiple, identical with \(F\), within the said \(V_q\). Moreover, with a view to showing that \(p_g\ge q\), we may assume that there is in \(V_q\) a simple surface identical with \(F\): for in the case of a multiple surface, recalling the \textit{remark} of \S~4, one sees that if this surface admits a continuous group of transformations into itself, while \(F\) does not possess such a group, the arithmetic genus of \(F\) turns out to be \(p_a\ge0\), and consequently, for \(F\) itself, \(p_g\ge q\).

So then, let there be in \(V_q\) a surface identical with \(F\) (which we again denote by \(F\)) that does not admit a continuous group of transformations into itself; by means of the \(\infty^q\) transformations of \(V_q\) our \(F\) is transformed into a series \(\infty^q\) of analogous surfaces, say \(\{F\}\). For simplicity let us first suppose \(q=3\). Since \(V_q=V_3\) has geometric genus one, the characteristic linear system \(\infty^2\) of \(F\), cut out on an \(F\) by the infinitely near surfaces, will consist of special curves (residual to the intersection with the canonical variety of \(V_3\) and, hence, canonical), so that for \(F\) there will result
\[
p_g\ge3. %
\]

% scan: s237R
% !PRINTED 449
\weberpage{449}
The reasoning extends to the case \(q>3\). In this case we shall take within the series \(\{F\}\) a series \(\infty^{q-3}\) of surfaces \(F\), which constitutes a variety \(V_{q-1}\), and we shall suppose in general that such a variety is contained in a linear system \(|V_{q-1}|\) of dimension \(r\) (as large as we like); then we shall construct the continuous system \(\infty^{r+q}\) \(\{V_{q-1}\}\) generated by \(|V_{q-1}|\) by means of the \(\infty^q\) transformations of \(V_q\). This continuous system has a characteristic linear system \(\infty^{r+q-1}\) of special or canonical varieties, cut out on a \(V_{q-1}\) by the infinitely near varieties; therefore the said \(V_{q-1}\) have genus
\[
P\ge r+q.
\]

Now let us evaluate the infinity of the \(V_{q-1}\) of \(\{V_{q-1}\}\) that contain an \(F\): there are \(\infty^q F\), and \(\infty^{r+q}V_{q-1}\) each containing \(\infty^{q-3}F\), so that the dimension of the system \(\{F\}\) turns out to be (in the most unfavourable hypothesis):
\[
r+2q-3-x=q
\]
where \(x\) denotes the infinity of the system of the \(V_{q-1}\) containing an \(F\): \(x=r+q-3\). It follows from this that the characteristic system of \(\{V_{q-1}\}\) in a \(V_q\) cuts out on an \(F\) contained in it a linear system (special or canonical) whose dimension is at least:
\[
r+q-1-(x-1)=3.
\]

And therefore one deduces that the genus of \(F\) is
\[
p_g\ge4.
\]

In conclusion we have proved that a surface \(F\) of genus \(p_g<3\), which does not contain an irrational pencil of curves, cannot have irregularity
\[
q>2;
\]
in particular an \(F\) of genus \(p_g=1\) will have arithmetic genus
\[
p_a\ge-1.
\]

The conclusion extends to the hypothesis that \(F\) contains an irrational pencil of curves \((C)\), because the computation of the Zeuthen-Segre invariant gives, in relation to this pencil,
\[
\begin{aligned}
\Delta-4&\ge12p_a-p^{(1)}+9,\\
12(p_a+1)-p^{(1)}&\le0,\\ %
p_a&\ge-1
\end{aligned}
\]
and for \(p_a=1\), \(p^{(1)}=1\). %

Theorem \textit{d}) is proved as follows.

Consider on \(F\) an (irreducible) linear system \(|C|\), of a certain genus \(\pi\).

Its adjoint system \(|C'|\) will have dimension
\[
p_a+\pi-1=\pi-1,
\]
\printersignature{\textsc{Enriques F.} - \textit{Superficie algebriche.}\quad 29}
% scan: s239L
% !PRINTED 450
\weberpage{450}
and will cut out on a \(C\) a series \(g^{\pi-2}_{2\pi-2}\), of deficiency one, contained in the canonical series. Moreover the said \(|C'|\) will be contained in a series \(\infty^1\) of para-adjoint linear systems \(|\overline{C}'|\), likewise cutting out on \(C\) a series \(g^{\pi-2}_{2\pi-2}\). Therefore there will be an elliptic paracanonical curve defined as the difference
\[
\overline{C}'-C,
\]
and the analogous paracanonical curves \(K\) will form on \(F\) an elliptic pencil: let us add that every curve \(2K\) will therefore belong to a linear system \(\infty^1\) (at least), and hence \(P_2>1\), \hfill Q.E.D.

\textit{Remark.} -- An example of a surface having genera \(p_a=0\), \(p_g=1\), \(p^{(1)}=1\), and possessing an elliptic pencil of elliptic curves, is given by the simplest paraelliptic surfaces, which we encountered in \S~11 of Chap. X.

It may be interesting to study and classify the surfaces with \(p_a=0\) and \(p_g=1\) having linear genus \(p^{(1)}=1\), or also \(p^{(1)}=2\ldotp\ldotp\ldotp\ldotp\)

The results obtained can be summarized by saying: \textit{The surfaces, not birationally equivalent to ruled surfaces \((P_{12}>0)\), that possess a canonical or pluricanonical curve of order zero} (and hence linear systems without base points, of genus \(\pi\) and degree \(n=2\pi-2\)) \textit{are defined by the condition}
\[
P_{12}=1,
\]
\textit{which entails}
\[
p^{(1)}=1
\]
\textit{and gives rise to the following families:}

1) \textit{\(p_g=0\) and \(p_a=0\), regular surfaces with bicanonical curve of order zero \((P_2=1)\), represented by the surface of the \(6^{\mathrm{th}}\) order that passes doubly through the edges of a tetrahedron} (Chap. VII, \S~1);

2) \textit{\(p_g=0\) and \(p_a=-1\), elliptic surface of the types} \(\mathrm{I}_a\ \mathrm{I}_b\ \mathrm{II}_a\ \mathrm{II}_b\ \mathrm{III}_a\ \mathrm{III}_b\ \mathrm{III}_e\) (Chap. X, \S\S~16-19); %

3) \textit{\(p_g=1\) and \(p_a=1\) \((P_2=1)\), regular surfaces with genera one} (Chap. VII, \S~2);

4) \textit{\(p_g=1\) and \(p_a=-1\) \((P_4=1)\) hyperelliptic surfaces} (\S\S~4 and 6 of this chapter).

\edsection{6. Note on the geometric theory of hyperelliptic surfaces.}{6. Note on the geometric theory of hyperelliptic surfaces.}{sec.XI.6}

Among the families of surfaces with \(P_{12}=1\), marked with nos. 1) 2) 3) 4) at the end of the preceding section, it is appropriate to distinguish: on the one hand the surfaces of genus \(p_g=1\), which are 3) and 4) (\(p_g=1\),
% scan: s239R
% !PRINTED 451
\weberpage{451}
\(p_a=1\) and \(p_a=-1\)), and on the other the surfaces of genus \(p_g=0\) (regular with \(P_2=1\), or elliptic), which are 1) and 2). The latter constitute a finite number of families, to each of which a certain number of moduli belongs: precisely: 10 moduli to 1), 2 moduli or 1 modulus to 2). On the other hand the surfaces 3) and 4) give rise to \textit{infinitely many families}, distinguished by the values of (at least) one arithmetic character, to each of which a certain number of moduli belongs.

What is asserted here has been proved for the surfaces with all genera one (the 3)) in Chap. XII, \S~2; but it also holds for the hyperelliptic surfaces 4), as must be explained here. %

The \textit{classification of hyperelliptic surfaces} results first of all, from the transcendental point of view, from their parametric representation by means of 4-fold periodic functions of two independent variables\origfootnote{1}{s239R-1}{Cf. e.g. \textsc{Enriques-Chisini}, \textit{Lezioni.} Book VI, chap. III, \S~39 (vol. IV, p. 242).}. Their primitive normal periods can be reduced to the table
\[
\begin{tabular}{cccc}
\(1\) & \(0\) & \(g\) & \(h\) \\
\(0\) & \(\dfrac{1}{\delta}\) & \(h\) & \(g',\)
\end{tabular}
\]
where \(\delta\) is an integer character (the so-called \textit{divisor} of the corresponding surfaces). Now for \(\delta=1\) one has the \textit{type surface}, which is the Jacobi surface, \(F_j\), representing the variety of the \textit{pairs of points of the curve of genus two}; whereas for \(\delta>1\) one has a \(\delta\)-fold \(F_j\) without branch curve, or else an involution of order \(\delta\) on the same surface \(F_j\).

The classification of hyperelliptic surfaces thus obtained has an algebro-geometric meaning, and so the problem arises of justifying it from this point of view. Let us say briefly to what order of considerations one is led in this way. %

First of all one can determine the characters of the Jacobi surface \(F_j\):
\[
p_g=1,\qquad p_a=-1
\]
and canonical curve of order zero (hence \(P_4=1\)). This is obtained by representing \(F_j\) on the double Kummer surface, of the \(4^{\mathrm{th}}\) order with 16 double points, as was seen in \S~8 of Chap. V\origfootnote{2}{s239R-2}{Cf. also \textsc{Enriques-Chisini}, \textit{Lezioni.} Bk. V. Chap. III, \S~40 (Vol. IV p. 24).}. This representation is based on the knowledge of the (involutory) transformations of the \(2^{\mathrm{nd}}\) kind of the variety of pairs of points of the curve \(C\) of genus two, such as are obtained by associating with each other the pairs that are resid%
% scan: s240L
% !PRINTED 452
\unskip\weberpage{452}%
\unskip ual with respect to a \(g^2_4\) taken on \(C\): an involution so defined is birationally identical with a Kummer surface, of the \(4^{\mathrm{th}}\) order, having 16 double points which correspond to the 16 pairs of points of \(C\) arising from the bisection of the \(g^2_4\).

Now on the Jacobi surface \(F_j\) the products of the said involutions \(I\) give rise to the transformations of a continuous abelian group \(\Gamma\); and there are cyclic transformations of \(\Gamma\) which generate finite groups \(G_\delta\) contained in it: such a \(G_\delta\) (let it be, e.g., a cyclic \(G_\delta\)) gives rise on \(F_j\) to an involution \(I_\delta\) of order \(\delta\), without double points, to which there corresponds in general a hyperelliptic surface \(F_\delta\) different from \(F_j\). One thus obtains the hyperelliptic surfaces of divisor \(\delta>1\). And one can see that they can also be produced starting from \(F_j\) taken \(\delta\)-fold, without branch curve; for this purpose one must recognize that with the groups of an involution \(I_\delta\) one can construct an involution of higher order, say \(r\delta\), which turns out to be birationally identical with \(F_j\).

In what has been said there is set forth the outline of a \textit{geometric theory of hyperelliptic surfaces}, based on the study of the cyclic transformations belonging to the group \(\Gamma\) of the Jacobi surface \(F_j\): such a study presents itself as an extension of the study of cyclic involutions on elliptic curves, which is found developed in the « Lezioni » of \textsc{Enriques-Chisini}\origfootnote{1}{s240L-1}{Book V, chap. III, \S\S~27-28 (vol. III, pp. 251-298).}.

But the consideration of hyperelliptic surfaces from the geometric point of view also confronts us with other interesting problems.

Indeed we have proved (Chap. IX, \S~9) that an abelian variety \(V_q\) can be regarded as identical with a simple or multiple variety contained imprimitively within the Jacobi variety \(V_p\) that corresponds to a curve of genus \(p\ge q\); but it does not hold in general that one can always take \(p=q\), because the number of moduli on which \(V_q\) depends is greater than that, \(3q-3\), of the curves of genus \(q\)\origfootnote{2}{s240L-2}{Cf. \textsc{Enriques-Chisini}, \textit{Lezioni.} Book V, chap. III, \S~33 (vol. III, p. 355) and Book VI, chap. III, \S~35 (vol. IV, p. 216).}. Now the transcendental theory shows us that hyperelliptic surfaces (with a given divisor) depend on \textit{three moduli}, and therefore they all reduce to the type of the Jacobi surface \(F_j\), simple or multiple, and, in the latter case, also to involutions (without double points) on \(F_j\).

The matter can be recognized from the geometric point of view in the way that we briefly indicate. On a hyperelliptic surface \(F\) (for which \(p_a=-1\)) a curve of genus \(\pi\) belongs, in general, to a linear system of dimension \(\pi-2\), in which there will be contained
% scan: s240R
% !PRINTED 453
\weberpage{453}
a finite number of curves \(C\), endowed with \(\pi-2\) double points, and therefore of genus two.

By means of the transformations of the first kind of \(F\), belonging to the group \(\Gamma\), one will then obtain on \(F\) a continuous system \(\infty^2\) \(\{C\}\) of curves of genus two, birationally identical with one another, whose \textit{three moduli} will furnish the moduli of \(F\). Now a \(C\) of this system is invariant under \textit{one} (involutory) transformation of the second kind\origfootnote{1}{s240R-1}{Referring to the parametric representation of the group \(\Gamma\), a transformation of the second kind can be made to correspond to that in which a \(\pi\) of \(\Gamma\) is changed into its inverse \(\pi^{-1}\). It is easy to convince oneself that, even in the case in which the curve \(C\) of genus two possesses singular transformations into itself, there cannot be several transformations of the \(2^{\mathrm{nd}}\) kind of \(F\) that leave it invariant.} of \(F\), which induces on \(C\) its \(g^1_2\); hence the group \(\Gamma\) is birationally identical both with the surface \(F\) (conceived as a locus of points) and with the system \(\{C\}\) regarded as a variety \(\infty^2\) of elements (curves \(C\)). It follows from this that \(F\) turns out to be birationally identical with the series \(\infty^2\) of the groups of points cut out by the curves of \(\{C\}\) on a particular \(C\): if the groups of this series are all non-equivalent, \(F\) turns out to be identical with the Jacobi surface that corresponds to the said \(C\); if instead for every group of the said series there is a certain number \(\delta\) of equivalent groups, the hyperelliptic \textit{surface} \(F\) turns out to be birationally \textit{identical with an involution \(I_\delta\) of order \(\delta\) on the Jacobi surface.} %

The theorem established takes on an interesting geometric meaning in relation to various constructions that can be given of hyperelliptic surfaces. Indeed one can construct a hyperelliptic surface starting from the Jacobi variety corresponding to a curve \(C_p\) of genus \(p\). When on this \(C_p\) there are two abelian integrals reducible to 4 periods (according to the geometric interpretation of the Picard-Poincaré theorem due to \textsc{Castelnuovo})\origfootnote{2}{s240R-2}{Cf. \textsc{Enriques-Chisini}, \textit{Lezioni.} Book VI, chap. III, \S~38 (vol. IV, p. 234)-} %
the corresponding Jacobi variety \(V_p\) will contain an algebraic subgroup \(\infty^2\) of the group \(\Gamma\) of the transformations of the first kind, and one will therefore have an abelian variety (that is, a hyperelliptic surface) contained imprimitively in \(V_p\). The preceding theorem tells us that the hyperelliptic surfaces constructed in this way can also be found again as simple or multiple Jacobi surfaces, or as involutions on the Jacobi surfaces, corresponding to the curve of genus \(p=2\).

It is useful to illustrate this property by referring to the simplest case, in which we are dealing with the hyperelliptic surfaces contained imprimitively within the Jacobi variety \(V_3\) that corresponds to a curve \(C_3\) of genus three. Here the group \(\Gamma\) of the transformations of
% scan: s241L
% !PRINTED 454
\weberpage{454}
the first kind of \(V_3\) will contain two complementary algebraic subgroups: a \(\Gamma_2\ \infty^2\), which gives rise to a family of imprimitivity consisting of \(\infty^1\) hyperelliptic surfaces \(F_2\), and a \(\Gamma_1\ \infty^1\), which gives rise to a family of imprimitivity consisting of \(\infty^2\) elliptic curves \(C_1\), the \(F_2\) and the \(C_1\) meeting in a certain number of points. And it is easy to prove that each \(C_1\) is a unirational transform of \(C_3\), that is, it corresponds to an elliptic involution \(\gamma_n^1\) (of a certain order \(n\)) belonging to \(C_3\).

Therefore one will be able to construct a hyperelliptic surface \(F_2\), contained imprimitively in a Jacobi variety \(V_3\), starting from a curve \(C_3\) of genus three to which there belongs an elliptic involution \(\gamma_n^1\); for this purpose it will suffice to construct the \(C_3\) that corresponds to a cubic \(K\) taken \(n\)-fold, on which 4 branch points are chosen arbitrarily, with the corresponding substitutions among the branches\origfootnote{1}{s241L-1}{Cf. the existence theorem for multiple curves in \textsc{Enriques-Chisini}, \textit{Lezioni.} Book V, chap. IV, \S~38 (vol. III, p. 427).}. It will be noted that the \(C_3\) containing an elliptic \(\gamma_n^1\) will depend on 4 moduli: the absolute invariant of \(K\) and the \textit{distances} of 3 branch points from the remaining one, on the same elliptic curve \(K\).

Setting aside the absolute invariant of \(K\), which corresponds to that of the \(C_1\) on \(V_3\), there remain 3 moduli belonging to the complementary surfaces \(F_2\).

Let us add that the Jacobi \(V_3\) constructed above will contain, in general, an involution consisting of the groups of points of intersection of the \(F_2\) and the \(C_1\): the variety \(\overline{V}_3\) representing these groups, which is a unirational transform of \(V_3\), will be the Jacobi variety of a curve of genus 3 degenerate into an elliptic curve and a curve of genus two; so then (in accordance with the theorem set forth above), the hyperelliptic surfaces \(F_2\), which form a system of imprimitivity for the group of the transformations of the first kind of \(V_3\), turn out to be birationally identical with an involution on one of the analogous \(\overline{F}_2\), which is the Jacobi surface corresponding to the aforesaid curve of genus two.

May the indications given above suffice to point out the interest, the beauty and the richness of the properties that belong to hyperelliptic surfaces.

We shall only add that the study of the irregular hyperelliptic involutions on a hyperelliptic surface \(F\) (and in particular on the Jacobi surface \(F_j\)) extends naturally into that of other involutions, irregular elliptic or regular with genera one, which can belong to \(F\) and lead to surfaces to which there still belongs a
% scan: s241R
% !PRINTED 455
\weberpage{455}
parametric representation by hyperelliptic functions \textit{of rank} \(r>1\)\origfootnote{1}{s241R-1}{Where to each point of the surface there correspond \(r\) points within the prismatoid of the periods.}.

The irregular elliptic involutions (which belong to \(F\) with particular moduli) give rise to the elliptic surfaces with \(p_a=-1\), \(p_g=0\) and \(P_{12}=1\), possessing a virtual canonical curve of order zero, as has already been mentioned in \S\S~16-19 of Chap. X.

The regular involutions of order 2 belonging to a hyperelliptic surface of divisor \(\delta\geq1\) lead in general to surfaces (hyperelliptic in an extended sense) of rank \(r=2\), and precisely: for \(\delta=1\), to the Kummer surfaces of the \(4^{\mathrm{th}}\) order with 16 double points (Chap. V, \S~8); and for \(\delta>1\), to a series of analogous surfaces of order \(n=2\pi-2\), with canonical hyperplane sections of genus \(\pi\), in the space \(S_\pi\) \((\pi=4,5,6,\ldotp\ldotp\ldotp\ldotp)\), endowed with 16 double points. But for particular moduli there also occur on an \(F\) \textit{singular regular involutions}, which are generated by a finite group of birational transformations of the said \(F\); and on the basis of this generation (proved by \textsc{Enriques-Severi}) they have been classified in their entirety by the memoirs already cited, which complement each other, of \textsc{Enriques} and \textsc{Severi}\origfootnote{2}{s241R-2}{\textit{Mémoire sur les surfaces hyperellyptiques.} Acta Mathemat., vols. 32 and 33 (1908).} %
and of \textsc{Bagnera} and \textsc{De Franchis}\origfootnote{3}{s241R-3}{\textit{Le superficie algebriche le quali ammettono una rappresentazione parametrica mediante funzioni iperellittiche di due argomenti.} Memorie della Soc. It. d. Scienze (dotta dei XL), 1908.}. %

Finally we shall recall that this order of considerations extends from hyperelliptic surfaces to regular surfaces with all genera one, on which it is likewise possible to find regular involutions of the same genus \(p_g=1\) (or with \(p_g=0\) and bigenus \(P_2=1\)), which are generated by a finite group of birational transformations. In this connection it will suffice to mention the interesting works of the Belgian school, and particularly of \textsc{Godeaux}.

\edsection{7. General classification of algebraic surfaces\origfootnote{4}{s241R-4}{Cf. \textsc{Enriques}, \textit{Sulla classificazione delle superficie algebriche in particolare di genere lineare \(p^{(1)}=1\).} Rendic. Acc. Lincei, February 1914.}.}{7. General classification of algebraic surfaces.}{sec.XI.7}

We are now in a position to gather into a synthesis the principal results of the theory of surfaces, by setting forth a general classification of algebraic surfaces according to the value of the twelve-genus \(P_{12}\): in what follows this classification can be narrowed down more closely, as we shall see in \S\S~8 and 9.

% scan: s242L
% !PRINTED 456
\weberpage{456}
We distinguish algebraic surfaces according to the values of \(P_{12}\) and — subordinately — of the (absolute) linear genus \(p^{(1)}\).

We shall have \textit{four great families of surfaces}:

\(A)\) \textit{for \(P_{12}=0\), the surfaces birationally equivalent to ruled surfaces} (rational or not);

\(B)\) \textit{for \(P_{12}=1\), the surface with canonical curve}, effective or virtual, \textit{of order zero};

\(C)\) \textit{for \(P_{12}>1\) and \(p^{(1)}=1\), the surfaces that contain infinite linear systems of canonical or pluricanonical curves composed of the elliptic curves of a pencil};

\(D)\) \textit{for \(P_{12}>1\) and \(p^{(1)}>1\), the surfaces that admit canonical or pluricanonical surfaces as projective model.}

It is appropriate to recall and explain the \textit{fundamental and characteristic properties} that distinguish these four families of surfaces:

\(A)\) \(P_{12}=0\). Surfaces birationally equivalent to ruled surfaces (rational or not): they possess linear systems of curves of genus \(\pi\) and degree \(n>2\pi-2\), exceptional curves, infinite in number, that cannot be eliminated, rational groups \(\infty^1\) of transformations, and continuous series of birational transformations into themselves not forming a group. They divide into

\(A')\) (regular) rational surfaces, characterized by the conditions
\[
p_a=P_2=0,
\]
and

\(A'')\) irrational ruled surfaces, for which
\[
p_a<0;
\]
the ruled surfaces of genus \(p=-p_a>1\) are characterized (as we shall see in \S~8) by the simple condition
\[
p_a<-1,
\]
while the ruled surfaces of genus \(p=1\) are distinguished from the properly elliptic surfaces with \(p_a=-1\) by having \(P_4=P_6=P_{12}=0\)). %

\(B)\) \(P_{12}=1,\ p^{(1)}=1\) (\(p_a\geq-1\)). Surfaces possessing a canonical or pluricanonical curve of order zero. For every system without base points, of genus \(\pi\) and degree \(n\), on the surface freed of exceptional curves, one has
\[
n=2\pi-2,
\]
so that one obtains, in general, as projective images of this type, surfaces \(F_n\) with plane or hyperplane sections of genus \(\pi\), of order \(n=2\pi-2\), in a space of \(r=\pi,\ \pi-1\) or \(\pi-2\) dimensions (the three values of \(r\) corresponding to the three cases: \(p_a=p_g=1\), \(p_a=p_g=0\), and \(p_a=-1\)).

The surfaces \(F_n\) of family \(B)\) divide into two categories:

\(B')\) surfaces with canonical (plane or hyperplane) sections, which have geometric genus \(p_g=1\) and arithmetic genus \(p_a=1\) or
% scan: s242R
% !PRINTED 457
\weberpage{457}
\(p_a=-1\), that is, regular surfaces with all genera one (Chap. VII, \S~2) and hyperelliptic surfaces;

\(B'')\) surfaces whose sections are curves that are not canonical but represent linear series arising from the division by \(s\) of the \(s\)-canonical series (\(s=2,3,4,6\)).

These are the regular surfaces of genus \(p_a=p_g=0\) and \(P_2=1\) (\(P_3=0\)), reducible to the sextic \(F_6\) that passes doubly through the edges of a tetrahedron (Chap. VII, \S~1), and the elliptic surfaces with \(p_a=-1\) and \(p_g=0\) containing a linear pencil of elliptic curves, which we have distinguished into seven types \(I_a\,I_b\,II_a\,II_b\,III_a\,III_b\,III_c\) in \S\S~16-19 of Chap. X. %

Between the surfaces \(B'\) and the \(B''\) there is this difference: that, while the \(B''\) give rise to a finite number of types, the \(B'\) give rise to an infinite number of families of surfaces, depending on (at least) one integer character, capable of taking infinitely many values.

This has been proved for the regular surfaces with all genera one in \S~2 of Chap. VII; (where it was also seen that 19 moduli belong to each family); but it is also true for hyperelliptic surfaces, where, alongside the Jacobi surface \(F_j\) corresponding to the curve of genus two, there are the infinitely many others which correspond to involutions of order \(\delta\) on \(F_j\) or to \(F_j\) taken \(\delta\)-fold, without branch curve: each class being defined within the respective family by 3 moduli (cf. \S~6).

On the other hand, the surfaces \(B')\) give rise, as has been said, to a finite number, and precisely to 8 families, sharply characterized and containing: the first (regular surfaces) 10 moduli (Chap. VII, \S~1) and the others (elliptic) 2 moduli (the \(I_a\,I_b\)) or 1 modulus (the remaining ones, of the harmonic or equianharmonic type). %

Let us add that the surfaces \(B')\) always contain pencils of elliptic curves and therefore come close to the surfaces \(C)\), in contrast with the \(B'')\), which do not in general contain such pencils. %

\(C)\) \(P_{12}>1\) and \(p^{(1)}=1\). Surfaces possessing a (rational or irrational) pencil of elliptic curves \(K\), components of canonical or pluricanonical curves. Together with the \(B'')\) these surfaces form the vast family of \textit{surfaces that contain a pencil of elliptic curves.}

These surfaces can be classified, as we have seen in the regular case (Chap. VII, \S~4), according to certain integer characters, namely:

1) the arithmetic genus \(p_a\);

2) the determinant \(n\), that is, the order of the smallest group of points that can be determined rationally on a \(K\);

% scan: s243L
% !PRINTED 458
\weberpage{458}
3) the number of curves of the pencil that reduce to a multiple curve; to which is added, in the irregular case,

4) the genus \(\varrho\) of the pencil, which is, in general (\textit{Remark}, Chap. IX, \S~8) %
\[
\begin{aligned}
&\varrho=p_g-p_a\\
&\quad(p_a\geq0)
\end{aligned}
\]
and, instead,
\[
\varrho=p_g-p_a-1
\]
in the case in which the surface contains, besides \((K)\), a second elliptic pencil of curves, a case which entails \(p_a=-1\) and leads to the elliptic surfaces.

For the regular surfaces \(C)\) we have recognized that their infinitely many families, depending on the characters indicated above, can be distinguished by means of the values of the plurigenera (Chap. VII, \S~4).

These results seem, in large part, to extend also to \textit{irregular surfaces}: \(p_a<p_g\). In any case there is here a \textit{subject of study}, which we point out to the attention of the young cultivators of geometry.

\(D)\) \(P_{12}>1,\ p^{(1)}>1\). These surfaces, which must be regarded as constituting the \textit{most general case}, unlike the preceding \(A)\ B)\) and \(C)\), give rise, for given values of the characters \(p_a,\ p_g,\ p^{(1)}\), to a \textit{finite number of families}, each of which contains a continuous system of classes, depending on a \textit{certain number} of \textit{parameters} or \textit{moduli}. Indeed to the surfaces \(D)\) there corresponds a type of canonical or pluricanonical surface, whose order is \(p^{(1)}-1\) or resp. a multiple of \(p^{(1)}-1\), which offers a projective model of the surfaces themselves, regarded with respect to birational transformations: the classification of the surfaces of a given order, with certain given characters, can lead only to a finite number of arithmetically distinct types.

We have already had occasion to study the \textit{general problem of the classification of surfaces with \(p^{(1)}>1\) in the regular case} (Chap. VIII). The considerations developed extend, in large part, to the case in which the hypothesis \(p_a=p_g\) is dropped; in any case there is here a \textit{subject of research} for scholars. It is noteworthy that, when one proceeds to construct the canonical or pluricanonical surfaces with given \(p_g\) and \(p^{(1)}\), one usually encounters regular surfaces, and only by way of exception does the possibility of irregular surfaces appear; thus the canonical surfaces for \(p_g\geq4\) which are constructed in the first cases \(p^{(1)}=5,6,7\ldotp\ldotp\ldotp\ldotp\), (Chap. VIII) turn out to be regular, and similarly it seems that the following also necessarily turn out to be regular: the tricanonical surfaces for \(p^{(1)}=2\) and the bicanonical ones for \(p^{(1)}=3\) which are constructed in
% scan: s243R
% !PRINTED 459
\weberpage{459}
correspondence with a \(p_a>0\), according to the procedure of \S\S~14-20 of Chap. VIII,\origfootnote{1}{s243R-1}{Considerations which are completed by the theorems on the simple integrals belonging to a surface lead in general to the assertion that the surfaces whose canonical system is not composed of the curves of a pencil, of irregularity \(q\geq3\) with \(p_a\geq0\), have linear genus \(p^{(1)}\geq6p_a+4\); if instead \(q=2\) and \(p_a\geq0\), then \(p^{(1)}\geq3p_a+4\).

Therefore there would remain to be examined only the surfaces of irregularity \(q=1\) possessing an elliptic pencil of curves (Communication kindly made to us by G. \textsc{Castelnuovo}, November 1942).}. %

Here it is appropriate to add that, by virtue of an inequality which we discuss in the following section, the surfaces \(D)\) (with \(p^{(1)}>1\)) all have arithmetic genus: \(p_a\geq0\).

\edsection{8. Surfaces of negative arithmetic genus.}{8. Surfaces of negative arithmetic genus.}{sec.XI.8}

The classification that we have set forth in the preceding section can be narrowed down more closely on the basis of an inequality between the irregularity and the geometric genus, which up to now it has not been possible to prove in an algebro-geometric way, and which must therefore be drawn from the transcendental theory (cf. Historical note, p. 463). The following theorem holds: \textit{For every surface of genus \(p_g\) and irregularity \(q=p_g-p_a\) that does not contain an irrational pencil of curves, one has} %
\[
p_g>2(q-2)
\]
that is
\[
p_g<2p_a+4.
\]
Hence: \textit{if \(p_g\geq2p_a+4\), the surface contains an irrational pencil of curves \((C)\)}.

The preceding inequality is certainly satisfied if the arithmetic genus is
\[
p_a<-1\quad(p_g\geq0).
\]

Therefore one can assert that the surfaces \(F\) of arithmetic genus \(p_a<-1\) contain an irrational pencil \((C)\): a pencil of genus \(p>0\) formed of curves \(C\) of a certain genus \(\pi\). But one easily recognizes that \(\pi=0\), and hence \(F\) can be transformed into a ruled surface (of genus \(p=-p_a\)).

To this end let us reason by contradiction: if \(\pi>0\), \(F\) can be reduced to a surface without exceptional curves, of absolute linear genus \(p^{(1)}\geq1\); then — computing its Zeuthen-Segre invariant in relation to \((C)\) — one finds (as in Chap. X, \S~4)
\[
\begin{aligned}
&\Delta-4=12p_a-p^{(1)}+9,\\
&\Delta\geq0\\
&p_a\geq-1\ !
\end{aligned}
\]

% scan: s244L
% !PRINTED 460
\weberpage{460}
We conclude, therefore, that \textit{the surfaces of arithmetic genus \(p_a<-1\) are all birationally equivalent to ruled surfaces} (of genus \(p=-p_a\)): in the theorem (Chap. X, \S~4) which assigns for this purpose the conditions \(p_a=-p\) and \(p_g=0\), the condition \(p_g=0\) is superfluous, since it results as a consequence of the other.

Let us take \(p_a=-1\). If the surface \(F\) has geometric genus
\[
p_g>1,
\]
the inequality
\[
p_g\geq2p_a+4
\]
is still satisfied, and hence \(F\) possesses an irrational pencil of curves \(C\) (of genus \(\pi>0\)). But, in relation to this pencil, one finds the Zeuthen-Segre invariant
\[
\begin{aligned}
&\Delta-4=12p_a-p^{(1)}+9\\
&\quad p_a=-1,\quad\Delta\geq0
\end{aligned}
\]
and hence
\[
p^{(1)}=1.
\]

Thus the condition \(p^{(1)}=1\), which — in addition to \(p_a=-1\) and \(p_g>1\) — characterizes elliptic surfaces (Chap. XI, \S~4), is a consequence of the other two conditions. And, since for \(p_g=0\) and \(p_a=-1\) one again has \(p^{(1)}=1\) and \(F\) turns out to be an elliptic surface (proper or improper), we can state that: \textit{the surfaces of arithmetic genus \(p_a=-1\) and of geometric genus \(p_g\ne1\) are elliptic} (proper elliptic surfaces or ruled surfaces of genus one).

Finally let us examine the case:
\[
p_a=-1,\quad p_g=1.
\]

In this case we have seen (Chap. X, \S~4) that \(F\) is hyperelliptic (with \(P_4=1\)) or else can be reduced to a hyperelliptic surface \(\Phi\), multiple with elliptic branch curve, and then it is itself an elliptic surface (with \(P_4>1\)).

Thus we can complete the theorems established in \S~4 of Chap. X, in the sense already announced there, by saying that:

\textit{Elliptic and hyperelliptic surfaces are characterized globally by the condition}
\[
p_a=-1. %
\]
(joined with those that exclude the improper elliptic surfaces, birationally equivalent to ruled surfaces: \(P_4+P_6=0\), or \(P_{12}=0\)).

\textit{Hyperelliptic surfaces} (which only for particular moduli are also — in two ways — elliptic) \textit{are distinguished from elliptic ones by the conditions}
\[
p_g=P_4=1
\]
(instead of
\[
p_g=1,\quad P_4>1).
\]

% scan: s244R
% !PRINTED 461
\weberpage{461}
The results obtained can be summarized by stating the theorem:

\textit{The surfaces of negative arithmetic genus} (\(p_a<0\)) \textit{are}:

1) \textit{birationally equivalent to ruled surfaces} (irrational).

2) \textit{or else elliptic or hyperelliptic surfaces} (in the latter case: \(p_g=1\)). %

Or also:

\textit{Negative arithmetic genus serves to define the whole family of surfaces, not rational, that admit a continuous group of birational transformations.}

\textit{Historical note.} – The fundamental inequality \(p_g<2p_a+4\), for surfaces that do not contain an irrational pencil of curves, derives — as has been said — from considerations of a transcendental order. Let us state beforehand that the genus \(p_g\) of a surface \(F\) is the number of linearly independent double integrals of the first kind that belong to \(F\), and on the other hand that the irregularity \(q\) is the number of simple integrals of the first kind of the same \(F\); and let us recall that \textsc{Noether} pointed out how one can « construct a double integral of the first kind starting from two simple integrals that are not functions of one another ».

Now the fact that the surface \(F\) possesses two simple integrals, \(I_1\) and \(I_2\), that are functions of one another, means that \(F\) contains an irrational pencil of algebraic curves \(C\), on which \(I_1\) and \(I_2\) remain constant.

This observation presented itself for the first time to \textsc{Castelnuovo}, with regard to surfaces of genus \(p_g=0\) (cf. Chap. IX, \S~6); later (1905), for surfaces of arbitrary genus \(p_g\), to \textsc{De Franchis}\origfootnote{1}{s244R-1}{\textit{Sulle superficie algebriche le quali contengono un fascio irrazionale di curve,} Rendic. Circolo Mat. di Palermo (25 April 1905).}, \textsc{Castelnuovo}\origfootnote{2}{s244R-2}{\textit{Sulle superficie aventi il genere aritmetico negativo.} Ib., 14 May 1905 (Memorie scelte, XXVII, p. 501).} and \textsc{Enriques}\origfootnote{3}{s244R-3}{\textit{Sulle superficie algebriche che ammettono un gruppo continuo di trasformazioni etc.} (Ibid., 14 May 1905).}, who saw here the appropriate means of study for classifying the surfaces of genus \(p_a<0\).

The relation between the genus \(p_g\) and the irregularity \(q\), or between \(p_g\) and \(p_a\), was made precise and proved in the simplest way by \textsc{Castelnuovo} with the aforesaid inequality
\[
p_g<2p_a+4,
\]
valid for the surfaces \(F\) that do not contain an irrational pencil of curves. Meanwhile, on the basis of this relation, \textsc{De Franchis}
% scan: s245L
% !PRINTED 462
\weberpage{462}
stated that the \(F\) of genus \(p_a<-1\) possess an irrational pencil of curves\origfootnote{1}{s245L-1}{Let us add that \textsc{De Franchis} (making precise an earlier result, of 1904) establishes that « \textit{every irregular double plane possesses an elliptic or hyperelliptic pencil of curves} ».}, while \textsc{Castelnuovo} and \textsc{Enriques} — making use of an argument already employed under the hypothesis \(p_g=0\) (Chap. X, \S~4) — recognized that these \textit{surfaces are birationally equivalent to ruled surfaces, of genus \(p=-p_a\) \((>1)\)}.

Finally \textsc{Enriques} deepened the analysis of the surfaces for which \(p_a=-1\), succeeding in recognizing that they are elliptic or hyperelliptic surfaces.

The developments that have been discussed had a sequel in the important works of A. \textsc{Rosenblatt}\origfootnote{2}{s245L-2}{\textit{Sur les surfaces irrégulières etc.} Rendic. Circolo Mat. di Palermo, 1913.} and A. \textsc{Comessatti}\origfootnote{3}{s245L-3}{Rendic. Acc. Lincei, 1919. and \textit{Intorno alle superficie algebriche irregolari etc.} Rendic. Circolo Mat. di Palermo, vol. XLVI, 1922.}, which yield the classification of the surfaces for which
\[
p_g\geq2p_a+4.
\]

Already \textsc{Rosenblatt} (1913) recognizes that they can be reduced to a few types. \textsc{Comessatti} (1919 and 1922) takes up again the analysis of \textsc{Castelnuovo}, treating the more general problem of « determining the relations in finite terms that hold between several functions of two variables when their mutual Jacobians satisfy a sufficient number of linear homogeneous equations with constant coefficients ».

From this he derives, in particular, the classification of the algebraic surfaces for which
\[
p_g\geq2p_a+4,
\]
proving that they belong to the following types:

1) ruled surfaces of genus \(p>1\);

2) elliptic surfaces with \(p_g>1\) and \(p_a=-1\);

3) surfaces representing the pairs of points of two curves of genus 2 and \(\pi\): \(p_g=2\pi,\ p_a=\pi-2\);

4) particular surfaces of linear genus \(p^{(1)}=1\) with \(p_a\geq0\).

For these brilliant results too, as for the inequality of \textsc{Castelnuovo}, it is desirable that one should succeed in giving an algebro-geometric proof.

% scan: s245R
% !PRINTED 463
\weberpage{463}
\edsection{9. Summary of the refined classification.}{9. Summary of the refined classification.}{sec.XI.9}

The theorems set forth on surfaces of negative arithmetic genus serve to narrow down more closely the general classification of surfaces, in particular implying that for \(p^{(1)}>1\) we must have \(p_a\geq0\).

This deduction has a significance of a general order: for the surfaces \(D\) (\(P_{12}>1,\ p^{(1)}>1\)), since \(p_a\geq0\), one obtains a bicanonical system of dimension \(P_2-1=p_a+p^{(1)}-1\), and hence — for \(p^{(1)}>3\) — one gets \(P_2-1\geq3\). Thus, apart from a doubt which must be critically examined\origfootnote{1}{s245R-1}{As has been done for regular surfaces in \S\S~18-20 of chap. IV.} concerning the reducibility of the bicanonical curves, the surfaces for which \(p^{(1)}>3\) admit as projective model a bicanonical surface (which, in particular cases, may be multiple).

And hence there appears the special interest of possessing a complete classification of the surfaces for which \(p^{(1)}=2\) and \(p^{(1)}=3\), to which the studies set forth by us in \S\S~14-20 of Chap. VIII refer; all the more so if it should happen to be found that for such surfaces (apart perhaps from some exception) there is verified a posteriori the hypothesis of regularity which in our studies was admitted a priori\origfootnote{2}{s245R-2}{Cf. the footnote at the end of \S~7 of this Chapter.}.

Let us summarize the more precise results obtained in the general classification of algebraic surfaces, making use also of some formulae whose meaning has been explained above, in \S~7.

\[
\begin{aligned}
&\begin{tabular}{@{}l@{}l@{}}
\begin{tabular}{@{}>{\raggedright\arraybackslash}p{0.33\linewidth}@{}}
\(P_{12}=0\): ruled surfaces \(n>2\pi-2\), exceptional curves that cannot be eliminated, continuous family of transformations not forming a group. %
\end{tabular}
&
\(\left\{
\begin{tabular}{@{}>{\raggedright\arraybackslash}p{0.61\linewidth}@{}}
\(p_a=0\) rational surfaces \((p_a=P_2=0)\)\\
\(p_a=-1\) elliptic ruled surfaces \((p_a=-1,\)\newline \(\qquad P_4=P_6=0)\)\\
\(p_a<-1\) ruled surfaces of genus \(p=-p_a>1\). %
\end{tabular}
\right.\)
\end{tabular}
\end{aligned}
\]

\[
\begin{aligned}
&\begin{tabular}{@{}l@{}l@{}}
\begin{tabular}{@{}>{\raggedright\arraybackslash}p{0.20\linewidth}@{}}
\(P_{12}=1\): virtual canonical curve of ord. 0,\\
\(\qquad n=2\pi-2\)\\
\((p^{(1)}=1)\)
\end{tabular}
&
\(\left\{
\begin{tabular}{@{}l@{}l@{}}
\(p_g=1\) &
\(\left\{
\begin{tabular}{@{}>{\raggedright\arraybackslash}p{0.62\linewidth}@{}}
\(p_a=1,\) infinitely many families of surfaces \((p_a=P_2=1)\)\\ %
\(p_a=-1,\) infinitely many families of hyperelliptic surfaces of divisor \(\delta=1,\ 2,\ \ldotp\ldotp\ldotp\ldotp\)\newline \(\qquad(p_a=-1,\ p_g=P_4=1).\)
\end{tabular}
\right.\)
\\
\(p_g=0\) &
\(\left\{
\begin{tabular}{@{}>{\raggedright\arraybackslash}p{0.62\linewidth}@{}}
\(p_a=0,\) surfaces of the \(6^{\mathrm{th}}\) order passing doubly through the edges of a tetrahedron \((p_a=P_3=0,\ P_2=1).\)\\
\(p_a=-1,\) elliptic surfaces \(I_a\,I_b\,II_a\,II_b\)\newline \(\qquad III_a\,III_b\,III_c\) \((P_2=0,1,\ P_4=0,1,\)\newline \(\qquad P_3=0,1,\ P_4=0,1).\) %
\end{tabular}
\right.\)
\end{tabular}
\right.\)
\end{tabular}
\end{aligned}
\]

\Needspace{110mm}
% scan: s246L
% !PRINTED 464
\weberpage{464}
\[
\begin{aligned}
\begin{gathered}
P_{12}>1\\
(n<2\pi-2)
\end{gathered}
\;
\left\{
\begin{aligned}
p^{(1)}=1\ &\left\{
\begin{aligned}
&p_a=-1,\ \text{elliptic surfaces with a pencil}\\
&\quad\text{of genus }p_g\ \text{of elliptic curves}\\
&\quad(p_a=-1,\ P_4\ne1).\\[1ex]
&p_a\ge0,\ \text{surfaces whose canonical and}\\
&\quad\text{pluricanonical curves are composed of the}\\
&\quad\text{elliptic curves of a pencil of genus}\\
&\quad p_g-p_a:\ \text{they depend on several arbitrary}\\
&\quad\text{integer characters.}
\end{aligned}
\right.\\[1ex]
p^{(1)}>1\ &\left\{
\begin{aligned}
&p_a\ge0,\ \text{surfaces that admit a projective}\\
&\quad\text{model, canonical }(p_g>3)\ \text{or}\\
&\quad\text{bicanonical }(p^{(1)}>3)\ \text{etc.: a finite}\\
&\quad\text{number of types for a given value of}\\
&\quad p^{(1)}.
\end{aligned}
\right.
\end{aligned}
\right.
\end{aligned}
\]

At this point may we be allowed to pause for a moment, as in an alpine ascent one loves to halt on the peak attained and from there contemplate the spectacle of Nature that offers itself to the view.

Fifty years ago the study of these theories began in Italy, theories barely sketched by the genius of a precursor (\textsc{Max Noether}); at that time, joking about the difficulties and the exceptions that were encountered on every side, people used to say that, while algebraic curves (already composed into a harmonious theory) are created by God, surfaces instead are the work of the Devil.

Now, on the contrary, it becomes manifest that it pleased God to create for surfaces an order of more recondite harmonies, in which a marvellous beauty shines forth, and that He willed in them — let us say with the Poet —
\begin{verse}
of His creating spirit\\
to stamp a vaster imprint.
\end{verse}

The richness of the properties and the beauty, long hidden, that are revealed here must not constitute a reason for vain pride for the Italian geometric school or for the foreign mathematicians who have collaborated in discovering them, but rather must arouse a sense of reverence for that marvellous order of mathematical entities which thought finds before itself and almost gathers, like the living species, from Mother Nature; and thus they must nourish the faith of young researchers that behind the difficulties, the exceptions, the apparent incongruities, there really is, in this world of entities, a divine harmony, which the concerted efforts of scholars will succeed ever better in bringing to light. %
% scan: s246R
% !PRINTED none
\weberblank{[s246R]}

\par\edFlush\clearpage

\end{document}
